\beginnumbering
\eledsection*{Adhyāya 1}
\pstart \ledsidenote{\passref{1.1.0}}\edtext{\skt{नमः}}{\lemma{\skt{नमः}}\Afootnote{N K; \skt{ओं नमो}\ H}}\skt{ }\edtext{\skt{कमलहस्ताय}}{\lemma{\skt{कमलहस्ताय}}\Afootnote{K; \skt{नागा\appbrk{}र्जु\appbrk{}न\appbrk{}पा\appbrk{}दाय}\ N; \skt{धन्वन्तरये}\ H}}\skt{~॥}\pend

\pstart \ledsidenote{\passref{1.1.1}}\skt{अथातो }\edtext{\skt{वेदो\-त्प\-त्ति\-म\-ध्यायं}}{\lemma{\skt{अध्यायं}}\Afootnote{A K\textsuperscript{pc}; \skt{आध्यायं}\ K\textsuperscript{ac} N; \skt{नामाध्यायं}\ H}}\skt{ व्याख्यास्यामः~।}\pend

\pstart \ledsidenote{\passref{1.1.3}}\skt{अथ खलु }\edtext{\skt{भग\-व\-न्त\-म\-म\-र\-व\-र\-मृ\-षि\-ग\-ण\-प\-रि\-वृ\-त्त\-मा\-श्र\-मस्थं}}{\lemma{\skt{अमरवरम्}}\Afootnote{A K H; \skt{अमरवर}\ N}}\skt{ काशिराजं }\edtext{\edtext{\skt{दिवो\-दा\-स\-मौ\-प\-धे\-न\-व\-वै\-त\-र\-णौ\-र\-भ्र\-पु\-ष्क\-ला\-व\-त\-क\-र\-वी\-र\-गो\-पु\-र\-र\-क्षि\-त\-भो\-ज\-सु\-श्रु\-त\-प्र\-भृ\-तय}}{\lemma{\skt{दिवोदासम्}}\Afootnote{A K\textsuperscript{pc} N H; \skt{दिवोदासम्म्}\ K\textsuperscript{ac}}}}{\lemma{\skt{औप\appbrk{}धे\appbrk{}न\appbrk{}व\appbrk{}वै\appbrk{}त\appbrk{}र\appbrk{}णौ\appbrk{}र\appbrk{}भ्र\appbrk{}पु\appbrk{}ष्क\appbrk{}ला\appbrk{}व\appbrk{}त\appbrk{}क\appbrk{}र\appbrk{}वी\appbrk{}र\appbrk{}गो\appbrk{}पु\appbrk{}र\appbrk{}र\appbrk{}क्षि\appbrk{}त\appbrk{}भो\appbrk{}ज\appbrk{}सु\appbrk{}श्रु\appbrk{}त\appbrk{}प्र\appbrk{}भृ\appbrk{}तय}}\Afootnote{\skt{धन्व\appbrk{}न्त\appbrk{}रि\appbrk{}मौ\appbrk{}प\appbrk{}धे\appbrk{}न\appbrk{}व\appbrk{}वै\appbrk{}त\appbrk{}र\appbrk{}णौ\appbrk{}र\appbrk{}भ्र\appbrk{}पौ\appbrk{}ष्क\appbrk{}ला\appbrk{}व\appbrk{}त\appbrk{}क\appbrk{}र\appbrk{}वी\appbrk{}र्य(\appbrk{}र)\appbrk{}गो\appbrk{}पु\appbrk{}र\appbrk{}र\appbrk{}क्षि\appbrk{}त\appbrk{}सु\appbrk{}श्रु\appbrk{}त\appbrk{}प्र\appbrk{}भृ\appbrk{}तय}\ A; \skt{औपधेनवः वैत\appbrk{}र\appbrk{}णो\appbrk{}र\appbrk{}भ्र\appbrk{}पु\appbrk{}ष्क\appbrk{}ला\appbrk{}व\appbrk{}त\appbrk{}क\appbrk{}र\appbrk{}वी\appbrk{}र\appbrk{}गो\appbrk{}पु\appbrk{}र\appbrk{}र\appbrk{}क्षि\appbrk{}त\appbrk{}भो\appbrk{}ज\appbrk{}सु\appbrk{}श्रु\appbrk{}त\appbrk{}प्र\appbrk{}भृ\appbrk{}तय}\ K\textsuperscript{ac}; \skt{धन्व\appbrk{}न्त\appbrk{}रि\appbrk{}मौ\appbrk{}प\appbrk{}धे\appbrk{}नवः वैत\appbrk{}र\appbrk{}णो\appbrk{}र\appbrk{}भ्र\appbrk{}पु\appbrk{}ष्क\appbrk{}ला\appbrk{}व\appbrk{}त\appbrk{}क\appbrk{}र\appbrk{}वी\appbrk{}र\appbrk{}गो\appbrk{}पु\appbrk{}र\appbrk{}र\appbrk{}क्षि\appbrk{}त\appbrk{}भो\appbrk{}ज\appbrk{}सु\appbrk{}श्रु\appbrk{}त\appbrk{}प्र\appbrk{}भृ\appbrk{}तय}\ K\textsuperscript{pc}; \skt{औप\appbrk{}धे\appbrk{}न\appbrk{}व\appbrk{}वै\appbrk{}स्त\appbrk{}र\appbrk{}णौ\appbrk{}र\appbrk{}भ्र\appbrk{}पौ\appbrk{}ष्क\appbrk{}ला\appbrk{}व\appbrk{}त\appbrk{}क\appbrk{}र\appbrk{}वी\appbrk{}र\appbrk{}गो\appbrk{}पु\appbrk{}र\appbrk{}र\appbrk{}क्षि\appbrk{}त\appbrk{}भो\appbrk{}ज\appbrk{}सु\appbrk{}श्रु\appbrk{}त\appbrk{}प्र\appbrk{}भृ\appbrk{}तव्य}\ N; \skt{धन्व\appbrk{}न्त\appbrk{}रि\appbrk{}मौ\appbrk{}प\appbrk{}धे\appbrk{}न\appbrk{}व\appbrk{}वै\appbrk{}र\appbrk{}णो\appbrk{}र\appbrk{}भ्र\appbrk{}पु\appbrk{}ष्क\appbrk{}ला\appbrk{}व\appbrk{}त\appbrk{}क\appbrk{}र\appbrk{}वी\appbrk{}र\appbrk{}गो\appbrk{}पु\appbrk{}र\appbrk{}र\appbrk{}क्षि\appbrk{}त\appbrk{}भो\appbrk{}ज\appbrk{}सु\appbrk{}श्रु\appbrk{}त\appbrk{}प्र\appbrk{}भृ\appbrk{}तय}\ H}}\footnoteA{The scribe of the earliest MS, K, added the name \emph{dhanvantari} in the margin; the scribe of N did not have the name; the scribe of H included it in the text.}\skt{ }\edtext{\skt{ऊचुः}}{\lemma{\skt{ऊचुः}}\Afootnote{A N H; \skt{ऊचु}\ K}}\skt{~।}\pend

\pstart \ledsidenote{\passref{1.1.4}}\edtext{\edtext{\edtext{\edtext{\skt{भग\-व\-ञ्शा\-री\-र\-मा\-न\-सा\-ग\-न्तु\-भि\-र्व्या\-धि\-भि\-र्वि\-वि\-ध\-वे\-द\-ना\-भि\-घा\-तो\-प\-द्रु\-ता\-न्स\-ना\-था\-न\-ना\-थ\-व\-द्वि\-चे\-ष्ट\-मा\-ना\-न्वि\-क्रो\-श\-तश्च}}{\lemma{\skt{भग\appbrk{}व\appbrk{}ञ्शा\appbrk{}री\appbrk{}र\appbrk{}मा\appbrk{}न\appbrk{}सा\appbrk{}ग\appbrk{}न्तु\appbrk{}भिर्}}\Afootnote{\skt{भग\appbrk{}व\appbrk{}न्शा\appbrk{}री\appbrk{}र\appbrk{}मा\appbrk{}न\appbrk{}सा\appbrk{}ग\appbrk{}न्तु\appbrk{}भिर्}\ A; \skt{भग\appbrk{}व\appbrk{}री\appbrk{}र\appbrk{}मा\appbrk{}न\appbrk{}सा\appbrk{}गन्तु भि…}\ K\textsuperscript{ac}; \skt{भग\appbrk{}व\appbrk{}च्छा\appbrk{}री\appbrk{}र\appbrk{}मा\appbrk{}न\appbrk{}सा\appbrk{}गन्तु भि…}\ K\textsuperscript{pc}; \skt{भग\appbrk{}व\appbrk{}ञ्छ\appbrk{}री\appbrk{}र\appbrk{}मा\appbrk{}न\appbrk{}सा\appbrk{}ग\appbrk{}न्तु\appbrk{}भिर्}\ N; \skt{भग\appbrk{}व\appbrk{}च्छा\appbrk{}री\appbrk{}र\appbrk{}मा\appbrk{}न\appbrk{}सा\appbrk{}ग\appbrk{}न्तु\appbrk{}भिर्}\ H}}}{\lemma{\skt{व्याधिभिर्}}\Afootnote{A K H; \skt{व्याधिभि}\ N}}}{\lemma{\skt{विवि\appbrk{}ध\appbrk{}वे\appbrk{}द\appbrk{}ना\appbrk{}भि\appbrk{}घा\appbrk{}तो\appbrk{}प\appbrk{}द्रु\appbrk{}ता\appbrk{}न्स\appbrk{}ना\appbrk{}था\appbrk{}न\appbrk{}ना\appbrk{}थ\appbrk{}व\appbrk{}द्वि\appbrk{}चे\appbrk{}ष्ट\appbrk{}मा\appbrk{}ना\appbrk{}न्वि\appbrk{}क्रो\appbrk{}श\appbrk{}तश्}}\Afootnote{\skt{विवि\appbrk{}ध\appbrk{}वे\appbrk{}द\appbrk{}ना\appbrk{}भि\appbrk{}घा\appbrk{}तो\appbrk{}प\appbrk{}द्रु\appbrk{}ता\appbrk{}न्स\appbrk{}ना\appbrk{}था\appbrk{}न\appbrk{}प्य\appbrk{}ना\appbrk{}थ\appbrk{}व\appbrk{}द्वि\appbrk{}चे\appbrk{}ष्ट\appbrk{}मा\appbrk{}ना\appbrk{}न्वि\appbrk{}क्रो\appbrk{}श\appbrk{}तश्}\ A; \skt{वि…ध\appbrk{}वे\appbrk{}द\appbrk{}ना\appbrk{}भि\appbrk{}घा\appbrk{}तो\appbrk{}प\appbrk{}द्रुतां सना\appbrk{}था\appbrk{}न\appbrk{}ना\appbrk{}थ\appbrk{}व\appbrk{}द्वि\appbrk{}चे\appbrk{}ष्ट\appbrk{}मा\appbrk{}ना\appbrk{}न्वि\appbrk{}क्रो\appbrk{}शत}\ K; \skt{विवि\appbrk{}ध\appbrk{}वे\appbrk{}द\appbrk{}ना\appbrk{}भि\appbrk{}घा\appbrk{}तो\appbrk{}प\appbrk{}द्रु\appbrk{}ता\appbrk{}ना\appbrk{}स\appbrk{}ना\appbrk{}था\appbrk{}न\appbrk{}ना\appbrk{}थ\appbrk{}व\appbrk{}द्वि\appbrk{}चे\appbrk{}ष्ट\appbrk{}मा\appbrk{}ना\appbrk{}न्वि\appbrk{}क्रो\appbrk{}श\appbrk{}तश्}\ N; \skt{व्विवि\appbrk{}ध\appbrk{}वे\appbrk{}द\appbrk{}ना\appbrk{}भि\appbrk{}घा\appbrk{}तो\appbrk{}प\appbrk{}द्रुतां सना\appbrk{}था\appbrk{}न\appbrk{}ना\appbrk{}थ\appbrk{}व\appbrk{}द्वि\appbrk{}चे\appbrk{}ष्ट\appbrk{}मा\appbrk{}ना\appbrk{}न्वि\appbrk{}क्रो\appbrk{}श\appbrk{}तश्}\ H}}}{\lemma{\skt{च}}\Afootnote{A N H; \skt{…श्च}\ K}}\skt{ }\edtext{\skt{मान\-वा\-न\-भि\-स\-मीक्ष्य}}{\lemma{\skt{मान\appbrk{}वा\appbrk{}न\appbrk{}भि\appbrk{}स\appbrk{}मीक्ष्य}}\Afootnote{A K\textsuperscript{ac} N; \skt{मान\appbrk{}वा\appbrk{}न\appbrk{}म\appbrk{}भि\appbrk{}स\appbrk{}मीक्ष्य}\ K\textsuperscript{pc}; \skt{मान\appbrk{}वा\appbrk{}ना\appbrk{}म\appbrk{}भि\appbrk{}स\appbrk{}मीक्ष्य}\ H}}\skt{ मनसि नः }\edtext{\skt{पीडाभवत्}}{\lemma{\skt{पीडाभवत्}}\Afootnote{K N H; \skt{पीडा भवति}\ A}}\skt{~।}\pend

\pstart \ledsidenote{\passref{1.1.4b}}\skt{तेषां सुखैषिणां }\edtext{\skt{रोगो\-प\-श\-मा\-र्थ\-मा\-त्म\-नश्च}}{\lemma{\skt{रोगो\appbrk{}प\appbrk{}श\appbrk{}मा\appbrk{}र्थम्}}\Afootnote{A K H; \skt{रोगो\appbrk{}प\appbrk{}ग\appbrk{}मा\appbrk{}र्थम्}\ N}}\skt{ }\edtext{\edtext{\skt{प्राण\-या\-त्रा\-र्थ\-मा\-यु\-र्वेदं}}{\lemma{\skt{प्राणयात्रार्थम्}}\Afootnote{A K H; \skt{प्राणयातार्थम्}\ N}}}{\lemma{\skt{आयुर्वेदं}}\Afootnote{N H; \skt{प्रजा\appbrk{}हि\appbrk{}त\appbrk{}हे\appbrk{}तो\appbrk{}रा\appbrk{}यु\appbrk{}र्वेदं}\ A; \skt{आयु}\ K}}\skt{ }\edtext{\skt{इच्छाम}}{\lemma{\skt{इच्छाम}}\Afootnote{N H; \skt{श्रोतुमिच्छाम}\ A; \skt{श्वेदमिच्छाम}\ K}}\skt{ }\edtext{\skt{उप\-दि\-श्य\-मा\-नम्}}{\lemma{\skt{उप\appbrk{}दि\appbrk{}श्य\appbrk{}मा\appbrk{}नम्}}\Afootnote{K N H; \skt{इहो\appbrk{}प\appbrk{}दि\appbrk{}श्य\appbrk{}मा\appbrk{}नम्}\ A}}\skt{~। }\edtext{\edtext{\skt{अत्रा\-य\-त्त\-मै\-हि\-क\-मा\-मु\-ष्मि\-कञ्च}}{\lemma{\skt{अत्रायत्तम्}}\Afootnote{A K N; \skt{अत्त्राय उत्तम्}\ H}}}{\lemma{\skt{ऐहिकम्}}\Afootnote{A K N; \skt{औहिकम्}\ H}}\skt{ }\edtext{\edtext{\skt{श्रेय\-स्त\-द्भ\-ग\-व\-न्त\-मु\-प\-सन्नाः}}{\lemma{\skt{तद्भगवन्तम्}}\Afootnote{K N H; \skt{तद्भगवन्तम्}\ A}}\skt{ स्मः}}{\lemma{\skt{उपसन्नाः स्मः}}\Afootnote{\skt{उपपन्नाः स्मः}\ A; \skt{उपसन्ना स्म}\ K N; \skt{उपसन्ना स्मः}\ H}}\skt{ }\edtext{\skt{शिष्यत्वेनेति}}{\lemma{\skt{शिष्यत्वेनेति}}\Afootnote{A N H; \skt{शिष्यत्वे}\ K}}\skt{~।}\pend

\pstart \ledsidenote{\passref{1.1.5}}\skt{तानुवाच भग\-वा\-न्स्वा\-ग\-तम्वः }\edtext{\skt{सर्व}}{\lemma{\skt{सर्व}}\Afootnote{A K H; \skt{सव}\ N}}\skt{ एवामीमांस्या }\edtext{\skt{अध्याप्याश्च}}{\lemma{\skt{अध्याप्याश्च}}\Afootnote{A K\textsuperscript{pc} N H; \skt{अ प्याश्च}\ K\textsuperscript{ac}}}\skt{ }\edtext{\skt{भवन्तो}}{\lemma{\skt{भवन्तो}}\Afootnote{A K N; \skt{भगवन्तो}\ H}}\skt{ वत्साः~।}\pend

\pstart \ledsidenote{\passref{1.1.6}}\skt{इह }\edtext{\skt{खल्वायुर्वेदो}}{\lemma{\skt{खल्वायुर्वेदो}}\Afootnote{K N H; \skt{खल्वायुर्वेदं}\ A}}\skt{ }\edtext{\skt{नाम}}{\lemma{\skt{नाम}}\Afootnote{K N H; \skt{नामो}\ A}}\skt{ }\edtext{\edtext{\edtext{\skt{यदु\-पा\-ङ्ग\-म\-थ\-र्व\-वे\-द\-स्यो\-क्त\-म\-नु\-त्पा\-द्यैव}}{\lemma{\skt{यदुपाङ्गम्}}\Afootnote{H; \skt{पाङ्गम्}\ A; \skt{यमुपाङ्गम्}\ K; \skt{यदुपागम्}\ N}}}{\lemma{\skt{अथ\appbrk{}र्व\appbrk{}वे\appbrk{}द\appbrk{}स्यो\appbrk{}क्तम्}}\Afootnote{K N H; \skt{अथर्ववेदस्या}\ A}}\skt{ च}}{\lemma{\skt{अनुत्पाद्यैव च}}\Afootnote{K N H; \skt{नुत्पाद्यैव}\ A}}\skt{ प्रजाः }\edtext{\skt{श्लोक\-श\-त\-स\-ह\-स्र\-म\-ध्या\-य\-स\-ह\-स्रञ्च}}{\lemma{\skt{श्लोक\appbrk{}श\appbrk{}त\appbrk{}स\appbrk{}ह\appbrk{}स्रम्}}\Afootnote{A K\textsuperscript{pc} N H; \skt{श्लोकसहस्रम्}\ K\textsuperscript{ac}}}\skt{ }\edtext{\edtext{\skt{कृत\-वा\-न्स्व\-य\-म्भू\-र\-ल्पा\-यु\-ष्का\-ल्प\-मे\-ध\-स्त्व\-ञ्चा\-लोक्य}}{\lemma{\skt{अल्पा\appbrk{}यु\appbrk{}ष्का\appbrk{}ल्प\appbrk{}मे\appbrk{}ध\appbrk{}स्त्वञ्}}\Afootnote{K N H; \skt{ततोऽ\appbrk{}ल्पा\appbrk{}यु\appbrk{}ष्ट्व\appbrk{}म\appbrk{}ल्प\appbrk{}मे\appbrk{}धस्त्वं}\ A}}}{\lemma{\skt{चालोक्य}}\Afootnote{A K N; \skt{च अलोक्य}\ H}}\skt{ नरा\-णा\-म्भू\-योऽ\-ष्टधा }\edtext{\skt{प्रणीतवान्}}{\lemma{\skt{प्रणीतवान्}}\Afootnote{A N H; \skt{प्रणी\appbrk{}त\appbrk{}वा\appbrk{}न्त\appbrk{}द्यथा}\ K}}\skt{~।}\pend

\pstart \ledsidenote{\passref{1.1.7}}\edtext{\skt{तद्यथा}}{\lemma{\skt{तद्यथा}}\Afootnote{A N H; \emph{om.}\ K}}\skt{ शल्यं शालाक्यं }\edtext{\skt{कायचिकित्सा}}{\lemma{\skt{कायचिकित्सा}}\Afootnote{A N H; \skt{का…चिकित्सा}\ K}}\skt{ भूतविद्या }\edtext{\skt{कौमा\-र\-भृ\-त्य\-म\-ग\-द\-तन्त्रं}}{\lemma{\skt{अगदतन्त्रं}}\Afootnote{A K N; \skt{अगदतन्त्र}\ H}}\skt{ रसायनतन्त्रं }\edtext{\edtext{\skt{वाजी\-क\-र\-ण\-त\-न्त्र\-मिति}}{\lemma{\skt{वाजी\appbrk{}क\appbrk{}र\appbrk{}ण\appbrk{}त\appbrk{}न्त्रम्}}\Afootnote{A N H; \skt{व्याजी\appbrk{}क\appbrk{}र\appbrk{}ण\appbrk{}त\appbrk{}न्त्रम्}\ K}}}{\lemma{\skt{इति}}\Afootnote{A N H; \skt{इति~। …}\ K}}\skt{~।}\pend

\pstart \ledsidenote{\passref{1.1.8.1}}\edtext{\skt{अथास्य}}{\lemma{\skt{अथास्य}}\Afootnote{A N H; \skt{स्य}\ K}}\skt{ }\edtext{\skt{प्रत्ये\-का\-ङ्ग\-ल\-क्ष\-ण\-स\-मासः}}{\lemma{\skt{प्रत्ये\appbrk{}का\appbrk{}ङ्ग\appbrk{}ल\appbrk{}क्ष\appbrk{}ण\appbrk{}स\appbrk{}मासः}}\Afootnote{K N H; \skt{प्रत्य\appbrk{}ङ्ग\appbrk{}ल\appbrk{}क्ष\appbrk{}ण\appbrk{}स\appbrk{}मासः}\ A}}\skt{~।}\pend

\pstart \ledsidenote{\passref{1.1.8.1a}}\skt{तत्र }\edtext{\skt{शल्यन्नाम}}{\lemma{\skt{शल्यन्}}\Afootnote{A H; \skt{शंल्यन्}\ K N}}\skt{ }\edtext{\skt{विवि\-ध\-तृ\-ण\-का\-ष्ठ\-पा\-षा\-ण\-पां\-सु\-लो\-ह\-लो\-ष्टा\-स्थि\-बा\-ल\-न\-ख\-पू\-या\-स्रा\-व\-दु\-ष्ट\-व्र\-णा\-न्त\-र्ग\-र्भ\-श\-ल्यो\-द्ध\-र\-णार्थं}}{\lemma{\skt{विवि\appbrk{}ध\appbrk{}तृ\appbrk{}ण\appbrk{}का\appbrk{}ष्ठ\appbrk{}पा\appbrk{}षा\appbrk{}ण\appbrk{}पां\appbrk{}सु\appbrk{}लो\appbrk{}ह\appbrk{}लो\appbrk{}ष्टा\appbrk{}स्थि\appbrk{}बा\appbrk{}ल\appbrk{}न\appbrk{}ख\appbrk{}पू\appbrk{}या\appbrk{}स्रा\appbrk{}व\appbrk{}दु\appbrk{}ष्ट\appbrk{}व्र\appbrk{}णा\appbrk{}न्त\appbrk{}र्ग\appbrk{}र्भ\appbrk{}श\appbrk{}ल्यो\appbrk{}द्ध\appbrk{}र\appbrk{}णार्थं}}\Afootnote{K; \skt{विवि\appbrk{}ध\appbrk{}तृ\appbrk{}ण\appbrk{}का\appbrk{}ष्ठ\appbrk{}पा\appbrk{}षा\appbrk{}ण\appbrk{}पां\appbrk{}शु\appbrk{}लो\appbrk{}ह\appbrk{}लो\appbrk{}ष्टा\appbrk{}स्थि\appbrk{}बा\appbrk{}ल\appbrk{}न\appbrk{}ख\appbrk{}पू\appbrk{}या\appbrk{}स्रा\appbrk{}व\appbrk{}दु\appbrk{}ष्ट\appbrk{}व्र\appbrk{}णा\appbrk{}न्त\appbrk{}र्ग\appbrk{}र्भ\appbrk{}श\appbrk{}ल्यो\appbrk{}द्ध\appbrk{}र\appbrk{}णार्थं}\ A; \skt{विवि\appbrk{}ध\appbrk{}तृ\appbrk{}ण\appbrk{}का\appbrk{}ष्ठ\appbrk{}पा\appbrk{}षा\appbrk{}ण\appbrk{}पां\appbrk{}सु\appbrk{}लो\appbrk{}ह\appbrk{}लो\appbrk{}ष्टा\appbrk{}स्थि\appbrk{}बा\appbrk{}ल\appbrk{}न\appbrk{}पू\appbrk{}या\appbrk{}स्रा\appbrk{}व\appbrk{}दु…\appbrk{}णा\appbrk{}न्…\appbrk{}र्ग\appbrk{}र्भ\appbrk{}श\appbrk{}ल्यो\appbrk{}द्ध\appbrk{}र\appbrk{}णार्थं}\ N; \skt{विवि\appbrk{}ध\appbrk{}तृ\appbrk{}ण\appbrk{}का\appbrk{}ष्ठ\appbrk{}पा\appbrk{}षा\appbrk{}ण\appbrk{}पां\appbrk{}शु\appbrk{}लो\appbrk{}ष्टा\appbrk{}स्थि\appbrk{}बा\appbrk{}ल\appbrk{}न\appbrk{}ख\appbrk{}पू\appbrk{}या\appbrk{}श्रा\appbrk{}व\appbrk{}दु\appbrk{}ष्ट\appbrk{}व्र\appbrk{}णा\appbrk{}न्त\appbrk{}र्ग\appbrk{}र्भ\appbrk{}श\appbrk{}ल्यो\appbrk{}द्ध\appbrk{}र\appbrk{}णार्थं}\ H}}\skt{ }\edtext{\skt{यन्त्र\-श\-स्त्र\-क्षा\-रा\-ग्नि\-प्र\-णि\-धा\-न\-वि\-नि\-श्च\-या\-र्थञ्च}}{\lemma{\skt{यन्त्र\appbrk{}श\appbrk{}स्त्र\appbrk{}क्षा\appbrk{}रा\appbrk{}ग्नि\appbrk{}प्र\appbrk{}णि\appbrk{}धा\appbrk{}न\appbrk{}वि\appbrk{}नि\appbrk{}श्च\appbrk{}या\appbrk{}र्थञ्च}}\Afootnote{K N H; \skt{यन्त्र\appbrk{}श\appbrk{}स्त्र\appbrk{}क्षा\appbrk{}रा\appbrk{}ग्नि\appbrk{}प्र\appbrk{}णि\appbrk{}धा\appbrk{}न\appbrk{}व्र\appbrk{}ण\appbrk{}वि\appbrk{}नि\appbrk{}श्च\appbrk{}यार्थं च (१)}\ A}}\skt{ }\edtext{\edtext{\skt{षष्ट्या\-भि\-धा\-नै\-रिति}}{\lemma{\skt{षष्ट्याभिधानैर्}}\Afootnote{N H; \emph{om.}\ A; \skt{षष्ट्याभिध्…नैर्}\ K}}}{\lemma{\skt{इति}}\Afootnote{K N H; \emph{om.}\ A}}\skt{~।}\pend

\pstart \ledsidenote{\passref{1.1.8.2}}\edtext{\edtext{\skt{शाला\-क्य\-त\-न्त्र\-न्ना\-मो\-र्द्ध्व\-ज\-त्रु\-ग\-तानां}}{\lemma{\skt{शालाक्यतन्त्रन्}}\Afootnote{K N H; \skt{शालाक्यं}\ A}}}{\lemma{\skt{नामो\appbrk{}र्द्ध्व\appbrk{}ज\appbrk{}त्रु\appbrk{}ग\appbrk{}तानां}}\Afootnote{A K N; \skt{नामो\appbrk{}र्ध्व\appbrk{}जा\appbrk{}त्रु\appbrk{}ग\appbrk{}तानां}\ H}}\skt{ }\edtext{\skt{विकाराणां}}{\lemma{\skt{विकाराणां}}\Afootnote{K H; \emph{om.}\ A; \skt{विकाराणा}\ N}}\skt{ }\edtext{\skt{श्रव\-ण\-न\-य\-न\-व\-द\-न\-घ्रा\-णा\-दि\-सं\-श्रि\-तानां}}{\lemma{\skt{श्रव\appbrk{}ण\appbrk{}न\appbrk{}य\appbrk{}न\appbrk{}व\appbrk{}द\appbrk{}न\appbrk{}घ्रा\appbrk{}णा\appbrk{}दि\appbrk{}सं\appbrk{}श्रि\appbrk{}तानां}}\Afootnote{A H; \skt{श्रव\appbrk{}ण\appbrk{}न\appbrk{}य\appbrk{}न\appbrk{}व\appbrk{}द\appbrk{}न\appbrk{}घ्रा\appbrk{}णा\appbrk{}दि\appbrk{}सं…\appbrk{}स्थि\appbrk{}तानां}\ K; \skt{श्रव\appbrk{}ण\appbrk{}न\appbrk{}य\appbrk{}न\appbrk{}व\appbrk{}द\appbrk{}न\appbrk{}घ्रा\appbrk{}णा\appbrk{}दि\appbrk{}सं\appbrk{}श्रि\appbrk{}ताना}\ N}}\skt{ }\edtext{\edtext{\skt{विका\-रा\-णा\-मु\-प\-श\-म\-क\-र\-णा\-र्थम्}}{\lemma{\skt{विकाराणाम्}}\Afootnote{K N H; \skt{व्याधीनाम्}\ A}}}{\lemma{\skt{उप\appbrk{}श\appbrk{}म\appbrk{}क\appbrk{}र\appbrk{}णा\appbrk{}र्थम्}}\Afootnote{K N H; \skt{उप\appbrk{}श\appbrk{}म\appbrk{}ना\appbrk{}र्थम् (२)}\ A}}\skt{~।}\pend

\pstart \ledsidenote{\passref{1.1.8.3}}\edtext{\skt{कायचिकीत्सा}}{\lemma{\skt{कायचिकीत्सा}}\Afootnote{K N H; \skt{कायचिकित्सा}\ A}}\skt{ नाम }\edtext{\skt{सर्व\-श\-री\-रा\-व\-स्थि\-तानां}}{\lemma{\skt{सर्व\appbrk{}श\appbrk{}री\appbrk{}रा\appbrk{}व\appbrk{}स्थि\appbrk{}तानां}}\Afootnote{H; \skt{सर्वा\appbrk{}ङ्ग\appbrk{}सं\appbrk{}श्रि\appbrk{}तानां}\ A; \skt{सर्व\appbrk{}श\appbrk{}री\appbrk{}रा\appbrk{}व\appbrk{}स्थि\appbrk{}ता\appbrk{}नाञ्च}\ K\textsuperscript{ac} N; \skt{सर्व\appbrk{}श\appbrk{}री\appbrk{}रा\appbrk{}व\appbrk{}स्थिता}\ K\textsuperscript{pc}}}\skt{ }\edtext{\edtext{\skt{व्याधी\-ना\-मु\-प\-श\-म\-क\-र\-णार्थं}}{\lemma{\skt{व्याधीनाम्}}\Afootnote{A K\textsuperscript{ac} N H; \skt{धीनाम्}\ K\textsuperscript{pc}}}}{\lemma{\skt{उप\appbrk{}श\appbrk{}म\appbrk{}क\appbrk{}र\appbrk{}णार्थं}}\Afootnote{K N H; \emph{om.}\ A}}\skt{ }\edtext{\edtext{\skt{ज्वर\-शो\-फ\-गु\-ल्म\-र\-क्त\-पि\-त्तो\-न्मा\-दा\-प\-स्मा\-र\-प्र\-मे\-हा\-ति\-सा\-रा\-दी\-नाञ्च}}{\lemma{\skt{ज्वर\appbrk{}शो\appbrk{}फ\appbrk{}गु\appbrk{}ल्म\appbrk{}र\appbrk{}क्त\appbrk{}पि\appbrk{}त्तो\appbrk{}न्मा\appbrk{}दा\appbrk{}प\appbrk{}स्मा\appbrk{}र\appbrk{}प्र\appbrk{}मे\appbrk{}हा\appbrk{}ति\appbrk{}सा\appbrk{}रा\appbrk{}दी\appbrk{}नाञ्}}\Afootnote{K; \skt{ज्वर\appbrk{}र\appbrk{}क्त\appbrk{}पि\appbrk{}त्त\appbrk{}शो\appbrk{}षो\appbrk{}न्मा\appbrk{}दा\appbrk{}प\appbrk{}स्मा\appbrk{}र\appbrk{}कु\appbrk{}ष्ठ\appbrk{}मे\appbrk{}हा\appbrk{}ति\appbrk{}सा\appbrk{}रा\appbrk{}दी\appbrk{}नाम्}\ A; \skt{ज्वर\appbrk{}शो\appbrk{}फ\appbrk{}गु\appbrk{}ल्म\appbrk{}र\appbrk{}क्त\appbrk{}पि\appbrk{}त्तो\appbrk{}न्मा\appbrk{}दा\appbrk{}प\appbrk{}स्मा\appbrk{}र\appbrk{}प्र\appbrk{}मे\appbrk{}हा\appbrk{}ती\appbrk{}सा\appbrk{}रा\appbrk{}दी\appbrk{}नाञ्}\ N H}}}{\lemma{\skt{च}}\Afootnote{K N H; \skt{उप\appbrk{}श\appbrk{}म\appbrk{}ना\appbrk{}र्थम् (३)}\ A}}\skt{~।}\pend

\pstart \ledsidenote{\passref{1.1.8.4}}\skt{भूतविद्या नाम }\edtext{\skt{देव\-ग\-न्ध\-र्व\-य\-क्ष\-रा\-क्ष\-स\-पि\-तृ\-पि\-शा\-च\-वि\-ना\-य\-क\-ना\-ग\-ग्र\-हो\-प\-सृ\-ष्ट\-चे\-तसां}}{\lemma{\skt{देव\appbrk{}ग\appbrk{}न्ध\appbrk{}र्व\appbrk{}य\appbrk{}क्ष\appbrk{}रा\appbrk{}क्ष\appbrk{}स\appbrk{}पि\appbrk{}तृ\appbrk{}पि\appbrk{}शा\appbrk{}च\appbrk{}वि\appbrk{}ना\appbrk{}य\appbrk{}क\appbrk{}ना\appbrk{}ग\appbrk{}ग्र\appbrk{}हो\appbrk{}प\appbrk{}सृ\appbrk{}ष्ट\appbrk{}चे\appbrk{}तसां}}\Afootnote{\skt{देवा\appbrk{}सु\appbrk{}र\appbrk{}ग\appbrk{}न्ध\appbrk{}र्व\appbrk{}य\appbrk{}क्ष\appbrk{}र\appbrk{}क्षः\appbrk{}पि\appbrk{}तृ\appbrk{}पि\appbrk{}शा\appbrk{}च\appbrk{}ना\appbrk{}ग\appbrk{}ग्र\appbrk{}हा\appbrk{}द्यु\appbrk{}प\appbrk{}सृ\appbrk{}ष्ट\appbrk{}चे\appbrk{}तसां}\ A; \skt{देव\appbrk{}ग\appbrk{}न्ध\appbrk{}र्व…\appbrk{}क्ष\appbrk{}रा\appbrk{}क्ष\appbrk{}स\appbrk{}पि\appbrk{}तृ\appbrk{}पि\appbrk{}शा\appbrk{}च\appbrk{}वि\appbrk{}ना\appbrk{}य\appbrk{}क\appbrk{}ना\appbrk{}ग\appbrk{}ग्र\appbrk{}हो\appbrk{}प\appbrk{}सृ\appbrk{}ष्ट\appbrk{}चे\appbrk{}तसां}\ K; \skt{देव\appbrk{}दी\appbrk{}न\appbrk{}व\appbrk{}ग\appbrk{}न्ध\appbrk{}र्व\appbrk{}य\appbrk{}क्ष\appbrk{}रा\appbrk{}क्ष\appbrk{}स\appbrk{}पि\appbrk{}तृ\appbrk{}पि\appbrk{}शा\appbrk{}च\appbrk{}वि\appbrk{}ना\appbrk{}य\appbrk{}क\appbrk{}ना\appbrk{}ग\appbrk{}ग्र\appbrk{}हो\appbrk{}प\appbrk{}सृ\appbrk{}ष्ट\appbrk{}चे\appbrk{}तसां}\ N; \skt{देव\appbrk{}दा\appbrk{}न\appbrk{}व\appbrk{}ग\appbrk{}न्ध\appbrk{}र्व्व\appbrk{}य\appbrk{}क्ष\appbrk{}रा\appbrk{}क्ष\appbrk{}स\appbrk{}पि\appbrk{}तृ\appbrk{}पि\appbrk{}शा\appbrk{}च\appbrk{}वि\appbrk{}ना\appbrk{}य\appbrk{}क\appbrk{}ना\appbrk{}ग\appbrk{}ग्र\appbrk{}हो\appbrk{}प\appbrk{}सृ\appbrk{}ष्ट\appbrk{}चे\appbrk{}तसां}\ H}}\skt{ शान्ति\-क\-र्म\-ब\-लि\-ह\-र\-णा\-दि\-ग्र\-हो\-प\-श\-म\-नार्थं~।}\pend

\pstart \ledsidenote{\passref{1.1.8.5}}\skt{कौमा\-र\-भृ\-त्य\-न्नाम }\edtext{\skt{कुमा\-र\-भ\-र\-ण\-धा\-त्री\-क्षी\-र\-दो\-ष\-सं\-शो\-ध\-नार्थं}}{\lemma{\skt{कुमा\appbrk{}र\appbrk{}भ\appbrk{}र\appbrk{}ण\appbrk{}धा\appbrk{}त्री\appbrk{}क्षी\appbrk{}र\appbrk{}दो\appbrk{}ष\appbrk{}सं\appbrk{}शो\appbrk{}ध\appbrk{}नार्थं}}\Afootnote{A K N; \skt{कुमा\appbrk{}रा\appbrk{}भ\appbrk{}र\appbrk{}ण\appbrk{}धा\appbrk{}त्री\appbrk{}क्षी\appbrk{}र\appbrk{}दो\appbrk{}ष\appbrk{}सं\appbrk{}शो\appbrk{}ध\appbrk{}नार्थं}\ H}}\skt{ }\edtext{\skt{दुष्ट\-स्त\-न्य\-ग्र\-ह\-स\-मु\-त्थि\-ता\-नाञ्च}}{\lemma{\skt{दुष्ट\appbrk{}स्त\appbrk{}न्य\appbrk{}ग्र\appbrk{}ह\appbrk{}स\appbrk{}मु\appbrk{}त्थि\appbrk{}ता\appbrk{}नाञ्च}}\Afootnote{K N H; \skt{दुष्ट\appbrk{}स्त\appbrk{}न्य\appbrk{}ग्र\appbrk{}ह\appbrk{}स\appbrk{}मु\appbrk{}त्थानां च}\ A}}\skt{ }\edtext{\skt{व्याधी\-ना\-मु\-प\-श\-म\-क\-र\-णा\-र्थम्}}{\lemma{\skt{उप\appbrk{}श\appbrk{}म\appbrk{}क\appbrk{}र\appbrk{}णा\appbrk{}र्थम्}}\Afootnote{K H; \skt{उप\appbrk{}श\appbrk{}म\appbrk{}ना\appbrk{}र्थम् (५)}\ A; \skt{उप\appbrk{}श\appbrk{}म\appbrk{}क\appbrk{}र\appbrk{}णार्थः}\ N}}\skt{~।}\pend

\pstart \ledsidenote{\passref{1.1.8.6}}\skt{अग\-द\-त\-न्त्र\-न्नाम }\edtext{\skt{सर्प\-की\-ट\-द\-ष्ट\-वि\-ष\-व्य\-ञ्ज\-नार्थं}}{\lemma{\skt{सर्प\appbrk{}की\appbrk{}ट\appbrk{}द\appbrk{}ष्ट\appbrk{}वि\appbrk{}ष\appbrk{}व्य\appbrk{}ञ्ज\appbrk{}नार्थं}}\Afootnote{K; \skt{सर्प\appbrk{}की\appbrk{}ट\appbrk{}लू\appbrk{}ता\appbrk{}मू\appbrk{}ष\appbrk{}का\appbrk{}दि\appbrk{}द\appbrk{}ष्ट\appbrk{}वि\appbrk{}ष\appbrk{}व्य\appbrk{}ञ्ज\appbrk{}नार्थं}\ A; \skt{सर्प\appbrk{}की\appbrk{}ट\appbrk{}लू\appbrk{}ता\appbrk{}द\appbrk{}ष्ट\appbrk{}स\appbrk{}री\appbrk{}सृ\appbrk{}प\appbrk{}वि\appbrk{}ष\appbrk{}व्य\appbrk{}ञ्ज\appbrk{}नार्थ}\ N; \skt{सर्प्प\appbrk{}की\appbrk{}ट\appbrk{}लू\appbrk{}ता\appbrk{}द\appbrk{}ष्ट\appbrk{}स\appbrk{}री\appbrk{}सृ\appbrk{}प\appbrk{}वि\appbrk{}ष\appbrk{}व्य\appbrk{}ञ्ज\appbrk{}ना\appbrk{}र्थम्}\ H}}\skt{ }\edtext{\skt{विवि\-ध\-वि\-ष\-वे\-गो\-प\-श\-म\-ना\-र्थञ्च}}{\lemma{\skt{विवि\appbrk{}ध\appbrk{}वि\appbrk{}ष\appbrk{}वे\appbrk{}गो\appbrk{}प\appbrk{}श\appbrk{}म\appbrk{}ना\appbrk{}र्थञ्}}\Afootnote{N H; \skt{विवि\appbrk{}ध\appbrk{}वि\appbrk{}ष\appbrk{}सं\appbrk{}यो\appbrk{}गो\appbrk{}प\appbrk{}श\appbrk{}म\appbrk{}नार्थं}\ A; \skt{विवि\appbrk{}ध\appbrk{}वि\appbrk{}ष\appbrk{}वेगो… शमनार्थञ्}\ K}}\skt{~।}\pend

\pstart \ledsidenote{\passref{1.1.8.7}}\skt{रसा\-य\-न\-त\-न्त्र\-न्नाम }\edtext{\edtext{\skt{वयः\-स्था\-प\-न\-मा\-यु\-र्मे\-धा\-क\-रणं}}{\lemma{\skt{वयःस्थापनम्}}\Afootnote{H; \skt{व}\ A; \skt{व…}\ K; \skt{वयस्थापनम्}\ N}}}{\lemma{\skt{आयु\appbrk{}र्मे\appbrk{}धा\appbrk{}क\appbrk{}रणं}}\Afootnote{N H; \skt{यःस्था\appbrk{}प\appbrk{}न\appbrk{}मा\appbrk{}यु\appbrk{}र्मे\appbrk{}धा\appbrk{}ब\appbrk{}ल\appbrk{}करं}\ A; \skt{…}\ K}}\skt{ }\edtext{\skt{व्याध्यु\-प\-श\-म\-क\-र\-णा\-र्थञ्च}}{\lemma{\skt{व्याध्यु\appbrk{}प\appbrk{}श\appbrk{}म\appbrk{}क\appbrk{}र\appbrk{}णा\appbrk{}र्थञ्}}\Afootnote{N H; \skt{रोगा\appbrk{}प\appbrk{}ह\appbrk{}र\appbrk{}ण\appbrk{}स\appbrk{}मर्थं}\ A; \skt{…प\appbrk{}श\appbrk{}म\appbrk{}क\appbrk{}र\appbrk{}णा\appbrk{}र्थञ्}\ K}}\skt{~।}\pend

\pstart \ledsidenote{\passref{1.1.8.8}}\edtext{\edtext{\skt{वाजी\-क\-र\-ण\-त\-न्त्र\-न्नाम}}{\lemma{\skt{वाजी\appbrk{}क\appbrk{}र\appbrk{}ण\appbrk{}त\appbrk{}न्त्रन्}}\Afootnote{A K\textsuperscript{ac} N H; \skt{व्याजी\appbrk{}क\appbrk{}र\appbrk{}ण\appbrk{}त\appbrk{}न्त्रन्}\ K\textsuperscript{pc}}}}{\lemma{\skt{नाम}}\Afootnote{K N H; \skt{नामा}\ A}}\skt{ }\edtext{\skt{स्वल्प\-दु\-ष्ट\-क्षी\-ण\-वि\-शु\-ष्क\-रे\-तसां}}{\lemma{\skt{स्वल्प\appbrk{}दु\appbrk{}ष्ट\appbrk{}क्षी\appbrk{}ण\appbrk{}वि\appbrk{}शु\appbrk{}ष्क\appbrk{}रे\appbrk{}तसां}}\Afootnote{K N H; \skt{ल्पदु\appbrk{}ष्ट\appbrk{}क्षी\appbrk{}ण\appbrk{}वि\appbrk{}शु\appbrk{}ष्क\appbrk{}रे\appbrk{}त\appbrk{}साम्}\ A}}\skt{ }\edtext{\skt{शुक्रा\-प्या\-य\-न\-प्र\-सा\-दो\-प\-च\-य\-ज\-न\-न\-नि\-मित्तं}}{\lemma{\skt{शुक्रा\appbrk{}प्या\appbrk{}य\appbrk{}न\appbrk{}प्र\appbrk{}सा\appbrk{}दो\appbrk{}प\appbrk{}च\appbrk{}य\appbrk{}ज\appbrk{}न\appbrk{}न\appbrk{}नि\appbrk{}मित्तं}}\Afootnote{K N; \skt{आप्या\appbrk{}य\appbrk{}न\appbrk{}प्र\appbrk{}सा\appbrk{}दो\appbrk{}प\appbrk{}च\appbrk{}य\appbrk{}ज\appbrk{}न\appbrk{}न\appbrk{}नि\appbrk{}मित्तं}\ A; \skt{शुक्रा\appbrk{}प्या\appbrk{}य\appbrk{}त\appbrk{}न\appbrk{}प्र\appbrk{}सा\appbrk{}दो\appbrk{}प\appbrk{}च\appbrk{}य\appbrk{}ज\appbrk{}न\appbrk{}न\appbrk{}नि\appbrk{}मित्तं}\ H}}\skt{ प्रह\-र्ष\-ज\-न\-ना\-र्थञ्च~।}\pend

\pstart \ledsidenote{\passref{1.1.9}}\skt{एव\-म\-य\-मा\-यु\-र्वे\-दोऽ\-ष्टाङ्ग उपदिश्यते~। अत्र }\edtext{\skt{कस्मै किं}}{\lemma{\skt{कस्मै किं}}\Afootnote{A N H; \skt{कस्मै कि}\ K}}\skt{ }\edtext{\edtext{\skt{वर्ण्यतामिति}}{\lemma{\skt{वर्ण्यताम्}}\Afootnote{N H; \skt{उच्यताम्}\ A; \skt{…}\ K}}}{\lemma{\skt{इति}}\Afootnote{A N H; \skt{…}\ K}}\skt{~।}\pend

\pstart \ledsidenote{\passref{1.1.10}}\edtext{\skt{त }\edtext{\skt{ऊचुरस्माकं}}{\lemma{\skt{अस्माकं}}\Afootnote{A N H; \emph{om.}\ K}}}{\lemma{\skt{त ऊचुर्}}\Afootnote{A N H; \emph{om.}\ K}}\skt{ }\edtext{\skt{सर्वमेव }\edtext{\skt{शल्य\-ज्ञा\-न\-म\-ल\-ङ्कृ\-त्वो\-प\-दि\-शतु}}{\lemma{\skt{अल\appbrk{}ङ्कृ\appbrk{}त्वो\appbrk{}प\appbrk{}दि\appbrk{}शतु}}\Afootnote{N H; \skt{मूलं कृत्वोपदिशतु}\ A; \skt{उपदिशतु}\ K\textsuperscript{ac}; \skt{अलंकृ+ उपदिशतु}\ K\textsuperscript{pc}}}}{\lemma{\skt{सर्वमेव शल्यज्ञानम्}}\Afootnote{H; \skt{सर्वेषामेव शल्यज्ञानं}\ A; \skt{शल्यज्ञानम्}\ K\textsuperscript{ac} K\textsuperscript{pc}; \skt{सर्वेषामेव ननमेव शल्यज्ञानम्}\ N\textsuperscript{ac}; \skt{सर्वेषामेव न शल्यज्ञानम्}\ N\textsuperscript{pc}}}\skt{ भगवानिति~।}\footnoteA{We follow the reading of MSS H and N, pace Harimoto 2013, in spite of the fact that \emph{alaṃkṛ[tvā]} is a marginal addition in K. K is badly damaged at this point, so it is not completely clear what happend to the manuscript, although it does look as if the word was not included in the text and was added by a different hand. But H is usually a faithful witness to the same text as K, and in this case even N concurs. The use of \emph{alaṅkṛtvā} in the sense "make accessible" is slightly unusual and probably explains the later banalization to \emph{mūla}.}\pend

\pstart \ledsidenote{\passref{1.1.11}}\edtext{\skt{स }\edtext{\skt{उवा\-चै\-व\-म\-स्त्विति}}{\lemma{\skt{अस्त्विति}}\Afootnote{A N H; \skt{…}\ K}}}{\lemma{\skt{स उवाचैवम्}}\Afootnote{A H; \emph{om.}\ K; \skt{स उवाच एवंम्}\ N}}\skt{~।}\pend

\pstart \ledsidenote{\passref{1.1.12}}\edtext{\skt{त }\edtext{\edtext{\skt{ऊचुर्भूयोऽस्माकं}}{\lemma{\skt{भूयो}}\Afootnote{A N H; \skt{{\symfont ‡}}\ K}}}{\lemma{\skt{ऽस्माकं}}\Afootnote{H; \skt{ऽपि भगवन्तम्}\ A; \skt{…कं}\ K; \skt{ऽस्माक स}\ N}}}{\lemma{\skt{त ऊचुर्}}\Afootnote{A N H; \skt{…}\ K}}\skt{ }\edtext{\edtext{\skt{सर्वे\-षा\-मे\-वै\-क\-म\-तीनां}}{\lemma{\skt{सर्वेषाम्}}\Afootnote{K H; \skt{अस्माकम्}\ A; \skt{र्वेषांम्}\ N}}}{\lemma{\skt{एवैकमतीनां}}\Afootnote{K H; \skt{एकमतीनां}\ A; \skt{एवै\appbrk{}क\appbrk{}म\appbrk{}ती\appbrk{}ना\appbrk{}नाम्}\ N}}\skt{ मत\-म\-भि\-स\-मीक्ष्य सुश्रुतो भगवन्तं }\edtext{\skt{प्रक्ष्यति}}{\lemma{\skt{प्रक्ष्यति}}\Afootnote{A K; \skt{प्रत्य}\ N H}}\skt{~। }\edtext{\skt{अस्यो\-प\-दि\-श्य\-मानं}}{\lemma{\skt{अस्यो\appbrk{}प\appbrk{}दि\appbrk{}श्य\appbrk{}मानं}}\Afootnote{K; \skt{अस्मै चोपदिश्यमानं}\ A; \skt{क्षस्यो\appbrk{}प\appbrk{}दि\appbrk{}श्य\appbrk{}मानं}\ N H}}\skt{ वय\-म\-प्यु\-प\-धा\-र\-यि\-ष्यामः~।}\pend

\pstart \ledsidenote{\passref{1.1.13}}\skt{स उवा\-चै\-व\-म\-स्त्विति~।}\pend

\pstart \ledsidenote{\passref{1.1.14-15}}\edtext{\skt{इह}}{\lemma{\skt{इह}}\Afootnote{K N H; \skt{वत्स सुश्रुत! इह}\ A}}\skt{ }\edtext{\skt{खल्वा\-यु\-र्वे\-द\-प्र\-यो\-जनं}}{\lemma{\skt{खल्वा\appbrk{}यु\appbrk{}र्वे\appbrk{}द\appbrk{}प्र\appbrk{}यो\appbrk{}जनं}}\Afootnote{A K N; \skt{खल्वायुर्व्वेदे प्रयोजनं}\ H}}\skt{ }\edtext{\skt{व्याध्युपसृष्टस्य}}{\lemma{\skt{व्याध्युपसृष्टस्य}}\Afootnote{K N H; \skt{व्याध्युपसृष्टानां}\ A}}\skt{ व्याधिपरिमोक्षः }\edtext{\edtext{\edtext{\skt{स्वस्थ\-र\-क्ष\-ण\-ञ्चा\-यु\-र\-स्मि\-न्वि\-द\-न्त्य\-नेन}}{\lemma{\skt{स्वस्थरक्षणञ्}}\Afootnote{K N H; \skt{स्वस्थस्य रक्षणं}\ A}}}{\lemma{\skt{चायुर्}}\Afootnote{K N H; \skt{च~॥१४॥ ‘आयुर्}\ A}}}{\lemma{\skt{अस्मिन्विदन्त्य्}}\Afootnote{K\textsuperscript{ac} N H; \skt{अस्मिन्विद्यते}\ A; \skt{अस्मिन्विंदन्त्य्}\ K\textsuperscript{pc}}}\skt{ }\edtext{\edtext{\skt{वायुर्विद्यत}}{\lemma{\skt{वायुर्}}\Afootnote{A K; \skt{वायु}\ N; \skt{चा}\ H\textsuperscript{ac}; \skt{चायु}\ H\textsuperscript{pc}}}}{\lemma{\skt{विद्यत}}\Afootnote{N; \skt{विन्दन्ति’}\ A; \skt{विन्द्यत}\ K; \emph{om.}\ H\textsuperscript{ac}; \skt{र्व्विद्यत}\ H\textsuperscript{pc}}}\skt{ }\edtext{\edtext{\skt{इत्यायुर्वेदः}}{\lemma{\skt{इत्य्}}\Afootnote{A K N H\textsuperscript{pc}; \emph{om.}\ H\textsuperscript{ac}}}}{\lemma{\skt{आयुर्वेदः}}\Afootnote{A K N H\textsuperscript{pc}; \skt{युर्व्वेदस्}\ H\textsuperscript{ac}}}\skt{~।}\pend

\pstart \ledsidenote{\passref{1.1.16}}\edtext{\edtext{\skt{तस्या\-ङ्ग\-व\-र\-मा\-ग\-म\-प्र\-त्य\-क्षा\-नु\-मा\-नो\-प\-मा\-नै\-र\-वि\-रु\-द्ध\-मु\-च्य\-मा\-न\-मु\-प\-धा\-र\-य\-ध्वम्}}{\lemma{\skt{आग\appbrk{}म\appbrk{}प्र\appbrk{}त्य\appbrk{}क्षा\appbrk{}नु\appbrk{}मा\appbrk{}नो\appbrk{}प\appbrk{}मा\appbrk{}नैर्}}\Afootnote{K\textsuperscript{ac} N H; \skt{आद्यं प्रत्य\appbrk{}क्षा\appbrk{}ग\appbrk{}मा\appbrk{}नु\appbrk{}मा\appbrk{}नो\appbrk{}प\appbrk{}मा\appbrk{}नैर्}\ A; \skt{आग\appbrk{}म\appbrk{}प्र\appbrk{}त्य\appbrk{}क्षा\appbrk{}नु\appbrk{}ना\appbrk{}मा\appbrk{}नो\appbrk{}प\appbrk{}मा\appbrk{}नैर्}\ K\textsuperscript{pc}}}}{\lemma{\skt{उप\appbrk{}धा\appbrk{}र\appbrk{}य\appbrk{}ध्वम्}}\Afootnote{K N H; \skt{उपधारय}\ A}}\skt{~।}\pend

\pstart \ledsidenote{\passref{1.1.17}}\skt{एतद्ध्यङ्गं प्रथमं }\edtext{\skt{प्रधानं}}{\lemma{\skt{प्रधानं}}\Afootnote{K N H; \emph{om.}\ A}}\skt{ }\edtext{\skt{प्राग\-भि\-हि\-त\-त्वा\-द्व्र\-ण\-सं\-रो\-ह\-ण\-क\-र\-त्वा\-द्य\-ज्ञ\-शि\-रः\-प्र\-धा\-न\-स\-न्धा\-नाच्च}}{\lemma{\skt{प्राग\appbrk{}भि\appbrk{}हि\appbrk{}त\appbrk{}त्वा\appbrk{}द्व्र\appbrk{}ण\appbrk{}सं\appbrk{}रो\appbrk{}ह\appbrk{}ण\appbrk{}क\appbrk{}र\appbrk{}त्वा\appbrk{}द्य\appbrk{}ज्ञ\appbrk{}शि\appbrk{}रः\appbrk{}प्र\appbrk{}धा\appbrk{}न\appbrk{}स\appbrk{}न्धा\appbrk{}नाच्च}}\Afootnote{K H; \skt{प्राग\appbrk{}भि\appbrk{}घा\appbrk{}त\appbrk{}व्र\appbrk{}ण\appbrk{}सं\appbrk{}रो\appbrk{}हा\appbrk{}द्य\appbrk{}ज्ञ\appbrk{}शि\appbrk{}रः\appbrk{}स\appbrk{}न्धा\appbrk{}नाच्च}\ A; \skt{प्राग\appbrk{}भि\appbrk{}धा\appbrk{}न\appbrk{}त्वा\appbrk{}द्व्र\appbrk{}ण\appbrk{}सं\appbrk{}रो\appbrk{}ह\appbrk{}ण\appbrk{}क\appbrk{}र\appbrk{}त्वा\appbrk{}द्य\appbrk{}ज्ञ\appbrk{}शि\appbrk{}र\appbrk{}प्र\appbrk{}धा\appbrk{}न\appbrk{}स\appbrk{}न्धा\appbrk{}नाच्च}\ N}}\skt{~। श्रूयते हि }\edtext{\skt{यथा}}{\lemma{\skt{यथा}}\Afootnote{A K\textsuperscript{pc} N H; \emph{om.}\ K\textsuperscript{ac}}}\skt{ }\edtext{\skt{पुरा}}{\lemma{\skt{पुरा}}\Afootnote{K N H; \emph{om.}\ A}}\skt{ रुद्रेण }\edtext{\edtext{\edtext{\skt{शिर\-श्छि\-न्न\-म\-श्विभ्यां}}{\lemma{\skt{शिरश्}}\Afootnote{K; \skt{यज्ञस्य शिरश्छिन्नमिति, ततो देवा अश्वि\appbrk{}ना\appbrk{}व\appbrk{}भि\appbrk{}ग\appbrk{}म्योचुः- भगवन्तौ नः श्रेष्ठतमौ युवां भविष्यथः, भवद्भ्यां यज्ञस्य शिरः सन्धातव्यमिति~। तावू\appbrk{}च\appbrk{}तु\appbrk{}रे\appbrk{}व\appbrk{}म\appbrk{}स्त्विति~। अथ तयोरर्थे देवा इन्द्रं यज्ञभागेन प्रासादयन्}\ A; \skt{यज्ञस्य शिरश्छिन्नमश्वि}\ N; \skt{शिरश्छिन्नम्}\ H\textsuperscript{ac}; \skt{यज्ञ शिरश्छिन्नम्}\ H\textsuperscript{pc}}}}{\lemma{\skt{छिन्नम्}}\Afootnote{K; \skt{ताभ्यां}\ A; \emph{om.}\ N H}}}{\lemma{\skt{अश्विभ्यां}}\Afootnote{K H; \skt{यज्ञस्य}\ A; \skt{भ्यां}\ N}}\skt{ }\edtext{\skt{पुनः}}{\lemma{\skt{पुनः}}\Afootnote{K N H; \skt{शिरः}\ A}}\skt{ }\edtext{\edtext{\skt{सन्धितमिति}}{\lemma{\skt{सन्धितम्}}\Afootnote{K N H; \skt{संहितम्”}\ A}}}{\lemma{\skt{इति}}\Afootnote{A N; \skt{इत्य्}\ K H}}\skt{~।}\pend

\pstart \ledsidenote{\passref{1.1.18}}\edtext{\edtext{\skt{अष्टानामपि}}{\lemma{\skt{अष्टानाम्}}\Afootnote{K N H; \skt{अष्टा}\ A}}}{\lemma{\skt{अपि}}\Afootnote{K N H; \skt{स्वपि}\ A}}\skt{ }\edtext{\edtext{\skt{चायु\-र्वे\-द\-त\-न्त्रा\-णा\-मे\-त\-दे\-वा\-धि\-क\-माशु}}{\lemma{\skt{चायु\appbrk{}र्वे\appbrk{}द\appbrk{}त\appbrk{}न्त्रा\appbrk{}णाम्}}\Afootnote{K N H; \skt{चायुर्वेदतन्त्रे}\ A}}}{\lemma{\skt{एत\appbrk{}दे\appbrk{}वा\appbrk{}धि\appbrk{}कम्}}\Afootnote{K N H; \skt{ष्वेत\appbrk{}दे\appbrk{}वा\appbrk{}धि\appbrk{}क\appbrk{}म\appbrk{}भि\appbrk{}म\appbrk{}तम्}\ A}}\skt{ }\edtext{\skt{क्रिया\-क\-र\-णा\-द्य\-न्त्र\-श\-स्त्र\-क्षा\-रा\-ग्नि\-प्र\-णि\-धा\-ना\-त्स\-र्व\-त\-न्त्र\-सा\-मा\-न्याच्च}}{\lemma{\skt{क्रिया\appbrk{}क\appbrk{}र\appbrk{}णा\appbrk{}द्य\appbrk{}न्त्र\appbrk{}श\appbrk{}स्त्र\appbrk{}क्षा\appbrk{}रा\appbrk{}ग्नि\appbrk{}प्र\appbrk{}णि\appbrk{}धा\appbrk{}ना\appbrk{}त्स\appbrk{}र्व\appbrk{}त\appbrk{}न्त्र\appbrk{}सा\appbrk{}मा\appbrk{}न्याच्च}}\Afootnote{K N H; \skt{क्रिया\appbrk{}क\appbrk{}र\appbrk{}णा\appbrk{}द्य\appbrk{}न्त्र\appbrk{}श\appbrk{}स्त्र\appbrk{}क्षा\appbrk{}रा\appbrk{}ग्नि\appbrk{}प्र\appbrk{}णि\appbrk{}धा\appbrk{}ना\appbrk{}त्स\appbrk{}र्व\appbrk{}त\appbrk{}न्त्र\appbrk{}सा\appbrk{}मा\appbrk{}न्याच्च}\ A}}\skt{~।}\pend

\pstart \ledsidenote{\passref{1.1.19}}\edtext{\skt{तदिदं}}{\lemma{\skt{तदिदं}}\Afootnote{A K H; \skt{तदिद}\ N}}\skt{ }\edtext{\skt{शाश्वतं}}{\lemma{\skt{शाश्वतं}}\Afootnote{A K H; \skt{शाश्वत}\ N}}\skt{ पुण्यं }\edtext{\skt{स्वर्ग्यं}}{\lemma{\skt{स्वर्ग्यं}}\Afootnote{A N; \skt{स्वर्गय}\ K\textsuperscript{ac}; \skt{स्वर्ग}\ K\textsuperscript{pc}; \skt{स्वर्ग्यां यं}\ H}}\skt{ }\edtext{\skt{यशस्यमायुष्यं}}{\lemma{\skt{यशस्यम्}}\Afootnote{A K N H\textsuperscript{pc}; \skt{शस्यम्}\ H\textsuperscript{ac}}}\skt{ }\edtext{\skt{वृत्तिकरञ्च}}{\lemma{\skt{वृत्तिकरञ्च}}\Afootnote{K N H; \skt{वृत्तिकरं चेति}\ A}}\skt{~।}\pend

\pstart \ledsidenote{\passref{1.1.20}}\edtext{\skt{तद्ब्रह्मा}}{\lemma{\skt{तद्ब्रह्मा}}\Afootnote{K; \skt{ब्रह्मा}\ A; \skt{त ब्रह्मा}\ N; \skt{तद्ब्रह्मो}\ H}}\skt{ }\edtext{\skt{प्रोवाच}}{\lemma{\skt{प्रोवाच}}\Afootnote{A K N; \skt{वाच}\ H}}\skt{ }\edtext{\edtext{\skt{तत्प्र\-जा\-प\-ति\-र\-धि\-जगे}}{\lemma{\skt{तत्प्रजापतिर्}}\Afootnote{K N H; \skt{ततः प्रजापतिर्}\ A}}}{\lemma{\skt{अधिजगे}}\Afootnote{A K H; \skt{अधिजचे}\ N}}\skt{ तस्मा\-द\-श्वि\-ना\-व\-श्वि\-भ्या\-मिन्द्र }\edtext{\skt{इन्द्रादहं}}{\lemma{\skt{इन्द्रादहं}}\Afootnote{A N; \skt{इन्द्रादःअं}\ K H}}\skt{ मया त्विह प्रदेयमर्थिभ्यः प्रजाहितहेतोः~।}\pend

\pstart\versepar \ledsidenote{\passref{1.1.21}}\edtext{\skt{भवति चात्र}}{\lemma{\skt{भवति चात्र}}\Afootnote{A; \skt{भ वति चात्र}\ K; \skt{भवति चात्र}\ N H}}\skt{~।}\footnoteA{The Nepalese witnesses commonly use the abbreviation "\emph{bha}" for "\emph{bhavati cātra}", introducing the next verse.}\skt{ }\pend
\pstart\versepar \skt{अहं हि धन्व\-न्त\-रि\-रा\-दि\-देवो जरा\-रु\-जा\-मृ\-त्यु\-हरो }\edtext{\skt{मराणां}}{\lemma{\skt{मराणां}}\Afootnote{K N; \skt{ऽमराणाम्}\ A H}}\skt{~। }\pend
\pstart\versepar \skt{}\edtext{\edtext{\skt{शल्य\-म्म\-ह\-च्छा\-स्त्र\-वरं}}{\lemma{\skt{शल्यम्}}\Afootnote{N H; \skt{शल्याङ्}\ A; \skt{…ल्यम्}\ K}}}{\lemma{\skt{महच्छास्त्रवरं}}\Afootnote{K H; \skt{गम\appbrk{}ङ्गै\appbrk{}र\appbrk{}प\appbrk{}रैर्}\ A; \skt{महच्छस्त्रवरं}\ N}}\skt{ }\edtext{\skt{गृहीत्वा}}{\lemma{\skt{गृहीत्वा}}\Afootnote{K N H; \skt{उपेतं}\ A}}\skt{ }\edtext{\skt{प्राप्तोऽस्मि गां}}{\lemma{\skt{ऽस्मि गां}}\Afootnote{A K; \skt{मि गां}\ N; \skt{ऽस्मि गम्}\ H}}\skt{ भूय इहोपदेष्टुं~।}\pend

\pstart \ledsidenote{\passref{1.1.22}}\edtext{\skt{तत्रास्मिञ्शास्त्रे}}{\lemma{\skt{तत्रास्मिञ्शास्त्रे}}\Afootnote{H; \skt{अस्मिञ्छास्त्रे}\ A; \skt{तत्रास्मिन्शास्त्रे}\ K; \skt{तत्रास्मिञ्छास्त्रे}\ N}}\skt{ पञ्च\-म\-हा\-भू\-त\-श\-री\-रि\-स\-म\-वायः पुरुष इत्युच्यते~। }\edtext{\skt{तस्मिन्क्रिया }\edtext{\skt{सोऽधिष्ठानं}}{\lemma{\skt{ऽधिष्ठानं}}\Afootnote{A K H; \skt{धिष्ठानं}\ N}}}{\lemma{\skt{तस्मिन्क्रिया सो}}\Afootnote{A K H; \skt{तस्मिं क्रिया सो}\ N}}\skt{~। }\edtext{\skt{कस्माल्~। लोकद्वैविध्याल्~। लोको हि}}{\lemma{\skt{कस्माल्~। लोकद्वैविध्याल्~। लोको हि}}\Afootnote{H; \skt{कस्मात्? लोकस्य द्वैविध्यात ; लोको हि}\ A; \skt{कस्मा\appbrk{}त्लो\appbrk{}क\appbrk{}द्वै\appbrk{}वि\appbrk{}ध्या\appbrk{}ल्लोको हि}\ K; \skt{कस्मा\appbrk{}ल्लौ\appbrk{}कि\appbrk{}क\appbrk{}वि\appbrk{}ध्या\appbrk{}ल्लोको हि}\ N}}\skt{ }\edtext{\skt{द्विविधः}}{\lemma{\skt{द्विविधः}}\Afootnote{A; \skt{द्विविध}\ K N H}}\footnoteA{The reading of \emph{dvividhan} might be a corruption of \emph{dvividhas}, which is a permitted sandhi.}\skt{ }\edtext{\skt{स्थावरो}}{\lemma{\skt{स्थावरो}}\Afootnote{A; \skt{न्स्थावरो}\ K H; \skt{स्थावर}\ N}}\skt{ }\edtext{\skt{जङ्गमश्च}}{\lemma{\skt{जङ्गमश्च}}\Afootnote{K N H; \skt{जङ्गमश्च}\ A}}\skt{~। द्विविधात्मक एवाग्नेयः }\edtext{\skt{सौम्यश्च}}{\lemma{\skt{सौम्यश्च}}\Afootnote{K N H; \skt{सौम्यश्च}\ A}}\skt{ }\edtext{\skt{तद्भू\-य\-स्त्वा\-त्प\-ञ्चा\-त्मको वा}}{\lemma{\skt{तद्भू\appbrk{}य\appbrk{}स्त्वा\appbrk{}त्प\appbrk{}ञ्चा\appbrk{}त्मको वा}}\Afootnote{A K H; \skt{तद्भू\appbrk{}य\appbrk{}स्त्वा\appbrk{}त्प\appbrk{}ञ्चा\appbrk{}तको वा}\ N}}\skt{~। तत्र चतुर्विधो भूतग्रामः }\edtext{\skt{संस्वे\-द\-जा\-द्रि\-ज\-ज\-रा\-यु\-जा\-ण्ड\-ज\-संज्ञः}}{\lemma{\skt{संस्वे\appbrk{}द\appbrk{}जा\appbrk{}द्रि\appbrk{}ज\appbrk{}ज\appbrk{}रा\appbrk{}यु\appbrk{}जा\appbrk{}ण्ड\appbrk{}ज\appbrk{}संज्ञः}}\Afootnote{\skt{संस्वे\appbrk{}द\appbrk{}ज\appbrk{}ज\appbrk{}रा\appbrk{}यु\appbrk{}जा\appbrk{}ण्ड\appbrk{}जो\appbrk{}द्भि\appbrk{}ज्ज\appbrk{}सञ्ज्ञः}\ A; \skt{संस्वे\appbrk{}द\appbrk{}जा\appbrk{}द्रि\appbrk{}ज\appbrk{}ज\appbrk{}रा\appbrk{}यु\appbrk{}जा\appbrk{}ण्ड\appbrk{}ज\appbrk{}सं\appbrk{}ज्ञास्}\ K; \skt{संस्वे\appbrk{}द\appbrk{}जा\appbrk{}द्रि\appbrk{}द\appbrk{}ज\appbrk{}ज\appbrk{}रा\appbrk{}यु\appbrk{}जा\appbrk{}ण्ड\appbrk{}ज\appbrk{}सं…स्}\ N; \skt{संस्वेदजो ह्रिद\appbrk{}ज\appbrk{}रा\appbrk{}यु\appbrk{}जा\appbrk{}ण्ड\appbrk{}ज\appbrk{}सं\appbrk{}ज्ञास्}\ H}}\skt{~। }\edtext{\skt{तस्मिन्पुरुषः}}{\lemma{\skt{तस्मिन्पुरुषः}}\Afootnote{H; \skt{तत्र पुरुषः}\ A; \skt{तस्मिस्पुरुषः}\ K; \skt{तस्मिं पुरुषः}\ N}}\skt{ }\edtext{\edtext{\edtext{\skt{प्रधा\-न\-स्त\-स्यो\-प\-क\-र\-ण\-म\-न्यत्}}{\lemma{\skt{प्रधानस्}}\Afootnote{K N H; \skt{प्रधानं}\ A}}}{\lemma{\skt{तस्यो\appbrk{}प\appbrk{}क\appbrk{}र\appbrk{}णम्}}\Afootnote{K N H; \skt{तस्योपकरण}\ A}}\skt{~। }\edtext{\skt{तस्मा\-त्पु\-रु\-षोऽ\-धि\-ष्ठानं}}{\lemma{\skt{ऽधिष्ठानं}}\Afootnote{A H; \skt{ऽधिष्ठान}\ K; \skt{धिष्ठानं}\ N}}}{\lemma{\skt{अन्यत्~। तस्मात्पुरुषो}}\Afootnote{K H; \skt{मन्यत्; तस्मात्पुरुषो}\ A; \skt{अन्यं तस्मात्पुरुषो}\ N}}\skt{~।}\pend

\pstart \ledsidenote{\passref{1.1.23}}\edtext{\skt{तद्दुःखसंयोगा}}{\lemma{\skt{तद्दुःखसंयोगा}}\Afootnote{A; \skt{तद्दुःखसंयोगा}\ K N H}}\skt{ व्याधय }\edtext{\skt{इत्युच्यन्ते}}{\lemma{\skt{इत्य्}}\Afootnote{K N H; \emph{om.}\ A}}\skt{~।}\pend

\pstart \ledsidenote{\passref{1.1.24}}\skt{ते चतुर्विधा आगन्तवः शारीरा }\edtext{\skt{मानसा}}{\lemma{\skt{मानसा}}\Afootnote{K N H; \skt{मानसाः}\ A}}\skt{ }\edtext{\edtext{\skt{स्वाभाविकाश्चेति}}{\lemma{\skt{स्वाभाविकाश्}}\Afootnote{K N H; \skt{स्वाभाविका}\ A}}}{\lemma{\skt{चेति}}\Afootnote{K N H; \skt{श्चेति}\ A}}\skt{~।}\pend

\pstart \ledsidenote{\passref{1.1.25.1}}\edtext{\edtext{\skt{तेष्वा\-ग\-न्त\-वोऽ\-भि\-घा\-त\-नि\-मित्ताः}}{\lemma{\skt{तेष्व्}}\Afootnote{K N H; \skt{तेषां}\ A}}}{\lemma{\skt{ऽभि\appbrk{}घा\appbrk{}त\appbrk{}नि\appbrk{}मित्ताः}}\Afootnote{A H; \skt{ऽभि\appbrk{}घा\appbrk{}त\appbrk{}नि\appbrk{}मित्ता}\ K; \skt{भिघातनिमित्ता}\ N}}\skt{~॥}\pend

\pstart \ledsidenote{\passref{1.1.25.2}}\edtext{\skt{शारी\-रा\-स्त्व\-न्न\-मूला}}{\lemma{\skt{त्वन्नमूला}}\Afootnote{K N H; \skt{त्वन्नपानमूला}\ A}}\skt{ }\edtext{\skt{वात\-पि\-त्त\-क\-फ\-शो\-णि\-त\-वै\-ष\-म्य\-नि\-मित्ताः}}{\lemma{\skt{वात\appbrk{}पि\appbrk{}त्त\appbrk{}क\appbrk{}फ\appbrk{}शो\appbrk{}णि\appbrk{}त\appbrk{}वै\appbrk{}ष\appbrk{}म्य\appbrk{}नि\appbrk{}मित्ताः}}\Afootnote{K N H; \skt{वात\appbrk{}पि\appbrk{}त्त\appbrk{}क\appbrk{}फ\appbrk{}शो\appbrk{}णि\appbrk{}त\appbrk{}स\appbrk{}न्नि\appbrk{}पा\appbrk{}त\appbrk{}वै\appbrk{}ष\appbrk{}म्य\appbrk{}नि\appbrk{}मित्ताः}\ A}}\skt{~।}\pend

\pstart \ledsidenote{\passref{1.1.25.3}}\skt{मानसास्तु }\edtext{\skt{क्रोध\-शो\-क\-दै\-न्य\-ह\-र्ष\-का\-म\-वि\-षा\-दे\-र्ष्या\-सू\-या\-मा\-त्स\-र्य\-लो\-भा\-दय}}{\lemma{\skt{क्रोध\appbrk{}शो\appbrk{}क\appbrk{}दै\appbrk{}न्य\appbrk{}ह\appbrk{}र्ष\appbrk{}का\appbrk{}म\appbrk{}वि\appbrk{}षा\appbrk{}दे\appbrk{}र्ष्या\appbrk{}सू\appbrk{}या\appbrk{}मा\appbrk{}त्स\appbrk{}र्य\appbrk{}लो\appbrk{}भा\appbrk{}दय}}\Afootnote{N H; \skt{क्रोध\appbrk{}शो\appbrk{}क\appbrk{}भ\appbrk{}य\appbrk{}ह\appbrk{}र्ष\appbrk{}वि\appbrk{}षा\appbrk{}दे\appbrk{}र्ष्या\appbrk{}भ्य\appbrk{}सू\appbrk{}या\appbrk{}दै\appbrk{}न्य\appbrk{}मा\appbrk{}त्स\appbrk{}र्य\appbrk{}का\appbrk{}म\appbrk{}लो\appbrk{}भ\appbrk{}प्र\appbrk{}भृ\appbrk{}तय}\ A; \skt{क्रोधा\appbrk{}शो\appbrk{}क\appbrk{}दै\appbrk{}न्य\appbrk{}ह\appbrk{}र्ष\appbrk{}का\appbrk{}म\appbrk{}वि\appbrk{}षा\appbrk{}दे\appbrk{}र्ष्या\appbrk{}सू\appbrk{}या\appbrk{}मा\appbrk{}त्स\appbrk{}र्य\appbrk{}लो\appbrk{}भा\appbrk{}दय}\ K}}\skt{ }\edtext{\skt{इच्छा\-द्वे\-ष\-नि\-मित्ताः}}{\lemma{\skt{इच्छा\appbrk{}द्वे\appbrk{}ष\appbrk{}नि\appbrk{}मित्ताः}}\Afootnote{K N H; \skt{इच्छा\appbrk{}द्वे\appbrk{}ष\appbrk{}भे\appbrk{}दै\appbrk{}र्भ\appbrk{}वन्ति}\ A}}\skt{~।}\pend

\pstart \ledsidenote{\passref{1.1.25.4}}\skt{स्वाभाविकास्तु }\edtext{\skt{क्षुत्पि\-पा\-सा\-ज\-रा\-मृ\-त्यु\-नि\-द्रा\-प्र\-कृ\-तयः}}{\lemma{\skt{क्षुत्पि\appbrk{}पा\appbrk{}सा\appbrk{}ज\appbrk{}रा\appbrk{}मृ\appbrk{}त्यु\appbrk{}नि\appbrk{}द्रा\appbrk{}प्र\appbrk{}कृ\appbrk{}तयः}}\Afootnote{K H; \skt{क्षुत्पि\appbrk{}पा\appbrk{}सा\appbrk{}ज\appbrk{}रा\appbrk{}मृ\appbrk{}त्यु\appbrk{}नि\appbrk{}द्रा\appbrk{}प्र\appbrk{}कृ(\appbrk{}भृ)\appbrk{}तयः}\ A; \skt{क्षुत्पि\appbrk{}पा\appbrk{}सा\appbrk{}ज\appbrk{}रा\appbrk{}मृ\appbrk{}त्यु\appbrk{}नि\appbrk{}द्रा\appbrk{}प्र\appbrk{}भृ\appbrk{}तयः}\ N}}\skt{~।}\pend

\pstart \ledsidenote{\passref{1.1.26}}\edtext{\skt{त एते}}{\lemma{\skt{त एते}}\Afootnote{A K H; \skt{एते}\ N}}\skt{ }\edtext{\skt{मनः\-श\-री\-रा\-धि\-ष्ठाना}}{\lemma{\skt{मनः\appbrk{}श\appbrk{}री\appbrk{}रा\appbrk{}धि\appbrk{}ष्ठाना}}\Afootnote{N H; \skt{मनः\appbrk{}श\appbrk{}री\appbrk{}रा\appbrk{}धि\appbrk{}ष्ठानाः}\ A; \skt{मनःशरीराधि…}\ K}}\skt{ }\edtext{\skt{भवन्ति}}{\lemma{\skt{भवन्ति}}\Afootnote{K N H; \emph{om.}\ A}}\skt{~।}\pend

\pstart \ledsidenote{\passref{1.1.27}}\skt{तेषां }\edtext{\skt{लेख\-न\-बृ\-म्ह\-ण\-सं\-शो\-ध\-न\-सं\-श\-म\-ना\-हा\-रा\-चाराः}}{\lemma{\skt{लेख\appbrk{}न\appbrk{}बृ\appbrk{}म्ह\appbrk{}ण\appbrk{}सं\appbrk{}शो\appbrk{}ध\appbrk{}न\appbrk{}सं\appbrk{}श\appbrk{}म\appbrk{}ना\appbrk{}हा\appbrk{}रा\appbrk{}चाराः}}\Afootnote{K N; \skt{संशो\appbrk{}ध\appbrk{}न\appbrk{}सं\appbrk{}श\appbrk{}म\appbrk{}ना\appbrk{}हा\appbrk{}रा\appbrk{}चाराः}\ A; \skt{लेख\appbrk{}न\appbrk{}बृ\appbrk{}म्ह\appbrk{}न\appbrk{}सं\appbrk{}शो\appbrk{}ध\appbrk{}न\appbrk{}सं\appbrk{}श\appbrk{}म\appbrk{}ना\appbrk{}हा\appbrk{}रा\appbrk{}चाराः}\ H}}\skt{ सम्यक्प्रयुक्ता निग्रहहेतवो }\edtext{\skt{भवन्ति}}{\lemma{\skt{भवन्ति}}\Afootnote{K N H; \emph{om.}\ A}}\skt{~।}\pend

\pstart \ledsidenote{\passref{1.1.28}}\skt{प्राणिनां }\edtext{\skt{पुन\-र्मू\-ल\-मा\-हारो}}{\lemma{\skt{पुनर्}}\Afootnote{A K H; \skt{पुन}\ N}}\skt{ बल\-व\-र्णौ\-ज\-साञ्च~। स षट्सु }\edtext{\skt{रसेष्वायत्तः~। रसाः}}{\lemma{\skt{आयत्तः~। रसाः}}\Afootnote{A; \skt{आयत्तरसाः}\ K N H}}\skt{ }\edtext{\skt{पुन\-र्द्र\-व्या\-श्र\-यिनः}}{\lemma{\skt{द्रव्याश्रयिनः}}\Afootnote{K; \skt{द्रव्याश्रयाः}\ A; \skt{दव्याश्रयिना}\ N; \skt{द्द्रव्याश्रयिणः}\ H}}\skt{~। द्रव्याणि }\edtext{\skt{पुनरोषध्यः~। ता}}{\lemma{\skt{ओषध्यः~। ता}}\Afootnote{K N; \skt{ओषधयः~। तास्तु}\ A; \skt{औषध्यः~। ता}\ H}}\skt{ }\edtext{\skt{द्विविधा}}{\lemma{\skt{द्विविधा}}\Afootnote{K N; \skt{द्विविधाः}\ A H}}\skt{~। स्थावरा जङ्गमाश्च~।}\pend

\pstart \ledsidenote{\passref{1.1.29}}\skt{तासां स्थाव\-रा\-श्च\-तु\-र्विधा वनस्पतयो वृक्षा }\edtext{\skt{ओषध्यो}}{\lemma{\skt{ओषध्यो}}\Afootnote{K H; \emph{om.}\ A; \skt{औषध्यो}\ N}}\skt{ }\edtext{\skt{वीरुध}}{\lemma{\skt{वीरुध}}\Afootnote{K N H; \skt{वीरुध, ओषधय}\ A}}\skt{ इति~। }\edtext{\skt{तास्वपुष्पा}}{\lemma{\skt{अपुष्पा}}\Afootnote{K N H; \skt{अपुष्पाः}\ A}}\skt{ }\edtext{\skt{फलवन्त्यो}}{\lemma{\skt{फलवन्त्यो}}\Afootnote{K H; \skt{फलवन्तो}\ A; \skt{…लवन्त्यो}\ N}}\skt{ वनस्पतयः~। }\edtext{\skt{पुष्पफलो ये ता}}{\lemma{\skt{पुष्पफलो ये ता}}\Afootnote{K H; \skt{पुष्पफलवन्तो}\ A; \skt{पुष्पफलो य्ये ता}\ N}}\skt{ वृक्षाः~। }\edtext{\edtext{\skt{फल\-पा\-क\-नि\-ष्ठा\-स्त्वो\-षध्यः}}{\lemma{\skt{फल\appbrk{}पा\appbrk{}क\appbrk{}नि\appbrk{}ष्ठास्}}\Afootnote{K N H; \skt{प्रतानवत्यः स्तम्बिन्यश्च वीरुधः, फलपाकनिष्ठा}\ A}}}{\lemma{\skt{त्वोषध्यः}}\Afootnote{K H; \skt{ओष}\ A; \skt{त्वौषध्यः}\ N}}\skt{~। }\edtext{\skt{प्रतानवत्यो}}{\lemma{\skt{प्रतानवत्यो}}\Afootnote{H; \skt{ध}\ A; \skt{प्रतानव…}\ K; \skt{प्रतानवन्त्यो}\ N}}\skt{ }\edtext{\skt{वीरुध}}{\lemma{\skt{वीरुध}}\Afootnote{K H; \skt{य}\ A; \skt{विरुध}\ N}}\skt{ इति~।}\pend

\pstart \ledsidenote{\passref{1.1.30}}\edtext{\skt{जङ्गमाः}}{\lemma{\skt{जङ्गमाः}}\Afootnote{K H; \skt{जङ्गमा}\ A N}}\skt{ खल्वपि }\edtext{\skt{चतुर्विधाः}}{\lemma{\skt{चतुर्विधाः}}\Afootnote{A N; \skt{चतुविधा}\ K H}}\skt{~। }\edtext{\skt{जरा\-यु\-जा\-ण्ड\-ज\-सं\-स्वे\-द\-जो\-द्भिदा}}{\lemma{\skt{जरा\appbrk{}यु\appbrk{}जा\appbrk{}ण्ड\appbrk{}ज\appbrk{}सं\appbrk{}स्वे\appbrk{}द\appbrk{}जो\appbrk{}द्भिदा}}\Afootnote{H; \skt{जरा\appbrk{}यु\appbrk{}जा\appbrk{}ण्ड\appbrk{}ज\appbrk{}स्वे\appbrk{}द\appbrk{}जोद्भि}\ A; \skt{जरा\appbrk{}यु\appbrk{}जा…\appbrk{}ज\appbrk{}स…\appbrk{}स्वे\appbrk{}द\appbrk{}जो\appbrk{}द्भिदाः}\ K; \skt{जरा\appbrk{}यु\appbrk{}जा\appbrk{}ण्ड\appbrk{}ज\appbrk{}स\appbrk{}स्वे\appbrk{}द\appbrk{}जो\appbrk{}द्भिदाः}\ N}}\skt{ }\edtext{\skt{इति}}{\lemma{\skt{इति}}\Afootnote{K N H; \skt{ज्जाः}\ A}}\skt{~। }\edtext{\skt{तेषा\-म्प\-शु\-म\-नु\-ष्य\-व्या\-ला\-दयो}}{\lemma{\skt{तेषाम्}}\Afootnote{K N H; \skt{तत्र}\ A}}\skt{ जरायुजाः~। }\edtext{\edtext{\skt{खग\-स\-री\-सृ\-प\-स\-र्पा\-स्त्व\-ण्डजाः}}{\lemma{\skt{खग\appbrk{}स\appbrk{}री\appbrk{}सृ\appbrk{}प\appbrk{}स\appbrk{}र्पास्}}\Afootnote{N H; \skt{खग\appbrk{}स\appbrk{}र्प\appbrk{}स\appbrk{}री\appbrk{}सृ\appbrk{}प\appbrk{}प्र\appbrk{}भृ\appbrk{}तयो}\ A; \skt{…सृप स्}\ K\textsuperscript{ac}; \skt{…सृप सर्पा स्}\ K\textsuperscript{pc}}}}{\lemma{\skt{त्वण्डजाः}}\Afootnote{K N H; \skt{ऽण्डजाः}\ A}}\skt{~। }\edtext{\skt{कृमि\-कु\-न्थ\-पि\-पी\-लि\-का\-प्र\-भृ\-तयः}}{\lemma{\skt{कृमि\appbrk{}कु\appbrk{}न्थ\appbrk{}पि\appbrk{}पी\appbrk{}लि\appbrk{}का\appbrk{}प्र\appbrk{}भृ\appbrk{}तयः}}\Afootnote{K; \skt{कृमि\appbrk{}की\appbrk{}ट\appbrk{}पि\appbrk{}पी\appbrk{}लि\appbrk{}का\appbrk{}प्र\appbrk{}भृ\appbrk{}तयः}\ A; \skt{कृमि\appbrk{}कु\appbrk{}ट्ठ\appbrk{}पि\appbrk{}पी\appbrk{}लि\appbrk{}का\appbrk{}प्र\appbrk{}भृ\appbrk{}तयः}\ N; \skt{क्रिमि\appbrk{}कु\appbrk{}ष्ठ\appbrk{}पि\appbrk{}पी\appbrk{}लि\appbrk{}का\appbrk{}प्र\appbrk{}भृ\appbrk{}तयः}\ H}}\footnoteA{\emph{kuṭṭha} is Prākṛta for "skin disease, leprosy," but this sense does not fit this passage; the word here must refer to an insect or similar creature. In our MSS, \emph{kuṭṭha} appears to be a transmitted form, possibly Sanskritized to \emph{kuṣṭha} by the scribe of MS H. It is not obvious how it could be a scribal error for the Vulgate's "\emph{kīṭa}." We read with K.}\skt{ }\edtext{\skt{संस्वेदजाः}}{\lemma{\skt{संस्वेदजाः}}\Afootnote{K N H; \skt{स्वेदजाः}\ A}}\skt{~। }\edtext{\edtext{\skt{इन्द्र\-गो\-प\-क\-म\-ण्डू\-क\-प्र\-भृ\-त\-य\-श्चो\-द्भिदाः}}{\lemma{\skt{इन्द्र\appbrk{}गो\appbrk{}प\appbrk{}क\appbrk{}म\appbrk{}ण्डू\appbrk{}क\appbrk{}प्र\appbrk{}भृ\appbrk{}त\appbrk{}यश्}}\Afootnote{K N H; \skt{इन्द्र\appbrk{}गो\appbrk{}प\appbrk{}म\appbrk{}ण्डू\appbrk{}क\appbrk{}प्र\appbrk{}भृ\appbrk{}तय}\ A}}}{\lemma{\skt{चोद्भिदाः}}\Afootnote{K N H; \skt{उद्भिज्जाः}\ A}}\skt{~।}\pend

\pstart \ledsidenote{\passref{1.1.31}}\skt{तत्र }\edtext{\skt{स्थाव\-रे\-भ्य\-स्त्व\-क्प\-त्र\-पु\-ष्प\-फ\-ल\-मू\-ल\-क\-न्द\-क्षी\-र\-नि\-र्या\-स\-सा\-र\-स्ने\-ह\-स्व\-रसाः}}{\lemma{\skt{त्वक्प\appbrk{}त्र\appbrk{}पु\appbrk{}ष्प\appbrk{}फ\appbrk{}ल\appbrk{}मू\appbrk{}ल\appbrk{}क\appbrk{}न्द\appbrk{}क्षी\appbrk{}र\appbrk{}नि\appbrk{}र्या\appbrk{}स\appbrk{}सा\appbrk{}र\appbrk{}स्ने\appbrk{}ह\appbrk{}स्व\appbrk{}रसाः}}\Afootnote{K H; \skt{त्वक्प\appbrk{}त्र\appbrk{}पु\appbrk{}ष्प\appbrk{}फ\appbrk{}ल\appbrk{}मू\appbrk{}ल\appbrk{}क\appbrk{}न्द\appbrk{}नि\appbrk{}र्या\appbrk{}स\appbrk{}स्व\appbrk{}र\appbrk{}सा\appbrk{}दयः}\ A; \skt{त्वक्प\appbrk{}त्र\appbrk{}पु\appbrk{}ष्प\appbrk{}फ\appbrk{}ल\appbrk{}मू\appbrk{}ल\appbrk{}क\appbrk{}न्द\appbrk{}क्षी\appbrk{}र\appbrk{}नि\appbrk{}र्या\appbrk{}स\appbrk{}सा\appbrk{}र\appbrk{}स्ने\appbrk{}ह\appbrk{}स्व\appbrk{}रसा}\ N}}\skt{ प्रयोजनवन्तः~। }\edtext{\skt{जङ्ग\-मे\-भ्य\-श्च\-र्म\-रो\-म\-न\-ख\-रु\-धि\-रा\-दयः}}{\lemma{\skt{चर्म\appbrk{}रो\appbrk{}म\appbrk{}न\appbrk{}ख\appbrk{}रु\appbrk{}धि\appbrk{}रा\appbrk{}दयः}}\Afootnote{K N H; \skt{चर्म\appbrk{}न\appbrk{}ख\appbrk{}रो\appbrk{}म\appbrk{}रु\appbrk{}धि\appbrk{}रा\appbrk{}दयः}\ A}}\skt{~।}\pend

\pstart \ledsidenote{\passref{1.1.32}}\edtext{\skt{पार्थिवस्तु}}{\lemma{\skt{पार्थिवस्तु}}\Afootnote{K N H; \skt{पार्थिवाः}\ A}}\skt{ }\edtext{\skt{सुव\-र्ण\-र\-ज\-ता\-दयः}}{\lemma{\skt{सुव\appbrk{}र्ण\appbrk{}र\appbrk{}ज\appbrk{}ता\appbrk{}दयः}}\Afootnote{K N H; \skt{सुव\appbrk{}र्ण\appbrk{}र\appbrk{}ज\appbrk{}त\appbrk{}म\appbrk{}णि\appbrk{}मु\appbrk{}क्ता\appbrk{}म\appbrk{}नः\appbrk{}शि\appbrk{}ला\appbrk{}मृ\appbrk{}त्क\appbrk{}पा\appbrk{}ला\appbrk{}दयः}\ A}}\skt{~।}\pend

\pstart \ledsidenote{\passref{1.1.33}}\edtext{\skt{कालकृतास्तु}}{\lemma{\skt{कालकृतास्तु}}\Afootnote{K N H; \skt{कालकृताः}\ A}}\skt{ }\edtext{\skt{प्रवा\-त\-नि\-वा\-ता\-त\-प\-च्छा\-या\-त\-मो\-ज्यो\-त्स्ना\-शी\-तो\-ष्ण\-व\-र्षासु}}{\lemma{\skt{प्रवा\appbrk{}त\appbrk{}नि\appbrk{}वा\appbrk{}ता\appbrk{}त\appbrk{}प\appbrk{}च्छा\appbrk{}या\appbrk{}त\appbrk{}मो\appbrk{}ज्यो\appbrk{}त्स्ना\appbrk{}शी\appbrk{}तो\appbrk{}ष्ण\appbrk{}व\appbrk{}र्षासु}}\Afootnote{K N H; \skt{प्रवा\appbrk{}त\appbrk{}नि\appbrk{}वा\appbrk{}ता\appbrk{}त\appbrk{}प\appbrk{}च्छा\appbrk{}या\appbrk{}ज्यो\appbrk{}त्स्ना\appbrk{}त\appbrk{}मः\appbrk{}शी\appbrk{}तो\appbrk{}ष्ण\appbrk{}वर्षा}\ A}}\skt{ }\edtext{\skt{सम्प्लवाः}}{\lemma{\skt{सम्प्लवाः}}\Afootnote{K H; \emph{om.}\ A; \skt{ससम्प्लवाः}\ N}}\skt{~। }\edtext{\skt{कालविशेषास्तु}}{\lemma{\skt{कालविशेषास्तु}}\Afootnote{K N H; \emph{om.}\ A}}\skt{ }\edtext{\skt{निमे\-ष\-का\-ष्ठा\-क\-ला\-मु\-हू\-र्ता\-हो\-रा\-त्र\-प\-क्ष\-मा\-स\-र्त्व\-य\-न\-स\-म्व\-त्स\-र\-यु\-ग\-वि\-शेषाः}}{\lemma{\skt{निमे\appbrk{}ष\appbrk{}का\appbrk{}ष्ठा\appbrk{}क\appbrk{}ला\appbrk{}मु\appbrk{}हू\appbrk{}र्ता\appbrk{}हो\appbrk{}रा\appbrk{}त्र\appbrk{}प\appbrk{}क्ष\appbrk{}मा\appbrk{}स\appbrk{}र्त्व\appbrk{}य\appbrk{}न\appbrk{}स\appbrk{}म्व\appbrk{}त्स\appbrk{}र\appbrk{}यु\appbrk{}ग\appbrk{}वि\appbrk{}शेषाः}}\Afootnote{K N H; \skt{होरा\appbrk{}त्र\appbrk{}प\appbrk{}क्ष\appbrk{}मा\appbrk{}स\appbrk{}र्त्व\appbrk{}य\appbrk{}ना\appbrk{}दयः संव\appbrk{}त्स\appbrk{}र\appbrk{}वि\appbrk{}शेषाः}\ A}}\skt{~।}\pend

\pstart \ledsidenote{\passref{1.1.34}}\edtext{\skt{स्वभावत}}{\lemma{\skt{स्वभावत}}\Afootnote{K N H; \skt{त एते स्वभावत}\ A}}\skt{ एव दोषाणां }\edtext{\skt{सञ्च\-य\-प्र\-को\-पो\-प\-श\-म\-प्र\-ती\-का\-र\-हे\-तवो}}{\lemma{\skt{सञ्च\appbrk{}य\appbrk{}प्र\appbrk{}को\appbrk{}पो\appbrk{}प\appbrk{}श\appbrk{}म\appbrk{}प्र\appbrk{}ती\appbrk{}का\appbrk{}र\appbrk{}हे\appbrk{}तवो}}\Afootnote{K H; \skt{सञ्च\appbrk{}य\appbrk{}प्र\appbrk{}को\appbrk{}प\appbrk{}प्र\appbrk{}श\appbrk{}म\appbrk{}प्र\appbrk{}ती\appbrk{}का\appbrk{}र\appbrk{}हे\appbrk{}तवः}\ A; \skt{सञ्च\appbrk{}य\appbrk{}प्र\appbrk{}को\appbrk{}प\appbrk{}श\appbrk{}म\appbrk{}प्र\appbrk{}ती\appbrk{}का\appbrk{}र\appbrk{}हे\appbrk{}तवो}\ N}}\skt{ }\edtext{\skt{भवन्ति}}{\lemma{\skt{भवन्ति}}\Afootnote{K N H; \emph{om.}\ A}}\skt{~। प्रयो\-ज\-न\-व\-न्तश्च~।}\pend

\pstart\versepar \ledsidenote{\passref{1.1.35}}\skt{भवन्ति चात्र~। }\pend
\pstart\versepar \skt{शारीराणां विकाराणामेष }\edtext{\edtext{\skt{वर्गश्चतुर्विधः}}{\lemma{\skt{वर्गश्}}\Afootnote{A K H; \skt{वर्गेश्}\ N}}}{\lemma{\skt{चतुर्विधः}}\Afootnote{A N H; \skt{चतुविधः}\ K}}\skt{~। }\pend
\pstart\versepar \skt{प्रकोपे प्रशमे चैव }\edtext{\edtext{\skt{हेतु\-रु\-क्त\-श्चि\-कि\-त्सकैः}}{\lemma{\skt{उक्तश्}}\Afootnote{K H; \skt{उक्त}\ A N}}}{\lemma{\skt{चिकित्सकैः}}\Afootnote{K N H; \skt{श्चिकित्सकैः}\ A}}\skt{~॥}\edtext{}{\lemma{\skt{भवन्ति चात्र~। शारीराणां}}\Afootnote{\skt{भवन्ति चात्र श्लोकाः- शारीराणां}\ A; \skt{भवन्ति चात्र॥ शारीराणां}\ K; \skt{भ शारीराणां}\ N H}}\pend

\pstart\versepar \ledsidenote{\passref{1.1.36}}\skt{आगन्तवस्तु ये रोगास्ते द्विधा }\edtext{\skt{निपतन्ति ह}}{\lemma{\skt{निपतन्ति ह}}\Afootnote{K N H; \skt{निपतन्ति हि}\ A}}\skt{~। }\pend
\pstart\versepar \skt{मनस्यन्ये }\edtext{\skt{शरीरे च}}{\lemma{\skt{शरीरे च}}\Afootnote{K N H; \skt{शरीरेऽन्ये}\ A}}\skt{ तेषान्तु द्विविधा क्रिया~॥}\pend

\pstart\versepar \ledsidenote{\passref{1.1.37}}\skt{शरी\-र\-प\-ति\-तानां तु }\edtext{\skt{शारी\-र\-व\-दु\-प\-क्रमः}}{\lemma{\skt{शारी\appbrk{}र\appbrk{}व\appbrk{}दु\appbrk{}प\appbrk{}क्रमः}}\Afootnote{K N H; \skt{शारी\appbrk{}र\appbrk{}व\appbrk{}दु\appbrk{}प\appbrk{}क्रमः}\ A}}\skt{~। }\pend
\pstart\versepar \skt{मानसानान्तु }\edtext{\edtext{\skt{शब्दादिरिष्टो}}{\lemma{\skt{शब्दादिर्}}\Afootnote{K N H; \skt{शब्दादि}\ A}}}{\lemma{\skt{इष्टो}}\Afootnote{K N H; \skt{रिष्टो}\ A}}\skt{ वर्गः सुखावहः~॥}\footnoteA{Note the good example of the inconsistent use by the scribe of witness H of anusvāra in 37ab (\emph{°n tu}) vs 37cd (\emph{°ṃ tu}).}\pend

\pstart \ledsidenote{\passref{1.1.38}}\skt{एव\-मे\-त\-त्पु\-रुषो व्याधिरौषधं क्रिया काल इति }\edtext{\skt{समासेन}}{\lemma{\skt{समासेन}}\Afootnote{K N H; \emph{om.}\ A}}\skt{ चतुष्टयं }\edtext{\edtext{\skt{व्याख्यातम्भवति}}{\lemma{\skt{व्याख्यातम्}}\Afootnote{K N H; \skt{समासेन व्याख्यातम्}\ A}}}{\lemma{\skt{भवति}}\Afootnote{K N H; \emph{om.}\ A}}\skt{~। तत्र पुरु\-ष\-ग्र\-ह\-णा\-त्त\-त्सं\-भ\-व\-द्र\-व्य\-स\-मूहो भूतादिरुक्तः }\edtext{\skt{तद\-ङ्ग\-प्र\-त्य\-ङ्ग\-वि\-क\-ल्पाश्च}}{\lemma{\skt{तद\appbrk{}ङ्ग\appbrk{}प्र\appbrk{}त्य\appbrk{}ङ्ग\appbrk{}वि\appbrk{}क\appbrk{}ल्पाश्च}}\Afootnote{A K; \skt{तद\appbrk{}ग\appbrk{}द्ग\appbrk{}प्र\appbrk{}त्य\appbrk{}ङ्ग\appbrk{}वि\appbrk{}क\appbrk{}ल्पाश्च}\ N; \skt{तदङ्गाप्रत्यङ्गा विकल्पाश्च}\ H\textsuperscript{ac}; \skt{तदङ्गप्रत्यङ्गा विकल्पाश्च}\ H\textsuperscript{pc}}}\skt{ }\edtext{\skt{त्वङ्मां\-सा\-सि\-रा\-स्ना\-य्व\-स्थि\-स\-न्धि\-प्र\-भृ\-तयः}}{\lemma{\skt{त्वङ्मां\appbrk{}सा\appbrk{}सि\appbrk{}रा\appbrk{}स्ना\appbrk{}य्व\appbrk{}स्थि\appbrk{}स\appbrk{}न्धि\appbrk{}प्र\appbrk{}भृ\appbrk{}तयः}}\Afootnote{K N H; \skt{त्वङ्मां\appbrk{}सा\appbrk{}स्थि\appbrk{}सि\appbrk{}रा\appbrk{}स्ना\appbrk{}यु\appbrk{}प्र\appbrk{}भृ\appbrk{}तयः}\ A}}\skt{~। }\edtext{\skt{व्याधि\-ग्र\-ह\-णा\-द्वा\-त\-पि\-त्त\-क\-फ\-शो\-णि\-त\-स\-न्नि\-पा\-ता\-ग\-न्तु\-स्व\-भा\-व\-नि\-मित्ताः}}{\lemma{\skt{व्याधि\appbrk{}ग्र\appbrk{}ह\appbrk{}णा\appbrk{}द्वा\appbrk{}त\appbrk{}पि\appbrk{}त्त\appbrk{}क\appbrk{}फ\appbrk{}शो\appbrk{}णि\appbrk{}त\appbrk{}स\appbrk{}न्नि\appbrk{}पा\appbrk{}ता\appbrk{}ग\appbrk{}न्तु\appbrk{}स्व\appbrk{}भा\appbrk{}व\appbrk{}नि\appbrk{}मित्ताः}}\Afootnote{K H; \skt{व्याधि\appbrk{}ग्र\appbrk{}ह\appbrk{}णा\appbrk{}द्वा\appbrk{}त\appbrk{}पि\appbrk{}त्त\appbrk{}क\appbrk{}फ\appbrk{}शो\appbrk{}णि\appbrk{}त\appbrk{}स\appbrk{}न्नि\appbrk{}पा\appbrk{}त\appbrk{}वै\appbrk{}ष\appbrk{}म्य\appbrk{}नि\appbrk{}मित्ताः}\ A; \skt{व्याधि\appbrk{}ग्रा\appbrk{}ह\appbrk{}णा\appbrk{}द्वा\appbrk{}त\appbrk{}पि\appbrk{}त्त\appbrk{}क\appbrk{}फ\appbrk{}शो\appbrk{}णि\appbrk{}त\appbrk{}स\appbrk{}न्नि\appbrk{}पा\appbrk{}ता\appbrk{}ग\appbrk{}न्तु\appbrk{}स्व\appbrk{}भा\appbrk{}व\appbrk{}नि\appbrk{}मित्ताः}\ N}}\skt{ सर्व एव व्याधयो }\edtext{\skt{व्याख्याता}}{\lemma{\skt{व्याख्याता}}\Afootnote{K N H; \skt{व्याख्याताः}\ A}}\skt{ }\edtext{\skt{भवन्ति}}{\lemma{\skt{भवन्ति}}\Afootnote{K H; \emph{om.}\ A; \skt{भवस्वन्ति}\ N}}\skt{~। }\edtext{\edtext{\skt{ओष\-ध\-ग्र\-ह\-णा\-द्द्र\-व्य\-र\-स\-गु\-ण\-वी\-र्य\-वि\-पा\-का\-ना\-मा\-देशः}}{\lemma{\skt{ओष\appbrk{}ध\appbrk{}ग्र\appbrk{}ह\appbrk{}णा\appbrk{}द्द्र\appbrk{}व्य\appbrk{}र\appbrk{}स\appbrk{}गु\appbrk{}ण\appbrk{}वी\appbrk{}र्य\appbrk{}वि\appbrk{}पा\appbrk{}का\appbrk{}नाम्}}\Afootnote{K N; \skt{औष\appbrk{}ध\appbrk{}ग्र\appbrk{}ह\appbrk{}णा\appbrk{}द्द्र\appbrk{}व्य\appbrk{}र\appbrk{}स\appbrk{}गु\appbrk{}ण\appbrk{}वी\appbrk{}र्य\appbrk{}वि\appbrk{}पा\appbrk{}काना}\ A; \skt{ओष\appbrk{}ध\appbrk{}ग्र\appbrk{}ह\appbrk{}णा\appbrk{}द्द्र\appbrk{}व्य\appbrk{}र\appbrk{}स\appbrk{}गु\appbrk{}ण\appbrk{}वी\appbrk{}र्य\appbrk{}वि\appbrk{}पाको ना}\ H}}}{\lemma{\skt{आदेशः}}\Afootnote{K N; \skt{मादेशः}\ A H}}\skt{~। }\edtext{\skt{क्रिया\-ग्र\-ह\-णा\-त्स्ने\-हा\-दीनि}}{\lemma{\skt{क्रिया\appbrk{}ग्र\appbrk{}ह\appbrk{}णा\appbrk{}त्स्ने\appbrk{}हा\appbrk{}दीनि}}\Afootnote{A K N H\textsuperscript{pc}; \skt{क्रिया\appbrk{}ग्र\appbrk{}ह\appbrk{}णा\appbrk{}त्स्नेहा}\ H\textsuperscript{ac}}}\skt{ }\edtext{\skt{च्छेद्यादीनि च}}{\lemma{\skt{च्छेद्यादीनि च}}\Afootnote{A K N H\textsuperscript{pc}; \skt{दीनि च}\ H\textsuperscript{ac}}}\skt{ }\edtext{\skt{कर्मा\-ण्यु\-प\-दि\-ष्टानि}}{\lemma{\skt{उपदिष्टानि}}\Afootnote{K N H; \skt{व्या}\ A}}\skt{ }\edtext{\skt{भवन्ति}}{\lemma{\skt{भवन्ति}}\Afootnote{K N H; \skt{ख्यातानि}\ A}}\skt{~। काल\-ग्र\-ह\-णा\-त्सर्व }\edtext{\skt{एव}}{\lemma{\skt{एव}}\Afootnote{K N H; \emph{om.}\ A}}\skt{ }\edtext{\skt{क्रियाकालादेशः}}{\lemma{\skt{क्रियाकालादेशः}}\Afootnote{K; \skt{क्रिया\appbrk{}का\appbrk{}ला\appbrk{}ना\appbrk{}मा\appbrk{}देशः}\ A; \skt{क्रिया\appbrk{}का\appbrk{}ल\appbrk{}नि\appbrk{}र्देशः}\ N; \skt{क्रिया\appbrk{}का\appbrk{}ल\appbrk{}का\appbrk{}ल\appbrk{}देशः}\ H}}\skt{~।}\pend

\pstart\versepar \ledsidenote{\passref{1.1.39}}\edtext{\skt{भवति चात्र}}{\lemma{\skt{भवति चात्र}}\Afootnote{A; \emph{om.}\ K; \skt{भवति चात्र}\ N H}}\skt{ }\pend
\pstart\versepar \skt{बीज\-ञ्चि\-कि\-त्सि\-त\-स्यै\-त\-त्स\-मा\-सेन }\edtext{\skt{प्रकीतम्}}{\lemma{\skt{प्रकीतम्}}\Afootnote{\skt{प्रकीर्तितम्}\ A K N H}}\skt{~। }\pend
\pstart\versepar \skt{}\edtext{\edtext{\edtext{\skt{सविं\-श\-म\-ध्या\-य\-श\-त\-मस्य}}{\lemma{\skt{सविंशम्}}\Afootnote{K N H; \skt{सविंश}\ A}}}{\lemma{\skt{अध्यायशतम्}}\Afootnote{K N H; \skt{मध्यायशत}\ A}}}{\lemma{\skt{अस्य}}\Afootnote{K N H; \skt{मस्य}\ A}}\skt{ }\edtext{\skt{व्याख्या}}{\lemma{\skt{व्याख्या}}\Afootnote{A K\textsuperscript{pc} N H; \emph{om.}\ K\textsuperscript{ac}}}\skt{ भविष्यति~॥}\pend

\pstart \ledsidenote{\passref{1.1.40}}\edtext{\skt{तच्च}}{\lemma{\skt{तच्च}}\Afootnote{A H; \skt{त}\ K N}}\skt{ }\edtext{\skt{विंश\-म\-ध्या\-य\-शतं}}{\lemma{\skt{विंशम्}}\Afootnote{H; \skt{सविंशम्}\ A; \skt{च्च विंशम्}\ K N}}\skt{ पञ्चसु स्थानेषु }\edtext{\skt{चेति}}{\lemma{\skt{चेति}}\Afootnote{K N H; \emph{om.}\ A}}\skt{~। }\edtext{\skt{तत्र}}{\lemma{\skt{तत्र}}\Afootnote{A K N; \skt{अत्र}\ H}}\skt{ }\edtext{\edtext{\skt{श्लोक\-स्था\-न\-नि\-दा\-न\-शा\-री\-र\-चि\-कि\-त्सि\-त\-क\-ल्पे\-ष्व\-र्थ\-व\-शा\-द्वि\-भज्य}}{\lemma{\skt{श्लोक\appbrk{}स्था\appbrk{}न\appbrk{}नि\appbrk{}दा\appbrk{}न\appbrk{}शा\appbrk{}री\appbrk{}र\appbrk{}चि\appbrk{}कि\appbrk{}त्सि\appbrk{}त\appbrk{}क\appbrk{}ल्पेष्व्}}\Afootnote{N H; \skt{सूत्र\appbrk{}नि\appbrk{}दा\appbrk{}न\appbrk{}शा\appbrk{}री\appbrk{}र\appbrk{}चि\appbrk{}कि\appbrk{}त्सि\appbrk{}त\appbrk{}क\appbrk{}ल्पेष्व्}\ A; \skt{श्लोक\appbrk{}स्था\appbrk{}न\appbrk{}नि\appbrk{}दा\appbrk{}न\appbrk{}शारी}\ K}}\skt{ उत्तरे}}{\lemma{\skt{अर्थ\appbrk{}व\appbrk{}शा\appbrk{}द्वि\appbrk{}भज्य उत्तरे}}\Afootnote{\skt{अर्थ\appbrk{}व\appbrk{}शा\appbrk{}त्सं\appbrk{}वि\appbrk{}भ\appbrk{}ज्यो\appbrk{}त्तरे}\ A; \skt{…}\ K; \skt{अर्थवशा}\ N; \skt{अर्थ\appbrk{}व\appbrk{}सा\appbrk{}द्वि\appbrk{}भज्य उत्तरे}\ H}}\skt{ }\edtext{\skt{वक्ष्यामः}}{\lemma{\skt{वक्ष्यामः}}\Afootnote{H; \skt{तन्त्रे शेषा\appbrk{}न\appbrk{}र्था\appbrk{}न्व\appbrk{}क्ष्यामः}\ A; \skt{…इभज्य उत्तरे वक्ष्यामः}\ K; \skt{द्विभज्योत्तरे वक्ष्यामः}\ N}}\skt{~।}\pend

\pstart\versepar \ledsidenote{\passref{1.1.41}}\edtext{\skt{भवति चात्र}}{\lemma{\skt{भवति चात्र}}\Afootnote{A; \emph{om.}\ K; \skt{भवति चात्र}\ N H}}\skt{~। }\pend
\pstart\versepar \skt{स्वयम्भुवा प्रोक्तमिदं सनातनम् }\pend
\pstart\versepar \skt{}\edtext{\skt{पठेत्तु यः}}{\lemma{\skt{पठेत्तु यः}}\Afootnote{K N H; \skt{पठेद्धि यः}\ A}}\skt{ काशि\-प\-ति\-प्र\-का\-शि\-तम्~। }\pend
\pstart\versepar \skt{स पुण्यकर्मा भुवि पूजितो नृपैर् }\pend
\pstart\versepar \skt{असुक्षये }\edtext{\skt{शक्र\-स\-लो\-क\-ता\-मि\-या\-दिति}}{\lemma{\skt{इयादिति}}\Afootnote{\skt{व्रजेत्}\ A; \skt{इयाद्}\ K N H}}\skt{~॥}\pend

\pstart \ledsidenote{\passref{1.1.trailer}}\edtext{\skt{॥ १~॥}}{\lemma{\skt{॥ १~॥}}\Afootnote{\skt{इति सुश्रु\appbrk{}त\appbrk{}सं\appbrk{}हि\appbrk{}तायां सूत्रस्थाने वेदोत्पत्तिर्नाम प्रथमोऽध्यायः}\ A; \skt{इति~॥{\symfont ❈}}\ K; \skt{इति}\ N H}}\pend

\endnumbering
\beginnumbering
\eledsection*{Adhyāya 2}
\pstart \ledsidenote{\passref{1.2.1}}\skt{अथातः शिष्यो\-प\-न\-य\-नी\-य\-म\-ध्यायं }\edtext{\skt{व्याख्यास्यामः}}{\lemma{\skt{व्याख्यास्यामः}}\Afootnote{A N H; \emph{om.}\ K}}\skt{~॥}\pend

\pstart \ledsidenote{\passref{1.2.3}}\edtext{\edtext{\skt{ब्राह्म\-ण\-क्ष\-त्रि\-य\-वै\-श्या\-ना\-म\-न्य\-त\-म\-म\-न्व\-य\-व\-यः\-शौ\-र्य\-शौ\-चा\-चा\-र\-वि\-न\-य\-श\-क्ति\-ब\-ल\-मे\-धा\-धृ\-ति\-स्मृ\-ति\-प्र\-ति\-प\-त्ति\-युक्तं}}{\lemma{\skt{अन्यतमम्}}\Afootnote{K N H; \skt{अन्यतम}\ A}}}{\lemma{\skt{अन्व\appbrk{}य\appbrk{}व\appbrk{}यः\appbrk{}शौ\appbrk{}र्य\appbrk{}शौ\appbrk{}चा\appbrk{}चा\appbrk{}र\appbrk{}वि\appbrk{}न\appbrk{}य\appbrk{}श\appbrk{}क्ति\appbrk{}ब\appbrk{}ल\appbrk{}मे\appbrk{}धा\appbrk{}धृ\appbrk{}ति\appbrk{}स्मृ\appbrk{}ति\appbrk{}प्र\appbrk{}ति\appbrk{}प\appbrk{}त्ति\appbrk{}युक्तं}}\Afootnote{\skt{मन्व\appbrk{}य\appbrk{}व\appbrk{}यः\appbrk{}शी\appbrk{}ल\appbrk{}शौ\appbrk{}र्य\appbrk{}शौ\appbrk{}चा\appbrk{}चा\appbrk{}र\appbrk{}वि\appbrk{}न\appbrk{}य\appbrk{}श\appbrk{}क्ति\appbrk{}ब\appbrk{}ल\appbrk{}मे\appbrk{}धा\appbrk{}धृ\appbrk{}ति\appbrk{}स्मृ\appbrk{}ति\appbrk{}म\appbrk{}ति\appbrk{}प्र\appbrk{}ति\appbrk{}प\appbrk{}त्ति\appbrk{}युक्तं}\ A; \skt{अन्वयः~। वयः…\appbrk{}र\appbrk{}वि\appbrk{}न\appbrk{}य\appbrk{}श\appbrk{}क्ति\appbrk{}व\appbrk{}ल\appbrk{}मे\appbrk{}धा\appbrk{}धृ\appbrk{}ति\appbrk{}स्मृ\appbrk{}ति\appbrk{}प्र\appbrk{}ति\appbrk{}प\appbrk{}त्ति\appbrk{}युक्तं}\ K; \skt{अन्वयः। वयः\appbrk{}शौ\appbrk{}र्य\appbrk{}शौ\appbrk{}चा\appbrk{}चा\appbrk{}र\appbrk{}वि\appbrk{}न\appbrk{}य\appbrk{}श\appbrk{}क्ति\appbrk{}व\appbrk{}ल\appbrk{}मे\appbrk{}धा\appbrk{}धृ\appbrk{}ति\appbrk{}स्मृ\appbrk{}ति\appbrk{}प्र\appbrk{}ति\appbrk{}प\appbrk{}त्ति\appbrk{}यु\appbrk{}क्तन्}\ N; \skt{अन्वयः। वय\appbrk{}सौ\appbrk{}र्य\appbrk{}शौ\appbrk{}चा\appbrk{}चा\appbrk{}र\appbrk{}वि\appbrk{}न\appbrk{}य\appbrk{}श\appbrk{}क्ति\appbrk{}ब\appbrk{}ल\appbrk{}मे\appbrk{}धा\appbrk{}धृ\appbrk{}ति\appbrk{}स्मृ\appbrk{}ति\appbrk{}प्र\appbrk{}ति\appbrk{}प\appbrk{}त्ति\appbrk{}युक्तं}\ H}}\footnoteA{Emending \emph{anvaya} against the Nepalese MSS is required to make syntactic sense of the passage. Dr. P. Maas pointed out the parallel with the mention of \emph{kula} {\symfont `}{\symfont `}family'' as one of the qualifiers for a student, in \emph{Carakasaṃhitā} 3.8.8.}\skt{ }\edtext{\edtext{\skt{तनु\-जि\-ह्वौ\-ष्ठ\-द\-न्ता\-ग्र\-मृ\-जु\-व\-क्त्रा\-क्षि\-नासं}}{\lemma{\skt{तनु\appbrk{}जि\appbrk{}ह्वौ\appbrk{}ष्ठ\appbrk{}द\appbrk{}न्ता\appbrk{}ग्रम्}}\Afootnote{K N H; \skt{तनु\appbrk{}जि\appbrk{}ह्वौ\appbrk{}ष्ठ\appbrk{}द\appbrk{}न्ताग्र}\ A}}}{\lemma{\skt{ऋजु\appbrk{}व\appbrk{}क्त्रा\appbrk{}क्षि\appbrk{}नासं}}\Afootnote{K N H; \skt{मृजु\appbrk{}व\appbrk{}क्त्रा\appbrk{}क्षि\appbrk{}नासं}\ A}}\skt{ प्रस\-न्न\-चि\-त्त\-वा\-क्चेष्टं क्लेशसहञ्च }\edtext{\edtext{\skt{शिष्य\-मु\-प\-न\-येत्}}{\lemma{\skt{शिष्यम्}}\Afootnote{K N H; \skt{भिषक्शिष्यम्}\ A}}\skt{~। स हि}}{\lemma{\skt{उपनयेत्~। स हि}}\Afootnote{K N H; \skt{उपनये}\ A}}\skt{ }\edtext{\edtext{\skt{गुणवान्तस्मै}}{\lemma{\skt{गुणवान्}}\Afootnote{K N H; \emph{om.}\ A}}}{\lemma{\skt{तस्मै}}\Afootnote{K N H; \skt{त्~। अ}\ A}}\skt{ }\edtext{\edtext{\skt{देयमतो॥}}{\lemma{\skt{देयम्}}\Afootnote{K N H; \emph{om.}\ A}}}{\lemma{\skt{अतो}}\Afootnote{K N H; \skt{तो}\ A}}\skt{ विपरीतगुणं नोपनयेत्~।}\pend

\pstart \ledsidenote{\passref{1.2.4}}\edtext{\edtext{\skt{शूद्रमपि}}{\lemma{\skt{शूद्रम्}}\Afootnote{K N H; \emph{om.}\ A}}}{\lemma{\skt{अपि}}\Afootnote{K N H; \emph{om.}\ A}}\skt{ }\edtext{\edtext{\edtext{\edtext{\skt{गुण\-व\-न्त\-म\-नु\-प\-नी\-त\-म\-ध्या\-प\-ये\-दि\-त्येके}}{\lemma{\skt{गुणवन्तम्}}\Afootnote{K N H; \emph{om.}\ A}}}{\lemma{\skt{अनुपनीतम्}}\Afootnote{K N H; \emph{om.}\ A}}}{\lemma{\skt{अध्यापयेदित्य्}}\Afootnote{K N H; \emph{om.}\ A}}}{\lemma{\skt{एके}}\Afootnote{K N H; \emph{om.}\ A}}\skt{~। उप\-न\-य\-नी\-यन्तु ब्राह्म\-ण\-म्प्र\-श\-स्तेषु }\edtext{\skt{तिथि\-क\-र\-ण\-मु\-हू\-र्त्त\-न\-क्ष\-त्रेषु}}{\lemma{\skt{तिथि\appbrk{}क\appbrk{}र\appbrk{}ण\appbrk{}मु\appbrk{}हू\appbrk{}र्त्त\appbrk{}न\appbrk{}क्ष\appbrk{}त्रेषु}}\Afootnote{A K; \skt{तिथि\appbrk{}क\appbrk{}र\appbrk{}ण\appbrk{}मू\appbrk{}हू\appbrk{}र्त्त\appbrk{}न\appbrk{}क्ष\appbrk{}त्रेषू}\ N; \skt{तिथि\appbrk{}क\appbrk{}र\appbrk{}ण\appbrk{}मु\appbrk{}हू\appbrk{}र्त्त\appbrk{}न\appbrk{}क्ष\appbrk{}त्रेषू}\ H}}\skt{ प्रशस्तायान्दिशि शुचौ }\edtext{\skt{देशे}}{\lemma{\skt{देशे}}\Afootnote{K N H; \skt{समे देशे}\ A}}\skt{ }\edtext{\skt{गोचर्ममात्रं}}{\lemma{\skt{गोचर्ममात्रं}}\Afootnote{K; \skt{चतुर्हस्तं चतुरस्रं}\ A; \skt{गो चर्म मात्रं}\ N; \skt{गो च स्म मात्रं}\ H}}\skt{ स्थण्डि\-ल\-मु\-प\-लिप्य }\edtext{\edtext{\skt{दर्भ\-सं\-स्त\-र\-ण\-म\-हितं}}{\lemma{\skt{दर्भ\appbrk{}सं\appbrk{}स्त\appbrk{}र\appbrk{}णम्}}\Afootnote{K N; \skt{गोमयेन, दर्भैः संस्तीर्य रत्न\appbrk{}पु\appbrk{}ष्प\appbrk{}ला\appbrk{}ज\appbrk{}भ\appbrk{}क्तै\appbrk{}र्दे\appbrk{}वताः पूजयित्वा विप्रान्भिषजश्च}\ A; \skt{दर्भ\appbrk{}सं\appbrk{}स्त\appbrk{}र\appbrk{}हितं कृत्वा पुष्पैर्श्च पूज}\ H}}}{\lemma{\skt{अहितं}}\Afootnote{K N; \skt{तत्रोल्लि}\ A; \emph{om.}\ H}}\skt{ }\edtext{\skt{कृत्वा}}{\lemma{\skt{कृत्वा}}\Afootnote{K N; \skt{ख्याभ्यु}\ A; \emph{om.}\ H}}\skt{ }\edtext{\edtext{\edtext{\skt{पुष्पै\-र्धू\-पै\-र्ज\-न्वै\-र्भ\-न्कैस्च}}{\lemma{\skt{पुष्पैर्}}\Afootnote{K; \skt{क्ष्य च}\ A; \skt{पुष्पै}\ N; \emph{om.}\ H}}}{\lemma{\skt{धूपैर्जन्वैर्}}\Afootnote{N; \skt{दक्षिणतो}\ A; \emph{om.}\ K H}}}{\lemma{\skt{भन्कैस्च}}\Afootnote{N; \skt{ब्रह्माणं}\ A; \skt{गन्}\ K; \emph{om.}\ H}}\skt{ }\edtext{\skt{पूजयित्वा}}{\lemma{\skt{पूजयित्वा}}\Afootnote{N; \skt{स्थापयित्वाऽग्निम्}\ A; \skt{धैर्}\ K; \skt{यित्वा}\ H}}\skt{ }\edtext{\skt{पला\-शो\-दु\-म्व\-र\-बि\-ल्वानां}}{\lemma{\skt{पला\appbrk{}शो\appbrk{}दु\appbrk{}म्व\appbrk{}र\appbrk{}बि\appbrk{}ल्वानां}}\Afootnote{N H; \skt{उपसमाधाय, खदि\appbrk{}र\appbrk{}प\appbrk{}ला\appbrk{}श\appbrk{}दे\appbrk{}व\appbrk{}दा\appbrk{}रु\appbrk{}बि\appbrk{}ल्वानां}\ A; \skt{धूपैर्भक्ष्यैश्च पूजयित्वा पला\appbrk{}शो\appbrk{}दु\appbrk{}म्व\appbrk{}र\appbrk{}बि\appbrk{}ल्वानां}\ K}}\skt{ }\edtext{\edtext{\edtext{\edtext{\edtext{\edtext{\skt{समि\-द्भि\-र्घृ\-त\-म\-क्ता\-भि\-र्दा\-र्वी\-हौ\-मि\-के\-ना\-ग्नि\-मु\-प\-स\-मा\-धा\-या\-ज्य\-ञ्जु\-हु\-यात्}}{\lemma{\skt{समिद्भिर्}}\Afootnote{K H; \skt{समिद्भि}\ A; \skt{समि}\ N}}}{\lemma{\skt{घृतम्}}\Afootnote{H; \skt{श्चतुर्णां वा क्षीरिवृक्षाणां (न्य\appbrk{}ग्रो\appbrk{}धो\appbrk{}दु\appbrk{}म्ब\appbrk{}रा\appbrk{}श्व\appbrk{}त्थ\appbrk{}म\appbrk{}धू\appbrk{}कानां) दधिमधुघृता}\ A; \skt{घृतमथक्तानि}\ K; \skt{द्भि घृतम्}\ N}}}{\lemma{\skt{अक्ताभिर्}}\Afootnote{N H; \skt{क्ताभिर्}\ A; \emph{om.}\ K}}}{\lemma{\skt{दार्वी\appbrk{}हौ\appbrk{}मि\appbrk{}के\appbrk{}ना\appbrk{}ग्निम्}}\Afootnote{N; \skt{दार्वीहौमिकेन विधिना सप्र\appbrk{}ण\appbrk{}वा\appbrk{}भि\appbrk{}र्म\appbrk{}हा\appbrk{}व्या\appbrk{}हृति}\ A; \skt{दार्वि\appbrk{}हौ\appbrk{}मि\appbrk{}के\appbrk{}ना\appbrk{}ग्निम्}\ K; \skt{द्दार्वी\appbrk{}हौ\appbrk{}मि\appbrk{}के\appbrk{}ना\appbrk{}ग्नि\appbrk{}मु\appbrk{}पस}\ H}}}{\lemma{\skt{उप\appbrk{}स\appbrk{}मा\appbrk{}धा\appbrk{}या\appbrk{}ज्यञ्}}\Afootnote{K N; \skt{भिः स्रुवे\appbrk{}णा\appbrk{}ज्या\appbrk{}हु\appbrk{}तीर्}\ A; \skt{माधायाज्यं}\ H}}\skt{~। }\edtext{\skt{प्रति\-दे\-व\-त\-मृ\-षिभ्यः}}{\lemma{\skt{ऋषिभ्यः}}\Afootnote{K N; \skt{मृषींश्}\ A; \skt{ऋषीभ्यः}\ H}}}{\lemma{\skt{जुहुयात्~। प्रतिदेवतम्}}\Afootnote{K N; \skt{जुहुयात्, ततः प्रतिदैवत}\ A; \skt{जुहुयात्~। प्रतिदैवतम्}\ H}}\skt{ }\edtext{\skt{शिष्यं}}{\lemma{\skt{शिष्यं}}\Afootnote{K N H; \skt{च}\ A}}\skt{ }\edtext{\skt{स्वाहा\-का\-रं\-ङ्का\-र\-येत्}}{\lemma{\skt{कारयेत्}}\Afootnote{K N H; \skt{कुर्यात् , शिष्यमपि कारयेत्}\ A}}\skt{~।}\pend

\pstart \ledsidenote{\passref{1.2.5}}\edtext{\skt{ब्राह्मणस्त्रयाणां}}{\lemma{\skt{ब्राह्मणस्}}\Afootnote{A K H; \skt{ब्राह्मणं}\ N}}\skt{ }\edtext{\skt{राजन्यो}}{\lemma{\skt{राजन्यो}}\Afootnote{K N H; \skt{वर्णा\appbrk{}ना\appbrk{}मु\appbrk{}प\appbrk{}न\appbrk{}यनं कर्तुमर्हति, (राजन्यो}\ A}}\skt{ }\edtext{\skt{द्वयस्या}}{\lemma{\skt{द्वयस्या}}\Afootnote{N; \skt{द्वयस्य}\ A K H}}\skt{ वैश्यो }\edtext{\skt{वैस्यस्यैव}}{\lemma{\skt{वैस्यस्यैव}}\Afootnote{K N; \skt{वैश्यस्यैवेति;) शूद्रमपि कुल\appbrk{}गु\appbrk{}ण\appbrk{}स\appbrk{}म्पन्नं मन्त्र\appbrk{}व\appbrk{}र्ज\appbrk{}म\appbrk{}नु\appbrk{}प\appbrk{}नी\appbrk{}त\appbrk{}म\appbrk{}ध्या\appbrk{}प\appbrk{}ये\appbrk{}दि\appbrk{}त्येके}\ A; \skt{वैश्यस्यैव}\ H}}\skt{~।}\pend

\pstart \ledsidenote{\passref{1.2.6}}\edtext{\skt{ततोऽग्निं}}{\lemma{\skt{ऽग्निं}}\Afootnote{K H; \skt{ऽग्निं त्रिः}\ A; \skt{ग्निं}\ N}}\skt{ परि\-णी\-या\-ग्नि\-सा\-क्षिकं शिष्यं }\edtext{\skt{ब्रूयात्~। काम\-क्रो\-ध\-लो\-भ\-मो\-ह\-मा\-ना\-हं\-का\-रे\-र्ष्या\-मा\-त्स\-र्य\-पा\-रू\-ष्य\-पै\-शु\-न्या\-नृ\-ता\-ल\-स्या\-स्य\-य\-श\-स्यानि}}{\lemma{\skt{ब्रूयात्~। काम\appbrk{}क्रो\appbrk{}ध\appbrk{}लो\appbrk{}भ\appbrk{}मो\appbrk{}ह\appbrk{}मा\appbrk{}ना\appbrk{}हं\appbrk{}का\appbrk{}रे\appbrk{}र्ष्या\appbrk{}मा\appbrk{}त्स\appbrk{}र्य\appbrk{}पा\appbrk{}रू\appbrk{}ष्य\appbrk{}पै\appbrk{}शु\appbrk{}न्या\appbrk{}नृ\appbrk{}ता\appbrk{}ल\appbrk{}स्या\appbrk{}स्य\appbrk{}य\appbrk{}श\appbrk{}स्यानि}}\Afootnote{\skt{ब्रूयात्- काम\appbrk{}क्रो\appbrk{}ध\appbrk{}लो\appbrk{}भ\appbrk{}मो\appbrk{}ह\appbrk{}मा\appbrk{}ना\appbrk{}ह\appbrk{}ङ्का\appbrk{}रे\appbrk{}र्ष्या\appbrk{}पा\appbrk{}रु\appbrk{}ष्य\appbrk{}पै\appbrk{}शु\appbrk{}न्या\appbrk{}नृ\appbrk{}ता\appbrk{}ल\appbrk{}स्या\appbrk{}य\appbrk{}श\appbrk{}स्यानि}\ A; \skt{ब्रूयात्~। काम\appbrk{}क्रो\appbrk{}ध\appbrk{}लो\appbrk{}भ\appbrk{}मो\appbrk{}ह\appbrk{}मा\appbrk{}ना\appbrk{}ह\appbrk{}ङ्क\appbrk{}रे\appbrk{}र्ष्या\appbrk{}मा\appbrk{}त्स\appbrk{}र्य\appbrk{}प\appbrk{}रु\appbrk{}ष\appbrk{}पै\appbrk{}शु\appbrk{}न्या\appbrk{}नृ\appbrk{}ता\appbrk{}ल\appbrk{}स्या\appbrk{}त्य\appbrk{}य\appbrk{}श\appbrk{}स्यानि}\ K; \skt{व्रूयात्। काम\appbrk{}क्रो\appbrk{}ध\appbrk{}लो\appbrk{}भ\appbrk{}मो\appbrk{}ह\appbrk{}मा\appbrk{}ना\appbrk{}हं\appbrk{}का\appbrk{}रे\appbrk{}र्ष्या\appbrk{}मा\appbrk{}त्स\appbrk{}र्य\appbrk{}रू\appbrk{}ष\appbrk{}पै\appbrk{}शु\appbrk{}न्या\appbrk{}नृ\appbrk{}ता\appbrk{}ल\appbrk{}स्या\appbrk{}स्य\appbrk{}य\appbrk{}श\appbrk{}स्यानि}\ N; \skt{ब्रूयात्। काम\appbrk{}क्रो\appbrk{}ध\appbrk{}लो\appbrk{}भ\appbrk{}मो\appbrk{}ह\appbrk{}मा\appbrk{}ना\appbrk{}हं\appbrk{}का\appbrk{}रे\appbrk{}र्ष्या\appbrk{}मा\appbrk{}त्स\appbrk{}र्य\appbrk{}पा\appbrk{}रु\appbrk{}ष्य\appbrk{}पै\appbrk{}शु\appbrk{}न्या\appbrk{}नृ\appbrk{}ता\appbrk{}ल\appbrk{}स्या\appbrk{}य\appbrk{}श\appbrk{}स्यानि}\ H}}\skt{ हित्वा }\edtext{\skt{काषायवाससा}}{\lemma{\skt{काषायवाससा}}\Afootnote{K N; \emph{om.}\ A; \skt{कषायवाससा}\ H}}\skt{ }\edtext{\skt{नीचनखरोम्णा}}{\lemma{\skt{नीचनखरोम्णा}}\Afootnote{A K N; \skt{नीचनखरोम्नां}\ H}}\skt{ }\edtext{\skt{त्रिरात्रं}}{\lemma{\skt{त्रिरात्रं}}\Afootnote{K N; \emph{om.}\ A; \skt{त्रिवारं}\ H}}\skt{ शुचिना }\edtext{\skt{सत्य\-ब्र\-ह्म\-च\-र्या\-भि\-वा\-द\-न\-प\-रेण}}{\lemma{\skt{सत्य\appbrk{}ब्र\appbrk{}ह्म\appbrk{}च\appbrk{}र्या\appbrk{}भि\appbrk{}वा\appbrk{}द\appbrk{}न\appbrk{}प\appbrk{}रेण}}\Afootnote{K N H; \skt{कषायवाससा सत्य\appbrk{}व्र\appbrk{}त\appbrk{}ब्र\appbrk{}ह्म\appbrk{}च\appbrk{}र्या\appbrk{}भि\appbrk{}वा\appbrk{}द\appbrk{}न\appbrk{}त\appbrk{}त्प\appbrk{}रे\appbrk{}णा\appbrk{}वश्यं}\ A}}\skt{ भवितव्यं~। }\edtext{\skt{ममा\-नु\-म\-त\-स्था\-न\-ग\-म\-न\-श\-य\-ना\-स\-न\-भो\-ज\-ना\-ध्य\-य\-न\-प\-रेण}}{\lemma{\skt{ममा\appbrk{}नु\appbrk{}म\appbrk{}त\appbrk{}स्था\appbrk{}न\appbrk{}ग\appbrk{}म\appbrk{}न\appbrk{}श\appbrk{}य\appbrk{}ना\appbrk{}स\appbrk{}न\appbrk{}भो\appbrk{}ज\appbrk{}ना\appbrk{}ध्य\appbrk{}य\appbrk{}न\appbrk{}प\appbrk{}रेण}}\Afootnote{N; \skt{मद\appbrk{}नु\appbrk{}म\appbrk{}त\appbrk{}स्था\appbrk{}न\appbrk{}ग\appbrk{}म\appbrk{}न\appbrk{}श\appbrk{}य\appbrk{}ना\appbrk{}स\appbrk{}न\appbrk{}भो\appbrk{}ज\appbrk{}ना\appbrk{}ध्य\appbrk{}य\appbrk{}न\appbrk{}प\appbrk{}रेण}\ A; \skt{ममा\appbrk{}नु\appbrk{}म\appbrk{}त\appbrk{}स्था\appbrk{}न\appbrk{}ग\appbrk{}म\appbrk{}न\appbrk{}श\appbrk{}य\appbrk{}ना\appbrk{}स\appbrk{}न\appbrk{}भो\appbrk{}ज\appbrk{}ना\appbrk{}ध्या\appbrk{}य\appbrk{}न\appbrk{}प\appbrk{}रेण}\ K; \skt{ममानुमते स्थान\appbrk{}ग\appbrk{}म\appbrk{}न\appbrk{}श\appbrk{}य\appbrk{}ना\appbrk{}श\appbrk{}न\appbrk{}भो\appbrk{}ज\appbrk{}ना\appbrk{}ध्य\appbrk{}य\appbrk{}न\appbrk{}प\appbrk{}रेण}\ H}}\skt{ भूत्वा मत्प्रियहितेषु }\edtext{\skt{वर्त्ति\-त\-व्य\-म\-तोऽ\-न्यथा}}{\lemma{\skt{ऽन्यथा}}\Afootnote{K; \skt{ऽन्यथा ते}\ A; \skt{न्यथा}\ N H}}\skt{ वर्त्त\-मा\-न\-स्या\-धर्म्मो }\edtext{\skt{भवत्यफला च}}{\lemma{\skt{अफला च}}\Afootnote{A N H; \skt{अफला}\ K}}\skt{ विद्या न च }\edtext{\skt{प्राकाश्यं}}{\lemma{\skt{प्राकाश्यं}}\Afootnote{A K N; \skt{प्राकास्यं}\ H}}\skt{ }\edtext{\skt{प्राप्नुयात्}}{\lemma{\skt{प्राप्नुयात्}}\Afootnote{K N H; \skt{प्राप्नोति}\ A}}\skt{~।}\pend

\pstart \ledsidenote{\passref{1.2.7}}\skt{अहम्वा त्वयि }\edtext{\skt{सम्यग्वर्त्तमाने}}{\lemma{\skt{सम्यग्वर्त्तमाने}}\Afootnote{A K N; \skt{सम्यक्वर्त्तमाने}\ H}}\skt{ }\edtext{\skt{यद्य\-न\-न्य\-था\-दर्शा}}{\lemma{\skt{अनन्यथादर्शा}}\Afootnote{N H; \skt{अन्यथादर्शी}\ A; \skt{अनन्यथादर्शी}\ K}}\skt{ }\edtext{\skt{स्यात्तदैव}}{\lemma{\skt{स्यात्तदैव}}\Afootnote{K\textsuperscript{pc} H; \skt{स्याम्}\ A; \skt{स्यात्तदै}\ K\textsuperscript{ac}; \skt{स्यात्तदेव}\ N}}\skt{ }\edtext{\skt{नासौ\-भा\-ग्य\-वि\-द्या\-फ\-ल\-भाक्च}}{\lemma{\skt{नासौ\appbrk{}भा\appbrk{}ग्य\appbrk{}वि\appbrk{}द्या\appbrk{}फ\appbrk{}ल\appbrk{}भाक्च}}\Afootnote{N; \skt{एनोभाग्भ}\ A; \skt{न सौभा\appbrk{}ग्य\appbrk{}वि\appbrk{}द्या\appbrk{}फ\appbrk{}ल\appbrk{}भाक्च}\ K H}}\skt{ }\edtext{\skt{भवेयं}}{\lemma{\skt{भवेयं}}\Afootnote{K N H; \skt{वेय\appbrk{}म\appbrk{}फ\appbrk{}ल\appbrk{}वि\appbrk{}द्यश्च}\ A}}\skt{~।}\pend

\pstart \ledsidenote{\passref{1.2.7.a}}\edtext{\skt{यस्मा\-द\-रो\-ग\-वता}}{\lemma{\skt{यस्मा\appbrk{}द\appbrk{}रो\appbrk{}ग\appbrk{}वता}}\Afootnote{N; \skt{यस्माद्रोगवता}\ K H}}\skt{ }\edtext{\skt{धर्म्मा\-र्थ\-का\-म\-मोक्षाः}}{\lemma{\skt{धर्म्मा\appbrk{}र्थ\appbrk{}का\appbrk{}म\appbrk{}मोक्षाः}}\Afootnote{K H; \skt{धर्म्मा\appbrk{}सौ\appbrk{}र्थ\appbrk{}का\appbrk{}म\appbrk{}मोक्षाः}\ N}}\skt{ }\edtext{\skt{प्राप्यन्ते}}{\lemma{\skt{प्राप्यन्ते}}\Afootnote{K N; \skt{प्राथ्यन्ते}\ H}}\skt{~।}\pend

\pstart \ledsidenote{\passref{1.2.8}}\edtext{\edtext{\edtext{\skt{तस्मा\-द्द्वि\-ज\-द\-रि\-द्र\-सा\-ध्व\-ना\-था\-भ्यु\-प\-ग\-त\-पा\-श\-ण्ड\-स्थि\-ता\-ना\-मा\-त्म\-बा\-न्ध\-वा\-ना\-मि\-वा\-त्म\-भे\-षजैः}}{\lemma{\skt{तस्मा\appbrk{}द्द्वि\appbrk{}ज\appbrk{}द\appbrk{}रि\appbrk{}द्र\appbrk{}सा\appbrk{}ध्व\appbrk{}ना\appbrk{}था\appbrk{}भ्यु\appbrk{}प\appbrk{}ग\appbrk{}त\appbrk{}पा\appbrk{}श\appbrk{}ण्ड\appbrk{}स्थि\appbrk{}ता\appbrk{}नाम्}}\Afootnote{N; \skt{द्विज\appbrk{}गु\appbrk{}रु\appbrk{}द\appbrk{}रि\appbrk{}द्र\appbrk{}मि\appbrk{}त्र\appbrk{}प्र\appbrk{}व्र\appbrk{}जि\appbrk{}तो\appbrk{}प\appbrk{}न\appbrk{}त\appbrk{}सा\appbrk{}ध्व\appbrk{}ना\appbrk{}था\appbrk{}भ्यु\appbrk{}प\appbrk{}ग\appbrk{}तानां}\ A; \skt{तस्मा\appbrk{}द्द्वि\appbrk{}ज\appbrk{}द\appbrk{}रि\appbrk{}द्र\appbrk{}सा\appbrk{}ध्व\appbrk{}ना\appbrk{}था\appbrk{}भ्यु\appbrk{}प\appbrk{}ग\appbrk{}त\appbrk{}पा\appbrk{}ष\appbrk{}ण्ड\appbrk{}स्थि\appbrk{}ता\appbrk{}नाम्}\ K; \skt{तस्मा द्विज\appbrk{}द\appbrk{}रि\appbrk{}द्र\appbrk{}सा\appbrk{}ध्व\appbrk{}ना\appbrk{}था\appbrk{}भ्यु\appbrk{}प\appbrk{}ग\appbrk{}त\appbrk{}पा\appbrk{}ष\appbrk{}ण्ड्स्थि\appbrk{}ता\appbrk{}नाम्}\ H}}}{\lemma{\skt{आत्म\appbrk{}बा\appbrk{}न्ध\appbrk{}वा\appbrk{}नाम्}}\Afootnote{K N H; \skt{चात्म\appbrk{}बा\appbrk{}न्ध\appbrk{}वा\appbrk{}नाम्}\ A}}}{\lemma{\skt{इवात्मभेषजैः}}\Afootnote{N; \skt{इव स्वभैषजैः}\ A; \skt{इव आत्मभेषजैः}\ K H}}\skt{~। प्रति\-क\-र्त\-व्य\-मेवं }\edtext{\skt{साधु}}{\lemma{\skt{साधु}}\Afootnote{A K N; \skt{साधुर्}\ H}}\skt{ भवति~। }\edtext{\skt{व्याध\-शा\-कु\-नि\-क\-प\-ति\-त\-पा\-प\-क\-र्त्तॄ\-णाञ्च न}}{\lemma{\skt{व्याध\appbrk{}शा\appbrk{}कु\appbrk{}नि\appbrk{}क\appbrk{}प\appbrk{}ति\appbrk{}त\appbrk{}पा\appbrk{}प\appbrk{}क\appbrk{}र्त्तॄ\appbrk{}णाञ्च न}}\Afootnote{N; \skt{व्याध\appbrk{}शा\appbrk{}कु\appbrk{}नि\appbrk{}क\appbrk{}प\appbrk{}ति\appbrk{}त\appbrk{}पा\appbrk{}प\appbrk{}का\appbrk{}रिणां च न}\ A; \skt{व्याध\appbrk{}सा\appbrk{}कु\appbrk{}नि\appbrk{}क\appbrk{}प\appbrk{}ति\appbrk{}त\appbrk{}या\appbrk{}प\appbrk{}क\appbrk{}र्त्तॄ\appbrk{}णाञ्च न}\ K H}}\skt{ प्रति\-क\-र्त्त\-व्य\-मेवं }\edtext{\skt{विद्या}}{\lemma{\skt{विद्या}}\Afootnote{A K N; \skt{विद्यां}\ H}}\skt{ प्रकाशते~। }\edtext{\skt{मित्र\-ध\-र्म\-का\-म\-य\-शांसि}}{\lemma{\skt{मित्र\appbrk{}ध\appbrk{}र्म\appbrk{}का\appbrk{}म\appbrk{}य\appbrk{}शांसि}}\Afootnote{N; \skt{मित्र\appbrk{}य\appbrk{}शो\appbrk{}ध\appbrk{}र्मा\appbrk{}र्थ\appbrk{}का\appbrk{}मांश्}\ A; \skt{मित्र\appbrk{}ध\appbrk{}र्म\appbrk{}का\appbrk{}म\appbrk{}य\appbrk{}शन्सि}\ K H}}\skt{ }\edtext{\skt{चावाप्नोति}}{\lemma{\skt{चावाप्नोति}}\Afootnote{K\textsuperscript{pc} N; \skt{च प्राप्नोति}\ A; \skt{वाप्नोति}\ K\textsuperscript{ac}; \skt{वाप्राप्नोति}\ H\textsuperscript{ac}; \skt{चावाप्राप्नोति}\ H\textsuperscript{pc}}}\skt{~॥}\pend

\pstart\versepar \ledsidenote{\passref{1.2.9}}\skt{भवतश्चात्र~॥ }\pend
\pstart\versepar \skt{कृष्णाष्टमी }\edtext{\edtext{\skt{तन्निधनेऽहनी}}{\lemma{\skt{तन्निधने}}\Afootnote{A H; \skt{त्तन्निधने}\ K N}}\skt{ द्वे}}{\lemma{\skt{ऽहनी द्वे}}\Afootnote{A K N; \skt{हनी द्वे …}\ H}}\skt{ }\pend
\pstart\versepar \skt{}\edtext{\edtext{\edtext{\edtext{\edtext{\skt{शुक्ला\-द\-येऽ\-प्ये\-व\-म\-ह\-र्द्वि\-सन्ध्यं}}{\lemma{\skt{शुक्लादये}}\Afootnote{K N; \skt{शुक्ले त}\ A; \skt{शुक्ले दये}\ H}}}{\lemma{\skt{ऽप्य्}}\Afootnote{K; \skt{थाऽ}\ A; \skt{प्य्}\ N H}}}{\lemma{\skt{एवम्}}\Afootnote{K N H; \skt{प्येव}\ A}}}{\lemma{\skt{अहर्}}\Afootnote{K N H; \skt{मह}\ A}}}{\lemma{\skt{द्विसन्ध्यं}}\Afootnote{K N H\textsuperscript{ac}; \skt{र्द्विसन्ध्यम्}\ A; \skt{द्विसयान्ध्यं}\ H\textsuperscript{pc}}}\skt{~। }\pend
\pstart\versepar \skt{अका\-ल\-वि\-द्यु\-त्स्त\-न\-यि\-त्नु\-घोषे }\pend
\pstart\versepar \skt{स्वत\-न्त्र\-रा\-ष्ट्र\-क्षि\-ति\-प\-व्य\-थासु~।}\edtext{}{\lemma{\skt{भवतश्चात्र~॥ कृष्णाष्टमी}}\Afootnote{\skt{भवतश्चात्र कृष्णेऽष्टमी}\ A; \skt{कृष्णाष्टमी}\ K N; \skt{भवति चात्र~॥ कृष्त्तषूमीन्}\ H}}\pend

\pstart\versepar \ledsidenote{\passref{1.2.10}}\edtext{\skt{श्मशा\-न\-या\-ना\-ध्व\-त\-ना\-ह\-वेषु}}{\lemma{\skt{श्मशा\appbrk{}न\appbrk{}या\appbrk{}ना\appbrk{}ध्व\appbrk{}त\appbrk{}ना\appbrk{}ह\appbrk{}वेषु}}\Afootnote{N; \skt{श्मशा\appbrk{}न\appbrk{}या\appbrk{}ना\appbrk{}द्य\appbrk{}त\appbrk{}ना\appbrk{}ह\appbrk{}वेषु}\ A H; \skt{श्मशा\appbrk{}न\appbrk{}या\appbrk{}ना\appbrk{}ध्य\appbrk{}त\appbrk{}ना\appbrk{}ह\appbrk{}वेषु}\ K}}\skt{ }\pend
\pstart\versepar \skt{}\edtext{\skt{तथो\-त्स\-वो\-त्पा\-ति\-क\-द\-र्श\-नेषु}}{\lemma{\skt{तथो\appbrk{}त्स\appbrk{}वो\appbrk{}त्पा\appbrk{}ति\appbrk{}क\appbrk{}द\appbrk{}र्श\appbrk{}नेषु}}\Afootnote{K; \skt{महो\appbrk{}त्स\appbrk{}वौ\appbrk{}त्पा\appbrk{}ति\appbrk{}क\appbrk{}द\appbrk{}र्श\appbrk{}नेषु}\ A; \skt{तथौ\appbrk{}त्स\appbrk{}व\appbrk{}ओ\appbrk{}त्पा\appbrk{}ति\appbrk{}क\appbrk{}द\appbrk{}र्श\appbrk{}नेषु}\ N; \skt{तथो\appbrk{}त्स\appbrk{}वो\appbrk{}त्या\appbrk{}ति\appbrk{}क\appbrk{}द\appbrk{}र्श\appbrk{}नेषु}\ H}}\skt{~। }\pend
\pstart\versepar \skt{}\edtext{\edtext{\skt{नाध्येयमन्येषु}}{\lemma{\skt{नाध्येयम्}}\Afootnote{K N; \skt{नाध्येय}\ A; \skt{नाध्योयम्}\ H}}\skt{ च}}{\lemma{\skt{अन्येषु च}}\Afootnote{N H; \skt{मन्येषु च}\ A; \skt{अन्ये}\ K}}\skt{ }\edtext{\skt{येषु}}{\lemma{\skt{येषु}}\Afootnote{A N H; \skt{षु}\ K}}\skt{ विप्राण्ण् }\pend
\pstart\versepar \skt{आधीयते नाशुचिना च नित्यम्~॥}\edtext{}{\lemma{\skt{विप्राण्णाधीयते}}\Afootnote{N; \skt{विप्रा नाधीयते}\ A H; \skt{विप्रा यते}\ K}}\pend

\pstart \ledsidenote{\passref{1.2.trailer}}\edtext{\skt{इति}}{\lemma{\skt{इति}}\Afootnote{N H; \skt{इति सुश्रु\appbrk{}त\appbrk{}सं\appbrk{}हि\appbrk{}तायां सूत्रस्थाने शिष्यो\appbrk{}प\appbrk{}न\appbrk{}य\appbrk{}नीयो नाम द्वितीयोऽध्यायः}\ A; \skt{इति~॥ {\symfont ❈}}\ K}}\skt{~॥}\pend

\endnumbering
\beginnumbering
\eledsection*{Adhyāya 3}
\pstart \ledsidenote{\passref{1.3.1}}\edtext{\edtext{\skt{अथा\-तोऽ\-ध्य\-य\-न\-स\-म्प्र\-दा\-नीयं}}{\lemma{\skt{अथातो}}\Afootnote{A K H; \skt{अथातोऽ}\ N}}}{\lemma{\skt{ऽध्य\appbrk{}य\appbrk{}न\appbrk{}स\appbrk{}म्प्र\appbrk{}दा\appbrk{}नीयं}}\Afootnote{A H; \skt{ध्यय\appbrk{}न\appbrk{}स\appbrk{}म्प्र\appbrk{}दा\appbrk{}नीयं}\ K N}}\skt{ }\edtext{\skt{व्याख्यास्यामः}}{\lemma{\skt{व्याख्यास्यामः}}\Afootnote{H; \skt{अध्यायं व्याख्यास्यामः}\ A; \emph{om.}\ K N}}\skt{~॥}\pend

\pstart \ledsidenote{\passref{1.3.3}}\skt{प्रागभिहितं }\edtext{\skt{सविं\-श\-म\-ध्या\-य\-शतं}}{\lemma{\skt{सविं\appbrk{}श\appbrk{}म\appbrk{}ध्या\appbrk{}य\appbrk{}शतं}}\Afootnote{\skt{सविं\appbrk{}श\appbrk{}म\appbrk{}ध्या\appbrk{}य\appbrk{}शतं}\ A K N H}}\skt{ पञ्चसु स्थानेषु }\edtext{\skt{चेति}}{\lemma{\skt{चेति}}\Afootnote{K N H; \emph{om.}\ A}}\skt{~॥ तत्र }\edtext{\skt{श्लोकस्थाने}}{\lemma{\skt{श्लोकस्थाने}}\Afootnote{K; \skt{सूत्रस्थान}\ A; \skt{श्लोकस्थने}\ N; \skt{…स्थाने}\ H}}\skt{ }\edtext{\skt{अध्यायाः}}{\lemma{\skt{अध्यायाः}}\Afootnote{K; \skt{मध्यायाः}\ A; \skt{ह्यध्याया}\ N; \skt{ध्यायाः}\ H}}\skt{ }\edtext{\skt{षट्चत्वारिंशत्~। षोड\-श\-नि\-दा\-नानि}}{\lemma{\skt{षट्चत्वारिंशत्~। षोड\appbrk{}श\appbrk{}नि\appbrk{}दा\appbrk{}नानि}}\Afootnote{A K N; \skt{षट्चत्वारिंशत्~। षोड्श निदानानि}\ H}}\skt{~। दश शारीराणि~। }\edtext{\skt{चत्वा\-रि\-श\-च्चि\-कि\-त्सि\-तानि}}{\lemma{\skt{चत्वा\appbrk{}रि\appbrk{}श\appbrk{}च्चि\appbrk{}कि\appbrk{}त्सि\appbrk{}तानि}}\Afootnote{K; \skt{चत्वा\appbrk{}रिं\appbrk{}श\appbrk{}च्चि\appbrk{}कि\appbrk{}त्सि\appbrk{}तानि}\ A; \skt{चत्वारिंश चिकित्सितानि}\ N; \skt{चत्वा\appbrk{}रिं\appbrk{}श\appbrk{}च्चि\appbrk{}कि\appbrk{}त्सि\appbrk{}तानि}\ H}}\skt{~। अष्टौ }\edtext{\skt{कल्पाः}}{\lemma{\skt{कल्पाः}}\Afootnote{K N; \skt{कल्पाः}\ A H}}\skt{~॥ }\edtext{\skt{भवन्ति}}{\lemma{\skt{भवन्ति}}\Afootnote{K N; \skt{त}\ A; \skt{… भवन्ति}\ H}}\skt{ }\edtext{\skt{चात्र}}{\lemma{\skt{चात्र}}\Afootnote{K N; \skt{दुत्तरं षट्षष्टिः}\ A; \skt{चात्र श्लोकाः}\ H}}\skt{~॥}\pend

\pstart\versepar \ledsidenote{\passref{1.3.4}}\skt{वेदोत्पत्तिः शिष्य\-न\-य\-स्त\-था\-ध्य\-य\-न\-दा\-निकः~। }\pend
\pstart\versepar \skt{}\edtext{\skt{प्रभा\-ष\-णा\-ग्र\-ह\-रणः}}{\lemma{\skt{प्रभा\appbrk{}ष\appbrk{}णा\appbrk{}ग्र\appbrk{}ह\appbrk{}रणः}}\Afootnote{K; \skt{प्रभा\appbrk{}ष\appbrk{}णा\appbrk{}ग्र\appbrk{}ह\appbrk{}रणा}\ A; \skt{प्रभा\appbrk{}ष\appbrk{}णा\appbrk{}ग्र\appbrk{}ह॥॥णः}\ N; \skt{प्रभा\appbrk{}ष\appbrk{}णा\appbrk{}ग्र\appbrk{}ह\appbrk{}रणः}\ H}}\skt{~। }\edtext{\skt{ऋतुचर्याथ}}{\lemma{\skt{ऋतुचर्याथ}}\Afootnote{K N; \skt{वृतुचर्याऽथ}\ A; \skt{चर्याथ}\ H}}\skt{ यान्त्रिकः~।}\pend

\pstart\versepar \ledsidenote{\passref{1.3.5}}\skt{शस्त्रावचारणं योग्या विशिखा क्षारकल्पनं~। }\pend
\pstart\versepar \skt{}\edtext{\skt{अग्नि\-क\-र्म\-ज\-लौ\-का\-ख्या\-व\-ध्यायौ}}{\lemma{\skt{अग्नि\appbrk{}क\appbrk{}र्म\appbrk{}ज\appbrk{}लौ\appbrk{}का\appbrk{}ख्या\appbrk{}व\appbrk{}ध्यायौ}}\Afootnote{\skt{अग्निकर्म जलौकाख्यो ह्यध्यायो}\ A; \skt{अग्नि\appbrk{}क\appbrk{}र्म\appbrk{}ज\appbrk{}लौ\appbrk{}का\appbrk{}ख्या\appbrk{}व\appbrk{}ध्यायौ}\ K; \skt{अग्नि\appbrk{}क\appbrk{}र्म्म\appbrk{}ज\appbrk{}लौ\appbrk{}का\appbrk{}ख्या\appbrk{}न\appbrk{}ध्यायो}\ N; \skt{अग्नि\appbrk{}क\appbrk{}र्म्म\appbrk{}ज\appbrk{}लौ\appbrk{}का\appbrk{}ख्या\appbrk{}व\appbrk{}ध्यायो}\ H}}\skt{ }\edtext{\skt{रक्तवर्ण्णनं}}{\lemma{\skt{रक्तवर्ण्णनं}}\Afootnote{A K H; \skt{रक्तवर्ण्णकम्}\ N}}\skt{~।}\pend

\pstart\versepar \ledsidenote{\passref{1.3.6}}\edtext{\skt{दोष\-धा\-तु\-म\-ला\-द्यानां}}{\lemma{\skt{दोष\appbrk{}धा\appbrk{}तु\appbrk{}म\appbrk{}ला\appbrk{}द्यानां}}\Afootnote{A K H; \skt{दोष\appbrk{}धा\appbrk{}तु\appbrk{}म\appbrk{}ला\appbrk{}द्याना}\ N}}\skt{ विज्ञानाध्याय एव च~। }\pend
\pstart\versepar \skt{}\edtext{\skt{कर्ण्णव्यथ}}{\lemma{\skt{कर्ण्णव्यथ}}\Afootnote{K; \skt{कर्णव्यधा}\ A; \skt{कर्ण्णव्यध}\ N H}}\skt{ }\edtext{\skt{आम\-प\-क्व\-मा\-लेपो}}{\lemma{\skt{आम\appbrk{}प\appbrk{}क्व\appbrk{}मा\appbrk{}लेपो}}\Afootnote{\skt{मप\appbrk{}क्वै\appbrk{}षा\appbrk{}वा\appbrk{}लेपो}\ A; \skt{आम\appbrk{}प\appbrk{}क्व\appbrk{}मा\appbrk{}लेपो}\ K; \skt{आम\appbrk{}प\appbrk{}क्व\appbrk{}मा\appbrk{}लयो}\ N H}}\skt{ }\edtext{\skt{व्रणितासनं}}{\lemma{\skt{व्रणितासनं}}\Afootnote{K N; \skt{व्रण्युपासनम्}\ A; \skt{व्रणितासनः}\ H}}\skt{~।}\pend

\pstart\versepar \ledsidenote{\passref{1.3.7}}\edtext{\skt{हिताहितो}}{\lemma{\skt{हिताहितो}}\Afootnote{A K N; \skt{हिताहितौ}\ H}}\skt{ व्रणप्रश्नो व्रणास्रावश्च यः }\edtext{\skt{स्मृतः}}{\lemma{\skt{स्मृतः}}\Afootnote{K N H; \skt{पृथ}\ A}}\skt{~। }\pend
\pstart\versepar \skt{}\edtext{\edtext{\skt{कृत्या\-कृ\-त्य\-वि\-धि\-र्व्या\-धि\-स\-मु\-द्दे\-शीय}}{\lemma{\skt{कृत्या\appbrk{}कृ\appbrk{}त्य\appbrk{}वि\appbrk{}धिर्}}\Afootnote{K N; \skt{क्~। कृत्याकृत्यविधि}\ A; \skt{कृत्याकृत्यविधि}\ H}}}{\lemma{\skt{व्याधि\appbrk{}स\appbrk{}मु\appbrk{}द्दे\appbrk{}शीय}}\Afootnote{K N; \skt{र्व्याधि\appbrk{}स\appbrk{}मु\appbrk{}द्दे\appbrk{}शीय}\ A H}}\skt{ }\edtext{\skt{एव}}{\lemma{\skt{एव}}\Afootnote{A K H; \skt{एव्र}\ N}}\skt{ च~।}\pend

\pstart\versepar \ledsidenote{\passref{1.3.8}}\skt{विनिश्चयः शस्त्रविधौ प्रन\-ष्ट\-ज्ञा\-नि\-क\-स्तथा~। }\pend
\pstart\versepar \skt{}\edtext{\edtext{\skt{शल्यो\-द्धृ\-ति\-र्व्र\-ण\-ज्ञानं}}{\lemma{\skt{शल्योद्धृतिर्}}\Afootnote{K N; \skt{शल्योद्धृति}\ A H}}}{\lemma{\skt{व्रणज्ञानं}}\Afootnote{K N; \skt{र्व्रणज्ञानं}\ A H}}\skt{ दूत\-स्व\-प्न\-नि\-द\-र्शनं~।}\pend

\pstart\versepar \ledsidenote{\passref{1.3.9}}\skt{पंचेन्द्रियन्तथा }\edtext{\skt{छाया}}{\lemma{\skt{छाया}}\Afootnote{A K N; \skt{चछाया}\ H}}\skt{ }\edtext{\skt{स्वभा\-वा\-द्वै\-कृ\-तोपि च}}{\lemma{\skt{स्वभा\appbrk{}वा\appbrk{}द्वै\appbrk{}कृ\appbrk{}तोपि च}}\Afootnote{H; \skt{स्वभावाद्वैकृतं तथा~। अ}\ A; \skt{स्वभाद्वैकृतोपि च}\ K\textsuperscript{ac}; \skt{स्वभा\appbrk{}वा\appbrk{}द्वै\appbrk{}कृ\appbrk{}तोपि च}\ K\textsuperscript{pc}; \skt{स्वाभा\appbrk{}वा\appbrk{}कृ\appbrk{}द्वै\appbrk{}कृतोऽ पि च}\ N}}\skt{~। }\pend
\pstart\versepar \skt{वारणो युक्तसेनीय }\edtext{\skt{आतु\-र\-क्र\-म\-मि\-श्रकौ}}{\lemma{\skt{आतु\appbrk{}र\appbrk{}क्र\appbrk{}म\appbrk{}मि\appbrk{}श्रकौ}}\Afootnote{K; \skt{आतु\appbrk{}र\appbrk{}क्र\appbrk{}म\appbrk{}भू\appbrk{}मिकौ}\ A; \skt{आतु\appbrk{}रा\appbrk{}क्र\appbrk{}म\appbrk{}मि\appbrk{}श्रकौ}\ N H}}\skt{~।}\pend

\pstart\versepar \ledsidenote{\passref{1.3.10}}\edtext{\skt{भूमिभागो}}{\lemma{\skt{भूमिभागो}}\Afootnote{K N; \skt{मिश्रकाख्यो}\ A; \skt{सूमिभा ग्यो}\ H\textsuperscript{ac}; \skt{सूमिभा गा ग्यो}\ H\textsuperscript{pc}}}\skt{ }\edtext{\skt{द्रव्यगणाः}}{\lemma{\skt{द्रव्यगणाः}}\Afootnote{K; \skt{द्रव्यगणः}\ A H N}}\skt{ संशुद्धौ }\edtext{\skt{शमने च यः}}{\lemma{\skt{शमने च यः}}\Afootnote{A K N; \skt{समने च यः}\ H}}\skt{~। }\pend
\pstart\versepar \skt{द्रव्यादीनाञ्च विज्ञानं }\edtext{\skt{विशेषे}}{\lemma{\skt{विशेषे}}\Afootnote{K N H; \skt{विशेषो}\ A}}\skt{ }\edtext{\skt{द्रव्यगेपरः}}{\lemma{\skt{द्रव्यगेपरः}}\Afootnote{K; \skt{द्रव्यगोऽपरः}\ A; \skt{द्रव्यगोपरः}\ N H}}\skt{~।}\pend

\pstart\versepar \ledsidenote{\passref{1.3.11}}\skt{रसज्ञानं }\edtext{\edtext{\skt{वम\-ना\-र्थ\-म\-ध्यायो}}{\lemma{\skt{वमनार्थम्}}\Afootnote{K N; \skt{वमनार्थ}\ A H}}}{\lemma{\skt{अध्यायो}}\Afootnote{K N; \skt{मध्यायो}\ A; \skt{मध्याये}\ H}}\skt{ }\edtext{\skt{रेचनस्य च}}{\lemma{\skt{रेचनस्य च}}\Afootnote{K N H; \skt{रेचनाय च}\ A}}\skt{~। }\pend
\pstart\versepar \skt{}\edtext{\edtext{\skt{द्रव\-द्र\-व्य\-वि\-धि\-स्त\-द्व\-द\-न्न\-पा\-न\-विधिः}}{\lemma{\skt{द्रव\appbrk{}द्र\appbrk{}व्य\appbrk{}वि\appbrk{}धिस्}}\Afootnote{K N H; \skt{द्रवद्रव्यविधि}\ A}}}{\lemma{\skt{तद्व\appbrk{}द\appbrk{}न्न\appbrk{}पा\appbrk{}न\appbrk{}विधिः}}\Afootnote{K N H; \skt{स्तद्व\appbrk{}द\appbrk{}न्न\appbrk{}पा\appbrk{}न\appbrk{}वि\appbrk{}धिस्}\ A}}\skt{ }\edtext{\skt{शुभः}}{\lemma{\skt{शुभः}}\Afootnote{K N H; \skt{तथा}\ A}}\skt{~।}\pend

\pstart\versepar \ledsidenote{\passref{1.3.12}}\edtext{\skt{सूच\-ना\-त्सू\-त्र\-णा\-च्चैव}}{\lemma{\skt{सूच\appbrk{}ना\appbrk{}त्सू\appbrk{}त्र\appbrk{}णा\appbrk{}च्चैव}}\Afootnote{K; \skt{सूच\appbrk{}ना\appbrk{}त्सू\appbrk{}त्र\appbrk{}णा\appbrk{}च्चैव}\ A H; \skt{सूचना सूत्रणाच्चैव}\ N}}\skt{ }\edtext{\skt{साध\-ना\-च्चा\-र्थ\-स\-न्ततेः}}{\lemma{\skt{साध\appbrk{}ना\appbrk{}च्चा\appbrk{}र्थ\appbrk{}स\appbrk{}न्ततेः}}\Afootnote{K N; \skt{सव\appbrk{}ना\appbrk{}च्चा\appbrk{}र्थ\appbrk{}स\appbrk{}न्ततेः}\ A H}}\skt{~। }\pend
\pstart\versepar \skt{}\edtext{\skt{षट्च\-त्वा\-रि\-म्श\-द\-ध्यायं}}{\lemma{\skt{षट्च\appbrk{}त्वा\appbrk{}रि\appbrk{}म्श\appbrk{}द\appbrk{}ध्यायं}}\Afootnote{K; \skt{षट्च\appbrk{}त्वा\appbrk{}रिं\appbrk{}श\appbrk{}द\appbrk{}ध्यायं}\ A N H}}\skt{ }\edtext{\edtext{\skt{सूत्र\-स्था\-न\-म्प्र\-च\-क्षते}}{\lemma{\skt{सूत्रस्थानम्}}\Afootnote{A K H; \skt{सूत्रस्थान्म्}\ N}}}{\lemma{\skt{प्रचक्षते}}\Afootnote{A K N; \skt{प्रचक्ष्यते}\ H}}\skt{~॥}\pend

\pstart\versepar \ledsidenote{\passref{1.3.13}}\edtext{\edtext{\skt{वात\-व्या\-धि\-क\-म\-र्शाम्शि}}{\lemma{\skt{वातव्याधिकम्}}\Afootnote{K N; \skt{वातव्याधिक}\ A H}}}{\lemma{\skt{अर्शाम्शि}}\Afootnote{K N; \skt{मर्शांसि}\ A H}}\skt{ साश्मरिश्च भगन्दरः }\pend
\pstart\versepar \skt{}\edtext{\skt{कुष्ठमेहोदरा}}{\lemma{\skt{कुष्ठमेहोदरा}}\Afootnote{K N H; \skt{कुष्ठमेहोदरं}\ A}}\skt{ }\edtext{\skt{मूढ}}{\lemma{\skt{मूढ}}\Afootnote{K N H; \skt{मूढो}\ A}}\skt{ }\edtext{\skt{विद्रध्यः}}{\lemma{\skt{विद्रध्यः}}\Afootnote{K N; \skt{विद्रधिः}\ A; \skt{विद्रधाः}\ H}}\skt{ परिसर्पणं~।}\pend

\pstart\versepar \ledsidenote{\passref{1.3.14}}\edtext{\skt{ग्रन्थि\-वृ\-द्धि\-क्षु\-द्र\-शू\-क\-भ\-ग्नाश्च}}{\lemma{\skt{ग्रन्थि\appbrk{}वृ\appbrk{}द्धि\appbrk{}क्षु\appbrk{}द्र\appbrk{}शू\appbrk{}क\appbrk{}भ\appbrk{}ग्नाश्च}}\Afootnote{A K N; \skt{ग्रन्थि\appbrk{}वृ\appbrk{}द्धि\appbrk{}क्षु\appbrk{}द्र\appbrk{}शू\appbrk{}क…\appbrk{}भ\appbrk{}ग्नाश्च}\ H}}\skt{ }\edtext{\skt{मुखरौगिकं}}{\lemma{\skt{मुखरौगिकं}}\Afootnote{K; \skt{मुखरोगिकम्}\ A H N}}\skt{~। }\pend
\pstart\versepar \skt{हेतु\-ल\-क्ष\-ण\-नि\-र्दे\-शा\-न्नि\-दा\-ना\-नीति षोडश~॥}\pend

\pstart\versepar \ledsidenote{\passref{1.3.15}}\edtext{\edtext{\edtext{\skt{भूत\-चि\-न्ता\-र\-जः\-शु\-द्धि\-र्ग\-र्भा\-व\-क्रा\-न्ति\-रेव}}{\lemma{\skt{भूत\appbrk{}चि\appbrk{}न्ता\appbrk{}र\appbrk{}जः\appbrk{}शु\appbrk{}द्धिर्}}\Afootnote{K; \skt{भूतचिन्ता रजःशुद्धि}\ A H; \skt{भूतचिन्ता राजशुद्धिर्}\ N}}}{\lemma{\skt{गर्भावक्रान्तिर्}}\Afootnote{K N; \skt{र्गर्भावक्रान्ति}\ A; \skt{र्ग्नर्भावक्रान्ति}\ H}}\skt{ च}}{\lemma{\skt{एव च}}\Afootnote{K N; \skt{रेव}\ A; \skt{रेव च}\ H}}\skt{~। }\pend
\pstart\versepar \skt{}\edtext{\skt{गर्भस्य च}}{\lemma{\skt{गर्भस्य च}}\Afootnote{K N H; \skt{च}\ A}}\skt{ }\edtext{\skt{व्याकरणं}}{\lemma{\skt{व्याकरणं}}\Afootnote{K N H; \skt{व्याकरणं च}\ A}}\skt{ }\edtext{\skt{शरीरस्य च}}{\lemma{\skt{शरीरस्य च}}\Afootnote{K N H; \skt{गर्भस्य शरीरस्य च}\ A}}\skt{ यत्स्मृतं~।}\pend

\pstart\versepar \ledsidenote{\passref{1.3.16}}\edtext{\edtext{\skt{प्रत्ये\-क\-म\-र्म\-नि\-र्देशः}}{\lemma{\skt{प्रत्येकम्}}\Afootnote{A K; \skt{प्रत्येक}\ N H}}}{\lemma{\skt{अर्मनिर्देशः}}\Afootnote{K; \skt{मर्मनिर्देशः}\ A N H}}\skt{ }\edtext{\edtext{\skt{सिरा\-व\-र्ण्ण\-न\-मेव}}{\lemma{\skt{सिरावर्ण्णनम्}}\Afootnote{A K N; \skt{सिवर्ण्णन}\ H\textsuperscript{ac}; \skt{सिरावर्ण्णन}\ H\textsuperscript{pc}}}\skt{ च}}{\lemma{\skt{एव च}}\Afootnote{A K N; \skt{मेव च}\ H}}\skt{~। }\pend
\pstart\versepar \skt{सिराव्यधो धमनीनां गर्भिण्या व्याकृतिस्तथा~।}\pend

\pstart\versepar \ledsidenote{\passref{1.3.17}}\skt{निर्दिष्टानि }\edtext{\skt{दशैतानि}}{\lemma{\skt{दशैतानि}}\Afootnote{A K N; \skt{दशौतानि}\ H}}\skt{ शारीराणि }\edtext{\skt{महर्षिणां}}{\lemma{\skt{महर्षिणां}}\Afootnote{K; \skt{महर्षिणा}\ A N H}}\skt{~। }\pend
\pstart\versepar \skt{विज्ञानार्थं शरीरस्य भिषजां }\edtext{\skt{योगिनामपि}}{\lemma{\skt{योगिनामपि}}\Afootnote{H; \skt{योगिनामपि}\ A N K}}\skt{~॥}\pend

\pstart\versepar \ledsidenote{\passref{1.3.18}}\skt{द्विव्रणीयो }\edtext{\skt{व्रणे}}{\lemma{\skt{व्रणे}}\Afootnote{K N; \skt{व्रणः}\ A; \skt{व्रण}\ H}}\skt{ सद्यो भग्नानां वातरोगिकं~। }\pend
\pstart\versepar \skt{}\edtext{\skt{महा\-वा\-ति\-क\-म\-र्शान्सि}}{\lemma{\skt{महा\appbrk{}वा\appbrk{}ति\appbrk{}क\appbrk{}म\appbrk{}र्शान्सि}}\Afootnote{A H; \skt{महा\appbrk{}वा\appbrk{}ति\appbrk{}क\appbrk{}म\appbrk{}र्शान्सि}\ K N}}\skt{ साश्मरिश्च भगन्दरः~।}\pend

\pstart\versepar \ledsidenote{\passref{1.3.19}}\skt{कुष्ठानां महताञ्चापि }\edtext{\skt{मैहि\-क\-म्पै\-डि\-क\-न्तथा}}{\lemma{\skt{मैहि\appbrk{}क\appbrk{}म्पै\appbrk{}डि\appbrk{}कन्}}\Afootnote{A K; \skt{मैहि\appbrk{}क\appbrk{}म्पै\appbrk{}हि\appbrk{}कन्}\ N; \skt{मैहिकं पैदिकन्}\ H}}\skt{~। }\pend
\pstart\versepar \skt{}\edtext{\skt{मधु\-मे\-हि\-चि\-कित्सा च}}{\lemma{\skt{मधु\appbrk{}मे\appbrk{}हि\appbrk{}चि\appbrk{}कित्सा च}}\Afootnote{K N; \skt{मधु\appbrk{}मे\appbrk{}ह\appbrk{}चि\appbrk{}कित्सा च}\ A H}}\skt{ तथा }\edtext{\skt{चोदरिणामपि}}{\lemma{\skt{चोदरिणामपि}}\Afootnote{H; \skt{चोदरिणामपि}\ A K N}}\skt{~।}\pend

\pstart\versepar \ledsidenote{\passref{1.3.20}}\skt{मूढ\-ग\-र्भ\-चि\-कित्सा च }\edtext{\skt{विद्र\-थी\-ना\-म्वि\-स\-र्पिणां}}{\lemma{\skt{विद्रथीनाम्}}\Afootnote{K; \skt{विद्रधीनां}\ A H; \skt{विद्रधी}\ N}}\skt{~। }\pend
\pstart\versepar \skt{}\edtext{\edtext{\skt{ग्रन्थी\-ना\-म्वृ\-द्ध्यु\-दं\-शानां}}{\lemma{\skt{ग्रन्थीनाम्}}\Afootnote{K N H; \skt{ग्रन्थि}\ A}}}{\lemma{\skt{वृद्ध्युदंशानां}}\Afootnote{K H; \skt{वृद्ध्युपदंशानां}\ A; \skt{वृद्ध्यपदंशानां}\ N}}\skt{ तथा च }\edtext{\skt{क्षुद्ररोगिकं}}{\lemma{\skt{क्षुद्ररोगिकं}}\Afootnote{K N H; \skt{क्षुद्ररोगिणाम्}\ A}}\skt{~।}\pend

\pstart\versepar \ledsidenote{\passref{1.3.21}}\skt{शूक\-दो\-ष\-चि\-कित्सा च तथा च मुखरोगिणां~। }\pend
\pstart\versepar \skt{शोफ\-स्या\-ना\-ग\-ता\-नाञ्च निषेधो मिश्रकन्तथा~।}\pend

\pstart\versepar \ledsidenote{\passref{1.3.22}}\edtext{\skt{व्याजीकरञ्च}}{\lemma{\skt{व्याजीकरञ्च}}\Afootnote{K N H; \skt{वाजीकरं च}\ A}}\skt{ यत्क्षीणे }\edtext{\skt{सर्वाबाधशमो}}{\lemma{\skt{सर्वाबाधशमो}}\Afootnote{A K H; \skt{सर्व्वाधसमो}\ N}}\skt{ पि च~। }\pend
\pstart\versepar \skt{}\edtext{\skt{मेधा\-यु\-ष्क\-र\-ण\-ञ्चापि}}{\lemma{\skt{चापि}}\Afootnote{A K N; \skt{चापि …}\ H}}\skt{ }\edtext{\skt{स्वभावाच्च}}{\lemma{\skt{स्वभावाच्च}}\Afootnote{K N H; \skt{स्वभावव्या}\ A}}\skt{ }\edtext{\skt{निवारणं}}{\lemma{\skt{निवारणं}}\Afootnote{K N H; \skt{धिवारणम्}\ A}}\skt{~।}\pend

\pstart\versepar \ledsidenote{\passref{1.3.23}}\skt{निवृ\-त्त\-स\-न्ता\-प\-करं कीर्तितञ्च }\edtext{\skt{रसायनं}}{\lemma{\skt{रसायनं}}\Afootnote{A K N; \skt{रसायनं~। म्}\ H}}\skt{~। }\pend
\pstart\versepar \skt{}\edtext{\skt{स्नेहोपयोगिकः}}{\lemma{\skt{स्नेहोपयोगिकः}}\Afootnote{K N; \skt{स्नेहोपयौगिकः}\ A; \skt{नेहोपयोगिकः}\ H}}\skt{ }\edtext{\skt{स्वेदाः}}{\lemma{\skt{स्वेदाः}}\Afootnote{K N H; \skt{स्वेदो}\ A}}\skt{ वमने }\edtext{\skt{सविरेचने}}{\lemma{\skt{सविरेचने}}\Afootnote{A K N H\textsuperscript{ac}; \skt{स विरेचने~॥ भ्र}\ H\textsuperscript{pc}}}\skt{~।}\pend

\pstart\versepar \ledsidenote{\passref{1.3.24}}\edtext{\edtext{\skt{तयो\-र्व्या\-प\-च्चि\-कित्सा}}{\lemma{\skt{तयोर्}}\Afootnote{A K H; \skt{ये योर्}\ N}}\skt{ च}}{\lemma{\skt{व्यापच्चिकित्सा च}}\Afootnote{A N H; \skt{व्यापच्चि}\ K}}\skt{ }\edtext{\skt{नेत्र\-ब\-स्ति\-वि\-भा\-गिकः}}{\lemma{\skt{नेत्र\appbrk{}ब\appbrk{}स्ति\appbrk{}वि\appbrk{}भा\appbrk{}गिकः}}\Afootnote{A H; \emph{om.}\ K N}}\skt{~। }\pend
\pstart\versepar \skt{}\edtext{\skt{नेत्र\-ब\-स्ति\-वि\-भा\-गिकं}}{\lemma{\skt{नेत्र\appbrk{}ब\appbrk{}स्ति\appbrk{}वि\appbrk{}भा\appbrk{}गिकं}}\Afootnote{N H; \emph{om.}\ A K}}\skt{~। }\edtext{\edtext{\skt{नेत्र\-ब\-स्ति\-वि\-प\-त्सि\-द्धि\-स्तथा}}{\lemma{\skt{नेत्र\appbrk{}ब\appbrk{}स्ति\appbrk{}वि\appbrk{}प\appbrk{}त्सि\appbrk{}द्धिस्}}\Afootnote{A N H; \emph{om.}\ K}}}{\lemma{\skt{तथा}}\Afootnote{A N H; \emph{om.}\ K}}\skt{ }\edtext{\skt{चोत्तरबस्तिकं}}{\lemma{\skt{चोत्तरबस्तिकं}}\Afootnote{N H; \skt{चोत्तरबस्तिकः}\ A; \emph{om.}\ K}}\skt{~॥}\pend

\pstart\versepar \ledsidenote{\passref{1.3.25}}\edtext{\skt{निरू\-ह\-क्र\-म\-सं\-ज्ञञ्च}}{\lemma{\skt{निरू\appbrk{}ह\appbrk{}क्र\appbrk{}म\appbrk{}सं\appbrk{}ज्ञञ्च}}\Afootnote{N H; \skt{निरू\appbrk{}ह\appbrk{}क्र\appbrk{}म\appbrk{}स\appbrk{}ञ्ज्ञश्च}\ A}}\skt{ }\edtext{\skt{तथै\-वा\-तु\-र\-सं\-ज्ञितं}}{\lemma{\skt{तथै\appbrk{}वा\appbrk{}तु\appbrk{}र\appbrk{}सं\appbrk{}ज्ञितं}}\Afootnote{N H; \skt{तथै\appbrk{}वा\appbrk{}तु\appbrk{}र\appbrk{}स\appbrk{}ञ्ज्ञकः}\ A}}\skt{~। }\pend
\pstart\versepar \skt{}\edtext{\edtext{\edtext{\skt{धूम\-न\-स्तो\-वि\-धि\-श्चो\-ग्र्य\-श्च\-त्वा\-रिं\-श\-दिति}}{\lemma{\skt{धूम\appbrk{}न\appbrk{}स्तो\appbrk{}वि\appbrk{}धिश्}}\Afootnote{N H; \skt{धूम\appbrk{}न\appbrk{}स्य\appbrk{}वि\appbrk{}धिश्}\ A}}}{\lemma{\skt{चोग्र्यश्}}\Afootnote{H; \skt{चान्त्य}\ A; \skt{चाग्र्यश्}\ N}}}{\lemma{\skt{चत्वारिंशदिति}}\Afootnote{N H; \skt{श्चत्वारिंशदिति}\ A}}\skt{ स्मृताः~॥}\pend

\pstart\versepar \ledsidenote{\passref{1.3.26}}\skt{प्रायश्चित्तं प्रशमनं चिकित्सा शान्तिकर्म्म च~। }\pend
\pstart\versepar \skt{}\edtext{\skt{पर्यायास्तस्य}}{\lemma{\skt{तस्य}}\Afootnote{A N; \skt{स्तस्य}\ H}}\skt{ }\edtext{\edtext{\skt{निर्दे\-शा\-च्चि\-कि\-त्सा\-स्था\-न\-मु\-च्यते}}{\lemma{\skt{निर्दे\appbrk{}शा\appbrk{}च्चि\appbrk{}कि\appbrk{}त्सा\appbrk{}स्था\appbrk{}नम्}}\Afootnote{N; \skt{निर्दे\appbrk{}शा\appbrk{}च्चि\appbrk{}कि\appbrk{}त्सा\appbrk{}स्थान}\ A H}}}{\lemma{\skt{उच्यते}}\Afootnote{N; \skt{मुच्यते}\ A H}}\skt{~॥}\pend

\pstart\versepar \ledsidenote{\passref{1.3.27}}\skt{अन्नस्य रक्षा विज्ञानं स्थाव\-र\-स्ये\-त\-रस्य च~। }\pend
\pstart\versepar \skt{}\edtext{\skt{सर्प्पदष्टविषं}}{\lemma{\skt{सर्प्पदष्टविषं}}\Afootnote{H; \skt{सर्पदष्टविष}\ A; \skt{सर्प्पदष्टं विष}\ N}}\skt{ ज्ञानं तस्यैव च चिकित्सितं~॥}\pend

\pstart\versepar \ledsidenote{\passref{1.3.28}}\edtext{\edtext{\skt{दुन्दु\-भे\-र्म्मू\-षि\-का\-णाञ्च}}{\lemma{\skt{दुन्दुभेर्}}\Afootnote{\skt{दुन्दुभे}\ A; \skt{दौन्दुभिर्}\ N H}}}{\lemma{\skt{म्मूषिकाणाञ्च}}\Afootnote{N H; \skt{र्मूषिकाणां च}\ A}}\skt{ कीटानाङ्कल्प एव च~। }\pend
\pstart\versepar \skt{अष्टौ कल्पाः समाख्याता }\edtext{\skt{विष\-भे\-ष\-ज\-क\-ल्प\-नात्}}{\lemma{\skt{विष\appbrk{}भे\appbrk{}ष\appbrk{}ज\appbrk{}क\appbrk{}ल्प\appbrk{}नात्}}\Afootnote{A H; \skt{विष\appbrk{}भे\appbrk{}ष\appbrk{}ज\appbrk{}क\appbrk{}ल्पा\appbrk{}नात्}\ N}}\skt{~॥}\pend

\pstart\versepar \ledsidenote{\passref{1.3.29}}\edtext{\skt{सविंशमध्याय}}{\lemma{\skt{सविंशमध्याय}}\Afootnote{H; \skt{अध्यायानां}\ A; \skt{सविंशमध्याय}\ N}}\skt{ }\edtext{\edtext{\skt{शत\-मे\-व\-मे\-त\-दु\-दी\-रितं}}{\lemma{\skt{शतम्}}\Afootnote{A N; \emph{om.}\ H}}}{\lemma{\skt{एवम्}}\Afootnote{N; \skt{विंशमेवम्}\ A; \skt{शतमेवम्}\ H}}\skt{~। }\pend
\pstart\versepar \skt{अतः परं }\edtext{\skt{स्वनाम्नान्तु}}{\lemma{\skt{स्वनाम्नान्तु}}\Afootnote{N H; \skt{स्वनाम्नैव}\ A}}\skt{ }\edtext{\edtext{\edtext{\skt{तन्त्र\-मु\-त्त\-र\-मु\-च्यते}}{\lemma{\skt{तन्त्रम्}}\Afootnote{A N; \skt{तन्त्र}\ H}}}{\lemma{\skt{उत्तरम्}}\Afootnote{A N; \skt{मुत्तर}\ H}}}{\lemma{\skt{उच्यते}}\Afootnote{A N; \skt{मुच्यते}\ H}}\skt{~॥}\pend

\pstart\versepar \ledsidenote{\passref{1.3.30}}\skt{अधिकृत्य कृतं }\edtext{\skt{यस्मा\-त्त\-न्त्र\-मे\-त\-दु\-प\-द्र\-वान्}}{\lemma{\skt{यस्मात्तन्त्रम्}}\Afootnote{N H; \skt{यस्मात्तन्त्र}\ A}}\skt{~। }\pend
\pstart\versepar \skt{औपद्रविक }\edtext{\skt{इत्येष}}{\lemma{\skt{इत्येष}}\Afootnote{H; \skt{इत्येष}\ A N}}\skt{ }\edtext{\skt{तस्या\-ग्र\-त्वा\-न्नि\-रु\-च्यते}}{\lemma{\skt{तस्याग्रत्वान्}}\Afootnote{N H; \skt{तस्याग्र्यत्वान्}\ A}}\skt{~॥}\edtext{}{\lemma{\skt{एत\appbrk{}दु\appbrk{}प\appbrk{}द्र\appbrk{}वान्~। औपद्रविक}}\Afootnote{N H; \skt{मेत\appbrk{}दु\appbrk{}प\appbrk{}द्र\appbrk{}वान्~। औपद्रविक}\ A}}\pend

\pstart\versepar \ledsidenote{\passref{1.3.31}}\skt{सन्धौ }\edtext{\skt{वर्त्मसु}}{\lemma{\skt{वर्त्मसु}}\Afootnote{N H; \skt{वर्त्मनि}\ A}}\skt{ शुक्ले च कृष्णे सर्वत्र दृष्टिषु~। }\pend
\pstart\versepar \skt{}\edtext{\edtext{\skt{संवि\-ज्ञा\-ना\-र्थ\-म\-ध्यायाः}}{\lemma{\skt{संविज्ञानार्थम्}}\Afootnote{N H; \skt{संविज्ञानार्थ}\ A}}}{\lemma{\skt{अध्यायाः}}\Afootnote{N H; \skt{मध्याया}\ A}}\skt{ गदानां तु प्रति प्रति~॥}\pend

\pstart\versepar \ledsidenote{\passref{1.3.32}}\edtext{\skt{चिकि\-त्सा\-प्र\-ति\-भा\-गीयो}}{\lemma{\skt{चिकि\appbrk{}त्सा\appbrk{}प्र\appbrk{}ति\appbrk{}भा\appbrk{}गीयो}}\Afootnote{N H; \skt{चिकि\appbrk{}त्सा\appbrk{}प्र\appbrk{}वि\appbrk{}भा\appbrk{}गीयो}\ A}}\skt{ }\edtext{\skt{वाता\-भि\-ष्य\-न्द\-वा\-रणः}}{\lemma{\skt{वाता\appbrk{}भि\appbrk{}ष्य\appbrk{}न्द\appbrk{}वा\appbrk{}रणः}}\Afootnote{A H; \skt{वाता\appbrk{}भि\appbrk{}ष्य\appbrk{}न्द\appbrk{}वा\appbrk{}रिणः}\ N}}\skt{~। }\pend
\pstart\versepar \skt{पैत्तस्य }\edtext{\skt{श्लेष्मलस्यापि}}{\lemma{\skt{श्लेष्मलस्यापि}}\Afootnote{N H; \skt{श्लैष्मिकस्यापि}\ A}}\skt{ }\edtext{\edtext{\skt{रौधिरस्य}}{\lemma{\skt{रौधिर्}}\Afootnote{H; \skt{रौधि}\ A; \skt{रोधि}\ N}}}{\lemma{\skt{अस्य}}\Afootnote{H; \skt{रस्य}\ A N}}\skt{ तथैव च~॥}\pend

\pstart\versepar \ledsidenote{\passref{1.3.33}}\edtext{\skt{लेख\-रो\-ग\-नि\-रो\-धश्च}}{\lemma{\skt{लेख\appbrk{}रो\appbrk{}ग\appbrk{}नि\appbrk{}रो\appbrk{}धश्च}}\Afootnote{H; \skt{लेख्य\appbrk{}भे\appbrk{}द्य\appbrk{}नि\appbrk{}षेधौ च}\ A; \skt{लेख्य\appbrk{}रो\appbrk{}ग\appbrk{}नि\appbrk{}षे\appbrk{}धश्च}\ N}}\skt{ छेद्यानां वर्त्मदृष्टिषु~। }\pend
\pstart\versepar \skt{}\edtext{\skt{क्रिया\-क\-ल्पो\-भि\-घा\-तश्च}}{\lemma{\skt{क्रिया\appbrk{}क\appbrk{}ल्पो\appbrk{}भि\appbrk{}घा\appbrk{}तश्च}}\Afootnote{H; \skt{क्रिया\appbrk{}क\appbrk{}ल्पोऽ\appbrk{}भि\appbrk{}घा\appbrk{}तश्च}\ A; \skt{क्रियाकल्पो भिघातश्च}\ N}}\skt{ }\edtext{\edtext{\skt{त्थास्तच्चिकित्सितं}}{\lemma{\skt{त्थास्}}\Afootnote{H; \skt{कर्णोत्था}\ A; \skt{कर्ण्णोत्थास्}\ N}}}{\lemma{\skt{तच्चिकित्सितं}}\Afootnote{H; \skt{स्तच्चिकित्सितम्}\ A; \skt{तच्चिकित्सितं}\ N}}\skt{~॥}\pend

\pstart\versepar \ledsidenote{\passref{1.3.34}}\skt{घ्राणोत्थानाञ्च विज्ञानं }\edtext{\skt{तद्ग\-द\-प्र\-ति\-षे\-धनं}}{\lemma{\skt{तद्ग\appbrk{}द\appbrk{}प्र\appbrk{}ति\appbrk{}षे\appbrk{}धनं}}\Afootnote{H; \skt{तद्ग\appbrk{}द\appbrk{}प्र\appbrk{}ति\appbrk{}षे\appbrk{}ध\appbrk{}नम्}\ A; \skt{तद्गदं प्रतिषेधनं}\ N}}\skt{~। }\pend
\pstart\versepar \skt{}\edtext{\skt{प्रति\-श्या\-य\-नि\-षे\-धश्च}}{\lemma{\skt{प्रति\appbrk{}श्या\appbrk{}य\appbrk{}नि\appbrk{}षे\appbrk{}धश्च}}\Afootnote{A H; \skt{प्रति\appbrk{}स्या\appbrk{}य\appbrk{}नि\appbrk{}षेध}\ N}}\skt{ }\edtext{\skt{शिरो\-ग\-त\-वि\-जा\-ननं}}{\lemma{\skt{शिरो\appbrk{}ग\appbrk{}त\appbrk{}वि\appbrk{}जा\appbrk{}ननं}}\Afootnote{N H; \skt{शिरो\appbrk{}ग\appbrk{}द\appbrk{}वि\appbrk{}वे\appbrk{}च\appbrk{}नम्}\ A}}\skt{~॥}\pend

\pstart\versepar \ledsidenote{\passref{1.3.35}}\skt{चिकित्सा }\edtext{\skt{तद्गदानाञ्च}}{\lemma{\skt{तद्गदानाञ्च}}\Afootnote{N H; \skt{तद्गदानां च}\ A}}\skt{ }\edtext{\skt{शालाक्ये}}{\lemma{\skt{शालाक्ये}}\Afootnote{H; \skt{शालाक्यं}\ A; \skt{शालाक्य}\ N}}\skt{ }\edtext{\skt{तन्त्र}}{\lemma{\skt{तन्त्र}}\Afootnote{N H; \skt{तन्त्रम्}\ A}}\skt{ उच्यते~। }\pend
\pstart\versepar \skt{नव\-ग्र\-हा\-कृ\-ति\-ज्ञानं स्कन्दस्य च निषेधनं~॥}\pend

\pstart\versepar \ledsidenote{\passref{1.3.36}}\skt{अप\-स्मा\-र\-श\-कु\-न्योश्च रेवत्याश्च पुनः पृथक्~। }\pend
\pstart\versepar \skt{}\edtext{\skt{पूत\-ना\-या\-स्त\-था\-न्धाया}}{\lemma{\skt{तथान्धाया}}\Afootnote{N H; \skt{तथाऽन्धायाः}\ A}}\skt{ }\edtext{\skt{मण्डिका}}{\lemma{\skt{मण्डिका}}\Afootnote{N H; \emph{om.}\ A}}\skt{ }\edtext{\skt{शीतपूतना}}{\lemma{\skt{शीतपूतना}}\Afootnote{N H; \skt{शीत\appbrk{}पू\appbrk{}त\appbrk{}न\appbrk{}म\appbrk{}ण्डिका}\ A}}\skt{~॥}\pend

\pstart\versepar \ledsidenote{\passref{1.3.37}}\skt{नैग\-मे\-ष\-चि\-कित्सा च ग्रहोत्पत्तिः }\edtext{\skt{सयोनिजाः}}{\lemma{\skt{सयोनिजाः}}\Afootnote{N H; \skt{सयोनिजा}\ A}}\skt{~। }\pend
\pstart\versepar \skt{कुमा\-र\-त\-न्त्र\-मि\-त्ये\-त\-च्छा\-री\-रेषु च कीर्त्तितं~॥}\pend

\pstart\versepar \ledsidenote{\passref{1.3.38}}\skt{ज्वरा\-ति\-सा\-र\-शो\-षाणां }\edtext{\skt{गुल्म\-हृ\-द्रो\-गि\-णा\-मपि}}{\lemma{\skt{गुल्म\appbrk{}हृ\appbrk{}द्रो\appbrk{}गि\appbrk{}णाम्}}\Afootnote{A H; \skt{गुल्म\appbrk{}हृ\appbrk{}द\appbrk{}रो\appbrk{}गि\appbrk{}णाम्}\ N}}\skt{~। }\pend
\pstart\versepar \skt{}\edtext{\skt{पाण्डूनां}}{\lemma{\skt{पाण्डूनां}}\Afootnote{A H; \skt{पाण्डूना}\ N}}\skt{ }\edtext{\skt{रक्तपित्तस्य}}{\lemma{\skt{रक्तपित्तस्य}}\Afootnote{A H; \skt{न्रक्तपित्तस्य}\ N}}\skt{ मूर्च्छायाः पानजाश्च ये~॥}\pend

\pstart\versepar \ledsidenote{\passref{1.3.39}}\skt{तृष्णा\-या\-श्छ\-र्द्दि\-हि\-क्कानां निषेधः श्वासकासयोः~। }\pend
\pstart\versepar \skt{स्वर\-भे\-द\-चि\-कित्सा च }\edtext{\skt{क्रिम्युदावर्तिनोः}}{\lemma{\skt{क्रिम्युदावर्तिनोः}}\Afootnote{N H; \skt{कृम्युदावर्तिनोः}\ A}}\skt{ पृथक्~॥}\pend

\pstart\versepar \ledsidenote{\passref{1.3.40}}\edtext{\skt{विमू\-चि\-का\-रा\-च\-क\-यो\-र्म्मू\-त्रा\-घा\-त\-वि\-कृ\-च्छ्रयोः}}{\lemma{\skt{विमू\appbrk{}चि\appbrk{}का\appbrk{}रा\appbrk{}च\appbrk{}क\appbrk{}यो\appbrk{}र्म्मू\appbrk{}त्रा\appbrk{}घा\appbrk{}त\appbrk{}वि\appbrk{}कृ\appbrk{}च्छ्रयोः}}\Afootnote{H; \skt{विसू\appbrk{}चि\appbrk{}का\appbrk{}रो\appbrk{}च\appbrk{}क\appbrk{}यो\appbrk{}र्मू\appbrk{}त्रा\appbrk{}घा\appbrk{}त\appbrk{}वि\appbrk{}कृ\appbrk{}च्छ्रयोः}\ A; \skt{विषू\appbrk{}चि\appbrk{}का\appbrk{}रो\appbrk{}च\appbrk{}कयो मूर्त्ता\appbrk{}घा\appbrk{}त\appbrk{}वि\appbrk{}कृ\appbrk{}च्छ्रयोः}\ N}}\skt{~। }\pend
\pstart\versepar \skt{इति कायचिकित्सायाः शेषमत्र प्रकीर्त्तितं॥}\pend

\pstart\versepar \ledsidenote{\passref{1.3.41}}\skt{अमा\-नु\-ष\-नि\-षे\-धश्च }\edtext{\skt{तथापस्मारिको}}{\lemma{\skt{तथापस्मारिको}}\Afootnote{A N; \skt{तथापस्मारिको}\ H}}\skt{ }\edtext{\skt{परः}}{\lemma{\skt{परः}}\Afootnote{A N; \skt{परः}\ H}}\skt{~। }\pend
\pstart\versepar \skt{उन्मा\-द\-प्र\-ति\-षे\-धश्च भूतविद्या निरुच्यते~॥}\pend

\pstart\versepar \ledsidenote{\passref{1.3.42}}\skt{रसभेदाः स्वस्थवृत्तं }\edtext{\edtext{\skt{युक्त\-य\-स्ता\-न्त्रि\-काश्च}}{\lemma{\skt{युक्तयस्}}\Afootnote{N H; \skt{युक्तय}\ A}}\skt{ याः}}{\lemma{\skt{तान्त्रिकाश्च याः}}\Afootnote{N H; \skt{स्तान्त्रिकाश्च याः}\ A}}\skt{~। }\pend
\pstart\versepar \skt{दोषभेदा इति ज्ञेया }\edtext{\edtext{\skt{अध्यायास्तत्र}}{\lemma{\skt{अध्यायास्}}\Afootnote{N H; \skt{अध्याया}\ A}}}{\lemma{\skt{तत्र}}\Afootnote{H; \skt{स्तन्त्र}\ A; \skt{तन्त्र}\ N}}\skt{ भूषणाः~॥}\pend

\pstart\versepar \ledsidenote{\passref{1.3.43}}\edtext{\skt{श्रेष्ठ\-त्वा\-दु\-त्त\-त्तरं }\edtext{\edtext{\skt{ह्येत\-त्त\-न्त्र\-मा\-हु\-र्म्म\-ह\-र्षयः}}{\lemma{\skt{आहुर्}}\Afootnote{A N; \skt{आहु}\ H}}}{\lemma{\skt{म्महर्षयः}}\Afootnote{A N; \skt{र्म्महर्षयः}\ H}}}{\lemma{\skt{श्रेष्ठ\appbrk{}त्वा\appbrk{}दु\appbrk{}त्त\appbrk{}त्तरं ह्य्}}\Afootnote{\skt{श्रेष्ठत्वादुत्तरं ह्य्}\ A; \skt{स्रेष्ठत्वादुत्तरं ह्य्}\ N; \skt{श्रेष्ठ\appbrk{}त्वा\appbrk{}दु\appbrk{}त्त\appbrk{}त्तरं ह्य्}\ H}}\skt{~। }\pend
\pstart\versepar \skt{}\edtext{\skt{बह्वर्थं}}{\lemma{\skt{बह्वर्थं}}\Afootnote{H; \skt{बह्वर्थ}\ A N}}\skt{ }\edtext{\skt{सङ्ग्र\-हा\-च्छ्रे\-ष्ठ\-मु\-त्तरं}}{\lemma{\skt{सङ्ग्रहाच्छ्रेष्ठम्}}\Afootnote{A H; \skt{स्रंग्रहाच्छ्रेष्ठम्}\ N}}\skt{ वापि पश्चिमं~॥}\pend

\pstart\versepar \ledsidenote{\passref{1.3.44}}\edtext{\skt{शालाक्यतन्त्रं}}{\lemma{\skt{शालाक्यतन्त्रं}}\Afootnote{A H; \skt{शालाक्यतन्त्र्कौ}\ N}}\skt{ }\edtext{\skt{कौमारं}}{\lemma{\skt{कौमारं}}\Afootnote{A H; \skt{मार}\ N}}\skt{ }\edtext{\skt{चिकित्सा}}{\lemma{\skt{चिकित्सा}}\Afootnote{A H; \skt{चिकित्स्का}\ N}}\skt{ }\edtext{\skt{कायिकी च याः}}{\lemma{\skt{कायिकी च याः}}\Afootnote{A H; \skt{यिकी च या}\ N}}\skt{~। }\pend
\pstart\versepar \skt{भूतविद्येति चत्वारि तन्त्रे तूत्तरसंज्ञिते~॥}\pend

\pstart\versepar \ledsidenote{\passref{1.3.45}}\edtext{\skt{व्याजीकरा}}{\lemma{\skt{व्याजीकरा}}\Afootnote{H; \skt{वाजीकरं}\ A; \skt{वाजीकराश्}\ N}}\skt{ }\edtext{\skt{चकित्सा च}}{\lemma{\skt{चकित्सा च}}\Afootnote{N; \skt{चिकित्सासु}\ A; \skt{चकित्सा च …}\ H}}\skt{ रसा\-य\-न\-वि\-धि\-स्तथा~। }\pend
\pstart\versepar \skt{विषतन्त्रं पुनः }\edtext{\skt{कल्पाः}}{\lemma{\skt{कल्पाः}}\Afootnote{A H; \skt{कल्पः}\ N}}\skt{ शल्यज्ञानं समन्ततः~॥}\pend

\pstart\versepar \ledsidenote{\passref{1.3.46}}\edtext{\edtext{\skt{इत्य\-ष्टा\-ङ्ग\-मि\-द\-न्त\-न्त्र\-मा\-दि\-दे\-वे\-प्र\-का\-शितं}}{\lemma{\skt{तन्त्रम्}}\Afootnote{N H; \skt{तन्त्र}\ A}}}{\lemma{\skt{आदि\appbrk{}दे\appbrk{}वे\appbrk{}प्र\appbrk{}का\appbrk{}शितं}}\Afootnote{H; \skt{मादि\appbrk{}दे\appbrk{}व\appbrk{}प्र\appbrk{}का\appbrk{}शि\appbrk{}तम्}\ A; \skt{आदि\appbrk{}दे\appbrk{}व\appbrk{}प्र\appbrk{}का\appbrk{}शितं}\ N}}\skt{~। }\pend
\pstart\versepar \skt{}\edtext{\skt{विधिनाधीत्य}}{\lemma{\skt{विधिनाधीत्य}}\Afootnote{A N H\textsuperscript{ac}; \skt{ञ्विधिनाधीत्य}\ H\textsuperscript{pc}}}\skt{ }\edtext{\skt{युञ्जाना}}{\lemma{\skt{युञ्जाना}}\Afootnote{A N; \skt{युञ्जाना …}\ H}}\skt{ भवन्ति प्राणदा भुवि~॥}\pend

\pstart \ledsidenote{\passref{1.3.47}}\edtext{\edtext{\skt{एत\-द\-व\-स्य\-म\-ध्येयं}}{\lemma{\skt{एतदवस्यम्}}\Afootnote{\skt{एतद्ध्यवश्यम्}\ A N; \skt{एतदवस्य}\ H}}}{\lemma{\skt{अध्येयं}}\Afootnote{A N; \skt{मध्येयं}\ H}}\skt{ अधीत्य च }\edtext{\edtext{\skt{कर्म्मा\-प्य\-व\-स्य\-मु\-पा\-सि\-तव्यं}}{\lemma{\skt{अवस्यम्}}\Afootnote{N H; \skt{अवश्यम्}\ A}}}{\lemma{\skt{उपासितव्यं}}\Afootnote{A H; \skt{उपा\appbrk{}या\appbrk{}सि\appbrk{}त\appbrk{}व्यम्}\ N}}\skt{ उभयज्ञो हि भिषग्राजार्हो }\edtext{\skt{भवति}}{\lemma{\skt{भवति}}\Afootnote{A H; \skt{भवति~॥ भ~॥}\ N}}\skt{~॥}\pend

\pstart\versepar \ledsidenote{\passref{1.3.48}}\edtext{\skt{भवन्ति}}{\lemma{\skt{भवन्ति}}\Afootnote{A; \emph{om.}\ N H}}\skt{ चात्र॥ }\pend
\pstart\versepar \skt{यस्तु केवलशास्त्रज्ञः कर्म्म\-स्व\-प\-रि\-नि\-ष्ठतः~। }\pend
\pstart\versepar \skt{स मुह्यत्यातुरं प्राप्य प्राप्य भीरुरिवाहवं~॥}\edtext{}{\lemma{\skt{चात्र॥ यस्तु}}\Afootnote{A; \skt{यस्तु}\ N H}}\edtext{}{\lemma{\skt{कर्म्म\appbrk{}स्व\appbrk{}प\appbrk{}रि\appbrk{}नि\appbrk{}ष्ठतः~। स}}\Afootnote{N H; \skt{कर्मस्व परिनिष्ठितः~। स}\ A}}\pend

\pstart\versepar \ledsidenote{\passref{1.3.49}}\edtext{\skt{यस्तु}}{\lemma{\skt{यस्तु}}\Afootnote{H; \skt{यस्तु}\ A N}}\skt{ कर्म्मसु निष्णातो धाष्ट्या\-च्छा\-स्त्र\-ब\-हि\-ष्कृतः~। }\pend
\pstart\versepar \skt{स सत्सु पूजान्नाप्नोति }\edtext{\skt{वधञ्चार्च्छति}}{\lemma{\skt{चार्च्छति}}\Afootnote{N H; \skt{चर्च्छति}\ A}}\skt{ राजतः~॥}\edtext{}{\lemma{\skt{धाष्ट्या\appbrk{}च्छा\appbrk{}स्त्र\appbrk{}ब\appbrk{}हि\appbrk{}ष्कृतः~। स}}\Afootnote{N H; \skt{धार्ष्ट्या\appbrk{}च्छा\appbrk{}स्त्र\appbrk{}ब\appbrk{}हि\appbrk{}ष्कृतः~। स}\ A}}\pend

\pstart\versepar \ledsidenote{\passref{1.3.50}}\edtext{\edtext{\skt{उभा\-वे\-ता\-व\-नि\-पु\-णा\-व\-स\-मर्थौ}}{\lemma{\skt{अनि\appbrk{}पु\appbrk{}णा\appbrk{}व\appbrk{}सम्}}\Afootnote{H; \skt{अनिपुणावस}\ A N}}}{\lemma{\skt{अर्थौ}}\Afootnote{H; \skt{मर्थौ}\ A; \skt{मर्थो}\ N}}\skt{ }\edtext{\skt{चिकित्सित्तं}}{\lemma{\skt{चिकित्सित्तं}}\Afootnote{H; \skt{स्वकर्मणि}\ A; \skt{चिकित्सित्तुं}\ N}}\skt{~। }\pend
\pstart\versepar \skt{}\edtext{\skt{अर्ध\-वे\-द\-ध\-रा\-वे\-ता\-वे\-क\-प\-क्षा\-विव}}{\lemma{\skt{अर्ध\appbrk{}वे\appbrk{}द\appbrk{}ध\appbrk{}रा\appbrk{}वे\appbrk{}ता\appbrk{}वे\appbrk{}क\appbrk{}प\appbrk{}क्षाव्}}\Afootnote{H; \skt{अर्ध\appbrk{}वे\appbrk{}द\appbrk{}ध\appbrk{}रा\appbrk{}वे\appbrk{}ता\appbrk{}वे\appbrk{}क\appbrk{}प\appbrk{}क्षाव्}\ A N}}\skt{ }\edtext{\skt{द्विजौ}}{\lemma{\skt{द्विजौ}}\Afootnote{A H; \skt{द्विजो}\ N}}\skt{~॥}\pend

\pstart\versepar \ledsidenote{\passref{1.3.51}}\edtext{\skt{ओष\-ध्योऽ\-मृ\-त\-क\-ल्पास्तु}}{\lemma{\skt{ओषध्यो}}\Afootnote{A H; \skt{औषध्यो}\ N}}\skt{ शस्त्रा\-श\-नि\-वि\-षो\-पमाः~। }\pend
\pstart\versepar \skt{}\edtext{\edtext{\edtext{\skt{अज्ञे\-नो\-प\-हि\-ता\-स्ता\-स्यु\-स्मा\-त्त\-म्प\-रि\-व\-र्ज\-येत्}}{\lemma{\skt{अज्ञे\appbrk{}नो\appbrk{}प\appbrk{}हि\appbrk{}ता\appbrk{}स्तास्}}\Afootnote{\skt{भव\appbrk{}न्त्य\appbrk{}ज्ञै\appbrk{}रु\appbrk{}प\appbrk{}हृ\appbrk{}तास्}\ A; \skt{अज्ञ\appbrk{}नो\appbrk{}प\appbrk{}हि\appbrk{}तास्ता}\ N; \skt{अज्ञे\appbrk{}नो\appbrk{}प\appbrk{}हि\appbrk{}तास्ता}\ H}}}{\lemma{\skt{युस्मात्तम्}}\Afootnote{\skt{तस्मादे}\ A; \skt{स्युस्तस्मात्तम्}\ N; \skt{स्युस्मात्तम्}\ H}}}{\lemma{\skt{परिवर्जयेत्}}\Afootnote{H; \skt{तान्विवर्जयेत्}\ A; \skt{प्ररिवर्जयेत्}\ N}}\skt{~॥}\pend

\pstart\versepar \ledsidenote{\passref{1.3.52}}\edtext{\edtext{\skt{स्नेहा\-दि\-ष्व\-न\-भि\-ज्ञास्तु}}{\lemma{\skt{स्नेहादिष्व्}}\Afootnote{N; \skt{स्नेहादि}\ A H}}}{\lemma{\skt{अनभिज्ञास्तु}}\Afootnote{N; \skt{ष्वनभिज्ञो य}\ A; \skt{ष्वनभिज्ञास्तु}\ H}}\skt{ }\edtext{\skt{छेद्यादिषु च}}{\lemma{\skt{छेद्यादिषु च}}\Afootnote{N H; \skt{श्छेद्यादिषु च}\ A}}\skt{ कर्म्मसु~। }\pend
\pstart\versepar \skt{}\edtext{\skt{मारयन्ति}}{\lemma{\skt{मारयन्ति}}\Afootnote{N H; \skt{स निहन्ति}\ A}}\skt{ }\edtext{\skt{नरं}}{\lemma{\skt{नरं}}\Afootnote{N H; \skt{जनं}\ A}}\skt{ }\edtext{\skt{लोभात्कुवैद्या}}{\lemma{\skt{लोभात्कुवैद्या}}\Afootnote{N H; \skt{लोभात्कुवैद्यो}\ A}}\skt{ नृपदोषतः~॥}\pend

\pstart\versepar \ledsidenote{\passref{1.3.53}}\skt{यस्तूभयज्ञो }\edtext{\skt{मतिमान्स}}{\lemma{\skt{मतिमान्स}}\Afootnote{A H; \skt{मतिमां स}\ N}}\skt{ }\edtext{\skt{सम\-र्थोऽ\-थ\-सा\-धने}}{\lemma{\skt{सम\appbrk{}र्थोऽ\appbrk{}थ\appbrk{}सा\appbrk{}धने}}\Afootnote{H; \skt{सम\appbrk{}र्थोऽ\appbrk{}र्थ\appbrk{}सा\appbrk{}धने}\ A; \skt{समर्थो र्थसाधने}\ N}}\skt{~। }\pend
\pstart\versepar \skt{आहवे कर्म्म निर्व्वोढुं }\edtext{\skt{द्विचक्र}}{\lemma{\skt{द्विचक्र}}\Afootnote{N H; \skt{द्विचक्रः}\ A}}\skt{ स्यन्दनो यथा~॥}\pend

\pstart \ledsidenote{\passref{1.3.54}}\edtext{\skt{वत्स}}{\lemma{\skt{वत्स}}\Afootnote{N H; \skt{अथ वत्स}\ A}}\skt{ }\edtext{\skt{यथा}}{\lemma{\skt{यथा}}\Afootnote{N H; \skt{तदे}\ A}}\skt{ }\edtext{\skt{अध्येयं}}{\lemma{\skt{अध्येयं}}\Afootnote{H; \skt{तदध्येयं}\ A; \emph{om.}\ N}}\skt{ }\edtext{\skt{तथोपधारय}}{\lemma{\skt{तथोपधारय}}\Afootnote{H; \skt{यथा तथोपधारय}\ A; \skt{ध्येयं तथोप}\ N}}\skt{ }\edtext{\skt{स्वोच्यमानं}}{\lemma{\skt{स्वोच्यमानं}}\Afootnote{H; \skt{मया प्रोच्यमानं}\ A; \skt{ध्यायस्वोच्यमानं}\ N}}\skt{~। अथ शुचये }\edtext{\skt{कृतं}}{\lemma{\skt{कृतं}}\Afootnote{N H; \emph{om.}\ A}}\skt{ }\edtext{\skt{अङ्गलाय}}{\lemma{\skt{अङ्गलाय}}\Afootnote{N H; \emph{om.}\ A}}\skt{ }\edtext{\skt{कृतो\-त्त\-रा\-स\-ङ्गा\-यो\-प\-स्थि\-ता\-या\-ध्य\-य\-न\-काले}}{\lemma{\skt{कृतो\appbrk{}त्त\appbrk{}रा\appbrk{}स\appbrk{}ङ्गा\appbrk{}यो\appbrk{}प\appbrk{}स्थि\appbrk{}ता\appbrk{}या\appbrk{}ध्य\appbrk{}य\appbrk{}न\appbrk{}काले}}\Afootnote{N H; \skt{कृतो\appbrk{}त्त\appbrk{}रा\appbrk{}स\appbrk{}ङ्गा\appbrk{}या\appbrk{}व्या\appbrk{}कु\appbrk{}ला\appbrk{}यो\appbrk{}प\appbrk{}स्थि\appbrk{}ता\appbrk{}या\appbrk{}ध्य\appbrk{}य\appbrk{}न\appbrk{}काले}\ A}}\skt{ शिष्याय~। }\edtext{\skt{यथाशक्तितो}}{\lemma{\skt{यथाशक्तितो}}\Afootnote{N H; \skt{यथाशक्ति}\ A}}\skt{ }\edtext{\edtext{\skt{गुरु\-रु\-प\-दि\-शे\-त्पदं}}{\lemma{\skt{गुरुर्}}\Afootnote{N; \skt{गुरु}\ A H}}}{\lemma{\skt{उपदिशेत्पदं}}\Afootnote{N; \skt{रुपदिशेत्पदं}\ A H}}\skt{ }\edtext{\skt{पदं}}{\lemma{\skt{पदं}}\Afootnote{H; \emph{om.}\ A; \skt{पम्}\ N}}\skt{ }\edtext{\skt{पादं}}{\lemma{\skt{पादं}}\Afootnote{N H; \emph{om.}\ A}}\skt{ पादं श्लोकं वा ते च }\edtext{\skt{क्रमे}}{\lemma{\skt{क्रमे}}\Afootnote{N H; \skt{प}\ A}}\skt{ }\edtext{\skt{यश}}{\lemma{\skt{यश}}\Afootnote{H; \skt{दपा}\ A; \emph{om.}\ N}}\skt{ }\edtext{\skt{भूयः}}{\lemma{\skt{भूयः}}\Afootnote{H; \skt{दश्लोका भूयः}\ A; \emph{om.}\ N}}\skt{ }\edtext{\skt{सन्धेयाः}}{\lemma{\skt{सन्धेयाः}}\Afootnote{H; \skt{क्रमे\appbrk{}णा\appbrk{}नु\appbrk{}स\appbrk{}न्धेयाः}\ A; \skt{ण भूयः सन्धेयाः}\ N}}\skt{~। }\edtext{\skt{एवमेकैकशो}}{\lemma{\skt{एकैकशो}}\Afootnote{A H; \skt{एकैकसोश्}\ N}}\skt{ }\edtext{\edtext{\skt{घट\-ये\-दा\-त्मा\-न\-ञ्चात्र}}{\lemma{\skt{घट\appbrk{}ये\appbrk{}दा\appbrk{}त्मा\appbrk{}नञ्}}\Afootnote{N H; \skt{घटयेदात्मना}\ A}}}{\lemma{\skt{चात्र}}\Afootnote{N H; \skt{चानु}\ A}}\skt{ }\edtext{\skt{पठे\-द\-द्रु\-त\-म\-वि\-लं\-बि\-त\-म\-न\-नु\-ना\-सिकं}}{\lemma{\skt{अविलंबितम्}}\Afootnote{N H; \skt{अवि\appbrk{}ल\appbrk{}म्बि\appbrk{}त\appbrk{}म\appbrk{}वि\appbrk{}श\appbrk{}ङ्कि\appbrk{}तम्}\ A}}\skt{ }\edtext{\edtext{\edtext{\skt{नाव्य\-क्ता\-ति\-नि\-पी\-डि\-त\-व\-र्ण्ण\-म\-क्षि\-भ्रु\-वौ\-ष्ठ\-ह\-स्त्तै\-र\-नतं}}{\lemma{\skt{नाव्य\appbrk{}क्ता\appbrk{}ति\appbrk{}नि\appbrk{}पी\appbrk{}डि\appbrk{}त\appbrk{}व\appbrk{}र्ण्णम्}}\Afootnote{H; \skt{सुव्य\appbrk{}क्ता\appbrk{}क्ष\appbrk{}र\appbrk{}म\appbrk{}पी\appbrk{}डि\appbrk{}त\appbrk{}व\appbrk{}र्णम्}\ A; \skt{नाव्य\appbrk{}क्ता\appbrk{}ति\appbrk{}नि\appbrk{}ष्पी\appbrk{}डि\appbrk{}त\appbrk{}व\appbrk{}र्ण्णम्}\ N}}}{\lemma{\skt{अक्षि\appbrk{}भ्रु\appbrk{}वौ\appbrk{}ष्ठ\appbrk{}ह\appbrk{}स्त्तैर्}}\Afootnote{H; \skt{अक्षि\appbrk{}भ्रु\appbrk{}वौ\appbrk{}ष्ठ\appbrk{}ह\appbrk{}स्तैर्}\ A; \skt{अक्षि\appbrk{}भ्रु\appbrk{}व\appbrk{}ष्ठ\appbrk{}ह\appbrk{}स्तेर्}\ N}}}{\lemma{\skt{अनतं}}\Afootnote{H; \skt{अनभिनीतं}\ A; \skt{अननुनीतं}\ N}}\skt{ सुसंस्कृतं }\edtext{\skt{नात्यु\-च्चै\-र्न्ना\-ति\-नीचैः}}{\lemma{\skt{न्नातिनीचैः}}\Afootnote{H; \skt{नातिनीचैश्च}\ A N}}\skt{ }\edtext{\skt{स्व}}{\lemma{\skt{स्व}}\Afootnote{N H; \skt{स्वरैः}\ A}}\skt{ }\edtext{\edtext{\skt{पठे\-त्त\-यो\-र\-धी\-या\-न\-योर्न्न}}{\lemma{\skt{पठेत्तयोर्}}\Afootnote{H; \skt{पठेत्~। न}\ A; \skt{पठे तयोर्}\ N}}}{\lemma{\skt{अधीयानयोर्न्न}}\Afootnote{H; \emph{om.}\ A; \skt{अधीयानयो न}\ N}}\skt{ चान्तरेण }\edtext{\skt{कश्चिद्व्रजेदिति}}{\lemma{\skt{कश्चिद्व्रजेदिति}}\Afootnote{N; \skt{कश्चि\appbrk{}द्व्र\appbrk{}जे\appbrk{}त्त\appbrk{}यो\appbrk{}र\appbrk{}धी\appbrk{}या\appbrk{}नयोः}\ A; \skt{कश्चिद्व्रजेदिति~॥ …}\ H}}\skt{~॥}\pend

\pstart\versepar \ledsidenote{\passref{1.3.55}}\edtext{\edtext{\skt{शुचिर्ग्गुरुपरो}}{\lemma{\skt{शुचिर्}}\Afootnote{N H; \skt{भवतश्चात्र~। शुचि}\ A}}}{\lemma{\skt{ग्गुरुपरो}}\Afootnote{N H; \skt{र्गुरुपरो}\ A}}\skt{ }\edtext{\edtext{\skt{दक्षा\-स्त\-न्द्री\-नि\-द्रा\-वि\-व\-र्ज्जितः}}{\lemma{\skt{दक्षास्}}\Afootnote{N H; \skt{दक्ष}\ A}}}{\lemma{\skt{तन्द्री\appbrk{}नि\appbrk{}द्रा\appbrk{}वि\appbrk{}व\appbrk{}र्ज्जितः}}\Afootnote{H; \skt{स्तन्द्रा\appbrk{}नि\appbrk{}द्रा\appbrk{}वि\appbrk{}व\appbrk{}र्जितः}\ A; \skt{तन्दी निद्राविवर्ज्जितः}\ N}}\skt{~। }\pend
\pstart\versepar \skt{पठन्नेतेन विधिना शिष्यः शास्त्रा\-न्त\-मा\-प्नु\-यात्~॥}\pend

\pstart\versepar \ledsidenote{\passref{1.3.56}}\edtext{\skt{वाक्सौ\-ष्ठ\-वे\-र्थ\-वि\-ज्ञाने}}{\lemma{\skt{वाक्सौ\appbrk{}ष्ठ\appbrk{}वे\appbrk{}र्थ\appbrk{}वि\appbrk{}ज्ञाने}}\Afootnote{N H; \skt{वाक्सौ\appbrk{}ष्ठ\appbrk{}वेऽ\appbrk{}र्थ\appbrk{}वि\appbrk{}ज्ञाने}\ A}}\skt{ }\edtext{\skt{प्रागल्भ्ये}}{\lemma{\skt{प्रागल्भ्ये}}\Afootnote{A H; \skt{प्राम्लैल्भे}\ N}}\skt{ }\edtext{\skt{कर्म्मनैपुणे}}{\lemma{\skt{कर्म्मनैपुणे}}\Afootnote{A H; \skt{कर्मनैपुरौ}\ N}}\skt{~। }\pend
\pstart\versepar \skt{}\edtext{\skt{तदभ्यासे च}}{\lemma{\skt{तदभ्यासे च}}\Afootnote{N H; \skt{तदभ्यासे च}\ A}}\skt{ सिद्धौ च }\edtext{\skt{यते\-ता\-ध्य\-य\-ना\-न्तग}}{\lemma{\skt{यते\appbrk{}ता\appbrk{}ध्य\appbrk{}य\appbrk{}ना\appbrk{}न्तग}}\Afootnote{\skt{यते\appbrk{}ता\appbrk{}ध्य\appbrk{}य\appbrk{}ना\appbrk{}न्तगः}\ A; \skt{यते\appbrk{}ता\appbrk{}ध्य\appbrk{}य\appbrk{}ना\appbrk{}न्तग इति}\ N H}}\skt{~॥}\pend

\pstart \ledsidenote{\passref{1.3.trailer}}\edtext{\skt{इति}}{\lemma{\skt{इति}}\Afootnote{\skt{इति सुश्रु\appbrk{}त\appbrk{}सं\appbrk{}हि\appbrk{}तायां सूत्र\appbrk{}स्था\appbrk{}नेऽ\appbrk{}ध्य\appbrk{}य\appbrk{}न\appbrk{}स\appbrk{}म्प्र\appbrk{}दा\appbrk{}नीयो नाम तृतीयोऽध्यायः}\ A; \skt{३~॥}\ N; \skt{३~॥}\ H}}\skt{~॥ ३~॥}\pend

\endnumbering
\beginnumbering
\eledsection*{Adhyāya 4}
\pstart \ledsidenote{\passref{1.4.1}}\skt{अथातः }\edtext{\skt{प्रभा\-ष\-णी\-य\-म\-ध्यायं}}{\lemma{\skt{प्रभाषणीयम्}}\Afootnote{A H; \skt{प्रवचीनीयम्}\ N}}\skt{ }\edtext{\skt{व्याख्यास्यामः}}{\lemma{\skt{व्याख्यास्यामः}}\Afootnote{A H; \emph{om.}\ N}}\skt{~॥}\pend

\pstart \ledsidenote{\passref{1.4.3}}\edtext{\edtext{\skt{अधि\-ग\-त\-म\-प्य\-ध्य\-य\-न\-म\-न\-नु\-भा\-षि\-त\-म\-र्थतः}}{\lemma{\skt{अन\appbrk{}नु\appbrk{}भा\appbrk{}षि\appbrk{}तम्}}\Afootnote{N H; \skt{अप्रभाषितम्}\ A}}}{\lemma{\skt{अर्थतः}}\Afootnote{N H; \skt{अर्थतः खरस्य चन्दनभार इव केवलं परिश्रमकरं भवति}\ A}}\pend

\pstart \ledsidenote{\passref{1.4.4}}\edtext{\skt{खरस्य}}{\lemma{\skt{खरस्य}}\Afootnote{N H; \skt{भवति चात्र यथा खरश्}\ A}}\skt{ }\edtext{\skt{चन्दनभार}}{\lemma{\skt{चन्दनभार}}\Afootnote{N H; \skt{चन्द\appbrk{}न\appbrk{}भा\appbrk{}र\appbrk{}वाही भारस्य}\ A}}\skt{ }\edtext{\skt{इव}}{\lemma{\skt{इव}}\Afootnote{N H; \skt{वेत्ता न तु चन्दनस्य~। एवं हि शास्त्राणि बहून्यधीत्य चार्थेषु मूढाः ख}\ A}}\skt{ }\edtext{\skt{केवलं}}{\lemma{\skt{केवलं}}\Afootnote{N H; \skt{रवद्व}\ A}}\skt{ }\edtext{\skt{परिश्रमकरो}}{\lemma{\skt{परिश्रमकरो}}\Afootnote{N H; \skt{हन्ति}\ A}}\skt{ }\edtext{\skt{भवति}}{\lemma{\skt{भवति}}\Afootnote{N H; \emph{om.}\ A}}\skt{~॥}\pend

\pstart \ledsidenote{\passref{1.4.5}}\edtext{\edtext{\edtext{\edtext{\skt{तस्मा\-त्स\-विं\-श\-म\-ध्या\-य\-श\-त\-म\-नु\-प\-द\-पा\-द\-श्लो\-क\-म\-व\-स्य\-म\-नु\-व\-र्णि\-तव्यं}}{\lemma{\skt{तस्मात्सविंशम्}}\Afootnote{N H; \skt{तस्मात्सविंश}\ A}}}{\lemma{\skt{अध्यायशतम्}}\Afootnote{N H; \skt{म ध्यायशतम्}\ A}}}{\lemma{\skt{अवस्यम्}}\Afootnote{N H; \skt{अनु\appbrk{}व\appbrk{}र्ण\appbrk{}यि\appbrk{}त\appbrk{}व्यम्}\ A}}}{\lemma{\skt{अनुवर्णितव्यं}}\Afootnote{N H; \skt{अनुश्रोतव्यं च}\ A}}\skt{~। }\edtext{\skt{कस्मात्सूक्ष्मा हि}}{\lemma{\skt{कस्मात्सूक्ष्मा हि}}\Afootnote{A H; \skt{कस्म सूक्ष्मा हि}\ N}}\skt{ }\edtext{\skt{द्रव्य\-र\-स\-गु\-ण\-वी\-र्य\-वि\-पा\-क\-दो\-ष\-धा\-तु\-म\-ला\-श\-य\-म\-र्म्म\-सि\-रा\-स्ना\-यु\-स\-न्ध्य\-स्थि\-वि\-भागाः}}{\lemma{\skt{द्रव्य\appbrk{}र\appbrk{}स\appbrk{}गु\appbrk{}ण\appbrk{}वी\appbrk{}र्य\appbrk{}वि\appbrk{}पा\appbrk{}क\appbrk{}दो\appbrk{}ष\appbrk{}धा\appbrk{}तु\appbrk{}म\appbrk{}ला\appbrk{}श\appbrk{}य\appbrk{}म\appbrk{}र्म्म\appbrk{}सि\appbrk{}रा\appbrk{}स्ना\appbrk{}यु\appbrk{}स\appbrk{}न्ध्य\appbrk{}स्थि\appbrk{}वि\appbrk{}भागाः}}\Afootnote{H; \skt{द्रव्य\appbrk{}र\appbrk{}स\appbrk{}गु\appbrk{}ण\appbrk{}वी\appbrk{}र्य\appbrk{}वि\appbrk{}पा\appbrk{}क\appbrk{}दो\appbrk{}ष\appbrk{}धा\appbrk{}तु\appbrk{}म\appbrk{}ला\appbrk{}श\appbrk{}य\appbrk{}म\appbrk{}र्म\appbrk{}सि\appbrk{}रा\appbrk{}स्ना\appbrk{}यु\appbrk{}स\appbrk{}न्ध्यस्थि}\ A; \skt{द्रव्य\appbrk{}र\appbrk{}स\appbrk{}गु\appbrk{}ण\appbrk{}वी\appbrk{}र्य\appbrk{}वि\appbrk{}पा\appbrk{}क\appbrk{}दो\appbrk{}ष\appbrk{}धा\appbrk{}तु\appbrk{}म\appbrk{}ला\appbrk{}स्र\appbrk{}य\appbrk{}म\appbrk{}र्म\appbrk{}सि\appbrk{}रा\appbrk{}स्ना\appbrk{}यु\appbrk{}स\appbrk{}न्ध्य\appbrk{}स्थि\appbrk{}वि\appbrk{}भागाः}\ N}}\skt{~। }\edtext{\skt{गर्भ\-सं\-भ\-व\-द्र\-व्य\-स\-मू\-ह\-वि\-भा\-गाश्च}}{\lemma{\skt{गर्भ\appbrk{}सं\appbrk{}भ\appbrk{}व\appbrk{}द्र\appbrk{}व्य\appbrk{}स\appbrk{}मू\appbrk{}ह\appbrk{}वि\appbrk{}भा\appbrk{}गाश्च}}\Afootnote{N H; \skt{गर्भ\appbrk{}स\appbrk{}म्भ\appbrk{}व\appbrk{}द्र\appbrk{}व्य\appbrk{}स\appbrk{}मू\appbrk{}ह\appbrk{}वि\appbrk{}भागा}\ A}}\skt{~। }\edtext{\skt{तथा}}{\lemma{\skt{तथा}}\Afootnote{N H; \skt{स्तथा}\ A}}\skt{ प्रन\-ष्ट\-श\-ल्यो\-द्ध\-र\-ण\-व्र\-ण\-वि\-नि\-श्च\-य\-भ\-ग्न\-वि\-कल्पाः~। }\edtext{\edtext{\skt{साध्य\-या\-प्य\-प्र\-त्या\-ख्ये\-यता}}{\lemma{\skt{साध्ययाप्य्}}\Afootnote{N H; \skt{साध्यया}\ A}}\skt{ च}}{\lemma{\skt{अप्रत्याख्येयता च}}\Afootnote{N H; \skt{प्यप्रत्याख्येयता च}\ A}}\skt{ }\edtext{\edtext{\edtext{\skt{विका\-रा\-णा\-मे\-व\-मा\-द\-यश्च}}{\lemma{\skt{विकाराणाम्}}\Afootnote{N H; \skt{विकाराणा}\ A}}}{\lemma{\skt{एवम्}}\Afootnote{N H; \skt{मेव}\ A}}}{\lemma{\skt{आदयश्च}}\Afootnote{N H; \skt{मादयश्चान्ये}\ A}}\skt{ }\edtext{\edtext{\skt{सहस्रशोऽन्ये}}{\lemma{\skt{सहस्रशो}}\Afootnote{N H; \skt{विशेषाः सहस्रशो}\ A}}}{\lemma{\skt{ऽन्ये}}\Afootnote{H; \skt{ये}\ A; \skt{न्ये}\ N}}\skt{ }\edtext{\skt{विशेषाः~। ये}}{\lemma{\skt{विशेषाः~। ये}}\Afootnote{N H; \skt{वि}\ A}}\skt{ चिन्त्यमाना }\edtext{\edtext{\skt{विम\-ल\-वि\-पु\-ल\-बु\-द्धे\-रपि}}{\lemma{\skt{विम\appbrk{}ल\appbrk{}वि\appbrk{}पु\appbrk{}ल\appbrk{}बु\appbrk{}द्धेर्}}\Afootnote{N H; \skt{विम\appbrk{}ल\appbrk{}वि\appbrk{}पु\appbrk{}ल\appbrk{}बुद्धे}\ A}}}{\lemma{\skt{अपि}}\Afootnote{N H; \skt{रपि}\ A}}\skt{ }\edtext{\edtext{\skt{बुद्धि\-मा\-कु\-ली\-कुर्युः}}{\lemma{\skt{बुद्धिम्}}\Afootnote{N H; \skt{बुद्धि}\ A}}\skt{ किं}}{\lemma{\skt{आकुलीकुर्युः किं}}\Afootnote{N H; \skt{माकुलीकुर्युः किं}\ A}}\skt{ }\edtext{\edtext{\edtext{\edtext{\skt{पुन\-र\-ल्प\-बु\-द्धे\-स्त\-स्मा\-द\-व\-श्य\-म\-नु\-व\-र्ण\-यि\-तव्यं}}{\lemma{\skt{पुनर्}}\Afootnote{N H; \skt{पुन}\ A}}}{\lemma{\skt{अल्पबुद्धेस्}}\Afootnote{N H; \skt{रल्पबुद्धेः}\ A}}}{\lemma{\skt{तस्मादवश्यम्}}\Afootnote{N; \skt{तस्मादवश्य}\ A; \skt{तस्मादवस्यम्}\ H}}}{\lemma{\skt{अनु\appbrk{}व\appbrk{}र्ण\appbrk{}यि\appbrk{}तव्यं}}\Afootnote{N H; \skt{मनु\appbrk{}प\appbrk{}द\appbrk{}पा\appbrk{}द\appbrk{}श्लो\appbrk{}क\appbrk{}म\appbrk{}नु\appbrk{}व\appbrk{}र्ण\appbrk{}यि\appbrk{}त\appbrk{}व्य\appbrk{}म\appbrk{}नु\appbrk{}श्रो\appbrk{}तव्यं च}\ A}}\skt{~॥}\pend

\pstart \ledsidenote{\passref{1.4.6}}\edtext{\edtext{\edtext{\edtext{\edtext{\skt{अन्य\-शा\-स्त्र\-वि\-ष\-यो\-प\-प\-न्ना\-ना\-म्चा\-र्था\-ना\-मि\-हो\-प\-नी\-ता\-ना\-म\-र्थ\-व\-शा\-त्ते\-षा\-न्त\-द्वि\-द्येभ्य}}{\lemma{\skt{अन्य\appbrk{}शा\appbrk{}स्त्र\appbrk{}वि\appbrk{}ष\appbrk{}यो\appbrk{}प\appbrk{}प\appbrk{}न्ना\appbrk{}नाम्}}\Afootnote{N H; \skt{अन्य\appbrk{}शा\appbrk{}स्त्रो\appbrk{}प\appbrk{}प\appbrk{}न्नानां}\ A; \emph{om.}\ K}}}{\lemma{\skt{चार्थानाम्}}\Afootnote{N H; \skt{चार्थाना}\ A; \emph{om.}\ K}}}{\lemma{\skt{इहो\appbrk{}प\appbrk{}नी\appbrk{}ता\appbrk{}नाम्}}\Afootnote{N H; \skt{मिहोपनीताना}\ A; \skt{…नाम्}\ K}}}{\lemma{\skt{अर्थ\appbrk{}व\appbrk{}शा\appbrk{}त्ते\appbrk{}षान्}}\Afootnote{K N H; \skt{मर्थवशात्तेषां}\ A}}}{\lemma{\skt{तद्विद्येभ्य}}\Afootnote{K N; \skt{तद्विद्येभ्य}\ A; \skt{तद्विद्यैभ्य}\ H}}\skt{ एव }\edtext{\edtext{\skt{व्याख्या\-न\-म\-नु\-श्रो\-तव्यं}}{\lemma{\skt{व्याख्यानम्}}\Afootnote{N H; \skt{व्याख्यान}\ A; \skt{व्याख्या}\ K}}\skt{~॥ न}}{\lemma{\skt{अनुश्रोतव्यं~॥ न}}\Afootnote{H; \skt{मनुश्रोतव्यं, क}\ A; \emph{om.}\ K; \skt{अनुस्रातव्यं~। न}\ N}}\skt{ }\edtext{\edtext{\edtext{\skt{ह्येकस्मिञ्च्छास्त्रे}}{\lemma{\skt{ह्य्}}\Afootnote{N H; \skt{स्मात्? न}\ A; \skt{न}\ K}}}{\lemma{\skt{एकस्मिञ्}}\Afootnote{N H; \skt{ह्येकस्मि}\ A; \emph{om.}\ K}}}{\lemma{\skt{च्छास्त्रे}}\Afootnote{N H; \skt{न्शास्त्रे}\ A; \emph{om.}\ K}}\skt{ }\edtext{\edtext{\skt{सर्व\-शा\-स्त्रा\-णा\-म\-व\-रोधः}}{\lemma{\skt{सर्वशास्त्राणाम्}}\Afootnote{N H\textsuperscript{pc}; \skt{शक्यः सर्वशास्त्राणा}\ A; \skt{…स्त्रे सर्वशास्त्राणाम्}\ K; \skt{मर्व्वशास्त्राणाम्}\ H\textsuperscript{ac}}}}{\lemma{\skt{अवरोधः}}\Afootnote{K N H; \skt{मवरोधः}\ A}}\skt{ कर्त्तुं }\edtext{\skt{शक्य}}{\lemma{\skt{शक्य}}\Afootnote{K N H; \emph{om.}\ A}}\skt{ }\edtext{\skt{इति}}{\lemma{\skt{इति}}\Afootnote{K N H; \emph{om.}\ A}}\skt{~॥ ६~॥}\pend

\pstart\versepar \ledsidenote{\passref{1.4.7}}\edtext{\edtext{\skt{शास्त्र\-मे\-क\-म\-धी\-यानो}}{\lemma{\skt{शास्त्रम्}}\Afootnote{K N H; \skt{भवन्ति चात्र}\ A}}\skt{ न}}{\lemma{\skt{अधीयानो न}}\Afootnote{K N H; \skt{शास्त्रमधीयानो न}\ A}}\skt{ }\edtext{\skt{विद्या\-च्छा\-स्त्र\-नि\-श्चयं}}{\lemma{\skt{विद्या\appbrk{}च्छा\appbrk{}स्त्र\appbrk{}नि\appbrk{}श्चयं}}\Afootnote{A K H; \skt{विद्या शास्त्रनिश्चयं}\ N}}\skt{~।}\footnoteA{Parallel to the Kātyāyanasmṛtisāroddhāraḥ, Rājaguṇāḥ v.66 (Kane 1933: 11).}\skt{ }\pend
\pstart\versepar \skt{तस्माद्बहुश्रुतः शास्त्रं }\edtext{\skt{विजा\-नी\-या\-च्चि\-कि\-त्सकः}}{\lemma{\skt{विजा\appbrk{}नी\appbrk{}या\appbrk{}च्चि\appbrk{}कि\appbrk{}त्सकः}}\Afootnote{K N H; \skt{विजा\appbrk{}नी\appbrk{}या\appbrk{}च्चि\appbrk{}कि\appbrk{}त्सकः}\ A}}\skt{~।}\pend

\pstart\versepar \ledsidenote{\passref{1.4.8}}\skt{शास्त्रं }\edtext{\edtext{\skt{गुरु\-मु\-खो\-द्गी\-र्ण्ण\-मा\-दा\-यो\-पास्य}}{\lemma{\skt{गुरु\appbrk{}मु\appbrk{}खो\appbrk{}द्गी\appbrk{}र्ण्णम्}}\Afootnote{K N H; \skt{गुरुमुखोद्गीर्ण}\ A}}}{\lemma{\skt{आदायोपास्य}}\Afootnote{K N H; \skt{मादायोपास्य}\ A}}\skt{ चासकृत्~। }\pend
\pstart\versepar \skt{यः कर्म कुरुते वैद्यः स }\edtext{\edtext{\skt{वैद्योऽन्ये}}{\lemma{\skt{वैद्यो}}\Afootnote{A K H; \skt{वेद्यो}\ N}}\skt{ तु}}{\lemma{\skt{ऽन्ये तु}}\Afootnote{A K H; \skt{न्ये तु}\ N}}\skt{ तस्कराः~॥}\pend

\pstart\versepar \ledsidenote{\passref{1.4.9}}\edtext{\skt{ओप\-धे\-न\-व\-मौ\-रभ्रं}}{\lemma{\skt{ओपधेनवम्}}\Afootnote{K H; \skt{औपधेनवम्}\ A N}}\skt{ }\edtext{\skt{सौश्रु\-त\-म्पौ\-ष्क\-ला\-वतं}}{\lemma{\skt{सौश्रुतम्}}\Afootnote{A K N H\textsuperscript{pc}; \skt{मौश्रुतं}\ H\textsuperscript{ac}}}\skt{~। }\pend
\pstart\versepar \skt{शेषाणां शल्यतन्त्राणां मूलान्येतानि }\edtext{\skt{निर्द्दिशेदिति}}{\lemma{\skt{निर्द्दिशेदिति}}\Afootnote{N H; \skt{निर्दिशेत्}\ A; \skt{निर्द्दिशेदिति~॥ {\symfont ❈}~॥ ४}\ K}}\skt{~॥}\pend

\pstart \ledsidenote{\passref{1.4.trailer}}\edtext{\skt{॥ ४~॥}}{\lemma{\skt{॥ ४~॥}}\Afootnote{\skt{इति सुश्रु\appbrk{}त\appbrk{}सं\appbrk{}हि\appbrk{}तायां सूत्रस्थाने प्रभाषणीयो नाम चतुर्थोऽध्यायः}\ A; \skt{॥ ४~॥}\ N H}}\pend

\endnumbering
\beginnumbering
\eledsection*{Adhyāya 5}
\pstart \ledsidenote{\passref{1.5.1}}\edtext{\skt{अथा\-तोऽ\-ग्रो\-प\-ह\-र\-णी\-य\-म\-ध्यायं}}{\lemma{\skt{ऽग्रो\appbrk{}प\appbrk{}ह\appbrk{}र\appbrk{}णी\appbrk{}यम्}}\Afootnote{A K H; \skt{ग्रोप\appbrk{}ह\appbrk{}र\appbrk{}णी\appbrk{}यम्}\ N}}\skt{ }\edtext{\skt{व्याख्यास्यामः}}{\lemma{\skt{व्याख्यास्यामः}}\Afootnote{A H; \emph{om.}\ K N}}\skt{~॥}\pend

\pstart \ledsidenote{\passref{1.5.3}}\skt{त्रिविधं कर्म्म पूर्वकर्म }\edtext{\skt{प्रधानकम्म}}{\lemma{\skt{प्रधानकम्म}}\Afootnote{K N; \skt{प्रधानकर्म}\ A H}}\skt{ }\edtext{\skt{पश्चात्कर्मेति}}{\lemma{\skt{पश्चात्कर्मेति}}\Afootnote{A K N; \skt{पश्चात्कर्म्म इति}\ H}}\skt{~। }\edtext{\skt{तद्व्या\-धि\-म्प्र\-त्यु\-प\-दे\-क्ष्यामः}}{\lemma{\skt{तद्व्याधिम्}}\Afootnote{A N; \skt{तद्व्याधिम्}\ K; \skt{तद्व्याधीं}\ H}}\skt{~।}\pend

\pstart \ledsidenote{\passref{1.5.4}}\skt{अस्य तु }\edtext{\skt{शास्त्रस्य}}{\lemma{\skt{शास्त्रस्य}}\Afootnote{A K H; \emph{om.}\ N}}\skt{ }\edtext{\skt{शस्त्र\-क\-र्म\-प्रा\-धा\-न्या\-त्पूर्वं}}{\lemma{\skt{शस्त्र\appbrk{}क\appbrk{}र्म\appbrk{}प्रा\appbrk{}धा\appbrk{}न्या\appbrk{}त्पूर्वं}}\Afootnote{K; \skt{शस्त्र\appbrk{}क\appbrk{}र्म\appbrk{}प्रा\appbrk{}धा\appbrk{}न्या\appbrk{}च्छ\appbrk{}स्त्र\appbrk{}क\appbrk{}र्मैव तावत्पूर्वम्}\ A; \skt{शस्त्र\appbrk{}क\appbrk{}र्म\appbrk{}प्रा\appbrk{}ध्या\appbrk{}न्यात्~। पूर्व्वं}\ N; \skt{शस्त्र\appbrk{}क\appbrk{}र्म्म\appbrk{}प्रा\appbrk{}धा\appbrk{}न्यात्~। पूर्व्व}\ H}}\skt{ }\edtext{\skt{शस्त्रसम्भाराने वोपदेक्ष्यामः}}{\lemma{\skt{शस्त्रसम्भाराने वोपदेक्ष्यामः}}\Afootnote{K N H; \skt{उप\appbrk{}दे\appbrk{}क्ष्या\appbrk{}म\appbrk{}स्त\appbrk{}त्स\appbrk{}म्भा\appbrk{}रांश्च}\ A}}\skt{~।}\pend

\pstart \ledsidenote{\passref{1.5.5}}\edtext{\skt{तच्च}}{\lemma{\skt{तच्च}}\Afootnote{A K N; \skt{तच्च}\ H}}\skt{ }\edtext{\skt{शस्त्र\-क\-र्मा\-ष्ट\-वि\-ध\-म्भ\-वति}}{\lemma{\skt{भवति}}\Afootnote{K N H; \emph{om.}\ A}}\skt{~। }\edtext{\skt{तद्यथा}}{\lemma{\skt{तद्यथा}}\Afootnote{A K N; \skt{तद्यथा}\ H}}\skt{~। छेद्यं भेद्यं लेख्यं }\edtext{\skt{वेध्य\-मे\-ष्य\-मा\-हार्यं}}{\lemma{\skt{एष्यम्}}\Afootnote{A K N; \skt{मेष्यं}\ H}}\skt{ }\edtext{\skt{विस्राव्यं}}{\lemma{\skt{विस्राव्यं}}\Afootnote{A K N; \skt{विश्राव्यं}\ H}}\skt{ }\edtext{\skt{सीव्यमिति}}{\lemma{\skt{सीव्यमिति}}\Afootnote{\skt{सीव्यमिति}\ A H; \skt{सीव्यमित्}\ K; \skt{सीव्यमित्य्}\ N}}\pend

\pstart \ledsidenote{\passref{1.5.6}}\edtext{\edtext{\skt{यतोऽ\-न्य\-त्क\-र्म\-चि\-की\-र्षुणा}}{\lemma{\skt{यतो}}\Afootnote{K; \skt{अतो}\ A; \emph{om.}\ N H}}}{\lemma{\skt{ऽन्य\appbrk{}त्क\appbrk{}र्म\appbrk{}चि\appbrk{}की\appbrk{}र्षुणा}}\Afootnote{K; \skt{ऽन्यतमं कर्म चिकीर्षता}\ A; \skt{अतो न्यत्कर्म चिकीर्षुणा}\ N; \skt{अतोऽन्यत्कर्म्म चिकीर्षुणा}\ H}}\skt{ }\edtext{\edtext{\skt{पूर्व\-मे\-वो\-प\-क\-ल्प\-यि\-त\-व्यानि}}{\lemma{\skt{पूर्वम्}}\Afootnote{K N; \skt{वैद्येन पूर्व}\ A; \skt{पूर्व्व}\ H}}}{\lemma{\skt{एवो\appbrk{}प\appbrk{}क\appbrk{}ल्प\appbrk{}यि\appbrk{}त\appbrk{}व्यानि}}\Afootnote{K N; \skt{मेवो\appbrk{}प\appbrk{}क\appbrk{}ल्प\appbrk{}यि\appbrk{}त\appbrk{}व्यानि}\ A H}}\skt{ भवन्ति~। }\edtext{\skt{तद्यथा}}{\lemma{\skt{तद्यथा}}\Afootnote{A K N; \skt{तद्यथा}\ H}}\skt{ }\edtext{\skt{यन्त्र\-श\-स्त्र\-क्षा\-रा\-ग्नि\-श\-ला\-का\-पि\-चु\-प्लो\-त\-प\-त्र\-सू\-त्र\-घृ\-त\-म\-धु\-प\-य\-स्तै\-ल\-त\-र्प\-ण\-क\-षा\-या\-ले\-प\-न\-क\-ल्क\-शी\-तो\-द\-क\-व्य\-ज\-न\-क\-टा\-हा\-दीनि}}{\lemma{\skt{यन्त्र\appbrk{}श\appbrk{}स्त्र\appbrk{}क्षा\appbrk{}रा\appbrk{}ग्नि\appbrk{}श\appbrk{}ला\appbrk{}का\appbrk{}पि\appbrk{}चु\appbrk{}प्लो\appbrk{}त\appbrk{}प\appbrk{}त्र\appbrk{}सू\appbrk{}त्र\appbrk{}घृ\appbrk{}त\appbrk{}म\appbrk{}धु\appbrk{}प\appbrk{}य\appbrk{}स्तै\appbrk{}ल\appbrk{}त\appbrk{}र्प\appbrk{}ण\appbrk{}क\appbrk{}षा\appbrk{}या\appbrk{}ले\appbrk{}प\appbrk{}न\appbrk{}क\appbrk{}ल्क\appbrk{}शी\appbrk{}तो\appbrk{}द\appbrk{}क\appbrk{}व्य\appbrk{}ज\appbrk{}न\appbrk{}क\appbrk{}टा\appbrk{}हा\appbrk{}दीनि}}\Afootnote{K; \skt{यन्त्र\appbrk{}श\appbrk{}स्त्र\appbrk{}क्षा\appbrk{}रा\appbrk{}ग्नि\appbrk{}श\appbrk{}ला\appbrk{}का\appbrk{}शृ\appbrk{}ङ्ग\appbrk{}ज\appbrk{}लौ\appbrk{}का\appbrk{}ला\appbrk{}बू\appbrk{}जा\appbrk{}म्ब\appbrk{}वौ\appbrk{}ष्ठ\appbrk{}पि\appbrk{}चु\appbrk{}प्रो\appbrk{}त\appbrk{}सू\appbrk{}त्र\appbrk{}प\appbrk{}त्र\appbrk{}प\appbrk{}ट्ट\appbrk{}म\appbrk{}धु\appbrk{}घृ\appbrk{}त\appbrk{}व\appbrk{}सा\appbrk{}प\appbrk{}य\appbrk{}स्तैल- तर्प\appbrk{}ण\appbrk{}क\appbrk{}षा\appbrk{}या\appbrk{}ले\appbrk{}प\appbrk{}न\appbrk{}क\appbrk{}ल्क\appbrk{}व्य\appbrk{}ज\appbrk{}न\appbrk{}शी\appbrk{}तो\appbrk{}ष्णो\appbrk{}द\appbrk{}क\appbrk{}क\appbrk{}टा\appbrk{}हा\appbrk{}दीनि}\ A; \skt{यन्त्र\appbrk{}श\appbrk{}स्त्र\appbrk{}क्षा\appbrk{}रा\appbrk{}ग्नि\appbrk{}श\appbrk{}ला\appbrk{}का\appbrk{}पि\appbrk{}चु\appbrk{}प्लो\appbrk{}त\appbrk{}प\appbrk{}त्र\appbrk{}सू\appbrk{}त\appbrk{}घृ\appbrk{}त\appbrk{}म\appbrk{}धु\appbrk{}प\appbrk{}य\appbrk{}स्तै\appbrk{}ल\appbrk{}त\appbrk{}र्प\appbrk{}ण\appbrk{}क\appbrk{}षा\appbrk{}या\appbrk{}ले\appbrk{}प\appbrk{}न\appbrk{}क\appbrk{}ल्क\appbrk{}शी\appbrk{}तो\appbrk{}द\appbrk{}क\appbrk{}व्य\appbrk{}ज\appbrk{}न\appbrk{}क\appbrk{}टा\appbrk{}हा\appbrk{}दीनि}\ N; \skt{यन्त्र\appbrk{}श\appbrk{}स्त्र\appbrk{}क्षा\appbrk{}रा\appbrk{}ग्नि\appbrk{}श\appbrk{}ला\appbrk{}का\appbrk{}पि\appbrk{}चु\appbrk{}प्लो\appbrk{}त\appbrk{}प\appbrk{}त्र\appbrk{}सू\appbrk{}त्र\appbrk{}घृ\appbrk{}त\appbrk{}म\appbrk{}धु\appbrk{}प\appbrk{}य\appbrk{}स्तै\appbrk{}ल\appbrk{}त\appbrk{}र्प्प\appbrk{}ण\appbrk{}क\appbrk{}षा\appbrk{}या\appbrk{}ल\appbrk{}प\appbrk{}न\appbrk{}क\appbrk{}ल्क\appbrk{}शी\appbrk{}तो\appbrk{}द\appbrk{}क\appbrk{}व्य\appbrk{}ज\appbrk{}न\appbrk{}क\appbrk{}टा\appbrk{}हा\appbrk{}दीनि}\ H}}\skt{ परिकर्मिणश्च }\edtext{\skt{स्निग्धा}}{\lemma{\skt{स्निग्धा}}\Afootnote{K N H; \skt{स्निग्धाः}\ A}}\skt{ स्थिरा बलवन्तः~।}\pend

\pstart \ledsidenote{\passref{1.5.7}}\skt{ततः प्रशस्तेषु }\edtext{\skt{तिथि\-क\-र\-ण\-मु\-हू\-र्त्त\-न\-क्ष\-त्रेषु}}{\lemma{\skt{तिथि\appbrk{}क\appbrk{}र\appbrk{}ण\appbrk{}मु\appbrk{}हू\appbrk{}र्त्त\appbrk{}न\appbrk{}क्ष\appbrk{}त्रेषु}}\Afootnote{A K H; \skt{तिथि\appbrk{}क\appbrk{}र\appbrk{}ण\appbrk{}मु\appbrk{}हू\appbrk{}र्त्त\appbrk{}न\appbrk{}क्षत्रे}\ N}}\skt{ }\edtext{\edtext{\edtext{\skt{दध्य\-क्ष\-ता\-न्न\-पा\-न\-र\-त्नै\-र्वि\-प्रां\-श्चा\-र्च\-यित्वा}}{\lemma{\skt{दध्य\appbrk{}क्ष\appbrk{}ता\appbrk{}न्न\appbrk{}पा\appbrk{}न\appbrk{}र\appbrk{}त्नैर्}}\Afootnote{K N H; \skt{दध्य\appbrk{}क्ष\appbrk{}ता\appbrk{}न्न\appbrk{}पा\appbrk{}न\appbrk{}रत्नै}\ A}}}{\lemma{\skt{विप्रांश्}}\Afootnote{K N; \skt{रग्निं विप्रा}\ A; \skt{व्विप्रा}\ H}}}{\lemma{\skt{चार्चयित्वा}}\Afootnote{K N; \skt{न्भिष\appbrk{}ज\appbrk{}श्चा\appbrk{}र्च\appbrk{}यित्वा}\ A; \skt{न्श्चार्च्चयित्वा}\ H}}\skt{ कृत\-ब\-लि\-म\-ङ्ग\-ल\-स्व\-स्ति\-वा\-चनं }\edtext{\skt{भुक्त\-व\-न्त\-मा\-तुरं}}{\lemma{\skt{भुक्त\appbrk{}व\appbrk{}न्त\appbrk{}मा\appbrk{}तुरं}}\Afootnote{H; \skt{लघुभुक्तवन्तं प्राङ्मु\appbrk{}ख\appbrk{}मा\appbrk{}तु\appbrk{}रम्}\ A; \skt{भुक्त\appbrk{}व\appbrk{}न्त\appbrk{}मा\appbrk{}तुरं}\ K N}}\skt{ }\edtext{\edtext{\skt{प्राङ्मु\-ख\-मु\-प\-वेश्य}}{\lemma{\skt{प्राङ्मुखम्}}\Afootnote{K N; \emph{om.}\ A; \skt{प्राग्मुख}\ H}}}{\lemma{\skt{उपवेश्य}}\Afootnote{A K N; \skt{मुपवेश्य}\ H}}\skt{ यन्त्रयित्वा }\edtext{\skt{मर्म\-सि\-रा\-स्ना\-यु\-स\-न्ध्य\-स्थि\-ध\-मनीः}}{\lemma{\skt{मर्म\appbrk{}सि\appbrk{}रा\appbrk{}स्ना\appbrk{}यु\appbrk{}स\appbrk{}न्ध्य\appbrk{}स्थि\appbrk{}ध\appbrk{}मनीः}}\Afootnote{K H\textsuperscript{pc}; \skt{प्रत्यङ्मुखो वैद्यो मर्म\appbrk{}सि\appbrk{}रा\appbrk{}स्ना\appbrk{}यु\appbrk{}स\appbrk{}न्ध्य\appbrk{}स्थि\appbrk{}ध\appbrk{}मनीः}\ A; \skt{मम\appbrk{}सि\appbrk{}रा\appbrk{}स्ना\appbrk{}यु\appbrk{}स\appbrk{}न्ध्य\appbrk{}स्थि\appbrk{}ध\appbrk{}मनीः}\ N; \skt{मर्म्म\appbrk{}मि\appbrk{}रा\appbrk{}स्ना\appbrk{}यु\appbrk{}म\appbrk{}न्ध्य\appbrk{}स्थि\appbrk{}ध\appbrk{}मनीः}\ H\textsuperscript{ac}}}\skt{ }\edtext{\skt{परिहरन्ननु लोमं}}{\lemma{\skt{परिहरन्ननु लोमं}}\Afootnote{K; \skt{परि\appbrk{}ह\appbrk{}र\appbrk{}न\appbrk{}नु\appbrk{}लोमं}\ A H; \skt{परि\appbrk{}ह\appbrk{}र\appbrk{}न्न\appbrk{}नु\appbrk{}लोम}\ N}}\skt{ }\edtext{\edtext{\skt{शस्त्र\-न्नि\-द\-ध्या\-दा\-पू\-य\-द\-र्श\-ना\-त्स\-कृ\-दे\-वो\-प\-ह\-रे\-च्छ\-स्त्र\-माशु}}{\lemma{\skt{शस्त्र\appbrk{}न्नि\appbrk{}द\appbrk{}ध्या\appbrk{}दा\appbrk{}पू\appbrk{}य\appbrk{}द\appbrk{}र्श\appbrk{}ना\appbrk{}त्स\appbrk{}कृ\appbrk{}दे\appbrk{}वो\appbrk{}प\appbrk{}ह\appbrk{}रे\appbrk{}च्छ\appbrk{}स्त्रम्}}\Afootnote{\skt{शस्त्रं निद\appbrk{}ध्या\appbrk{}दा\appbrk{}पू\appbrk{}य\appbrk{}द\appbrk{}र्श\appbrk{}नात् , सकृ\appbrk{}दे\appbrk{}वा\appbrk{}प\appbrk{}ह\appbrk{}रे\appbrk{}च्छस्त्र}\ A; \skt{शस्त्र\appbrk{}न्नि\appbrk{}द\appbrk{}ध्या\appbrk{}दा\appbrk{}पू\appbrk{}य\appbrk{}द\appbrk{}र्श\appbrk{}ना\appbrk{}त्स\appbrk{}कृ\appbrk{}दे\appbrk{}वो\appbrk{}प\appbrk{}ह\appbrk{}रे\appbrk{}च्छ\appbrk{}स्त्रम्}\ K; \skt{शस्त्रं निद\appbrk{}ध्या\appbrk{}दा\appbrk{}पू\appbrk{}य\appbrk{}द\appbrk{}र्शना सकृ\appbrk{}दे\appbrk{}वो\appbrk{}प\appbrk{}ह\appbrk{}रे\appbrk{}च्छ\appbrk{}स्त्रम्}\ N; \skt{शस्त्र\appbrk{}न्नि\appbrk{}द\appbrk{}ध्या\appbrk{}दा\appbrk{}पू\appbrk{}य\appbrk{}द\appbrk{}र्श\appbrk{}ना\appbrk{}त्स\appbrk{}कृ\appbrk{}दे\appbrk{}वो\appbrk{}प\appbrk{}ह\appbrk{}रे\appbrk{}च्छस्त्र}\ H}}\skt{ च}}{\lemma{\skt{आशु च}}\Afootnote{K N; \skt{माशु च}\ A H}}\skt{~। }\edtext{\edtext{\skt{महत्स्वपि}}{\lemma{\skt{महत्स्व्}}\Afootnote{K N; \skt{महत्स्व}\ A H}}\skt{ च}}{\lemma{\skt{अपि च}}\Afootnote{K; \skt{पि च}\ A H; \skt{अपि}\ N}}\skt{ पाकेषु }\edtext{\skt{द्व्यङ्गु\-ल\-त्र्य\-ङ्गु\-ला\-न्त\-रम्वा}}{\lemma{\skt{द्व्यङ्गु\appbrk{}ल\appbrk{}त्र्य\appbrk{}ङ्गु\appbrk{}ला\appbrk{}न्त\appbrk{}रम्वा}}\Afootnote{K N; \skt{द्व्यङ्गुलान्तरं त्र्यङ्गुलान्तरं वा}\ A; \skt{द्व्यङ्गु\appbrk{}ल\appbrk{}न्त्र्य\appbrk{}ङ्गु\appbrk{}ङ्गु\appbrk{}ला\appbrk{}न्तरं वा}\ H}}\skt{ }\edtext{\edtext{\skt{शस्त्रपदमुक्तं}}{\lemma{\skt{शस्त्रपदम्}}\Afootnote{A K N\textsuperscript{pc}; \skt{पशस्त्रपदम्}\ N\textsuperscript{ac}; \skt{शस्त्रपद}\ H}}}{\lemma{\skt{उक्तं}}\Afootnote{A K N; \skt{मुक्तं}\ H}}\skt{~॥}\pend

\pstart \ledsidenote{\passref{1.5.8}}\skt{तत्रायतो }\edtext{\skt{विशालः}}{\lemma{\skt{विशालः}}\Afootnote{A K N; \skt{विषाल}\ H}}\skt{ }\edtext{\skt{समः}}{\lemma{\skt{समः}}\Afootnote{A K N; \skt{स्ममस्}\ H}}\skt{ }\edtext{\skt{सुविभक्त}}{\lemma{\skt{सुविभक्त}}\Afootnote{K N H; \skt{सुविभक्तो}\ A}}\skt{ }\edtext{\skt{इति}}{\lemma{\skt{इति}}\Afootnote{H; \skt{निराश्रय इति}\ A; \skt{इति व्रणाः}\ K N}}\skt{ }\edtext{\skt{व्रणाः}}{\lemma{\skt{व्रणाः}}\Afootnote{H; \skt{व्रणगुणाः}\ A; \emph{om.}\ K N}}\skt{~॥}\pend

\pstart \ledsidenote{\passref{1.5.11}}\skt{एकेन वा }\edtext{\skt{व्रणेन न}}{\lemma{\skt{व्रणेन न}}\Afootnote{K N H; \skt{व्रणेना}\ A}}\skt{ }\edtext{\skt{विशुद्ध्यति}}{\lemma{\skt{विशुद्ध्यति}}\Afootnote{K N H; \skt{शुध्यमा}\ A}}\skt{ }\edtext{\edtext{\skt{ततोऽपरां}}{\lemma{\skt{ततो}}\Afootnote{K N; \skt{ने}\ A; \skt{ततोऽ}\ H}}}{\lemma{\skt{ऽपरां}}\Afootnote{K N; \skt{नाऽन्तरा}\ A; \skt{परां}\ H}}\skt{ }\edtext{\skt{बुद्ध्यापेक्षान्तरं}}{\lemma{\skt{बुद्ध्यापेक्षान्तरं}}\Afootnote{K; \skt{बुद्ध्याऽवेक्ष्यापरा}\ A; \skt{बुद्ध्यापेक्ष्यान्तरं}\ N; \skt{बुद्व्यापेक्षान्तरं}\ H}}\skt{ }\edtext{\edtext{\skt{व्रणाङ्कुर्यात्}}{\lemma{\skt{व्रणाङ्}}\Afootnote{K N; \skt{न्व्रणा}\ A; \skt{व्रणङ्}\ H}}}{\lemma{\skt{कुर्यात्}}\Afootnote{K N H; \skt{न्कुर्यात्}\ A}}\skt{~॥}\pend

\pstart\versepar \ledsidenote{\passref{1.5.9}}\edtext{\skt{भ}}{\lemma{\skt{भ}}\Afootnote{N H; \skt{भवतश्चात्र}\ A; \skt{भ वन्ति चात्र श्लोकाः।}\ K}}\skt{ }\pend
\pstart\versepar \skt{आयतश्च विशालश्च }\edtext{\skt{सुविभक्तो}}{\lemma{\skt{सुविभक्तो}}\Afootnote{A K N; \skt{सूविभक्तो}\ H}}\skt{ निराश्रयः~। }\pend
\pstart\versepar \skt{}\edtext{\skt{प्राप्त\-का\-ल\-कृ\-त\-श्चैव}}{\lemma{\skt{चैव}}\Afootnote{K N H; \skt{चापि}\ A}}\skt{ }\edtext{\skt{व्रणाः}}{\lemma{\skt{व्रणाः}}\Afootnote{K; \skt{व्रणः}\ A N H}}\skt{ कर्मणि }\edtext{\skt{शस्यते}}{\lemma{\skt{शस्यते}}\Afootnote{A K N; \skt{सश्यते}\ H}}\skt{~॥}\pend

\pstart\versepar \ledsidenote{\passref{1.5.10}}\edtext{\edtext{\skt{शौर्य\-मा\-शु\-क्रि\-या\-तीक्ष्णं}}{\lemma{\skt{शौर्यम्}}\Afootnote{K N; \skt{शौर्य}\ A; \skt{सौर्य}\ H}}}{\lemma{\skt{आशुक्रियातीक्ष्णं}}\Afootnote{K N; \skt{माशुक्रिया}\ A; \skt{माशुक्रियातीक्ष्णं}\ H}}\skt{ }\edtext{\edtext{\skt{शस्त्र\-म\-स्वे\-द\-वे\-पथुः}}{\lemma{\skt{शस्त्रम्}}\Afootnote{K N; \skt{शस्त्र}\ A H}}}{\lemma{\skt{अस्वेदवेपथुः}}\Afootnote{K N; \skt{तैक्ष्ण्य\appbrk{}म\appbrk{}स्वे\appbrk{}द\appbrk{}वे\appbrk{}पथु}\ A; \skt{मस्वेदवेपथुः}\ H}}\skt{~। }\pend
\pstart\versepar \skt{}\edtext{\skt{असम्मोहश्च}}{\lemma{\skt{असम्मोहश्च}}\Afootnote{A K N H\textsuperscript{pc}; \skt{अमंमोहश्च}\ H\textsuperscript{ac}}}\skt{ }\edtext{\skt{वैद्यस्य}}{\lemma{\skt{वैद्यस्य}}\Afootnote{A K N; \skt{वैद्यस्य…}\ H}}\skt{ शस्त्रकर्मणि }\edtext{\skt{पूज्यते}}{\lemma{\skt{पूज्यते}}\Afootnote{K N H; \skt{शस्यते}\ A}}\skt{~॥}\pend

\pstart\versepar \ledsidenote{\passref{1.5.12}}\edtext{\skt{यतो}}{\lemma{\skt{यतो}}\Afootnote{K N H; \skt{भवति चात्र यतो}\ A}}\skt{ यतो गतिं }\edtext{\skt{विद्यादुत्सङ्गो}}{\lemma{\skt{विद्यादुत्सङ्गो}}\Afootnote{K N; \skt{विद्यादुत्सङ्गो}\ A H}}\skt{ यत्र }\edtext{\skt{यत्र च}}{\lemma{\skt{यत्र च}}\Afootnote{A K H; \skt{य च}\ N}}\skt{~। }\pend
\pstart\versepar \skt{तत्र तत्र }\edtext{\skt{वर्णं}}{\lemma{\skt{वर्णं}}\Afootnote{K N; \skt{व्रणं}\ A H}}\skt{ }\edtext{\skt{कुर्याद्यथा}}{\lemma{\skt{कुर्याद्यथा}}\Afootnote{A K N; \skt{कुर्याद्यथा}\ H}}\skt{ दोषो न तिष्ठति~॥}\pend

\pstart \ledsidenote{\passref{1.5.13}}\skt{तत्र }\edtext{\skt{भ्रूग\-ण्ड\-ल\-ला\-टा\-क्षि\-पु\-ट\-क\-क्षा\-वं\-क्ष्णेषु}}{\lemma{\skt{भ्रूग\appbrk{}ण्ड\appbrk{}ल\appbrk{}ला\appbrk{}टा\appbrk{}क्षि\appbrk{}पु\appbrk{}ट\appbrk{}क\appbrk{}क्षा\appbrk{}वं\appbrk{}क्ष्णेषु}}\Afootnote{K; \skt{भ्रूग\appbrk{}ण्ड\appbrk{}श\appbrk{}ङ्ख\appbrk{}ल\appbrk{}ला\appbrk{}टा\appbrk{}क्षि\appbrk{}पु\appbrk{}टौ\appbrk{}ष्ठ\appbrk{}द\appbrk{}न्त\appbrk{}वे\appbrk{}ष्ट\appbrk{}क\appbrk{}क्षा\appbrk{}कु\appbrk{}क्षि\appbrk{}व\appbrk{}ङ्क्ष\appbrk{}णेषु}\ A; \skt{भ्रूग\appbrk{}ण्ड\appbrk{}ल\appbrk{}ला\appbrk{}टो\appbrk{}क्षि\appbrk{}कु\appbrk{}ट\appbrk{}क\appbrk{}क्षा\appbrk{}वं\appbrk{}क्ष\appbrk{}णेषु}\ N; \skt{भ्रूग\appbrk{}णु\appbrk{}ल\appbrk{}ला\appbrk{}टा\appbrk{}क्षि\appbrk{}कू\appbrk{}ट\appbrk{}क\appbrk{}क्षा\appbrk{}वं\appbrk{}क्ष\appbrk{}णेषु}\ H}}\skt{ तिर्यक्छेद उक्तः~।}\pend

\pstart \ledsidenote{\passref{1.5.15}}\edtext{\skt{अन्यथा तु}}{\lemma{\skt{अन्यथा तु}}\Afootnote{A K N H\textsuperscript{ac}; \skt{डेअन्यथा तु}\ H\textsuperscript{pc}}}\skt{ }\edtext{\skt{सिरा\-स्ना\-यु\-क्ष\-ण\-ना\-द\-ति\-मात्रं}}{\lemma{\skt{सिरा\appbrk{}स्ना\appbrk{}यु\appbrk{}क्ष\appbrk{}ण\appbrk{}ना\appbrk{}द\appbrk{}ति\appbrk{}मात्रं}}\Afootnote{K N H; \skt{सिरा\appbrk{}स्ना\appbrk{}यु\appbrk{}च्छे\appbrk{}द\appbrk{}नम्, अतिमात्रं}\ A}}\skt{~। }\edtext{\skt{वेदनाचिराच्च}}{\lemma{\skt{वेदनाचिराच्च}}\Afootnote{H; \skt{वेदना, चिरा}\ A; \skt{वेदनाचिराच्च}\ K N}}\skt{ }\edtext{\skt{व्रणसंरोहो}}{\lemma{\skt{व्रणसंरोहो}}\Afootnote{K N H; \skt{द्व्रणसंरोहो}\ A}}\skt{ }\edtext{\skt{मांस\-क\-न्दी\-प्रा\-दु\-र्भा\-वश्च}}{\lemma{\skt{मांस\appbrk{}क\appbrk{}न्दी\appbrk{}प्रा\appbrk{}दु\appbrk{}र्भा\appbrk{}वश्च}}\Afootnote{K N H; \skt{मांस\appbrk{}क\appbrk{}न्दी\appbrk{}प्रा\appbrk{}दु\appbrk{}र्भा\appbrk{}वश्चे}\ A}}\skt{ }\edtext{\skt{भवति}}{\lemma{\skt{भवति}}\Afootnote{K N H; \skt{ति}\ A}}\skt{~।}\pend

\pstart \ledsidenote{\passref{1.5.16}}\edtext{\skt{मूढ\-ग\-र्भो\-द\-रा\-श्म\-री\-भ\-ग\-न्द\-र\-मु\-ख\-रा\-गे\-ष्व\-भु\-क्त\-वतां}}{\lemma{\skt{मूढ\appbrk{}ग\appbrk{}र्भो\appbrk{}द\appbrk{}रा\appbrk{}श्म\appbrk{}री\appbrk{}भ\appbrk{}ग\appbrk{}न्द\appbrk{}र\appbrk{}मु\appbrk{}ख\appbrk{}रा\appbrk{}गे\appbrk{}ष्व\appbrk{}भु\appbrk{}क्त\appbrk{}वतां}}\Afootnote{\skt{मूढ\appbrk{}ग\appbrk{}र्भो\appbrk{}द\appbrk{}रा\appbrk{}र्शोऽ\appbrk{}श्म\appbrk{}री\appbrk{}भ\appbrk{}ग\appbrk{}न्द\appbrk{}र\appbrk{}मु\appbrk{}ख\appbrk{}रो\appbrk{}गे\appbrk{}ष्व\appbrk{}भु\appbrk{}क्त\appbrk{}वतः}\ A; \skt{मूढ\appbrk{}ग\appbrk{}र्भो\appbrk{}द\appbrk{}रा\appbrk{}श्म\appbrk{}री\appbrk{}भ\appbrk{}ग\appbrk{}न्द\appbrk{}र\appbrk{}मु\appbrk{}ख\appbrk{}रा\appbrk{}गे\appbrk{}ष्व\appbrk{}भु\appbrk{}क्त\appbrk{}वतां}\ K; \skt{मूढ\appbrk{}ग\appbrk{}र्भो\appbrk{}द\appbrk{}रा\appbrk{}श्म\appbrk{}भ\appbrk{}ग\appbrk{}न्दं\appbrk{}र\appbrk{}मु\appbrk{}ख\appbrk{}रा\appbrk{}गे\appbrk{}ष्व\appbrk{}भु\appbrk{}क्त\appbrk{}वतां}\ N; \skt{मूढ\appbrk{}ग\appbrk{}र्भौ\appbrk{}द\appbrk{}रा\appbrk{}श्म\appbrk{}री\appbrk{}भ\appbrk{}ग\appbrk{}न्द\appbrk{}र\appbrk{}मु\appbrk{}ख\appbrk{}रा\appbrk{}गे\appbrk{}ष्व\appbrk{}भु\appbrk{}क्त\appbrk{}वतां}\ H}}\skt{ }\edtext{\skt{कुर्वीतः}}{\lemma{\skt{कुर्वीतः}}\Afootnote{K N; \skt{कर्म कुर्वीत}\ A; \skt{कुर्व्वीत…}\ H}}\pend

\pstart \ledsidenote{\passref{1.5.17}}\skt{ततः }\edtext{\edtext{\skt{शस्त्रमवचार्य}}{\lemma{\skt{शस्त्रम्}}\Afootnote{K N H; \skt{शस्त्र}\ A}}}{\lemma{\skt{अवचार्य}}\Afootnote{K N H; \skt{मवचार्य}\ A}}\skt{ }\edtext{\skt{शीताभिरद्भि}}{\lemma{\skt{अद्भि}}\Afootnote{K\textsuperscript{pc} N; \skt{अद्भिर्}\ A K\textsuperscript{ac} H}}\skt{ }\edtext{\edtext{\edtext{\skt{परि\-षि\-च्य\-रा\-तु\-र\-मा\-श्वास्य}}{\lemma{\skt{परिषिच्यर्}}\Afootnote{K\textsuperscript{pc}; \emph{om.}\ A K\textsuperscript{ac}; \skt{परिषिच्य}\ N; \skt{परिषिंच्य}\ H}}}{\lemma{\skt{आतुरम्}}\Afootnote{\skt{आतुर}\ A; \skt{आतुर}\ K; \skt{चातुरम्}\ N H}}\skt{ च}}{\lemma{\skt{आश्वास्य च}}\Afootnote{N H; \skt{माश्वास्य}\ A K\textsuperscript{ac}; \skt{माश्वास्य च}\ K\textsuperscript{pc}}}\skt{ }\edtext{\skt{सम\-न्ता\-त्प\-रि\-पी\-ड्यां\-गुल्या}}{\lemma{\skt{सम\appbrk{}न्ता\appbrk{}त्प\appbrk{}रि\appbrk{}पी\appbrk{}ड्यां\appbrk{}गुल्या}}\Afootnote{A K N H\textsuperscript{pc}; \skt{मम\appbrk{}न्ता\appbrk{}त्प\appbrk{}रि\appbrk{}पी\appbrk{}ड्यां\appbrk{}गुल्या}\ H\textsuperscript{ac}}}\skt{ }\edtext{\edtext{\skt{व्रणमभिमृज्य}}{\lemma{\skt{व्रणम्}}\Afootnote{K N H; \skt{व्रण}\ A}}}{\lemma{\skt{अभिमृज्य}}\Afootnote{K N H\textsuperscript{ac}; \skt{मभिमृद्य(ज्य)}\ A; \skt{अभिसृज्य}\ H\textsuperscript{pc}}}\skt{ प्रक्षाल्य कषायेण }\edtext{\edtext{\skt{प्लोते\-नो\-द\-क\-मा\-दाय}}{\lemma{\skt{प्लोतेनोदकम्}}\Afootnote{K N H; \skt{प्रोतेनोदक}\ A}}}{\lemma{\skt{आदाय}}\Afootnote{K N H; \skt{मादाय}\ A}}\skt{ }\edtext{\skt{तिल\-क\-ल्क\-म\-धु\-गाढां}}{\lemma{\skt{तिल\appbrk{}क\appbrk{}ल्क\appbrk{}म\appbrk{}धु\appbrk{}गाढां}}\Afootnote{K; \skt{तिल\appbrk{}क\appbrk{}ल्क\appbrk{}म\appbrk{}धु\appbrk{}स\appbrk{}र्पिः\appbrk{}प्र\appbrk{}गा\appbrk{}ढा\appbrk{}मौ\appbrk{}ष\appbrk{}ध\appbrk{}युक्तां नातिस्निग्धां नातिरूक्षां}\ A; \skt{तिल\appbrk{}क\appbrk{}ल्क\appbrk{}म\appbrk{}धु\appbrk{}स\appbrk{}र्प्पि\appbrk{}प्र\appbrk{}गाढं}\ N; \skt{तिल\appbrk{}क\appbrk{}ल्क\appbrk{}म\appbrk{}धु\appbrk{}सर्प्पिः प्रगां}\ H\textsuperscript{ac}; \skt{तिल\appbrk{}क\appbrk{}ल्क\appbrk{}म\appbrk{}धु\appbrk{}सर्प्पिः प्रगाढां}\ H\textsuperscript{pc}}}\skt{ }\edtext{\skt{वर्त्तिम्प्रणिधाय}}{\lemma{\skt{प्रणिधाय}}\Afootnote{K N H; \skt{प्रणिदध्यात्; ततः}\ A}}\skt{ }\edtext{\skt{पत्रेणाच्छाद्य}}{\lemma{\skt{पत्रेणाच्छाद्य}}\Afootnote{K N H; \skt{कल्केनाच्छाद्य}\ A}}\skt{ }\edtext{\skt{कवलिकान्दत्वा}}{\lemma{\skt{कवलिकान्}}\Afootnote{K N H; \skt{घनां कवलिकां}\ A}}\skt{ }\edtext{\skt{बन्ध\-नो\-प\-पा\-द\-येत्~। }\edtext{\skt{वेद\-न\-आ\-र\-क्षो\-घ्नै\-र्धू\-प\-यित्वा}}{\lemma{\skt{धूपयित्वा}}\Afootnote{K; \skt{रक्षां कुर्वीत}\ A; \emph{om.}\ N H}}}{\lemma{\skt{बन्ध\appbrk{}नो\appbrk{}प\appbrk{}पा\appbrk{}द\appbrk{}येत्~। वेद\appbrk{}न\appbrk{}आ\appbrk{}र\appbrk{}क्षो\appbrk{}घ्नैर्}}\Afootnote{K\textsuperscript{pc}; \skt{वस्त्रपट्टेन बध्नीयात्; वेद\appbrk{}ना\appbrk{}र\appbrk{}क्षो\appbrk{}घ्नै\appbrk{}र्धू\appbrk{}पै\appbrk{}र्धू\appbrk{}प\appbrk{}येत्, रक्षोघ्नैश्च मन्त्रै}\ A; \skt{बन्ध\appbrk{}नो\appbrk{}प\appbrk{}पा\appbrk{}द\appbrk{}येत्~। वेद\appbrk{}न\appbrk{}र\appbrk{}क्षो\appbrk{}घ्नैर्}\ K\textsuperscript{ac}; \skt{बन्धे\appbrk{}नो\appbrk{}प\appbrk{}पा\appbrk{}द\appbrk{}येत्~। वेदनारक्षघ्नै धूपयित्वा}\ N; \skt{बन्ध\appbrk{}नो\appbrk{}प\appbrk{}पा\appbrk{}द\appbrk{}येत्~॥ वेद\appbrk{}ना\appbrk{}र\appbrk{}क्षो\appbrk{}घ्नै\appbrk{}र्द्धू\appbrk{}पैत्वा}\ H}}\pend

\pstart \ledsidenote{\passref{1.5.18}}\edtext{\skt{गुग्गु\-ल्व\-ग\-रु\-स\-र्ज्ज\-र\-स\-व\-चा\-गौ\-र\-स\-र्प\-ल\-व\-ण\-नि\-म्ब\-प\-त्रा\-ज्य\-शे\-षेण}}{\lemma{\skt{गुग्गु\appbrk{}ल्व\appbrk{}ग\appbrk{}रु\appbrk{}स\appbrk{}र्ज्ज\appbrk{}र\appbrk{}स\appbrk{}व\appbrk{}चा\appbrk{}गौ\appbrk{}र\appbrk{}स\appbrk{}र्प\appbrk{}ल\appbrk{}व\appbrk{}ण\appbrk{}नि\appbrk{}म्ब\appbrk{}प\appbrk{}त्रा\appbrk{}ज्य\appbrk{}शे\appbrk{}षेण}}\Afootnote{K; \skt{ततो गुग्गु\appbrk{}ल्व\appbrk{}गु\appbrk{}रु\appbrk{}स\appbrk{}र्ज\appbrk{}र\appbrk{}स\appbrk{}व\appbrk{}चा\appbrk{}गौ\appbrk{}र\appbrk{}स\appbrk{}र्ष\appbrk{}प\appbrk{}चू\appbrk{}र्णै\appbrk{}र्ल\appbrk{}व\appbrk{}ण\appbrk{}नि\appbrk{}म्ब\appbrk{}प\appbrk{}त्र\appbrk{}वि\appbrk{}मि\appbrk{}श्रै\appbrk{}रा\appbrk{}ज्य\appbrk{}यु\appbrk{}क्तै\appbrk{}र्धू\appbrk{}प\appbrk{}येत्~। आज्यशेषेण}\ A; \skt{गुग्गु\appbrk{}ल्व\appbrk{}ग\appbrk{}रु\appbrk{}स\appbrk{}ज्ज\appbrk{}र\appbrk{}स\appbrk{}वा\appbrk{}चा\appbrk{}गौ\appbrk{}र\appbrk{}स\appbrk{}र्ष\appbrk{}प\appbrk{}ल\appbrk{}व\appbrk{}ण\appbrk{}नि\appbrk{}म्ब\appbrk{}प\appbrk{}त्रा\appbrk{}ज्य\appbrk{}मि\appbrk{}श्रै\appbrk{}रा\appbrk{}ज्य\appbrk{}शे\appbrk{}षेण}\ N\textsuperscript{ac}; \skt{गुग्गु\appbrk{}ल्व\appbrk{}ग\appbrk{}रु\appbrk{}स\appbrk{}ज्ज\appbrk{}र\appbrk{}स\appbrk{}व\appbrk{}चा\appbrk{}गौ\appbrk{}र\appbrk{}स\appbrk{}र्ष\appbrk{}प\appbrk{}ल\appbrk{}व\appbrk{}ण\appbrk{}नि\appbrk{}म्ब\appbrk{}प\appbrk{}त्रा\appbrk{}ज्य\appbrk{}मि\appbrk{}श्रै\appbrk{}रा\appbrk{}ज्य\appbrk{}शे\appbrk{}षेण}\ N\textsuperscript{pc}; \skt{गुग्गु\appbrk{}ल्व\appbrk{}गु\appbrk{}रु\appbrk{}म\appbrk{}र्ज्ज\appbrk{}र\appbrk{}स\appbrk{}व\appbrk{}चा\appbrk{}गौ\appbrk{}र\appbrk{}स\appbrk{}र्ष\appbrk{}प\appbrk{}ल\appbrk{}व\appbrk{}ण\appbrk{}नि\appbrk{}म्ब\appbrk{}प\appbrk{}त्रा\appbrk{}ज्य\appbrk{}मिश्रै राज्यशेषेण}\ H\textsuperscript{ac}; \skt{गुग्गु\appbrk{}ल्व\appbrk{}गु\appbrk{}रु\appbrk{}स\appbrk{}र्ज्ज\appbrk{}र\appbrk{}स\appbrk{}व\appbrk{}चा\appbrk{}गौ\appbrk{}र\appbrk{}स\appbrk{}र्ष\appbrk{}प\appbrk{}ल\appbrk{}व\appbrk{}ण\appbrk{}नि\appbrk{}म्ब\appbrk{}प\appbrk{}त्रा\appbrk{}ज्य\appbrk{}मिश्रै राज्यशेषेण}\ H\textsuperscript{pc}}}\skt{ चास्य }\edtext{\skt{प्राणां}}{\lemma{\skt{प्राणां}}\Afootnote{K N; \skt{प्राणा}\ A H}}\skt{ }\edtext{\skt{समालभेत}}{\lemma{\skt{समालभेत}}\Afootnote{K N; \skt{न्समालभेत}\ A H}}\skt{~।}\pend

\pstart \ledsidenote{\passref{1.5.19}}\edtext{\skt{उदकुम्भाच्चापो}}{\lemma{\skt{उदकुम्भाच्चापो}}\Afootnote{K N; \skt{उदकुम्भाच्चापो}\ A H}}\skt{ गृहीत्वा प्रोक्ष\-य\-न्र\-क्षा\-कर्म्म }\edtext{\skt{कुर्यात्}}{\lemma{\skt{कुर्यात्}}\Afootnote{K N H; \skt{कुर्यात् ; तद्वक्ष्यामः}\ A}}\skt{~॥}\pend

\pstart\versepar \ledsidenote{\passref{1.5.20}}\edtext{\skt{कृत्या\-ना\-म्प\-रि\-र\-क्षार्थ}}{\lemma{\skt{परिरक्षार्थ}}\Afootnote{K N; \skt{प्रतिघातार्थं}\ A; \skt{परिरक्षार्थन्}\ H}}\skt{ तथा रक्षोभयस्य च }\pend
\pstart\versepar \skt{रक्षाकर्म्म }\edtext{\skt{करिष्यामि}}{\lemma{\skt{करिष्यामि}}\Afootnote{A K N; \skt{करिष्यामि …}\ H}}\skt{ ब्रह्मा }\edtext{\skt{तदनुमन्यतां}}{\lemma{\skt{तदनुमन्यतां}}\Afootnote{K N H; \skt{तद\appbrk{}नु\appbrk{}म\appbrk{}न्य\appbrk{}ताम्}\ A}}\skt{~।}\pend

\pstart\versepar \ledsidenote{\passref{1.5.21}}\edtext{\skt{नागा}}{\lemma{\skt{नागा}}\Afootnote{K N H; \skt{नागाः}\ A}}\skt{ पिशाचा }\edtext{\skt{गन्धर्वा}}{\lemma{\skt{गन्धर्वा}}\Afootnote{K N H; \skt{गन्धर्वाः}\ A}}\skt{ }\edtext{\skt{यक्ष\-र\-क्षं\-स्य\-थ\-ग्रहाः}}{\lemma{\skt{यक्ष\appbrk{}र\appbrk{}क्षं\appbrk{}स्य\appbrk{}थ\appbrk{}ग्रहाः}}\Afootnote{\skt{पितरो यक्षराक्षसाः}\ A; \skt{यक्षरक्षंस्यथ ग्रहाः}\ K; \skt{यक्षा रक्षांस्यथ ग्रहाः~।}\ N; \skt{यक्षा\appbrk{}र्क्षां\appbrk{}स्य\appbrk{}थ\appbrk{}ग्रहाः}\ H}}\skt{~। }\pend
\pstart\versepar \skt{}\edtext{\skt{अभिद्रवन्ति ये ये त्वा}}{\lemma{\skt{अभिद्रवन्ति ये ये त्वा}}\Afootnote{K N H; \skt{अभिद्रवन्ति ये ये त्वां}\ A}}\skt{ ब्रह्माद्याः घ्नन्तु }\edtext{\skt{तां सदा}}{\lemma{\skt{तां सदा}}\Afootnote{K; \skt{तान्सदा}\ A; \skt{तांन्सदा}\ N; \skt{तान्मदा}\ H}}\skt{~।}\pend

\pstart\versepar \ledsidenote{\passref{1.5.22}}\edtext{\skt{पृथि\-व्या\-म\-न्त\-रिक्षे च ये}}{\lemma{\skt{पृथि\appbrk{}व्या\appbrk{}म\appbrk{}न्त\appbrk{}रिक्षे च ये}}\Afootnote{\skt{पृथि\appbrk{}व्या\appbrk{}म\appbrk{}न्त\appbrk{}रीक्षे च ये}\ A; \skt{पृथि\appbrk{}व्या\appbrk{}म\appbrk{}न्त\appbrk{}रिक्षे च ये}\ K N; \skt{पृथि\appbrk{}व्या\appbrk{}म\appbrk{}न्त\appbrk{}रीक्षे च ये}\ H}}\skt{ चरन्ति निशाचराः~। }\pend
\pstart\versepar \skt{दिक्षु }\edtext{\skt{वास्तुनि ये}}{\lemma{\skt{वास्तुनि ये}}\Afootnote{K N H; \skt{वास्तुनिवा}\ A}}\skt{ }\edtext{\skt{चान्ये}}{\lemma{\skt{चान्ये}}\Afootnote{K N H; \skt{साश्च}\ A}}\skt{ }\edtext{\skt{पान्तु त्वा ते}}{\lemma{\skt{पान्तु त्वा ते}}\Afootnote{K N H; \skt{पान्तु त्वां ते}\ A}}\skt{ नमस्कृताः~।}\pend

\pstart\versepar \ledsidenote{\passref{1.5.23}}\edtext{\skt{पान्तु}}{\lemma{\skt{पान्तु}}\Afootnote{A K H; \skt{पां}\ N}}\skt{ }\edtext{\skt{त्वामृषयो}}{\lemma{\skt{त्वामृषयो}}\Afootnote{N H; \skt{त्वां मुनयो}\ A; \skt{त्वामृषयो}\ K}}\skt{ }\edtext{\skt{ब्रह्मविद्या}}{\lemma{\skt{ब्रह्मविद्या}}\Afootnote{K; \skt{ब्राह्मया दिव्या}\ A; \skt{ब्रह्मा विद्या}\ N; \skt{ब्राहमा विद्या}\ H}}\skt{ }\edtext{\edtext{\skt{राजर्षयस्तथा}}{\lemma{\skt{राजर्षयस्}}\Afootnote{K; \skt{राजर्षय}\ A H; \skt{राजार्षयस्}\ N}}}{\lemma{\skt{तथा}}\Afootnote{K N; \skt{स्तथा}\ A H}}\skt{~। }\pend
\pstart\versepar \skt{पर्वताश्चैव नद्यश्च सर्वाः सर्वे च सागराः~।}\pend

\pstart\versepar \ledsidenote{\passref{1.5.24}}\skt{अग्नी रक्षतु ते जिह्वां }\edtext{\skt{प्राणं}}{\lemma{\skt{प्राणं}}\Afootnote{K N H; \skt{प्राणा}\ A}}\skt{ }\edtext{\edtext{\skt{वायुस्तथैव}}{\lemma{\skt{वायुस्}}\Afootnote{K N; \skt{न्वायुस्}\ A; \skt{वायु}\ H}}\skt{ च}}{\lemma{\skt{तथैव च}}\Afootnote{A K N; \skt{स्तथैव च}\ H}}\skt{~। }\pend
\pstart\versepar \skt{सोमो }\edtext{\edtext{\skt{व्यानमपानन्ते}}{\lemma{\skt{व्यानम्}}\Afootnote{K N; \skt{व्यान}\ A H}}}{\lemma{\skt{अपानन्ते}}\Afootnote{K N; \skt{मपानं ते}\ A H}}\skt{ }\edtext{\skt{पर्जन्यः}}{\lemma{\skt{पर्जन्यः}}\Afootnote{A K H; \skt{पज्जन्यः}\ N}}\skt{ परिरक्षतु~।}\pend

\pstart\versepar \ledsidenote{\passref{1.5.25}}\skt{उदानम्विद्युतः पान्तु }\edtext{\skt{समानं}}{\lemma{\skt{समानं}}\Afootnote{A K N H\textsuperscript{pc}; \skt{समानं म्}\ H\textsuperscript{ac}}}\skt{ }\edtext{\skt{स्तनयित्नवः}}{\lemma{\skt{स्तनयित्नवः}}\Afootnote{A K H\textsuperscript{pc}; \skt{स्तनयिन्नवः}\ N; \skt{तनयिन्नवः}\ H\textsuperscript{ac}}}\skt{~। }\pend
\pstart\versepar \skt{}\edtext{\skt{बलमिन्द्रो}}{\lemma{\skt{बलमिन्द्रो}}\Afootnote{H; \skt{बलमिन्द्रो}\ A K N}}\skt{ }\edtext{\edtext{\edtext{\edtext{\skt{बल\-प\-ति\-र्म\-ति\-म्वा\-च\-स्प\-ति\-स्तथा}}{\lemma{\skt{बलपतिर्}}\Afootnote{A K N; \skt{बलपति}\ H}}}{\lemma{\skt{मतिम्}}\Afootnote{K N; \skt{मनुर्}\ A; \skt{र्म्मतिं}\ H}}}{\lemma{\skt{वाचस्पतिस्}}\Afootnote{K N; \skt{मन्ये मतिं}\ A; \skt{वाचस्पति}\ H}}}{\lemma{\skt{तथा}}\Afootnote{A K N; \skt{स्तथा}\ H}}\skt{~।}\pend

\pstart\versepar \ledsidenote{\passref{1.5.26}}\edtext{\skt{कामान्ते}}{\lemma{\skt{कामान्ते}}\Afootnote{K N H; \skt{कामांस्ते}\ A}}\skt{ }\edtext{\skt{पान्तु}}{\lemma{\skt{पान्तु}}\Afootnote{A K H; \emph{om.}\ N}}\skt{ }\edtext{\skt{गन्ध\-र्वा\-स्स\-त्व\-मि\-न्द्रो\-भि\-र\-क्षतु}}{\lemma{\skt{गन्ध\appbrk{}र्वा\appbrk{}स्स\appbrk{}त्व\appbrk{}मि\appbrk{}न्द्रो\appbrk{}भि\appbrk{}र\appbrk{}क्षतु}}\Afootnote{H; \skt{गन्धर्वाः सत्त्व\appbrk{}मि\appbrk{}न्द्रोऽ\appbrk{}भि\appbrk{}र\appbrk{}क्षतु}\ A; \skt{गन्ध\appbrk{}र्वा\appbrk{}स्स\appbrk{}त्व\appbrk{}मि\appbrk{}न्द्रोऽ\appbrk{}भि\appbrk{}र\appbrk{}क्षतु}\ K; \skt{गन्ध\appbrk{}र्वा\appbrk{}स\appbrk{}त्व\appbrk{}मिन्द्रो भिरक्षतु}\ N}}\skt{~। }\pend
\pstart\versepar \skt{प्रज्ञान्ते वरुणो राजा समुद्रो नाभिमण्डलं~।}\pend

\pstart\versepar \ledsidenote{\passref{1.5.27}}\skt{चक्षुः सूर्यो दिशः }\edtext{\skt{श्रोत्रं}}{\lemma{\skt{श्रोत्रं}}\Afootnote{K N H; \skt{श्रोत्रे}\ A}}\skt{ }\edtext{\skt{चन्द्रमा}}{\lemma{\skt{चन्द्रमा}}\Afootnote{K; \skt{चन्द्रमाः}\ A H; \skt{चममा}\ N}}\skt{ पातु ते मनः~। }\pend
\pstart\versepar \skt{नक्षत्राणि सदा रूपं }\edtext{\skt{छायाम्पातु}}{\lemma{\skt{पातु}}\Afootnote{K N H; \skt{पान्तु}\ A}}\skt{ निशा }\edtext{\skt{तव}}{\lemma{\skt{तव}}\Afootnote{K N H; \skt{स्तव}\ A}}\skt{~।}\pend

\pstart\versepar \ledsidenote{\passref{1.5.28}}\edtext{\skt{रेत\-स्त्वा\-प्या\-य\-यं\-त्वापो}}{\lemma{\skt{रेत\appbrk{}स्त्वा\appbrk{}प्या\appbrk{}य\appbrk{}यं\appbrk{}त्वापो}}\Afootnote{\skt{रेत\appbrk{}स्त्वा\appbrk{}प्या\appbrk{}य\appbrk{}य\appbrk{}न्त्वापो}\ A H; \skt{रेतस्त्वाप्याययं त्वापो}\ K; \skt{रेत\appbrk{}स्त्वा\appbrk{}प्या\appbrk{}य\appbrk{}यं\appbrk{}त्वापो}\ N}}\skt{ }\edtext{\edtext{\skt{रोमा\-ण्यौ\-ष\-ध\-य\-स्तथा}}{\lemma{\skt{रोमाण्य्}}\Afootnote{K\textsuperscript{pc} N; \skt{रोमाण्यो}\ A; \skt{रोण्य्}\ K\textsuperscript{ac}; \skt{रोमाण्यौ}\ H}}}{\lemma{\skt{औषधयस्}}\Afootnote{K; \skt{षधयस्}\ A H; \skt{ओषधयस्}\ N}}\skt{~। }\pend
\pstart\versepar \skt{आकाशङ्खानि ते }\edtext{\skt{पातु}}{\lemma{\skt{पातु}}\Afootnote{K N H; \skt{पान्तु}\ A}}\skt{ देहन्तव वसुन्धरा~।}\pend

\pstart\versepar \ledsidenote{\passref{1.5.29}}\skt{वैश्वानरः शिरः पातु }\edtext{\edtext{\skt{विष्णुस्तव}}{\lemma{\skt{विष्णुस्}}\Afootnote{K N; \skt{विष्णु}\ A H}}}{\lemma{\skt{तव}}\Afootnote{K N; \skt{स्तव}\ A H}}\skt{ पराक्रमं~। }\pend
\pstart\versepar \skt{}\edtext{\skt{पौरु\-ष\-म्पु\-रु\-ष\-श्रेष्ठो}}{\lemma{\skt{पौरुषम्}}\Afootnote{A K N; \skt{पौरुशं}\ H}}\skt{ ब्रह्मात्मानं }\edtext{\skt{भ्रुवौ}}{\lemma{\skt{भ्रुवौ}}\Afootnote{K N; \emph{om.}\ A; \skt{भुवौ}\ H}}\skt{ }\edtext{\skt{ध्रुवः}}{\lemma{\skt{ध्रुवः}}\Afootnote{K N H; \skt{ध्रुवो भ्रुवौ}\ A}}\skt{~।}\pend

\pstart\versepar \ledsidenote{\passref{1.5.30}}\skt{एता देहे विशेषेण तव नित्या हि देवताः~। }\pend
\pstart\versepar \skt{}\edtext{\skt{एतास्ताः}}{\lemma{\skt{एतास्ताः}}\Afootnote{K N; \skt{एतास्त्वां}\ A; \skt{एतास्ताः}\ H}}\skt{ सततं पान्तु }\edtext{\skt{दिशन्तु च}}{\lemma{\skt{दिशन्तु च}}\Afootnote{K N; \skt{दीर्घ\appbrk{}मा\appbrk{}यु\appbrk{}र\appbrk{}वा\appbrk{}प्नुहि}\ A; \skt{… दीशन्तु च}\ H}}\skt{ }\edtext{\skt{निरामयं}}{\lemma{\skt{निरामयं}}\Afootnote{K N H; \emph{om.}\ A}}\skt{~।}\pend

\pstart\versepar \ledsidenote{\passref{1.5.33}}\edtext{\skt{एतै\-र्वे\-दा\-त्म\-कै\-र्मन्त्रैः}}{\lemma{\skt{एतैर्}}\Afootnote{K N H; \skt{इति स्वाहा~। एतैर्}\ A}}\skt{~। }\edtext{\skt{कृत्य\-व्या\-धि\-वि\-ना\-शनैः}}{\lemma{\skt{कृत्य\appbrk{}व्या\appbrk{}धि\appbrk{}वि\appbrk{}ना\appbrk{}शनैः}}\Afootnote{K N; \skt{कृत्या\appbrk{}व्या\appbrk{}धि\appbrk{}वि\appbrk{}ना\appbrk{}शनैः}\ A H}}\skt{~। }\pend
\pstart\versepar \skt{मयैवं }\edtext{\skt{कृत\-र\-क्ष\-स्त्व\-न्दी\-र्घ\-मा\-यु\-र\-वा\-प्नुहि}}{\lemma{\skt{कृत\appbrk{}र\appbrk{}क्ष\appbrk{}स्त्व\appbrk{}न्दी\appbrk{}र्घ\appbrk{}मा\appbrk{}यु\appbrk{}र\appbrk{}वा\appbrk{}प्नुहि}}\Afootnote{A H; \skt{कृत\appbrk{}र\appbrk{}क्ष\appbrk{}स्त्व\appbrk{}न्दी\appbrk{}र्घ\appbrk{}मा\appbrk{}यु\appbrk{}र\appbrk{}वा\appbrk{}प्नुहि}\ K N}}\skt{~॥}\pend

\pstart \ledsidenote{\passref{1.5.34}}\skt{ततः }\edtext{\edtext{\edtext{\skt{कृत\-र\-क्ष\-क\-र्म\-मा\-तु\-र\-म\-गा\-र\-म्प्र\-वे\-श्या\-चा\-रि\-क\-मु\-प\-दि\-शेत्}}{\lemma{\skt{कृत\appbrk{}र\appbrk{}क्ष\appbrk{}क\appbrk{}र्म\appbrk{}मा\appbrk{}तु\appbrk{}र\appbrk{}म\appbrk{}गा\appbrk{}रम्}}\Afootnote{\skt{कृत\appbrk{}र\appbrk{}क्ष\appbrk{}मा\appbrk{}तु\appbrk{}र\appbrk{}मा\appbrk{}गारं}\ A; \skt{कृत\appbrk{}र\appbrk{}क्ष\appbrk{}क\appbrk{}र्म\appbrk{}मा\appbrk{}तु\appbrk{}र\appbrk{}म\appbrk{}गा\appbrk{}रम्}\ K; \skt{कृत\appbrk{}र\appbrk{}क्षा\appbrk{}क\appbrk{}र्म\appbrk{}मा\appbrk{}तु\appbrk{}र\appbrk{}म\appbrk{}गा\appbrk{}रम्}\ N; \skt{कृत\appbrk{}र\appbrk{}क्षा\appbrk{}क\appbrk{}र्म्म\appbrk{}मा\appbrk{}तु\appbrk{}र\appbrk{}म\appbrk{}गारं}\ H}}}{\lemma{\skt{प्रवे\appbrk{}श्या\appbrk{}चा\appbrk{}रि\appbrk{}कम्}}\Afootnote{K N; \skt{प्रवेश्य, आचारिकम्}\ A; \skt{प्रवेश्याचारिक}\ H}}}{\lemma{\skt{उपदिशेत्}}\Afootnote{K N; \skt{आदिशेत्}\ A; \skt{मुपदिशेत्}\ H}}\pend

\pstart \ledsidenote{\passref{1.5.35}}\edtext{\skt{तत\-स्तृ\-ती\-ये\-हनि}}{\lemma{\skt{तत\appbrk{}स्तृ\appbrk{}ती\appbrk{}ये\appbrk{}हनि}}\Afootnote{H; \skt{तत\appbrk{}स्तृ\appbrk{}ती\appbrk{}येऽ\appbrk{}हनि}\ A; \skt{तत\appbrk{}स्तृ\appbrk{}ती\appbrk{}येऽ\appbrk{}हनि}\ K; \skt{ततस्तृतीये हनि}\ N}}\skt{ }\edtext{\skt{विमुच्यैवमेव}}{\lemma{\skt{विमुच्यैवमेव}}\Afootnote{N H; \skt{विमुच्यैवमेव}\ A K}}\skt{ }\edtext{\skt{बध्नीयान्न}}{\lemma{\skt{बध्नीयान्न}}\Afootnote{K N H; \skt{बध्नी\appbrk{}या\appbrk{}द्व\appbrk{}स्त्र\appbrk{}प\appbrk{}ट्टेन; न}\ A}}\skt{ चैनं }\edtext{\skt{त्वर\-मा\-णो\-प\-रे\-द्यु\-र्मो\-क्ष\-येत्}}{\lemma{\skt{त्वर\appbrk{}मा\appbrk{}णो\appbrk{}प\appbrk{}रे\appbrk{}द्यु\appbrk{}र्मो\appbrk{}क्ष\appbrk{}येत्}}\Afootnote{\skt{त्वर\appbrk{}मा\appbrk{}णोऽ\appbrk{}प\appbrk{}रे\appbrk{}द्यु\appbrk{}र्मो\appbrk{}क्ष\appbrk{}येत्}\ A K N; \skt{त्वरमाणै र्परे\appbrk{}द्यु\appbrk{}र्म्मो\appbrk{}क्ष\appbrk{}येत्}\ H}}\pend

\pstart \ledsidenote{\passref{1.5.36}}\edtext{\skt{द्विती\-य\-दि\-व\-स\-मो\-क्ष\-णा\-द्वि\-ग्र\-थितो}}{\lemma{\skt{द्विती\appbrk{}य\appbrk{}दि\appbrk{}व\appbrk{}स\appbrk{}मो\appbrk{}क्ष\appbrk{}णा\appbrk{}द्वि\appbrk{}ग्र\appbrk{}थितो}}\Afootnote{K; \skt{द्विती\appbrk{}य\appbrk{}दि\appbrk{}व\appbrk{}स\appbrk{}प\appbrk{}रि\appbrk{}मो\appbrk{}क्ष\appbrk{}णा\appbrk{}द्वि\appbrk{}ग्र\appbrk{}थितो}\ A; \skt{द्विती\appbrk{}य\appbrk{}दि\appbrk{}व\appbrk{}स\appbrk{}मो\appbrk{}क्ष\appbrk{}णा\appbrk{}द्वि\appbrk{}ग्र\appbrk{}न्थितो}\ N; \skt{द्विती\appbrk{}य\appbrk{}दि\appbrk{}व\appbrk{}स\appbrk{}मो\appbrk{}क्ष\appbrk{}णा\appbrk{}द्वि\appbrk{}ग्र\appbrk{}थितो}\ H}}\skt{ व्रणः }\edtext{\skt{चिरा\-दु\-प\-सं\-रो\-ह\-त्यु\-ग्र\-रुक्त}}{\lemma{\skt{चिरा\appbrk{}दु\appbrk{}प\appbrk{}सं\appbrk{}रो\appbrk{}ह\appbrk{}त्यु\appbrk{}ग्र\appbrk{}रुक्त}}\Afootnote{\skt{चिरा\appbrk{}दु\appbrk{}प\appbrk{}सं\appbrk{}रो\appbrk{}हति, तीव्ररुजश्च}\ A; \skt{चिरा\appbrk{}दु\appbrk{}प\appbrk{}सं\appbrk{}रो\appbrk{}ह\appbrk{}त्यु\appbrk{}ग्र\appbrk{}रुक्त}\ K; \skt{चिरा\appbrk{}दु\appbrk{}प\appbrk{}सं\appbrk{}रो\appbrk{}ह\appbrk{}त्यु\appbrk{}ग्र\appbrk{}रुक्ष}\ N; \skt{चिरा\appbrk{}दु\appbrk{}प\appbrk{}सं\appbrk{}रो\appbrk{}ह\appbrk{}त्यु\appbrk{}ग्र\appbrk{}रुक्भ}\ H}}\skt{ }\edtext{\skt{भवति}}{\lemma{\skt{भवति}}\Afootnote{A K N; \skt{वति}\ H}}\skt{~॥}\pend

\pstart \ledsidenote{\passref{1.5.37}}\skt{अत }\edtext{\skt{उर्ध्वं}}{\lemma{\skt{उर्ध्वं}}\Afootnote{K; \skt{ऊर्ध्वं}\ A; \skt{उर्ध्व}\ N; \skt{ऊर्द्धन्}\ H}}\skt{ }\edtext{\skt{दोष\-का\-ल\-ब\-ला\-दी\-न\-वेक्ष्य}}{\lemma{\skt{दोष\appbrk{}का\appbrk{}ल\appbrk{}ब\appbrk{}ला\appbrk{}दी\appbrk{}न\appbrk{}वेक्ष्य}}\Afootnote{A K; \skt{देष\appbrk{}का\appbrk{}ल\appbrk{}ब\appbrk{}ला\appbrk{}दी\appbrk{}न्य\appbrk{}वेक्ष्य}\ N; \skt{देश\appbrk{}का\appbrk{}ल\appbrk{}ब\appbrk{}ला\appbrk{}दी\appbrk{}न्य\appbrk{}वेक्ष्य}\ H}}\skt{ }\edtext{\skt{कषा\-या\-ले\-प\-न\-ब\-न्धा\-हा\-रा\-द्वि\-द\-ध्यान्}}{\lemma{\skt{कषा\appbrk{}या\appbrk{}ले\appbrk{}प\appbrk{}न\appbrk{}ब\appbrk{}न्धा\appbrk{}हा\appbrk{}रा\appbrk{}द्वि\appbrk{}द\appbrk{}ध्यान्}}\Afootnote{K N; \skt{कषा\appbrk{}या\appbrk{}ले\appbrk{}प\appbrk{}न\appbrk{}ब\appbrk{}न्धा\appbrk{}हा\appbrk{}रा\appbrk{}चा\appbrk{}रा\appbrk{}न्वि\appbrk{}द\appbrk{}ध्यात्}\ A; \skt{कषा\appbrk{}या\appbrk{}ले\appbrk{}प\appbrk{}न\appbrk{}ब\appbrk{}न्धा\appbrk{}हा\appbrk{}रा\appbrk{}द्वि\appbrk{}द\appbrk{}ध्यान्}\ H}}\pend

\pstart \ledsidenote{\passref{1.5.38}}\skt{न चैनं }\edtext{\skt{त्वरमारणः}}{\lemma{\skt{त्वरमारणः}}\Afootnote{N H\textsuperscript{pc}; \skt{त्वरमाणः}\ A H\textsuperscript{ac}; \skt{त्वरमारणः}\ K}}\skt{ }\edtext{\skt{सान्तर्दोषं}}{\lemma{\skt{सान्तर्दोषं}}\Afootnote{A K; \skt{सान्तर्दोष}\ N; \skt{सान्तदोषं}\ H}}\skt{ }\edtext{\skt{रोहयेत्~। स }\edtext{\skt{ह्यल्पे\-ना\-प्य\-प\-चा\-रे\-णा\-भ्य\-न्त\-र\-मु\-त्स\-ङ्ग\-ङ्कृत्वा}}{\lemma{\skt{अल्पे\appbrk{}ना\appbrk{}प्य\appbrk{}प\appbrk{}चा\appbrk{}रे\appbrk{}णा\appbrk{}भ्य\appbrk{}न्त\appbrk{}र\appbrk{}मु\appbrk{}त्स\appbrk{}ङ्ग\appbrk{}ङ्कृत्वा}}\Afootnote{\skt{अल्पे\appbrk{}ना\appbrk{}प्य\appbrk{}प\appbrk{}चा\appbrk{}रे\appbrk{}णा\appbrk{}भ्य\appbrk{}न्त\appbrk{}र\appbrk{}मु\appbrk{}त्सङ्गं कृत्वा}\ A; \skt{अल्पे\appbrk{}ना\appbrk{}प्य\appbrk{}प\appbrk{}चा\appbrk{}रे\appbrk{}णा\appbrk{}भ्य\appbrk{}न्त\appbrk{}र\appbrk{}मु\appbrk{}त्स\appbrk{}ङ्ग\appbrk{}ङ्कृत्वा}\ K N; \skt{ह्यल्पे\appbrk{}ना\appbrk{}प्य\appbrk{}प\appbrk{}चा\appbrk{}रे\appbrk{}णा\appbrk{}भ्य\appbrk{}न्त\appbrk{}र\appbrk{}मु\appbrk{}त्स\appbrk{}ङ्ग\appbrk{}ङ्कृत्वा}\ H\textsuperscript{ac}; \skt{ह्यल्पे\appbrk{}ओ\appbrk{}ना\appbrk{}प्य\appbrk{}प\appbrk{}चा\appbrk{}रे\appbrk{}णा\appbrk{}भ्य\appbrk{}न्त\appbrk{}र\appbrk{}मु\appbrk{}त्स\appbrk{}ङ्ग\appbrk{}ङ्कृत्वा}\ H\textsuperscript{pc}}}}{\lemma{\skt{रोहयेत्~। स ह्य्}}\Afootnote{K N; \skt{रोपयेत्; स ह्य्}\ A; \skt{रोहयेत्~। स}\ H}}\skt{ भूयो }\edtext{\skt{विकरोति}}{\lemma{\skt{विकरोति}}\Afootnote{K N H; \skt{ऽपि विकरोति}\ A}}\skt{~।}\pend

\pstart\versepar \ledsidenote{\passref{1.5.39abcd}}\edtext{\skt{तस्मात्सुशुद्धं}}{\lemma{\skt{तस्मात्सुशुद्धं}}\Afootnote{K N H; \skt{भवन्ति चात्र तस्मा}\ A}}\skt{ }\edtext{\skt{बहि\-र\-भ्य\-न्त\-र\-तश्च}}{\lemma{\skt{बहि\appbrk{}र\appbrk{}भ्य\appbrk{}न्त\appbrk{}र\appbrk{}तश्च}}\Afootnote{K H; \skt{दन्तर्बहिश्चै}\ A; \skt{बहि\appbrk{}र\appbrk{}भ्य\appbrk{}न्त\appbrk{}र\appbrk{}तश्च}\ N}}\skt{ }\edtext{\skt{व्रणं}}{\lemma{\skt{व्रणं}}\Afootnote{K N H; \skt{व सुशुद्धं}\ A}}\skt{ संरोहयेत्~। }\pend
\pstart\versepar \skt{}\edtext{\edtext{\skt{रूढेऽ\-प्य\-जी\-र्ण्ण\-व्या\-या\-म\-व्य\-वा\-या\-दी\-न्वि\-व\-र्ज्ज\-येत्}}{\lemma{\skt{ऽप्य्}}\Afootnote{A; \skt{प्य्}\ K N H}}}{\lemma{\skt{अजी\appbrk{}र्ण्ण\appbrk{}व्या\appbrk{}या\appbrk{}म\appbrk{}व्य\appbrk{}वा\appbrk{}या\appbrk{}दी\appbrk{}न्वि\appbrk{}व\appbrk{}र्ज्ज\appbrk{}येत्}}\Afootnote{A H; \skt{अजी\appbrk{}र्ण्ण\appbrk{}व्या\appbrk{}या\appbrk{}म\appbrk{}व्य\appbrk{}वा\appbrk{}या\appbrk{}दी\appbrk{}न्वि\appbrk{}व\appbrk{}र्ज्ज\appbrk{}येत्}\ K; \skt{अजीर्ण्णव्यायामे व्यवा\appbrk{}या\appbrk{}दी\appbrk{}न्वि\appbrk{}व\appbrk{}र्ज्ज\appbrk{}येत्}\ N}}\skt{~॥}\edtext{}{\lemma{\skt{संरोहयेत्~। रूढे}}\Afootnote{K H\textsuperscript{pc}; \skt{रोपयेद्व्रणम्~। रूढे}\ A; \skt{स रोहयेत्~। रूढे}\ N; \skt{रोहयेत्। रूढे}\ H\textsuperscript{ac}}}\pend

\pstart\versepar \ledsidenote{\passref{1.5.40}}\edtext{\skt{भवन्ति चात्र॥}}{\lemma{\skt{भवन्ति चात्र}}\Afootnote{\skt{हेमन्ते}\ A K N; \emph{om.}\ H}}\skt{ }\pend
\pstart\versepar \skt{}\edtext{\skt{हेमन्ते}}{\lemma{\skt{हेमन्ते}}\Afootnote{H; \skt{शिशिरे}\ A; \emph{om.}\ K N}}\skt{ }\edtext{\skt{च}}{\lemma{\skt{च}}\Afootnote{K N H; \skt{चैव}\ A}}\skt{ }\edtext{\skt{वसन्ते च}}{\lemma{\skt{वसन्ते च}}\Afootnote{K N H; \skt{वसन्ते}\ A}}\skt{ }\edtext{\skt{शिशिरे}}{\lemma{\skt{शिशिरे}}\Afootnote{K N H; \emph{om.}\ A}}\skt{ चापि मोक्षयेत्~। }\pend
\pstart\versepar \skt{त्र्यहा }\edtext{\skt{द्व्यहा\-च्छ\-र\-द्ग्री\-ष्म\-व\-र्षा\-स्वपि}}{\lemma{\skt{द्व्यहा\appbrk{}च्छ\appbrk{}र\appbrk{}द्ग्री\appbrk{}ष्म\appbrk{}व\appbrk{}र्षास्व्}}\Afootnote{K H; \skt{द्द्व्यहा\appbrk{}च्छ\appbrk{}र\appbrk{}द्ग्री\appbrk{}ष्म\appbrk{}व\appbrk{}र्षास्व्}\ A; \skt{द्व्यहा\appbrk{}च्छ\appbrk{}र\appbrk{}ग्री\appbrk{}ष्म\appbrk{}व\appbrk{}र्षास्व्}\ N}}\skt{ च }\edtext{\skt{बुद्धिमां}}{\lemma{\skt{बुद्धिमां}}\Afootnote{K N; \skt{बुद्धिमान्}\ A H}}\skt{~।}\pend

\pstart\versepar \ledsidenote{\passref{1.5.41}}\edtext{\skt{अतिपातिषु}}{\lemma{\skt{अतिपातिषु}}\Afootnote{A K N; \skt{अतियातिषु}\ H}}\skt{ रोगेषु }\edtext{\skt{नेक्षेतैनं}}{\lemma{\skt{नेक्षेतैनं}}\Afootnote{K N; \skt{नेच्छेद्विधिम्}\ A; \skt{नेक्षतैनं}\ H}}\skt{ }\edtext{\skt{विधिं}}{\lemma{\skt{विधिं}}\Afootnote{K N H; \skt{इमं}\ A}}\skt{ भिषक्~। }\pend
\pstart\versepar \skt{प्रदी\-प्ता\-द्गा\-र\-व\-च्छीघ्रं }\edtext{\skt{तत्र}}{\lemma{\skt{तत्र}}\Afootnote{A K H; \skt{तत}\ N}}\skt{ कुर्यात्प्रतिक्रियां~।}\edtext{}{\lemma{\skt{भिषक्~। प्रदी\appbrk{}प्ता\appbrk{}द्गा\appbrk{}र\appbrk{}व\appbrk{}च्छीघ्रं}}\Afootnote{K; \skt{भिषक्~। प्रदी\appbrk{}प्ता\appbrk{}गा\appbrk{}र\appbrk{}व\appbrk{}च्छीघ्रं}\ A; \skt{भिषक्~॥ प्रदी\appbrk{}प्त\appbrk{}द्गा\appbrk{}रं\appbrk{}व\appbrk{}च्छीघ्रं}\ N; \skt{भिषक्~। प्रदी\appbrk{}प्ता\appbrk{}ङ्गा\appbrk{}र\appbrk{}व\appbrk{}च्छी\appbrk{}घ्रन्}\ H}}\pend

\pstart\versepar \ledsidenote{\passref{1.5.42}}\skt{या वेद\-ना\-श\-स्त्र\-नि\-पा\-त\-जाता तीव्रा }\edtext{\skt{शरीरे}}{\lemma{\skt{शरीरे}}\Afootnote{K H; \skt{शरीरं}\ A; \skt{च्छरीरे}\ N}}\skt{ }\edtext{\skt{प्रतनोति}}{\lemma{\skt{प्रतनोति}}\Afootnote{K N H; \skt{प्रदुनोति}\ A}}\skt{ }\edtext{\skt{जन्तोः}}{\lemma{\skt{जन्तोः}}\Afootnote{A K H; \skt{…}\ N}}\skt{~। }\pend
\pstart\versepar \skt{}\edtext{\skt{घृतेन सा}}{\lemma{\skt{घृतेन सा}}\Afootnote{K N H\textsuperscript{pc}; \skt{धृतेन सा}\ A; \skt{घृतेन मा}\ H\textsuperscript{ac}}}\skt{ }\edtext{\edtext{\skt{शान्तिमुपैति}}{\lemma{\skt{शान्तिम्}}\Afootnote{A K N; \skt{शान्ति}\ H}}}{\lemma{\skt{उपैति}}\Afootnote{A K N; \skt{मुपैति}\ H}}\skt{ }\edtext{\skt{नेति}}{\lemma{\skt{नेति}}\Afootnote{K; \skt{सिक्ता कोष्णेन यष्टी\appbrk{}म\appbrk{}धु\appbrk{}का\appbrk{}न्वि\appbrk{}तेना}\ A; \skt{सिक्ता कोष्णेन यष्टी\appbrk{}म\appbrk{}धु\appbrk{}का\appbrk{}न्वि\appbrk{}ते\appbrk{}नेयि~॥ ह्य५}\ N; \skt{सिक्ता कोष्णेन यष्टी\appbrk{}म\appbrk{}धु\appbrk{}का\appbrk{}न्विते नेति}\ H}}\skt{~॥}\pend

\pstart \ledsidenote{\passref{1.5.trailer}}\edtext{\skt{ह्य}}{\lemma{\skt{ह्य}}\Afootnote{K H; \skt{इति सुश्रु\appbrk{}त\appbrk{}सं\appbrk{}हि\appbrk{}तायां सूत्र\appbrk{}स्था\appbrk{}नेऽ\appbrk{}ग्रो\appbrk{}प\appbrk{}ह\appbrk{}र\appbrk{}णीयो नाम पञ्चमोऽध्यायः}\ A}}\skt{~॥}\pend

\endnumbering
\beginnumbering
\eledsection*{Adhyāya 6}
\pstart \ledsidenote{\passref{1.6.1}}\edtext{\skt{अथातो}}{\lemma{\skt{अथातो}}\Afootnote{K N H; \skt{अथात}\ A}}\skt{ }\edtext{\skt{ऋतुचर्यां}}{\lemma{\skt{ऋतुचर्यां}}\Afootnote{N H; \skt{ऋतुचर्यम्}\ A K}}\skt{ }\edtext{\skt{व्याख्यास्यामः}}{\lemma{\skt{व्याख्यास्यामः}}\Afootnote{H; \skt{अध्यायं व्याख्यास्यामः}\ A; \emph{om.}\ K N}}\skt{~॥१॥}\pend

\pstart \ledsidenote{\passref{1.6.3}}\skt{कालो हि }\edtext{\skt{भग\-वा\-न्स्व\-यं\-भु\-र\-ना\-दि\-म\-ध्य\-नि\-धनः}}{\lemma{\skt{भग\appbrk{}वा\appbrk{}न्स्व\appbrk{}यं\appbrk{}भुर्}}\Afootnote{N; \skt{नाम (भगवान्) स्वयम्भुर्}\ A; \skt{भगवां स्वयंभुर्}\ K; \skt{भग\appbrk{}वा\appbrk{}न्स्व\appbrk{}य\appbrk{}म्भूर्}\ H}}\skt{~। अत्र }\edtext{\skt{रस\-व्या\-प\-त्स\-म्प\-त्तयो}}{\lemma{\skt{रस\appbrk{}व्या\appbrk{}प\appbrk{}त्स\appbrk{}म्प\appbrk{}त्तयो}}\Afootnote{K N H; \skt{रस\appbrk{}व्या\appbrk{}प\appbrk{}त्स\appbrk{}म्पत्ती}\ A}}\skt{ जीवितमरणे च }\edtext{\skt{मनु\-ष्या\-णा\-मा\-यत्ते~। स}}{\lemma{\skt{आयत्ते~। स}}\Afootnote{A; \skt{आयत्ते स्स}\ K; \skt{आयन्ते स्म}\ N; \skt{आयन्ते स्स}\ H}}\skt{ }\edtext{\skt{सूक्ष्मामपि}}{\lemma{\skt{सूक्ष्माम्}}\Afootnote{A K H; \skt{सूक्ष्मम्}\ N}}\skt{ कलान्न लीयत इति कालः~। }\edtext{\skt{सङ्कलयति}}{\lemma{\skt{सङ्कलयति}}\Afootnote{A K; \skt{सङ्कालयति}\ N H}}\skt{ }\edtext{\skt{कलयति वा}}{\lemma{\skt{कलयति वा}}\Afootnote{K N H; \skt{कालयति वा}\ A}}\skt{ भूतानीति कालः~॥२॥}\pend

\pstart \ledsidenote{\passref{1.6.4}}\skt{तस्य सम्वत्सरात्मनो भगवानादित्यो }\edtext{\skt{गतिविशेषेण}}{\lemma{\skt{गतिविशेषेण}}\Afootnote{K N H; \skt{गतिविशेषेणा}\ A}}\skt{ }\edtext{\skt{निमे\-ष\-का\-ष्ठा\-क\-ला\-मु\-हू\-र्त्ता\-हो\-रा\-त्र\-प\-क्ष\-मा\-स\-र्त्व\-य\-न\-स\-म्व\-त्स\-र\-यु\-ग\-प्र\-वि\-भागं}}{\lemma{\skt{निमे\appbrk{}ष\appbrk{}का\appbrk{}ष्ठा\appbrk{}क\appbrk{}ला\appbrk{}मु\appbrk{}हू\appbrk{}र्त्ता\appbrk{}हो\appbrk{}रा\appbrk{}त्र\appbrk{}प\appbrk{}क्ष\appbrk{}मा\appbrk{}स\appbrk{}र्त्व\appbrk{}य\appbrk{}न\appbrk{}स\appbrk{}म्व\appbrk{}त्स\appbrk{}र\appbrk{}यु\appbrk{}ग\appbrk{}प्र\appbrk{}वि\appbrk{}भागं}}\Afootnote{K; \skt{क्षिनि\appbrk{}मे\appbrk{}ष\appbrk{}का\appbrk{}ष्ठा\appbrk{}क\appbrk{}ला\appbrk{}मु\appbrk{}हू\appbrk{}र्ता\appbrk{}हो\appbrk{}रा\appbrk{}त्र\appbrk{}प\appbrk{}क्ष\appbrk{}मा\appbrk{}स\appbrk{}र्त्व\appbrk{}य\appbrk{}न\appbrk{}सं\appbrk{}व\appbrk{}त्स\appbrk{}र\appbrk{}यु\appbrk{}ग\appbrk{}प्र\appbrk{}वि\appbrk{}भागं}\ A; \skt{निमे\appbrk{}ष\appbrk{}का\appbrk{}ष्ठा\appbrk{}क\appbrk{}ला\appbrk{}मु\appbrk{}हू\appbrk{}र्ता\appbrk{}हो\appbrk{}रा\appbrk{}त्र\appbrk{}प\appbrk{}क्ष\appbrk{}मा\appbrk{}स\appbrk{}र्त्व\appbrk{}य\appbrk{}न\appbrk{}सं\appbrk{}व\appbrk{}र\appbrk{}यु\appbrk{}ग\appbrk{}प्र\appbrk{}वि\appbrk{}भागं}\ N; \skt{निमे\appbrk{}ष\appbrk{}का\appbrk{}ष्ठा\appbrk{}क\appbrk{}ला\appbrk{}मु\appbrk{}हू\appbrk{}र्त्ता\appbrk{}हो\appbrk{}रा\appbrk{}त्र\appbrk{}प\appbrk{}क्ष\appbrk{}मा\appbrk{}स\appbrk{}र्त्त\appbrk{}य\appbrk{}न\appbrk{}सं\appbrk{}व\appbrk{}त्स\appbrk{}र\appbrk{}यु\appbrk{}ग\appbrk{}प्र\appbrk{}ति\appbrk{}भा\appbrk{}गङ्}\ H}}\skt{ करोति~॥३॥}\pend

\pstart \ledsidenote{\passref{1.6.5}}\skt{तत्र }\edtext{\skt{लघ्व\-क्षि\-नि\-पा\-त\-मात्रो}}{\lemma{\skt{लघ्व\appbrk{}क्षि\appbrk{}नि\appbrk{}पा\appbrk{}त\appbrk{}मात्रो}}\Afootnote{K H; \skt{लघ्व\appbrk{}क्ष\appbrk{}रो\appbrk{}च्चा\appbrk{}र\appbrk{}ण\appbrk{}मात्रो}\ A; \skt{लघ्वक्षिनिपात}\ N}}\skt{ }\edtext{\skt{निमेषः}}{\lemma{\skt{निमेषः}}\Afootnote{K H; \skt{ऽक्षिनिमेषः}\ A; \skt{मात्रो निमेषः}\ N}}\skt{~। }\edtext{\skt{पञ्च\-द\-श\-नि\-मेषा}}{\lemma{\skt{पञ्च\appbrk{}द\appbrk{}श\appbrk{}नि\appbrk{}मेषा}}\Afootnote{K; \skt{पञ्च\appbrk{}द\appbrk{}शाऽ\appbrk{}क्षि\appbrk{}नि\appbrk{}मेषाः}\ A; \skt{पञ्च\appbrk{}द\appbrk{}श\appbrk{}नि\appbrk{}मे\appbrk{}ष\appbrk{}मात्रा}\ N; \skt{पञ्च\appbrk{}द\appbrk{}श\appbrk{}नि\appbrk{}मेषाः}\ H}}\skt{ काष्ठा~। }\edtext{\skt{त्रिङ्शत्काष्ठा}}{\lemma{\skt{त्रिङ्शत्काष्ठा}}\Afootnote{K N H; \skt{त्रिंशत्काष्ठाः}\ A}}\skt{ कला~। }\edtext{\skt{विङ्शतिकलो}}{\lemma{\skt{विङ्शतिकलो}}\Afootnote{A K; \skt{विंशतिकला}\ N H}}\skt{ मुहूर्त्तः }\edtext{\skt{कलायाः}}{\lemma{\skt{कलायाः}}\Afootnote{K N H; \skt{कला}\ A}}\skt{ दशभागश्च~। }\edtext{\edtext{\skt{त्रिंश\-न्मु\-हू\-र्त्त\-म\-हो\-रा\-त्रम्}}{\lemma{\skt{त्रिंशन्मुहूर्त्तम्}}\Afootnote{K; \skt{त्रिंशन्मुहूर्त}\ A; \skt{त्रिंषन्मुहूर्त्तम्}\ N; \skt{त्रिंशत्मुहूर्त्तम्}\ H}}}{\lemma{\skt{अहोरात्रम्}}\Afootnote{K N H; \skt{महोरात्रं}\ A}}\skt{~। }\edtext{\skt{पञ्च\-द\-शा\-हो\-रात्रः}}{\lemma{\skt{पञ्च\appbrk{}द\appbrk{}शा\appbrk{}हो\appbrk{}रात्रः}}\Afootnote{K N H; \skt{पञ्च\appbrk{}द\appbrk{}शा\appbrk{}हो\appbrk{}रा\appbrk{}त्राणि}\ A}}\skt{ }\edtext{\skt{पक्षः~। स}}{\lemma{\skt{पक्षः~। स}}\Afootnote{K N H; \skt{पक्षः, स च}\ A}}\skt{ द्विविधः }\edtext{\skt{कृष्णः}}{\lemma{\skt{कृष्णः}}\Afootnote{K N H; \skt{शुक्लः}\ A}}\skt{ }\edtext{\skt{शुक्लश्च~। तौ}}{\lemma{\skt{शुक्लश्च~। तौ}}\Afootnote{K N H; \skt{कृष्णश्च, तौ}\ A}}\skt{ मासः~॥४॥}\pend

\pstart \ledsidenote{\passref{1.6.6}}\skt{तत्र माघादयो द्वादश }\edtext{\skt{मासाः}}{\lemma{\skt{मासाः}}\Afootnote{K N H; \skt{मासाः}\ A}}\skt{ }\edtext{\skt{सम्वत्सरः}}{\lemma{\skt{सम्वत्सरः}}\Afootnote{K N H; \emph{om.}\ A}}\skt{~। }\edtext{\skt{द्विमा\-सि\-क\-मृ\-तु\-ङ्कृत्वा}}{\lemma{\skt{द्विमासिकम्}}\Afootnote{A K; \skt{द्विमासिकाम्}\ N; \skt{Sप्रि\appbrk{}न्ग्द्वि\appbrk{}मा\appbrk{}सि\appbrk{}कम्}\ H}}\skt{ षडृतवो }\edtext{\skt{भवन्ति~। ते च}}{\lemma{\skt{भवन्ति~। ते च}}\Afootnote{K N H; \skt{भवन्ति; ते}\ A}}\skt{ शिशि\-र\-व\-स\-न्त\-ग्री\-ष्म\-व\-र्षा\-श\-र\-द्धे\-मन्ताः~। }\edtext{\skt{तेषा\-न्त\-प\-स्त\-पस्यौ}}{\lemma{\skt{तेषान्}}\Afootnote{A K; \skt{तेषन्}\ N H}}\skt{ शिशिरः~। मधुमाधवौ वसन्तः~। शुचिशुक्रौ ग्रीष्मः नभोनभस्यौ वर्षा~। }\edtext{\skt{इषोर्जौ}}{\lemma{\skt{इषोर्जौ}}\Afootnote{A K H; \skt{इषोजौ}\ N}}\skt{ शरत्~। सहस्सहस्यौ }\edtext{\skt{हेमन्त}}{\lemma{\skt{हेमन्त}}\Afootnote{A K N H\textsuperscript{pc}; \skt{हेमन्त ः}\ H\textsuperscript{ac}}}\skt{ }\edtext{\skt{इति}}{\lemma{\skt{इति}}\Afootnote{A K N H\textsuperscript{pc}; \skt{ति}\ H\textsuperscript{ac}}}\skt{~॥५॥}\pend

\pstart \ledsidenote{\passref{1.6.7}}\skt{त एते }\edtext{\skt{शीतो\-ष्ण\-व\-र्ष\-वा\-त\-ल\-क्षणाः}}{\lemma{\skt{शीतो\appbrk{}ष्ण\appbrk{}व\appbrk{}र्ष\appbrk{}वा\appbrk{}त\appbrk{}ल\appbrk{}क्षणाः}}\Afootnote{K\textsuperscript{pc} H; \skt{शीतो\appbrk{}ष्ण\appbrk{}व\appbrk{}र्ष\appbrk{}ल\appbrk{}क्ष\appbrk{}णाश्}\ A K\textsuperscript{ac}; \skt{शीतो\appbrk{}ष्ण\appbrk{}व\appbrk{}र्षा\appbrk{}वा\appbrk{}त\appbrk{}ल\appbrk{}क्षणाः}\ N}}\skt{~। चन्द्रादित्ययोः }\edtext{\skt{काल\-वि\-भा\-ग\-क\-र\-त्वा\-द\-यने द्वे}}{\lemma{\skt{काल\appbrk{}वि\appbrk{}भा\appbrk{}ग\appbrk{}क\appbrk{}र\appbrk{}त्वा\appbrk{}द\appbrk{}यने द्वे}}\Afootnote{A K H; \skt{काल\appbrk{}वि\appbrk{}भा\appbrk{}ग\appbrk{}क\appbrk{}र\appbrk{}त्वा\appbrk{}द्व\appbrk{}यने द्वे}\ N}}\skt{ भवतः~। दक्षि\-ण\-मु\-त्त\-रञ्च~। तयो\-र्द\-क्षि\-ण\-म्व\-र्षा\-श\-र\-द्धे\-मन्ताः~। तेषु }\edtext{\skt{भग\-वा\-ना\-प्या\-यते}}{\lemma{\skt{भग\appbrk{}वा\appbrk{}ना\appbrk{}प्या\appbrk{}यते}}\Afootnote{A K; \skt{भग\appbrk{}वा\appbrk{}न्ना\appbrk{}प्या\appbrk{}य्यते}\ N; \skt{भग\appbrk{}वा\appbrk{}न्ना\appbrk{}प्या\appbrk{}य्यते}\ H}}\skt{ }\edtext{\skt{सोमः}}{\lemma{\skt{सोमः}}\Afootnote{A H; \skt{सोमः}\ K N}}\skt{~। }\edtext{\skt{अम्ल\-ल\-व\-ण\-म\-धु\-राश्च}}{\lemma{\skt{अम्ल\appbrk{}ल\appbrk{}व\appbrk{}ण\appbrk{}म\appbrk{}धु\appbrk{}राश्च}}\Afootnote{A H; \skt{आम्ल्ब\appbrk{}ल\appbrk{}व\appbrk{}ण\appbrk{}म\appbrk{}धु\appbrk{}राश्च}\ K; \skt{आम्ब्ल\appbrk{}ल\appbrk{}व\appbrk{}ण\appbrk{}म\appbrk{}धु\appbrk{}राश्च}\ N}}\skt{ रसा बलवन्तो भवन्ति~। }\edtext{\skt{उत्तरोत्तरञ्च}}{\lemma{\skt{उत्तरोत्तरञ्च}}\Afootnote{A K H; \skt{उत्तरोत्तरश्च}\ N}}\skt{ सर्वप्राणिनां बल\-म\-भि\-व\-र्द्धते~। }\edtext{\skt{उत्तरं}}{\lemma{\skt{उत्तरं}}\Afootnote{K N H; \skt{उत्तरं च}\ A}}\skt{ शिशि\-र\-व\-स\-न्त\-ग्रीष्माः~। तेषु }\edtext{\skt{भग\-वा\-ना\-प्या\-य\-तेऽर्कः}}{\lemma{\skt{भग\appbrk{}वा\appbrk{}ना\appbrk{}प्या\appbrk{}यते}}\Afootnote{A K; \skt{भग\appbrk{}वा\appbrk{}न्ना\appbrk{}प्या\appbrk{}य्यते}\ N H}}\skt{~। }\edtext{\skt{कटु\-ति\-क्त\-क\-षा\-याश्च}}{\lemma{\skt{कटु\appbrk{}ति\appbrk{}क्त\appbrk{}क\appbrk{}षा\appbrk{}याश्च}}\Afootnote{K N\textsuperscript{pc} H; \skt{तिक्त\appbrk{}क\appbrk{}षा\appbrk{}य\appbrk{}क\appbrk{}टु\appbrk{}काश्च}\ A; \skt{कटु\appbrk{}ति\appbrk{}क्त\appbrk{}का\appbrk{}षा\appbrk{}याश्च}\ N\textsuperscript{ac}}}\skt{ रसा }\edtext{\skt{बलवत्तरा}}{\lemma{\skt{बलवत्तरा}}\Afootnote{K\textsuperscript{pc} H; \skt{बलवन्तो}\ A N; \skt{बलत्तरा}\ K\textsuperscript{ac}}}\skt{ भवन्ति~। }\edtext{\edtext{\skt{उत्तरोत्तरञ्च}}{\lemma{\skt{उत्तरोत्तरञ्}}\Afootnote{A K\textsuperscript{pc} H; \skt{उत्तरोत्त}\ K\textsuperscript{ac}; \skt{उत्तरोत्तरश्}\ N}}}{\lemma{\skt{च}}\Afootnote{K\textsuperscript{pc} N H; \skt{च सर्व}\ A; \skt{रआञ्च}\ K\textsuperscript{ac}}}\skt{ }\edtext{\skt{प्राणि\-ना\-म्ब\-ल\-म्प\-रि\-ही\-यते}}{\lemma{\skt{परिहीयते}}\Afootnote{K N H; \skt{उपहीयते}\ A}}\skt{~॥६॥}\pend

\pstart\versepar \ledsidenote{\passref{1.6.8}}\edtext{\skt{भवति}}{\lemma{\skt{भवति}}\Afootnote{A N H; \emph{om.}\ K}}\skt{ }\edtext{\skt{चात्र}}{\lemma{\skt{चात्र}}\Afootnote{A N H; \emph{om.}\ K}}\skt{~॥ }\pend
\pstart\versepar \skt{}\edtext{\skt{सोमः}}{\lemma{\skt{सोमः}}\Afootnote{K N\textsuperscript{pc} H; \skt{शीतांशुः}\ A; \skt{सोम}\ N\textsuperscript{ac}}}\skt{ }\edtext{\skt{क्लेदयते}}{\lemma{\skt{क्लेदयते}}\Afootnote{K N H; \skt{क्लेदयत्यु}\ A}}\skt{ }\edtext{\skt{भूमिं}}{\lemma{\skt{भूमिं}}\Afootnote{K; \skt{र्वीं वि}\ A; \skt{भूमि}\ N; \skt{भूमीं}\ H}}\skt{ }\edtext{\skt{सूर्यः}}{\lemma{\skt{सूर्यः}}\Afootnote{K N H; \skt{वस्वा}\ A}}\skt{ }\edtext{\skt{शोषयते}}{\lemma{\skt{शोषयते}}\Afootnote{K N H; \skt{न्शोषयत्य}\ A}}\skt{ पुनः~। }\pend
\pstart\versepar \skt{तावुभावपि संश्रित्य वायुः पालयति }\edtext{\skt{प्रजाः}}{\lemma{\skt{प्रजाः}}\Afootnote{A K H; \skt{पजः}\ N}}\skt{~॥७॥}\edtext{}{\lemma{\skt{पुनः~। ताव्}}\Afootnote{K N H; \skt{पि~। ताव्}\ A}}\pend

\pstart \ledsidenote{\passref{1.6.9}}\skt{अथ }\edtext{\skt{खल्वयने}}{\lemma{\skt{खल्वयने}}\Afootnote{K N H; \skt{खल्वयने द्वे}\ A}}\skt{ }\edtext{\skt{युग\-प\-त्सं\-व\-त्सरो}}{\lemma{\skt{युग\appbrk{}प\appbrk{}त्सं\appbrk{}व\appbrk{}त्सरो}}\Afootnote{A H; \skt{युग\appbrk{}प\appbrk{}त्स\appbrk{}म्व\appbrk{}सरो}\ K\textsuperscript{ac}; \skt{युग\appbrk{}प\appbrk{}त्स\appbrk{}म्व\appbrk{}त्स\appbrk{}सरो}\ K\textsuperscript{pc}; \skt{युगपत्सम्वरो}\ N}}\skt{ }\edtext{\skt{भवति~। ते द्वे}}{\lemma{\skt{भवति~। ते द्वे}}\Afootnote{K\textsuperscript{pc} N H; \skt{भवति~। ते}\ A K\textsuperscript{ac}}}\skt{ }\edtext{\skt{अयने}}{\lemma{\skt{अयने}}\Afootnote{K\textsuperscript{pc} N H; \emph{om.}\ A K\textsuperscript{ac}}}\skt{ }\edtext{\skt{वर्षः}}{\lemma{\skt{वर्षः}}\Afootnote{K\textsuperscript{pc} H; \emph{om.}\ A K\textsuperscript{ac}; \skt{वर्ष}\ N}}\skt{ }\edtext{\skt{संवत्सरः}}{\lemma{\skt{संवत्सरः}}\Afootnote{K\textsuperscript{pc} H; \emph{om.}\ A K\textsuperscript{ac}; \skt{सम्वरः}\ N}}\skt{ }\edtext{\skt{परिवत्सरः}}{\lemma{\skt{परिवत्सरः}}\Afootnote{K\textsuperscript{pc} N H; \emph{om.}\ A K\textsuperscript{ac}}}\skt{ }\edtext{\skt{इडावत्सरः}}{\lemma{\skt{इडावत्सरः}}\Afootnote{K\textsuperscript{pc}; \emph{om.}\ A K\textsuperscript{ac}; \skt{इडावत्सरः इडवत्सरः}\ N; \skt{इडावत्सरः~। इद्वत्सरः}\ H}}\skt{ }\edtext{\skt{वत्सर}}{\lemma{\skt{वत्सर}}\Afootnote{K\textsuperscript{pc} N H; \emph{om.}\ A K\textsuperscript{ac}}}\skt{ }\edtext{\edtext{\skt{इत्येवं}}{\lemma{\skt{इत्य्}}\Afootnote{K\textsuperscript{pc} N H; \emph{om.}\ A K\textsuperscript{ac}}}}{\lemma{\skt{एवं}}\Afootnote{K\textsuperscript{pc} N H; \emph{om.}\ A K\textsuperscript{ac}}}\skt{ }\edtext{\skt{पञ्च}}{\lemma{\skt{पञ्च}}\Afootnote{K\textsuperscript{pc} N H; \emph{om.}\ A K\textsuperscript{ac}}}\skt{ }\edtext{\skt{पञ्च}}{\lemma{\skt{पञ्च}}\Afootnote{K\textsuperscript{pc} H; \emph{om.}\ A K\textsuperscript{ac}; \skt{वञ्च}\ N}}\skt{ }\edtext{\skt{वर्षाणि~॥ ते}}{\lemma{\skt{वर्षाणि~॥ ते}}\Afootnote{N H; \skt{तु}\ A; \emph{om.}\ K\textsuperscript{ac}; \skt{वर्षानि~॥ ते}\ K\textsuperscript{pc}}}\skt{ पञ्च युगमिति }\edtext{\skt{संज्ञां}}{\lemma{\skt{संज्ञां}}\Afootnote{A N; \skt{संज्ञा}\ K H}}\skt{ लभन्ते~। स एष }\edtext{\skt{निमे\-षा\-दि\-र्यु\-ग\-प\-र्यन्तः}}{\lemma{\skt{निमेषादिर्}}\Afootnote{A K H; \skt{निमेषादि}\ N}}\skt{ काल\-श्च\-क्र\-व\-त्प\-रि\-व\-र्त्त\-मानः }\edtext{\skt{कालचक्र}}{\lemma{\skt{कालचक्र}}\Afootnote{K N H; \skt{कालचक्रम्}\ A}}\skt{ }\edtext{\skt{इत्युच्यते}}{\lemma{\skt{उच्यते}}\Afootnote{K N H; \skt{उच्यत इत्येके}\ A}}\skt{~॥८॥}\pend

\pstart \ledsidenote{\passref{1.6.9.add}}\skt{एवं दक्षिणायने }\edtext{\skt{रात्रिर्व्याख्याता}}{\lemma{\skt{रात्रिर्}}\Afootnote{K H; \skt{रात्रि}\ N}}\skt{~॥९॥}\pend

\pstart \ledsidenote{\passref{1.6.10}}\skt{इह तु वर्षा\-श\-र\-द्धे\-म\-न्त\-व\-स\-न्त\-ग्री\-ष्म\-प्रा\-वृषः षडृतवो भवन्ति }\edtext{\skt{दोषो\-प\-च\-य\-प्र\-को\-प\-प्र\-श\-म\-नि\-मि\-त्तम्~। ते तु}}{\lemma{\skt{दोषो\appbrk{}प\appbrk{}च\appbrk{}य\appbrk{}प्र\appbrk{}को\appbrk{}प\appbrk{}प्र\appbrk{}श\appbrk{}म\appbrk{}नि\appbrk{}मि\appbrk{}त्तम्~। ते तु}}\Afootnote{K N H; \skt{दोषो\appbrk{}प\appbrk{}च\appbrk{}य\appbrk{}प्र\appbrk{}को\appbrk{}पो\appbrk{}प\appbrk{}श\appbrk{}म\appbrk{}नि\appbrk{}मित्तं; ते तु}\ A}}\skt{ }\edtext{\edtext{\skt{भाद्र\-प\-दा\-द्यै\-र्द्वि\-मा\-सि\-के\-नैवं}}{\lemma{\skt{भाद्रपदाद्यैर्}}\Afootnote{H; \skt{भाद्रपदाद्येन}\ A; \skt{भद्रपदाद्यैर्}\ K; \skt{भद्रपदाद्यै}\ N}}}{\lemma{\skt{द्विमासिकेनैवं}}\Afootnote{K N H; \skt{द्विमासिकेन}\ A}}\skt{ व्याख्याताः~। तद्यथा~। }\edtext{\skt{भद्र\-प\-दा\-श्व\-युजौ}}{\lemma{\skt{भद्र\appbrk{}प\appbrk{}दा\appbrk{}श्व\appbrk{}युजौ}}\Afootnote{K H; \skt{भाद्र\appbrk{}प\appbrk{}दा\appbrk{}श्व\appbrk{}युजौ}\ A; \skt{भद्रपदास्च युजौ}\ N}}\skt{ वर्षाः~। कार्त्ति\-क\-मा\-र्ग\-शीर्षौ शरत्~। पौषमाघौ हेमन्तः~। फाल्गुनचैत्रौ }\edtext{\skt{वसन्तः}}{\lemma{\skt{वसन्तः}}\Afootnote{A K H; \emph{om.}\ N}}\skt{~। }\edtext{\skt{वैशाखज्येष्ठौ}}{\lemma{\skt{वैशाखज्येष्ठौ}}\Afootnote{A K H; \skt{वैशाखजेष्ठ}\ N}}\skt{ }\edtext{\skt{ग्रीष्मः}}{\lemma{\skt{ग्रीष्मः}}\Afootnote{A K H; \skt{व्…}\ N}}\skt{~। }\edtext{\skt{आषाढश्रावणौ}}{\lemma{\skt{आषाढश्रावणौ}}\Afootnote{A K H; \emph{om.}\ N}}\skt{ }\edtext{\skt{प्रावृडिति}}{\lemma{\skt{प्रावृडिति}}\Afootnote{A K H; \emph{om.}\ N}}\skt{~॥१०॥}\pend

\pstart \ledsidenote{\passref{1.6.11}}\edtext{\skt{तत्र}}{\lemma{\skt{तत्र}}\Afootnote{A K H; \emph{om.}\ N}}\skt{ }\edtext{\edtext{\edtext{\skt{वर्षा\-स्वौ\-ष\-ध्य\-स्त\-रु\-ण्योऽ\-ल्प\-वीर्या}}{\lemma{\skt{वर्षास्व्}}\Afootnote{A K H; \emph{om.}\ N}}}{\lemma{\skt{औषध्यस्}}\Afootnote{K H; \skt{ओषधयस्}\ A; \skt{…औषध्यस्}\ N}}}{\lemma{\skt{ऽल्पवीर्या}}\Afootnote{A H; \skt{ल्पवीर्या}\ K N}}\skt{ }\edtext{\edtext{\skt{आहा\-र\-त्व\-मु\-प\-गता}}{\lemma{\skt{आहारत्वम्}}\Afootnote{K N H; \emph{om.}\ A}}}{\lemma{\skt{उपगता}}\Afootnote{K N H; \emph{om.}\ A}}\skt{ }\edtext{\skt{विदह्यन्ते}}{\lemma{\skt{विदह्यन्ते}}\Afootnote{K N H; \emph{om.}\ A}}\skt{~। आपश्चाप्रशान्ताः }\edtext{\skt{क्षितिमलप्रायाः~। }\edtext{\skt{तास्तू\-प\-यु\-ज्य\-माना}}{\lemma{\skt{तूपयुज्यमाना}}\Afootnote{K N H; \skt{उपयुज्यमाना}\ A}}}{\lemma{\skt{क्षितिमलप्रायाः~। तास्}}\Afootnote{K N H; \skt{क्षितिमलप्रायाः, ता}\ A}}\skt{ नभसि मेघावतते जलप्रक्लिन्नायां भूमौ }\edtext{\skt{क्लिन्नदेहानां}}{\lemma{\skt{क्लिन्नदेहानां}}\Afootnote{A K N; \skt{…… क्लिन्नदेहानान्}\ H}}\skt{ }\edtext{\skt{देहिनां}}{\lemma{\skt{देहिनां}}\Afootnote{K N H; \skt{प्राणिनां}\ A}}\skt{ }\edtext{\skt{शीत\-वा\-त\-व\-र्ष\-वि\-ष्ट\-म्भि\-ता\-ग्नीनां}}{\lemma{\skt{शीत\appbrk{}वा\appbrk{}त\appbrk{}व\appbrk{}र्ष\appbrk{}वि\appbrk{}ष्ट\appbrk{}म्भि\appbrk{}ता\appbrk{}ग्नीनां}}\Afootnote{K H; \skt{शीत\appbrk{}वा\appbrk{}त\appbrk{}वि\appbrk{}ष्ट\appbrk{}म्भि\appbrk{}ता\appbrk{}ग्नीनां}\ A; \skt{शीत\appbrk{}वा\appbrk{}त\appbrk{}व\appbrk{}र्ष\appbrk{}वि\appbrk{}स्तं\appbrk{}भि\appbrk{}ता\appbrk{}ग्नी।नां}\ N}}\skt{ विदह्यन्ते~। }\edtext{\skt{विदा\-हा\-त्पि\-त्त\-स\-ञ्च\-य\-मा\-पा\-द\-यन्ति}}{\lemma{\skt{विदा\appbrk{}हा\appbrk{}त्पि\appbrk{}त्त\appbrk{}स\appbrk{}ञ्च\appbrk{}यम्}}\Afootnote{A; \skt{सवि\appbrk{}दा\appbrk{}हा\appbrk{}त्पि\appbrk{}त्त\appbrk{}स\appbrk{}ञ्च\appbrk{}यम्}\ K H N}}\skt{~। स सञ्चयः }\edtext{\skt{शरदि}}{\lemma{\skt{शरदि}}\Afootnote{A K N H\textsuperscript{pc}; \skt{शरैदि}\ H\textsuperscript{ac}}}\skt{ }\edtext{\skt{प्रविरलमेघे}}{\lemma{\skt{प्रविरलमेघे}}\Afootnote{A K N; \skt{प्रविरलमेण}\ H}}\skt{ }\edtext{\edtext{\skt{विय\-त्यु\-प\-शु\-ष्यति}}{\lemma{\skt{वियत्य्}}\Afootnote{A K N; \skt{विलयत्य्}\ H}}}{\lemma{\skt{उपशुष्यति}}\Afootnote{A K H; \skt{उपशुश्यति}\ N}}\skt{ }\edtext{\edtext{\skt{पङ्केऽ\-र्क\-कि\-र\-ण\-प्र\-वि\-ला\-यितः}}{\lemma{\skt{पङ्के}}\Afootnote{A K N; \skt{पङ्का}\ H}}}{\lemma{\skt{ऽर्क\appbrk{}कि\appbrk{}र\appbrk{}ण\appbrk{}प्र\appbrk{}वि\appbrk{}ला\appbrk{}यितः}}\Afootnote{A N; \skt{र्ककि\appbrk{}र\appbrk{}ण\appbrk{}प्र\appbrk{}वि\appbrk{}ला\appbrk{}यितः}\ K H}}\skt{ }\edtext{\edtext{\skt{पैत्ति\-का\-न्व्या\-धी\-ञ्ज\-न\-यति}}{\lemma{\skt{पैत्तिकान्व्याधीञ्}}\Afootnote{A; \skt{पैत्तिकां व्याधीं}\ K; \skt{पित्तैकां व्याधीं}\ N; \skt{पैत्तिकान्व्याधी}\ H}}\skt{~। ता}}{\lemma{\skt{जनयति~। ता}}\Afootnote{A K N; \skt{न्जनयति~। ता}\ H}}\skt{ }\edtext{\skt{एवौषध्यः}}{\lemma{\skt{एवौषध्यः}}\Afootnote{K; \skt{एवौषधयः}\ A N H}}\skt{ काल\-प\-रि\-णा\-मा\-त्प\-रि\-ण\-त\-वीर्या बलवत्यो हेमन्ते }\edtext{\skt{भवन्ति}}{\lemma{\skt{भवन्ति}}\Afootnote{A K H; \skt{भवति}\ N}}\skt{~। आपश्च प्रशान्ताः स्निग्धा अत्यर्थं }\edtext{\skt{गुर्व्यः~। ता}}{\lemma{\skt{गुर्व्यः~। ता}}\Afootnote{K H; \skt{गुर्व्याश्च, ता}\ A; \skt{गुरुव्यस्ता}\ N}}\skt{ उपयुज्यमानाः }\edtext{\skt{मन्द\-कि\-र\-ण\-त्वा\-द्भा\-नो\-स्स\-तु\-षा\-रो\-प\-ष्ट\-म्भि\-त\-दे\-हानां}}{\lemma{\skt{सतु\appbrk{}षा\appbrk{}रो\appbrk{}प\appbrk{}ष्ट\appbrk{}म्भि\appbrk{}त\appbrk{}दे\appbrk{}हानां}}\Afootnote{N H; \skt{सतु\appbrk{}षा\appbrk{}र\appbrk{}प\appbrk{}व\appbrk{}नो\appbrk{}प\appbrk{}स्त\appbrk{}म्भि\appbrk{}त\appbrk{}दे\appbrk{}हानां}\ A; \skt{सतु\appbrk{}षा\appbrk{}रौ\appbrk{}प\appbrk{}ष्ट\appbrk{}म्भि\appbrk{}त\appbrk{}दे\appbrk{}हानां}\ K}}\skt{ देहि\-ना\-म\-वि\-दग्धाः }\edtext{\skt{स्नेहा\-द्गौ\-र\-वा\-दु\-प\-ले\-पि\-त्वाच्च}}{\lemma{\skt{स्नेहा\appbrk{}द्गौ\appbrk{}र\appbrk{}वा\appbrk{}दु\appbrk{}प\appbrk{}ले\appbrk{}पि\appbrk{}त्वाच्च}}\Afootnote{K N H; \skt{स्नेहा\appbrk{}च्छै\appbrk{}त्या\appbrk{}द्गौ\appbrk{}र\appbrk{}वा\appbrk{}दु\appbrk{}प\appbrk{}ले\appbrk{}पाच्च}\ A}}\skt{ }\edtext{\skt{श्लेष्म\-ण\-स्स\-ञ्च\-य\-मा\-पा\-द\-यन्ति}}{\lemma{\skt{श्लेष्मणस्}}\Afootnote{K N H; \skt{श्लेष्म}\ A}}\skt{~। स संचयो }\edtext{\skt{वस\-न्तेऽ\-र्क\-कि\-र\-ण\-प्र\-वि\-ला\-यितः}}{\lemma{\skt{ऽर्क\appbrk{}कि\appbrk{}र\appbrk{}ण\appbrk{}प्र\appbrk{}वि\appbrk{}ला\appbrk{}यितः}}\Afootnote{H; \skt{ऽर्क\appbrk{}र\appbrk{}श्मि\appbrk{}प्र\appbrk{}वि\appbrk{}ला\appbrk{}यित}\ A; \skt{र्ककि\appbrk{}र\appbrk{}ण\appbrk{}प्र\appbrk{}वि\appbrk{}ला\appbrk{}यितः}\ K; \skt{र्ककिरणप्र}\ N}}\skt{ }\edtext{\skt{श्लेष्मि\-का\-न्व्या\-धी\-न्ज\-न\-यति~। ता}}{\lemma{\skt{श्लेष्मि\appbrk{}का\appbrk{}न्व्या\appbrk{}धी\appbrk{}न्ज\appbrk{}न\appbrk{}यति~। ता}}\Afootnote{H; \skt{ईष\appbrk{}त्स्त\appbrk{}ब्ध\appbrk{}दे\appbrk{}हानां देहिनां श्लैष्मि\appbrk{}का\appbrk{}न्व्या\appbrk{}धी\appbrk{}ञ्ज\appbrk{}न\appbrk{}यति~। ता}\ A; \skt{श्लेष्मिकां व्याधीं जनयति~। ता}\ K; \skt{लायितः श्लेष्मिकान्व्याधीं जनयति~। ता}\ N}}\skt{ }\edtext{\skt{एवौषध्यो}}{\lemma{\skt{एवौषध्यो}}\Afootnote{K N H; \skt{एवौषधयो}\ A}}\skt{ }\edtext{\skt{ग्रीष्मे}}{\lemma{\skt{ग्रीष्मे}}\Afootnote{K H; \skt{निदाघे}\ A; \skt{ग्रीष्म}\ N}}\skt{ निःसारा रूक्षा }\edtext{\skt{अतिमात्रलघ्व्यो}}{\lemma{\skt{अतिमात्रलघ्व्यो}}\Afootnote{K N H; \skt{अतिमात्रं लघ्व्यो}\ A}}\skt{ }\edtext{\edtext{\skt{भवन्त्यापश्च}}{\lemma{\skt{भवन्त्य्}}\Afootnote{A K N; \skt{भवत्य्}\ H}}\skt{ ता}}{\lemma{\skt{आपश्च ता}}\Afootnote{A K H; \skt{आपश्च}\ N}}\skt{ उपयुज्यमानाः}\footnoteA{At this point, K adds: \emph{(tatra śaraddhemantayor madhye samam ahorātraṃ yugapac cādhimāsakau bhavataḥ viṃśatikāṣṭhāś caikanimeṣārdhaś ca dinaparivṛddhir uttarāyaṇe)}. The round brackets appear to have been added by another hand.}\skt{ }\edtext{\edtext{\skt{सूर्य\-प्र\-ता\-पो\-प\-शो\-षि\-त\-दे\-हा\-ना\-न्दे\-हिनां}}{\lemma{\skt{सूर्य\appbrk{}प्र\appbrk{}ता\appbrk{}पो\appbrk{}प\appbrk{}शो\appbrk{}षि\appbrk{}त\appbrk{}दे\appbrk{}हा\appbrk{}नान्}}\Afootnote{A H; \skt{तत्र शर\appbrk{}द्धे\appbrk{}म\appbrk{}न्त\appbrk{}यो\appbrk{}र्मध्ये सममहोरात्रं युग\appbrk{}प\appbrk{}च्चा\appbrk{}धि\appbrk{}मा\appbrk{}सकौ भवतः विंश\appbrk{}ति\appbrk{}का\appbrk{}ष्ठा\appbrk{}श्चै\appbrk{}क\appbrk{}नि\appbrk{}मे\appbrk{}षा\appbrk{}र्धश्च दिन\appbrk{}प\appbrk{}रि\appbrk{}वृ\appbrk{}द्धि\appbrk{}रु\appbrk{}त्त\appbrk{}रा\appbrk{}यणे~। सूर्य\appbrk{}प्र\appbrk{}ता\appbrk{}पो\appbrk{}प\appbrk{}शो\appbrk{}षि\appbrk{}त\appbrk{}दे\appbrk{}हा\appbrk{}नान्}\ K; \skt{सूर्य\appbrk{}प्र\appbrk{}ता\appbrk{}पो\appbrk{}य\appbrk{}शो\appbrk{}षि\appbrk{}त\appbrk{}दे\appbrk{}हानां}\ N}}}{\lemma{\skt{देहिनां}}\Afootnote{A K H; \skt{देहिना}\ N}}\skt{ }\edtext{\skt{रौक्ष्याल्लाघवाच्च}}{\lemma{\skt{रौक्ष्याल्लाघवाच्च}}\Afootnote{K H; \skt{रौक्ष्याल्लघुत्वाच्च}\ A; \skt{रोक्ष्या लाघवाच्च}\ N}}\skt{ }\edtext{\skt{वायोः}}{\lemma{\skt{वायोः}}\Afootnote{A K H; \skt{योः}\ N}}\skt{ सञ्च\-य\-मा\-पा\-द\-यन्ति~। स सञ्चयः प्रावृषि }\edtext{\skt{चात्य\-र्थ\-ञ्ज\-लो\-प\-क्लि\-न्नायां}}{\lemma{\skt{चात्यर्थञ्}}\Afootnote{A H; \skt{चात्यर्थंञ्}\ K; \skt{चात्यर्थ}\ N}}\skt{ भूमौ }\edtext{\skt{चाति\-क्लि\-न्न\-दे\-हानां}}{\lemma{\skt{चाति\appbrk{}क्लि\appbrk{}न्न\appbrk{}दे\appbrk{}हानां}}\Afootnote{K N H; \skt{क्लिन्नदेहानां}\ A}}\skt{ देहिनां शीत\-वा\-त\-व\-र्षे\-रितो }\edtext{\edtext{\skt{वाति\-का\-न्व्या\-धी\-ञ्ज\-न\-यति}}{\lemma{\skt{वातिकान्व्याधीञ्}}\Afootnote{A; \skt{वातिकां व्याधींञ्}\ K; \skt{वातिकां व्याधी}\ N; \skt{वातिकान्व्याधी}\ H}}}{\lemma{\skt{जनयति}}\Afootnote{A K; \skt{ण्जनयति}\ N; \skt{न्जनयति}\ H}}\skt{~। }\edtext{\skt{एवमेषां}}{\lemma{\skt{एषां}}\Afootnote{K N; \skt{एष}\ A; \skt{एषन्}\ H}}\skt{ दोषाणां }\edtext{\skt{सञ्च\-य\-प्र\-को\-प\-हे\-तु\-रुक्तः}}{\lemma{\skt{उक्तः}}\Afootnote{A K N; \skt{उक्ताः}\ H}}\skt{~॥११॥}\pend

\pstart \ledsidenote{\passref{1.6.12}}\skt{तत्र वर्षा\-हे\-म\-न्त\-ग्री\-ष्मेषु सञ्चि\-ता\-ना\-न्दो\-षाणां }\edtext{\skt{शर\-द्व\-स\-न्त\-प्रा\-वृट्सु च}}{\lemma{\skt{शर\appbrk{}द्व\appbrk{}स\appbrk{}न्त\appbrk{}प्रा\appbrk{}वृट्सु च}}\Afootnote{A K H; \skt{शर\appbrk{}द्धे\appbrk{}म\appbrk{}न्त\appbrk{}प्रा\appbrk{}वृट्सु च}\ N}}\skt{ }\edtext{\skt{प्रकुपितानां}}{\lemma{\skt{प्रकुपितानां}}\Afootnote{A K H; \skt{प्रकुपिताना}\ N}}\skt{ }\edtext{\skt{निर्हरणम्}}{\lemma{\skt{निर्हरणम्}}\Afootnote{K N H; \skt{निर्हरणं कर्तव्यम्}\ A}}\skt{~॥१२॥}\footnoteA{In the 1931 ed. of the vulgate, the text numbering skips from 1.61.11 to 1.6.13, omitting the number 1.6.12. In the 1938 ed. the vulgate \emph{tatra varṣāhemanta ... nirharaṇaṃ kartavyam} is numbered 1.6.12 (not 1.6.13) and the numbers are one more than the 1931 ed. from there to the end of the \emph{adhyāya}.}\pend

\pstart \ledsidenote{\passref{1.6.13}}\skt{तत्र }\edtext{\skt{पैत्तिकानां}}{\lemma{\skt{पैत्तिकानां}}\Afootnote{A K H; \skt{पैतिकानां}\ N}}\skt{ व्याधीनां }\edtext{\skt{हेमन्ते}}{\lemma{\skt{हेमन्ते}}\Afootnote{K N H; \skt{उपशमो हेमन्ते}\ A}}\skt{ }\edtext{\skt{व्युपशमः}}{\lemma{\skt{व्युपशमः}}\Afootnote{K N H; \emph{om.}\ A}}\skt{~। }\edtext{\skt{श्लैष्मि\-का\-णा\-न्नि\-दाघे}}{\lemma{\skt{श्लैष्मिकाणान्}}\Afootnote{A K N; \skt{श्लैष्मिकानां}\ H}}\skt{ वातिकानां शरदि~। }\edtext{\edtext{\skt{स्वभावतस्त्वेते}}{\lemma{\skt{स्वभावतस्}}\Afootnote{K N H; \skt{स्वभावत एव; त}\ A}}}{\lemma{\skt{त्वेते}}\Afootnote{K N H; \skt{एते}\ A}}\skt{ }\edtext{\skt{सञ्च\-य\-प्र\-को\-पो\-प\-शमा}}{\lemma{\skt{सञ्च\appbrk{}य\appbrk{}प्र\appbrk{}को\appbrk{}पो\appbrk{}प\appbrk{}शमा}}\Afootnote{A K N; \skt{सञ्च\appbrk{}य\appbrk{}प्र\appbrk{}को\appbrk{}पो\appbrk{}प\appbrk{}शम्}\ H}}\skt{ }\edtext{\skt{व्याख्याताः}}{\lemma{\skt{व्याख्याताः}}\Afootnote{A K N; \skt{आख्याताः}\ H}}\skt{~॥१३॥}\pend

\pstart \ledsidenote{\passref{1.6.14}}\skt{तत्र }\edtext{\skt{दिवसपूर्वाह्णे}}{\lemma{\skt{दिवसपूर्वाह्णे}}\Afootnote{K N H; \skt{पूर्वाह्णे}\ A}}\skt{ }\edtext{\skt{वसन्तलिङ्गं}}{\lemma{\skt{वसन्तलिङ्गं}}\Afootnote{K N H; \skt{वसन्तस्य लिङ्गं}\ A}}\skt{ मध्याह्ने ग्रीष्मस्य अपराह्णे }\edtext{\skt{प्रावृड्लिङ्गं}}{\lemma{\skt{प्रावृड्लिङ्गं}}\Afootnote{K H; \skt{प्रावृषः}\ A; \skt{प्रावृड्लिंग}\ N}}\skt{ }\edtext{\skt{प्रदोषे}}{\lemma{\skt{प्रदोषे}}\Afootnote{A K H; \skt{महौषे}\ N}}\skt{ वार्षिकं शार\-द\-म\-र्ध\-रात्रे प्रत्युषसि }\edtext{\skt{हैम\-न\-मु\-प\-ल\-क्ष\-येत्}}{\lemma{\skt{हैमनम्}}\Afootnote{K; \skt{हैमन्तम्}\ A; \skt{हेमनम्}\ N; \skt{हेमन्तम्}\ H}}\skt{~। एव\-म\-हो\-रा\-त्र\-मपि वर्षमिव }\edtext{\edtext{\skt{शीतो\-ष्ण\-व\-र्ष\-दो\-षो\-प\-च\-य\-प्र\-को\-पो\-प\-श\-मै\-र्जा\-नी\-यात्}}{\lemma{\skt{शीतो\appbrk{}ष्ण\appbrk{}व\appbrk{}र्ष\appbrk{}दो\appbrk{}षो\appbrk{}प\appbrk{}च\appbrk{}य\appbrk{}प्र\appbrk{}को\appbrk{}पो\appbrk{}प\appbrk{}श\appbrk{}मैर्}}\Afootnote{K N H; \skt{शीतो\appbrk{}ष्ण\appbrk{}व\appbrk{}र्ष\appbrk{}ल\appbrk{}क्षणं दोषो\appbrk{}प\appbrk{}च\appbrk{}य\appbrk{}प्र\appbrk{}को\appbrk{}पो\appbrk{}प\appbrk{}श\appbrk{}मैर्}\ A}}}{\lemma{\skt{जानीयात्}}\Afootnote{A K H; \skt{जनीयात्}\ N}}\skt{~॥१४॥}\pend

\pstart \ledsidenote{\passref{1.6.15}}\edtext{\skt{तत्रा\-व्या\-प\-न्ने\-ष्वृ\-तु\-ष्व\-व्या\-पन्ना}}{\lemma{\skt{तत्रा\appbrk{}व्या\appbrk{}प\appbrk{}न्ने\appbrk{}ष्वृ\appbrk{}तुष्व्}}\Afootnote{K N; \skt{तत्र, अव्यापन्नेषु ऋतुष्व्}\ A; \skt{तत्र व्यापन्नेष्वृतुष्व्}\ H}}\skt{ ओषधयो }\edtext{\skt{भवन्त्यापश्च~। }\edtext{\skt{तास्तू\-प\-यु\-ज्य\-मानाः}}{\lemma{\skt{तूपयुज्यमानाः}}\Afootnote{K N H; \skt{उपयुज्यमानाः}\ A}}}{\lemma{\skt{आपश्च~। तास्}}\Afootnote{K N H; \skt{आपश्च; ता}\ A}}\skt{ }\edtext{\skt{प्राणा\-यु\-र्ब\-ल\-वी\-र्यौ\-ज\-स्कर्यो}}{\lemma{\skt{प्राणा\appbrk{}यु\appbrk{}र्ब\appbrk{}ल\appbrk{}वी\appbrk{}र्यौ\appbrk{}ज\appbrk{}स्कर्यो}}\Afootnote{A K H; \skt{प्राणा\appbrk{}यु\appbrk{}ब\appbrk{}ल\appbrk{}वी\appbrk{}र्यौ\appbrk{}ज\appbrk{}स्कर्यो}\ N}}\skt{ भवन्ति~॥१५॥}\pend

\pstart \ledsidenote{\passref{1.6.16}}\edtext{\edtext{\skt{तासा\-म्पु\-न\-र्व्या\-प\-दोऽ\-दृ\-ष्ट\-का\-रि\-तानि}}{\lemma{\skt{तासाम्}}\Afootnote{K N H; \skt{तेषां}\ A}}}{\lemma{\skt{ऽदृ\appbrk{}ष्ट\appbrk{}का\appbrk{}रि\appbrk{}तानि}}\Afootnote{\skt{ऽदृष्टकारिताः}\ A; \skt{दृष्टकारितानि}\ K N H}}\skt{ }\edtext{\skt{शीतो\-ष्ण\-वा\-त\-व\-र्षाणि}}{\lemma{\skt{शीतो\appbrk{}ष्ण\appbrk{}वा\appbrk{}त\appbrk{}व\appbrk{}र्षाणि}}\Afootnote{A K\textsuperscript{pc} N H; \skt{शीतोष्णवर्षाणि}\ K\textsuperscript{ac}}}\skt{~। खलु }\edtext{\edtext{\skt{विप\-री\-ता\-न्यौ\-ष\-धी\-र्व्या\-पा\-द\-य\-न्त्या\-पश्च}}{\lemma{\skt{औषधीर्}}\Afootnote{K H; \skt{ओषधीर्}\ A; \skt{औषधी}\ N}}}{\lemma{\skt{आपश्}}\Afootnote{K N H; \skt{अपश्}\ A}}\skt{~॥१६॥}\pend

\pstart \ledsidenote{\passref{1.6.17}}\skt{तासा\-मु\-प\-यो\-गा\-द्वि\-वि\-ध\-रो\-ग\-प्रा\-दु\-र्भावो मरको वा }\edtext{\skt{भवति}}{\lemma{\skt{भवति}}\Afootnote{K N H; \skt{भवेदिति}\ A}}\skt{~॥१७॥}\pend

\pstart \ledsidenote{\passref{1.6.19}}\edtext{\edtext{\skt{कदा\-चि\-द\-व्या\-प\-न्ने\-ष्वपि}}{\lemma{\skt{कदा\appbrk{}चि\appbrk{}द\appbrk{}व्या\appbrk{}प\appbrk{}न्नेष्व्}}\Afootnote{A N H; \emph{om.}\ K}}}{\lemma{\skt{अपि}}\Afootnote{A N H; \emph{om.}\ K}}\skt{ }\edtext{\edtext{\skt{कृत्या\-भि\-शा\-प\-र\-क्षः\-क्रो\-धा\-द\-ध\-र्मै\-रु\-ध्व\-स्यन्ते}}{\lemma{\skt{कृत्या\appbrk{}भि\appbrk{}शा\appbrk{}प\appbrk{}र\appbrk{}क्षः\appbrk{}क्रो\appbrk{}धा\appbrk{}द\appbrk{}ध\appbrk{}र्मैर्}}\Afootnote{N H; \skt{ऋतुषु कृत्या\appbrk{}भि\appbrk{}शा\appbrk{}प\appbrk{}र\appbrk{}क्षः\appbrk{}क्रो\appbrk{}धा\appbrk{}ध\appbrk{}र्मैर्}\ A; \skt{…शा\appbrk{}प\appbrk{}र\appbrk{}क्षः\appbrk{}क्रो\appbrk{}धा\appbrk{}द\appbrk{}ध\appbrk{}र्मैर्}\ K}}}{\lemma{\skt{उध्वस्यन्ते}}\Afootnote{H; \skt{उपध्वस्यन्ते}\ A; \skt{उध्वंस्यन्ते}\ K N}}\skt{ जनपदाः~। विषौ\-ष\-धि\-पु\-ष्प\-ग\-न्धेन वा }\edtext{\skt{वायु\-नो\-प\-नी\-तेन}}{\lemma{\skt{वायु\appbrk{}नो\appbrk{}प\appbrk{}नी\appbrk{}तेन}}\Afootnote{K N H; \skt{वायु\appbrk{}नो\appbrk{}प\appbrk{}नी\appbrk{}ते\appbrk{}ना\appbrk{}क्र\appbrk{}म्यते यो देशस्तत्र दोष\appbrk{}प्र\appbrk{}कृ\appbrk{}त्य\appbrk{}वि\appbrk{}शे\appbrk{}षेण}\ A}}\skt{~। }\edtext{\skt{कास\-श्वा\-स\-प्र\-ति\-श्या\-य\-शि\-रो\-रो\-ग\-ज्व\-रै\-रु\-प\-त\-प्यन्ते}}{\lemma{\skt{कास\appbrk{}श्वा\appbrk{}स\appbrk{}प्र\appbrk{}ति\appbrk{}श्या\appbrk{}य\appbrk{}शि\appbrk{}रो\appbrk{}रो\appbrk{}ग\appbrk{}ज्व\appbrk{}रैर्}}\Afootnote{K N; \skt{कास\appbrk{}श्वा\appbrk{}स\appbrk{}व\appbrk{}म\appbrk{}थु\appbrk{}प्र\appbrk{}ति\appbrk{}श्या\appbrk{}य\appbrk{}शि\appbrk{}रो\appbrk{}रु\appbrk{}ग्ज्व\appbrk{}रैर्}\ A; \skt{कास\appbrk{}श्वा\appbrk{}स\appbrk{}प्र\appbrk{}ति\appbrk{}स्या\appbrk{}य\appbrk{}शि\appbrk{}रो\appbrk{}रो\appbrk{}ग\appbrk{}ज्व\appbrk{}रैर्}\ H}}\skt{ }\edtext{\skt{प्रजाः}}{\lemma{\skt{प्रजाः}}\Afootnote{K H; \emph{om.}\ A; \skt{प्रजः}\ N}}\skt{ ग्रह\-न\-क्ष\-त्र\-च\-रि\-तैर्वा~। }\edtext{\skt{शय\-ना\-स\-न\-या\-न\-वा\-ह\-न\-म\-णि\-र\-त्नो\-प\-क\-र\-ण\-ग\-र्हि\-त\-ल\-क्ष\-ण\-प्रा\-दु\-र्भा\-वैर्वा}}{\lemma{\skt{शय\appbrk{}ना\appbrk{}स\appbrk{}न\appbrk{}या\appbrk{}न\appbrk{}वा\appbrk{}ह\appbrk{}न\appbrk{}म\appbrk{}णि\appbrk{}र\appbrk{}त्नो\appbrk{}प\appbrk{}क\appbrk{}र\appbrk{}ण\appbrk{}ग\appbrk{}र्हि\appbrk{}त\appbrk{}ल\appbrk{}क्ष\appbrk{}ण\appbrk{}प्रा\appbrk{}दु\appbrk{}र्भा\appbrk{}वैर्}}\Afootnote{K N H\textsuperscript{pc}; \skt{गृह\appbrk{}दा\appbrk{}र\appbrk{}श\appbrk{}य\appbrk{}ना\appbrk{}स\appbrk{}न\appbrk{}या\appbrk{}न\appbrk{}वा\appbrk{}ह\appbrk{}न\appbrk{}म\appbrk{}णि\appbrk{}र\appbrk{}त्नो\appbrk{}प\appbrk{}क\appbrk{}र\appbrk{}ण\appbrk{}ग\appbrk{}र्हि\appbrk{}त\appbrk{}ल\appbrk{}क्ष\appbrk{}ण\appbrk{}नि\appbrk{}मि\appbrk{}त्त\appbrk{}प्रा\appbrk{}दु\appbrk{}र्भा\appbrk{}वैर्}\ A; \skt{शय\appbrk{}ना\appbrk{}म\appbrk{}न\appbrk{}या\appbrk{}न\appbrk{}वा\appbrk{}ह\appbrk{}न\appbrk{}म\appbrk{}णि\appbrk{}र\appbrk{}त्नो\appbrk{}प\appbrk{}क\appbrk{}र\appbrk{}ण\appbrk{}ग\appbrk{}र्हि\appbrk{}त\appbrk{}ल\appbrk{}क्ष\appbrk{}ण\appbrk{}प्रा\appbrk{}दु\appbrk{}र्भा\appbrk{}वैर्}\ H\textsuperscript{ac}}}\skt{~॥१८॥}\pend

\pstart \ledsidenote{\passref{1.6.20}}\edtext{\edtext{\skt{जातोऽत्र}}{\lemma{\skt{जातो}}\Afootnote{K N H; \skt{त}\ A}}}{\lemma{\skt{ऽत्र}}\Afootnote{K N H; \skt{त्र, स्थान}\ A}}\skt{ }\edtext{\skt{परि\-त्या\-ग\-शा\-न्ति\-प्रा\-य\-श्चि\-त्त\-ब\-लि\-म\-ङ्ग\-ल\-ज\-प\-हो\-मो\-प\-हा\-रे\-ज्या\-ञ्ज\-लि\-न\-म\-स्का\-र\-त\-पो\-नि\-य\-म\-द\-या\-दा\-न\-दी\-क्षा\-भ्यु\-प\-ग\-म\-दे\-व\-ता\-ब्रा\-ह्म\-ण\-गु\-रु\-प\-रै\-र्भ\-वि\-त\-व्य\-मेवं}}{\lemma{\skt{परि\appbrk{}त्या\appbrk{}ग\appbrk{}शा\appbrk{}न्ति\appbrk{}प्रा\appbrk{}य\appbrk{}श्चि\appbrk{}त्त\appbrk{}ब\appbrk{}लि\appbrk{}म\appbrk{}ङ्ग\appbrk{}ल\appbrk{}ज\appbrk{}प\appbrk{}हो\appbrk{}मो\appbrk{}प\appbrk{}हा\appbrk{}रे\appbrk{}ज्या\appbrk{}ञ्ज\appbrk{}लि\appbrk{}न\appbrk{}म\appbrk{}स्का\appbrk{}र\appbrk{}त\appbrk{}पो\appbrk{}नि\appbrk{}य\appbrk{}म\appbrk{}द\appbrk{}या\appbrk{}दा\appbrk{}न\appbrk{}दी\appbrk{}क्षा\appbrk{}भ्यु\appbrk{}प\appbrk{}ग\appbrk{}म\appbrk{}दे\appbrk{}व\appbrk{}ता\appbrk{}ब्रा\appbrk{}ह्म\appbrk{}ण\appbrk{}गु\appbrk{}रु\appbrk{}प\appbrk{}रैर्}}\Afootnote{K N H; \skt{परि\appbrk{}त्या\appbrk{}ग\appbrk{}शा\appbrk{}न्ति\appbrk{}क\appbrk{}र्म\appbrk{}प्रा\appbrk{}य\appbrk{}श्चि\appbrk{}त्त\appbrk{}म\appbrk{}ङ्ग\appbrk{}ल\appbrk{}ज\appbrk{}प\appbrk{}हो\appbrk{}मो\appbrk{}प\appbrk{}हा\appbrk{}रे\appbrk{}ज्या\appbrk{}ञ्ज\appbrk{}लि\appbrk{}न\appbrk{}म\appbrk{}स्का\appbrk{}र\appbrk{}त\appbrk{}पो\appbrk{}नि\appbrk{}य\appbrk{}म\appbrk{}द\appbrk{}या\appbrk{}दा\appbrk{}न\appbrk{}दी\appbrk{}क्षा\appbrk{}भ्यु\appbrk{}प\appbrk{}ग\appbrk{}म\appbrk{}दे\appbrk{}व\appbrk{}ता\appbrk{}ब्रा\appbrk{}ह्म\appbrk{}ण\appbrk{}गु\appbrk{}रु\appbrk{}प\appbrk{}रैर्}\ A}}\skt{ }\edtext{\skt{साधु}}{\lemma{\skt{साधु}}\Afootnote{A K; \skt{साधुर्}\ N H}}\skt{ भवति~॥१९॥}\pend

\pstart\versepar \ledsidenote{\passref{1.6.37}}\skt{भवति चात्र~। }\pend
\pstart\versepar \skt{}\edtext{\skt{स्वगुणैरपि}}{\lemma{\skt{अपि}}\Afootnote{K N H; \skt{अति}\ A}}\skt{ }\edtext{\skt{युक्तेषु}}{\lemma{\skt{युक्तेषु}}\Afootnote{A K H; \skt{भूते}\ N}}\skt{ विपरीतेषु वा पुनः~। }\pend
\pstart\versepar \skt{विषमेष्वपि वा }\edtext{\skt{दोषाः}}{\lemma{\skt{दोषाः}}\Afootnote{A; \skt{दोषा}\ K N H}}\skt{ कुप्यन्त्यृतुषु देहिनाम्~। }\edtext{\skt{इति}}{\lemma{\skt{इति}}\Afootnote{K N H; \emph{om.}\ A}}\skt{~॥२०॥}\edtext{}{\lemma{\skt{भवति चात्र~। स्वगुणैर्}}\Afootnote{\skt{स्वगुणैर्}\ A N K H}}\pend

\pstart \ledsidenote{\passref{1.6.trailer}}\edtext{\skt{॥ ६~॥}}{\lemma{\skt{॥ ६~॥}}\Afootnote{\skt{इति सुश्रु\appbrk{}त\appbrk{}सं\appbrk{}हि\appbrk{}तायां सूत्रस्थाने ऋतुचर्या नाम षष्ठोऽध्यायः}\ A; \skt{भ्र}\ K H; \skt{॥ {\symfont ❈}}\ N}}\pend

\endnumbering
\beginnumbering
\eledsection*{Adhyāya 7}
\pstart \ledsidenote{\passref{1.7.1}}\skt{अथातो यन्त्र\-वि\-धि\-म\-ध्यायं }\edtext{\skt{व्याख्यास्यामः}}{\lemma{\skt{व्याख्यास्यामः}}\Afootnote{A N H; \skt{व्याख्या}\ K}}\skt{~॥}\pend

\pstart \ledsidenote{\passref{1.7.2.add}}\edtext{\skt{त्वह्यं}}{\lemma{\skt{त्वह्यं}}\Afootnote{K; \skt{न्वह्मा}\ N; \skt{त्वह्या}\ H}}\skt{ }\edtext{\skt{ससन्धिश्रोतः}}{\lemma{\skt{ससन्धिश्रोतः}}\Afootnote{K; \skt{त्समन्विस्रोत्रः}\ N; \skt{त्ससन्धिश्रोतः}\ H}}\skt{ }\edtext{\edtext{\skt{स्नाय्व\-स्थि\-को\-ष्ठ\-ग\-त\-श\-ल्या\-द्ध\-र\-णा\-र्थ\-मु\-प\-दि\-श्यते}}{\lemma{\skt{स्नाय्व\appbrk{}स्थि\appbrk{}को\appbrk{}ष्ठ\appbrk{}ग\appbrk{}त\appbrk{}श\appbrk{}ल्या\appbrk{}द्ध\appbrk{}र\appbrk{}णा\appbrk{}र्थम्}}\Afootnote{K H; \skt{श्नाद्व\appbrk{}स्छि\appbrk{}को\appbrk{}ष्ठ\appbrk{}ग\appbrk{}त\appbrk{}श\appbrk{}ल्यो\appbrk{}द्व\appbrk{}र\appbrk{}णाव्व}\ N}}}{\lemma{\skt{उपदिश्यते}}\Afootnote{K H; \skt{म्रपदिश्यते}\ N}}\skt{~।}\pend

\pstart \ledsidenote{\passref{1.7.3}}\edtext{\edtext{\edtext{\skt{यन्त्र\-श\-त\-मे\-को\-त्त\-र\-मत्र}}{\lemma{\skt{यन्त्रशतम्}}\Afootnote{A K H; \skt{टानुशतम्}\ N}}}{\lemma{\skt{एकोत्तरम्}}\Afootnote{A K H; \skt{एका}\ N}}}{\lemma{\skt{अत्र}}\Afootnote{A K H; \emph{om.}\ N}}\skt{ }\edtext{\skt{हस्तयन्त्रमेव}}{\lemma{\skt{हस्तयन्त्रम्}}\Afootnote{K H; \skt{हस्तम्}\ A; \skt{यन्त्रम्}\ N}}\skt{ प्रधानतमं }\edtext{\skt{यन्त्रा\-णा\-म\-व\-गच्छ}}{\lemma{\skt{अवगच्छ}}\Afootnote{K N H; \skt{अवगच्छ, (किं}\ A}}\skt{ }\edtext{\skt{तद\-धी\-न\-त्वा\-द्य\-न्त्र\-क\-र्मणां}}{\lemma{\skt{तद\appbrk{}धी\appbrk{}न\appbrk{}त्वा\appbrk{}द्य\appbrk{}न्त्र\appbrk{}क\appbrk{}र्मणां}}\Afootnote{K H; \skt{कारणं ? यस्माद्धस्तादृते यन्त्रा\appbrk{}णा\appbrk{}म\appbrk{}प्र\appbrk{}वृ\appbrk{}त्ति\appbrk{}रेव) तद\appbrk{}धी\appbrk{}न\appbrk{}त्वा\appbrk{}द्य\appbrk{}न्त्र\appbrk{}क\appbrk{}र्म\appbrk{}णाम्}\ A; \skt{तद\appbrk{}धी\appbrk{}न\appbrk{}त्वा\appbrk{}द्य\appbrk{}न्त्र\appbrk{}क\appbrk{}र्मणां}\ N}}\skt{~।}\pend

\pstart \ledsidenote{\passref{1.7.4}}\skt{तत्र मनः शरी\-रा\-बा\-ध\-क\-राणि शल्यानि }\edtext{\skt{तेषा\-मा\-ह\-र\-णो\-पायो}}{\lemma{\skt{तेषा\appbrk{}मा\appbrk{}ह\appbrk{}र\appbrk{}णो\appbrk{}पायो}}\Afootnote{K N; \skt{तेषां आहरणोपायो}\ A; \skt{तेषा\appbrk{}मा\appbrk{}भ\appbrk{}र\appbrk{}णो\appbrk{}पायो}\ H}}\skt{ यन्त्राणि}\pend

\pstart \ledsidenote{\passref{1.7.5}}\skt{तानि षट्प्रकाराणि }\edtext{\skt{भवन्ति}}{\lemma{\skt{भवन्ति}}\Afootnote{K N H; \emph{om.}\ A}}\skt{~। }\edtext{\skt{तद्त्यथा}}{\lemma{\skt{तद्त्यथा}}\Afootnote{\skt{तद्यथा}\ A H; \skt{तद्यथा}\ K N}}\skt{~। स्वस्तिकयन्त्राणि~। सन्दम्शयन्त्राणि~। }\edtext{\skt{ताडयन्त्राणि}}{\lemma{\skt{ताडयन्त्राणि}}\Afootnote{K H; \skt{तालयन्त्राणि}\ A; \skt{तातयन्त्राणि}\ N}}\skt{~। नाडीयन्त्राणि~। }\edtext{\skt{शलाकायन्त्राणि}}{\lemma{\skt{शलाकायन्त्राणि}}\Afootnote{A K H; \skt{शलाकयन्त्राणि}\ N}}\skt{~। उपयन्त्राणि चेति~॥}\pend

\pstart \ledsidenote{\passref{1.7.6}}\skt{तत्र }\edtext{\skt{चतुर्विम्शति}}{\lemma{\skt{चतुर्विम्शति}}\Afootnote{K N; \skt{चतुर्विंशतिः}\ A H}}\skt{ स्वस्तिकयन्त्राणि~। द्वे }\edtext{\skt{सन्दम्शयन्त्रे~। द्वे}}{\lemma{\skt{सन्दम्शयन्त्रे~। द्वे}}\Afootnote{A K H; \skt{सन्दशयन्त्रे~। द्वे}\ N}}\skt{ एव }\edtext{\skt{ताडययन्त्रे}}{\lemma{\skt{ताडययन्त्रे}}\Afootnote{K; \skt{तालयन्त्रे}\ A; \skt{तातययन्त्रे}\ N; \skt{ताडयन्त्रे}\ H}}\skt{~॥ }\edtext{\skt{विंशतिर्नाड्यः}}{\lemma{\skt{विंशतिर्नाड्यः}}\Afootnote{A K; \skt{विंसति नाड्यः}\ N; \skt{विंशतिनाड्यः}\ H}}\skt{ }\edtext{\skt{अष्टाविंशति}}{\lemma{\skt{अष्टाविंशति}}\Afootnote{K N H; \skt{अष्टाविंशतिः}\ A}}\skt{ शलाकाः~। }\edtext{\skt{पञ्चविंशति}}{\lemma{\skt{पञ्चविंशति}}\Afootnote{A K; \skt{पञ्चविंशतिर्}\ N H}}\skt{ }\edtext{\skt{रुपयन्त्राणीति}}{\lemma{\skt{रुपयन्त्राणीति}}\Afootnote{K; \skt{रुपयन्त्राणि}\ A; \skt{उपयन्त्रानीति}\ N; \skt{उपयन्त्राणिति}\ H\textsuperscript{ac}; \skt{उपयन्त्राणीति}\ H\textsuperscript{pc}}}\skt{~।}\pend

\pstart \ledsidenote{\passref{1.7.7}}\skt{तानि प्रायशो }\edtext{\skt{लोहानि}}{\lemma{\skt{लोहानि}}\Afootnote{K N H; \skt{लौहानि}\ A}}\skt{ भवन्ति~। }\edtext{\skt{तत्प्र\-ति\-रू\-प\-काणि वा}}{\lemma{\skt{तत्प्र\appbrk{}ति\appbrk{}रू\appbrk{}प\appbrk{}काणि वा}}\Afootnote{A K H; \skt{तन्प्र\appbrk{}ति\appbrk{}रू\appbrk{}प\appbrk{}काणि वा}\ N}}\skt{ }\edtext{\skt{तदलाभे}}{\lemma{\skt{तदलाभे}}\Afootnote{K; \skt{तदलाभे}\ A H; \emph{om.}\ N}}\pend

\pstart \ledsidenote{\passref{1.7.8}}\edtext{\skt{तत्र}}{\lemma{\skt{तत्र}}\Afootnote{A K H; \skt{तदा लाभे तत्र}\ N}}\skt{ नानाप्रकाराणां }\edtext{\edtext{\skt{व्याडा\-ना\-म्मृ\-ग\-प\-क्षि\-णा\-म्मु\-खै\-र्मु\-खानि}}{\lemma{\skt{व्याडानाम्}}\Afootnote{K N H\textsuperscript{ac}; \skt{व्यालानां}\ A; \skt{व्याडोडानां}\ H\textsuperscript{pc}}}}{\lemma{\skt{मुखानि}}\Afootnote{A K H; \skt{मुस्यानि}\ N}}\skt{ }\edtext{\skt{यन्त्राणां}}{\lemma{\skt{यन्त्राणां}}\Afootnote{A K H; \skt{यन्त्रानिणां}\ N}}\skt{ प्रायशः सदृशानि }\edtext{\skt{भवन्ति}}{\lemma{\skt{भवन्ति}}\Afootnote{K N H; \emph{om.}\ A}}\skt{~। }\edtext{\skt{तस्मा\-त्सा\-रू\-प्या\-दा\-ग\-मा\-दु\-प\-दे\-शा\-द\-न्यत्र दर्शनात्~। युक्तितश्च}}{\lemma{\skt{तस्मा\appbrk{}त्सा\appbrk{}रू\appbrk{}प्या\appbrk{}दा\appbrk{}ग\appbrk{}मा\appbrk{}दु\appbrk{}प\appbrk{}दे\appbrk{}शा\appbrk{}द\appbrk{}न्यत्र दर्शनात्~। युक्तितश्च}}\Afootnote{K; \skt{तस्मा\appbrk{}त्त\appbrk{}त्सा\appbrk{}रू\appbrk{}प्या\appbrk{}दा\appbrk{}ग\appbrk{}मा\appbrk{}दु\appbrk{}प\appbrk{}दे\appbrk{}शा\appbrk{}द\appbrk{}न्य\appbrk{}य\appbrk{}न्त्र\appbrk{}द\appbrk{}र्श\appbrk{}ना\appbrk{}द्यु\appbrk{}क्ति\appbrk{}तश्च}\ A; \skt{तस्मा\appbrk{}त्सा\appbrk{}रू\appbrk{}प्या\appbrk{}दा\appbrk{}ग\appbrk{}मा\appbrk{}दु\appbrk{}प\appbrk{}दे\appbrk{}शा\appbrk{}द\appbrk{}न्य\appbrk{}त्र\appbrk{}द\appbrk{}र्श\appbrk{}ना\appbrk{}त\appbrk{}र\appbrk{}यु\appbrk{}क्ति\appbrk{}तस्च}\ N; \skt{तस्मा\appbrk{}त्त्सा\appbrk{}रू\appbrk{}प्या\appbrk{}दा\appbrk{}ग\appbrk{}मा\appbrk{}दु\appbrk{}प\appbrk{}दे\appbrk{}शा\appbrk{}द\appbrk{}न्य\appbrk{}त्र\appbrk{}द\appbrk{}र्श\appbrk{}नात्~। युक्तितश्च}\ H\textsuperscript{ac}; \skt{तस्मा\appbrk{}त्त्सा\appbrk{}रू\appbrk{}प्या\appbrk{}दा\appbrk{}ग\appbrk{}मा\appbrk{}दु\appbrk{}प\appbrk{}दे\appbrk{}शा\appbrk{}द\appbrk{}न्य\appbrk{}त्र\appbrk{}द\appbrk{}र्श\appbrk{}नात्~। तृ युक्तितश्च}\ H\textsuperscript{pc}}}\skt{ कारयेत्}\pend

\pstart\versepar \ledsidenote{\passref{1.7.9}}\skt{समाहितानि यन्त्राणि~। }\edtext{\skt{खर\-श्ल\-क्ष्ण\-मु\-खानि च}}{\lemma{\skt{खर\appbrk{}श्ल\appbrk{}क्ष्ण\appbrk{}मु\appbrk{}खानि च}}\Afootnote{A K H; \skt{खरश्लमुखानि च}\ N}}\skt{~। }\pend
\pstart\versepar \skt{सुदृढानि सुरूपाणि~। }\edtext{\skt{सुग्रहाणि च}}{\lemma{\skt{सुग्रहाणि च}}\Afootnote{A K\textsuperscript{pc} N H; \skt{सुग्रआहाणि च}\ K\textsuperscript{ac}}}\skt{ कारयेत्~॥}\pend

\pstart \ledsidenote{\passref{1.7.10}}\edtext{\skt{स्वस्ति\-क\-य\-न्त्रा\-ण्य\-ष्टा\-द\-शा\-ङ्गु\-लानि}}{\lemma{\skt{स्वस्ति\appbrk{}क\appbrk{}य\appbrk{}न्त्रा\appbrk{}ण्य\appbrk{}ष्टा\appbrk{}द\appbrk{}शा\appbrk{}ङ्गु\appbrk{}लानि}}\Afootnote{K; \skt{तत्र स्वस्तिकयन्त्राणि- अष्टा\appbrk{}द\appbrk{}शा\appbrk{}ङ्गु\appbrk{}ल\appbrk{}प्र\appbrk{}मा\appbrk{}णानि}\ A; \skt{स्वस्ति\appbrk{}क\appbrk{}य\appbrk{}न्त्रा\appbrk{}ण्य\appbrk{}ष्टा\appbrk{}द\appbrk{}शा\appbrk{}गु\appbrk{}लानि शि}\ N; \skt{स्वस्ति\appbrk{}क\appbrk{}य\appbrk{}न्त्रा\appbrk{}ण्य\appbrk{}ष्टा\appbrk{}द\appbrk{}शा\appbrk{}ङ्गु\appbrk{}लानि}\ H}}\skt{~। }\edtext{\skt{सिंह\-व्या\-घ्र\-त\-र\-क्ष्वृ\-क्ष\-क\-द्वी\-पि\-मा\-र्ज्जा\-र\-शृ\-गा\-ल\-मृ\-गे\-र्वा\-रु\-का\-क\-क\-ङ्क\-कु\-र\-र\-चा\-स\-भास शश\-घा\-त्यु\-लू\-क\-ची\-रि\-ल्ल\-श्ये\-न\-गृ\-द्ध्रो\-त्क्रो\-स\-भृ\-ङ्ग\-रा\-जा\-ञ्ज\-लि\-क\-क\-र्ण्णा\-व\-भ\-ञ्ज\-न\-न\-न्दी\-मु\-ख\-मु\-खानि}}{\lemma{\skt{सिंह\appbrk{}व्या\appbrk{}घ्र\appbrk{}त\appbrk{}र\appbrk{}क्ष्वृ\appbrk{}क्ष\appbrk{}क\appbrk{}द्वी\appbrk{}पि\appbrk{}मा\appbrk{}र्ज्जा\appbrk{}र\appbrk{}शृ\appbrk{}गा\appbrk{}ल\appbrk{}मृ\appbrk{}गे\appbrk{}र्वा\appbrk{}रु\appbrk{}का\appbrk{}क\appbrk{}क\appbrk{}ङ्क\appbrk{}कु\appbrk{}र\appbrk{}र\appbrk{}चा\appbrk{}स\appbrk{}भास शश\appbrk{}घा\appbrk{}त्यु\appbrk{}लू\appbrk{}क\appbrk{}ची\appbrk{}रि\appbrk{}ल्ल\appbrk{}श्ये\appbrk{}न\appbrk{}गृ\appbrk{}द्ध्रो\appbrk{}त्क्रो\appbrk{}स\appbrk{}भृ\appbrk{}ङ्ग\appbrk{}रा\appbrk{}जा\appbrk{}ञ्ज\appbrk{}लि\appbrk{}क\appbrk{}क\appbrk{}र्ण्णा\appbrk{}व\appbrk{}भ\appbrk{}ञ्ज\appbrk{}न\appbrk{}न\appbrk{}न्दी\appbrk{}मु\appbrk{}ख\appbrk{}मु\appbrk{}खानि}}\Afootnote{K\textsuperscript{pc}; \skt{सिंह\appbrk{}व्या\appbrk{}घ्र\appbrk{}वृ\appbrk{}क\appbrk{}त\appbrk{}र\appbrk{}क्ष्वृ\appbrk{}क्ष\appbrk{}द्वी\appbrk{}पि\appbrk{}मा\appbrk{}र्जा\appbrk{}र\appbrk{}शृ\appbrk{}गा\appbrk{}ल\appbrk{}मृ\appbrk{}गै\appbrk{}र्वा\appbrk{}रु\appbrk{}क\appbrk{}का\appbrk{}क\appbrk{}क\appbrk{}ङ्क\appbrk{}कु\appbrk{}र\appbrk{}र\appbrk{}चा\appbrk{}स\appbrk{}भा\appbrk{}स\appbrk{}श\appbrk{}श\appbrk{}घा\appbrk{}त्यु\appbrk{}लू\appbrk{}क\appbrk{}चि\appbrk{}ल्लि\appbrk{}श्येन- गृध्र\appbrk{}क्रौ\appbrk{}ञ्च\appbrk{}भृ\appbrk{}ङ्ग\appbrk{}रा\appbrk{}जा\appbrk{}ञ्ज\appbrk{}लि\appbrk{}क\appbrk{}र्णा\appbrk{}व\appbrk{}भ\appbrk{}ञ्ज\appbrk{}न\appbrk{}न\appbrk{}न्दी\appbrk{}मु\appbrk{}ख\appbrk{}मु\appbrk{}खानि}\ A; \skt{सिंह\appbrk{}व्या\appbrk{}घ्र\appbrk{}त\appbrk{}र\appbrk{}क्ष्वृ\appbrk{}क्ष\appbrk{}क\appbrk{}द्वी\appbrk{}पि\appbrk{}मा\appbrk{}र्ज्जा\appbrk{}र\appbrk{}शृ\appbrk{}गा\appbrk{}ल\appbrk{}मृ\appbrk{}गे\appbrk{}र्वा\appbrk{}रु\appbrk{}का\appbrk{}क\appbrk{}क\appbrk{}ङ्क\appbrk{}कु\appbrk{}र\appbrk{}र\appbrk{}चास शश\appbrk{}घा\appbrk{}त्यु\appbrk{}लू\appbrk{}क\appbrk{}ची\appbrk{}रि\appbrk{}ल्ल\appbrk{}श्ये\appbrk{}न\appbrk{}गृ\appbrk{}द्ध्रो\appbrk{}त्क्रो\appbrk{}स\appbrk{}भृ\appbrk{}ङ्ग\appbrk{}रा\appbrk{}जा\appbrk{}ञ्ज\appbrk{}लि\appbrk{}क\appbrk{}क\appbrk{}र्ण्णा\appbrk{}व\appbrk{}भ\appbrk{}ञ्ज\appbrk{}न\appbrk{}न\appbrk{}न्दी\appbrk{}मु\appbrk{}ख\appbrk{}मु\appbrk{}खानि}\ K\textsuperscript{ac}; \skt{न्हव्या\appbrk{}घ्र\appbrk{}त\appbrk{}र\appbrk{}क्वृ\appbrk{}क्ष\appbrk{}क\appbrk{}द्वी\appbrk{}षि\appbrk{}मा\appbrk{}र्ज\appbrk{}र\appbrk{}शृ\appbrk{}गा\appbrk{}ल\appbrk{}मृ\appbrk{}गे\appbrk{}र्वा\appbrk{}रु\appbrk{}का\appbrk{}क\appbrk{}कं\appbrk{}का\appbrk{}क्त\appbrk{}र\appbrk{}र\appbrk{}चा\appbrk{}स\appbrk{}भा\appbrk{}स\appbrk{}श\appbrk{}म\appbrk{}घा\appbrk{}त्यु\appbrk{}लू\appbrk{}वी\appbrk{}रि\appbrk{}ल्ल\appbrk{}स्ये\appbrk{}न\appbrk{}गृध्रो क्रोश\appbrk{}भृ\appbrk{}ङ्ग\appbrk{}रा\appbrk{}ज\appbrk{}ञ्ज\appbrk{}लि\appbrk{}क\appbrk{}क\appbrk{}ण्ला\appbrk{}व\appbrk{}भ\appbrk{}ञ्ज\appbrk{}न\appbrk{}न\appbrk{}दी\appbrk{}मु\appbrk{}ख\appbrk{}मु\appbrk{}खानि}\ N; \skt{सिंह\appbrk{}व्या\appbrk{}घ्र\appbrk{}त\appbrk{}र\appbrk{}क्ष्वृ\appbrk{}क्ष\appbrk{}क\appbrk{}द्वी\appbrk{}पि\appbrk{}मा\appbrk{}र्ज्जा\appbrk{}र\appbrk{}शृ\appbrk{}गा\appbrk{}ल\appbrk{}मृ\appbrk{}गे\appbrk{}र्व्वा\appbrk{}रु\appbrk{}का\appbrk{}क\appbrk{}क\appbrk{}ङ्क\appbrk{}कु\appbrk{}र\appbrk{}र\appbrk{}चा\appbrk{}स\appbrk{}भा\appbrk{}स\appbrk{}श\appbrk{}श\appbrk{}द्या\appbrk{}त्यु\appbrk{}लू\appbrk{}क\appbrk{}ची\appbrk{}रि\appbrk{}ल्ल\appbrk{}श्ये\appbrk{}न\appbrk{}गृ\appbrk{}द्ध्रो\appbrk{}त्क्रो\appbrk{}स\appbrk{}भृ\appbrk{}ङ्ग\appbrk{}रा\appbrk{}जा\appbrk{}ञ्ज\appbrk{}लि\appbrk{}क\appbrk{}र्ण्णा\appbrk{}व\appbrk{}भ\appbrk{}ञ्ज\appbrk{}न\appbrk{}न\appbrk{}न्दि\appbrk{}मु\appbrk{}ख\appbrk{}मु\appbrk{}खानि}\ H}}\skt{~। }\edtext{\skt{मसूराकृतिभिः}}{\lemma{\skt{मसूराकृतिभिः}}\Afootnote{A K H\textsuperscript{ac}; \skt{मसृराकृतिभिः}\ N; \skt{मसू\appbrk{}यू\appbrk{}रा\appbrk{}कृ\appbrk{}तिभिः}\ H\textsuperscript{pc}}}\skt{ }\edtext{\skt{कीलै\-र\-व\-ब\-द्धानि}}{\lemma{\skt{अवबद्धानि}}\Afootnote{A K; \skt{अबद्धानि}\ N; \skt{अवद्धानि}\ H}}\skt{~। }\edtext{\skt{मूले\-ङ्कु\-श\-व\-दा\-वृ\-त्त\-वा\-र\-ङ्गानि}}{\lemma{\skt{मूले\appbrk{}ङ्कु\appbrk{}श\appbrk{}व\appbrk{}दा\appbrk{}वृ\appbrk{}त्त\appbrk{}वा\appbrk{}र\appbrk{}ङ्गानि}}\Afootnote{K; \skt{मूलेऽ\appbrk{}ङ्कु\appbrk{}श\appbrk{}व\appbrk{}दा\appbrk{}वृ\appbrk{}त्त\appbrk{}वा\appbrk{}र\appbrk{}ङ्गाणि}\ A; \skt{मृल\appbrk{}ङ्कु\appbrk{}श\appbrk{}व\appbrk{}ढ\appbrk{}वृ\appbrk{}त्त\appbrk{}वा\appbrk{}र\appbrk{}ङ्गानि}\ N; \skt{मूले\appbrk{}ङ्कु\appbrk{}श\appbrk{}व\appbrk{}ट\appbrk{}वृ\appbrk{}त्त\appbrk{}वा\appbrk{}र\appbrk{}ङ्गाणि}\ H}}\skt{~। }\edtext{\skt{अस्थि\-वि\-द\-ष्ट\-श\-ल्यो\-द्ध\-र\-णा\-र्थ\-मु\-प\-दि\-श्यन्ते}}{\lemma{\skt{अस्थि\appbrk{}वि\appbrk{}द\appbrk{}ष्ट\appbrk{}श\appbrk{}ल्यो\appbrk{}द्ध\appbrk{}र\appbrk{}णा\appbrk{}र्थम्}}\Afootnote{A K H; \skt{अस्थि\appbrk{}वि\appbrk{}द\appbrk{}ष्ट\appbrk{}श\appbrk{}ल्यो\appbrk{}द्व\appbrk{}र\appbrk{}णा\appbrk{}र्थम्}\ N}}\skt{~।}\pend

\pstart \ledsidenote{\passref{1.7.11}}\edtext{\skt{सनि\-ग्र\-हा\-नु\-ग्रहौ}}{\lemma{\skt{सनि\appbrk{}ग्र\appbrk{}हा\appbrk{}नु\appbrk{}ग्रहौ}}\Afootnote{K H; \skt{सनि\appbrk{}ग्र\appbrk{}होऽ\appbrk{}नि\appbrk{}ग्र\appbrk{}हश्च}\ A; \skt{स निग्रहानुग्रहो}\ N}}\skt{ }\edtext{\skt{सन्दंशौ}}{\lemma{\skt{सन्दंशौ}}\Afootnote{A K H; \skt{सन्दत्सौ}\ N}}\skt{ }\edtext{\skt{षोडशाङ्गुलौ}}{\lemma{\skt{षोडशाङ्गुलौ}}\Afootnote{A K H; \skt{षोडशागुलौ}\ N}}\skt{~। }\edtext{\skt{त्वङ्मा\-न्स\-सि\-रा\-स्ना\-यु\-ग\-त\-श\-ल्यो\-द्ध\-र\-णा\-र्थ\-मु\-प\-दि\-श्येते}}{\lemma{\skt{त्वङ्मा\appbrk{}न्स\appbrk{}सि\appbrk{}रा\appbrk{}स्ना\appbrk{}यु\appbrk{}ग\appbrk{}त\appbrk{}श\appbrk{}ल्यो\appbrk{}द्ध\appbrk{}र\appbrk{}णा\appbrk{}र्थम्}}\Afootnote{K; \skt{भवतः, त्वङ्मां\appbrk{}स\appbrk{}सि\appbrk{}रा\appbrk{}स्ना\appbrk{}यु\appbrk{}ग\appbrk{}त\appbrk{}श\appbrk{}ल्यो\appbrk{}द्ध\appbrk{}र\appbrk{}णा\appbrk{}र्थम्}\ A; \skt{त्वद्मां\appbrk{}स\appbrk{}सि\appbrk{}रा\appbrk{}स्ना\appbrk{}यु\appbrk{}ग\appbrk{}त\appbrk{}श\appbrk{}ल्यो\appbrk{}द्व\appbrk{}र\appbrk{}णा\appbrk{}र्थम्}\ N; \skt{त्वग्मां\appbrk{}स\appbrk{}सि\appbrk{}रा\appbrk{}स्ना\appbrk{}यु\appbrk{}ग\appbrk{}त\appbrk{}श\appbrk{}ल्यो\appbrk{}द्ध\appbrk{}र\appbrk{}णा\appbrk{}र्थम्}\ H}}\skt{~।}\pend

\pstart \ledsidenote{\passref{1.7.12}}\edtext{\skt{ताडयन्त्रे}}{\lemma{\skt{ताडयन्त्रे}}\Afootnote{K H; \skt{तालयन्त्रे}\ A; \skt{तातयन्त्रे}\ N}}\skt{ }\edtext{\skt{द्वे~।द्वादशाङ्गुले}}{\lemma{\skt{द्वे~।द्वादशाङ्गुले}}\Afootnote{K H; \skt{द्वादशाङ्गुले}\ A; \skt{द्वादशाङ्गुल}\ N}}\skt{ }\edtext{\skt{मत्स्य\-ना\-ल\-क\-व\-दे\-क\-ना\-ल\-द्विके}}{\lemma{\skt{मत्स्य\appbrk{}ना\appbrk{}ल\appbrk{}क\appbrk{}व\appbrk{}दे\appbrk{}क\appbrk{}ना\appbrk{}ल\appbrk{}द्विके}}\Afootnote{K; \skt{मत्स्य\appbrk{}ता\appbrk{}ल\appbrk{}व\appbrk{}दे\appbrk{}क\appbrk{}ता\appbrk{}ल\appbrk{}द्वि\appbrk{}ता\appbrk{}लके}\ A; \skt{मत्स्ये\appbrk{}ना\appbrk{}ल\appbrk{}क\appbrk{}व\appbrk{}दे\appbrk{}क\appbrk{}ना\appbrk{}ल\appbrk{}द्वि\appbrk{}ना\appbrk{}लके}\ N; \skt{मत्स्य\appbrk{}ना\appbrk{}ल\appbrk{}क\appbrk{}व\appbrk{}दे\appbrk{}क\appbrk{}ना\appbrk{}ल…के}\ H}}\skt{ }\edtext{\skt{कर्ण्ण\-ना\-सा\-स्रो\-तो\-ग\-त\-श\-ल्यो\-द्ध\-र\-णा\-र्थम्}}{\lemma{\skt{कर्ण्ण\appbrk{}ना\appbrk{}सा\appbrk{}स्रो\appbrk{}तो\appbrk{}ग\appbrk{}त\appbrk{}श\appbrk{}ल्यो\appbrk{}द्ध\appbrk{}र\appbrk{}णा\appbrk{}र्थम्}}\Afootnote{K; \skt{कर्ण\appbrk{}ना\appbrk{}सा\appbrk{}ना\appbrk{}डी\appbrk{}श\appbrk{}ल्या\appbrk{}ना\appbrk{}मा\appbrk{}ह\appbrk{}र\appbrk{}णा\appbrk{}र्थम्}\ A; \skt{वर्ण्ल\appbrk{}ना\appbrk{}सा\appbrk{}स्रो\appbrk{}त्रो\appbrk{}ग\appbrk{}त\appbrk{}श\appbrk{}ल्यो\appbrk{}द्व\appbrk{}र\appbrk{}ण्यर्थ}\ N; \skt{कर्ण्ण\appbrk{}ना\appbrk{}सा\appbrk{}श्रो\appbrk{}त्रो\appbrk{}ग\appbrk{}त\appbrk{}श\appbrk{}ल्यो\appbrk{}द्ध\appbrk{}र\appbrk{}णार्थं}\ H}}\skt{~।}\pend

\pstart \ledsidenote{\passref{1.7.13}}\edtext{\edtext{\edtext{\edtext{\skt{नाडी\-य\-न्त्रा\-ण्य\-ने\-क\-प्र\-का\-रा\-ण्य\-ने\-क\-प्र\-यो\-ज\-ना\-न्य\-ने\-क\-तो\-मु\-खा\-न्यु\-भ\-य\-तो\-मु\-खानि}}{\lemma{\skt{नाडीयन्त्राण्य्}}\Afootnote{A K; \skt{नाडीयन्त्रा}\ N H}}}{\lemma{\skt{अने\appbrk{}क\appbrk{}प्र\appbrk{}का\appbrk{}राण्य्}}\Afootnote{A K; \skt{न्यने\appbrk{}क\appbrk{}प्र\appbrk{}का\appbrk{}राण्य्}\ N; \skt{ण्यनेकप्रकारा}\ H}}}{\lemma{\skt{अने\appbrk{}क\appbrk{}प्र\appbrk{}यो\appbrk{}ज\appbrk{}नान्य्}}\Afootnote{A K; \skt{अने\appbrk{}क\appbrk{}प्र\appbrk{}यो\appbrk{}जना}\ N; \skt{ण्य नेकप्रयोजना}\ H}}}{\lemma{\skt{अने\appbrk{}क\appbrk{}तो\appbrk{}मु\appbrk{}खा\appbrk{}न्यु\appbrk{}भ\appbrk{}य\appbrk{}तो\appbrk{}मु\appbrk{}खानि}}\Afootnote{K; \skt{एक\appbrk{}तो\appbrk{}मु\appbrk{}खा\appbrk{}न्यु\appbrk{}भ\appbrk{}य\appbrk{}तो\appbrk{}मु\appbrk{}खानि च; तानि}\ A; \skt{न्यैक\appbrk{}तो\appbrk{}मु\appbrk{}खो\appbrk{}भ\appbrk{}य\appbrk{}तो\appbrk{}मु\appbrk{}खानि}\ N; \skt{न्यने\appbrk{}क\appbrk{}तो\appbrk{}मु\appbrk{}खा\appbrk{}न्यु\appbrk{}भ\appbrk{}य\appbrk{}तो\appbrk{}मु\appbrk{}खानि}\ H}}\skt{~। }\edtext{\skt{स्रोतो\-ग\-त\-ग\-ल\-श\-ल्यो\-द्ध\-र\-णार्थं}}{\lemma{\skt{स्रोतो\appbrk{}ग\appbrk{}त\appbrk{}ग\appbrk{}ल\appbrk{}श\appbrk{}ल्यो\appbrk{}द्ध\appbrk{}र\appbrk{}णार्थं}}\Afootnote{K; \skt{स्रोतो\appbrk{}ग\appbrk{}त\appbrk{}श\appbrk{}ल्यो\appbrk{}द्ध\appbrk{}र\appbrk{}णार्थं}\ A; \skt{स्रोतो\appbrk{}ग\appbrk{}त\appbrk{}ग\appbrk{}ल\appbrk{}श\appbrk{}ल्यो\appbrk{}द्व\appbrk{}र\appbrk{}णार्थ}\ N; \skt{श्रोतो\appbrk{}ग\appbrk{}त\appbrk{}ग्ल\appbrk{}श\appbrk{}ल्यो\appbrk{}द्ध\appbrk{}र\appbrk{}णार्थं}\ H}}\skt{ }\edtext{\edtext{\skt{क्रिया\-सौ\-क\-र्या\-र्थ\-मा\-चू\-ष\-णार्थं}}{\lemma{\skt{क्रिया\appbrk{}सौ\appbrk{}क\appbrk{}र्या\appbrk{}र्थम्}}\Afootnote{K; \skt{रोग\appbrk{}द\appbrk{}र्श\appbrk{}ना\appbrk{}र्थम्}\ A; \skt{क्रियासौकर्यार्थ}\ N H}}}{\lemma{\skt{आचूषणार्थं}}\Afootnote{A K; \skt{मावृषनार्थ}\ N; \skt{माचूषणार्थं}\ H}}\skt{ }\edtext{\skt{रोग\-द\-र्श\-ना\-र्थञ्च}}{\lemma{\skt{रोग\appbrk{}द\appbrk{}र्श\appbrk{}ना\appbrk{}र्थञ्च}}\Afootnote{K N H; \skt{क्रियासौकर्यार्थं चेति}\ A}}\skt{~। तानि }\edtext{\skt{स्रोतो\-द्वा\-र\-प\-रि\-णा\-हानि}}{\lemma{\skt{स्रोतो\appbrk{}द्वा\appbrk{}र\appbrk{}प\appbrk{}रि\appbrk{}णा\appbrk{}हानि}}\Afootnote{A K H; \skt{स्रोत्र\appbrk{}द्वा\appbrk{}प\appbrk{}रि\appbrk{}णा\appbrk{}हानि}\ N}}\skt{ }\edtext{\skt{यथा\-यो\-ग\-दी\-र्घाणि}}{\lemma{\skt{यथा\appbrk{}यो\appbrk{}ग\appbrk{}दी\appbrk{}र्घाणि}}\Afootnote{A K; \skt{यथा\appbrk{}यो\appbrk{}ग\appbrk{}प्र\appbrk{}दी\appbrk{}र्घाणि}\ N H}}\skt{ }\edtext{\skt{भवन्ति}}{\lemma{\skt{भवन्ति}}\Afootnote{K N H; \skt{च}\ A}}\skt{~। तत्र }\edtext{\edtext{\edtext{\skt{भग\-न्द\-रा\-र्शो\-र्बु\-द\-व्र\-ण\-ब\-स्त्यु\-त्त\-र\-ब\-स्ति\-मू\-त्र\-वृ\-द्धि\-द\-को\-द\-र\-धू\-म\-नि\-रु\-द्ध\-प्र\-क\-श\-स\-न्नि\-रु\-द्ध\-गु\-द\-य\-न्त्रा\-ण्य\-ला\-बू\-शृ\-ङ्ग\-य\-न्त्राणि}}{\lemma{\skt{भगन्दरार्शोर्}}\Afootnote{K; \skt{भगन्दरार्शो}\ A; \skt{भगन्दरार्सो}\ N; \skt{भगन्दरार्णोऽ}\ H}}}{\lemma{\skt{बुद\appbrk{}व्र\appbrk{}ण\appbrk{}ब\appbrk{}स्त्यु\appbrk{}त्त\appbrk{}र\appbrk{}ब\appbrk{}स्ति\appbrk{}मू\appbrk{}त्र\appbrk{}वृ\appbrk{}द्धि\appbrk{}द\appbrk{}को\appbrk{}द\appbrk{}र\appbrk{}धू\appbrk{}म\appbrk{}नि\appbrk{}रु\appbrk{}द्ध\appbrk{}प्र\appbrk{}क\appbrk{}श\appbrk{}स\appbrk{}न्नि\appbrk{}रु\appbrk{}द्ध\appbrk{}गु\appbrk{}द\appbrk{}य\appbrk{}न्त्राण्य्}}\Afootnote{K; \skt{व्रण\appbrk{}ब\appbrk{}स्त्यु\appbrk{}त्त\appbrk{}र\appbrk{}ब\appbrk{}स्ति\appbrk{}मू\appbrk{}त्र\appbrk{}वृ\appbrk{}द्धि\appbrk{}द\appbrk{}को\appbrk{}द\appbrk{}र\appbrk{}धू\appbrk{}म\appbrk{}नि\appbrk{}रु\appbrk{}द्ध\appbrk{}प्र\appbrk{}क\appbrk{}श\appbrk{}स\appbrk{}न्नि\appbrk{}रु\appbrk{}द्ध\appbrk{}गु\appbrk{}द\appbrk{}यन्त्रा}\ A; \skt{र्बुद\appbrk{}व्र\appbrk{}ण\appbrk{}ब\appbrk{}स्त्यु\appbrk{}त्त\appbrk{}र\appbrk{}ब\appbrk{}स्ति\appbrk{}मृ\appbrk{}र्त्त\appbrk{}वृ\appbrk{}द्वि\appbrk{}द\appbrk{}को\appbrk{}द\appbrk{}र\appbrk{}धृ\appbrk{}म\appbrk{}नि\appbrk{}रु\appbrk{}द्व\appbrk{}प्र\appbrk{}क\appbrk{}श\appbrk{}स\appbrk{}न्नि\appbrk{}रु\appbrk{}द्व\appbrk{}गु\appbrk{}द\appbrk{}य\appbrk{}न्त्रण्य्}\ N; \skt{र्ब्बुद\appbrk{}व्र\appbrk{}ण\appbrk{}ब\appbrk{}स्त्यु\appbrk{}त्त\appbrk{}र\appbrk{}ब\appbrk{}स्ति\appbrk{}मू\appbrk{}त्र\appbrk{}वृ\appbrk{}द्धि\appbrk{}द\appbrk{}को\appbrk{}द\appbrk{}र\appbrk{}धू\appbrk{}म\appbrk{}नि\appbrk{}रु\appbrk{}द्ध\appbrk{}प्र\appbrk{}का\appbrk{}श\appbrk{}स\appbrk{}न्नि\appbrk{}रु\appbrk{}द्ध\appbrk{}गु\appbrk{}द\appbrk{}य\appbrk{}न्त्राण्य्}\ H}}}{\lemma{\skt{अला\appbrk{}बू\appbrk{}शृ\appbrk{}ङ्ग\appbrk{}य\appbrk{}न्त्राणि}}\Afootnote{K H; \skt{ण्यला\appbrk{}बू\appbrk{}शृ\appbrk{}ङ्ग\appbrk{}य\appbrk{}न्त्राणि}\ A; \skt{अला\appbrk{}वृ\appbrk{}सृ\appbrk{}ङ्ग\appbrk{}य\appbrk{}न्त्राणि}\ N}}\skt{ }\edtext{\skt{चोप\-रि\-ष्टा\-द्व\-क्ष्यामः}}{\lemma{\skt{चोप\appbrk{}रि\appbrk{}ष्टा\appbrk{}द्व\appbrk{}क्ष्यामः}}\Afootnote{K H; \skt{चोप\appbrk{}रि\appbrk{}ष्टा\appbrk{}द्व\appbrk{}क्ष्यामः}\ A; \skt{चोप\appbrk{}रि\appbrk{}ष्टा\appbrk{}द्व\appbrk{}क्वृमः}\ N}}\skt{~।}\pend

\pstart \ledsidenote{\passref{1.7.14}}\edtext{\skt{शला\-का\-य\-न्त्रा\-ण्यपि}}{\lemma{\skt{शला\appbrk{}का\appbrk{}य\appbrk{}न्त्रा\appbrk{}ण्यपि}}\Afootnote{K; \skt{शला\appbrk{}का\appbrk{}य\appbrk{}न्त्रा\appbrk{}ण्यपि}\ A N H}}\skt{ नानाप्रकाराणि नाना\-प्र\-यो\-ज\-नानि }\edtext{\skt{यथा\-यो\-ग\-प\-रि\-णा\-ह\-प्र\-दी\-र्घाणि}}{\lemma{\skt{यथा\appbrk{}यो\appbrk{}ग\appbrk{}प\appbrk{}रि\appbrk{}णा\appbrk{}ह\appbrk{}प्र\appbrk{}दी\appbrk{}र्घाणि}}\Afootnote{K N H; \skt{यथा\appbrk{}यो\appbrk{}ग\appbrk{}प\appbrk{}रि\appbrk{}णा\appbrk{}ह\appbrk{}दी\appbrk{}र्घाणि}\ A}}\skt{ }\edtext{\skt{भवन्ति}}{\lemma{\skt{भवन्ति}}\Afootnote{K N H; \skt{च}\ A}}\skt{~। }\edtext{\edtext{\skt{तेषा\-ङ्ग\-ण्डू\-प\-द\-श\-र\-पुङ्ख}}{\lemma{\skt{तेषाङ्}}\Afootnote{A K H; \skt{तेषं}\ N}}}{\lemma{\skt{गण्डू\appbrk{}प\appbrk{}द\appbrk{}श\appbrk{}र\appbrk{}पुङ्ख}}\Afootnote{K H; \skt{गण्डू\appbrk{}प\appbrk{}द\appbrk{}स\appbrk{}र्प\appbrk{}फ\appbrk{}णश}\ A; \skt{गण्डु\appbrk{}प\appbrk{}द\appbrk{}स\appbrk{}र\appbrk{}पु\appbrk{}खार}\ N}}\skt{~। }\edtext{\skt{सर्प्पहनु}}{\lemma{\skt{सर्प्पहनु}}\Afootnote{K N H; \skt{रपुङ्ख}\ A}}\skt{~। बडिशमुखे द्वे द्वे~। }\edtext{\edtext{\skt{एष\-ण\-व्यू\-ह\-न\-चा\-ल\-ना\-ह\-र\-णा\-र्थ\-मु\-प\-दि\-श्यन्ते}}{\lemma{\skt{एष\appbrk{}ण\appbrk{}व्यू\appbrk{}ह\appbrk{}न\appbrk{}चा\appbrk{}ल\appbrk{}ना\appbrk{}ह\appbrk{}र\appbrk{}णा\appbrk{}र्थम्}}\Afootnote{A K; \skt{एष\appbrk{}ण\appbrk{}व्या\appbrk{}ह\appbrk{}न\appbrk{}चा\appbrk{}ल\appbrk{}ना\appbrk{}ह\appbrk{}र\appbrk{}णा\appbrk{}र्थम्}\ N; \skt{एष\appbrk{}ण\appbrk{}द\appbrk{}व्यू\appbrk{}ह\appbrk{}न\appbrk{}चा\appbrk{}ल\appbrk{}ना\appbrk{}ह\appbrk{}र\appbrk{}णा\appbrk{}र्थम्}\ H}}}{\lemma{\skt{उपदिश्यन्ते}}\Afootnote{K N H; \skt{उपदिश्येते}\ A}}\skt{~। }\edtext{\skt{मसू\-र\-द\-ल\-मा\-त्र\-मुखे द्वे}}{\lemma{\skt{मसू\appbrk{}र\appbrk{}द\appbrk{}ल\appbrk{}मा\appbrk{}त्र\appbrk{}मुखे द्वे}}\Afootnote{A K H; \skt{मसृ\appbrk{}र\appbrk{}द\appbrk{}ल\appbrk{}मा\appbrk{}त्र\appbrk{}मुखे द्वे}\ N}}\skt{ }\edtext{\skt{किञ्चिदानताग्रे}}{\lemma{\skt{किञ्चिदानताग्रे}}\Afootnote{K N H; \skt{किञ्चिदानताग्रे}\ A}}\skt{ }\edtext{\edtext{\skt{स्रोतो\-ग\-त\-श\-ल्यो\-द्ध\-र\-णा\-र्थ\-म्ष\-ट्का\-र्प्पा\-स\-कृ\-तो\-ष्णी\-ष्णाणि}}{\lemma{\skt{स्रोतो\appbrk{}ग\appbrk{}त\appbrk{}श\appbrk{}ल्यो\appbrk{}द्ध\appbrk{}र\appbrk{}णा\appbrk{}र्थम्}}\Afootnote{A K; \skt{स्रोतो\appbrk{}ग\appbrk{}त\appbrk{}श\appbrk{}ल्यो\appbrk{}द्व\appbrk{}र\appbrk{}णार्थं}\ N; \skt{स्रोतो\appbrk{}ग\appbrk{}त\appbrk{}श\appbrk{}ल्या\appbrk{}द्ध\appbrk{}र\appbrk{}णार्थं}\ H}}}{\lemma{\skt{षट्का\appbrk{}र्प्पा\appbrk{}स\appbrk{}कृ\appbrk{}तो\appbrk{}ष्णी\appbrk{}ष्णाणि}}\Afootnote{K; \skt{षट्का\appbrk{}र्पा\appbrk{}स\appbrk{}कृ\appbrk{}तो\appbrk{}ष्णी\appbrk{}षाणि}\ A; \skt{षट्क\appbrk{}र्प्या\appbrk{}स\appbrk{}कृ\appbrk{}तो\appbrk{}ष्णी\appbrk{}षाणि}\ N; \skt{षट्का\appbrk{}र्प्या\appbrk{}म\appbrk{}कृ\appbrk{}तो\appbrk{}ष्णी\appbrk{}ष्णाणि}\ H\textsuperscript{ac}; \skt{षट्का\appbrk{}र्प्या\appbrk{}स\appbrk{}कृ\appbrk{}तो\appbrk{}ष्णी\appbrk{}ष्णाणि}\ H\textsuperscript{pc}}}\skt{~। प्रमा\-र्ज\-न\-क्रि\-यासु~। }\edtext{\skt{क्षारौ\-ष\-ध\-प्र\-णि\-धा\-ना\-र्थ\-न्त्रीणि}}{\lemma{\skt{क्षारौ\appbrk{}ष\appbrk{}ध\appbrk{}प्र\appbrk{}णि\appbrk{}धा\appbrk{}ना\appbrk{}र्थ\appbrk{}न्त्रीणि}}\Afootnote{K H; \skt{त्रीणि}\ A; \skt{क्षारौ\appbrk{}ष\appbrk{}ध\appbrk{}प्र\appbrk{}णि\appbrk{}धा\appbrk{}नार्थ त्रीणि}\ N}}\skt{ }\edtext{\skt{दर्व्याकृतीनि}}{\lemma{\skt{दर्व्याकृतीनि}}\Afootnote{A K H; \skt{द्रर्व्याकृतीणि}\ N}}\skt{ }\edtext{\skt{खल्लमुखानि}}{\lemma{\skt{खल्लमुखानि}}\Afootnote{K H\textsuperscript{ac}; \skt{खल्लमुखानि, क्षारौ\appbrk{}ष\appbrk{}ध\appbrk{}प्र\appbrk{}णि\appbrk{}धा\appbrk{}नार्थं ; त्रीण्यन्यानि}\ A; \skt{खल्लमुखा}\ N; \skt{खल्लमुखानि~। जंब्ब\appbrk{}ष्ट\appbrk{}ज\appbrk{}म्बु\appbrk{}जस}\ H\textsuperscript{pc}}}\skt{~। }\edtext{\edtext{\skt{जाम्ब्वं\-कु\-श\-व\-द\-ना\-न्य\-ग्नि\-क\-र्माणि}}{\lemma{\skt{जाम्ब्वं\appbrk{}कु\appbrk{}श\appbrk{}व\appbrk{}द\appbrk{}नान्य्}}\Afootnote{K; \skt{जाम्ब\appbrk{}व\appbrk{}व\appbrk{}द\appbrk{}नानि, त्रीण्य\appbrk{}ङ्कु\appbrk{}श\appbrk{}व\appbrk{}द\appbrk{}नानि}\ A; \skt{नि~। जव्कं\appbrk{}कु\appbrk{}श\appbrk{}व\appbrk{}द\appbrk{}नान्य्}\ N; \skt{जंब्बष्टवदना}\ H\textsuperscript{ac}; \skt{वदना}\ H\textsuperscript{pc}}}}{\lemma{\skt{अग्निकर्माणि}}\Afootnote{K N; \skt{षडे\appbrk{}वा\appbrk{}ग्नि\appbrk{}क\appbrk{}र्मस्व}\ A; \skt{न्यग्निकर्म्माणि}\ H}}\skt{ }\edtext{\skt{त्रीणि}}{\lemma{\skt{त्रीणि}}\Afootnote{K H; \skt{भिप्रे}\ A; \skt{त्रीणी}\ N}}\skt{ }\edtext{\skt{त्रिणि}}{\lemma{\skt{त्रिणि}}\Afootnote{K; \skt{तानि}\ A; \skt{त्रीणी}\ N; \skt{त्रीणि}\ H}}\skt{~। }\edtext{\skt{नासा\-र्बु\-द\-ह\-र\-णा\-र्थ\-मेकं}}{\lemma{\skt{नासा\appbrk{}र्बु\appbrk{}द\appbrk{}ह\appbrk{}र\appbrk{}णा\appbrk{}र्थ\appbrk{}मेकं}}\Afootnote{A K; \skt{नासा\appbrk{}र्ब्बु\appbrk{}द\appbrk{}ह\appbrk{}र\appbrk{}णा\appbrk{}र्थ\appbrk{}मेकं}\ N; \skt{नामा\appbrk{}र्ब्बु\appbrk{}द\appbrk{}ह\appbrk{}र\appbrk{}णा\appbrk{}र्थ\appbrk{}मेकं}\ H}}\skt{ }\edtext{\skt{कोला\-स्थि\-द\-ल\-मात्रं}}{\lemma{\skt{कोला\appbrk{}स्थि\appbrk{}द\appbrk{}ल\appbrk{}मात्रं}}\Afootnote{K N H\textsuperscript{ac}; \skt{कोला\appbrk{}स्थि\appbrk{}द\appbrk{}ल\appbrk{}मा\appbrk{}त्र\appbrk{}मुखं}\ A; \skt{कोला\appbrk{}स्थि\appbrk{}द\appbrk{}ल\appbrk{}भ\appbrk{}य३\appbrk{}मात्रं}\ H\textsuperscript{pc}}}\skt{ }\edtext{\skt{खल्लतीक्ष्णोष्ठं}}{\lemma{\skt{खल्लतीक्ष्णोष्ठं}}\Afootnote{K H; \skt{खल्लतीक्ष्णौष्ठं}\ A; \skt{खल्वतीक्ष्णोष्ठ}\ N}}\skt{~। }\edtext{\skt{अञ्जनार्थमेकं}}{\lemma{\skt{अञ्जनार्थमेकं}}\Afootnote{A K; \skt{अञ्जनार्थमेक}\ N; \skt{अञ्जनार्थमेकं}\ H}}\skt{ }\edtext{\skt{कला\-य\-प\-रि\-म\-ण्ड\-ल\-मु\-भ\-यतो}}{\lemma{\skt{कला\appbrk{}य\appbrk{}प\appbrk{}रि\appbrk{}म\appbrk{}ण्ड\appbrk{}ल\appbrk{}मु\appbrk{}भ\appbrk{}यतो}}\Afootnote{A K; \skt{कला\appbrk{}य\appbrk{}प\appbrk{}रि\appbrk{}म\appbrk{}ण्ड\appbrk{}ल\appbrk{}मु\appbrk{}भु\appbrk{}यतो}\ N; \skt{कला\appbrk{}य\appbrk{}प\appbrk{}रि\appbrk{}म\appbrk{}ण्ड\appbrk{}ल\appbrk{}मु\appbrk{}भ\appbrk{}यतो}\ H}}\skt{ मुकुलाग्रं~। }\edtext{\skt{मूत्र\-मा\-र्ग\-वि\-शु\-द्ध्य\-र्थ\-मेकं}}{\lemma{\skt{मूत्र\appbrk{}मा\appbrk{}र्ग\appbrk{}वि\appbrk{}शु\appbrk{}द्ध्य\appbrk{}र्थ\appbrk{}मेकं}}\Afootnote{K; \skt{मूत्र\appbrk{}मा\appbrk{}र्ग\appbrk{}वि\appbrk{}शो\appbrk{}ध\appbrk{}ना\appbrk{}र्थ\appbrk{}मेकं}\ A; \skt{मृत्र\appbrk{}मो\appbrk{}ज\appbrk{}वि\appbrk{}शु\appbrk{}द्ध्य\appbrk{}र्थ\appbrk{}मेकं}\ N; \skt{मूत्र\appbrk{}मा\appbrk{}र्ग्ग\appbrk{}वि\appbrk{}शु\appbrk{}ध्य\appbrk{}र्थ\appbrk{}मेकं}\ H\textsuperscript{ac}; \skt{मूडो\appbrk{}त्र\appbrk{}मा\appbrk{}र्ग्ग\appbrk{}वि\appbrk{}शु\appbrk{}ध्य\appbrk{}र्थ\appbrk{}मेकं}\ H\textsuperscript{pc}}}\skt{ }\edtext{\skt{माल\-ती\-पु\-ष्प\-वृ\-न्ता\-ग्र\-प्र\-मा\-ण\-प\-रि\-म\-ण्ड\-ल\-मिति}}{\lemma{\skt{माल\appbrk{}ती\appbrk{}पु\appbrk{}ष्प\appbrk{}वृ\appbrk{}न्ता\appbrk{}ग्र\appbrk{}प्र\appbrk{}मा\appbrk{}ण\appbrk{}प\appbrk{}रि\appbrk{}म\appbrk{}ण्ड\appbrk{}ल\appbrk{}मिति}}\Afootnote{K; \skt{माल\appbrk{}ती\appbrk{}पु\appbrk{}ष्प\appbrk{}वृ\appbrk{}न्ता\appbrk{}ग्र\appbrk{}प्र\appbrk{}मा\appbrk{}ण\appbrk{}प\appbrk{}रि\appbrk{}म\appbrk{}ण्ड\appbrk{}ल\appbrk{}मिति}\ A; \skt{माल\appbrk{}ती\appbrk{}पु\appbrk{}ष्प\appbrk{}वृ\appbrk{}त्ता\appbrk{}ग्र\appbrk{}प्र\appbrk{}मा\appbrk{}ण\appbrk{}प\appbrk{}रि\appbrk{}म\appbrk{}ण्ड\appbrk{}ल\appbrk{}मिति}\ N; \skt{माल\appbrk{}ती\appbrk{}पु\appbrk{}ष्प…\appbrk{}वृ\appbrk{}त्ता\appbrk{}ग्र\appbrk{}प्र\appbrk{}मा\appbrk{}ण\appbrk{}प\appbrk{}रि\appbrk{}म\appbrk{}ण्ड\appbrk{}ल\appbrk{}मिति}\ H\textsuperscript{ac}; \skt{माल\appbrk{}ती\appbrk{}पु\appbrk{}ष्प…\appbrk{}वृ\appbrk{}त्ता\appbrk{}ग्र\appbrk{}प्र\appbrk{}मा\appbrk{}ण\appbrk{}प\appbrk{}रि\appbrk{}म\appbrk{}ल\appbrk{}मिति}\ H\textsuperscript{pc}}}\skt{~॥}\pend

\pstart \ledsidenote{\passref{1.7.15}}\edtext{\skt{उपयन्त्राण्यपि}}{\lemma{\skt{उपयन्त्राण्यपि}}\Afootnote{K; \skt{उपयन्त्राण्यपि}\ A N H}}\skt{ }\edtext{\skt{रज्जु\-वे\-णि\-का\-च\-र्मा\-न्त\-र्व\-ल्क\-ल\-ल\-ता\-व\-स्त्रा\-ष्ठी\-ला\-ष्म\-न्त\-क\-मु\-द्ग\-र\-पा\-णि\-पा\-द\-त\-ला\-ङ्गु\-लि\-जि\-ह्वा\-द\-न्त\-न\-ख\-मु\-ख\-बा\-ला\-श्म\-शाखा ष्ठीव\-न\-प्र\-वा\-ह\-ण\-ह\-र्षा\-म\-स्का\-र्त्त\-ग\-तानि}}{\lemma{\skt{रज्जु\appbrk{}वे\appbrk{}णि\appbrk{}का\appbrk{}च\appbrk{}र्मा\appbrk{}न्त\appbrk{}र्व\appbrk{}ल्क\appbrk{}ल\appbrk{}ल\appbrk{}ता\appbrk{}व\appbrk{}स्त्रा\appbrk{}ष्ठी\appbrk{}ला\appbrk{}ष्म\appbrk{}न्त\appbrk{}क\appbrk{}मु\appbrk{}द्ग\appbrk{}र\appbrk{}पा\appbrk{}णि\appbrk{}पा\appbrk{}द\appbrk{}त\appbrk{}ला\appbrk{}ङ्गु\appbrk{}लि\appbrk{}जि\appbrk{}ह्वा\appbrk{}द\appbrk{}न्त\appbrk{}न\appbrk{}ख\appbrk{}मु\appbrk{}ख\appbrk{}बा\appbrk{}ला\appbrk{}श्म\appbrk{}शाखा ष्ठीव\appbrk{}न\appbrk{}प्र\appbrk{}वा\appbrk{}ह\appbrk{}ण\appbrk{}ह\appbrk{}र्षा\appbrk{}म\appbrk{}स्का\appbrk{}र्त्त\appbrk{}ग\appbrk{}तानि}}\Afootnote{K\textsuperscript{pc}; \skt{रज्जु\appbrk{}वे\appbrk{}णि\appbrk{}का\appbrk{}प\appbrk{}ट्ट\appbrk{}च\appbrk{}र्मा\appbrk{}न्त\appbrk{}व\appbrk{}ल्क\appbrk{}ल\appbrk{}ल\appbrk{}ता\appbrk{}व\appbrk{}स्त्रा\appbrk{}ष्ठी\appbrk{}ला\appbrk{}श्म\appbrk{}मु\appbrk{}द्ग\appbrk{}र\appbrk{}पा\appbrk{}णि\appbrk{}पा\appbrk{}द\appbrk{}त\appbrk{}ला\appbrk{}ङ्गु\appbrk{}लि\appbrk{}जिह्वा- दन्त\appbrk{}न\appbrk{}ख\appbrk{}मु\appbrk{}ख\appbrk{}बा\appbrk{}ला\appbrk{}श्व\appbrk{}क\appbrk{}ट\appbrk{}क\appbrk{}शा\appbrk{}खा\appbrk{}ष्ठी\appbrk{}व\appbrk{}न\appbrk{}प्र\appbrk{}वा\appbrk{}ह\appbrk{}ण\appbrk{}ह\appbrk{}र्षा\appbrk{}य\appbrk{}स्का\appbrk{}न्त\appbrk{}म\appbrk{}यानि}\ A; \skt{रज्जु\appbrk{}वे\appbrk{}णि\appbrk{}का\appbrk{}च\appbrk{}र्मा\appbrk{}न्त\appbrk{}र्व\appbrk{}ल्क\appbrk{}ल\appbrk{}ल\appbrk{}ता\appbrk{}व\appbrk{}स्त्रा\appbrk{}ष्ठी\appbrk{}ला\appbrk{}ष्म\appbrk{}न्त\appbrk{}क\appbrk{}मु\appbrk{}द्ग\appbrk{}र\appbrk{}पा\appbrk{}णि\appbrk{}पा\appbrk{}द\appbrk{}त\appbrk{}ला\appbrk{}ङ्गु\appbrk{}लि\appbrk{}जि\appbrk{}ह्वा\appbrk{}द\appbrk{}न्त\appbrk{}न\appbrk{}ख\appbrk{}मु\appbrk{}ख\appbrk{}बा\appbrk{}ला\appbrk{}श्मशा ष्ठीव\appbrk{}न\appbrk{}प्र\appbrk{}वा\appbrk{}ह\appbrk{}ण\appbrk{}ह\appbrk{}र्षा\appbrk{}म\appbrk{}स्का\appbrk{}र्त्त\appbrk{}ग\appbrk{}तानि}\ K\textsuperscript{ac}; \skt{रज्जु\appbrk{}वे\appbrk{}णि\appbrk{}का\appbrk{}च\appbrk{}र्मा\appbrk{}न्त\appbrk{}न्व\appbrk{}ल्क\appbrk{}ला\appbrk{}ल\appbrk{}ता\appbrk{}व\appbrk{}स्त्रा\appbrk{}ष्ठी\appbrk{}ला\appbrk{}ष्म\appbrk{}न्त\appbrk{}क\appbrk{}मु\appbrk{}द्ग\appbrk{}र\appbrk{}पा\appbrk{}णि\appbrk{}पा\appbrk{}द\appbrk{}त\appbrk{}लां\appbrk{}गु\appbrk{}लि\appbrk{}जि\appbrk{}ह्वा\appbrk{}द\appbrk{}न्त\appbrk{}न\appbrk{}ख\appbrk{}मु\appbrk{}ख\appbrk{}बा\appbrk{}ला\appbrk{}श्म\appbrk{}शा\appbrk{}खा\appbrk{}ष्ठी\appbrk{}व\appbrk{}न\appbrk{}प्र\appbrk{}वा\appbrk{}ह\appbrk{}ण\appbrk{}ह\appbrk{}र्षा\appbrk{}म\appbrk{}स्का\appbrk{}र्त्त\appbrk{}ग\appbrk{}तानि}\ N; \skt{रज्जु\appbrk{}वे\appbrk{}णि\appbrk{}का\appbrk{}च\appbrk{}र्म्मा\appbrk{}न्त\appbrk{}व\appbrk{}ल्क\appbrk{}ल\appbrk{}ल\appbrk{}ता\appbrk{}व\appbrk{}स्त्रा\appbrk{}ष्ठी\appbrk{}ला\appbrk{}म्मा\appbrk{}न्त\appbrk{}मु\appbrk{}द्ग\appbrk{}र\appbrk{}पा\appbrk{}णि\appbrk{}पा\appbrk{}द\appbrk{}त\appbrk{}ला\appbrk{}ङ्गु\appbrk{}लि\appbrk{}जि(\appbrk{}ह्वा)\appbrk{}द\appbrk{}न्त\appbrk{}न\appbrk{}ख\appbrk{}मु\appbrk{}ख\appbrk{}बा\appbrk{}ला\appbrk{}श्म\appbrk{}शा\appbrk{}खा\appbrk{}ष्ठी\appbrk{}व\appbrk{}न\appbrk{}प्र\appbrk{}वा\appbrk{}ह\appbrk{}न\appbrk{}ह\appbrk{}र्षा\appbrk{}म\appbrk{}स्का\appbrk{}र्त्त\appbrk{}ग\appbrk{}तनि}\ H\textsuperscript{ac}; \skt{रज्जु\appbrk{}वे\appbrk{}णि\appbrk{}का\appbrk{}च\appbrk{}र्म्मा\appbrk{}न्त\appbrk{}व\appbrk{}ल्क\appbrk{}ल\appbrk{}ल\appbrk{}ता\appbrk{}व\appbrk{}स्त्रा\appbrk{}ष्ठी\appbrk{}ला\appbrk{}स्मा\appbrk{}न्त\appbrk{}मु\appbrk{}द्ग\appbrk{}र\appbrk{}पा\appbrk{}णि\appbrk{}पा\appbrk{}द\appbrk{}त\appbrk{}ला\appbrk{}ङ्गु\appbrk{}लि\appbrk{}जि\appbrk{}ह्व(\appbrk{}ह्वा)\appbrk{}द\appbrk{}न्त\appbrk{}न\appbrk{}ख\appbrk{}मु\appbrk{}ख\appbrk{}बा\appbrk{}ला\appbrk{}श्म\appbrk{}शा\appbrk{}खा\appbrk{}ष्ठी\appbrk{}व\appbrk{}न\appbrk{}प्र\appbrk{}वा\appbrk{}ह\appbrk{}न\appbrk{}ह\appbrk{}र्षा\appbrk{}म\appbrk{}स्का\appbrk{}र्त्त\appbrk{}ग\appbrk{}तनि}\ H\textsuperscript{pc}}}\skt{ }\edtext{\skt{क्षारा\-ग्नि\-भे\-ष\-जानि}}{\lemma{\skt{क्षारा\appbrk{}ग्नि\appbrk{}भे\appbrk{}ष\appbrk{}जानि}}\Afootnote{A K H; \skt{खोरा\appbrk{}ग्नि\appbrk{}भे\appbrk{}ष\appbrk{}जानि}\ N}}\skt{ चेति~॥}\pend

\pstart\versepar \ledsidenote{\passref{1.7.16}}\skt{एतानि देहे सर्वस्मिं }\edtext{\skt{देहस्यावयवे}}{\lemma{\skt{देहस्यावयवे}}\Afootnote{A K N; \skt{देहम्स्यावयवे}\ H\textsuperscript{ac}; \skt{देहस्स्यावयवे}\ H\textsuperscript{pc}}}\skt{ तथा }\pend
\pstart\versepar \skt{}\edtext{\skt{सन्धौ}}{\lemma{\skt{सन्धौ}}\Afootnote{A K H; \skt{सन्धो}\ N}}\skt{ }\edtext{\skt{कोष्ठे}}{\lemma{\skt{कोष्ठे}}\Afootnote{A K H; \skt{कोष्ठ}\ N}}\skt{ }\edtext{\skt{धमन्या च}}{\lemma{\skt{धमन्या च}}\Afootnote{K; \skt{धमन्यां च}\ A; \skt{धमन्यो च}\ N; \skt{धमन्या च …}\ H}}\skt{ यथायोगं प्रयोजयेत्~॥}\pend

\pstart \ledsidenote{\passref{1.7.17}}\skt{यन्त्रकर्माणि तु }\edtext{\skt{दुष्ट\-व्र\-ण\-ब\-न्ध\-न\-व्य\-ञ्ज\-न\-व\-र्त्त\-न\-भ्रा\-म\-ण\-पी\-ड\-न\-ता\-ड\-न\-नि\-व\-र\-ण\-वि\-व\-र\-ण\-सी\-व\-ण\-मा\-र्ग\-शो\-ध\-न\-वि\-क\-र्ष\-णा\-ह\-र\-णा\-ञ्छ\-नो\-न्न\-म\-न\-वि\-रे\-च\-भ\-ञ्ज\-नो\-न्म\-थ\-न\-चू\-ष\-णै\-ष\-ण\-दा\-र\-ण\-ऋ\-जू\-क\-र\-ण\-प्र\-क्षा\-ल\-न\-प्र\-ध\-म\-न\-प्र\-मा\-र्ज\-नानि}}{\lemma{\skt{दुष्ट\appbrk{}व्र\appbrk{}ण\appbrk{}ब\appbrk{}न्ध\appbrk{}न\appbrk{}व्य\appbrk{}ञ्ज\appbrk{}न\appbrk{}व\appbrk{}र्त्त\appbrk{}न\appbrk{}भ्रा\appbrk{}म\appbrk{}ण\appbrk{}पी\appbrk{}ड\appbrk{}न\appbrk{}ता\appbrk{}ड\appbrk{}न\appbrk{}नि\appbrk{}व\appbrk{}र\appbrk{}ण\appbrk{}वि\appbrk{}व\appbrk{}र\appbrk{}ण\appbrk{}सी\appbrk{}व\appbrk{}ण\appbrk{}मा\appbrk{}र्ग\appbrk{}शो\appbrk{}ध\appbrk{}न\appbrk{}वि\appbrk{}क\appbrk{}र्ष\appbrk{}णा\appbrk{}ह\appbrk{}र\appbrk{}णा\appbrk{}ञ्छ\appbrk{}नो\appbrk{}न्न\appbrk{}म\appbrk{}न\appbrk{}वि\appbrk{}रे\appbrk{}च\appbrk{}भ\appbrk{}ञ्ज\appbrk{}नो\appbrk{}न्म\appbrk{}थ\appbrk{}न\appbrk{}चू\appbrk{}ष\appbrk{}णै\appbrk{}ष\appbrk{}ण\appbrk{}दा\appbrk{}र\appbrk{}ण\appbrk{}ऋ\appbrk{}जू\appbrk{}क\appbrk{}र\appbrk{}ण\appbrk{}प्र\appbrk{}क्षा\appbrk{}ल\appbrk{}न\appbrk{}प्र\appbrk{}ध\appbrk{}म\appbrk{}न\appbrk{}प्र\appbrk{}मा\appbrk{}र्ज\appbrk{}नानि}}\Afootnote{K; \skt{निर्घा\appbrk{}त\appbrk{}न\appbrk{}पू\appbrk{}र\appbrk{}ण\appbrk{}ब\appbrk{}न्ध\appbrk{}न\appbrk{}व्यू\appbrk{}ह\appbrk{}न\appbrk{}व\appbrk{}र्त\appbrk{}न\appbrk{}चा\appbrk{}ल\appbrk{}न\appbrk{}वि\appbrk{}व\appbrk{}र्त\appbrk{}न\appbrk{}वि\appbrk{}व\appbrk{}र\appbrk{}ण\appbrk{}पी\appbrk{}ड\appbrk{}न\appbrk{}मा\appbrk{}र्ग\appbrk{}वि\appbrk{}शो\appbrk{}ध\appbrk{}न\appbrk{}वि\appbrk{}क\appbrk{}र्ष\appbrk{}णा\appbrk{}ह\appbrk{}र\appbrk{}णा\appbrk{}ञ्छ\appbrk{}नो\appbrk{}न्न\appbrk{}मन- विन\appbrk{}म\appbrk{}न\appbrk{}भ\appbrk{}ञ्ज\appbrk{}नो\appbrk{}न्म\appbrk{}थ\appbrk{}ना\appbrk{}चू\appbrk{}ष\appbrk{}णै\appbrk{}ष\appbrk{}ण\appbrk{}दा\appbrk{}र\appbrk{}ण\appbrk{}र्जू\appbrk{}क\appbrk{}र\appbrk{}ण\appbrk{}प्र\appbrk{}क्षा\appbrk{}ल\appbrk{}न\appbrk{}प्र\appbrk{}ध\appbrk{}म\appbrk{}न\appbrk{}प्र\appbrk{}मा\appbrk{}र्ज\appbrk{}नानि}\ A; \skt{दुष्ट\appbrk{}प्र\appbrk{}ण\appbrk{}ब\appbrk{}न्ध\appbrk{}न\appbrk{}व्य\appbrk{}ज\appbrk{}न\appbrk{}व\appbrk{}र्त्त\appbrk{}न\appbrk{}भ्रा\appbrk{}म\appbrk{}ण\appbrk{}पी\appbrk{}ड\appbrk{}न\appbrk{}ता\appbrk{}ड\appbrk{}न\appbrk{}नि\appbrk{}व\appbrk{}र\appbrk{}ण\appbrk{}वि\appbrk{}व\appbrk{}र\appbrk{}ण\appbrk{}सी\appbrk{}व\appbrk{}न\appbrk{}मा\appbrk{}र्ग\appbrk{}सो\appbrk{}ध\appbrk{}न\appbrk{}वि\appbrk{}क\appbrk{}र्ष\appbrk{}णा\appbrk{}ह\appbrk{}र\appbrk{}णा\appbrk{}ञ्छ\appbrk{}नो\appbrk{}न्ना\appbrk{}म\appbrk{}न\appbrk{}वि\appbrk{}र\appbrk{}म\appbrk{}न\appbrk{}भ\appbrk{}ञ्ज\appbrk{}नो\appbrk{}न्म\appbrk{}थ\appbrk{}न\appbrk{}वृ\appbrk{}ष\appbrk{}णै\appbrk{}ष\appbrk{}ण\appbrk{}दा\appbrk{}र\appbrk{}ण\appbrk{}ह्जृ\appbrk{}क\appbrk{}र\appbrk{}ण\appbrk{}प्र\appbrk{}क्षा\appbrk{}ल\appbrk{}ण\appbrk{}प्र\appbrk{}ध\appbrk{}म\appbrk{}न\appbrk{}प्र\appbrk{}मा\appbrk{}र्ज्ज\appbrk{}नानि}\ N; \skt{दुष्ट\appbrk{}व्र\appbrk{}ण\appbrk{}ब\appbrk{}न्ध\appbrk{}न\appbrk{}व्य\appbrk{}ञ्ज\appbrk{}न\appbrk{}व\appbrk{}र्त्त\appbrk{}न\appbrk{}भ्रा\appbrk{}म\appbrk{}ण\appbrk{}पी\appbrk{}ड\appbrk{}न\appbrk{}ता\appbrk{}ड\appbrk{}न\appbrk{}नि\appbrk{}व\appbrk{}र\appbrk{}ण\appbrk{}वि\appbrk{}व\appbrk{}र\appbrk{}ण\appbrk{}सी\appbrk{}व\appbrk{}न\appbrk{}मा\appbrk{}र्ग्ग\appbrk{}शो\appbrk{}ध\appbrk{}न\appbrk{}वि\appbrk{}क\appbrk{}र्ष\appbrk{}णा\appbrk{}ह\appbrk{}र\appbrk{}णा\appbrk{}ञ्छ\appbrk{}नो\appbrk{}न्न\appbrk{}म\appbrk{}न\appbrk{}वि\appbrk{}र\appbrk{}म\appbrk{}ण\appbrk{}भ\appbrk{}ञ्ज\appbrk{}नो\appbrk{}न्म\appbrk{}थ\appbrk{}चू\appbrk{}ष\appbrk{}णै\appbrk{}ष\appbrk{}ण\appbrk{}दा\appbrk{}र\appbrk{}ण\appbrk{}ऋ\appbrk{}जू\appbrk{}क\appbrk{}र\appbrk{}ण\appbrk{}प्र\appbrk{}क्षा\appbrk{}ल\appbrk{}न\appbrk{}प्र\appbrk{}ध\appbrk{}म\appbrk{}न\appbrk{}प्र\appbrk{}मा\appbrk{}र्ज्ज\appbrk{}नानि}\ H\textsuperscript{ac}; \skt{दुष्ट\appbrk{}व्र\appbrk{}ण\appbrk{}ब\appbrk{}न्ध\appbrk{}न\appbrk{}व्य\appbrk{}ञ्ज\appbrk{}न\appbrk{}व\appbrk{}र्त्त\appbrk{}न\appbrk{}भ्रा\appbrk{}म\appbrk{}ण\appbrk{}पी\appbrk{}ड\appbrk{}न\appbrk{}ता\appbrk{}ड\appbrk{}न\appbrk{}नि\appbrk{}व\appbrk{}र\appbrk{}ण\appbrk{}वि\appbrk{}व\appbrk{}र\appbrk{}ण\appbrk{}सी\appbrk{}व\appbrk{}न\appbrk{}मा\appbrk{}र्ग्ग\appbrk{}शो\appbrk{}ध\appbrk{}न\appbrk{}वि\appbrk{}क\appbrk{}र्ष\appbrk{}णा\appbrk{}ह\appbrk{}र\appbrk{}णा\appbrk{}ञ्छ\appbrk{}नो\appbrk{}न्न\appbrk{}म\appbrk{}न\appbrk{}वि\appbrk{}र\appbrk{}म\appbrk{}ण\appbrk{}भ\appbrk{}ञ्ज\appbrk{}नो\appbrk{}न्म\appbrk{}थ\appbrk{}न\appbrk{}चू\appbrk{}ष\appbrk{}णै\appbrk{}ष\appbrk{}ण\appbrk{}दा\appbrk{}र\appbrk{}ण\appbrk{}ऋ\appbrk{}जू\appbrk{}क\appbrk{}र\appbrk{}ण\appbrk{}प्र\appbrk{}क्षा\appbrk{}ल\appbrk{}न\appbrk{}प्र\appbrk{}ध\appbrk{}म\appbrk{}न\appbrk{}प्र\appbrk{}मा\appbrk{}र्ज्ज\appbrk{}नानि}\ H\textsuperscript{pc}}}\skt{ }\edtext{\skt{चतुर्विंशति}}{\lemma{\skt{चतुर्विंशति}}\Afootnote{K; \skt{चतुर्विंशतिः}\ A H; \skt{चतुविंसति}\ N}}\skt{ }\edtext{\skt{भवन्ति}}{\lemma{\skt{भवन्ति}}\Afootnote{K N; \emph{om.}\ A; \skt{भवति}\ H}}\skt{~॥}\pend

\pstart\versepar \ledsidenote{\passref{1.7.18}}\edtext{\skt{स्वबुध्या}}{\lemma{\skt{स्वबुध्या}}\Afootnote{A K N; \skt{स्वबु}\ H}}\skt{ }\edtext{\skt{विभ\-जे\-द्य\-न्त्र\-य\-न्त्र\-क\-र्माणि}}{\lemma{\skt{विभ\appbrk{}जे\appbrk{}द्य\appbrk{}न्त्र\appbrk{}य\appbrk{}न्त्र\appbrk{}क\appbrk{}र्माणि}}\Afootnote{K; \skt{चापि विभ\appbrk{}जे\appbrk{}द्य\appbrk{}न्त्र\appbrk{}क\appbrk{}र्माणि}\ A; \skt{विभजेद्युक्ता यन्त्रकर्माणि}\ N; \skt{र्ध्या विभ\appbrk{}जे\appbrk{}द्यु\appbrk{}ज्या\appbrk{}य\appbrk{}न्त्र\appbrk{}क\appbrk{}र्म्माणि}\ H}}\skt{ बुद्धिमान्~। }\pend
\pstart\versepar \skt{अस\-ङ्ख्ये\-य\-वि\-क\-ल्प\-त्वा\-च्छ\-ल्या\-ना\-मिति निश्चयः~॥}\pend

\pstart \ledsidenote{\passref{1.7.19}}\edtext{\edtext{\edtext{\edtext{\edtext{\skt{तत्रा\-ति\-स्थू\-ल\-म\-सा\-र\-म\-ति\-दी\-र्घ\-म\-ति\-ह्र\-स्व\-म\-ग्रा\-हि\-वक्रं}}{\lemma{\skt{तत्रातिस्थूलम्}}\Afootnote{K H; \skt{तत्र, अतिस्थूलम्}\ A; \skt{तत्रातिस्थृल}\ N}}}{\lemma{\skt{असारम्}}\Afootnote{A K H; \skt{मसारम्}\ N}}}{\lemma{\skt{अतिदीर्घम्}}\Afootnote{A K N; \skt{अतिदी}\ H}}}{\lemma{\skt{अतिह्रस्वम्}}\Afootnote{A K N; \emph{om.}\ H}}}{\lemma{\skt{अग्राहिवक्रं}}\Afootnote{K; \skt{अग्राहि, विषमग्राहि, वक्रं}\ A; \skt{अग्राहि वक्त्रं}\ N; \skt{र्घ मतिह्र्स्वमग्राहि वक्रं}\ H}}\skt{ }\edtext{\skt{शिथि\-ल\-म\-त्यु\-न्नतं}}{\lemma{\skt{शिथिलम्}}\Afootnote{A K N; \skt{सिथिलम्}\ H}}\skt{ मृदुकीलं }\edtext{\skt{मृदुपाशं}}{\lemma{\skt{मृदुपाशं}}\Afootnote{K; \skt{मृदुमुखं}\ A; \skt{मृदुपोशं}\ N; \skt{मृदुपासं}\ H}}\skt{ }\edtext{\skt{मृदुमुखमिति}}{\lemma{\skt{मृदुमुखम्}}\Afootnote{K N H; \skt{मृदुपाशम्}\ A}}\skt{ }\edtext{\skt{द्वाद\-श\-य\-न्त्र\-दोषाः}}{\lemma{\skt{द्वाद\appbrk{}श\appbrk{}य\appbrk{}न्त्र\appbrk{}दोषाः}}\Afootnote{A K N H\textsuperscript{pc}; \skt{द्वादश यन्त्रन्त्रादोषाः}\ H\textsuperscript{ac}}}\skt{~॥}\pend

\pstart\versepar \ledsidenote{\passref{1.7.20}}\edtext{\edtext{\edtext{\skt{एतै\-र्दो\-षे\-र्वि\-मुक्तं}}{\lemma{\skt{एतैर्}}\Afootnote{A K H; \skt{एतै}\ N}}}{\lemma{\skt{दोषेर्}}\Afootnote{K; \skt{दोषैर्}\ A N H}}\skt{ तु}}{\lemma{\skt{विमुक्तं तु}}\Afootnote{K N H; \skt{विनिर्मुक्तं}\ A}}\skt{ }\edtext{\edtext{\skt{यन्त्र\-म\-ष्टा\-द\-शां\-गुलं}}{\lemma{\skt{यन्त्रम्}}\Afootnote{K H; \skt{यन्त्र}\ A N}}}{\lemma{\skt{अष्टादशांगुलं}}\Afootnote{K H; \skt{मष्टा\appbrk{}द\appbrk{}शा\appbrk{}ङ्गु\appbrk{}लम्}\ A; \skt{मष्टा\appbrk{}द\appbrk{}शा\appbrk{}गु\appbrk{}लम्}\ N}}\skt{~। }\pend
\pstart\versepar \skt{प्रशस्तं भिषजा ज्ञेयं }\edtext{\skt{तद्धि}}{\lemma{\skt{तद्धि}}\Afootnote{A K; \skt{तद्वि}\ N H}}\skt{ कर्मसु }\edtext{\skt{योजयेत्}}{\lemma{\skt{योजयेत्}}\Afootnote{A K H; \skt{योजयोत्}\ N}}\skt{~॥}\pend

\pstart\versepar \ledsidenote{\passref{1.7.21}}\skt{दृश्यं }\edtext{\skt{सिंहमुखाद्यैस्तु}}{\lemma{\skt{सिंहमुखाद्यैस्तु}}\Afootnote{A K N; \skt{सिंहमुखाद्यैस्त्र}\ H}}\skt{ }\edtext{\skt{गूढं}}{\lemma{\skt{गूढं}}\Afootnote{A K H; \skt{गृढं}\ N}}\skt{ }\edtext{\skt{कङ्कमुखादिभिः}}{\lemma{\skt{कङ्कमुखादिभिः}}\Afootnote{A K H; \skt{कद्कमुखादिभिः}\ N}}\skt{~। }\pend
\pstart\versepar \skt{}\edtext{\skt{शल्यं}}{\lemma{\skt{शल्यं}}\Afootnote{K H; \emph{om.}\ A; \skt{शल्य}\ N}}\skt{ }\edtext{\skt{स्वस्तिकयन्त्रैस्तु}}{\lemma{\skt{स्वस्तिकयन्त्रैस्तु}}\Afootnote{K N; \emph{om.}\ A; \skt{स्वस्तिक्यन्त्रै}\ H}}\skt{ }\edtext{\skt{निर्ह\-रे\-त्त\-द्भि\-ष\-क्छनैः}}{\lemma{\skt{निर्ह\appbrk{}रे\appbrk{}त्त\appbrk{}द्भि\appbrk{}ष\appbrk{}क्छनैः}}\Afootnote{K; \skt{निर्हरेत शनैः शल्यं शास्त्र\appbrk{}यु\appbrk{}क्ति\appbrk{}व्य\appbrk{}पे\appbrk{}क्षया}\ A; \skt{निरहरैत्त द्भिषक्छनैः}\ N; \skt{स्त्र निर्हरेत्त भिषक्छनैः}\ H}}\skt{~॥}\pend

\pstart\versepar \ledsidenote{\passref{1.7.22}}\edtext{\skt{विवर्त्तते}}{\lemma{\skt{विवर्त्तते}}\Afootnote{K H; \skt{नि(वि)वर्तते}\ A; \skt{विवर्त्तन}\ N}}\skt{ साध्ववगाहते च }\edtext{\skt{गृह्णाति}}{\lemma{\skt{गृह्णाति}}\Afootnote{K H; \skt{शल्यं नि}\ A; \emph{om.}\ N}}\skt{ }\edtext{\skt{गृह्योद्धरते च}}{\lemma{\skt{गृह्योद्धरते च}}\Afootnote{A K H; \skt{गृहोद्वरते च}\ N}}\skt{ यस्मात्~। }\pend
\pstart\versepar \skt{}\edtext{\skt{तस्मा\-त्स्मृ\-त\-ङ्का\-ङ्क\-मुखं}}{\lemma{\skt{काङ्कमुखं}}\Afootnote{K; \skt{कङ्कमुखं}\ A H; \skt{काङ्कामुखं}\ N}}\skt{ प्रधानं स्थानेषु }\edtext{\edtext{\skt{सर्वेष्वविकारि}}{\lemma{\skt{सर्वेष्व्}}\Afootnote{A K H; \skt{सर्वे}\ N}}}{\lemma{\skt{अविकारि}}\Afootnote{K H; \skt{अधि(वि)कारि}\ A; \skt{ष्वविकारि}\ N}}\skt{ }\edtext{\skt{यच्चेति}}{\lemma{\skt{यच्चेति}}\Afootnote{K N H; \skt{चैव}\ A}}\skt{~॥}\edtext{}{\lemma{\skt{यस्मात्~। तस्मात्स्मृतङ्}}\Afootnote{K N; \skt{यस्मात्~। यन्त्रेष्वतः}\ A; \skt{यस्मात्~। यस्मात्स्मृतं}\ H}}\pend

\pstart \ledsidenote{\passref{1.7.trailer}}\edtext{\skt{॥ ङ्}}{\lemma{\skt{॥ ङ्}}\Afootnote{H; \skt{इति सुश्रु\appbrk{}त\appbrk{}सं\appbrk{}हि\appbrk{}तायां सूत्रस्थाने यन्त्रविधिर्नाम सप्तमोऽध्यायः}\ A; \skt{{\symfont ❈}~॥ ७}\ K}}\skt{~॥}\pend

\endnumbering
\beginnumbering
\eledsection*{Adhyāya 8}
\pstart \ledsidenote{\passref{1.8.1}}\skt{अथातः }\edtext{\edtext{\skt{शस्त्रा\-व\-चा\-र\-णी\-य\-म\-ध्यायं}}{\lemma{\skt{शस्त्रा\appbrk{}व\appbrk{}चा\appbrk{}र\appbrk{}णी\appbrk{}यम्}}\Afootnote{A K H; \skt{शस्त्रा\appbrk{}वि\appbrk{}चा\appbrk{}र\appbrk{}णी\appbrk{}यम्}\ N}}}{\lemma{\skt{अध्यायं}}\Afootnote{A N H; \emph{om.}\ K}}\skt{ }\edtext{\skt{व्याख्यास्यामः}}{\lemma{\skt{व्याख्यास्यामः}}\Afootnote{A N H; \emph{om.}\ K}}\skt{~।}\pend

\pstart \ledsidenote{\passref{1.8.3}}\edtext{\skt{एकविंशतिः}}{\lemma{\skt{एकविंशतिः}}\Afootnote{H; \skt{विंशतिः}\ A; \skt{एकविंशति}\ K N}}\skt{ शस्त्राणि }\edtext{\skt{भवन्ति}}{\lemma{\skt{भवन्ति}}\Afootnote{K N H; \emph{om.}\ A}}\skt{~। }\edtext{\skt{तद्यथा}}{\lemma{\skt{तद्यथा}}\Afootnote{A K H; \skt{तद्यथा}\ N}}\skt{~। }\edtext{\skt{मण्ड\-ला\-ग्रा\-र्ध\-म\-ण्ड\-ला\-ग्र\-क\-र\-प\-त्र\-वृ\-द्धि\-प\-त्र\-न\-ख\-श\-स्त्र\-मु\-द्रि\-को\-त्प\-ल\-प\-त्र\-का\-ध्य\-र्ध\-धा\-र\-सू\-ची\-कु\-श\-प\-त्रा\-टा\-मु\-ख\-श\-रा\-री\-मु\-खा\-न्त\-र्मु\-ख\-त्रि\-कु\-र्च\-क\-कु\-ठा\-रि\-का\-व्री\-हि\-मु\-खा\-रा\-वे\-त\-स\-प\-त्र\-ब\-डि\-श\-द\-न्त\-श\-ङ्क्वे\-षण्य}}{\lemma{\skt{मण्ड\appbrk{}ला\appbrk{}ग्रा\appbrk{}र्ध\appbrk{}म\appbrk{}ण्ड\appbrk{}ला\appbrk{}ग्र\appbrk{}क\appbrk{}र\appbrk{}प\appbrk{}त्र\appbrk{}वृ\appbrk{}द्धि\appbrk{}प\appbrk{}त्र\appbrk{}न\appbrk{}ख\appbrk{}श\appbrk{}स्त्र\appbrk{}मु\appbrk{}द्रि\appbrk{}को\appbrk{}त्प\appbrk{}ल\appbrk{}प\appbrk{}त्र\appbrk{}का\appbrk{}ध्य\appbrk{}र्ध\appbrk{}धा\appbrk{}र\appbrk{}सू\appbrk{}ची\appbrk{}कु\appbrk{}श\appbrk{}प\appbrk{}त्रा\appbrk{}टा\appbrk{}मु\appbrk{}ख\appbrk{}श\appbrk{}रा\appbrk{}री\appbrk{}मु\appbrk{}खा\appbrk{}न्त\appbrk{}र्मु\appbrk{}ख\appbrk{}त्रि\appbrk{}कु\appbrk{}र्च\appbrk{}क\appbrk{}कु\appbrk{}ठा\appbrk{}रि\appbrk{}का\appbrk{}व्री\appbrk{}हि\appbrk{}मु\appbrk{}खा\appbrk{}रा\appbrk{}वे\appbrk{}त\appbrk{}स\appbrk{}प\appbrk{}त्र\appbrk{}ब\appbrk{}डि\appbrk{}श\appbrk{}द\appbrk{}न्त\appbrk{}श\appbrk{}ङ्क्वे\appbrk{}षण्य}}\Afootnote{\skt{मण्ड\appbrk{}ला\appbrk{}ग्र\appbrk{}क\appbrk{}र\appbrk{}प\appbrk{}त्र\appbrk{}वृ\appbrk{}द्धि\appbrk{}प\appbrk{}त्र\appbrk{}न\appbrk{}ख\appbrk{}श\appbrk{}स्त्र\appbrk{}मु\appbrk{}द्रि\appbrk{}को\appbrk{}त्प\appbrk{}ल\appbrk{}प\appbrk{}त्र\appbrk{}का\appbrk{}र्ध\appbrk{}धा\appbrk{}र\appbrk{}सू\appbrk{}ची\appbrk{}कु\appbrk{}श\appbrk{}प\appbrk{}त्रा\appbrk{}टी\appbrk{}मुख- शरा\appbrk{}रि\appbrk{}मु\appbrk{}खा\appbrk{}न्त\appbrk{}र्मु\appbrk{}ख\appbrk{}त्रि\appbrk{}कू\appbrk{}र्च\appbrk{}क\appbrk{}कु\appbrk{}ठा\appbrk{}रि\appbrk{}का\appbrk{}व्री\appbrk{}हि\appbrk{}मु\appbrk{}खा\appbrk{}रा\appbrk{}वे\appbrk{}त\appbrk{}स\appbrk{}प\appbrk{}त्र\appbrk{}क\appbrk{}ब\appbrk{}डि\appbrk{}श\appbrk{}द\appbrk{}न्त\appbrk{}श\appbrk{}ङ्क्वे\appbrk{}षण्य}\ A; \skt{मण्ड\appbrk{}ला\appbrk{}ग्रा\appbrk{}र्ध\appbrk{}म\appbrk{}ण्ड\appbrk{}ला\appbrk{}ग्र\appbrk{}क\appbrk{}र\appbrk{}प\appbrk{}त्र\appbrk{}वृ\appbrk{}द्धि\appbrk{}प\appbrk{}त्र\appbrk{}न\appbrk{}ख\appbrk{}श\appbrk{}स्त्र\appbrk{}मु\appbrk{}द्रि\appbrk{}को\appbrk{}त्प\appbrk{}ल\appbrk{}प\appbrk{}त्र\appbrk{}का\appbrk{}ध्य\appbrk{}र्ध\appbrk{}धा\appbrk{}र\appbrk{}सू\appbrk{}ची\appbrk{}कु\appbrk{}श\appbrk{}प\appbrk{}त्रा\appbrk{}टा\appbrk{}मु\appbrk{}ख\appbrk{}श\appbrk{}रा\appbrk{}री\appbrk{}मु\appbrk{}खा\appbrk{}न्त\appbrk{}र्मु\appbrk{}ख\appbrk{}त्रि\appbrk{}कु\appbrk{}र्च\appbrk{}क\appbrk{}कु\appbrk{}ठा\appbrk{}रि\appbrk{}का\appbrk{}व्री\appbrk{}हि\appbrk{}मु\appbrk{}खा\appbrk{}रा\appbrk{}वे\appbrk{}त्र\appbrk{}स\appbrk{}प\appbrk{}त्र\appbrk{}ब\appbrk{}डि\appbrk{}श\appbrk{}द\appbrk{}न्त\appbrk{}श\appbrk{}ङ्क्वे\appbrk{}षण्य}\ K N; \skt{मण्ड\appbrk{}ला\appbrk{}ग्रा\appbrk{}र्द्ध\appbrk{}म\appbrk{}ण्ड\appbrk{}ल्ग्र।\appbrk{}क\appbrk{}र\appbrk{}प\appbrk{}त्र\appbrk{}वृ\appbrk{}द्धि\appbrk{}प\appbrk{}त्र\appbrk{}न\appbrk{}ख\appbrk{}श\appbrk{}स्त्र\appbrk{}मु\appbrk{}द्रि\appbrk{}को\appbrk{}त्प\appbrk{}ल\appbrk{}प\appbrk{}त्र\appbrk{}का\appbrk{}ध्य\appbrk{}र्द्ध\appbrk{}धा\appbrk{}रा\appbrk{}सू\appbrk{}ची\appbrk{}कु\appbrk{}श\appbrk{}प\appbrk{}त्रा\appbrk{}टा\appbrk{}मु\appbrk{}ख\appbrk{}श\appbrk{}रा\appbrk{}री\appbrk{}मु\appbrk{}खा\appbrk{}न्त\appbrk{}र्म्मु\appbrk{}ख\appbrk{}त्रि\appbrk{}कु\appbrk{}र्च्च\appbrk{}क\appbrk{}कु\appbrk{}ठा\appbrk{}रि\appbrk{}का।\appbrk{}व्री\appbrk{}हि\appbrk{}मु\appbrk{}खा\appbrk{}रा।\appbrk{}वे\appbrk{}त\appbrk{}स\appbrk{}प\appbrk{}त्र\appbrk{}ब\appbrk{}डि\appbrk{}श\appbrk{}द\appbrk{}न्त\appbrk{}शं\appbrk{}क्वे\appbrk{}षण्य}\ H\textsuperscript{ac}; \skt{मण्ड\appbrk{}ला\appbrk{}ग्रा\appbrk{}र्द्ध\appbrk{}म\appbrk{}ण्ड\appbrk{}ला\appbrk{}ग्र।\appbrk{}क\appbrk{}र\appbrk{}प\appbrk{}त्र\appbrk{}वृ\appbrk{}द्धि\appbrk{}प\appbrk{}त्र\appbrk{}न\appbrk{}ख\appbrk{}श\appbrk{}स्त्र\appbrk{}मु\appbrk{}द्रि\appbrk{}को\appbrk{}त्प\appbrk{}ल\appbrk{}प\appbrk{}त्र\appbrk{}का\appbrk{}ध्य\appbrk{}र्द्ध\appbrk{}धा\appbrk{}रा\appbrk{}सू\appbrk{}ची\appbrk{}कु\appbrk{}श\appbrk{}प\appbrk{}त्रा\appbrk{}टा\appbrk{}मु\appbrk{}ख\appbrk{}श\appbrk{}रा\appbrk{}री\appbrk{}मु\appbrk{}खा\appbrk{}न्त\appbrk{}र्म्मु\appbrk{}ख\appbrk{}त्रि\appbrk{}कु\appbrk{}र्च्च\appbrk{}क\appbrk{}कु\appbrk{}ठा\appbrk{}रि\appbrk{}का।\appbrk{}व्री\appbrk{}हि\appbrk{}मु\appbrk{}खा\appbrk{}रा।\appbrk{}वे\appbrk{}त\appbrk{}स\appbrk{}प\appbrk{}त्र\appbrk{}ब\appbrk{}डि\appbrk{}श\appbrk{}द\appbrk{}न्त\appbrk{}शं\appbrk{}क्वे\appbrk{}षण्य}\ H\textsuperscript{pc}}}\skt{ इति~।}\pend

\pstart \ledsidenote{\passref{1.8.4}}\skt{तत्र }\edtext{\edtext{\skt{मण्ड\-ला\-ग्र\-म\-र्ध\-म\-ण्ड\-ला\-ग्र\-क\-र\-प\-त्राणि}}{\lemma{\skt{मण्डलाग्रम्}}\Afootnote{K N H; \emph{om.}\ A}}}{\lemma{\skt{अर्ध\appbrk{}म\appbrk{}ण्ड\appbrk{}ला\appbrk{}ग्र\appbrk{}क\appbrk{}र\appbrk{}प\appbrk{}त्राणि}}\Afootnote{K N H\textsuperscript{pc}; \skt{मण्ड\appbrk{}ला\appbrk{}ग्र\appbrk{}क\appbrk{}र\appbrk{}पत्रे स्यातां}\ A; \skt{अर्द्ध\appbrk{}म\appbrk{}ण्ड\appbrk{}ला\appbrk{}ग्र\appbrk{}क\appbrk{}र\appbrk{}त्राणि}\ H\textsuperscript{ac}}}\skt{ च्छेदने लेखने }\edtext{\skt{चोपदिश्यन्ते}}{\lemma{\skt{चोपदिश्यन्ते}}\Afootnote{K N H; \skt{च}\ A}}\skt{~। }\edtext{\skt{वृद्धि\-प\-त्र\-न\-ख\-श\-स्त्र\-मु\-द्रि\-को\-त्प\-ल\-प\-त्र\-का\-ध्य\-र्ध\-धा\-राणि}}{\lemma{\skt{वृद्धि\appbrk{}प\appbrk{}त्र\appbrk{}न\appbrk{}ख\appbrk{}श\appbrk{}स्त्र\appbrk{}मु\appbrk{}द्रि\appbrk{}को\appbrk{}त्प\appbrk{}ल\appbrk{}प\appbrk{}त्र\appbrk{}का\appbrk{}ध्य\appbrk{}र्ध\appbrk{}धा\appbrk{}राणि}}\Afootnote{K H; \skt{वृद्धि\appbrk{}प\appbrk{}त्र\appbrk{}न\appbrk{}ख\appbrk{}श\appbrk{}स्त्र\appbrk{}मु\appbrk{}द्रि\appbrk{}को\appbrk{}त्प\appbrk{}ल\appbrk{}प\appbrk{}त्र\appbrk{}का\appbrk{}र्ध\appbrk{}धा\appbrk{}राणि}\ A; \skt{वृद्धि\appbrk{}प\appbrk{}त्र।\appbrk{}न\appbrk{}ख\appbrk{}श\appbrk{}स्त्र।\appbrk{}मु\appbrk{}दि\appbrk{}को\appbrk{}त्प\appbrk{}ल\appbrk{}प\appbrk{}त्रा\appbrk{}ध्य\appbrk{}र्ध\appbrk{}ण\appbrk{}राणि}\ N\textsuperscript{ac}; \skt{वृद्धि\appbrk{}प\appbrk{}त्र।\appbrk{}न\appbrk{}ख\appbrk{}श\appbrk{}स्त्र।\appbrk{}मु\appbrk{}दि\appbrk{}को\appbrk{}त्प\appbrk{}ल\appbrk{}प\appbrk{}त्रा\appbrk{}ध्य\appbrk{}र्ध\appbrk{}धा\appbrk{}राणि}\ N\textsuperscript{pc}}}\skt{ }\edtext{\skt{भेदने}}{\lemma{\skt{भेदने}}\Afootnote{K N H; \skt{छेदने}\ A}}\skt{ }\edtext{\skt{च्छेदने}}{\lemma{\skt{च्छेदने}}\Afootnote{K N H; \skt{भेदने}\ A}}\skt{ }\edtext{\skt{चोपयुज्यन्ते}}{\lemma{\skt{चोपयुज्यन्ते}}\Afootnote{K N H; \skt{च}\ A}}\skt{~। }\edtext{\skt{सूची\-कु\-ठा\-रि\-का\-व्री\-हि\-मु\-खा\-रा\-वे\-त\-स\-प\-त्राणि}}{\lemma{\skt{सूची\appbrk{}कु\appbrk{}ठा\appbrk{}रि\appbrk{}का\appbrk{}व्री\appbrk{}हि\appbrk{}मु\appbrk{}खा\appbrk{}रा\appbrk{}वे\appbrk{}त\appbrk{}स\appbrk{}प\appbrk{}त्राणि}}\Afootnote{K H; \skt{सूची\appbrk{}कु\appbrk{}श\appbrk{}प\appbrk{}त्रा\appbrk{}टी(\appbrk{}टा)\appbrk{}मु\appbrk{}ख\appbrk{}श\appbrk{}रा\appbrk{}रि\appbrk{}मु\appbrk{}खा\appbrk{}न्त\appbrk{}र्मु\appbrk{}ख\appbrk{}त्रि\appbrk{}कू\appbrk{}र्च\appbrk{}कानि विस्रावणे, कुठा\appbrk{}रि\appbrk{}का\appbrk{}व्री\appbrk{}हि\appbrk{}मु\appbrk{}खा\appbrk{}रा\appbrk{}वे\appbrk{}त\appbrk{}स\appbrk{}प\appbrk{}त्र\appbrk{}काणि}\ A; \skt{सूची\appbrk{}कु\appbrk{}ठा\appbrk{}रि\appbrk{}का\appbrk{}व्री\appbrk{}हि\appbrk{}मु\appbrk{}खा\appbrk{}रा\appbrk{}वे\appbrk{}त\appbrk{}स\appbrk{}प्र\appbrk{}त्राणि}\ N}}\skt{ }\edtext{\skt{वेधने}}{\lemma{\skt{वेधने}}\Afootnote{K N H; \skt{व्यधने}\ A}}\skt{~। }\edtext{\edtext{\skt{एषण्येषणे}}{\lemma{\skt{एषण्य्}}\Afootnote{K N H; \emph{om.}\ A}}}{\lemma{\skt{एषणे}}\Afootnote{K N H; \emph{om.}\ A}}\skt{~। }\edtext{\skt{अनुलोमाः}}{\lemma{\skt{अनुलोमाः}}\Afootnote{K N H; \emph{om.}\ A}}\skt{ }\edtext{\skt{करीराः}}{\lemma{\skt{करीराः}}\Afootnote{K N H; \emph{om.}\ A}}\skt{ }\edtext{\skt{शस्त्रवृत्ताश्च}}{\lemma{\skt{शस्त्रवृत्ताश्च}}\Afootnote{K N H; \emph{om.}\ A}}\skt{~। }\edtext{\skt{सूची}}{\lemma{\skt{सूची}}\Afootnote{K N H; \skt{सूची च}\ A}}\skt{ }\edtext{\skt{बडि\-श\-द\-न्त\-श\-ङ्कु\-श्चा\-ह\-रणे}}{\lemma{\skt{बडि\appbrk{}श\appbrk{}द\appbrk{}न्त\appbrk{}श\appbrk{}ङ्कुश्}}\Afootnote{K N; \skt{बडिशं दन्तशङ्कुश्}\ A; \skt{बडि\appbrk{}श\appbrk{}द\appbrk{}न्त\appbrk{}श\appbrk{}कुश्}\ H}}\skt{~। }\edtext{\skt{कुश\-प\-त्रा\-टा\-मु\-ख\-श\-रा\-रि\-मु\-खा\-न्त\-र्मु\-ख\-त्रि\-कु\-र्च\-कानि}}{\lemma{\skt{कुश\appbrk{}प\appbrk{}त्रा\appbrk{}टा\appbrk{}मु\appbrk{}ख\appbrk{}श\appbrk{}रा\appbrk{}रि\appbrk{}मु\appbrk{}खा\appbrk{}न्त\appbrk{}र्मु\appbrk{}ख\appbrk{}त्रि\appbrk{}कु\appbrk{}र्च\appbrk{}कानि}}\Afootnote{K N\textsuperscript{pc} H; \skt{एष}\ A; \skt{कुश\appbrk{}प\appbrk{}त्रा\appbrk{}टा\appbrk{}ष्टा\appbrk{}मु\appbrk{}ख\appbrk{}श\appbrk{}रा\appbrk{}रि\appbrk{}मु\appbrk{}खा\appbrk{}न्त\appbrk{}र्मु\appbrk{}ख\appbrk{}त्रि\appbrk{}कु\appbrk{}र्च\appbrk{}कानि}\ N\textsuperscript{ac}}}\skt{ }\edtext{\skt{विस्रावणे}}{\lemma{\skt{विस्रावणे}}\Afootnote{K N; \skt{ण्येषणे}\ A; \skt{विश्रावणे}\ H}}\skt{~। }\edtext{\skt{सूची}}{\lemma{\skt{सूची}}\Afootnote{K N H; \skt{आनुलोम्ये च, सूच्यः}\ A}}\skt{ }\edtext{\skt{सीव्यन}}{\lemma{\skt{सीव्यन}}\Afootnote{K N H; \skt{सीवने}\ A}}\skt{ }\edtext{\skt{इत्यष्टविधः}}{\lemma{\skt{अष्टविधः}}\Afootnote{K H; \skt{अष्टविधे कर्मण्युपयोगः}\ A; \skt{अविधः}\ N}}\skt{ शस्त्राणां }\edtext{\edtext{\skt{कर्मण्युपयोगो}}{\lemma{\skt{कर्मण्य्}}\Afootnote{K N; \emph{om.}\ A; \skt{कर्म्माण्य्}\ H}}}{\lemma{\skt{उपयोगो}}\Afootnote{K N H; \emph{om.}\ A}}\skt{ व्याख्यातः~।}\pend

\pstart \ledsidenote{\passref{1.8.5}}\edtext{\skt{तेषां}}{\lemma{\skt{तेषां}}\Afootnote{K N H; \skt{तेषां अथ}\ A}}\skt{ यथायोगं }\edtext{\skt{ग्रहणम्}}{\lemma{\skt{ग्रहणम्}}\Afootnote{K N H; \skt{ग्रहण}\ A}}\skt{~। }\edtext{\edtext{\skt{कर्मस्वेष}}{\lemma{\skt{कर्मस्व्}}\Afootnote{K N H; \skt{समासोपायः कर्मसु}\ A}}}{\lemma{\skt{एष}}\Afootnote{K N H; \emph{om.}\ A}}\skt{ }\edtext{\skt{शस्त्र\-ग्र\-ह\-ण\-स\-मासः}}{\lemma{\skt{शस्त्र\appbrk{}ग्र\appbrk{}ह\appbrk{}ण\appbrk{}स\appbrk{}मासः}}\Afootnote{K N H; \skt{वक्ष्यते- तत्र}\ A}}\skt{~। }\edtext{\skt{वृद्धिपत्रन्तु}}{\lemma{\skt{वृद्धिपत्रन्तु}}\Afootnote{K N H; \skt{वृद्धिपत्रं}\ A}}\skt{ }\edtext{\skt{वृन्त\-फ\-ल\-सा\-धा\-रणे}}{\lemma{\skt{वृन्त\appbrk{}फ\appbrk{}ल\appbrk{}सा\appbrk{}धा\appbrk{}रणे}}\Afootnote{A K; \skt{वृन्त\appbrk{}फ\appbrk{}ल\appbrk{}सा\appbrk{}रणे}\ N; \skt{वृत्त\appbrk{}फ\appbrk{}ल\appbrk{}सा\appbrk{}धा\appbrk{}रणे}\ H}}\skt{ भागे }\edtext{\skt{गृह्णीयात्~। }\edtext{\skt{भेदनान्येवं}}{\lemma{\skt{एवं}}\Afootnote{K H; \skt{न्येवं}\ A; \skt{एव}\ N}}}{\lemma{\skt{गृह्णीयात्~। भेदनान्य्}}\Afootnote{K N H; \skt{गृह्णीयात्, भेदना}\ A}}\skt{ सर्वाणि~। }\edtext{\skt{वृद्धि\-प\-त्र\-व\-द\-र्ध\-म\-ण्ड\-लाग्रं}}{\lemma{\skt{वृद्धि\appbrk{}प\appbrk{}त्र\appbrk{}व\appbrk{}द\appbrk{}र्ध\appbrk{}म\appbrk{}ण्ड\appbrk{}लाग्रं}}\Afootnote{K H; \skt{वृद्धिपत्रं मण्डलाग्रं च}\ A; \skt{वृद्धि\appbrk{}प\appbrk{}प\appbrk{}त्र\appbrk{}व\appbrk{}द\appbrk{}र्ध\appbrk{}म\appbrk{}ण्ड\appbrk{}लाग्रं}\ N}}\skt{ }\edtext{\skt{किञ्चि\-दु\-त्ता\-न\-पा\-णिना}}{\lemma{\skt{किञ्चि\appbrk{}दु\appbrk{}त्ता\appbrk{}न\appbrk{}पा\appbrk{}णिना}}\Afootnote{K H; \skt{किञ्चिदुत्तानेन पाणिना}\ A; \skt{किंचि\appbrk{}द\appbrk{}त्ता\appbrk{}न\appbrk{}पा\appbrk{}णिना}\ N}}\skt{ लेखने }\edtext{\skt{बहुशोऽवचार्य}}{\lemma{\skt{ऽवचार्य}}\Afootnote{\skt{ऽवचार्यं}\ A; \skt{वचार्य}\ K N H}}\skt{ }\edtext{\skt{वृन्ताग्रेण}}{\lemma{\skt{वृन्ताग्रेण}}\Afootnote{K N; \skt{वृन्ताग्रे}\ A; \skt{वृत्ताग्रेण}\ H}}\skt{ }\edtext{\skt{विस्रावणानि}}{\lemma{\skt{विस्रावणानि}}\Afootnote{A K N; \skt{विश्रावणानि}\ H}}\skt{~। विशेषेण तु }\edtext{\skt{बाल\-वृ\-द्ध\-सु\-कु\-मा\-र\-भी\-रु\-ना\-रीणां}}{\lemma{\skt{बाल\appbrk{}वृ\appbrk{}द्ध\appbrk{}सु\appbrk{}कु\appbrk{}मा\appbrk{}र\appbrk{}भी\appbrk{}रु\appbrk{}ना\appbrk{}रीणां}}\Afootnote{A K\textsuperscript{pc} N H; \skt{बाल\appbrk{}वृ\appbrk{}द्ध\appbrk{}सु\appbrk{}कु\appbrk{}मा\appbrk{}र\appbrk{}भी\appbrk{}रु\appbrk{}नारी}\ K\textsuperscript{ac}}}\skt{ राज्ञां }\edtext{\skt{राजमात्राणां वा}}{\lemma{\skt{राजमात्राणां वा}}\Afootnote{K N H; \skt{राजमा(पु)त्राणां च}\ A}}\skt{ }\edtext{\skt{त्रिकुर्चकेन}}{\lemma{\skt{त्रिकुर्चकेन}}\Afootnote{K N H\textsuperscript{pc}; \skt{त्रिकूर्चकेन}\ A; \skt{त्रिकुर्च्चके}\ H\textsuperscript{ac}}}\skt{ }\edtext{\skt{विस्रावयेत्~। }\edtext{\skt{तल\-प्र\-च्छा\-दि\-त\-वृ\-न्ता\-ग्र\-म\-ङ्गु\-ष्ठ\-प्र\-दे\-शि\-नीभ्यां}}{\lemma{\skt{अङ्गु\appbrk{}ष्ठ\appbrk{}प्र\appbrk{}दे\appbrk{}शि\appbrk{}नीभ्यां}}\Afootnote{K N H; \skt{मङ्गु\appbrk{}ष्ठ\appbrk{}प्र\appbrk{}दे\appbrk{}शि\appbrk{}नीभ्यां}\ A}}}{\lemma{\skt{विस्रावयेत्~। तल\appbrk{}प्र\appbrk{}च्छा\appbrk{}दि\appbrk{}त\appbrk{}वृ\appbrk{}न्ता\appbrk{}ग्रम्}}\Afootnote{K N; \skt{विस्रावयेत्, तल\appbrk{}प्र\appbrk{}च्छा\appbrk{}दि\appbrk{}त\appbrk{}वृन्त}\ A; \skt{विश्रावयेत्~। तल\appbrk{}प्र\appbrk{}च्छा\appbrk{}दि\appbrk{}त\appbrk{}वृ\appbrk{}त्ता\appbrk{}ग्रम्}\ H}}\skt{ व्रीहिमुखम्~। कुठारिकां }\edtext{\skt{वाम\-ह\-स्त\-गृ\-ही\-त\-पुच्छां}}{\lemma{\skt{वाम\appbrk{}ह\appbrk{}स्त\appbrk{}गृ\appbrk{}ही\appbrk{}त\appbrk{}पुच्छां}}\Afootnote{K N H; \skt{वाम\appbrk{}ह\appbrk{}स्त\appbrk{}न्य\appbrk{}स्ता\appbrk{}मित}\ A}}\skt{ }\edtext{\skt{दक्षि\-ण\-ह\-स्ता\-ङ्गु\-ष्ठा\-व\-ष्ट\-ब्धया}}{\lemma{\skt{दक्षि\appbrk{}ण\appbrk{}ह\appbrk{}स्ता\appbrk{}ङ्गु\appbrk{}ष्ठा\appbrk{}व\appbrk{}ष्ट\appbrk{}ब्धया}}\Afootnote{K\textsuperscript{pc} H; \skt{रह\appbrk{}स्त\appbrk{}म\appbrk{}ध्य\appbrk{}मा\appbrk{}ङ्गु\appbrk{}ल्याऽ\appbrk{}ङ्गु\appbrk{}ष्ठ\appbrk{}वि\appbrk{}ष्ट\appbrk{}ब्धयाऽ}\ A; \skt{दक्षि\appbrk{}ण\appbrk{}ह\appbrk{}स्तां\appbrk{}गुष्ठा}\ K\textsuperscript{ac}; \skt{दक्षि\appbrk{}ण\appbrk{}ह\appbrk{}स्तां\appbrk{}गु\appbrk{}ष्ठा\appbrk{}व\appbrk{}स्क\appbrk{}न्धया}\ N}}\skt{ }\edtext{\skt{मध्यमयाङ्गुल्या}}{\lemma{\skt{मध्यमयाङ्गुल्या}}\Afootnote{K N H; \emph{om.}\ A}}\skt{ }\edtext{\skt{निहन्यात्~। तत्र}}{\lemma{\skt{निहन्यात्~। तत्र}}\Afootnote{K N H; \skt{भिहन्यात् , आरा}\ A}}\skt{ }\edtext{\skt{कर\-प\-त्रा\-रा\-वे\-त\-स\-प\-त्र\-ब\-डि\-श\-द\-न्त\-श\-ङ्क्वे\-ष\-णी\-र्मूले}}{\lemma{\skt{कर\appbrk{}प\appbrk{}त्रा\appbrk{}रा\appbrk{}वे\appbrk{}त\appbrk{}स\appbrk{}प\appbrk{}त्र\appbrk{}ब\appbrk{}डि\appbrk{}श\appbrk{}द\appbrk{}न्त\appbrk{}श\appbrk{}ङ्क्वे\appbrk{}ष\appbrk{}णीर्}}\Afootnote{H; \skt{करपत्रैषण्यो}\ A; \skt{कर\appbrk{}प\appbrk{}त्रा\appbrk{}रा\appbrk{}वे\appbrk{}त\appbrk{}स\appbrk{}प\appbrk{}त्र\appbrk{}ब\appbrk{}डि\appbrk{}श\appbrk{}द\appbrk{}न्त\appbrk{}श\appbrk{}क्वे\appbrk{}ष\appbrk{}णीर्}\ K; \skt{कर\appbrk{}प\appbrk{}त्रा\appbrk{}रा\appbrk{}वे\appbrk{}त\appbrk{}स\appbrk{}प\appbrk{}त्र\appbrk{}ब\appbrk{}डि\appbrk{}ष\appbrk{}द\appbrk{}न्त\appbrk{}स\appbrk{}क्वे\appbrk{}ष\appbrk{}णीर्}\ N}}\skt{ }\edtext{\skt{प्रदे\-शि\-नी\-प्र\-यु\-क्तम्}}{\lemma{\skt{प्रदे\appbrk{}शि\appbrk{}नी\appbrk{}प्र\appbrk{}यु\appbrk{}क्तम्}}\Afootnote{K N H; \skt{शेषाणि तु यथायोगं गृह्णीयात्}\ A}}\skt{~।}\pend

\pstart \ledsidenote{\passref{1.8.6}}\edtext{\skt{मुद्रिकासदृशं}}{\lemma{\skt{मुद्रिकासदृशं}}\Afootnote{K N H; \emph{om.}\ A}}\skt{ }\edtext{\edtext{\skt{नखा\-का\-र\-श\-स्त्र\-मु\-ख\-ञ्च\-तु\-र्वे\-धनं}}{\lemma{\skt{नखा\appbrk{}का\appbrk{}र\appbrk{}श\appbrk{}स्त्र\appbrk{}मु\appbrk{}खञ्}}\Afootnote{K N H; \emph{om.}\ A}}}{\lemma{\skt{चतुर्वेधनं}}\Afootnote{K N H; \emph{om.}\ A}}\skt{ }\edtext{\skt{सूक्ष्म\-डो\-रा\-व\-बद्धं}}{\lemma{\skt{सूक्ष्म\appbrk{}डो\appbrk{}रा\appbrk{}व\appbrk{}बद्धं}}\Afootnote{K N; \emph{om.}\ A; \skt{सूक्ष्म\appbrk{}तो\appbrk{}रा\appbrk{}व\appbrk{}बद्धं}\ H}}\skt{ }\edtext{\skt{मुद्रिकाशस्त्रं}}{\lemma{\skt{मुद्रिकाशस्त्रं}}\Afootnote{K N H; \emph{om.}\ A}}\skt{ }\edtext{\skt{तेषा\-न्ना\-म\-भि\-रे\-वा\-कृ\-तयः}}{\lemma{\skt{तेषान्}}\Afootnote{A K N; \skt{तेषान्तेषान्}\ H}}\skt{ प्रायेण व्याख्याताः~।}\pend

\pstart \ledsidenote{\passref{1.8.7-8}}\skt{तत्र }\edtext{\edtext{\skt{नख\-व\-र्ध\-नै\-ष\-ण्या\-व\-ष्टा\-ङ्गुल्यौ}}{\lemma{\skt{नख\appbrk{}व\appbrk{}र्ध\appbrk{}नै\appbrk{}ष\appbrk{}ण्याव्}}\Afootnote{K N H; \skt{नखशस्त्रैषण्या}\ A}}}{\lemma{\skt{अष्टाङ्गुल्यौ}}\Afootnote{K N H; \skt{वष्टाङ्गुले}\ A}}\skt{ }\edtext{\skt{सूच्यो}}{\lemma{\skt{सूच्यो}}\Afootnote{A K H; \skt{सूच्ये}\ N}}\skt{ }\edtext{\skt{वक्ष्यन्ते}}{\lemma{\skt{वक्ष्यन्ते}}\Afootnote{K N H; \skt{वक्ष्यन्ते, (प्र\appbrk{}दे\appbrk{}शि\appbrk{}न्य\appbrk{}ग्र\appbrk{}प\appbrk{}र्व\appbrk{}प्र\appbrk{}दे\appbrk{}श\appbrk{}प्र\appbrk{}माणा मुद्रिका, दशाङ्गुला शरारिमुखी सा च (या सा) कर्तरीति कथ्यते)}\ A}}\skt{~। शेषाणि तु षडङ्गुलानि तानि सुग्रहाणि }\edtext{\skt{सुधाराणि}}{\lemma{\skt{सुधाराणि}}\Afootnote{K N H; \skt{सुलोहानि, सुधाराणि}\ A}}\skt{ सुरूपाणि }\edtext{\skt{सुलोहानि}}{\lemma{\skt{सुलोहानि}}\Afootnote{K\textsuperscript{pc} N H; \skt{सु}\ A; \emph{om.}\ K\textsuperscript{ac}}}\skt{ }\edtext{\skt{सुधौतानि}}{\lemma{\skt{सुधौतानि}}\Afootnote{K N H; \emph{om.}\ A}}\skt{ }\edtext{\skt{समा\-ञ्चि\-त\-मु\-खानि}}{\lemma{\skt{समा\appbrk{}ञ्चि\appbrk{}त\appbrk{}मु\appbrk{}खानि}}\Afootnote{N H; \skt{समा\appbrk{}हि\appbrk{}त\appbrk{}मु\appbrk{}खा\appbrk{}ग्राणि, अकरालानि}\ A; \skt{समां\appbrk{}चि\appbrk{}त\appbrk{}मु\appbrk{}खा\appbrk{}निणि}\ K}}\skt{ चेति शस्त्रसम्पत्~।}\pend

\pstart \ledsidenote{\passref{1.8.10}}\skt{तत्र धारा भेदनानां }\edtext{\skt{मासूरी}}{\lemma{\skt{मासूरी}}\Afootnote{A K N; \skt{मासुरी}\ H}}\skt{~। }\edtext{\skt{लेख\-ना\-ना\-म\-र्ध\-मा\-सूरी}}{\lemma{\skt{अर्धमासूरी}}\Afootnote{A K N; \skt{अर्द्धमा सूरी}\ H\textsuperscript{ac}; \skt{अर्द्धम सूरी}\ H\textsuperscript{pc}}}\skt{~। }\edtext{\skt{वेध्यानां}}{\lemma{\skt{वेध्यानां}}\Afootnote{K N H; \skt{व्यधनानां}\ A}}\skt{ }\edtext{\skt{विस्रावणानाञ्च}}{\lemma{\skt{विस्रावणानाञ्च}}\Afootnote{A K N; \skt{विश्रावणाञ्च}\ H}}\skt{ }\edtext{\skt{कैशिकी}}{\lemma{\skt{कैशिकी}}\Afootnote{A K H; \skt{केशिकी}\ N}}\skt{~। }\edtext{\skt{छेद\-ना\-ना\-म\-र्ध\-कै\-शिकी}}{\lemma{\skt{अर्धकैशिकी}}\Afootnote{K H; \skt{अर्धकैशिकीति}\ A; \skt{अर्धकेशिकी}\ N}}\skt{~।}\pend

\pstart \ledsidenote{\passref{1.8.11}}\edtext{\skt{बडि\-श\-द\-न्त\-शङ्कु}}{\lemma{\skt{बडि\appbrk{}श\appbrk{}द\appbrk{}न्त\appbrk{}शङ्कु}}\Afootnote{K N H; \skt{बडिशं दन्तशङ्कुश्}\ A}}\skt{ }\edtext{\skt{चानताग्रे}}{\lemma{\skt{चानताग्रे}}\Afootnote{A K\textsuperscript{pc} N H; \skt{नताग्रे~।}\ K\textsuperscript{ac}}}\skt{~। }\edtext{\skt{तीक्ष्ण\-क\-ण्ट\-क\-प्र\-थ\-म\-य\-व\-प\-त्र\-मुखी}}{\lemma{\skt{तीक्ष्ण\appbrk{}क\appbrk{}ण्ट\appbrk{}क\appbrk{}प्र\appbrk{}थ\appbrk{}म\appbrk{}य\appbrk{}व\appbrk{}प\appbrk{}त्र\appbrk{}मुखी}}\Afootnote{K H; \skt{तीक्ष्ण\appbrk{}क\appbrk{}ण्ट\appbrk{}क\appbrk{}प्र\appbrk{}थ\appbrk{}म\appbrk{}य\appbrk{}व\appbrk{}प\appbrk{}त्रमु}\ A; \skt{तीक्ष्ण\appbrk{}क\appbrk{}ण्ट\appbrk{}क\appbrk{}प्र\appbrk{}थ\appbrk{}म\appbrk{}य\appbrk{}व\appbrk{}प्र\appbrk{}त्र\appbrk{}मुखी}\ N}}\skt{ }\edtext{\skt{यवपत्रा}}{\lemma{\skt{यवपत्रा}}\Afootnote{K H; \emph{om.}\ A; \skt{यवप्रत्रा}\ N}}\skt{ }\edtext{\skt{एषणी}}{\lemma{\skt{एषणी}}\Afootnote{K N H; \skt{ख्येषणी}\ A}}\skt{ गण्डू\-प\-दा\-का\-र\-मुखी }\edtext{\skt{चेति}}{\lemma{\skt{चेति}}\Afootnote{K N H; \skt{च)}\ A}}\skt{~।}\pend

\pstart \ledsidenote{\passref{1.8.9}}\skt{तत्र वक्रं }\edtext{\skt{कुण्ठं}}{\lemma{\skt{कुण्ठं}}\Afootnote{A K H; \skt{कुण्ड}\ N}}\skt{ }\edtext{\skt{खण्डं}}{\lemma{\skt{खण्डं}}\Afootnote{A K\textsuperscript{pc}; \skt{खण्ड}\ K\textsuperscript{ac} H; \skt{खंण्ड}\ N}}\skt{ }\edtext{\edtext{\skt{खर\-धा\-रा\-ति\-स्थू\-ल\-म\-त्य\-ल्प\-म\-ति\-दी\-र्घ\-म\-ति\-ह्र\-स्व\-मि\-त्यष्टौ}}{\lemma{\skt{खर\appbrk{}धा\appbrk{}रा\appbrk{}ति\appbrk{}स्थू\appbrk{}लम्}}\Afootnote{K N H; \skt{खरधारम्, अतिस्थूलम्}\ A}}}{\lemma{\skt{अतिह्रस्वम्}}\Afootnote{A K H; \skt{अतिह्रष्वम्}\ N}}\skt{ }\edtext{\skt{शस्त्रदोषाः}}{\lemma{\skt{शस्त्रदोषाः}}\Afootnote{A K N\textsuperscript{pc} H; \skt{शस्त्रदोषाषाः}\ N\textsuperscript{ac}}}\skt{~। अतो }\edtext{\skt{विप\-री\-त\-गु\-ण\-मा\-द\-द्या\-द\-न्यत्र}}{\lemma{\skt{आदद्यादन्यत्र}}\Afootnote{K N H; \skt{आददीत, अन्यत्र}\ A}}\skt{ }\edtext{\skt{करपत्रात्~। तद्धि}}{\lemma{\skt{करपत्रात्~। तद्धि}}\Afootnote{A K H; \skt{करप्रत्रात्~॥ तद्धि}\ N}}\skt{ }\edtext{\skt{खर\-धा\-र\-म\-स्थि\-च्छे\-द\-ना\-र्थम्}}{\lemma{\skt{खरधारम्}}\Afootnote{A K H; \skt{खुरधारम्}\ N}}\skt{~।}\pend

\pstart \ledsidenote{\passref{1.8.13}}\edtext{\skt{तेषान्निशानी}}{\lemma{\skt{निशानी}}\Afootnote{K N; \skt{निशानार्थं}\ A; \skt{निमानी}\ H}}\skt{ }\edtext{\skt{श्लक्ष्णशिलिका}}{\lemma{\skt{श्लक्ष्णशिलिका}}\Afootnote{K N; \skt{श्लक्ष्णशिला मा}\ A; \skt{श्लक्ष्णसिलिका}\ H}}\skt{ }\edtext{\skt{धारा\-स\-म्पा\-द\-नार्थं}}{\lemma{\skt{धारा\appbrk{}स\appbrk{}म्पा\appbrk{}द\appbrk{}नार्थं}}\Afootnote{K N H; \skt{षवर्णा; धारा\appbrk{}सं\appbrk{}स्था\appbrk{}प\appbrk{}नार्थं}\ A}}\skt{ शाल्मलीफलकं }\edtext{\skt{चेति}}{\lemma{\skt{चेति}}\Afootnote{K N H; \skt{इति}\ A}}\skt{~॥}\pend

\pstart\versepar \ledsidenote{\passref{1.8.14}}\skt{भवति चात्र~॥ }\pend
\pstart\versepar \skt{यदा }\edtext{\skt{सुनिशितं}}{\lemma{\skt{सुनिशितं}}\Afootnote{A K H; \skt{सुनिषितं}\ N}}\skt{ शस्त्रं }\edtext{\skt{रोम\-वा\-हि\-सु\-सं\-स्थि\-तम्}}{\lemma{\skt{रोम\appbrk{}वा\appbrk{}हि\appbrk{}सु\appbrk{}सं\appbrk{}स्थि\appbrk{}तम्}}\Afootnote{K N; \skt{रोमच्छेदि सुसंस्थितम्}\ A; \skt{रोम\appbrk{}वा\appbrk{}हि\appbrk{}सु\appbrk{}सं\appbrk{}स्थितः}\ H}}\skt{~। }\pend
\pstart\versepar \skt{}\edtext{\skt{सुगृ\-ही\-त\-म्प्र\-मा\-णेन}}{\lemma{\skt{प्रमाणेन}}\Afootnote{A K H; \skt{प्रमाणेनेन}\ N}}\skt{ तदा }\edtext{\skt{शस्त्रं}}{\lemma{\skt{शस्त्रं}}\Afootnote{K N H; \skt{कर्म}\ A}}\skt{ }\edtext{\skt{निपातयेत्}}{\lemma{\skt{निपातयेत्}}\Afootnote{K H; \skt{सु योजयेत्}\ A; \skt{निधारयेत्}\ N}}\skt{~॥}\edtext{}{\lemma{\skt{भवति चात्र~॥ यदा}}\Afootnote{A; \skt{यदा}\ K N H}}\pend

\pstart \ledsidenote{\passref{1.8.15}}\edtext{\skt{अनुशस्त्राणि}}{\lemma{\skt{अनुशस्त्राणि}}\Afootnote{K N H; \skt{अनुशस्त्राणि तु}\ A}}\skt{ }\edtext{\skt{त्वक्क्षा\-र\-स्फ\-टि\-क\-का\-च\-कु\-र\-वि\-न्द\-ज\-लौ\-का\-ग्नि\-न\-ख\-प\-त्राणि}}{\lemma{\skt{त्वक्क्षा\appbrk{}र\appbrk{}स्फ\appbrk{}टि\appbrk{}क\appbrk{}का\appbrk{}च\appbrk{}कु\appbrk{}र\appbrk{}वि\appbrk{}न्द\appbrk{}ज\appbrk{}लौ\appbrk{}का\appbrk{}ग्नि\appbrk{}न\appbrk{}ख\appbrk{}प\appbrk{}त्राणि}}\Afootnote{K H; \skt{त्वक्सा\appbrk{}र\appbrk{}स्फ\appbrk{}टि\appbrk{}क\appbrk{}का\appbrk{}च\appbrk{}कु\appbrk{}रु\appbrk{}वि\appbrk{}न्द\appbrk{}ज\appbrk{}लौ\appbrk{}को\appbrk{}ग्नि\appbrk{}क्षा\appbrk{}र\appbrk{}न\appbrk{}ख\appbrk{}गो\appbrk{}जी\appbrk{}शे\appbrk{}फा\appbrk{}लि\appbrk{}का\appbrk{}शा\appbrk{}क\appbrk{}प\appbrk{}त्र\appbrk{}क\appbrk{}री\appbrk{}र\appbrk{}बा\appbrk{}ला\appbrk{}ङ्गु\appbrk{}लय इति}\ A; \skt{त्वक्क्षा\appbrk{}र\appbrk{}स्फ\appbrk{}टि\appbrk{}क\appbrk{}का\appbrk{}च\appbrk{}कु\appbrk{}र\appbrk{}वि\appbrk{}न्द\appbrk{}ज\appbrk{}लौ\appbrk{}का\appbrk{}ग्नि\appbrk{}न\appbrk{}ख\appbrk{}पत्रा}\ N}}\skt{~॥}\pend

\pstart\versepar \ledsidenote{\passref{1.8.16}}\skt{शिशूनां }\edtext{\edtext{\skt{शस्त्र\-भी\-रू\-णा\-म\-नु\-श\-स्त्राणि}}{\lemma{\skt{शस्त्रभीरूणाम्}}\Afootnote{A K H; \skt{शस्त्रभीरुणाम्}\ N}}}{\lemma{\skt{अनुशस्त्राणि}}\Afootnote{K N H; \skt{शस्त्राभावे च}\ A}}\skt{ योजयेत्~। }\pend
\pstart\versepar \skt{त्वक्क्षा\-रा\-दि\-च\-तु\-र्वर्गं }\edtext{\skt{भेद्ये}}{\lemma{\skt{भेद्ये}}\Afootnote{K N H; \skt{छेद्ये}\ A}}\skt{ }\edtext{\skt{च्छेद्ये च}}{\lemma{\skt{च्छेद्ये च}}\Afootnote{K N H; \skt{भेद्ये च}\ A}}\skt{ }\edtext{\skt{बुद्धिमान्}}{\lemma{\skt{बुद्धिमान्}}\Afootnote{A N H; \skt{बुद्धिमां}\ K}}\skt{~।}\edtext{}{\lemma{\skt{योजयेत्~। त्वक्क्षा\appbrk{}रा\appbrk{}दि\appbrk{}च\appbrk{}तु\appbrk{}र्वर्गं}}\Afootnote{K H; \skt{योजयेत्~। त्वक्सा\appbrk{}रा\appbrk{}दि\appbrk{}च\appbrk{}तु\appbrk{}र्वर्गं}\ A; \skt{योजयेत्~॥ त्वक्क्षा\appbrk{}रा\appbrk{}दि\appbrk{}च\appbrk{}तु\appbrk{}व\appbrk{}र्गम्}\ N}}\pend

\pstart\versepar \ledsidenote{\passref{1.8.17}}\edtext{\skt{आहार्ये}}{\lemma{\skt{आहार्ये}}\Afootnote{K H; \skt{आहार्य}\ A; \skt{आहार}\ N}}\skt{ }\edtext{\skt{च्छेद्ये}}{\lemma{\skt{च्छेद्ये}}\Afootnote{K H; \skt{च्छेद्य}\ A; \skt{च्छेद्ये}\ N}}\skt{ }\edtext{\skt{भेद्ये च}}{\lemma{\skt{भेद्ये च}}\Afootnote{K N; \skt{भेद्येषु}\ A; \skt{भेदे च}\ H}}\skt{ नखं शक्येषु योजयेत्~। }\pend
\pstart\versepar \skt{विधिः प्रवक्ष्यते }\edtext{\skt{पश्चा\-द\-ग्नि\-क्षा\-र\-ज\-लौ\-कसां}}{\lemma{\skt{पश्चा\appbrk{}द\appbrk{}ग्नि\appbrk{}क्षा\appbrk{}र\appbrk{}ज\appbrk{}लौ\appbrk{}कसां}}\Afootnote{K N H; \skt{पश्चा\appbrk{}त्क्षा\appbrk{}र\appbrk{}व\appbrk{}ह्नि\appbrk{}ज\appbrk{}लौ\appbrk{}क\appbrk{}साम्}\ A}}\skt{~।}\edtext{}{\lemma{\skt{योजयेत्~। विधिः}}\Afootnote{A K H; \skt{योजयेत्~। विधि}\ N}}\pend

\pstart\versepar \ledsidenote{\passref{1.8.18}}\edtext{\skt{ये }\edtext{\skt{स्युर्मुखगता}}{\lemma{\skt{मुखगता}}\Afootnote{A K H; \skt{मुखगतां}\ N}}}{\lemma{\skt{ये स्युर्}}\Afootnote{A K H; \skt{ये सु}\ N}}\skt{ रोगा }\edtext{\skt{नेत्र\-व\-र्म\-ग\-ताश्च ये}}{\lemma{\skt{नेत्र\appbrk{}व\appbrk{}र्म\appbrk{}ग\appbrk{}ताश्च ये}}\Afootnote{K N H; \skt{नेत्र\appbrk{}व\appbrk{}र्त्म\appbrk{}ग\appbrk{}ताश्च ये}\ A}}\skt{~। }\pend
\pstart\versepar \skt{}\edtext{\edtext{\skt{गोजी\-शे\-फा\-लि\-का\-शा\-क\-प\-त्रै\-र्वि\-स्रा\-व\-येत्तु}}{\lemma{\skt{गोजी\appbrk{}शे\appbrk{}फा\appbrk{}लि\appbrk{}का\appbrk{}शा\appbrk{}क\appbrk{}प\appbrk{}त्रैर्}}\Afootnote{K H; \skt{गोजी\appbrk{}शे\appbrk{}फा\appbrk{}लि\appbrk{}का\appbrk{}शा\appbrk{}क\appbrk{}पत्रै}\ A; \skt{गोजी\appbrk{}से\appbrk{}हा\appbrk{}लि\appbrk{}का\appbrk{}शा\appbrk{}क\appbrk{}पत्रै}\ N}}}{\lemma{\skt{विस्रावयेत्तु}}\Afootnote{K N; \skt{र्विस्रावयेत्तु}\ A; \skt{व्विश्रावयेत्तु}\ H}}\skt{ }\edtext{\skt{तान्}}{\lemma{\skt{तान्}}\Afootnote{A K; \skt{तां}\ N H}}\skt{~॥}\pend

\pstart\versepar \ledsidenote{\passref{1.8.19}}\edtext{\skt{शस्त्राण्येतानि}}{\lemma{\skt{शस्त्राण्य्}}\Afootnote{K N; \skt{एष्ये\appbrk{}ष्वे\appbrk{}ष\appbrk{}ण्य\appbrk{}लाभे तु बालाङ्गुल्यङ्कुरा हिताः~॥ शस्त्राण्य्}\ A; \skt{शस्त्रान्य्}\ H}}\skt{ }\edtext{\skt{मति\-मा\-न्शु\-द्ध\-शै\-क्या\-य\-सानि तु}}{\lemma{\skt{मति\appbrk{}मा\appbrk{}न्शु\appbrk{}द्ध\appbrk{}शै\appbrk{}क्या\appbrk{}य\appbrk{}सानि तु}}\Afootnote{A; \skt{मतिमां शुद्ध\appbrk{}शै\appbrk{}क्या\appbrk{}य\appbrk{}सानि तु}\ K; \skt{मतिमां शुद्ध\appbrk{}शै\appbrk{}क्या\appbrk{}य\appbrk{}सानि तुं}\ N; \skt{मति\appbrk{}मा\appbrk{}न्शु\appbrk{}द्ध\appbrk{}सै\appbrk{}क्या\appbrk{}य\appbrk{}सानि तु}\ H}}\skt{~। }\pend
\pstart\versepar \skt{}\edtext{\skt{कार\-ये\-त्क\-र\-ण\-प्राप्तः}}{\lemma{\skt{कार\appbrk{}ये\appbrk{}त्क\appbrk{}र\appbrk{}ण\appbrk{}प्राप्तः}}\Afootnote{K\textsuperscript{pc} H; \skt{कार\appbrk{}ये\appbrk{}त्क\appbrk{}र\appbrk{}ण\appbrk{}प्राप्तं}\ A; \skt{कार\appbrk{}ये\appbrk{}त्क\appbrk{}र\appbrk{}ण\appbrk{}प्राप्त}\ K\textsuperscript{ac}; \skt{कार\appbrk{}ये\appbrk{}त्क\appbrk{}ण\appbrk{}प्राप्तः}\ N}}\skt{ }\edtext{\skt{कर्मा\-र\-क\-र्म\-को\-विद}}{\lemma{\skt{कर्मा\appbrk{}र\appbrk{}क\appbrk{}र्म\appbrk{}को\appbrk{}विद}}\Afootnote{K\textsuperscript{pc} N; \skt{कर्मारं कर्मकोविदम्}\ A; \skt{कर्मा\appbrk{}रं\appbrk{}क\appbrk{}र्म\appbrk{}को\appbrk{}विद}\ K\textsuperscript{ac}; \skt{कर्म्मारः कर्म्मकोविद}\ H}}\skt{ }\edtext{\skt{इति}}{\lemma{\skt{इति}}\Afootnote{K N H; \emph{om.}\ A}}\skt{~॥}\pend

\pstart \ledsidenote{\passref{1.8.trailer}}\edtext{\skt{॥ ८~॥}}{\lemma{\skt{॥ ८~॥}}\Afootnote{\skt{इति सुश्रु\appbrk{}त\appbrk{}सं\appbrk{}हि\appbrk{}तायां सूत्रस्थाने शस्त्रा\appbrk{}व\appbrk{}चा\appbrk{}र\appbrk{}णीयो नामाष्टमोऽध्यायः}\ A; \skt{॥ {\symfont ❈}~॥ ८}\ K; \skt{८}\ H}}\pend

\endnumbering
\beginnumbering
\eledsection*{Adhyāya 9}
\pstart \ledsidenote{\passref{1.9.1}}\skt{अथातो योग्या\-सू\-त्री\-य\-म\-ध्यायं }\edtext{\skt{व्याख्यास्यामः}}{\lemma{\skt{व्याख्यास्यामः}}\Afootnote{A N H; \emph{om.}\ K}}\skt{~॥}\pend

\pstart \ledsidenote{\passref{1.9.3}}\edtext{\skt{अधि\-ग\-त\-स\-र्व\-शा\-स्त्र\-मपि}}{\lemma{\skt{अधि\appbrk{}ग\appbrk{}त\appbrk{}स\appbrk{}र्व\appbrk{}शा\appbrk{}स्त्रम्}}\Afootnote{K N H; \skt{अधि\appbrk{}ग\appbrk{}त\appbrk{}स\appbrk{}र्व\appbrk{}शा\appbrk{}स्त्रा\appbrk{}र्थम्}\ A}}\skt{ शिष्यं योग्यां }\edtext{\skt{कारयेत्~। च्छेद्यादिषु}}{\lemma{\skt{कारयेत्~। च्छेद्यादिषु}}\Afootnote{K N; \skt{कारयेत्~। स्नेहादिषु}\ A; \skt{कारयेत्~। च्छेद्यादीषु}\ H}}\skt{ }\edtext{\skt{स्नेहादिषु}}{\lemma{\skt{स्नेहादिषु}}\Afootnote{K; \skt{छेद्यादिषु च}\ A; \skt{स्नेहादीषु च}\ N; \skt{स्नेहादि च}\ H}}\skt{ }\edtext{\skt{कर्म\-प\-थ\-मु\-प\-दि\-शेत्~। बहुश्रुतो}}{\lemma{\skt{उपदिशेत्~। बहुश्रुतो}}\Afootnote{N; \skt{उपदिशेत्~। सुबहुश्रुतो}\ A; \skt{उपदिशेत्~। बहुश्रुतो}\ K; \skt{उपदिश्येत्~। बहुश्रुतोऽ}\ H}}\skt{ }\edtext{\skt{प्यकृतयोगः}}{\lemma{\skt{प्यकृतयोगः}}\Afootnote{\skt{ऽप्यकृतयोग्यः}\ A N; \skt{प्यकृतयोगः}\ K; \skt{प्यकृतयोग्यः}\ H}}\skt{ }\edtext{\skt{कर्मस्वयोग्यो}}{\lemma{\skt{कर्मस्वयोग्यो}}\Afootnote{\skt{कर्मस्वयोग्यो}\ A K H; \skt{कर्मश्चयोग्यो}\ N}}\skt{ भवति~।}\pend

\pstart \ledsidenote{\passref{1.9.4}}\skt{तत्र }\edtext{\skt{पुष्प\-फ\-ला\-ला\-वु\-त्र\-पु\-से\-र्वा\-रु\-क\-प्र\-भृ\-तिषु}}{\lemma{\skt{पुष्प\appbrk{}फ\appbrk{}ला\appbrk{}ला\appbrk{}वु\appbrk{}त्र\appbrk{}पु\appbrk{}से\appbrk{}र्वा\appbrk{}रु\appbrk{}क\appbrk{}प्र\appbrk{}भृ\appbrk{}तिषु}}\Afootnote{K; \skt{पुष्प\appbrk{}फ\appbrk{}ला\appbrk{}ला\appbrk{}बू\appbrk{}का\appbrk{}लि\appbrk{}न्द\appbrk{}क\appbrk{}त्र\appbrk{}पु\appbrk{}सै(\appbrk{}सो)\appbrk{}र्वा\appbrk{}रु\appbrk{}क\appbrk{}र्का\appbrk{}रु\appbrk{}क\appbrk{}प्र\appbrk{}भृ\appbrk{}तिषु}\ A; \skt{पुष्प\appbrk{}ह\appbrk{}ला\appbrk{}ला\appbrk{}वू\appbrk{}त्र\appbrk{}पु\appbrk{}मे\appbrk{}वा\appbrk{}रु\appbrk{}क\appbrk{}प्र\appbrk{}भृ\appbrk{}तिषु}\ N; \skt{पुष्प\appbrk{}फ\appbrk{}ला\appbrk{}ला\appbrk{}वु\appbrk{}त्र\appbrk{}पु\appbrk{}षे\appbrk{}र्व्वा\appbrk{}रु\appbrk{}क\appbrk{}प्र\appbrk{}भृ\appbrk{}तिषु}\ H}}\skt{ }\edtext{\skt{च्छेद्यशेषां}}{\lemma{\skt{च्छेद्यशेषां}}\Afootnote{K\textsuperscript{pc}; \skt{छेद्यविशेषान्}\ A K\textsuperscript{ac} H; \skt{छेद्यविषौषान्}\ N}}\skt{ }\edtext{\skt{दर्शयेत्~। उत्क\-र्त\-ना\-प\-क\-र्त्त\-नानि}}{\lemma{\skt{दर्शयेत्~। उत्क\appbrk{}र्त\appbrk{}ना\appbrk{}प\appbrk{}क\appbrk{}र्त्त\appbrk{}नानि}}\Afootnote{A K H; \skt{दर्शयेत्~॥ उत्क\appbrk{}र्त्त\appbrk{}ना\appbrk{}य\appbrk{}क\appbrk{}र्त्त\appbrk{}नानि}\ N}}\skt{ }\edtext{\skt{चोपदिशेत्~। दृति\-ब\-स्ति\-प्र\-से\-व\-क\-पू\-र्णेषु}}{\lemma{\skt{चोपदिशेत्~। दृति\appbrk{}ब\appbrk{}स्ति\appbrk{}प्र\appbrk{}से\appbrk{}व\appbrk{}क\appbrk{}पू\appbrk{}र्णेषु}}\Afootnote{K N; \skt{चोपदिशेत्; दृति\appbrk{}ब\appbrk{}स्ति\appbrk{}प्र\appbrk{}से\appbrk{}व\appbrk{}क\appbrk{}प्र\appbrk{}भृ\appbrk{}ति\appbrk{}षू\appbrk{}द\appbrk{}क\appbrk{}प\appbrk{}ङ्क\appbrk{}पू\appbrk{}र्णेषु}\ A; \skt{चोपदिश्येत्~। दृति\appbrk{}ब\appbrk{}स्ति\appbrk{}प्र\appbrk{}से\appbrk{}व\appbrk{}क\appbrk{}पू\appbrk{}र्ण्णेषु}\ H}}\skt{ भेद्ययोग्यां~। }\edtext{\skt{सरोम्णि}}{\lemma{\skt{सरोम्णि}}\Afootnote{A K H; \skt{सरोस्ति}\ N}}\skt{ }\edtext{\skt{चर्मातते}}{\lemma{\skt{चर्मातते}}\Afootnote{K; \skt{चर्मण्यातते}\ A; \skt{चर्मतते}\ N; \skt{चर्म्मातने}\ H}}\skt{ लेख्यस्य~। }\edtext{\skt{मृत\-प\-शु\-म\-हि\-षा\-स्व\-शि\-रा\-सू\-त्प\-ल\-ना\-लेषु}}{\lemma{\skt{मृत\appbrk{}प\appbrk{}शु\appbrk{}म\appbrk{}हि\appbrk{}षा\appbrk{}स्व\appbrk{}शि\appbrk{}रा\appbrk{}सू\appbrk{}त्प\appbrk{}ल\appbrk{}ना\appbrk{}लेषु}}\Afootnote{K; \skt{मृत\appbrk{}प\appbrk{}शु\appbrk{}सि\appbrk{}रा\appbrk{}सू\appbrk{}त्प\appbrk{}ल\appbrk{}ना\appbrk{}लेषु च}\ A; \skt{मृत\appbrk{}प\appbrk{}शु\appbrk{}म\appbrk{}हि\appbrk{}षा\appbrk{}श्व\appbrk{}शि\appbrk{}वा\appbrk{}सू\appbrk{}त्य\appbrk{}ल\appbrk{}ना\appbrk{}लेषु}\ N; \skt{मृत\appbrk{}प\appbrk{}शु\appbrk{}म\appbrk{}हि\appbrk{}षा\appbrk{}श्व\appbrk{}शि\appbrk{}रा\appbrk{}सू\appbrk{}त्प\appbrk{}ल\appbrk{}ना\appbrk{}लेषु}\ H}}\skt{ वेध्यस्य~। }\edtext{\skt{घुणो\-प\-ह\-त\-का\-ष्ठ\-वे\-णु\-न\-ल\-ना\-डी\-शु\-ष्का\-ला\-बु\-मु\-खे\-ष्वे\-ष्यस्य}}{\lemma{\skt{घुणो\appbrk{}प\appbrk{}ह\appbrk{}त\appbrk{}का\appbrk{}ष्ठ\appbrk{}वे\appbrk{}णु\appbrk{}न\appbrk{}ल\appbrk{}ना\appbrk{}डी\appbrk{}शु\appbrk{}ष्का\appbrk{}ला\appbrk{}बु\appbrk{}मु\appbrk{}खे\appbrk{}ष्वे\appbrk{}ष्यस्य}}\Afootnote{K\textsuperscript{pc} H\textsuperscript{ac}; \skt{घुणो\appbrk{}प\appbrk{}ह\appbrk{}त\appbrk{}का\appbrk{}ष्ठ\appbrk{}वे\appbrk{}णु\appbrk{}न\appbrk{}ल\appbrk{}ना\appbrk{}ली\appbrk{}शु\appbrk{}ष्का\appbrk{}ला\appbrk{}बू\appbrk{}मु\appbrk{}खे\appbrk{}ष्वे\appbrk{}ष्यस्य}\ A; \skt{घुणो\appbrk{}प\appbrk{}ह\appbrk{}त\appbrk{}का\appbrk{}ष्ठ\appbrk{}त\appbrk{}वे\appbrk{}णु\appbrk{}न\appbrk{}ल\appbrk{}ना\appbrk{}डी\appbrk{}शु\appbrk{}ष्का\appbrk{}ला\appbrk{}बु\appbrk{}मु\appbrk{}खे\appbrk{}ष्वे\appbrk{}ष्यस्य}\ K\textsuperscript{ac}; \skt{घुणो\appbrk{}प\appbrk{}ह\appbrk{}त\appbrk{}का\appbrk{}ष्ठ\appbrk{}वे\appbrk{}णु\appbrk{}न\appbrk{}ल\appbrk{}ना\appbrk{}डी\appbrk{}शु\appbrk{}ष्का\appbrk{}ला\appbrk{}वू\appbrk{}मु\appbrk{}खे\appbrk{}ष्ठे\appbrk{}ष्यस्य}\ N; \skt{घुणो\appbrk{}प\appbrk{}ह\appbrk{}त\appbrk{}का\appbrk{}ष्ठ\appbrk{}भ्र\appbrk{}वे\appbrk{}णु\appbrk{}न\appbrk{}ल\appbrk{}ना\appbrk{}डी\appbrk{}शु\appbrk{}ष्का\appbrk{}ला\appbrk{}बु\appbrk{}मु\appbrk{}खे\appbrk{}ष्वे\appbrk{}ष्यस्य}\ H\textsuperscript{pc}}}\skt{~। }\edtext{\skt{पन\-स\-बि\-म्बी\-फ\-ल\-म\-ज्जा\-मृ\-त\-प\-शु\-द\-न्ते\-ष्वा\-हा\-र्यस्य}}{\lemma{\skt{पन\appbrk{}स\appbrk{}बि\appbrk{}म्बी\appbrk{}फ\appbrk{}ल\appbrk{}म\appbrk{}ज्जा\appbrk{}मृ\appbrk{}त\appbrk{}प\appbrk{}शु\appbrk{}द\appbrk{}न्ते\appbrk{}ष्वा\appbrk{}हा\appbrk{}र्यस्य}}\Afootnote{K H; \skt{पन\appbrk{}स\appbrk{}बि\appbrk{}म्बी\appbrk{}बि\appbrk{}ल्व\appbrk{}फ\appbrk{}ल\appbrk{}म\appbrk{}ज्ज\appbrk{}मृ\appbrk{}त\appbrk{}प\appbrk{}शु\appbrk{}द\appbrk{}न्ते\appbrk{}ष्वा\appbrk{}हा\appbrk{}र्यस्य; मधू\appbrk{}च्छि\appbrk{}ष्टो\appbrk{}प\appbrk{}लिप्ते शाल्मलीफलके विस्राव्यस्य}\ A; \skt{पन\appbrk{}स\appbrk{}विं\appbrk{}बी\appbrk{}फ\appbrk{}ल\appbrk{}मु\appbrk{}ज्जा\appbrk{}मृ\appbrk{}भप पशु\appbrk{}द\appbrk{}न्ते\appbrk{}ष्ठ\appbrk{}हा\appbrk{}र्यस्य}\ N}}\skt{ }\edtext{\skt{सूक्ष्म\-घ\-न\-व\-स्त्रा\-न्त\-मृ\-दु\-च\-र्मा\-न्तयोः}}{\lemma{\skt{सूक्ष्म\appbrk{}घ\appbrk{}न\appbrk{}व\appbrk{}स्त्रा\appbrk{}न्त\appbrk{}मृ\appbrk{}दु\appbrk{}च\appbrk{}र्मा\appbrk{}न्तयोः}}\Afootnote{K H; \skt{सूक्ष्म\appbrk{}घ\appbrk{}न\appbrk{}व\appbrk{}स्त्रा\appbrk{}न्त\appbrk{}यो\appbrk{}र्मृ\appbrk{}दु\appbrk{}च\appbrk{}र्मा\appbrk{}न्त\appbrk{}योश्च}\ A; \skt{सूक्ष्म\appbrk{}म\appbrk{}न\appbrk{}व\appbrk{}स्त्रा\appbrk{}न्ठ\appbrk{}मृ\appbrk{}दु\appbrk{}च\appbrk{}र्मा\appbrk{}न्तयोः}\ N}}\skt{ }\edtext{\skt{सीव्यस्य~। मृदु\-व\-ध्र\-मां\-स\-पे\-स्यु\-त्प\-ल\-ना\-लेषु च}}{\lemma{\skt{सीव्यस्य~। मृदु\appbrk{}व\appbrk{}ध्र\appbrk{}मां\appbrk{}स\appbrk{}पे\appbrk{}स्यु\appbrk{}त्प\appbrk{}ल\appbrk{}ना\appbrk{}लेषु च}}\Afootnote{K H; \skt{सीव्यस्य; पुस्त\appbrk{}म\appbrk{}य\appbrk{}पु\appbrk{}रु\appbrk{}षा\appbrk{}ङ्ग\appbrk{}प्र\appbrk{}त्य\appbrk{}ङ्ग\appbrk{}वि\appbrk{}शे\appbrk{}षेषु बन्धनयोग्यां; मृदुषु मांस\appbrk{}ख\appbrk{}ण्डे\appbrk{}ष्व\appbrk{}ग्नि\appbrk{}क्षा\appbrk{}र\appbrk{}योग्यां; मृदु\appbrk{}च\appbrk{}र्म\appbrk{}मां\appbrk{}स\appbrk{}पे\appbrk{}शी\appbrk{}षू\appbrk{}त्प\appbrk{}ल\appbrk{}ना\appbrk{}लेषु च}\ A; \skt{मीव्यस्य मृदु\appbrk{}च\appbrk{}प्र\appbrk{}मां\appbrk{}स\appbrk{}पे\appbrk{}स्यु\appbrk{}त्प\appbrk{}ल\appbrk{}ना\appbrk{}लेषु च}\ N}}\skt{ }\edtext{\skt{कर्ण्ण\-स\-न्धि\-ब\-न्ध\-न\-यो\-ग्याम्}}{\lemma{\skt{कर्ण्ण\appbrk{}स\appbrk{}न्धि\appbrk{}ब\appbrk{}न्ध\appbrk{}न\appbrk{}यो\appbrk{}ग्याम्}}\Afootnote{K; \skt{कर्ण\appbrk{}स\appbrk{}न्धि\appbrk{}ब\appbrk{}न्ध\appbrk{}यो\appbrk{}ग्याम्}\ A H; \skt{कर्ण्ण\appbrk{}स\appbrk{}च्चि\appbrk{}व\appbrk{}न्व\appbrk{}न\appbrk{}योग्यां}\ N}}\skt{~। }\edtext{\skt{पुस्त\-म\-य\-पु\-रु\-षां\-ग\-प्र\-त्य\-ङ्गेषु}}{\lemma{\skt{पुस्त\appbrk{}म\appbrk{}य\appbrk{}पु\appbrk{}रु\appbrk{}षां\appbrk{}ग\appbrk{}प्र\appbrk{}त्य\appbrk{}ङ्गेषु}}\Afootnote{K; \skt{उ}\ A; \skt{पुस्त\appbrk{}म\appbrk{}य\appbrk{}षु\appbrk{}रु\appbrk{}षां\appbrk{}ग\appbrk{}प्र\appbrk{}त्प\appbrk{}द्गेषु}\ N; \skt{पुम्त\appbrk{}म\appbrk{}य\appbrk{}पु\appbrk{}रु\appbrk{}शा\appbrk{}ङ्ग\appbrk{}प्र\appbrk{}त्य\appbrk{}ङ्गेषु}\ H\textsuperscript{ac}; \skt{पुस्त\appbrk{}म\appbrk{}य\appbrk{}पु\appbrk{}रु\appbrk{}शा\appbrk{}ङ्ग\appbrk{}प्र\appbrk{}त्य\appbrk{}ङ्गेषु}\ H\textsuperscript{pc}}}\skt{ }\edtext{\skt{बन्धनयोग्यां}}{\lemma{\skt{बन्धनयोग्यां}}\Afootnote{K N H; \skt{दक}\ A}}\skt{~। }\edtext{\skt{घटा\-ला\-वु\-मु\-खेषु}}{\lemma{\skt{घटा\appbrk{}ला\appbrk{}वु\appbrk{}मु\appbrk{}खेषु}}\Afootnote{K H; \skt{पूर्ण\appbrk{}घ\appbrk{}ट\appbrk{}पा\appbrk{}र्श्व\appbrk{}स्रो\appbrk{}त\appbrk{}स्य\appbrk{}ला\appbrk{}बू\appbrk{}मु\appbrk{}खादि}\ A; \skt{घठा\appbrk{}ला\appbrk{}वु\appbrk{}मु\appbrk{}खेषु}\ N}}\skt{ }\edtext{\skt{बस्ति\-व्र\-ण\-ब\-स्ति\-पी\-ड\-न\-योग्यां}}{\lemma{\skt{बस्ति\appbrk{}व्र\appbrk{}ण\appbrk{}ब\appbrk{}स्ति\appbrk{}पी\appbrk{}ड\appbrk{}न\appbrk{}योग्यां}}\Afootnote{K H; \skt{षु च}\ A; \skt{बस्ति\appbrk{}पी\appbrk{}ड\appbrk{}न\appbrk{}योग्यां}\ N}}\skt{~। }\edtext{\edtext{\skt{नेत्र\-प्र\-णि\-धा\-न\-ब\-स्ति\-पी\-ड\-न\-यो\-रिति}}{\lemma{\skt{नेत्र\appbrk{}प्र\appbrk{}णि\appbrk{}धा\appbrk{}न\appbrk{}ब\appbrk{}स्ति\appbrk{}पी\appbrk{}ड\appbrk{}न\appbrk{}योर्}}\Afootnote{K H; \skt{नेत्र\appbrk{}प्र\appbrk{}णि\appbrk{}धा\appbrk{}न\appbrk{}ब\appbrk{}स्ति\appbrk{}व्र\appbrk{}ण\appbrk{}ब\appbrk{}स्ति\appbrk{}पी\appbrk{}ड\appbrk{}नयो}\ A; \skt{नेत्र\appbrk{}प्र\appbrk{}णि\appbrk{}धा\appbrk{}न\appbrk{}ब\appbrk{}स्ति\appbrk{}पी\appbrk{}ड\appbrk{}नयो}\ N}}}{\lemma{\skt{इति}}\Afootnote{K H; \skt{ग्यामिति}\ A; \skt{रिति}\ N}}\skt{~॥}\pend

\pstart\versepar \ledsidenote{\passref{1.9.5ab}}\skt{भवति चात्र~॥ }\pend
\pstart\versepar \skt{एवमादिषु }\edtext{\skt{मेधावी}}{\lemma{\skt{मेधावी}}\Afootnote{A K H; \skt{मोधावी}\ N}}\skt{ }\edtext{\skt{योग्या }\edtext{\skt{कर्मण्यशेषतः}}{\lemma{\skt{अशेषतः}}\Afootnote{K H; \skt{र्हेषु यथाविधि}\ A; \skt{असेसतः}\ N}}}{\lemma{\skt{योग्या कर्मण्य्}}\Afootnote{K\textsuperscript{pc} H; \skt{योग्या}\ A; \skt{योकर्मण्य्}\ K\textsuperscript{ac}; \skt{योजूकर्मण्य्}\ N}}\skt{~।}\edtext{}{\lemma{\skt{भवति चात्र~॥ एवम्}}\Afootnote{\skt{भवतश्चात्र एवम्}\ A; \skt{एवम्}\ K N H}}\pend

\pstart \ledsidenote{\passref{1.9.6cd}}\edtext{\skt{यस्य}}{\lemma{\skt{यस्य}}\Afootnote{A K H; \skt{यस्तु}\ N}}\skt{ }\edtext{\skt{यस्येह}}{\lemma{\skt{यस्येह}}\Afootnote{K H; \skt{यत्रेह}\ A; \skt{यस्तु व्ह्र}\ N}}\skt{ }\edtext{\skt{साधर्म्यं}}{\lemma{\skt{साधर्म्यं}}\Afootnote{A K H; \skt{साधर्म्य}\ N}}\skt{ तत्र }\edtext{\skt{योग्याञ्च}}{\lemma{\skt{योग्याञ्च}}\Afootnote{K H; \skt{योग्यां स}\ A; \skt{योग्याव}\ N}}\skt{ }\edtext{\skt{कारयेत्}}{\lemma{\skt{कारयेत्}}\Afootnote{K N H; \skt{माचरेत्}\ A}}\skt{~॥}\pend

\pstart \ledsidenote{\passref{1.9.trailer}}\edtext{\skt{॥ ९~॥}}{\lemma{\skt{॥ ९~॥}}\Afootnote{\skt{इति सुश्रु\appbrk{}त\appbrk{}सं\appbrk{}हि\appbrk{}तायां सूत्रस्थाने योग्यासूत्रीयो नाम नवमोऽध्यायः}\ A; \skt{{\symfont ❈}~॥ ९}\ K; \skt{डे}\ N H}}\pend

\endnumbering
\beginnumbering
\eledsection*{Adhyāya 10}
\pstart \ledsidenote{\passref{1.10.1}}\skt{अथातो }\edtext{\skt{विशि\-खा\-नु\-प्र\-वे\-श\-नीयं}}{\lemma{\skt{विशि\appbrk{}खा\appbrk{}नु\appbrk{}प्र\appbrk{}वे\appbrk{}श\appbrk{}नीयं}}\Afootnote{K H; \skt{विशि\appbrk{}खा\appbrk{}नु\appbrk{}प्र\appbrk{}वे\appbrk{}श\appbrk{}नी\appbrk{}य\appbrk{}म\appbrk{}ध्यायं}\ A; \skt{विशि\appbrk{}खा\appbrk{}नु\appbrk{}प्र\appbrk{}वे\appbrk{}शि\appbrk{}नीयं}\ N}}\skt{ }\edtext{\skt{व्याख्यास्यामः}}{\lemma{\skt{व्याख्यास्यामः}}\Afootnote{A N H; \emph{om.}\ K}}\skt{~॥}\pend

\pstart \ledsidenote{\passref{1.10.3}}\edtext{\skt{अधि\-ग\-त\-त\-न्त्रे\-णो\-पा\-सि\-त\-त\-न्त्रा\-र्थेन}}{\lemma{\skt{अधि\appbrk{}ग\appbrk{}त\appbrk{}त\appbrk{}न्त्रे\appbrk{}णो\appbrk{}पा\appbrk{}सि\appbrk{}त\appbrk{}त\appbrk{}न्त्रा\appbrk{}र्थेन}}\Afootnote{A K; \skt{अधि\appbrk{}ग\appbrk{}त\appbrk{}न्त्रे\appbrk{}नो\appbrk{}पा\appbrk{}सि\appbrk{}त\appbrk{}त\appbrk{}न्त्रा\appbrk{}र्थिषु}\ N; \skt{अधि\appbrk{}ग\appbrk{}त\appbrk{}त\appbrk{}न्त्रे\appbrk{}णौ\appbrk{}पा\appbrk{}सि\appbrk{}त\appbrk{}त\appbrk{}न्त्रा\appbrk{}र्थेन}\ H}}\skt{ }\edtext{\skt{दृष्टकर्मणा}}{\lemma{\skt{दृष्टकर्मणा}}\Afootnote{A K H; \skt{दृष्टकर्मण}\ N}}\skt{ कृतयोग्येन }\edtext{\skt{शास्त्रा\-र्थ\-न्नि\-ग\-दता}}{\lemma{\skt{शास्त्रा\appbrk{}र्थ\appbrk{}न्नि\appbrk{}ग\appbrk{}दता}}\Afootnote{K H; \skt{शास्त्रं निगदता}\ A; \skt{शास्त्रा\appbrk{}र्थ\appbrk{}नि\appbrk{}ग\appbrk{}दता}\ N}}\skt{ }\edtext{\skt{राजानुज्ञातेन}}{\lemma{\skt{राजानुज्ञातेन}}\Afootnote{K H; \skt{राजानुज्ञाते}\ A; \skt{राजातुज्ञातेन}\ N}}\skt{ }\edtext{\skt{वैद्येन}}{\lemma{\skt{वैद्येन}}\Afootnote{K H; \skt{न}\ A; \skt{वेद्येन}\ N}}\skt{ }\edtext{\skt{विशि\-खा\-च\-रि\-तव्या}}{\lemma{\skt{विशि\appbrk{}खा\appbrk{}च\appbrk{}रि\appbrk{}तव्या}}\Afootnote{K H; \emph{om.}\ A; \skt{शिखाचरितव्या}\ N}}\skt{~। }\edtext{\skt{नीचनखरोम्णा}}{\lemma{\skt{नीचनखरोम्णा}}\Afootnote{A K H; \skt{नीचनखरोष्णा}\ N}}\skt{ }\edtext{\skt{शुचिना}}{\lemma{\skt{शुचिना}}\Afootnote{A K N; \skt{षुचिना}\ H}}\skt{ }\edtext{\skt{शुचि\-व\-स्त्र\-प\-रि\-हि\-तेन}}{\lemma{\skt{शुचि\appbrk{}व\appbrk{}स्त्र\appbrk{}प\appbrk{}रि\appbrk{}हि\appbrk{}तेन}}\Afootnote{K N H; \skt{शुक्ल\appbrk{}व\appbrk{}स्त्र\appbrk{}प\appbrk{}रि\appbrk{}हि\appbrk{}तेन}\ A}}\skt{ छत्रवता }\edtext{\skt{सोपा\-न\-त्के\-ना\-नु\-द्ध\-त\-वे\-षेण}}{\lemma{\skt{सोपा\appbrk{}न\appbrk{}त्के\appbrk{}ना\appbrk{}नु\appbrk{}द्ध\appbrk{}त\appbrk{}वे\appbrk{}षेण}}\Afootnote{K N; \skt{दण्डहस्तेन सोपा\appbrk{}न\appbrk{}त्के\appbrk{}ना\appbrk{}नु\appbrk{}द्ध\appbrk{}त\appbrk{}वे\appbrk{}शेन}\ A; \skt{सोपा\appbrk{}न\appbrk{}त्के\appbrk{}ना\appbrk{}नु\appbrk{}द्ध\appbrk{}त\appbrk{}वे\appbrk{}षेन}\ H}}\skt{ सुमनसा }\edtext{\skt{कल्या\-णा\-भि\-व्या\-हा\-रे\-णा\-कु\-ह\-केन}}{\lemma{\skt{कल्या\appbrk{}णा\appbrk{}भि\appbrk{}व्या\appbrk{}हा\appbrk{}रे\appbrk{}णा\appbrk{}कु\appbrk{}ह\appbrk{}केन}}\Afootnote{A K H; \skt{कल्या\appbrk{}णा\appbrk{}भि\appbrk{}व्या\appbrk{}हा\appbrk{}रे\appbrk{}णा\appbrk{}कु\appbrk{}हके}\ N}}\skt{ बन्धुभूतेन भूतानां }\edtext{\skt{सहायवता}}{\lemma{\skt{सहायवता}}\Afootnote{K N H; \skt{सुसहायवता वैद्येन विशि\appbrk{}खाऽ\appbrk{}नु\appbrk{}प्र\appbrk{}वे\appbrk{}ष्टव्या}\ A}}\skt{~।}\pend

\pstart \ledsidenote{\passref{1.10.4.1}}\skt{ततो }\edtext{\edtext{\edtext{\skt{दूत\-नि\-मि\-त्त\-श\-कु\-न\-म\-ङ्ग\-ला\-नु\-लो\-म्ये\-ना\-तु\-र\-गृ\-ह\-मा\-ग\-म्यो\-प\-वि\-श्या\-तु\-र\-म\-भि\-प\-श्ये\-त्स्पृ\-शे\-त्पृ\-च्छेच्च}}{\lemma{\skt{दूत\appbrk{}नि\appbrk{}मि\appbrk{}त्त\appbrk{}श\appbrk{}कु\appbrk{}न\appbrk{}म\appbrk{}ङ्ग\appbrk{}ला\appbrk{}नु\appbrk{}लो\appbrk{}म्ये\appbrk{}ना\appbrk{}तु\appbrk{}र\appbrk{}गृ\appbrk{}हम्}}\Afootnote{K H; \skt{दूत\appbrk{}नि\appbrk{}मि\appbrk{}त्त\appbrk{}श\appbrk{}कु\appbrk{}न\appbrk{}म\appbrk{}ङ्ग\appbrk{}ला\appbrk{}नु\appbrk{}लो\appbrk{}म्ये\appbrk{}ना\appbrk{}तु\appbrk{}र\appbrk{}गृह}\ A; \skt{दूत\appbrk{}नि\appbrk{}मि\appbrk{}त्त\appbrk{}श\appbrk{}कु\appbrk{}न\appbrk{}म\appbrk{}ङ्ग\appbrk{}ला\appbrk{}नु\appbrk{}लो\appbrk{}म्ये\appbrk{}ना\appbrk{}तु\appbrk{}र\appbrk{}गृ\appbrk{}हम्}\ N}}}{\lemma{\skt{आग\appbrk{}म्यो\appbrk{}प\appbrk{}वि\appbrk{}श्या\appbrk{}तु\appbrk{}रम्}}\Afootnote{K H; \skt{मभिगम्य, उपविश्य, आतुर}\ A; \skt{आगम्यातुरम्}\ N}}}{\lemma{\skt{अभि\appbrk{}प\appbrk{}श्ये\appbrk{}त्स्पृ\appbrk{}शे\appbrk{}त्पृ\appbrk{}च्छेच्च}}\Afootnote{K N H; \skt{मभि\appbrk{}प\appbrk{}श्ये\appbrk{}त्स्पृ\appbrk{}शे\appbrk{}त्पृ\appbrk{}च्छेच्च}\ A}}\skt{ }\edtext{\edtext{\skt{त्रिभि\-रे\-तै\-र्वि\-ज्ञा\-नो\-पायैः}}{\lemma{\skt{एतैर्}}\Afootnote{K N H; \skt{एतै}\ A}}}{\lemma{\skt{विज्ञानोपायैः}}\Afootnote{K N H; \skt{र्विज्ञानोपायै}\ A}}\skt{~। }\edtext{\edtext{\edtext{\skt{दीर्घ\-मा\-यु\-षोऽ\-ल्पा\-युषो}}{\lemma{\skt{दीर्घम्}}\Afootnote{K N H; \emph{om.}\ A}}}{\lemma{\skt{आयुषो}}\Afootnote{K H; \skt{रो}\ A; \skt{आयुष्यो}\ N}}}{\lemma{\skt{ऽल्पायुषो}}\Afootnote{\skt{गाः प्रायशो}\ A; \skt{ल्पायुषो}\ K N; \skt{ऽल्पायुशो}\ H}}\skt{ }\edtext{\skt{वेदितव्याः}}{\lemma{\skt{वेदितव्याः}}\Afootnote{K; \skt{वेदितव्या}\ A N H}}\skt{~।}\footnoteA{MS K seems to have \emph{veditavyāḥ} although the final visarga isn't clear and the manuscript is damaged at this point. It was read as a visarga by Ācārya (1939: 77, note 2).}\footnoteA{The vulgate here adds a passage that is taken to be a reference to the doctrines of the Carakasaṃhitā. See HIML 1A, 351-352. The commentary Bhānumatī by Cakrapāṇidatta does not support this added passage (\emph{ayaṃ pāṭhaś cakrāsaṃmataḥ}: Ācārya 1939: 77, note 2).}\pend

\pstart \ledsidenote{\passref{1.10.5}}\edtext{\skt{तत्र दृष्ट्वा यौ}}{\lemma{\skt{तत्र दृष्ट्वा यौ}}\Afootnote{K\textsuperscript{pc}; \skt{तत्र श्रोत्रे\appbrk{}न्द्रि\appbrk{}य\appbrk{}वि\appbrk{}ज्ञेया विशेषा रोगेषु व्रणा\appbrk{}स्रा\appbrk{}व\appbrk{}वि\appbrk{}ज्ञा\appbrk{}नी\appbrk{}या\appbrk{}दिषु वक्ष्यन्ते- ‘तत्र सफेनं रक्त\appbrk{}मी\appbrk{}र\appbrk{}य\appbrk{}न्न\appbrk{}निलः सशब्दो निर्गच्छति’(सू. अ.२६) इत्येवमादयः, स्पर्श\appbrk{}ने\appbrk{}न्द्रि\appbrk{}य\appbrk{}वि\appbrk{}ज्ञेयाः शीतो\appbrk{}ष्ण\appbrk{}श्ल\appbrk{}क्ष्ण\appbrk{}क\appbrk{}र्क\appbrk{}श\appbrk{}मृ\appbrk{}दु\appbrk{}क\appbrk{}ठि\appbrk{}न\appbrk{}त्वा\appbrk{}दयः (स्पर्शविशेषा ) ज्वरशोफादिषु, चक्षु\appbrk{}रि\appbrk{}न्द्रि\appbrk{}य\appbrk{}वि\appbrk{}ज्ञेयाः शरी\appbrk{}रो\appbrk{}प\appbrk{}च\appbrk{}या\appbrk{}प\appbrk{}च\appbrk{}या\appbrk{}यु\appbrk{}र्ल\appbrk{}क्ष\appbrk{}ण\appbrk{}ब\appbrk{}ल\appbrk{}व\appbrk{}र्ण\appbrk{}वि\appbrk{}का\appbrk{}रा\appbrk{}दयः, रस\appbrk{}ने\appbrk{}न्द्रि\appbrk{}य\appbrk{}वि\appbrk{}ज्ञेयाः प्रमेहादिषु रसविशेषाः, घ्राणेन्द्रिय}\ A; \skt{तत्र दृष्ट्वायौ वर्ण्णवै}\ K\textsuperscript{ac}; \skt{तत्र दृष्ट्वा शरी\appbrk{}रो\appbrk{}प\appbrk{}च\appbrk{}या\appbrk{}प\appbrk{}चयो वर्ण\appbrk{}वै\appbrk{}कृ\appbrk{}ति\appbrk{}च्छायां चातुरस्य}\ N; \skt{तत्र दृष्ट्वा शरी\appbrk{}रो\appbrk{}प\appbrk{}च\appbrk{}या\appbrk{}प\appbrk{}यौ\appbrk{}व\appbrk{}ण्ण\appbrk{}वैकृ}\ H}}\skt{ }\edtext{\skt{वर्ण्ण\-वै\-कृ\-ति\-च्छायां}}{\lemma{\skt{वर्ण्ण\appbrk{}वै\appbrk{}कृ\appbrk{}ति\appbrk{}च्छायां}}\Afootnote{K\textsuperscript{pc}; \skt{विज्ञेया अरिष्टलिङ्}\ A; \skt{तिच्छायां}\ K\textsuperscript{ac} H; \emph{om.}\ N}}\skt{ }\edtext{\skt{चातुरस्य}}{\lemma{\skt{चातुरस्य}}\Afootnote{K H; \skt{गादिषु व्रणा}\ A; \skt{भि}\ N}}\skt{ }\edtext{\skt{भिषग्जानीयात्~। स्पृष्ट्वा}}{\lemma{\skt{भिषग्जानीयात्~। स्पृष्ट्वा}}\Afootnote{K H; \skt{नामव्रणानां च गन्ध}\ A; \skt{षग्जानीयात्~॥ स्पृष्टा}\ N}}\skt{ }\edtext{\skt{शीतोष्णादीं}}{\lemma{\skt{शीतोष्णादीं}}\Afootnote{K; \emph{om.}\ A; \skt{शीतोष्णादी}\ N H}}\skt{ }\edtext{\skt{स्पर्शविशेषा}}{\lemma{\skt{स्पर्शविशेषा}}\Afootnote{K; \skt{विशेषाः, प्रश्नेन च}\ A; \skt{स्पर्शविशेषो}\ N; \skt{न्स्पर्शविशेषा}\ H}}\skt{ }\edtext{\skt{विप\-री\-ता\-वि\-प\-रीतां}}{\lemma{\skt{विप\appbrk{}री\appbrk{}ता\appbrk{}वि\appbrk{}प\appbrk{}रीतां}}\Afootnote{K N; \skt{विजानीया}\ A; \skt{विपरीता विपरीता}\ H}}\skt{ }\edtext{\skt{ज्वरशोफादीं}}{\lemma{\skt{ज्वरशोफादीं}}\Afootnote{K; \emph{om.}\ A; \skt{सोहादीं}\ N; \skt{न्ज्वरशोफादी}\ H}}\skt{~। }\edtext{\skt{पृष्ट्वा}}{\lemma{\skt{पृष्ट्वा}}\Afootnote{K N; \emph{om.}\ A; \skt{न्~॥ पृष्ट्वा}\ H}}\skt{ }\edtext{\skt{देशं}}{\lemma{\skt{देशं}}\Afootnote{K N H; \skt{द्देशं}\ A}}\skt{ कालं }\edtext{\edtext{\edtext{\skt{जाति\-सा\-म्य\-मा\-त\-ङ्क\-मु\-त्पत्ति}}{\lemma{\skt{जातिसाम्यम्}}\Afootnote{K; \skt{जातिं सात्म्य}\ A; \skt{जातिमार्म्य्}\ N; \skt{चातिसात्म्यम्}\ H}}}{\lemma{\skt{आतङ्कम्}}\Afootnote{K H; \skt{मातङ्कस}\ A; \skt{आमातङ्काम्}\ N}}}{\lemma{\skt{उत्पत्ति}}\Afootnote{K N H; \skt{मुत्पत्तिं}\ A}}\skt{ }\edtext{\skt{वेद\-ना\-स\-मु\-च्छ्रायो}}{\lemma{\skt{वेद\appbrk{}ना\appbrk{}स\appbrk{}मु\appbrk{}च्छ्रायो}}\Afootnote{K H; \skt{वेद\appbrk{}ना\appbrk{}स\appbrk{}मु\appbrk{}च्छ्रायं}\ A; \skt{वेद\appbrk{}ना\appbrk{}स\appbrk{}मु\appbrk{}च्छायो}\ N}}\skt{ }\edtext{\edtext{\skt{बलाबलमग्निं}}{\lemma{\skt{बलाबलम्}}\Afootnote{K N H; \skt{बलमन्त}\ A}}}{\lemma{\skt{अग्निं}}\Afootnote{K N H; \skt{रग्निं}\ A}}\skt{ }\edtext{\skt{वात\-मू\-त्र\-पु\-री\-षा\-णा\-म\-प्र\-वृ\-त्ति\-प्र\-वृत्तिं}}{\lemma{\skt{अप्र\appbrk{}वृ\appbrk{}त्ति\appbrk{}प्र\appbrk{}वृत्तिं}}\Afootnote{K N; \skt{प्रवृ\appbrk{}त्ति\appbrk{}म\appbrk{}प्र\appbrk{}वृत्तिं}\ A; \skt{अप्र\appbrk{}वृ\appbrk{}त्ति\appbrk{}प्र\appbrk{}वृ\appbrk{}त्तीञ्}\ H}}\skt{ }\edtext{\skt{चेति}}{\lemma{\skt{चेति}}\Afootnote{K N H; \skt{काल\appbrk{}प्र\appbrk{}क\appbrk{}र्षा\appbrk{}दींश्च विशेषान्~। आत्मसदृशेषु विज्ञा\appbrk{}ना\appbrk{}भ्यु\appbrk{}पा\appbrk{}येषु तत्स्था\appbrk{}नी\appbrk{}यै\appbrk{}र्जा\appbrk{}नी\appbrk{}यात्}\ A}}\skt{~॥}\pend

\pstart\versepar \ledsidenote{\passref{1.10.7}}\edtext{\skt{भवति चात्र}}{\lemma{\skt{भवति चात्र}}\Afootnote{A; \emph{om.}\ K N; \skt{भवति चात्र}\ H}}\skt{~॥ }\pend
\pstart\versepar \skt{}\edtext{\skt{मिथ्यादृष्ट्या}}{\lemma{\skt{मिथ्यादृष्ट्या}}\Afootnote{N H; \skt{मिथ्यादृष्टा}\ A; \skt{भवति चात्र~॥ मिथ्यादृष्ट्वा}\ K}}\skt{ विकारा हि }\edtext{\edtext{\skt{दुरा\-ख्या\-ता\-स्त\-थैव}}{\lemma{\skt{दुराख्यातास्}}\Afootnote{A H; \emph{om.}\ K; \skt{द्वराकृता}\ N}}\skt{ च}}{\lemma{\skt{तथैव च}}\Afootnote{A N H; \emph{om.}\ K}}\skt{~। }\pend
\pstart\versepar \skt{}\edtext{\skt{तथा}}{\lemma{\skt{तथा}}\Afootnote{A N H; \emph{om.}\ K}}\skt{ }\edtext{\skt{दुष्परिस्पृष्टाश्च}}{\lemma{\skt{दुष्परिस्पृष्टाश्च}}\Afootnote{H; \skt{दुष्परिमृष्टाश्च}\ A; \skt{दू…श्च}\ K; \skt{दुष्प\appbrk{}रि\appbrk{}पृ\appbrk{}ष्ट्या\appbrk{}वैश्च}\ N}}\skt{ }\edtext{\edtext{\skt{मोह\-ये\-यु\-श्चि\-कि\-त्सकं}}{\lemma{\skt{मोहयेयुश्}}\Afootnote{K N H; \skt{मोहयेयु}\ A}}}{\lemma{\skt{चिकित्सकं}}\Afootnote{K N H; \skt{श्चिकित्सकम्}\ A}}\skt{~॥}\pend

\pstart\versepar \ledsidenote{\passref{1.10.7.add}}\edtext{\skt{तस्मात्परीक्ष्यः}}{\lemma{\skt{तस्मात्परीक्ष्यः}}\Afootnote{K H; \skt{तस्मात्चरीक्षाः}\ N}}\skt{ सततं भिषजा }\edtext{\skt{सिद्धिमिच्छता}}{\lemma{\skt{सिद्धिमिच्छता}}\Afootnote{K H; \skt{सिद्धिमिच्छता}\ N}}\skt{~। }\pend
\pstart\versepar \skt{}\edtext{\skt{युक्तितो}}{\lemma{\skt{युक्तितो}}\Afootnote{K H; \skt{युक्तिना}\ N}}\skt{ }\edtext{\skt{व्याधयः}}{\lemma{\skt{व्याधयः}}\Afootnote{K H; \skt{व्यापयः}\ N}}\skt{ }\edtext{\skt{सर्वे}}{\lemma{\skt{सर्वे}}\Afootnote{K H; \skt{सर्व्वै}\ N}}\skt{ }\edtext{\edtext{\skt{प्रमा\-णै\-र्द\-र्श\-ना\-दिभिः}}{\lemma{\skt{प्रमाणैर्}}\Afootnote{N H; \skt{प्रमाणै}\ K}}}{\lemma{\skt{दर्शनादिभिः}}\Afootnote{N H; \skt{र्दर्शनादिभिः}\ K}}\skt{~॥}\pend

\pstart \ledsidenote{\passref{1.10.6}}\edtext{\skt{एव\-म\-भि\-स\-मीक्ष्य}}{\lemma{\skt{अभिसमीक्ष्य}}\Afootnote{A K H; \skt{अभिसमिक्ष्य}\ N}}\skt{ }\edtext{\skt{साध्यां}}{\lemma{\skt{साध्यां}}\Afootnote{K H; \skt{साध्या}\ A N}}\skt{ }\edtext{\skt{साधयेद्याप्यां}}{\lemma{\skt{साधयेद्याप्यां}}\Afootnote{K; \skt{न्साधयेत्, याप्या}\ A; \skt{साधयद्याप्यां}\ N; \skt{साधयेद्याप्या}\ H}}\skt{ }\edtext{\skt{याप\-ये\-द\-सा\-ध्या\-न्नो\-प\-क्रा\-मेत्~। परि\-व\-त्स\-रो\-षि\-तांश्च}}{\lemma{\skt{याप\appbrk{}ये\appbrk{}द\appbrk{}सा\appbrk{}ध्या\appbrk{}न्नो\appbrk{}प\appbrk{}क्रा\appbrk{}मेत्~। परि\appbrk{}व\appbrk{}त्स\appbrk{}रो\appbrk{}षि\appbrk{}तांश्च}}\Afootnote{K; \skt{न्यापयेत्, असा\appbrk{}ध्या\appbrk{}न्नै\appbrk{}वो\appbrk{}प\appbrk{}क्र\appbrk{}मेत, परि\appbrk{}सं\appbrk{}व\appbrk{}त्स\appbrk{}रो\appbrk{}त्थि\appbrk{}तांश्च}\ A; \skt{पाय\appbrk{}ये\appbrk{}द\appbrk{}सा\appbrk{}ध्या\appbrk{}न्ना\appbrk{}प\appbrk{}क्रा\appbrk{}मेत्~॥ परि\appbrk{}व\appbrk{}त\appbrk{}स\appbrk{}रो\appbrk{}षितां स्व}\ N; \skt{न्याप\appbrk{}ये\appbrk{}द\appbrk{}सा\appbrk{}ध्या\appbrk{}न्नो\appbrk{}प\appbrk{}क्र\appbrk{}मेत्~। परि\appbrk{}सं\appbrk{}व\appbrk{}त्स\appbrk{}रो\appbrk{}षि\appbrk{}तान्श्च}\ H}}\skt{ }\edtext{\edtext{\skt{विकाराम्प्रायशः}}{\lemma{\skt{विकाराम्}}\Afootnote{K N; \skt{विकारा}\ A H}}}{\lemma{\skt{प्रायशः}}\Afootnote{K N; \skt{न्प्रायशो}\ A H}}\skt{ }\edtext{\skt{परिवर्ज्जयेत्}}{\lemma{\skt{परिवर्ज्जयेत्}}\Afootnote{K N H; \skt{वर्जयेत्}\ A}}\skt{~।}\pend

\pstart \ledsidenote{\passref{1.10.8}}\skt{तत्र साध्या }\edtext{\skt{अपि}}{\lemma{\skt{अपि}}\Afootnote{A K H\textsuperscript{ac}; \skt{त्चषि}\ N; \skt{अपि डो}\ H\textsuperscript{pc}}}\skt{ व्याधयः }\edtext{\skt{प्रायशो}}{\lemma{\skt{प्रायशो}}\Afootnote{K N H; \skt{प्रायेणैषां}\ A}}\skt{ }\edtext{\skt{दुश्चिकित्सा}}{\lemma{\skt{दुश्चिकित्सा}}\Afootnote{K N H; \skt{दुश्चिकित्स्य}\ A}}\skt{ }\edtext{\skt{भवन्ति}}{\lemma{\skt{भवन्ति}}\Afootnote{K N H; \skt{तमा भवन्ति~। तद्यथा}\ A}}\skt{~। }\edtext{\skt{श्रोत्रि\-य\-नृ\-प\-ति\-स्त्री\-बा\-ल\-वृ\-द्ध\-भी\-रु\-दु\-र्ब\-ल\-वै\-द्य\-वि\-द\-ग्ध\-व्या\-धि\-गु\-ह\-क\-द\-रि\-द्र\-कृ\-प\-ण\-क्रो\-ध\-ना\-ना\-त्म\-व\-तात्}}{\lemma{\skt{श्रोत्रि\appbrk{}य\appbrk{}नृ\appbrk{}प\appbrk{}ति\appbrk{}स्त्री\appbrk{}बा\appbrk{}ल\appbrk{}वृ\appbrk{}द्ध\appbrk{}भी\appbrk{}रु\appbrk{}दु\appbrk{}र्ब\appbrk{}ल\appbrk{}वै\appbrk{}द्य\appbrk{}वि\appbrk{}द\appbrk{}ग्ध\appbrk{}व्या\appbrk{}धि\appbrk{}गु\appbrk{}ह\appbrk{}क\appbrk{}द\appbrk{}रि\appbrk{}द्र\appbrk{}कृ\appbrk{}प\appbrk{}ण\appbrk{}क्रो\appbrk{}ध\appbrk{}ना\appbrk{}ना\appbrk{}त्म\appbrk{}व\appbrk{}तात्}}\Afootnote{\skt{श्रोत्रि\appbrk{}य\appbrk{}नृ\appbrk{}प\appbrk{}ति\appbrk{}स्त्री\appbrk{}बा\appbrk{}ल\appbrk{}वृ\appbrk{}द्ध\appbrk{}भी\appbrk{}रु\appbrk{}रा\appbrk{}ज\appbrk{}से\appbrk{}व\appbrk{}क\appbrk{}कि\appbrk{}त\appbrk{}व\appbrk{}दु\appbrk{}र्ब\appbrk{}ल\appbrk{}वै\appbrk{}द्य\appbrk{}वि\appbrk{}द\appbrk{}ग्ध\appbrk{}व्या\appbrk{}धि\appbrk{}गो\appbrk{}प\appbrk{}क\appbrk{}द\appbrk{}रि\appbrk{}द्र\appbrk{}कृ\appbrk{}प\appbrk{}ण\appbrk{}क्रो\appbrk{}ध\appbrk{}ना\appbrk{}ना\appbrk{}म\appbrk{}ना\appbrk{}त्म\appbrk{}व\appbrk{}ता\appbrk{}म\appbrk{}ना\appbrk{}थानां च; एवं निरूप्य चिकित्सां कुर्व\appbrk{}न्ध\appbrk{}र्मा\appbrk{}र्थ\appbrk{}का\appbrk{}म\appbrk{}य\appbrk{}शांसि प्राप्नोति}\ A; \skt{श्रोत्रि\appbrk{}य\appbrk{}नृ\appbrk{}प\appbrk{}ति\appbrk{}स्त्री\appbrk{}बा\appbrk{}ल\appbrk{}वृ\appbrk{}द्ध\appbrk{}भी\appbrk{}रु\appbrk{}दु\appbrk{}र्ब\appbrk{}ल\appbrk{}वै\appbrk{}द्य\appbrk{}वि\appbrk{}द\appbrk{}ग्ध\appbrk{}व्या\appbrk{}धि\appbrk{}गृ\appbrk{}ह\appbrk{}क\appbrk{}द\appbrk{}रि\appbrk{}द्र\appbrk{}कृ\appbrk{}प\appbrk{}ण\appbrk{}क्रो\appbrk{}ध\appbrk{}ना\appbrk{}ना\appbrk{}त्म\appbrk{}व\appbrk{}तात्}\ K; \skt{स्रोत्रि\appbrk{}य\appbrk{}नृ\appbrk{}प\appbrk{}ति\appbrk{}स्त्री\appbrk{}बा\appbrk{}ल\appbrk{}वृ\appbrk{}द्ध\appbrk{}भी\appbrk{}रु\appbrk{}दु\appbrk{}र्ब्ब\appbrk{}ल\appbrk{}वै\appbrk{}द्य\appbrk{}विं\appbrk{}द\appbrk{}ग्ध\appbrk{}व्या\appbrk{}धि\appbrk{}गू\appbrk{}ह\appbrk{}क\appbrk{}द\appbrk{}रि\appbrk{}द्र\appbrk{}कृ\appbrk{}प\appbrk{}ण\appbrk{}क्रो\appbrk{}ध\appbrk{}ना\appbrk{}त्म\appbrk{}वतां}\ N; \skt{श्रोत्रि\appbrk{}य\appbrk{}नृ\appbrk{}प\appbrk{}ति\appbrk{}स्त्री\appbrk{}बा\appbrk{}ल\appbrk{}वृ\appbrk{}द्ध\appbrk{}भी\appbrk{}रु\appbrk{}दु\appbrk{}र्ब्ब\appbrk{}ल\appbrk{}वै\appbrk{}द्य\appbrk{}वि\appbrk{}द\appbrk{}ग्ध\appbrk{}व्या\appbrk{}धि\appbrk{}गू\appbrk{}ह\appbrk{}क\appbrk{}द\appbrk{}रि\appbrk{}द्र\appbrk{}कृ\appbrk{}प\appbrk{}ण\appbrk{}क्रो\appbrk{}ध\appbrk{}ना\appbrk{}ना\appbrk{}त्म\appbrk{}व\appbrk{}तात्}\ H}}\skt{~॥}\pend

\pstart\versepar \ledsidenote{\passref{1.10.9}}\edtext{\skt{भवन्ति}}{\lemma{\skt{भवन्ति}}\Afootnote{K N; \skt{भवति}\ A H}}\skt{ चात्र~॥ }\pend
\pstart\versepar \skt{स्त्रीभिः }\edtext{\skt{सहास्यं}}{\lemma{\skt{सहास्यं}}\Afootnote{K; \skt{सहास्यां}\ A N; \skt{सहाम्यं}\ H\textsuperscript{ac}; \skt{ञ्सहास्यं}\ H\textsuperscript{pc}}}\skt{ }\edtext{\skt{संवास}}{\lemma{\skt{संवास}}\Afootnote{K; \skt{संवासं}\ A N; \skt{संवादं}\ H}}\skt{ परिहासञ्च वर्जयेत्~। }\pend
\pstart\versepar \skt{दत्तं }\edtext{\skt{ताभिर्न}}{\lemma{\skt{ताभिर्न}}\Afootnote{K N H; \skt{ताभ्यो ना}\ A}}\skt{ }\edtext{\skt{गृह्णी\-या\-द\-न्ना\-द\-न्य\-द्भिष देति}}{\lemma{\skt{गृह्णी\appbrk{}या\appbrk{}द\appbrk{}न्ना\appbrk{}द\appbrk{}न्य\appbrk{}द्भिष देति}}\Afootnote{\skt{देय\appbrk{}म\appbrk{}न्ना\appbrk{}द\appbrk{}न्य\appbrk{}द्भि\appbrk{}ष\appbrk{}ग्वरैः}\ A; \skt{गृह्णी\appbrk{}या\appbrk{}द\appbrk{}न्ना\appbrk{}द\appbrk{}न्य\appbrk{}द्भि\appbrk{}ष\appbrk{}देति~॥ {\symfont ❈}~॥ १०}\ K; \skt{मृह्णी\appbrk{}या\appbrk{}द\appbrk{}न्या\appbrk{}द\appbrk{}न्य\appbrk{}द्भि\appbrk{}ष\appbrk{}ग्भ्यु\appbrk{}देति~॥ डो}\ N; \skt{गृह्णी\appbrk{}या\appbrk{}द\appbrk{}न्ना\appbrk{}द\appbrk{}न्य\appbrk{}द्भि\appbrk{}ष\appbrk{}देति~॥ डो}\ H}}\skt{~॥}\edtext{}{\lemma{\skt{वर्जयेत्~। दत्तं}}\Afootnote{K N H; \skt{वर्जयेत्~। दत्तं च}\ A}}\pend

\pstart\versepar \ledsidenote{\passref{1.10.9.add}}\edtext{\skt{वेदो\-त्प\-त्ति\-शि\-ष्य\-दी\-क्षा\-दा\-न\-म\-ध्य\-य\-नस्य}}{\lemma{\skt{वेदो\appbrk{}त्प\appbrk{}त्ति\appbrk{}शि\appbrk{}ष्य\appbrk{}दी\appbrk{}क्षा\appbrk{}दा\appbrk{}नम्}}\Afootnote{K H; \skt{वेदो\appbrk{}त्प\appbrk{}त्ति\appbrk{}शि\appbrk{}ष्या\appbrk{}दी\appbrk{}क्षा\appbrk{}दा\appbrk{}नम्}\ N}}\skt{ च~। }\pend
\pstart\versepar \skt{}\edtext{\skt{प्रभा\-ष\-ण\-ञ्चा\-ग्र\-हरं}}{\lemma{\skt{प्रभाषणञ्}}\Afootnote{K H; \skt{प्रभाषणश्}\ N}}\skt{ }\edtext{\skt{ऋतुचर्या}}{\lemma{\skt{ऋतुचर्या}}\Afootnote{K\textsuperscript{pc} N; \skt{णं ऋतुचर्या}\ K\textsuperscript{ac}; \skt{मृतुचर्या}\ H}}\skt{ तथैव च~। }\pend
\pstart\versepar \skt{यन्त्रं शस्त्रावचारञ्च योग्या\-सू\-त्री\-य\-मेव च~। }\pend
\pstart\versepar \skt{}\edtext{\skt{विशि\-खा\-नु\-प्र\-वे\-शञ्च}}{\lemma{\skt{विशि\appbrk{}खा\appbrk{}नु\appbrk{}प्र\appbrk{}वे\appbrk{}शञ्च}}\Afootnote{K H; \skt{विशि\appbrk{}क्षा\appbrk{}नु\appbrk{}प्र\appbrk{}वे\appbrk{}शश्च}\ N}}\skt{ प्रोक्तं वै प्रथमो दश~॥}\pend

\pstart \ledsidenote{\passref{1.10.trailer}}\edtext{\skt{॥ १०~॥}}{\lemma{\skt{॥ १०~॥}}\Afootnote{\skt{इति सुश्रु\appbrk{}त\appbrk{}सं\appbrk{}हि\appbrk{}तायां सूत्रस्थाने विशि\appbrk{}खा\appbrk{}नु\appbrk{}प्र\appbrk{}वे\appbrk{}श\appbrk{}नीयो नाम दशमोऽध्यायः}\ A; \skt{{\symfont ❈}}\ K; \skt{डो}\ H}}\pend

\endnumbering
\beginnumbering
\eledsection*{Adhyāya 11}
\pstart \ledsidenote{\passref{1.11.1}}\skt{अथातः }\edtext{\skt{क्षार\-पा\-क\-वि\-धि\-म\-ध्यायं}}{\lemma{\skt{क्षार\appbrk{}पा\appbrk{}क\appbrk{}वि\appbrk{}धि\appbrk{}म\appbrk{}ध्यायं}}\Afootnote{K; \skt{क्षार\appbrk{}पा\appbrk{}क\appbrk{}वि\appbrk{}धि\appbrk{}म\appbrk{}ध्यायं व्याख्यास्यामः}\ A H; \skt{स्वार\appbrk{}पा\appbrk{}क\appbrk{}वि\appbrk{}धि\appbrk{}म\appbrk{}ध्यायं व्याख्यास्यामः}\ N}}\skt{~॥}\pend

\pstart \ledsidenote{\passref{1.11.3}}\edtext{\skt{अनुशस्त्रेभ्यः}}{\lemma{\skt{अनुशस्त्रेभ्यः}}\Afootnote{K H; \skt{शस्त्रानुशस्त्रेभ्यः}\ A; \skt{अनुशस्त्रेभ्य}\ N}}\skt{ }\edtext{\skt{क्षारः}}{\lemma{\skt{क्षारः}}\Afootnote{A K H; \skt{क्षारं}\ N}}\skt{ प्रधानतमो }\edtext{\skt{भवति}}{\lemma{\skt{भवति}}\Afootnote{K N H; \emph{om.}\ A}}\skt{~। }\edtext{\skt{च्छेद्य\-भे\-द्य\-ले\-ख्य\-क\-र\-णा\-द्वि\-शे\-ष\-क्रि\-या\-व\-चा\-र\-णाच्च}}{\lemma{\skt{च्छेद्य\appbrk{}भे\appbrk{}द्य\appbrk{}ले\appbrk{}ख्य\appbrk{}क\appbrk{}र\appbrk{}णा\appbrk{}द्वि\appbrk{}शे\appbrk{}ष\appbrk{}क्रि\appbrk{}या\appbrk{}व\appbrk{}चा\appbrk{}र\appbrk{}णाच्च}}\Afootnote{K; \skt{छेद्य\appbrk{}भे\appbrk{}द्य\appbrk{}ले\appbrk{}ख्य\appbrk{}क\appbrk{}र\appbrk{}णा\appbrk{}त्त्रि\appbrk{}दो\appbrk{}ष\appbrk{}घ्न\appbrk{}त्वा\appbrk{}द्वि\appbrk{}शे\appbrk{}ष\appbrk{}क्रि\appbrk{}या\appbrk{}व\appbrk{}चा\appbrk{}र\appbrk{}णाच्च}\ A; \skt{छेद्य\appbrk{}भे\appbrk{}द्य\appbrk{}ले\appbrk{}ख्य\appbrk{}क\appbrk{}र\appbrk{}णात्~॥ विशे\appbrk{}ष\appbrk{}क्रि\appbrk{}या\appbrk{}व\appbrk{}च\appbrk{}ण्णाच्चं}\ N; \skt{च्छेद्य\appbrk{}भे\appbrk{}द्य\appbrk{}ले\appbrk{}ख्य\appbrk{}क\appbrk{}र\appbrk{}णा\appbrk{}द्वि\appbrk{}शे\appbrk{}ष\appbrk{}क्रि\appbrk{}या\appbrk{}व\appbrk{}चा\appbrk{}र\appbrk{}णाच्च}\ H}}\skt{~।}\pend

\pstart \ledsidenote{\passref{1.11.4}}\skt{तत्र }\edtext{\skt{क्षर\-णा\-त्क्ष\-ण\-णाद्वा}}{\lemma{\skt{क्षर\appbrk{}णा\appbrk{}त्क्ष\appbrk{}ण\appbrk{}णाद्वा}}\Afootnote{K; \skt{क्षर\appbrk{}णा\appbrk{}त्क्ष\appbrk{}ण\appbrk{}नाद्वा}\ A N; \skt{क्षर\appbrk{}णा\appbrk{}त्क्ष\appbrk{}ण\appbrk{}नाद्वा}\ H}}\skt{ क्षारः~।}\pend

\pstart \ledsidenote{\passref{1.11.5}}\edtext{\skt{नानौ\-ष\-धि\-स\-म\-वा\-या\-त्त्रि\-दो\-षघ्नः}}{\lemma{\skt{नानौ\appbrk{}ष\appbrk{}धि\appbrk{}स\appbrk{}म\appbrk{}वा\appbrk{}या\appbrk{}त्त्रि\appbrk{}दो\appbrk{}षघ्नः}}\Afootnote{K H; \skt{नानौ\appbrk{}ष\appbrk{}धि\appbrk{}स\appbrk{}म\appbrk{}वा\appbrk{}या\appbrk{}त्त्रि\appbrk{}दो\appbrk{}षघ्नः}\ A; \skt{नानौ\appbrk{}ष\appbrk{}धि\appbrk{}स\appbrk{}म\appbrk{}वा\appbrk{}या\appbrk{}न्त्रि\appbrk{}दो\appbrk{}षघ्नः}\ N}}\skt{~। }\edtext{\edtext{\skt{शुक्ल\-त्वा\-त्सौ\-म्य\-स्तस्य}}{\lemma{\skt{शुक्लत्वात्सौम्यस्}}\Afootnote{A K; \skt{शुक्लत्वात्सौम्य}\ N; \skt{शुक्लत्वा… त्सौम्यस्}\ H}}}{\lemma{\skt{तस्य}}\Afootnote{A K H; \skt{स्तस्य}\ N}}\skt{ सौम्यस्यापि }\edtext{\skt{सतो}}{\lemma{\skt{सतो}}\Afootnote{A K H; \skt{स}\ N}}\skt{ }\edtext{\skt{दह\-न\-प\-च\-न\-दा\-र\-ण\-श\-क्ति\-र\-वि\-रुद्धा}}{\lemma{\skt{दह\appbrk{}न\appbrk{}प\appbrk{}च\appbrk{}न\appbrk{}दा\appbrk{}र\appbrk{}ण\appbrk{}श\appbrk{}क्ति\appbrk{}र\appbrk{}वि\appbrk{}रुद्धा}}\Afootnote{K; \skt{दह\appbrk{}न\appbrk{}प\appbrk{}च\appbrk{}न\appbrk{}दा\appbrk{}र\appbrk{}णा\appbrk{}दि\appbrk{}श\appbrk{}क्ति\appbrk{}र\appbrk{}वि\appbrk{}रुद्धा}\ A; \skt{त्वोद\appbrk{}ह\appbrk{}न\appbrk{}प\appbrk{}च\appbrk{}न\appbrk{}दा\appbrk{}र\appbrk{}ण\appbrk{}श\appbrk{}क्ति\appbrk{}रि\appbrk{}रु\appbrk{}रुद्धा}\ N; \skt{दद\appbrk{}ह\appbrk{}न\appbrk{}प\appbrk{}च\appbrk{}न\appbrk{}दा\appbrk{}र\appbrk{}ण\appbrk{}श\appbrk{}क्ति\appbrk{}र\appbrk{}वि\appbrk{}रुद्धा}\ H}}\skt{~। }\edtext{\skt{सख\-ल्वा\-ग्ने\-यौ\-ष\-धि\-भूय कः}}{\lemma{\skt{सख\appbrk{}ल्वा\appbrk{}ग्ने\appbrk{}यौ\appbrk{}ष\appbrk{}धि\appbrk{}भूय कः}}\Afootnote{K; \skt{आग्ने\appbrk{}यौ\appbrk{}ष\appbrk{}धि\appbrk{}गु\appbrk{}ण\appbrk{}भू\appbrk{}यि\appbrk{}ष्ठ\appbrk{}त्वा\appbrk{}त्क\appbrk{}टुक}\ A; \skt{स खल्वा\appbrk{}र्मे\appbrk{}यौ\appbrk{}ष\appbrk{}धि\appbrk{}भू\appbrk{}य\appbrk{}स्त्वा\appbrk{}त्क\appbrk{}ठ्रकं}\ N; \skt{स खल्वा\appbrk{}ग्ने\appbrk{}यौ\appbrk{}ष\appbrk{}धि\appbrk{}भू\appbrk{}यि\appbrk{}ष्ठा\appbrk{}त्क\appbrk{}टुकः}\ H}}\skt{~। उष्णस्तीक्ष्णः पाचनो }\edtext{\skt{विलायनः}}{\lemma{\skt{विलायनः}}\Afootnote{K N H; \skt{विलयनः}\ A}}\skt{~। शोधनो }\edtext{\skt{रोपणः}}{\lemma{\skt{रोपणः}}\Afootnote{K; \skt{रोपणः शोषणः}\ A H; \skt{रोषणः}\ N}}\skt{ }\edtext{\skt{स्तम्भ\-नो\-ल्ले\-ख\-न\-क्रि\-म्या\-म\-क\-फ\-वि\-ष\-मे\-द\-सा\-मु\-प\-हन्ता}}{\lemma{\skt{स्तम्भ\appbrk{}नो\appbrk{}ल्ले\appbrk{}ख\appbrk{}न\appbrk{}क्रि\appbrk{}म्या\appbrk{}म\appbrk{}क\appbrk{}फ\appbrk{}वि\appbrk{}ष\appbrk{}मे\appbrk{}द\appbrk{}साम्}}\Afootnote{K H; \skt{स्तम्भनो लेखनः कृम्या\appbrk{}म\appbrk{}क\appbrk{}फ\appbrk{}कु\appbrk{}ष्ठ\appbrk{}वि\appbrk{}ष\appbrk{}मे\appbrk{}द\appbrk{}साम्}\ A; \skt{स्तम्भ\appbrk{}नों\appbrk{}ले\appbrk{}ख\appbrk{}न\appbrk{}क्रि\appbrk{}म्या\appbrk{}म\appbrk{}क\appbrk{}ह\appbrk{}वि\appbrk{}ष\appbrk{}मे\appbrk{}द\appbrk{}साम्}\ N}}\skt{ }\edtext{\skt{पुंस्त्वस्य}}{\lemma{\skt{पुंस्त्वस्य}}\Afootnote{A K H; \skt{पुस्त्वस्य}\ N}}\skt{ चातिसेवितः~।}\pend

\pstart \ledsidenote{\passref{1.11.6}}\skt{स द्विविधः प्रतिसारणीयः~। पानीयश्च}\pend

\pstart \ledsidenote{\passref{1.11.7}}\skt{तत्र प्रतिसारणीयः~। }\edtext{\skt{कुष्ठ\-कि\-टि\-भ\-द\-र्द्रु\-म\-ण्ड\-ल\-कि\-ला\-स\-भ\-ग\-न्द\-रा\-र्शो\-र्बु\-द\-दु\-ष्ट\-व्र\-ण\-ना\-डी\-च\-र्म\-की\-ल\-ति\-ल\-का\-ल\-क\-न\-च्छ\-व्य\-ङ्ग\-बा\-ह्य\-क्रि\-मि\-वि\-षा\-दिषु}}{\lemma{\skt{कुष्ठ\appbrk{}कि\appbrk{}टि\appbrk{}भ\appbrk{}द\appbrk{}र्द्रु\appbrk{}म\appbrk{}ण्ड\appbrk{}ल\appbrk{}कि\appbrk{}ला\appbrk{}स\appbrk{}भ\appbrk{}ग\appbrk{}न्द\appbrk{}रा\appbrk{}र्शो\appbrk{}र्बु\appbrk{}द\appbrk{}दु\appbrk{}ष्ट\appbrk{}व्र\appbrk{}ण\appbrk{}ना\appbrk{}डी\appbrk{}च\appbrk{}र्म\appbrk{}की\appbrk{}ल\appbrk{}ति\appbrk{}ल\appbrk{}का\appbrk{}ल\appbrk{}क\appbrk{}न\appbrk{}च्छ\appbrk{}व्य\appbrk{}ङ्ग\appbrk{}बा\appbrk{}ह्य\appbrk{}क्रि\appbrk{}मि\appbrk{}वि\appbrk{}षा\appbrk{}दिषु}}\Afootnote{K H; \skt{कुष्ठ\appbrk{}कि\appbrk{}टि\appbrk{}भ\appbrk{}द\appbrk{}द्रु\appbrk{}म\appbrk{}ण्ड\appbrk{}ल\appbrk{}कि\appbrk{}ला\appbrk{}स\appbrk{}भ\appbrk{}ग\appbrk{}न्द\appbrk{}रा\appbrk{}र्बु\appbrk{}दा\appbrk{}र्शो\appbrk{}दु\appbrk{}ष्ट\appbrk{}व्र\appbrk{}ण\appbrk{}ना\appbrk{}डी\appbrk{}च\appbrk{}र्म\appbrk{}कील- तिल\appbrk{}का\appbrk{}ल\appbrk{}क\appbrk{}न्य\appbrk{}च्छ\appbrk{}व्य\appbrk{}ङ्ग\appbrk{}म\appbrk{}श\appbrk{}क\appbrk{}बा\appbrk{}ह्य\appbrk{}वि\appbrk{}द्र\appbrk{}धि\appbrk{}कृ\appbrk{}मि\appbrk{}वि\appbrk{}षा\appbrk{}दिषू}\ A; \skt{कुष्ट\appbrk{}कि\appbrk{}टि\appbrk{}भ\appbrk{}द\appbrk{}द्रु\appbrk{}म\appbrk{}ण्ड\appbrk{}ल\appbrk{}कि\appbrk{}ला\appbrk{}स\appbrk{}भ\appbrk{}ग\appbrk{}न्द\appbrk{}रा\appbrk{}र्सो\appbrk{}र्बु\appbrk{}द\appbrk{}दु\appbrk{}ष्ट\appbrk{}व्र\appbrk{}ण\appbrk{}ना\appbrk{}डी\appbrk{}क\appbrk{}र्म\appbrk{}की\appbrk{}ल\appbrk{}ति\appbrk{}ल\appbrk{}का\appbrk{}ल\appbrk{}क\appbrk{}न\appbrk{}छ\appbrk{}व्य\appbrk{}ङ्ग\appbrk{}बा\appbrk{}हु\appbrk{}कृ\appbrk{}मि\appbrk{}वि\appbrk{}षा\appbrk{}दिषु}\ N}}\skt{ }\edtext{\skt{चोपदिश्यते}}{\lemma{\skt{चोपदिश्यते}}\Afootnote{K N H; \skt{पदिश्यते}\ A}}\skt{~। सप्तसु च }\edtext{\skt{मुख\-रो\-गे\-षू\-प\-जि\-ह्वो\-प\-कु\-श\-द\-न्त\-वै\-द\-र्भ\-मे\-द\-जो\-ष्ठ\-प्र\-को\-पेषु}}{\lemma{\skt{मुख\appbrk{}रो\appbrk{}गे\appbrk{}षू\appbrk{}प\appbrk{}जि\appbrk{}ह्वो\appbrk{}प\appbrk{}कु\appbrk{}श\appbrk{}द\appbrk{}न्त\appbrk{}वै\appbrk{}द\appbrk{}र्भ\appbrk{}मे\appbrk{}द\appbrk{}जो\appbrk{}ष्ठ\appbrk{}प्र\appbrk{}को\appbrk{}पेषु}}\Afootnote{K H; \skt{मुख\appbrk{}रो\appbrk{}गे\appbrk{}षू\appbrk{}प\appbrk{}जि\appbrk{}ह्वा\appbrk{}धि\appbrk{}जि\appbrk{}ह्वो\appbrk{}प\appbrk{}कु\appbrk{}श\appbrk{}द\appbrk{}न्त\appbrk{}वै\appbrk{}द\appbrk{}र्भेषु}\ A; \skt{मुख\appbrk{}रो\appbrk{}गे\appbrk{}सू\appbrk{}प\appbrk{}जि\appbrk{}ह्वो\appbrk{}प\appbrk{}कु\appbrk{}श\appbrk{}द\appbrk{}न्त\appbrk{}वै\appbrk{}द\appbrk{}र्भ\appbrk{}मे\appbrk{}द\appbrk{}जे\appbrk{}ष्ट\appbrk{}प्र\appbrk{}को\appbrk{}पेषु}\ N}}\skt{ }\edtext{\skt{त्रिषु च}}{\lemma{\skt{त्रिषु च}}\Afootnote{\skt{तिसृषु च}\ A; \skt{त्रिसृषु च}\ K H; \skt{तिसृसुष च}\ N}}\skt{ }\edtext{\skt{रोहिणीषु}}{\lemma{\skt{रोहिणीषु}}\Afootnote{A K H; \skt{रोहणीषु}\ N}}\skt{ }\edtext{\edtext{\skt{एते\-ष्वे\-वा\-नु\-श\-स्त्र\-पा\-त\-न\-मुक्तं}}{\lemma{\skt{एतेष्व्}}\Afootnote{A K N; \skt{एते}\ H}}}{\lemma{\skt{एवा\appbrk{}नु\appbrk{}श\appbrk{}स्त्र\appbrk{}पा\appbrk{}त\appbrk{}नम्}}\Afootnote{K N; \skt{एवा\appbrk{}नु\appbrk{}श\appbrk{}स्त्र\appbrk{}प्र\appbrk{}णि\appbrk{}धा\appbrk{}नम्}\ A; \skt{ष्वेवा\appbrk{}नु\appbrk{}श\appbrk{}स्त्र\appbrk{}पा\appbrk{}त\appbrk{}नम्}\ H}}\skt{~॥}\footnoteA{Ḍalhaṇa reported that some people read \emph{evānuśastrapātanam} as here, as opposed to the vulgate's \emph{anuśastrapraṇidhānam}.}\pend

\pstart \ledsidenote{\passref{1.11.8}}\edtext{\skt{पानीयस्तु}}{\lemma{\skt{पानीयस्तु}}\Afootnote{A K N; \skt{पानी}\ H}}\skt{ }\edtext{\skt{गुल्मा\-द\-रा\-ग्नि\-सं\-गा\-जी\-र्ण्णा\-ना\-ह\-श\-र्क\-रा\-श्म\-र्य\-भ्य\-न्त\-र\-क्रि\-मि\-वि\-षा\-शस्सु}}{\lemma{\skt{गुल्मा\appbrk{}द\appbrk{}रा\appbrk{}ग्नि\appbrk{}सं\appbrk{}गा\appbrk{}जी\appbrk{}र्ण्णा\appbrk{}ना\appbrk{}ह\appbrk{}श\appbrk{}र्क\appbrk{}रा\appbrk{}श्म\appbrk{}र्य\appbrk{}भ्य\appbrk{}न्त\appbrk{}र\appbrk{}क्रि\appbrk{}मि\appbrk{}वि\appbrk{}षा\appbrk{}शस्सु}}\Afootnote{K; \skt{गर\appbrk{}गु\appbrk{}ल्मो\appbrk{}द\appbrk{}रा\appbrk{}ग्नि\appbrk{}स\appbrk{}ङ्गा\appbrk{}जी\appbrk{}र्णा\appbrk{}रो\appbrk{}च\appbrk{}का\appbrk{}नाह- शर्क\appbrk{}रा\appbrk{}श्म\appbrk{}र्या\appbrk{}भ्य\appbrk{}न्त\appbrk{}र\appbrk{}वि\appbrk{}द्र\appbrk{}धि\appbrk{}कृ\appbrk{}मि\appbrk{}वि\appbrk{}षा\appbrk{}र्शःसू}\ A; \skt{गुल्मो\appbrk{}द\appbrk{}रा\appbrk{}ग्नि\appbrk{}सं\appbrk{}शा\appbrk{}जी\appbrk{}र्ण्णा\appbrk{}ना\appbrk{}ह\appbrk{}स\appbrk{}र्क्क\appbrk{}रा\appbrk{}स्म\appbrk{}र्ग्यी\appbrk{}र्य\appbrk{}भ\appbrk{}द\appbrk{}न्त\appbrk{}र\appbrk{}क्रि\appbrk{}मि\appbrk{}वि\appbrk{}षा\appbrk{}र्सस्सु}\ N; \skt{यस्तु गुल्मो\appbrk{}द\appbrk{}रा\appbrk{}ग्नि\appbrk{}स\appbrk{}ङ्गा\appbrk{}जी\appbrk{}र्ण्णा\appbrk{}ना\appbrk{}ह\appbrk{}श\appbrk{}र्क्क\appbrk{}रा\appbrk{}श्म\appbrk{}र्य\appbrk{}भ्य\appbrk{}न्त\appbrk{}र\appbrk{}क्रि\appbrk{}मि\appbrk{}वि\appbrk{}षा\appbrk{}र्शःमू}\ H\textsuperscript{ac}; \skt{यस्तु गुल्मो\appbrk{}द\appbrk{}रा\appbrk{}ग्नि\appbrk{}स\appbrk{}ङ्गा\appbrk{}जी\appbrk{}र्ण्णा\appbrk{}ना\appbrk{}ह\appbrk{}श\appbrk{}र्क्क\appbrk{}रा\appbrk{}श्म\appbrk{}र्य\appbrk{}भ्य\appbrk{}न्त\appbrk{}र\appbrk{}क्रि\appbrk{}मि\appbrk{}वि\appbrk{}षा\appbrk{}र्शःसू}\ H\textsuperscript{pc}}}\skt{ }\edtext{\skt{चोपयुज्यते}}{\lemma{\skt{चोपयुज्यते}}\Afootnote{K N H; \skt{पयुज्यते}\ A}}\skt{~॥}\pend

\pstart \ledsidenote{\passref{1.11.10}}\edtext{\skt{तस्य}}{\lemma{\skt{तस्य}}\Afootnote{K H; \skt{तं चेत\appbrk{}र\appbrk{}क्षा\appbrk{}र\appbrk{}व\appbrk{}द्दग्ध्वा परिस्रावयेत्; तस्य}\ A; \skt{तस्य विस्तरो न्यता\appbrk{}र्थे\appbrk{}त\appbrk{}र\appbrk{}चि\appbrk{}कीर्षुः}\ N}}\skt{ }\edtext{\edtext{\skt{विस्तारोऽन्यत्र}}{\lemma{\skt{विस्तारो}}\Afootnote{K; \skt{विस्तरो}\ A; \emph{om.}\ N; \skt{विस्तरोऽ}\ H}}}{\lemma{\skt{ऽन्यत्र}}\Afootnote{A; \skt{न्यत्र}\ K H; \emph{om.}\ N}}\pend

\pstart \ledsidenote{\passref{1.11.11}}\edtext{\skt{अथेतरं}}{\lemma{\skt{अथेतरं}}\Afootnote{\skt{अथेतरस्}\ A; \emph{om.}\ K N; \skt{अर्थेतर}\ H}}\skt{ }\edtext{\skt{चिकीर्षुः}}{\lemma{\skt{चिकीर्षुः}}\Afootnote{H; \skt{त्रिविधो मृदु\appbrk{}र्म\appbrk{}ध्य\appbrk{}स्ती\appbrk{}क्ष्णश्च~। तं चिकीर्षुः}\ A; \skt{अर्थतर चिकीर्षुः}\ K; \emph{om.}\ N}}\skt{~। }\edtext{\skt{शरदि}}{\lemma{\skt{शरदि}}\Afootnote{A K H; \emph{om.}\ N}}\skt{ }\edtext{\edtext{\skt{शुचि\-रु\-प\-व\-स\-न्प्र\-श\-स्त\-दे\-श\-जा\-त\-म\-नु\-प\-हतं}}{\lemma{\skt{शुचिर्}}\Afootnote{K H; \skt{गिरिसानुजं शुचिर्}\ A; \skt{शरद शुचिर्}\ N}}}{\lemma{\skt{उप\appbrk{}व\appbrk{}स\appbrk{}न्प्र\appbrk{}श\appbrk{}स्त\appbrk{}दे\appbrk{}श\appbrk{}जा\appbrk{}तम्}}\Afootnote{K N H; \skt{उपोष्य प्रशस्तेऽहनि प्रश\appbrk{}स्त\appbrk{}दे\appbrk{}श\appbrk{}जा\appbrk{}तम्}\ A}}\skt{ मध्यमवयसं }\edtext{\edtext{\skt{काल\-मु\-ष्क\-क\-म\-धि\-वा\-स्या\-प\-रेद्युः}}{\lemma{\skt{कालमुष्ककम्}}\Afootnote{K H; \skt{महा\appbrk{}न्त\appbrk{}म\appbrk{}सि\appbrk{}त\appbrk{}मु\appbrk{}ष्क\appbrk{}कम्}\ A; \skt{कालमुस्ककम}\ N}}}{\lemma{\skt{अधि\appbrk{}वा\appbrk{}स्या\appbrk{}प\appbrk{}रेद्युः}}\Afootnote{A K H; \skt{धिवास्यापरेद्युः}\ N}}\skt{ }\edtext{\skt{पातयित्वा}}{\lemma{\skt{पातयित्वा}}\Afootnote{K H; \skt{पाटयित्वा}\ A N}}\skt{ }\edtext{\skt{कण्डसः}}{\lemma{\skt{कण्डसः}}\Afootnote{K; \skt{खण्डशः}\ A; \skt{काण्डशः}\ N; \skt{काण्डसः}\ H}}\skt{ }\edtext{\skt{प्रकल्प्य}}{\lemma{\skt{प्रकल्प्य}}\Afootnote{K H; \skt{प्रकल्प्या}\ A N}}\skt{ }\edtext{\skt{निवातदेशे}}{\lemma{\skt{निवातदेशे}}\Afootnote{H; \skt{वपाट्य निवाते देशे}\ A; \skt{निवात}\ K N}}\skt{ }\edtext{\edtext{\skt{चितिङ्कृत्वा}}{\lemma{\skt{चितिङ्}}\Afootnote{H; \skt{निचितिं}\ A; \emph{om.}\ K N}}\skt{ तिल\-ना\-लै\-रा\-दी\-प\-येत्}}{\lemma{\skt{कृत्वा तिलनालैर्}}\Afootnote{H; \skt{कृत्वा सुधाशर्कराश्च प्रक्षिप्य तिलनालैर्}\ A; \skt{देशे चितिङ्कृत्वा तिलनालैर्}\ K; \skt{दोशे चितिं कृत्वा तिलनालैर्}\ N}}\skt{~। अथोपशान्तेऽग्नौ }\edtext{\skt{तद्भस्म}}{\lemma{\skt{तद्भस्म}}\Afootnote{K N; \skt{तद्भस्म}\ A H}}\skt{ }\edtext{\skt{पृथ\-ग्गृ\-ह्णी\-या\-त्भ\-स्म\-श\-र्क\-राश्च}}{\lemma{\skt{पृथ\appbrk{}ग्गृ\appbrk{}ह्णी\appbrk{}या\appbrk{}त्भ\appbrk{}स्म\appbrk{}श\appbrk{}र्क\appbrk{}राश्च}}\Afootnote{K N H; \skt{पृथ\appbrk{}ग्गृ\appbrk{}ह्णी\appbrk{}या\appbrk{}द्भ\appbrk{}स्म\appbrk{}श\appbrk{}र्क\appbrk{}राश्च}\ A}}\skt{~॥ अथानेनैव }\edtext{\skt{कल्पेन}}{\lemma{\skt{कल्पेन}}\Afootnote{K H; \skt{विधानेन}\ A; \skt{कर्ष्केन्}\ N}}\skt{ }\edtext{\skt{कुट\-ज\-प\-ला\-शा\-श्व\-क\-र्ण्ण\-पा\-रि\-भ\-द्र\-क\-बि\-भी\-त\-का\-र\-ग्व\-ध\-ति\-ल्व\-का\-र्क\-स्नु\-ह्या\-पा\-मा\-र्ग\-न\-क्त\-मा\-ल\-वृ\-ष\-क\-द\-ली\-चि\-त्र\-के\-न्द्र\-य\-वृ\-क्षा\-स्फो\-ता\-श्व\-मा\-र\-क\-स\-प्त\-च्छ\-दा\-ग्नि\-मन्थाः}}{\lemma{\skt{कुट\appbrk{}ज\appbrk{}प\appbrk{}ला\appbrk{}शा\appbrk{}श्व\appbrk{}क\appbrk{}र्ण्ण\appbrk{}पा\appbrk{}रि\appbrk{}भ\appbrk{}द्र\appbrk{}क\appbrk{}बि\appbrk{}भी\appbrk{}त\appbrk{}का\appbrk{}र\appbrk{}ग्व\appbrk{}ध\appbrk{}ति\appbrk{}ल्व\appbrk{}का\appbrk{}र्क\appbrk{}स्नु\appbrk{}ह्या\appbrk{}पा\appbrk{}मा\appbrk{}र्ग\appbrk{}न\appbrk{}क्त\appbrk{}मा\appbrk{}ल\appbrk{}वृ\appbrk{}ष\appbrk{}क\appbrk{}द\appbrk{}ली\appbrk{}चि\appbrk{}त्र\appbrk{}के\appbrk{}न्द्र\appbrk{}य\appbrk{}वृ\appbrk{}क्षा\appbrk{}स्फो\appbrk{}ता\appbrk{}श्व\appbrk{}मा\appbrk{}र\appbrk{}क\appbrk{}स\appbrk{}प्त\appbrk{}च्छ\appbrk{}दा\appbrk{}ग्नि\appbrk{}मन्थाः}}\Afootnote{\skt{कुट\appbrk{}ज\appbrk{}प\appbrk{}ला\appbrk{}शा\appbrk{}श्व\appbrk{}क\appbrk{}र्ण\appbrk{}पा\appbrk{}रि\appbrk{}भ\appbrk{}द्र\appbrk{}क\appbrk{}बि\appbrk{}भी\appbrk{}त\appbrk{}का\appbrk{}र\appbrk{}ग्व\appbrk{}ध\appbrk{}ति\appbrk{}ल्व\appbrk{}का\appbrk{}र्क\appbrk{}स्नु\appbrk{}ह्य\appbrk{}पा\appbrk{}मा\appbrk{}र्ग\appbrk{}पा\appbrk{}ट\appbrk{}ला\appbrk{}न\appbrk{}क्त\appbrk{}मा\appbrk{}ल\appbrk{}वृ\appbrk{}ष\appbrk{}क\appbrk{}द\appbrk{}ली\appbrk{}चि\appbrk{}त्र\appbrk{}क\appbrk{}पू\appbrk{}ती\appbrk{}के\appbrk{}न्द्र\appbrk{}वृ\appbrk{}क्षा\appbrk{}स्फो\appbrk{}ता\appbrk{}श्व\appbrk{}मा\appbrk{}र\appbrk{}क\appbrk{}स\appbrk{}प्त\appbrk{}च्छ\appbrk{}दा\appbrk{}ग्नि\appbrk{}म\appbrk{}न्थ\appbrk{}गु\appbrk{}ञ्जाश्}\ A; \skt{कुट\appbrk{}ज\appbrk{}प\appbrk{}ला\appbrk{}शा\appbrk{}श्व\appbrk{}क\appbrk{}र्ण्ण\appbrk{}पा\appbrk{}रि\appbrk{}भ\appbrk{}द्र\appbrk{}क\appbrk{}बि\appbrk{}भी\appbrk{}त\appbrk{}का\appbrk{}र\appbrk{}ग्व\appbrk{}ध\appbrk{}ति\appbrk{}ल्व\appbrk{}का\appbrk{}र्क\appbrk{}स्नु\appbrk{}हा\appbrk{}पा\appbrk{}मा\appbrk{}र्ग\appbrk{}न\appbrk{}क्त\appbrk{}मा\appbrk{}ल\appbrk{}वृ\appbrk{}ष\appbrk{}क\appbrk{}द\appbrk{}ली\appbrk{}चि\appbrk{}त्र\appbrk{}के\appbrk{}न्द्र\appbrk{}वृ\appbrk{}क्षा\appbrk{}स्फो\appbrk{}ता\appbrk{}श्व\appbrk{}मा\appbrk{}र\appbrk{}क\appbrk{}स\appbrk{}प्त\appbrk{}च्छ\appbrk{}दा\appbrk{}ग्नि\appbrk{}मन्थाः}\ K\textsuperscript{ac}; \skt{कुट\appbrk{}ज\appbrk{}प\appbrk{}ला\appbrk{}शा\appbrk{}श्व\appbrk{}क\appbrk{}र्ण्ण\appbrk{}पा\appbrk{}रि\appbrk{}भ\appbrk{}द्र\appbrk{}क\appbrk{}बि\appbrk{}भी\appbrk{}त\appbrk{}का\appbrk{}र\appbrk{}ग्व\appbrk{}ध\appbrk{}ति\appbrk{}ल्व\appbrk{}का\appbrk{}र्क\appbrk{}स्नु\appbrk{}हा\appbrk{}पा\appbrk{}मा\appbrk{}र्ग\appbrk{}न\appbrk{}क्त\appbrk{}मा\appbrk{}ल\appbrk{}वृ\appbrk{}ष\appbrk{}क\appbrk{}द\appbrk{}ली\appbrk{}चि\appbrk{}त्र\appbrk{}के\appbrk{}न्द्र\appbrk{}य\appbrk{}वृ\appbrk{}क्षा\appbrk{}स्फो\appbrk{}ता\appbrk{}श्व\appbrk{}मा\appbrk{}र\appbrk{}क\appbrk{}स\appbrk{}प्त\appbrk{}च्छ\appbrk{}दा\appbrk{}ग्नि\appbrk{}मन्थाः}\ K\textsuperscript{pc}; \skt{कुट\appbrk{}ज\appbrk{}पा\appbrk{}ला\appbrk{}शाश्च कर्ण्ण\appbrk{}पा\appbrk{}रि\appbrk{}भ\appbrk{}द्र\appbrk{}क\appbrk{}वि\appbrk{}नी\appbrk{}त\appbrk{}का\appbrk{}र\appbrk{}ग्व\appbrk{}ध\appbrk{}ति\appbrk{}ण्व\appbrk{}का\appbrk{}र्क\appbrk{}स्नु\appbrk{}हा\appbrk{}या\appbrk{}मा\appbrk{}र्ज\appbrk{}न\appbrk{}क्त\appbrk{}मा\appbrk{}ला\appbrk{}वृ\appbrk{}ष\appbrk{}क\appbrk{}द\appbrk{}ली\appbrk{}चि\appbrk{}त्र\appbrk{}के\appbrk{}न्द्र\appbrk{}वृ\appbrk{}क्षा\appbrk{}स्फो\appbrk{}ता\appbrk{}स्व\appbrk{}मा\appbrk{}र\appbrk{}क\appbrk{}स\appbrk{}प्त\appbrk{}च्छ\appbrk{}दा\appbrk{}ग्नि\appbrk{}मन्थः}\ N; \skt{कुट\appbrk{}ज\appbrk{}प\appbrk{}ला\appbrk{}शा\appbrk{}श्व\appbrk{}क\appbrk{}र्ण्ण\appbrk{}पा\appbrk{}रि\appbrk{}भ\appbrk{}द्र\appbrk{}क\appbrk{}बि\appbrk{}भी\appbrk{}त\appbrk{}का\appbrk{}र\appbrk{}ग्व\appbrk{}ध\appbrk{}ति\appbrk{}ल्व\appbrk{}का\appbrk{}र्क्क\appbrk{}स्नु\appbrk{}हा\appbrk{}पा\appbrk{}मा\appbrk{}र्ग्ग\appbrk{}न\appbrk{}क्त\appbrk{}मा\appbrk{}ल\appbrk{}वृ\appbrk{}ष\appbrk{}क\appbrk{}द\appbrk{}ली\appbrk{}चि\appbrk{}त्र\appbrk{}के\appbrk{}न्द्र\appbrk{}य\appbrk{}वृ\appbrk{}क्षा\appbrk{}स्फो\appbrk{}ता\appbrk{}श्व\appbrk{}मा\appbrk{}र\appbrk{}क\appbrk{}स\appbrk{}प्त\appbrk{}च्छ\appbrk{}दा\appbrk{}ग्नि\appbrk{}मन्ः}\ H\textsuperscript{ac}; \skt{कुट\appbrk{}ज\appbrk{}प\appbrk{}ला\appbrk{}शा\appbrk{}श्व\appbrk{}क\appbrk{}र्ण्ण\appbrk{}पा\appbrk{}रि\appbrk{}भ\appbrk{}द्र\appbrk{}क\appbrk{}बि\appbrk{}भी\appbrk{}त\appbrk{}का\appbrk{}र\appbrk{}ग्व\appbrk{}ध\appbrk{}ति\appbrk{}ल्व\appbrk{}का\appbrk{}र्क्क\appbrk{}स्नु\appbrk{}हा\appbrk{}पा\appbrk{}मा\appbrk{}र्ग्ग\appbrk{}न\appbrk{}क्त\appbrk{}मा\appbrk{}ल\appbrk{}वृ\appbrk{}ष\appbrk{}क\appbrk{}द\appbrk{}ली\appbrk{}चि\appbrk{}त्र\appbrk{}के\appbrk{}न्द्र\appbrk{}न्दु\appbrk{}य\appbrk{}वृ\appbrk{}क्षा\appbrk{}स्फो\appbrk{}ता\appbrk{}श्व\appbrk{}मा\appbrk{}र\appbrk{}क\appbrk{}स\appbrk{}प्त\appbrk{}च्छ\appbrk{}दा\appbrk{}ग्नि\appbrk{}मन्ः}\ H\textsuperscript{pc}}}\skt{~। }\edtext{\skt{चतस्रः}}{\lemma{\skt{चतस्रः}}\Afootnote{A K H; \skt{चत}\ N}}\skt{ }\edtext{\skt{कोशातक्यः}}{\lemma{\skt{कोशातक्यः}}\Afootnote{K H; \skt{च कोशातकीः}\ A; \skt{स्रकोशातकृः}\ N}}\skt{ }\edtext{\skt{समू\-ल\-फ\-ल\-शा\-खा\-प\-त्रा\-न्द\-हे\-त्ततः}}{\lemma{\skt{समू\appbrk{}ल\appbrk{}फ\appbrk{}ल\appbrk{}शा\appbrk{}खा\appbrk{}प\appbrk{}त्रा\appbrk{}न्द\appbrk{}हे\appbrk{}त्ततः}}\Afootnote{K H; \skt{समू\appbrk{}ल\appbrk{}फ\appbrk{}ल\appbrk{}प\appbrk{}त्र\appbrk{}शाखा दहेत्~। ततः}\ A; \skt{समू\appbrk{}ल\appbrk{}फ\appbrk{}ल\appbrk{}शा\appbrk{}प\appbrk{}त्रा\appbrk{}न्दे\appbrk{}ह\appbrk{}न्ततः}\ N}}\skt{ }\edtext{\edtext{\skt{क्षार\-द्रो\-ण\-मु\-द\-क\-द्रोणैः}}{\lemma{\skt{क्षारद्रोणम्}}\Afootnote{A H; \skt{क्षारद्रोण}\ K N}}}{\lemma{\skt{उदकद्रोणैः}}\Afootnote{A H; \skt{मुदकद्रोणैः}\ K; \skt{मुदकद्रोण्णैः}\ N}}\skt{ }\edtext{\edtext{\skt{षड्भिरालोड्य}}{\lemma{\skt{षड्भिर्}}\Afootnote{A N H; \skt{षड्भि}\ K}}}{\lemma{\skt{आलोड्य}}\Afootnote{A N H; \skt{रालोड्य}\ K}}\skt{ }\edtext{\skt{मूत्रैश्च}}{\lemma{\skt{मूत्रैश्च}}\Afootnote{K N H; \skt{मूत्रैर्वा}\ A}}\skt{ }\edtext{\edtext{\skt{यथोक्तैर्महति}}{\lemma{\skt{यथोक्तैर्}}\Afootnote{A N H; \skt{यथो}\ K}}}{\lemma{\skt{महति}}\Afootnote{N H; \skt{एक\appbrk{}विं\appbrk{}श\appbrk{}ति\appbrk{}कृत्वः परिस्राव्य, महति}\ A; \skt{क्तैर्महति}\ K}}\skt{ }\edtext{\skt{कटाहे}}{\lemma{\skt{कटाहे}}\Afootnote{A K H; \skt{काटाहे}\ N}}\skt{ }\edtext{\edtext{\edtext{\skt{शनै\-श्श\-नै\-र्द\-व्या\-व\-घ\-ट्ट\-य\-न्वि\-प\-चेत्स}}{\lemma{\skt{शनैश्}}\Afootnote{N H; \skt{श}\ A; \skt{शनै}\ K}}}{\lemma{\skt{शनैर्}}\Afootnote{N H; \skt{नैर्}\ A; \skt{श्शनै}\ K}}}{\lemma{\skt{दव्या\appbrk{}व\appbrk{}घ\appbrk{}ट्ट\appbrk{}य\appbrk{}न्वि\appbrk{}प\appbrk{}चेत्स}}\Afootnote{\skt{दर्व्याऽ\appbrk{}व\appbrk{}घ\appbrk{}ट्ट\appbrk{}य\appbrk{}न्वि\appbrk{}प\appbrk{}चेत्स}\ A H; \skt{र्दव्या\appbrk{}व\appbrk{}घ\appbrk{}ट्ट\appbrk{}य\appbrk{}न्वि\appbrk{}प\appbrk{}चेत्स}\ K; \skt{दव्या\appbrk{}व\appbrk{}घ\appbrk{}ट्ट\appbrk{}य\appbrk{}न्वि\appbrk{}प\appbrk{}चे\appbrk{}त्स\appbrk{}च्चि\appbrk{}प\appbrk{}चेत्स}\ N}}\skt{ यदा }\edtext{\skt{भव\-त्य\-च्छो\-र\-क्त\-स्तीक्ष्णः}}{\lemma{\skt{भव\appbrk{}त्य\appbrk{}च्छो\appbrk{}र\appbrk{}क्त\appbrk{}स्तीक्ष्णः}}\Afootnote{N; \skt{भवत्यच्छो रक्तस्तीक्ष्णः}\ A; \skt{भव\appbrk{}त्य\appbrk{}च्छो\appbrk{}र\appbrk{}क्त\appbrk{}स्तीक्ष्णः}\ K; \skt{भवत्यच्छो रक्तस्तीक्ष्णाः}\ H}}\skt{ }\edtext{\skt{पिच्छिलश्च }\edtext{\skt{तमादायेतरं}}{\lemma{\skt{आदायेतरं}}\Afootnote{H; \skt{आदाय महति वस्त्रे परिस्राव्येतरं}\ A; \skt{तमादायेतरं}\ K; \skt{लस्च तमादायेतरं}\ N}}}{\lemma{\skt{पिच्छिलश्च तम्}}\Afootnote{A H; \skt{पिच्छिलश्च}\ K; \skt{पिच्छि}\ N}}\skt{ }\edtext{\skt{संसृज्य}}{\lemma{\skt{संसृज्य}}\Afootnote{K H; \skt{विभज्य}\ A; \skt{संसृकृ}\ N}}\skt{ }\edtext{\edtext{\skt{पुनरपि}}{\lemma{\skt{पुनर्}}\Afootnote{A N H; \skt{पुन}\ K}}}{\lemma{\skt{अपि}}\Afootnote{N H; \skt{अ}\ A; \skt{रपि}\ K}}\skt{ }\edtext{\skt{पाका\-या\-धि\-श्र\-ये\-त्तत}}{\lemma{\skt{पाका\appbrk{}या\appbrk{}धि\appbrk{}श्र\appbrk{}ये\appbrk{}त्तत}}\Afootnote{K N; \skt{ग्नावधिश्रयेत्~। तत}\ A; \skt{पाक याधिश्रयेत्तत}\ H}}\skt{ }\edtext{\skt{एव च}}{\lemma{\skt{एव च}}\Afootnote{K N H; \skt{एव}\ A}}\skt{ }\edtext{\skt{क्षारो\-द\-क\-कु\-ड\-व\-म\-ध्यर्ध}}{\lemma{\skt{क्षारो\appbrk{}द\appbrk{}क\appbrk{}कु\appbrk{}ड\appbrk{}व\appbrk{}म\appbrk{}ध्यर्ध}}\Afootnote{K; \skt{क्षारो\appbrk{}द\appbrk{}का\appbrk{}त्कु\appbrk{}ड\appbrk{}व\appbrk{}म\appbrk{}ध्यर्धं}\ A; \skt{क्षारोदं कुडवमध्यर्ध}\ N; \skt{क्षारोदक कुड\appbrk{}व\appbrk{}म\appbrk{}ध्य\appbrk{}र्द्धङ्}\ H}}\skt{ }\edtext{\skt{कृत्वा\-प\-न\-ये\-त्ततः}}{\lemma{\skt{कृत्वा\appbrk{}प\appbrk{}न\appbrk{}ये\appbrk{}त्ततः}}\Afootnote{K N H; \skt{वाऽपनयेत्~। ततः}\ A}}\skt{ }\edtext{\skt{कट\-श\-र्क\-रा\-भ\-स्म\-श\-र्क\-राश्च}}{\lemma{\skt{कट\appbrk{}श\appbrk{}र्क\appbrk{}रा\appbrk{}भ\appbrk{}स्म\appbrk{}श\appbrk{}र्क\appbrk{}राश्च}}\Afootnote{K N H\textsuperscript{ac}; \skt{कट\appbrk{}श\appbrk{}र्क\appbrk{}रा\appbrk{}भ\appbrk{}स्म\appbrk{}शर्क}\ A; \skt{कट\appbrk{}श\appbrk{}र्क्क\appbrk{}रा\appbrk{}भ\appbrk{}स्म\appbrk{}श\appbrk{}र्क्करा श्च~। डो}\ H\textsuperscript{pc}}}\skt{~। }\edtext{\skt{क्षीर\-प\-क\-शं\-ख\-ना\-भी\-र\-ग्नि\-वर्ण्णाः}}{\lemma{\skt{क्षीर\appbrk{}प\appbrk{}क\appbrk{}शं\appbrk{}ख\appbrk{}ना\appbrk{}भी\appbrk{}र\appbrk{}ग्नि\appbrk{}वर्ण्णाः}}\Afootnote{\skt{राक्षी\appbrk{}र\appbrk{}पा\appbrk{}क\appbrk{}श\appbrk{}ङ्ख\appbrk{}ना\appbrk{}भी\appbrk{}र\appbrk{}ग्नि\appbrk{}वर्णाः}\ A; \skt{क्षीर\appbrk{}प\appbrk{}क\appbrk{}शं\appbrk{}ख\appbrk{}ना\appbrk{}भी\appbrk{}र\appbrk{}ग्नि\appbrk{}वर्ण्णाः}\ K; \skt{क्षीर\appbrk{}ष\appbrk{}क\appbrk{}शं\appbrk{}ख\appbrk{}ना\appbrk{}भी\appbrk{}र\appbrk{}क्त\appbrk{}वल्लीः}\ N; \skt{क्षीरे पक\appbrk{}शं\appbrk{}ख\appbrk{}ना\appbrk{}भी\appbrk{}र\appbrk{}ग्नि\appbrk{}वर्ण्णाः}\ H\textsuperscript{ac}; \skt{क्षीरे पक\appbrk{}शं\appbrk{}ख\appbrk{}ना\appbrk{}भी\appbrk{}र\appbrk{}ग्नि\appbrk{}वर्ण्णाः}\ H\textsuperscript{pc}}}\skt{ }\edtext{\skt{कृत्वायसे}}{\lemma{\skt{कृत्वायसे}}\Afootnote{A K N; \skt{कृत्यायसे}\ H}}\skt{ पात्रे }\edtext{\skt{तस्मिं}}{\lemma{\skt{तस्मिं}}\Afootnote{K N; \skt{तस्मि}\ A H}}\skt{ }\edtext{\skt{क्षारोदके}}{\lemma{\skt{क्षारोदके}}\Afootnote{K N; \skt{न्नेव क्षारोदके}\ A; \skt{न्क्षारोदके}\ H}}\skt{ }\edtext{\skt{निषिच्य}}{\lemma{\skt{निषिच्य}}\Afootnote{A N H; \skt{भिषिच्य}\ K}}\skt{~। पिष्ट्वा }\edtext{\skt{तथैव च}}{\lemma{\skt{तथैव च}}\Afootnote{K N H; \skt{तेनैव}\ A}}\skt{ }\edtext{\skt{प्रतीवापो}}{\lemma{\skt{प्रतीवापो}}\Afootnote{K N H; \emph{om.}\ A}}\skt{ }\edtext{\skt{यथालाभं}}{\lemma{\skt{यथालाभं}}\Afootnote{K N H; \emph{om.}\ A}}\skt{ }\edtext{\skt{दन्ती\-चि\-त्र\-क\-ला\-ङ्ग\-ली\-पू\-ती\-क\-प्र\-वा\-ड\-ता\-ल\-प\-त्री\-वि\-ड\-सौ\-व\-र्चि\-का\-क\-न\-क\-क्षी\-री\-हिं\-गु\-व\-चा\-ति\-वि\-षा\-शुक्तीः }\edtext{\skt{श्लक्ष्णचूर्ण्णङ्कृत्वा}}{\lemma{\skt{कृत्वा}}\Afootnote{K N H; \emph{om.}\ A}}}{\lemma{\skt{दन्ती\appbrk{}चि\appbrk{}त्र\appbrk{}क\appbrk{}ला\appbrk{}ङ्ग\appbrk{}ली\appbrk{}पू\appbrk{}ती\appbrk{}क\appbrk{}प्र\appbrk{}वा\appbrk{}ड\appbrk{}ता\appbrk{}ल\appbrk{}प\appbrk{}त्री\appbrk{}वि\appbrk{}ड\appbrk{}सौ\appbrk{}व\appbrk{}र्चि\appbrk{}का\appbrk{}क\appbrk{}न\appbrk{}क\appbrk{}क्षी\appbrk{}री\appbrk{}हिं\appbrk{}गु\appbrk{}व\appbrk{}चा\appbrk{}ति\appbrk{}वि\appbrk{}षा\appbrk{}शुक्तीः श्लक्ष्णचूर्ण्णङ्}}\Afootnote{K; \skt{द्विव्र\appbrk{}णेऽ\appbrk{}ष्ट\appbrk{}प\appbrk{}ल\appbrk{}सं\appbrk{}मितं शङ्खनाभ्यादीनां प्रमाणं प्रतिवाप्य}\ A; \skt{वन्ती चित्र\appbrk{}क\appbrk{}ला\appbrk{}ङ्ग\appbrk{}ली\appbrk{}पू\appbrk{}ति\appbrk{}क\appbrk{}प्र\appbrk{}तु\appbrk{}ता\appbrk{}ल\appbrk{}प\appbrk{}त्री\appbrk{}वि\appbrk{}ड\appbrk{}सौ\appbrk{}व\appbrk{}र्च्चि\appbrk{}का\appbrk{}क\appbrk{}न\appbrk{}क\appbrk{}क्षी\appbrk{}री\appbrk{}हिं\appbrk{}गु\appbrk{}व\appbrk{}चा\appbrk{}ति\appbrk{}वि\appbrk{}षा\appbrk{}शु\appbrk{}क्ती॥\appbrk{}श्ल\appbrk{}क्ष्ण\appbrk{}चूर्ण्णं}\ N; \skt{दन्ती\appbrk{}चि\appbrk{}त्र\appbrk{}क\appbrk{}ला\appbrk{}ङ्ग\appbrk{}ली\appbrk{}पू\appbrk{}ती\appbrk{}क\appbrk{}प्र\appbrk{}वा\appbrk{}ल\appbrk{}ता\appbrk{}ल\appbrk{}पत्त्री विड\appbrk{}सौ\appbrk{}व\appbrk{}र्च्चि\appbrk{}का\appbrk{}क\appbrk{}न\appbrk{}क\appbrk{}क्षी\appbrk{}री\appbrk{}हिं\appbrk{}गु\appbrk{}व\appbrk{}चा\appbrk{}ति\appbrk{}वि\appbrk{}षा\appbrk{}शु\appbrk{}क्तीः\appbrk{}श्ल\appbrk{}क्ष्ण\appbrk{}चू\appbrk{}र्ण्णाङ्}\ H\textsuperscript{ac} H\textsuperscript{pc}}}\skt{ }\edtext{\skt{निदध्यात्~। }\edtext{\skt{सत\-त\-म\-प्र\-म\-त्तश्च}}{\lemma{\skt{अप्रमत्तश्च}}\Afootnote{K N H\textsuperscript{pc}; \skt{अप्रमत्तश्चै}\ A; \skt{अप्रत्तश्च}\ H\textsuperscript{ac}}}}{\lemma{\skt{निदध्यात्~। सततम्}}\Afootnote{K N H; \skt{सततम्}\ A}}\skt{ }\edtext{\skt{दर्व्यावघटायं}}{\lemma{\skt{दर्व्यावघटायं}}\Afootnote{K; \skt{नमवघट्टय}\ A; \skt{दप्यावघत्तयं}\ N; \skt{दर्व्व्यावघटय}\ H}}\skt{ }\edtext{\skt{विपचेत्~। स}}{\lemma{\skt{विपचेत्~। स}}\Afootnote{K N; \skt{न्विपचेत्~। स}\ A H}}\skt{ यथा }\edtext{\skt{नातिसान्द्रो}}{\lemma{\skt{नातिसान्द्रो}}\Afootnote{A K H; \skt{नाति…सान्द्रो}\ N}}\skt{ नातिद्रवश्च भवति तथा }\edtext{\skt{प्रयतेत}}{\lemma{\skt{प्रयतेत}}\Afootnote{A K H; \skt{पयतेत}\ N}}\skt{~। }\edtext{\edtext{\edtext{\skt{अथै\-न\-मा\-ग\-त\-पा\-क\-म\-व\-ता\-र्या\-नु\-गु\-प्त\-मा\-यसे}}{\lemma{\skt{अथैनम्}}\Afootnote{A K H; \skt{अथेनम्}\ N}}}{\lemma{\skt{आगतपाकम्}}\Afootnote{A K H; \skt{आगतया क}\ N}}}{\lemma{\skt{अव\appbrk{}ता\appbrk{}र्या\appbrk{}नु\appbrk{}गु\appbrk{}प्तम्}}\Afootnote{A K H; \skt{र्मव\appbrk{}ता\appbrk{}र्पा\appbrk{}नु\appbrk{}गु\appbrk{}प्तम्}\ N}}\skt{ कुम्भे }\edtext{\skt{निदध्यात्}}{\lemma{\skt{निदध्यात्}}\Afootnote{K N H; \skt{संवृतमुखे निदध्यादेष मध्यमः}\ A}}\skt{~।}\edtext{}{\lemma{\skt{आदीपयेत्~। अथोपशान्तेऽग्नौ}}\Afootnote{A; \skt{आदी\appbrk{}प\appbrk{}ये\appbrk{}त्य\appbrk{}थो\appbrk{}प\appbrk{}शा\appbrk{}न्तेग्नौ}\ K H; \skt{आदीपयेत्~॥ यथोपशान्तर्मौ}\ N}}\pend

\pstart \ledsidenote{\passref{1.11.15}}\edtext{\skt{क्षीणबले च}}{\lemma{\skt{क्षीणबले च}}\Afootnote{K N H; \skt{क्षीणबले तु}\ A}}\skt{ }\edtext{\skt{क्षारो\-द\-क\-मा\-व\-पे\-द्ब\-ल\-क\-र\-णार्थं}}{\lemma{\skt{आव\appbrk{}पे\appbrk{}द्ब\appbrk{}ल\appbrk{}क\appbrk{}र\appbrk{}णार्थं}}\Afootnote{A K H; \skt{आव\appbrk{}ये\appbrk{}द्ब\appbrk{}ल\appbrk{}क\appbrk{}र\appbrk{}णार्थं}\ N}}\skt{~॥}\pend

\pstart\versepar \ledsidenote{\passref{1.11.16}}\edtext{\skt{भवति}}{\lemma{\skt{भवति}}\Afootnote{N H; \skt{भवतश्}\ A; \emph{om.}\ K}}\skt{ }\edtext{\skt{चात्र}}{\lemma{\skt{चात्र}}\Afootnote{A N; \emph{om.}\ K H}}\skt{~॥ }\pend
\pstart\versepar \skt{नैवातितीक्ष्णो न मृदुः शुक्लः }\edtext{\skt{श्लक्ष्णोति}}{\lemma{\skt{श्लक्ष्णोति}}\Afootnote{K H; \skt{श्लक्ष्णोऽथ}\ A; \skt{क्ष्णोति}\ N}}\skt{ पिच्छिलः~। }\pend
\pstart\versepar \skt{अविष्यन्दी शिवः शीघ्रः }\edtext{\skt{क्षारो}}{\lemma{\skt{क्षारो}}\Afootnote{A K H; \skt{क्षा}\ N}}\skt{ }\edtext{\skt{ह्यष्टगुणः}}{\lemma{\skt{ह्यष्टगुणः}}\Afootnote{K; \skt{ह्यष्टगुणः}\ A H; \skt{रोह्यष्टगुणः}\ N}}\skt{ स्मृतः~।}\pend

\pstart\versepar \ledsidenote{\passref{1.11.17}}\edtext{\edtext{\skt{अत्यु\-ष्ण\-म\-ति\-पै\-च्छि\-ल्य\-म\-ति\-तै\-क्ष्ण्य\-वि\-स\-र्पिता}}{\lemma{\skt{अत्यु\appbrk{}ष्ण\appbrk{}म\appbrk{}ति\appbrk{}पै\appbrk{}च्छि\appbrk{}ल्यम्}}\Afootnote{K N; \skt{अति\appbrk{}मा\appbrk{}र्द\appbrk{}व\appbrk{}श्वै\appbrk{}त्यौ\appbrk{}ष्ण्य\appbrk{}तै\appbrk{}क्ष्ण्य\appbrk{}पै\appbrk{}च्छिल्य}\ A; \skt{अत्यौ\appbrk{}ष्ण\appbrk{}म\appbrk{}ति\appbrk{}पै\appbrk{}च्छि\appbrk{}ल्यम्}\ H}}}{\lemma{\skt{अति\appbrk{}तै\appbrk{}क्ष्ण्य\appbrk{}वि\appbrk{}स\appbrk{}र्पिता}}\Afootnote{K H; \skt{सर्पि}\ A; \skt{अति\appbrk{}ती\appbrk{}क्ष्ण्य\appbrk{}वि\appbrk{}स\appbrk{}र्प्पिता}\ N}}\skt{~। }\pend
\pstart\versepar \skt{}\edtext{\skt{अत्यर्थं}}{\lemma{\skt{अत्यर्थं}}\Afootnote{K N H; \emph{om.}\ A}}\skt{ }\edtext{\skt{मार्दवं}}{\lemma{\skt{मार्दवं}}\Afootnote{K H; \emph{om.}\ A; \skt{मार्द्रवं}\ N}}\skt{ }\edtext{\edtext{\skt{शैत्यमत्यर्थं}}{\lemma{\skt{शैत्यम्}}\Afootnote{K N H; \emph{om.}\ A}}}{\lemma{\skt{अत्यर्थं}}\Afootnote{K H; \emph{om.}\ A; \skt{अत्यर्थ}\ N}}\skt{ }\edtext{\edtext{\skt{सान्द्रमेव}}{\lemma{\skt{सान्द्रम्}}\Afootnote{K N H; \emph{om.}\ A}}\skt{ च}}{\lemma{\skt{एव च}}\Afootnote{K H; \emph{om.}\ A; \skt{एव च~॥ ः}\ N}}\skt{~। }\pend
\pstart\versepar \skt{}\edtext{\skt{हीनौ\-ष\-ध्य\-वि\-प\-क्वत्वं}}{\lemma{\skt{हीनौ\appbrk{}ष\appbrk{}ध्य\appbrk{}वि\appbrk{}प\appbrk{}क्वत्वं}}\Afootnote{K H; \emph{om.}\ A; \skt{इनौ\appbrk{}ष\appbrk{}ध्य\appbrk{}वि\appbrk{}प\appbrk{}क्वत्वं}\ N}}\skt{ }\edtext{\skt{क्षारदोषा}}{\lemma{\skt{क्षारदोषा}}\Afootnote{K; \emph{om.}\ A; \skt{त्तार दोषा}\ N; \skt{क्षार दोष}\ H}}\skt{ }\edtext{\skt{न च}}{\lemma{\skt{न च}}\Afootnote{K; \emph{om.}\ A; \skt{नव}\ N H}}\skt{ }\edtext{\skt{स्मृताः}}{\lemma{\skt{स्मृताः}}\Afootnote{K N H; \skt{ताः~। सान्द्रताऽपक्वता हीनद्रव्यता दोष उच्यते}\ A}}\skt{~॥}\pend

\pstart \ledsidenote{\passref{1.11.18}}\skt{तत्र }\edtext{\skt{क्षारसाध्यव्याधिं}}{\lemma{\skt{क्षारसाध्यव्याधिं}}\Afootnote{K N H; \skt{क्षारसाध्यव्याधि}\ A}}\skt{ व्याधि\-त\-मु\-प\-वेश्य }\edtext{\skt{निवातासम्बाधे}}{\lemma{\skt{निवातासम्बाधे}}\Afootnote{K N H; \skt{निवातातपे देशेऽसम्बाधे}\ A}}\skt{ }\edtext{\skt{देशे}}{\lemma{\skt{देशे}}\Afootnote{K N H; \emph{om.}\ A}}\skt{ }\edtext{\skt{अग्रो\-प\-ह\-र\-णी\-यो\-क्तो\-प\-सं\-भा\-र\-सं\-भृ\-त\-तनोः}}{\lemma{\skt{अग्रो\appbrk{}प\appbrk{}ह\appbrk{}र\appbrk{}णी\appbrk{}यो\appbrk{}क्तो\appbrk{}प\appbrk{}सं\appbrk{}भा\appbrk{}र\appbrk{}सं\appbrk{}भृ\appbrk{}त\appbrk{}तनोः}}\Afootnote{K H; \skt{ऽग्रो\appbrk{}प\appbrk{}ह\appbrk{}र\appbrk{}णी\appbrk{}यो\appbrk{}क्तेन विधा\appbrk{}ने\appbrk{}नो\appbrk{}प\appbrk{}स\appbrk{}म्भृ\appbrk{}त\appbrk{}स\appbrk{}म्भारं, ततो}\ A; \skt{अग्रो\appbrk{}प\appbrk{}ह\appbrk{}र\appbrk{}णी\appbrk{}यो\appbrk{}क्ते\appbrk{}प\appbrk{}सं\appbrk{}भा\appbrk{}र\appbrk{}सं\appbrk{}भृ\appbrk{}त\appbrk{}तनोः}\ N}}\skt{~। }\edtext{\skt{अन्य\-त\-म\-म\-व\-घृ\-ष्या\-व\-लिख्य}}{\lemma{\skt{अन्य\appbrk{}त\appbrk{}म\appbrk{}म\appbrk{}व\appbrk{}घृ\appbrk{}ष्या\appbrk{}व\appbrk{}लिख्य}}\Afootnote{\skt{ऽस्य तमवकाशं निरी\appbrk{}क्ष्या\appbrk{}व\appbrk{}घृ\appbrk{}ष्या\appbrk{}व\appbrk{}लिख्य}\ A; \skt{अन्य\appbrk{}त\appbrk{}म\appbrk{}म\appbrk{}व\appbrk{}घृ\appbrk{}ष्या\appbrk{}व\appbrk{}लिख्य}\ K H; \skt{अन्य\appbrk{}त\appbrk{}म\appbrk{}म\appbrk{}व\appbrk{}घृ\appbrk{}व्या\appbrk{}वलि ख्य}\ N}}\skt{ }\edtext{\skt{प्रच्छयित्वा वा}}{\lemma{\skt{प्रच्छयित्वा वा}}\Afootnote{K H; \skt{प्रच्छयित्वा}\ A; \skt{प्रच्छयिव्वा}\ N}}\skt{ }\edtext{\skt{शलाकया}}{\lemma{\skt{शलाकया}}\Afootnote{A K N; \skt{शलाक…या}\ H}}\skt{ क्षारं }\edtext{\skt{प्रतिसार्य}}{\lemma{\skt{प्रतिसार्य}}\Afootnote{K H; \skt{प्रतिसार}\ A; \skt{प्रतिसा}\ N}}\skt{ }\edtext{\edtext{\skt{वाक्छ\-त\-मा\-त्र\-मु\-पै\-क्षेत}}{\lemma{\skt{वाक्छतमात्रम्}}\Afootnote{K H; \skt{येत्, दत्त्वा वाक्शतमात्रम्}\ A; \skt{यवा\appbrk{}क्छ\appbrk{}त\appbrk{}मा\appbrk{}त्रम्}\ N}}}{\lemma{\skt{उपैक्षेत}}\Afootnote{K H; \skt{उपेक्षेत}\ A N}}\skt{~।}\pend

\pstart\versepar \ledsidenote{\passref{1.11.19}}\skt{तस्मिं निपतिते व्याधौ }\edtext{\skt{कृष्णता}}{\lemma{\skt{कृष्णता}}\Afootnote{A K H; \skt{क्षणता}\ N}}\skt{ दग्धलक्षणं~। }\pend
\pstart\versepar \skt{}\edtext{\skt{तत्राम्लवर्गः}}{\lemma{\skt{तत्राम्लवर्गः}}\Afootnote{A K; \skt{तत्राम्ब्लवर्गः}\ N; \skt{तत्रावर्ग्गः}\ H}}\skt{ }\edtext{\skt{समनः}}{\lemma{\skt{समनः}}\Afootnote{K N H; \skt{शमनः}\ A}}\skt{ }\edtext{\skt{सर्पि\-र्म\-धु\-स\-मा\-युतः}}{\lemma{\skt{सर्पि\appbrk{}र्म\appbrk{}धु\appbrk{}स\appbrk{}मा\appbrk{}युतः}}\Afootnote{K H; \skt{सर्पि\appbrk{}र्म\appbrk{}धु\appbrk{}क\appbrk{}सं\appbrk{}युतः}\ A; \skt{सर्प्पि\appbrk{}म\appbrk{}धु\appbrk{}स\appbrk{}मा\appbrk{}युतः}\ N}}\skt{~।}\pend

\pstart\versepar \ledsidenote{\passref{1.11.20}}\edtext{\skt{अथ च}}{\lemma{\skt{अथ च}}\Afootnote{K; \skt{अथ चे}\ A H N}}\skt{ }\edtext{\skt{स्थिर\-मू\-ल\-त्वा\-त्क्षा\-र\-द\-ग्धन्न}}{\lemma{\skt{स्थिर\appbrk{}मू\appbrk{}ल\appbrk{}त्वा\appbrk{}त्क्षा\appbrk{}र\appbrk{}द\appbrk{}ग्धन्न}}\Afootnote{K N; \skt{त्स्थिर\appbrk{}मू\appbrk{}ल\appbrk{}त्वा\appbrk{}त्क्षा\appbrk{}र\appbrk{}दग्धं न}\ A H}}\skt{ }\edtext{\skt{दीर्यते}}{\lemma{\skt{दीर्यते}}\Afootnote{K H; \skt{शीर्यते}\ A; \skt{दीयते}\ N}}\skt{~। }\pend
\pstart\versepar \skt{}\edtext{\skt{इद\-मा\-ले\-प\-न\-न्तत्र}}{\lemma{\skt{इद\appbrk{}मा\appbrk{}ले\appbrk{}प\appbrk{}न\appbrk{}न्तत्र}}\Afootnote{A K H; \skt{इदमालेपनं तत्र}\ N}}\skt{ }\edtext{\skt{शीघ्र}}{\lemma{\skt{शीघ्र}}\Afootnote{K; \skt{सम}\ A; \skt{सीघ्रं}\ N; \skt{शीघ्रं}\ H}}\skt{ }\edtext{\skt{सम\-व\-चा\-र\-येत्}}{\lemma{\skt{सम\appbrk{}व\appbrk{}चा\appbrk{}र\appbrk{}येत्}}\Afootnote{K N H; \skt{ग्रम\appbrk{}व\appbrk{}चा\appbrk{}र\appbrk{}येत्}\ A}}\skt{~।}\pend

\pstart\versepar \ledsidenote{\passref{1.11.21}}\edtext{\skt{अम्ल\-का\-ञ्जि\-क\-बी\-जानां}}{\lemma{\skt{अम्ल\appbrk{}का\appbrk{}ञ्जि\appbrk{}क\appbrk{}बी\appbrk{}जानां}}\Afootnote{K H; \skt{अम्ल\appbrk{}का\appbrk{}ञ्जि\appbrk{}क\appbrk{}बी\appbrk{}जानि}\ A; \skt{अंब्ल\appbrk{}कं\appbrk{}ञ्जि\appbrk{}क\appbrk{}बी\appbrk{}जानं}\ N}}\skt{ }\edtext{\edtext{\skt{तिला\-म्म\-धु\-क\-मेव}}{\lemma{\skt{तिलाम्}}\Afootnote{K N H; \skt{तिला}\ A}}\skt{ च}}{\lemma{\skt{मधुकमेव च}}\Afootnote{K H; \skt{न्मधुकमेव च}\ A; \skt{मधुकमेव च}\ N}}\skt{~। }\pend
\pstart\versepar \skt{}\edtext{\skt{प्रपिष्य}}{\lemma{\skt{प्रपिष्य}}\Afootnote{K H; \skt{प्रपेष्य}\ A; \skt{प्रसिष्य}\ N}}\skt{ समभागानि }\edtext{\skt{तेनै\-व\-म\-नु\-ले\-प\-येत्}}{\lemma{\skt{तेनै\appbrk{}व\appbrk{}म\appbrk{}नु\appbrk{}ले\appbrk{}प\appbrk{}येत्}}\Afootnote{N H; \skt{तेनै\appbrk{}न\appbrk{}म\appbrk{}नु\appbrk{}ले\appbrk{}प\appbrk{}येत्}\ A; \skt{तेनै\appbrk{}व\appbrk{}म\appbrk{}नु\appbrk{}ले\appbrk{}प\appbrk{}येत्}\ K}}\skt{~।}\pend

\pstart\versepar \ledsidenote{\passref{1.11.22}}\skt{तिलकल्कः समधुको }\edtext{\skt{घृताक्तो}}{\lemma{\skt{घृताक्तो}}\Afootnote{A K H; \skt{घृताक्ते}\ N}}\skt{ }\edtext{\skt{व्रणलेपनः}}{\lemma{\skt{व्रणलेपनः}}\Afootnote{K N; \skt{व्रणरोपणः}\ A; \skt{व्रणलेपणः}\ H}}\skt{~॥ }\pend
\pstart\versepar \skt{}\edtext{\skt{रसेनाम्लेन}}{\lemma{\skt{रसेनाम्लेन}}\Afootnote{A K H; \skt{रसेनाम्ब्लेन}\ N}}\skt{ तीक्ष्णेन }\edtext{\skt{वीर्योष्णेन}}{\lemma{\skt{वीर्योष्णेन}}\Afootnote{A K H; \skt{वीयाष्णेन}\ N}}\skt{ }\edtext{\skt{तथैव च}}{\lemma{\skt{तथैव च}}\Afootnote{K N H; \skt{च योजितः}\ A}}\skt{~।}\pend

\pstart\versepar \ledsidenote{\passref{1.11.23}}\edtext{\skt{आग्नेयेननाग्नि}}{\lemma{\skt{आग्नेयेननाग्नि}}\Afootnote{K; \skt{आग्नेयेनाग्नि}\ A N; \skt{आग्नेयै…नाग्नि}\ H}}\skt{ }\edtext{\skt{सदृशः}}{\lemma{\skt{सदृशः}}\Afootnote{K N H; \skt{ना तुल्यः}\ A}}\skt{ कथं }\edtext{\skt{क्षार}}{\lemma{\skt{क्षार}}\Afootnote{K N; \skt{क्षारः}\ A H}}\skt{ प्रशाम्यति~। }\pend
\pstart\versepar \skt{एवं }\edtext{\skt{चेन्मन्यसे}}{\lemma{\skt{चेन्मन्यसे}}\Afootnote{A K H; \skt{चेत्मन्यसे}\ N}}\skt{ वत्स प्रोच्यमानं }\edtext{\skt{शृणुष्व}}{\lemma{\skt{शृणुष्व}}\Afootnote{K N H; \skt{निबोध}\ A}}\skt{ मे~।}\pend

\pstart\versepar \ledsidenote{\passref{1.11.24}}\edtext{\skt{अम्लवर्ज्या रसां}}{\lemma{\skt{अम्लवर्ज्या रसां}}\Afootnote{K H; \skt{अम्लवर्जान्रसा}\ A; \skt{अंब्लवर्ज्यारसां}\ N}}\skt{ }\edtext{\skt{क्षारे}}{\lemma{\skt{क्षारे}}\Afootnote{K H; \skt{न्क्षारे}\ A; \skt{त्लारे}\ N}}\skt{ सर्वानेव विभावयेत्~। }\pend
\pstart\versepar \skt{कटुक था~। }\pend
\pstart\versepar \skt{}\edtext{\skt{अम्लेन}}{\lemma{\skt{अम्लेन}}\Afootnote{A K H; \skt{अंब्लेन}\ N}}\skt{ }\edtext{\skt{सह}}{\lemma{\skt{सह}}\Afootnote{A K N H\textsuperscript{pc}; \skt{मह}\ H\textsuperscript{ac}}}\skt{ संयुक्तः }\edtext{\skt{सुतीक्ष्णो लवणो}}{\lemma{\skt{सुतीक्ष्णो लवणो}}\Afootnote{K; \skt{सतीक्ष्णलवणो}\ A; \skt{सुतीक्ष्णलवणो}\ N; \skt{सुतीक्ष्णोलवणो}\ H}}\skt{ }\edtext{\skt{रसः}}{\lemma{\skt{रसः}}\Afootnote{K N H; \skt{रसे}\ A}}\skt{~।}\edtext{}{\lemma{\skt{विभावयेत्~। कटुक था}}\Afootnote{K; \skt{विभावयेत्~। कटुकस्तत्र भूयिष्ठो लव\appbrk{}णोऽ\appbrk{}नु\appbrk{}र\appbrk{}स\appbrk{}स्तथा}\ A; \skt{विभावयेत्~॥ कटुकस्तत्र भूयिष्ठो लवणेन रसस्तथा}\ N; \skt{विभावयेत्~। कटुकस्तत्र भूयिष्ठो लव\appbrk{}णो\appbrk{}नु\appbrk{}र\appbrk{}स\appbrk{}स्तथा}\ H}}\pend

\pstart\versepar \ledsidenote{\passref{1.11.25}}\edtext{\edtext{\skt{माधूर्यमाशु}}{\lemma{\skt{माधूर्यम्}}\Afootnote{K; \skt{माधुर्यं}\ A; \skt{माधूय}\ N; \skt{माधूर्य}\ H}}}{\lemma{\skt{आशु}}\Afootnote{K; \skt{भज}\ A; \skt{मास्त}\ N; \skt{माशु}\ H}}\skt{ }\edtext{\skt{व्रजति}}{\lemma{\skt{व्रजति}}\Afootnote{K H; \skt{तेऽत्यर्थं}\ A; \skt{क्रेजति}\ N}}\skt{ }\edtext{\skt{तीक्ष्णभावञ्च}}{\lemma{\skt{तीक्ष्णभावञ्च}}\Afootnote{K N H; \skt{तीक्ष्णभावं वि}\ A}}\skt{ मुञ्चति~। }\pend
\pstart\versepar \skt{}\edtext{\skt{माधूर्ययोगान्न}}{\lemma{\skt{माधूर्ययोगान्न}}\Afootnote{K H; \skt{माधुर्याच्छममा}\ A; \skt{माधूययोगान्न}\ N}}\skt{ }\edtext{\skt{दहे\-द\-ग्नि\-र\-ग्नि\-रि\-वा\-प्लुतः}}{\lemma{\skt{दहे\appbrk{}द\appbrk{}ग्नि\appbrk{}र\appbrk{}ग्नि\appbrk{}रि\appbrk{}वा\appbrk{}प्लुतः}}\Afootnote{K; \skt{प्नोति वह्नि\appbrk{}र\appbrk{}द्भि\appbrk{}रि\appbrk{}वा\appbrk{}प्लुतः}\ A; \skt{दहे\appbrk{}द\appbrk{}ग्नि\appbrk{}र\appbrk{}द्भि\appbrk{}रि\appbrk{}वा\appbrk{}प्लुतः}\ N; \skt{दहे\appbrk{}द\appbrk{}ग्नि\appbrk{}र\appbrk{}ग्नि\appbrk{}रि\appbrk{}वा\appbrk{}प्लुतः~॥ …}\ H}}\skt{~॥}\pend

\pstart \ledsidenote{\passref{1.11.26}}\skt{तत्र }\edtext{\skt{सम्यग्दग्धे}}{\lemma{\skt{सम्यग्दग्धे}}\Afootnote{A K H; \skt{सम्यग्दश्चै}\ N}}\skt{ }\edtext{\skt{विकारोपशमो}}{\lemma{\skt{विकारोपशमो}}\Afootnote{A K N; \skt{विकारोपसमो}\ H}}\skt{ }\edtext{\skt{लाघ\-व\-म\-ना\-स्रा\-वश्च}}{\lemma{\skt{लाघ\appbrk{}व\appbrk{}म\appbrk{}ना\appbrk{}स्रा\appbrk{}वश्च}}\Afootnote{A K N; \skt{लाघ\appbrk{}व\appbrk{}म\appbrk{}ना\appbrk{}श्रा\appbrk{}वश्च}\ H}}\skt{~। }\edtext{\skt{हीने तोद\-क\-ण्डू\-जा\-ड्यानि}}{\lemma{\skt{हीने तोद\appbrk{}क\appbrk{}ण्डू\appbrk{}जा\appbrk{}ड्यानि}}\Afootnote{K; \skt{हीनदग्धे तोद\appbrk{}क\appbrk{}ण्डु\appbrk{}जा\appbrk{}ड्यानि}\ A; \skt{हीने\appbrk{}तो\appbrk{}द\appbrk{}क\appbrk{}ण्डु\appbrk{}जा\appbrk{}र्ड्यानि}\ N; \skt{हीने\appbrk{}तो\appbrk{}द\appbrk{}द\appbrk{}क\appbrk{}ण्डू\appbrk{}जा\appbrk{}ड्यादि}\ H}}\skt{ व्याधिवृद्धिश्च~। }\edtext{\skt{अतिदग्धे}}{\lemma{\skt{अतिदग्धे}}\Afootnote{A K N; \skt{… अतिदग्धे}\ H\textsuperscript{ac}; \skt{…डो अतिदग्धे}\ H\textsuperscript{pc}}}\skt{ }\edtext{\skt{दाह\-पा\-क\-श्र\-मा\-ङ्ग\-म\-द\-क्लमाः}}{\lemma{\skt{दाह\appbrk{}पा\appbrk{}क\appbrk{}श्र\appbrk{}मा\appbrk{}ङ्ग\appbrk{}म\appbrk{}द\appbrk{}क्लमाः}}\Afootnote{K; \skt{दाह\appbrk{}पा\appbrk{}क\appbrk{}रा\appbrk{}ग\appbrk{}स्रा\appbrk{}वा\appbrk{}ङ्ग\appbrk{}म\appbrk{}र्द\appbrk{}क्लम}\ A; \skt{दाह\appbrk{}पा\appbrk{}क\appbrk{}स्र\appbrk{}मा\appbrk{}ङ्ग\appbrk{}म\appbrk{}र्द्द\appbrk{}क्लमाः}\ N; \skt{दाह\appbrk{}पा\appbrk{}क\appbrk{}श्रा\appbrk{}मा\appbrk{}ङ्ग\appbrk{}म\appbrk{}र्द्द\appbrk{}क्लमाः}\ H}}\skt{~। }\edtext{\edtext{\skt{पिपा\-सा\-म\-र\-णा\-ञ्चेति}}{\lemma{\skt{पिपा\appbrk{}सा\appbrk{}म\appbrk{}र\appbrk{}णाञ्}}\Afootnote{K; \skt{पिपासामूर्च्छाः स्युर्मरणं}\ A; \skt{पिपासामरणं}\ N H}}}{\lemma{\skt{चेति}}\Afootnote{K N H; \skt{वा}\ A}}\skt{~।}\pend

\pstart \ledsidenote{\passref{1.11.27}}\edtext{\skt{क्षार\-द\-ग्ध\-व्र\-णञ्च}}{\lemma{\skt{क्षार\appbrk{}द\appbrk{}ग्ध\appbrk{}व्र\appbrk{}णञ्च}}\Afootnote{K H; \skt{क्षारदग्धव्रणं तु}\ A; \skt{क्षार\appbrk{}द\appbrk{}ग्ध\appbrk{}व्र\appbrk{}णञ्च तु}\ N}}\skt{ यथादोषं }\edtext{\skt{यथाव्याधिं}}{\lemma{\skt{यथाव्याधिं}}\Afootnote{K H; \skt{यथाव्याधि}\ A; \skt{यथाव्याधिञ्ञ्}\ N}}\skt{ }\edtext{\skt{चोपक्रमेत}}{\lemma{\skt{चोपक्रमेत}}\Afootnote{K H; \skt{चोपक्रमेत्}\ A; \skt{चापक्रमेत्}\ N}}\skt{~।}\pend

\pstart \ledsidenote{\passref{1.11.28}}\edtext{\skt{अथा}}{\lemma{\skt{अथा}}\Afootnote{K; \skt{अथ नैते}\ A; \skt{अथ}\ N H}}\skt{ }\edtext{\skt{क्षारकृत्या}}{\lemma{\skt{क्षारकृत्या}}\Afootnote{K H; \skt{क्षारकृत्याः}\ A; \skt{क्षारक्षड्या}\ N}}\skt{ }\edtext{\skt{भवन्ति}}{\lemma{\skt{भवन्ति}}\Afootnote{K H; \skt{तद्यथा}\ A; \skt{भवन्तिं}\ N}}\skt{~। }\edtext{\skt{दुर्ब\-ल\-बा\-ल\-स्थि\-वि\-र\-भी\-रु\-स\-र्वा\-ङ्ग\-शू\-नो\-द\-री\-ग\-र्भि\-णी\-ऋ\-तु\-म\-ती\-प्र\-वृ\-द्ध\-ज्व\-री\-प्र\-मे\-ही\-रू\-क्ष\-क्ष\-त\-क्षी\-ण\-तृ\-ष्णा\-मू\-र्च्छो\-प\-द्रु\-त\-क्लीबा}}{\lemma{\skt{दुर्ब\appbrk{}ल\appbrk{}बा\appbrk{}ल\appbrk{}स्थि\appbrk{}वि\appbrk{}र\appbrk{}भी\appbrk{}रु\appbrk{}स\appbrk{}र्वा\appbrk{}ङ्ग\appbrk{}शू\appbrk{}नो\appbrk{}द\appbrk{}री\appbrk{}ग\appbrk{}र्भि\appbrk{}णी\appbrk{}ऋ\appbrk{}तु\appbrk{}म\appbrk{}ती\appbrk{}प्र\appbrk{}वृ\appbrk{}द्ध\appbrk{}ज्व\appbrk{}री\appbrk{}प्र\appbrk{}मे\appbrk{}ही\appbrk{}रू\appbrk{}क्ष\appbrk{}क्ष\appbrk{}त\appbrk{}क्षी\appbrk{}ण\appbrk{}तृ\appbrk{}ष्णा\appbrk{}मू\appbrk{}र्च्छो\appbrk{}प\appbrk{}द्रु\appbrk{}त\appbrk{}क्लीबा}}\Afootnote{K; \skt{दुर्ब\appbrk{}ल\appbrk{}बा\appbrk{}ल\appbrk{}स्थ\appbrk{}वि\appbrk{}र\appbrk{}भी\appbrk{}रु\appbrk{}स\appbrk{}र्वा\appbrk{}ङ्ग\appbrk{}शू\appbrk{}नो\appbrk{}द\appbrk{}रि\appbrk{}र\appbrk{}क्त\appbrk{}पि\appbrk{}त्ति\appbrk{}ग\appbrk{}र्भि\appbrk{}ण्यृ\appbrk{}तु\appbrk{}म\appbrk{}ती\appbrk{}प्र\appbrk{}वृ\appbrk{}द्ध\appbrk{}ज्व\appbrk{}रि\appbrk{}प्र\appbrk{}मे\appbrk{}हि\appbrk{}रू\appbrk{}क्ष\appbrk{}क्ष\appbrk{}त\appbrk{}क्षीण- तृष्णा\appbrk{}मू\appbrk{}र्च्छो\appbrk{}प\appbrk{}द्रु\appbrk{}त\appbrk{}क्लीबा}\ A; \skt{दुर्ब\appbrk{}ल\appbrk{}बा\appbrk{}ल\appbrk{}छ\appbrk{}वि\appbrk{}र\appbrk{}भी\appbrk{}रु\appbrk{}स\appbrk{}र्वा\appbrk{}ङ्ग\appbrk{}सू\appbrk{}नो\appbrk{}द\appbrk{}री\appbrk{}ग\appbrk{}र्भि\appbrk{}णी\appbrk{}द\appbrk{}न्त\appbrk{}म\appbrk{}ती\appbrk{}प्र\appbrk{}वृ\appbrk{}द्ध\appbrk{}ज्व\appbrk{}री\appbrk{}प्र\appbrk{}मे\appbrk{}ही\appbrk{}रू\appbrk{}क्ष\appbrk{}क्षे\appbrk{}त\appbrk{}क्षी\appbrk{}ण\appbrk{}तृ\appbrk{}ष्णा\appbrk{}मू\appbrk{}र्च्छो\appbrk{}प\appbrk{}द्रु\appbrk{}त\appbrk{}क्लीबा}\ N; \skt{दुर्ब्ब\appbrk{}ल\appbrk{}बा\appbrk{}ल\appbrk{}स्थ\appbrk{}वि\appbrk{}र\appbrk{}भी\appbrk{}रु\appbrk{}स\appbrk{}र्व्वा\appbrk{}ङ्ग\appbrk{}शू\appbrk{}नो\appbrk{}द\appbrk{}री\appbrk{}ग\appbrk{}र्भ्भि\appbrk{}णी\appbrk{}ऋ\appbrk{}तु\appbrk{}म\appbrk{}ती\appbrk{}प्र\appbrk{}वृ\appbrk{}द्ध\appbrk{}ज्व\appbrk{}री\appbrk{}प्र\appbrk{}मे\appbrk{}ही\appbrk{}रू\appbrk{}क्ष\appbrk{}क्ष\appbrk{}त\appbrk{}क्षी\appbrk{}ण\appbrk{}तृ\appbrk{}ष्णा\appbrk{}मू\appbrk{}र्च्छो\appbrk{}प\appbrk{}द्रु\appbrk{}त\appbrk{}क्लीबा}\ H\textsuperscript{ac}; \skt{दुर्ब्ब\appbrk{}ल\appbrk{}डे\appbrk{}बा\appbrk{}ल\appbrk{}स्थ\appbrk{}वि\appbrk{}र\appbrk{}भी\appbrk{}रु\appbrk{}स\appbrk{}र्व्वा\appbrk{}ङ्ग\appbrk{}शू\appbrk{}नो\appbrk{}द\appbrk{}री\appbrk{}ग\appbrk{}र्भ्भि\appbrk{}णी\appbrk{}ऋ\appbrk{}तु\appbrk{}म\appbrk{}ती\appbrk{}प्र\appbrk{}वृ\appbrk{}द्ध\appbrk{}ज्व\appbrk{}री\appbrk{}प्र\appbrk{}मे\appbrk{}ही\appbrk{}रू\appbrk{}क्ष\appbrk{}क्ष\appbrk{}त\appbrk{}क्षी\appbrk{}ण\appbrk{}तृ\appbrk{}ष्णा\appbrk{}मू\appbrk{}र्च्छो\appbrk{}प\appbrk{}द्रु\appbrk{}त\appbrk{}क्लीबा}\ H\textsuperscript{pc}}}\skt{ }\edtext{\skt{उद्धृ\-तो\-द्वृ\-त्त\-फ\-ल\-यो\-न्यश्च}}{\lemma{\skt{उद्धृ\appbrk{}तो\appbrk{}द्वृ\appbrk{}त्त\appbrk{}फ\appbrk{}ल\appbrk{}यो\appbrk{}न्यश्च}}\Afootnote{\skt{पवृ\appbrk{}त्तो\appbrk{}द्वृ\appbrk{}त्त\appbrk{}फ\appbrk{}ल\appbrk{}यो\appbrk{}नयः}\ A; \skt{उद्धृ\appbrk{}तो\appbrk{}द्वृ\appbrk{}त्त\appbrk{}फ\appbrk{}ल\appbrk{}यो\appbrk{}न्यश्च}\ K; \skt{द्रद्वृ\appbrk{}तो\appbrk{}द्वृ\appbrk{}त्ते\appbrk{}फ\appbrk{}ल\appbrk{}यो\appbrk{}न्य\appbrk{}स्च\appbrk{}स्तथा}\ N; \skt{उद्धृ\appbrk{}तो\appbrk{}द्वृ\appbrk{}त्त\appbrk{}फ\appbrk{}ल\appbrk{}यो\appbrk{}न\appbrk{}यश्च}\ H}}\pend

\pstart \ledsidenote{\passref{1.11.29}}\edtext{\skt{तथा}}{\lemma{\skt{तथा}}\Afootnote{A K H; \emph{om.}\ N}}\skt{ }\edtext{\skt{मर्म\-सि\-रा\-स्ना\-यु\-स\-न्धि\-त\-रु\-णा\-स्थि\-से\-व\-नी\-ग\-ल\-ना\-भि\-न\-खा\-न्त\-र\-से\-प\-स्रोतः}}{\lemma{\skt{मर्म\appbrk{}सि\appbrk{}रा\appbrk{}स्ना\appbrk{}यु\appbrk{}स\appbrk{}न्धि\appbrk{}त\appbrk{}रु\appbrk{}णा\appbrk{}स्थि\appbrk{}से\appbrk{}व\appbrk{}नी\appbrk{}ग\appbrk{}ल\appbrk{}ना\appbrk{}भि\appbrk{}न\appbrk{}खा\appbrk{}न्त\appbrk{}र\appbrk{}से\appbrk{}प\appbrk{}स्रोतः}}\Afootnote{K N; \skt{मर्म\appbrk{}सि\appbrk{}रा\appbrk{}स्ना\appbrk{}यु\appbrk{}स\appbrk{}न्धि\appbrk{}त\appbrk{}रु\appbrk{}णा\appbrk{}स्थि\appbrk{}से\appbrk{}व\appbrk{}नी\appbrk{}ध\appbrk{}म\appbrk{}नी\appbrk{}ग\appbrk{}ल\appbrk{}ना\appbrk{}भि\appbrk{}न\appbrk{}खा\appbrk{}न्तः\appbrk{}शे\appbrk{}फः\appbrk{}स्रोतः}\ A; \skt{मर्म्म\appbrk{}सि\appbrk{}रा\appbrk{}स्ना\appbrk{}यु\appbrk{}स\appbrk{}न्धि\appbrk{}त\appbrk{}रु\appbrk{}णा\appbrk{}स्थि\appbrk{}से\appbrk{}व\appbrk{}नी\appbrk{}ग\appbrk{}ल\appbrk{}ना\appbrk{}भि\appbrk{}न\appbrk{}खा\appbrk{}न्त\appbrk{}र\appbrk{}से\appbrk{}फः\appbrk{}श्रोतः}\ H}}\skt{ }\edtext{\edtext{\skt{स्वल्प\-मां\-स\-प्र\-दे\-शे\-ष्व\-क्ष्णोश्च}}{\lemma{\skt{स्वल्प\appbrk{}मां\appbrk{}स\appbrk{}प्र\appbrk{}दे\appbrk{}शेष्व्}}\Afootnote{K N H; \skt{स्वल्पमांसेषु च देशेष्व्}\ A}}\skt{ न}}{\lemma{\skt{अक्ष्णोश्च न}}\Afootnote{A K H; \skt{अक्षांस्च न}\ N}}\skt{ दद्यादन्यत्र }\edtext{\skt{वर्मरोगात्}}{\lemma{\skt{वर्मरोगात्}}\Afootnote{K H; \skt{वर्त्मरोगात्}\ A; \skt{चर्मरोगात्}\ N}}\skt{~।}\pend

\pstart \ledsidenote{\passref{1.11.30}}\skt{तत्र क्षारसाध्येष्वपि व्याधिषु }\edtext{\edtext{\skt{शून\-गा\-त्र\-म\-स्थि\-शू\-लि\-न\-म\-न्न\-द्वे\-षिणं}}{\lemma{\skt{शूनगात्रम्}}\Afootnote{A K N; \skt{सूनगात्रम्}\ H}}}{\lemma{\skt{अन्नद्वेषिणं}}\Afootnote{A K H; \skt{अन्नद्वेसिनं}\ N}}\skt{ हृद\-य\-स\-न्धि\-पी\-डो\-प\-द्रु\-तञ्च क्षारो न }\edtext{\skt{साधयति}}{\lemma{\skt{साधयति}}\Afootnote{K N H; \skt{साधयति (दु\appbrk{}र्ब\appbrk{}ल\appbrk{}बा\appbrk{}ल\appbrk{}स्थ\appbrk{}वि\appbrk{}रा\appbrk{}दी\appbrk{}न्प्र\appbrk{}ति\appbrk{}सा\appbrk{}र\appbrk{}णीय इति)}\ A}}\skt{~॥}\pend

\pstart\versepar \ledsidenote{\passref{1.11.31}}\skt{भवति चात्र~॥ }\pend
\pstart\versepar \skt{विषा\-ग्नि\-श\-स्त्रा\-श\-नि\-मृ\-त्यु\-तुल्यः क्षारो भव\-त्य\-ल्प\-म\-ति\-प्र\-युक्तः~। }\pend
\pstart\versepar \skt{}\edtext{\skt{सत्वप्रमत्तेन}}{\lemma{\skt{सत्वप्रमत्तेन}}\Afootnote{K H; \skt{स धीमता सम्य}\ A; \skt{सन्धप्रमत्तेन}\ N}}\skt{ }\edtext{\skt{सदा}}{\lemma{\skt{सदा}}\Afootnote{K N H; \skt{गनु}\ A}}\skt{ }\edtext{\skt{प्रयुक्तो}}{\lemma{\skt{प्रयुक्तो}}\Afootnote{A K H; \skt{प्रयुक्ते}\ N}}\skt{ }\edtext{\skt{रोगा\-न्नि\-ह\-न्या\-द\-चि\-रेण}}{\lemma{\skt{रोगान्}}\Afootnote{A K H; \skt{रोगा}\ N}}\skt{ }\edtext{\skt{घोरानिति}}{\lemma{\skt{घोरानिति}}\Afootnote{K N H; \skt{घोरान्}\ A}}\skt{~॥}\edtext{}{\lemma{\skt{भवति चात्र~॥ विषा\appbrk{}ग्नि\appbrk{}श\appbrk{}स्त्रा\appbrk{}श\appbrk{}नि\appbrk{}मृ\appbrk{}त्यु\appbrk{}तुल्यः}}\Afootnote{\skt{भवति चात्र विषा\appbrk{}ग्नि\appbrk{}श\appbrk{}स्त्रा\appbrk{}श\appbrk{}नि\appbrk{}मृ\appbrk{}त्यु\appbrk{}कल्पः}\ A; \skt{विषा\appbrk{}ग्नि\appbrk{}श\appbrk{}स्त्रा\appbrk{}श\appbrk{}नि\appbrk{}मृ\appbrk{}त्यु\appbrk{}तुल्यः}\ K H; \skt{भवति चात्र~॥ विषा\appbrk{}ग्नि\appbrk{}श\appbrk{}स्त्रा\appbrk{}स\appbrk{}नि\appbrk{}मृ\appbrk{}त्यु\appbrk{}तुल्यः}\ N}}\pend

\pstart \ledsidenote{\passref{1.11.trailer}}\edtext{\skt{॥ ११~॥}}{\lemma{\skt{॥ ११~॥}}\Afootnote{\skt{इति सुश्रु\appbrk{}त\appbrk{}सं\appbrk{}हि\appbrk{}तायां सूत्रस्थाने क्षार\appbrk{}पा\appbrk{}क\appbrk{}वि\appbrk{}धि\appbrk{}र्ना\appbrk{}मै\appbrk{}का\appbrk{}द\appbrk{}शोऽ\appbrk{}ध्यायः}\ A; \skt{{\symfont ❈}~॥ ११}\ K; \skt{डों}\ N; \skt{डो ए}\ H}}\pend

\endnumbering
\beginnumbering
\eledsection*{Adhyāya 12}
\pstart \ledsidenote{\passref{1.12.1}}\edtext{\edtext{\edtext{\skt{अथा\-तोऽ\-ग्नि\-क\-र्म\-वि\-धि\-म\-ध्यायां}}{\lemma{\skt{अथातो}}\Afootnote{A L N H; \skt{अथा}\ K}}}{\lemma{\skt{ऽग्नि\appbrk{}क\appbrk{}र्म\appbrk{}वि\appbrk{}धिम्}}\Afootnote{A L N H; \skt{ग्निकर्मविधिं}\ K}}}{\lemma{\skt{अध्यायां}}\Afootnote{N; \skt{अध्यायं}\ A L; \emph{om.}\ K H}}\skt{ }\edtext{\skt{व्याख्यास्यामः}}{\lemma{\skt{व्याख्यास्यामः}}\Afootnote{A N H; \skt{व्याख्यास्यामः~॥…}\ L; \emph{om.}\ K}}\skt{~॥}\pend

\pstart \ledsidenote{\passref{1.12.3}}\edtext{\edtext{\skt{क्षारा\-द\-ग्नि\-र्ग\-रीयां}}{\lemma{\skt{क्षारादग्निर्}}\Afootnote{K N H; \skt{क्षारादग्नि}\ A; \skt{…क्षारादग्निर्}\ L}}}{\lemma{\skt{गरीयां}}\Afootnote{K N; \skt{र्गरीया}\ A; \skt{गरीया}\ L H}}\skt{ }\edtext{\skt{क्रियासु}}{\lemma{\skt{क्रियासु}}\Afootnote{K N; \skt{न्क्रियासु}\ A L H}}\skt{ }\edtext{\edtext{\skt{व्याख्या\-त\-स्त\-द्द\-ग्धानां}}{\lemma{\skt{व्याख्यातस्}}\Afootnote{A K\textsuperscript{pc}; \skt{व्याख्यातः~॥ …}\ L; \skt{व्याख्यातआस्}\ K\textsuperscript{ac}; \skt{व्याख्यातास्}\ N H}}}{\lemma{\skt{तद्दग्धानां}}\Afootnote{A K N H; \skt{तद्दग्धानां}\ L}}\skt{ }\edtext{\skt{रोगा\-णा\-म\-पु\-न\-र्भ\-वात्~। }\edtext{\skt{स्वेद\-श\-स्त्र\-क्षा\-रै\-र\-श\-क्यानां}}{\lemma{\skt{अशक्यानां}}\Afootnote{L K H; \skt{असाध्यानां}\ A; \skt{असक्षानां}\ N}}}{\lemma{\skt{अपुनर्भवात्~। स्वेद\appbrk{}श\appbrk{}स्त्र\appbrk{}क्षा\appbrk{}रैर्}}\Afootnote{K H\textsuperscript{ac}; \skt{अपु\appbrk{}न\appbrk{}र्भा\appbrk{}वा\appbrk{}द्भे\appbrk{}ष\appbrk{}ज\appbrk{}श\appbrk{}स्त्र\appbrk{}क्षा\appbrk{}रैर्}\ A; \skt{अपुनर्भवात्~॥ … भेष\appbrk{}ज\appbrk{}श\appbrk{}स्त्र\appbrk{}क्षा\appbrk{}रैर्}\ L; \skt{अपु\appbrk{}न\appbrk{}र्भा\appbrk{}वा\appbrk{}त्स्वे\appbrk{}द\appbrk{}स\appbrk{}स्त्र\appbrk{}क्षा\appbrk{}रैर्}\ N; \skt{अपुनर्भवात्~। स्वेद\appbrk{}श\appbrk{}स्त्र\appbrk{}त\appbrk{}क्षा\appbrk{}रैर्}\ H\textsuperscript{pc}}}\skt{ }\edtext{\skt{तत्साधनाच्च}}{\lemma{\skt{तत्साधनाच्च}}\Afootnote{K N H; \skt{तत्साध्यत्वाच्च}\ A L}}\skt{~।}\pend

\pstart \ledsidenote{\passref{1.12.4}}\edtext{\skt{अथैमानि}}{\lemma{\skt{अथैमानि}}\Afootnote{K N H; \skt{अथेमानि}\ A; \skt{…अथैमानि}\ L}}\skt{ दह\-नो\-प\-क\-र\-णानि }\edtext{\skt{भवन्ति}}{\lemma{\skt{भवन्ति}}\Afootnote{A K H; \emph{om.}\ L; \skt{भव्न्ति}\ N}}\skt{~। }\edtext{\skt{पिप्प\-ल्य\-जा\-श\-कृ\-द्गो\-द\-न्त\-श\-र\-श\-ला\-का\-ज\-म्व्वो\-ष्ठे\-त\-र\-लो\-ह\-क्षौ\-द्र\-गु\-ड\-स्ने\-हा\-दीनि}}{\lemma{\skt{पिप्प\appbrk{}ल्य\appbrk{}जा\appbrk{}श\appbrk{}कृ\appbrk{}द्गो\appbrk{}द\appbrk{}न्त\appbrk{}श\appbrk{}र\appbrk{}श\appbrk{}ला\appbrk{}का\appbrk{}ज\appbrk{}म्व्वो\appbrk{}ष्ठे\appbrk{}त\appbrk{}र\appbrk{}लो\appbrk{}ह\appbrk{}क्षौ\appbrk{}द्र\appbrk{}गु\appbrk{}ड\appbrk{}स्ने\appbrk{}हा\appbrk{}दीनि}}\Afootnote{K; \skt{तद्यथा- पिप्प\appbrk{}ल्य\appbrk{}जा\appbrk{}श\appbrk{}कृ\appbrk{}द्गो\appbrk{}द\appbrk{}न्त\appbrk{}श\appbrk{}र\appbrk{}श\appbrk{}ला\appbrk{}का\appbrk{}जा\appbrk{}म्ब\appbrk{}वौ\appbrk{}ष्ठे\appbrk{}त\appbrk{}र\appbrk{}लौहाः क्षौद्र\appbrk{}गु\appbrk{}ड\appbrk{}स्ने\appbrk{}हाश्च}\ A; \skt{तद्यथा~॥ पिप्प\appbrk{}ल्य\appbrk{}जा\appbrk{}श\appbrk{}कृ\appbrk{}द्गो\appbrk{}दं\appbrk{}त\appbrk{}श\appbrk{}र\appbrk{}श\appbrk{}ला\appbrk{}का\appbrk{}जं\appbrk{}आ\appbrk{}ब\appbrk{}वौ\appbrk{}ष्ट्वे\appbrk{}त\appbrk{}त\appbrk{}र\appbrk{}लौहाः क्षौद्र\appbrk{}गु\appbrk{}ड\appbrk{}स्ने\appbrk{}हाश्च~॥ …}\ L\textsuperscript{ac}; \skt{तद्यथा~॥ पिप्प\appbrk{}ल्य\appbrk{}जा\appbrk{}श\appbrk{}कृ\appbrk{}द्गो\appbrk{}दं\appbrk{}त\appbrk{}श\appbrk{}र\appbrk{}श\appbrk{}ला\appbrk{}का\appbrk{}जं\appbrk{}ब\appbrk{}वौ\appbrk{}ष्ट्वे\appbrk{}त\appbrk{}त\appbrk{}र\appbrk{}लौहाः क्षौद्र\appbrk{}गु\appbrk{}ड\appbrk{}स्ने\appbrk{}हाश्च~॥ …}\ L\textsuperscript{pc}; \skt{पिप्प\appbrk{}ल्य\appbrk{}जा\appbrk{}स\appbrk{}कृ\appbrk{}द्गो\appbrk{}द\appbrk{}न्त\appbrk{}स\appbrk{}र\appbrk{}स\appbrk{}ला\appbrk{}का\appbrk{}जां\appbrk{}बो\appbrk{}ष्ठो\appbrk{}त\appbrk{}र\appbrk{}लो\appbrk{}ह\appbrk{}क्षौ\appbrk{}द्र\appbrk{}गु\appbrk{}ड\appbrk{}स्ने\appbrk{}हा\appbrk{}दीनि}\ N; \skt{पिप्प\appbrk{}ल्य\appbrk{}जा\appbrk{}स\appbrk{}कृ\appbrk{}द्गो\appbrk{}द\appbrk{}न्त\appbrk{}श\appbrk{}र\appbrk{}श\appbrk{}ला\appbrk{}का\appbrk{}जं\appbrk{}ब्बो\appbrk{}ष्ठे\appbrk{}त\appbrk{}र\appbrk{}ला\appbrk{}ह\appbrk{}क्षौ\appbrk{}द्र\appbrk{}गु\appbrk{}ड\appbrk{}स्ने\appbrk{}हा\appbrk{}दीनि}\ H}}\skt{~। तत्र }\edtext{\skt{पिप्प\-ल्य\-जा\-श\-कृ\-द्गो\-द\-न्त\-श\-र\-श\-ला\-का\-स्त्व\-ग्ग\-तानां}}{\lemma{\skt{त्वग्गतानां}}\Afootnote{A L N H; \skt{स्त्वग्गतानां}\ K}}\skt{~। }\edtext{\skt{जम्ब्वो\-ष्ठे\-त\-र\-लोहाः}}{\lemma{\skt{जम्ब्वो\appbrk{}ष्ठे\appbrk{}त\appbrk{}र\appbrk{}लोहाः}}\Afootnote{K N; \skt{जाम्ब\appbrk{}वौ\appbrk{}ष्ठे\appbrk{}त\appbrk{}र\appbrk{}लौहा}\ A; \skt{… जांवौ\appbrk{}ष्ठे\appbrk{}त\appbrk{}र\appbrk{}लौहा}\ L\textsuperscript{ac}; \skt{… जांब\appbrk{}वौ\appbrk{}ष्ठे\appbrk{}त\appbrk{}र\appbrk{}लौहा}\ L\textsuperscript{pc}; \skt{जंब्बो\appbrk{}ष्ठे\appbrk{}त\appbrk{}र\appbrk{}लाहा}\ H}}\skt{ }\edtext{\skt{मांसगतानां}}{\lemma{\skt{मांसगतानां}}\Afootnote{A K N H; \skt{मांसगतानां~॥ …}\ L}}\skt{~। क्षौद्रगुडस्नेहाः }\edtext{\skt{सिरा\-स्ना\-यु\-स\-न्ध्य\-स्थि\-ग\-ता\-नाम्}}{\lemma{\skt{सिरा\appbrk{}स्ना\appbrk{}यु\appbrk{}स\appbrk{}न्ध्य\appbrk{}स्थि\appbrk{}ग\appbrk{}ता\appbrk{}नाम्}}\Afootnote{A N H; \skt{शिरा\appbrk{}स्ना\appbrk{}यु\appbrk{}सं\appbrk{}ध्य\appbrk{}स्थि\appbrk{}ग\appbrk{}तानां~॥…}\ L; \skt{सिरा\appbrk{}स्ना\appbrk{}यु\appbrk{}स\appbrk{}न्ध्य\appbrk{}स्थि\appbrk{}ग\appbrk{}ताना}\ K}}\skt{~।}\pend

\pstart \ledsidenote{\passref{1.12.5}}\edtext{\skt{तत्राग्निकर्म}}{\lemma{\skt{तत्राग्निकर्म}}\Afootnote{A K N H; \emph{om.}\ L}}\skt{ }\edtext{\skt{सर्वर्तुषु}}{\lemma{\skt{सर्वर्तुषु}}\Afootnote{A K N H; \emph{om.}\ L}}\skt{ }\edtext{\skt{कुर्यादन्यत्र}}{\lemma{\skt{कुर्यादन्यत्र}}\Afootnote{A K N H; \emph{om.}\ L}}\skt{ }\edtext{\skt{शरद्ग्रीष्माभ्यां}}{\lemma{\skt{शरद्ग्रीष्माभ्यां}}\Afootnote{A K N H; \skt{…}\ L}}\skt{~। }\edtext{\edtext{\skt{तत्रा\-प्या\-त्य\-यि\-केऽ\-ग्नि\-साध्ये}}{\lemma{\skt{तत्राप्यात्ययिके}}\Afootnote{A L K H; \skt{तत्राप्यात्ययिके}\ N}}}{\lemma{\skt{ऽग्निसाध्ये}}\Afootnote{H; \skt{ऽग्निकर्मसाध्ये}\ A L; \skt{ग्निसाध्ये}\ K N}}\skt{ व्याधौ तत्प्रत्यनीकं }\edtext{\skt{विधिं}}{\lemma{\skt{विधिं}}\Afootnote{A K N H; \skt{विधि}\ L}}\skt{ }\edtext{\skt{कृत्वा}}{\lemma{\skt{कृत्वा}}\Afootnote{A K N H; \skt{कृत्वेति~॥ …}\ L}}\pend

\pstart \ledsidenote{\passref{1.12.6}}\edtext{\skt{सर्व\-व्या\-धि\-ष्वृ\-तुषु च}}{\lemma{\skt{सर्व\appbrk{}व्या\appbrk{}धि\appbrk{}ष्वृ\appbrk{}तुषु च}}\Afootnote{A K N H; \skt{… सर्वव्याधिषु च}\ L}}\skt{ }\edtext{\skt{पिच्छि\-ल\-म\-न्न\-म्भु\-क्त\-वतः}}{\lemma{\skt{पिच्छि\appbrk{}ल\appbrk{}म\appbrk{}न्न\appbrk{}म्भु\appbrk{}क्त\appbrk{}वतः}}\Afootnote{H; \skt{पिच्छिलमन्नं भुक्तवतः; मूढ\appbrk{}ग\appbrk{}र्भा\appbrk{}श्म\appbrk{}री\appbrk{}भ\appbrk{}ग\appbrk{}न्द\appbrk{}रो\appbrk{}द\appbrk{}रा\appbrk{}र्शो\appbrk{}मु\appbrk{}ख\appbrk{}रो\appbrk{}गे\appbrk{}ष्व\appbrk{}भु\appbrk{}क्त\appbrk{}वतः कर्म कुर्वीत}\ A; \skt{पिछलमन्नं भुक्तवतः~॥ … मूढ\appbrk{}ग\appbrk{}र्भा\appbrk{}श्म\appbrk{}री\appbrk{}भ\appbrk{}गं\appbrk{}द\appbrk{}रो\appbrk{}द\appbrk{}रा\appbrk{}र्शो\appbrk{}मु\appbrk{}ख\appbrk{}रो\appbrk{}गे\appbrk{}ष्व\appbrk{}भु\appbrk{}क्त\appbrk{}वतः~॥ …}\ L; \skt{पिच्छि\appbrk{}ल\appbrk{}म\appbrk{}न्न\appbrk{}म्भु\appbrk{}क्त\appbrk{}वतः}\ K N}}\skt{~।}\pend

\pstart \ledsidenote{\passref{1.12.7}}\edtext{\skt{तत्र}}{\lemma{\skt{तत्र}}\Afootnote{A K N H; \skt{… तत्र}\ L}}\skt{ }\edtext{\edtext{\skt{द्विवि\-ध\-म\-ग्नि\-द\-ग्ध\-मा\-हु\-रेके}}{\lemma{\skt{द्विविधम्}}\Afootnote{N; \skt{द्विविध}\ A K H; \skt{द्विधम्}\ L}}}{\lemma{\skt{अग्नि\appbrk{}द\appbrk{}ग्ध\appbrk{}मा\appbrk{}हु\appbrk{}रेके}}\Afootnote{\skt{मग्नि\appbrk{}क\appbrk{}र्मा\appbrk{}हु\appbrk{}रेके}\ A; \skt{अग्नि\appbrk{}क\appbrk{}र्मा\appbrk{}हु\appbrk{}रेके}\ L; \skt{मग्नि\appbrk{}द\appbrk{}ग्ध\appbrk{}मा\appbrk{}हु\appbrk{}रेके}\ K H; \skt{अग्नि\appbrk{}द\appbrk{}ग्ध\appbrk{}मा\appbrk{}हु\appbrk{}रेके}\ N}}\skt{~। त्वग्दग्धं मांसदग्धञ्च~। इह तु }\edtext{\skt{सिरास्नायुषु}}{\lemma{\skt{सिरास्नायुषु}}\Afootnote{K H; \skt{सिरास्नायु}\ A L; \skt{शिरास्नायु}\ N}}\skt{ }\edtext{\edtext{\skt{सन्ध्यस्थिष्वपि}}{\lemma{\skt{सन्ध्यस्थिष्व्}}\Afootnote{A K N H; \skt{संध्यस्थि}\ L}}\skt{ न}}{\lemma{\skt{अपि न}}\Afootnote{A K N H; \skt{ष्वपि न}\ L}}\skt{ }\edtext{\skt{प्रतिषिद्धोऽग्निः}}{\lemma{\skt{ऽग्निः}}\Afootnote{A H; \skt{ऽग्निः~॥ …}\ L; \skt{ग्निः}\ K N}}\skt{~।}\pend

\pstart \ledsidenote{\passref{1.12.8}}\edtext{\skt{तत्र}}{\lemma{\skt{तत्र}}\Afootnote{A K N H; \skt{… तत्र}\ L}}\skt{ }\edtext{\skt{शब्दप्रादुर्भावो}}{\lemma{\skt{शब्दप्रादुर्भावो}}\Afootnote{A K N H; \skt{शब्दप्रादुर्भावा}\ L}}\skt{ दुर्गन्धता }\edtext{\skt{त्वक्संकोचश्च}}{\lemma{\skt{त्वक्संकोचश्च}}\Afootnote{A K H; \skt{त्वग्संकोचश्च}\ L; \skt{त्वक्संकोचस्}\ N}}\skt{ त्वग्दग्धे~। }\edtext{\skt{कपो\-त\-व\-र्ण्ण\-ता\-ल्प\-श्व\-य\-थु\-वे\-द\-नता}}{\lemma{\skt{कपो\appbrk{}त\appbrk{}व\appbrk{}र्ण्ण\appbrk{}ता\appbrk{}ल्प\appbrk{}श्व\appbrk{}य\appbrk{}थु\appbrk{}वे\appbrk{}द\appbrk{}नता}}\Afootnote{K N H; \skt{कपो\appbrk{}त\appbrk{}व\appbrk{}र्ण\appbrk{}ताऽ\appbrk{}ल्प\appbrk{}श्व\appbrk{}य\appbrk{}थु\appbrk{}वे\appbrk{}दना}\ A; \skt{कपोतवर्णता अल्प\appbrk{}श्व\appbrk{}य\appbrk{}थु\appbrk{}वे\appbrk{}द\appbrk{}नता}\ L}}\skt{ }\edtext{\skt{शुष्क\-सं\-कु\-चि\-त\-व्र\-णता च}}{\lemma{\skt{शुष्क\appbrk{}सं\appbrk{}कु\appbrk{}चि\appbrk{}त\appbrk{}व्र\appbrk{}णता च}}\Afootnote{A K; \skt{शुष्क\appbrk{}सं\appbrk{}कु\appbrk{}चि\appbrk{}त\appbrk{}व्र\appbrk{}णता}\ L; \skt{शुष्क\appbrk{}सं\appbrk{}कु\appbrk{}चि\appbrk{}त\appbrk{}व्र\appbrk{}णाता च}\ N; \skt{शुष्क\appbrk{}सं\appbrk{}कु\appbrk{}ञ्चि\appbrk{}त\appbrk{}व्र\appbrk{}णता च}\ H\textsuperscript{ac}; \skt{शुष्क\appbrk{}सं\appbrk{}कु\appbrk{}कु\appbrk{}ञ्चि\appbrk{}चि\appbrk{}त\appbrk{}व्र\appbrk{}णता च}\ H\textsuperscript{pc}}}\skt{ }\edtext{\skt{मांसदग्धे}}{\lemma{\skt{मांसदग्धे}}\Afootnote{A L K N H\textsuperscript{pc}; \skt{मांमदग्धे}\ H\textsuperscript{ac}}}\skt{~। }\edtext{\skt{कृष्णो\-न्न\-त\-व्र\-णता}}{\lemma{\skt{कृष्णो\appbrk{}न्न\appbrk{}त\appbrk{}व्र\appbrk{}णता}}\Afootnote{A L K H; \skt{कृष्णो\appbrk{}न्न\appbrk{}त\appbrk{}व्र\appbrk{}णाता}\ N}}\skt{ }\edtext{\skt{स्राव\-स\-न्नि\-रो\-धश्च}}{\lemma{\skt{स्राव\appbrk{}स\appbrk{}न्नि\appbrk{}रो\appbrk{}धश्च}}\Afootnote{A L K N H\textsuperscript{pc}; \skt{स्राव\appbrk{}म\appbrk{}न्नि\appbrk{}रो\appbrk{}धश्च}\ H\textsuperscript{ac}}}\skt{ }\edtext{\skt{सिरास्नायुदग्धे}}{\lemma{\skt{सिरास्नायुदग्धे}}\Afootnote{A K N H\textsuperscript{pc}; \skt{शिरास्नायुदग्धे}\ L; \skt{मिरास्नायुदग्धे}\ H\textsuperscript{ac}}}\skt{ }\edtext{\skt{रूक्ष\-ता\-रु\-ण\-ता\-क\-र्क\-श\-स्थि\-र\-व्र\-णता च}}{\lemma{\skt{रूक्ष\appbrk{}ता\appbrk{}रु\appbrk{}ण\appbrk{}ता\appbrk{}क\appbrk{}र्क\appbrk{}श\appbrk{}स्थि\appbrk{}र\appbrk{}व्र\appbrk{}णता च}}\Afootnote{K N H; \skt{रूक्षारुणता कर्क\appbrk{}श\appbrk{}स्थि\appbrk{}र\appbrk{}व्र\appbrk{}णता च}\ A; \skt{रुक्षारुणता कर्क्क\appbrk{}श\appbrk{}स्थि\appbrk{}र\appbrk{}व्र\appbrk{}णता च}\ L}}\skt{ }\edtext{\skt{सन्ध्यस्थिदग्धे}}{\lemma{\skt{सन्ध्यस्थिदग्धे}}\Afootnote{A K N H; \skt{संध्यस्थिदग्धे~॥ …}\ L}}\skt{~॥}\pend

\pstart \ledsidenote{\passref{1.12.9}}\edtext{\skt{तत्र}}{\lemma{\skt{तत्र}}\Afootnote{A K N H; \skt{… तत्र}\ L}}\skt{ }\edtext{\skt{शिरो\-रो\-गा\-धि\-म\-न्थ\-यो\-र्भ्रू\-ल\-ला\-ट\-श\-ङ्ख\-दे\-शेषु}}{\lemma{\skt{शिरो\appbrk{}रो\appbrk{}गा\appbrk{}धि\appbrk{}म\appbrk{}न्थ\appbrk{}यो\appbrk{}र्भ्रू\appbrk{}ल\appbrk{}ला\appbrk{}ट\appbrk{}श\appbrk{}ङ्ख\appbrk{}दे\appbrk{}शेषु}}\Afootnote{K; \skt{शिरो\appbrk{}रो\appbrk{}गा\appbrk{}धि\appbrk{}म\appbrk{}न्थ\appbrk{}यो\appbrk{}र्भ्रू\appbrk{}ल\appbrk{}ला\appbrk{}ट\appbrk{}श\appbrk{}ङ्ख\appbrk{}प्र\appbrk{}दे\appbrk{}शेषु}\ A; \skt{शिरोरोगाधि अक्षिमंथयो भ्रूल\appbrk{}ला\appbrk{}ट\appbrk{}शं\appbrk{}ख\appbrk{}प्र\appbrk{}दे\appbrk{}षे\appbrk{}शेषु}\ L\textsuperscript{ac}; \skt{शिरोरोगाधि अक्षिमंथयो भ्रूल\appbrk{}ला\appbrk{}ट\appbrk{}शं\appbrk{}ख\appbrk{}प्र\appbrk{}दे\appbrk{}शेषु}\ L\textsuperscript{pc}; \skt{शिरो\appbrk{}रो\appbrk{}गा\appbrk{}थि\appbrk{}म\appbrk{}न्थ\appbrk{}यो\appbrk{}र्भ्रू\appbrk{}ल\appbrk{}ला\appbrk{}टा\appbrk{}शं\appbrk{}ख\appbrk{}दे\appbrk{}शेषु}\ N; \skt{शिरो\appbrk{}रो\appbrk{}गा\appbrk{}धि\appbrk{}म\appbrk{}न्थ\appbrk{}यो\appbrk{}र्ल\appbrk{}ला\appbrk{}ट\appbrk{}शं\appbrk{}ख\appbrk{}दे\appbrk{}शेषु}\ H}}\skt{ }\edtext{\edtext{\skt{दहे\-द्व\-र्त्म\-रो\-गे\-ष्वा\-र्द्र\-न\-क्त\-क\-प्र\-ति\-च्छ\-न्ना\-न्दृष्टिं}}{\lemma{\skt{दहे\appbrk{}द्व\appbrk{}र्त्म\appbrk{}रो\appbrk{}गेष्व्}}\Afootnote{A K; \skt{दहेत्~॥ … वर्त्मरोगेष्व्}\ L; \skt{दहे\appbrk{}द्च\appbrk{}र्म\appbrk{}रो\appbrk{}गेष्व्}\ N; \skt{दहे\appbrk{}द्व\appbrk{}र्त्म्म\appbrk{}रो\appbrk{}गेष्व्}\ H}}}{\lemma{\skt{आर्द्र\appbrk{}न\appbrk{}क्त\appbrk{}क\appbrk{}प्र\appbrk{}ति\appbrk{}च्छ\appbrk{}न्नान्}}\Afootnote{L K N H; \skt{आर्द्रा\appbrk{}ल\appbrk{}क्त\appbrk{}क\appbrk{}प्र\appbrk{}ति\appbrk{}च्छन्नां}\ A}}\skt{ कृत्वा }\edtext{\skt{वर्त्मरोमकूपां}}{\lemma{\skt{वर्त्मरोमकूपां}}\Afootnote{K; \skt{वर्त्म\appbrk{}रो\appbrk{}म\appbrk{}कू\appbrk{}पान्}\ A; \skt{वर्म\appbrk{}रो\appbrk{}म\appbrk{}कु\appbrk{}पान् …}\ L; \skt{वर्मरोमकूपां}\ N; \skt{वर्त्म्मरोमकूपां}\ H}}\skt{~।}\pend

\pstart \ledsidenote{\passref{1.12.10}}\edtext{\skt{त्वङ्मां\-स\-सि\-रा\-स्ना\-यु\-स\-न्ध्य\-स्थि\-ग\-त\-मु\-ग्र\-रुजौ}}{\lemma{\skt{त्वङ्मां\appbrk{}स\appbrk{}सि\appbrk{}रा\appbrk{}स्ना\appbrk{}यु\appbrk{}स\appbrk{}न्ध्य\appbrk{}स्थि\appbrk{}ग\appbrk{}त\appbrk{}मु\appbrk{}ग्र\appbrk{}रुजौ}}\Afootnote{\skt{त्वङ्मां\appbrk{}स\appbrk{}सि\appbrk{}रा\appbrk{}स्ना\appbrk{}यु\appbrk{}स\appbrk{}न्ध्य\appbrk{}स्थि\appbrk{}स्थि\appbrk{}तेऽ\appbrk{}त्यु\appbrk{}ग्र\appbrk{}रुजि}\ A; \skt{… त्वग्मां\appbrk{}स\appbrk{}शि\appbrk{}रा\appbrk{}स्ना\appbrk{}यु\appbrk{}सं\appbrk{}ध्य\appbrk{}स्थि\appbrk{}स्थिते त्युग्ररुजि}\ L; \skt{त्वङ्मां\appbrk{}स\appbrk{}सि\appbrk{}रा\appbrk{}स्ना\appbrk{}यु\appbrk{}स\appbrk{}न्ध्य\appbrk{}स्थि\appbrk{}ग\appbrk{}त\appbrk{}मु\appbrk{}ग्र\appbrk{}रुजौ}\ K; \skt{त्वङ्म\appbrk{}न्स\appbrk{}सि\appbrk{}रा\appbrk{}स्ना\appbrk{}यु\appbrk{}स\appbrk{}न्ध्य\appbrk{}स्थि\appbrk{}ग\appbrk{}त\appbrk{}मु\appbrk{}ग्रु\appbrk{}रोजौ}\ N; \skt{त्वद्यां\appbrk{}स\appbrk{}सि\appbrk{}रा\appbrk{}स्ना\appbrk{}यु\appbrk{}स\appbrk{}न्ध्य\appbrk{}स्थि\appbrk{}ग\appbrk{}त\appbrk{}मु\appbrk{}ग्र\appbrk{}रुजं}\ H\textsuperscript{ac}; \skt{त्वद्यां\appbrk{}स\appbrk{}सि\appbrk{}रा\appbrk{}स्ना\appbrk{}यु\appbrk{}स\appbrk{}न्ध्य\appbrk{}स्थि\appbrk{}ग\appbrk{}त\appbrk{}मु\appbrk{}ग्र\appbrk{}रु\appbrk{}जंजे वायुं}\ H\textsuperscript{pc}}}\skt{ }\edtext{\skt{वायौ}}{\lemma{\skt{वायौ}}\Afootnote{A L K N H\textsuperscript{pc}; \skt{वायुं}\ H\textsuperscript{ac}}}\skt{~। }\edtext{\skt{दुष्ट\-व्र\-ण\-मु\-च्छि\-त\-क\-ठि\-न\-मां\-स\-ग\-ग्र\-न्थ्य\-र्बु\-दा\-प\-ची\-ग\-ल\-ग\-ण्ड\-गृ\-द्ध्र\-सी\-म\-स\-क\-गु\-ल्मो\-दर भग\-न्द\-रा\-र्श\-स\-न्धि\-श्ली\-प\-द\-च\-र्म\-की\-ल\-ति\-ल\-का\-ल\-क\-सि\-रा\-च्छे\-द\-ना\-डी\-शो\-णि\-ता\-ति\-प्र\-वृ\-त्तिषु}}{\lemma{\skt{दुष्ट\appbrk{}व्र\appbrk{}ण\appbrk{}मु\appbrk{}च्छि\appbrk{}त\appbrk{}क\appbrk{}ठि\appbrk{}न\appbrk{}मां\appbrk{}स\appbrk{}ग\appbrk{}ग्र\appbrk{}न्थ्य\appbrk{}र्बु\appbrk{}दा\appbrk{}प\appbrk{}ची\appbrk{}ग\appbrk{}ल\appbrk{}ग\appbrk{}ण्ड\appbrk{}गृ\appbrk{}द्ध्र\appbrk{}सी\appbrk{}म\appbrk{}स\appbrk{}क\appbrk{}गु\appbrk{}ल्मो\appbrk{}दर भग\appbrk{}न्द\appbrk{}रा\appbrk{}र्श\appbrk{}स\appbrk{}न्धि\appbrk{}श्ली\appbrk{}प\appbrk{}द\appbrk{}च\appbrk{}र्म\appbrk{}की\appbrk{}ल\appbrk{}ति\appbrk{}ल\appbrk{}का\appbrk{}ल\appbrk{}क\appbrk{}सि\appbrk{}रा\appbrk{}च्छे\appbrk{}द\appbrk{}ना\appbrk{}डी\appbrk{}शो\appbrk{}णि\appbrk{}ता\appbrk{}ति\appbrk{}प्र\appbrk{}वृ\appbrk{}त्तिषु}}\Afootnote{K; \skt{उच्छ्रि\appbrk{}त\appbrk{}क\appbrk{}ठि\appbrk{}न\appbrk{}सु\appbrk{}प्त\appbrk{}मांसे व्रणे ग्रन्थ्य\appbrk{}र्शोऽ\appbrk{}र्बु\appbrk{}द\appbrk{}भ\appbrk{}ग\appbrk{}न्द\appbrk{}रा\appbrk{}प\appbrk{}ची\appbrk{}श्ली\appbrk{}प\appbrk{}द\appbrk{}च\appbrk{}र्म\appbrk{}की\appbrk{}ल\appbrk{}ति\appbrk{}ल\appbrk{}का\appbrk{}ल\appbrk{}का\appbrk{}न्त्र\appbrk{}वृ\appbrk{}द्धि\appbrk{}स\appbrk{}न्धि\appbrk{}सि\appbrk{}रा\appbrk{}च्छे\appbrk{}द\appbrk{}ना\appbrk{}दिषु नाडी\appbrk{}शो\appbrk{}णि\appbrk{}ता\appbrk{}ति\appbrk{}प्र\appbrk{}वृ\appbrk{}त्तिषु}\ A; \skt{उछ्रि\appbrk{}त\appbrk{}क\appbrk{}ठि\appbrk{}न\appbrk{}सु\appbrk{}प्त\appbrk{}मांसे व्रद्रणे ग्रंथ्य\appbrk{}र्शो\appbrk{}र्बु\appbrk{}द\appbrk{}भ\appbrk{}गं\appbrk{}दरा पची\appbrk{}श्ली\appbrk{}प\appbrk{}द\appbrk{}च\appbrk{}र्म\appbrk{}की\appbrk{}ल\appbrk{}ति\appbrk{}ल\appbrk{}का\appbrk{}ल\appbrk{}कां\appbrk{}त्र\appbrk{}वृ\appbrk{}द्धि\appbrk{}सं\appbrk{}धि\appbrk{}सि\appbrk{}रा\appbrk{}छे\appbrk{}द\appbrk{}ना\appbrk{}दिषु नाडै\appbrk{}ई\appbrk{}शो\appbrk{}णि\appbrk{}ता\appbrk{}ति\appbrk{}प्र\appbrk{}त्तिषु}\ L\textsuperscript{ac}; \skt{उछ्रि\appbrk{}त\appbrk{}क\appbrk{}ठि\appbrk{}न\appbrk{}सु\appbrk{}प्त\appbrk{}मांसे व्रणे ग्रंथ्य\appbrk{}र्शो\appbrk{}र्बु\appbrk{}द\appbrk{}भ\appbrk{}गं\appbrk{}दरा पची\appbrk{}श्ली\appbrk{}प\appbrk{}द\appbrk{}च\appbrk{}र्म\appbrk{}की\appbrk{}ल\appbrk{}ति\appbrk{}ल\appbrk{}का\appbrk{}ल\appbrk{}कां\appbrk{}त्र\appbrk{}वृ\appbrk{}द्धि\appbrk{}सं\appbrk{}धि\appbrk{}सि\appbrk{}रा\appbrk{}छे\appbrk{}द\appbrk{}ना\appbrk{}दिषु नाड\appbrk{}ई\appbrk{}शो\appbrk{}णि\appbrk{}ता\appbrk{}ति\appbrk{}प्र\appbrk{}वृ\appbrk{}त्तिषु}\ L\textsuperscript{pc}; \skt{दुष्ट\appbrk{}व्र\appbrk{}ण\appbrk{}मु\appbrk{}च्छ्रि\appbrk{}त\appbrk{}क\appbrk{}ठि\appbrk{}न\appbrk{}मां\appbrk{}स\appbrk{}ग्र\appbrk{}न्थ्या\appbrk{}र्बु\appbrk{}दा\appbrk{}प\appbrk{}ची\appbrk{}ग\appbrk{}ल\appbrk{}ग\appbrk{}ण्ड\appbrk{}गृ\appbrk{}द्ध्र\appbrk{}सी\appbrk{}म\appbrk{}स\appbrk{}क\appbrk{}गु\appbrk{}ल्मो\appbrk{}द\appbrk{}र\appbrk{}भ\appbrk{}ग\appbrk{}क्ष\appbrk{}रा\appbrk{}र्स\appbrk{}म\appbrk{}न्वि\appbrk{}स्ली\appbrk{}प\appbrk{}द\appbrk{}च\appbrk{}र्म\appbrk{}की\appbrk{}ल\appbrk{}डि\appbrk{}ल\appbrk{}का\appbrk{}ल\appbrk{}क\appbrk{}सि\appbrk{}रा\appbrk{}च्छे\appbrk{}द\appbrk{}ना\appbrk{}डी\appbrk{}शो\appbrk{}णि\appbrk{}ता\appbrk{}ति\appbrk{}प्र\appbrk{}वृ\appbrk{}त्तिषु}\ N; \skt{दुष्ट\appbrk{}व्र\appbrk{}ण\appbrk{}मु\appbrk{}च्छ्रि\appbrk{}त\appbrk{}क\appbrk{}ठि\appbrk{}नमां सग्र\appbrk{}न्थ्य\appbrk{}र्ब्बु\appbrk{}दा\appbrk{}प\appbrk{}ची\appbrk{}ग\appbrk{}ल\appbrk{}ग\appbrk{}णु\appbrk{}गृ\appbrk{}द्ध्र\appbrk{}सी\appbrk{}म\appbrk{}स\appbrk{}क\appbrk{}गु\appbrk{}ल्मो\appbrk{}द\appbrk{}र\appbrk{}भ\appbrk{}ग\appbrk{}न्द\appbrk{}रा\appbrk{}र्शः\appbrk{}स\appbrk{}न्धि\appbrk{}श्ली\appbrk{}प\appbrk{}द\appbrk{}च\appbrk{}र्म्म\appbrk{}की\appbrk{}ल\appbrk{}ति\appbrk{}ल\appbrk{}का\appbrk{}ल\appbrk{}क\appbrk{}सि\appbrk{}रा\appbrk{}च्छेद नाडी\appbrk{}शो\appbrk{}णि\appbrk{}ता\appbrk{}ति\appbrk{}प्र\appbrk{}वृ\appbrk{}त्तिषु}\ H\textsuperscript{ac}; \skt{दुष्ट\appbrk{}व्र\appbrk{}ण\appbrk{}मु\appbrk{}च्छ्रि\appbrk{}त\appbrk{}क\appbrk{}ठि\appbrk{}नमां सग्र\appbrk{}न्थ्य\appbrk{}र्ब्बु\appbrk{}दा\appbrk{}प\appbrk{}ची\appbrk{}ग\appbrk{}ल\appbrk{}ग\appbrk{}णु\appbrk{}गृ\appbrk{}द्ध्र\appbrk{}सी\appbrk{}म\appbrk{}स\appbrk{}क\appbrk{}गु\appbrk{}ल्मो\appbrk{}द\appbrk{}र\appbrk{}भ\appbrk{}ग\appbrk{}न्द\appbrk{}रा\appbrk{}र्शः\appbrk{}स\appbrk{}न्धि\appbrk{}श्ली\appbrk{}प\appbrk{}द\appbrk{}च\appbrk{}र्म्म\appbrk{}की\appbrk{}ल\appbrk{}ति\appbrk{}ल\appbrk{}का\appbrk{}ल\appbrk{}क\appbrk{}सि\appbrk{}थं\appbrk{}रा\appbrk{}च्छेद नाडी\appbrk{}शो\appbrk{}णि\appbrk{}ता\appbrk{}ति\appbrk{}प्र\appbrk{}वृ\appbrk{}त्तिषु}\ H\textsuperscript{pc}}}\skt{ चाग्निकर्म }\edtext{\skt{कुर्यात्}}{\lemma{\skt{कुर्यात्}}\Afootnote{A K H; \skt{कुर्यात्~॥ …}\ L; \skt{कुयात्}\ N}}\skt{~।}\pend

\pstart \ledsidenote{\passref{1.12.11}}\edtext{\skt{तत्र}}{\lemma{\skt{तत्र}}\Afootnote{A K N H; \skt{… तत्र}\ L}}\skt{ }\edtext{\edtext{\skt{वल\-य\-बि\-न्दु\-रे\-रे\-खा\-प्र\-ति\-सा\-र\-णा\-ञ्चेति}}{\lemma{\skt{वल\appbrk{}य\appbrk{}बि\appbrk{}न्दु\appbrk{}रे\appbrk{}रे\appbrk{}खा\appbrk{}प्र\appbrk{}ति\appbrk{}सा\appbrk{}र\appbrk{}णाञ्}}\Afootnote{K; \skt{वल\appbrk{}य\appbrk{}बि\appbrk{}न्दु\appbrk{}वि\appbrk{}ले\appbrk{}खा\appbrk{}प्र\appbrk{}ति\appbrk{}सा\appbrk{}रणा}\ A L; \skt{वल\appbrk{}य\appbrk{}बिं\appbrk{}दु\appbrk{}रे\appbrk{}खा\appbrk{}प्र\appbrk{}ति\appbrk{}सा\appbrk{}रणं}\ N H}}}{\lemma{\skt{चेति}}\Afootnote{K N H; \skt{नीति}\ A L}}\skt{ }\edtext{\skt{दहणविशेषाः}}{\lemma{\skt{दहणविशेषाः}}\Afootnote{K; \skt{दहनविशेषाः}\ A H; \skt{दहनविशेषाः~॥ …}\ L; \skt{दहनविशेषा}\ N}}\skt{~॥}\pend

\pstart\versepar \ledsidenote{\passref{1.12.12}}\edtext{\skt{भवति चात्र}}{\lemma{\skt{भवति चात्र}}\Afootnote{A; \skt{… भवति चात्र}\ L; \emph{om.}\ K N; \skt{भवति चात्र}\ H}}\skt{~॥ }\pend
\pstart\versepar \skt{रोगस्य }\edtext{\edtext{\skt{संस्थानमवेक्ष्य}}{\lemma{\skt{संस्थानम्}}\Afootnote{A L K; \skt{संस्थान}\ N H}}}{\lemma{\skt{अवेक्ष्य}}\Afootnote{A L K; \skt{मवेक्ष्य}\ N H}}\skt{ }\edtext{\skt{धीमां}}{\lemma{\skt{धीमां}}\Afootnote{K N H; \skt{सम्यङ्}\ A; \skt{सम्य}\ L}}\skt{ }\edtext{\skt{नरस्य}}{\lemma{\skt{नरस्य}}\Afootnote{A K N H; \skt{क्नर्मस्य}\ L}}\skt{ मर्माणि }\edtext{\skt{बलाबलञ्च}}{\lemma{\skt{बलाबलञ्च}}\Afootnote{A K N H; \skt{वलावलं च}\ L}}\skt{~। }\pend
\pstart\versepar \skt{}\edtext{\skt{व्याधिन्तथर्तुञ्च}}{\lemma{\skt{तथर्तुञ्च}}\Afootnote{A K N H; \skt{तथार्त्तुं च}\ L}}\skt{ }\edtext{\skt{समीक्ष्य}}{\lemma{\skt{समीक्ष्य}}\Afootnote{A K N H; \emph{om.}\ L}}\skt{ }\edtext{\skt{सम्यक्ततो}}{\lemma{\skt{सम्यक्ततो}}\Afootnote{A K H; \skt{सम्यक्त}\ L N}}\skt{ }\edtext{\skt{व्यव\-स्ये\-द्भि\-ष\-ग\-ग्नि\-कर्म}}{\lemma{\skt{व्यव\appbrk{}स्ये\appbrk{}द्भि\appbrk{}ष\appbrk{}ग\appbrk{}ग्नि\appbrk{}कर्म}}\Afootnote{H; \skt{ऽव्य\appbrk{}व\appbrk{}स्ये\appbrk{}द्भि\appbrk{}ष\appbrk{}ग\appbrk{}ग्नि\appbrk{}कर्म}\ A; \skt{तोध्य\appbrk{}व\appbrk{}स्ये\appbrk{}द्भि\appbrk{}ष\appbrk{}ग्नि\appbrk{}कर्मः~॥ …}\ L; \skt{व्यव\appbrk{}स्ये\appbrk{}द्भि\appbrk{}ष\appbrk{}ग्नि\appbrk{}कर्म}\ K\textsuperscript{ac}; \skt{व्यव\appbrk{}स्ये\appbrk{}द्भि\appbrk{}ष\appbrk{}ग\appbrk{}ग्नि\appbrk{}कर्म}\ K\textsuperscript{pc}; \skt{तोध्य\appbrk{}व\appbrk{}स्ये\appbrk{}त्भि\appbrk{}ष\appbrk{}ग\appbrk{}ग्नि\appbrk{}कर्म}\ N}}\skt{~॥}\pend

\pstart \ledsidenote{\passref{1.12.13}}\edtext{\skt{तत्र}}{\lemma{\skt{तत्र}}\Afootnote{A K N H; \skt{… तत्र}\ L}}\skt{ }\edtext{\skt{सम्यग्दग्धेषु}}{\lemma{\skt{सम्यग्दग्धेषु}}\Afootnote{K H; \skt{सम्यग्दग्धे}\ A L; \skt{संम्यग्दग्धे}\ N}}\skt{ }\edtext{\skt{मधु\-स\-र्पि\-र\-भ्यङ्गः}}{\lemma{\skt{मधु\appbrk{}स\appbrk{}र्पि\appbrk{}र\appbrk{}भ्यङ्गः}}\Afootnote{K N; \skt{मधु\appbrk{}स\appbrk{}र्पि\appbrk{}र्भ्या\appbrk{}म\appbrk{}भ्यङ्गः}\ A; \skt{मध्यु\appbrk{}स\appbrk{}र्प्पि\appbrk{}भ्यां\appbrk{}म\appbrk{}भ्यंगः~॥ …}\ L; \skt{मधु\appbrk{}स\appbrk{}र्प्पि\appbrk{}र\appbrk{}भ्यङ्गोः}\ H}}\skt{~॥}\pend

\pstart \ledsidenote{\passref{1.12.14}}\edtext{\skt{अथेमानि}}{\lemma{\skt{अथेमानि}}\Afootnote{K N H; \skt{अथेमानग्निना}\ A; \skt{… अथेमानग्निना}\ L}}\skt{ }\edtext{\edtext{\skt{परि\-ह\-रे\-त्पि\-त्त\-प्र\-कृ\-ति\-म\-न्तः\-शो\-णितं}}{\lemma{\skt{परि\appbrk{}ह\appbrk{}रे\appbrk{}त्पि\appbrk{}त्त\appbrk{}प्र\appbrk{}कृ\appbrk{}तिम्}}\Afootnote{A L K H; \skt{परि\appbrk{}ह\appbrk{}रे\appbrk{}त्पि\appbrk{}त्त\appbrk{}प्र\appbrk{}कृति}\ N}}}{\lemma{\skt{अन्तःशोणितं}}\Afootnote{A L K H; \skt{मन्तः सोणितं}\ N}}\skt{~। }\edtext{\edtext{\skt{भिन्न\-को\-ष्ठ\-म\-नु\-द्धृ\-त\-शल्यं}}{\lemma{\skt{भिन्नकोष्ठम्}}\Afootnote{L K N H; \skt{भिन्नकोष्ठ}\ A}}}{\lemma{\skt{अनुद्धृतशल्यं}}\Afootnote{L K H; \skt{मनुद्धृतशल्यं}\ A; \skt{अनुद्धृत}\ N}}\skt{ }\edtext{\skt{बाल\-वृ\-द्ध\-भी\-रु\-दु\-र्ब\-ल\-म\-ने\-क\-व्या\-धि\-पी\-डि\-त\-म\-स्वे\-द्यांश्च}}{\lemma{\skt{बाल\appbrk{}वृ\appbrk{}द्ध\appbrk{}भी\appbrk{}रु\appbrk{}दु\appbrk{}र्ब\appbrk{}ल\appbrk{}म\appbrk{}ने\appbrk{}क\appbrk{}व्या\appbrk{}धि\appbrk{}पी\appbrk{}डि\appbrk{}त\appbrk{}म\appbrk{}स्वे\appbrk{}द्यांश्च}}\Afootnote{K H; \skt{दुर्बलं बालं वृद्धं भीरु\appbrk{}म\appbrk{}ने\appbrk{}क\appbrk{}व्र\appbrk{}ण\appbrk{}पी\appbrk{}डि\appbrk{}त\appbrk{}म\appbrk{}स्वे\appbrk{}द्यां\appbrk{}श्चेति}\ A; \skt{बालं वृद्धं भीरु\appbrk{}म\appbrk{}ने\appbrk{}क\appbrk{}व्र\appbrk{}ण\appbrk{}पी\appbrk{}डि\appbrk{}त\appbrk{}म\appbrk{}स्वे\appbrk{}द्यां\appbrk{}श्चेति~॥ …}\ L; \skt{सल्यं बाल\appbrk{}वृ\appbrk{}द्ध\appbrk{}भी\appbrk{}रु\appbrk{}दु\appbrk{}र्ब\appbrk{}ल\appbrk{}म\appbrk{}ने\appbrk{}क\appbrk{}व्या\appbrk{}धि\appbrk{}पी\appbrk{}डि\appbrk{}त\appbrk{}म\appbrk{}स्वे\appbrk{}द्यांस्च}\ N}}\skt{~॥}\pend

\pstart \ledsidenote{\passref{1.12.15}}\edtext{\skt{अत}}{\lemma{\skt{अत}}\Afootnote{A K N H; \skt{… अत}\ L}}\skt{ }\edtext{\edtext{\skt{ऊर्द्धमितरथा}}{\lemma{\skt{ऊर्द्धम्}}\Afootnote{K N; \skt{ऊर्ध्व}\ A; \skt{ऊर्ध्वम्}\ L H}}\skt{ दग्धं}}{\lemma{\skt{इतरथा दग्धं}}\Afootnote{L K N H; \skt{मितरथादग्ध}\ A}}\skt{ }\edtext{\skt{वक्ष्यामः}}{\lemma{\skt{वक्ष्यामः}}\Afootnote{L K N H; \skt{लक्षणं वक्ष्यामः}\ A}}\skt{ तत्र स्निग्धं }\edtext{\edtext{\edtext{\edtext{\skt{रूक्ष\-ञ्चा\-श्रि\-त्य\-द्र\-व्य\-म\-ग्नि\-र्द\-हति}}{\lemma{\skt{रूक्षञ्}}\Afootnote{A K N H; \skt{रुक्षं}\ L}}}{\lemma{\skt{चाश्रित्यद्रव्यम्}}\Afootnote{K N\textsuperscript{pc} H; \skt{वाऽऽ(चा)श्रित्य द्रव्य}\ A; \skt{वाश्रित्य द्रव्यम्}\ L; \skt{चाक्षञ्चाश्रित्य द्रव्यम्}\ N\textsuperscript{ac}}}}{\lemma{\skt{अग्निर्}}\Afootnote{K N H; \skt{मग्निर्}\ A; \skt{अग्नि}\ L}}}{\lemma{\skt{दहति}}\Afootnote{A K N H; \skt{दहंति}\ L}}\skt{~। }\edtext{\skt{अतिसन्तप्तो हि}}{\lemma{\skt{अतिसन्तप्तो हि}}\Afootnote{K N H; \skt{अग्निसन्तप्तो हि}\ A L}}\skt{ स्नेहः }\edtext{\skt{सूक्ष्म\-मा\-र्गा\-नु\-सा\-रि\-त्वा\-त्त्व\-गा\-दी\-न\-नु\-प्र\-वि\-श्याशु}}{\lemma{\skt{सूक्ष्म\appbrk{}मा\appbrk{}र्गा\appbrk{}नु\appbrk{}सा\appbrk{}रि\appbrk{}त्वा\appbrk{}त्त्व\appbrk{}गा\appbrk{}दी\appbrk{}न\appbrk{}नु\appbrk{}प्र\appbrk{}वि\appbrk{}श्याशु}}\Afootnote{K N; \skt{सूक्ष्म\appbrk{}सि\appbrk{}रा\appbrk{}नु\appbrk{}सा\appbrk{}रि\appbrk{}त्वा\appbrk{}त्त्व\appbrk{}गा\appbrk{}दी\appbrk{}न\appbrk{}नु\appbrk{}प्र\appbrk{}वि\appbrk{}श्याशु}\ A; \skt{सूक्ष्म\appbrk{}सि\appbrk{}रा\appbrk{}नु\appbrk{}सा\appbrk{}रि\appbrk{}त्वा\appbrk{}त्त्व\appbrk{}गा\appbrk{}दी\appbrk{}न\appbrk{}नु\appbrk{}प्र\appbrk{}विश्य आशु}\ L; \skt{सूक्ष्म\appbrk{}मा\appbrk{}र्ग्गा\appbrk{}नु\appbrk{}सा\appbrk{}रि\appbrk{}त्वा\appbrk{}त्त्व\appbrk{}गा\appbrk{}दी\appbrk{}न\appbrk{}नु\appbrk{}प्र\appbrk{}वि\appbrk{}श्याशु}\ H}}\skt{ }\edtext{\skt{दहति}}{\lemma{\skt{दहति}}\Afootnote{A L K H; \skt{ददति}\ N}}\skt{~। }\edtext{\skt{तस्मा\-त्स्ने\-ह\-द\-ग्धेऽ\-धिका}}{\lemma{\skt{ऽधिका}}\Afootnote{A; \skt{धिका}\ L K N H}}\skt{ }\edtext{\skt{रुजा}}{\lemma{\skt{रुजा}}\Afootnote{K N H; \skt{रुजो}\ A L}}\skt{ }\edtext{\skt{भवति}}{\lemma{\skt{भवति}}\Afootnote{K N H; \skt{भवन्ति}\ A; \skt{भवतंति~॥ …}\ L}}\skt{~।}\pend

\pstart \ledsidenote{\passref{1.12.16}}\edtext{\skt{तत्र}}{\lemma{\skt{तत्र}}\Afootnote{A K N H; \skt{… तत्र}\ L}}\skt{ प्लुष्टं }\edtext{\skt{दुर्दग्धं}}{\lemma{\skt{दुर्दग्धं}}\Afootnote{A L K H; \skt{दुदग्धं}\ N}}\skt{ }\edtext{\edtext{\edtext{\skt{सम्य\-ग्द\-ग्ध\-म\-ति\-द\-ग्ध\-मिति}}{\lemma{\skt{सम्यग्दग्धम्}}\Afootnote{K N H; \skt{सम्यग्दग्ध}\ A; \skt{सम्यग्दग}\ L}}}{\lemma{\skt{अतिदग्धम्}}\Afootnote{K H; \skt{मतिदग्धं}\ A; \skt{द्धमतिदग्धं}\ L; \skt{इति दग्धम्}\ N}}}{\lemma{\skt{इति}}\Afootnote{K N H; \skt{चेति}\ A L}}\skt{ }\edtext{\edtext{\edtext{\skt{चतु\-र्वि\-ध\-म\-ग्नि\-द\-ग्ध\-म्भ\-वति}}{\lemma{\skt{चतुर्विधम्}}\Afootnote{L K H; \skt{चतुर्विध}\ A; \skt{चतुविदम्}\ N}}}{\lemma{\skt{अग्निदग्धम्}}\Afootnote{L K N H; \skt{मग्निदग्धम्}\ A}}}{\lemma{\skt{भवति}}\Afootnote{K N H; \emph{om.}\ A; \skt{…}\ L}}\skt{~। तत्र }\edtext{\edtext{\edtext{\skt{यद्वि\-व\-ण्ण\-मु\-ष्य\-तेऽ\-ति\-मात्रं}}{\lemma{\skt{यद्विवण्णम्}}\Afootnote{K; \skt{यद्विवर्णं}\ A L N H}}}{\lemma{\skt{उष्यते}}\Afootnote{K N H; \skt{प्लुष्यते}\ A L}}}{\lemma{\skt{ऽतिमात्रं}}\Afootnote{A L H; \skt{तिमात्रं}\ K N}}\skt{ }\edtext{\skt{तत्प्लुष्टं}}{\lemma{\skt{तत्प्लुष्टं}}\Afootnote{A K H; \skt{तदप्लुष्टं …}\ L; \skt{प्लुष्टं}\ N}}\skt{~। }\edtext{\skt{यत्रो\-त्प\-त्ति\-स्फो\-टना तीव्र\-दा\-ह\-दो\-ष\-वे\-दना}}{\lemma{\skt{यत्रो\appbrk{}त्प\appbrk{}त्ति\appbrk{}स्फो\appbrk{}टना तीव्र\appbrk{}दा\appbrk{}ह\appbrk{}दो\appbrk{}ष\appbrk{}वे\appbrk{}दना}}\Afootnote{K\textsuperscript{pc}; \skt{यत्रोत्तिष्ठन्ति स्फोटा\appbrk{}स्ती\appbrk{}व्रा\appbrk{}श्चो\appbrk{}ष\appbrk{}दा\appbrk{}ह\appbrk{}रा\appbrk{}ग\appbrk{}पा\appbrk{}क\appbrk{}वे\appbrk{}द\appbrk{}नाश्}\ A L; \skt{यत्रोत्पत्तिस्फोट तीव्र\appbrk{}दा\appbrk{}ह\appbrk{}दो\appbrk{}ष\appbrk{}वे\appbrk{}दना}\ K\textsuperscript{ac}; \skt{यत्रोत्पत्ति स्फोटानां स्त्रीव्र\appbrk{}दा\appbrk{}ह\appbrk{}चो\appbrk{}ष\appbrk{}वे\appbrk{}दना}\ N; \skt{यत्रोत्पत्ति स्फोट\appbrk{}ना\appbrk{}स्ती\appbrk{}व्र\appbrk{}दा\appbrk{}ह\appbrk{}दो\appbrk{}ष\appbrk{}वे\appbrk{}दना}\ H}}\skt{ }\edtext{\skt{चिरा\-च्चो\-प\-शा\-म्यति}}{\lemma{\skt{चिरा\appbrk{}च्चो\appbrk{}प\appbrk{}शा\appbrk{}म्यति}}\Afootnote{K H; \skt{चिरा\appbrk{}च्चो\appbrk{}प\appbrk{}शा\appbrk{}म्यन्ति}\ A L; \skt{चिरा चोपशाम्यति}\ N}}\skt{ }\edtext{\skt{तद्दुर्दग्धं}}{\lemma{\skt{तद्दुर्दग्धं}}\Afootnote{A K; \skt{तद्दुर्द्दग्धं~॥ …}\ L; \skt{तद्दुदग्धं}\ N; \skt{तद्दुरुद्दग्धः}\ H}}\skt{~। }\edtext{\edtext{\skt{सम्य\-ग्द\-ग्ध\-म\-न\-व\-गाढं}}{\lemma{\skt{सम्यग्दग्धम्}}\Afootnote{A L K N; \skt{सम्यक्दग्धम्}\ H}}}{\lemma{\skt{अनवगाढं}}\Afootnote{A L K H; \skt{अवगाढं}\ N}}\skt{ }\edtext{\skt{पक्वतालवर्ण्णं}}{\lemma{\skt{पक्वतालवर्ण्णं}}\Afootnote{K H; \skt{तालवर्णं}\ A L; \skt{पक्वतालवर्ण्ण}\ N}}\skt{ }\edtext{\skt{सुसंस्थितं}}{\lemma{\skt{सुसंस्थितं}}\Afootnote{A L K N; \skt{सुसंस्थि}\ H\textsuperscript{ac}; \skt{सुसंस्थित}\ H\textsuperscript{pc}}}\skt{ }\edtext{\skt{पूर्व\-ल\-क्ष\-ण\-सं\-युक्तं च}}{\lemma{\skt{पूर्व\appbrk{}ल\appbrk{}क्ष\appbrk{}ण\appbrk{}सं\appbrk{}युक्तं च}}\Afootnote{K N H; \skt{पूर्व\appbrk{}ल\appbrk{}क्ष\appbrk{}ण\appbrk{}युक्तं च}\ A; \skt{पूर्व\appbrk{}ल\appbrk{}क्ष\appbrk{}ण\appbrk{}युक्तं च~॥ …}\ L}}\skt{~। }\edtext{\skt{अतिदग्धे तु}}{\lemma{\skt{अतिदग्धे तु}}\Afootnote{K N H; \skt{अतिदग्धे}\ A L}}\skt{ }\edtext{\skt{मांसावलंबनं}}{\lemma{\skt{मांसावलंबनं}}\Afootnote{A L K H; \skt{मांसावलवनं}\ N}}\skt{ गात्रविश्लेषः }\edtext{\skt{सन्धिवैगुण्यं}}{\lemma{\skt{सन्धिवैगुण्यं}}\Afootnote{K H; \emph{om.}\ A L; \skt{सन्धिर्वैगुण्यं}\ N}}\skt{ }\edtext{\skt{सिरा\-स्ना\-यु\-स\-न्ध्य\-स्थि\-व्या\-पा\-दनं}}{\lemma{\skt{सिरा\appbrk{}स्ना\appbrk{}यु\appbrk{}स\appbrk{}न्ध्य\appbrk{}स्थि\appbrk{}व्या\appbrk{}पा\appbrk{}दनं}}\Afootnote{K H; \skt{सिरा\appbrk{}स्ना\appbrk{}यु\appbrk{}स\appbrk{}न्ध्य\appbrk{}स्थि\appbrk{}व्या\appbrk{}पा\appbrk{}दन}\ A; \skt{सिरा\appbrk{}स्ना\appbrk{}यु\appbrk{}सं\appbrk{}ध्य\appbrk{}स्थि\appbrk{}व्या\appbrk{}पा\appbrk{}दा\appbrk{}नम्}\ L; \skt{मर्म\appbrk{}शि\appbrk{}रा\appbrk{}स्ना\appbrk{}यु\appbrk{}स\appbrk{}न्ध्य\appbrk{}स्थि\appbrk{}व्या\appbrk{}पा\appbrk{}दनं}\ N}}\skt{ }\edtext{\edtext{\skt{ज्वर\-दा\-ह\-पि\-पा\-सा\-मू\-र्च्छा\-श्चो\-प\-द्रवा}}{\lemma{\skt{ज्वर\appbrk{}दा\appbrk{}ह\appbrk{}पि\appbrk{}पा\appbrk{}सा\appbrk{}मू\appbrk{}र्च्छाश्}}\Afootnote{K; \skt{मतिमात्रं ज्वर\appbrk{}दा\appbrk{}ह\appbrk{}पि\appbrk{}पा\appbrk{}सा\appbrk{}मूर्च्छा}\ A; \skt{अतिमात्रं ज्वरदाहपिपा सामूर्छाश्}\ L; \skt{ज्वर\appbrk{}दा\appbrk{}ह\appbrk{}पि\appbrk{}पा\appbrk{}सा\appbrk{}मूर्च्छा}\ N H}}}{\lemma{\skt{चोपद्रवा}}\Afootnote{L K N H; \skt{श्चोपद्रवा}\ A}}\skt{ भवन्ति~। }\edtext{\edtext{\edtext{\skt{सक्रि\-मि\-श्चे\-त्व्र\-ण\-श्चास्य}}{\lemma{\skt{सक्रिमिश्}}\Afootnote{K\textsuperscript{pc} H; \emph{om.}\ A L K\textsuperscript{ac}; \skt{सत्क्रिमिश्}\ N}}}{\lemma{\skt{चेत्व्रणश्}}\Afootnote{K\textsuperscript{pc} N H; \skt{व्रणश्}\ A L K\textsuperscript{ac}}}}{\lemma{\skt{चास्य}}\Afootnote{A K N H; \skt{च अस्य}\ L}}\skt{ }\edtext{\edtext{\skt{चिरे\-णो\-प\-रो\-ह\-त्यु\-प\-रू\-ढश्च}}{\lemma{\skt{चिरे\appbrk{}णो\appbrk{}प\appbrk{}रो\appbrk{}हत्य्}}\Afootnote{K; \skt{चिरेण रोहति}\ A L; \skt{चिरे\appbrk{}णो\appbrk{}प\appbrk{}ह\appbrk{}रो\appbrk{}हत्य्}\ N; \skt{चिरे\appbrk{}णो\appbrk{}प\appbrk{}ह\appbrk{}रोत्य्}\ H}}}{\lemma{\skt{उपरूढश्च}}\Afootnote{K N H; \skt{रूढश्च}\ A; \skt{रुढश्च}\ L}}\skt{ विवर्ण्णो भवति~। }\edtext{\edtext{\edtext{\skt{तदे\-त\-च्च\-तु\-र्वि\-ध\-म\-ग्नि\-द\-ग्ध\-ल\-क्ष\-ण\-मा\-नु\-पू\-र्व्योक्तं}}{\lemma{\skt{तदे\appbrk{}त\appbrk{}च्च\appbrk{}तु\appbrk{}र्वि\appbrk{}धम्}}\Afootnote{A K H; \emph{om.}\ L; \skt{तद्दे\appbrk{}त\appbrk{}च्च\appbrk{}तु\appbrk{}र्वि\appbrk{}दम्}\ N}}}{\lemma{\skt{अग्नि\appbrk{}द\appbrk{}ग्ध\appbrk{}ल\appbrk{}क्ष\appbrk{}णम्}}\Afootnote{A K N H; \emph{om.}\ L}}}{\lemma{\skt{आनुपूर्व्योक्तं}}\Afootnote{K; \skt{आ}\ A; \emph{om.}\ L; \skt{आनुपूक्तं}\ N; \skt{आनुपूर्व्वोक्तं}\ H}}\skt{ }\edtext{\skt{पूर्व\-क\-र्म\-प्र\-सा\-धकं}}{\lemma{\skt{पूर्व\appbrk{}क\appbrk{}र्म\appbrk{}प्र\appbrk{}सा\appbrk{}धकं}}\Afootnote{K N H\textsuperscript{pc}; \skt{त्मक\appbrk{}र्म\appbrk{}प्र\appbrk{}सा\appbrk{}धकं}\ A; \emph{om.}\ L; \skt{पूर्व्वकर्म्म प्र}\ H\textsuperscript{ac}}}\skt{ }\edtext{\skt{भवति}}{\lemma{\skt{भवति}}\Afootnote{A K N; \emph{om.}\ L; \skt{धानं भवति}\ H\textsuperscript{ac}; \skt{प्रधानं भवति}\ H\textsuperscript{pc}}}\skt{~॥}\pend

\pstart\versepar \ledsidenote{\passref{1.12.17}}\edtext{\skt{भवति}}{\lemma{\skt{भवति}}\Afootnote{\skt{भवन्ति}\ A; \skt{… भवति}\ L; \emph{om.}\ K N H}}\skt{ }\edtext{\skt{चात्र}}{\lemma{\skt{चात्र}}\Afootnote{A L; \emph{om.}\ K N; \skt{भवति चात्र}\ H}}\skt{~॥ }\pend
\pstart\versepar \skt{अग्निना कोपितं }\edtext{\skt{पित्तं}}{\lemma{\skt{पित्तं}}\Afootnote{K N H; \skt{रक्तं}\ A L}}\skt{ भृशञ्जन्तोः }\edtext{\skt{प्रधावति}}{\lemma{\skt{प्रधावति}}\Afootnote{K N H; \skt{प्रकुप्यति}\ A L}}\skt{~। }\pend
\pstart\versepar \skt{}\edtext{\edtext{\skt{ततस्तेनैव}}{\lemma{\skt{ततस्}}\Afootnote{A L N; \skt{तत}\ K H}}}{\lemma{\skt{तेनैव}}\Afootnote{A L N; \skt{स्तेनैव}\ K H}}\skt{ वेगेन }\edtext{\edtext{\skt{रक्त\-ञ्चा\-प्यु\-प\-दी\-र्यते}}{\lemma{\skt{रक्तञ्}}\Afootnote{K N H; \emph{om.}\ A L}}}{\lemma{\skt{चाप्युपदीर्यते}}\Afootnote{\skt{पित्त\appbrk{}म\appbrk{}स्या\appbrk{}भ्यु\appbrk{}दी\appbrk{}र्यते}\ A; \skt{पित्त\appbrk{}म\appbrk{}स्या\appbrk{}भ्यु\appbrk{}दी\appbrk{}र्यते~॥ …}\ L; \skt{चाप्युपदीर्यते}\ K H; \skt{चाप्यपदीर्यते}\ N}}\skt{~।}\pend

\pstart\versepar \ledsidenote{\passref{1.12.18}}\edtext{\edtext{\edtext{\skt{तुल्यवीर्येऽप्युभे}}{\lemma{\skt{तुल्यवीर्ये}}\Afootnote{A N H; \skt{… तुल्यवीर्ये}\ L; \skt{तुल्यविर्ये}\ K}}}{\lemma{\skt{ऽप्य्}}\Afootnote{\emph{om.}\ A L; \skt{प्य्}\ K N; \skt{ऽप्यु}\ H}}\skt{ }\edtext{\skt{ह्येते}}{\lemma{\skt{एते}}\Afootnote{L K N; \skt{ह्येते}\ A H}}}{\lemma{\skt{उभे ह्य्}}\Afootnote{L K\textsuperscript{pc} N; \skt{उभे}\ A; \skt{उ ह्य्}\ K\textsuperscript{ac}; \skt{भे}\ H}}\skt{ रसतो द्रव्यतस्तथा~। }\pend
\pstart\versepar \skt{तेनास्य }\edtext{\edtext{\skt{वेदनास्तीव्रा}}{\lemma{\skt{वेदनास्}}\Afootnote{K N; \skt{वेदना}\ A L H}}}{\lemma{\skt{तीव्रा}}\Afootnote{L K H; \skt{स्तीव्राः}\ A; \skt{त्रीव्रा}\ N}}\skt{ प्रकृत्या च }\edtext{\skt{विदह्यते}}{\lemma{\skt{विदह्यते}}\Afootnote{A K N H; \skt{विदह्यते~॥ …}\ L}}\skt{~।}\pend

\pstart\versepar \ledsidenote{\passref{1.12.19}}\skt{स्फोटाः शीघ्रं प्रजायन्ते }\edtext{\edtext{\skt{ज्वरस्तृष्णा}}{\lemma{\skt{ज्वरस्}}\Afootnote{A L K H; \skt{ज्वर}\ N}}\skt{ च}}{\lemma{\skt{तृष्णा च}}\Afootnote{A K N H; \skt{तृषा च}\ L}}\skt{ }\edtext{\skt{बाधते}}{\lemma{\skt{बाधते}}\Afootnote{A K N H; \skt{बाधते~॥ …}\ L}}\skt{~॥ }\pend
\pstart\versepar \skt{}\edtext{\skt{दग्ध\-स्यो\-प\-श\-मा\-र्थाय}}{\lemma{\skt{दग्ध\appbrk{}स्यो\appbrk{}प\appbrk{}श\appbrk{}मा\appbrk{}र्थाय}}\Afootnote{A L K N; \skt{दग्ध\appbrk{}स्यो\appbrk{}प\appbrk{}श\appbrk{}मा\appbrk{}र्थाय …}\ H}}\skt{ चिकित्सा }\edtext{\skt{सम्प्रवक्ष्यते}}{\lemma{\skt{सम्प्रवक्ष्यते}}\Afootnote{A K N H; \skt{संप्रक्ष्यते~॥ …}\ L}}\skt{~।}\pend

\pstart\versepar \ledsidenote{\passref{1.12.20}}\edtext{\skt{प्लुष्ट\-स्या\-ग्नि\-प्र\-त\-पनं}}{\lemma{\skt{प्लुष्ट\appbrk{}स्या\appbrk{}ग्नि\appbrk{}प्र\appbrk{}त\appbrk{}पनं}}\Afootnote{A L K N; \skt{प्लुष्ट\appbrk{}स्या\appbrk{}ग्नि\appbrk{}प्र\appbrk{}त\appbrk{}पनं…}\ H}}\skt{ }\edtext{\edtext{\edtext{\skt{कार्य\-मु\-ष्ण\-न्त\-थौ\-षधं}}{\lemma{\skt{कार्यम्}}\Afootnote{L N; \skt{कार्य}\ A K H}}}{\lemma{\skt{उष्णन्}}\Afootnote{N; \skt{मुष्णं}\ A K H; \skt{उष्ण}\ L}}}{\lemma{\skt{तथौषधं}}\Afootnote{A K N H; \skt{तथौषधं~॥ …}\ L}}\skt{~। }\pend
\pstart\versepar \skt{शरीरे स्विन्नभूयिष्ठे }\edtext{\skt{स्विनं}}{\lemma{\skt{स्विनं}}\Afootnote{K; \skt{स्विन्नं}\ A L H; \skt{स्विन्न}\ N}}\skt{ भवति }\edtext{\skt{शोणितं}}{\lemma{\skt{शोणितं}}\Afootnote{A K N H; \skt{शोणितं~। …}\ L}}\skt{~।}\pend

\pstart\versepar \ledsidenote{\passref{1.12.21}}\edtext{\skt{प्रकृत्या }\edtext{\skt{ह्युदकं}}{\lemma{\skt{उदकं}}\Afootnote{A K N H; \skt{ह्युदकं}\ L}}}{\lemma{\skt{प्रकृत्या ह्य्}}\Afootnote{A K N H; \skt{प्रकृत्या}\ L}}\skt{ शीतं }\edtext{\edtext{\skt{स्कन्दयत्यथ}}{\lemma{\skt{स्कन्दयत्य्}}\Afootnote{A K H; \skt{स्कंदय}\ L; \skt{स्कन्दयस्त}\ N}}}{\lemma{\skt{अथ}}\Afootnote{K H; \skt{अति}\ A; \skt{त्य}\ L; \skt{थ}\ N}}\skt{ }\edtext{\skt{शोणितं}}{\lemma{\skt{शोणितं}}\Afootnote{A K N H\textsuperscript{ac}; \skt{तिशोणितं}\ L; \skt{शोणिथं एतं}\ H\textsuperscript{pc}}}\skt{~। }\pend
\pstart\versepar \skt{}\edtext{\skt{तस्मात्सुखयति ह्युष्ण}}{\lemma{\skt{तस्मात्सुखयति ह्य्}}\Afootnote{A K N H; \skt{तस्मात्सुखयति}\ L}}\skt{ न तु शीतं }\edtext{\skt{कथञ्चनः}}{\lemma{\skt{कथञ्चनः}}\Afootnote{K N; \skt{कथञ्चन}\ A H; \skt{कथंचन~॥१॥ …}\ L}}\skt{~।}\edtext{}{\lemma{\skt{उष्ण न तु}}\Afootnote{K\textsuperscript{pc}; \skt{उष्णं ननु}\ A; \skt{ह्युष्णं म तु}\ L; \skt{उष्ण न}\ K\textsuperscript{ac}; \skt{उष्णं न तु}\ N H}}\pend

\pstart\versepar \ledsidenote{\passref{1.12.22}}\edtext{\edtext{\skt{शीतामुष्णाञ्च}}{\lemma{\skt{शीताम्}}\Afootnote{N; \skt{शीता}\ A K H; \skt{… शीताम्}\ L}}\skt{ च}}{\lemma{\skt{उष्णाञ्च च}}\Afootnote{\skt{मुष्णां च}\ A H; \skt{उष्णां च}\ L N; \skt{मुष्णाञ्च च}\ K}}\skt{ दुर्दग्धे }\edtext{\skt{क्रियां}}{\lemma{\skt{क्रियां}}\Afootnote{A L K H; \skt{क्रि}\ N}}\skt{ }\edtext{\skt{कुर्यात्ततः}}{\lemma{\skt{कुर्यात्ततः}}\Afootnote{K H; \skt{कुर्या}\ A; \skt{कुर्यात्भिष}\ L; \emph{om.}\ N}}\skt{ }\edtext{\skt{पुनः}}{\lemma{\skt{पुनः}}\Afootnote{K H; \emph{om.}\ A N; \skt{क्पुनः}\ L}}\skt{~। }\pend
\pstart\versepar \skt{}\edtext{\skt{घृता\-ले\-प\-न\-से\-कांस्तु}}{\lemma{\skt{घृता\appbrk{}ले\appbrk{}प\appbrk{}न\appbrk{}से\appbrk{}कांस्तु}}\Afootnote{\skt{द्भिषक्पुनः~। घृता\appbrk{}ले\appbrk{}प\appbrk{}न\appbrk{}से\appbrk{}कांस्तु}\ A; \skt{घृता\appbrk{}ले\appbrk{}प\appbrk{}न\appbrk{}से\appbrk{}कांस्तु शीतामेवास्य कारयेत्~॥ घृता\appbrk{}ले\appbrk{}प\appbrk{}न\appbrk{}से\appbrk{}कांस्तु}\ L; \skt{घृता\appbrk{}ले\appbrk{}प\appbrk{}न\appbrk{}से\appbrk{}कांस्तु}\ K H; \skt{या कुयां ततः पुनः~॥ घृता\appbrk{}ले\appbrk{}प\appbrk{}न\appbrk{}से\appbrk{}कांस्तु}\ N}}\skt{ शीतानेवास्य }\edtext{\skt{कारयेत्}}{\lemma{\skt{कारयेत्}}\Afootnote{A K N H; \skt{कारयेत्~॥ …}\ L}}\skt{~॥}\pend

\pstart\versepar \ledsidenote{\passref{1.12.23}}\edtext{\skt{सम्यग्दग्धे}}{\lemma{\skt{सम्यग्दग्धे}}\Afootnote{A K H; \skt{… सम्यक्दग्धे}\ L; \skt{संम्यग्दग्धे}\ N}}\skt{ }\edtext{\skt{तु काक्षी\-री\-प्ल\-क्ष\-च\-न्द\-न\-गै\-रिकैः}}{\lemma{\skt{तु काक्षी\appbrk{}री\appbrk{}प्ल\appbrk{}क्ष\appbrk{}च\appbrk{}न्द\appbrk{}न\appbrk{}गै\appbrk{}रिकैः}}\Afootnote{K H; \skt{तुगा\appbrk{}क्षी\appbrk{}री\appbrk{}प्ल\appbrk{}क्ष\appbrk{}च\appbrk{}न्द\appbrk{}न\appbrk{}गै\appbrk{}रिकैः}\ A L N}}\skt{~। }\pend
\pstart\versepar \skt{सामृतैः सर्पिषा }\edtext{\edtext{\skt{युक्तैरालेपं}}{\lemma{\skt{युक्तैर्}}\Afootnote{N; \skt{स्निग्धै}\ A; \skt{स्निग्धैर्}\ L; \skt{युक्तै}\ K H}}}{\lemma{\skt{आलेपं}}\Afootnote{L N; \skt{रालेपं}\ A K H}}\skt{ कारयेद्भिषक्~॥}\pend

\pstart\versepar \ledsidenote{\passref{1.12.24}}\edtext{\skt{ग्राम्यानूपैः}}{\lemma{\skt{ग्राम्यानूपैः}}\Afootnote{K N H; \skt{ग्राम्यानूपौ}\ A L}}\skt{ }\edtext{\skt{सजलजैः}}{\lemma{\skt{सजलजैः}}\Afootnote{K N H; \skt{दकैश्चैनं}\ A; \skt{दकैश्चैनं}\ L}}\skt{ }\edtext{\edtext{\skt{पिष्टैर्मांसैश्च}}{\lemma{\skt{पिष्टैर्}}\Afootnote{L K N H; \skt{पिष्टै}\ A}}}{\lemma{\skt{मांसैश्च}}\Afootnote{K N H; \skt{र्मांसैः प्र}\ A; \skt{मासैः प्र}\ L}}\skt{ लेपयेत्~। }\pend
\pstart\versepar \skt{पित्त\-वि\-द्र\-धि\-व\-च्चैनं }\edtext{\edtext{\skt{प्रशा\-न्त्यो\-ष्मा\-ण\-मा\-च\-रेत्}}{\lemma{\skt{प्रशान्त्योष्माणम्}}\Afootnote{\skt{सन्ततोष्माण}\ A; \skt{संततोष्माणम्}\ L; \skt{प्रशान्त्योष्माण}\ K; \skt{प्रशान्तोष्माणम्}\ N; \skt{प्रशांत्योष्मान}\ H}}}{\lemma{\skt{आचरेत्}}\Afootnote{N; \skt{माचरेत्}\ A K H; \skt{आचरेत्~॥ …}\ L}}\skt{~।}\edtext{}{\lemma{\skt{लेपयेत्~। पित्त\appbrk{}वि\appbrk{}द्र\appbrk{}धि\appbrk{}व\appbrk{}च्चैनं}}\Afootnote{L\textsuperscript{pc}; \skt{लेप\appbrk{}ये\appbrk{}त्पि\appbrk{}त्त\appbrk{}वि\appbrk{}द्र\appbrk{}धि\appbrk{}व\appbrk{}च्चैनं}\ A H K N; \skt{लेपलयेत्~॥ पित्त\appbrk{}वि\appbrk{}द्र\appbrk{}धि\appbrk{}व\appbrk{}च्चैनं}\ L\textsuperscript{ac}}}\pend

\pstart\versepar \ledsidenote{\passref{1.12.25}}\edtext{\skt{अतिदग्धे तु}}{\lemma{\skt{अतिदग्धे तु}}\Afootnote{K N H; \skt{अतिदग्धे वि}\ A; \skt{… अतिदग्धे वि}\ L}}\skt{ शीर्ण्णानि }\edtext{\edtext{\skt{मांसान्युद्धृत्य}}{\lemma{\skt{मांसान्य्}}\Afootnote{K N; \skt{मांसा}\ A H; \skt{मांसानु}\ L}}}{\lemma{\skt{उद्धृत्य}}\Afootnote{K N; \skt{न्युद्धृत्य}\ A H; \skt{धृत्य}\ L}}\skt{ शीतलां }\pend
\pstart\versepar \skt{}\edtext{\skt{क्रिया\-ङ्कु\-र्या\-च्चू\-र्ण्ण\-काले}}{\lemma{\skt{कुर्याच्चूर्ण्णकाले}}\Afootnote{K H; \skt{कुर्याद्भिषक्प}\ A; \skt{कुर्याद्भिषक्पश्चा}\ L; \skt{कुर्या चूर्ण्णकाले}\ N}}\skt{ }\edtext{\skt{शालि\-त\-ण्डु\-ल\-क\-ण्डनैः}}{\lemma{\skt{शालि\appbrk{}त\appbrk{}ण्डु\appbrk{}ल\appbrk{}क\appbrk{}ण्डनैः}}\Afootnote{K; \skt{श्चाच्छा\appbrk{}लि\appbrk{}त\appbrk{}ण्डु\appbrk{}ल\appbrk{}क\appbrk{}ण्डनैः}\ A; \skt{शालि\appbrk{}द\appbrk{}तुं\appbrk{}दु\appbrk{}ल\appbrk{}कं\appbrk{}डनैः}\ L; \skt{शालि\appbrk{}त\appbrk{}ण्डु\appbrk{}ल\appbrk{}का\appbrk{}ण्डनैः}\ N; \skt{… शालि\appbrk{}त\appbrk{}ण्डु\appbrk{}ल\appbrk{}क\appbrk{}ण्डनैः}\ H}}\skt{~।}\pend

\pstart\versepar \ledsidenote{\passref{1.12.26}}\edtext{\edtext{\skt{तिन्दु\-क्या\-स्त्व\-क्क\-षा\-यैर्वा}}{\lemma{\skt{तिन्दुक्यास्}}\Afootnote{K N; \skt{तिन्दुकी}\ A L; \skt{तिन्दुक्या}\ H}}}{\lemma{\skt{त्वक्कषायैर्वा}}\Afootnote{K H; \skt{त्वक्कपालैर्वा}\ A L; \skt{त्वक्कशायैर्वा}\ N}}\skt{ }\edtext{\skt{मृदु\-भृ\-ष्टै\-रु\-पा\-च\-रेत्}}{\lemma{\skt{मृदुभृष्टैर्}}\Afootnote{K N; \skt{घृतमिश्रैः}\ A L; \skt{मृदुभृष्ठैर्}\ H}}\skt{~। }\pend
\pstart\versepar \skt{}\edtext{\skt{व्रण\-ङ्गु\-डू\-ची\-प\-त्रैर्वा}}{\lemma{\skt{गुडूचीपत्रैर्वा}}\Afootnote{A K N H; \skt{गडूचीपत्रैर्वा}\ L}}\skt{ च्छादयेदथ }\edtext{\skt{चोदकैः}}{\lemma{\skt{चोदकैः}}\Afootnote{K; \skt{वौदकैः}\ A; \skt{वैदकैः}\ L; \skt{चोडकैः}\ N H}}\skt{~।}\edtext{}{\lemma{\skt{उपाचरेत्~। व्रणङ्}}\Afootnote{K H; \skt{प्रलेपयेत्व्रणं}\ A L; \skt{उपाचरेत्~। व्रणां}\ N}}\pend

\pstart \ledsidenote{\passref{1.12.27}}\edtext{\edtext{\skt{क्रियाङ्कुर्याच्च}}{\lemma{\skt{क्रियाङ्}}\Afootnote{K N H; \skt{क्रियां च निखिलां}\ A; \skt{क्रियां च}\ L}}}{\lemma{\skt{कुर्याच्च}}\Afootnote{K N H; \skt{कुर्या}\ A; \skt{कुर्याद}\ L}}\skt{ }\edtext{\skt{निखिलां}}{\lemma{\skt{निखिलां}}\Afootnote{K N H; \emph{om.}\ A; \skt{खिलां}\ L}}\skt{ }\edtext{\skt{भिष\-क्पि\-त्त\-वि\-स\-र्प\-वत्}}{\lemma{\skt{भिष\appbrk{}क्पि\appbrk{}त्त\appbrk{}वि\appbrk{}स\appbrk{}र्प\appbrk{}वत्}}\Afootnote{K N H; \skt{द्भिष\appbrk{}क्पि\appbrk{}त्त\appbrk{}वि\appbrk{}स\appbrk{}र्प\appbrk{}वत्~॥ मधूच्छिष्टं समधुकं रोध्रं सर्जरसं तथा}\ A; \skt{भिष\appbrk{}क्पि\appbrk{}त्त\appbrk{}वि\appbrk{}स\appbrk{}र्प्प\appbrk{}वत्~॥ … मधुछिष्टं समधुकं रोध्रं सर्जरसं तथा}\ L}}\skt{~॥}\pend

\pstart \ledsidenote{\passref{1.12.29cd}}\edtext{\skt{अथेदानीं}}{\lemma{\skt{अथेदानीं}}\Afootnote{K N H; \skt{अत ऊर्ध्वं प्रवक्ष्यामि}\ A; \skt{… अत उर्द्ध्वं प्रवक्ष्यामि}\ L}}\skt{ }\edtext{\skt{धूमो\-प\-ह\-त\-वि\-धानं}}{\lemma{\skt{धूमो\appbrk{}प\appbrk{}ह\appbrk{}त\appbrk{}वि\appbrk{}धानं}}\Afootnote{K N H; \skt{धूमो\appbrk{}प\appbrk{}ह\appbrk{}त\appbrk{}ल\appbrk{}क्ष\appbrk{}णम्}\ A L}}\skt{ }\edtext{\skt{वक्ष्यामः}}{\lemma{\skt{वक्ष्यामः}}\Afootnote{K N H; \emph{om.}\ A L}}\skt{~।}\pend

\pstart\versepar \ledsidenote{\passref{1.12.30}}\edtext{\edtext{\skt{श्वसित्याध्माति}}{\lemma{\skt{श्वसित्य्}}\Afootnote{A L K N; \skt{श्वसि}\ H}}}{\lemma{\skt{आध्माति}}\Afootnote{K; \skt{क्षौति}\ A L; \skt{आग्नोति}\ N; \skt{त्याध्माति}\ H}}\skt{ चात्यर्थं }\edtext{\skt{कासते}}{\lemma{\skt{कासते}}\Afootnote{K N H; \skt{अप्याधमति}\ A L}}\skt{ }\edtext{\skt{क्षवते}}{\lemma{\skt{क्षवते}}\Afootnote{K N H; \skt{कासते}\ A; \skt{काशत}\ L}}\skt{ }\edtext{\skt{भृशं}}{\lemma{\skt{भृशं}}\Afootnote{K H; \emph{om.}\ A L; \skt{भृषं}\ N}}\skt{~। }\pend
\pstart\versepar \skt{}\edtext{\skt{चक्षुषः}}{\lemma{\skt{चक्षुषः}}\Afootnote{K N H; \skt{चक्षुषोः}\ A L}}\skt{ परिदाहश्च राग\-श्चै\-वो\-प\-जा\-यते~।}\pend

\pstart\versepar \ledsidenote{\passref{1.12.31}}\edtext{\skt{सधूमकं}}{\lemma{\skt{सधूमकं}}\Afootnote{A L K N; \skt{समधूकन्}\ H}}\skt{ }\edtext{\skt{निश्वसिति}}{\lemma{\skt{निश्वसिति}}\Afootnote{A K N H; \skt{निःश्वसति}\ L}}\skt{ }\edtext{\edtext{\skt{घ्रेयमन्यन्न}}{\lemma{\skt{घ्रेयम्}}\Afootnote{L K N H; \skt{घ्रेय}\ A}}}{\lemma{\skt{अन्यन्न}}\Afootnote{L K H; \skt{मन्यन्न}\ A; \skt{अत्यन्न}\ N}}\skt{ वेत्ति च~। }\pend
\pstart\versepar \skt{तथैव च }\edtext{\skt{रसां}}{\lemma{\skt{रसां}}\Afootnote{L K N; \skt{रसा}\ A H}}\skt{ }\edtext{\skt{सर्वां}}{\lemma{\skt{सर्वां}}\Afootnote{K N; \skt{न्सर्वा}\ A H; \skt{सर्वा}\ L}}\skt{ }\edtext{\edtext{\skt{स्मृति\-श्चा\-सो\-प\-ह\-न्यते}}{\lemma{\skt{स्मृतिश्}}\Afootnote{K; \skt{न्श्रुतिश्}\ A L; \skt{स्मृति}\ N; \skt{न्स्मृतिश्}\ H}}}{\lemma{\skt{चासोपहन्यते}}\Afootnote{K; \skt{चास्योपहन्यते}\ A L N H}}\skt{~।}\pend

\pstart\versepar \ledsidenote{\passref{1.12.32}}\edtext{\skt{तृष्णा\-दा\-ह\-ज्व\-र\-श्चेतः}}{\lemma{\skt{तृष्णा\appbrk{}दा\appbrk{}ह\appbrk{}ज्व\appbrk{}र\appbrk{}श्चेतः}}\Afootnote{K; \skt{तृष्णा\appbrk{}दा\appbrk{}ह\appbrk{}ज्व\appbrk{}र\appbrk{}युतः}\ A L N H}}\skt{ }\edtext{\edtext{\skt{सीदत्यर्थञ्च}}{\lemma{\skt{सीदत्य्}}\Afootnote{A L N; \skt{सीद}\ K H}}}{\lemma{\skt{अर्थञ्च}}\Afootnote{\skt{अथ च}\ A L N; \emph{om.}\ K; \skt{त्यर्थञ्च}\ H}}\skt{ }\edtext{\skt{मूर्च्छति}}{\lemma{\skt{मूर्च्छति}}\Afootnote{A L N H; \emph{om.}\ K}}\skt{~। }\pend
\pstart\versepar \skt{}\edtext{\skt{धूमोपहत}}{\lemma{\skt{धूमोपहत}}\Afootnote{A L N H; \emph{om.}\ K}}\skt{ }\edtext{\edtext{\skt{इत्येष}}{\lemma{\skt{इत्य्}}\Afootnote{A L N H; \emph{om.}\ K}}}{\lemma{\skt{एष}}\Afootnote{N H; \skt{एषः}\ A; \skt{एष~॥ …}\ L; \emph{om.}\ K}}\skt{ }\edtext{\skt{शृणु}}{\lemma{\skt{शृणु}}\Afootnote{A L N H; \emph{om.}\ K}}\skt{ }\edtext{\skt{तस्य}}{\lemma{\skt{तस्य}}\Afootnote{A L N H; \emph{om.}\ K}}\skt{ }\edtext{\skt{चिकित्सितं}}{\lemma{\skt{चिकित्सितं}}\Afootnote{A L N H; \emph{om.}\ K}}\skt{~॥}\pend

\pstart\versepar \ledsidenote{\passref{1.12.33}}\edtext{\edtext{\skt{सर्पिरिक्षुरसो}}{\lemma{\skt{सर्पिर्}}\Afootnote{L N H; \skt{सर्पि}\ A}}}{\lemma{\skt{इक्षुरसो}}\Afootnote{N H; \skt{रिक्षुरसं}\ A; \skt{इक्षुरसं}\ L}}\skt{ }\edtext{\skt{द्राक्षा}}{\lemma{\skt{द्राक्षा}}\Afootnote{N H; \skt{द्राक्षां}\ A L}}\skt{ पयो वा शर्कराम्बु वा~। }\pend
\pstart\versepar \skt{}\edtext{\edtext{\skt{मधु\-रा\-म्ला\-न्र\-सां\-श्चापि}}{\lemma{\skt{मधु\appbrk{}रा\appbrk{}म्ला\appbrk{}न्र\appbrk{}सांश्}}\Afootnote{\skt{मधुराम्लौ रसौ}\ A L; \skt{मधुरांब्लां रसांश्}\ N; \skt{मधु\appbrk{}रा\appbrk{}म्ला\appbrk{}न्र\appbrk{}साश्}\ H}}}{\lemma{\skt{चापि}}\Afootnote{N H; \skt{वाऽपि}\ A L}}\skt{ }\edtext{\skt{वमनार्थाय}}{\lemma{\skt{वमनार्थाय}}\Afootnote{N H; \skt{वमनाय प्र}\ A L\textsuperscript{pc}; \skt{वमनायआ प्र}\ L\textsuperscript{ac}}}\skt{ दापयेत्~॥}\pend

\pstart\versepar \ledsidenote{\passref{1.12.34}}\skt{वमतः }\edtext{\skt{सुध्यते}}{\lemma{\skt{सुध्यते}}\Afootnote{N; \emph{om.}\ A L; \skt{शुध्यते}\ H}}\skt{ }\edtext{\skt{कोष्ठं धूमगन्धस्च}}{\lemma{\skt{कोष्ठं धूमगन्धस्च}}\Afootnote{N; \skt{कोष्ठशुद्धिः स्याद्धूमगन्धश्च}\ A; \skt{कोष्ठशुद्धिः स्यात्धूमगंधिश्च}\ L; \skt{कोष्ठधूम गन्धश्च}\ H}}\skt{ नश्यति~। }\pend
\pstart\versepar \skt{}\edtext{\skt{अनेन}}{\lemma{\skt{अनेन}}\Afootnote{N H; \emph{om.}\ A L}}\skt{ विधिना }\edtext{\skt{तस्य}}{\lemma{\skt{तस्य}}\Afootnote{N H; \skt{नेन}\ A L}}\skt{ }\edtext{\skt{कासस्वासौ}}{\lemma{\skt{कासस्वासौ}}\Afootnote{N; \skt{शाम्यन्ति स}\ A L; \skt{कासश्चासौ}\ H}}\skt{ }\edtext{\skt{प्रशस्यतः}}{\lemma{\skt{प्रशस्यतः}}\Afootnote{N; \skt{दन\appbrk{}क्ष\appbrk{}व\appbrk{}थु\appbrk{}ज्वराः}\ A L; \skt{प्रणस्यतः}\ H}}\skt{~॥}\pend

\pstart\versepar \ledsidenote{\passref{1.12.35}}\edtext{\skt{आत्मानञ्च}}{\lemma{\skt{आत्मानञ्च}}\Afootnote{N; \emph{om.}\ A L; \skt{आध्मानञ्च}\ H}}\skt{ }\edtext{\edtext{\skt{ज्वरश्चैव}}{\lemma{\skt{ज्वरश्}}\Afootnote{N H; \emph{om.}\ A L}}}{\lemma{\skt{चैव}}\Afootnote{N H; \emph{om.}\ A L}}\skt{ }\edtext{\skt{तृष्णा}}{\lemma{\skt{तृष्णा}}\Afootnote{N H; \emph{om.}\ A L}}\skt{ }\edtext{\edtext{\skt{दाहस्तथैव}}{\lemma{\skt{दाहस्}}\Afootnote{H; \skt{दाह\appbrk{}मू\appbrk{}र्च्छा\appbrk{}तृ\appbrk{}डा\appbrk{}ध्मा\appbrk{}न\appbrk{}श्वा\appbrk{}स\appbrk{}का\appbrk{}साश्}\ A; \skt{दाह\appbrk{}मू\appbrk{}र्छा\appbrk{}तृ\appbrk{}डा\appbrk{}ध्मा\appbrk{}न\appbrk{}स्वा\appbrk{}सा\appbrk{}का\appbrk{}साश्}\ L; \skt{दाह}\ N}}\skt{ च}}{\lemma{\skt{तथैव च}}\Afootnote{N H; \skt{च दारुणाः}\ A L}}\skt{~। }\pend
\pstart\versepar \skt{}\edtext{\edtext{\skt{मधु\-रै\-र्ल\-व\-णा\-म्लैश्च}}{\lemma{\skt{मधुरैर्}}\Afootnote{L N H; \skt{मधुरै}\ A}}}{\lemma{\skt{लवणाम्लैश्च}}\Afootnote{L H; \skt{र्लवणाम्लैश्च}\ A; \skt{लवणाम्ब्लैश्च}\ N}}\skt{ कटुकैः }\edtext{\skt{कवडग्रहैः}}{\lemma{\skt{कवडग्रहैः}}\Afootnote{N H; \skt{कवलग्रहैः}\ A; \skt{क्वलग्रहैः}\ L}}\skt{~॥}\pend

\pstart \ledsidenote{\passref{1.12.35ef}}\skt{वान्तस्य कण्ठशुद्धिः स्याद्धूमगन्धश्च नश्यति~।}\pend

\pstart \ledsidenote{\passref{1.12.36ab}}\edtext{\skt{सम्यग्गृह्णाती}}{\lemma{\skt{सम्यग्गृह्णाती}}\Afootnote{H; \skt{सम्यग्गृह्णातीन्}\ A L; \skt{संम्यग्गृह्णाति}\ N}}\skt{ }\edtext{\edtext{\skt{विष\-या\-न्म\-न\-श्चात्र}}{\lemma{\skt{विषयान्मनश्}}\Afootnote{N H; \skt{द्रियार्थान्मनश्}\ A L}}}{\lemma{\skt{चात्र}}\Afootnote{N H; \skt{चास्य}\ A L}}\skt{ प्रसीदति~॥}\pend

\pstart \ledsidenote{\passref{1.12.36cd}}\edtext{\skt{शिरोविरेचनं}}{\lemma{\skt{शिरोविरेचनं}}\Afootnote{A N H; \skt{शिरौविरेचनं}\ L}}\skt{ चास्मै }\edtext{\skt{दद्याद्योगेन}}{\lemma{\skt{दद्याद्योगेन}}\Afootnote{A L H; \skt{दद्याद्योगेअन}\ N}}\skt{ शास्त्रवित्~।}\pend

\pstart\versepar \ledsidenote{\passref{1.12.37}}\edtext{\skt{तेनास्य}}{\lemma{\skt{तेनास्य}}\Afootnote{N H; \skt{दृष्टिर्विशुध्यते चास्य}\ A; \skt{दृष्टि विशुध्यते चास्य}\ L}}\skt{ }\edtext{\skt{शुध्यते}}{\lemma{\skt{शुध्यते}}\Afootnote{N H; \emph{om.}\ A L}}\skt{ }\edtext{\skt{दृष्टिः}}{\lemma{\skt{दृष्टिः}}\Afootnote{N H; \emph{om.}\ A L}}\skt{ }\edtext{\edtext{\skt{सिरश्चैवास्य}}{\lemma{\skt{सिरश्}}\Afootnote{N; \skt{शिरो}\ A L; \skt{शिरश्}\ H}}}{\lemma{\skt{चैवास्य}}\Afootnote{N H\textsuperscript{pc}; \skt{ग्रीवं च}\ A L; \skt{चैवा}\ H\textsuperscript{ac}}}\skt{ देहिनः~। }\pend
\pstart\versepar \skt{}\edtext{\skt{अविदाहि}}{\lemma{\skt{अविदाहि}}\Afootnote{A N H; \skt{अविदाह}\ L}}\skt{ लघु }\edtext{\edtext{\skt{स्निग्धमाहारं}}{\lemma{\skt{स्निग्धम्}}\Afootnote{L N H; \skt{स्निग्ध}\ A}}}{\lemma{\skt{आहारं}}\Afootnote{L N H; \skt{माहारं}\ A}}\skt{ चास्य }\edtext{\skt{कल्पयेत्}}{\lemma{\skt{कल्पयेत्}}\Afootnote{A N H; \skt{कल्पयेत्~॥ …}\ L}}\skt{~।}\pend

\pstart\versepar \ledsidenote{\passref{1.12.38}}\edtext{\edtext{\skt{उष्ण\-वा\-ता\-त\-पै\-र्दग्धे}}{\lemma{\skt{उष्ण\appbrk{}वा\appbrk{}ता\appbrk{}त\appbrk{}पैर्}}\Afootnote{N H; \skt{उष्णवातातपै}\ A; \skt{… उष्ण वातातपैर्}\ L}}}{\lemma{\skt{दग्धे}}\Afootnote{L H; \skt{र्दग्धे}\ A; \skt{दग्ध}\ N}}\skt{ शीतः }\edtext{\skt{कार्यो}}{\lemma{\skt{कार्यो}}\Afootnote{A L H; \skt{कायो}\ N}}\skt{ विधिः }\edtext{\skt{सदा}}{\lemma{\skt{सदा}}\Afootnote{A N H; \skt{सदा~॥ …}\ L}}\skt{~॥ }\pend
\pstart\versepar \skt{}\edtext{\skt{शीत\-व\-र्षा\-नि\-ल\-हते}}{\lemma{\skt{शीत\appbrk{}व\appbrk{}र्षा\appbrk{}नि\appbrk{}ल\appbrk{}हते}}\Afootnote{N H; \skt{शीत\appbrk{}व\appbrk{}र्षा\appbrk{}नि\appbrk{}लै\appbrk{}र्दग्धे}\ A L}}\skt{ }\edtext{\skt{उष्णः}}{\lemma{\skt{उष्णः}}\Afootnote{N; \skt{स्निग्ध}\ A; \skt{स्निग्धम्}\ L; \skt{उष्ण}\ H}}\skt{ }\edtext{\skt{स्निग्धश्च}}{\lemma{\skt{स्निग्धश्च}}\Afootnote{N H; \skt{मुष्णं च}\ A; \skt{उष्णश्च}\ L}}\skt{ }\edtext{\skt{शस्यते}}{\lemma{\skt{शस्यते}}\Afootnote{A N H; \skt{शस्यते~॥ …}\ L}}\skt{~॥}\pend

\pstart \ledsidenote{\passref{1.12.39}}\edtext{\skt{तथातितेजसा}}{\lemma{\skt{तथातितेजसा}}\Afootnote{A N H; \skt{तथापि तेजसा}\ L}}\skt{ दग्धे }\edtext{\edtext{\skt{सिद्धिर्नैकान्तिकी}}{\lemma{\skt{सिद्धिर्}}\Afootnote{L H; \skt{सिद्ध}\ A; \skt{सि}\ N}}}{\lemma{\skt{नैकान्तिकी}}\Afootnote{H; \skt{र्नास्ति कथञ्चन~। इन्द्र\appbrk{}व\appbrk{}ज्रा\appbrk{}ग्नि\appbrk{}द\appbrk{}ग्धेऽपि जीवति प्रतिका}\ A; \skt{नास्ति कथंचनेति~॥ …}\ L; \skt{द्धि नैकान्तिकी}\ N}}\skt{ }\edtext{\skt{भवेदिति}}{\lemma{\skt{भवेदिति}}\Afootnote{N H; \skt{रयेत्~। स्नेहा\appbrk{}भ्य\appbrk{}ङ्ग\appbrk{}प\appbrk{}री\appbrk{}षेकैः प्रदेहैश्च तथा भिषक्}\ A; \emph{om.}\ L}}\skt{~॥}\pend

\pstart \ledsidenote{\passref{1.12.trailer}}\edtext{\skt{॥ १२~॥}}{\lemma{\skt{॥ १२~॥}}\Afootnote{\skt{इति सुश्रु\appbrk{}त\appbrk{}सं\appbrk{}हि\appbrk{}तायां सूत्र\appbrk{}स्था\appbrk{}नेऽ\appbrk{}ग्नि\appbrk{}क\appbrk{}र्म\appbrk{}वि\appbrk{}धि\appbrk{}र्नाम द्वादशोऽध्यायः}\ A; \skt{इति सूत्र\appbrk{}स्था\appbrk{}न\appbrk{}नि\appbrk{}बं\appbrk{}ध\appbrk{}सं\appbrk{}ग्रहे द्वादशोऽध्यायः}\ L; \skt{१२}\ N; \skt{१२~॥ ०~॥}\ H}}\pend

\endnumbering
\beginnumbering
\eledsection*{Adhyāya 13}
\pstart \ledsidenote{\passref{1.13.1}}\skt{अथातो }\edtext{\skt{जलायुकाध्यायं}}{\lemma{\skt{जलायुकाध्यायं}}\Afootnote{N; \skt{जलौ\appbrk{}का\appbrk{}व\appbrk{}चा\appbrk{}र\appbrk{}णी\appbrk{}य\appbrk{}म\appbrk{}ध्यायं}\ A; \skt{जला\appbrk{}यु\appbrk{}ष्का\appbrk{}ध्या\appbrk{}यम्}\ H}}\skt{ व्याख्यास्यामः~॥}\pend

\pstart \ledsidenote{\passref{1.13.3}}\edtext{\skt{नृपा\-ढ्य\-सु\-कु\-मा\-र\-बा\-ल\-स्थ\-वि\-र\-भी\-रु\-ना\-री\-णा\-म\-नु\-ग्र\-हार्थं}}{\lemma{\skt{नृपा\appbrk{}ढ्य\appbrk{}सु\appbrk{}कु\appbrk{}मा\appbrk{}र\appbrk{}बा\appbrk{}ल\appbrk{}स्थ\appbrk{}वि\appbrk{}र\appbrk{}भी\appbrk{}रु\appbrk{}ना\appbrk{}री\appbrk{}णाम्}}\Afootnote{N H; \skt{नृपा\appbrk{}ढ्य\appbrk{}बा\appbrk{}ल\appbrk{}स्थ\appbrk{}वि\appbrk{}र\appbrk{}भी\appbrk{}रु\appbrk{}दु\appbrk{}र्ब\appbrk{}ल\appbrk{}ना\appbrk{}री\appbrk{}सु\appbrk{}कु\appbrk{}मा\appbrk{}रा\appbrk{}णाम्}\ A}}\skt{ }\edtext{\skt{पर\-म\-सु\-कु\-मा\-रोऽयं}}{\lemma{\skt{ऽयं}}\Afootnote{A N; \skt{यं}\ H}}\skt{ }\edtext{\edtext{\skt{शोणि\-ता\-व\-से\-च\-नो\-पा\-योऽ\-भि\-हितो}}{\lemma{\skt{शोणि\appbrk{}ता\appbrk{}व\appbrk{}से\appbrk{}च\appbrk{}नो\appbrk{}पायो}}\Afootnote{A; \skt{शोणि\appbrk{}ता\appbrk{}व\appbrk{}से\appbrk{}च\appbrk{}नो\appbrk{}पाया}\ N; \skt{शोणि\appbrk{}ता\appbrk{}व\appbrk{}से\appbrk{}नो\appbrk{}पायो}\ H}}}{\lemma{\skt{ऽभिहितो}}\Afootnote{A H; \skt{भिहितो}\ N}}\skt{ जलौकसः~॥}\pend

\pstart \ledsidenote{\passref{1.13.4}}\skt{तत्र वात\-पि\-त्त\-क\-फ\-दु\-ष्ट\-शो\-णितं }\edtext{\skt{यथासंख्यं}}{\lemma{\skt{यथासंख्यं}}\Afootnote{A H; \skt{यथासख्यं}\ N}}\skt{ शृङ्ग\-ज\-लौ\-का\-ला\-बु\-भि\-र\-व\-से\-च\-येत्~। सर्वाणि सर्वैर्वा }\edtext{\skt{विस्राव्यं}}{\lemma{\skt{विस्राव्यं}}\Afootnote{\skt{(विशेषस्तु विस्राव्यं शृङ्ग\appbrk{}ज\appbrk{}लौ\appbrk{}का\appbrk{}ला\appbrk{}बु\appbrk{}भि\appbrk{}र्गृ\appbrk{}ह्णी\appbrk{}यात्)}\ A; \skt{विशेषस्तु विश्राव्यं}\ N; \skt{विशेषस्तु विस्राव्यं}\ H}}\skt{~॥}\footnoteA{The 1931 and 1938 editions of the vulgate report an insertion at this point, although they report it differently. Ḍalhaṇa reported variant readings at this point in the text.}\pend

\pstart\versepar \ledsidenote{\passref{1.13.5}}\skt{भवन्ति चात्र~॥ }\pend
\pstart\versepar \skt{स्निग्धं }\edtext{\skt{श्लक्ष्णं}}{\lemma{\skt{श्लक्ष्णं}}\Afootnote{N H; \skt{उष्णं}\ A}}\skt{ }\edtext{\skt{समधुरं}}{\lemma{\skt{समधुरं}}\Afootnote{A N\textsuperscript{pc}; \skt{समधुराङ्}\ N\textsuperscript{ac}; \skt{सुमधुरं}\ H}}\skt{ }\edtext{\skt{गवां}}{\lemma{\skt{गवां}}\Afootnote{H; \skt{स्निग्धं गवां}\ A; \skt{गवाङ्ग}\ N}}\skt{ }\edtext{\skt{शृङ्गं}}{\lemma{\skt{शृङ्गं}}\Afootnote{A H; \skt{गं}\ N}}\skt{ प्रकीर्त्तितं~। }\pend
\pstart\versepar \skt{तस्मा\-द्वा\-तो\-प\-सृष्टे तु हित\-न्त\-द\-व\-से\-चने~॥}\edtext{}{\lemma{\skt{भवन्ति चात्र~॥ स्निग्धं}}\Afootnote{\skt{भवन्ति चात्र श्लोकाः}\ A; \skt{भवन्ति चात्र~॥ स्निग्धं}\ N H}}\pend

\pstart\versepar \ledsidenote{\passref{1.13.5.1}}\skt{अर्द्धचन्द्राकृति महत्तनु सप्ताङ्गुलायतं~। }\pend
\pstart\versepar \skt{प्रच्छिते }\edtext{\skt{दाप\-ये\-त्पू\-र्व्व\-मा\-स्ये\-ना\-चू\-ष\-ये\-द्बली}}{\lemma{\skt{आस्ये\appbrk{}ना\appbrk{}चू\appbrk{}ष\appbrk{}ये\appbrk{}द्बली}}\Afootnote{H; \skt{आस्ये\appbrk{}ना\appbrk{}चू\appbrk{}ष\appbrk{}य\appbrk{}न्वली}\ N}}\skt{~॥ }\footnoteA{The 1938 edition of the vulgate (p. 55, note 4) reports the insertion of this verse after 1.13.5 at this point. Ḍalhaṇa also reported this reading.}\pend

\pstart\versepar \ledsidenote{\passref{1.13.6}}\skt{शीताधिवासा मधुरा जलौका वारिसंभवा~। }\pend
\pstart\versepar \skt{तस्मा\-त्पि\-त्तो\-प\-सृष्टे तु }\edtext{\skt{हितास्ता}}{\lemma{\skt{हितास्ता}}\Afootnote{N H; \skt{हिता सा}\ A}}\skt{ }\edtext{\skt{अवसेचने}}{\lemma{\skt{अवसेचने}}\Afootnote{N H; \skt{त्ववसेचने}\ A}}\skt{~॥}\pend

\pstart\versepar \ledsidenote{\passref{1.13.7}}\edtext{\skt{कटु}}{\lemma{\skt{कटु}}\Afootnote{N H; \skt{अलाबु कटुकं}\ A}}\skt{ }\edtext{\skt{रूक्षञ्च}}{\lemma{\skt{रूक्षञ्च}}\Afootnote{N H; \skt{रूक्षं}\ A}}\skt{ तीक्ष्णञ्च }\edtext{\skt{अलाबु}}{\lemma{\skt{अलाबु}}\Afootnote{H; \emph{om.}\ A; \skt{अलाबुः}\ N}}\skt{ परिकीर्त्तितं~।}\footnoteA{The hiatus between \emph{ca} and \emph{alābu} is required for correct metre.}\skt{ }\pend
\pstart\versepar \skt{तस्मा\-च्छ्ले\-ष्मो\-प\-सृष्टे तु हित\-न्त\-द\-व\-से\-चने~॥}\pend

\pstart \ledsidenote{\passref{1.13.8}}\skt{तत्र प्रच्छिते }\edtext{\skt{तनु\-ब\-स्ति\-प\-ट\-ला\-व\-न\-द्धेन}}{\lemma{\skt{तनु\appbrk{}ब\appbrk{}स्ति\appbrk{}प\appbrk{}ट\appbrk{}ला\appbrk{}व\appbrk{}न\appbrk{}द्धेन}}\Afootnote{A H; \skt{तनु\appbrk{}ब\appbrk{}स्थि\appbrk{}प\appbrk{}ट\appbrk{}ला\appbrk{}बु\appbrk{}न\appbrk{}द्धेन}\ N}}\skt{ शृङ्गेण }\edtext{\skt{शोणि\-त\-म\-व\-से\-च\-येत्~। आचू\-ष\-णा\-द\-न्त\-र्द्दी\-प्ते\-ना\-ला\-बुना}}{\lemma{\skt{अवसेचयेत्~। आचू\appbrk{}ष\appbrk{}णा\appbrk{}द\appbrk{}न्त\appbrk{}र्द्दी\appbrk{}प्ते\appbrk{}ना\appbrk{}ला\appbrk{}बुना}}\Afootnote{H; \skt{अवसेचये, आचूषणात्, सान्त\appbrk{}र्दी\appbrk{}प\appbrk{}याऽ\appbrk{}लब्धा~। (जलायुका वक्ष्यन्ते)}\ A; \skt{असे\appbrk{}च\appbrk{}ये\appbrk{}त\appbrk{}चू\appbrk{}ष\appbrk{}णा\appbrk{}द\appbrk{}न्त\appbrk{}र्दी\appbrk{}प्ते\appbrk{}ना\appbrk{}ला\appbrk{}बुना}\ N\textsuperscript{ac}; \skt{अव\appbrk{}से\appbrk{}च\appbrk{}ये\appbrk{}त\appbrk{}चू\appbrk{}ष\appbrk{}णा\appbrk{}द\appbrk{}न्त\appbrk{}र्दी\appbrk{}प्ते\appbrk{}ना\appbrk{}ला\appbrk{}बुना}\ N\textsuperscript{pc}}}\skt{~।}\pend

\pstart \ledsidenote{\passref{1.13.9}}\edtext{\edtext{\skt{जल\-मा\-सा\-मा\-यु\-रि\-त्यतो}}{\lemma{\skt{इत्य्}}\Afootnote{N H; \skt{इ}\ A}}}{\lemma{\skt{अतो}}\Afootnote{N H; \emph{om.}\ A}}\skt{ }\edtext{\skt{जलायुकाः}}{\lemma{\skt{जलायुकाः}}\Afootnote{N H; \skt{त्य्जलायुकाः}\ A}}\skt{~। }\edtext{\skt{ओको}}{\lemma{\skt{ओको}}\Afootnote{N H; \emph{om.}\ A}}\skt{ }\edtext{\skt{निवासो}}{\lemma{\skt{निवासो}}\Afootnote{N H; \emph{om.}\ A}}\skt{ जलमासामोक }\edtext{\edtext{\skt{इत्यतो}}{\lemma{\skt{इत्य्}}\Afootnote{N H; \skt{इ}\ A}}}{\lemma{\skt{अतो}}\Afootnote{N H; \emph{om.}\ A}}\skt{ }\edtext{\skt{जलौकसः}}{\lemma{\skt{जलौकसः}}\Afootnote{N H; \skt{त्य्जलौकसः}\ A}}\skt{~॥}\pend

\pstart \ledsidenote{\passref{1.13.10}}\edtext{\skt{ता द्वादश}}{\lemma{\skt{ता द्वादश}}\Afootnote{A H; \skt{ताः द्वाद}\ N}}\skt{~। }\edtext{\skt{सविषाः}}{\lemma{\skt{सविषाः}}\Afootnote{H; \skt{तासां सविषाः}\ A; \skt{शः सविषाः}\ N}}\skt{ }\edtext{\skt{षट्~। तावन्त्य}}{\lemma{\skt{षट्~। तावन्त्य}}\Afootnote{N H; \skt{षट्, तावत्य}\ A}}\skt{ }\edtext{\skt{एव}}{\lemma{\skt{एव}}\Afootnote{A H; \skt{एव निव}\ N}}\skt{ निर्विषाः~।}\pend

\pstart \ledsidenote{\passref{1.13.11}}\edtext{\skt{कृष्णा}}{\lemma{\skt{कृष्णा}}\Afootnote{N H; \skt{तत्र सविषाः - कृष्णा}\ A}}\skt{ कर्बुरा अलगर्दा इन्द्रायुधा सामुद्रिका गोवन्दना चेति~। }\edtext{\skt{तास्वञ्जनवर्ण्णा}}{\lemma{\skt{अञ्जनवर्ण्णा}}\Afootnote{N H; \skt{अञ्ज\appbrk{}न\appbrk{}चू\appbrk{}र्ण\appbrk{}वर्णा}\ A}}\skt{ }\edtext{\skt{पृथुशीर्षा}}{\lemma{\skt{पृथुशीर्षा}}\Afootnote{N H; \skt{पृथुशीराः}\ A}}\skt{ कृष्णा }\edtext{\skt{नाम}}{\lemma{\skt{नाम}}\Afootnote{N H; \emph{om.}\ A}}\skt{~। }\edtext{\skt{वर्मि\-म\-त्स्य\-व\-दा\-यता}}{\lemma{\skt{वर्मि\appbrk{}म\appbrk{}त्स्य\appbrk{}व\appbrk{}दा\appbrk{}यता}}\Afootnote{A; \skt{वर्मि\appbrk{}व\appbrk{}त्स्य\appbrk{}व\appbrk{}दा\appbrk{}यता}\ N; \skt{चर्म्मि\appbrk{}म\appbrk{}त्स्य\appbrk{}व\appbrk{}दा\appbrk{}यता}\ H}}\skt{ च्छिन्नोन्नतकुक्षिः कर्बुरा }\edtext{\skt{नाम}}{\lemma{\skt{नाम}}\Afootnote{N H; \emph{om.}\ A}}\skt{~। }\edtext{\skt{रोमशा}}{\lemma{\skt{रोमशा}}\Afootnote{A N; \skt{रोमषा}\ H}}\skt{ महापार्श्वा }\edtext{\skt{कृष्णमुखा}}{\lemma{\skt{कृष्णमुखा}}\Afootnote{N H; \skt{कृष्णमुखी}\ A}}\skt{ अलगर्दा }\edtext{\skt{नाम}}{\lemma{\skt{नाम}}\Afootnote{N H; \emph{om.}\ A}}\skt{~। }\edtext{\skt{इन्द्रा\-यु\-ध\-व\-दू\-र्द्ध्व\-राजी}}{\lemma{\skt{इन्द्रा\appbrk{}यु\appbrk{}ध\appbrk{}व\appbrk{}दू\appbrk{}र्द्ध्व\appbrk{}राजी}}\Afootnote{N\textsuperscript{pc} H; \skt{इन्द्रा\appbrk{}यु\appbrk{}ध\appbrk{}व\appbrk{}दू\appbrk{}र्द्ध्व\appbrk{}रा\appbrk{}ज्भिश्}\ A; \skt{इन्द्रा\appbrk{}यु\appbrk{}धा\appbrk{}व\appbrk{}दू\appbrk{}र्ध्व\appbrk{}राजी}\ N\textsuperscript{ac}}}\skt{ चित्रा }\edtext{\skt{इन्द्रायुधा}}{\lemma{\skt{इन्द्रायुधा}}\Afootnote{A H; \skt{इन्द्रयुधा}\ N}}\skt{ }\edtext{\skt{नाम~। ईष\-द\-सि\-त\-पी\-तिका}}{\lemma{\skt{नाम~। ईष\appbrk{}द\appbrk{}सि\appbrk{}त\appbrk{}पी\appbrk{}तिका}}\Afootnote{H; \skt{ईष\appbrk{}द\appbrk{}सि\appbrk{}त\appbrk{}पी\appbrk{}तिका}\ A; \skt{ईष\appbrk{}द\appbrk{}षि\appbrk{}त\appbrk{}पी\appbrk{}तिका}\ N}}\skt{ }\edtext{\skt{विचि\-त्र\-पु\-ष्पा\-कृ\-ति\-चिता}}{\lemma{\skt{विचि\appbrk{}त्र\appbrk{}पु\appbrk{}ष्पा\appbrk{}कृ\appbrk{}ति\appbrk{}चिता}}\Afootnote{N H; \skt{विचि\appbrk{}त्र\appbrk{}पु\appbrk{}ष्पा\appbrk{}कृ\appbrk{}ति\appbrk{}चित्रा}\ A}}\skt{ सामुद्रिका }\edtext{\skt{नाम}}{\lemma{\skt{नाम}}\Afootnote{N H; \emph{om.}\ A}}\skt{~। गोवृ\-ष\-ण\-व\-द\-धो\-भागे }\edtext{\skt{द्विधा\-भू\-ता\-कृ\-ति\-र\-णु\-मुखी}}{\lemma{\skt{अणुमुखी}}\Afootnote{A H; \skt{अणमुखी}\ N}}\skt{ }\edtext{\skt{गोवन्दना}}{\lemma{\skt{गोवन्दना}}\Afootnote{N H; \skt{गोवन्द}\ A}}\skt{ }\edtext{\skt{नाम}}{\lemma{\skt{नाम}}\Afootnote{N H; \skt{नेति}\ A}}\skt{~। ताभिर्दष्टे }\edtext{\skt{दंशे}}{\lemma{\skt{दंशे}}\Afootnote{N H; \skt{पुरुषे दंशे}\ A}}\skt{ श्वय\-थु\-र\-ति\-मात्रं कण्डूमूर्च्छा ज्वरो }\edtext{\edtext{\edtext{\skt{दाहश्छर्द्दिरिति}}{\lemma{\skt{दाहश्}}\Afootnote{A H; \skt{दाह}\ N}}}{\lemma{\skt{छर्द्दिर्}}\Afootnote{A H; \skt{च्छर्द्दिर्}\ N}}}{\lemma{\skt{इति}}\Afootnote{N H; \skt{मदः सदनमिति}\ A}}\skt{ लिङ्गानि भवन्ति~। तत्र महागदः }\edtext{\skt{पाना\-ले\-प\-ना\-दि\-षू\-प\-योज्यः}}{\lemma{\skt{पाना\appbrk{}ले\appbrk{}प\appbrk{}ना\appbrk{}दि\appbrk{}षू\appbrk{}प\appbrk{}योज्यः}}\Afootnote{N; \skt{पाना\appbrk{}ले\appbrk{}प\appbrk{}न\appbrk{}न\appbrk{}स्य्क\appbrk{}र्मा\appbrk{}दि\appbrk{}षू\appbrk{}प\appbrk{}योज्यः}\ A; \skt{पाना\appbrk{}ले\appbrk{}प\appbrk{}ना\appbrk{}दि\appbrk{}षू\appbrk{}प\appbrk{}योज्याः}\ H}}\skt{~। इन्द्रा\-यु\-धा\-द\-ष्ट\-म\-सा\-ध्य\-मि\-त्येताः सविषाः सचिकित्सिता व्याख्याताः~॥}\pend

\pstart \ledsidenote{\passref{1.13.12}}\skt{अथ }\edtext{\skt{निर्विषाः}}{\lemma{\skt{निर्विषाः}}\Afootnote{A H; \skt{निर्विषा}\ N}}\skt{~। कपिला पिङ्गला शङ्कुमुखी }\edtext{\skt{मूषिका}}{\lemma{\skt{मूषिका}}\Afootnote{A H; \skt{मूपिका}\ N}}\skt{ पुण्डरीकमुखी सावरिका चेति~। तत्र मनः\-शि\-ला\-र\-ञ्जि\-ता\-भ्या\-मिव पार्श्वाभ्यां }\edtext{\skt{पृष्ठे}}{\lemma{\skt{पृष्ठे}}\Afootnote{A H; \skt{पृष्टे}\ N}}\skt{ }\edtext{\skt{स्निग्धमुद्गवर्ण्णा}}{\lemma{\skt{स्निग्धमुद्गवर्ण्णा}}\Afootnote{H; \skt{स्निग्धा मुद्गवर्णा}\ A; \skt{स्निग्धवर्ण्णा}\ N}}\skt{ कपिला }\edtext{\skt{नाम}}{\lemma{\skt{नाम}}\Afootnote{N H; \emph{om.}\ A}}\skt{~। }\edtext{\skt{किञ्चिद्रक्ता}}{\lemma{\skt{किञ्चिद्रक्ता}}\Afootnote{A N; \skt{…किञ्चिद्रक्ता}\ H}}\skt{ वृत्तकाया }\edtext{\skt{पिङ्गल्याशुगा}}{\lemma{\skt{पिङ्गल्याशुगा}}\Afootnote{N H; \skt{पिङ्गाऽऽशुगा च}\ A}}\skt{ पिङ्गला }\edtext{\skt{नाम~। यकृद्वर्ण्णा}}{\lemma{\skt{नाम~। यकृद्वर्ण्णा}}\Afootnote{H; \skt{यकृद्वर्णा}\ A N}}\skt{ शीघ्रपायिनी }\edtext{\skt{दीर्घमुखी}}{\lemma{\skt{दीर्घमुखी}}\Afootnote{N H; \skt{दीर्घतीक्ष्णमुखी}\ A}}\skt{ शङ्कुमुखी }\edtext{\skt{नाम~। मूषि\-का\-कृ\-ति\-व\-र्ण्णा\-नि\-ष्ट\-गन्धा}}{\lemma{\skt{नाम~। मूषि\appbrk{}का\appbrk{}कृ\appbrk{}ति\appbrk{}व\appbrk{}र्ण्णा\appbrk{}नि\appbrk{}ष्ट\appbrk{}गन्धा}}\Afootnote{\skt{मूषि\appbrk{}का\appbrk{}कृ\appbrk{}ति\appbrk{}व\appbrk{}र्णाऽ\appbrk{}नि\appbrk{}ष्ट\appbrk{}गन्धा च}\ A; \skt{दूषि\appbrk{}का\appbrk{}कृ\appbrk{}ति\appbrk{}व\appbrk{}र्ण्णा\appbrk{}नि\appbrk{}ष्ट\appbrk{}गन्धा}\ N\textsuperscript{ac}; \skt{विदू\appbrk{}षि\appbrk{}का\appbrk{}कृ\appbrk{}ति\appbrk{}व\appbrk{}र्ण्णा\appbrk{}नि\appbrk{}ष्ट\appbrk{}गन्धा}\ N\textsuperscript{pc}; \skt{नाम।… मूषि\appbrk{}का\appbrk{}कृ\appbrk{}ति\appbrk{}व\appbrk{}र्ण्णा\appbrk{}नि\appbrk{}ष्ट\appbrk{}गन्धा}\ H}}\skt{ मूषिका }\edtext{\skt{नाम}}{\lemma{\skt{नाम}}\Afootnote{N H; \emph{om.}\ A}}\skt{~। मुद्गवर्ण्णा }\edtext{\skt{पुण्ड\-री\-क\-तु\-ल्य\-वक्त्रा}}{\lemma{\skt{पुण्ड\appbrk{}री\appbrk{}क\appbrk{}तु\appbrk{}ल्य\appbrk{}वक्त्रा}}\Afootnote{A H; \skt{पुण्ड\appbrk{}री\appbrk{}क\appbrk{}तु\appbrk{}ल्य\appbrk{}वक्ता}\ N}}\skt{ }\edtext{\skt{पुण्डरीका}}{\lemma{\skt{पुण्डरीका}}\Afootnote{N H; \skt{पुण्डरीक}\ A}}\skt{ }\edtext{\skt{नाम}}{\lemma{\skt{नाम}}\Afootnote{N H; \skt{मुखी}\ A}}\skt{~। }\edtext{\skt{पद्म\-प\-त्र\-व\-र्ण्णा\-ष्टा\-द\-शा\-ङ्गु\-ल\-प्र\-माणा}}{\lemma{\skt{पद्म\appbrk{}प\appbrk{}त्र\appbrk{}व\appbrk{}र्ण्णा\appbrk{}ष्टा\appbrk{}द\appbrk{}शा\appbrk{}ङ्गु\appbrk{}ल\appbrk{}प्र\appbrk{}माणा}}\Afootnote{N\textsuperscript{pc} H; \skt{स्निग्धा पद्म\appbrk{}प\appbrk{}त्र\appbrk{}व\appbrk{}र्णाऽ\appbrk{}ष्टा\appbrk{}द\appbrk{}शा\appbrk{}ङ्गु\appbrk{}ल\appbrk{}प्र\appbrk{}माणा}\ A; \skt{पद्म\appbrk{}प\appbrk{}त्र\appbrk{}व\appbrk{}र्ण्णा\appbrk{}ष्टा\appbrk{}द\appbrk{}शां\appbrk{}द\appbrk{}गु\appbrk{}ल\appbrk{}प्र\appbrk{}माणा}\ N\textsuperscript{ac}}}\skt{ सावरिका }\edtext{\skt{नाम~। सा}}{\lemma{\skt{नाम~। सा}}\Afootnote{N H; \skt{सा च}\ A}}\skt{ पश्वर्थे }\edtext{\skt{त्वविषा}}{\lemma{\skt{त्वविषा}}\Afootnote{N H; \skt{इत्येता अविषा}\ A}}\skt{ }\edtext{\skt{व्याख्याताः}}{\lemma{\skt{व्याख्याताः}}\Afootnote{A H; \skt{व्याख्याता}\ N}}\skt{~॥}\pend

\pstart \ledsidenote{\passref{1.13.13}}\skt{तासां }\edtext{\skt{यव\-न\-पा\-ण्ड्य\-स\-ह्य\-पो\-त\-ना\-दीनि}}{\lemma{\skt{यव\appbrk{}न\appbrk{}पा\appbrk{}ण्ड्य\appbrk{}स\appbrk{}ह्य\appbrk{}पो\appbrk{}त\appbrk{}ना\appbrk{}दीनि}}\Afootnote{N H; \skt{यव\appbrk{}न\appbrk{}पा\appbrk{}ण्ड्य\appbrk{}स\appbrk{}ह्य\appbrk{}पौ\appbrk{}त\appbrk{}ना\appbrk{}दीनि}\ A}}\skt{ क्षेत्राणि }\edtext{\skt{भवन्ति}}{\lemma{\skt{भवन्ति}}\Afootnote{N H; \emph{om.}\ A}}\skt{~। }\edtext{\skt{तासां}}{\lemma{\skt{तासां}}\Afootnote{N H; \skt{तेषु}\ A}}\skt{ महाशरीरा बलवत्यः शीघ्रपायिन्यो महाशना निर्विषाश्च विशेषेण भवन्ति~।}\pend

\pstart \ledsidenote{\passref{1.13.14}}\skt{तत्र }\edtext{\skt{सवि\-ष\-की\-ट\-द\-र्द्दु\-र\-मू\-त्र\-पु\-री\-ष\-को\-थ\-जाताः}}{\lemma{\skt{सवि\appbrk{}ष\appbrk{}की\appbrk{}ट\appbrk{}द\appbrk{}र्द्दु\appbrk{}र\appbrk{}मू\appbrk{}त्र\appbrk{}पु\appbrk{}री\appbrk{}ष\appbrk{}को\appbrk{}थ\appbrk{}जाताः}}\Afootnote{H; \skt{सवि\appbrk{}ष\appbrk{}म\appbrk{}त्स्य\appbrk{}की\appbrk{}ट\appbrk{}द\appbrk{}र्दु\appbrk{}र\appbrk{}मू\appbrk{}त्र\appbrk{}पु\appbrk{}री\appbrk{}ष\appbrk{}को\appbrk{}थ\appbrk{}जाताः}\ A; \skt{सवि\appbrk{}ष\appbrk{}की\appbrk{}टा\appbrk{}द\appbrk{}र्द्दु\appbrk{}र\appbrk{}मू\appbrk{}त्र\appbrk{}पु\appbrk{}री\appbrk{}ष\appbrk{}को\appbrk{}थ\appbrk{}जाताः}\ N}}\skt{ }\edtext{\edtext{\skt{कलुषेष्वम्भःसु}}{\lemma{\skt{कलुषेष्व्}}\Afootnote{A H; \skt{कलुष्वेष्व्}\ N}}\skt{ च}}{\lemma{\skt{अम्भःसु च}}\Afootnote{N H; \skt{अम्भसु च}\ A}}\skt{ सविषाः~। }\edtext{\skt{पद्मो\-त्प\-ल\-कु\-मु\-द\-सौ\-ग\-न्धि\-क\-शै\-वा\-ल\-को\-थ\-जाता}}{\lemma{\skt{पद्मो\appbrk{}त्प\appbrk{}ल\appbrk{}कु\appbrk{}मु\appbrk{}द\appbrk{}सौ\appbrk{}ग\appbrk{}न्धि\appbrk{}क\appbrk{}शै\appbrk{}वा\appbrk{}ल\appbrk{}को\appbrk{}थ\appbrk{}जाता}}\Afootnote{\skt{पद्मो\appbrk{}त्प\appbrk{}ल\appbrk{}न\appbrk{}लि\appbrk{}न\appbrk{}कु\appbrk{}मु\appbrk{}द\appbrk{}सौ\appbrk{}ग\appbrk{}न्धि\appbrk{}क\appbrk{}कु\appbrk{}व\appbrk{}ल\appbrk{}य\appbrk{}पु\appbrk{}ण्ड\appbrk{}री\appbrk{}क\appbrk{}शै\appbrk{}व\appbrk{}ल\appbrk{}को\appbrk{}थ\appbrk{}जाता}\ A; \skt{पद्मो\appbrk{}प\appbrk{}ल\appbrk{}कु\appbrk{}मु\appbrk{}द\appbrk{}सौ\appbrk{}ग\appbrk{}न्धि\appbrk{}क\appbrk{}शै\appbrk{}वा\appbrk{}ल\appbrk{}को\appbrk{}थ\appbrk{}जाता}\ N; \skt{पद्मो\appbrk{}त्प\appbrk{}ल\appbrk{}कु\appbrk{}मु\appbrk{}द\appbrk{}सौ\appbrk{}ग\appbrk{}न्धि\appbrk{}क\appbrk{}शे\appbrk{}वा\appbrk{}ल\appbrk{}को\appbrk{}थ\appbrk{}जाता}\ H}}\skt{ }\edtext{\skt{विमलेष्वम्भःसु}}{\lemma{\skt{विमलेष्व्}}\Afootnote{A N; \skt{विमले}\ H}}\skt{ च निर्व्विषाः~॥}\pend

\pstart\versepar \ledsidenote{\passref{1.13.15}}\edtext{\skt{भवति}}{\lemma{\skt{भवति}}\Afootnote{\skt{भवन्ति}\ A; \emph{om.}\ N H}}\skt{ }\edtext{\skt{चात्र}}{\lemma{\skt{चात्र}}\Afootnote{A; \emph{om.}\ N; \skt{भवति चात्र}\ H}}\skt{~॥ }\pend
\pstart\versepar \skt{क्षेत्रेषु विचरन्त्येताः सलि\-ला\-ढ्य\-सु\-ग\-न्धिषु~। }\pend
\pstart\versepar \skt{न च सङ्कीर्ण्णचारिण्यो न च }\edtext{\skt{पङ्केशयाः}}{\lemma{\skt{पङ्केशयाः}}\Afootnote{A H; \skt{पंकोशयाः}\ N}}\skt{ }\edtext{\skt{स्मृताः}}{\lemma{\skt{स्मृताः}}\Afootnote{H; \skt{सुखाः}\ A; \skt{स्मृता}\ N}}\skt{~॥}\pend

\pstart \ledsidenote{\passref{1.13.16}}\skt{तासां }\edtext{\skt{ग्रह\-ण\-मा\-र्द्र\-च\-र्म\-णा\-न्यैर्वा}}{\lemma{\skt{आर्द्र\appbrk{}च\appbrk{}र्म\appbrk{}णा\appbrk{}न्यैर्वा}}\Afootnote{N H; \skt{आर्द्रचर्मणा, अन्यैर्वा}\ A}}\skt{ }\edtext{\skt{प्रयो\-गै\-र्गृ\-ह्णी\-यात्}}{\lemma{\skt{गृह्णीयात्}}\Afootnote{A; \skt{गृहीत्वा}\ N H}}\footnoteA{The transmitted Nepalese reading \emph{gṛhītvā} is hard to construe unless taken over to the start of the next paragraph. However, then it would sit uneasily before \emph{atha}, which seems to demarcate the start of different procedural stages.}\skt{~।}\pend

\pstart \ledsidenote{\passref{1.13.17}}\edtext{\skt{अथैता}}{\lemma{\skt{अथैता}}\Afootnote{N H; \skt{अथैनां}\ A}}\skt{ नवे महति }\edtext{\skt{घटे}}{\lemma{\skt{घटे}}\Afootnote{A H; \skt{घटो}\ N}}\skt{ }\edtext{\skt{सर\-स्त\-डा\-को\-द\-क\-प\-ङ्का\-ना\-वाप्य}}{\lemma{\skt{सर\appbrk{}स्त\appbrk{}डा\appbrk{}को\appbrk{}द\appbrk{}क\appbrk{}प\appbrk{}ङ्का\appbrk{}ना\appbrk{}वाप्य}}\Afootnote{N H; \skt{सर\appbrk{}स्त\appbrk{}डा\appbrk{}गो\appbrk{}द\appbrk{}क\appbrk{}प\appbrk{}ङ्क\appbrk{}मा\appbrk{}वाप्य}\ A}}\skt{ }\edtext{\skt{निदध्यात्~। भक्ता\-र्थ\-ञ्चा\-सा\-मु\-प\-ह\-रेत्}}{\lemma{\skt{निदध्यात्~। भक्तार्थञ्}}\Afootnote{H; \skt{निदध्यात्; भक्ष्यार्थे}\ A; \skt{निदद्यात्~। भक्तार्थञ्}\ N}}\skt{~। शेवालं }\edtext{\skt{वल्लू\-र\-मो\-द\-कांश्च}}{\lemma{\skt{ओदकांश्च}}\Afootnote{H; \skt{औदकांश्च}\ A; \skt{ओदकाश्च}\ N}}\skt{ }\edtext{\skt{कन्दांश्चूर्णीकृत्य}}{\lemma{\skt{कन्दांश्}}\Afootnote{A H; \skt{कन्दां}\ N}}\skt{ }\edtext{\skt{शय्यार्थे}}{\lemma{\skt{शय्यार्थे}}\Afootnote{N H; \skt{शय्यार्थं}\ A}}\skt{ }\edtext{\skt{तृणमोदकानि}}{\lemma{\skt{ओदकानि}}\Afootnote{N H; \skt{औदकानि च}\ A}}\skt{ पत्राणि }\edtext{\skt{त्र्यहात्त्र्यहाच्चासां}}{\lemma{\skt{त्र्यहात्त्र्यहाच्चासां}}\Afootnote{H; \skt{त्र्यहा\appbrk{}त्त्र्य\appbrk{}हा\appbrk{}च्चा\appbrk{}भ्योऽन्य}\ A; \skt{त्र्यहा त्र्यहाच्चासां}\ N}}\skt{ }\edtext{\edtext{\skt{जल\-भ\-क्त\-न्द\-द्यात्}}{\lemma{\skt{जलभक्तन्}}\Afootnote{N\textsuperscript{pc} H; \skt{ज्जलं भक्ष्यं च}\ A; \skt{जभक्तं}\ N\textsuperscript{ac}}}\skt{~। सप्त\-रा\-त्रा\-त्स\-प्त\-रा\-त्रा\-द्घ\-ट\-मन्यं}}{\lemma{\skt{दद्यात्~। सप्त\appbrk{}रा\appbrk{}त्रा\appbrk{}त्स\appbrk{}प्त\appbrk{}रा\appbrk{}त्रा\appbrk{}द्घ\appbrk{}टम्}}\Afootnote{N H\textsuperscript{pc}; \skt{दद्यात्; सप्त\appbrk{}रा\appbrk{}त्रा\appbrk{}त्स\appbrk{}प्त\appbrk{}रा\appbrk{}त्राच्च घटम्}\ A; \skt{दद्यात्~। सप्त\appbrk{}रा\appbrk{}त्रा\appbrk{}त्स\appbrk{}प्त\appbrk{}रा\appbrk{}त्रा\appbrk{}द्घ\appbrk{}टाम्}\ H\textsuperscript{ac}}}\skt{ सङ्क्रामयेत्~॥}\edtext{}{\lemma{\skt{उपहरेत्~। शेवालं}}\Afootnote{N H; \skt{उप\appbrk{}ह\appbrk{}रे\appbrk{}च्छै\appbrk{}वलं}\ A}}\pend

\pstart\versepar \ledsidenote{\passref{1.13.18}}\edtext{\skt{भवति}}{\lemma{\skt{भवति}}\Afootnote{N H; \emph{om.}\ A}}\skt{ }\edtext{\skt{चात्र}}{\lemma{\skt{चात्र}}\Afootnote{N H; \emph{om.}\ A}}\skt{~। }\pend
\pstart\versepar \skt{}\edtext{\skt{स्थूलमध्याः}}{\lemma{\skt{स्थूलमध्याः}}\Afootnote{A; \skt{स्थूलमध्या}\ N; \skt{श्लोकः~॥ स्थूलमध्या}\ H}}\skt{ }\edtext{\edtext{\edtext{\skt{परि\-क्लि\-ष्टा\-स्त\-न्व्य\-श्चा\-क्षे\-त्र\-जाश्च}}{\lemma{\skt{परिक्लिष्टास्}}\Afootnote{A; \skt{परिक्लिष्टा}\ N H}}}{\lemma{\skt{तन्व्यश्}}\Afootnote{N H; \skt{पृथ्व्यो}\ A}}}{\lemma{\skt{चाक्षेत्रजाश्च}}\Afootnote{N H; \skt{मन्दविचे}\ A}}\skt{ }\edtext{\skt{याः}}{\lemma{\skt{याः}}\Afootnote{N H; \skt{ष्टिताः}\ A}}\skt{~। }\pend
\pstart\versepar \skt{}\edtext{\skt{अग्रा\-हि\-ण्योऽ\-ल्प\-पा\-यिन्यः}}{\lemma{\skt{ऽल्पपायिन्यः}}\Afootnote{A H; \skt{ल्पपायिन्यः}\ N}}\skt{ सविषाश्च न }\edtext{\skt{पोषयेत्}}{\lemma{\skt{पोषयेत्}}\Afootnote{N H; \skt{पूजिताः}\ A}}\skt{~॥}\footnoteA{The vulgate edition notes "other readings" that correspond to those in the the Nepalese witnesses for \emph{tanvyaś cākṣetrajāś ca yāḥ} and \emph{poṣayet}.}\pend

\pstart \ledsidenote{\passref{1.13.19}}\skt{अथ }\edtext{\skt{जलौ\-का\-व\-से\-क\-सा\-ध्य\-व्याधिं}}{\lemma{\skt{जलौ\appbrk{}का\appbrk{}व\appbrk{}से\appbrk{}क\appbrk{}सा\appbrk{}ध्य\appbrk{}व्याधिं}}\Afootnote{N; \skt{जलौ\appbrk{}को\appbrk{}व\appbrk{}से\appbrk{}क\appbrk{}साध्य}\ A; \skt{जलौ\appbrk{}का\appbrk{}व\appbrk{}से\appbrk{}क\appbrk{}साध्यं व्याधिं}\ H}}\skt{ व्याधि\-त\-मु\-प\-वेश्य }\edtext{\skt{संवेश्य वा}}{\lemma{\skt{संवेश्य वा}}\Afootnote{A H; \skt{संवेष्य वा}\ N}}\skt{ विरूक्ष्य }\edtext{\skt{तमवकाशं}}{\lemma{\skt{तम्}}\Afootnote{N H; \skt{चास्य तम्}\ A}}\skt{ }\edtext{\skt{मृद्गो\-म\-य\-चू\-र्ण्णै\-र्य\-द्य\-त्स\-रुजं}}{\lemma{\skt{यद्यत्सरुजं}}\Afootnote{N; \skt{यद्यरुजः}\ A; \skt{यद्यत्सरुजा}\ H}}\skt{ }\edtext{\skt{स्याद्~। अथ}}{\lemma{\skt{स्याद्~। अथ}}\Afootnote{N H; \skt{स्यात्~। गृ}\ A}}\skt{ }\edtext{\skt{जलौकसः}}{\lemma{\skt{जलौकसः}}\Afootnote{N H; \skt{हीताश्च ताः}\ A}}\skt{ }\edtext{\skt{सर्ष\-प\-र\-ज\-नी\-प्र\-दि\-ग्ध\-गात्र्यः}}{\lemma{\skt{सर्ष\appbrk{}प\appbrk{}र\appbrk{}ज\appbrk{}नी\appbrk{}प्र\appbrk{}दि\appbrk{}ग्ध\appbrk{}गात्र्यः}}\Afootnote{N H; \skt{सर्ष\appbrk{}प\appbrk{}र\appbrk{}ज\appbrk{}नी\appbrk{}क\appbrk{}ल्को\appbrk{}द\appbrk{}क\appbrk{}प्र\appbrk{}दि\appbrk{}ग्ध\appbrk{}गात्रीः}\ A}}\skt{ }\edtext{\skt{सलि\-ल\-स\-र\-क\-म\-ध्य\-स\-ञ्चा\-रिणी}}{\lemma{\skt{सलि\appbrk{}ल\appbrk{}स\appbrk{}र\appbrk{}क\appbrk{}म\appbrk{}ध्य\appbrk{}स\appbrk{}ञ्चा\appbrk{}रिणी}}\Afootnote{N H; \skt{सलि\appbrk{}ल\appbrk{}स\appbrk{}र\appbrk{}क\appbrk{}मध्ये मुहूर्तस्थिता}\ A}}\skt{ }\edtext{\skt{विगतमलाः}}{\lemma{\skt{विगतमलाः}}\Afootnote{N H; \skt{विगतक्लमा}\ A}}\skt{ }\edtext{\skt{कृत्वा}}{\lemma{\skt{कृत्वा}}\Afootnote{N H; \skt{ज्ञात्वा}\ A}}\skt{ }\edtext{\edtext{\skt{रोगङ्ग्राहयेत्}}{\lemma{\skt{रोगङ्}}\Afootnote{\skt{ताभी रोगं}\ A; \skt{रोग}\ N; \skt{…~। रोगङ्}\ H}}\skt{~।}\footnoteA{Note the irregular sandhi of f. pl. \emph{-sañcāriṇī}.}\skt{ अथ}}{\lemma{\skt{ग्राहयेत्~। अथ}}\Afootnote{N H; \skt{ग्राहयेत्~। श्लक्ष्ण\appbrk{}शु\appbrk{}क्ला\appbrk{}र्द्र\appbrk{}पि\appbrk{}चु\appbrk{}प्रो\appbrk{}ता\appbrk{}व\appbrk{}च्छन्नां कृत्वा मुखम}\ A}}\skt{ }\edtext{\skt{न}}{\lemma{\skt{न}}\Afootnote{N H; \skt{पावृणुयात्; अ}\ A}}\skt{ }\edtext{\skt{गृह्णत्याः}}{\lemma{\skt{गृह्णत्याः}}\Afootnote{N H; \skt{गृह्णन्त्यै}\ A}}\skt{ }\edtext{\skt{क्षीरबिन्दुं}}{\lemma{\skt{क्षीरबिन्दुं}}\Afootnote{N H; \skt{क्षीरविन्दुं}\ A}}\skt{ }\edtext{\skt{शोणितबिन्दुं वा}}{\lemma{\skt{शोणितबिन्दुं वा}}\Afootnote{A H; \skt{शोणितबिदुम्वा}\ N}}\skt{ }\edtext{\skt{निदध्यात्~। शस्त्रपदानि वा}}{\lemma{\skt{निदध्यात्~। शस्त्रपदानि वा}}\Afootnote{N H; \skt{दद्यात्, शस्त्रपदानि वा}\ A}}\skt{ कुर्वीत~। }\edtext{\skt{अथैवमपि}}{\lemma{\skt{अथैवम्}}\Afootnote{H; \skt{यद्येवम्}\ A; \skt{अथेवम्}\ N}}\skt{ न }\edtext{\skt{गृह्णी\-या\-त\-न्या\-ङ्ग्रा\-ह\-येत्}}{\lemma{\skt{गृह्णीयातन्याङ्}}\Afootnote{N H; \skt{गृह्णीयात्तदाऽन्यां}\ A}}\skt{~।}\pend

\pstart \ledsidenote{\passref{1.13.20}}\edtext{\skt{यदा}}{\lemma{\skt{यदा}}\Afootnote{N H; \skt{यदा च}\ A}}\skt{ }\edtext{\edtext{\skt{निवि\-श\-तेऽ\-श्व\-खु\-र\-व\-दा\-न\-न\-ङ्कृ\-त्वो\-न्नाम्य}}{\lemma{\skt{ऽश्व\appbrk{}खु\appbrk{}र\appbrk{}व\appbrk{}दा\appbrk{}न\appbrk{}नङ्}}\Afootnote{A; \skt{श्वखु\appbrk{}र\appbrk{}व\appbrk{}दा\appbrk{}ननं}\ N H}}\skt{ च}}{\lemma{\skt{कृत्वोन्नाम्य च}}\Afootnote{H\textsuperscript{ac}; \skt{कृत्वोन्नम्य च}\ A; \skt{कृत्वोनाम्य च}\ N; \skt{कृत्वोन्नाम्य च च}\ H\textsuperscript{pc}}}\skt{ }\edtext{\edtext{\edtext{\skt{स्कन्ध\-मे\-व\-ञ्जा\-नी\-या\-द्गृ\-ह्णा\-तीति}}{\lemma{\skt{स्कन्धम्}}\Afootnote{A H\textsuperscript{pc}; \skt{स्कन्दम्}\ N; \skt{कन्धं}\ H\textsuperscript{ac}}}}{\lemma{\skt{एवञ्}}\Afootnote{N H\textsuperscript{ac}; \skt{तदा}\ A; \skt{एव}\ H\textsuperscript{pc}}}}{\lemma{\skt{जानी\appbrk{}या\appbrk{}द्गृ\appbrk{}ह्णा\appbrk{}तीति}}\Afootnote{A N H\textsuperscript{pc}; \skt{जानी\appbrk{}या\appbrk{}त्द्गृ\appbrk{}ह्णा\appbrk{}तीति}\ H\textsuperscript{ac}}}\skt{~। }\edtext{\edtext{\skt{अथै\-ना\-मा\-र्द्र\-प्लो\-ता\-व\-च्छ\-न्ना\-ङ्कृत्वा}}{\lemma{\skt{अथैनाम्}}\Afootnote{N H; \skt{गृह्णन्तीं}\ A}}}{\lemma{\skt{आर्द्र\appbrk{}प्लो\appbrk{}ता\appbrk{}व\appbrk{}च्छ\appbrk{}न्नाङ्}}\Afootnote{\skt{चार्द्र\appbrk{}व\appbrk{}स्त्रा\appbrk{}व\appbrk{}च्छन्नां}\ A; \skt{आर्द्र\appbrk{}प्लो\appbrk{}ता\appbrk{}व\appbrk{}च्छन्न}\ N; \skt{आर्द्द्र\appbrk{}प्लो\appbrk{}ता\appbrk{}व\appbrk{}च्छ\appbrk{}न्नङ्}\ H}}\skt{ धारयेत्~॥}\pend

\pstart \ledsidenote{\passref{1.13.21}}\edtext{\skt{अथ}}{\lemma{\skt{अथ}}\Afootnote{N H; \emph{om.}\ A}}\skt{ दंशे }\edtext{\edtext{\skt{तोद\-क\-ण्डू\-प्रा\-दु\-र्भा\-वा\-ज्जा\-नी\-या\-च्छु\-द्ध\-मा\-द\-दा\-तीति}}{\lemma{\skt{तोद\appbrk{}क\appbrk{}ण्डू\appbrk{}प्रा\appbrk{}दु\appbrk{}र्भा\appbrk{}वा\appbrk{}ज्जा\appbrk{}नी\appbrk{}या\appbrk{}च्छु\appbrk{}द्धम्}}\Afootnote{\skt{तोद\appbrk{}क\appbrk{}ण्डु\appbrk{}प्रा\appbrk{}दु\appbrk{}र्भा\appbrk{}वै\appbrk{}र्जा\appbrk{}नी\appbrk{}या\appbrk{}च्छु\appbrk{}द्ध\appbrk{}मि\appbrk{}य\appbrk{}मा\appbrk{}दत्त इति; शुद्धम्}\ A; \skt{तोद\appbrk{}क\appbrk{}ण्डू\appbrk{}प्रा\appbrk{}दु\appbrk{}र्भावो जानीयाच्छुद्धम्}\ N H}}\skt{~। तामपनयेत्}}{\lemma{\skt{आददातीति~। ताम्}}\Afootnote{N H; \skt{आददानाम्}\ A}}\skt{~। अथ शोणितगन्धेन न }\edtext{\skt{मुञ्चेत्~। मुखमस्याः}}{\lemma{\skt{मुञ्चेत्~। मुखम्}}\Afootnote{N H; \skt{मुञ्चेन्मुखम्}\ A}}\skt{ सैन्ध\-व\-चू\-र्णे\-ना\-व\-कि\-रेत्~।}\pend

\pstart \ledsidenote{\passref{1.13.22}}\edtext{\skt{अथैनां}}{\lemma{\skt{अथैनां}}\Afootnote{H; \skt{अथ प}\ A; \emph{om.}\ K; \skt{अथैना}\ N}}\skt{ }\edtext{\edtext{\skt{शालि\-त\-ण्डु\-ल\-का\-ण्ड\-न\-प्र\-लि\-प्ता\-न्तै\-ल\-ल\-व\-णा\-भ्य\-क्त\-मुखीं}}{\lemma{\skt{शालि\appbrk{}त\appbrk{}ण्डु\appbrk{}ल\appbrk{}का\appbrk{}ण्ड\appbrk{}न\appbrk{}प्र\appbrk{}लि\appbrk{}प्तान्}}\Afootnote{H; \skt{तितां तण्डु\appbrk{}ल\appbrk{}क\appbrk{}ण्ड\appbrk{}न\appbrk{}प्र\appbrk{}दि\appbrk{}ग्ध\appbrk{}गात्रीं}\ A; \emph{om.}\ K; \skt{शालि\appbrk{}त\appbrk{}ण्डु\appbrk{}ल\appbrk{}का\appbrk{}ण्ड\appbrk{}न\appbrk{}प्र\appbrk{}लि\appbrk{}प्तान्तौ}\ N}}}{\lemma{\skt{तैल\appbrk{}ल\appbrk{}व\appbrk{}णा\appbrk{}भ्य\appbrk{}क्त\appbrk{}मुखीं}}\Afootnote{A; \emph{om.}\ K; \skt{लल\appbrk{}व\appbrk{}णा\appbrk{}भ्य\appbrk{}क्त\appbrk{}मुखीं}\ N; \skt{तैल\appbrk{}ल\appbrk{}व\appbrk{}णो\appbrk{}भ्य\appbrk{}क्त\appbrk{}मुखीं}\ H}}\skt{ }\edtext{\edtext{\skt{वाम\-ह\-स्त\-गृ\-ही\-त\-पु\-च्छा\-न्द\-क्षि\-ण\-ह\-स्ता\-ङ्गु\-ष्ठा\-ङ्गु\-लीभ्यां}}{\lemma{\skt{वाम\appbrk{}ह\appbrk{}स्त\appbrk{}गृ\appbrk{}ही\appbrk{}त\appbrk{}पु\appbrk{}च्छान्}}\Afootnote{N H; \skt{वाम\appbrk{}ह\appbrk{}स्ता\appbrk{}ङ्गु\appbrk{}ष्ठा\appbrk{}ङ्गु\appbrk{}लीभ्यां गृहीतपुच्छां}\ A; \emph{om.}\ K}}}{\lemma{\skt{दक्षि\appbrk{}ण\appbrk{}ह\appbrk{}स्ता\appbrk{}ङ्गु\appbrk{}ष्ठा\appbrk{}ङ्गु\appbrk{}लीभ्यां}}\Afootnote{N H; \skt{दक्षि\appbrk{}ण\appbrk{}ह\appbrk{}स्ता\appbrk{}ङ्गु\appbrk{}ष्ठा\appbrk{}ङ्गु\appbrk{}लिभ्यां}\ A; \emph{om.}\ K}}\skt{ }\edtext{\skt{शनैः}}{\lemma{\skt{शनैः}}\Afootnote{A N H; \emph{om.}\ K}}\skt{ }\edtext{\edtext{\skt{शनै\-र\-नु\-लो\-म\-मा\-र्ज्ज\-यन्ना}}{\lemma{\skt{शनैर्}}\Afootnote{A N H; \emph{om.}\ K}}}{\lemma{\skt{अनु\appbrk{}लो\appbrk{}म\appbrk{}मा\appbrk{}र्ज्ज\appbrk{}यन्ना}}\Afootnote{N H; \skt{अनु\appbrk{}लो\appbrk{}म\appbrk{}म\appbrk{}नु\appbrk{}मा\appbrk{}र्ज\appbrk{}येदा}\ A; \emph{om.}\ K}}\skt{ }\edtext{\skt{मुखा\-द्वा\-म\-ये\-द्या\-व\-त्स\-म्य\-ग्वा\-न्तेति}}{\lemma{\skt{मुखा\appbrk{}द्वा\appbrk{}म\appbrk{}ये\appbrk{}द्या\appbrk{}व\appbrk{}त्स\appbrk{}म्य\appbrk{}ग्वा\appbrk{}न्तेति}}\Afootnote{H; \skt{मुखात्, वाम\appbrk{}ये\appbrk{}त्ता\appbrk{}व\appbrk{}द्या\appbrk{}व\appbrk{}त्स\appbrk{}म्य\appbrk{}ग्वा\appbrk{}न्त\appbrk{}लि\appbrk{}ङ्गा\appbrk{}नीति}\ A; \skt{वत्सम्यग्वान्तेति}\ K; \skt{मुखाद्वामये वत्सम्यग्वान्तेति}\ N}}\skt{~। सम्यग्वान्ता सलिलसरके न्यस्ता भोक्तुकामा }\edtext{\skt{सती}}{\lemma{\skt{सती}}\Afootnote{A K N H\textsuperscript{pc}; \skt{मती}\ H\textsuperscript{ac}}}\skt{ चरेद्~। या }\edtext{\skt{सीदति न}}{\lemma{\skt{सीदति न}}\Afootnote{K H; \skt{सीदती न}\ A; \skt{शीति न}\ N}}\skt{ चेष्टते सा }\edtext{\skt{दुर्वान्ता}}{\lemma{\skt{दुर्वान्ता}}\Afootnote{\skt{दुर्वान्ता, तां}\ A; \skt{दुर्वान्तां}\ K H; \skt{दुर्त्वान्ता}\ N}}\skt{ पुनः }\edtext{\skt{सम्यग्वामयेत्~। }\edtext{\skt{दुर्वान्तायास्तु}}{\lemma{\skt{तु}}\Afootnote{K N H; \skt{व्याधिरसाध्य}\ A}}}{\lemma{\skt{सम्यग्वामयेत्~। दुर्वान्तायास्}}\Afootnote{H; \skt{सम्यग्वामयेत्~। दुर्वान्ताया}\ A; \skt{सम्यक्वामयेत्~। दुर्वान्तायास्}\ K N}}\skt{ }\edtext{\skt{इन्द्रमदो}}{\lemma{\skt{इन्द्रमदो}}\Afootnote{A; \skt{इन्द्रपदो}\ K H; \skt{इन्द्रापदो}\ N}}\skt{ नाम }\edtext{\edtext{\skt{व्याधिरसाध्यो}}{\lemma{\skt{व्याधिर्}}\Afootnote{K N H; \emph{om.}\ A}}}{\lemma{\skt{असाध्यो}}\Afootnote{K N H; \emph{om.}\ A}}\skt{ }\edtext{\skt{भवति}}{\lemma{\skt{भवति}}\Afootnote{K N H; \skt{भवति~। अथ सुवान्तां पूर्व\appbrk{}व\appbrk{}त्स\appbrk{}न्नि\appbrk{}द\appbrk{}ध्यात्}\ A}}\skt{~।}\pend

\pstart\versepar \ledsidenote{\passref{1.13.22.1}}\skt{अप्रहृष्टशिरः }\edtext{\skt{पाति}}{\lemma{\skt{पाति}}\Afootnote{N; \skt{पात्य}\ K H}}\skt{ कायेनोद्वेष्टते सकृत्~। }\pend
\pstart\versepar \skt{या चोष्णं कुरुते }\edtext{\skt{तोयं}}{\lemma{\skt{तोयं}}\Afootnote{K N\textsuperscript{pc} H; \skt{तोष्टयन्}\ N\textsuperscript{ac}}}\skt{ }\edtext{\edtext{\skt{तस्यामिन्द्रमदः}}{\lemma{\skt{तस्याम्}}\Afootnote{K H; \skt{तस्या}\ N}}}{\lemma{\skt{इन्द्रमदः}}\Afootnote{\skt{इदमदः}\ K; \skt{निद्रमदः}\ N; \skt{इन्द्रपदः}\ H}}\skt{ स्मृतः~॥}\footnoteA{The vulgate, third edition, reports the reading of this verse in a Nepalese witness, probably H: \emph{aprahṛṣṭaśiraḥpādakāyenodveṣṭate sakṛt | yā coṣṇaṃ kurute tāpaṃ tasyām indramadaḥ smṛtaḥ} (Ācārya 1938: 58, n. 4).}\pend

\pstart \ledsidenote{\passref{1.13.22b}}\skt{अथैनां पूर्व\-व\-त्स\-न्नि\-द\-ध्यात्~।}\pend

\pstart \ledsidenote{\passref{1.13.23}}\edtext{\skt{शोणितस्य च}}{\lemma{\skt{शोणितस्य च}}\Afootnote{K N H; \skt{शोणितस्य}\ A}}\skt{ }\edtext{\edtext{\skt{योगा\-यो\-ग\-म\-वेक्ष्य}}{\lemma{\skt{योगायोगम्}}\Afootnote{K N H; \skt{योगायोगा}\ A}}}{\lemma{\skt{अवेक्ष्य}}\Afootnote{K N H; \skt{नवेक्ष्य}\ A}}\skt{ }\edtext{\skt{जलौकामुखं}}{\lemma{\skt{जलौकामुखं}}\Afootnote{K N H; \skt{शत\appbrk{}धौ\appbrk{}त\appbrk{}घृ\appbrk{}ता\appbrk{}भ्यङ्गः, तत्पिचुधारणं वा; जलौकोव्रणा}\ A}}\skt{ }\edtext{\skt{मधु\-ना\-व\-घ\-ट्ट\-येत्~। बध्नीत वा}}{\lemma{\skt{मधु\appbrk{}ना\appbrk{}व\appbrk{}घ\appbrk{}ट्ट\appbrk{}येत्~। बध्नीत वा}}\Afootnote{K N; \skt{न्मधु\appbrk{}नाऽ\appbrk{}व\appbrk{}घ\appbrk{}ट्ट\appbrk{}येत्, शीताभिरद्भिः परि\appbrk{}षे\appbrk{}च\appbrk{}ये\appbrk{}द्ब\appbrk{}ध्नीत वा}\ A; \skt{मधु\appbrk{}ना\appbrk{}व\appbrk{}घ\appbrk{}ट्ट\appbrk{}येत्…। बध्नीत वा}\ H}}\skt{ }\edtext{\skt{कषा\-य\-म\-धु\-र\-स्नि\-ग्ध\-शी\-तैश्च}}{\lemma{\skt{कषा\appbrk{}य\appbrk{}म\appbrk{}धु\appbrk{}र\appbrk{}स्नि\appbrk{}ग्ध\appbrk{}शी\appbrk{}तैश्च}}\Afootnote{A H; \skt{कशा\appbrk{}य\appbrk{}म\appbrk{}धु\appbrk{}र\appbrk{}स्नि\appbrk{}ग्ध\appbrk{}शी\appbrk{}तैश्च}\ K N}}\skt{ प्रदेहैः प्रदिह्यादिति~॥}\pend

\pstart\versepar \ledsidenote{\passref{1.13.23.1}}\edtext{\skt{भवतश्चात्र}}{\lemma{\skt{भवतश्चात्र}}\Afootnote{\emph{om.}\ K N; \skt{भवतश्चात्र}\ H}}\skt{~॥ }\pend
\pstart\versepar \skt{पीतमात्रे जलौ\-का\-भि\-र्घृ\-तेन परिषेचयेत्~। }\pend
\pstart\versepar \skt{शोणि\-त\-स्था\-प\-नी\-यैश्च शोणितं परिषेचयेत्}\pend

\pstart\versepar \ledsidenote{\passref{1.13.24}}\skt{क्षेत्राणि }\edtext{\skt{ग्रहणञ्चापि}}{\lemma{\skt{चापि}}\Afootnote{K N H; \skt{जातीः}\ A}}\skt{ पोषणं सावचारणम्~। }\pend
\pstart\versepar \skt{}\edtext{\skt{जानीयाद्यो}}{\lemma{\skt{जानीयाद्यो}}\Afootnote{K N H; \emph{om.}\ A}}\skt{ }\edtext{\skt{जलौकानां}}{\lemma{\skt{जलौकानां}}\Afootnote{K N H; \skt{जलौकसां च यो}\ A}}\skt{ }\edtext{\skt{स}}{\lemma{\skt{स}}\Afootnote{K N H; \skt{वेत्ति तत्साध्यान्स ज}\ A}}\skt{ }\edtext{\skt{राज्ञः}}{\lemma{\skt{राज्ञः}}\Afootnote{K N H; \skt{येद्गदान्}\ A}}\skt{ }\edtext{\edtext{\skt{कर्त्तुमर्हति}}{\lemma{\skt{कर्त्तुम्}}\Afootnote{K N H; \emph{om.}\ A}}}{\lemma{\skt{अर्हति}}\Afootnote{K N H; \emph{om.}\ A}}\skt{~॥}\pend

\pstart \ledsidenote{\passref{1.13.trailer}}\edtext{\skt{१३}}{\lemma{\skt{१३}}\Afootnote{\skt{इति सुश्रु\appbrk{}त\appbrk{}सं\appbrk{}हि\appbrk{}तायां सूत्रस्थाने जलौ\appbrk{}का\appbrk{}व\appbrk{}चा\appbrk{}र\appbrk{}णीयो नाम त्रयोदशोऽध्यायः}\ A; \skt{१३}\ K; \skt{१० ३}\ N; \skt{१३~॥}\ H}}\pend

\endnumbering
\beginnumbering
\eledsection*{Adhyāya 14}
\pstart \ledsidenote{\passref{1.14.1}}\edtext{\skt{अथातः}}{\lemma{\skt{अथातः}}\Afootnote{A K H; \skt{अथात}\ N}}\skt{ शोणि\-त\-व\-र्ण्ण\-नीयं }\edtext{\skt{व्याख्यास्यामः}}{\lemma{\skt{व्याख्यास्यामः}}\Afootnote{H; \skt{अध्यायं व्याख्यास्यामः}\ A; \emph{om.}\ K; \skt{अध्यमाया व्याख्यास्यामः}\ N\textsuperscript{ac}; \skt{अध्या\appbrk{}या\appbrk{}म्व्या\appbrk{}ख्या\appbrk{}स्यामः}\ N\textsuperscript{pc}}}\skt{~॥}\pend

\pstart \ledsidenote{\passref{1.14.3}}\edtext{\skt{पाञ्चभौतिकस्य}}{\lemma{\skt{पाञ्चभौतिकस्य}}\Afootnote{K N H; \skt{तत्र पाञ्चभौतिकस्य}\ A}}\skt{ }\edtext{\skt{चतु\-र्वि\-ध\-स्या\-हा\-रस्य}}{\lemma{\skt{चतु\appbrk{}र्वि\appbrk{}ध\appbrk{}स्या\appbrk{}हा\appbrk{}रस्य}}\Afootnote{K N H; \skt{चतुर्विधस्य}\ A}}\skt{ }\edtext{\skt{षड्रसोपेतस्य}}{\lemma{\skt{षड्रसोपेतस्य}}\Afootnote{K N H; \skt{षड्रसस्य}\ A}}\skt{ द्विवि\-ध\-वी\-र्य\-स्या\-ष्ट\-वि\-ध\-वी\-र्यस्य वा }\edtext{\skt{अने\-क\-गु\-णो\-प\-यु\-क्त\-स्या\-हा\-रस्य}}{\lemma{\skt{अने\appbrk{}क\appbrk{}गु\appbrk{}णो\appbrk{}प\appbrk{}यु\appbrk{}क्त\appbrk{}स्या\appbrk{}हा\appbrk{}रस्य}}\Afootnote{K N H; \skt{नेक\appbrk{}गु\appbrk{}ण\appbrk{}स्यो\appbrk{}प\appbrk{}यु\appbrk{}क्त\appbrk{}स्या\appbrk{}हा\appbrk{}रस्य}\ A}}\skt{ }\edtext{\skt{सम्य\-क्प\-रि\-ण\-तस्य }\edtext{\skt{यस्ते\-जो\-गु\-ण\-भूतः}}{\lemma{\skt{तेजोगुणभूतः}}\Afootnote{N H; \skt{तेजोभूतः}\ A K\textsuperscript{ac}; \skt{तेजोणभूतः}\ K\textsuperscript{pc}}}}{\lemma{\skt{सम्य\appbrk{}क्प\appbrk{}रि\appbrk{}ण\appbrk{}तस्य यस्}}\Afootnote{A K H; \skt{सम्य\appbrk{}क्ष\appbrk{}रि\appbrk{}ण\appbrk{}तस्य यस्}\ N}}\skt{ सारः }\edtext{\skt{परमसूक्ष्मः स}}{\lemma{\skt{परमसूक्ष्मः स}}\Afootnote{K N; \skt{परमसूक्ष्मः स {\symfont `}}\ A; \skt{परमशूक्ष्मः स}\ H}}\skt{ रस इत्युच्यते~। तस्य हृदयं }\edtext{\skt{स्थानं~। स}}{\lemma{\skt{स्थानं~। स}}\Afootnote{A K N H\textsuperscript{pc}; \skt{स्थानं~। म}\ H\textsuperscript{ac}}}\skt{ }\edtext{\edtext{\skt{हृद\-या\-च्च\-तु\-र्विं\-श\-ति\-र्ध\-म\-नी\-र\-नु\-प्र\-वि\-श्यो\-र्ध्वगा}}{\lemma{\skt{हृद\appbrk{}या\appbrk{}च्च\appbrk{}तु\appbrk{}र्विं\appbrk{}श\appbrk{}तिर्}}\Afootnote{K N H; \skt{हृद\appbrk{}या\appbrk{}च्च\appbrk{}तु\appbrk{}र्विं\appbrk{}शति}\ A}}}{\lemma{\skt{अनु\appbrk{}प्र\appbrk{}वि\appbrk{}श्यो\appbrk{}र्ध्वगा}}\Afootnote{K N H\textsuperscript{pc}; \skt{अनु\appbrk{}प्र\appbrk{}वि\appbrk{}श्यो\appbrk{}र्ध्वगा}\ A; \skt{अनुप्रविर्द्ध्वगा}\ H\textsuperscript{ac}}}\skt{ दश }\edtext{\skt{दश}}{\lemma{\skt{दश}}\Afootnote{K N H; \skt{द}\ A}}\skt{ }\edtext{\edtext{\skt{चाधो\-गा\-मि\-न्य\-श्च\-त\-स्र\-स्ति\-र्यग्गाः}}{\lemma{\skt{चाधोगामिन्यश्}}\Afootnote{K N H; \skt{शाधोगामिन्यश्}\ A}}}{\lemma{\skt{चतस्रस्}}\Afootnote{K N H; \skt{चतस्रश्च}\ A}}\skt{ कृत्स्नं }\edtext{\skt{शरी\-र\-म\-ह\-र\-ह\-स्त\-र्प\-यति}}{\lemma{\skt{अहरहस्}}\Afootnote{A K H; \skt{अहरहस्}\ N}}\skt{ }\edtext{\skt{जीवयति}}{\lemma{\skt{जीवयति}}\Afootnote{N H; \skt{वर्धयति}\ A; \skt{जीवति}\ K\textsuperscript{ac}; \skt{जीपवति}\ K\textsuperscript{pc}}}\skt{ }\edtext{\skt{यापयति}}{\lemma{\skt{यापयति}}\Afootnote{K N H; \skt{धारयति}\ A}}\skt{ }\edtext{\skt{वर्धयति}}{\lemma{\skt{वर्धयति}}\Afootnote{K N H; \skt{यापयति}\ A}}\skt{ }\edtext{\skt{चादृ\-ष्ट\-है\-तु\-केन}}{\lemma{\skt{चादृ\appbrk{}ष्ट\appbrk{}है\appbrk{}तु\appbrk{}केन}}\Afootnote{H; \skt{चादृ\appbrk{}ष्ट\appbrk{}हे\appbrk{}तु\appbrk{}केन}\ A N K}}\skt{ कर्मणा~। तस्य }\edtext{\edtext{\edtext{\skt{शरी\-र\-म\-नु\-स\-र\-तोऽ\-नु\-मा\-ना\-द्ग\-ति\-रु\-प\-ल\-क्ष\-यि\-तव्या}}{\lemma{\skt{अनुसरतो}}\Afootnote{A K N; \skt{अनुसर}\ H}}}{\lemma{\skt{ऽनु\appbrk{}मा\appbrk{}ना\appbrk{}द्ग\appbrk{}तिर्}}\Afootnote{A; \skt{नुमानाद्गतिर्}\ K; \skt{नुमानद्गतिर्}\ N; \skt{तोनु\appbrk{}मा\appbrk{}ना\appbrk{}ङ्ग\appbrk{}तिर्}\ H}}}{\lemma{\skt{उप\appbrk{}ल\appbrk{}क्ष\appbrk{}यि\appbrk{}तव्या}}\Afootnote{A H; \skt{उपलक्षयित्वा}\ K N}}\skt{ }\edtext{\skt{क्षय\-वृ\-द्धि\-हे\-तुकी}}{\lemma{\skt{क्षय\appbrk{}वृ\appbrk{}द्धि\appbrk{}हे\appbrk{}तुकी}}\Afootnote{H; \skt{क्षयवृद्धिवै}\ A; \skt{क्षय\appbrk{}वृ\appbrk{}द्धी\appbrk{}हे\appbrk{}तुकी}\ K\textsuperscript{ac} K\textsuperscript{pc} N}}\skt{~॥ }\edtext{\skt{तस्मिं}}{\lemma{\skt{तस्मिं}}\Afootnote{K\textsuperscript{ac} N; \skt{कृतैः~। तस्मि}\ A; \skt{गुण तस्मिं}\ K\textsuperscript{pc}; \skt{तस्मि}\ H}}\skt{ }\edtext{\skt{सर्व\-श\-री\-रा\-व\-य\-व\-दो\-ष\-धा\-तु\-म\-ला\-श\-या\-नु\-सा\-रिणि}}{\lemma{\skt{सर्व\appbrk{}श\appbrk{}री\appbrk{}रा\appbrk{}व\appbrk{}य\appbrk{}व\appbrk{}दो\appbrk{}ष\appbrk{}धा\appbrk{}तु\appbrk{}म\appbrk{}ला\appbrk{}श\appbrk{}या\appbrk{}नु\appbrk{}सा\appbrk{}रिणि}}\Afootnote{K; \skt{न्सर्व\appbrk{}श\appbrk{}री\appbrk{}रा\appbrk{}व\appbrk{}य\appbrk{}व\appbrk{}दो\appbrk{}ष\appbrk{}धा\appbrk{}तु\appbrk{}म\appbrk{}ला\appbrk{}श\appbrk{}या\appbrk{}नु\appbrk{}सा\appbrk{}रिणि}\ A; \skt{शरी\appbrk{}रा\appbrk{}व\appbrk{}य\appbrk{}व\appbrk{}दो\appbrk{}ष\appbrk{}धा\appbrk{}तु\appbrk{}म\appbrk{}ला\appbrk{}श\appbrk{}या\appbrk{}नु\appbrk{}सा\appbrk{}रिणि}\ N; \skt{न्सर्व्व\appbrk{}श\appbrk{}री\appbrk{}रा\appbrk{}व\appbrk{}य\appbrk{}व\appbrk{}दो\appbrk{}ष\appbrk{}धा\appbrk{}तु\appbrk{}म\appbrk{}ला\appbrk{}स\appbrk{}या\appbrk{}नु\appbrk{}सा\appbrk{}रिणि}\ H}}\skt{ रसे जिज्ञासा किमयं }\edtext{\skt{सौम्यस्तैजस}}{\lemma{\skt{तैजस}}\Afootnote{A N; \skt{तेजस}\ K H}}\skt{ इति~। }\edtext{\skt{स}}{\lemma{\skt{स}}\Afootnote{K N H; \skt{अत्रोच्यते स}\ A}}\skt{ खलु }\edtext{\skt{द्रव\-त्वा\-द\-नु\-स\-रणे}}{\lemma{\skt{द्रव\appbrk{}त्वा\appbrk{}द\appbrk{}नु\appbrk{}स\appbrk{}रणे}}\Afootnote{K N; \skt{द्रवानुसारी}\ A; \skt{द्रवा\appbrk{}द\appbrk{}नु\appbrk{}स\appbrk{}रणे}\ H}}\skt{ स्नेह\-न\-जी\-व\-न\-त\-र्प\-ण\-धा\-र\-णा\-दि\-भि\-र्वि\-शेषैः सौम्य इत्यवगम्यते~।}\pend

\pstart \ledsidenote{\passref{1.14.4}}\skt{स खल्वाप्यो रसो यकृत्प्लीहानौ प्राप्य रागमुपैति~॥}\pend

\pstart\versepar \ledsidenote{\passref{1.14.5}}\skt{भवन्ति चात्र~॥ }\pend
\pstart\versepar \skt{रञ्जितास्तेजसा }\edtext{\skt{त्वापः}}{\lemma{\skt{त्वापः}}\Afootnote{A K N; \skt{त्वाप्पः}\ H}}\skt{ शरीरस्थेन देहिनां~। }\pend
\pstart\versepar \skt{अव्यापन्नाः }\edtext{\skt{प्रसन्नेन}}{\lemma{\skt{प्रसन्नेन}}\Afootnote{A K N; \skt{प्रसन्नेन…}\ H}}\skt{ }\edtext{\skt{रक्तमित्येव}}{\lemma{\skt{एव}}\Afootnote{K N H; \skt{अभिधीय}\ A}}\skt{ }\edtext{\skt{तद्विदुः}}{\lemma{\skt{तद्विदुः}}\Afootnote{K N H; \skt{ते}\ A}}\skt{~॥}\edtext{}{\lemma{\skt{भवन्ति चात्र~॥ रञ्जितास्}}\Afootnote{\skt{श्लोकौ चात्र भवतः रञ्जितास्}\ A; \skt{रञ्जितास्}\ K H N}}\pend

\pstart\versepar \ledsidenote{\passref{1.14.6}}\skt{रसादेव }\edtext{\skt{स्त्रियां}}{\lemma{\skt{स्त्रियां}}\Afootnote{\skt{स्त्रिया}\ A K H; \skt{स्त्रियो}\ N}}\skt{ }\edtext{\skt{रक्तमृतुसंज्ञं}}{\lemma{\skt{ऋतुसंज्ञं}}\Afootnote{K N H; \skt{रजःसञ्ज्ञं}\ A}}\skt{ प्रवर्तते }\pend
\pstart\versepar \skt{}\edtext{\skt{द्वादशाद्वर्धते}}{\lemma{\skt{द्वादशाद्वर्धते}}\Afootnote{K N H; \skt{तद्व\appbrk{}र्षा\appbrk{}द्द्वा\appbrk{}द\appbrk{}शादू}\ A}}\skt{ }\edtext{\skt{वर्षाद्याति}}{\lemma{\skt{वर्षाद्याति}}\Afootnote{K N; \skt{र्ध्वं याति}\ A; \skt{वर्शाद्याति}\ H}}\skt{ पञ्चाशतः क्षयं~।}\pend

\pstart \ledsidenote{\passref{1.14.7}}\skt{आर्त्तवं शोणितं }\edtext{\skt{त्वाग्नेयमाहुः}}{\lemma{\skt{आहुः}}\Afootnote{K N H; \emph{om.}\ A}}\skt{~॥ }\edtext{\skt{अग्नी\-षो\-मी\-य\-त्वा\-द्ग\-र्भस्य}}{\lemma{\skt{अग्नी\appbrk{}षो\appbrk{}मी\appbrk{}य\appbrk{}त्वा\appbrk{}द्ग\appbrk{}र्भस्य}}\Afootnote{A K N; \skt{अग्नी\appbrk{}षो\appbrk{}मी\appbrk{}य\appbrk{}त्वा\appbrk{}द्ग\appbrk{}र्भ्भस्य}\ H}}\skt{~।}\pend

\pstart \ledsidenote{\passref{1.14.8}}\edtext{\edtext{\skt{पाञ्च\-भौ\-ति\-क\-त्व\-म\-परे}}{\lemma{\skt{पाञ्च\appbrk{}भौ\appbrk{}ति\appbrk{}क\appbrk{}त्वम्}}\Afootnote{K N H; \skt{पाञ्चभौतिकं त्व}\ A}}}{\lemma{\skt{अपरे}}\Afootnote{K N H; \skt{परे}\ A}}\skt{ }\edtext{\skt{जीवं}}{\lemma{\skt{जीवं}}\Afootnote{K N H; \skt{जीव}\ A}}\skt{ }\edtext{\skt{रक्त\-मा\-हु\-रा\-चार्याः}}{\lemma{\skt{आहुर्}}\Afootnote{A K H; \emph{om.}\ N}}\skt{~॥}\pend

\pstart\versepar \ledsidenote{\passref{1.14.9}}\skt{भवन्ति चात्र~॥ }\pend
\pstart\versepar \skt{विस्रता द्रवता रागः }\edtext{\skt{स्पन्दनन्तनुता}}{\lemma{\skt{तनुता}}\Afootnote{K N H; \skt{लघुता}\ A}}\skt{ तथा~। }\pend
\pstart\versepar \skt{}\edtext{\edtext{\skt{पृथि\-व्या\-दि\-गु\-णा\-स्त्वेते}}{\lemma{\skt{पृथि\appbrk{}व्या\appbrk{}दि\appbrk{}गु\appbrk{}णास्}}\Afootnote{K N H; \skt{भूम्यादीनां गुणा}\ A}}}{\lemma{\skt{त्वेते}}\Afootnote{N H; \skt{ह्येते}\ A; \skt{चेते}\ K}}\skt{ दृश्यन्ते }\edtext{\skt{शोणिते}}{\lemma{\skt{शोणिते}}\Afootnote{K N H; \skt{चात्र शोणिते}\ A}}\skt{ }\edtext{\skt{यतः}}{\lemma{\skt{यतः}}\Afootnote{K N; \emph{om.}\ A; \skt{यतः~॥ …}\ H}}\skt{~।}\edtext{}{\lemma{\skt{भवन्ति चात्र~॥ विस्रता}}\Afootnote{\skt{विस्रता}\ A K N; \skt{भवति चात्र~॥ विश्रता}\ H}}\pend

\pstart\versepar \ledsidenote{\passref{1.14.10}}\skt{रसाद्रक्तं ततो मांसं }\edtext{\skt{मांसान्मेदः}}{\lemma{\skt{मांसान्मेदः}}\Afootnote{A K; \skt{मांसान्मेद}\ N; \skt{मांसात्मेदः}\ H}}\skt{ }\edtext{\skt{प्रवर्त्तते}}{\lemma{\skt{प्रवर्त्तते}}\Afootnote{K H; \skt{प्र}\ A N}}\skt{ }\pend
\pstart\versepar \skt{}\edtext{\skt{मेदसोऽस्थि}}{\lemma{\skt{मेदसो}}\Afootnote{K H\textsuperscript{ac}; \skt{जायते~। मेदसो}\ A; \skt{वर्त्तते मेदसो}\ N; \skt{पञ्चशत्रानि~। मेदसो}\ H\textsuperscript{pc}}}\skt{ ततो मज्जा }\edtext{\skt{मज्जा}}{\lemma{\skt{मज्जा}}\Afootnote{K N H; \skt{मज्ज्ञः}\ A}}\skt{ }\edtext{\skt{शुक्रं}}{\lemma{\skt{शुक्रं}}\Afootnote{A K N; \skt{च्छुक्रन्}\ H}}\skt{ }\edtext{\skt{ततः}}{\lemma{\skt{ततः}}\Afootnote{K N H; \skt{तु}\ A}}\skt{ }\edtext{\skt{प्रजाः}}{\lemma{\skt{प्रजाः}}\Afootnote{K N H; \skt{जायते}\ A}}\skt{~॥}\pend

\pstart \ledsidenote{\passref{1.14.11}}\edtext{\edtext{\skt{तत्रै\-षा\-न्धा\-तू\-ना\-म\-न्न\-पा\-न\-रसः}}{\lemma{\skt{तत्रैषान्}}\Afootnote{K N H; \skt{तत्रैतेषां}\ A}}}{\lemma{\skt{अन्नपानरसः}}\Afootnote{A K N; \skt{अन्नपानरसः…}\ H}}\skt{ }\edtext{\skt{प्रीणयिता}}{\lemma{\skt{प्रीणयिता}}\Afootnote{A H; \skt{प्रीणायिता}\ K N}}\skt{ }\edtext{\skt{भवति}}{\lemma{\skt{भवति}}\Afootnote{K N H; \emph{om.}\ A}}\skt{~॥}\pend

\pstart\versepar \ledsidenote{\passref{1.14.12}}\skt{भवति चात्र~॥ }\pend
\pstart\versepar \skt{रसजं पुरुषं }\edtext{\skt{विद्याद्रसं}}{\lemma{\skt{विद्याद्रसं}}\Afootnote{A K H; \skt{विद्या रसं}\ N}}\skt{ }\edtext{\skt{रक्षेत}}{\lemma{\skt{रक्षेत}}\Afootnote{K N H\textsuperscript{pc}; \skt{रक्षेत्प्र}\ A; \skt{रक्ते तक्षेत}\ H\textsuperscript{ac}}}\skt{ यत्नतः~। }\edtext{\skt{अन्न\-पा\-न\-प्र\-यो\-गेण}}{\lemma{\skt{अन्न\appbrk{}पा\appbrk{}न\appbrk{}प्र\appbrk{}यो\appbrk{}गेण}}\Afootnote{N; \skt{अन्नात्पानाञ्च मतिमा}\ A; \skt{अन्न\appbrk{}पा\appbrk{}न\appbrk{}प्र\appbrk{}यो\appbrk{}गेन}\ K H}}\skt{ }\edtext{\skt{आहारेण}}{\lemma{\skt{आहारेण}}\Afootnote{K N H; \skt{नाचाराञ्चा}\ A}}\skt{ }\edtext{\skt{सुयन्त्रितः}}{\lemma{\skt{सुयन्त्रितः}}\Afootnote{K N H; \skt{प्यतन्द्रितः}\ A}}\edtext{}{\lemma{\skt{भवति चात्र~॥ रसजं}}\Afootnote{\skt{रसजं}\ A K H; \skt{रजसं}\ N}}\pend

\pstart \ledsidenote{\passref{1.14.13}}\skt{तत्र रस }\edtext{\skt{गतौ}}{\lemma{\skt{गतौ}}\Afootnote{A; \skt{गतो}\ K N H}}\skt{ }\edtext{\skt{धातु\-र\-ह\-र\-ह\-र्ग\-च्छ\-तीति}}{\lemma{\skt{गच्छतीति}}\Afootnote{A K N H\textsuperscript{ac}; \skt{ग्गथंच्छतीति}\ H\textsuperscript{pc}}}\skt{ }\edtext{\skt{रसः}}{\lemma{\skt{रसः}}\Afootnote{K N H; \skt{अतो रसः}\ A}}\skt{~।}\pend

\pstart \ledsidenote{\passref{1.14.14}}\edtext{\skt{पञ्चविंशति}}{\lemma{\skt{पञ्चविंशति}}\Afootnote{K N; \skt{स खलु त्रीणि त्रीणि}\ A; \skt{स त्रीणि}\ H}}\skt{ }\edtext{\skt{कलाशतानि}}{\lemma{\skt{कलाशतानि}}\Afootnote{\skt{कला}\ A H; \skt{कलामतानि}\ K N}}\footnoteA{The reading of H here resembles the vulgate.}\skt{~। }\edtext{\skt{चतुरशीतिश्च}}{\lemma{\skt{चतुरशीतिश्च}}\Afootnote{N; \skt{सहस्राणि पञ्च}\ A H; \skt{चतुरशीतिम्च}\ K}}\skt{ }\edtext{\skt{नव}}{\lemma{\skt{नव}}\Afootnote{N; \skt{दश}\ A H; \emph{om.}\ K}}\skt{ }\edtext{\skt{च}}{\lemma{\skt{च}}\Afootnote{A N; \skt{नव च}\ K; \skt{कलाश्च}\ H}}\skt{ }\edtext{\skt{काष्ठा}}{\lemma{\skt{काष्ठा}}\Afootnote{K N; \skt{कला}\ A; \skt{ता}\ H}}\skt{ }\edtext{\edtext{\skt{एकै\-क\-स्मि\-न्धा\-ता\-व\-व\-ति\-ष्ठन्ते}}{\lemma{\skt{धाताव्}}\Afootnote{A K H; \skt{धातावाताव्}\ N}}}{\lemma{\skt{अवतिष्ठन्ते}}\Afootnote{K H; \skt{अवतिष्ठते}\ A; \skt{आवतिष्ठन्ते}\ N}}\skt{~। एवम्मासेन रसः }\edtext{\skt{शुक्रीभवति}}{\lemma{\skt{शुक्रीभवति}}\Afootnote{K N H; \skt{शुक्रं}\ A}}\skt{ }\edtext{\edtext{\skt{स्त्रीणा\-ञ्चा\-र्त्त\-व\-मिति}}{\lemma{\skt{चार्त्तवम्}}\Afootnote{A K N H\textsuperscript{ac}; \skt{चार्त्तव}\ H\textsuperscript{pc}}}}{\lemma{\skt{इति}}\Afootnote{K N H\textsuperscript{ac}; \skt{भवति}\ A; \skt{हृमिति}\ H\textsuperscript{pc}}}\skt{~॥}\pend

\pstart \ledsidenote{\passref{1.14.15}}\edtext{\edtext{\skt{भवतश्चात्र}}{\lemma{\skt{भवतश्}}\Afootnote{\skt{भवति}\ A H; \emph{om.}\ K N}}}{\lemma{\skt{चात्र}}\Afootnote{A H; \emph{om.}\ K N}}\skt{~॥ अष्टा\-द\-श\-स\-ह\-स्राणि }\edtext{\skt{संख्या }\edtext{\skt{ह्यस्मिन्समुच्चये}}{\lemma{\skt{अस्मिन्समुच्चये}}\Afootnote{A; \skt{अस्मिं समुच्चये}\ K N; \skt{ह्यस्मिन्समुच्चये}\ H}}}{\lemma{\skt{संख्या ह्य्}}\Afootnote{A K N; \skt{संख्या}\ H}}\skt{~। }\edtext{\skt{कला\-ना\-न्न\-व\-ति\-श्चापि}}{\lemma{\skt{चापि}}\Afootnote{K N H; \skt{प्रोक्ता}\ A}}\skt{ }\edtext{\skt{स्वत\-न्त्र\-प\-र\-त\-न्त्रतः}}{\lemma{\skt{स्वत\appbrk{}न्त्र\appbrk{}प\appbrk{}र\appbrk{}त\appbrk{}न्त्रतः}}\Afootnote{K N H; \skt{स्वत\appbrk{}न्त्र\appbrk{}प\appbrk{}र\appbrk{}त\appbrk{}न्त्रयोः}\ A}}\skt{~॥ }\edtext{\skt{रसे}}{\lemma{\skt{रसे}}\Afootnote{K N H; \emph{om.}\ A}}\skt{ }\edtext{\edtext{\skt{गतिविशेषोऽयं}}{\lemma{\skt{गतिविशेषो}}\Afootnote{K N H; \emph{om.}\ A}}}{\lemma{\skt{ऽयं}}\Afootnote{K N H; \emph{om.}\ A}}\skt{ }\edtext{\edtext{\skt{मन्दा\-ग्ने\-र\-व\-सा\-निकः}}{\lemma{\skt{मन्दाग्नेर्}}\Afootnote{K N H; \emph{om.}\ A}}}{\lemma{\skt{अवसानिकः}}\Afootnote{N; \emph{om.}\ A; \skt{एवमानिकः}\ K; \skt{अवमानिकः}\ H}}\skt{~। }\edtext{\skt{अन\-यै\-वो\-दि\-ता\-ग्नेश्च}}{\lemma{\skt{अन\appbrk{}यै\appbrk{}वो\appbrk{}दि\appbrk{}ता\appbrk{}ग्नेश्च}}\Afootnote{K N; \emph{om.}\ A; \skt{अन\appbrk{}यै\appbrk{}वो\appbrk{}दि\appbrk{}ता\appbrk{}ग्नेश्च…}\ H}}\skt{ }\edtext{\skt{विज्ञेयः}}{\lemma{\skt{विज्ञेयः}}\Afootnote{K N H; \emph{om.}\ A}}\skt{ }\edtext{\skt{कालसंख्यया}}{\lemma{\skt{कालसंख्यया}}\Afootnote{K N H; \emph{om.}\ A}}\skt{~॥}\pend

\pstart \ledsidenote{\passref{1.14.16}}\edtext{\skt{स शब्दा\-र्चि\-जा\-ल\-स\-न्ता\-न\-व\-द\-णुना}}{\lemma{\skt{स शब्दा\appbrk{}र्चि\appbrk{}जा\appbrk{}ल\appbrk{}स\appbrk{}न्ता\appbrk{}न\appbrk{}व\appbrk{}द\appbrk{}णुना}}\Afootnote{\skt{स शब्दा\appbrk{}र्चि\appbrk{}र्ज\appbrk{}ल\appbrk{}स\appbrk{}न्ता\appbrk{}न\appbrk{}व\appbrk{}द\appbrk{}णुना}\ A; \skt{स शब्दा\appbrk{}र्चि\appbrk{}जा\appbrk{}ल…\appbrk{}स\appbrk{}न्ता\appbrk{}न\appbrk{}व\appbrk{}द\appbrk{}णुना}\ K; \skt{स सब्दा\appbrk{}च्चि\appbrk{}ज\appbrk{}ल\appbrk{}स\appbrk{}न्ता\appbrk{}न\appbrk{}व\appbrk{}द\appbrk{}णुना}\ N; \skt{स शब्दा\appbrk{}र्च्चि\appbrk{}ज\appbrk{}ल\appbrk{}स\appbrk{}न्ता\appbrk{}न\appbrk{}व\appbrk{}द\appbrk{}नुना}\ H}}\skt{ }\edtext{\skt{विशे\-षे\-णा\-नु\-स\-र\-त्येवं}}{\lemma{\skt{विशे\appbrk{}षे\appbrk{}णा\appbrk{}नु\appbrk{}स\appbrk{}रत्य्}}\Afootnote{K N H; \skt{विशे\appbrk{}षे\appbrk{}णा\appbrk{}नु\appbrk{}धा\appbrk{}वत्य्}\ A}}\skt{ शरीरं }\edtext{\skt{केवलं}}{\lemma{\skt{केवलं}}\Afootnote{A K\textsuperscript{pc} N H; \emph{om.}\ K\textsuperscript{ac}}}\skt{~।}\pend

\pstart \ledsidenote{\passref{1.14.17}}\edtext{\edtext{\skt{व्याजी\-क\-र\-ण्य\-स्त्वौ\-षध्यः}}{\lemma{\skt{व्याजीकरण्यस्}}\Afootnote{K N; \skt{वाजीकरण्यस्}\ A; \skt{व्याजीकेरण्यस्}\ H}}}{\lemma{\skt{त्वौषध्यः}}\Afootnote{K N; \skt{त्वोषधयः}\ A; \skt{त्वौषधयः}\ H}}\skt{ }\edtext{\skt{स्वगु\-ण\-ब\-लो\-त्क\-र्षा\-द्वि\-रे\-च\-न\-व\-दु\-प\-युक्ताः}}{\lemma{\skt{स्वगु\appbrk{}ण\appbrk{}ब\appbrk{}लो\appbrk{}त्क\appbrk{}र्षा\appbrk{}द्वि\appbrk{}रे\appbrk{}च\appbrk{}न\appbrk{}व\appbrk{}दु\appbrk{}प\appbrk{}युक्ताः}}\Afootnote{H; \skt{स्वब\appbrk{}ल\appbrk{}गु\appbrk{}णो\appbrk{}त्क\appbrk{}र्षा\appbrk{}द्वि\appbrk{}रे\appbrk{}च\appbrk{}न\appbrk{}व\appbrk{}दु\appbrk{}प\appbrk{}युक्ताः}\ A; \skt{स्वगु\appbrk{}ण\appbrk{}ब\appbrk{}ला\appbrk{}त्क\appbrk{}र्षा\appbrk{}द्वि\appbrk{}रे\appbrk{}च\appbrk{}न\appbrk{}दु\appbrk{}प\appbrk{}युक्ताः}\ K\textsuperscript{ac}; \skt{स्वगु\appbrk{}ण\appbrk{}ब\appbrk{}ला\appbrk{}त्क\appbrk{}र्षा\appbrk{}द्वि\appbrk{}रे\appbrk{}च\appbrk{}न\appbrk{}व\appbrk{}दु\appbrk{}प\appbrk{}युक्ताः}\ K\textsuperscript{pc}; \skt{स्वगु\appbrk{}ण\appbrk{}ब\appbrk{}ला\appbrk{}त्क\appbrk{}र्षा\appbrk{}द्वि\appbrk{}रे\appbrk{}च\appbrk{}न\appbrk{}दु\appbrk{}प\appbrk{}युक्तः}\ N}}\skt{ शुक्रं }\edtext{\skt{विरेचयन्ति}}{\lemma{\skt{विरेचयन्ति}}\Afootnote{K N H; \skt{शीघ्रं विरेचयन्ति}\ A}}\skt{~॥}\pend

\pstart \ledsidenote{\passref{1.14.18}}\skt{यथा हि पुष्पमुकुलस्थो }\edtext{\skt{गन्धो न}}{\lemma{\skt{गन्धो न}}\Afootnote{A K H; \skt{गन्धो}\ N}}\skt{ शक्यमिहास्तीति वक्तुं~। नैव नास्तीति~। अथ चास्ति }\edtext{\edtext{\skt{सता\-म्भा\-वा\-ना\-मु\-त्प\-त्ति\-रिति}}{\lemma{\skt{सताम्}}\Afootnote{A K N; \skt{शताम्}\ H}}}{\lemma{\skt{उत्पत्तिर्}}\Afootnote{K N H; \skt{अभिव्यक्तिर्}\ A}}\skt{ }\edtext{\skt{कृत्वा}}{\lemma{\skt{कृत्वा}}\Afootnote{K N H; \skt{ज्ञात्वा}\ A}}\skt{ }\edtext{\skt{केवलं तु}}{\lemma{\skt{केवलं तु}}\Afootnote{K N H; \skt{केवलं}\ A}}\skt{ }\edtext{\skt{सौक्ष्म्या\-न्ना\-भि\-व्य\-ज्यते}}{\lemma{\skt{सौक्ष्म्यान्}}\Afootnote{A K H; \skt{सौक्ष्म्यं}\ N}}\skt{~। स एव }\edtext{\skt{विवृतकेसरे}}{\lemma{\skt{विवृतकेसरे}}\Afootnote{H; \skt{विवृ\appbrk{}त\appbrk{}प\appbrk{}त्र\appbrk{}के\appbrk{}शरे}\ A; \skt{विवृते केशरे}\ K; \skt{विवृत्केसरे}\ N}}\skt{ }\edtext{\skt{पुष्पे}}{\lemma{\skt{पुष्पे}}\Afootnote{A H; \skt{पुष्प}\ K N}}\skt{ }\edtext{\skt{काला\-न्त\-रे\-णा\-भि\-व्यक्तो}}{\lemma{\skt{काला\appbrk{}न्त\appbrk{}रे\appbrk{}णा\appbrk{}भि\appbrk{}व्यक्तो}}\Afootnote{K N H; \skt{काला\appbrk{}न्त\appbrk{}रे\appbrk{}णा\appbrk{}भि\appbrk{}व्यक्तिं}\ A}}\skt{ }\edtext{\skt{भवति}}{\lemma{\skt{भवति}}\Afootnote{K N H; \skt{गच्छति}\ A}}\skt{~। एव\-म्बा\-ला\-ना\-मपि }\edtext{\skt{वयः\-प\-रि\-णा\-मा\-च्छु\-क्र\-प्रा\-दु\-र्भावो}}{\lemma{\skt{वयः\appbrk{}प\appbrk{}रि\appbrk{}णा\appbrk{}मा\appbrk{}च्छु\appbrk{}क्र\appbrk{}प्रा\appbrk{}दु\appbrk{}र्भावो}}\Afootnote{A K N; \skt{वयः\appbrk{}प\appbrk{}रि\appbrk{}णा\appbrk{}मा\appbrk{}च्छु\appbrk{}क्र\appbrk{}प्रा\appbrk{}दु\appbrk{}र्भ्भावो}\ H}}\skt{ भवति~। }\edtext{\skt{रोम\-रा\-ज्या\-र्त्त\-वा\-दिश्च}}{\lemma{\skt{रोम\appbrk{}रा\appbrk{}ज्या\appbrk{}र्त्त\appbrk{}वा\appbrk{}दिश्च}}\Afootnote{K H; \skt{रोम\appbrk{}रा\appbrk{}ज्या\appbrk{}द\appbrk{}यश्च}\ A; \skt{रोम\appbrk{}रा\appbrk{}ज्या\appbrk{}र्त्त\appbrk{}वा\appbrk{}दिश्च}\ N}}\skt{ }\edtext{\skt{विशेषो}}{\lemma{\skt{विशेषो}}\Afootnote{K N; \skt{विशेषा}\ A; \skt{विशेषौ}\ H}}\skt{ नारीणाम्~।}\pend

\pstart \ledsidenote{\passref{1.14.19}}\skt{स }\edtext{\skt{एवा\-न्न\-र\-सोऽ\-भि\-वृ\-द्धानां}}{\lemma{\skt{ऽभिवृद्धानां}}\Afootnote{K N H; \skt{वृद्धानां}\ A}}\skt{ }\edtext{\skt{परि\-प\-क्व\-श\-री\-र\-त्वा\-द\-प्री\-णनो}}{\lemma{\skt{परि\appbrk{}प\appbrk{}क्व\appbrk{}श\appbrk{}री\appbrk{}र\appbrk{}त्वा\appbrk{}द\appbrk{}प्री\appbrk{}णनो}}\Afootnote{H; \skt{(जरा) परि\appbrk{}प\appbrk{}क्व\appbrk{}श\appbrk{}री\appbrk{}र\appbrk{}त्वा\appbrk{}द\appbrk{}प्री\appbrk{}णनो}\ A; \skt{परि\appbrk{}प\appbrk{}क्व\appbrk{}श\appbrk{}री\appbrk{}र\appbrk{}त्वा\appbrk{}द\appbrk{}प्री\appbrk{}णानो}\ K N}}\skt{ भवति~।}\pend

\pstart \ledsidenote{\passref{1.14.20}}\skt{त एते }\edtext{\skt{शरी\-र\-धा\-र\-णा\-द्धा\-तव}}{\lemma{\skt{शरी\appbrk{}र\appbrk{}धा\appbrk{}र\appbrk{}णा\appbrk{}द्धा\appbrk{}तव}}\Afootnote{A H; \skt{शरी\appbrk{}धा\appbrk{}णा\appbrk{}र\appbrk{}णा\appbrk{}द्धा\appbrk{}तव}\ K; \skt{शरी\appbrk{}र\appbrk{}धा\appbrk{}रि\appbrk{}णा\appbrk{}द्धा\appbrk{}तव}\ N}}\skt{ इत्युच्यन्ते~।}\pend

\pstart \ledsidenote{\passref{1.14.21}}\skt{तेषां क्षयवृद्धी शोणितनिमित्ते~। }\edtext{\skt{तस्मा\-त्त\-द\-धि\-कृत्य}}{\lemma{\skt{तस्मा\appbrk{}त्त\appbrk{}द\appbrk{}धि\appbrk{}कृत्य}}\Afootnote{A K N H\textsuperscript{pc}; \skt{तम्मा\appbrk{}त्त\appbrk{}द\appbrk{}धि\appbrk{}कृत्य}\ H\textsuperscript{ac}}}\skt{ वक्ष्यामः~॥ तत्र }\edtext{\edtext{\skt{सफे\-नि\-ल\-म\-रुणं}}{\lemma{\skt{सफेनिलम्}}\Afootnote{K N H; \skt{फेनिल}\ A}}}{\lemma{\skt{अरुणं}}\Afootnote{K N H; \skt{मरुणं}\ A}}\skt{ }\edtext{\skt{कृष्ण\-म्प\-रु\-ष\-न्तनु}}{\lemma{\skt{परुषन्}}\Afootnote{A K N; \skt{परूषं}\ H}}\skt{ }\edtext{\skt{शीघ्रगमस्कन्दि च}}{\lemma{\skt{शीघ्रगमस्कन्दि च}}\Afootnote{A; \skt{शीघ्रगमस्कन्दि च}\ K H; \skt{शीघ्रमस्कन्दि च}\ N}}\skt{ }\edtext{\skt{वातदुष्टं}}{\lemma{\skt{वातदुष्टं}}\Afootnote{K N H; \skt{वातेन दुष्टं}\ A}}\skt{~। नीलम्पीतं हरितं }\edtext{\skt{श्यावं}}{\lemma{\skt{श्यावं}}\Afootnote{A K H; \skt{श्यामम्}\ N}}\skt{ }\edtext{\edtext{\edtext{\skt{विस्र\-म\-नि\-ष्ट\-म्पि\-पी\-लि\-का\-म\-क्षि\-का\-णाञ्च}}{\lemma{\skt{विस्रमनिष्टम्}}\Afootnote{A; \skt{विश्रमनिष्टं}\ K H N}}}{\lemma{\skt{पिपी\appbrk{}लि\appbrk{}का\appbrk{}म\appbrk{}क्षि\appbrk{}का\appbrk{}णाञ्}}\Afootnote{K N H; \skt{पिपी\appbrk{}लि\appbrk{}का\appbrk{}म\appbrk{}क्षि\appbrk{}काणा}\ A}}}{\lemma{\skt{च}}\Afootnote{K N H; \skt{मस्कन्दि च}\ A}}\skt{ पित्तेन~। }\edtext{\skt{गैरि\-को\-द\-क\-प्र\-काशं}}{\lemma{\skt{गैरि\appbrk{}को\appbrk{}द\appbrk{}क\appbrk{}प्र\appbrk{}काशं}}\Afootnote{K N H; \skt{दुष्टं; गैरि\appbrk{}को\appbrk{}द\appbrk{}क\appbrk{}प्र\appbrk{}ती\appbrk{}काशं}\ A}}\skt{ }\edtext{\skt{स्निग्धं}}{\lemma{\skt{स्निग्धं}}\Afootnote{A K H; \skt{सिग्धं}\ N}}\skt{ }\edtext{\skt{शीतं}}{\lemma{\skt{शीतं}}\Afootnote{K N H; \skt{शीतलं}\ A}}\skt{ }\edtext{\skt{बहलं}}{\lemma{\skt{बहलं}}\Afootnote{A K H; \skt{वहल}\ N}}\skt{ पिच्छिलं }\edtext{\skt{विस्रावि}}{\lemma{\skt{विस्रावि}}\Afootnote{K N; \skt{चिरस्रावि}\ A; \skt{विश्रावि}\ H}}\skt{ }\edtext{\skt{मांस\-पे\-शी\-स\-म\-प्र\-भञ्च}}{\lemma{\skt{मांस\appbrk{}पे\appbrk{}शी\appbrk{}स\appbrk{}म\appbrk{}प्र\appbrk{}भञ्च}}\Afootnote{K N H; \skt{मांसपेशीप्रभं च}\ A}}\skt{ }\edtext{\skt{श्लेष्मणा}}{\lemma{\skt{श्लेष्मणा}}\Afootnote{K N H; \skt{श्लेष्मदुष्टं}\ A}}\skt{~। }\edtext{\skt{सर्व\-ल\-क्ष\-ण\-युक्तं}}{\lemma{\skt{सर्व\appbrk{}ल\appbrk{}क्ष\appbrk{}ण\appbrk{}युक्तं}}\Afootnote{K N; \skt{सर्व\appbrk{}ल\appbrk{}क्ष\appbrk{}ण\appbrk{}सं\appbrk{}युक्तं काञ्जिकाभं विशेषतो दुर्गन्धि च}\ A; \skt{सर्व्व\appbrk{}ल\appbrk{}क्ष\appbrk{}ण\appbrk{}सं\appbrk{}युक्तं}\ H}}\skt{ }\edtext{\skt{सन्निपातेन}}{\lemma{\skt{सन्निपातेन}}\Afootnote{K N H; \skt{सन्निपा}\ A}}\skt{~। }\edtext{\skt{पित्त\-व\-द्र\-क्ते\-ना\-ति\-कृ\-ष्णञ्च}}{\lemma{\skt{पित्त\appbrk{}व\appbrk{}द्र\appbrk{}क्ते\appbrk{}ना\appbrk{}ति\appbrk{}कृ\appbrk{}ष्णञ्च}}\Afootnote{K N; \skt{तदुष्टं}\ A; \skt{पित्त\appbrk{}व\appbrk{}द्र\appbrk{}क्ते\appbrk{}ना\appbrk{}ति\appbrk{}कृ\appbrk{}ष्णश्च}\ H}}\skt{~। }\edtext{\skt{द्विदो\-ष\-लि\-ङ्ग\-सं\-सृष्टं}}{\lemma{\skt{द्विदो\appbrk{}ष\appbrk{}लि\appbrk{}ङ्ग\appbrk{}सं\appbrk{}सृष्टं}}\Afootnote{K N H; \skt{द्विदोषलिङ्गं संसृष्टम्}\ A}}\skt{ }\edtext{\skt{द्विदोषं}}{\lemma{\skt{द्विदोषं}}\Afootnote{K N H; \emph{om.}\ A}}\skt{~।}\pend

\pstart \ledsidenote{\passref{1.14.22}}\edtext{\edtext{\edtext{\skt{इन्द्र\-गो\-प\-क\-प्र\-का\-श\-म\-सं\-ह\-त\-म\-वि\-व\-र्ण्णञ्च}}{\lemma{\skt{इन्द्र\appbrk{}गो\appbrk{}प\appbrk{}क\appbrk{}प्र\appbrk{}का\appbrk{}शम्}}\Afootnote{K N H; \skt{इन्द्र\appbrk{}गो\appbrk{}प\appbrk{}क\appbrk{}प्र\appbrk{}ती\appbrk{}काश}\ A}}}{\lemma{\skt{असंहतम्}}\Afootnote{K N H; \skt{मसंहत}\ A}}}{\lemma{\skt{अविवर्ण्णञ्च}}\Afootnote{K N H; \skt{मविवर्णं च}\ A}}\skt{ }\edtext{\skt{प्रकृतिस्थमिति}}{\lemma{\skt{इति}}\Afootnote{K N H; \emph{om.}\ A}}\skt{ जानीयात्~॥}\pend

\pstart \ledsidenote{\passref{1.14.23}}\edtext{\skt{विस्राव्यानन्यत्र}}{\lemma{\skt{विस्राव्यानन्यत्र}}\Afootnote{K N; \skt{विस्राव्याण्यन्यत्र}\ A; \skt{विश्राव्यानन्यत्र}\ H}}\skt{ }\edtext{\skt{वक्ष्यामः}}{\lemma{\skt{वक्ष्यामः}}\Afootnote{A K N\textsuperscript{pc} H; \skt{वक्ष्यामिअः}\ N\textsuperscript{ac}}}\skt{~॥}\pend

\pstart \ledsidenote{\passref{1.14.24}}\edtext{\skt{अथाविस्राव्याः}}{\lemma{\skt{अथाविस्राव्याः}}\Afootnote{A K N; \skt{अथाविश्राव्याः}\ H}}\skt{~। सर्वाङ्गशोफः }\edtext{\skt{क्षीणस्य}}{\lemma{\skt{क्षीणस्य}}\Afootnote{A N; \skt{क्षीणस्य वा}\ K H\textsuperscript{pc}; \skt{क्षीण वा}\ H\textsuperscript{ac}}}\skt{ }\edtext{\skt{चाम्ल\-भो\-ज\-न\-नि\-मित्तः}}{\lemma{\skt{चाम्ल\appbrk{}भो\appbrk{}ज\appbrk{}न\appbrk{}नि\appbrk{}मित्तः}}\Afootnote{A K H; \skt{चाम्व्ल\appbrk{}भो\appbrk{}ज\appbrk{}न\appbrk{}नि\appbrk{}मित्तः}\ N}}\skt{~। }\edtext{\skt{पाण्डु\-रो\-ग्य\-र्श\-स्यु\-द\-रि\-शो\-षि\-ग\-र्भि\-णी\-नाञ्च}}{\lemma{\skt{पाण्डु\appbrk{}रो\appbrk{}ग्य\appbrk{}र्श\appbrk{}स्यु\appbrk{}द\appbrk{}रि\appbrk{}शो\appbrk{}षि\appbrk{}ग\appbrk{}र्भि\appbrk{}णी\appbrk{}नाञ्च}}\Afootnote{K N; \skt{पाण्डु\appbrk{}रो\appbrk{}ग्य\appbrk{}र्श\appbrk{}सो\appbrk{}द\appbrk{}रि\appbrk{}शो\appbrk{}षि\appbrk{}ग\appbrk{}र्भि\appbrk{}णीनां च}\ A; \skt{पाण्दु\appbrk{}रो\appbrk{}ग्य\appbrk{}र्श\appbrk{}स्यु\appbrk{}द\appbrk{}रि\appbrk{}शो\appbrk{}षि\appbrk{}ग\appbrk{}र्भ्भि\appbrk{}णी\appbrk{}नाञ्च}\ H}}\skt{ श्वयथवः~।}\pend

\pstart \ledsidenote{\passref{1.14.26}}\skt{तत्र ऋज्वसंकीर्ण्णं सूक्ष्मं }\edtext{\skt{सम\-म\-न\-व\-गा\-ढ\-म\-नु\-त्ता\-न\-माशु}}{\lemma{\skt{अनुत्तानम्}}\Afootnote{A K H; \skt{अन्त्रत्तानम्}\ N}}\skt{ }\edtext{\skt{शस्त्रञ्च}}{\lemma{\skt{शस्त्रञ्च}}\Afootnote{K N H; \skt{च शस्त्रं}\ A}}\skt{ }\edtext{\skt{पातयेत्}}{\lemma{\skt{पातयेत्}}\Afootnote{K N H; \skt{पात\appbrk{}ये\appbrk{}न्म\appbrk{}र्म\appbrk{}सि\appbrk{}रा\appbrk{}स्ना\appbrk{}यु\appbrk{}स\appbrk{}न्धीनां चानुपघाति}\ A}}\skt{~।}\pend

\pstart \ledsidenote{\passref{1.14.26a}}\edtext{\skt{हृद\-य\-ब\-स्ति\-गु\-द\-ना\-भि\-क\-क्ष्य\-व\-ङ्क्ष\-णा\-क्षि\-कू\-ट\-पा\-णि\-पा\-द\-त\-लांश्च}}{\lemma{\skt{हृद\appbrk{}य\appbrk{}ब\appbrk{}स्ति\appbrk{}गु\appbrk{}द\appbrk{}ना\appbrk{}भि\appbrk{}क\appbrk{}क्ष्य\appbrk{}व\appbrk{}ङ्क्ष\appbrk{}णा\appbrk{}क्षि\appbrk{}कू\appbrk{}ट\appbrk{}पा\appbrk{}णि\appbrk{}पा\appbrk{}द\appbrk{}त\appbrk{}लांश्च}}\Afootnote{K N; \skt{हृद\appbrk{}य\appbrk{}ब\appbrk{}स्ति\appbrk{}गु\appbrk{}द\appbrk{}ना\appbrk{}भि\appbrk{}क\appbrk{}क्ष\appbrk{}वं\appbrk{}क्ष\appbrk{}णा\appbrk{}क्षि\appbrk{}कू\appbrk{}ट\appbrk{}पा\appbrk{}णि\appbrk{}पा\appbrk{}द\appbrk{}त\appbrk{}लाश्च}\ H}}\skt{ वर्जयेत्~।}\pend

\pstart \ledsidenote{\passref{1.14.26b}}\edtext{\skt{पूयगर्भां}}{\lemma{\skt{पूयगर्भां}}\Afootnote{K N; \skt{पूयगर्भ्भां}\ H}}\skt{ पुन\-र्य\-थो\-क्तै\-रे\-वो\-प\-च\-रेत्~।}\pend

\pstart \ledsidenote{\passref{1.14.27}}\edtext{\skt{तत्र}}{\lemma{\skt{तत्र}}\Afootnote{K N H; \skt{तत्र दुर्दिने}\ A}}\skt{ }\edtext{\skt{दुर्विद्धे}}{\lemma{\skt{दुर्विद्धे}}\Afootnote{A N H; \skt{दुर्वद्धे}\ K}}\skt{ शीत\-वा\-त\-यो\-र\-स्विन्ने }\edtext{\skt{भुक्ते च}}{\lemma{\skt{भुक्ते च}}\Afootnote{K N H; \skt{भुक्तमात्रे}\ A}}\skt{ }\edtext{\skt{स्कन्न\-त्वा\-च्छो\-णि\-तन्न}}{\lemma{\skt{स्कन्न\appbrk{}त्वा\appbrk{}च्छो\appbrk{}णि\appbrk{}तन्न}}\Afootnote{K N H; \skt{स्कन्दत्वाच्छोणितं न}\ A}}\skt{ }\edtext{\skt{स्रवति}}{\lemma{\skt{स्रवति}}\Afootnote{A K N; \skt{श्रवति}\ H}}\skt{~। अल्पं वा }\edtext{\skt{स्रवति}}{\lemma{\skt{स्रवति}}\Afootnote{A K N; \skt{श्रवति}\ H}}\skt{~॥}\pend

\pstart\versepar \ledsidenote{\passref{1.14.28}}\edtext{\skt{भवति चात्र}}{\lemma{\skt{भवति चात्र}}\Afootnote{H; \skt{मद\appbrk{}मू\appbrk{}र्च्छा\appbrk{}श्र\appbrk{}मा\appbrk{}र्तानां}\ A; \emph{om.}\ K\textsuperscript{ac} N; \skt{भ~॥ चात्र}\ K\textsuperscript{pc}}}\skt{~। }\pend
\pstart\versepar \skt{}\edtext{\skt{वात\-वि\-ण्मू\-त्र\-सं\-गेषु}}{\lemma{\skt{वात\appbrk{}वि\appbrk{}ण्मू\appbrk{}त्र\appbrk{}सं\appbrk{}गेषु}}\Afootnote{K; \skt{वात\appbrk{}वि\appbrk{}ण्मू\appbrk{}त्र\appbrk{}सङ्गि}\ A; \skt{॥ वात\appbrk{}वि\appbrk{}ड्मू\appbrk{}त्र\appbrk{}सं\appbrk{}गेषु}\ N; \skt{वात\appbrk{}वि\appbrk{}मू\appbrk{}त्र\appbrk{}स\appbrk{}ङ्गेषु}\ H}}\skt{ }\edtext{\skt{मद\-मू\-र्च्छा\-श्र\-मेषु च}}{\lemma{\skt{मद\appbrk{}मू\appbrk{}र्च्छा\appbrk{}श्र\appbrk{}मेषु च}}\Afootnote{K N H; \skt{नाम्}\ A}}\skt{~। }\pend
\pstart\versepar \skt{}\edtext{\skt{निद्राभिभूते}}{\lemma{\skt{निद्राभिभूते}}\Afootnote{K N H; \skt{निद्राभिभूत}\ A}}\skt{ }\edtext{\skt{शीते वा}}{\lemma{\skt{शीते वा}}\Afootnote{K N H; \skt{भीतानां}\ A}}\skt{ }\edtext{\edtext{\skt{नृणा\-न्ना\-सृ\-क्स्र\-वे\-दिति}}{\lemma{\skt{नृणान्}}\Afootnote{A K N; \skt{नॄणान्}\ H}}}{\lemma{\skt{नासृक्स्रवेदिति}}\Afootnote{K N; \skt{नासृक्प्रवर्तते}\ A; \skt{नासृवेदिति}\ H\textsuperscript{ac}; \skt{नासृवेथंदिति}\ H\textsuperscript{pc}}}\skt{~॥}\pend

\pstart \ledsidenote{\passref{1.14.29}}\edtext{\edtext{\skt{तद्दु\-ष्ट\-शो\-णि\-त\-म\-नि\-र्ह्रि\-य\-माणं}}{\lemma{\skt{तद्दु\appbrk{}ष्ट\appbrk{}शो\appbrk{}णि\appbrk{}तम्}}\Afootnote{K N H; \skt{तद्दुष्टं शोणित}\ A}}}{\lemma{\skt{अनिर्ह्रियमाणं}}\Afootnote{K N H; \skt{मनिर्ह्रियमाणं}\ A}}\skt{ }\edtext{\skt{व्याधिवृद्धिं}}{\lemma{\skt{व्याधिवृद्धिं}}\Afootnote{K N H; \emph{om.}\ A}}\skt{ }\edtext{\skt{करोति}}{\lemma{\skt{करोति}}\Afootnote{K N H; \skt{शोफ\appbrk{}दा\appbrk{}ह\appbrk{}रा\appbrk{}ग\appbrk{}पा\appbrk{}क\appbrk{}वे\appbrk{}दना जनयेत्}\ A}}\skt{~।}\pend

\pstart \ledsidenote{\passref{1.14.30}}\edtext{\edtext{\edtext{\skt{अत्यु\-ष्णा\-ति\-स्वि\-न्ना\-ति\-वि\-द्धे\-ष्व\-ज्ञ\-स्रा\-वि\-त\-म\-ति\-प्र\-व\-र्त्तते}}{\lemma{\skt{अत्यु\appbrk{}ष्णा\appbrk{}ति\appbrk{}स्वि\appbrk{}न्ना\appbrk{}ति\appbrk{}वि\appbrk{}द्धेष्व्}}\Afootnote{K N H; \skt{अत्यु\appbrk{}ष्णेऽ\appbrk{}ति\appbrk{}स्वि\appbrk{}न्नेऽ\appbrk{}ति\appbrk{}विद्धे}\ A}}}{\lemma{\skt{अज्ञस्रावितम्}}\Afootnote{K N H; \skt{ऽज्ञैर्विस्रावित}\ A}}}{\lemma{\skt{अतिप्रवर्त्तते}}\Afootnote{K N H; \skt{मतिप्रवर्तते}\ A}}\skt{~। }\edtext{\skt{तदतिप्रवृत्तं}}{\lemma{\skt{तदतिप्रवृत्तं}}\Afootnote{A K H; \skt{तदतिप्रवृत्त}\ N}}\skt{ }\edtext{\edtext{\edtext{\skt{शिरो\-भि\-ता\-प\-मा\-न्ध्य\-ति\-मि\-र\-प्रा\-दु\-र्भा\-व\-न्चा\-नु\-क्ष\-या\-क्षे\-प\-क\-म्प\-प\-क्षा\-घा\-त\-मे\-का\-ङ्ग\-वि\-कारं}}{\lemma{\skt{शिरोभितापम्}}\Afootnote{K N; \skt{शिरोऽभिताप}\ A; \skt{शिरोऽभितापम्}\ H}}}{\lemma{\skt{आन्ध्य\appbrk{}ति\appbrk{}मि\appbrk{}र\appbrk{}प्रा\appbrk{}दु\appbrk{}र्भा\appbrk{}व\appbrk{}न्चा\appbrk{}नु\appbrk{}क्ष\appbrk{}या\appbrk{}क्षे\appbrk{}प\appbrk{}क\appbrk{}म्प\appbrk{}प\appbrk{}क्षा\appbrk{}घा\appbrk{}तम्}}\Afootnote{\skt{मान्ध्य\appbrk{}म\appbrk{}धि\appbrk{}म\appbrk{}न्थ\appbrk{}ति\appbrk{}मि\appbrk{}र\appbrk{}प्रा\appbrk{}दु\appbrk{}र्भावं धातु\appbrk{}क्ष\appbrk{}य\appbrk{}मा\appbrk{}क्षे\appbrk{}पकं दाहं पक्षाघात}\ A; \skt{आन्ध्य\appbrk{}ति\appbrk{}मि\appbrk{}र\appbrk{}प्रा\appbrk{}दु\appbrk{}र्भा\appbrk{}व\appbrk{}म्धा\appbrk{}तु\appbrk{}क्ष\appbrk{}या\appbrk{}क्षे\appbrk{}प\appbrk{}कस्य पक्षाघातम्}\ K N; \skt{आन्ध्य\appbrk{}ति\appbrk{}मि\appbrk{}र\appbrk{}प्रा\appbrk{}दु\appbrk{}र्भ्भा\appbrk{}व\appbrk{}न्धा\appbrk{}तु\appbrk{}क्ष\appbrk{}या\appbrk{}क्षे\appbrk{}पकं पक्षाघातम्}\ H\textsuperscript{ac}; \skt{आन्ध्य\appbrk{}ति\appbrk{}मि\appbrk{}र\appbrk{}भ्र\appbrk{}प्रा\appbrk{}दु\appbrk{}र्भ्भा\appbrk{}व\appbrk{}न्धा\appbrk{}तु\appbrk{}क्ष\appbrk{}या\appbrk{}क्षे\appbrk{}पकं पक्षाघातम्}\ H\textsuperscript{pc}}}}{\lemma{\skt{एकाङ्गविकारं}}\Afootnote{K N H; \skt{मेकाङ्गविकारं}\ A}}\skt{ }\edtext{\skt{हिक्काकासं}}{\lemma{\skt{हिक्काकासं}}\Afootnote{N; \skt{हिक्कां श्वास}\ A; \skt{हिक्कां कासं}\ K H}}\skt{ }\edtext{\skt{श्वासं}}{\lemma{\skt{श्वासं}}\Afootnote{K N H; \skt{कासौ}\ A}}\skt{ पाण्डुरोगं }\edtext{\skt{मरणं}}{\lemma{\skt{मरणं}}\Afootnote{A N H; \emph{om.}\ K}}\skt{ }\edtext{\skt{वाशु}}{\lemma{\skt{वाशु}}\Afootnote{K N H; \skt{चापा}\ A}}\skt{ }\edtext{\skt{करोति}}{\lemma{\skt{करोति}}\Afootnote{K N H; \skt{दयति}\ A}}\skt{~।}\pend

\pstart \ledsidenote{\passref{1.14.31}}\edtext{\skt{तन्नातिशीते}}{\lemma{\skt{तन्नातिशीते}}\Afootnote{K N H; \skt{तस्मान्न शीते}\ A}}\skt{ नात्युष्णे }\edtext{\skt{नास्विन्ने}}{\lemma{\skt{नास्विन्ने}}\Afootnote{A K N; \emph{om.}\ H}}\skt{ नातितापिते }\edtext{\skt{यवागूं}}{\lemma{\skt{यवागूं}}\Afootnote{A K N; \skt{यवागूस्}\ H}}\skt{ प्रतिपीतस्य }\edtext{\skt{शोणितं}}{\lemma{\skt{शोणितं}}\Afootnote{A K H; \skt{शौणितं}\ N}}\skt{ }\edtext{\skt{मोक्षयेद्भिषक्}}{\lemma{\skt{मोक्षयेद्भिषक्}}\Afootnote{A K H; \skt{मोक्षये भिषक्}\ N}}\skt{~॥}\pend

\pstart\versepar \ledsidenote{\passref{1.14.32}}\skt{सम्यग्गत्वा यदा }\edtext{\skt{रक्तं}}{\lemma{\skt{रक्तं}}\Afootnote{A K N H\textsuperscript{pc}; \skt{रक्तं म्}\ H\textsuperscript{ac}}}\skt{ }\edtext{\skt{स्वय\-मे\-वा\-व\-ति\-ष्ठते}}{\lemma{\skt{स्वयम्}}\Afootnote{A K N H\textsuperscript{pc}; \skt{वयम्}\ H\textsuperscript{ac}}}\skt{ }\pend
\pstart\versepar \skt{}\edtext{\skt{शुद्ध\-मे\-व\-म्वि\-जा\-नी\-या\-त्स\-म्य\-ग्वि\-स्रा\-वि\-तञ्च}}{\lemma{\skt{एवम्}}\Afootnote{K N H; \skt{तदा}\ A}}\skt{ तत्~।}\pend

\pstart\versepar \ledsidenote{\passref{1.14.33}}\skt{लाघवं }\edtext{\skt{वेद\-ना\-शा\-न्ति\-र्व्या\-धि\-वे\-ग\-प\-रि\-क्षयः}}{\lemma{\skt{व्याधि\appbrk{}वे\appbrk{}ग\appbrk{}प\appbrk{}रि\appbrk{}क्षयः}}\Afootnote{K N H; \skt{व्याधे\appbrk{}र्वे\appbrk{}ग\appbrk{}प\appbrk{}रि\appbrk{}क्षयः}\ A}}\skt{~। }\pend
\pstart\versepar \skt{सम्यग्विस्राविते लिङ्गं प्रसादो मनसस्तथा~॥}\pend

\pstart\versepar \ledsidenote{\passref{1.14.34}}\skt{त्वग्दोषा ग्रन्थयः शोफा रोगाः }\edtext{\skt{शोणितजाश्च ये}}{\lemma{\skt{शोणितजाश्च ये}}\Afootnote{A K H; \skt{शोणितजाश्च त्ये~।}\ N}}\skt{~। }\pend
\pstart\versepar \skt{रक्त\-मो\-क्ष\-ण\-शी\-लानां न भवन्ति कदाचन~॥}\pend

\pstart \ledsidenote{\passref{1.14.35}}\skt{अथ खल्व\-प्र\-व\-र्त्त\-माने~। }\edtext{\skt{एलाशीतशिवः कुष्ठ\-त\-ग\-र\-पा\-ठा\-भ\-द्र\-दा\-रु\-वि\-डं\-ग\-चि\-त्र\-क\-त्रि\-क\-टु\-का\-गा\-र\-धू\-म\-ह\-रि\-द्रा\-र्क्कां\-कु\-र\-न\-क्त\-मा\-ल\-फ\-लै\-र्य\-था\-लाभं}}{\lemma{\skt{एलाशीतशिवः कुष्ठ\appbrk{}त\appbrk{}ग\appbrk{}र\appbrk{}पा\appbrk{}ठा\appbrk{}भ\appbrk{}द्र\appbrk{}दा\appbrk{}रु\appbrk{}वि\appbrk{}डं\appbrk{}ग\appbrk{}चि\appbrk{}त्र\appbrk{}क\appbrk{}त्रि\appbrk{}क\appbrk{}टु\appbrk{}का\appbrk{}गा\appbrk{}र\appbrk{}धू\appbrk{}म\appbrk{}ह\appbrk{}रि\appbrk{}द्रा\appbrk{}र्क्कां\appbrk{}कु\appbrk{}र\appbrk{}न\appbrk{}क्त\appbrk{}मा\appbrk{}ल\appbrk{}फ\appbrk{}लैर्}}\Afootnote{H; \skt{रक्ते एला\appbrk{}शी\appbrk{}त\appbrk{}शि\appbrk{}व\appbrk{}कु\appbrk{}ष्ठ\appbrk{}त\appbrk{}ग\appbrk{}र\appbrk{}पा\appbrk{}ठा\appbrk{}भ\appbrk{}द्र\appbrk{}दा\appbrk{}रु\appbrk{}वि\appbrk{}ड\appbrk{}ङ्ग\appbrk{}चि\appbrk{}त्र\appbrk{}क\appbrk{}त्रि\appbrk{}क\appbrk{}टु\appbrk{}का\appbrk{}गा\appbrk{}र\appbrk{}धू\appbrk{}म\appbrk{}ह\appbrk{}रि\appbrk{}द्रा\appbrk{}र्का\appbrk{}ङ्कु\appbrk{}र\appbrk{}न\appbrk{}क्त\appbrk{}मा\appbrk{}ल\appbrk{}फ\appbrk{}लैर्}\ A; \skt{एला शीतशिवः कुष्ठ\appbrk{}त\appbrk{}ग\appbrk{}र\appbrk{}पा\appbrk{}ठा\appbrk{}भ\appbrk{}द्र\appbrk{}दा\appbrk{}रु\appbrk{}वि\appbrk{}डं\appbrk{}ग\appbrk{}चि\appbrk{}त्र\appbrk{}क\appbrk{}त्रि\appbrk{}क\appbrk{}टु\appbrk{}का\appbrk{}गा\appbrk{}र\appbrk{}धू\appbrk{}म\appbrk{}ह\appbrk{}रि\appbrk{}द्रा\appbrk{}र्का\appbrk{}ङ्कु\appbrk{}न\appbrk{}क्त\appbrk{}मा\appbrk{}ल\appbrk{}फलै}\ K; \skt{एला\appbrk{}शी\appbrk{}त\appbrk{}शि\appbrk{}व\appbrk{}कु\appbrk{}ष्ठ\appbrk{}त\appbrk{}ग\appbrk{}र\appbrk{}पा\appbrk{}ठा\appbrk{}भ\appbrk{}द्र\appbrk{}दा\appbrk{}रु\appbrk{}वि\appbrk{}डि\appbrk{}ङ्ग\appbrk{}चि\appbrk{}त्र\appbrk{}त\appbrk{}क\appbrk{}टु\appbrk{}का\appbrk{}गा\appbrk{}र\appbrk{}धू\appbrk{}म\appbrk{}ह\appbrk{}रि\appbrk{}द्रा\appbrk{}र्क्का\appbrk{}ङ्कु\appbrk{}र\appbrk{}न\appbrk{}क्त\appbrk{}मा\appbrk{}ल\appbrk{}फ\appbrk{}लैर्}\ N}}\skt{ }\edtext{\skt{त्रिभिश्चतुर्भिः}}{\lemma{\skt{चतुर्भिः}}\Afootnote{A K N; \skt{चतुर्भ्भिः}\ H}}\skt{~। }\edtext{\skt{समस्तैर्वा}}{\lemma{\skt{समस्तैर्वा}}\Afootnote{A K N; \skt{…समस्तैर्व्वा}\ H}}\skt{ }\edtext{\edtext{\edtext{\skt{लव\-ण\-प्र\-गा\-ढै\-र्व्र\-ण\-मु\-ख\-म\-व\-घ\-र्ष\-ये\-देवं}}{\lemma{\skt{लव\appbrk{}ण\appbrk{}प्र\appbrk{}गा\appbrk{}ढैर्}}\Afootnote{K N H; \skt{चूर्णी\appbrk{}कृ\appbrk{}तै\appbrk{}र्ल\appbrk{}व\appbrk{}ण\appbrk{}तै\appbrk{}ल\appbrk{}प्र\appbrk{}गा\appbrk{}ढैर्}\ A}}}{\lemma{\skt{व्रणमुखम्}}\Afootnote{A N H; \skt{व्रणमु}\ K}}}{\lemma{\skt{अव\appbrk{}घ\appbrk{}र्ष\appbrk{}ये\appbrk{}देवं}}\Afootnote{A N H; \skt{ख घर्षयेदेवं}\ K}}\skt{ }\edtext{\skt{साधु}}{\lemma{\skt{साधु}}\Afootnote{K N H; \skt{सम्य}\ A}}\skt{ }\edtext{\skt{वहति}}{\lemma{\skt{वहति}}\Afootnote{\skt{क्प्रवर्तते}\ A; \skt{भवति}\ K H; \skt{…ति}\ N\textsuperscript{ac}; \skt{वदति}\ N\textsuperscript{pc}}}\skt{~॥}\pend

\pstart \ledsidenote{\passref{1.14.36}}\skt{अथातिप्रवृत्ते }\edtext{\edtext{\edtext{\skt{लोध्र\-म\-धु\-क\-प्रि\-य\-ङ्गु\-प\-त्ता\-ङ्ग\-गै\-रि\-क\-र\-सा\-ञ्ज\-न\-श\-ङ्ख\-शु\-क्ति\-य\-व\-मा\-ष\-गो\-धू\-म\-स\-र्ज\-र\-स\-चू\-र्ण्णै\-र\-ना\-र्द्रै\-र्व्र\-ण\-मु\-ख\-म\-व\-चू\-र्ण्ण्या\-ङ्गु\-ल्य\-ग्रे\-णा\-व\-पी\-ड\-येत्}}{\lemma{\skt{लोध्र\appbrk{}म\appbrk{}धु\appbrk{}क\appbrk{}प्रि\appbrk{}य\appbrk{}ङ्गु\appbrk{}प\appbrk{}त्ता\appbrk{}ङ्ग\appbrk{}गै\appbrk{}रि\appbrk{}क\appbrk{}र\appbrk{}सा\appbrk{}ञ्ज\appbrk{}न\appbrk{}श\appbrk{}ङ्ख\appbrk{}शु\appbrk{}क्ति\appbrk{}य\appbrk{}व\appbrk{}मा\appbrk{}ष\appbrk{}गो\appbrk{}धू\appbrk{}म\appbrk{}स\appbrk{}र्ज\appbrk{}र\appbrk{}स\appbrk{}चू\appbrk{}र्ण्णैर्}}\Afootnote{K N H; \skt{रोध्र\appbrk{}म\appbrk{}धु\appbrk{}क\appbrk{}प्रि\appbrk{}य\appbrk{}ङ्गु\appbrk{}प\appbrk{}त्त\appbrk{}ङ्गु\appbrk{}प\appbrk{}त्त\appbrk{}ङ्ग\appbrk{}गै\appbrk{}रि\appbrk{}क\appbrk{}स\appbrk{}र्ज\appbrk{}र\appbrk{}स\appbrk{}र\appbrk{}सा\appbrk{}ञ्ज\appbrk{}न\appbrk{}शा\appbrk{}ल्म\appbrk{}ली\appbrk{}पु\appbrk{}ष्प\appbrk{}श\appbrk{}ङ्ख\appbrk{}शु\appbrk{}क्ति\appbrk{}मा\appbrk{}ष\appbrk{}य\appbrk{}व\appbrk{}गो\appbrk{}धू\appbrk{}म\appbrk{}चूर्णैः}\ A}}}{\lemma{\skt{अनार्द्रैर्}}\Afootnote{K N H; \skt{शनैः शनैर्}\ A}}\skt{~। साल\-स\-र्ज\-र्जु\-ना\-रि\-मे\-द\-ग्र\-न्थि\-ध\-व\-ध\-न्व\-न\-त्व\-ग्भिर्वा}}{\lemma{\skt{अव\appbrk{}चू\appbrk{}र्ण्ण्या\appbrk{}ङ्गु\appbrk{}ल्य\appbrk{}ग्रे\appbrk{}णा\appbrk{}व\appbrk{}पी\appbrk{}ड\appbrk{}येत्~। साल\appbrk{}स\appbrk{}र्ज\appbrk{}र्जु\appbrk{}ना\appbrk{}रि\appbrk{}मे\appbrk{}द\appbrk{}ग्र\appbrk{}न्थि\appbrk{}ध\appbrk{}व\appbrk{}ध\appbrk{}न्व\appbrk{}न\appbrk{}त्व\appbrk{}ग्भिर्वा}}\Afootnote{\skt{अव\appbrk{}चू\appbrk{}र्ण्या\appbrk{}ङ्गु\appbrk{}ल्य\appbrk{}ग्रे\appbrk{}णा\appbrk{}व\appbrk{}पी\appbrk{}ड\appbrk{}ये\appbrk{}त्सा\appbrk{}ल\appbrk{}स\appbrk{}र्जा\appbrk{}र्जु\appbrk{}ना\appbrk{}रि\appbrk{}मे\appbrk{}द\appbrk{}मे\appbrk{}ष\appbrk{}शृ\appbrk{}ङ्ग\appbrk{}ध\appbrk{}व\appbrk{}ध\appbrk{}न्व\appbrk{}न\appbrk{}त्व\appbrk{}ग्भिर्वा}\ A; \skt{अव\appbrk{}चू\appbrk{}र्ण्णा\appbrk{}ङ्गु\appbrk{}ल्य\appbrk{}ग्रे\appbrk{}णा\appbrk{}व\appbrk{}पी\appbrk{}ड\appbrk{}येत्~। साल\appbrk{}स\appbrk{}र्जा\appbrk{}र्जु\appbrk{}ना\appbrk{}रि\appbrk{}मे\appbrk{}द\appbrk{}ग्र\appbrk{}न्थि\appbrk{}ध\appbrk{}व\appbrk{}ध\appbrk{}न्व\appbrk{}न\appbrk{}त्वग्भि वा}\ K; \skt{अव\appbrk{}चू\appbrk{}र्ण्णा\appbrk{}ङ्गु\appbrk{}ल्य\appbrk{}ग्रे\appbrk{}णा\appbrk{}व\appbrk{}पी\appbrk{}ड\appbrk{}येत्~॥ साल\appbrk{}स\appbrk{}र्ज\appbrk{}र्जु\appbrk{}ना\appbrk{}रि\appbrk{}मे\appbrk{}द\appbrk{}ग्र\appbrk{}न्थि\appbrk{}ध\appbrk{}व\appbrk{}ध\appbrk{}न्व\appbrk{}न\appbrk{}त्व\appbrk{}ग्भिर्वा}\ N; \skt{अव\appbrk{}चू\appbrk{}र्ण्या\appbrk{}ङ्गु\appbrk{}ल्य\appbrk{}ग्रे\appbrk{}णा\appbrk{}व\appbrk{}पी\appbrk{}ड\appbrk{}येत्~। शाल\appbrk{}स\appbrk{}र्ज्जा\appbrk{}र्ज्जु\appbrk{}ना\appbrk{}रि\appbrk{}मे\appbrk{}द\appbrk{}ग्र\appbrk{}न्थि\appbrk{}ध\appbrk{}च\appbrk{}ध\appbrk{}न्व\appbrk{}न\appbrk{}त्व\appbrk{}ग्भिर्व्वा}\ H}}\skt{ }\edtext{\skt{चूर्ण्णीकृताभिः}}{\lemma{\skt{चूर्ण्णीकृताभिः}}\Afootnote{K N H; \skt{चूर्णिताभिः}\ A}}\skt{ }\edtext{\skt{क्षौमेण}}{\lemma{\skt{क्षौमेण}}\Afootnote{A K H; \skt{क्षौमेणा}\ N}}\skt{ }\edtext{\skt{वध्यासितेन}}{\lemma{\skt{वध्यासितेन}}\Afootnote{K N H; \skt{वा ध्मापितेन}\ A}}\skt{ }\edtext{\skt{समुद्रफेनेन}}{\lemma{\skt{समुद्रफेनेन}}\Afootnote{K N H; \skt{समुद्रफेन}\ A}}\skt{ लाक्षाचूर्ण्णैर्वा }\edtext{\edtext{\skt{यथो\-क्तै\-र्ब\-न्ध\-न\-द्र\-व्यै\-र्गाढं}}{\lemma{\skt{यथोक्तैर्}}\Afootnote{K N H; \skt{यथोक्त}\ A}}}{\lemma{\skt{बन्धनद्रव्यैर्}}\Afootnote{K N H; \skt{र्व्रण\appbrk{}ब\appbrk{}न्ध\appbrk{}न\appbrk{}द्र\appbrk{}व्यैर्}\ A}}\skt{ }\edtext{\skt{बध्नीयात्~। व्यधानन्तरं}}{\lemma{\skt{बध्नीयात्~। व्यधानन्तरं}}\Afootnote{N H; \skt{बध्नीया}\ A; \skt{बन्धीयात्~। व्यधानन्तरं}\ K}}\skt{ }\edtext{\edtext{\skt{पुनर्व्यधयेत्}}{\lemma{\skt{पुनर्}}\Afootnote{K N H; \emph{om.}\ A}}\skt{~। शीता\-च्छा\-द\-न\-भो\-ज\-ना\-गा\-र\-प\-रि\-षेकैः}}{\lemma{\skt{व्यधयेत्~। शीता\appbrk{}च्छा\appbrk{}द\appbrk{}न\appbrk{}भो\appbrk{}ज\appbrk{}ना\appbrk{}गा\appbrk{}र\appbrk{}प\appbrk{}रि\appbrk{}षेकैः}}\Afootnote{K N H; \skt{त्शीता\appbrk{}च्छा\appbrk{}द\appbrk{}न\appbrk{}भो\appbrk{}ज\appbrk{}ना\appbrk{}गारैः शीतैः प्रदे\appbrk{}ह\appbrk{}प\appbrk{}रि\appbrk{}षे\appbrk{}कैश्}\ A}}\skt{~। }\edtext{\edtext{\skt{शीतैरालेपैः}}{\lemma{\skt{शीतैर्}}\Afootnote{K N H; \emph{om.}\ A}}}{\lemma{\skt{आलेपैः}}\Afootnote{K N H; \emph{om.}\ A}}\skt{ }\edtext{\edtext{\skt{प्रदे\-है\-र्वो\-प\-च\-रे\-द\-ग्निना}}{\lemma{\skt{प्रदेहैर्}}\Afootnote{N H; \emph{om.}\ A; \skt{प्रदहैर्}\ K}}\skt{ वा}}{\lemma{\skt{वोप\appbrk{}च\appbrk{}रे\appbrk{}द\appbrk{}ग्निना वा}}\Afootnote{\skt{चोपचरेत्, क्षारेणाग्निना वा}\ A; \skt{वाप\appbrk{}च\appbrk{}रे\appbrk{}द\appbrk{}ग्निना वा}\ K N H}}\skt{ }\edtext{\skt{दहेद्यथोक्तं}}{\lemma{\skt{दहेद्यथोक्तं}}\Afootnote{K N H; \skt{दहेद्यथोक्तं}\ A}}\skt{ }\edtext{\skt{काकोल्यादिक्वाथं वा}}{\lemma{\skt{काकोल्यादिक्वाथं वा}}\Afootnote{K N H; \skt{व्यधादनन्तरं तामे\appbrk{}वा\appbrk{}ति\appbrk{}प्र\appbrk{}वृत्तां सिरां विध्ये\appbrk{}त्का\appbrk{}को\appbrk{}ल्या\appbrk{}दि\appbrk{}क्वाथं वा}\ A}}\skt{ }\edtext{\skt{शर्क\-रा\-म\-धु\-म\-धुरं}}{\lemma{\skt{शर्क\appbrk{}रा\appbrk{}म\appbrk{}धु\appbrk{}म\appbrk{}धुरं}}\Afootnote{A K H; \skt{शर्क\appbrk{}रा\appbrk{}म\appbrk{}धु\appbrk{}म\appbrk{}धुर}\ N}}\skt{ }\edtext{\skt{पाययेत्~। एण\-ह\-रि\-णो\-र\-भ्रं\-म\-हि\-ष\-श\-श\-व\-रा\-हाणां वा}}{\lemma{\skt{पाययेत्~। एण\appbrk{}ह\appbrk{}रि\appbrk{}णो\appbrk{}र\appbrk{}भ्रं\appbrk{}म\appbrk{}हि\appbrk{}ष\appbrk{}श\appbrk{}श\appbrk{}व\appbrk{}रा\appbrk{}हाणां वा}}\Afootnote{\skt{पाय\appbrk{}ये\appbrk{}ते\appbrk{}ण\appbrk{}ह\appbrk{}रि\appbrk{}णो\appbrk{}र\appbrk{}भ्र\appbrk{}श\appbrk{}श\appbrk{}म\appbrk{}हि\appbrk{}ष\appbrk{}व\appbrk{}रा\appbrk{}हाणां वा}\ A; \skt{पाचयेत्~। एण\appbrk{}ह\appbrk{}रि\appbrk{}णो\appbrk{}रभ्रं महि\appbrk{}ष\appbrk{}श\appbrk{}श\appbrk{}व\appbrk{}रा\appbrk{}हांणां वा}\ K; \skt{पाययेत्~। एण\appbrk{}ह\appbrk{}रि\appbrk{}णो\appbrk{}र\appbrk{}भ्र\appbrk{}म\appbrk{}हि\appbrk{}ष\appbrk{}श\appbrk{}श\appbrk{}व\appbrk{}रा\appbrk{}हानां वा}\ N; \skt{पाय\appbrk{}ये\appbrk{}ते\appbrk{}ण\appbrk{}ह\appbrk{}रि\appbrk{}णो\appbrk{}र\appbrk{}भ्र\appbrk{}म\appbrk{}हि\appbrk{}ष\appbrk{}श\appbrk{}श\appbrk{}व\appbrk{}रा\appbrk{}हाणां वा}\ H}}\skt{ रुधिरं }\edtext{\skt{क्षीरयूषं}}{\lemma{\skt{क्षीरयूषं}}\Afootnote{K H; \skt{क्षीरयूष}\ A N}}\skt{ }\edtext{\skt{रसैश्चाश्नीयात्~। उपद्रवांश्च}}{\lemma{\skt{चाश्नीयात्~। उपद्रवांश्च}}\Afootnote{K N H; \skt{सुस्नि\appbrk{}ग्धै\appbrk{}श्चा\appbrk{}श्नी\appbrk{}या\appbrk{}तु\appbrk{}प\appbrk{}द्र\appbrk{}वांश्च}\ A}}\skt{ }\edtext{\skt{यथो\-क्ता\-नु\-प\-च\-रेत्}}{\lemma{\skt{यथो\appbrk{}क्ता\appbrk{}नु\appbrk{}प\appbrk{}च\appbrk{}रेत्}}\Afootnote{K N H; \skt{यथा\appbrk{}स्व\appbrk{}मु\appbrk{}प\appbrk{}च\appbrk{}रेत्}\ A}}\skt{~॥}\pend

\pstart\versepar \ledsidenote{\passref{1.14.37}}\skt{भवति चात्र~॥ }\pend
\pstart\versepar \skt{धातुक्षयात्स्रुते रक्ते मन्दः सञ्जायते }\edtext{\skt{नलः}}{\lemma{\skt{नलः}}\Afootnote{N; \skt{ऽनलः}\ A K H}}\skt{~। }\pend
\pstart\versepar \skt{पवनश्च परं कोपं याति तस्मात्प्रयत्नतः~।}\edtext{}{\lemma{\skt{भवति चात्र~॥ धातुक्षयात्स्रुते}}\Afootnote{\skt{भवन्ति चात्र धातुक्षयात्स्रुते}\ A; \skt{धातुक्षयात्स्रुते}\ K H; \skt{भव~॥ धातुक्षयात्स्रुते}\ N}}\pend

\pstart\versepar \ledsidenote{\passref{1.14.38}}\edtext{\skt{तन्ना\-ति\-शी\-तै\-र्ल\-घुभिः}}{\lemma{\skt{तन्नातिशीतैर्}}\Afootnote{A K H; \skt{तन्नातिशी}\ N}}\skt{ स्निग्धैः }\edtext{\skt{शोणितवर्धनैः}}{\lemma{\skt{शोणितवर्धनैः}}\Afootnote{A K N; \skt{शोणितेवर्द्धनैः}\ H}}\skt{~। }\pend
\pstart\versepar \skt{}\edtext{\edtext{\skt{ईष\-द\-म्लै\-र\-न\-म्लैर्वा}}{\lemma{\skt{ईषदम्लैर्}}\Afootnote{A K H; \skt{ईषदब्लैर्}\ N}}}{\lemma{\skt{अनम्लैर्वा}}\Afootnote{A K H; \skt{अनब्लैर्वा}\ N}}\skt{ भोजनैः समुपाचरेत्~॥}\pend

\pstart\versepar \ledsidenote{\passref{1.14.39}}\skt{चतुर्विधं यदेतद्धि }\edtext{\skt{रुधिरस्य}}{\lemma{\skt{रुधिरस्य}}\Afootnote{A K N H\textsuperscript{pc}; \skt{रुधिरम्य}\ H\textsuperscript{ac}}}\skt{ निवारणं~। }\pend
\pstart\versepar \skt{}\edtext{\skt{सन्धानं}}{\lemma{\skt{सन्धानं}}\Afootnote{A K N; \skt{सन्धान}\ H}}\skt{ स्कन्दनं चैव पाचनं दहनं तथा~॥}\pend

\pstart\versepar \ledsidenote{\passref{1.14.40}}\edtext{\skt{व्रणं}}{\lemma{\skt{व्रणं}}\Afootnote{A K H; \skt{व्रण}\ N}}\skt{ }\edtext{\skt{कषायः}}{\lemma{\skt{कषायः}}\Afootnote{A K N\textsuperscript{pc} H; \skt{कषाय}\ N\textsuperscript{ac}}}\skt{ सन्धत्ते रक्तं स्कन्दयते हिमं~। }\pend
\pstart\versepar \skt{तथा सम्पाचयेद्भस्म दाहः }\edtext{\skt{सङ्कोचयेत्सिरां}}{\lemma{\skt{सङ्कोचयेत्सिरां}}\Afootnote{K N H; \skt{सङ्कोचयेत्सिराः}\ A}}\skt{~॥}\pend

\pstart\versepar \ledsidenote{\passref{1.14.41}}\skt{अस्कन्दमाने रुधिरे सन्धानानि प्रयोजयेत्~। }\pend
\pstart\versepar \skt{}\edtext{\skt{सन्धा\-नै\-र्भ्रा\-श्य\-माने तु}}{\lemma{\skt{भ्राश्यमाने तु}}\Afootnote{\skt{भ्रश्यमाने तु}\ A N H K}}\skt{ पाचनैः समुपाचरेत्~॥}\edtext{}{\lemma{\skt{प्रयोजयेत्~। सन्धानैर्}}\Afootnote{K N H; \skt{प्रयोजयेत्~। सन्धाने}\ A}}\pend

\pstart\versepar \ledsidenote{\passref{1.14.42}}\edtext{\skt{कल्पै\-रे\-भि\-स्त्रि\-भि\-र्वैद्यः}}{\lemma{\skt{एभिस्}}\Afootnote{K N H; \skt{एतैस्}\ A}}\skt{ }\edtext{\skt{प्रयतेत}}{\lemma{\skt{प्रयतेत}}\Afootnote{A K N H\textsuperscript{ac}; \skt{प्रयथंतेत}\ H\textsuperscript{pc}}}\skt{ }\edtext{\skt{यथाविधिः}}{\lemma{\skt{यथाविधिः}}\Afootnote{K N H; \skt{यथाविधि}\ A}}\skt{~। }\pend
\pstart\versepar \skt{}\edtext{\skt{असिद्धिमत्सु}}{\lemma{\skt{असिद्धिमत्सु}}\Afootnote{A K H; \skt{असिद्धिम…}\ N}}\skt{ }\edtext{\skt{चैतेषु}}{\lemma{\skt{चैतेषु}}\Afootnote{A K H; \skt{… षु}\ N}}\skt{ }\edtext{\skt{दाहपरम}}{\lemma{\skt{दाहपरम}}\Afootnote{K N; \skt{दाहः परम}\ A H}}\skt{ इष्यते~।}\pend

\pstart\versepar \ledsidenote{\passref{1.14.43}}\edtext{\skt{सशेषदोषे}}{\lemma{\skt{सशेषदोषे}}\Afootnote{K N H; \skt{शेषदोषे}\ A}}\skt{ }\edtext{\skt{रुधिरे न}}{\lemma{\skt{रुधिरे न}}\Afootnote{K N H; \skt{यतो रक्ते न}\ A}}\skt{ }\edtext{\edtext{\skt{व्याधि\-र\-ति\-व\-र्त्तते}}{\lemma{\skt{व्याधिर्}}\Afootnote{K N H; \skt{व्याधि}\ A}}}{\lemma{\skt{अतिवर्त्तते}}\Afootnote{K N H; \skt{रतिवर्तते}\ A}}\skt{~। }\pend
\pstart\versepar \skt{}\edtext{\skt{न शेषं}}{\lemma{\skt{न शेषं}}\Afootnote{K; \skt{सावशेषे}\ A; \skt{सशेषं}\ N H}}\skt{ }\edtext{\skt{स्थापयेत्तस्मान्न च}}{\lemma{\skt{स्थापयेत्तस्मान्न च}}\Afootnote{K N H; \skt{ततः स्थेयं न तु}\ A}}\skt{ }\edtext{\skt{कुर्यादतिक्रियां}}{\lemma{\skt{कुर्यादतिक्रियां}}\Afootnote{K N H; \skt{कुर्या\appbrk{}द\appbrk{}ति\appbrk{}क्र\appbrk{}मम्}\ A}}\skt{~।}\pend

\pstart\versepar \ledsidenote{\passref{1.14.44}}\skt{देहस्य }\edtext{\skt{रुधिरं}}{\lemma{\skt{रुधिरं}}\Afootnote{A K N H\textsuperscript{ac}; \skt{ञ्रुधिरं}\ H\textsuperscript{pc}}}\skt{ मूलं रुधिरेणैव धार्यते~। }\pend
\pstart\versepar \skt{}\edtext{\skt{तस्माद्रक्षेद्धि}}{\lemma{\skt{तस्माद्रक्षेद्धि}}\Afootnote{K N H; \skt{तस्मा}\ A}}\skt{ }\edtext{\skt{रुधिरं}}{\lemma{\skt{रुधिरं}}\Afootnote{K\textsuperscript{pc} N H; \emph{om.}\ A; \skt{रुरं}\ K\textsuperscript{ac}}}\skt{ }\edtext{\skt{रुधिरं}}{\lemma{\skt{रुधिरं}}\Afootnote{K N H; \emph{om.}\ A}}\skt{ }\edtext{\skt{जीव}}{\lemma{\skt{जीव}}\Afootnote{K N H; \emph{om.}\ A}}\skt{ }\edtext{\skt{उच्यते}}{\lemma{\skt{उच्यते}}\Afootnote{K N H; \skt{द्यत्नेन संरक्ष्यं रक्तं जीव इति स्थितिः}\ A}}\skt{~॥}\pend

\pstart\versepar \ledsidenote{\passref{1.14.45}}\edtext{\skt{स्रुतरक्तस्य}}{\lemma{\skt{स्रुतरक्तस्य}}\Afootnote{A K; \skt{च्युतरक्तस्य}\ N; \skt{श्रुतरक्तस्य}\ H}}\skt{ सेकाद्यैः शीतैः प्रकुपितेऽनिले }\pend
\pstart\versepar \skt{शोफं }\edtext{\skt{सतोदं}}{\lemma{\skt{सतोदं}}\Afootnote{A K H; \skt{सतोद}\ N}}\skt{ कोष्णेन सर्पिषा }\edtext{\skt{परि\-षे\-च\-ये\-दिति}}{\lemma{\skt{परि\appbrk{}षे\appbrk{}च\appbrk{}ये\appbrk{}दिति}}\Afootnote{\skt{परिषेचत्}\ A; \skt{परि\appbrk{}षे\appbrk{}च\appbrk{}ये\appbrk{}दिति~॥ १४~॥ {\symfont ❈}}\ K; \skt{परि\appbrk{}से\appbrk{}च\appbrk{}ये\appbrk{}दिति~॥ १४}\ N; \skt{परि\appbrk{}षे\appbrk{}च\appbrk{}ये\appbrk{}दिति~॥ डोह्व}\ H}}\skt{~॥ १४~॥}\pend

\endnumbering
\beginnumbering
\eledsection*{Adhyāya 15}
\pstart \ledsidenote{\passref{1.15.1}}\skt{अथातो }\edtext{\skt{दोष\-धा\-तु\-म\-ल\-क्ष\-य\-वृद्धिं}}{\lemma{\skt{दोष\appbrk{}धा\appbrk{}तु\appbrk{}म\appbrk{}ल\appbrk{}क्ष\appbrk{}य\appbrk{}वृद्धिं}}\Afootnote{K N H; \skt{दोष\appbrk{}धा\appbrk{}तु\appbrk{}म\appbrk{}ल\appbrk{}क्ष\appbrk{}य\appbrk{}वृद्धि}\ A}}\skt{ }\edtext{\skt{व्याख्यास्यामः}}{\lemma{\skt{व्याख्यास्यामः}}\Afootnote{H; \skt{विज्ञा\appbrk{}नी\appbrk{}य\appbrk{}म\appbrk{}ध्यायं व्याख्यास्यामः}\ A; \emph{om.}\ K; \skt{अध्यायं व्याख्यास्यामः}\ N}}\skt{~॥}\pend

\pstart \ledsidenote{\passref{1.15.3}}\skt{दोष\-धा\-तु\-म\-ल\-मूलं हि शरीरं~। }\edtext{\edtext{\edtext{\skt{तस्मा\-त्फ\-ल\-ल\-क्ष\-ण\-मे\-ते\-षा\-मु\-प\-धा\-र\-यश्च}}{\lemma{\skt{तस्मा\appbrk{}त्फ\appbrk{}ल\appbrk{}ल\appbrk{}क्ष\appbrk{}ण\appbrk{}मे\appbrk{}ते\appbrk{}षाम्}}\Afootnote{K; \skt{तस्मादेतेषां}\ A; \skt{तस्मा\appbrk{}त्फ\appbrk{}ल\appbrk{}ल\appbrk{}क्ष\appbrk{}ण\appbrk{}मे\appbrk{}ते\appbrk{}षाम्}\ N H}}}{\lemma{\skt{उपधारयश्}}\Afootnote{K; \skt{लक्ष\appbrk{}ण\appbrk{}मु\appbrk{}च्य\appbrk{}मा\appbrk{}न\appbrk{}मु\appbrk{}प\appbrk{}धा\appbrk{}रय}\ A; \skt{उपधारय}\ N H}}}{\lemma{\skt{च}}\Afootnote{K; \emph{om.}\ A; \skt{स्व}\ N H}}\skt{~।}\pend

\pstart \ledsidenote{\passref{1.15.4.1}}\edtext{\skt{तत्र}}{\lemma{\skt{तत्र}}\Afootnote{K N H; \skt{तत्र प्र}\ A}}\skt{ }\edtext{\skt{स्पन्द\-नो\-द्व\-ह\-न\-पू\-र\-ण\-वि\-वे\-क\-धा\-र\-ण\-ल\-क्षणो}}{\lemma{\skt{स्पन्द\appbrk{}नो\appbrk{}द्व\appbrk{}ह\appbrk{}न\appbrk{}पू\appbrk{}र\appbrk{}ण\appbrk{}वि\appbrk{}वे\appbrk{}क\appbrk{}धा\appbrk{}र\appbrk{}ण\appbrk{}ल\appbrk{}क्षणो}}\Afootnote{A K N; \skt{स्पन्द\appbrk{}मो\appbrk{}द्व\appbrk{}ह\appbrk{}न\appbrk{}पू\appbrk{}र\appbrk{}ण\appbrk{}वि\appbrk{}वे\appbrk{}क\appbrk{}धा\appbrk{}र\appbrk{}ण\appbrk{}ल\appbrk{}क्षणो}\ H}}\skt{ वायुः पञ्चधा प्रविभक्तः }\edtext{\skt{शरी\-र\-न्त\-न्त्र\-यति}}{\lemma{\skt{शरी\appbrk{}र\appbrk{}न्त\appbrk{}न्त्र\appbrk{}यति}}\Afootnote{K N; \skt{शरीरं धारयति (१)}\ A; \skt{शरीरन्तन्त्रय्ति~। …}\ H}}\skt{~।}\pend

\pstart \ledsidenote{\passref{1.15.4.2}}\edtext{\skt{रागः पक्ति\-स्ते\-ज\-ऊ\-ष्म\-कृ\-त्पित्तं}}{\lemma{\skt{रागः पक्ति\appbrk{}स्ते\appbrk{}ज\appbrk{}ऊ\appbrk{}ष्म\appbrk{}कृ\appbrk{}त्पित्तं}}\Afootnote{K H; \skt{राग\appbrk{}प\appbrk{}क्त्यो\appbrk{}ज\appbrk{}स्ते\appbrk{}जो\appbrk{}मे\appbrk{}धो\appbrk{}ष्म\appbrk{}कृ\appbrk{}त्पित्तं पञ्चधा प्रवि\appbrk{}भ\appbrk{}क्त\appbrk{}म\appbrk{}ग्नि\appbrk{}क\appbrk{}र्म\appbrk{}णाऽ\appbrk{}नु\appbrk{}ग्रहं करोति; (२)}\ A; \skt{रागः पक्तिस्तेज ऊष्माकृत्पित्तं}\ N}}\skt{~।}\pend

\pstart \ledsidenote{\passref{1.15.4.3}}\edtext{\skt{सन्धि\-सं\-श्ले\-ष\-ण\-स्ने\-ह\-न\-रो\-प\-ण\-बृं\-हणः}}{\lemma{\skt{सन्धि\appbrk{}सं\appbrk{}श्ले\appbrk{}ष\appbrk{}ण\appbrk{}स्ने\appbrk{}ह\appbrk{}न\appbrk{}रो\appbrk{}प\appbrk{}ण\appbrk{}बृं\appbrk{}हणः}}\Afootnote{K N H; \skt{सन्धि\appbrk{}सं\appbrk{}श्ले\appbrk{}ष\appbrk{}ण\appbrk{}स्ने\appbrk{}ह\appbrk{}न\appbrk{}रो\appbrk{}प\appbrk{}ण\appbrk{}पू\appbrk{}रण}\ A}}\skt{ }\edtext{\skt{श्लेष्मा}}{\lemma{\skt{श्लेष्मा}}\Afootnote{K N H; \skt{बल\appbrk{}स्थै\appbrk{}र्य\appbrk{}कृ\appbrk{}च्छ्लेष्मा पञ्चधा प्रविभक्त उद\appbrk{}क\appbrk{}क\appbrk{}र्म\appbrk{}णाऽ\appbrk{}नु\appbrk{}ग्रहं करोति}\ A}}\skt{~।}\pend

\pstart \ledsidenote{\passref{1.15.5.1}}\edtext{\skt{रसः}}{\lemma{\skt{रसः}}\Afootnote{K N\textsuperscript{pc}; \skt{रसस्तुष्टिं प्रीणनं रक्तपुष्टिं च करोति, रक्तं वर्ण}\ A; \skt{रस}\ N\textsuperscript{ac}; \skt{रसः प्रीणयति~। रक्तं}\ H}}\skt{ }\edtext{\skt{प्रीणयति}}{\lemma{\skt{प्रीणयति}}\Afootnote{K N; \skt{प्रसादं मां}\ A; \emph{om.}\ H}}\skt{~। }\edtext{\skt{रक्तञ्जीवयति}}{\lemma{\skt{रक्तञ्जीवयति}}\Afootnote{K N; \skt{सपुष्टिं जीवयति च}\ A; \skt{जीवयति}\ H}}\skt{~। }\edtext{\skt{मांसं}}{\lemma{\skt{मांसं}}\Afootnote{K N H; \skt{मांसं शरीरपुष्टिं}\ A}}\skt{ }\edtext{\skt{लेपयति}}{\lemma{\skt{लेपयति}}\Afootnote{K N H; \skt{मेदसश्च}\ A}}\skt{ मेदः }\edtext{\skt{स्नेहयति}}{\lemma{\skt{स्नेहयति}}\Afootnote{K N; \skt{स्नेहस्वेदौ दृढत्वं पुष्टिमस्थ्नां च, अस्थीनि देहधारणं मज्ज्ञः पुष्टिं च, मज्जा स्नेहं बलं शुक्रपुष्टिं पूरणमस्थ्नां च करोति}\ A; \skt{स्नेहयति~। अस्थि धारयति~। मज्जा पूर}\ H}}\skt{~। }\edtext{\skt{अस्थि}}{\lemma{\skt{अस्थि}}\Afootnote{K N; \skt{शुक्रं}\ A; \emph{om.}\ H}}\skt{ }\edtext{\skt{धारयति}}{\lemma{\skt{धारयति}}\Afootnote{K N; \skt{धैर्यं च्यव}\ A; \emph{om.}\ H}}\skt{~। }\edtext{\skt{मज्जा}}{\lemma{\skt{मज्जा}}\Afootnote{K N; \skt{नं प्री}\ A; \emph{om.}\ H}}\skt{ }\edtext{\skt{पूरयति}}{\lemma{\skt{पूरयति}}\Afootnote{K N; \skt{तिं देहबलं हर्षं}\ A; \skt{यति}\ H}}\skt{~। }\edtext{\skt{बीजा\-र्थ\-ह\-र्ष\-कृ\-च्छुक्रं}}{\lemma{\skt{बीजा\appbrk{}र्थ\appbrk{}ह\appbrk{}र्ष\appbrk{}कृ\appbrk{}च्छुक्रं}}\Afootnote{K N; \skt{बीजार्थं च; (१)}\ A; \skt{बीजा\appbrk{}र्थ\appbrk{}ह\appbrk{}र्ष\appbrk{}कृ\appbrk{}च्छुक्रं}\ H}}\skt{ }\edtext{\skt{क्लेदयति}}{\lemma{\skt{क्लेदयति}}\Afootnote{K N H; \emph{om.}\ A}}\skt{~।}\pend

\pstart \ledsidenote{\passref{1.15.5.2}}\edtext{\skt{बस्ति\-क्ले\-द\-कृ\-न्मूत्रं}}{\lemma{\skt{बस्ति\appbrk{}क्ले\appbrk{}द\appbrk{}कृ\appbrk{}न्मूत्रं}}\Afootnote{K\textsuperscript{pc} N; \skt{पुरी\appbrk{}ष\appbrk{}मु\appbrk{}प\appbrk{}स्तम्भं वाय्वग्निधारणं च, बस्ति\appbrk{}पू\appbrk{}र\appbrk{}ण\appbrk{}वि\appbrk{}क्ले\appbrk{}द\appbrk{}कृ\appbrk{}न्मूत्रं}\ A; \skt{बस्तिक्ले कृन्मूत्रं}\ K\textsuperscript{ac}; \skt{बस्ति\appbrk{}पू\appbrk{}र\appbrk{}ण\appbrk{}कृ\appbrk{}त्मूत्रं}\ H}}\skt{~। }\edtext{\skt{प्राण\-वा\-य्व\-ग्नि\-धा\-र\-णा\-व\-ष्ट\-म्भ\-कृ\-त्पु\-रीशं}}{\lemma{\skt{प्राण\appbrk{}वा\appbrk{}य्व\appbrk{}ग्नि\appbrk{}धा\appbrk{}र\appbrk{}णा\appbrk{}व\appbrk{}ष्ट\appbrk{}म्भ\appbrk{}कृ\appbrk{}त्पु\appbrk{}रीशं}}\Afootnote{K; \emph{om.}\ A; \skt{प्राण\appbrk{}वा\appbrk{}य्व\appbrk{}ग्नि\appbrk{}धा\appbrk{}र\appbrk{}णा\appbrk{}व\appbrk{}ष्ट\appbrk{}म्भ\appbrk{}कृ\appbrk{}त्पु\appbrk{}रीषं}\ N H}}\skt{~। }\edtext{\skt{स्वेदः}}{\lemma{\skt{स्वेदः}}\Afootnote{A K N\textsuperscript{pc} H; \skt{स्वेद}\ N\textsuperscript{ac}}}\skt{ }\edtext{\skt{क्लेदयति}}{\lemma{\skt{क्लेदयति}}\Afootnote{K N H; \skt{क्लेद\appbrk{}त्व\appbrk{}क्सौ\appbrk{}कु\appbrk{}मा\appbrk{}र्य\appbrk{}कृत्; (२)}\ A}}\skt{~।}\pend

\pstart \ledsidenote{\passref{1.15.5.3}}\edtext{\skt{गर्भ\-ल\-क्ष\-ण\-मा\-र्त्तवं}}{\lemma{\skt{गर्भ\appbrk{}ल\appbrk{}क्ष\appbrk{}ण\appbrk{}मा\appbrk{}र्त्तवं}}\Afootnote{K; \skt{रक्त\appbrk{}ल\appbrk{}क्ष\appbrk{}ण\appbrk{}मा\appbrk{}र्तवं}\ A; \emph{om.}\ N H}}\skt{~। }\edtext{\skt{स्तन्यं}}{\lemma{\skt{स्तन्यं}}\Afootnote{K; \skt{गर्भकृच्च, गर्भो गर्भलक्षणं, स्तन्यं}\ A; \skt{गर्भ\appbrk{}ल\appbrk{}क्ष\appbrk{}ण\appbrk{}मा\appbrk{}र्त्त\appbrk{}व\appbrk{}म्स्तन्यं}\ N H}}\skt{ }\edtext{\edtext{\skt{स्तना\-पी\-न\-ज\-न\-न\-जी\-व\-न\-मिति}}{\lemma{\skt{स्तना\appbrk{}पी\appbrk{}न\appbrk{}ज\appbrk{}न\appbrk{}न\appbrk{}जी\appbrk{}व\appbrk{}नम्}}\Afootnote{K N; \skt{स्तन\appbrk{}यो\appbrk{}रा\appbrk{}पी\appbrk{}न\appbrk{}त्व\appbrk{}ज\appbrk{}ननं जीवनं}\ A; \skt{स्तना\appbrk{}पी\appbrk{}न\appbrk{}ज\appbrk{}नन जीवन}\ H}}}{\lemma{\skt{इति}}\Afootnote{K N; \skt{चेति}\ A; \skt{मिति}\ H}}\skt{~।}\pend

\pstart \ledsidenote{\passref{1.15.6}}\skt{तत्र }\edtext{\skt{विधि\-व\-त्प\-रि\-र\-क्षणं}}{\lemma{\skt{विधि\appbrk{}व\appbrk{}त्प\appbrk{}रि\appbrk{}र\appbrk{}क्षणं}}\Afootnote{A K H; \skt{विधि\appbrk{}व\appbrk{}त्प\appbrk{}री\appbrk{}क्षणं}\ N}}\skt{ }\edtext{\skt{कुर्वीत}}{\lemma{\skt{कुर्वीत}}\Afootnote{A K H; \skt{कुर्वीतः}\ N}}\skt{~॥}\pend

\pstart \ledsidenote{\passref{1.15.7}}\edtext{\skt{अतः}}{\lemma{\skt{अतः}}\Afootnote{K N H; \skt{अत}\ A}}\skt{ }\edtext{\skt{सर्वेषां}}{\lemma{\skt{सर्वेषां}}\Afootnote{K N H; \skt{ऊर्ध्वमेषां}\ A}}\skt{ }\edtext{\skt{क्षयलक्षणं}}{\lemma{\skt{क्षयलक्षणं}}\Afootnote{K N H; \skt{क्षीणलक्षणं}\ A}}\skt{ }\edtext{\skt{व्याख्यास्यामः}}{\lemma{\skt{व्याख्यास्यामः}}\Afootnote{K N H; \emph{om.}\ A}}\skt{~। }\edtext{\skt{तत्र}}{\lemma{\skt{तत्र}}\Afootnote{K N H; \emph{om.}\ A}}\skt{ }\edtext{\skt{वातक्षये}}{\lemma{\skt{वातक्षये}}\Afootnote{K H; \skt{वक्ष्यामः~। तत्र, वातक्षये}\ A; \skt{वात…}\ N}}\skt{ }\edtext{\skt{मन्दचेष्टता}}{\lemma{\skt{मन्दचेष्टता}}\Afootnote{A K H; \skt{…}\ N}}\skt{ }\edtext{\skt{अल्प\-वा\-क्त्व\-म\-ल्प\-प्र\-हर्षो}}{\lemma{\skt{अल्प\appbrk{}वा\appbrk{}क्त्व\appbrk{}म\appbrk{}ल्प\appbrk{}प्र\appbrk{}हर्षो}}\Afootnote{K; \skt{ल्पवा\appbrk{}क्त्व\appbrk{}म\appbrk{}प्र\appbrk{}हर्षो}\ A; \skt{… अल्पप्रहर्षो}\ N; \skt{अल्प\appbrk{}वा\appbrk{}क्त्व\appbrk{}म\appbrk{}ल्प\appbrk{}प्र\appbrk{}हर्षो}\ H}}\skt{ मूढसञ्ज्ञता च~। पित्तक्षये मन्दो\-ष्मा\-ग्नि\-ता\-नि\-ष्प्र\-भता च~। श्लेष्मक्षये }\edtext{\skt{रूक्ष\-ता\-न्त\-र्दा\-ह\-आ\-मा\-श\-ये\-त\-रा\-श\-य\-शू\-न्य\-ता\-शि\-र\-सश्च}}{\lemma{\skt{रूक्ष\appbrk{}ता\appbrk{}न्त\appbrk{}र्दा\appbrk{}ह\appbrk{}आ\appbrk{}मा\appbrk{}श\appbrk{}ये\appbrk{}त\appbrk{}रा\appbrk{}श\appbrk{}य\appbrk{}शू\appbrk{}न्य\appbrk{}ता\appbrk{}शि\appbrk{}र\appbrk{}सश्}}\Afootnote{K; \skt{रूक्षताऽन्तर्दाह आमा\appbrk{}श\appbrk{}ये\appbrk{}त\appbrk{}र\appbrk{}श्ले\appbrk{}ष्मा\appbrk{}श\appbrk{}य\appbrk{}शू\appbrk{}न्यता सन्धिशैथिल्यं (तृष्णा दौर्बल्यं प्रजागरणं)}\ A; \skt{रूक्षान्तर्दाह आमा\appbrk{}श\appbrk{}ये\appbrk{}त\appbrk{}रा\appbrk{}स\appbrk{}य\appbrk{}शू\appbrk{}न्यता शिरसश्}\ N; \skt{रूक्ष\appbrk{}ता\appbrk{}न्त\appbrk{}र्द्दा\appbrk{}ह\appbrk{}आ\appbrk{}मा\appbrk{}स\appbrk{}ये\appbrk{}त\appbrk{}रा\appbrk{}स\appbrk{}य\appbrk{}शू\appbrk{}न्यता शिरसश्}\ H}}\skt{~।}\pend

\pstart \ledsidenote{\passref{1.15.8}}\skt{तत्र }\edtext{\edtext{\skt{स्वयो\-नि\-व\-र्ध\-ना\-न्येव}}{\lemma{\skt{स्वयो\appbrk{}नि\appbrk{}व\appbrk{}र्ध\appbrk{}नान्य्}}\Afootnote{K N H; \skt{स्वयोनिवर्धना}\ A}}}{\lemma{\skt{एव}}\Afootnote{K N H; \skt{न्येव}\ A}}\skt{ प्रतीकारः~॥}\pend

\pstart \ledsidenote{\passref{1.15.9}}\skt{रसक्षये }\edtext{\skt{हृद\-य\-पी\-डा\-कम्पः}}{\lemma{\skt{हृद\appbrk{}य\appbrk{}पी\appbrk{}डा\appbrk{}कम्पः}}\Afootnote{K N H; \skt{हृत्पीडाकम्प}\ A}}\skt{ }\edtext{\skt{शोषः}}{\lemma{\skt{शोषः}}\Afootnote{K N H; \emph{om.}\ A}}\skt{ }\edtext{\skt{शून्यता}}{\lemma{\skt{शून्यता}}\Afootnote{A K N; \skt{शून्याता}\ H}}\skt{ }\edtext{\skt{तृष्णा च}}{\lemma{\skt{तृष्णा च}}\Afootnote{K N H; \skt{स्तृष्णा च}\ A}}\skt{~। शोणितक्षये }\edtext{\edtext{\skt{त्वक्पा\-रु\-ष्य\-म\-म्ल\-शी\-त\-प्रा\-र्थना}}{\lemma{\skt{त्वक्पारुष्यम्}}\Afootnote{K N H; \skt{त्वक्पारुष्य}\ A}}}{\lemma{\skt{अम्ल\appbrk{}शी\appbrk{}त\appbrk{}प्रा\appbrk{}र्थना}}\Afootnote{K N H; \skt{मम्ल\appbrk{}शी\appbrk{}त\appbrk{}प्रा\appbrk{}र्थना}\ A}}\skt{ सिराशैथिल्यञ्च~॥ मांसक्षये }\edtext{\skt{स्फिग्ग\-ण्ण्डौ\-ष्ठो\-प\-स्थो\-रु\-व\-क्षः\-क\-क्ष\-शु\-ष्कता}}{\lemma{\skt{स्फिग्ग\appbrk{}ण्ण्डौ\appbrk{}ष्ठो\appbrk{}प\appbrk{}स्थो\appbrk{}रु\appbrk{}व\appbrk{}क्षः\appbrk{}क\appbrk{}क्ष\appbrk{}शु\appbrk{}ष्कता}}\Afootnote{K; \skt{स्फिग्ग\appbrk{}ण्डौ\appbrk{}ष्ठो\appbrk{}प\appbrk{}स्थो\appbrk{}रु\appbrk{}व\appbrk{}क्षः\appbrk{}क\appbrk{}क्षा\appbrk{}पि\appbrk{}ण्डि\appbrk{}को\appbrk{}द\appbrk{}र\appbrk{}ग्री\appbrk{}वा\appbrk{}शु\appbrk{}ष्कता}\ A; \skt{हिग्गण्डौष्ठोप}\ N; \skt{स्फिग्ग\appbrk{}ण्डो\appbrk{}ष्ठो\appbrk{}प\appbrk{}स्थो\appbrk{}रु\appbrk{}व\appbrk{}क्षः\appbrk{}क\appbrk{}क्ष\appbrk{}शुष्क}\ H}}\skt{ }\edtext{\skt{धमनीनाञ्च}}{\lemma{\skt{धमनीनाञ्च}}\Afootnote{K; \skt{रौक्ष्यतोदौ गात्राणां सदनं धमनी}\ A; \skt{ष्ठो… धमनीनाञ्च}\ N; \skt{ता धमनीनाञ्च}\ H}}\skt{ }\edtext{\skt{शैथिल्यं}}{\lemma{\skt{शैथिल्यं}}\Afootnote{K N H; \skt{शैथिल्यं च}\ A}}\skt{~। }\edtext{\skt{मेदक्षये}}{\lemma{\skt{मेदक्षये}}\Afootnote{K N; \skt{मेदःक्षये}\ A H}}\skt{ प्लीहाभिवृद्धिः सन्धिशून्यता रौक्ष्यं मेदु\-र\-मां\-स\-प्रा\-र्थना च~। अस्थिक्षये }\edtext{\skt{अस्थिशूलो}}{\lemma{\skt{अस्थिशूलो}}\Afootnote{K\textsuperscript{pc} N H; \skt{ऽस्थिशूलं}\ A; \skt{शूलो}\ K\textsuperscript{ac}}}\skt{ }\edtext{\skt{दन्त\-न\-ख\-भं\-गा\-रौ\-क्ष्यञ्च}}{\lemma{\skt{दन्त\appbrk{}न\appbrk{}ख\appbrk{}भं\appbrk{}गा\appbrk{}रौ\appbrk{}क्ष्यञ्च}}\Afootnote{K; \skt{दन्तनखभङ्गो रौक्ष्यं च}\ A N H}}\skt{~। }\edtext{\skt{मज्जाक्षये}}{\lemma{\skt{मज्जाक्षये}}\Afootnote{K; \skt{मज्जक्षये}\ A N H}}\skt{ }\edtext{\skt{ल्पशुक्रता}}{\lemma{\skt{ल्पशुक्रता}}\Afootnote{K; \skt{ऽल्पशुक्रता}\ A N H}}\skt{ पर्वभेदो }\edtext{\skt{स्थितोदः}}{\lemma{\skt{स्थितोदः}}\Afootnote{K; \skt{ऽस्थिनिस्तोदो}\ A; \skt{स्थि}\ N; \skt{ऽस्थिनो}\ H}}\skt{ }\edtext{\skt{शून्यता}}{\lemma{\skt{शून्यता}}\Afootnote{K; \skt{ऽस्थिशून्यता च}\ A; \skt{तोदः शून्यता}\ N; \skt{दःशून्यता}\ H}}\skt{~। शुक्रक्षये }\edtext{\edtext{\skt{मेढ्र\-वृ\-ष\-ण\-वे\-द\-ना\-श\-क्ति\-र्मै\-थुने}}{\lemma{\skt{मेढ्र\appbrk{}वृ\appbrk{}ष\appbrk{}ण\appbrk{}वे\appbrk{}द\appbrk{}ना\appbrk{}श\appbrk{}क्तिर्}}\Afootnote{K N H; \skt{मेढ्र\appbrk{}वृ\appbrk{}ष\appbrk{}ण\appbrk{}वे\appbrk{}द\appbrk{}नाऽ\appbrk{}शक्ति}\ A}}}{\lemma{\skt{मैथुने}}\Afootnote{K N H; \skt{र्मैथुने}\ A}}\skt{ }\edtext{\skt{चिराद्वा}}{\lemma{\skt{चिराद्वा}}\Afootnote{A K H; \skt{प्रचाराद्वा}\ N}}\skt{ प्रसेकः प्रसेके }\edtext{\skt{चाल्प\-शु\-क्र\-र\-क्त\-द\-र्शनं}}{\lemma{\skt{चाल्प\appbrk{}शु\appbrk{}क्र\appbrk{}र\appbrk{}क्त\appbrk{}द\appbrk{}र्शनं}}\Afootnote{K N H; \skt{चाल्प\appbrk{}र\appbrk{}क्त\appbrk{}शु\appbrk{}क्र\appbrk{}द\appbrk{}र्श\appbrk{}नम्}\ A}}\skt{~।}\pend

\pstart \ledsidenote{\passref{1.15.11}}\skt{पुरीषक्षये हृद\-य\-पा\-र्श्व\-पीडा सशब्दस्य च }\edtext{\skt{वायो\-रू\-ध्व\-ग\-मनं}}{\lemma{\skt{ऊध्वगमनं}}\Afootnote{K; \skt{ऊर्ध्वगमनं}\ A N H}}\skt{ कुक्षौ सञ्चरणं च मूत्रक्षये }\edtext{\skt{बस्तितोदो}}{\lemma{\skt{बस्तितोदो}}\Afootnote{A K N; \skt{बस्तितोदो}\ H\textsuperscript{ac}; \skt{बस्तितो}\ H\textsuperscript{pc}}}\skt{ }\edtext{\skt{ल्पमूत्रता च}}{\lemma{\skt{ल्पमूत्रता च}}\Afootnote{K N H\textsuperscript{ac}; \skt{ऽल्पमूत्रता च}\ A; \skt{दोल्पमूत्रता च}\ H\textsuperscript{pc}}}\skt{~॥ }\edtext{\skt{अत्रापि}}{\lemma{\skt{अत्रापि}}\Afootnote{A K H; \skt{अतो पि}\ N}}\skt{ }\edtext{\skt{स्वयो\-नि\-व\-र्ध\-न\-द्र\-व्यो\-प\-योगः}}{\lemma{\skt{स्वयो\appbrk{}नि\appbrk{}व\appbrk{}र्ध\appbrk{}न\appbrk{}द्र\appbrk{}व्यो\appbrk{}प\appbrk{}योगः}}\Afootnote{A K H; \skt{स्वयो\appbrk{}नि\appbrk{}व\appbrk{}र्ध\appbrk{}न\appbrk{}द्र\appbrk{}व्यो\appbrk{}प्र\appbrk{}योगः}\ N}}\skt{~॥ स्वेदक्षये स्तब्ध\-रो\-म\-कू\-पता }\edtext{\edtext{\skt{स्पर्शवैगुण्यश्च}}{\lemma{\skt{स्पर्शवैगुण्यश्}}\Afootnote{K; \skt{त्वक्शोषः स्पर्शवैगुण्यं}\ A; \skt{स्पर्शवैगुण्यञ्}\ N H\textsuperscript{ac}; \skt{स्पर्शवैगु}\ H\textsuperscript{pc}}}}{\lemma{\skt{च}}\Afootnote{K N H\textsuperscript{ac}; \skt{स्वेदनाशश्च}\ A; \skt{ण्यथं द्रञ्च}\ H\textsuperscript{pc}}}\skt{ }\edtext{\skt{तत्राभ्यङ्ग}}{\lemma{\skt{तत्राभ्यङ्ग}}\Afootnote{K N H; \skt{तत्राभ्यङ्गः}\ A}}\skt{ स्वेदोपयोगश्च~॥}\pend

\pstart \ledsidenote{\passref{1.15.12}}\skt{आर्त्तवक्षये }\edtext{\skt{यथो\-चि\-त\-का\-ला\-द\-र्श\-न\-म\-ल्पता वा}}{\lemma{\skt{यथो\appbrk{}चि\appbrk{}त\appbrk{}का\appbrk{}ला\appbrk{}द\appbrk{}र्श\appbrk{}न\appbrk{}म\appbrk{}ल्पता वा}}\Afootnote{A K; \skt{यथो\appbrk{}चि\appbrk{}त\appbrk{}का\appbrk{}ला\appbrk{}द\appbrk{}र्श\appbrk{}न\appbrk{}म\appbrk{}ल्पता वा}\ N H}}\skt{ योनिवेदना च~। तत्र }\edtext{\skt{संशो\-ध\-न\-मा\-ग्ने\-या\-नाञ्च}}{\lemma{\skt{संशो\appbrk{}ध\appbrk{}न\appbrk{}मा\appbrk{}ग्ने\appbrk{}या\appbrk{}नाञ्च}}\Afootnote{A K; \skt{संशो\appbrk{}ध\appbrk{}न\appbrk{}मा\appbrk{}ग्ने\appbrk{}यानां च}\ N H}}\skt{ }\edtext{\edtext{\skt{द्रव्या\-ना\-मु\-प\-योगः}}{\lemma{\skt{द्रव्यानाम्}}\Afootnote{K N H; \skt{द्रव्याणां}\ A}}}{\lemma{\skt{उपयोगः}}\Afootnote{K N; \skt{विधि\appbrk{}व\appbrk{}दु\appbrk{}प\appbrk{}योगः}\ A; \skt{उपयोगः…}\ H}}\skt{~॥ }\edtext{\skt{स्तन्यक्षये}}{\lemma{\skt{स्तन्यक्षये}}\Afootnote{K H; \skt{स्तनक्षये}\ A; \skt{स्त…क्षये}\ N}}\skt{ }\edtext{\skt{स्तनयो}}{\lemma{\skt{स्तनयो}}\Afootnote{K; \skt{स्तनयोर्}\ A H N}}\skt{ म्लानता }\edtext{\skt{स्तन्यासम्भवश्च}}{\lemma{\skt{स्तन्यासम्भवश्च}}\Afootnote{K H; \skt{स्तन्यासम्भवो}\ A; \skt{स्तन्य्}\ N}}\skt{~। }\edtext{\skt{तत्र श्लेष्म\-व\-र्ध\-न\-द्र\-व्यो\-प\-योगः}}{\lemma{\skt{तत्र श्लेष्म\appbrk{}व\appbrk{}र्ध\appbrk{}न\appbrk{}द्र\appbrk{}व्यो\appbrk{}प\appbrk{}योगः}}\Afootnote{K H; \skt{ऽल्पता वा; तत्र श्लेष्म\appbrk{}व\appbrk{}र्ध\appbrk{}न\appbrk{}द्र\appbrk{}व्यो\appbrk{}प\appbrk{}योगः}\ A; \skt{…द्रव्योपयोगः}\ N}}\skt{~॥ }\edtext{\skt{गर्भक्षये}}{\lemma{\skt{गर्भक्षये}}\Afootnote{A K N; \skt{गर्भ्भक्षये}\ H}}\skt{ }\edtext{\skt{गर्भा\-स्प\-न्द\-न\-म\-नु\-न्न\-त\-कु\-क्षिता च}}{\lemma{\skt{गर्भा\appbrk{}स्प\appbrk{}न्द\appbrk{}न\appbrk{}म\appbrk{}नु\appbrk{}न्न\appbrk{}त\appbrk{}कु\appbrk{}क्षिता च}}\Afootnote{A K; \skt{गर्भा\appbrk{}स्प\appbrk{}न्द\appbrk{}न\appbrk{}म\appbrk{}नु\appbrk{}न्न\appbrk{}त\appbrk{}कुक्षि… च}\ N; \skt{गर्भ्भा\appbrk{}स्प\appbrk{}न्द\appbrk{}न\appbrk{}म\appbrk{}नु\appbrk{}न्न\appbrk{}त\appbrk{}कु\appbrk{}क्षिता च}\ H}}\skt{~। तत्र प्राप्त\-ब\-स्ति\-का\-लायाः क्षीर\-ब\-स्ति\-प्र\-योगो }\edtext{\skt{मेध्यान्न }\edtext{\skt{प्रयोगश्च}}{\lemma{\skt{च}}\Afootnote{K N H; \skt{चेति}\ A}}}{\lemma{\skt{मेध्यान्न प्रयोगश्}}\Afootnote{K N H; \skt{मेद्या\appbrk{}न्नो\appbrk{}प\appbrk{}यो\appbrk{}गश्}\ A}}\skt{~॥}\pend

\pstart \ledsidenote{\passref{1.15.13}}\skt{अत }\edtext{\edtext{\skt{ऊर्ध्व\-म\-ति\-प्र\-वृ\-द्धानां}}{\lemma{\skt{ऊर्ध्वम्}}\Afootnote{K N H; \skt{ऊर्ध्व}\ A}}}{\lemma{\skt{अतिप्रवृद्धानां}}\Afootnote{K N H; \skt{मतिवृद्धानां}\ A}}\skt{ }\edtext{\skt{दोषधातूनां}}{\lemma{\skt{दोषधातूनां}}\Afootnote{K N H; \skt{दोष\appbrk{}धा\appbrk{}तु\appbrk{}म\appbrk{}लानां}\ A}}\skt{ }\edtext{\skt{लक्ष\-ण\-मु\-प\-दे\-क्ष्यामः}}{\lemma{\skt{उपदेक्ष्यामः}}\Afootnote{K N H; \skt{वक्ष्यामः}\ A}}\skt{ }\edtext{\skt{तत्र}}{\lemma{\skt{तत्र}}\Afootnote{K N; \skt{वृद्धिः पुनरेषां स्वयो\appbrk{}नि\appbrk{}व\appbrk{}र्ध\appbrk{}ना\appbrk{}त्यु\appbrk{}प\appbrk{}से\appbrk{}व\appbrk{}ना\appbrk{}द्भ\appbrk{}वति~। तत्र}\ A; \skt{… तत्र}\ H}}\skt{ वातवृद्धौ }\edtext{\skt{कार्श्यं}}{\lemma{\skt{कार्श्यं}}\Afootnote{K H; \skt{वाक्पारुष्यं कार्श्यं}\ A; \skt{कार्श्य}\ N}}\skt{ }\edtext{\skt{कार्ष्ण्यं}}{\lemma{\skt{कार्ष्ण्यं}}\Afootnote{A K H; \skt{कार्ष्ण्य}\ N}}\skt{ }\edtext{\skt{गात्रस्फुरणता}}{\lemma{\skt{गात्रस्फुरणता}}\Afootnote{K N H; \skt{गात्रस्फुरण}\ A}}\skt{ }\edtext{\skt{उष्णकामता}}{\lemma{\skt{उष्णकामता}}\Afootnote{K N H; \skt{मुष्णकामि(म)ता}\ A}}\skt{ }\edtext{\skt{निद्रा\-ना\-शो\-ल्प\-ब\-लत्वं}}{\lemma{\skt{निद्रा\appbrk{}ना\appbrk{}शो\appbrk{}ल्प\appbrk{}ब\appbrk{}लत्वं}}\Afootnote{K\textsuperscript{pc}; \skt{निद्रा\appbrk{}ना\appbrk{}शोऽ\appbrk{}ल्प\appbrk{}ब\appbrk{}लत्वं}\ A H; \skt{निद्रा\appbrk{}ना\appbrk{}शो\appbrk{}ल्प\appbrk{}बत्वं}\ K\textsuperscript{ac}; \skt{निद्रा\appbrk{}ना\appbrk{}शो’\appbrk{}ल्प\appbrk{}ब\appbrk{}लत्वं}\ N}}\skt{ }\edtext{\skt{गाढवर्चस्वता च}}{\lemma{\skt{गाढवर्चस्वता च}}\Afootnote{K; \skt{गाढवर्चस्त्वं च}\ A; \skt{गाढवर्च्चस्कता च}\ N H}}\skt{~॥ पित्तवृद्धौ पीतावभासता }\edtext{\skt{सन्ता\-पः\-शी\-त\-का\-मि\-त्व\-म\-ल्प\-नि\-द्रता}}{\lemma{\skt{सन्ता\appbrk{}पः\appbrk{}शी\appbrk{}त\appbrk{}का\appbrk{}मि\appbrk{}त्व\appbrk{}म\appbrk{}ल्प\appbrk{}नि\appbrk{}द्रता}}\Afootnote{A H; \skt{सन्ता\appbrk{}पः\appbrk{}शी\appbrk{}त\appbrk{}का\appbrk{}मि\appbrk{}त्व\appbrk{}म\appbrk{}ल्प\appbrk{}नि\appbrk{}द्रता}\ K N}}\skt{ मूर्च्छा }\edtext{\skt{बलहानिः}}{\lemma{\skt{बलहानिः}}\Afootnote{K N H; \skt{बलहानि}\ A}}\skt{ }\edtext{\skt{पीत\-वि\-ण्मू\-त्र\-त्वञ्च}}{\lemma{\skt{पीत\appbrk{}वि\appbrk{}ण्मू\appbrk{}त्र\appbrk{}त्वञ्च}}\Afootnote{K N H; \skt{रिन्द्रियदौर्बल्यं पीत\appbrk{}वि\appbrk{}ण्मू\appbrk{}त्र\appbrk{}ने\appbrk{}त्रत्वं च}\ A}}\skt{~॥ श्लेष्मवृद्धौ शौक्ल्यं }\edtext{\skt{शैथ्यं}}{\lemma{\skt{शैथ्यं}}\Afootnote{K; \skt{शैत्यं}\ A; \skt{शैत्यं}\ N H}}\skt{ स्थैर्यं }\edtext{\edtext{\edtext{\skt{गौर\-व\-म\-ग्नि\-सा\-द\-स्तन्द्रा}}{\lemma{\skt{गौरवम्}}\Afootnote{K\textsuperscript{pc} N H; \skt{गौरव}\ A K\textsuperscript{ac}}}}{\lemma{\skt{अग्निसादस्}}\Afootnote{K\textsuperscript{pc} N H; \skt{मवसाद}\ A; \skt{सादस्}\ K\textsuperscript{ac}}}}{\lemma{\skt{तन्द्रा}}\Afootnote{K N H; \skt{स्तन्द्रा}\ A}}\skt{ निद्रा }\edtext{\skt{सन्ध्यति}}{\lemma{\skt{सन्ध्यति}}\Afootnote{K N H; \skt{सन्धिवि}\ A}}\skt{ }\edtext{\skt{श्लिष्टता}}{\lemma{\skt{श्लिष्टता}}\Afootnote{K N; \skt{श्लेषश्}\ A; \skt{श्लिष्तता}\ H}}\skt{ च~॥}\pend

\pstart \ledsidenote{\passref{1.15.14}}\edtext{\skt{रसोऽतिप्रवृद्धो}}{\lemma{\skt{ऽतिप्रवृद्धो}}\Afootnote{K N H; \skt{ऽतिवृद्धो}\ A}}\skt{ }\edtext{\skt{हृदये}}{\lemma{\skt{हृदये}}\Afootnote{K N H; \skt{हृदयो}\ A}}\skt{ }\edtext{\skt{क्लेदं}}{\lemma{\skt{क्लेदं}}\Afootnote{K N H; \skt{त्क्लेदं}\ A}}\skt{ प्रसे\-क\-ञ्चा\-पा\-द\-यति~। रक्तं रक्ताङ्गाक्षितां~॥ }\edtext{\skt{मांसं}}{\lemma{\skt{मांसं}}\Afootnote{K N H; \skt{सिरापूर्णत्वं च; मांसं}\ A}}\skt{ }\edtext{\skt{स्फिग्ग\-ण्डौ\-ष्ठो\-प\-स्थो\-रु\-बा\-हु\-ज\-ङ्घासु}}{\lemma{\skt{स्फिग्ग\appbrk{}ण्डौ\appbrk{}ष्ठो\appbrk{}प\appbrk{}स्थो\appbrk{}रु\appbrk{}बा\appbrk{}हु\appbrk{}ज\appbrk{}ङ्घासु}}\Afootnote{A K N; \skt{स्फिग्ग\appbrk{}ण्डो\appbrk{}ष्ठो\appbrk{}प\appbrk{}स्थो\appbrk{}रु\appbrk{}बा\appbrk{}हु\appbrk{}जं\appbrk{}घासु}\ H}}\skt{ }\edtext{\skt{वृद्धिं}}{\lemma{\skt{वृद्धिं}}\Afootnote{A K H; \skt{वृद्धि}\ N}}\skt{ गुरुगात्रताञ्च~। मेदः }\edtext{\skt{स्निग्धा\-ङ्ग\-ता\-मु\-द\-र\-पा\-र्श्व\-वृद्धिं}}{\lemma{\skt{स्निग्धा\appbrk{}ङ्ग\appbrk{}ता\appbrk{}मु\appbrk{}द\appbrk{}र\appbrk{}पा\appbrk{}र्श्व\appbrk{}वृद्धिं}}\Afootnote{A K; \skt{स्निग्धा\appbrk{}ङ्ग\appbrk{}ता\appbrk{}मु\appbrk{}द\appbrk{}र\appbrk{}पा\appbrk{}र्श्व\appbrk{}वृ\appbrk{}द्धिङ्}\ N H}}\skt{ }\edtext{\skt{कासश्वासो दौर्गन्ध्यञ्च}}{\lemma{\skt{कासश्वासो दौर्गन्ध्यञ्च}}\Afootnote{K; \skt{कास\appbrk{}श्वा\appbrk{}सा\appbrk{}दी\appbrk{}न्दौ\appbrk{}र्गन्ध्यं च}\ A; \skt{कासश्वासौ दौर्गन्ध्यञ्च}\ N H}}\skt{~॥ }\edtext{\edtext{\skt{अस्थ्य\-ध्य\-स्थी\-न्य\-धि\-द\-न्तांश्च}}{\lemma{\skt{अस्थ्यध्यस्थीन्य्}}\Afootnote{K; \skt{अस्थ्यध्यस्थी}\ A H; \skt{अस्थि अध्यस्थीन्य्}\ N}}}{\lemma{\skt{अधिदन्तांश्च}}\Afootnote{K N; \skt{न्यधिदन्तांश्च}\ A; \skt{त्यधिदन्तांश्च}\ H}}\skt{~॥ मज्जा }\edtext{\skt{सर्वा\-ङ्ग\-ने\-त्र\-गौ\-रवं}}{\lemma{\skt{सर्वा\appbrk{}ङ्ग\appbrk{}ने\appbrk{}त्र\appbrk{}गौ\appbrk{}रवं}}\Afootnote{K H; \skt{सर्वा\appbrk{}ङ्ग\appbrk{}ने\appbrk{}त्र\appbrk{}गौ\appbrk{}रवं च}\ A; \skt{सर्वा\appbrk{}ङ्गा\appbrk{}ने\appbrk{}त्र\appbrk{}गौ\appbrk{}रवं}\ N}}\skt{~॥ शुक्रं }\edtext{\skt{शुक्राश्मरिति प्रादुर्भावं}}{\lemma{\skt{शुक्राश्मरिति प्रादुर्भावं}}\Afootnote{\skt{शुक्रा\appbrk{}श्म\appbrk{}री\appbrk{}म\appbrk{}ति\appbrk{}प्रा\appbrk{}दु\appbrk{}र्भावं च}\ A; \skt{शुक्राश्मरि…ति प्रादुर्भावं}\ K; \skt{शुक्रा\appbrk{}श्म\appbrk{}र्य\appbrk{}ति\appbrk{}प्रा\appbrk{}दु\appbrk{}र्भा\appbrk{}वम्}\ N; \skt{शुक्रा\appbrk{}स्म\appbrk{}र्य\appbrk{}ति\appbrk{}प्रा\appbrk{}दु\appbrk{}र्भ्भावं}\ H}}\skt{~॥}\pend

\pstart \ledsidenote{\passref{1.15.15}}\edtext{\edtext{\skt{पुरीशमाटोपः}}{\lemma{\skt{पुरीशम्}}\Afootnote{K; \skt{पुरीषम्}\ A N H}}}{\lemma{\skt{आटोपः}}\Afootnote{K N H; \skt{आटोपं}\ A}}\skt{ कुक्षौ शूलञ्च~॥ मूत्रं }\edtext{\skt{मुहुर्मुहुः}}{\lemma{\skt{मुहुर्मुहुः}}\Afootnote{K N H; \skt{मूत्रवृद्धिं मुहुर्मुहुः}\ A}}\skt{ }\edtext{\skt{प्रवृत्तिं}}{\lemma{\skt{प्रवृत्तिं}}\Afootnote{A K H; \skt{प्रवृ}\ N}}\skt{ }\edtext{\edtext{\skt{तोदञ्च}}{\lemma{\skt{तोदञ्}}\Afootnote{K; \skt{बस्तितोदम्}\ A; \skt{त्ति तोदञ्}\ N; \skt{तोद}\ H}}}{\lemma{\skt{च}}\Afootnote{K N; \skt{आध्मानं च}\ A; \skt{मञ्च}\ H}}\skt{~॥ स्वेदः }\edtext{\skt{कण्डू}}{\lemma{\skt{कण्डू}}\Afootnote{K N H; \skt{त्वचो}\ A}}\skt{ }\edtext{\skt{दौर्गन्ध्यञ्च}}{\lemma{\skt{च}}\Afootnote{K N H; \skt{कण्डूं च}\ A}}\skt{~॥}\pend

\pstart \ledsidenote{\passref{1.15.16}}\edtext{\skt{स्तन्यं}}{\lemma{\skt{स्तन्यं}}\Afootnote{K N H; \skt{आर्त\appbrk{}व\appbrk{}म\appbrk{}ङ्ग\appbrk{}म\appbrk{}र्द\appbrk{}म\appbrk{}ति\appbrk{}प्र\appbrk{}वृत्तिं दौर्गन्ध्यं च; स्तन्यं}\ A}}\skt{ }\edtext{\edtext{\skt{स्तन\-यो\-र\-ति\-पी\-नत्वं}}{\lemma{\skt{स्तनयोर्}}\Afootnote{A K N; \skt{स्तनयो}\ H}}}{\lemma{\skt{अतिपीनत्वं}}\Afootnote{K N; \skt{आपीनत्वं}\ A; \skt{र्रतिपीनत्वं}\ H}}\skt{ }\edtext{\skt{मुहुर्मुहुः}}{\lemma{\skt{मुहुर्मुहुः}}\Afootnote{A K; \skt{मुहुर्मुहुः}\ N; \skt{मुहुर्म्मुहु}\ H}}\skt{ }\edtext{\edtext{\skt{प्रवृ\-त्ति\-म\-ति\-तो\-दञ्च}}{\lemma{\skt{प्रवृत्तिम्}}\Afootnote{A K N; \skt{प्रवृत्ति}\ H}}}{\lemma{\skt{अतितोदञ्च}}\Afootnote{K N; \skt{तोदं}\ A; \skt{मति तोदञ्च}\ H}}\skt{~॥ }\edtext{\skt{आर्त्त\-व\-म\-ङ्ग\-मर्दो}}{\lemma{\skt{आर्त्त\appbrk{}व\appbrk{}म\appbrk{}ङ्ग\appbrk{}मर्दो}}\Afootnote{\emph{om.}\ A; \skt{आर्त्त\appbrk{}व\appbrk{}म\appbrk{}ङ्ग\appbrk{}मर्दो}\ K H; \skt{आर्त्त\appbrk{}व\appbrk{}म\appbrk{}ङ्ग\appbrk{}मर्दं}\ N}}\skt{ }\edtext{\skt{दौर्बल्यञ्च}}{\lemma{\skt{दौर्बल्यञ्च}}\Afootnote{K; \skt{च; गर्भो जठराभिवृद्धिं स्वेदं च}\ A; \skt{दौर्बल्यञ्च~। गर्भो जठ\appbrk{}रा\appbrk{}भि\appbrk{}वृ\appbrk{}द्धि\appbrk{}ङ्क\appbrk{}रोति}\ N; \skt{दौर्ब्बल्यञ्च~॥ गर्भ्भो जठ\appbrk{}रा\appbrk{}भि\appbrk{}वृ\appbrk{}द्धि\appbrk{}ङ्क\appbrk{}रोति}\ H}}\pend

\pstart \ledsidenote{\passref{1.15.17}}\edtext{\skt{तेषां}}{\lemma{\skt{तेषां}}\Afootnote{K N H; \skt{तेषां यथास्वं संशोधनं}\ A}}\skt{ }\edtext{\edtext{\skt{क्षप\-ण\-म\-वि\-रुद्धैः}}{\lemma{\skt{क्षपणम्}}\Afootnote{K N H; \skt{क्षपणं च}\ A}}}{\lemma{\skt{अविरुद्धैः}}\Afootnote{K N H; \skt{क्षयादविरुद्धैः}\ A}}\skt{ क्रियाविशेषैः }\edtext{\skt{कुर्वीत}}{\lemma{\skt{कुर्वीत}}\Afootnote{N H; \skt{प्रकुर्वीत}\ A; \skt{कुर्वीत बलक्षयं}\ K}}\skt{~।}\pend

\pstart \ledsidenote{\passref{1.15.19}}\edtext{\skt{अत}}{\lemma{\skt{अत}}\Afootnote{K; \skt{बललक्षणं बल\appbrk{}क्ष\appbrk{}य\appbrk{}ल\appbrk{}क्षणं चात}\ A; \skt{बलक्षयमत}\ N H}}\skt{ }\edtext{\skt{ऊर्ध्व\-म\-नु\-व्या\-ख्या\-स्यामः}}{\lemma{\skt{अनुव्याख्यास्यामः}}\Afootnote{K N; \skt{उपदेक्ष्यामः~। तत्र}\ A; \skt{अन्व्याख्यास्यामः}\ H}}\skt{~॥ रसादीनां शुक्रान्तानां धातूनां }\edtext{\edtext{\edtext{\skt{यत्प\-र\-न्ते\-ज\-स्त\-त्ख\-ल्वो\-ज\-स्त\-देव}}{\lemma{\skt{तेजस्}}\Afootnote{K N H; \skt{तेज}\ A}}}{\lemma{\skt{तत्खल्वोजस्}}\Afootnote{K N H; \skt{स्तत्खल्वोज}\ A}}}{\lemma{\skt{तदेव}}\Afootnote{K N H; \skt{स्तदेव}\ A}}\skt{ बलमित्युच्यते~। }\edtext{\skt{शास्त्रसिद्धान्तात्}}{\lemma{\skt{शास्त्रसिद्धान्तात्}}\Afootnote{K N H; \skt{स्वशा\appbrk{}स्त्र\appbrk{}सि\appbrk{}द्धा\appbrk{}न्तात्}\ A}}\pend

\pstart \ledsidenote{\passref{1.15.20}}\skt{तत्र बलेन स्थिरो\-प\-चि\-त\-मां\-सता }\edtext{\skt{सर्व\-चे\-ष्टा\-स्व\-प्र\-ति\-घातः}}{\lemma{\skt{सर्व\appbrk{}चे\appbrk{}ष्टा\appbrk{}स्व\appbrk{}प्र\appbrk{}ति\appbrk{}घातः}}\Afootnote{A K; \skt{सर्व\appbrk{}चे\appbrk{}ष्टा\appbrk{}स्व\appbrk{}प्र\appbrk{}ती\appbrk{}घातः}\ N; \skt{सर्व्व\appbrk{}चे\appbrk{}ष्टा\appbrk{}स्व\appbrk{}प्र\appbrk{}ती\appbrk{}इ\appbrk{}घा…तः}\ H}}\skt{ स्वर\-व\-र्ण्ण\-प्र\-सादो }\edtext{\skt{बाह्याभ्यन्तराणां च}}{\lemma{\skt{बाह्याभ्यन्तराणां च}}\Afootnote{K N H; \skt{बाह्या\appbrk{}ना\appbrk{}मा\appbrk{}भ्य\appbrk{}न्त\appbrk{}राणां च}\ A}}\skt{ }\edtext{\edtext{\edtext{\skt{कर\-णा\-ना\-मा\-त्म\-का\-र्य\-प्र\-ति\-प\-त्ति\-र्भ\-वति}}{\lemma{\skt{करणानाम्}}\Afootnote{K H; \skt{करणाना}\ A N}}}{\lemma{\skt{आत्म\appbrk{}का\appbrk{}र्य\appbrk{}प्र\appbrk{}ति\appbrk{}प\appbrk{}त्तिर्}}\Afootnote{K; \skt{मात्म\appbrk{}का\appbrk{}र्य\appbrk{}प्र\appbrk{}ति\appbrk{}प\appbrk{}त्तिर्}\ A; \skt{त्मका\appbrk{}र्य\appbrk{}प्र\appbrk{}ति\appbrk{}प\appbrk{}त्तिर्}\ N; \skt{आत्म\appbrk{}का\appbrk{}र्य\appbrk{}प्र\appbrk{}त्ति\appbrk{}प\appbrk{}त्तिर्}\ H}}}{\lemma{\skt{भवति}}\Afootnote{A K N; \skt{भ्भवति}\ H}}\skt{~॥}\pend

\pstart\versepar \ledsidenote{\passref{1.15.21}}\edtext{\skt{भवति चात्र}}{\lemma{\skt{भवति चात्र}}\Afootnote{\skt{भवन्ति चात्र}\ A; \emph{om.}\ K N; \skt{भवन्ति चात्र}\ H}}\skt{~॥ }\pend
\pstart\versepar \skt{ओजः सोमात्मकं स्निग्धं }\edtext{\skt{शीतं}}{\lemma{\skt{शीतं}}\Afootnote{K N H; \skt{शुक्लं}\ A}}\skt{ }\edtext{\skt{श्लक्ष्णं}}{\lemma{\skt{श्लक्ष्णं}}\Afootnote{K N H; \skt{शीतं}\ A}}\skt{ स्थिरं सरं~। }\pend
\pstart\versepar \skt{विविक्तं मृदु मृत्स्नञ्च }\edtext{\skt{प्राणा\-य\-त\-न\-मु\-त्तमं}}{\lemma{\skt{उत्तमं}}\Afootnote{A K N; \skt{उत्त्तमं}\ H}}\skt{~॥}\pend

\pstart \ledsidenote{\passref{1.15.22ab}}\skt{देहः }\edtext{\edtext{\skt{सावयवस्तेन}}{\lemma{\skt{सावयवस्}}\Afootnote{K N H; \skt{सावयव}\ A}}}{\lemma{\skt{तेन}}\Afootnote{K N H; \skt{स्तेन}\ A}}\skt{ व्याप्तो भवति }\edtext{\skt{देहिनाम्}}{\lemma{\skt{देहिनाम्}}\Afootnote{K N H; \skt{देहिनः}\ A}}\skt{~।}\pend

\pstart\versepar \ledsidenote{\passref{1.15.23abcd}}\edtext{\skt{अभि\-घा\-ता\-त्क्ष\-या\-त्को\-पा\-द्ध्या\-ना\-च्छो\-का\-च्छ्र\-मा\-त्क्षुधः}}{\lemma{\skt{अभि\appbrk{}घा\appbrk{}ता\appbrk{}त्क्ष\appbrk{}या\appbrk{}त्को\appbrk{}पा\appbrk{}द्ध्या\appbrk{}ना\appbrk{}च्छो\appbrk{}का\appbrk{}च्छ्र\appbrk{}मा\appbrk{}त्क्षुधः}}\Afootnote{K; \skt{अभि\appbrk{}घा\appbrk{}ता\appbrk{}त्क्ष\appbrk{}या\appbrk{}त्को\appbrk{}पा\appbrk{}च्छो\appbrk{}का\appbrk{}द्ध्या\appbrk{}ना\appbrk{}च्छ्र\appbrk{}मा\appbrk{}त्क्षुधः}\ A; \skt{अभि\appbrk{}घा\appbrk{}ता\appbrk{}त्क्ष\appbrk{}या\appbrk{}त्को\appbrk{}पा\appbrk{}द्ध्या\appbrk{}ना\appbrk{}च्छो\appbrk{}का\appbrk{}च्छ्र\appbrk{}मा\appbrk{}त्क्षुधः}\ N; \skt{अभि\appbrk{}घा\appbrk{}ता\appbrk{}त्क्ष\appbrk{}या\appbrk{}त्को\appbrk{}पा\appbrk{}द्व्या\appbrk{}ना\appbrk{}च्छो\appbrk{}का\appbrk{}च्छ्र\appbrk{}मा\appbrk{}त्क्षुधः}\ H}}\skt{~। }\pend
\pstart\versepar \skt{ओजः संक्षीयते }\edtext{\skt{देहे}}{\lemma{\skt{देहे}}\Afootnote{K N H; \skt{ह्येभ्यो}\ A}}\skt{ धातु\-ग्र\-ह\-ण\-निः\-सृतं~॥}\pend

\pstart \ledsidenote{\passref{1.15.24}}\edtext{\skt{तत्र}}{\lemma{\skt{तत्र}}\Afootnote{K N H; \skt{तस्य}\ A}}\skt{ विस्रंसो व्यापत्क्षय इति }\edtext{\skt{लिङ्गानि}}{\lemma{\skt{लिङ्गानि}}\Afootnote{K N H; \skt{(त्रयो दोषाः;) लिङ्गानि}\ A}}\skt{ भवन्ति~। सन्धिविश्लेषो गात्राणां सदनं }\edtext{\skt{दोष\-च्य\-व\-न\-क्रि\-या\-स\-न्नि\-रो\-धश्च}}{\lemma{\skt{दोष\appbrk{}च्य\appbrk{}व\appbrk{}न\appbrk{}क्रि\appbrk{}या\appbrk{}स\appbrk{}न्नि\appbrk{}रो\appbrk{}धश्च}}\Afootnote{K; \skt{दोषच्यवनं क्रिया\appbrk{}स\appbrk{}न्नि\appbrk{}रो\appbrk{}धश्च}\ A N; \skt{दोषच्यवनं क्रियो सन्निरोधश्च}\ H\textsuperscript{ac}; \skt{दोथंषच्यवनं क्रियो सन्निरोधश्च}\ H\textsuperscript{pc}}}\skt{ विस्रंसे~। }\edtext{\skt{स्तब्धता}}{\lemma{\skt{स्तब्धता}}\Afootnote{K N H; \skt{स्तब्ध}\ A}}\skt{ गुरुगात्र ता }\edtext{\skt{शोफो}}{\lemma{\skt{शोफो}}\Afootnote{K N H; \skt{वातशोफो}\ A}}\skt{ वर्ण्णभेदो }\edtext{\edtext{\skt{ग्लानिस्तन्द्रा}}{\lemma{\skt{ग्लानिस्}}\Afootnote{K N H; \skt{ग्लानि}\ A}}}{\lemma{\skt{तन्द्रा}}\Afootnote{K N H; \skt{स्तन्द्रा}\ A}}\skt{ }\edtext{\skt{निद्रा}}{\lemma{\skt{निद्रा}}\Afootnote{K N H; \skt{निद्रा च}\ A}}\skt{ व्यापन्ने~। }\edtext{\skt{मांसक्षयो}}{\lemma{\skt{मांसक्षयो}}\Afootnote{K N H; \skt{मूर्च्छा मांसक्षयो}\ A}}\skt{ मोहः प्रलापो मरणमिति च }\edtext{\skt{क्षीणो}}{\lemma{\skt{क्षीणो}}\Afootnote{K; \skt{क्षये}\ A; \skt{क्षीणे}\ N H}}\skt{~।}\pend

\pstart \ledsidenote{\passref{1.15.28}}\edtext{\skt{तत्र}}{\lemma{\skt{तत्र}}\Afootnote{K N H; \skt{पूर्वोक्तानि च लिङ्गानि मरणं च बलक्षये~। तत्र}\ A}}\skt{ विस्रंसे व्यापन्ने च }\edtext{\edtext{\skt{क्रिया\-वि\-शे\-षै\-र\-वि\-रु\-द्धै\-र्ब\-ल\-म\-ध्या\-य\-ये\-त्मू\-ढ\-सं\-ज्ञ\-मि\-त\-रञ्च}}{\lemma{\skt{अध्या\appbrk{}य\appbrk{}ये\appbrk{}त्मू\appbrk{}ढ\appbrk{}सं\appbrk{}ज्ञम्}}\Afootnote{K; \skt{आप्याययेत् ; इतरं तु मूढसञ्ज्ञं}\ A; \skt{अध्या\appbrk{}य\appbrk{}ये\appbrk{}न्मू\appbrk{}ढ\appbrk{}सं\appbrk{}ज्ञम्}\ N; \skt{अध्या\appbrk{}य\appbrk{}ये\appbrk{}त्मू\appbrk{}ढ\appbrk{}सं\appbrk{}ज्ञाम्}\ H}}}{\lemma{\skt{इतरञ्च}}\Afootnote{K N H; \emph{om.}\ A}}\skt{ वर्ज्जयेत्~।}\pend

\pstart\versepar \ledsidenote{\passref{1.15.31}}\skt{यस्य धातुक्षयाद्वायुः }\edtext{\skt{सञ्ज्ञाकर्म्म}}{\lemma{\skt{सञ्ज्ञाकर्म्म}}\Afootnote{K N H; \skt{सञ्ज्ञां कर्म}\ A}}\skt{ विनाशयेत् }\pend
\pstart\versepar \skt{प्रक्षीणञ्च बलं }\edtext{\skt{यस्य तौ न}}{\lemma{\skt{यस्य तौ न}}\Afootnote{K N H; \skt{यस्य नासौ}\ A}}\skt{ }\edtext{\skt{शक्तौ}}{\lemma{\skt{शक्तौ}}\Afootnote{K; \skt{शक्यश्}\ A N; \skt{शक्त्यौ}\ H}}\skt{ चिकित्सितुं~॥}\edtext{}{\lemma{\skt{विना\appbrk{}श\appbrk{}ये\appbrk{}त्प्र\appbrk{}क्षी\appbrk{}णञ्च}}\Afootnote{K N; \skt{च नाशयेत्~। प्रक्षीणं च}\ A; \skt{विनाशयेत्~। प्रक्षिणं च}\ H}}\pend

\pstart \ledsidenote{\passref{1.15.32}}\skt{रस\-नि\-मि\-त्त\-मेव स्थौल्यङ्कार्श्यञ्च~। तत्र }\edtext{\skt{श्लेष्म\-ला\-हा\-र\-से\-वि\-तो\-ध्य\-श\-न\-शी\-ल\-स्या\-व्या\-या\-मिनो}}{\lemma{\skt{श्लेष्म\appbrk{}ला\appbrk{}हा\appbrk{}र\appbrk{}से\appbrk{}वि\appbrk{}तो\appbrk{}ध्य\appbrk{}श\appbrk{}न\appbrk{}शी\appbrk{}ल\appbrk{}स्या\appbrk{}व्या\appbrk{}या\appbrk{}मिनो}}\Afootnote{K; \skt{श्लेष्म\appbrk{}ला\appbrk{}हा\appbrk{}र\appbrk{}से\appbrk{}वि\appbrk{}नोऽ\appbrk{}ध्य\appbrk{}श\appbrk{}न\appbrk{}शी\appbrk{}ल\appbrk{}स्या\appbrk{}व्या\appbrk{}या\appbrk{}मिनो}\ A; \skt{श्लेष्म\appbrk{}ला\appbrk{}हा\appbrk{}र\appbrk{}से\appbrk{}विनो ध्यश\appbrk{}न\appbrk{}शी\appbrk{}ल\appbrk{}स्या\appbrk{}व्या\appbrk{}या\appbrk{}मिनो}\ N; \skt{श्लेष्म\appbrk{}ला\appbrk{}हा\appbrk{}र\appbrk{}से\appbrk{}वि\appbrk{}नोऽ\appbrk{}ध्य\appbrk{}स\appbrk{}न\appbrk{}शी\appbrk{}ल\appbrk{}स्या\appbrk{}व्या\appbrk{}या\appbrk{}मिनो}\ H}}\skt{ दिवा\-स्व\-प्न\-र\-तस्य }\edtext{\skt{आम}}{\lemma{\skt{आम}}\Afootnote{K N; \skt{चाम}\ A; \skt{म}\ H}}\skt{ एवान्नरसो मधुरतरश्च }\edtext{\edtext{\skt{शरी\-र\-म\-नु\-क्र\-म\-मा\-णो\-ति\-स्ने\-हा\-न्मेदो}}{\lemma{\skt{शरीरम्}}\Afootnote{K N H; \skt{शरीर}\ A}}}{\lemma{\skt{अनु\appbrk{}क्र\appbrk{}म\appbrk{}मा\appbrk{}णो\appbrk{}ति\appbrk{}स्ने\appbrk{}हा\appbrk{}न्मेदो}}\Afootnote{K H; \skt{मनु\appbrk{}क्रा\appbrk{}म\appbrk{}न्न\appbrk{}ति\appbrk{}स्ने\appbrk{}हा\appbrk{}न्मेदो}\ A; \skt{अनुक्रममाणो तिस्नेहान्मेदो}\ N}}\skt{ जनयति~। }\edtext{\skt{मेदसो}}{\lemma{\skt{मेदसो}}\Afootnote{K N H; \skt{तद}\ A}}\skt{ }\edtext{\edtext{\skt{तिप्र\-वृ\-द्ध\-त्वा\-द्वा\-ठ\-र्य\-मा\-पा\-द\-यति}}{\lemma{\skt{तिप्र\appbrk{}वृ\appbrk{}द्ध\appbrk{}त्वा\appbrk{}द्वा\appbrk{}ठ\appbrk{}र्यम्}}\Afootnote{K N H; \skt{तिस्थौल्य}\ A}}\skt{~। }\edtext{\skt{तमतिवठरं}}{\lemma{\skt{अतिवठरं}}\Afootnote{K H; \skt{मतिस्थूलं}\ A; \skt{अतिपठरं}\ N}}}{\lemma{\skt{आपादयति~। तम्}}\Afootnote{K N H; \skt{मापादयति ; त}\ A}}\skt{ }\edtext{\skt{क्षुद्र\-स्वा\-स\-पि\-पा\-सा\-क्षु\-त्स्व\-प्न\-स्वे\-द\-दौ\-र्ग\-न्ध्य\-क्र\-थ\-न\-गा\-त्र\-सा\-द\-ग\-ग\-द\-त्वानि}}{\lemma{\skt{क्षुद्र\appbrk{}स्वा\appbrk{}स\appbrk{}पि\appbrk{}पा\appbrk{}सा\appbrk{}क्षु\appbrk{}त्स्व\appbrk{}प्न\appbrk{}स्वे\appbrk{}द\appbrk{}दौ\appbrk{}र्ग\appbrk{}न्ध्य\appbrk{}क्र\appbrk{}थ\appbrk{}न\appbrk{}गा\appbrk{}त्र\appbrk{}सा\appbrk{}द\appbrk{}ग\appbrk{}ग\appbrk{}द\appbrk{}त्वानि}}\Afootnote{K; \skt{क्षुद्र\appbrk{}श्वा\appbrk{}स\appbrk{}पि\appbrk{}पा\appbrk{}सा\appbrk{}क्षु\appbrk{}त्स्व\appbrk{}प्न\appbrk{}स्वे\appbrk{}द\appbrk{}गा\appbrk{}त्र\appbrk{}दौ\appbrk{}र्ग\appbrk{}न्ध्य\appbrk{}क्र\appbrk{}थ\appbrk{}न\appbrk{}गा\appbrk{}त्र\appbrk{}सा\appbrk{}द\appbrk{}ग\appbrk{}द्ग\appbrk{}द\appbrk{}त्वानि}\ A; \skt{क्षुद्र\appbrk{}श्वा\appbrk{}स\appbrk{}पि\appbrk{}पा\appbrk{}सा\appbrk{}क्षु\appbrk{}त्स्व\appbrk{}प्न\appbrk{}स्वे\appbrk{}द\appbrk{}दौ\appbrk{}र्ग\appbrk{}न्ध्य\appbrk{}ग्र\appbrk{}थ\appbrk{}गा\appbrk{}त्र\appbrk{}सा\appbrk{}द\appbrk{}ग\appbrk{}द्ग\appbrk{}द\appbrk{}त्वानि}\ N; \skt{क्षुद्र\appbrk{}स्वा\appbrk{}स\appbrk{}पि\appbrk{}पा\appbrk{}सा\appbrk{}क्षु\appbrk{}त्स्व\appbrk{}प्न\appbrk{}स्वे\appbrk{}द\appbrk{}दौ\appbrk{}र्ग्ग\appbrk{}न्ध्य\appbrk{}क्र\appbrk{}थ\appbrk{}न\appbrk{}गा\appbrk{}त्र\appbrk{}सा\appbrk{}द\appbrk{}ग\appbrk{}द्ग\appbrk{}द\appbrk{}त्वानि}\ H}}\skt{ }\edtext{\edtext{\skt{क्षिप्र\-मे\-वा\-वि\-शन्ति}}{\lemma{\skt{क्षिप्रम्}}\Afootnote{K N H; \skt{क्षिप्र}\ A}}}{\lemma{\skt{एवाविशन्ति}}\Afootnote{K N H; \skt{मेवाविशन्ति}\ A}}\skt{~। }\edtext{\skt{सौकु\-मा\-र्या\-त्मे\-दसः}}{\lemma{\skt{सौकु\appbrk{}मा\appbrk{}र्या\appbrk{}त्मे\appbrk{}दसः}}\Afootnote{K H; \skt{सौकु\appbrk{}मा\appbrk{}र्या\appbrk{}न्मे\appbrk{}दसः}\ A; \skt{सौकु\appbrk{}मा\appbrk{}र्या\appbrk{}न्मे\appbrk{}दसः}\ N}}\skt{ }\edtext{\skt{सर्व\-क्रि\-या\-स्व\-स\-म\-र्थत्वं}}{\lemma{\skt{सर्व\appbrk{}क्रि\appbrk{}या\appbrk{}स्व\appbrk{}स\appbrk{}म\appbrk{}र्थत्वं}}\Afootnote{K\textsuperscript{pc} H; \skt{सर्व\appbrk{}क्रि\appbrk{}या\appbrk{}स्व\appbrk{}स\appbrk{}मर्थः}\ A; \skt{सर्व\appbrk{}क्रि\appbrk{}स्व\appbrk{}स\appbrk{}म\appbrk{}र्थत्वं}\ K\textsuperscript{ac}; \skt{सर्व\appbrk{}क्रि\appbrk{}या\appbrk{}स्व\appbrk{}स\appbrk{}म\appbrk{}र्थ\appbrk{}त्वम्}\ N}}\skt{ }\edtext{\skt{भवति}}{\lemma{\skt{भवति}}\Afootnote{K N H; \emph{om.}\ A}}\skt{~। }\edtext{\skt{कफफमेदो}}{\lemma{\skt{कफफमेदो}}\Afootnote{K; \skt{कफमेदो}\ A; \skt{कफमेदो}\ N H}}\skt{ }\edtext{\skt{निरु\-द्ध\-मा\-र्ग\-त्वा\-च्चा\-ल्प\-व्य\-वायो}}{\lemma{\skt{निरु\appbrk{}द्ध\appbrk{}मा\appbrk{}र्ग\appbrk{}त्वा\appbrk{}च्चा\appbrk{}ल्प\appbrk{}व्य\appbrk{}वायो}}\Afootnote{K N; \skt{निरु\appbrk{}द्ध\appbrk{}मा\appbrk{}र्ग\appbrk{}त्वा\appbrk{}च्चा\appbrk{}ल्प\appbrk{}व्य\appbrk{}वायो}\ A; \skt{नुरु\appbrk{}द्ध\appbrk{}मा\appbrk{}र्ग्ग\appbrk{}त्वा…\appbrk{}च्चा\appbrk{}ल्प\appbrk{}व्य\appbrk{}वायो}\ H}}\skt{ भवति~। }\edtext{\skt{आवृ\-त\-मा\-र्ग\-त्वा\-देव च}}{\lemma{\skt{आवृ\appbrk{}त\appbrk{}मा\appbrk{}र्ग\appbrk{}त्वा\appbrk{}देव च}}\Afootnote{K N H; \skt{आवृ\appbrk{}त\appbrk{}मा\appbrk{}र्ग\appbrk{}त्वा\appbrk{}देव}\ A}}\skt{ शेषा धातवो }\edtext{\edtext{\skt{नाप्या\-य\-न्तेऽ\-त्य\-र्थ\-मतो}}{\lemma{\skt{ऽत्यर्थम्}}\Afootnote{K H; \skt{ऽत्यर्थ}\ A; \skt{त्यर्थम्}\ N}}}{\lemma{\skt{अतो}}\Afootnote{K N H; \skt{मतो}\ A}}\skt{ }\edtext{\skt{ल्पप्राणो}}{\lemma{\skt{ल्पप्राणो}}\Afootnote{K N; \skt{ऽल्पप्राणो}\ A H}}\skt{ भवति~। }\edtext{\edtext{\skt{प्रमे\-ह\-पि\-ड\-का\-ज्व\-र\-भ\-ग\-न्द\-र\-वि\-द्र\-धि\-वा\-त\-वि\-का\-णा\-रा\-म\-न्य\-तमं}}{\lemma{\skt{प्रमे\appbrk{}ह\appbrk{}पि\appbrk{}ड\appbrk{}का\appbrk{}ज्व\appbrk{}र\appbrk{}भ\appbrk{}ग\appbrk{}न्द\appbrk{}र\appbrk{}वि\appbrk{}द्र\appbrk{}धि\appbrk{}वा\appbrk{}त\appbrk{}वि\appbrk{}का\appbrk{}णा\appbrk{}राम्}}\Afootnote{\skt{प्रमे\appbrk{}ह\appbrk{}पि\appbrk{}ड\appbrk{}का\appbrk{}ज्व\appbrk{}र\appbrk{}भ\appbrk{}ग\appbrk{}न्द\appbrk{}र\appbrk{}वि\appbrk{}द्र\appbrk{}धि\appbrk{}वा\appbrk{}त\appbrk{}वि\appbrk{}का\appbrk{}राणा}\ A; \skt{प्रमे\appbrk{}ह\appbrk{}पि\appbrk{}ड\appbrk{}का\appbrk{}ज्व\appbrk{}र\appbrk{}भ\appbrk{}ग\appbrk{}न्द\appbrk{}र\appbrk{}वि\appbrk{}द्र\appbrk{}धि\appbrk{}वा\appbrk{}त\appbrk{}वि\appbrk{}काणा}\ K\textsuperscript{ac}; \skt{प्रमे\appbrk{}ह\appbrk{}पि\appbrk{}ड\appbrk{}का\appbrk{}ज्व\appbrk{}र\appbrk{}भ\appbrk{}ग\appbrk{}न्द\appbrk{}र\appbrk{}वि\appbrk{}द्र\appbrk{}धि\appbrk{}वा\appbrk{}त\appbrk{}वि\appbrk{}का\appbrk{}णारा}\ K\textsuperscript{pc}; \skt{प्रमे\appbrk{}ह\appbrk{}पि\appbrk{}ट\appbrk{}का\appbrk{}ज्व\appbrk{}र\appbrk{}भ\appbrk{}ग\appbrk{}न्द\appbrk{}र\appbrk{}वि\appbrk{}द्र\appbrk{}धि\appbrk{}वा\appbrk{}त\appbrk{}वि\appbrk{}रू\appbrk{}प\appbrk{}णा\appbrk{}नाम्}\ N; \skt{प्रमे\appbrk{}ह\appbrk{}पी\appbrk{}ट\appbrk{}का\appbrk{}ज्व\appbrk{}र\appbrk{}भ\appbrk{}ग\appbrk{}न्द\appbrk{}र\appbrk{}वि\appbrk{}द्र\appbrk{}धि\appbrk{}वा\appbrk{}त\appbrk{}वि\appbrk{}का\appbrk{}रा\appbrk{}णाम्}\ H}}}{\lemma{\skt{अन्यतमं}}\Afootnote{N H; \skt{मन्यतमं}\ A K}}\skt{ प्राप्य }\edtext{\skt{मर\-ण\-मु\-प\-याति}}{\lemma{\skt{मरणम्}}\Afootnote{K N H; \skt{पञ्चत्वम्}\ A}}\skt{~। सर्व एव चास्य रोगा बलवन्तो }\edtext{\skt{भवन्ति}}{\lemma{\skt{भवन्ति}}\Afootnote{K H; \skt{भवन्}\ A; \skt{भवति}\ N}}\skt{~। }\edtext{\skt{कस्मा\-दा\-वृ\-त\-मा\-र्ग\-त्वा\-त्स्रो\-त\-सा\-म\-त\-स्त\-स्यो\-त्प\-त्ति\-हेतुं}}{\lemma{\skt{कस्मा\appbrk{}दा\appbrk{}वृ\appbrk{}त\appbrk{}मा\appbrk{}र्ग\appbrk{}त्वा\appbrk{}त्स्रो\appbrk{}त\appbrk{}सा\appbrk{}म\appbrk{}तस्}}\Afootnote{K; \skt{त्यावृ\appbrk{}त\appbrk{}मा\appbrk{}र्ग\appbrk{}त्वा\appbrk{}त्स्रो\appbrk{}त\appbrk{}साम्; अतस्}\ A; \skt{कस्मा\appbrk{}दा\appbrk{}वृ\appbrk{}त\appbrk{}मा\appbrk{}र्ग\appbrk{}त्वा\appbrk{}त्स्रो\appbrk{}त\appbrk{}सा\appbrk{}म\appbrk{}तस्}\ N H}}\skt{ परि\-ह\-रे\-दु\-त्पन्ने तु }\edtext{\skt{शिला\-ज\-तु\-गु\-ग्गु\-ल\-मू\-त्र\-त्रि\-फ\-ला\-लो\-ह\-र\-जो\-र\-सा\-ञ्ज\-न\-म\-धु\-य\-व\-मु\-द्ग\-को\-र\-दू\-षो\-द्दा\-ल\-का\-दीनां}}{\lemma{\skt{शिला\appbrk{}ज\appbrk{}तु\appbrk{}गु\appbrk{}ग्गु\appbrk{}ल\appbrk{}मू\appbrk{}त्र\appbrk{}त्रि\appbrk{}फ\appbrk{}ला\appbrk{}लो\appbrk{}ह\appbrk{}र\appbrk{}जो\appbrk{}र\appbrk{}सा\appbrk{}ञ्ज\appbrk{}न\appbrk{}म\appbrk{}धु\appbrk{}य\appbrk{}व\appbrk{}मु\appbrk{}द्ग\appbrk{}को\appbrk{}र\appbrk{}दू\appbrk{}षो\appbrk{}द्दा\appbrk{}ल\appbrk{}का\appbrk{}दीनां}}\Afootnote{K; \skt{शिला\appbrk{}ज\appbrk{}तु\appbrk{}गु\appbrk{}ग्गु\appbrk{}लु\appbrk{}गो\appbrk{}मू\appbrk{}त्र\appbrk{}त्रि\appbrk{}फ\appbrk{}ला\appbrk{}लो\appbrk{}ह\appbrk{}र\appbrk{}जो\appbrk{}र\appbrk{}सा\appbrk{}ञ्ज\appbrk{}न\appbrk{}म\appbrk{}धु\appbrk{}य\appbrk{}व\appbrk{}मु\appbrk{}द्ग\appbrk{}को\appbrk{}र\appbrk{}दू\appbrk{}ष\appbrk{}क\appbrk{}श्या\appbrk{}मा\appbrk{}को\appbrk{}द्दा\appbrk{}ल\appbrk{}का\appbrk{}दीनां}\ A; \skt{शिला\appbrk{}ज\appbrk{}तु\appbrk{}गु\appbrk{}ग्गु\appbrk{}लु\appbrk{}मू\appbrk{}त्र\appbrk{}त्रि\appbrk{}फ\appbrk{}ला\appbrk{}लो\appbrk{}ह\appbrk{}र\appbrk{}जो\appbrk{}र\appbrk{}सा\appbrk{}ञ्ज\appbrk{}न\appbrk{}म\appbrk{}धु\appbrk{}य\appbrk{}व\appbrk{}मु\appbrk{}द्ग\appbrk{}को\appbrk{}र\appbrk{}दू\appbrk{}षो\appbrk{}द्दा\appbrk{}ल\appbrk{}का\appbrk{}दीनां}\ N H}}\skt{ }\edtext{\skt{विरू\-क्ष\-ण\-च्छे\-द\-नी\-यानां}}{\lemma{\skt{विरू\appbrk{}क्ष\appbrk{}ण\appbrk{}च्छे\appbrk{}द\appbrk{}नी\appbrk{}यानां}}\Afootnote{K N H; \skt{विरू\appbrk{}क्ष\appbrk{}ण\appbrk{}च्छे\appbrk{}द\appbrk{}नी\appbrk{}यानां च}\ A}}\skt{ }\edtext{\skt{द्रव्याणां}}{\lemma{\skt{द्रव्याणां}}\Afootnote{A K N; \skt{द्रव्यानां}\ H}}\skt{ }\edtext{\skt{विधि\-व\-दु\-प\-योगो}}{\lemma{\skt{विधि\appbrk{}व\appbrk{}दु\appbrk{}प\appbrk{}योगो}}\Afootnote{K N H; \skt{विधि\appbrk{}व\appbrk{}दु\appbrk{}प\appbrk{}योगो}\ A}}\skt{ }\edtext{\skt{व्याया\-म\-ले\-ख\-न\-ब\-स्त्यु\-प\-यो\-ग\-श्चेति}}{\lemma{\skt{व्याया\appbrk{}म\appbrk{}ले\appbrk{}ख\appbrk{}न\appbrk{}ब\appbrk{}स्त्यु\appbrk{}प\appbrk{}यो\appbrk{}गश्}}\Afootnote{K N H; \skt{व्यायामो लेख\appbrk{}न\appbrk{}ब\appbrk{}स्त्यु\appbrk{}प\appbrk{}यो\appbrk{}गश्}\ A}}\skt{~॥}\pend

\pstart \ledsidenote{\passref{1.15.33}}\skt{तत्र }\edtext{\skt{पुन\-र्वा\-त\-ला\-हा\-र\-से\-वि\-नो\-ति\-व्य\-वा\-य\-व्या\-या\-मा\-ध्य\-य\-न\-चि\-न्ता\-भ\-य\-शो\-क\-रा\-त्रि\-जा\-ग\-र\-ण\-पि\-पा\-सा\-क्षु\-त्क्ष\-या\-ल्पा\-स\-न\-प्र\-भृ\-ति\-भि\-रु\-प\-शो\-षितो}}{\lemma{\skt{वात\appbrk{}ला\appbrk{}हा\appbrk{}र\appbrk{}से\appbrk{}वि\appbrk{}नो\appbrk{}ति\appbrk{}व्य\appbrk{}वा\appbrk{}य\appbrk{}व्या\appbrk{}या\appbrk{}मा\appbrk{}ध्य\appbrk{}य\appbrk{}न\appbrk{}चि\appbrk{}न्ता\appbrk{}भ\appbrk{}य\appbrk{}शो\appbrk{}क\appbrk{}रा\appbrk{}त्रि\appbrk{}जा\appbrk{}ग\appbrk{}र\appbrk{}ण\appbrk{}पि\appbrk{}पा\appbrk{}सा\appbrk{}क्षु\appbrk{}त्क्ष\appbrk{}या\appbrk{}ल्पा\appbrk{}स\appbrk{}न\appbrk{}प्र\appbrk{}भृ\appbrk{}ति\appbrk{}भिर्}}\Afootnote{K; \skt{वात\appbrk{}ला\appbrk{}हा\appbrk{}र\appbrk{}से\appbrk{}वि\appbrk{}नोऽ\appbrk{}ति\appbrk{}व्या\appbrk{}या\appbrk{}म\appbrk{}व्य\appbrk{}वा\appbrk{}या\appbrk{}ध्य\appbrk{}य\appbrk{}न\appbrk{}भ\appbrk{}य\appbrk{}शो\appbrk{}क\appbrk{}ध्या\appbrk{}न\appbrk{}रा\appbrk{}त्रि\appbrk{}जा\appbrk{}ग\appbrk{}र\appbrk{}ण\appbrk{}पि\appbrk{}पा\appbrk{}साक्षु त्कषा\appbrk{}या\appbrk{}ल्पा\appbrk{}श\appbrk{}न\appbrk{}प्र\appbrk{}भृ\appbrk{}ति\appbrk{}भिर्}\ A; \skt{वात\appbrk{}ला\appbrk{}हा\appbrk{}र\appbrk{}से\appbrk{}विनो तिव्य\appbrk{}वा\appbrk{}य\appbrk{}व्या\appbrk{}या\appbrk{}मा\appbrk{}ध्य\appbrk{}य\appbrk{}न\appbrk{}चि\appbrk{}न्ता\appbrk{}भ\appbrk{}य\appbrk{}शो\appbrk{}क\appbrk{}रा\appbrk{}त्रि\appbrk{}जा\appbrk{}ग\appbrk{}र\appbrk{}ण\appbrk{}पि\appbrk{}पा\appbrk{}सा\appbrk{}क्षु\appbrk{}त्क्ष\appbrk{}या\appbrk{}ल्पा\appbrk{}श\appbrk{}न\appbrk{}प्र\appbrk{}भृ\appbrk{}ति\appbrk{}भिर्}\ N; \skt{व्वात\appbrk{}ला\appbrk{}हा\appbrk{}र\appbrk{}से\appbrk{}वि\appbrk{}नोऽ\appbrk{}ति\appbrk{}व्य\appbrk{}वा\appbrk{}य\appbrk{}व्या\appbrk{}या\appbrk{}मा\appbrk{}ध्य\appbrk{}य\appbrk{}न\appbrk{}चि\appbrk{}स्ता\appbrk{}भ\appbrk{}य\appbrk{}शो\appbrk{}क\appbrk{}रा\appbrk{}त्रि\appbrk{}जा\appbrk{}ग\appbrk{}र\appbrk{}ण\appbrk{}पि\appbrk{}पा\appbrk{}सा\appbrk{}क्षु\appbrk{}त्क्षा\appbrk{}या\appbrk{}ल्पा\appbrk{}स\appbrk{}न\appbrk{}प्र\appbrk{}भृ\appbrk{}ति\appbrk{}भिर्}\ H}}\skt{ रसधातुः }\edtext{\skt{शरी\-र\-म\-नु\-क्र\-म\-मा\-णौ\-ल्प\-त्वान्न}}{\lemma{\skt{शरी\appbrk{}र\appbrk{}म\appbrk{}नु\appbrk{}क्र\appbrk{}म\appbrk{}मा\appbrk{}णौ\appbrk{}ल्प\appbrk{}त्वान्न}}\Afootnote{K; \skt{शरी\appbrk{}र\appbrk{}म\appbrk{}न\appbrk{}नु\appbrk{}क्रा\appbrk{}म\appbrk{}न्न\appbrk{}ल्प\appbrk{}त्वान्न}\ A; \skt{शरीर …}\ N; \skt{शरी\appbrk{}र\appbrk{}म\appbrk{}नु\appbrk{}क्र\appbrk{}म\appbrk{}मा\appbrk{}णोऽ\appbrk{}ल्प\appbrk{}त्वान्न}\ H}}\skt{ }\edtext{\skt{प्रीणायति}}{\lemma{\skt{प्रीणायति}}\Afootnote{K; \skt{प्रीणाति}\ A; \skt{…ति}\ N; \skt{प्रीणयति}\ H}}\skt{ तस्मादतिकार्श्यं भवति~। }\edtext{\skt{सोतिकृशः}}{\lemma{\skt{सोतिकृशः}}\Afootnote{K; \skt{सोऽतिकृशः}\ A N; \skt{सोऽतिकृषः}\ H}}\skt{ }\edtext{\edtext{\skt{क्षुत्पि\-पा\-सा\-शी\-तो\-ष्ण\-वा\-त\-व\-र्षा\-भा\-रा\-दा\-ने\-ष्व\-स\-हि\-ष्णु\-र्वा\-त\-रो\-ग\-प्रा\-यो\-ल्प\-प्राणः}}{\lemma{\skt{क्षुत्पि\appbrk{}पा\appbrk{}सा\appbrk{}शी\appbrk{}तो\appbrk{}ष्ण\appbrk{}वा\appbrk{}त\appbrk{}व\appbrk{}र्षा\appbrk{}भा\appbrk{}रा\appbrk{}दा\appbrk{}नेष्व्}}\Afootnote{K N H; \skt{क्षुत्पि\appbrk{}पा\appbrk{}सा\appbrk{}शी\appbrk{}तो\appbrk{}ष्ण\appbrk{}वा\appbrk{}त\appbrk{}व\appbrk{}र्ष\appbrk{}भा\appbrk{}रा\appbrk{}दा\appbrk{}नेष्व्}\ A}}}{\lemma{\skt{अस\appbrk{}हि\appbrk{}ष्णु\appbrk{}र्वा\appbrk{}त\appbrk{}रो\appbrk{}ग\appbrk{}प्रा\appbrk{}यो\appbrk{}ल्प\appbrk{}प्राणः}}\Afootnote{K; \skt{अस\appbrk{}हि\appbrk{}ष्णु\appbrk{}र्वा\appbrk{}त\appbrk{}रो\appbrk{}ग\appbrk{}प्रा\appbrk{}योऽ\appbrk{}ल्प\appbrk{}प्रा\appbrk{}णश्च}\ A; \skt{अस\appbrk{}हि\appbrk{}ष्णु\appbrk{}र्वा\appbrk{}त\appbrk{}रो\appbrk{}ग\appbrk{}प्रायो ल्पप्राणः}\ N; \skt{अस\appbrk{}हि\appbrk{}ष्णु\appbrk{}र्व्वा\appbrk{}त\appbrk{}रो\appbrk{}ग\appbrk{}प्रा\appbrk{}योऽ\appbrk{}ल्प\appbrk{}प्राणः}\ H}}\skt{ }\edtext{\skt{क्रियासु च}}{\lemma{\skt{क्रियासु च}}\Afootnote{K N H; \skt{क्रियासु}\ A}}\skt{ भवति~। }\edtext{\skt{का }\edtext{\skt{प्लीहो\-द\-रा\-ग्नि\-सा\-द\-गु\-ल्म\-र\-क्त\-पि\-त्ता\-ना\-म\-न्य\-तमं}}{\lemma{\skt{अन्यतमं}}\Afootnote{K N; \skt{अन्यतम}\ A; \skt{अन्यतस}\ H}}}{\lemma{\skt{का प्लीहो\appbrk{}द\appbrk{}रा\appbrk{}ग्नि\appbrk{}सा\appbrk{}द\appbrk{}गु\appbrk{}ल्म\appbrk{}र\appbrk{}क्त\appbrk{}पि\appbrk{}त्ता\appbrk{}नाम्}}\Afootnote{K; \skt{श्वास\appbrk{}का\appbrk{}स\appbrk{}शो\appbrk{}ष\appbrk{}प्ली\appbrk{}हो\appbrk{}द\appbrk{}रा\appbrk{}ग्नि\appbrk{}सा\appbrk{}द\appbrk{}गु\appbrk{}ल्म\appbrk{}र\appbrk{}क्त\appbrk{}पि\appbrk{}त्ता\appbrk{}नाम्}\ A; \skt{कास\appbrk{}श्वा\appbrk{}स\appbrk{}प्ली\appbrk{}हो\appbrk{}द\appbrk{}रा\appbrk{}ग्नि\appbrk{}सा\appbrk{}द\appbrk{}गु\appbrk{}ल्म\appbrk{}र\appbrk{}क्त\appbrk{}पि\appbrk{}त्ता\appbrk{}नाम्}\ N H}}\skt{ }\edtext{\skt{प्राप्य}}{\lemma{\skt{प्राप्य}}\Afootnote{K N; \skt{मासाद्य}\ A; \skt{मप्राप्य}\ H}}\skt{ मर\-ण\-मु\-प\-याति सर्व एव चास्य रोगा बलवन्तो }\edtext{\skt{भवन्ति}}{\lemma{\skt{भवन्ति}}\Afootnote{K N H; \skt{भवन्}\ A}}\skt{~। }\edtext{\skt{कस्मा\-द\-ल्प\-प्रा\-ण\-त्वा\-द\-त\-स्त\-स्यो\-त्प\-त्ति\-हे\-तु\-म्प\-रि\-ह\-रेत्}}{\lemma{\skt{कस्मा\appbrk{}द\appbrk{}ल्प\appbrk{}प्रा\appbrk{}ण\appbrk{}त्वा\appbrk{}द\appbrk{}तस्}}\Afootnote{K N H; \skt{त्यल्पप्राणत्वात् ; अतस्}\ A}}\skt{~। उत्पन्ने तु }\edtext{\skt{पय\-स्या\-श्व\-ग\-न्धा\-वि\-दा\-री\-वि\-दा\-रि\-ग\-न्धा\-श\-ता\-व\-री\-ना\-ग\-ब\-लानां}}{\lemma{\skt{पय\appbrk{}स्या\appbrk{}श्व\appbrk{}ग\appbrk{}न्धा\appbrk{}वि\appbrk{}दा\appbrk{}री\appbrk{}वि\appbrk{}दा\appbrk{}रि\appbrk{}ग\appbrk{}न्धा\appbrk{}श\appbrk{}ता\appbrk{}व\appbrk{}री\appbrk{}ना\appbrk{}ग\appbrk{}ब\appbrk{}लानां}}\Afootnote{K H; \skt{पय\appbrk{}स्या\appbrk{}श्व\appbrk{}ग\appbrk{}न्धा\appbrk{}वि\appbrk{}दा\appbrk{}रि\appbrk{}ग\appbrk{}न्धा\appbrk{}श\appbrk{}ता\appbrk{}व\appbrk{}री\appbrk{}ब\appbrk{}ला\appbrk{}ति\appbrk{}ब\appbrk{}ला\appbrk{}ना\appbrk{}ग\appbrk{}ब\appbrk{}लानां}\ A; \skt{पय\appbrk{}स्या\appbrk{}श्व\appbrk{}ग\appbrk{}न्धा\appbrk{}दि\appbrk{}वि\appbrk{}दा\appbrk{}री\appbrk{}वि\appbrk{}दा\appbrk{}रि\appbrk{}ग\appbrk{}न्ध\appbrk{}श\appbrk{}ता\appbrk{}व\appbrk{}री\appbrk{}ना\appbrk{}ग\appbrk{}ब\appbrk{}ला\appbrk{}नाम्}\ N}}\skt{ }\edtext{\skt{मधुराणां}}{\lemma{\skt{मधुराणां}}\Afootnote{K N H; \skt{मधुराणा}\ A}}\skt{ }\edtext{\skt{मन्यासां}}{\lemma{\skt{मन्यासां}}\Afootnote{A K; \skt{अन्यासां}\ N; \skt{मन्यासञ्}\ H}}\skt{ }\edtext{\skt{चौषधीनां}}{\lemma{\skt{चौषधीनां}}\Afootnote{A K N H\textsuperscript{ac}; \skt{चौषधीनां ल}\ H\textsuperscript{pc}}}\skt{ }\edtext{\skt{विधि\-व\-दु\-प\-योगः}}{\lemma{\skt{विधि\appbrk{}व\appbrk{}दु\appbrk{}प\appbrk{}योगः}}\Afootnote{K; \skt{उपयोगः}\ A; \skt{विधि\appbrk{}व\appbrk{}दु\appbrk{}प\appbrk{}योगः}\ N H\textsuperscript{ac} H\textsuperscript{pc}}}\skt{~। }\edtext{\skt{क्षीर\-द\-धि\-घृ\-त\-मां\-स\-शा\-लि\-ष\-ष्टि\-क\-य\-व\-गो\-धू\-मानां च}}{\lemma{\skt{क्षीर\appbrk{}द\appbrk{}धि\appbrk{}घृ\appbrk{}त\appbrk{}मां\appbrk{}स\appbrk{}शा\appbrk{}लि\appbrk{}ष\appbrk{}ष्टि\appbrk{}क\appbrk{}य\appbrk{}व\appbrk{}गो\appbrk{}धू\appbrk{}मानां च}}\Afootnote{A K H; \skt{क्षीर\appbrk{}द\appbrk{}धि\appbrk{}घृ\appbrk{}त\appbrk{}मा\appbrk{}न्स\appbrk{}शा\appbrk{}लि\appbrk{}ष\appbrk{}ष्टि\appbrk{}य\appbrk{}व\appbrk{}गो\appbrk{}धू\appbrk{}मा\appbrk{}नाञ्च}\ N}}\skt{ }\edtext{\edtext{\skt{दिवा\-स्व\-प्न\-ब्र\-ह्म\-च\-र्या\-व्या\-या\-म\-बृं\-ह\-ण\-ब\-स्त्यु\-प\-यो\-ग\-श्चेति}}{\lemma{\skt{दिवा\appbrk{}स्व\appbrk{}प्न\appbrk{}ब्र\appbrk{}ह्म\appbrk{}च\appbrk{}र्या\appbrk{}व्या\appbrk{}या\appbrk{}म\appbrk{}बृं\appbrk{}ह\appbrk{}ण\appbrk{}ब\appbrk{}स्त्यु\appbrk{}प\appbrk{}यो\appbrk{}गश्}}\Afootnote{K N; \skt{दिवा\appbrk{}स्व\appbrk{}प्न\appbrk{}ब्र\appbrk{}ह्म\appbrk{}च\appbrk{}र्या\appbrk{}व्या\appbrk{}या\appbrk{}म\appbrk{}बृं\appbrk{}ह\appbrk{}ण\appbrk{}ब\appbrk{}स्त्यु\appbrk{}प\appbrk{}योग}\ A; \skt{दिवा\appbrk{}स्व\appbrk{}प्न\appbrk{}ब्र\appbrk{}ह्म\appbrk{}च\appbrk{}र्या\appbrk{}व्या\appbrk{}या\appbrk{}म\appbrk{}बृ\appbrk{}ह\appbrk{}ण\appbrk{}ब\appbrk{}स्त्यु\appbrk{}प\appbrk{}यो\appbrk{}गश्}\ H}}}{\lemma{\skt{चेति}}\Afootnote{K N H; \skt{श्चेति}\ A}}\skt{~।}\pend

\pstart \ledsidenote{\passref{1.15.34}}\skt{यः }\edtext{\edtext{\skt{पुन\-रु\-भ\-य\-सा\-धा\-र\-णा\-न्यु\-प\-से\-वते}}{\lemma{\skt{उभ\appbrk{}य\appbrk{}सा\appbrk{}धा\appbrk{}र\appbrk{}णान्य्}}\Afootnote{K N H; \skt{उभ\appbrk{}य\appbrk{}सा\appbrk{}धा\appbrk{}र\appbrk{}णान्या}\ A}}}{\lemma{\skt{उपसेवते}}\Afootnote{K N; \skt{सेवेत}\ A; \skt{उपसेवेत}\ H}}\skt{ तस्यान्नरसः }\edtext{\skt{शरी\-र\-म\-नु\-क्र\-म\-माणः}}{\lemma{\skt{अनुक्रममाणः}}\Afootnote{K N H; \skt{अनुक्राम}\ A}}\skt{ }\edtext{\skt{समा\-न्धा\-तू\-नु\-प\-चि\-नोति}}{\lemma{\skt{समा\appbrk{}न्धा\appbrk{}तू\appbrk{}नु\appbrk{}प\appbrk{}चि\appbrk{}नोति}}\Afootnote{K N H; \skt{न्समा\appbrk{}न्धा\appbrk{}तू\appbrk{}नु\appbrk{}प\appbrk{}चि\appbrk{}नोति}\ A}}\skt{ }\edtext{\skt{सम\-धा\-तु\-त्वा\-न्म\-ध्य\-श\-रीरो}}{\lemma{\skt{सम\appbrk{}धा\appbrk{}तु\appbrk{}त्वा\appbrk{}न्म\appbrk{}ध्य\appbrk{}श\appbrk{}रीरो}}\Afootnote{K; \skt{सम\appbrk{}धा\appbrk{}तु\appbrk{}त्वा\appbrk{}न्म\appbrk{}ध्य\appbrk{}श\appbrk{}रीरो}\ A; \skt{सधा\appbrk{}तु\appbrk{}त्वा\appbrk{}न्म\appbrk{}ध्य\appbrk{}श\appbrk{}रीरो}\ N; \skt{सम\appbrk{}धा\appbrk{}तु\appbrk{}त्वा\appbrk{}त्म\appbrk{}ध्य\appbrk{}श\appbrk{}रीरो}\ H}}\skt{ भवति~। }\edtext{\skt{सर्वक्रियासु च}}{\lemma{\skt{सर्वक्रियासु च}}\Afootnote{K N H; \skt{सर्वक्रियासु}\ A}}\skt{ समर्थः क्षुत्पि\-पा\-सा\-शी\-तो\-ष्ण\-वा\-त\-व\-र्षा\-त\-प\-सहो }\edtext{\skt{बलवांश्च}}{\lemma{\skt{बलवांश्च}}\Afootnote{K N H; \skt{बलवांश्च}\ A}}\skt{ }\edtext{\skt{भवति सः}}{\lemma{\skt{भवति सः}}\Afootnote{K N H; \skt{स}\ A}}\skt{~। }\edtext{\edtext{\skt{सत\-त\-म\-नु\-पा\-ल\-यि\-तव्य}}{\lemma{\skt{सततम्}}\Afootnote{A N H; \skt{सतत}\ K}}}{\lemma{\skt{अनु\appbrk{}पा\appbrk{}ल\appbrk{}यि\appbrk{}तव्य}}\Afootnote{A N H; \skt{मनु\appbrk{}पा\appbrk{}ल\appbrk{}यि\appbrk{}तव्य}\ K}}\skt{ इति~॥}\pend

\pstart\versepar \ledsidenote{\passref{1.15.35}}\edtext{\skt{भवति}}{\lemma{\skt{भवति}}\Afootnote{N; \skt{भवन्ति}\ A; \emph{om.}\ K H}}\skt{ चात्र~॥ }\pend
\pstart\versepar \skt{}\edtext{\edtext{\skt{द्वावप्येतौ}}{\lemma{\skt{अप्य्}}\Afootnote{K N H; \emph{om.}\ A}}}{\lemma{\skt{एतौ}}\Afootnote{K N H; \skt{अत्यन्}\ A}}\skt{ }\edtext{\skt{विगहितौ}}{\lemma{\skt{विगहितौ}}\Afootnote{K; \skt{तगर्हितावेतौ}\ A; \skt{विगर्हितौ}\ N H}}\skt{ सदा }\edtext{\skt{स्थूलकृशौ}}{\lemma{\skt{स्थूलकृशौ}}\Afootnote{A K N; \skt{स्थूलकृषौ}\ H}}\skt{ नरौ~। }\pend
\pstart\versepar \skt{श्रेष्ठो मध्यशरीरस्तु कृशः }\edtext{\skt{स्थूला तु}}{\lemma{\skt{स्थूला तु}}\Afootnote{K; \skt{स्थूलात्तु}\ A N H}}\skt{ पूजितः~।}\edtext{}{\lemma{\skt{चात्र~॥ द्वाव्}}\Afootnote{N; \skt{चात्र}\ A; \skt{भवति चात्र~॥ द्वाव्}\ K; \skt{भवति चात्र~॥ द्वाव्}\ H}}\pend

\pstart\versepar \ledsidenote{\passref{1.15.36}}\skt{दोषः प्रकुपितो }\edtext{\skt{धातूं}}{\lemma{\skt{धातूं}}\Afootnote{K N; \skt{धातू}\ A H}}\skt{ }\edtext{\edtext{\skt{क्षप\-य\-त्या\-म\-ते\-जसा}}{\lemma{\skt{क्षपयत्य्}}\Afootnote{K; \skt{न्क्षपयत्या}\ A; \skt{क्षपयतित्य्}\ N; \skt{न्क्षपयत्य्}\ H}}}{\lemma{\skt{आमतेजसा}}\Afootnote{K; \skt{त्मतेजसा}\ A; \skt{आत्मतेजसा}\ N H}}\skt{~। }\pend
\pstart\versepar \skt{इद्धः स्वतेजसा }\edtext{\edtext{\edtext{\skt{वह्नि\-रु\-खा\-ग\-त\-मि\-वो\-दकं}}{\lemma{\skt{वह्निर्}}\Afootnote{K N; \skt{वह्नि}\ A H}}}{\lemma{\skt{उखागतम्}}\Afootnote{K N; \skt{रुखागत}\ A; \skt{रुग्वागतम्}\ H}}}{\lemma{\skt{इवोदकं}}\Afootnote{K N H; \skt{मिवोदकम्}\ A}}\skt{~॥}\pend

\pstart\versepar \ledsidenote{\passref{1.15.37}}\edtext{\edtext{\skt{वैल\-क्ष\-ण्या\-च्छ\-री\-रा\-णा\-म\-स्था\-यि\-त्वा\-त्त\-थैव}}{\lemma{\skt{वैल\appbrk{}क्ष\appbrk{}ण्या\appbrk{}च्छ\appbrk{}री\appbrk{}रा\appbrk{}णाम्}}\Afootnote{K N H; \skt{वैल\appbrk{}क्ष\appbrk{}ण्या\appbrk{}च्छ\appbrk{}री\appbrk{}राणा}\ A}}\skt{ च}}{\lemma{\skt{अस्था\appbrk{}यि\appbrk{}त्वा\appbrk{}त्त\appbrk{}थैव च}}\Afootnote{K N H; \skt{मस्था\appbrk{}यि\appbrk{}त्वा\appbrk{}त्त\appbrk{}थैव च}\ A}}\skt{~। }\pend
\pstart\versepar \skt{दोष\-धा\-तु\-म\-लानां तु }\edtext{\skt{परिमाणान्न}}{\lemma{\skt{परिमाणान्न}}\Afootnote{K; \skt{परिमाणं न}\ A N H}}\skt{ विद्यते~।}\pend

\pstart\versepar \ledsidenote{\passref{1.15.38}}\skt{एषां समत्वं यच्चापि भिष\-ग्भि\-र\-भि\-धी\-यते~। }\pend
\pstart\versepar \skt{न तत्स्वास्थ्यादृते शक्यं }\edtext{\edtext{\skt{वक्तुमन्येन}}{\lemma{\skt{वक्तुम्}}\Afootnote{K N H; \skt{वक्तु}\ A}}}{\lemma{\skt{अन्येन}}\Afootnote{K N H; \skt{मन्येन}\ A}}\skt{ हेतुना~।}\edtext{}{\lemma{\skt{अभिधीयते~। न}}\Afootnote{K N H; \skt{अवधार्यते~। न}\ A}}\pend

\pstart\versepar \ledsidenote{\passref{1.15.39}}\edtext{\skt{दोषादीनान्तु}}{\lemma{\skt{दोषादीनान्तु}}\Afootnote{K N H; \skt{दोषादीनां त्व}\ A}}\skt{ सम\-ता\-म\-नु\-मा\-नेन लक्षयेत्~। }\pend
\pstart\versepar \skt{प्रसन्नात्मेन्द्रियं }\edtext{\skt{ज्ञात्वा}}{\lemma{\skt{ज्ञात्वा}}\Afootnote{K N; \skt{वीक्ष्य}\ A; \emph{om.}\ H}}\skt{ }\edtext{\edtext{\skt{पुरुषन्तत्र}}{\lemma{\skt{पुरुषन्}}\Afootnote{A K N; \emph{om.}\ H}}}{\lemma{\skt{तत्र}}\Afootnote{K N; \skt{कुश}\ A; \emph{om.}\ H}}\skt{ }\edtext{\skt{बुद्धिमां}}{\lemma{\skt{बुद्धिमां}}\Afootnote{K N; \skt{लो भिषक्}\ A; \emph{om.}\ H}}\skt{~।}\edtext{}{\lemma{\skt{लक्षयेत्~। प्रसन्नात्मेन्द्रियं}}\Afootnote{K; \skt{लक्षयेत्~। अप्रसन्नेन्द्रियं}\ A; \skt{लक्षयेत्~। प्रसन्नामेन्द्रियं}\ N; \skt{लक्षयेत्}\ H}}\pend

\pstart\versepar \ledsidenote{\passref{1.15.40cdef}}\skt{क्षप\-ये\-द्बृं\-ह\-ये\-च्चापि }\edtext{\skt{दोषधातुमलां}}{\lemma{\skt{दोषधातुमलां}}\Afootnote{K N; \skt{दोषधातुमला}\ A H}}\skt{ भिषक्~। }\pend
\pstart\versepar \skt{ताव\-द्या\-व\-द\-रोगः }\edtext{\skt{स्यादेतत्साम्यस्य}}{\lemma{\skt{स्यादेतत्साम्यस्य}}\Afootnote{A K N; \skt{स्यादेतत्सात्म्यस्य}\ H}}\skt{ लक्षणं~।}\edtext{}{\lemma{\skt{भिषक्~। ताव\appbrk{}द्या\appbrk{}व\appbrk{}द\appbrk{}रोगः}}\Afootnote{K N; \skt{न्भिष\appbrk{}क्ता\appbrk{}व\appbrk{}द्या\appbrk{}व\appbrk{}द\appbrk{}रोगः}\ A H}}\pend

\pstart\versepar \ledsidenote{\passref{1.15.41}}\skt{समदोषः समाग्निश्च सम\-धा\-तु\-म\-ल\-क्रियः~। }\pend
\pstart\versepar \skt{प्रस\-न्ना\-त्मे\-न्द्रि\-य\-मनाः स्वस्थ }\edtext{\skt{इत्यभिधीयते}}{\lemma{\skt{अभिधीयते}}\Afootnote{A K N; \skt{उपदिश्यते}\ H\textsuperscript{ac}; \skt{उप\appbrk{}दि\appbrk{}श्य\appbrk{}ते\appbrk{}त्य\appbrk{}भि\appbrk{}धी\appbrk{}यते}\ H\textsuperscript{pc}}}\skt{~॥}\pend

\pstart \ledsidenote{\passref{1.15.trailer}}\edtext{\skt{१५}}{\lemma{\skt{१५}}\Afootnote{\skt{इति सुश्रु\appbrk{}त\appbrk{}सं\appbrk{}हि\appbrk{}तायां सूत्रस्थाने दोष\appbrk{}धा\appbrk{}तु\appbrk{}म\appbrk{}ल\appbrk{}क्ष\appbrk{}य\appbrk{}वृ\appbrk{}द्धि\appbrk{}वि\appbrk{}ज्ञा\appbrk{}नीयो नाम पञ्चदशोऽध्यायः}\ A; \skt{ण्डोहू~॥ {\symfont ❈}~॥ १५}\ K; \skt{ण्ट ःऋ~॥ ॰}\ N; \skt{॥ डोहू}\ H}}\pend

\endnumbering
\beginnumbering
\eledsection*{Adhyāya 16}
\pstart \ledsidenote{\passref{1.16.1}}\skt{अथातः }\edtext{\skt{कर्ण्णव्यधविधिं}}{\lemma{\skt{कर्ण्णव्यधविधिं}}\Afootnote{K; \skt{कर्ण\appbrk{}व्य\appbrk{}ध\appbrk{}ब\appbrk{}न्ध\appbrk{}वि\appbrk{}धिम्}\ A; \skt{कर्णव्यध}\ N H}}\skt{ }\edtext{\skt{व्याख्यास्यामः}}{\lemma{\skt{व्याख्यास्यामः}}\Afootnote{\skt{अध्यायं व्याख्यास्यामः}\ A; \emph{om.}\ K; \skt{विधिम्व्याख्यास्यामः}\ N H}}\skt{~॥१॥}\pend

\pstart \ledsidenote{\passref{1.16.3}}\skt{रक्षा\-भू\-ष\-ण\-नि\-मि\-त्त\-म्बा\-लस्य कर्ण्णौ }\edtext{\skt{व्यधयेत्~। तौ}}{\lemma{\skt{व्यधयेत्~। तौ}}\Afootnote{K N H; \skt{विध्येते~। तौ}\ A}}\skt{ षष्ठे }\edtext{\skt{मासे}}{\lemma{\skt{मासे}}\Afootnote{K N H; \skt{मासि}\ A}}\skt{ }\edtext{\skt{सप्तमे वा}}{\lemma{\skt{सप्तमे वा}}\Afootnote{A K H; \skt{वा}\ N}}\skt{ शुक्लपक्षे प्रशस्तेषु तिथि\-क\-र\-ण\-मु\-हू\-र्त्त\-न\-क्ष\-त्रेषु }\edtext{\skt{कृत\-म\-ङ्ग\-ल\-स्व\-स्ति\-वा\-चनं}}{\lemma{\skt{कृत\appbrk{}म\appbrk{}ङ्ग\appbrk{}ल\appbrk{}स्व\appbrk{}स्ति\appbrk{}वा\appbrk{}चनं}}\Afootnote{A; \skt{कृतमङ्गलं स्वस्तिवाचनं}\ K N H}}\footnoteA{The compound \emph{kṛtamaṅgalasvastivācanaṃ} is an emendation based on the similar text at Su.śā.3.2.25.}\skt{ }\edtext{\skt{धात्र्यङ्के}}{\lemma{\skt{धात्र्यङ्के}}\Afootnote{A N H; \skt{धात्र्यङ्}\ K}}\skt{ }\edtext{\skt{कुमा\-र\-मु\-प\-वे\-श्या\-भि\-सा\-न्त्व\-य\-मानो}}{\lemma{\skt{कुमा\appbrk{}र\appbrk{}मु\appbrk{}प\appbrk{}वे\appbrk{}श्या\appbrk{}भि\appbrk{}सा\appbrk{}न्त्व\appbrk{}य\appbrk{}मानो}}\Afootnote{H; \skt{कुमारधराङ्के वा कुमा\appbrk{}र\appbrk{}मु\appbrk{}प\appbrk{}वेश्य बालक्रीडनकैः प्रलो\appbrk{}भ्या\appbrk{}भि\appbrk{}सा\appbrk{}न्त्वय}\ A; \skt{को कुमा\appbrk{}र\appbrk{}मु\appbrk{}प\appbrk{}वे\appbrk{}श्या\appbrk{}भि\appbrk{}सा\appbrk{}न्त्व\appbrk{}य\appbrk{}मानः}\ K; \skt{कुमा\appbrk{}र\appbrk{}क\appbrk{}मु\appbrk{}प\appbrk{}वे\appbrk{}श्या\appbrk{}भि\appbrk{}सां\appbrk{}त्व\appbrk{}य\appbrk{}मानः}\ N}}\skt{ }\edtext{\skt{भिष\-ग्वा\-म\-ह\-स्ते\-ना\-कृष्य}}{\lemma{\skt{भिष\appbrk{}ग्वा\appbrk{}म\appbrk{}ह\appbrk{}स्ते\appbrk{}ना\appbrk{}कृष्य}}\Afootnote{K N H; \skt{न्भिष\appbrk{}ग्वा\appbrk{}म\appbrk{}ह\appbrk{}स्ते\appbrk{}ना\appbrk{}कृष्य}\ A}}\skt{ कर्ण्णन्दैवकृते }\edtext{\skt{छिद्रे}}{\lemma{\skt{छिद्रे}}\Afootnote{K N H; \skt{छिद्र}\ A}}\skt{ }\edtext{\skt{दक्षिणहस्तेन}}{\lemma{\skt{दक्षिणहस्तेन}}\Afootnote{K N H; \skt{आदि\appbrk{}त्य\appbrk{}क\appbrk{}रा\appbrk{}व\appbrk{}भा\appbrk{}सिते शनैः शनै\appbrk{}र्द\appbrk{}क्षि\appbrk{}ण\appbrk{}ह\appbrk{}स्तेन}\ A}}\skt{ }\edtext{\skt{ऋजु}}{\lemma{\skt{ऋजु}}\Afootnote{K; \skt{र्जु}\ A; \skt{ऋजुम्}\ N H}}\skt{ }\edtext{\skt{विध्येत्~। पूर्व्वन्दक्षिणं}}{\lemma{\skt{विध्येत्~। पूर्व्वन्}}\Afootnote{K N H; \skt{विध्येत्प्रतनुकं सूच्या बहलमारया पूर्वं}\ A}}\skt{ कुमारस्य }\edtext{\skt{वामङ्कन्यायाः}}{\lemma{\skt{कन्यायाः}}\Afootnote{K\textsuperscript{pc} N H; \skt{कुमार्याः}\ A; \skt{कन्या}\ K\textsuperscript{ac}}}\skt{~। }\edtext{\skt{प्रतनुं}}{\lemma{\skt{प्रतनुं}}\Afootnote{K\textsuperscript{pc}; \emph{om.}\ A K\textsuperscript{ac}; \skt{प्रतनू}\ N H}}\skt{ }\edtext{\skt{सूच्या}}{\lemma{\skt{सूच्या}}\Afootnote{K\textsuperscript{pc} N H; \emph{om.}\ A; \skt{याः}\ K\textsuperscript{ac}}}\skt{ }\edtext{\edtext{\skt{बहलमारया}}{\lemma{\skt{बहलम्}}\Afootnote{K N H; \skt{ततः}\ A}}}{\lemma{\skt{आरया}}\Afootnote{K N H; \skt{पिचुवर्तिं प्रवेशयेत्}\ A}}\skt{~॥२॥}\footnoteA{Ḍalhaṇa recorded the alternative reading \emph{bhakṣyaviśeṣair vā} before \emph{bālakrīḍanakaiḥ pralobhya} in the vulgate.}\pend

\pstart \ledsidenote{\passref{1.16.4}}\edtext{\skt{शोणि\-त\-ब\-हु\-त्वेऽ\-ति\-वे\-द\-नायां}}{\lemma{\skt{ऽतिवेदनायां}}\Afootnote{\skt{न वेदनया}\ A; \skt{तिवेदनायां}\ K N H}}\skt{ चान्य\-दे\-श\-वि\-द्ध\-मिति }\edtext{\skt{जानीयात्~। निरुपद्रवता}}{\lemma{\skt{जानीयात्~। निरुपद्रवता}}\Afootnote{K N H; \skt{जानी\appbrk{}या\appbrk{}त्नि\appbrk{}रु\appbrk{}प\appbrk{}द्र\appbrk{}व\appbrk{}तया}\ A}}\skt{ }\edtext{\skt{तद्दे\-श\-वि\-द्ध\-लि\-ङ्गम्}}{\lemma{\skt{तद्दे\appbrk{}श\appbrk{}वि\appbrk{}द्ध\appbrk{}लि\appbrk{}ङ्गम्}}\Afootnote{K N H; \skt{तद्दे\appbrk{}श\appbrk{}वि\appbrk{}द्ध\appbrk{}मिति}\ A}}\skt{~॥३॥}\footnoteA{At this point, witness K is missing a folio, so the rest of this chapter is constructed on the basis of witnesses N and H.}\pend

\pstart \ledsidenote{\passref{1.16.5}}\edtext{\skt{तत्र यदृच्छाविद्धायां}}{\lemma{\skt{तत्र यदृच्छाविद्धायां}}\Afootnote{N H; \skt{तत्राज्ञेन यदृच्छया विद्धासु}\ A}}\skt{ }\edtext{\edtext{\skt{सिरायामज्ञेन}}{\lemma{\skt{सिरायाम्}}\Afootnote{N H; \skt{सिरासु कालि\appbrk{}का\appbrk{}म\appbrk{}र्म\appbrk{}रि\appbrk{}का\appbrk{}लो\appbrk{}हि\appbrk{}ति\appbrk{}का\appbrk{}सू\appbrk{}प\appbrk{}द्रवा भवन्ति~। तत्र कालिकायां}\ A}}}{\lemma{\skt{अज्ञेन}}\Afootnote{N H; \emph{om.}\ A}}\skt{ }\edtext{\skt{ज्वर\-दा\-ह\-श्व\-य\-थु\-वे\-द\-ना\-ग्र\-न्थि\-म\-न्या\-स्त\-म्भा\-प\-ता\-न\-क\-शि\-रो\-ग्र\-ह\-क\-र्ण्ण\-शू\-लानि}}{\lemma{\skt{ज्वर\appbrk{}दा\appbrk{}ह\appbrk{}श्व\appbrk{}य\appbrk{}थु\appbrk{}वे\appbrk{}द\appbrk{}ना\appbrk{}ग्र\appbrk{}न्थि\appbrk{}म\appbrk{}न्या\appbrk{}स्त\appbrk{}म्भा\appbrk{}प\appbrk{}ता\appbrk{}न\appbrk{}क\appbrk{}शि\appbrk{}रो\appbrk{}ग्र\appbrk{}ह\appbrk{}क\appbrk{}र्ण्ण\appbrk{}शू\appbrk{}लानि}}\Afootnote{\skt{ज्वरो दाहः श्वयथुर्वेदना च भवति मर्मरिकायां वेदना ज्वरो ग्रन्थयश्च लोहितिकायां मन्या\appbrk{}स्त\appbrk{}म्भा\appbrk{}प\appbrk{}ता\appbrk{}न\appbrk{}क\appbrk{}शि\appbrk{}रो\appbrk{}ग्र\appbrk{}ह\appbrk{}क\appbrk{}र्ण\appbrk{}शू\appbrk{}लानि}\ A; \skt{ज्वर\appbrk{}दा\appbrk{}ह\appbrk{}श्व\appbrk{}य\appbrk{}थु\appbrk{}र्वे\appbrk{}द\appbrk{}ना\appbrk{}ग्र\appbrk{}न्थि\appbrk{}म\appbrk{}न्या\appbrk{}स्त\appbrk{}म्भा\appbrk{}प\appbrk{}ता\appbrk{}न\appbrk{}क\appbrk{}शि\appbrk{}रो\appbrk{}ग्र\appbrk{}ह\appbrk{}क\appbrk{}र्ण्ण\appbrk{}शू\appbrk{}लानि}\ N H}}\skt{ }\edtext{\skt{भवन्ति}}{\lemma{\skt{भवन्ति}}\Afootnote{N H; \skt{भवन्ति~। तेषु यथास्वं प्रतिकुर्वीत}\ A}}\skt{~॥४॥}\pend

\pstart \ledsidenote{\passref{1.16.6}}\edtext{\skt{दोष\-स\-मु\-द\-या\-द\-प्र\-श\-स्त\-व्य\-धाद्वा}}{\lemma{\skt{दोष\appbrk{}स\appbrk{}मु\appbrk{}द\appbrk{}या\appbrk{}द\appbrk{}प्र\appbrk{}श\appbrk{}स्त\appbrk{}व्य\appbrk{}धाद्वा}}\Afootnote{N H; \skt{क्लिष्ट\appbrk{}जि\appbrk{}ह्मा\appbrk{}प्र\appbrk{}श\appbrk{}स्त\appbrk{}सू\appbrk{}ची\appbrk{}व्य\appbrk{}धा\appbrk{}द्गा\appbrk{}ढ\appbrk{}त\appbrk{}र\appbrk{}व\appbrk{}र्ति\appbrk{}त्वा\appbrk{}द्दो\appbrk{}ष\appbrk{}स\appbrk{}मु\appbrk{}दा\appbrk{}या\appbrk{}द\appbrk{}प्र\appbrk{}श\appbrk{}स्त\appbrk{}व्य\appbrk{}धाद्वा}\ A}}\skt{ }\edtext{\skt{तत्र}}{\lemma{\skt{तत्र}}\Afootnote{N H; \skt{यत्र संरम्भो वेदना वा भवति तत्र}\ A}}\skt{ }\edtext{\skt{वर्त्तिमपहृत्य}}{\lemma{\skt{अपहृत्य}}\Afootnote{N H; \skt{उपहृत्या}\ A}}\skt{ }\edtext{\edtext{\skt{यव\-म\-धु\-क\-म\-ञ्जि\-ष्ठा\-ग\-न्ध\-र्व्व\-ह\-स्त\-मू\-लै\-र्म्म\-धु\-घृ\-त\-प्र\-गा\-ढै\-रा\-ले\-प\-येत्}}{\lemma{\skt{यव\appbrk{}म\appbrk{}धु\appbrk{}क\appbrk{}म\appbrk{}ञ्जि\appbrk{}ष्ठा\appbrk{}ग\appbrk{}न्ध\appbrk{}र्व्व\appbrk{}ह\appbrk{}स्त\appbrk{}मू\appbrk{}लैर्}}\Afootnote{H; \skt{शु मधु\appbrk{}कै\appbrk{}र\appbrk{}ण्ड\appbrk{}मू\appbrk{}ल\appbrk{}म\appbrk{}ञ्जि\appbrk{}ष्ठा\appbrk{}य\appbrk{}व\appbrk{}ति\appbrk{}ल\appbrk{}क\appbrk{}ल्कैर्}\ A; \skt{यव\appbrk{}म\appbrk{}धु\appbrk{}क\appbrk{}म\appbrk{}ञ्जि\appbrk{}ष्ठा\appbrk{}ग\appbrk{}न्द\appbrk{}र्व\appbrk{}ह\appbrk{}स्त\appbrk{}मूलै}\ N}}\skt{~। }\edtext{\skt{सुरू\-ढ\-ञ्चै\-न\-म्पु\-न\-र्व्वि\-ध्येत्}}{\lemma{\skt{व्विध्येत्}}\Afootnote{N H; \skt{विध्येत्विधानं तु पूर्वोक्तमेव}\ A}}}{\lemma{\skt{आलेपयेत्~। सुरूढञ्}}\Afootnote{N H; \skt{आले\appbrk{}प\appbrk{}ये\appbrk{}त्ता\appbrk{}व\appbrk{}द्या\appbrk{}व\appbrk{}त्सु\appbrk{}रूढ इति सुरूढं}\ A}}\skt{~॥५॥}\footnoteA{Ḍalhaṇa (1.16.6) stated that some do not read \emph{surūḍhañ cainam punar vidhyet}.}\pend

\pstart \ledsidenote{\passref{1.16.7}}\edtext{\edtext{\skt{सम्य\-ग्वि\-द्ध\-मा\-म\-तै\-ल\-प\-रि\-षे\-के\-णो\-प\-च\-रेत्}}{\lemma{\skt{सम्यग्विद्धम्}}\Afootnote{N H; \skt{तत्र सम्यग्विद्धम्}\ A}}\skt{~। त्र्यहात्त्र्यहाद्वर्त्तिं}}{\lemma{\skt{आम\appbrk{}तै\appbrk{}ल\appbrk{}प\appbrk{}रि\appbrk{}षे\appbrk{}के\appbrk{}णो\appbrk{}प\appbrk{}च\appbrk{}रेत्~। त्र्यहात्त्र्यहाद्वर्त्तिं}}\Afootnote{N; \skt{आमतैलेन परि\appbrk{}षे\appbrk{}च\appbrk{}ये\appbrk{}त्त्र्य\appbrk{}हा\appbrk{}त्त्र्य\appbrk{}हाच्च वर्तिं}\ A; \skt{आम\appbrk{}तै\appbrk{}ल\appbrk{}प\appbrk{}रि\appbrk{}षे\appbrk{}क\appbrk{}णो\appbrk{}प\appbrk{}च\appbrk{}रेत्~। त्र्यहात्त्र्यहाद्वर्त्तिं}\ H}}\skt{ }\edtext{\skt{स्थूलतरीं}}{\lemma{\skt{स्थूलतरीं}}\Afootnote{N H; \skt{स्थूलतरां}\ A}}\footnoteA{The unusual form \emph{sthūlatarīṃ} is supported by both manuscripts and we have retained it in spite of only meagre evidence for the form in epic Sanskrit.}\skt{ }\edtext{\skt{कुर्व्वीत}}{\lemma{\skt{कुर्व्वीत}}\Afootnote{N H; \skt{दद्या}\ A}}\skt{ }\edtext{\skt{परिषेकञ्च तमेव}}{\lemma{\skt{परिषेकञ्च तम्}}\Afootnote{N H; \skt{त्परिषेकं च तम्}\ A}}\skt{~॥६॥}\pend

\pstart \ledsidenote{\passref{1.16.8}}\skt{अथ व्यप\-ग\-त\-दो\-षो\-प\-द्रवे }\edtext{\skt{कर्ण्णेऽ\-लं\-प्र\-व\-र्द्ध\-नार्थं}}{\lemma{\skt{ऽलं\appbrk{}प्र\appbrk{}व\appbrk{}र्द्ध\appbrk{}नार्थं}}\Afootnote{\skt{वर्धनार्थं}\ A; \skt{लंप्रवर्धनार्थं}\ N H}}\skt{ }\edtext{\edtext{\skt{लघु\-प्र\-व\-र्द्ध\-न\-क\-मा\-मु\-ञ्चेत्}}{\lemma{\skt{लघु\appbrk{}प्र\appbrk{}व\appbrk{}र्द्ध\appbrk{}न\appbrk{}कम्}}\Afootnote{\skt{लघुवर्धनकं}\ A; \skt{ल प्रवर्धनका}\ N; \skt{लघु\appbrk{}प्र\appbrk{}व\appbrk{}र्द्ध\appbrk{}नका}\ H}}}{\lemma{\skt{आमुञ्चेत्}}\Afootnote{\skt{कुर्यात्}\ A; \skt{मो मुञ्चेत्}\ N; \skt{मा मुंचेत्}\ H\textsuperscript{ac}; \skt{मै मुंचेत्}\ H\textsuperscript{pc}}}\skt{~॥७॥}\pend

\pstart\versepar \ledsidenote{\passref{1.16.9}}\skt{एवं }\edtext{\skt{सम्वर्द्धितः}}{\lemma{\skt{सम्वर्द्धितः}}\Afootnote{N H; \skt{विवर्धितः}\ A}}\skt{ कर्ण्णश्छिद्यते तु द्विधा }\edtext{\skt{नृणाम्}}{\lemma{\skt{नृणाम्}}\Afootnote{A H; \skt{नृणा}\ N}}\skt{~। }\pend
\pstart\versepar \skt{}\edtext{\skt{दोषतो}}{\lemma{\skt{दोषतो}}\Afootnote{A; \skt{दोषटो}\ N H}}\skt{ वाभिघाताद्वा }\edtext{\skt{सन्धानान्तस्य}}{\lemma{\skt{सन्धानान्}}\Afootnote{N H; \skt{सन्धानं}\ A}}\skt{ मे शृणु~॥८॥}\pend

\pstart \ledsidenote{\passref{1.16.10}}\skt{तत्र समासेन }\edtext{\skt{पञ्च\-द\-श\-स\-न्धा\-ना\-कृ\-तयो}}{\lemma{\skt{पञ्च\appbrk{}द\appbrk{}श\appbrk{}स\appbrk{}न्धा\appbrk{}ना\appbrk{}कृ\appbrk{}तयो}}\Afootnote{H; \skt{पञ्च\appbrk{}द\appbrk{}श\appbrk{}क\appbrk{}र्ण\appbrk{}ब\appbrk{}न्धा\appbrk{}कृ\appbrk{}तयः}\ A; \skt{पञ्च\appbrk{}द\appbrk{}श\appbrk{}स\appbrk{}न्धा\appbrk{}कृ\appbrk{}तयो}\ N}}\skt{ }\edtext{\skt{भवन्ति}}{\lemma{\skt{भवन्ति}}\Afootnote{N H; \emph{om.}\ A}}\skt{~।}\footnoteA{Cakrapāṇi (1.16.9–13) and Ḍalhaṇa (1.16.10) pointed out that others read \emph{pañcadaśakarṇakṛtayaḥ} (instead of \emph{pañcadaśasandhānākṛtayaḥ}). Ḍalhaṇa (1.16.10) also mentioned that some read \emph{samunnatasamobhayapāliḥ} (instead of \emph{vṛttāyatasamobhayapālir}) and others do not read \emph{saṃkṣiptādayaḥ pañcāsādhyāḥ}.}\skt{ तद्यथा~। }\edtext{\skt{नेमीसन्धानकः}}{\lemma{\skt{नेमीसन्धानकः}}\Afootnote{N H; \skt{नेमिसन्धानक}\ A}}\skt{~। उत्पलभेद्यकः~। वल्लूरकः~। आसङ्गिमः~। गण्डकर्ण्णः~। आहार्यः~। निर्व्वेधिमः~। व्यायोजिमः~। }\edtext{\skt{कपाटसन्धिकः}}{\lemma{\skt{कपाटसन्धिकः}}\Afootnote{A H; \emph{om.}\ N}}\skt{~। }\edtext{\skt{अर्द्ध\-क\-पा\-ट\-स\-न्धिकः}}{\lemma{\skt{अर्द्ध\appbrk{}क\appbrk{}पा\appbrk{}ट\appbrk{}स\appbrk{}न्धिकः}}\Afootnote{A H; \skt{कपाटसन्धिकः}\ N}}\skt{~। सङ्क्षिप्तः~। हीनकर्ण्णः~। वल्लीकर्ण्णः~। }\edtext{\skt{यष्टीकर्ण्णः}}{\lemma{\skt{यष्टीकर्ण्णः}}\Afootnote{N H; \skt{यष्टिकर्णः}\ A}}\skt{~। }\edtext{\skt{काकौष्ठः}}{\lemma{\skt{काकौष्ठः}}\Afootnote{N; \skt{काकौष्ठक}\ A; \skt{काकौष्ठभः}\ H}}\skt{~। }\edtext{\skt{इति}}{\lemma{\skt{इति}}\Afootnote{A N; \skt{ति}\ H}}\skt{~। तेषु }\edtext{\skt{तत्र}}{\lemma{\skt{तत्र}}\Afootnote{N H; \emph{om.}\ A}}\skt{ }\edtext{\edtext{\skt{पृथु\-ला\-य\-त\-स\-मो\-भ\-य\-पा\-लि\-र्ने\-मी\-स\-न्धा\-नकः}}{\lemma{\skt{पृथु\appbrk{}ला\appbrk{}य\appbrk{}त\appbrk{}स\appbrk{}मो\appbrk{}भ\appbrk{}य\appbrk{}पा\appbrk{}लिर्}}\Afootnote{A; \skt{पृथुलातसमो}\ N; \skt{पृथु\appbrk{}ला\appbrk{}य\appbrk{}स\appbrk{}मो\appbrk{}भ\appbrk{}य\appbrk{}पा\appbrk{}लिर्}\ H}}}{\lemma{\skt{नेमीसन्धानकः}}\Afootnote{N H; \skt{नेमिसन्धानकः}\ A}}\skt{~। }\edtext{\skt{वृत्ता\-य\-त\-स\-मो\-भ\-य\-पा\-लि\-रु\-त्प\-ल\-भे\-द्यकः}}{\lemma{\skt{उत्पलभेद्यकः}}\Afootnote{A; \skt{उत्पलभेद्यः}\ N; \skt{उत्पलभेदकः}\ H}}\skt{~। }\edtext{\skt{ह्रस्व\-वृ\-त्त\-स\-मो\-भ\-य\-पा\-लि\-र्व\-ल्लू\-र\-क\-र्ण्णकः}}{\lemma{\skt{वल्लूरकर्ण्णकः}}\Afootnote{H; \skt{वल्लूरकः}\ A; \skt{वलूरकर्ण्णकः}\ N}}\skt{~। अभ्य\-न्त\-र\-दी\-र्घै\-क\-पा\-लि\-रा\-स\-ङ्गिमः~। }\edtext{\edtext{\skt{बाह्य\-दी\-र्घै\-क\-पा\-लि\-र्ग्ग\-ण्ड\-क\-र्ण्णकः}}{\lemma{\skt{बाह्य\appbrk{}दी\appbrk{}र्घै\appbrk{}क\appbrk{}पा\appbrk{}लिर्}}\Afootnote{A; \skt{बाह्यै\appbrk{}क\appbrk{}दी\appbrk{}र्घै\appbrk{}क\appbrk{}पा\appbrk{}लिर्}\ N H}}}{\lemma{\skt{ग्गण्डकर्ण्णकः}}\Afootnote{N H; \skt{गण्डकर्णः}\ A}}\skt{~। }\edtext{\edtext{\skt{अपा\-लि\-रु\-भ\-य\-तोऽ\-प्या\-हार्यः}}{\lemma{\skt{उभयतो}}\Afootnote{A N; \skt{उभयतोऽ}\ H}}}{\lemma{\skt{ऽप्य्}}\Afootnote{A N; \skt{प्य्}\ H}}\skt{~। }\edtext{\skt{पीठो\-प\-म\-पा\-लि\-र्नि\-र्व्वे\-धिमः}}{\lemma{\skt{निर्व्वेधिमः}}\Afootnote{N H; \skt{उभयतः क्षीण\appbrk{}पु\appbrk{}त्रि\appbrk{}का\appbrk{}श्रितो निर्वेधिमः}\ A}}\skt{~। }\edtext{\skt{अणु\-स्थू\-ल\-स\-म\-वि\-ष\-म\-पा\-लि\-र्व्या\-यो\-जिमः}}{\lemma{\skt{अणु\appbrk{}स्थू\appbrk{}ल\appbrk{}स\appbrk{}म\appbrk{}वि\appbrk{}ष\appbrk{}म\appbrk{}पा\appbrk{}लिर्}}\Afootnote{N; \skt{स्थूला\appbrk{}णु\appbrk{}स\appbrk{}म\appbrk{}वि\appbrk{}ष\appbrk{}म\appbrk{}पा\appbrk{}लिर्}\ A; \skt{अशु\appbrk{}स्थू\appbrk{}ल\appbrk{}स\appbrk{}म\appbrk{}वि\appbrk{}ष\appbrk{}म\appbrk{}पा\appbrk{}लिर्}\ H}}\skt{~। अभ्य\-न्त\-र\-दी\-र्घै\-क\-पा\-लि\-रि\-त\-रा\-ल्प\-पालिः }\edtext{\skt{कपाटसन्धिकः}}{\lemma{\skt{कपाटसन्धिकः}}\Afootnote{A N; \skt{कवाटासन्धिकः}\ H}}\skt{~। }\edtext{\skt{बाह्य\-दी\-र्घै\-क\-पा\-लि\-रि\-त\-रा\-ल्प\-पा\-लि\-श्चा\-र्द्ध\-क\-पा\-ट\-स\-न्धिकः}}{\lemma{\skt{चार्द्ध\appbrk{}क\appbrk{}पा\appbrk{}ट\appbrk{}स\appbrk{}न्धिकः}}\Afootnote{\skt{अर्ध\appbrk{}क\appbrk{}पा\appbrk{}ट\appbrk{}स\appbrk{}न्धिकः}\ A; \skt{चार्द्ध\appbrk{}क\appbrk{}वा\appbrk{}प\appbrk{}स\appbrk{}न्धिकः}\ N; \skt{वार्द्ध\appbrk{}क\appbrk{}वा\appbrk{}ट\appbrk{}स\appbrk{}न्धिकः}\ H}}\skt{~। }\edtext{\skt{तत्रैते}}{\lemma{\skt{तत्रैते}}\Afootnote{N H; \skt{तत्र द}\ A}}\skt{ }\edtext{\skt{दश\-क\-र्ण्ण\-स\-न्धि\-वि\-कल्पा}}{\lemma{\skt{दश\appbrk{}क\appbrk{}र्ण्ण\appbrk{}स\appbrk{}न्धि\appbrk{}वि\appbrk{}कल्पा}}\Afootnote{N H; \skt{शैते कर्ण\appbrk{}ब\appbrk{}न्ध\appbrk{}वि\appbrk{}कल्पाः}\ A}}\skt{ }\edtext{\skt{बन्ध्या}}{\lemma{\skt{बन्ध्या}}\Afootnote{N H; \skt{साध्याः}\ A}}\skt{ }\edtext{\skt{भवन्ति}}{\lemma{\skt{भवन्ति}}\Afootnote{N H; \emph{om.}\ A}}\skt{~। }\edtext{\skt{तेषा\-न्ना\-म\-भि\-रे\-वा\-कृ\-तयः}}{\lemma{\skt{नामभिर्}}\Afootnote{N H; \skt{स्वनामभिर्}\ A}}\skt{ प्रायेण व्याख्याताः~। संक्षिप्तादयः पञ्चासाध्याः~। तत्र }\edtext{\skt{शुष्क\-श\-ष्कु\-लि\-रि\-त\-रा\-ल्प\-पालिः}}{\lemma{\skt{शुष्कशष्कुलिर्}}\Afootnote{N H; \skt{शुष्क\appbrk{}श\appbrk{}ष्कु\appbrk{}लि\appbrk{}रु\appbrk{}त्स\appbrk{}न्न\appbrk{}पा\appbrk{}लिर्}\ A}}\skt{ संक्षिप्तः~। अन\-धि\-ष्ठा\-न\-पालिः }\edtext{\skt{पर्यन्तयोश्च}}{\lemma{\skt{पर्यन्तयोश्च}}\Afootnote{H; \skt{पर्यन्तयोः}\ A; \emph{om.}\ N}}\skt{ क्षीणमांसो हीनकर्ण्णः~। }\edtext{\skt{तनु\-वि\-ष\-म\-पा\-लि\-र्व\-ल्ली\-कर्ण्णः}}{\lemma{\skt{तनु\appbrk{}वि\appbrk{}ष\appbrk{}म\appbrk{}पा\appbrk{}लिर्}}\Afootnote{N H; \skt{तनु\appbrk{}वि\appbrk{}ष\appbrk{}मा\appbrk{}ल्प\appbrk{}पा\appbrk{}लिर्}\ A}}\skt{~। }\edtext{\skt{ग्रन्थितमांसः}}{\lemma{\skt{ग्रन्थितमांसः}}\Afootnote{N H; \skt{ग्रथितमांस}\ A}}\skt{ }\edtext{\edtext{\skt{स्तब्ध\-सि\-रा\-त\-त\-सू\-क्ष्म\-पा\-लि\-र्य\-ष्टी\-कर्ण्णः}}{\lemma{\skt{स्तब्ध\appbrk{}सि\appbrk{}रा\appbrk{}त\appbrk{}त\appbrk{}सू\appbrk{}क्ष्म\appbrk{}पा\appbrk{}लिर्}}\Afootnote{N H; \skt{स्तब्ध\appbrk{}सि\appbrk{}रा\appbrk{}सं\appbrk{}त\appbrk{}त\appbrk{}सू\appbrk{}क्ष्म\appbrk{}पा\appbrk{}लिर्}\ A}}}{\lemma{\skt{यष्टीकर्ण्णः}}\Afootnote{N H; \skt{यष्टिकर्णः}\ A}}\skt{~। }\edtext{\skt{निर्म्मां\-स\-सं\-क्षि\-प्ता\-ग्रा\-ल्प\-शो\-णि\-त\-पालिः}}{\lemma{\skt{निर्म्मां\appbrk{}स\appbrk{}सं\appbrk{}क्षि\appbrk{}प्ता\appbrk{}ग्रा\appbrk{}ल्प\appbrk{}शो\appbrk{}णि\appbrk{}त\appbrk{}पालिः}}\Afootnote{A H; \skt{निमा\appbrk{}स\appbrk{}सं\appbrk{}क्षि\appbrk{}प्ता\appbrk{}ग्रा\appbrk{}ल्प\appbrk{}शो\appbrk{}णि\appbrk{}त\appbrk{}पालिः}\ N}}\skt{ }\edtext{\skt{काकौष्ठ}}{\lemma{\skt{काकौष्ठ}}\Afootnote{N H; \skt{काकौष्ठक}\ A}}\skt{ इति~। }\edtext{\skt{बद्धेष्वपि}}{\lemma{\skt{अपि}}\Afootnote{N H; \skt{अपि तु}\ A}}\skt{ }\edtext{\skt{दाह\-पा\-क\-स्रा\-व\-शो\-फ\-युक्ता न}}{\lemma{\skt{दाह\appbrk{}पा\appbrk{}क\appbrk{}स्रा\appbrk{}व\appbrk{}शो\appbrk{}फ\appbrk{}युक्ता न}}\Afootnote{N; \skt{शोफ\appbrk{}दा\appbrk{}ह\appbrk{}रा\appbrk{}ग\appbrk{}पा\appbrk{}क\appbrk{}पि\appbrk{}ड\appbrk{}का\appbrk{}स्रा\appbrk{}व\appbrk{}युक्ता न}\ A; \skt{दाह\appbrk{}पा\appbrk{}क\appbrk{}श्रा\appbrk{}व\appbrk{}शो\appbrk{}फ\appbrk{}युक्ता न}\ H}}\skt{ सिद्धिमुपयान्ति~॥९॥}\pend

\pstart \ledsidenote{\passref{1.16.15}}\edtext{\edtext{\skt{अतोऽन्यतमस्य}}{\lemma{\skt{अतो}}\Afootnote{A H; \skt{त}\ N}}}{\lemma{\skt{ऽन्यतमस्य}}\Afootnote{\skt{ऽन्यतमं}\ A; \skt{तो न्यतमस्य}\ N; \skt{न्यतमस्य}\ H}}\skt{ बन्धञ्चिकीर्षुः }\edtext{\skt{अग्रो\-प\-ह\-र\-णी\-यो\-क्तो\-प\-स\-म्भृ\-त\-स\-म्भारः}}{\lemma{\skt{अग्रो\appbrk{}प\appbrk{}ह\appbrk{}र\appbrk{}णी\appbrk{}यो\appbrk{}क्तो\appbrk{}प\appbrk{}स\appbrk{}म्भृ\appbrk{}त\appbrk{}स\appbrk{}म्भारः}}\Afootnote{H; \skt{अग्रो\appbrk{}प\appbrk{}ह\appbrk{}र\appbrk{}णी\appbrk{}यो\appbrk{}क्तो\appbrk{}प\appbrk{}सं\appbrk{}भृ\appbrk{}त\appbrk{}सं\appbrk{}भारं}\ A; \skt{अग्रो\appbrk{}प\appbrk{}सं\appbrk{}ह\appbrk{}र\appbrk{}णी\appbrk{}यो\appbrk{}क्तो\appbrk{}प\appbrk{}स\appbrk{}म्भृ\appbrk{}त\appbrk{}स\appbrk{}म्भारः}\ N}}\skt{ }\edtext{\skt{विशे\-ष\-त\-श्चा\-त्रो\-प\-ह\-रे\-त्सु\-रा\-म\-ण्ड\-क्षी\-र\-मु\-दकं}}{\lemma{\skt{चात्रो\appbrk{}प\appbrk{}ह\appbrk{}रे\appbrk{}त्सु\appbrk{}रा\appbrk{}म\appbrk{}ण्ड\appbrk{}क्षी\appbrk{}रम्}}\Afootnote{\skt{चात्रो\appbrk{}प\appbrk{}ह\appbrk{}रे\appbrk{}त्सु\appbrk{}रा\appbrk{}मण्डं क्षीरम्}\ A; \skt{चाग्रो\appbrk{}प\appbrk{}ह\appbrk{}र\appbrk{}णी\appbrk{}या\appbrk{}त्सु\appbrk{}रा\appbrk{}म\appbrk{}ण्ड\appbrk{}क्षी\appbrk{}रम्}\ N H}}\footnoteA{The MSS reading \emph{viśeṣataś cāgropaharaṇīyāt} has been emended to \emph{viśeṣataś cātropaharet} to make sense of the list of ingredients, which is in the accusative case. Also, the repetition of \emph{agropaharaṇīyāt} in the MSS suggests that its second occurrence, which does not make good sense here, is a dittographic error.}\skt{ }\edtext{\skt{धान्या\-म्ल\-क\-पा\-ल\-चू\-र्ण्ण\-ञ्चेति}}{\lemma{\skt{धान्या\appbrk{}म्ल\appbrk{}क\appbrk{}पा\appbrk{}ल\appbrk{}चू\appbrk{}र्ण्णञ्}}\Afootnote{N H; \skt{धान्याम्लं कपालचूर्णं}\ A}}\skt{~। ततोऽङ्गनां }\edtext{\skt{पुरुषम्वा}}{\lemma{\skt{पुरुषम्वा}}\Afootnote{A H; \skt{पुरुषञ्वा}\ N}}\skt{ ग्रथितकेशान्तं लघु\-भु\-क्त\-व\-न्त\-माप्तैः सुपरिगृहीतं }\edtext{\skt{कृत्वा च}}{\lemma{\skt{कृत्वा च}}\Afootnote{\skt{च कृत्वा}\ A; \emph{om.}\ N H}}\skt{ }\edtext{\skt{बन्धानुपधार्य}}{\lemma{\skt{बन्धानुपधार्य}}\Afootnote{N; \skt{बन्धमुपधार्य}\ A; \skt{बन्धानुपपाद्य}\ H}}\skt{ }\edtext{\skt{छेद्य\-भे\-द्य\-ले\-ख्य\-व्य\-ध\-नै\-रु\-प\-पाद्य}}{\lemma{\skt{छेद्य\appbrk{}भे\appbrk{}द्य\appbrk{}ले\appbrk{}ख्य\appbrk{}व्य\appbrk{}ध\appbrk{}नैर्}}\Afootnote{N H; \skt{छेद्य\appbrk{}भे\appbrk{}द्य\appbrk{}ले\appbrk{}ख्य\appbrk{}व्य\appbrk{}ध\appbrk{}नै\appbrk{}रु\appbrk{}प\appbrk{}प\appbrk{}न्नैर्}\ A}}\skt{ }\edtext{\edtext{\edtext{\edtext{\skt{कर्ण्ण\-शो\-णि\-त\-म\-वे\-क्ष्यै\-त\-द्दु\-ष्ट\-म\-दु\-ष्ट\-म्वेति}}{\lemma{\skt{कर्ण्णशोणितम्}}\Afootnote{A H; \skt{कर्ण्णशोणितत}\ N}}}{\lemma{\skt{अवे\appbrk{}क्ष्यै\appbrk{}त\appbrk{}द्दु\appbrk{}ष्टम्}}\Afootnote{H; \skt{अवेक्ष्य दुष्टम्}\ A; \skt{अवे\appbrk{}क्ष्ये\appbrk{}त\appbrk{}द्दु\appbrk{}ष्टम्}\ N}}}{\lemma{\skt{अदुष्टम्}}\Afootnote{A H; \skt{अदुष्टश्}\ N}}}{\lemma{\skt{वेति}}\Afootnote{A; \skt{चेति}\ N H}}\skt{ }\edtext{\skt{ततो}}{\lemma{\skt{ततो}}\Afootnote{N H; \skt{तत्र}\ A}}\skt{ वातदुष्टे }\edtext{\skt{धान्याम्लोदकाभ्यां}}{\lemma{\skt{धान्याम्लोदकाभ्यां}}\Afootnote{H; \skt{धान्या\appbrk{}म्लो\appbrk{}ष्णो\appbrk{}द\appbrk{}काभ्यां}\ A; \skt{धान्याव्लोदकाभ्यां}\ N}}\skt{ पित्तदुष्टे }\edtext{\skt{शीतो\-द\-क\-प\-योभ्यां}}{\lemma{\skt{शीतो\appbrk{}द\appbrk{}क\appbrk{}प\appbrk{}योभ्यां}}\Afootnote{A H; \skt{शीतो\appbrk{}द\appbrk{}को\appbrk{}प\appbrk{}योभ्यां}\ N}}\skt{ श्लेष्मदुष्टे }\edtext{\skt{सुरा\-म\-ण्डो\-द\-काभ्यां}}{\lemma{\skt{सुरा\appbrk{}म\appbrk{}ण्डो\appbrk{}द\appbrk{}काभ्यां}}\Afootnote{N H; \skt{सुरा\appbrk{}म\appbrk{}ण्डो\appbrk{}ष्णो\appbrk{}द\appbrk{}काभ्यां}\ A}}\skt{ प्रक्षाल्य }\edtext{\edtext{\skt{कर्ण्ण\-म्पु\-न\-र\-व\-लि\-खेत्}}{\lemma{\skt{कर्ण्णम्}}\Afootnote{N H; \skt{कर्णौ}\ A}}\skt{~। अनु\-न्न\-त\-म\-ही\-न\-म\-वि\-ष\-मञ्च}}{\lemma{\skt{अवलिखेत्~। अनुन्नतम्}}\Afootnote{H; \skt{अव\appbrk{}लि\appbrk{}ख्या\appbrk{}नु\appbrk{}न्न\appbrk{}तम्}\ A; \skt{अवलिखेत~। अनुन्नतम्}\ N}}\skt{ }\edtext{\skt{कर्ण्ण\-स\-न्धि\-न्नि\-वेश्य}}{\lemma{\skt{निवेश्य}}\Afootnote{N H; \skt{सन्निवेश्य}\ A}}\skt{ स्थितरक्तं }\edtext{\skt{सन्दर्श्य}}{\lemma{\skt{सन्दर्श्य}}\Afootnote{N H; \skt{सन्दध्यात्~। ततो}\ A}}\skt{ मधु\-घृ\-ते\-ना\-भ्यज्य }\edtext{\edtext{\skt{पिचु\-प्लो\-त\-यो\-र\-न्य\-त\-रे\-णा\-व\-गुण्ठ्य}}{\lemma{\skt{पिचुप्लोतयोर्}}\Afootnote{N H; \skt{पिचुप्रोतयोर्}\ A}}}{\lemma{\skt{अन्य\appbrk{}त\appbrk{}रे\appbrk{}णा\appbrk{}व\appbrk{}गुण्ठ्य}}\Afootnote{A N; \skt{अन्य\appbrk{}त\appbrk{}रे\appbrk{}णा\appbrk{}व\appbrk{}गुण्ठ्यो}\ H}}\skt{ }\edtext{\edtext{\skt{नाति\-गा\-ढ\-न्ना\-ति\-शि\-थिलं}}{\lemma{\skt{नातिगाढन्}}\Afootnote{N H; \skt{सूत्रे\appbrk{}णा\appbrk{}न\appbrk{}व\appbrk{}गा\appbrk{}ढम}\ A}}}{\lemma{\skt{नातिशिथिलं}}\Afootnote{N H; \skt{नतिशिथिलं}\ A}}\skt{ }\edtext{\skt{सूत्रेणावबध्य}}{\lemma{\skt{सूत्रेणावबध्य}}\Afootnote{H; \skt{च बद्ध्वा}\ A; \skt{सूत्रेणावबद्ध}\ N}}\skt{ }\edtext{\skt{कपा\-ल\-चू\-र्ण्णे\-ना\-व\-की\-र्या\-चा\-रि\-क\-मु\-प\-दि\-शेत्~।}\footnoteA{\emph{Aṣṭāṅgahṛdayasaṃhitā} 1.18.52--53: \emph{atha grathitvā keśāntaṃ kṛtvā chedanalekhanam | niveśya sandhiṃ suṣamaṃ na nimnaṃ na samunnatam || 52 || abhyajya madhusarpirbhyāṃ picuplotāvaguṇṭhitam | sūtreṇāgāḍhaśithilaṃ baddhvā cūrṇair avākiret || 53 ||}}\skt{ द्विव्रणीयोक्तेन}}{\lemma{\skt{उपदिशेत्~। द्विव्रणीयोक्तेन}}\Afootnote{A H; \skt{उप\appbrk{}दि\appbrk{}शे\appbrk{}त्द्वि\appbrk{}व\appbrk{}णी\appbrk{}यो\appbrk{}क्तेन}\ N}}\skt{ }\edtext{\skt{चान्ने\-नो\-प\-च\-रेत्}}{\lemma{\skt{चान्ने\appbrk{}नो\appbrk{}प\appbrk{}च\appbrk{}रेत्}}\Afootnote{H; \skt{च विधा\appbrk{}ने\appbrk{}नो\appbrk{}प\appbrk{}च\appbrk{}रेत्}\ A; \skt{उपपोचरेत्}\ N}}\skt{~॥१०॥}\footnoteA{The additional verses in A (from \emph{bhavanti cātra to śāstravit}) were probably also absent in the version of the Suśrutasaṃhitā commented on by Cakrapāṇi, who cited them in his commentary as being "read by some" in regard to the joins (\emph{sandhāna}) they describe.}\pend

\pstart\versepar \ledsidenote{\passref{1.16.16}}\edtext{\skt{विघ\-ट्ट\-न\-न्दि\-वा\-स्वप्नं}}{\lemma{\skt{विघट्टनन्}}\Afootnote{N H; \skt{भवति चात्र~। विघट्टनं}\ A}}\skt{ व्याया\-म\-म\-ति\-भो\-ज\-नम्~। }\pend
\pstart\versepar \skt{}\edtext{\skt{व्यवा\-य\-म\-ग्नि\-स\-न्ता\-प\-म्वा\-क्श्र\-मञ्च}}{\lemma{\skt{अग्निसन्तापम्}}\Afootnote{A H; \skt{आग्निसन्ताप}\ N}}\skt{ }\edtext{\skt{विवर्ज्जयेत्}}{\lemma{\skt{विवर्ज्जयेत्}}\Afootnote{A H; \skt{वर्जयेत्}\ N}}\skt{~॥११॥}\pend

\pstart \ledsidenote{\passref{1.16.17}}\edtext{\edtext{\skt{नाति\-शु\-द्ध\-र\-क्त\-म\-ति\-प्र\-वृ\-त्त\-रक्तं}}{\lemma{\skt{नाति\appbrk{}शु\appbrk{}द्ध\appbrk{}र\appbrk{}क्तम्}}\Afootnote{N H; \skt{न चाशुद्धरक्तम्}\ A}}}{\lemma{\skt{अति\appbrk{}प्र\appbrk{}वृ\appbrk{}त्त\appbrk{}रक्तं}}\Afootnote{A H; \skt{अतिप्रवृत्तं रक्तं}\ N}}\skt{ क्षीणरक्तं वा सन्दध्यात्~। स हि वातदुष्टे }\edtext{\edtext{\skt{रक्तबद्धोऽरूढो}}{\lemma{\skt{रक्तबद्धो}}\Afootnote{N H; \skt{रक्ते}\ A}}}{\lemma{\skt{ऽरूढो}}\Afootnote{\skt{रूढो}\ A H; \skt{रुढो}\ N}}\skt{ }\edtext{\skt{परि\-पु\-ट\-न\-वा\-न्भ\-वति}}{\lemma{\skt{परि\appbrk{}पु\appbrk{}ट\appbrk{}न\appbrk{}वा\appbrk{}न्भ\appbrk{}वति}}\Afootnote{\skt{ऽपि परिपुटनवा}\ A; \skt{परि\appbrk{}पु\appbrk{}ट\appbrk{}वा\appbrk{}म्भ\appbrk{}वति}\ N; \skt{परि\appbrk{}पु\appbrk{}ट\appbrk{}न\appbrk{}वा\appbrk{}म्भ\appbrk{}वति}\ H}}\skt{~। }\edtext{\skt{पित्तदुष्टे}}{\lemma{\skt{पित्तदुष्टे}}\Afootnote{H; \skt{न्पित्तदुष्टे}\ A; \skt{पित्तदुष्टै}\ N}}\skt{ }\edtext{\skt{गाढ\-पा\-क\-रा\-ग\-वान्~। श्लेष्मदुष्टे}}{\lemma{\skt{गाढ\appbrk{}पा\appbrk{}क\appbrk{}रा\appbrk{}ग\appbrk{}वान्~। श्लेष्मदुष्टे}}\Afootnote{H; \skt{दाह\appbrk{}पा\appbrk{}क\appbrk{}रा\appbrk{}ग\appbrk{}वे\appbrk{}द\appbrk{}ना\appbrk{}वा\appbrk{}न्श्ले\appbrk{}ष्म\appbrk{}दुष्टे}\ A; \skt{गाढ\appbrk{}पा\appbrk{}क\appbrk{}रा\appbrk{}ग\appbrk{}वान्~। श्लेषदुष्टे}\ N}}\skt{ }\edtext{\skt{स्तब्धकर्ण्णः}}{\lemma{\skt{स्तब्धकर्ण्णः}}\Afootnote{H; \skt{स्तब्धः}\ A; \skt{स्तब्धवर्ण्णः}\ N}}\skt{ }\edtext{\skt{कण्डू\-मा\-न\-ति\-प्र\-वृ\-त्त\-स्रावः}}{\lemma{\skt{कण्डू\appbrk{}मा\appbrk{}न\appbrk{}ति\appbrk{}प्र\appbrk{}वृ\appbrk{}त्त\appbrk{}स्रावः}}\Afootnote{N; \skt{कण्डू\appbrk{}मा\appbrk{}न\appbrk{}ति\appbrk{}प्र\appbrk{}वृ\appbrk{}त्त\appbrk{}रक्ते श्याव}\ A; \skt{कण्डूमान्~। अति\appbrk{}प्र\appbrk{}वृ\appbrk{}त्त\appbrk{}श्रावः}\ H}}\skt{ }\edtext{\skt{शोफ\-वा\-न्क्षी\-णा\-ल्प\-मांसो न}}{\lemma{\skt{शोफ\appbrk{}वा\appbrk{}न्क्षी\appbrk{}णा\appbrk{}ल्प\appbrk{}मांसो न}}\Afootnote{H; \skt{शोफ\appbrk{}वा\appbrk{}न्क्षी\appbrk{}णोऽ\appbrk{}ल्प\appbrk{}मांसो न}\ A; \skt{शोफवान्~। क्षीणो ल्पमांसो न}\ N}}\skt{ वृद्धिमुपैति~॥१२॥}\pend

\pstart \ledsidenote{\passref{1.16.18}}\edtext{\skt{स}}{\lemma{\skt{स}}\Afootnote{N H; \skt{आमतैलेन त्रिरात्रं परि\appbrk{}षे\appbrk{}च\appbrk{}ये\appbrk{}त्त्रि\appbrk{}रा\appbrk{}त्राच्च पिचुं परिवर्तयेत्~। स}\ A}}\skt{ यदा }\edtext{\skt{रूढो}}{\lemma{\skt{रूढो}}\Afootnote{H; \skt{सुरूढो}\ A; \skt{रुढो}\ N}}\skt{ निरुपद्रवः }\edtext{\skt{कर्ण्णो}}{\lemma{\skt{कर्ण्णो}}\Afootnote{N H; \skt{सवर्णो}\ A}}\skt{ भवति तदैनं शनैः }\edtext{\skt{शनै\-र\-भि\-व\-र्द्ध\-येत्~। अन्यथा}}{\lemma{\skt{अभिवर्द्धयेत्~। अन्यथा}}\Afootnote{N H; \skt{अभिवर्धयेत्~। अतोऽन्यथा}\ A}}\skt{ }\edtext{\skt{संर\-म्भ\-दा\-ह\-पा\-क\-वे\-द\-ना\-वा\-न्भ\-वति}}{\lemma{\skt{संर\appbrk{}म्भ\appbrk{}दा\appbrk{}ह\appbrk{}पा\appbrk{}क\appbrk{}वे\appbrk{}द\appbrk{}ना\appbrk{}वा\appbrk{}न्भ\appbrk{}वति}}\Afootnote{\skt{संर\appbrk{}म्भ\appbrk{}दा\appbrk{}ह\appbrk{}पा\appbrk{}क\appbrk{}रा\appbrk{}ग\appbrk{}वे\appbrk{}द\appbrk{}नावा}\ A; \skt{संर\appbrk{}म्भ\appbrk{}दा\appbrk{}ह\appbrk{}पा\appbrk{}क\appbrk{}रा\appbrk{}ग\appbrk{}वे\appbrk{}द\appbrk{}ना\appbrk{}वा\appbrk{}न्भ\appbrk{}वति}\ N; \skt{संर\appbrk{}म्भ\appbrk{}दा\appbrk{}ह\appbrk{}पा\appbrk{}क\appbrk{}वे\appbrk{}द\appbrk{}ना\appbrk{}वा\appbrk{}म्भ\appbrk{}वति}\ H}}\skt{~। }\edtext{\edtext{\skt{पुनरपि}}{\lemma{\skt{पुनर्}}\Afootnote{N H; \skt{न्पुनश्}\ A}}}{\lemma{\skt{अपि}}\Afootnote{N H; \emph{om.}\ A}}\skt{ }\edtext{\skt{छिद्येत}}{\lemma{\skt{छिद्येत}}\Afootnote{N H; \skt{छिद्यते वा}\ A}}\skt{~॥१३॥}\pend

\pstart \ledsidenote{\passref{1.16.19}}\edtext{\skt{अथा\-प्र\-दु\-ष्ट\-स्या\-भि\-व\-र्द्ध\-ना\-र्थ\-म\-भ्यङ्गः}}{\lemma{\skt{अथा\appbrk{}प्र\appbrk{}दु\appbrk{}ष्ट\appbrk{}स्या\appbrk{}भि\appbrk{}व\appbrk{}र्द्ध\appbrk{}ना\appbrk{}र्थम्}}\Afootnote{\skt{अथा\appbrk{}स्या\appbrk{}प्र\appbrk{}दु\appbrk{}ष्ट\appbrk{}स्या\appbrk{}भि\appbrk{}व\appbrk{}र्ध\appbrk{}ना\appbrk{}र्थम्}\ A; \skt{अथा\appbrk{}प्र\appbrk{}दु\appbrk{}ष्ट\appbrk{}स्या\appbrk{}वि\appbrk{}व\appbrk{}र्ध\appbrk{}ना\appbrk{}र्थम्}\ N; \skt{अथास्याः प्रदु\appbrk{}ष्ट\appbrk{}स्या\appbrk{}भि\appbrk{}व\appbrk{}र्द्ध\appbrk{}ना\appbrk{}र्थम्}\ H}}\skt{~। }\edtext{\skt{गोधा\-प्र\-तु\-द\-वि\-ष्कि\-रा\-नू\-पौ\-द\-क\-व\-सा\-म\-ज्जा\-प\-य\-स्तैलं}}{\lemma{\skt{गोधा\appbrk{}प्र\appbrk{}तु\appbrk{}द\appbrk{}वि\appbrk{}ष्कि\appbrk{}रा\appbrk{}नू\appbrk{}पौ\appbrk{}द\appbrk{}क\appbrk{}व\appbrk{}सा\appbrk{}म\appbrk{}ज्जा\appbrk{}प\appbrk{}य\appbrk{}स्तैलं}}\Afootnote{N H; \skt{तद्यथा गोधा\appbrk{}प्र\appbrk{}तु\appbrk{}द\appbrk{}वि\appbrk{}ष्कि\appbrk{}रा\appbrk{}नू\appbrk{}पौ\appbrk{}द\appbrk{}क\appbrk{}व\appbrk{}सा\appbrk{}म\appbrk{}ज्जानौ पयः सर्पिस्तैलं}\ A}}\skt{ }\edtext{\skt{गौर\-स\-र्ष\-प\-जञ्च}}{\lemma{\skt{गौर\appbrk{}स\appbrk{}र्ष\appbrk{}प\appbrk{}जञ्च}}\Afootnote{A H; \skt{सर्षपजञ्च}\ N}}\skt{ यथालाभं }\edtext{\edtext{\skt{संभृ\-त्या\-र्का\-ल\-र्क\-ब\-ला\-ति\-ब\-ला\-न\-न्ता\-वि\-दा\-री\-म\-धु\-क\-ज\-ल\-शू\-क\-प्र\-ति\-वा\-प\-न्तै\-ल\-म्पा\-च\-यित्वा}}{\lemma{\skt{संभृ\appbrk{}त्या\appbrk{}र्का\appbrk{}ल\appbrk{}र्क\appbrk{}ब\appbrk{}ला\appbrk{}ति\appbrk{}ब\appbrk{}ला\appbrk{}न\appbrk{}न्ता\appbrk{}वि\appbrk{}दा\appbrk{}री\appbrk{}म\appbrk{}धु\appbrk{}क\appbrk{}ज\appbrk{}ल\appbrk{}शू\appbrk{}क\appbrk{}प्र\appbrk{}ति\appbrk{}वा\appbrk{}पन्}}\Afootnote{H; \skt{संभृ\appbrk{}त्या\appbrk{}र्का\appbrk{}ल\appbrk{}र्क\appbrk{}ब\appbrk{}ला\appbrk{}ति\appbrk{}ब\appbrk{}ला\appbrk{}न\appbrk{}न्ता\appbrk{}पा\appbrk{}मा\appbrk{}र्गा\appbrk{}श्व\appbrk{}ग\appbrk{}न्धा\appbrk{}वि\appbrk{}दा\appbrk{}रि\appbrk{}ग\appbrk{}न्धा\appbrk{}क्षी\appbrk{}र\appbrk{}शु\appbrk{}क्ला\appbrk{}ज\appbrk{}ल\appbrk{}शू\appbrk{}क\appbrk{}म\appbrk{}धु\appbrk{}र\appbrk{}व\appbrk{}र्ग\appbrk{}प\appbrk{}य\appbrk{}स्या\appbrk{}प्र\appbrk{}ति\appbrk{}वापं}\ A; \skt{सम्भृ\appbrk{}त्या\appbrk{}र्का\appbrk{}ल\appbrk{}क\appbrk{}व\appbrk{}ला\appbrk{}ति\appbrk{}ब\appbrk{}ला\appbrk{}न\appbrk{}न्ता\appbrk{}वि\appbrk{}दा\appbrk{}री\appbrk{}म\appbrk{}धु\appbrk{}क\appbrk{}ज\appbrk{}ल\appbrk{}शू\appbrk{}क\appbrk{}प्र\appbrk{}ति\appbrk{}वापं}\ N}}}{\lemma{\skt{तैलम्}}\Afootnote{N H; \skt{तैलं वा}\ A}}\skt{ }\edtext{\skt{स्वनु\-गु\-प्त\-न्नि\-द\-ध्यात्}}{\lemma{\skt{निदध्यात्}}\Afootnote{A H; \skt{निदद्यात्}\ N}}\skt{~॥१४॥}\footnoteA{Ḍalhaṇa (1.16.18) noted that some read \emph{rājasarṣapajaṃ} in the place of \emph{gaurasarṣapajaṃ}. This reading appears to have been accepted by Cakrapāṇi (1.16.18–20), who glossed \emph{rājasarṣapaja} as \emph{śvetasarṣapa}. Cakrapāṇi also said that some read sarpis in the place of \emph{payas}. In the compound beginning with \emph{arka}, Ḍalhaṇa noted that some read \emph{arkapuṣpī}.}\pend

\pstart\versepar \ledsidenote{\passref{1.16.20}}\edtext{\skt{स्वेदितो}}{\lemma{\skt{स्वेदितो}}\Afootnote{H; \skt{स्वेदितोन्}\ A; \skt{॥ भवति चात्र~॥ स्वदितो}\ N}}\skt{ }\edtext{\skt{मर्दि\-त\-ङ्क\-र्ण्ण\-म\-नेन}}{\lemma{\skt{अनेन}}\Afootnote{H; \skt{स्नेहेनै}\ A; \skt{अने}\ N}}\skt{ }\edtext{\skt{म्रक्षयेद्बुधः}}{\lemma{\skt{म्रक्षयेद्बुधः}}\Afootnote{N H; \skt{तेन योजये}\ A}}\skt{~। }\pend
\pstart\versepar \skt{}\edtext{\skt{ततोऽनुपद्रवः}}{\lemma{\skt{ऽनुपद्रवः}}\Afootnote{A; \skt{नुपद्रवम्}\ N; \skt{नुपद्रवः}\ H}}\skt{ सम्यग्बलवांश्च विवर्द्धते~॥१५॥}\footnoteA{N has a \emph{kākapāda} after \emph{ane}, but the missing letter (one would expect '\emph{na}') has not been supplied in a margin or elsewhere.}\pend

\pstart \ledsidenote{\passref{1.16.22}}\edtext{\skt{ये तु}}{\lemma{\skt{ये तु}}\Afootnote{N H; \skt{तैलं विपक्वं सक्षी\appbrk{}र\appbrk{}म\appbrk{}भ्य\appbrk{}ङ्गा\appbrk{}त्पा\appbrk{}लि\appbrk{}व\appbrk{}र्ध\appbrk{}नम्~। ये तु}\ A}}\skt{ कर्ण्णा न वर्द्धन्ते }\edtext{\skt{स्नेह\-स्वे\-दो\-प\-पा\-दिताः}}{\lemma{\skt{स्नेह\appbrk{}स्वे\appbrk{}दो\appbrk{}प\appbrk{}पा\appbrk{}दिताः}}\Afootnote{N H; \skt{स्वेद\appbrk{}स्ने\appbrk{}हो\appbrk{}प\appbrk{}पा\appbrk{}दिताः}\ A}}\skt{~।}\pend

\pstart \ledsidenote{\passref{1.16.23}}\edtext{\skt{तेषामपाङ्गे}}{\lemma{\skt{अपाङ्गे}}\Afootnote{N H; \skt{अपाङ्ग}\ A}}\skt{ }\edtext{\skt{त्वबहिः}}{\lemma{\skt{त्वबहिः}}\Afootnote{H; \skt{देशे तु}\ A; \skt{त्ववर्हि}\ N}}\skt{ }\edtext{\skt{कुर्या\-त्प्र\-छा\-न\-मेव}}{\lemma{\skt{कुर्यात्प्रछानम्}}\Afootnote{A; \skt{कुयात्प्रछानम्}\ N; \skt{कुर्यात्प्रछन्नम्}\ H}}\skt{ }\edtext{\skt{च}}{\lemma{\skt{च}}\Afootnote{N H; \skt{तु}\ A}}\skt{~॥१६॥}\footnoteA{Ḍalhaṇa (1.16.23) noted that some read \emph{teṣām apāṅgacchedyaṃ hi kāryam ābhyantaraṃ bhavet}.}\pend

\pstart\versepar \ledsidenote{\passref{1.16.26.0}}\skt{अमिताः }\edtext{\skt{कर्ण्णबन्धास्तु}}{\lemma{\skt{कर्ण्णबन्धास्तु}}\Afootnote{A N; \skt{कर्ण्णबन्धोस्तु}\ H}}\skt{ विज्ञेयाः कुशलैरिह~। }\pend
\pstart\versepar \skt{यो यथा }\edtext{\skt{सुनिविष्टः}}{\lemma{\skt{सुनिविष्टः}}\Afootnote{N H; \skt{सुविशिष्टः}\ A}}\skt{ }\edtext{\skt{स्यात्तत्तथा}}{\lemma{\skt{स्यात्तत्तथा}}\Afootnote{N H; \skt{स्यात्तं तथा}\ A}}\skt{ }\edtext{\skt{योजयेद्भिषक्}}{\lemma{\skt{योजयेद्भिषक्}}\Afootnote{H; \skt{विनियोजयेत्}\ A; \skt{योजये भिषक्}\ N}}\skt{~॥१७॥}\footnoteA{Ḍalhaṇa (1.16.26) stated that some read \emph{suniviṣṭaḥ} (the reading of the Nepalese version) instead of \emph{suviśiṣṭaḥ}.}\pend

\pstart\versepar \ledsidenote{\passref{1.16.25}}\skt{जातरोमा }\edtext{\skt{सुवर्त्मा च}}{\lemma{\skt{सुवर्त्मा च}}\Afootnote{A; \skt{सुपर्मा च}\ N; \skt{सुवर्म्मा च}\ H}}\skt{ }\edtext{\skt{श्लिष्टसन्धिः}}{\lemma{\skt{श्लिष्टसन्धिः}}\Afootnote{A H; \skt{श्लिष्टसन्धिम्}\ N}}\skt{ समः स्थिरः~। }\pend
\pstart\versepar \skt{}\edtext{\skt{सुरूढोऽवेदनो यस्तु तं}}{\lemma{\skt{ऽवेदनो यस्तु तं}}\Afootnote{\skt{ऽवेदनो यश्च तं}\ A; \skt{वेदनो यस्तु त}\ N H}}\skt{ }\edtext{\skt{कर्ण्णं}}{\lemma{\skt{कर्ण्णं}}\Afootnote{A; \skt{त्कर्णं}\ N H}}\skt{ वर्द्धयेच्छनैः~॥१८॥}\pend

\pstart\versepar \ledsidenote{\passref{1.16.27}}\edtext{\edtext{\skt{विश्ले\-षि\-ता\-या\-मथ}}{\lemma{\skt{विश्लेषितायाम्}}\Afootnote{N H; \skt{विश्लेषितायास्}\ A}}}{\lemma{\skt{अथ}}\Afootnote{N H; \skt{त्वथ}\ A}}\skt{ }\edtext{\skt{नासिकायां}}{\lemma{\skt{नासिकायां}}\Afootnote{H; \skt{नासिकाया}\ A N}}\skt{ }\pend
\pstart\versepar \skt{वक्ष्यामि सन्धानविधिं यथावत्~। }\pend
\pstart\versepar \skt{नासाप्रमाणं }\edtext{\skt{पृथिवीरुहाणां}}{\lemma{\skt{पृथिवीरुहाणां}}\Afootnote{A H; \skt{पृथि\appbrk{}वी\appbrk{}रु\appbrk{}हा\appbrk{}णम्}\ N}}\skt{ }\pend
\pstart\versepar \skt{}\edtext{\skt{पत्रं}}{\lemma{\skt{पत्रं}}\Afootnote{A H; \skt{पत्र}\ N}}\skt{ गृहीत्वा त्ववलम्बि तस्य~॥१९॥}\footnoteA{Cakrapāṇidatta said that others read \emph{nāsāsandhānavidhim} here. Ḍalhaṇa (1.16.27–31) stated that some read, \emph{chinnāṃ tu nāsikāṃ dṛṣṭvā vayaḥsthasya śarīriṇaḥ | nāsānurūpaṃ saṃcchidya patraṃ gaṇḍe niveśayet ||}}\edtext{}{\lemma{\skt{यथावत्~। नासाप्रमाणं}}\Afootnote{A H; \skt{यथावत्~। नासाप्रमाण}\ N}}\pend

\pstart\versepar \ledsidenote{\passref{1.16.28}}\skt{तेन प्रमाणेन हि गण्डपार्श्वाद् }\pend
\pstart\versepar \skt{उत्कृत्य }\edtext{\skt{वध्रं}}{\lemma{\skt{वध्रं}}\Afootnote{N; \skt{बद्धं}\ A; \skt{वन्ध्र}\ H}}\skt{ त्वथ नासिकाग्रम्~। }\pend
\pstart\versepar \skt{विलिख्य चाशु प्रतिसन्दधीत }\pend
\pstart\versepar \skt{तं }\edtext{\skt{साधु\-ब\-द्ध\-म्भि\-ष\-ग\-प्र\-मत्तः}}{\lemma{\skt{साधुबद्धम्}}\Afootnote{H; \skt{त्साधुबन्धैर्}\ A; \skt{साधुवद्ध}\ N}}\skt{~॥२०॥}\edtext{}{\lemma{\skt{प्रतिसन्दधीत तं}}\Afootnote{N H; \skt{प्रतिसंदधीत त}\ A}}\pend

\pstart\versepar \ledsidenote{\passref{1.16.29}}\edtext{\skt{सुसीवितं}}{\lemma{\skt{सुसीवितं}}\Afootnote{\skt{सुसंहितं}\ A; \skt{सुसीवित}\ N; \skt{सुशीवितं}\ H}}\skt{ सम्यगतो यथावन् }\pend
\pstart\versepar \skt{नाडी\-द्व\-ये\-ना\-भि\-स\-मीक्ष्य नह्येत्~। }\pend
\pstart\versepar \skt{उन्नामयित्वा }\edtext{\skt{त्ववचूर्ण्णयीत}}{\lemma{\skt{त्ववचूर्ण्णयीत}}\Afootnote{N H; \skt{अवचूर्णयेत्तु}\ A}}\skt{ }\pend
\pstart\versepar \skt{}\edtext{\skt{पत्ता\-ङ्ग\-य\-ष्टी\-म\-धु\-का\-ञ्ज\-नैश्च}}{\lemma{\skt{पत्ता\appbrk{}ङ्ग\appbrk{}य\appbrk{}ष्टी\appbrk{}म\appbrk{}धु\appbrk{}का\appbrk{}ञ्ज\appbrk{}नैश्}}\Afootnote{\skt{पत\appbrk{}ङ्ग\appbrk{}य\appbrk{}ष्टी\appbrk{}म\appbrk{}धु\appbrk{}का\appbrk{}ञ्ज\appbrk{}नैश्}\ A N; \skt{पत्त्रा\appbrk{}ङ्ग\appbrk{}य\appbrk{}ष्टी\appbrk{}म\appbrk{}धु\appbrk{}का\appbrk{}ञ्ज\appbrk{}नैश्}\ H}}\skt{~॥२१॥}\edtext{}{\lemma{\skt{नह्येत्~। उन्नामयित्वा}}\Afootnote{N H; \skt{बद्ध्वा~। प्रोन्नम्य चैनाम्}\ A}}\pend

\pstart\versepar \ledsidenote{\passref{1.16.30}}\skt{संछाद्य सम्यक्पिचुना }\edtext{\skt{व्रणन्तु}}{\lemma{\skt{व्रणन्तु}}\Afootnote{H; \skt{सितेन}\ A; \skt{व्रण तुन्}\ N}}\skt{ }\pend
\pstart\versepar \skt{तैलेन सिञ्चे\-द\-स\-कृ\-त्ति\-ला\-नाम्~। }\pend
\pstart\versepar \skt{घृतञ्च पाय्यः स नरः सुजीर्ण्णे }\pend
\pstart\versepar \skt{स्निग्धो }\edtext{\skt{विरेच्यः}}{\lemma{\skt{विरेच्यः}}\Afootnote{A; \skt{विरेच्य}\ N H}}\skt{ }\edtext{\skt{स्वय\-थो\-प\-दे\-शम्}}{\lemma{\skt{स्वय\appbrk{}थो\appbrk{}प\appbrk{}दे\appbrk{}शम्}}\Afootnote{H; \skt{स यथोपदेशम्}\ A; \skt{स्वयथोपदेशः}\ N}}\skt{~॥२२॥}\pend

\pstart\versepar \ledsidenote{\passref{1.16.31}}\skt{रूढञ्च }\edtext{\edtext{\skt{सन्धा\-न\-मु\-पा\-गतं}}{\lemma{\skt{सन्धानम्}}\Afootnote{A H; \skt{सन्धाम्}\ N}}\skt{ वै}}{\lemma{\skt{उपागतं वै}}\Afootnote{N; \skt{उपागतं स्या}\ A; \skt{उपागतश्चै}\ H}}\skt{ }\pend
\pstart\versepar \skt{}\edtext{\skt{तद्वध्रशेषं तु}}{\lemma{\skt{तद्वध्रशेषं तु}}\Afootnote{N H; \skt{त्तदर्धशेषं तु}\ A}}\skt{ पुनर्निकृन्तेत्~। }\pend
\pstart\versepar \skt{हीन\-म्पु\-न\-र्व\-र्द्ध\-यितुं }\edtext{\skt{यतेत}}{\lemma{\skt{यतेत}}\Afootnote{A H; \skt{यतेतः}\ N}}\skt{ }\pend
\pstart\versepar \skt{}\edtext{\skt{समञ्च}}{\lemma{\skt{समञ्च}}\Afootnote{N H; \skt{समां च}\ A}}\skt{ }\edtext{\skt{कुर्या\-द\-ति\-वृ\-द्ध\-मां\-सम्}}{\lemma{\skt{कुर्या\appbrk{}द\appbrk{}ति\appbrk{}वृ\appbrk{}द्ध\appbrk{}मां\appbrk{}सम्}}\Afootnote{N H; \skt{कुर्या\appbrk{}द\appbrk{}ति\appbrk{}वृ\appbrk{}द्ध\appbrk{}मां\appbrk{}साम्}\ A}}\skt{~॥२३॥ }\pend
\pstart\versepar \skt{}\edtext{\skt{इति}}{\lemma{\skt{इति}}\Afootnote{N H; \emph{om.}\ A}}\skt{~॥ }\edtext{\skt{ओम्}}{\lemma{\skt{ओम्}}\Afootnote{H; \emph{om.}\ A N}}\skt{~॥}\edtext{}{\lemma{\skt{निकृन्तेत्~। हीनम्}}\Afootnote{N H; \skt{निकृन्तेत्~। हीनां}\ A}}\pend

\endnumbering
\beginnumbering
\eledsection*{Adhyāya 17}
\pstart \ledsidenote{\passref{1.17.1}}\skt{अथात }\edtext{\skt{आम\-प\-क्वै\-ष\-णी\-य\-म्व्या\-ख्या\-स्यामः}}{\lemma{\skt{आम\appbrk{}प\appbrk{}क्वै\appbrk{}ष\appbrk{}णी\appbrk{}यम्}}\Afootnote{\skt{आम\appbrk{}प\appbrk{}क्वै\appbrk{}ष\appbrk{}णी\appbrk{}य\appbrk{}म\appbrk{}ध्यायं}\ A; \skt{आम\appbrk{}प\appbrk{}क्वे\appbrk{}स\appbrk{}ष\appbrk{}नी\appbrk{}य\appbrk{}म\appbrk{}ध्यायं}\ N; \skt{आम\appbrk{}प\appbrk{}क्वे\appbrk{}ष\appbrk{}णी\appbrk{}यम्}\ H}}\skt{~॥}\pend

\pstart \ledsidenote{\passref{1.17.3}}\edtext{\skt{अथ}}{\lemma{\skt{अथ}}\Afootnote{N H; \emph{om.}\ A}}\skt{ }\edtext{\skt{शोफसमुत्थाना}}{\lemma{\skt{शोफसमुत्थाना}}\Afootnote{A H; \skt{सौफसमुत्थाना}\ N}}\skt{ }\edtext{\skt{ग्रन्थि\-वि\-द्र\-ध्य\-ल\-जी\-प्र\-भृ\-तयः}}{\lemma{\skt{ग्रन्थि\appbrk{}वि\appbrk{}द्र\appbrk{}ध्य\appbrk{}ल\appbrk{}जी\appbrk{}प्र\appbrk{}भृ\appbrk{}तयः}}\Afootnote{A N; \skt{ग्र…न्थि\appbrk{}वि\appbrk{}द्र\appbrk{}ध्य\appbrk{}ल\appbrk{}जी\appbrk{}प्र\appbrk{}भृ\appbrk{}तयः}\ H}}\skt{~। प्रायेण }\edtext{\skt{व्याधयोऽभिहिता}}{\lemma{\skt{ऽभिहिता}}\Afootnote{A H; \skt{भिहिता}\ N}}\skt{ }\edtext{\skt{अनेकाकृतयः~। }\edtext{\skt{तैर्व्विलक्षणः}}{\lemma{\skt{व्विलक्षणः}}\Afootnote{A H; \skt{र्विलक्षणः}\ N}}}{\lemma{\skt{अनेकाकृतयः~। तैर्}}\Afootnote{A H; \skt{अनेकाकृतयः~। तै}\ N}}\skt{ }\edtext{\edtext{\skt{पृथुर्ग्रथितः}}{\lemma{\skt{पृथुर्}}\Afootnote{N H; \skt{पृथु}\ A}}}{\lemma{\skt{ग्रथितः}}\Afootnote{N; \skt{र्ग्रथितः}\ A; \skt{ग्ग्र…थितः}\ H}}\skt{ समो }\edtext{\skt{विषमो वा}}{\lemma{\skt{विषमो वा}}\Afootnote{A H; \skt{विसमो वा}\ N}}\skt{ }\edtext{\skt{त्वग्मांसस्थायी}}{\lemma{\skt{त्वग्मांसस्थायी}}\Afootnote{\skt{त्वङ्मांसस्थायी}\ A; \skt{त्वद्मात्सस्थायी}\ N; \skt{त्वद्मांसस्थायी}\ H}}\skt{ }\edtext{\skt{संघातः}}{\lemma{\skt{संघातः}}\Afootnote{H; \skt{दोषसङ्घातः}\ A; \skt{सङ्व्मातः}\ N}}\skt{ }\edtext{\skt{शरी\-रै\-क\-दे\-शो\-त्थितः}}{\lemma{\skt{शरी\appbrk{}रै\appbrk{}क\appbrk{}दे\appbrk{}शो\appbrk{}त्थितः}}\Afootnote{A H; \skt{शरी\appbrk{}रै\appbrk{}क\appbrk{}दे\appbrk{}शो\appbrk{}स्थितः}\ N}}\skt{ शोफ इत्युच्यते~।}\pend

\pstart \ledsidenote{\passref{1.17.4}}\edtext{\skt{वात\-पि\-त्त\-क\-फ\-शो\-णि\-त\-स\-न्नि\-पा\-ता\-ग\-न्तु\-कनि ररुणः}}{\lemma{\skt{वात\appbrk{}पि\appbrk{}त्त\appbrk{}क\appbrk{}फ\appbrk{}शो\appbrk{}णि\appbrk{}त\appbrk{}स\appbrk{}न्नि\appbrk{}पा\appbrk{}ता\appbrk{}ग\appbrk{}न्तु\appbrk{}कनि ररुणः}}\Afootnote{K; \skt{स षड्विधो वात\appbrk{}पि\appbrk{}त्त\appbrk{}क\appbrk{}फ\appbrk{}शो\appbrk{}णि\appbrk{}त\appbrk{}स\appbrk{}न्नि\appbrk{}पा\appbrk{}ता\appbrk{}ग\appbrk{}न्तु\appbrk{}नि\appbrk{}मित्तः~। तस्य दोष\appbrk{}रू\appbrk{}प\appbrk{}व्य\appbrk{}ञ्ज\appbrk{}नै\appbrk{}र्ल\appbrk{}क्ष\appbrk{}णानि व्याख्यास्यामः~। तत्र वातशोफः कृष्णोऽरुणो}\ A; \skt{स षट्विधौ भवति~। वात\appbrk{}पि\appbrk{}त्त\appbrk{}क\appbrk{}फ\appbrk{}सो\appbrk{}णि\appbrk{}त\appbrk{}स\appbrk{}न्नि\appbrk{}पा\appbrk{}ता\appbrk{}गं\appbrk{}तु\appbrk{}क\appbrk{}नि\appbrk{}मित्तः तत्र वातस्च यथुररुणः}\ N; \skt{स षड्विधो भवति~। वात\appbrk{}पि\appbrk{}त्त\appbrk{}क\appbrk{}फ\appbrk{}स\appbrk{}न्नि\appbrk{}पा\appbrk{}ता\appbrk{}ग\appbrk{}न्तु\appbrk{}क\appbrk{}नि\appbrk{}मित्तः~। तत्र वात\appbrk{}श्व\appbrk{}य\appbrk{}थु\appbrk{}र\appbrk{}रुणः}\ H\textsuperscript{ac}; \skt{स षड्विधो भवति~। वात\appbrk{}पि\appbrk{}त्त\appbrk{}क\appbrk{}फ\appbrk{}शो\appbrk{}णि\appbrk{}त\appbrk{}स\appbrk{}न्नि\appbrk{}पा\appbrk{}ता\appbrk{}ग\appbrk{}न्तु\appbrk{}क\appbrk{}नि\appbrk{}मित्तः~। तत्र वात\appbrk{}श्व\appbrk{}य\appbrk{}थु\appbrk{}र\appbrk{}रुणः}\ H\textsuperscript{pc}}}\skt{ }\edtext{\skt{कृष्णो वा}}{\lemma{\skt{कृष्णो वा}}\Afootnote{K N H; \skt{वा}\ A}}\skt{ परुषो }\edtext{\skt{मृदु\-र\-न\-व\-स्थि\-त\-स्तो\-दा\-द\-य\-श्चात्र}}{\lemma{\skt{मृदु\appbrk{}र\appbrk{}न\appbrk{}व\appbrk{}स्थि\appbrk{}त\appbrk{}स्तो\appbrk{}दा\appbrk{}द\appbrk{}यश्}}\Afootnote{K; \skt{मृदु\appbrk{}र\appbrk{}न\appbrk{}व\appbrk{}स्थि\appbrk{}ता\appbrk{}स्तो\appbrk{}दा\appbrk{}द\appbrk{}यश्}\ A; \skt{मृदु\appbrk{}र\appbrk{}न\appbrk{}व\appbrk{}स्थि\appbrk{}ता\appbrk{}स्तो\appbrk{}दा\appbrk{}यस्}\ N; \skt{मृदु\appbrk{}र\appbrk{}न\appbrk{}व\appbrk{}स्थि\appbrk{}त\appbrk{}स्तो\appbrk{}दा\appbrk{}द\appbrk{}यश्}\ H}}\skt{ }\edtext{\skt{वेदनाविशेषो}}{\lemma{\skt{वेदनाविशेषो}}\Afootnote{K; \skt{वेदनाविशेषा}\ A H; \skt{वेदनाविसेषा}\ N}}\skt{ भवन्ति~। }\edtext{\skt{पित्तश्वयथुः}}{\lemma{\skt{पित्तश्वयथुः}}\Afootnote{K H; \skt{पित्तशोफः}\ A; \skt{पित्तस्च यथु}\ N}}\skt{ पीतः }\edtext{\skt{सरक्तो वा}}{\lemma{\skt{सरक्तो वा}}\Afootnote{K N H; \skt{मृदुः सरक्तो वा}\ A}}\skt{ }\edtext{\edtext{\skt{शीघ्रा\-नु\-सा\-री\-मृ\-दु\-र्दा\-हा\-द\-य\-श्चात्र}}{\lemma{\skt{शीघ्रा\appbrk{}नु\appbrk{}सा\appbrk{}री\appbrk{}मृ\appbrk{}दुर्}}\Afootnote{K H; \skt{शीघ्रानुसार्यो}\ A; \skt{शीघ्रानुसारी मृदु}\ N}}}{\lemma{\skt{दाहादयश्}}\Afootnote{K; \skt{षादयश्}\ A; \skt{दार्हादयश्}\ N; \skt{द्दादयश्}\ H}}\skt{ }\edtext{\skt{वेदनाविशेषो}}{\lemma{\skt{वेदनाविशेषो}}\Afootnote{K; \skt{वेदनाविशेषा}\ A H; \skt{वेदनाविसेषा}\ N}}\skt{ भवन्ति~। }\edtext{\skt{श्लेष्मश्वयथुः}}{\lemma{\skt{श्लेष्मश्वयथुः}}\Afootnote{A K H; \skt{स्लेष्मश्च यथुः}\ N}}\skt{ पाण्डुः }\edtext{\skt{शुक्लो वा}}{\lemma{\skt{शुक्लो वा}}\Afootnote{K N H; \emph{om.}\ A}}\skt{ कठिनः }\edtext{\skt{स्निग्धो}}{\lemma{\skt{स्निग्धो}}\Afootnote{K H; \skt{स्निग्धः शीतो स्निग्धो}\ A; \skt{शितः स्निग्धो}\ N}}\skt{ मन्दानुसारी कण्ड्वादयश्चात्र }\edtext{\skt{वेदनाविशेषो}}{\lemma{\skt{वेदनाविशेषो}}\Afootnote{K; \skt{वेदनाविशेषा}\ A N H}}\skt{ भवन्ति~। }\edtext{\skt{सन्नि\-पा\-त\-श्व\-यथुः}}{\lemma{\skt{सन्नि\appbrk{}पा\appbrk{}त\appbrk{}श्व\appbrk{}यथुः}}\Afootnote{K; \skt{सर्व\appbrk{}व\appbrk{}र्ण\appbrk{}वे\appbrk{}दनः सन्नि\appbrk{}पा\appbrk{}त\appbrk{}श्व\appbrk{}यथुः}\ A; \skt{सन्निपातश्च यथुः}\ N; \skt{सन्नि\appbrk{}पा\appbrk{}त\appbrk{}श्व\appbrk{}य\appbrk{}तुअ}\ H\textsuperscript{ac}; \skt{सन्नि\appbrk{}पा\appbrk{}त\appbrk{}श्व\appbrk{}य\appbrk{}तुअः}\ H\textsuperscript{pc}}}\skt{ }\edtext{\skt{सर्व\-दो\-ष\-लि\-ङ्ग\-वि\-शे\-षो\-पेतः}}{\lemma{\skt{सर्व\appbrk{}दो\appbrk{}ष\appbrk{}लि\appbrk{}ङ्ग\appbrk{}वि\appbrk{}शे\appbrk{}षो\appbrk{}पेतः}}\Afootnote{K N; \emph{om.}\ A; \skt{सर्व्व\appbrk{}दो\appbrk{}ष\appbrk{}लि\appbrk{}ङ्ग\appbrk{}वि\appbrk{}शे\appbrk{}षो\appbrk{}पेतेः}\ H}}\skt{~। }\edtext{\skt{पित्त\-व\-च्छो\-णि\-त\-जो\-ति\-कृ\-ष्णश्च}}{\lemma{\skt{पित्त\appbrk{}व\appbrk{}च्छो\appbrk{}णि\appbrk{}त\appbrk{}जो\appbrk{}ति\appbrk{}कृ\appbrk{}ष्णश्च}}\Afootnote{K H; \skt{पित्त\appbrk{}व\appbrk{}च्छो\appbrk{}णि\appbrk{}त\appbrk{}जोऽ\appbrk{}ति\appbrk{}कृ\appbrk{}ष्णश्च}\ A; \skt{पित्त\appbrk{}व\appbrk{}च्छो\appbrk{}णि\appbrk{}त\appbrk{}जो\appbrk{}ति\appbrk{}कृ\appbrk{}त्सश्चा}\ N}}\skt{~। }\edtext{\edtext{\skt{पित्त\-र\-क्त\-ल\-क्ष\-ण\-श्चा\-ग\-न्तु\-र्लो\-हि\-ता\-व\-भा\-सश्च}}{\lemma{\skt{पित्त\appbrk{}र\appbrk{}क्त\appbrk{}ल\appbrk{}क्ष\appbrk{}णश्}}\Afootnote{K N H; \skt{पित्त\appbrk{}र\appbrk{}क्त\appbrk{}ल\appbrk{}क्षण}\ A}}}{\lemma{\skt{चाग\appbrk{}न्तु\appbrk{}र्लो\appbrk{}हि\appbrk{}ता\appbrk{}व\appbrk{}भा\appbrk{}सश्}}\Afootnote{K; \skt{आग\appbrk{}न्तु\appbrk{}र्लो\appbrk{}हि\appbrk{}ता\appbrk{}व\appbrk{}भा\appbrk{}सश्}\ A; \skt{चाग\appbrk{}न्तु\appbrk{}लो\appbrk{}हि\appbrk{}ता\appbrk{}व\appbrk{}भा\appbrk{}सश्}\ N; \skt{चाग\appbrk{}न्तु\appbrk{}र्ला\appbrk{}हि\appbrk{}ता\appbrk{}भा\appbrk{}सश्}\ H}}\skt{~।}\pend

\pstart \ledsidenote{\passref{1.17.5}}\skt{स यदा बाह्याभ्यन्तरैः }\edtext{\skt{क्रियाविशेषैर्न}}{\lemma{\skt{क्रियाविशेषैर्न}}\Afootnote{K N; \skt{क्रियाविशेषैर्न}\ A; \skt{क्रियाविशेषैर्न्}\ H}}\skt{ }\edtext{\skt{शक्यते}}{\lemma{\skt{शक्यते}}\Afootnote{K; \skt{सम्भावितः}\ A; \skt{सकृते}\ N; \skt{न शक्यते}\ H}}\skt{ }\edtext{\skt{प्रशमयितुं}}{\lemma{\skt{प्रशमयितुं}}\Afootnote{A K H; \skt{प्रभेमयितुं}\ N}}\skt{ }\edtext{\skt{क्रिया\-वि\-प\-र्य\-या\-द्ब\-हु\-त्वाद्वा}}{\lemma{\skt{क्रिया\appbrk{}वि\appbrk{}प\appbrk{}र्य\appbrk{}या\appbrk{}द्ब\appbrk{}हु\appbrk{}त्वाद्वा}}\Afootnote{K H; \skt{क्रिया\appbrk{}वि\appbrk{}प\appbrk{}र्य\appbrk{}या\appbrk{}द्ब\appbrk{}हु\appbrk{}त्वाद्वा}\ A N}}\skt{ दोषाणां }\edtext{\skt{पाकयाभिमुखो}}{\lemma{\skt{पाकयाभिमुखो}}\Afootnote{K\textsuperscript{pc} H; \skt{तदा पाकाभिमुखो}\ A; \skt{पाक\appbrk{}आ\appbrk{}या\appbrk{}भि\appbrk{}मुखो}\ K\textsuperscript{ac}; \skt{पाकायाभिमुखो}\ N}}\skt{ भवति }\edtext{\skt{तस्यामस्य}}{\lemma{\skt{तस्यामस्य}}\Afootnote{A K H; \skt{तस्याम्स्य}\ N}}\skt{ पच्यमानस्य }\edtext{\skt{पक्वस्य च}}{\lemma{\skt{पक्वस्य च}}\Afootnote{A K N; \skt{पक्वस्य}\ H}}\skt{ }\edtext{\edtext{\edtext{\skt{लक्ष\-ण\-मु\-च्य\-मा\-न\-मु\-प\-धा\-र\-यश्च}}{\lemma{\skt{लक्षणम्}}\Afootnote{K H; \skt{लक्षण}\ A N}}}{\lemma{\skt{उच्यमानम्}}\Afootnote{K H; \skt{मुच्यमानम्}\ A N}}}{\lemma{\skt{उपधारयश्च}}\Afootnote{K N; \skt{उपधारय}\ A; \skt{उपधारयस्व}\ H}}\skt{~। तत्र }\edtext{\skt{मन्दोष्मता}}{\lemma{\skt{मन्दोष्मता}}\Afootnote{A K H; \skt{मन्दोस्मता}\ N}}\skt{ }\edtext{\edtext{\skt{त्वक्सा\-व\-र्ण्य\-न्स्थै\-र्यं\-म\-ल्प\-रु\-ज\-ता\-ल्प\-शो\-फता}}{\lemma{\skt{त्वक्सा\appbrk{}व\appbrk{}र्ण्य\appbrk{}न्स्थै\appbrk{}र्यंम्}}\Afootnote{K; \skt{त्वक्सवर्णता शीतशोफता स्थैर्यं मन्}\ A; \skt{त्वन्सा\appbrk{}व\appbrk{}ल्ली\appbrk{}स्थै\appbrk{}र्यर्य}\ N; \skt{त्वक्साव्वर्ण्य स्थैर्यम्}\ H}}}{\lemma{\skt{अल्प\appbrk{}रु\appbrk{}ज\appbrk{}ता\appbrk{}ल्प\appbrk{}शो\appbrk{}फता}}\Afootnote{K H; \skt{दवे\appbrk{}द\appbrk{}न\appbrk{}ताऽ\appbrk{}ल्प\appbrk{}शो\appbrk{}फता}\ A; \skt{मल्प\appbrk{}रु\appbrk{}ज\appbrk{}ता\appbrk{}श्प\appbrk{}शो\appbrk{}फता}\ N}}\skt{ }\edtext{\edtext{\skt{चाम\-ल\-क्ष\-ण\-मु\-द्दिष्टं}}{\lemma{\skt{चामलक्षणम्}}\Afootnote{K H; \skt{चामलक्षण}\ A N}}}{\lemma{\skt{उद्दिष्टं}}\Afootnote{K H; \skt{मुद्दिष्टं}\ A N}}\skt{~॥ }\edtext{\skt{सूचीभिरिव}}{\lemma{\skt{सूचीभिर्}}\Afootnote{K H; \skt{सूचिभिर्}\ A N}}\skt{ }\edtext{\skt{निस्तुद्यते}}{\lemma{\skt{निस्तुद्यते}}\Afootnote{A K H; \skt{निस्तूद्यते}\ N}}\skt{ }\edtext{\skt{दश्यत}}{\lemma{\skt{दश्यत}}\Afootnote{A K H; \emph{om.}\ N}}\skt{ }\edtext{\skt{इव च}}{\lemma{\skt{इव च}}\Afootnote{K N H; \skt{इव}\ A}}\skt{ }\edtext{\edtext{\skt{पिपी\-लि\-का\-भि\-श्छि\-द्यते}}{\lemma{\skt{पिपीलिकाभिश्}}\Afootnote{K N H; \skt{पिपीलिकाभिः, ताभिश्च}\ A}}}{\lemma{\skt{छिद्यते}}\Afootnote{K H; \skt{संसर्प्यत इव, छिद्यत इव श}\ A; \skt{छि}\ N}}\skt{ }\edtext{\skt{भिद्यत}}{\lemma{\skt{भिद्यत}}\Afootnote{K H; \skt{स्त्रेण, भिद्यत}\ A; \skt{द्यते भिद्यते}\ N}}\skt{ }\edtext{\skt{इव च}}{\lemma{\skt{इव च}}\Afootnote{K H; \skt{इव}\ A; \emph{om.}\ N}}\skt{ }\edtext{\skt{शस्त्रेण}}{\lemma{\skt{शस्त्रेण}}\Afootnote{K H; \skt{शक्तिभिः}\ A; \emph{om.}\ N}}\skt{ }\edtext{\skt{ताड्यत}}{\lemma{\skt{ताड्यत}}\Afootnote{A K H; \emph{om.}\ N}}\skt{ }\edtext{\skt{इव च}}{\lemma{\skt{इव च}}\Afootnote{K N H; \skt{इव}\ A}}\skt{ }\edtext{\skt{दण्डेनभिः पीड्यत}}{\lemma{\skt{दण्डेनभिः पीड्यत}}\Afootnote{K; \skt{दण्डेन, पीड्यत}\ A; \skt{दण्डेन निःपीड्यत}\ N H}}\skt{ }\edtext{\skt{इव च}}{\lemma{\skt{इव च}}\Afootnote{K N H; \skt{इव}\ A}}\skt{ पाणिना घट्यत इव }\edtext{\skt{चांगुल्या}}{\lemma{\skt{चांगुल्या}}\Afootnote{A K H; \skt{चागुल्या}\ N}}\skt{ दह्यते पच्यत इव }\edtext{\skt{चाग्निक्षाराभ्यां}}{\lemma{\skt{चाग्निक्षाराभ्यां}}\Afootnote{A K H; \skt{चाग्निक्षाराभ्या}\ N}}\skt{ }\edtext{\skt{मूषा}}{\lemma{\skt{मूषा}}\Afootnote{K N; \skt{ओष}\ A; \skt{मूष}\ H}}\skt{ }\edtext{\skt{चोष\-प\-रि\-दा\-हाश्च}}{\lemma{\skt{चोष\appbrk{}प\appbrk{}रि\appbrk{}दा\appbrk{}हाश्च}}\Afootnote{K N; \skt{चोष\appbrk{}प\appbrk{}री\appbrk{}दा\appbrk{}हाश्च}\ A H}}\skt{ भवन्ति~। वृश्चिकविद्ध इव च स्थाना\-स\-न\-श\-य\-नेषु न शान्तिमुपैति~। }\edtext{\edtext{\skt{आध्मा\-त\-ब\-स्ति\-रि\-वा\-त\-तश्च}}{\lemma{\skt{आध्मातबस्तिर्}}\Afootnote{K H\textsuperscript{pc}; \skt{आध्मातबस्ति}\ A N; \skt{आध्मातबम्तिर्}\ H\textsuperscript{ac}}}}{\lemma{\skt{इवाततश्च}}\Afootnote{K H; \skt{रिवाततश्च}\ A; \skt{रिवाततस्च}\ N}}\skt{ शोफो भवति }\edtext{\skt{त्वग्वैवर्ण्यं}}{\lemma{\skt{त्वग्वैवर्ण्यं}}\Afootnote{A K; \skt{त्वगैवर्ण्ण्यं}\ N; \skt{त्वक्वैवर्ण्ण्यं}\ H}}\skt{ }\edtext{\skt{शोफा\-भि\-वृ\-द्धि\-र्ज्वरो दाहः}}{\lemma{\skt{शोफा\appbrk{}भि\appbrk{}वृ\appbrk{}द्धि\appbrk{}र्ज्वरो दाहः}}\Afootnote{K; \skt{शोफा\appbrk{}भि\appbrk{}वृ\appbrk{}द्धि\appbrk{}र्ज्व\appbrk{}र\appbrk{}दाह}\ A; \skt{शोफा\appbrk{}ति\appbrk{}वृ\appbrk{}द्धि\appbrk{}ज्व\appbrk{}रो\appbrk{}दाहः}\ N; \skt{शोफा\appbrk{}ति\appbrk{}वृ\appbrk{}द्धि\appbrk{}र्ज्ज्व\appbrk{}रो\appbrk{}दाहः}\ H}}\skt{ पिपासा भक्ता रुचिश्च पच्यमानलिङ्गं~॥ }\edtext{\skt{वेद\-नो\-प\-शा\-न्ति\-र्नि\-र्लो\-हि\-ता\-ल्प\-शो\-फता च}}{\lemma{\skt{वेद\appbrk{}नो\appbrk{}प\appbrk{}शा\appbrk{}न्ति\appbrk{}र्नि\appbrk{}र्लो\appbrk{}हि\appbrk{}ता\appbrk{}ल्प\appbrk{}शो\appbrk{}फता च}}\Afootnote{K\textsuperscript{pc}; \skt{वेदनोपशान्तिः पाण्डु\appbrk{}ताऽ\appbrk{}ल्प\appbrk{}शो\appbrk{}फता}\ A; \skt{वेद\appbrk{}नो\appbrk{}प\appbrk{}शा\appbrk{}न्ति\appbrk{}र्लो\appbrk{}हि\appbrk{}ता\appbrk{}ल्प\appbrk{}शो\appbrk{}फता च}\ K\textsuperscript{ac}; \skt{वेद\appbrk{}नो\appbrk{}प\appbrk{}शा\appbrk{}न्ति\appbrk{}नि\appbrk{}र्लौ\appbrk{}हि\appbrk{}ता\appbrk{}आ\appbrk{}ल्प\appbrk{}शो\appbrk{}फता च}\ N; \skt{वेद\appbrk{}नो\appbrk{}प\appbrk{}शा\appbrk{}न्ति\appbrk{}र्न्नि\appbrk{}र्ल्लो\appbrk{}हि\appbrk{}ता\appbrk{}ल्प\appbrk{}शो\appbrk{}फता च}\ H}}\skt{ }\edtext{\skt{बलीप्रादुर्भावः}}{\lemma{\skt{बलीप्रादुर्भावः}}\Afootnote{K N; \skt{वलीप्रादुर्भाव}\ A; \skt{वलीप्रादुर्भ्भाव}\ H}}\skt{ }\edtext{\skt{त्वक्परिपोटनं}}{\lemma{\skt{त्वक्परिपोटनं}}\Afootnote{K N H; \skt{स्त्वक्परिपुटनं}\ A}}\skt{ }\edtext{\edtext{\skt{निम्न\-द\-र्श\-न\-म\-ङ्गु\-ल्या\-व\-पी\-डिते}}{\lemma{\skt{निम्नदर्शनम्}}\Afootnote{K H; \skt{निम्नदर्शन}\ A N}}}{\lemma{\skt{अङ्गु\appbrk{}ल्या\appbrk{}व\appbrk{}पी\appbrk{}डिते}}\Afootnote{K; \skt{मङ्गु\appbrk{}ल्याऽ\appbrk{}व\appbrk{}पी\appbrk{}डिते}\ A N; \skt{अङ्गुल्या च पीडिते}\ H}}\skt{ }\edtext{\skt{बस्ता\-वि\-व\-चो\-द\-क\-स\-ञ्च\-रणं}}{\lemma{\skt{बस्ता\appbrk{}वि\appbrk{}व\appbrk{}चो\appbrk{}द\appbrk{}क\appbrk{}स\appbrk{}ञ्च\appbrk{}रणं}}\Afootnote{K H; \skt{प्रत्युन्नमनं, बस्ता\appbrk{}वि\appbrk{}वो\appbrk{}द\appbrk{}क\appbrk{}स\appbrk{}ञ्च\appbrk{}रणं}\ A; \skt{बस्ताविव चोद\appbrk{}क\appbrk{}स\appbrk{}ञ्च\appbrk{}रण}\ N}}\skt{ पूयस्य }\edtext{\edtext{\edtext{\edtext{\skt{प्रपी\-ड\-य\-त्ये\-क\-म\-न्त\-मन्ते}}{\lemma{\skt{प्रपीडयत्य्}}\Afootnote{K N H; \skt{प्रपीडय}\ A}}}{\lemma{\skt{एकम्}}\Afootnote{K H; \skt{त्येक}\ A; \skt{एक}\ N}}}{\lemma{\skt{अन्तम्}}\Afootnote{K H; \skt{मन्त}\ A; \skt{मत्त}\ N}}}{\lemma{\skt{अन्ते}}\Afootnote{K H; \skt{मन्ते}\ A N}}\skt{ चावपीडिते }\edtext{\edtext{\skt{मुहुर्मुहुस्तोदः}}{\lemma{\skt{मुहुर्मुहुस्}}\Afootnote{K; \skt{मुहुर्मुहु}\ A N; \skt{मुहुर्मुहुस्}\ H}}}{\lemma{\skt{तोदः}}\Afootnote{K H\textsuperscript{ac}; \skt{स्तोदः}\ A; \skt{स्तादः}\ N; \skt{तोदः ल}\ H\textsuperscript{pc}}}\skt{ }\edtext{\skt{कण्डू\-र\-नु\-न्न\-तता}}{\lemma{\skt{कण्डू\appbrk{}र\appbrk{}नु\appbrk{}न्न\appbrk{}तता}}\Afootnote{K H; \skt{कण्डूरुन्नतता}\ A; \skt{कण्डू\appbrk{}र\appbrk{}णु\appbrk{}न्न\appbrk{}नता}\ N}}\skt{ }\edtext{\edtext{\edtext{\skt{व्याधे\-रु\-प\-द्र\-व\-शा\-न्ति\-र्भ\-क्ता\-भि\-काङ्क्षा}}{\lemma{\skt{व्याधेर्}}\Afootnote{K H; \skt{व्याधे}\ A N}}}{\lemma{\skt{उप\appbrk{}द्र\appbrk{}व\appbrk{}शा\appbrk{}न्तिर्}}\Afootnote{K H; \skt{रुपद्रवशान्ति}\ A N}}\skt{ च}}{\lemma{\skt{भक्ताभिकाङ्क्षा च}}\Afootnote{K N; \skt{र्भक्ताभिकाङ्क्षा च}\ A; \skt{भ्भक्ताभिकांक्षा च}\ H}}\skt{ }\edtext{\skt{परिपक्वलिङ्गं}}{\lemma{\skt{परिपक्वलिङ्गं}}\Afootnote{K N H; \skt{पक्वलिङ्गम्}\ A}}\skt{~॥ कफजेषु }\edtext{\skt{खलु}}{\lemma{\skt{खलु}}\Afootnote{K N H; \skt{तु}\ A}}\skt{ रोगेषु }\edtext{\skt{गम्भी\-रा\-नु\-ग\-त\-त्वा\-द\-भि\-घा\-त\-जेषु च}}{\lemma{\skt{गम्भी\appbrk{}रा\appbrk{}नु\appbrk{}ग\appbrk{}त\appbrk{}त्वा\appbrk{}द\appbrk{}भि\appbrk{}घा\appbrk{}त\appbrk{}जेषु च}}\Afootnote{K H; \skt{गम्भी\appbrk{}र\appbrk{}ग\appbrk{}ति\appbrk{}त्वा\appbrk{}द\appbrk{}भि\appbrk{}घा\appbrk{}त\appbrk{}जेषु वा}\ A; \skt{गम्भी\appbrk{}रा\appbrk{}नु\appbrk{}ग\appbrk{}त\appbrk{}त्वा\appbrk{}द\appbrk{}भि\appbrk{}घा\appbrk{}त\appbrk{}त्वा\appbrk{}द\appbrk{}भि\appbrk{}घा\appbrk{}त\appbrk{}जेषु च}\ N}}\skt{ }\edtext{\edtext{\skt{केषु\-चि\-द\-स\-म\-स्त\-म्प\-क्व\-ल\-क्षणं}}{\lemma{\skt{केषु\appbrk{}चि\appbrk{}द\appbrk{}स\appbrk{}म\appbrk{}स्तम्}}\Afootnote{A K H; \skt{केशु चदसमस्तम्}\ N}}}{\lemma{\skt{पक्वलक्षणं}}\Afootnote{A K N H\textsuperscript{ac}; \skt{पक्व\appbrk{}क्व\appbrk{}ल\appbrk{}क्ष\appbrk{}णन्}\ H\textsuperscript{pc}}}\skt{ दृष्ट्वा }\edtext{\edtext{\skt{पक्व\-म\-प\-क्व\-मिति}}{\lemma{\skt{पक्वम्}}\Afootnote{K N H; \skt{पक्व}\ A}}}{\lemma{\skt{अपक्वम्}}\Afootnote{K N H; \skt{मपक्वम्}\ A}}\skt{ मन्यमानो }\edtext{\skt{भिष\-ङ्मो\-ह\-मु\-पैति}}{\lemma{\skt{भिषङ्मोहम्}}\Afootnote{A K; \skt{भिषच्मोहम्}\ N; \skt{भिषहम्}\ H}}\skt{~। तत्र हि त्वक्सवर्ण्णता }\edtext{\skt{शीत\-शो\-फ\-ता\-ल्प\-रु\-ज\-ता\-श्म\-वच्च}}{\lemma{\skt{शीत\appbrk{}शो\appbrk{}फ\appbrk{}ता\appbrk{}ल्प\appbrk{}रु\appbrk{}ज\appbrk{}ता\appbrk{}श्म\appbrk{}वच्च}}\Afootnote{K N; \skt{शीतशोफता स्थैर्य\appbrk{}म\appbrk{}ल्प\appbrk{}रु\appbrk{}ज\appbrk{}ताऽ\appbrk{}श्म\appbrk{}वच्च}\ A; \skt{शीत\appbrk{}शो\appbrk{}फ\appbrk{}ता\appbrk{}ल्प\appbrk{}रु\appbrk{}ज\appbrk{}ता\appbrk{}स्प\appbrk{}वच्च}\ H}}\skt{ }\edtext{\skt{घनता न}}{\lemma{\skt{घनता न}}\Afootnote{A K H; \skt{पनता न}\ N}}\skt{ }\edtext{\skt{तत्र}}{\lemma{\skt{तत्र}}\Afootnote{K N H; \emph{om.}\ A}}\skt{ }\edtext{\edtext{\skt{मोहमुपेयात्}}{\lemma{\skt{मोहम्}}\Afootnote{A K H; \skt{मोह}\ N}}}{\lemma{\skt{उपेयात्}}\Afootnote{K H; \skt{उपेयादिति}\ A; \skt{मुपैयात्}\ N}}\skt{~।}\pend

\pstart\versepar \ledsidenote{\passref{1.17.6}}\edtext{\skt{भवति}}{\lemma{\skt{भवति}}\Afootnote{\skt{भवन्ति}\ A N; \emph{om.}\ K H}}\skt{ }\edtext{\skt{चात्र}}{\lemma{\skt{चात्र}}\Afootnote{A N; \emph{om.}\ K; \skt{भवति चात्र}\ H}}\skt{~॥ }\pend
\pstart\versepar \skt{}\edtext{\edtext{\skt{आम\-म्वि\-द\-ह्य\-मा\-नञ्च}}{\lemma{\skt{आमम्}}\Afootnote{A K H; \skt{आम}\ N}}}{\lemma{\skt{विदह्यमानञ्च}}\Afootnote{K H; \skt{विपच्यमानं च}\ A; \skt{ष्विदह्यमान}\ N}}\skt{ }\edtext{\skt{सम्यक्पक्वञ्च यो}}{\lemma{\skt{सम्यक्पक्वञ्च यो}}\Afootnote{A K H; \skt{सम्यक्पक्वञ्च यो}\ N}}\skt{ भिषक्~। }\pend
\pstart\versepar \skt{जानीयात्स भवेद्वैद्यः }\edtext{\edtext{\skt{शेषा\-स्त\-स्क\-र\-वृ\-त्तयः}}{\lemma{\skt{शेषास्}}\Afootnote{A K H; \skt{शेषा}\ N}}}{\lemma{\skt{तस्करवृत्तयः}}\Afootnote{A K H; \skt{स्तस्करवृत्तयः}\ N}}\skt{~॥}\pend

\pstart\versepar \ledsidenote{\passref{1.17.7}}\skt{वातादृते नास्ति रुजा न पाकः पित्तादृते }\edtext{\skt{नास्ति}}{\lemma{\skt{नास्ति}}\Afootnote{A K H; \skt{नस्ति}\ N}}\skt{ }\edtext{\skt{कफाच्च}}{\lemma{\skt{कफाच्च}}\Afootnote{A K H; \skt{कफा च}\ N}}\skt{ पूयः~। }\pend
\pstart\versepar \skt{}\edtext{\skt{तस्मात्समस्ताः}}{\lemma{\skt{तस्मात्समस्ताः}}\Afootnote{K; \skt{तस्मात्समस्ता}\ A; \skt{तस्मस्मात्समस्ताः}\ N; \skt{…~। तस्मात्समस्ताः}\ H}}\skt{ }\edtext{\skt{परिपाककाले}}{\lemma{\skt{परिपाककाले}}\Afootnote{K N H; \skt{न्परिपाककाले}\ A}}\skt{ पचन्ति }\edtext{\skt{शोफं}}{\lemma{\skt{शोफं}}\Afootnote{K N H; \skt{शोफांस्}\ A}}\skt{ त्रय एव दोषाः~॥}\pend

\pstart\versepar \ledsidenote{\passref{1.17.7.add1}}\edtext{\skt{नर्ते}}{\lemma{\skt{नर्ते}}\Afootnote{K\textsuperscript{pc} N H; \emph{om.}\ K\textsuperscript{ac}}}\skt{ }\edtext{\skt{रुजा}}{\lemma{\skt{रुजा}}\Afootnote{K\textsuperscript{pc} N; \emph{om.}\ K\textsuperscript{ac}; \skt{रूजा}\ H}}\skt{ }\edtext{\edtext{\skt{वातमृते}}{\lemma{\skt{वातम्}}\Afootnote{K\textsuperscript{pc}; \emph{om.}\ K\textsuperscript{ac}; \skt{वात}\ N H}}\skt{ च}}{\lemma{\skt{ऋते च}}\Afootnote{K\textsuperscript{pc}; \emph{om.}\ K\textsuperscript{ac}; \skt{षृते च}\ N; \skt{मृते च}\ H}}\skt{ }\edtext{\skt{पित्तं}}{\lemma{\skt{पित्तं}}\Afootnote{K\textsuperscript{pc} N H; \emph{om.}\ K\textsuperscript{ac}}}\skt{ }\pend
\pstart\versepar \skt{}\edtext{\skt{पाकः}}{\lemma{\skt{पाकः}}\Afootnote{K\textsuperscript{pc} N H; \emph{om.}\ K\textsuperscript{ac}}}\skt{ }\edtext{\edtext{\skt{कफश्चापिविना}}{\lemma{\skt{कफश्}}\Afootnote{K\textsuperscript{pc}; \emph{om.}\ K\textsuperscript{ac}; \skt{कफञ्}\ N H}}\skt{ न}}{\lemma{\skt{चापिविना न}}\Afootnote{K\textsuperscript{pc} H; \emph{om.}\ K\textsuperscript{ac}; \skt{चात्मविना न}\ N}}\skt{ }\edtext{\skt{पूयः}}{\lemma{\skt{पूयः}}\Afootnote{K\textsuperscript{pc} N H; \emph{om.}\ K\textsuperscript{ac}}}\skt{~। }\pend
\pstart\versepar \skt{}\edtext{\skt{तस्माद्विपाकं}}{\lemma{\skt{तस्माद्विपाकं}}\Afootnote{K\textsuperscript{pc} H; \emph{om.}\ K\textsuperscript{ac}; \skt{तस्माद्विपाक}\ N}}\skt{ }\edtext{\skt{परिपाककाले}}{\lemma{\skt{परिपाककाले}}\Afootnote{K\textsuperscript{pc} N H; \emph{om.}\ K\textsuperscript{ac}}}\skt{ }\pend
\pstart\versepar \skt{}\edtext{\skt{प्रयान्ति}}{\lemma{\skt{प्रयान्ति}}\Afootnote{K\textsuperscript{pc} N H; \emph{om.}\ K\textsuperscript{ac}}}\skt{ }\edtext{\edtext{\edtext{\skt{शोफास्त्रिभिरेव}}{\lemma{\skt{शोफास्}}\Afootnote{K\textsuperscript{pc} N H; \emph{om.}\ K\textsuperscript{ac}}}}{\lemma{\skt{त्रिभिर्}}\Afootnote{K\textsuperscript{pc} N H; \emph{om.}\ K\textsuperscript{ac}}}}{\lemma{\skt{एव}}\Afootnote{K\textsuperscript{pc} N; \emph{om.}\ K\textsuperscript{ac}; \skt{ऐव}\ H}}\skt{ }\edtext{\skt{दोषैः}}{\lemma{\skt{दोषैः}}\Afootnote{K\textsuperscript{pc} N H; \emph{om.}\ K\textsuperscript{ac}}}\skt{~॥}\pend

\pstart\versepar \ledsidenote{\passref{1.17.8}}\edtext{\skt{काला\-न्त\-रे\-णा\-भ्यु\-दि\-तन्तु}}{\lemma{\skt{काला\appbrk{}न्त\appbrk{}रे\appbrk{}णा\appbrk{}भ्यु\appbrk{}दि\appbrk{}तन्तु}}\Afootnote{A K H; \skt{काला\appbrk{}न्त\appbrk{}ना\appbrk{}भ्यु\appbrk{}दितं तु}\ N}}\skt{ पित्तं कृत्वा }\edtext{\skt{वशे}}{\lemma{\skt{वशे}}\Afootnote{A K N; \skt{वसे}\ H}}\skt{ वातकफौ }\edtext{\skt{प्रसह्य}}{\lemma{\skt{प्रसह्य}}\Afootnote{A K H; \skt{प्रसह्यः}\ N}}\skt{~। }\pend
\pstart\versepar \skt{}\edtext{\edtext{\skt{पचत्यतः}}{\lemma{\skt{पचत्य्}}\Afootnote{A K N; \skt{…~। पच}\ H}}}{\lemma{\skt{अतः}}\Afootnote{A K N; \skt{त्यतः}\ H}}\skt{ शोणितमेव पाको }\edtext{\skt{मतोपरेषां}}{\lemma{\skt{मतोपरेषां}}\Afootnote{K N H; \skt{मतोऽपरेषां}\ A}}\skt{ विदुषां द्वितीयः~॥}\pend

\pstart\versepar \ledsidenote{\passref{1.17.8.add1}}\edtext{\edtext{\skt{द्रव्या\-णा\-ञ्च\-न्द\-ना\-दीनां}}{\lemma{\skt{द्रव्याणाञ्}}\Afootnote{K\textsuperscript{pc} N; \emph{om.}\ K\textsuperscript{ac}; \skt{द्रव्यानाञ्}\ H}}}{\lemma{\skt{चन्दनादीनां}}\Afootnote{K\textsuperscript{pc} N H; \emph{om.}\ K\textsuperscript{ac}}}\skt{ }\edtext{\skt{दग्धानां}}{\lemma{\skt{दग्धानां}}\Afootnote{K\textsuperscript{pc} N H; \emph{om.}\ K\textsuperscript{ac}}}\skt{ }\edtext{\skt{श्वेतता}}{\lemma{\skt{श्वेतता}}\Afootnote{K\textsuperscript{pc}; \emph{om.}\ K\textsuperscript{ac}; \skt{स्वेततां}\ N; \skt{स्वेतता}\ H}}\skt{ }\edtext{\skt{यथा}}{\lemma{\skt{यथा}}\Afootnote{K\textsuperscript{pc} N H; \emph{om.}\ K\textsuperscript{ac}}}\skt{~। }\pend
\pstart\versepar \skt{}\edtext{\skt{तद्वत्पित्तोष्मणा}}{\lemma{\skt{तद्वत्पित्तोष्मणा}}\Afootnote{K\textsuperscript{pc} N H; \emph{om.}\ K\textsuperscript{ac}}}\skt{ }\edtext{\skt{दग्धं}}{\lemma{\skt{दग्धं}}\Afootnote{K\textsuperscript{pc} N H; \emph{om.}\ K\textsuperscript{ac}}}\skt{ }\edtext{\skt{रक्त}}{\lemma{\skt{रक्त}}\Afootnote{K\textsuperscript{pc}; \emph{om.}\ K\textsuperscript{ac}; \skt{रक्तं}\ N H}}\skt{ }\edtext{\edtext{\skt{पूयमिहोच्यते}}{\lemma{\skt{पूयम्}}\Afootnote{K\textsuperscript{pc} N H; \emph{om.}\ K\textsuperscript{ac}}}}{\lemma{\skt{इहोच्यते}}\Afootnote{K\textsuperscript{pc} N H; \emph{om.}\ K\textsuperscript{ac}}}\skt{~॥}\pend

\pstart \ledsidenote{\passref{1.17.9}}\edtext{\skt{तत्रामच्छेदे}}{\lemma{\skt{तत्रामच्छेदे}}\Afootnote{K N H; \skt{तत्र, आमच्छेदे}\ A}}\skt{ }\edtext{\skt{सिरा\-स्ना\-यु\-व्या\-पा\-दनं}}{\lemma{\skt{सिरा\appbrk{}स्ना\appbrk{}यु\appbrk{}व्या\appbrk{}पा\appbrk{}दनं}}\Afootnote{K N H; \skt{मांस\appbrk{}सि\appbrk{}रा\appbrk{}स्ना\appbrk{}यु\appbrk{}स\appbrk{}न्ध्य\appbrk{}स्थि\appbrk{}व्या\appbrk{}पा\appbrk{}द\appbrk{}नम्}\ A}}\skt{ }\edtext{\skt{शोणि\-ता\-ति\-प्र\-वृत्तिः}}{\lemma{\skt{शोणि\appbrk{}ता\appbrk{}ति\appbrk{}प्र\appbrk{}वृत्तिः}}\Afootnote{K N H; \skt{अतिमात्रं शोणि\appbrk{}ता\appbrk{}ति\appbrk{}प्र\appbrk{}वृ\appbrk{}त्तिर्}\ A}}\skt{ }\edtext{\edtext{\skt{वेद\-ना\-प्रा\-दु\-र्भा\-वो\-व\-द\-र\-ण\-म\-ने\-को\-प\-द्र\-व\-द\-र्शनं}}{\lemma{\skt{वेद\appbrk{}ना\appbrk{}प्रा\appbrk{}दु\appbrk{}र्भा\appbrk{}वो\appbrk{}व\appbrk{}द\appbrk{}र\appbrk{}णम्}}\Afootnote{K; \skt{वेद\appbrk{}ना\appbrk{}प्रा\appbrk{}दु\appbrk{}र्भा\appbrk{}वोऽ\appbrk{}व\appbrk{}द\appbrk{}र\appbrk{}णम्}\ A; \skt{वेद\appbrk{}ना\appbrk{}प्रा\appbrk{}दु\appbrk{}र्भावो वदरण}\ N; \skt{व्वेद\appbrk{}ना\appbrk{}प्रा\appbrk{}दु\appbrk{}र्भ्भा\appbrk{}वो\appbrk{}व\appbrk{}द\appbrk{}र\appbrk{}णम्}\ H}}}{\lemma{\skt{अने\appbrk{}को\appbrk{}प\appbrk{}द्र\appbrk{}व\appbrk{}द\appbrk{}र्शनं}}\Afootnote{A K H; \skt{मनेको पद्रवदर्शनं}\ N}}\skt{ }\edtext{\skt{क्षतविद्रधिर्वा}}{\lemma{\skt{क्षतविद्रधिर्वा}}\Afootnote{K; \skt{क्षतविद्रधिर्वा}\ A H; \skt{क्वनविद्रधिर्वा}\ N}}\skt{ भवति~। स }\edtext{\skt{यदा तु}}{\lemma{\skt{यदा तु}}\Afootnote{K N H; \skt{यदा}\ A}}\skt{ भयमोहाभ्यां }\edtext{\skt{पक्व\-म\-प\-क्व\-मिति}}{\lemma{\skt{अपक्वम्}}\Afootnote{K N H; \skt{अप्यपक्वम्}\ A}}\skt{ मन्यमानः चिरमुपेक्षते }\edtext{\skt{व्याधिं}}{\lemma{\skt{व्याधिं}}\Afootnote{A K H; \skt{व्याधि}\ N}}\skt{ }\edtext{\skt{वैद्यः~। स}}{\lemma{\skt{वैद्यः~। स}}\Afootnote{K N H; \skt{वैद्यस्तदा}\ A}}\skt{ }\edtext{\skt{गम्भीरानुगतो}}{\lemma{\skt{गम्भीरानुगतो}}\Afootnote{A K H; \skt{गम्भीरानुगतो}\ N}}\skt{ द्वार\-म\-ल\-भ\-मानः पूयः }\edtext{\edtext{\edtext{\skt{स्वमा\-श\-य\-म\-व\-दा\-यो\-त्स\-ङ्ग\-ङ्कृत्वा}}{\lemma{\skt{स्वमाशयम्}}\Afootnote{K N; \skt{स्वमाश्रयम्}\ A; \skt{स्वमासय}\ H}}}{\lemma{\skt{अव\appbrk{}दा\appbrk{}यो\appbrk{}त्स\appbrk{}ङ्गङ्}}\Afootnote{K; \skt{अवदार्योत्सङ्गं}\ A; \skt{अवदार्यो}\ N; \skt{मव\appbrk{}दा\appbrk{}र्यो\appbrk{}त्स\appbrk{}ङ्गङ्}\ H}}}{\lemma{\skt{कृत्वा}}\Afootnote{K H; \skt{महा\appbrk{}न्त\appbrk{}म\appbrk{}व\appbrk{}काशं कृत्वा}\ A; \skt{तसङ्गं कृत्वा}\ N}}\skt{ नाडीञ्जनयित्वा }\edtext{\edtext{\skt{भवत्यसाध्यः}}{\lemma{\skt{भवत्य्}}\Afootnote{K N H; \skt{कृच्छ्रसाध्यो भवत्य्}\ A}}}{\lemma{\skt{असाध्यः}}\Afootnote{K N H; \skt{असाध्यो वेति}\ A}}\skt{~॥}\pend

\pstart\versepar \ledsidenote{\passref{1.17.10}}\edtext{\skt{भवति}}{\lemma{\skt{भवति}}\Afootnote{A; \emph{om.}\ K H; \skt{भवन्ति}\ N}}\skt{ चात्र~॥ }\pend
\pstart\versepar \skt{}\edtext{\skt{यश्छि\-न\-त्या\-म\-म\-ज्ञा\-ना\-द्यश्च}}{\lemma{\skt{छिन\appbrk{}त्या\appbrk{}म\appbrk{}म\appbrk{}ज्ञा\appbrk{}ना\appbrk{}द्यश्च}}\Afootnote{\skt{छिन\appbrk{}त्त्या\appbrk{}म\appbrk{}म\appbrk{}ज्ञा\appbrk{}ना\appbrk{}द्यश्च}\ A; \skt{श्छि नत्या\appbrk{}म\appbrk{}म\appbrk{}ज्ञा\appbrk{}ना\appbrk{}द्यश्च}\ K; \skt{छिन\appbrk{}त्या\appbrk{}म\appbrk{}म\appbrk{}ज्ञा\appbrk{}ना\appbrk{}द्यश्च}\ N; \skt{छिन\appbrk{}त्त्या\appbrk{}म\appbrk{}म\appbrk{}ज्ञा\appbrk{}ना\appbrk{}द्यश्च}\ H}}\skt{ }\edtext{\edtext{\skt{पक्वमुपेक्षते}}{\lemma{\skt{पक्वम्}}\Afootnote{A N H; \skt{पक्व}\ K}}}{\lemma{\skt{उपेक्षते}}\Afootnote{A N H; \skt{मुपेक्षते}\ K}}\skt{~। }\pend
\pstart\versepar \skt{}\edtext{\edtext{\skt{श्वपचाविव}}{\lemma{\skt{श्वपचाव्}}\Afootnote{A K N; \skt{श्वपचा}\ H}}}{\lemma{\skt{इव}}\Afootnote{A K N; \skt{विव}\ H}}\skt{ }\edtext{\skt{सस्य\-श्चे\-त्ता\-व\-नि\-श्चि\-त\-का\-रिणौ}}{\lemma{\skt{सस्य\appbrk{}श्चे\appbrk{}त्ता\appbrk{}व\appbrk{}नि\appbrk{}श्चि\appbrk{}त\appbrk{}का\appbrk{}रिणौ}}\Afootnote{K; \skt{मन्तव्यौ ताव\appbrk{}नि\appbrk{}श्चि\appbrk{}त\appbrk{}का\appbrk{}रिणौ}\ A; \skt{सन्य\appbrk{}स्ये\appbrk{}त्ता\appbrk{}वि\appbrk{}नि\appbrk{}श्चि\appbrk{}त\appbrk{}का\appbrk{}रिणौ}\ N; \skt{सम्प\appbrk{}श्ये\appbrk{}त्ता\appbrk{}व\appbrk{}नि\appbrk{}श्चि\appbrk{}त\appbrk{}का\appbrk{}रिणौ}\ H}}\skt{~॥}\edtext{}{\lemma{\skt{चात्र~॥ यश्}}\Afootnote{A N; \skt{भवति चात्र~॥ य}\ K; \skt{भवति चात्र~॥ यश्}\ H}}\pend

\pstart\versepar \ledsidenote{\passref{1.17.11}}\edtext{\skt{प्राक्छ\-स्त्र\-क\-र्म\-ण\-श्चेष्टं}}{\lemma{\skt{प्राक्छ\appbrk{}स्त्र\appbrk{}क\appbrk{}र्म\appbrk{}णश्}}\Afootnote{A K H; \skt{प्राच्छस्त्रकर्मण}\ N}}\skt{ भोजयेदातुरं भिषक्~। }\pend
\pstart\versepar \skt{पानेयं }\edtext{\skt{पाययेत्मद्यं}}{\lemma{\skt{पाययेत्मद्यं}}\Afootnote{K H; \skt{पाययेन्मद्यं}\ A; \skt{पाययेन्मद्यंन्}\ N}}\skt{ तीक्ष्णं }\edtext{\skt{योऽवेदनासहं}}{\lemma{\skt{ऽवेदनासहं}}\Afootnote{\skt{वेदनाऽसहः}\ A; \skt{वेदनासहं}\ K H N}}\skt{~॥}\edtext{}{\lemma{\skt{भिषक्~। पानेयं}}\Afootnote{K H; \skt{भिषक्~। मद्यपं}\ A; \skt{भिषक्~। पानपं}\ N}}\pend

\pstart\versepar \ledsidenote{\passref{1.17.12}}\skt{न }\edtext{\skt{मूर्छ\-त्य\-न्न\-सं\-यो\-गा\-न्मत्तः}}{\lemma{\skt{अन्न\appbrk{}सं\appbrk{}यो\appbrk{}गा\appbrk{}न्मत्तः}}\Afootnote{A K; \skt{अन्न\appbrk{}सं\appbrk{}यो\appbrk{}गा\appbrk{}न्मत्त}\ N; \skt{अन्न\appbrk{}सं\appbrk{}यो\appbrk{}गा\appbrk{}त्मत्तः}\ H}}\skt{ शस्त्रं न बुध्यते~। }\pend
\pstart\versepar \skt{}\edtext{\skt{तस्मा\-द\-व\-श्य\-म्भो\-क्तव्यं}}{\lemma{\skt{तस्मादवश्यम्}}\Afootnote{A K H; \skt{तस्मादवश्य}\ N}}\skt{ रोगेषूक्तेषु कर्मणि~॥}\pend

\pstart\versepar \ledsidenote{\passref{1.17.13}}\skt{प्राणो }\edtext{\skt{ह्याभ्यन्तरो}}{\lemma{\skt{आभ्यन्तरो}}\Afootnote{A K H; \skt{आभ्यन्तरे}\ N}}\skt{ }\edtext{\skt{नृणा\-म्बा\-ह्य\-प्रा\-ण\-गु\-णा\-न्वितः}}{\lemma{\skt{नृणाम्}}\Afootnote{K; \skt{नॄणां}\ A; \skt{नॄणां}\ N H}}\skt{~। }\pend
\pstart\versepar \skt{धार\-य\-त्य\-वि\-रो\-धेन }\edtext{\skt{उच्छ्र\-य\-म्पा\-ञ्च\-भौ\-तिकं}}{\lemma{\skt{उच्छ्रयम्}}\Afootnote{K H; \skt{शरीरं}\ A; \skt{द्रच्छ्रायं}\ N}}\skt{~॥}\pend

\pstart\versepar \ledsidenote{\passref{1.17.14}}\edtext{\skt{यो }\edtext{\edtext{\skt{ह्युत्थितोऽल्पो}}{\lemma{\skt{उत्थितो}}\Afootnote{K H; \emph{om.}\ A; \skt{उक्तश्छि}\ N}}}{\lemma{\skt{ऽल्पो}}\Afootnote{A H; \skt{ल्पो}\ K; \skt{तो ल्पो}\ N}}}{\lemma{\skt{यो ह्य्}}\Afootnote{K N H; \emph{om.}\ A}}\skt{ }\edtext{\skt{यदि वा}}{\lemma{\skt{यदि वा}}\Afootnote{K N H; \emph{om.}\ A}}\skt{ }\edtext{\skt{महांभ्यात्क्रियां}}{\lemma{\skt{महांभ्यात्क्रियां}}\Afootnote{K; \skt{महान्वा क्रियया}\ A; \skt{महां स्यात्क्रियां}\ N; \skt{महान्स्यात्क्रियां}\ H}}\skt{ }\edtext{\skt{विना}}{\lemma{\skt{विना}}\Afootnote{K N; \skt{विना यः समुच्छ्रितः}\ A; \skt{विना पाकमु}\ H}}\skt{ }\edtext{\edtext{\skt{पाकमुपैति}}{\lemma{\skt{पाकम्}}\Afootnote{A K N; \emph{om.}\ H}}}{\lemma{\skt{उपैति}}\Afootnote{A K N; \skt{पैति}\ H}}\skt{ }\edtext{\skt{शोफः}}{\lemma{\skt{शोफः}}\Afootnote{A K N H\textsuperscript{ac}; \skt{शोफः~। ल}\ H\textsuperscript{pc}}}\skt{~। }\pend
\pstart\versepar \skt{विशालमूलो विषमं विदग्धः स }\edtext{\skt{कृच्छ्रतां}}{\lemma{\skt{कृच्छ्रतां}}\Afootnote{A K H; \skt{कृच्छ्रता}\ N}}\skt{ यात्य\-व\-गा\-ढ\-दोषः~॥}\pend

\pstart\versepar \ledsidenote{\passref{1.17.15}}\edtext{\skt{आले\-प\-वि\-स्रा\-व\-ण\-शो\-ध\-नैस्तु}}{\lemma{\skt{आले\appbrk{}प\appbrk{}वि\appbrk{}स्रा\appbrk{}व\appbrk{}ण\appbrk{}शो\appbrk{}ध\appbrk{}नैस्तु}}\Afootnote{K; \skt{आले\appbrk{}प\appbrk{}वि\appbrk{}स्रा\appbrk{}व\appbrk{}णा\appbrk{}शो\appbrk{}ध\appbrk{}नैस्तु}\ A; \skt{आले\appbrk{}प\appbrk{}वि\appbrk{}श्रा\appbrk{}व\appbrk{}ण\appbrk{}सौ\appbrk{}ध\appbrk{}नैस्तु}\ N; \skt{आले\appbrk{}प\appbrk{}वि\appbrk{}श्रा\appbrk{}व\appbrk{}न\appbrk{}शो\appbrk{}ध\appbrk{}नैस्तु}\ H}}\skt{ }\edtext{\skt{सम्य\-क्प्र\-यु\-क्तै\-र्यदि}}{\lemma{\skt{सम्यक्प्रयुक्तैर्}}\Afootnote{A K H; \skt{सम्यक्युकैर्}\ N}}\skt{ नोपशाम्येत्~। }\pend
\pstart\versepar \skt{शीघ्रं }\edtext{\edtext{\skt{विप\-च्ये\-त्स\-म\-म\-ल्प\-मूलः}}{\lemma{\skt{विपच्येत्समम्}}\Afootnote{K; \skt{शीघ्रं समम्}\ A; \skt{विपस्यत्समम्}\ N; \skt{विपच्येत्सम}\ H}}}{\lemma{\skt{अल्पमूलः}}\Afootnote{A K N; \skt{मल्पमूलः}\ H}}\skt{ }\edtext{\edtext{\skt{सम्पि\-ण्डि\-त\-श्चो\-प\-रि\-चा\-ल्प\-दोषः}}{\lemma{\skt{सम्पिण्डितश्}}\Afootnote{K; \skt{स पिण्डितश्}\ A H\textsuperscript{pc}; \skt{स पिण्डित}\ N; \skt{सन्पिण्डितश्}\ H\textsuperscript{ac}}}}{\lemma{\skt{चोप\appbrk{}रि\appbrk{}चा\appbrk{}ल्प\appbrk{}दोषः}}\Afootnote{K; \skt{चोपरि चोन्नतः स्यात्}\ A; \skt{श्चोप\appbrk{}रि\appbrk{}चा\appbrk{}ल्प\appbrk{}दोषः}\ N; \skt{चोपरि वाल्पऐदोषः}\ H\textsuperscript{ac}; \skt{चोपरि वाल्पदोषः}\ H\textsuperscript{pc}}}\edtext{}{\lemma{\skt{नोपशाम्येत्~। शीघ्रं}}\Afootnote{K N H\textsuperscript{ac}; \skt{नोपशाम्येत्~। पच्येत}\ A; \skt{नोपशाम्येत्~। हु शीघ्रं}\ H\textsuperscript{pc}}}\pend

\pstart\versepar \ledsidenote{\passref{1.17.16}}\edtext{\skt{कक्षं}}{\lemma{\skt{कक्षं}}\Afootnote{A K N; \skt{कक्षा}\ H}}\skt{ }\edtext{\skt{समासाद्य}}{\lemma{\skt{समासाद्य}}\Afootnote{A K N; \skt{न्समासाद्य}\ H}}\skt{ }\edtext{\skt{यथा च}}{\lemma{\skt{यथा च}}\Afootnote{K N H; \skt{यथैव}\ A}}\skt{ }\edtext{\edtext{\skt{वह्निर्वाय्वीरितः}}{\lemma{\skt{वह्निर्}}\Afootnote{K N H; \skt{वह्नि}\ A}}}{\lemma{\skt{वाय्वीरितः}}\Afootnote{K; \skt{र्वाय्वीरितः}\ A; \skt{बाष्पीरितः}\ N; \skt{वाचितः}\ H}}\skt{ सन्दहति }\edtext{\skt{प्रसह्य}}{\lemma{\skt{प्रसह्य}}\Afootnote{A K H; \skt{प्रसह्यः}\ N}}\skt{~। }\pend
\pstart\versepar \skt{तथैव }\edtext{\skt{पूयो\-प्य\-वि\-निः\-सृतो हि}}{\lemma{\skt{पूयो\appbrk{}प्य\appbrk{}वि\appbrk{}निः\appbrk{}सृतो हि}}\Afootnote{K; \skt{पूयोऽ\appbrk{}प्य\appbrk{}वि\appbrk{}निः\appbrk{}सृतो हि}\ A; \skt{पूयो पि विनिग्मृतो हि}\ N; \skt{पूयोऽ\appbrk{}प्य\appbrk{}वि\appbrk{}निः\appbrk{}सृतो हि}\ H}}\skt{ }\edtext{\skt{मान्सं}}{\lemma{\skt{मान्सं}}\Afootnote{A K H; \skt{मान्स}\ N}}\skt{ }\edtext{\skt{सिरास्नायु च}}{\lemma{\skt{सिरास्नायु च}}\Afootnote{K N H; \skt{सिराः स्नायु च}\ A}}\skt{ खादतीह~॥}\pend

\pstart\versepar \ledsidenote{\passref{1.17.17}}\skt{आदौ }\edtext{\skt{विम्लापनं}}{\lemma{\skt{विम्लापनं}}\Afootnote{A K H; \skt{विस्लापनं}\ N}}\skt{ }\edtext{\edtext{\skt{कुर्या\-त्द्वि\-ती\-य\-म\-व\-से\-चनं}}{\lemma{\skt{कुर्यात्द्वितीयम्}}\Afootnote{A K; \skt{कुर्या द्वितीयम्}\ N; \skt{कुर्या द्वितीय}\ H}}}{\lemma{\skt{अवसेचनं}}\Afootnote{A K N; \skt{मवसेचनं}\ H}}\skt{~। }\pend
\pstart\versepar \skt{तृती\-य\-मु\-प\-नाहं तु चतुर्थीं }\edtext{\skt{पाटनक्रियां}}{\lemma{\skt{पाटनक्रियां}}\Afootnote{A K H; \skt{पाठनक्रियां}\ N}}\skt{~।}\pend

\pstart\versepar \ledsidenote{\passref{1.17.18}}\edtext{\skt{पञ्चमं}}{\lemma{\skt{पञ्चमं}}\Afootnote{A K H; \skt{पञ्चम}\ N}}\skt{ शोधनं }\edtext{\skt{विद्यात्षष्ठं}}{\lemma{\skt{विद्यात्षष्ठं}}\Afootnote{K H; \skt{कुर्यात्षष्ठं}\ A; \skt{स्विद्यात्~। षष्ठं}\ N}}\skt{ }\edtext{\edtext{\skt{रोपणमिष्यते}}{\lemma{\skt{रोपणम्}}\Afootnote{A K; \skt{रोपण}\ N H}}}{\lemma{\skt{इष्यते}}\Afootnote{A K; \skt{मिष्यते}\ N H}}\skt{~। }\pend
\pstart\versepar \skt{एते क्रमा }\edtext{\skt{व्रण\-स्यो\-क्ता\-स्स\-प्तमं}}{\lemma{\skt{व्रणस्योक्तास्}}\Afootnote{A K H; \skt{क्रमाव्रणस्योक्ताः}\ N}}\skt{ }\edtext{\skt{वैकृ\-ता\-प\-ह\-मिति}}{\lemma{\skt{वैकृ\appbrk{}ता\appbrk{}प\appbrk{}ह\appbrk{}मिति}}\Afootnote{\skt{वैकृतापहम्}\ A N H K}}\skt{~॥}\pend

\pstart \ledsidenote{\passref{1.17.trailer}}\edtext{\skt{१७~॥}}{\lemma{\skt{१७~॥}}\Afootnote{\skt{इति सुश्रु\appbrk{}त\appbrk{}सं\appbrk{}हि\appbrk{}तायां सूत्रस्थाने आम\appbrk{}प\appbrk{}क्वै\appbrk{}ष\appbrk{}णीयो नाम सप्तदशोऽध्यायः}\ A; \skt{इति १७~॥ {\symfont ❈}}\ K; \skt{इति}\ N; \skt{इति~॥ डोञ्}\ H}}\pend

\endnumbering
\beginnumbering
\eledsection*{Adhyāya 18}
\pstart \ledsidenote{\passref{1.18.1}}\edtext{\skt{अथात}}{\lemma{\skt{अथात}}\Afootnote{K N H; \skt{अथातो}\ A}}\skt{ }\edtext{\skt{आले\-प\-व्र\-ण\-ब\-न्ध\-विधि}}{\lemma{\skt{आले\appbrk{}प\appbrk{}व्र\appbrk{}ण\appbrk{}ब\appbrk{}न्ध\appbrk{}विधि}}\Afootnote{K; \skt{व्रणा\appbrk{}ले\appbrk{}प\appbrk{}न\appbrk{}ब\appbrk{}न्ध\appbrk{}वि\appbrk{}धिम्}\ A; \skt{आत्म\appbrk{}ष\appbrk{}व्र\appbrk{}ण\appbrk{}ब\appbrk{}न्ध\appbrk{}वि\appbrk{}धिम्}\ N; \skt{आले\appbrk{}प\appbrk{}न\appbrk{}व्र\appbrk{}ण\appbrk{}ब\appbrk{}न्ध\appbrk{}वि\appbrk{}धिम्}\ H}}\skt{ }\edtext{\skt{व्याख्यास्यामः}}{\lemma{\skt{व्याख्यास्यामः}}\Afootnote{H; \skt{अध्यायं व्याख्यास्यामः}\ A N; \emph{om.}\ K}}\skt{~॥}\pend

\pstart \ledsidenote{\passref{1.18.3}}\edtext{\skt{आलेपन}}{\lemma{\skt{आलेपन}}\Afootnote{K N H; \skt{आलेप}\ A}}\skt{ आद्य उपक्रम एष सर्वशोफानां }\edtext{\skt{सामान्यतः}}{\lemma{\skt{सामान्यतः}}\Afootnote{K N H; \skt{सामान्यः}\ A}}\skt{ }\edtext{\skt{प्रधानतमश्च}}{\lemma{\skt{प्रधानतमश्च}}\Afootnote{A K H; \skt{प्रधानतमच्च}\ N}}\skt{ }\edtext{\skt{तम्प्रति}}{\lemma{\skt{तम्प्रति}}\Afootnote{K N H; \skt{तं च}\ A}}\skt{ }\edtext{\skt{प्रति\-रो\-ग\-म्व\-क्ष्यामः}}{\lemma{\skt{वक्ष्यामः}}\Afootnote{K N H; \skt{वक्ष्यामः ततो बन्धः प्रधानं तेन शुद्धि\appbrk{}र्व्र\appbrk{}ण\appbrk{}रो\appbrk{}प\appbrk{}ण\appbrk{}म\appbrk{}स्थि\appbrk{}सं\appbrk{}धि\appbrk{}स्थैर्यं च}\ A}}\skt{~।}\pend

\pstart \ledsidenote{\passref{1.18.4}}\skt{तत्र }\edtext{\edtext{\skt{प्रति\-लो\-म\-मा\-लि\-म्पे\-न्ना\-नु\-लोम}}{\lemma{\skt{प्रतिलोमम्}}\Afootnote{A K; \skt{प्रतिलोम}\ N H}}\skt{~। प्रतिलोमे हि}}{\lemma{\skt{आलि\appbrk{}म्पे\appbrk{}न्ना\appbrk{}नु\appbrk{}लोम~। प्रतिलोमे हि}}\Afootnote{K; \skt{आलिम्पेत्~। प्रतिलोमे हि}\ A; \skt{मालि\appbrk{}ष्पे\appbrk{}न्ना\appbrk{}नु\appbrk{}लोमं~॥ प्रतिलोमे हि}\ N; \skt{मालि\appbrk{}म्पे\appbrk{}न्ना\appbrk{}नु\appbrk{}लो\appbrk{}ममः~। प्रतिलोमे हि}\ H}}\skt{ }\edtext{\skt{सम्य\-गौ\-ष\-ध\-म\-व\-ति\-ष्ठते}}{\lemma{\skt{सम्य\appbrk{}गौ\appbrk{}ष\appbrk{}ध\appbrk{}म\appbrk{}व\appbrk{}ति\appbrk{}ष्ठते}}\Afootnote{K; \skt{सम्य\appbrk{}गौ\appbrk{}ष\appbrk{}ध\appbrk{}म\appbrk{}व\appbrk{}ति\appbrk{}ष्ठते}\ A N H}}\skt{~। }\edtext{\skt{प्रविशति च}}{\lemma{\skt{प्रविशति च}}\Afootnote{K N H; \skt{ऽनुप्रविशति}\ A}}\skt{ }\edtext{\edtext{\skt{रोमकूपैस्तस्य}}{\lemma{\skt{रोमकूपैस्}}\Afootnote{K N; \skt{रोम\appbrk{}कू\appbrk{}पा\appbrk{}न्स्वे\appbrk{}द\appbrk{}वा\appbrk{}हि\appbrk{}भिश्च सिरामुखैर्वीर्यं प्राप्नो}\ A; \skt{रोमकूपैस्त}\ H}}}{\lemma{\skt{तस्य}}\Afootnote{K N; \skt{ति}\ A; \skt{स्य}\ H}}\skt{ }\edtext{\skt{प्रमाणं}}{\lemma{\skt{प्रमाणं}}\Afootnote{K N H; \emph{om.}\ A}}\skt{ }\edtext{\edtext{\skt{माहि\-षा\-र्द्र\-च\-र्मा\-त्से\-ध\-मु\-प\-दि\-शन्ति}}{\lemma{\skt{माहि\appbrk{}षा\appbrk{}र्द्र\appbrk{}च\appbrk{}र्मा\appbrk{}त्से\appbrk{}धम्}}\Afootnote{K; \emph{om.}\ A; \skt{माहि\appbrk{}षा\appbrk{}र्द्र\appbrk{}च\appbrk{}र्मो\appbrk{}त्से\appbrk{}धम्}\ N; \skt{माभि\appbrk{}षा\appbrk{}र्द्र\appbrk{}च\appbrk{}र्म्मो\appbrk{}त्से\appbrk{}धम्}\ H}}}{\lemma{\skt{उपदिशन्ति}}\Afootnote{K N H; \emph{om.}\ A}}\skt{~।}\pend

\pstart \ledsidenote{\passref{1.18.5}}\skt{न च }\edtext{\edtext{\skt{शुष्क\-मु\-पे\-क्षि\-त\-व्य\-म\-न्यत्र}}{\lemma{\skt{शुष्कम्}}\Afootnote{K N H; \skt{शुष्यमाण}\ A}}}{\lemma{\skt{उपेक्षितव्यम्}}\Afootnote{K N H; \skt{मुपेक्षेत}\ A}}\skt{ }\edtext{\skt{पीडयितव्यात्~। }\edtext{\skt{शुष्कमपार्थकं}}{\lemma{\skt{अपार्थकं}}\Afootnote{K H; \skt{ह्यपार्थको}\ A; \skt{मषर्थक}\ N}}}{\lemma{\skt{पीडयितव्यात्~। शुष्कम्}}\Afootnote{K H; \skt{पीडयितव्यात्~। शुष्को}\ A; \skt{पीडयितव्यात्~। शुष्क}\ N}}\skt{ }\edtext{\skt{रुजाकरश्च}}{\lemma{\skt{रुजाकरश्च}}\Afootnote{K; \skt{रुक्करश्च}\ A; \skt{रुजाकरञ्च}\ N H}}\skt{ }\edtext{\skt{भवति}}{\lemma{\skt{भवति}}\Afootnote{K N H; \emph{om.}\ A}}\skt{~।}\pend

\pstart \ledsidenote{\passref{1.18.6}}\edtext{\edtext{\skt{आले\-प\-प्र\-दे\-ह\-यो\-र\-न्त\-र\-मा\-लेपः}}{\lemma{\skt{आले\appbrk{}प\appbrk{}प्र\appbrk{}दे\appbrk{}ह\appbrk{}योर्}}\Afootnote{K N H; \skt{स त्रिविधः प्रलेपः, प्रदेह, आलेपश्च~। प्रले\appbrk{}प\appbrk{}प्र\appbrk{}दे\appbrk{}ह\appbrk{}योर्}\ A}}}{\lemma{\skt{आलेपः}}\Afootnote{K N H; \skt{तत्र प्रलेपः}\ A}}\skt{ }\edtext{\edtext{\edtext{\skt{शीत\-स्त\-नु\-र\-वि\-शोषी}}{\lemma{\skt{शीतस्}}\Afootnote{K N; \skt{शीत}\ A H}}}{\lemma{\skt{तनुर्}}\Afootnote{K; \skt{स्तनुर्}\ A H; \skt{ततु}\ N}}}{\lemma{\skt{अविशोषी}}\Afootnote{A K; \skt{रविसोषी}\ N; \skt{अ}\ H}}\skt{ }\edtext{\skt{विशोषी वा}}{\lemma{\skt{विशोषी वा}}\Afootnote{A K H; \skt{विशोषी च}\ N}}\skt{~। }\edtext{\edtext{\skt{प्रदेहस्तूष्णः}}{\lemma{\skt{प्रदेहस्}}\Afootnote{K N; \skt{प्रदेह}\ A H}}}{\lemma{\skt{तूष्णः}}\Afootnote{K N; \skt{स्तूष्णः}\ A H}}\skt{ }\edtext{\skt{शीतो वा}}{\lemma{\skt{शीतो वा}}\Afootnote{A K H; \skt{शीतो}\ N}}\skt{ }\edtext{\skt{बह\-लो\-ब\-हु\-वि\-शोषी च}}{\lemma{\skt{बह\appbrk{}लो\appbrk{}ब\appbrk{}हु\appbrk{}वि\appbrk{}शोषी च}}\Afootnote{K; \skt{बह\appbrk{}लोऽ\appbrk{}ब\appbrk{}हु\appbrk{}र\appbrk{}वि\appbrk{}शोषी च}\ A; \skt{वह\appbrk{}लो\appbrk{}ब\appbrk{}ह्व\appbrk{}वि\appbrk{}शोषी च}\ N; \skt{बहलो}\ H}}\skt{~। }\edtext{\skt{तत्र}}{\lemma{\skt{तत्र}}\Afootnote{K; \skt{मध्यमोऽत्रालेपः~। तत्र}\ A; \skt{तत्र पित्त}\ N; \skt{बह्वविशोषी च~। तत्र}\ H}}\skt{ }\edtext{\skt{रक्त\-पि\-त्त\-प्र\-सा\-द\-कृ\-दा\-लेपः}}{\lemma{\skt{रक्त\appbrk{}पि\appbrk{}त्त\appbrk{}प्र\appbrk{}सा\appbrk{}द\appbrk{}कृ\appbrk{}दा\appbrk{}लेपः}}\Afootnote{A K; \skt{रक्त\appbrk{}प्र\appbrk{}सा\appbrk{}द\appbrk{}कृ\appbrk{}दा\appbrk{}लेपः}\ N; \skt{रक्त\appbrk{}पि\appbrk{}त्त\appbrk{}प्र\appbrk{}सा\appbrk{}द\appbrk{}कृ\appbrk{}दा\appbrk{}लेपः~।}\ H}}\skt{~। }\edtext{\skt{शोधनो}}{\lemma{\skt{शोधनो}}\Afootnote{K N H; \skt{प्रदेहो वात\appbrk{}श्ले\appbrk{}ष्म\appbrk{}प्र\appbrk{}श\appbrk{}मनः शोधनो}\ A}}\skt{ }\edtext{\skt{रोपणः}}{\lemma{\skt{रोपणः}}\Afootnote{A K N; \skt{रोपण}\ H}}\skt{ }\edtext{\skt{शोफ\-वे\-द\-ना\-प\-ग\-मश्च}}{\lemma{\skt{शोफ\appbrk{}वे\appbrk{}द\appbrk{}ना\appbrk{}प\appbrk{}ग\appbrk{}मश्च}}\Afootnote{K N H; \skt{शोफ\appbrk{}वे\appbrk{}द\appbrk{}ना\appbrk{}प\appbrk{}हश्च}\ A}}\skt{ }\edtext{\skt{तस्योपयोगः}}{\lemma{\skt{तस्योपयोगः}}\Afootnote{A K N; \skt{तस्योपभोगः}\ H}}\skt{ }\edtext{\skt{क्षताक्षतेषु}}{\lemma{\skt{क्षताक्षतेषु}}\Afootnote{A K H; \skt{क्षताक्षताषु}\ N}}\skt{ }\edtext{\skt{यस्तु}}{\lemma{\skt{यस्तु}}\Afootnote{K; \skt{यस्तु}\ A N H}}\skt{ क्षते\-षू\-प\-यु\-ज्यते स भूयः कल्क इति }\edtext{\skt{संज्ञां}}{\lemma{\skt{संज्ञां}}\Afootnote{A K H; \skt{संज्ञा}\ N}}\skt{ लभते~। }\edtext{\skt{निरु\-ध्या\-ले\-प\-न\-संज्ञः}}{\lemma{\skt{निरु\appbrk{}ध्या\appbrk{}ले\appbrk{}प\appbrk{}न\appbrk{}संज्ञः}}\Afootnote{K; \skt{निरु\appbrk{}द्धा\appbrk{}ले\appbrk{}प\appbrk{}न\appbrk{}सञ्ज्ञः}\ A; \skt{निरु\appbrk{}ध्या\appbrk{}ले\appbrk{}प\appbrk{}न\appbrk{}सं\appbrk{}ज्ञान्}\ N; \skt{निरु\appbrk{}ध्या\appbrk{}ले\appbrk{}प\appbrk{}न\appbrk{}संज्ञा}\ H}}\skt{ }\edtext{\skt{तेना\-स्रा\-व\-म\-नि\-रोधो}}{\lemma{\skt{तेना\appbrk{}स्रा\appbrk{}व\appbrk{}म\appbrk{}नि\appbrk{}रोधो}}\Afootnote{K; \skt{तेना\appbrk{}स्रा\appbrk{}व\appbrk{}स\appbrk{}न्नि\appbrk{}रोधो}\ A N; \skt{तेना\appbrk{}श्रा\appbrk{}व\appbrk{}स\appbrk{}न्नि\appbrk{}रोधौ}\ H}}\skt{ }\edtext{\edtext{\skt{मृदु\-पू\-ति\-मां\-सा\-प\-क\-र्ष\-ण\-म\-न्त\-र्दो\-षता}}{\lemma{\skt{मृदु\appbrk{}पू\appbrk{}ति\appbrk{}मां\appbrk{}सा\appbrk{}प\appbrk{}क\appbrk{}र्ष\appbrk{}णम्}}\Afootnote{K N H; \skt{मृदुता पूति\appbrk{}मां\appbrk{}सा\appbrk{}प\appbrk{}क\appbrk{}र्षण}\ A}}}{\lemma{\skt{अन्तर्दोषता}}\Afootnote{K H; \skt{मनन्तर्दोषता}\ A; \skt{अनन्तर्दोषता}\ N}}\skt{ }\edtext{\skt{शुर्व्रणशुद्धिश्च}}{\lemma{\skt{शुर्व्रणशुद्धिश्च}}\Afootnote{K; \skt{व्रणशुद्धिश्च}\ A; \skt{शुव्रणशुद्धिश्च}\ N H\textsuperscript{pc}; \skt{शुर्व्रणशुद्धिश्च}\ H\textsuperscript{ac}}}\skt{ भवति~।}\pend

\pstart \ledsidenote{\passref{1.18.12}}\edtext{\skt{न चालेपनं}}{\lemma{\skt{न चालेपनं}}\Afootnote{K H; \skt{न चालेपं}\ A; \skt{न चालेपन}\ N}}\skt{ }\edtext{\skt{रात्रौ}}{\lemma{\skt{रात्रौ}}\Afootnote{A K H; \skt{शंत्रौ}\ N}}\skt{ }\edtext{\skt{प्रयुञ्जीत}}{\lemma{\skt{प्रयुञ्जीत}}\Afootnote{K N H; \skt{प्रयुञ्जी}\ A}}\skt{ }\edtext{\skt{शैत्यात्तु श्लेष्मणः}}{\lemma{\skt{शैत्यात्तु श्लेष्मणः}}\Afootnote{K; \emph{om.}\ A; \skt{शैत्यां श्लेष्मणः}\ N; \skt{शैत्यानुश्लेष्मणः}\ H}}\skt{ }\edtext{\skt{ऊर्ध्व\-वि\-वृ\-त\-रो\-म\-कू\-प\-त्वा\-दू\-ष्मा\-नि\-रेति}}{\lemma{\skt{ऊर्ध्व\appbrk{}वि\appbrk{}वृ\appbrk{}त\appbrk{}रो\appbrk{}म\appbrk{}कू\appbrk{}प\appbrk{}त्वा\appbrk{}दू\appbrk{}ष्मा\appbrk{}नि\appbrk{}रेति}}\Afootnote{K; \skt{त , मा भूच्छै\appbrk{}त्य\appbrk{}वि\appbrk{}ह\appbrk{}तो\appbrk{}ष्म\appbrk{}ण\appbrk{}स्त\appbrk{}द\appbrk{}नि\appbrk{}र्ग\appbrk{}मा\appbrk{}द्वि\appbrk{}का\appbrk{}र\appbrk{}प्र\appbrk{}वृ\appbrk{}त्ति\appbrk{}रिति}\ A; \skt{ऊर्ध्व\appbrk{}वि\appbrk{}वृ\appbrk{}त\appbrk{}रो\appbrk{}म\appbrk{}कू\appbrk{}प\appbrk{}त्वा\appbrk{}दू\appbrk{}ष्मा\appbrk{}नि\appbrk{}रेति}\ N; \skt{ऊर्द्ध्व\appbrk{}वि\appbrk{}वृ\appbrk{}त\appbrk{}रा\appbrk{}म\appbrk{}कू\appbrk{}प\appbrk{}त्वा\appbrk{}दू\appbrk{}ष्मा\appbrk{}न\appbrk{}नि\appbrk{}रेति}\ H}}\skt{~॥}\pend

\pstart\versepar \ledsidenote{\passref{1.18.7}}\skt{अविदग्धेषु शोफेषु }\edtext{\edtext{\skt{हित\-मा\-ले\-प\-म\-न\-म्भ\-वेत्}}{\lemma{\skt{हितम्}}\Afootnote{K N H; \skt{हित}\ A}}}{\lemma{\skt{आलेपमनम्}}\Afootnote{K; \skt{मालेपनं}\ A; \skt{आलेपनं}\ N H}}\skt{~। }\pend
\pstart\versepar \skt{यथा }\edtext{\skt{स्वदोषशमनं}}{\lemma{\skt{स्वदोषशमनं}}\Afootnote{K; \skt{स्वं दोषशमनं}\ A; \skt{च दोशमनन्}\ N; \skt{स्वं दोषशमनं…}\ H}}\skt{ दाह\-क\-ण्डू\-रु\-जा\-पहं~॥}\edtext{}{\lemma{\skt{भवेत्~। यथा}}\Afootnote{A K H; \skt{भवेत्~। यथाश्}\ N}}\pend

\pstart\versepar \ledsidenote{\passref{1.18.9}}\skt{मर्मदेशेषु ये रोगा गुह्येष्वपि सदा नृणां~। }\pend
\pstart\versepar \skt{संशोधनाय तेषान्तु कुर्यादालेपनं भिषक्~॥}\pend

\pstart \ledsidenote{\passref{1.18.16}}\skt{अत }\edtext{\skt{ऊर्ध्व\-व्र\-ण\-ब\-न्ध\-न\-द्र\-व्या\-ण्यु\-प\-दे\-क्ष्यामः}}{\lemma{\skt{ऊर्ध्व\appbrk{}व्र\appbrk{}ण\appbrk{}ब\appbrk{}न्ध\appbrk{}न\appbrk{}द्र\appbrk{}व्याण्य्}}\Afootnote{K; \skt{ऊर्ध्वं व्रण\appbrk{}ब\appbrk{}न्ध\appbrk{}न\appbrk{}द्र\appbrk{}व्याण्य्}\ A H; \skt{ऊर्ध्वं व्रण\appbrk{}ब\appbrk{}न्ध\appbrk{}न\appbrk{}द्र\appbrk{}व्यान्य्}\ N}}\skt{~। }\edtext{\edtext{\skt{क्षौम\-का\-र्पा\-सि\-का\-वि\-क\-दु\-कू\-ल\-कौ\-शे\-य\-प\-त्रो\-र्ण्ण\-ची\-न\-प\-ट्ट\-च\-र्मा\-न्त\-र्ब\-ल्क\-ल\-ल\-ता\-वि\-द\-ल\-र\-ज्जा\-बा\-ल\-तू\-ल\-स\-न्ता\-नि\-का\-ला\-हा\-दी\-न्य\-था\-व्याधिं}}{\lemma{\skt{क्षौम\appbrk{}का\appbrk{}र्पा\appbrk{}सि\appbrk{}का\appbrk{}वि\appbrk{}क\appbrk{}दु\appbrk{}कू\appbrk{}ल\appbrk{}कौ\appbrk{}शे\appbrk{}य\appbrk{}प\appbrk{}त्रो\appbrk{}र्ण्ण\appbrk{}ची\appbrk{}न\appbrk{}प\appbrk{}ट्ट\appbrk{}च\appbrk{}र्मा\appbrk{}न्त\appbrk{}र्ब\appbrk{}ल्क\appbrk{}ल\appbrk{}ल\appbrk{}ता\appbrk{}वि\appbrk{}द\appbrk{}ल\appbrk{}र\appbrk{}ज्जा\appbrk{}बा\appbrk{}ल\appbrk{}तू\appbrk{}ल\appbrk{}स\appbrk{}न्ता\appbrk{}नि\appbrk{}का\appbrk{}ला\appbrk{}हा\appbrk{}दीन्य्}}\Afootnote{K; \skt{तद्यथा- क्षौम\appbrk{}का\appbrk{}र्पा\appbrk{}सा\appbrk{}वि\appbrk{}क\appbrk{}दु\appbrk{}कू\appbrk{}ल\appbrk{}कौ\appbrk{}शे\appbrk{}य\appbrk{}प\appbrk{}त्रो\appbrk{}र्ण\appbrk{}ची\appbrk{}न\appbrk{}प\appbrk{}ट्ट\appbrk{}च\appbrk{}र्मा\appbrk{}न्त\appbrk{}र्व\appbrk{}ल्क\appbrk{}ला\appbrk{}ला\appbrk{}बू\appbrk{}श\appbrk{}क\appbrk{}ल\appbrk{}ल\appbrk{}ता\appbrk{}वि\appbrk{}द\appbrk{}ल\appbrk{}र\appbrk{}ज्जु\appbrk{}तू\appbrk{}ल\appbrk{}फ\appbrk{}ल\appbrk{}स\appbrk{}न्ता\appbrk{}नि\appbrk{}का\appbrk{}लौ\appbrk{}हा\appbrk{}नीति}\ A; \skt{क्षौम\appbrk{}का\appbrk{}र्पा\appbrk{}सि\appbrk{}का\appbrk{}वि\appbrk{}क\appbrk{}दु\appbrk{}कू\appbrk{}ल\appbrk{}कौ\appbrk{}से\appbrk{}य\appbrk{}प\appbrk{}त्तौ\appbrk{}र्ण्णा\appbrk{}ची\appbrk{}न\appbrk{}प\appbrk{}ट्ट\appbrk{}च\appbrk{}र्मा\appbrk{}न्त\appbrk{}र्व\appbrk{}ल्क\appbrk{}ल\appbrk{}ल\appbrk{}ता\appbrk{}दि\appbrk{}द\appbrk{}ल\appbrk{}र\appbrk{}ज्जु\appbrk{}वा\appbrk{}ल\appbrk{}तू\appbrk{}ल\appbrk{}स\appbrk{}न्ता\appbrk{}नि\appbrk{}का\appbrk{}लो\appbrk{}हादी}\ N; \skt{क्षौम\appbrk{}का\appbrk{}र्प्पा\appbrk{}सि\appbrk{}का\appbrk{}वि\appbrk{}क\appbrk{}दु\appbrk{}कू\appbrk{}ल\appbrk{}कौ\appbrk{}शे\appbrk{}य\appbrk{}प\appbrk{}त्तौ\appbrk{}र्ण्णा\appbrk{}ची\appbrk{}न\appbrk{}प\appbrk{}ट्ट\appbrk{}च\appbrk{}र्म्मा\appbrk{}न्त\appbrk{}र्व्व\appbrk{}ल्क\appbrk{}ल\appbrk{}ल\appbrk{}ता\appbrk{}वि\appbrk{}द\appbrk{}ल\appbrk{}रज्जु वा लगू\appbrk{}ल\appbrk{}स\appbrk{}न्ता\appbrk{}नि\appbrk{}का\appbrk{}लो\appbrk{}हा\appbrk{}दीन्य्}\ H}}}{\lemma{\skt{अथाव्याधिं}}\Afootnote{K; \skt{तेषां व्याधिं}\ A; \skt{न्येषं व्याधि}\ N; \skt{एषां व्याधिङ्}\ H}}\skt{ }\edtext{\skt{काल\-श्चा\-वे\-क्ष्यो\-प\-योगः}}{\lemma{\skt{कालश्}}\Afootnote{K; \skt{कालं}\ A N H}}\skt{~। }\edtext{\edtext{\edtext{\skt{प्रमा\-ण\-त\-श्चे\-षा\-मा\-देशः}}{\lemma{\skt{प्रमाणतश्}}\Afootnote{K H; \skt{प्रकरणतश्}\ A; \skt{प्रमाणत}\ N}}}{\lemma{\skt{चेषाम्}}\Afootnote{K H; \skt{चैषाम्}\ A; \skt{श्चेषाम्}\ N}}}{\lemma{\skt{आदेशः}}\Afootnote{A K; \skt{आदेशः~॥ कोश\appbrk{}दा\appbrk{}म\appbrk{}स्च\appbrk{}स्ति\appbrk{}का\appbrk{}तु\appbrk{}वे\appbrk{}ल्लितं}\ N; \skt{आदेशः…}\ H}}\skt{~।}\pend

\pstart \ledsidenote{\passref{1.18.17}}\edtext{\skt{कोशदाम शाखासु}}{\lemma{\skt{कोशदाम शाखासु}}\Afootnote{K; \skt{तत्र कोश\appbrk{}दा\appbrk{}म\appbrk{}स्व\appbrk{}स्ति\appbrk{}का\appbrk{}नु\appbrk{}वे\appbrk{}ल्लित}\ A; \skt{शाखासु}\ N; \skt{कोश\appbrk{}दा\appbrk{}म\appbrk{}स्व\appbrk{}स्ति\appbrk{}का\appbrk{}नु\appbrk{}वे\appbrk{}ल्लितं शाखासु}\ H\textsuperscript{ac}; \skt{ल भ्र कोश\appbrk{}दा\appbrk{}म\appbrk{}स्व\appbrk{}स्ति\appbrk{}का\appbrk{}नु\appbrk{}वे\appbrk{}ल्लितं शाखासु}\ H\textsuperscript{pc}}}\skt{ }\edtext{\skt{ग्रीवा\-मे\-ढ्र\-मू\-तो\-ली\-म\-ण्ड\-ल\-स्थ\-वि\-का\-य\-म\-ल\-क\-ख\-ट्वा\-ची\-न\-वि\-वन्ध वितानगोफणाः}}{\lemma{\skt{ग्रीवा\appbrk{}मे\appbrk{}ढ्र\appbrk{}मू\appbrk{}तो\appbrk{}ली\appbrk{}म\appbrk{}ण्ड\appbrk{}ल\appbrk{}स्थ\appbrk{}वि\appbrk{}का\appbrk{}य\appbrk{}म\appbrk{}ल\appbrk{}क\appbrk{}ख\appbrk{}ट्वा\appbrk{}ची\appbrk{}न\appbrk{}वि\appbrk{}वन्ध वितानगोफणाः}}\Afootnote{K; \skt{मु(प्र)\appbrk{}तो\appbrk{}ली\appbrk{}म\appbrk{}ण्ड\appbrk{}ल\appbrk{}स्थ\appbrk{}गि\appbrk{}का\appbrk{}य\appbrk{}म\appbrk{}क\appbrk{}ख\appbrk{}ट्वा\appbrk{}ची\appbrk{}न\appbrk{}वि\appbrk{}ब\appbrk{}न्ध\appbrk{}वि\appbrk{}ता\appbrk{}न\appbrk{}गो\appbrk{}फणाः}\ A; \skt{ग्रीवा\appbrk{}मे\appbrk{}ढ्र\appbrk{}मू\appbrk{}तो\appbrk{}ली\appbrk{}म\appbrk{}ण्ड\appbrk{}ल\appbrk{}स्थ\appbrk{}वि\appbrk{}का\appbrk{}य\appbrk{}म\appbrk{}ल\appbrk{}क\appbrk{}ख\appbrk{}ट्वा\appbrk{}ची\appbrk{}न\appbrk{}वि\appbrk{}ब\appbrk{}न्ध\appbrk{}वि\appbrk{}ता\appbrk{}न\appbrk{}गो\appbrk{}ष्फणाः}\ N; \skt{ग्रीवामेढ्र मूतो\appbrk{}ली\appbrk{}म\appbrk{}ण्ड\appbrk{}ल\appbrk{}म्थ\appbrk{}वि\appbrk{}का\appbrk{}य\appbrk{}म\appbrk{}ल\appbrk{}क\appbrk{}ख\appbrk{}ट्वा\appbrk{}ची\appbrk{}न\appbrk{}वि\appbrk{}ब\appbrk{}न्ध\appbrk{}वि\appbrk{}ता\appbrk{}न\appbrk{}गो\appbrk{}स्फणाः}\ H\textsuperscript{ac}; \skt{ग्रीवामेढ्र मूतो\appbrk{}ली\appbrk{}म\appbrk{}ण्ड\appbrk{}ल\appbrk{}स्थ\appbrk{}वि\appbrk{}का\appbrk{}य\appbrk{}म\appbrk{}ल\appbrk{}क\appbrk{}ख\appbrk{}ट्वा\appbrk{}ची\appbrk{}न\appbrk{}वि\appbrk{}ब\appbrk{}न्ध\appbrk{}वि\appbrk{}ता\appbrk{}न\appbrk{}गो\appbrk{}स्फणाः}\ H\textsuperscript{pc}}}\skt{ पञ्चाङ्गी चेति चतु\-र्द\-श\-ब\-न्ध\-वि\-शेषाः~। तेषा\-न्ना\-म\-भि\-रे\-वा\-कृ\-तयः प्रायेण व्याख्याताः~।}\pend

\pstart \ledsidenote{\passref{1.18.18}}\skt{तत्र }\edtext{\edtext{\skt{कोश\-ञ्ज\-ङ्घा\-ङ्गु\-लि\-प\-र्व्वसु}}{\lemma{\skt{कोशञ्}}\Afootnote{K N H; \skt{कोश}\ A}}}{\lemma{\skt{जङ्घा\appbrk{}ङ्गु\appbrk{}लि\appbrk{}प\appbrk{}र्व्वसु}}\Afootnote{K N H; \skt{मङ्गु\appbrk{}ष्ठा\appbrk{}ङ्गु\appbrk{}लि\appbrk{}प\appbrk{}र्वसु}\ A}}\skt{ }\edtext{\skt{विदध्यात्~। }\edtext{\skt{दाम\-म\-सं\-बा\-धेऽङ्गे}}{\lemma{\skt{असंबाधेऽङ्गे}}\Afootnote{\skt{सम्बाधेऽङ्गे}\ A; \skt{असम्}\ K; \skt{मसन्धोर्धेङ्गे}\ N; \skt{मसंबाधेङ्गे}\ H}}}{\lemma{\skt{विदध्यात्~। दामम्}}\Afootnote{K; \skt{विदध्यात्दाम}\ A N H}}\skt{ }\edtext{\skt{सान्धि\-कू\-र्च्च\-क\-स्त\-ना\-न्त\-र\-क\-र्ण्णेषु}}{\lemma{\skt{सान्धि\appbrk{}कू\appbrk{}र्च्च\appbrk{}क\appbrk{}स्त\appbrk{}ना\appbrk{}न्त\appbrk{}र\appbrk{}क\appbrk{}र्ण्णेषु}}\Afootnote{H; \skt{सन्धि\appbrk{}कू\appbrk{}र्च\appbrk{}क\appbrk{}भ्रू\appbrk{}स्त\appbrk{}ना\appbrk{}न्त\appbrk{}र\appbrk{}त\appbrk{}ल\appbrk{}क\appbrk{}र्णेषु}\ A; \emph{om.}\ K; \skt{सन्धि\appbrk{}कू\appbrk{}र्च\appbrk{}क\appbrk{}स्त\appbrk{}ना\appbrk{}न्त\appbrk{}र\appbrk{}क\appbrk{}र्ण्णेषु}\ N}}\skt{ }\edtext{\edtext{\skt{स्वस्ति\-क\-म\-नु\-वे\-ल्लितं}}{\lemma{\skt{स्वस्तिकम्}}\Afootnote{A N; \emph{om.}\ K; \skt{स्वस्तिक}\ H}}}{\lemma{\skt{अनुवेल्लितं}}\Afootnote{A; \emph{om.}\ K; \skt{अनुवेलितं}\ N; \skt{म वेल्लितं}\ H\textsuperscript{ac}; \skt{म नुवेल्लितं}\ H\textsuperscript{pc}}}\skt{ }\edtext{\skt{शाखासु}}{\lemma{\skt{शाखासु}}\Afootnote{A N H; \emph{om.}\ K}}\skt{ }\edtext{\skt{ग्रीवा\-मे\-ढ्र\-यो\-न्मू\-तोलीं}}{\lemma{\skt{ग्रीवा\appbrk{}मे\appbrk{}ढ्र\appbrk{}यो\appbrk{}न्मू\appbrk{}तोलीं}}\Afootnote{H; \skt{ग्रीवामेढ्रयोः मु(प्र)तोलीं}\ A; \emph{om.}\ K; \skt{ग्रीवा\appbrk{}मे\appbrk{}ढ्रा\appbrk{}यो\appbrk{}र्मू\appbrk{}तोली}\ N}}\skt{ }\edtext{\skt{वृत्तेङ्गे}}{\lemma{\skt{वृत्तेङ्गे}}\Afootnote{N H; \skt{वृत्तेऽङ्गे}\ A; \skt{वा}\ K}}\skt{ }\edtext{\edtext{\skt{मण्ड\-ल\-म\-ङ्गु\-ष्ठा\-ङ्गु\-लि\-मे\-ढ्रा\-ग्रेषु}}{\lemma{\skt{मण्डलम्}}\Afootnote{A N H; \emph{om.}\ K}}}{\lemma{\skt{अङ्गु\appbrk{}ष्ठा\appbrk{}ङ्गु\appbrk{}लि\appbrk{}मे\appbrk{}ढ्रा\appbrk{}ग्रेषु}}\Afootnote{A H; \emph{om.}\ K; \skt{अष्टा\appbrk{}ङ्गु\appbrk{}लि\appbrk{}मे\appbrk{}ढ्रा\appbrk{}ग्रेषु}\ N}}\skt{ }\edtext{\skt{स्थविका}}{\lemma{\skt{स्थविका}}\Afootnote{H; \skt{स्थगिकां}\ A; \skt{धेङ्}\ K; \skt{सुस्थविका}\ N}}\skt{~। }\edtext{\skt{यमलव्रणाभ्यां}}{\lemma{\skt{यमलव्रणाभ्यां}}\Afootnote{N H; \skt{यम\appbrk{}ल\appbrk{}व्र\appbrk{}ण\appbrk{}योर्}\ A; \emph{om.}\ K}}\skt{ }\edtext{\skt{यमलकं}}{\lemma{\skt{यमलकं}}\Afootnote{H; \skt{यमकं}\ A; \emph{om.}\ K; \skt{यमलका}\ N}}\skt{~। }\edtext{\skt{हनु\-ग…\-ण्ड\-शं\-खेषु }\edtext{\edtext{\skt{खट्वा\-म\-पा\-ङ्ग\-यो\-श्चीनं}}{\lemma{\skt{खट्वाम्}}\Afootnote{A N H; \emph{om.}\ K}}}{\lemma{\skt{अपाङ्गयोश्}}\Afootnote{A N H; \emph{om.}\ K}}}{\lemma{\skt{हनुग{\symfont ‡}ण्डशंखेषु}}\Afootnote{H; \skt{हनु\appbrk{}श\appbrk{}ङ्ख\appbrk{}ग\appbrk{}ण्डेषु}\ A; \emph{om.}\ K; \skt{हव\appbrk{}ग\appbrk{}ण्ड\appbrk{}श\appbrk{}ङ्खेषु}\ N}}\skt{~। }\edtext{\skt{पृष्ठोदरोरःसु विबन्धं}}{\lemma{\skt{पृष्ठोदरोरःसु}}\Afootnote{A N H; \emph{om.}\ K}}\skt{ }\edtext{\skt{मूर्द्ध्नि वितामं}}{\lemma{\skt{मूर्द्ध्नि}}\Afootnote{H; \skt{मूर्ध}\ A; \emph{om.}\ K; \skt{मू}\ N}}\skt{~। }\edtext{\skt{चिपु\-क\-ना\-मौ\-ष्ठां\-स\-ब\-स्तिषु गोष्फणा}}{\lemma{\skt{चिपु\appbrk{}क\appbrk{}ना\appbrk{}मौ\appbrk{}ष्ठां\appbrk{}स\appbrk{}ब\appbrk{}स्तिषु}}\Afootnote{H; \skt{चिबु\appbrk{}क\appbrk{}ना\appbrk{}सौ\appbrk{}ष्ठां\appbrk{}स\appbrk{}ब\appbrk{}स्तिषु}\ A; \emph{om.}\ K; \skt{चिपु\appbrk{}क\appbrk{}ना\appbrk{}सौ\appbrk{}ष्ठां\appbrk{}स\appbrk{}ब\appbrk{}स्तिषु}\ N}}\skt{ }\edtext{\edtext{\skt{जत्रु\-णि\-प\-ञ्चा\-ङ्गी\-मिति}}{\lemma{\skt{जत्रु\appbrk{}णि\appbrk{}प\appbrk{}ञ्चा\appbrk{}ङ्गीम्}}\Afootnote{N H; \skt{जत्रुण ऊर्ध्वं पञ्चाङ्गीम्}\ A; \skt{गे स}\ K}}\skt{~। यो वा यस्मिन्प्रदेशे}}{\lemma{\skt{इति~। यो वा}}\Afootnote{A N H; \emph{om.}\ K}}\skt{ }\edtext{\skt{सुनिविष्टो भवति}}{\lemma{\skt{सुनिविष्टो}}\Afootnote{A H; \emph{om.}\ K; \skt{विनिविष्टो}\ N}}\skt{ }\edtext{\edtext{\skt{तत्त\-स्मि\-न्नि\-द\-ध्यात्}}{\lemma{\skt{तत्तस्मिन्}}\Afootnote{N H; \skt{तं तस्मि}\ A; \emph{om.}\ K}}\skt{~।}}{\lemma{\skt{निदध्यात्}}\Afootnote{N H; \skt{न्विदध्यात्}\ A; \emph{om.}\ K}}\edtext{}{\lemma{\skt{चीनं}}\Afootnote{A N H; \emph{om.}\ K}}\edtext{}{\lemma{\skt{विबन्धं}}\Afootnote{A H; \emph{om.}\ K; \skt{विबन्ध}\ N}}\edtext{}{\lemma{\skt{वितामं}}\Afootnote{H; \skt{नि वितानं}\ A; \emph{om.}\ K; \skt{ध्निवितानं}\ N}}\edtext{}{\lemma{\skt{गोष्फणा}}\Afootnote{H; \skt{गोफणां}\ A; \emph{om.}\ K; \skt{शोष्फण}\ N}}\edtext{}{\lemma{\skt{यस्मिन्प्रदेशे}}\Afootnote{H; \skt{यस्मि\appbrk{}न्श\appbrk{}री\appbrk{}र\appbrk{}प्र\appbrk{}देशे}\ A; \emph{om.}\ K; \skt{यस्मिन्प्रदेसेष}\ N}}\edtext{}{\lemma{\skt{भवति}}\Afootnote{A N H; \emph{om.}\ K}}\pend

\pstart \ledsidenote{\passref{1.18.19}}\edtext{\edtext{\skt{यन्त्र\-न\-मू\-र्द्ध\-म\-ध\-स्ति\-र्यक्च}}{\lemma{\skt{यन्त्रनम्}}\Afootnote{H; \skt{यन्त्रणम्}\ A; \skt{यन्त्रण}\ N}}}{\lemma{\skt{ऊर्द्ध\appbrk{}म\appbrk{}ध\appbrk{}स्ति\appbrk{}र्यक्च}}\Afootnote{H; \skt{ऊर्ध्व\appbrk{}म\appbrk{}ध\appbrk{}स्ति\appbrk{}र्यक्च}\ A; \skt{मूर्ध\appbrk{}म\appbrk{}ध\appbrk{}स्ति\appbrk{}र्यक्च}\ N}}\skt{ }\edtext{\skt{भवति}}{\lemma{\skt{भवति}}\Afootnote{N H; \emph{om.}\ A}}\skt{~।}\pend

\pstart \ledsidenote{\passref{1.18.20}}\skt{तत्र }\edtext{\skt{घना\-ङ्का\-व\-लि\-का\-न्दत्वा}}{\lemma{\skt{कावलिकान्}}\Afootnote{H; \skt{कवलिकां}\ A N}}\skt{ }\edtext{\edtext{\skt{चाम\-प\-रि\-क्षे\-प\-मृ\-जु\-म\-ना\-वि\-द्ध\-म\-स\-ङ्कु\-चितं}}{\lemma{\skt{चाम\appbrk{}प\appbrk{}रि\appbrk{}क्षे\appbrk{}पम्}}\Afootnote{H; \skt{वाम\appbrk{}ह\appbrk{}स्त\appbrk{}प\appbrk{}रि\appbrk{}क्षे\appbrk{}पम्}\ A; \skt{वामपरिक्षेप}\ N}}}{\lemma{\skt{ऋजुम्}}\Afootnote{A H; \skt{मृजुम्}\ N}}\skt{ मृदु पट्टन्निवेश्य बध्नीयात्~। न च व्रणस्योपरि कुर्यात्~। }\edtext{\edtext{\skt{ग्रन्थि\-मा\-बा\-ध\-ङ्क\-रोति}}{\lemma{\skt{आबाधङ्}}\Afootnote{N H; \skt{आबाध}\ A}}}{\lemma{\skt{करोति}}\Afootnote{N H; \skt{करं च}\ A}}\skt{~।}\pend

\pstart \ledsidenote{\passref{1.18.21}}\edtext{\skt{अथास्य न च}}{\lemma{\skt{अथास्य न च}}\Afootnote{N H; \skt{न च}\ A}}\skt{ }\edtext{\skt{विके\-शि\-कौ\-ष\-धे\-ति\-स्नि\-ग्धे\-ति\-रूक्षे}}{\lemma{\skt{विके\appbrk{}शि\appbrk{}कौ\appbrk{}ष\appbrk{}धे\appbrk{}ति\appbrk{}स्नि\appbrk{}ग्धे\appbrk{}ति\appbrk{}रूक्षे}}\Afootnote{H; \skt{विकेशिकौषधे अतिस्निग्धे अतिरूक्षे}\ A; \skt{विके\appbrk{}शि\appbrk{}कौ\appbrk{}ष\appbrk{}धे\appbrk{}ति\appbrk{}रूक्षे}\ N}}\skt{ विषमे वा कुर्व्वीत~। }\edtext{\skt{कस्मा\-द\-ति\-स्ने\-हा\-त्क्ले\-द\-यति}}{\lemma{\skt{कस्मा\appbrk{}द\appbrk{}ति\appbrk{}स्ने\appbrk{}हा\appbrk{}त्क्ले\appbrk{}द\appbrk{}यति}}\Afootnote{N H; \skt{यस्मा\appbrk{}द\appbrk{}ति\appbrk{}स्ने\appbrk{}हा\appbrk{}त्क्लेदो}\ A}}\skt{~। }\edtext{\skt{रौक्ष्याच्छिनत्ति}}{\lemma{\skt{रौक्ष्याच्छिनत्ति}}\Afootnote{H; \skt{रौक्ष्याच्छेदो}\ A; \skt{रौक्ष्यच्छिंनत्ति}\ N}}\skt{ }\edtext{\skt{दुर्न्या\-स्त\-व्र\-ण\-च\-र्म्मा\-व\-घ\-र्षणं}}{\lemma{\skt{दुर्न्या\appbrk{}स्त\appbrk{}व्र\appbrk{}ण\appbrk{}च\appbrk{}र्म्मा\appbrk{}व\appbrk{}घ\appbrk{}र्षणं}}\Afootnote{H; \skt{दुर्न्या\appbrk{}सा\appbrk{}द्व्र\appbrk{}ण\appbrk{}व\appbrk{}र्त्मा\appbrk{}व\appbrk{}घ\appbrk{}र्ष\appbrk{}णम्}\ A; \skt{दुन्य\appbrk{}स्च\appbrk{}व्र\appbrk{}ण\appbrk{}व\appbrk{}र्मा\appbrk{}व\appbrk{}घ\appbrk{}र्ष\appbrk{}णङ्}\ N}}\skt{ }\edtext{\skt{करोति}}{\lemma{\skt{करोति}}\Afootnote{N H; \emph{om.}\ A}}\skt{~। }\edtext{\skt{युक्त}}{\lemma{\skt{युक्त}}\Afootnote{N H; \emph{om.}\ A}}\skt{ }\edtext{\skt{स्नेहत्यादाशु}}{\lemma{\skt{स्नेहत्यादाशु}}\Afootnote{H; \emph{om.}\ A; \skt{स्नेहत्वादाशु}\ N}}\skt{ }\edtext{\skt{रोहति}}{\lemma{\skt{रोहति}}\Afootnote{N H; \skt{इति}\ A}}\skt{~।}\pend

\pstart \ledsidenote{\passref{1.18.22}}\skt{तत्र }\edtext{\edtext{\skt{व्रणा\-य\-त\-न\-वि\-शे\-षा\-द्ब\-न्ध\-स्त्रि\-विधो}}{\lemma{\skt{व्रणा\appbrk{}य\appbrk{}त\appbrk{}न\appbrk{}वि\appbrk{}शे\appbrk{}षा\appbrk{}द्ब\appbrk{}न्धस्}}\Afootnote{\skt{व्रणा\appbrk{}य\appbrk{}त\appbrk{}न\appbrk{}वि\appbrk{}शे\appbrk{}षा\appbrk{}द्ब\appbrk{}न्ध\appbrk{}वि\appbrk{}शे\appbrk{}षस्}\ A; \skt{व्रणा\appbrk{}य\appbrk{}त\appbrk{}न\appbrk{}वि\appbrk{}से\appbrk{}षा\appbrk{}द्बन्ध}\ N; \skt{व्रणा\appbrk{}य\appbrk{}त\appbrk{}न\appbrk{}वि\appbrk{}शे\appbrk{}षा\appbrk{}द्बन्ध}\ H}}}{\lemma{\skt{त्रिविधो}}\Afootnote{A; \skt{स्त्रिविधो}\ N H}}\skt{ भवति~। गाढस्समः शिथिल इति~।}\pend

\pstart \ledsidenote{\passref{1.18.24}}\skt{तत्र }\edtext{\skt{स्फिक्कु\-क्षि\-वं\-क्ष\-ण\-शिरः}}{\lemma{\skt{स्फिक्कु\appbrk{}क्षि\appbrk{}वं\appbrk{}क्ष\appbrk{}ण\appbrk{}शिरः}}\Afootnote{H; \skt{स्फिक्कु\appbrk{}क्षि\appbrk{}क\appbrk{}क्षा\appbrk{}व\appbrk{}ङ्क्ष\appbrk{}णो\appbrk{}रु\appbrk{}शिरः}\ A; \skt{स्फिक्कु\appbrk{}क्षि\appbrk{}च\appbrk{}ङ्क्ष\appbrk{}ण\appbrk{}शि\appbrk{}रस्}\ N}}\skt{ सुगाढः~। }\edtext{\skt{शाखा\-व\-द\-न\-क\-र्ण्ण\-क\-ण्ठ\-मे\-ढ्र\-मु\-ष्क\-पा\-र्श्वो\-द\-रो\-रस्मु}}{\lemma{\skt{शाखा\appbrk{}व\appbrk{}द\appbrk{}न\appbrk{}क\appbrk{}र्ण्ण\appbrk{}क\appbrk{}ण्ठ\appbrk{}मे\appbrk{}ढ्र\appbrk{}मु\appbrk{}ष्क\appbrk{}पा\appbrk{}र्श्वो\appbrk{}द\appbrk{}रो\appbrk{}रस्मु}}\Afootnote{H; \skt{शाखा\appbrk{}व\appbrk{}द\appbrk{}न\appbrk{}क\appbrk{}र्ण\appbrk{}क\appbrk{}ण्ठ\appbrk{}मे\appbrk{}ढ्र\appbrk{}मु\appbrk{}ष्क\appbrk{}पृ\appbrk{}ष्ठ\appbrk{}पा\appbrk{}र्श्वो\appbrk{}द\appbrk{}रो\appbrk{}रःसु}\ A; \skt{शाखा\appbrk{}व\appbrk{}द\appbrk{}न\appbrk{}क\appbrk{}र्ण्ण\appbrk{}क\appbrk{}ण्ठ\appbrk{}मे\appbrk{}ढ्र\appbrk{}सु\appbrk{}ष्को\appbrk{}पा\appbrk{}र्स्वो\appbrk{}द\appbrk{}रो\appbrk{}रःसु}\ N}}\skt{ समः~। अक्ष्णोः }\edtext{\skt{सन्धिषु च}}{\lemma{\skt{सन्धिषु च}}\Afootnote{A H; \skt{सन्धिसु च}\ N}}\skt{ }\edtext{\skt{शिथिलः}}{\lemma{\skt{शिथिलः}}\Afootnote{N H; \skt{शिथिल इति}\ A}}\skt{~।}\pend

\pstart \ledsidenote{\passref{1.18.25}}\skt{तत्र पैत्तिकं गाढस्थाने समं बध्नीयात्~। समस्थाने शिथिलं~। शिथिलस्थाने नैव }\edtext{\skt{शोणितदुष्टञ्च}}{\lemma{\skt{शोणितदुष्टञ्च}}\Afootnote{N H; \skt{एवं शोणितदुष्टं च}\ A}}\skt{~॥ }\edtext{\skt{श्लैष्मिकं}}{\lemma{\skt{श्लैष्मिकं}}\Afootnote{A H; \skt{श्लेष्मिकं}\ N}}\skt{ शिथिलस्थाने समं~। समस्थाने गाढं~। गाढस्थाने गाढतरं~। }\edtext{\skt{वातदुष्टञ्च}}{\lemma{\skt{वातदुष्टञ्}}\Afootnote{N H; \skt{एवं वातदुष्टं}\ A}}\skt{~॥}\pend

\pstart \ledsidenote{\passref{1.18.26}}\skt{तत्र }\edtext{\skt{पैत्तिकं}}{\lemma{\skt{पैत्तिकं}}\Afootnote{A H; \skt{पैत्ति}\ N}}\skt{ }\edtext{\skt{ग्रीष्मे}}{\lemma{\skt{ग्रीष्मे}}\Afootnote{H; \skt{शरदि ग्रीष्मे}\ A; \skt{क}\ N}}\skt{ }\edtext{\skt{द्विरह्ने}}{\lemma{\skt{द्विरह्ने}}\Afootnote{H; \skt{द्विरह्नो}\ A N}}\skt{ }\edtext{\skt{बध्नीयात्~। रक्तो\-प\-द्रु\-त\-म\-प्येवं}}{\lemma{\skt{बध्नीयात्~। रक्तो\appbrk{}प\appbrk{}द्रु\appbrk{}त\appbrk{}म\appbrk{}प्येवं}}\Afootnote{\skt{बध्नी\appbrk{}या\appbrk{}त्र\appbrk{}क्तो\appbrk{}प\appbrk{}द्रु\appbrk{}त\appbrk{}म\appbrk{}प्येवं}\ A N H}}\skt{~। }\edtext{\skt{श्लैष्मिकं}}{\lemma{\skt{श्लैष्मिकं}}\Afootnote{A H; \skt{स्लेष्मिकं}\ N}}\skt{ }\edtext{\skt{हेमन्ते}}{\lemma{\skt{हेमन्ते}}\Afootnote{N H; \skt{हेमन्त}\ A}}\skt{ }\edtext{\edtext{\skt{तृतीयेऽह्नि}}{\lemma{\skt{तृतीये}}\Afootnote{N H; \skt{वसन्तयो}\ A}}}{\lemma{\skt{ऽह्नि}}\Afootnote{H; \skt{स्त्र्यहा}\ A; \skt{हनि}\ N}}\skt{~। }\edtext{\skt{वातो\-प\-द्रु\-त\-म\-प्येवं}}{\lemma{\skt{वातो\appbrk{}प\appbrk{}द्रु\appbrk{}त\appbrk{}म\appbrk{}प्येवं}}\Afootnote{\skt{त् , वातो\appbrk{}प\appbrk{}द्रु\appbrk{}त\appbrk{}म\appbrk{}प्ये\appbrk{}वम्}\ A; \skt{वातो\appbrk{}प\appbrk{}द्रु\appbrk{}त\appbrk{}म\appbrk{}प्येवं}\ N H}}\skt{~। }\edtext{\skt{एवं}}{\lemma{\skt{एवं}}\Afootnote{H; \skt{एव}\ A N}}\skt{ }\edtext{\skt{मूह्य}}{\lemma{\skt{मूह्य}}\Afootnote{N H; \skt{मभ्यूह्य}\ A}}\skt{ }\edtext{\skt{बंध\-म्वि\-प\-र्य\-यञ्च}}{\lemma{\skt{बंध\appbrk{}म्वि\appbrk{}प\appbrk{}र्य\appbrk{}यञ्च}}\Afootnote{N H; \skt{बन्धविपर्ययं च}\ A}}\skt{ }\edtext{\skt{कुर्व्वीत}}{\lemma{\skt{कुर्व्वीत}}\Afootnote{H; \skt{कुर्यात्}\ A; \skt{कुर्वीतः}\ N}}\skt{~।}\pend

\pstart \ledsidenote{\passref{1.18.27}}\edtext{\skt{सम\-शि\-थि\-ल\-स्था\-नेषु}}{\lemma{\skt{सम\appbrk{}शि\appbrk{}थि\appbrk{}ल\appbrk{}स्था\appbrk{}नेषु}}\Afootnote{N H; \skt{तत्र, सम\appbrk{}शि\appbrk{}थि\appbrk{}ल\appbrk{}स्था\appbrk{}नेषु}\ A}}\skt{ }\edtext{\skt{गाढ}}{\lemma{\skt{गाढ}}\Afootnote{H; \skt{गा}\ A; \emph{om.}\ N}}\skt{ }\edtext{\skt{बंधे}}{\lemma{\skt{बंधे}}\Afootnote{H; \skt{ढं}\ A; \emph{om.}\ N}}\skt{ }\edtext{\edtext{\skt{विके\-सि\-कौ\-ष\-ध\-नै\-र\-र्थक्यं}}{\lemma{\skt{विके\appbrk{}सि\appbrk{}कौ\appbrk{}ष\appbrk{}ध\appbrk{}नैर्}}\Afootnote{H; \emph{om.}\ A N}}}{\lemma{\skt{अर्थक्यं}}\Afootnote{H; \emph{om.}\ A N}}\skt{ }\edtext{\skt{शोफ\-वे\-द\-ना\-प्रा\-दु\-र्भ्भा\-वश्च}}{\lemma{\skt{शोफ\appbrk{}वे\appbrk{}द\appbrk{}ना\appbrk{}प्रा\appbrk{}दु\appbrk{}र्भ्भा\appbrk{}वश्च}}\Afootnote{H; \emph{om.}\ A N}}\skt{~। }\edtext{\skt{गाढ\-स\-म\-स्था\-नेषु}}{\lemma{\skt{गाढ\appbrk{}स\appbrk{}म\appbrk{}स्था\appbrk{}नेषु}}\Afootnote{H; \emph{om.}\ A; \skt{गाढ}\ N}}\skt{ }\edtext{\skt{सिथिल}}{\lemma{\skt{सिथिल}}\Afootnote{H; \skt{ब}\ A; \emph{om.}\ N}}\skt{ }\edtext{\skt{बन्धे}}{\lemma{\skt{बन्धे}}\Afootnote{N H; \skt{द्धे}\ A}}\skt{ }\edtext{\edtext{\skt{विके\-शि\-कौ\-ष\-ध\-नै\-र\-र्थक्यं}}{\lemma{\skt{विके\appbrk{}शि\appbrk{}कौ\appbrk{}ष\appbrk{}ध\appbrk{}नैर्}}\Afootnote{A N; \skt{विकेसिकौषध}\ H}}\skt{ शोफ\-वे\-द\-ना\-प्रा\-दु\-र्भा\-वश्च}}{\lemma{\skt{अर्थक्यं शोफ\appbrk{}वे\appbrk{}द\appbrk{}ना\appbrk{}प्रा\appbrk{}दु\appbrk{}र्भा\appbrk{}वश्च}}\Afootnote{A; \skt{अर्थ\appbrk{}क्यै\appbrk{}शो\appbrk{}फ\appbrk{}वे\appbrk{}द\appbrk{}ना\appbrk{}प्रा\appbrk{}दु\appbrk{}र्भा\appbrk{}वश्च}\ N; \emph{om.}\ H}}\skt{ }\edtext{\skt{गाढ\-स\-म\-स्था\-नेषु}}{\lemma{\skt{गाढ\appbrk{}स\appbrk{}म\appbrk{}स्था\appbrk{}नेषु}}\Afootnote{A N; \emph{om.}\ H}}\skt{ }\edtext{\skt{शिथिलं बद्धे विके\-शि\-कौ\-ष\-ध\-प\-तनं}}{\lemma{\skt{शिथिलं बद्धे विके\appbrk{}शि\appbrk{}कौ\appbrk{}ष\appbrk{}ध\appbrk{}प\appbrk{}तनं}}\Afootnote{A; \skt{शिथि\appbrk{}ल\appbrk{}व\appbrk{}र्धे\appbrk{}वि\appbrk{}के\appbrk{}शि\appbrk{}कौ\appbrk{}ष\appbrk{}ध\appbrk{}प\appbrk{}तनं}\ N; \skt{पतनम्}\ H\textsuperscript{ac}; \skt{पल ञ्तनम्}\ H\textsuperscript{pc}}}\skt{ }\edtext{\edtext{\skt{पट्ट\-स\-ञ्चा\-रा\-द्व्र\-ण\-व\-क्त्र\-घ\-र्ष\-णञ्च}}{\lemma{\skt{पट्ट\appbrk{}स\appbrk{}ञ्चा\appbrk{}रा\appbrk{}द्व्र\appbrk{}ण\appbrk{}व\appbrk{}क्त्र\appbrk{}घ\appbrk{}र्ष\appbrk{}णञ्}}\Afootnote{N; \skt{पट्ट\appbrk{}स\appbrk{}ञ्चा\appbrk{}रा\appbrk{}द्व्र\appbrk{}ण\appbrk{}व\appbrk{}र्त्मा\appbrk{}व\appbrk{}घ\appbrk{}र्ष\appbrk{}ण\appbrk{}मिति; गाढ\appbrk{}शि\appbrk{}थि\appbrk{}ल\appbrk{}स्था\appbrk{}नेषु समं}\ A; \skt{पस\appbrk{}ञ्चा\appbrk{}रा\appbrk{}द्व्र\appbrk{}ण\appbrk{}व\appbrk{}घ\appbrk{}र्ष\appbrk{}णञ्च}\ H\textsuperscript{ac}; \skt{पट्ट\appbrk{}स\appbrk{}ञ्चा\appbrk{}रा\appbrk{}द्व्र\appbrk{}ण\appbrk{}व\appbrk{}घ\appbrk{}र्ष\appbrk{}णञ्च}\ H\textsuperscript{pc}}}}{\lemma{\skt{च}}\Afootnote{N; \skt{बद्धे च गुणाभाव इति}\ A; \emph{om.}\ H}}\skt{~।}\pend

\pstart \ledsidenote{\passref{1.18.28}}\edtext{\skt{अवि\-प\-री\-त\-बन्धे}}{\lemma{\skt{अवि\appbrk{}प\appbrk{}री\appbrk{}त\appbrk{}बन्धे}}\Afootnote{A H; \skt{अवि\appbrk{}प\appbrk{}री\appbrk{}त\appbrk{}बन्ध}\ N}}\skt{ }\edtext{\skt{वेद\-ना\-शा\-न्ति\-र\-सृ\-क्प्र\-सादो}}{\lemma{\skt{वेदनाशान्तिर्}}\Afootnote{N H; \skt{वेद\appbrk{}नो\appbrk{}प\appbrk{}शा\appbrk{}न्तिर्}\ A}}\skt{ मार्द्दवञ्च~।}\pend

\pstart \ledsidenote{\passref{1.18.29}}\edtext{\skt{अबध्यमाने}}{\lemma{\skt{अबध्यमाने}}\Afootnote{N H; \skt{अबध्यमानो}\ A}}\skt{ }\edtext{\edtext{\edtext{\skt{शीत\-वा\-ता\-त\-प\-र\-जो\-व\-र्ष\-दं\-श\-म\-स\-क\-म\-क्षि\-का\-प्र\-भृ\-ति\-भि\-र\-भि\-घा\-त\-वि\-शे\-षै\-र्ह\-न्यते}}{\lemma{\skt{शीत\appbrk{}वा\appbrk{}ता\appbrk{}त\appbrk{}प\appbrk{}र\appbrk{}जो\appbrk{}व\appbrk{}र्ष\appbrk{}दं\appbrk{}श\appbrk{}म\appbrk{}स\appbrk{}क\appbrk{}म\appbrk{}क्षि\appbrk{}का\appbrk{}प्र\appbrk{}भृ\appbrk{}ति\appbrk{}भिर्}}\Afootnote{N H; \skt{दंश\appbrk{}म\appbrk{}श\appbrk{}क\appbrk{}तृ\appbrk{}ण\appbrk{}का\appbrk{}ष्ठो\appbrk{}प\appbrk{}ल\appbrk{}पां\appbrk{}शु\appbrk{}शी\appbrk{}त\appbrk{}वा\appbrk{}ता\appbrk{}त\appbrk{}प\appbrk{}प्र\appbrk{}भृ\appbrk{}ति\appbrk{}भिर्}\ A}}}{\lemma{\skt{अभि\appbrk{}घा\appbrk{}त\appbrk{}वि\appbrk{}शे\appbrk{}षैर्}}\Afootnote{H; \skt{विशेषैर्}\ A; \skt{अति\appbrk{}घा\appbrk{}त\appbrk{}वि\appbrk{}शे\appbrk{}षैर्}\ N}}}{\lemma{\skt{हन्यते}}\Afootnote{N H; \skt{अभिहन्यते}\ A}}\skt{ व्रणः~। विवि\-ध\-वे\-द\-नो\-प\-द्रु\-तश्च दुष्टतामुपैति~। आलेपनादीनि चास्य }\edtext{\skt{विशो\-प\-मु\-प\-यान्ति}}{\lemma{\skt{विशोपम्}}\Afootnote{H; \skt{विशोषम्}\ A N}}\skt{~॥}\pend

\pstart\versepar \ledsidenote{\passref{1.18.30}}\edtext{\skt{चूर्ण्णि\-त\-म्म\-थि\-त\-म्भग्नं}}{\lemma{\skt{भग्नं}}\Afootnote{A H; \skt{भग्न}\ N}}\skt{ }\edtext{\skt{विश्लि\-ष्ट\-म\-ति\-पा\-दितं}}{\lemma{\skt{अतिपादितं}}\Afootnote{H; \skt{अतिपातितम्}\ A N}}\skt{~। }\pend
\pstart\versepar \skt{}\edtext{\skt{अस्थि\-स्ना\-यु\-सि\-रा\-च्छि\-न्न\-माशु}}{\lemma{\skt{अस्थि\appbrk{}स्ना\appbrk{}यु\appbrk{}सि\appbrk{}रा\appbrk{}च्छि\appbrk{}न्नम्}}\Afootnote{A H; \skt{अस्थि\appbrk{}श्ना\appbrk{}यु\appbrk{}च्छि\appbrk{}रा\appbrk{}च्छि\appbrk{}न्नम्}\ N}}\skt{ बन्धेन रोहति~॥}\pend

\pstart\versepar \ledsidenote{\passref{1.18.31}}\edtext{\skt{सुखञ्च}}{\lemma{\skt{सुखञ्च}}\Afootnote{N H; \skt{सुखमे}\ A}}\skt{ }\edtext{\skt{व्रणितः}}{\lemma{\skt{व्रणितः}}\Afootnote{N H; \skt{वं व्रणी}\ A}}\skt{ शेते सुखङ्गच्छति तिष्ठति~। }\pend
\pstart\versepar \skt{}\edtext{\skt{सुख}}{\lemma{\skt{सुख}}\Afootnote{N H; \skt{सुखं}\ A}}\skt{ शय्यासनस्थस्य क्षिप्रं संरोहति व्रणः~॥}\pend

\pstart \ledsidenote{\passref{1.18.32}}\edtext{\skt{अबन्ध्या}}{\lemma{\skt{अबन्ध्या}}\Afootnote{N H; \skt{अबन्ध्याः}\ A}}\skt{ }\edtext{\skt{पित्त\-र\-क्ता\-भि\-घा\-त\-नि\-मित्ता}}{\lemma{\skt{पित्त\appbrk{}र\appbrk{}क्ता\appbrk{}भि\appbrk{}घा\appbrk{}त\appbrk{}नि\appbrk{}मित्ता}}\Afootnote{H; \skt{पित्त\appbrk{}र\appbrk{}क्ता\appbrk{}भि\appbrk{}घा\appbrk{}त\appbrk{}वि\appbrk{}ष\appbrk{}नि\appbrk{}मित्ता}\ A; \skt{पित्त\appbrk{}र\appbrk{}क्ता\appbrk{}ख्या\appbrk{}भि\appbrk{}घा\appbrk{}त\appbrk{}नि\appbrk{}मित्ता}\ N}}\skt{ }\edtext{\skt{यदा\-तो\-द\-दा\-ह\-वे\-द\-ना\-भि\-भूताः}}{\lemma{\skt{यदा\appbrk{}तो\appbrk{}द\appbrk{}दा\appbrk{}ह\appbrk{}वे\appbrk{}द\appbrk{}ना\appbrk{}भि\appbrk{}भूताः}}\Afootnote{H; \skt{यदा शोफ\appbrk{}दा\appbrk{}ह\appbrk{}पा\appbrk{}क\appbrk{}रा\appbrk{}ग\appbrk{}तो\appbrk{}द\appbrk{}वे\appbrk{}द\appbrk{}ना\appbrk{}भि\appbrk{}भूताः}\ A; \skt{यदा तोद\appbrk{}दा\appbrk{}ह\appbrk{}वे\appbrk{}द\appbrk{}ना\appbrk{}भि\appbrk{}भृताः}\ N}}\skt{ क्षाराग्निदग्धाः }\edtext{\skt{पाका\-त्प्र\-कु\-थि\-त\-प्र\-सी\-र्ण्ण\-मां\-साश्च}}{\lemma{\skt{पाका\appbrk{}त्प्र\appbrk{}कु\appbrk{}थि\appbrk{}त\appbrk{}प्र\appbrk{}सी\appbrk{}र्ण्ण\appbrk{}मां\appbrk{}साश्च}}\Afootnote{N H; \skt{पाका\appbrk{}त्प्र\appbrk{}कु\appbrk{}थि\appbrk{}त\appbrk{}प्र\appbrk{}शी\appbrk{}र्ण\appbrk{}मां\appbrk{}साश्च}\ A}}\skt{ भवन्ति~॥}\pend

\pstart\versepar \ledsidenote{\passref{1.18.33}}\edtext{\skt{कुष्ठि\-ना\-म\-ग्नि\-द\-ग्धानां}}{\lemma{\skt{कुष्ठिनाम्}}\Afootnote{A H; \skt{भवन्ति चात्र~॥ कुष्टिनाम्}\ N}}\skt{ }\edtext{\skt{पिटका}}{\lemma{\skt{पिटका}}\Afootnote{N H; \skt{पिडका}\ A}}\skt{ मधुमेहिनां~। }\pend
\pstart\versepar \skt{}\edtext{\skt{कर्ण्णि\-का\-श्चो\-न्दु\-रु\-विषे}}{\lemma{\skt{चोन्दुरुविषे}}\Afootnote{A H; \skt{चोन्दुरुविषैर्}\ N}}\skt{ }\edtext{\skt{विष\-जु\-ष्ट\-व्र\-णाश्च ये}}{\lemma{\skt{विष\appbrk{}जु\appbrk{}ष्ट\appbrk{}व्र\appbrk{}णाश्च ये}}\Afootnote{H; \skt{विषजुष्टाश्च ये व्रणाः}\ A; \skt{विष\appbrk{}ज\appbrk{}ष्ट\appbrk{}व्र\appbrk{}णाश्च ये}\ N}}\skt{~॥}\pend

\pstart\versepar \ledsidenote{\passref{1.18.34}}\edtext{\skt{मांसपाके न}}{\lemma{\skt{मांसपाके न}}\Afootnote{A H; \skt{मान्सापाकेन}\ N}}\skt{ बध्यन्ते गुदपाके च दारुणे~। }\pend
\pstart\versepar \skt{}\edtext{\edtext{\skt{स्वबुद्ध्यार्वपि}}{\lemma{\skt{स्वबुद्ध्यार्}}\Afootnote{H; \skt{स्वबुद्ध्या}\ A N}}}{\lemma{\skt{वपि}}\Afootnote{H; \skt{चापि}\ A N}}\skt{ }\edtext{\skt{विभ\-जे\-त्कृ\-त्या\-कृ\-त्येषु}}{\lemma{\skt{विभ\appbrk{}जे\appbrk{}त्कृ\appbrk{}त्या\appbrk{}कृ\appbrk{}त्येषु}}\Afootnote{H; \skt{विभ\appbrk{}जे\appbrk{}त्कृ\appbrk{}त्या\appbrk{}कृ\appbrk{}त्यांश्च}\ A; \skt{विभ\appbrk{}जो\appbrk{}त्कृ\appbrk{}त्या\appbrk{}कृ\appbrk{}त्येषु}\ N}}\skt{ बुद्धिमान्~॥}\pend

\pstart\versepar \ledsidenote{\passref{1.18.35}}\edtext{\edtext{\skt{दोषन्देहञ्च}}{\lemma{\skt{दोषन्}}\Afootnote{H; \skt{देशं}\ A; \skt{दोष}\ N}}}{\lemma{\skt{देहञ्च}}\Afootnote{H; \skt{दोषं च}\ A; \skt{ञ्हञ्च}\ N\textsuperscript{ac}; \skt{ञ्देहञ्च}\ N\textsuperscript{pc}}}\skt{ विज्ञाय }\edtext{\skt{व्रणन्श्च}}{\lemma{\skt{व्रणन्श्च}}\Afootnote{H; \skt{व्रणं च}\ A; \skt{व्रणासु}\ N}}\skt{ व्रणकोविदः~। }\pend
\pstart\versepar \skt{}\edtext{\skt{ऋतूञ्च}}{\lemma{\skt{ऋतूञ्च}}\Afootnote{H; \skt{ऋतूंश्च}\ A; \skt{ऋतुञ्च}\ N}}\skt{ }\edtext{\skt{परिसंख्याय… ततो }\edtext{\skt{बन्धी\-न्नि\-वे\-श\-ये\-दिति}}{\lemma{\skt{बन्धीन्}}\Afootnote{H; \skt{बन्धान्}\ A; \skt{बन्धा}\ N}}}{\lemma{\skt{परिसंख्याय{\symfont ‡} ततो}}\Afootnote{H; \skt{परिसङ्ख्याय ततो}\ A N}}\skt{~॥}\edtext{}{\lemma{\skt{निवेशयेदिति}}\Afootnote{H; \skt{निवेशयेत्}\ A N}}\pend

\pstart \ledsidenote{\passref{1.18.trailer}}\edtext{\skt{डोड}}{\lemma{\skt{डोड}}\Afootnote{H; \skt{इति सुश्रु\appbrk{}त\appbrk{}सं\appbrk{}हि\appbrk{}तायां सूत्रस्थाने व्रणा\appbrk{}ले\appbrk{}प\appbrk{}न\appbrk{}ब\appbrk{}न्ध\appbrk{}वि\appbrk{}धि\appbrk{}र्ना\appbrk{}मा\appbrk{}ष्टा\appbrk{}द\appbrk{}शोऽ\appbrk{}ध्यायः}\ A; \skt{२५~॥ {\symfont ❈}}\ K; \emph{om.}\ N}}\skt{~॥}\pend

\endnumbering
\beginnumbering
\eledsection*{Adhyāya 19}
\pstart \ledsidenote{\passref{1.19.1}}\skt{॥ अथातो }\edtext{\skt{व्रणि\-तो\-पा\-स\-नी\-य\-म्व्या\-ख्या\-स्यामः}}{\lemma{\skt{व्रणि\appbrk{}तो\appbrk{}पा\appbrk{}स\appbrk{}नी\appbrk{}यम्}}\Afootnote{H; \skt{व्रणि\appbrk{}तो\appbrk{}पा\appbrk{}स\appbrk{}नी\appbrk{}य\appbrk{}म\appbrk{}ध्यायं}\ A N}}\skt{~॥}\pend

\pstart \ledsidenote{\passref{1.19.3}}\edtext{\skt{अथ}}{\lemma{\skt{अथ}}\Afootnote{N H; \emph{om.}\ A}}\skt{ व्रणितस्य }\edtext{\edtext{\edtext{\skt{प्रथ\-म\-मे\-वा\-गा\-र\-म\-न्वि\-च्छेत्}}{\lemma{\skt{प्रथमम्}}\Afootnote{H; \skt{प्रथम}\ A; \skt{प्रथ}\ N}}}{\lemma{\skt{एवागारम्}}\Afootnote{H; \skt{मेवागार}\ A; \skt{मेवागारम्}\ N}}\skt{~॥ प्रश\-स्त\-वा\-स्तु\-सु\-नि\-विष्टं}}{\lemma{\skt{अन्विच्छेत्~॥ प्रश\appbrk{}स्त\appbrk{}वा\appbrk{}स्तु\appbrk{}सु\appbrk{}नि\appbrk{}विष्टं}}\Afootnote{H; \skt{मन्वि\appbrk{}च्छे\appbrk{}त्त\appbrk{}च्चा\appbrk{}गारं प्रशस्त}\ A; \skt{अन्विच्छेत्~। प्रस\appbrk{}स्त\appbrk{}वा\appbrk{}श्तु\appbrk{}सु\appbrk{}वि\appbrk{}निष्टं}\ N}}\skt{ }\edtext{\skt{शुच्य\-वा\-ता\-त\-पञ्च}}{\lemma{\skt{शुच्य\appbrk{}वा\appbrk{}ता\appbrk{}त\appbrk{}पञ्च}}\Afootnote{H; \skt{वास्त्वादिकं कार्यम्}\ A; \skt{शुच्यवार्तपञ्च}\ N}}\skt{~।}\pend

\pstart\versepar \ledsidenote{\passref{1.19.4}}\edtext{\skt{अप्र\-श\-स्त\-वा\-स्तु\-नि\-गृहे}}{\lemma{\skt{अप्र\appbrk{}श\appbrk{}स्त\appbrk{}वा\appbrk{}स्तु\appbrk{}नि\appbrk{}गृहे}}\Afootnote{N H; \skt{प्रशस्तवास्तुनि गृहे}\ A}}\skt{ }\edtext{\edtext{\skt{सम्बा\-धेऽ\-शु\-चि\-न्या\-तपे}}{\lemma{\skt{सम्बाधे}}\Afootnote{N H; \emph{om.}\ A}}}{\lemma{\skt{ऽशुचिन्यातपे}}\Afootnote{H; \skt{शुचावातप}\ A; \skt{शुचिन्यातपे}\ N}}\skt{ }\pend
\pstart\versepar \skt{}\edtext{\skt{चातिवाते च}}{\lemma{\skt{चातिवाते च}}\Afootnote{N H; \skt{वर्जिते~। निवाते न च}\ A}}\skt{ रोगाः }\edtext{\skt{स्युश्चिरं}}{\lemma{\skt{स्युश्चिरं}}\Afootnote{H; \skt{स्युः}\ A; \skt{स्युश्चिर}\ N}}\skt{ }\edtext{\skt{शारीरमानसाः}}{\lemma{\skt{शारीरमानसाः}}\Afootnote{N H; \skt{शारी\appbrk{}रा\appbrk{}ग\appbrk{}न्तु\appbrk{}मा\appbrk{}नसाः}\ A}}\skt{~।}\pend

\pstart \ledsidenote{\passref{1.19.5}}\edtext{\edtext{\edtext{\skt{तस्मि\-ञ्छ\-य\-न\-म\-स\-म्बाधं}}{\lemma{\skt{तस्मिञ्}}\Afootnote{N H; \skt{तस्मि}\ A}}}{\lemma{\skt{छयनम्}}\Afootnote{H; \skt{ञ्शयन}\ A; \skt{च्छयन}\ N}}}{\lemma{\skt{असम्बाधं}}\Afootnote{H; \skt{मसम्बाधं}\ A N}}\skt{ }\edtext{\skt{स्वास्तीर्णअणं}}{\lemma{\skt{स्वास्तीर्णअणं}}\Afootnote{H; \skt{स्वास्तीर्णं}\ A; \skt{वातीर्ण्णं}\ N}}\skt{ }\edtext{\skt{मनो\-ज्ञ\-म्प्रा\-क्छी\-र्ष\-न्स\-श\-स्त्र\-ङ्कु\-र्व्वीत}}{\lemma{\skt{प्राक्छी\appbrk{}र्ष\appbrk{}न्स\appbrk{}श\appbrk{}स्त्रङ्}}\Afootnote{H; \skt{प्राक्शिरस्कं सशस्त्रं}\ A; \skt{प्राक्छींर्षससस्त्रं}\ N}}\skt{~।}\pend

\pstart\versepar \ledsidenote{\passref{1.19.6}}\skt{सुख\-चे\-ष्टा\-प्र\-चारः स्यात्~। स्वास्तीर्ण्णे शयने व्रणी }\pend
\pstart\versepar \skt{}\edtext{\skt{प्राच्यादिशि}}{\lemma{\skt{प्राच्यादिशि}}\Afootnote{N H; \skt{प्राच्यां दिशि}\ A}}\skt{ }\edtext{\skt{स्थितो}}{\lemma{\skt{स्थितो}}\Afootnote{H; \skt{स्थिता}\ A N}}\skt{ }\edtext{\edtext{\edtext{\skt{देवा\-स्त\-त्पू\-ज\-ना\-र्थ\-न्त\-तः\-शि\-रस्}}{\lemma{\skt{देवास्}}\Afootnote{N H; \skt{देवा}\ A}}}{\lemma{\skt{तत्पूजनार्थन्}}\Afootnote{N H; \skt{स्तत्पूजार्थं}\ A}}}{\lemma{\skt{ततःशिरस्}}\Afootnote{N H; \skt{च तच्छिरः}\ A}}\pend

\pstart \ledsidenote{\passref{1.19.7}}\skt{तस्मि\-न्सु\-हृ\-द्भि\-र\-नु\-कूलैः प्रिय\-म्व\-दै\-रु\-पा\-स्य\-मानो }\edtext{\skt{यथे\-ष्ट\-मा\-शी\-तेति}}{\lemma{\skt{यथे\appbrk{}ष्ट\appbrk{}मा\appbrk{}शी\appbrk{}तेति}}\Afootnote{H; \skt{यथेष्टमासीत}\ A; \skt{यथे\appbrk{}ष्ट\appbrk{}मा\appbrk{}शी\appbrk{}तेडि}\ N}}\skt{~॥}\pend

\pstart \ledsidenote{\passref{1.19.8}}\edtext{\skt{भवन्ति}}{\lemma{\skt{भवन्ति}}\Afootnote{N; \skt{सुहृदो विक्षिपन्त्याशु कथा\appbrk{}भि\appbrk{}र्व्र\appbrk{}ण\appbrk{}वे\appbrk{}दनाः~। आश्वासयन्तो}\ A; \skt{भ वति चात्र}\ H}}\skt{ }\edtext{\skt{चात्र॥}}{\lemma{\skt{चात्र}}\Afootnote{N; \skt{बहुशः स्वनुकूलाः प्रियंवदाः}\ A; \emph{om.}\ H}}\pend

\pstart\versepar \ledsidenote{\passref{1.19.8a}}\skt{सुहृदो }\edtext{\skt{विक्षिपन्त्याशु}}{\lemma{\skt{आशु}}\Afootnote{H; \skt{आसुः}\ N}}\skt{ }\edtext{\edtext{\skt{कथा\-भि\-र्व्व्र\-ण\-वे\-दनाः}}{\lemma{\skt{कथाभिर्}}\Afootnote{H; \skt{कथाभि}\ N}}}{\lemma{\skt{व्व्रणवेदनाः}}\Afootnote{H; \skt{र्व्रणवेदनाः}\ N}}\skt{~। }\pend
\pstart\versepar \skt{आश्वासयन्तो }\edtext{\skt{बहुशः}}{\lemma{\skt{बहुशः}}\Afootnote{H; \skt{बहुषः}\ N}}\skt{ }\edtext{\skt{स्वनुकूलाः}}{\lemma{\skt{स्वनुकूलाः}}\Afootnote{H; \skt{स्वनुकूला}\ N}}\skt{ प्रियम्वदाः~॥}\pend

\pstart \ledsidenote{\passref{1.19.9}}\skt{न च }\edtext{\skt{दिवा\-नि\-द्रा\-व\-शगः}}{\lemma{\skt{दिवा\appbrk{}नि\appbrk{}द्रा\appbrk{}व\appbrk{}शगः}}\Afootnote{A H; \skt{निद्रावसगः}\ N}}\skt{ स्याद्}\pend

\pstart \ledsidenote{\passref{1.19.11}}\edtext{\skt{उत्था\-स\-म्वे\-स\-न\-प\-रि\-मा\-र्ज्ज\-ना\-दिषु}}{\lemma{\skt{उत्था\appbrk{}स\appbrk{}म्वे\appbrk{}स\appbrk{}न\appbrk{}प\appbrk{}रि\appbrk{}मा\appbrk{}र्ज्ज\appbrk{}ना\appbrk{}दिषु}}\Afootnote{H; \skt{उत्था\appbrk{}न\appbrk{}सं\appbrk{}वे\appbrk{}श\appbrk{}न\appbrk{}प\appbrk{}रि\appbrk{}व\appbrk{}र्त\appbrk{}न\appbrk{}च\appbrk{}ङ्क्र\appbrk{}म\appbrk{}णो\appbrk{}च्चै\appbrk{}र्भा\appbrk{}षणा}\ A; \skt{उत्था\appbrk{}रा\appbrk{}सं\appbrk{}वे\appbrk{}स\appbrk{}न\appbrk{}प\appbrk{}रि\appbrk{}मा\appbrk{}र्ज्ज\appbrk{}ना\appbrk{}दिषु}\ N}}\skt{ }\edtext{\skt{चात्म\-चे\-ष्टा\-स्व\-प्र\-मत्तो}}{\lemma{\skt{चात्म\appbrk{}चे\appbrk{}ष्टा\appbrk{}स्व\appbrk{}प्र\appbrk{}मत्तो}}\Afootnote{N; \skt{द्यास्वा\appbrk{}त्म\appbrk{}चे\appbrk{}ष्टा\appbrk{}स्व\appbrk{}प्र\appbrk{}मत्तो}\ A; \skt{चात्म\appbrk{}चे\appbrk{}ष्टा\appbrk{}स्व\appbrk{}प्र\appbrk{}मत्तो}\ H}}\skt{ व्रणं }\edtext{\skt{रक्षेत्~॥ भवन्ति}}{\lemma{\skt{रक्षेत्~॥ भवन्ति}}\Afootnote{N; \skt{संरक्षेत्}\ A; \skt{रक्षेत्~॥ भवति चात्र}\ H}}\skt{ }\edtext{\skt{चात्र॥}}{\lemma{\skt{चात्र}}\Afootnote{N; \emph{om.}\ A H}}\pend

\pstart\versepar \ledsidenote{\passref{1.19.12}}\edtext{\skt{स्थाना\-स\-न\-ञ्च\-ङ्क्र\-मणं}}{\lemma{\skt{स्थाना\appbrk{}स\appbrk{}न\appbrk{}ञ्च\appbrk{}ङ्क्र\appbrk{}मणं}}\Afootnote{A H; \skt{स्थाना\appbrk{}श\appbrk{}न\appbrk{}ञ्च\appbrk{}क्र\appbrk{}म\appbrk{}नन्}\ N}}\skt{ दिवा\-स्व\-प्न\-न्त\-थैव च~। }\pend
\pstart\versepar \skt{व्रणितो न }\edtext{\skt{विनिषेवेत}}{\lemma{\skt{विनिषेवेत}}\Afootnote{H\textsuperscript{pc}; \skt{निषेवेत}\ A; \skt{निसेवेत}\ N; \skt{विषेवेत}\ H\textsuperscript{ac}}}\skt{ शक्तिमानपि मानवः~॥}\pend

\pstart \ledsidenote{\passref{1.19.14}}\skt{गम्यानाञ्च }\edtext{\skt{स्त्रीणां}}{\lemma{\skt{स्त्रीणां}}\Afootnote{A H; \skt{स्त्रीणा}\ N}}\skt{ सन्द\-र्श\-न\-स\-म्भा\-ष\-ण\-सं\-स्प\-र्श\-नानि दूरत एव }\edtext{\skt{परिहरेत्~॥ भवन्ति}}{\lemma{\skt{परिहरेत्~॥ भवन्ति}}\Afootnote{\skt{परिहरेत्}\ A; \skt{परिहरेत्~॥ भवति}\ N; \skt{परिहरेत्~॥ भ}\ H}}\skt{ }\edtext{\skt{चात्र॥}}{\lemma{\skt{चात्र}}\Afootnote{N; \emph{om.}\ A H}}\pend

\pstart\versepar \ledsidenote{\passref{1.19.15}}\edtext{\skt{स्त्रीणां}}{\lemma{\skt{स्त्रीणां}}\Afootnote{H; \emph{om.}\ A; \skt{स्त्रीणं}\ N}}\skt{ }\edtext{\edtext{\skt{सन्द\-र्श\-ना\-च्छु\-क्र\-ङ्क\-दा\-चि\-द्ब\-लितं}}{\lemma{\skt{सन्द\appbrk{}र्श\appbrk{}ना\appbrk{}च्छु\appbrk{}क्रङ्}}\Afootnote{N H; \skt{स्त्रीदर्शनादिभिः शुक्रं}\ A}}}{\lemma{\skt{कदाचिद्बलितं}}\Afootnote{H; \skt{कदाचिच्चलितं}\ A N}}\skt{ स्रवेत्~। }\pend
\pstart\versepar \skt{}\edtext{\edtext{\edtext{\skt{ग्राम्य\-ध\-र्म्म\-कृ\-ता\-न्दो\-षा\-न्सोऽ\-सं\-स\-र्ग्गेऽ\-प्य\-वा\-प्नु\-यात्}}{\lemma{\skt{दोषान्सो}}\Afootnote{A H; \skt{दोषांसो}\ N}}}{\lemma{\skt{ऽसंसर्ग्गे}}\Afootnote{A H; \skt{संसर्गे}\ N}}}{\lemma{\skt{ऽप्यवाप्नुयात्}}\Afootnote{A H; \skt{प्यवाप्नुयात्}\ N}}\skt{~॥}\edtext{}{\lemma{\skt{स्रवेत्~। ग्राम्य\appbrk{}ध\appbrk{}र्म्म\appbrk{}कृ\appbrk{}तान्}}\Afootnote{A H; \skt{श्रवेत्~। ग्राम्यधर्मकृता}\ N}}\pend

\pstart \ledsidenote{\passref{1.19.16}}\edtext{\skt{नव\-धा\-न्य\-मा\-ष\-ति\-ल\-क\-शा\-य\-कु\-ल\-त्थ\-नि\-ष्पा\-व\-क\-ह\-रि\-त\-शो\-का\-म्ल\-ल\-व\-ण\-गु\-ड\-पि\-ष्ट\-वि\-कृ\-त\-शु\-ष्क\-शा\-का\-जा\-वि\-मां\-स\-शीतो}}{\lemma{\skt{नव\appbrk{}धा\appbrk{}न्य\appbrk{}मा\appbrk{}ष\appbrk{}ति\appbrk{}ल\appbrk{}क\appbrk{}शा\appbrk{}य\appbrk{}कु\appbrk{}ल\appbrk{}त्थ\appbrk{}नि\appbrk{}ष्पा\appbrk{}व\appbrk{}क\appbrk{}ह\appbrk{}रि\appbrk{}त\appbrk{}शो\appbrk{}का\appbrk{}म्ल\appbrk{}ल\appbrk{}व\appbrk{}ण\appbrk{}गु\appbrk{}ड\appbrk{}पि\appbrk{}ष्ट\appbrk{}वि\appbrk{}कृ\appbrk{}त\appbrk{}शु\appbrk{}ष्क\appbrk{}शा\appbrk{}का\appbrk{}जा\appbrk{}वि\appbrk{}मां\appbrk{}स\appbrk{}शीतो}}\Afootnote{H; \skt{नव\appbrk{}धा\appbrk{}न्य\appbrk{}मा\appbrk{}ष\appbrk{}ति\appbrk{}ल\appbrk{}क\appbrk{}ला\appbrk{}य\appbrk{}कु\appbrk{}ल\appbrk{}त्थ\appbrk{}नि\appbrk{}ष्पा\appbrk{}व\appbrk{}ह\appbrk{}रि\appbrk{}त\appbrk{}क\appbrk{}शा\appbrk{}का\appbrk{}म्ल\appbrk{}ल\appbrk{}व\appbrk{}ण\appbrk{}क\appbrk{}टु\appbrk{}क\appbrk{}गु\appbrk{}ड\appbrk{}पि\appbrk{}ष्ट\appbrk{}वि\appbrk{}कृ\appbrk{}ति\appbrk{}व\appbrk{}ल्लू\appbrk{}र\appbrk{}शु\appbrk{}ष्क\appbrk{}शा\appbrk{}का\appbrk{}जा\appbrk{}वि\appbrk{}का\appbrk{}नू\appbrk{}पौ\appbrk{}दक- मांस\appbrk{}व\appbrk{}सा\appbrk{}शी\appbrk{}तो\appbrk{}द\appbrk{}क\appbrk{}कृ\appbrk{}श\appbrk{}रापा}\ A; \skt{नव\appbrk{}धा\appbrk{}न्य\appbrk{}मा\appbrk{}ष\appbrk{}ति\appbrk{}ल\appbrk{}क\appbrk{}षा\appbrk{}य\appbrk{}कु\appbrk{}ल\appbrk{}त्थ\appbrk{}नि\appbrk{}ष्पा\appbrk{}व\appbrk{}ह\appbrk{}रि\appbrk{}त\appbrk{}शा\appbrk{}का\appbrk{}व्ल\appbrk{}ल\appbrk{}व\appbrk{}ण\appbrk{}गु\appbrk{}ड\appbrk{}पि\appbrk{}ष्ट\appbrk{}वि\appbrk{}कृ\appbrk{}त\appbrk{}शु\appbrk{}ष्क\appbrk{}शा\appbrk{}का\appbrk{}जा\appbrk{}वि\appbrk{}मां\appbrk{}स\appbrk{}सी\appbrk{}तो\appbrk{}दक}\ N}}\skt{ }\edtext{\skt{दक\-द\-धि\-दु\-ग्ध\-त\-क्र\-प्र\-भृ\-तीनि}}{\lemma{\skt{दक\appbrk{}द\appbrk{}धि\appbrk{}दु\appbrk{}ग्ध\appbrk{}त\appbrk{}क्र\appbrk{}प्र\appbrk{}भृ\appbrk{}तीनि}}\Afootnote{H; \skt{यस\appbrk{}द\appbrk{}धि\appbrk{}दु\appbrk{}ग्ध\appbrk{}त\appbrk{}क्र\appbrk{}प्र\appbrk{}भृ\appbrk{}तीनि}\ A; \skt{दधि\appbrk{}दु\appbrk{}ग्ध\appbrk{}त\appbrk{}क्र\appbrk{}प्र\appbrk{}भृ\appbrk{}तीनि}\ N}}\skt{ }\edtext{\skt{परिहरेत्~॥ भवन्ति}}{\lemma{\skt{परिहरेत्~॥ भवन्ति}}\Afootnote{\skt{परिहरेत्}\ A H; \skt{परिहरेत्~॥ भवति}\ N}}\skt{ }\edtext{\skt{चात्र॥}}{\lemma{\skt{चात्र}}\Afootnote{N; \emph{om.}\ A H}}\pend

\pstart\versepar \ledsidenote{\passref{1.19.17}}\skt{तक्रान्तो }\edtext{\skt{नव\-धा\-न्या\-दि\-र्यो\-य\-म्वर्ग्ग}}{\lemma{\skt{नव\appbrk{}धा\appbrk{}न्या\appbrk{}दि\appbrk{}र्यो\appbrk{}य\appbrk{}म्वर्ग्ग}}\Afootnote{H; \skt{नव\appbrk{}धा\appbrk{}न्या\appbrk{}दि\appbrk{}र्योऽयं वर्ग}\ A; \skt{नवधान्यादिर्यो यवर्ग}\ N}}\skt{ उदाहृतः~। }\pend
\pstart\versepar \skt{}\edtext{\skt{दोषसञ्जननो}}{\lemma{\skt{दोषसञ्जननो}}\Afootnote{N H; \skt{दोषसञ्जननो ह्य्}\ A}}\skt{ }\edtext{\skt{ह्येष}}{\lemma{\skt{ह्येष}}\Afootnote{N H; \skt{एष}\ A}}\skt{ }\edtext{\skt{विज्ञेयः}}{\lemma{\skt{विज्ञेयः}}\Afootnote{A H; \skt{विज्ञेय}\ N}}\skt{ पूयवर्द्धनः~॥}\pend

\pstart \ledsidenote{\passref{1.19.18}}\edtext{\edtext{\skt{मद्य\-प\-ञ्य\-मै\-रे\-या\-रि\-ष्टा\-स\-व\-म\-धु\-सु\-रा\-वि\-हा\-रा\-म्प\-रि\-ह\-रेत्}}{\lemma{\skt{मद्य\appbrk{}प\appbrk{}ञ्य\appbrk{}मै\appbrk{}रे\appbrk{}या\appbrk{}रि\appbrk{}ष्टा\appbrk{}स\appbrk{}व\appbrk{}म\appbrk{}धु\appbrk{}सु\appbrk{}रा\appbrk{}वि\appbrk{}हा\appbrk{}राम्}}\Afootnote{H; \skt{मद्यपश्च मैरे\appbrk{}या\appbrk{}रि\appbrk{}ष्टा\appbrk{}स\appbrk{}व\appbrk{}सी\appbrk{}धु\appbrk{}सु\appbrk{}रा\appbrk{}वि\appbrk{}कारा}\ A; \skt{मद्यपश्च मैरि\appbrk{}या\appbrk{}रि\appbrk{}ष्टा\appbrk{}श\appbrk{}व\appbrk{}स\appbrk{}धु\appbrk{}र्थ\appbrk{}रा\appbrk{}वि\appbrk{}हारां}\ N}}}{\lemma{\skt{परिहरेत्}}\Afootnote{N H; \skt{न्परिहरेत्}\ A}}\skt{~॥}\pend

\pstart\versepar \ledsidenote{\passref{1.19.19}}\edtext{\skt{मद्यमम्लञ्च}}{\lemma{\skt{मद्यमम्लञ्च}}\Afootnote{H; \skt{मद्यमम्लं त}\ A; \skt{मद्यमव्लञ्च}\ N}}\skt{ }\edtext{\skt{रूक्षञ्च}}{\lemma{\skt{रूक्षञ्च}}\Afootnote{H; \skt{था रूक्षं}\ A; \skt{श्मक्षञ्च}\ N}}\skt{ तीक्ष्णमुष्णञ्च वीर्यतः~। }\pend
\pstart\versepar \skt{आशुकारि च तत्पीतं }\edtext{\skt{क्षिप्र\-म्व्या\-पा\-द\-ये\-द्व्रणं}}{\lemma{\skt{क्षिप्रम्}}\Afootnote{A H; \skt{क्षिप्र}\ N}}\skt{~॥}\pend

\pstart \ledsidenote{\passref{1.19.20}}\edtext{\edtext{\skt{वाता\-त\-प\-र\-जो\-धू\-मा\-ति\-भो\-ज\-ना\-नि\-ष्ट\-श्र\-व\-ण\-द\-र्श\-ना\-म\-र्ष\-शो\-क\-वि\-ष\-म\-श\-य\-ना\-स\-न\-रा\-त्रि\-जा\-ग\-र\-णा\-त्म\-क्षि\-का\-दि\-भि\-श्चा\-बा\-धा\-म्प\-रि\-ह\-रेत्}}{\lemma{\skt{वाता\appbrk{}त\appbrk{}प\appbrk{}र\appbrk{}जो\appbrk{}धू\appbrk{}मा\appbrk{}ति\appbrk{}भो\appbrk{}ज\appbrk{}ना\appbrk{}नि\appbrk{}ष्ट\appbrk{}श्र\appbrk{}व\appbrk{}ण\appbrk{}द\appbrk{}र्श\appbrk{}ना\appbrk{}म\appbrk{}र्ष\appbrk{}शो\appbrk{}क\appbrk{}वि\appbrk{}ष\appbrk{}म\appbrk{}श\appbrk{}य\appbrk{}ना\appbrk{}स\appbrk{}न\appbrk{}रा\appbrk{}त्रि\appbrk{}जा\appbrk{}ग\appbrk{}र\appbrk{}णा\appbrk{}त्म\appbrk{}क्षि\appbrk{}का\appbrk{}दि\appbrk{}भिश्}}\Afootnote{H; \skt{वाता\appbrk{}त\appbrk{}प\appbrk{}र\appbrk{}जो\appbrk{}धू\appbrk{}मा\appbrk{}व\appbrk{}श्या\appbrk{}या\appbrk{}ति\appbrk{}से\appbrk{}व\appbrk{}ना\appbrk{}ति\appbrk{}भो\appbrk{}ज\appbrk{}ना\appbrk{}नि\appbrk{}ष्ट\appbrk{}भो\appbrk{}ज\appbrk{}न\appbrk{}श्र\appbrk{}व\appbrk{}ण\appbrk{}द\appbrk{}र्श\appbrk{}ने\appbrk{}र्ष्या\appbrk{}म\appbrk{}र्ष\appbrk{}भ\appbrk{}य\appbrk{}शो\appbrk{}क\appbrk{}ध्या\appbrk{}न\appbrk{}रा\appbrk{}त्रि\appbrk{}जा\appbrk{}ग\appbrk{}र\appbrk{}ण\appbrk{}वि\appbrk{}ष\appbrk{}मा\appbrk{}शन- शय\appbrk{}नो\appbrk{}प\appbrk{}वा\appbrk{}स\appbrk{}वा\appbrk{}ग्व्या\appbrk{}या\appbrk{}म\appbrk{}स्था\appbrk{}न\appbrk{}च\appbrk{}ङ्क्र\appbrk{}म\appbrk{}ण\appbrk{}शी\appbrk{}त\appbrk{}वा\appbrk{}त\appbrk{}वि\appbrk{}रु\appbrk{}द्धा\appbrk{}ध्य\appbrk{}श\appbrk{}ना\appbrk{}जी\appbrk{}र्ण\appbrk{}म\appbrk{}क्षिका}\ A; \skt{वाता\appbrk{}त\appbrk{}प\appbrk{}र\appbrk{}जो\appbrk{}धू\appbrk{}मा\appbrk{}ति\appbrk{}भो\appbrk{}ज\appbrk{}ना\appbrk{}नि\appbrk{}ष्ट\appbrk{}श्र\appbrk{}व\appbrk{}ण\appbrk{}द\appbrk{}र्श\appbrk{}ना\appbrk{}म\appbrk{}प\appbrk{}शो\appbrk{}क\appbrk{}वि\appbrk{}ष\appbrk{}म\appbrk{}श\appbrk{}य\appbrk{}ना\appbrk{}स\appbrk{}न\appbrk{}रा\appbrk{}त्रि\appbrk{}जा\appbrk{}ग\appbrk{}र\appbrk{}णा\appbrk{}त्म\appbrk{}क्षि\appbrk{}कादि}\ N}}}{\lemma{\skt{चाबाधाम्}}\Afootnote{H; \skt{द्या बाधाः}\ A; \skt{ष्ठावाधं}\ N}}\skt{~॥}\pend

\pstart\versepar \ledsidenote{\passref{1.19.21}}\edtext{\skt{व्रणि\-त\-स्यो\-प\-त\-प्तस्य}}{\lemma{\skt{व्रणि\appbrk{}त\appbrk{}स्यो\appbrk{}प\appbrk{}त\appbrk{}प्तस्य}}\Afootnote{N H; \skt{व्रणिनः सम्प्रतप्तस्य}\ A}}\skt{ }\edtext{\skt{कारणैरेव}}{\lemma{\skt{एव}}\Afootnote{A H; \skt{एवम्}\ N}}\skt{ }\edtext{\skt{मादिभिः}}{\lemma{\skt{मादिभिः}}\Afootnote{A H; \skt{आदिभिः}\ N}}\skt{~। }\pend
\pstart\versepar \skt{क्षीण\-शो\-णि\-त\-मां\-सस्य भुक्तं }\edtext{\skt{सम्यग्न}}{\lemma{\skt{सम्यग्न}}\Afootnote{H; \skt{सम्यङ्न}\ A; \skt{संम्यग्न}\ N}}\skt{ }\edtext{\skt{जीर्यति}}{\lemma{\skt{जीर्यति}}\Afootnote{A H; \skt{जीग्यति}\ N}}\skt{~॥}\pend

\pstart\versepar \ledsidenote{\passref{1.19.22}}\edtext{\skt{अजी\-र्ण\-अ\-णा\-त्प\-व\-ना\-दीनां}}{\lemma{\skt{अजी\appbrk{}र्ण\appbrk{}अ\appbrk{}णा\appbrk{}त्प\appbrk{}व\appbrk{}ना\appbrk{}दीनां}}\Afootnote{H; \skt{अजी\appbrk{}र्णा\appbrk{}त्प\appbrk{}व\appbrk{}ना\appbrk{}दीनां}\ A N}}\skt{ विभ्र\-मो\-ब\-ल\-वा\-न्भ\-वेत्~। }\pend
\pstart\versepar \skt{ततः }\edtext{\skt{शोफ\-रु\-जा\-श्रा\-व\-दा\-ह\-पा\-का\-न\-वा\-प्नु\-यात्}}{\lemma{\skt{शोफ\appbrk{}रु\appbrk{}जा\appbrk{}श्रा\appbrk{}व\appbrk{}दा\appbrk{}ह\appbrk{}पा\appbrk{}का\appbrk{}न\appbrk{}वा\appbrk{}प्नु\appbrk{}यात्}}\Afootnote{N H; \skt{शोफ\appbrk{}रु\appbrk{}जा\appbrk{}स्रा\appbrk{}व\appbrk{}दा\appbrk{}ह\appbrk{}पा\appbrk{}का\appbrk{}न\appbrk{}वा\appbrk{}प्नु\appbrk{}यात्}\ A}}\skt{~॥}\edtext{}{\lemma{\skt{विभ्र\appbrk{}मो\appbrk{}ब\appbrk{}ल\appbrk{}वा\appbrk{}न्भ\appbrk{}वेत्~। ततः}}\Afootnote{H; \skt{विभ्रमो बलवान्भवेत्~। ततः}\ A; \skt{विभ्रमो बलवान्भवेत्~। ततः}\ N}}\pend

\pstart \ledsidenote{\passref{1.19.23}}\edtext{\skt{सदा च}}{\lemma{\skt{सदा च}}\Afootnote{N H; \skt{सदा}\ A}}\skt{ }\edtext{\skt{नीच\-न\-ख\-रो\-म्णा\-शु\-चि\-ना\-शु\-चि\-वा\-स\-सा\-शा\-न्ति\-म\-ङ्ग\-ल\-दे\-व\-ता\-ब्रा\-ह्म\-ण\-गु\-रु\-प\-रे\-ण\-भ\-वि\-तव्यं}}{\lemma{\skt{नीच\appbrk{}न\appbrk{}ख\appbrk{}रो\appbrk{}म्णा\appbrk{}शु\appbrk{}चि\appbrk{}ना\appbrk{}शु\appbrk{}चि\appbrk{}वा\appbrk{}स\appbrk{}सा\appbrk{}शा\appbrk{}न्ति\appbrk{}म\appbrk{}ङ्ग\appbrk{}ल\appbrk{}दे\appbrk{}व\appbrk{}ता\appbrk{}ब्रा\appbrk{}ह्म\appbrk{}ण\appbrk{}गु\appbrk{}रु\appbrk{}प\appbrk{}रे\appbrk{}ण\appbrk{}भ\appbrk{}वि\appbrk{}तव्यं}}\Afootnote{H; \skt{नीचनखरोम्णा शुचिना शुक्लवाससा शान्ति\appbrk{}म\appbrk{}ङ्ग\appbrk{}ल\appbrk{}दे\appbrk{}व\appbrk{}ता\appbrk{}ब्रा\appbrk{}ह्म\appbrk{}ण\appbrk{}गु\appbrk{}रु\appbrk{}प\appbrk{}रेण भवितव्यम्}\ A; \skt{नीचनखरोष्णा शुचिना शुचिवाससा शान्ति\appbrk{}म\appbrk{}ङ्ग\appbrk{}ल\appbrk{}दे\appbrk{}व\appbrk{}ता\appbrk{}ब्रा\appbrk{}ह्म\appbrk{}ण\appbrk{}गु\appbrk{}रु\appbrk{}प\appbrk{}रेण भवितव्यम्}\ N}}\skt{~॥ }\edtext{\skt{तत्कस्य}}{\lemma{\skt{तत्कस्य}}\Afootnote{N H; \skt{इति~। तत्कस्य}\ A}}\skt{ हेतोः }\edtext{\skt{हिंसाविहाराणि}}{\lemma{\skt{हिंसाविहाराणि}}\Afootnote{N H; \skt{हिंसाविहाराणि हि}\ A}}\skt{ }\edtext{\skt{तु}}{\lemma{\skt{तु}}\Afootnote{N H; \skt{महावीर्याणि}\ A}}\skt{ रक्षा\-न्सि\-प\-शु\-प\-ति\-कु\-बे\-र\-कु\-मा\-रा\-नु\-च\-राणि }\edtext{\edtext{\skt{मान्स\-शो\-णि\-त\-प्र\-त्वा\-त्क्ष\-त\-ज\-नि\-मि\-त्त\-म्व्र\-णि\-त\-मुप}}{\lemma{\skt{मान्स\appbrk{}शो\appbrk{}णि\appbrk{}त\appbrk{}प्र\appbrk{}त्वा\appbrk{}त्क्ष\appbrk{}त\appbrk{}ज\appbrk{}नि\appbrk{}मि\appbrk{}त्तम्}}\Afootnote{H; \skt{मांस\appbrk{}शो\appbrk{}णि\appbrk{}त\appbrk{}प्रि\appbrk{}य\appbrk{}त्वा\appbrk{}त्क्ष\appbrk{}त\appbrk{}ज\appbrk{}नि\appbrk{}मित्तं}\ A N}}}{\lemma{\skt{व्रणितम्}}\Afootnote{N H; \skt{व्रणिनम्}\ A}}\skt{ सर्प्पन्ति सत्का\-रा\-र्थ\-ञ्जि\-घां\-सूनि वा }\edtext{\skt{कदाचिद्भवन्ति}}{\lemma{\skt{कदाचिद्भवन्ति}}\Afootnote{N H; \skt{कदाचित्}\ A}}\skt{~॥}\pend

\pstart\versepar \ledsidenote{\passref{1.19.24}}\edtext{\skt{तेषां}}{\lemma{\skt{तेषां}}\Afootnote{N H; \skt{भवति चात्र~। तेषां}\ A}}\skt{ सत्का\-र\-का\-मा\-ना\-म्प्र\-य\-ते\-ता\-न्त\-रा\-त्मना~। }\pend
\pstart\versepar \skt{}\edtext{\skt{धूप\-मा\-ल्यो\-प\-हा\-राश्च}}{\lemma{\skt{धूप\appbrk{}मा\appbrk{}ल्यो\appbrk{}प\appbrk{}हा\appbrk{}राश्च}}\Afootnote{N H; \skt{धूप\appbrk{}ब\appbrk{}ल्यु\appbrk{}प\appbrk{}हा\appbrk{}रांश्च}\ A}}\skt{ }\edtext{\edtext{\skt{भक्षा\-श्चै\-वो\-प\-हा\-र\-येत्}}{\lemma{\skt{भक्षाश्}}\Afootnote{H; \skt{भक्ष्यां}\ A; \skt{भक्षां}\ N}}}{\lemma{\skt{चैवो\appbrk{}प\appbrk{}हा\appbrk{}र\appbrk{}येत्}}\Afootnote{N H; \skt{श्चैवो\appbrk{}प\appbrk{}हा\appbrk{}र\appbrk{}येत्}\ A}}\skt{~॥}\pend

\pstart \ledsidenote{\passref{1.19.25}}\skt{ते तु सन्तर्प्पिता }\edtext{\skt{आत्मवन्न}}{\lemma{\skt{आत्मवन्न}}\Afootnote{H; \skt{आत्मवन्तं न}\ A N}}\skt{ }\edtext{\edtext{\edtext{\skt{हिंस्यु\-स्त\-स्मा\-त्स\-त\-त\-म\-त\-न्द्रि\-त\-ज\-न\-प\-रि\-वृतो}}{\lemma{\skt{हिंस्युस्}}\Afootnote{A H; \skt{हि स्यु}\ N}}}{\lemma{\skt{तस्मात्सततम्}}\Afootnote{N H; \skt{तस्मात्सतत}\ A}}}{\lemma{\skt{अत\appbrk{}न्द्रि\appbrk{}त\appbrk{}ज\appbrk{}न\appbrk{}प\appbrk{}रि\appbrk{}वृतो}}\Afootnote{N H; \skt{मतन्द्रितो जनपरिवृतो}\ A}}\skt{ }\edtext{\skt{नित्य\-दी\-पो\-द\-क\-श\-स्त्र\-स्र\-ग्दा\-मा\-ल\-ङ्कृत}}{\lemma{\skt{नित्य\appbrk{}दी\appbrk{}पो\appbrk{}द\appbrk{}क\appbrk{}श\appbrk{}स्त्र\appbrk{}स्र\appbrk{}ग्दा\appbrk{}मा\appbrk{}ल\appbrk{}ङ्कृत}}\Afootnote{H; \skt{नित्यं दीपो\appbrk{}द\appbrk{}क\appbrk{}श\appbrk{}स्त्र\appbrk{}स्र\appbrk{}ग्दा\appbrk{}म\appbrk{}पु\appbrk{}ष्प\appbrk{}ला\appbrk{}जा\appbrk{}द्य\appbrk{}ल\appbrk{}ङ्कृते}\ A; \skt{नित्य\appbrk{}दी\appbrk{}पो\appbrk{}द\appbrk{}क\appbrk{}श\appbrk{}स्त्र\appbrk{}स्र\appbrk{}ग्दा\appbrk{}मा\appbrk{}न\appbrk{}ङ्कृत}\ N}}\skt{ वेश्मनि }\edtext{\edtext{\skt{सम्प\-न्म\-ङ्ग\-ल\-म\-नोऽ\-नु\-कूलाः}}{\lemma{\skt{सम्प\appbrk{}न्म\appbrk{}ङ्ग\appbrk{}ल\appbrk{}मनो}}\Afootnote{H; \skt{सम्प\appbrk{}न्म\appbrk{}ङ्ग\appbrk{}ल\appbrk{}मनो}\ A; \skt{सांप\appbrk{}न्म\appbrk{}ङ्ग\appbrk{}ल\appbrk{}मनो}\ N}}}{\lemma{\skt{ऽनुकूलाः}}\Afootnote{A H; \skt{नुकूलाः}\ N}}\skt{ कथाः }\edtext{\skt{शृण्वन्नासीत}}{\lemma{\skt{शृण्वन्नासीत}}\Afootnote{A H; \skt{शृण्वन्नासीत्}\ N}}\skt{~॥}\pend

\pstart\versepar \ledsidenote{\passref{1.19.26}}\edtext{\skt{सम्प\-त्म\-ङ्ग\-ल\-यु\-क्ताभिः}}{\lemma{\skt{सम्प\appbrk{}त्म\appbrk{}ङ्ग\appbrk{}ल\appbrk{}यु\appbrk{}क्ताभिः}}\Afootnote{H; \skt{सम्प\appbrk{}दा\appbrk{}द्य\appbrk{}नु\appbrk{}कू\appbrk{}लाभिः}\ A; \skt{सम्प\appbrk{}दा\appbrk{}द्भा\appbrk{}ल\appbrk{}युक्ता}\ N}}\skt{ }\edtext{\skt{कथाभिः}}{\lemma{\skt{कथाभिः}}\Afootnote{A H; \skt{भिः}\ N}}\skt{ प्रीतमानसः~। }\pend
\pstart\versepar \skt{}\edtext{\skt{आसा\-वा\-न्व्या\-धि\-मो\-क्षाय}}{\lemma{\skt{आसा\appbrk{}वा\appbrk{}न्व्या\appbrk{}धि\appbrk{}मो\appbrk{}क्षाय}}\Afootnote{H; \skt{आशा\appbrk{}वा\appbrk{}न्व्या\appbrk{}धि\appbrk{}मो\appbrk{}क्षाय}\ A; \skt{आषां\appbrk{}वा\appbrk{}व्या\appbrk{}धि\appbrk{}मो\appbrk{}क्षाय}\ N}}\skt{ क्षिप्रं }\edtext{\edtext{\skt{सुख\-म\-वा\-प्नु\-यात्}}{\lemma{\skt{सुखम्}}\Afootnote{A H; \skt{सुख}\ N}}}{\lemma{\skt{अवाप्नुयात्}}\Afootnote{A H; \skt{मवाप्नुयात्}\ N}}\skt{~॥}\pend

\pstart \ledsidenote{\passref{1.19.27}}\edtext{\edtext{\edtext{\edtext{\skt{ऋक्सा\-म\-य\-जु\-र्भि\-र्म्म\-न्त्रै\-र\-प\-रै\-श्चा\-शी\-र्व्वा\-दै\-रु\-पा\-ध्या\-य\-भि\-ष\-जाश्च}}{\lemma{\skt{ऋक्सा\appbrk{}म\appbrk{}य\appbrk{}जु\appbrk{}र्भिर्}}\Afootnote{H; \skt{ऋग्य\appbrk{}जुः\appbrk{}सा\appbrk{}मा\appbrk{}थ\appbrk{}र्व\appbrk{}वे\appbrk{}दाभि}\ A; \skt{ऋत्स्क्ता\appbrk{}म\appbrk{}य\appbrk{}जु\appbrk{}भिर्}\ N}}}{\lemma{\skt{म्मन्त्रैरपरैश्}}\Afootnote{H; \skt{हितैरपरैश्}\ A; \skt{मन्त्रैरपरैश्}\ N}}}{\lemma{\skt{चाशीर्व्वादैर्}}\Afootnote{N H; \skt{चाशीर्विधानैर्}\ A}}}{\lemma{\skt{उपा\appbrk{}ध्या\appbrk{}य\appbrk{}भि\appbrk{}ष\appbrk{}जाश्च}}\Afootnote{N H; \skt{उपाध्याया भिषजश्च}\ A}}\skt{ सन्ध्ययो रक्षाङ्कुर्युः~॥}\pend

\pstart\versepar \ledsidenote{\passref{1.19.28}}\skt{सर्ष\-पा\-रि\-ष्ट\-प\-त्राभ्यां सर्प्पि\-षा\-ल\-व\-णेन च~। }\pend
\pstart\versepar \skt{}\edtext{\skt{द्विरह्नः}}{\lemma{\skt{द्विरह्नः}}\Afootnote{A H; \skt{द्विरह्न}\ N}}\skt{ }\edtext{\skt{कारयेद्रूपं}}{\lemma{\skt{कारयेद्रूपं}}\Afootnote{H; \skt{कारयेद्धूपं}\ A N}}\skt{ }\edtext{\skt{सप्त\-रा\-त्र\-म\-त\-न्द्रितः}}{\lemma{\skt{सप्त\appbrk{}रा\appbrk{}त्र\appbrk{}म\appbrk{}त\appbrk{}न्द्रितः}}\Afootnote{H; \skt{दश\appbrk{}रा\appbrk{}त्र\appbrk{}म\appbrk{}त\appbrk{}न्द्रितः}\ A; \skt{सप्त\appbrk{}रा\appbrk{}त्र\appbrk{}म\appbrk{}त\appbrk{}न्द्रितः}\ N}}\skt{~॥}\pend

\pstart \ledsidenote{\passref{1.19.29}}\edtext{\edtext{\skt{छत्त्रा\-ति\-च्छ\-त्त्र\-ला\-ङ्गु\-ली\-ञ्ज\-टिलां}}{\lemma{\skt{छत्त्रा\appbrk{}ति\appbrk{}च्छ\appbrk{}त्त्र\appbrk{}ला\appbrk{}ङ्गु\appbrk{}लीञ्}}\Afootnote{N H; \skt{छत्रामतिच्छत्रां लाङ्गू(ङ्ग)लीं}\ A}}}{\lemma{\skt{जटिलां}}\Afootnote{A H; \skt{जटिला}\ N}}\skt{ ब्रह्मचारिणीं~। }\edtext{\skt{लक्ष्मी\-गु\-हा\-म\-ति\-गु\-हा\-व\-चा\-म\-ति\-वि\-षा\-न्तथा}}{\lemma{\skt{लक्ष्मी\appbrk{}गु\appbrk{}हा\appbrk{}म\appbrk{}ति\appbrk{}गु\appbrk{}हा\appbrk{}व\appbrk{}चा\appbrk{}म\appbrk{}ति\appbrk{}वि\appbrk{}षा\appbrk{}न्तथा}}\Afootnote{N H; \skt{लक्ष्मीं गुहामतिगुहां वचामतिविषां}\ A}}\skt{~॥ }\edtext{\skt{शतवीर्यां}}{\lemma{\skt{शतवीर्यां}}\Afootnote{A H; \skt{शतवीर्या}\ N}}\skt{ सह }\edtext{\skt{स्रवीर्यां}}{\lemma{\skt{स्रवीर्यां}}\Afootnote{A H; \skt{स्रवीर्या}\ N}}\skt{ }\edtext{\edtext{\skt{सिद्धार्थाश्चापि}}{\lemma{\skt{सिद्धार्थाश्}}\Afootnote{H; \skt{सिद्धार्थकांश्}\ A; \skt{सिद्धार्थाश्}\ N}}}{\lemma{\skt{चापि}}\Afootnote{H; \skt{च शि}\ A; \emph{om.}\ N}}\skt{ }\edtext{\skt{धारयेत्~॥ भवन्ति चात्र॥}}{\lemma{\skt{धारयेत्~॥ भवन्ति चात्र}}\Afootnote{\skt{रसा धारयेत्}\ A; \skt{चापि धारयेत्}\ N; \skt{धारयेत्~॥ भवति चात्र}\ H}}\pend

\pstart\versepar \ledsidenote{\passref{1.19.31}}\skt{अनेन }\edtext{\skt{विधाना}}{\lemma{\skt{विधाना}}\Afootnote{H; \skt{विधिना}\ A N}}\skt{ }\edtext{\edtext{\skt{युक्तमारादेव}}{\lemma{\skt{युक्तम्}}\Afootnote{A H; \skt{युक्ता}\ N}}}{\lemma{\skt{आरादेव}}\Afootnote{H; \skt{आदावेव}\ A; \skt{मारादेव}\ N}}\skt{ निशाचराः~। }\pend
\pstart\versepar \skt{वनं }\edtext{\skt{केसरिणाक्रान्तं}}{\lemma{\skt{केसरिणाक्रान्तं}}\Afootnote{H; \skt{केशरिणऽऽक्रान्तं}\ A; \skt{केशारिणाक्रान्त}\ N}}\skt{ वर्ज्जयन्ति मृगा इव~॥}\pend

\pstart\versepar \ledsidenote{\passref{1.19.33}}\edtext{\edtext{\skt{बालो\-शी\-रै\-र्व्र\-ण\-म्बी\-जैर्न्न}}{\lemma{\skt{बालोशीरैर्}}\Afootnote{H; \emph{om.}\ A N}}\skt{ च नम्प\-रि\-घ\-टू\-येत्}}{\lemma{\skt{व्रणम्बीजैर्न्न च नम्}}\Afootnote{H; \emph{om.}\ A N}}\skt{~। }\pend
\pstart\versepar \skt{न }\edtext{\skt{तुदेन्न च}}{\lemma{\skt{तुदेन्न च}}\Afootnote{H; \emph{om.}\ A N}}\skt{ }\edtext{\skt{कण्डूयाच्छयानः}}{\lemma{\skt{कण्डूयाच्छयानः}}\Afootnote{H; \emph{om.}\ A N}}\skt{ परिपालयेत्~॥ }\pend
\pstart\versepar \skt{जीर्णशाल्योदनं }\edtext{\edtext{\skt{स्निग्ध\-म\-ल्य\-मु\-ष्क\-न्द्र\-वो\-त्तरं}}{\lemma{\skt{स्निग्ध\appbrk{}म\appbrk{}ल्य\appbrk{}मु\appbrk{}ष्कन्}}\Afootnote{H; \emph{om.}\ A N}}}{\lemma{\skt{द्रवोत्तरं}}\Afootnote{H; \emph{om.}\ A N}}\skt{~। }\pend
\pstart\versepar \skt{}\edtext{\skt{भुञ्जानो}}{\lemma{\skt{भुञ्जानो}}\Afootnote{H; \emph{om.}\ A N}}\skt{ }\edtext{\edtext{\skt{जाङ्गलैर्म्मासैः}}{\lemma{\skt{जाङ्गलैर्}}\Afootnote{H; \emph{om.}\ A N}}}{\lemma{\skt{म्मासैः}}\Afootnote{H; \emph{om.}\ A N}}\skt{ }\edtext{\skt{शीघ्रं}}{\lemma{\skt{शीघ्रं}}\Afootnote{H; \emph{om.}\ A N}}\skt{ }\edtext{\edtext{\skt{व्रणमपोहति}}{\lemma{\skt{व्रणम्}}\Afootnote{H; \emph{om.}\ A N}}}{\lemma{\skt{अपोहति}}\Afootnote{H; \emph{om.}\ A N}}\skt{~॥ }\pend
\pstart\versepar \skt{}\edtext{\skt{तण्डु\-ली\-य\-क\-जी\-वन्ती}}{\lemma{\skt{तण्डु\appbrk{}ली\appbrk{}य\appbrk{}क\appbrk{}जी\appbrk{}वन्ती}}\Afootnote{A H; \emph{om.}\ N}}\skt{ }\edtext{\skt{सुनि\-ष\-र्ण्ण\-क\-वा\-स्तुकैः}}{\lemma{\skt{सुनि\appbrk{}ष\appbrk{}र्ण्ण\appbrk{}क\appbrk{}वा\appbrk{}स्तुकैः}}\Afootnote{H; \skt{सुनि\appbrk{}ष\appbrk{}ण्ण\appbrk{}क\appbrk{}वा\appbrk{}स्तुकैः}\ A; \emph{om.}\ N}}\skt{~। }\pend
\pstart\versepar \skt{}\edtext{\skt{बाल\-मू\-ल\-क\-वा\-र्त्ताकी}}{\lemma{\skt{बाल\appbrk{}मू\appbrk{}ल\appbrk{}क\appbrk{}वा\appbrk{}र्त्ताकी}}\Afootnote{N H; \skt{बाल\appbrk{}मू\appbrk{}ल\appbrk{}क\appbrk{}वा\appbrk{}र्ताक}\ A}}\skt{ पटोलैः }\edtext{\skt{कारवल्लकैः}}{\lemma{\skt{कारवल्लकैः}}\Afootnote{H; \skt{कारवेल्लकैः}\ A N}}\skt{~॥}\edtext{}{\lemma{\skt{परिघटूयेत्~। न}}\Afootnote{H; \emph{om.}\ A N}}\edtext{}{\lemma{\skt{परिपालयेत्~॥ जीर्णशाल्योदनं}}\Afootnote{H; \emph{om.}\ A N}}\pend

\pstart\versepar \ledsidenote{\passref{1.19.34}}\skt{सदाडिमैः }\edtext{\edtext{\skt{साम\-ल\-कै\-र्घृ\-त\-भृष्टैः}}{\lemma{\skt{सामलकैर्}}\Afootnote{N H; \skt{सामलकै}\ A}}}{\lemma{\skt{घृतभृष्टैः}}\Afootnote{H; \skt{र्घृतभृष्टैः}\ A; \skt{घृतपृष्टैः}\ N}}\skt{ }\edtext{\skt{ससैन्धवैःः}}{\lemma{\skt{ससैन्धवैःः}}\Afootnote{H; \skt{ससैन्धवैः}\ A; \skt{ससेन्धवैः}\ N}}\skt{~। }\pend
\pstart\versepar \skt{}\edtext{\skt{अन्यैरेव }\edtext{\skt{ङ्गुणैर्व्वापि}}{\lemma{\skt{व्वापि}}\Afootnote{A H; \skt{चापि}\ N}}}{\lemma{\skt{एव ङ्गुणैर्}}\Afootnote{A H; \skt{एव गुणैश्}\ N}}\skt{ }\edtext{\skt{मुद्गादीनां}}{\lemma{\skt{मुद्गादीनां}}\Afootnote{A H; \skt{शुङ्गहीनं}\ N}}\skt{ }\edtext{\skt{रसेन वा}}{\lemma{\skt{रसेन वा}}\Afootnote{N H; \skt{रसेन वा शक्तून्विलेपीं कुल्माषाञ्जलं चापि शृतं पिबेत्}\ A}}\skt{~॥}\pend

\pstart\versepar \ledsidenote{\passref{1.19.35}}\skt{दिवा न निद्रा वशगो निवा\-त\-गृ\-ह\-गो\-चरः~। }\pend
\pstart\versepar \skt{व्रणी }\edtext{\skt{वैद्यवसे}}{\lemma{\skt{वैद्यवसे}}\Afootnote{N H; \skt{वैद्यवशे}\ A}}\skt{ }\edtext{\skt{तिष्ठच्छीघ्रं}}{\lemma{\skt{तिष्ठच्छीघ्रं}}\Afootnote{H; \skt{तिष्ठञ्शीघ्रं}\ A; \skt{तिष्ठेच्छीघ्रं}\ N}}\skt{ }\edtext{\edtext{\skt{व्रणमपोहति}}{\lemma{\skt{व्रणम्}}\Afootnote{H; \skt{व्रण}\ A N}}}{\lemma{\skt{अपोहति}}\Afootnote{H; \skt{मपोहति}\ A N}}\skt{~॥}\pend

\pstart \ledsidenote{\passref{1.19.36}}\edtext{\skt{तौ च}}{\lemma{\skt{तौ च}}\Afootnote{N; \skt{(व्रणे श्वय\appbrk{}थु\appbrk{}रा\appbrk{}या\appbrk{}सात्स च रागश्च जागारात्~। तौ च}\ A}}\skt{ रुक्च }\edtext{\skt{दिवास्वापात्ताश्च}}{\lemma{\skt{दिवास्वापात्ताश्च}}\Afootnote{N; \skt{दिवास्वापात्ताश्च}\ A}}\skt{ मृत्युश्च मैथुनात्~॥}\pend

\pstart\versepar \ledsidenote{\passref{1.19.37}}\skt{एवं वृत्तसमाचारो व्रणी सम्पद्यते सुखी~। }\pend
\pstart\versepar \skt{आयुश्च }\edtext{\edtext{\skt{दीर्घमाप्नोति}}{\lemma{\skt{दीर्घम्}}\Afootnote{N H; \skt{दीर्घ}\ A}}}{\lemma{\skt{आप्नोति}}\Afootnote{N H; \skt{माप्नोति}\ A}}\skt{ }\edtext{\skt{धन्वन्तरिवचो}}{\lemma{\skt{धन्वन्तरिवचो}}\Afootnote{A N; \skt{धन्वतरिवचो}\ H}}\skt{ यथा~॥}\pend

\endnumbering
\beginnumbering
\eledsection*{Adhyāya 20}
\pstart \ledsidenote{\passref{1.20.1}}\edtext{\skt{॥ अथातो}}{\lemma{\skt{॥ अथातो}}\Afootnote{N H; \skt{अथतो}\ A}}\skt{ }\edtext{\skt{हिताहितीयं}}{\lemma{\skt{हिताहितीयं}}\Afootnote{H; \skt{हिता\appbrk{}हि\appbrk{}ती\appbrk{}य\appbrk{}म\appbrk{}ध्यायं}\ A; \skt{हिती\appbrk{}हि\appbrk{}ती\appbrk{}य\appbrk{}म\appbrk{}ध्यायं}\ N}}\skt{ व्याख्यास्यामः~॥}\pend

\pstart \ledsidenote{\passref{1.20.3}}\edtext{\skt{यद्वायोः}}{\lemma{\skt{यद्वायोः}}\Afootnote{N H; \skt{यद्वायोः}\ A}}\skt{ पथ्य\-न्त\-त्पि\-त्त\-स्या\-प\-थ्य\-मि\-त्य\-नेन हेतुना न }\edtext{\edtext{\skt{किञ्चि\-द्द्र\-व्य\-मे\-का\-न्तेन}}{\lemma{\skt{किञ्चिद्द्रव्यम्}}\Afootnote{N; \skt{किञ्चिद्द्रव्य}\ A; \skt{किञ्चि द्रव्यम्}\ H}}}{\lemma{\skt{एकान्तेन}}\Afootnote{N H; \skt{मेकान्तेन}\ A}}\skt{ }\edtext{\edtext{\skt{हित\-म\-हि\-त\-म्वा\-स्तीति}}{\lemma{\skt{हितम्}}\Afootnote{N H; \skt{हित}\ A}}}{\lemma{\skt{अहितम्}}\Afootnote{N H; \skt{महितं}\ A}}\skt{ }\edtext{\skt{केचिदाचार्या}}{\lemma{\skt{केचिदाचार्या}}\Afootnote{H; \skt{केचिदाचार्या}\ A; \skt{केचिदाचार्य}\ N}}\skt{ }\edtext{\skt{ब्रुवते~॥ तं तु न}}{\lemma{\skt{ब्रुवते~॥ तं तु न}}\Afootnote{N H; \skt{ब्रुवते~। तत्तु न}\ A}}\skt{ सम्यक्~। इह खलु }\edtext{\skt{द्रव्याणि}}{\lemma{\skt{द्रव्याणि}}\Afootnote{N H; \skt{यस्माद्द्रव्याणि}\ A}}\skt{ स्वभावतः }\edtext{\skt{संयोगतश्च}}{\lemma{\skt{संयोगतश्च}}\Afootnote{N H; \skt{संयोगतश्चै}\ A}}\skt{ }\edtext{\skt{एकान्तहितानि}}{\lemma{\skt{एकान्तहितानि}}\Afootnote{N H; \skt{कान्तहिता}\ A}}\skt{ }\edtext{\skt{एकान्ताहितानि}}{\lemma{\skt{एकान्ताहितानि}}\Afootnote{N H; \skt{न्येकान्ताहितानि}\ A}}\skt{ हिताहितानि च भवन्ति~॥}\pend

\pstart \ledsidenote{\passref{1.20.4}}\edtext{\skt{तत्रै\-का\-न्ता\-हि\-तानि}}{\lemma{\skt{तत्रै\appbrk{}का\appbrk{}न्ता\appbrk{}हि\appbrk{}तानि}}\Afootnote{H; \skt{तत्र एकान्तहितानि}\ A N}}\skt{ }\edtext{\skt{जातिसात्म्यत्वात्~। सलि\-ल\-घृ\-त\-दु\-ग्ध्वौ\-द\-न\-प्र\-भृ\-तीनि}}{\lemma{\skt{जातिसात्म्यत्वात्~। सलि\appbrk{}ल\appbrk{}घृ\appbrk{}त\appbrk{}दु\appbrk{}ग्ध्वौ\appbrk{}द\appbrk{}न\appbrk{}प्र\appbrk{}भृ\appbrk{}तीनि}}\Afootnote{H; \skt{जाति\appbrk{}सा\appbrk{}त्म्या\appbrk{}त्स\appbrk{}लि\appbrk{}ल\appbrk{}घृ\appbrk{}त\appbrk{}दु\appbrk{}ग्धौ\appbrk{}द\appbrk{}न\appbrk{}प्र\appbrk{}भृ\appbrk{}तीनि}\ A; \skt{जातिसात्म्यत्वात्~। सलि\appbrk{}ल\appbrk{}घृ\appbrk{}त\appbrk{}दु\appbrk{}ग्धो\appbrk{}द\appbrk{}न\appbrk{}प्र\appbrk{}भृ\appbrk{}तीनि}\ N}}\skt{~। एकान्ताहितानि तु दह\-न\-प\-च\-न\-मा\-र\-णा\-दिषु }\edtext{\edtext{\skt{प्रवृ\-त्त्ता\-न्य\-ग्नि\-क्षा\-र\-वि\-षा\-दीनि}}{\lemma{\skt{प्रवृत्त्तान्य्}}\Afootnote{N H; \skt{प्रवृत्तान्य}\ A}}}{\lemma{\skt{अग्नि\appbrk{}क्षा\appbrk{}र\appbrk{}वि\appbrk{}षा\appbrk{}दीनि}}\Afootnote{N H; \skt{नल\appbrk{}क्षा\appbrk{}र\appbrk{}वि\appbrk{}षा\appbrk{}दीनि}\ A}}\skt{ }\edtext{\edtext{\skt{संयो\-ग\-त\-स्त्व\-प\-राणि}}{\lemma{\skt{संयोगतस्}}\Afootnote{N H; \skt{संयोगाद}\ A}}}{\lemma{\skt{त्वपराणि}}\Afootnote{H; \skt{पराणि}\ A; \skt{मृपराणि}\ N}}\skt{ विषतुल्यानि भवन्ति~। हिताहितानि तु }\edtext{\skt{यद्वायोः}}{\lemma{\skt{यद्वायोः}}\Afootnote{N H; \skt{यद्वायोः}\ A}}\skt{ }\edtext{\skt{पथ्य\-न्त\-त्पि\-त्त\-स्या\-प\-थ्य\-मिति}}{\lemma{\skt{इति}}\Afootnote{A H; \skt{इत्य्}\ N}}\pend

\pstart \ledsidenote{\passref{1.20.5}}\edtext{\skt{एतद्द्वित्वान्न}}{\lemma{\skt{एतद्द्वित्वान्न}}\Afootnote{H; \skt{अतः}\ A; \skt{एतद्विर्वान}\ N}}\skt{ }\edtext{\edtext{\edtext{\skt{सर्व्व\-व्र\-णि\-ना\-म\-य\-मा\-हा\-रार्थे}}{\lemma{\skt{सर्व्वव्रणिनाम्}}\Afootnote{N; \skt{सर्वप्राणिना}\ A; \skt{सर्व्वव्रणिना}\ H}}}{\lemma{\skt{अयम्}}\Afootnote{N; \skt{मय}\ A H}}}{\lemma{\skt{आहारार्थे}}\Afootnote{\skt{माहारार्थं}\ A; \skt{आहार्थै}\ N; \skt{माहारार्थे}\ H}}\skt{ वर्ग्ग उपदिश्यते~। }\edtext{\skt{रक्त\-शा\-लि\-ष\-ष्टि\-का\-ङ्गु\-क\-मु\-कु\-न्द\-क\-पा\-ण्डु\-क\-क\-लम}}{\lemma{\skt{रक्त\appbrk{}शा\appbrk{}लि\appbrk{}ष\appbrk{}ष्टि\appbrk{}का\appbrk{}ङ्गु\appbrk{}क\appbrk{}मु\appbrk{}कु\appbrk{}न्द\appbrk{}क\appbrk{}पा\appbrk{}ण्डु\appbrk{}क\appbrk{}क\appbrk{}लम}}\Afootnote{H; \skt{तद्यथा- रक्त\appbrk{}शा\appbrk{}लि\appbrk{}ष\appbrk{}ष्टि\appbrk{}क\appbrk{}क\appbrk{}ङ्गु\appbrk{}क्र\appbrk{}मु\appbrk{}कु\appbrk{}न्द\appbrk{}क\appbrk{}पा\appbrk{}ण्डु\appbrk{}क\appbrk{}पी\appbrk{}त\appbrk{}क\appbrk{}प्र\appbrk{}मो\appbrk{}द\appbrk{}क\appbrk{}का\appbrk{}ल\appbrk{}का\appbrk{}स\appbrk{}न\appbrk{}क\appbrk{}पु\appbrk{}ष्प\appbrk{}क\appbrk{}क\appbrk{}र्द\appbrk{}म\appbrk{}क\appbrk{}श\appbrk{}कु\appbrk{}ना\appbrk{}हृ\appbrk{}त\appbrk{}सु\appbrk{}ग\appbrk{}न्ध\appbrk{}क\appbrk{}क\appbrk{}लम}\ A; \skt{रक्त\appbrk{}शा\appbrk{}लि\appbrk{}ष\appbrk{}ष्ठि\appbrk{}का\appbrk{}ङ्ग\appbrk{}क\appbrk{}मु\appbrk{}कु\appbrk{}न्द\appbrk{}क\appbrk{}पा\appbrk{}ण्डु\appbrk{}क\appbrk{}क\appbrk{}लम}\ N}}\skt{ }\edtext{\skt{नीवा\-रो\-द्र\-वो\-द्दा\-ल\-क\-श्या\-मा\-कवाः}}{\lemma{\skt{नीवा\appbrk{}रो\appbrk{}द्र\appbrk{}वो\appbrk{}द्दा\appbrk{}ल\appbrk{}क\appbrk{}श्या\appbrk{}मा\appbrk{}कवाः}}\Afootnote{H; \skt{नीवा\appbrk{}र\appbrk{}को\appbrk{}द्र\appbrk{}वो\appbrk{}द्दा\appbrk{}ल\appbrk{}क\appbrk{}श्या\appbrk{}मा\appbrk{}क\appbrk{}गो\appbrk{}धू\appbrk{}म\appbrk{}य\appbrk{}व\appbrk{}वै\appbrk{}णवै}\ A; \skt{नीवा\appbrk{}र\appbrk{}को\appbrk{}द्द्र\appbrk{}वो\appbrk{}द्दा\appbrk{}ल\appbrk{}क\appbrk{}स्या\appbrk{}मा\appbrk{}क\appbrk{}यवाः}\ N}}\skt{~॥ }\edtext{\skt{एण\-ह\-रि\-ण\-कु\-र\-ङ्गा\-मृ\-ग\-मा\-तृ\-का\-स्व\-दं\-ष्ट्री\-क्र\-क\-र\-ला\-व\-ति\-त्ति\-रि\-क\-पि\-ञ्ज\-ल\-व\-र्त्ती\-र\-क\-व\-र्त्तकाः}}{\lemma{\skt{एण\appbrk{}ह\appbrk{}रि\appbrk{}ण\appbrk{}कु\appbrk{}र\appbrk{}ङ्गा\appbrk{}मृ\appbrk{}ग\appbrk{}मा\appbrk{}तृ\appbrk{}का\appbrk{}स्व\appbrk{}दं\appbrk{}ष्ट्री\appbrk{}क्र\appbrk{}क\appbrk{}र\appbrk{}ला\appbrk{}व\appbrk{}ति\appbrk{}त्ति\appbrk{}रि\appbrk{}क\appbrk{}पि\appbrk{}ञ्ज\appbrk{}ल\appbrk{}व\appbrk{}र्त्ती\appbrk{}र\appbrk{}क\appbrk{}व\appbrk{}र्त्तकाः}}\Afootnote{H; \skt{णह\appbrk{}रि\appbrk{}ण\appbrk{}कु\appbrk{}र\appbrk{}ङ्ग\appbrk{}मृ\appbrk{}ग\appbrk{}मा\appbrk{}तृ\appbrk{}का\appbrk{}श्व\appbrk{}दं\appbrk{}ष्ट्रा\appbrk{}क\appbrk{}रा\appbrk{}ल\appbrk{}क्र\appbrk{}क\appbrk{}र\appbrk{}क\appbrk{}पो\appbrk{}त\appbrk{}ला\appbrk{}व\appbrk{}ति\appbrk{}त्ति\appbrk{}रि\appbrk{}क\appbrk{}पि\appbrk{}ञ्ज\appbrk{}ल\appbrk{}व\appbrk{}र्ती\appbrk{}र\appbrk{}व\appbrk{}र्तिका}\ A; \skt{एण\appbrk{}ह\appbrk{}रि\appbrk{}ण\appbrk{}कु\appbrk{}र\appbrk{}ङ्ग\appbrk{}मृ\appbrk{}ग\appbrk{}मा\appbrk{}तृ\appbrk{}काश्च दंष्ट्री\appbrk{}क्र\appbrk{}क\appbrk{}र\appbrk{}ला\appbrk{}व\appbrk{}ति\appbrk{}त्ति\appbrk{}रि\appbrk{}क\appbrk{}पि\appbrk{}ञ्ज\appbrk{}ल\appbrk{}व\appbrk{}र्त्ती\appbrk{}क\appbrk{}व\appbrk{}र्त्तकाः}\ N}}\skt{~॥ }\edtext{\skt{मुद्ग\-व\-न\-म\-सू\-र\-मु\-कु\-ष्ठ\-ह\-रे\-ण\-अ\-वा\-ढ\-की\-स\-तीनाः}}{\lemma{\skt{मुद्ग\appbrk{}व\appbrk{}न\appbrk{}म\appbrk{}सू\appbrk{}र\appbrk{}मु\appbrk{}कु\appbrk{}ष्ठ\appbrk{}ह\appbrk{}रे\appbrk{}ण\appbrk{}अ\appbrk{}वा\appbrk{}ढ\appbrk{}की\appbrk{}स\appbrk{}तीनाः}}\Afootnote{H; \skt{मुद्ग\appbrk{}व\appbrk{}न\appbrk{}मु\appbrk{}द्ग\appbrk{}म\appbrk{}कु\appbrk{}ष्ठ\appbrk{}क\appbrk{}ला\appbrk{}य\appbrk{}म\appbrk{}सू\appbrk{}र\appbrk{}म\appbrk{}ङ्ग\appbrk{}ल्य\appbrk{}च\appbrk{}ण\appbrk{}क\appbrk{}ह\appbrk{}रे\appbrk{}ण्वा\appbrk{}ढ\appbrk{}की\appbrk{}स\appbrk{}तीना}\ A; \skt{मुद्ग\appbrk{}व\appbrk{}न\appbrk{}मु\appbrk{}द्ग\appbrk{}म\appbrk{}शू\appbrk{}र\appbrk{}म\appbrk{}कु\appbrk{}ष्ट\appbrk{}ह\appbrk{}रे\appbrk{}ण्वा\appbrk{}ढ\appbrk{}की\appbrk{}श\appbrk{}तीनाः}\ N}}\skt{~॥ }\edtext{\skt{चिल्ली\-वा\-स्तू\-क\-सु\-नि\-ष\-र्ण्ण\-क\-जी\-वन्ती}}{\lemma{\skt{चिल्ली\appbrk{}वा\appbrk{}स्तू\appbrk{}क\appbrk{}सु\appbrk{}नि\appbrk{}ष\appbrk{}र्ण्ण\appbrk{}क\appbrk{}जी\appbrk{}वन्ती}}\Afootnote{H; \skt{श्चिल्लि\appbrk{}वा\appbrk{}स्तु\appbrk{}क\appbrk{}सु\appbrk{}नि\appbrk{}ष\appbrk{}ण्ण\appbrk{}क\appbrk{}जी\appbrk{}वन्ती}\ A; \skt{चिल्ली\appbrk{}वा\appbrk{}स्तू\appbrk{}क\appbrk{}सु\appbrk{}नि\appbrk{}ष\appbrk{}ण्ण\appbrk{}क\appbrk{}जी\appbrk{}वन्ती}\ N}}\skt{ }\edtext{\skt{तण्डु\-ली\-य\-क\-म\-ण्डू\-क\-पर्ण्णाः}}{\lemma{\skt{तण्डु\appbrk{}ली\appbrk{}य\appbrk{}क\appbrk{}म\appbrk{}ण्डू\appbrk{}क\appbrk{}पर्ण्णाः}}\Afootnote{H; \skt{तण्डु\appbrk{}ली\appbrk{}य\appbrk{}क\appbrk{}म\appbrk{}ण्डू\appbrk{}क\appbrk{}पर्ण्यः}\ A; \skt{तण्डु\appbrk{}ली\appbrk{}य\appbrk{}क\appbrk{}म\appbrk{}ण्डू\appbrk{}क\appbrk{}पर्ण्णिः}\ N}}\skt{~। गव्यं घृतं }\edtext{\skt{सैन्ध\-व\-दा\-डि\-मा\-म\-ल\-क\-मि\-त्येष}}{\lemma{\skt{सैन्ध\appbrk{}व\appbrk{}दा\appbrk{}डि\appbrk{}मा\appbrk{}म\appbrk{}ल\appbrk{}कम्}}\Afootnote{N H; \skt{सैन्धवं, दाडि\appbrk{}मा\appbrk{}म\appbrk{}ल\appbrk{}कम्}\ A}}\skt{ }\edtext{\skt{वर्ग्गः}}{\lemma{\skt{वर्ग्गः}}\Afootnote{A H; \skt{वर्गाः}\ N}}\skt{ }\edtext{\skt{सर्व्वव्रणिन्यः}}{\lemma{\skt{सर्व्वव्रणिन्यः}}\Afootnote{H; \skt{सर्वप्राणिनां}\ A; \skt{सर्वव्रणिभ्यः}\ N}}\skt{ सामान्यतः }\edtext{\skt{पथ्यतमः}}{\lemma{\skt{पथ्यतमः}}\Afootnote{A H; \skt{पथ्यतमाः}\ N}}\skt{~।}\pend

\pstart \ledsidenote{\passref{1.20.6}}\skt{तथा }\edtext{\skt{ब्रह्म\-च\-र्य\-नि\-वा\-त\-श\-य\-णो\-ष्णो\-द\-का\-दि\-वा\-स्व\-प्ना\-व्या\-या\-मा\-धू\-म\-सेवा च}}{\lemma{\skt{ब्रह्म\appbrk{}च\appbrk{}र्य\appbrk{}नि\appbrk{}वा\appbrk{}त\appbrk{}श\appbrk{}य\appbrk{}णो\appbrk{}ष्णो\appbrk{}द\appbrk{}का\appbrk{}दि\appbrk{}वा\appbrk{}स्व\appbrk{}प्ना\appbrk{}व्या\appbrk{}या\appbrk{}मा\appbrk{}धू\appbrk{}म\appbrk{}सेवा च}}\Afootnote{H; \skt{ब्रह्म\appbrk{}च\appbrk{}र्य\appbrk{}नि\appbrk{}वा\appbrk{}त\appbrk{}श\appbrk{}य\appbrk{}नो\appbrk{}ष्णो\appbrk{}द\appbrk{}क\appbrk{}स्ना\appbrk{}न\appbrk{}नि\appbrk{}शा\appbrk{}स्व\appbrk{}प्न\appbrk{}व्या\appbrk{}या\appbrk{}माश्चै}\ A; \skt{ब्रह्म\appbrk{}च\appbrk{}र्य\appbrk{}नि\appbrk{}वा\appbrk{}त\appbrk{}श\appbrk{}र\appbrk{}णो\appbrk{}ष्णो\appbrk{}द\appbrk{}का\appbrk{}दि\appbrk{}वा\appbrk{}स्व\appbrk{}प्ना\appbrk{}व्या\appbrk{}या\appbrk{}मा\appbrk{}धू\appbrk{}म\appbrk{}सेवा च}\ N}}\skt{ }\edtext{\skt{एकान्तः}}{\lemma{\skt{एकान्तः}}\Afootnote{N H; \skt{कान्ततः}\ A}}\skt{ }\edtext{\skt{पथ्यतमानि}}{\lemma{\skt{पथ्यतमानि}}\Afootnote{N H; \skt{पथ्यतमाः}\ A}}\skt{~।}\pend

\pstart \ledsidenote{\passref{1.20.7}}\edtext{\skt{अहितानि}}{\lemma{\skt{अहितानि}}\Afootnote{N H; \skt{एका\appbrk{}न्त\appbrk{}हि\appbrk{}ता\appbrk{}न्ये\appbrk{}का\appbrk{}न्ता\appbrk{}हि\appbrk{}तानि तु}\ A}}\skt{ प्रागुपदिष्टानि~। हिताहितानि तु }\edtext{\skt{यद्वायोः}}{\lemma{\skt{यद्वायोः}}\Afootnote{N; \skt{यद्वायोः}\ A H}}\skt{ पथ्य\-न्त\-त्पि\-त्त\-स्या\-प\-थ्य\-मिति~॥}\pend

\pstart \ledsidenote{\passref{1.20.8}}\edtext{\edtext{\skt{संयो\-ग\-त\-स्त्व\-प\-राणि}}{\lemma{\skt{संयोगतस्}}\Afootnote{N H; \skt{संयोगत}\ A}}}{\lemma{\skt{त्वपराणि}}\Afootnote{N H; \skt{स्त्वपराणि}\ A}}\skt{ विषतुल्यानि भवन्ति~। }\edtext{\skt{वल्ली\-फ\-ल\-ज\-व\-क\-क\-री\-र\-आ\-म्ल\-फ\-ल\-ल\-व\-ण\-कु\-ल\-त्थ\-पि\-ण्याक}}{\lemma{\skt{वल्ली\appbrk{}फ\appbrk{}ल\appbrk{}ज\appbrk{}व\appbrk{}क\appbrk{}क\appbrk{}री\appbrk{}र\appbrk{}आ\appbrk{}म्ल\appbrk{}फ\appbrk{}ल\appbrk{}ल\appbrk{}व\appbrk{}ण\appbrk{}कु\appbrk{}ल\appbrk{}त्थ\appbrk{}पि\appbrk{}ण्याक}}\Afootnote{H; \skt{तद्यथा- वल्ली\appbrk{}फ\appbrk{}ल\appbrk{}क\appbrk{}व\appbrk{}क\appbrk{}क\appbrk{}री\appbrk{}रा\appbrk{}म्ल\appbrk{}फ\appbrk{}ल\appbrk{}ल\appbrk{}व\appbrk{}ण\appbrk{}कु\appbrk{}ल\appbrk{}त्थ\appbrk{}पि\appbrk{}ण्याक}\ A; \skt{वल्ली\appbrk{}फ\appbrk{}ल\appbrk{}क\appbrk{}वक~। करीर~। आम्लफल~। लवण~। कुलत्थ~। पिण्या}\ N}}\skt{ }\edtext{\skt{तैल\-शु\-ष्क\-शा\-क\-म\-द्य\-जा\-म्व\-जा\-ङ्ग\-ल\-चि\-लि\-मि\-लि\-म\-त्स्य\-गो\-धा\-व\-रा\-हा\-न\-चै\-क\-ध्य\-म\-श्नी\-यात्~। पयसा}}{\lemma{\skt{तैल\appbrk{}शु\appbrk{}ष्क\appbrk{}शा\appbrk{}क\appbrk{}म\appbrk{}द्य\appbrk{}जा\appbrk{}म्व\appbrk{}जा\appbrk{}ङ्ग\appbrk{}ल\appbrk{}चि\appbrk{}लि\appbrk{}मि\appbrk{}लि\appbrk{}म\appbrk{}त्स्य\appbrk{}गो\appbrk{}धा\appbrk{}व\appbrk{}रा\appbrk{}हा\appbrk{}न\appbrk{}चै\appbrk{}क\appbrk{}ध्य\appbrk{}म\appbrk{}श्नी\appbrk{}यात्~। पयसा}}\Afootnote{\skt{दधि\appbrk{}तै\appbrk{}ल\appbrk{}वि\appbrk{}रो\appbrk{}हि\appbrk{}पि\appbrk{}ष्ट\appbrk{}शु\appbrk{}ष्क\appbrk{}शा\appbrk{}का\appbrk{}जा\appbrk{}वि\appbrk{}क\appbrk{}मां\appbrk{}स\appbrk{}म\appbrk{}द्य\appbrk{}जा\appbrk{}म्ब\appbrk{}व\appbrk{}चि\appbrk{}लि\appbrk{}चि\appbrk{}म\appbrk{}म\appbrk{}त्स्य\appbrk{}गो\appbrk{}धा\appbrk{}व\appbrk{}रा\appbrk{}हांश्च नैक\appbrk{}ध्य\appbrk{}म\appbrk{}श्नी\appbrk{}या\appbrk{}त्प\appbrk{}यसा}\ A; \skt{क~। तैल~। शुष्कशाक~। मद्यजाम्ब~। जङ्गल~। चिलिचिमिमत्स्य~। गोधा\appbrk{}व\appbrk{}रा\appbrk{}हा\appbrk{}न\appbrk{}चै\appbrk{}क\appbrk{}ध्य\appbrk{}म\appbrk{}श्नी\appbrk{}यात्~। पयसा}\ N; \skt{तैल\appbrk{}शु\appbrk{}ष्क\appbrk{}शा\appbrk{}क\appbrk{}म\appbrk{}द्य\appbrk{}जा\appbrk{}म्व\appbrk{}जा\appbrk{}ङ्ग\appbrk{}ल\appbrk{}चि\appbrk{}लि\appbrk{}मि\appbrk{}लि\appbrk{}म\appbrk{}त्स्य\appbrk{}गो\appbrk{}धा\appbrk{}व\appbrk{}रा\appbrk{}हा\appbrk{}न\appbrk{}चै\appbrk{}क\appbrk{}ध्य\appbrk{}म\appbrk{}श्नी\appbrk{}यात्~। पयसा}\ H}}\skt{ }\edtext{\edtext{\skt{प्राक्पयसोऽन्ते}}{\lemma{\skt{प्राक्पयसो}}\Afootnote{H; \emph{om.}\ A; \skt{प्रात्प्रूयसोन्}\ N}}\skt{ वा}}{\lemma{\skt{ऽन्ते वा}}\Afootnote{H; \emph{om.}\ A; \skt{ते वा}\ N}}\skt{ }\edtext{\skt{पयसः}}{\lemma{\skt{पयसः}}\Afootnote{N H; \emph{om.}\ A}}\skt{~।}\pend

\pstart \ledsidenote{\passref{1.20.14}}\edtext{\skt{कपोताश्च}}{\lemma{\skt{कपोताश्च}}\Afootnote{N H; \skt{अतः कर्म\appbrk{}वि\appbrk{}रु\appbrk{}द्धा\appbrk{}न्व\appbrk{}क्ष्यामः- कपोता}\ A}}\skt{ }\edtext{\skt{सर्ष\-प\-तै\-ल\-सि\-द्धाश्च}}{\lemma{\skt{सर्ष\appbrk{}प\appbrk{}तै\appbrk{}ल\appbrk{}सि\appbrk{}द्धाश्च}}\Afootnote{N H; \skt{न्सर्ष\appbrk{}प\appbrk{}तै\appbrk{}ल\appbrk{}भृ\appbrk{}ष्टान्}\ A}}\skt{ }\edtext{\skt{नाश्नीयात्~। }\edtext{\skt{कपि\-ञ्ज\-ल\-म\-यू\-र\-ला\-व\-ति\-त्ति\-रि\-गो\-धा\-श्चै\-र\-ण्ड\-का\-ष्ठ\-सिद्धा}}{\lemma{\skt{चैर\appbrk{}ण्ड\appbrk{}का\appbrk{}ष्ठ\appbrk{}सिद्धा}}\Afootnote{N H; \skt{श्चैर\appbrk{}ण्ड\appbrk{}दा\appbrk{}र्व्य\appbrk{}ग्नि\appbrk{}सिद्धा}\ A}}}{\lemma{\skt{नाश्नीयात्~। कपि\appbrk{}ञ्ज\appbrk{}ल\appbrk{}म\appbrk{}यू\appbrk{}र\appbrk{}ला\appbrk{}व\appbrk{}ति\appbrk{}त्ति\appbrk{}रि\appbrk{}गो\appbrk{}धाश्}}\Afootnote{H; \skt{नाद्यात्; कपि\appbrk{}ञ्ज\appbrk{}ल\appbrk{}म\appbrk{}यू\appbrk{}र\appbrk{}ला\appbrk{}व\appbrk{}ति\appbrk{}त्ति\appbrk{}रि\appbrk{}गोधा}\ A; \skt{नाश्नीयात्~। कपि\appbrk{}ञ्ज\appbrk{}ल\appbrk{}म\appbrk{}पू\appbrk{}र\appbrk{}ला\appbrk{}व\appbrk{}ति\appbrk{}त्ति\appbrk{}रि\appbrk{}गो\appbrk{}धाश्}\ N}}\skt{ }\edtext{\skt{एरण्डतैलन}}{\lemma{\skt{एरण्डतैलन}}\Afootnote{H; \skt{एरण्डतैल}\ A; \skt{एरण्डतैलेन}\ N}}\skt{ }\edtext{\skt{नाद्यात्~। मधु\-घृ\-त\-स\-म\-घृ\-त\-म\-धु\-चा\-न्त\-रि\-क्षौ\-द\-का\-नु\-पानं}}{\lemma{\skt{नाद्यात्~। मधु\appbrk{}घृ\appbrk{}त\appbrk{}स\appbrk{}म\appbrk{}घृ\appbrk{}त\appbrk{}म\appbrk{}धु\appbrk{}चा\appbrk{}न्त\appbrk{}रि\appbrk{}क्षौ\appbrk{}द\appbrk{}का\appbrk{}नु\appbrk{}पानं}}\Afootnote{H; \skt{सिद्धा वा नाद्या}\ A; \skt{नाद्यात्~। मधुघृते समधृते मधु\appbrk{}चा\appbrk{}न्त\appbrk{}रि\appbrk{}क्षो\appbrk{}द\appbrk{}पानं}\ N}}\skt{~। }\edtext{\skt{कांसभाजने}}{\lemma{\skt{कांसभाजने}}\Afootnote{N H; \skt{त्; कांस्यभाजने}\ A}}\skt{ दश\-रा\-त्र\-प\-र्यु\-षितं सर्प्पिः~। }\edtext{\skt{मधुनोष्णेन वा}}{\lemma{\skt{मधुनोष्णेन वा}}\Afootnote{H; \skt{मधु चोष्णैर्}\ A; \skt{मधुनोम्लेन वा}\ N}}\skt{ }\edtext{\skt{दधिमद्यञ्च}}{\lemma{\skt{दधिमद्यञ्च}}\Afootnote{N H; \emph{om.}\ A}}\skt{~। }\edtext{\skt{पित्तेन}}{\lemma{\skt{पित्तेन}}\Afootnote{N H; \emph{om.}\ A}}\skt{ }\edtext{\skt{चाममांसानि}}{\lemma{\skt{चाममांसानि}}\Afootnote{N H; \emph{om.}\ A}}\skt{~। }\edtext{\skt{सुरा\-कृ\-स\-र\-पा\-य\-साश्च}}{\lemma{\skt{सुरा\appbrk{}कृ\appbrk{}स\appbrk{}र\appbrk{}पा\appbrk{}य\appbrk{}साश्च}}\Afootnote{H; \skt{उष्णे}\ A; \skt{सुर\appbrk{}कृ\appbrk{}स\appbrk{}रा\appbrk{}या\appbrk{}पा\appbrk{}य\appbrk{}सांश्च}\ N}}\skt{ }\edtext{\skt{नैक\-ध्य\-म\-श्नी\-यात्~। सौवीरेण}}{\lemma{\skt{नैक\appbrk{}ध्य\appbrk{}म\appbrk{}श्नी\appbrk{}यात्~। सौवीरेण}}\Afootnote{H; \emph{om.}\ A; \skt{नैके\appbrk{}ध्य\appbrk{}म\appbrk{}श्नी\appbrk{}यात्~। सौवीरेण}\ N}}\skt{ }\edtext{\skt{सह\-ति\-ल\-स\-ष्कुली}}{\lemma{\skt{सह\appbrk{}ति\appbrk{}ल\appbrk{}स\appbrk{}ष्कुली}}\Afootnote{H; \emph{om.}\ A; \skt{सह तिल संस्कुली}\ N}}\skt{~। }\edtext{\skt{तक्रेण}}{\lemma{\skt{तक्रेण}}\Afootnote{N H; \emph{om.}\ A}}\skt{ }\edtext{\skt{मधु\-घृ\-त\-धा\-ना\-पृ\-ष\-त\-मां\-सानि}}{\lemma{\skt{मधु\appbrk{}घृ\appbrk{}त\appbrk{}धा\appbrk{}ना\appbrk{}पृ\appbrk{}ष\appbrk{}त\appbrk{}मां\appbrk{}सानि}}\Afootnote{N H; \emph{om.}\ A}}\skt{~। }\edtext{\skt{गोधाम्मधुना}}{\lemma{\skt{गोधाम्मधुना}}\Afootnote{H; \emph{om.}\ A; \skt{शोधाम्मधुना}\ N}}\skt{~। }\edtext{\skt{क्षौद्रासवेन}}{\lemma{\skt{क्षौद्रासवेन}}\Afootnote{N H; \emph{om.}\ A}}\skt{ }\edtext{\skt{मत्स्यां}}{\lemma{\skt{मत्स्यां}}\Afootnote{N H; \emph{om.}\ A}}\skt{~। }\edtext{\skt{पृषतमांसम्वा}}{\lemma{\skt{पृषतमांसम्वा}}\Afootnote{N H; \skt{वा}\ A}}\skt{ }\edtext{\skt{मैरे\-य\-मा\-ध्वी\-काभ्यां}}{\lemma{\skt{मैरे\appbrk{}य\appbrk{}मा\appbrk{}ध्वी\appbrk{}काभ्यां}}\Afootnote{N H; \skt{मत्स्यपरिप}\ A}}\skt{~। }\edtext{\skt{मत्स्या\-नु\-पो\-त\-कया}}{\lemma{\skt{मत्स्या\appbrk{}नु\appbrk{}पो\appbrk{}त\appbrk{}कया}}\Afootnote{H; \skt{चने शृङ्गवेरपरि}\ A; \skt{मत्स्या\appbrk{}तु\appbrk{}पो\appbrk{}त\appbrk{}कया}\ N}}\skt{ }\edtext{\skt{सहै\-क्षु\-वि\-कृ\-तीश्च}}{\lemma{\skt{सहै\appbrk{}क्षु\appbrk{}वि\appbrk{}कृ\appbrk{}तीश्च}}\Afootnote{N H; \skt{पच}\ A}}\skt{ }\edtext{\skt{सर्व्वाः}}{\lemma{\skt{सर्व्वाः}}\Afootnote{H; \skt{ने वा}\ A; \skt{सन्धाः}\ N}}\skt{~। }\edtext{\skt{गुड\-ङ्का\-क\-मा\-च्या\-म\-धुना}}{\lemma{\skt{गुड\appbrk{}ङ्का\appbrk{}क\appbrk{}मा\appbrk{}च्या\appbrk{}म\appbrk{}धुना}}\Afootnote{H; \skt{सिद्धां काकमाचीं; तिल\appbrk{}क\appbrk{}ल्क\appbrk{}सि\appbrk{}द्धमु}\ A; \skt{गुड\appbrk{}का\appbrk{}क\appbrk{}शा\appbrk{}च्या\appbrk{}म\appbrk{}धुना}\ N}}\skt{ }\edtext{\skt{मूलकानि च}}{\lemma{\skt{मूलकानि च}}\Afootnote{N H; \skt{पोदिकाशाकं}\ A}}\skt{~। }\edtext{\skt{नैकध्यमधुना}}{\lemma{\skt{नैकध्यमधुना}}\Afootnote{N H; \skt{नारिके}\ A}}\skt{ }\edtext{\skt{वराहं}}{\lemma{\skt{वराहं}}\Afootnote{H; \skt{लेन वराह}\ A; \skt{वशहं}\ N}}\skt{~। }\edtext{\skt{दध्ना}}{\lemma{\skt{दध्ना}}\Afootnote{N H; \skt{वसा}\ A}}\skt{ }\edtext{\skt{कुक्कुटं}}{\lemma{\skt{कुक्कुटं}}\Afootnote{H; \emph{om.}\ A; \skt{कुर्कुढं}\ N}}\skt{~। }\edtext{\skt{मद्येन}}{\lemma{\skt{मद्येन}}\Afootnote{N H; \skt{परिभृ}\ A}}\skt{ }\edtext{\skt{बलाकां}}{\lemma{\skt{बलाकां}}\Afootnote{H; \skt{ष्टां बलाकां; भास\appbrk{}म\appbrk{}ङ्गा\appbrk{}र\appbrk{}शूल्यं नाश्नीयादिति}\ A; \skt{वलाकं}\ N}}\skt{~।}\pend

\pstart \ledsidenote{\passref{1.20.17}}\edtext{\skt{तर\-त\-म\-यो\-ग\-यु\-क्तांश्च}}{\lemma{\skt{तर\appbrk{}त\appbrk{}म\appbrk{}यो\appbrk{}ग\appbrk{}यु\appbrk{}क्तांश्च}}\Afootnote{N H; \skt{तर\appbrk{}त\appbrk{}म\appbrk{}यो\appbrk{}ग\appbrk{}यु\appbrk{}क्तांश्च}\ A}}\skt{~। }\edtext{\edtext{\edtext{\edtext{\edtext{\skt{भावा\-न\-ति\-रू\-क्ष\-म\-ति\-स्नि\-ग्ध\-म\-त्यु\-ष्ण\-म\-ति\-शी\-त\-मे\-व\-मा\-दी\-न्वि\-व\-र्ज्ज\-येत्}}{\lemma{\skt{भावा\appbrk{}न\appbrk{}ति\appbrk{}रू\appbrk{}क्षम्}}\Afootnote{H; \skt{भावानतिस्निग्धा}\ A; \skt{भावा\appbrk{}न\appbrk{}ति\appbrk{}स्ल\appbrk{}क्षम्}\ N}}}{\lemma{\skt{अतिस्निग्धम्}}\Afootnote{N H; \skt{नतिरूक्षा}\ A}}}{\lemma{\skt{अत्युष्णम्}}\Afootnote{N H; \skt{नत्युष्णा}\ A}}}{\lemma{\skt{अतिशीतम्}}\Afootnote{N H; \skt{नतिशीतानित्य्}\ A}}\skt{~। न}}{\lemma{\skt{एव\appbrk{}मा\appbrk{}दी\appbrk{}न्वि\appbrk{}व\appbrk{}र्ज्ज\appbrk{}येत्~। न}}\Afootnote{N H; \skt{एव\appbrk{}मा\appbrk{}दी\appbrk{}न्वि\appbrk{}व\appbrk{}र्ज\appbrk{}येत्}\ A}}\skt{ }\edtext{\skt{मधुनोष्णं}}{\lemma{\skt{मधुनोष्णं}}\Afootnote{H; \emph{om.}\ A; \skt{मधुना उष्णं}\ N}}\skt{~। }\edtext{\skt{नोष्ण}}{\lemma{\skt{नोष्ण}}\Afootnote{H; \emph{om.}\ A; \skt{नोष्णो}\ N}}\skt{ }\edtext{\skt{नोष्णार्त्तो~। न}}{\lemma{\skt{नोष्णार्त्तो~। न}}\Afootnote{N H; \emph{om.}\ A}}\skt{ }\edtext{\skt{मधु\-ना\-ति\-ल\-क\-ल्के\-नो\-पो\-तकां~। न}}{\lemma{\skt{मधु\appbrk{}ना\appbrk{}ति\appbrk{}ल\appbrk{}क\appbrk{}ल्के\appbrk{}नो\appbrk{}पो\appbrk{}तकां~। न}}\Afootnote{H; \emph{om.}\ A; \skt{मधुना तिल\appbrk{}क\appbrk{}न्के\appbrk{}नौ\appbrk{}पे\appbrk{}तकां~॥ न}\ N}}\skt{ }\edtext{\skt{पिप्प\-ली\-म\-त्स्य\-वे\-सया}}{\lemma{\skt{पिप्प\appbrk{}ली\appbrk{}म\appbrk{}त्स्य\appbrk{}वे\appbrk{}सया}}\Afootnote{H; \emph{om.}\ A; \skt{पिप्प\appbrk{}ली\appbrk{}म\appbrk{}त्स्य\appbrk{}व\appbrk{}श\appbrk{}सया}\ N}}\skt{~। }\edtext{\skt{प्रिय\-ङ्ग्व\-नु\-लि\-प्तेन}}{\lemma{\skt{प्रिय\appbrk{}ङ्ग्व\appbrk{}नु\appbrk{}लि\appbrk{}प्तेन}}\Afootnote{H; \emph{om.}\ A; \skt{प्रियं\appbrk{}श्चा\appbrk{}नु\appbrk{}लि\appbrk{}प्तेन}\ N}}\skt{ }\edtext{\skt{पाय\-स\-म\-श्नी\-यात्}}{\lemma{\skt{पाय\appbrk{}स\appbrk{}म\appbrk{}श्नी\appbrk{}यात्}}\Afootnote{H; \emph{om.}\ A; \skt{पाय\appbrk{}स\appbrk{}म\appbrk{}श्नी\appbrk{}यात्}\ N}}\skt{~॥}\pend

\pstart\versepar \ledsidenote{\passref{1.20.18}}\edtext{\skt{भवन्ति}}{\lemma{\skt{भवन्ति}}\Afootnote{A N; \emph{om.}\ H}}\skt{ }\edtext{\skt{चात्र॥}}{\lemma{\skt{चात्र}}\Afootnote{A N; \emph{om.}\ H}}\skt{ }\pend
\pstart\versepar \skt{विरु\-द्धा\-न्ये\-व\-मा\-दीनि वीर्यतो यानि कानिचित्~। }\pend
\pstart\versepar \skt{}\edtext{\skt{तान्येकान्तान्येव}}{\lemma{\skt{एकान्तान्य्}}\Afootnote{H; \skt{एकान्ताहितान्य्}\ A N}}\skt{ }\edtext{\skt{शेष\-म्वि\-द्या\-धि\-ता\-हितं}}{\lemma{\skt{विद्याधिताहितं}}\Afootnote{H; \skt{विद्या\appbrk{}द्धि\appbrk{}ता\appbrk{}हि\appbrk{}तम्}\ A N}}\skt{~॥ }\edtext{\skt{भवन्ति चात्र॥}}{\lemma{\skt{भवन्ति चात्र}}\Afootnote{\emph{om.}\ A N H}}\pend

\pstart\versepar \ledsidenote{\passref{1.20.19}}\edtext{\edtext{\skt{व्याधि\-मि\-न्द्रि\-य\-दौ\-र्ब्ब\-ल्य\-म्म\-रण्वा}}{\lemma{\skt{व्याधि\appbrk{}मि\appbrk{}न्द्रि\appbrk{}य\appbrk{}दौ\appbrk{}र्ब्ब\appbrk{}ल्यम्}}\Afootnote{A H; \skt{व्याधि\appbrk{}मि\appbrk{}न्द्रि\appbrk{}य\appbrk{}दौ\appbrk{}र्बल्यं}\ N}}}{\lemma{\skt{मरण्वा}}\Afootnote{H; \skt{मरणं}\ A N}}\skt{ }\edtext{\skt{नियच्छति}}{\lemma{\skt{नियच्छति}}\Afootnote{H; \skt{चाधिगच्छति}\ A; \skt{वा नियच्छति}\ N}}\skt{~। }\pend
\pstart\versepar \skt{विरु\-द्ध\-र\-स\-वी\-र्याणि }\edtext{\skt{भुञ्जा\-नोऽ\-ना\-त्म\-वा\-न्नरः}}{\lemma{\skt{ऽनात्मवान्}}\Afootnote{A H; \skt{नात्मभवान्}\ N}}\skt{~॥}\pend

\pstart\versepar \ledsidenote{\passref{1.20.20}}\edtext{\skt{यत्कि\-ञ्चि\-द्दो\-ष\-मु\-त्क्लेश्य}}{\lemma{\skt{यत्कि\appbrk{}ञ्चि\appbrk{}द्दो\appbrk{}ष\appbrk{}मु\appbrk{}त्क्लेश्य}}\Afootnote{\skt{यत्कि\appbrk{}ञ्चि\appbrk{}द्दो\appbrk{}ष\appbrk{}मु\appbrk{}त्क्लेश्य}\ A; \skt{यत्किं\appbrk{}चि\appbrk{}द्दो\appbrk{}ष\appbrk{}मु\appbrk{}त्क्लेश्यं}\ N; \skt{यत्कि\appbrk{}ञ्चि\appbrk{}द्दो\appbrk{}ष\appbrk{}मु\appbrk{}त्क्लेश्य}\ H}}\skt{ भुक्तं कायान्न निर्हरेत्~। }\pend
\pstart\versepar \skt{रसादिषु }\edtext{\skt{रसा\-र्थ\-त्वा\-त्त\-द्वि\-का\-राय}}{\lemma{\skt{रसा\appbrk{}र्थ\appbrk{}त्वा\appbrk{}त्त\appbrk{}द्वि\appbrk{}का\appbrk{}राय}}\Afootnote{H; \skt{ष्वयथार्थं वा तद्विकाराय}\ A; \skt{रसादिषु रसा\appbrk{}र्थ\appbrk{}त्वा\appbrk{}त्त\appbrk{}द्वि\appbrk{}का\appbrk{}राय}\ N}}\skt{ कल्पते~॥}\edtext{}{\lemma{\skt{निर्हरेत्~। रसादिषु}}\Afootnote{N H; \skt{निर्हरेत्~। रसादि}\ A}}\pend

\pstart\versepar \ledsidenote{\passref{1.20.21}}\edtext{\skt{विरु\-द्धा\-श\-न\-जा\-न्रो\-गा\-न्प्र\-ति\-हन्ति}}{\lemma{\skt{विरु\appbrk{}द्धा\appbrk{}श\appbrk{}न\appbrk{}जा\appbrk{}न्रो\appbrk{}गा\appbrk{}न्प्र\appbrk{}ति\appbrk{}हन्ति}}\Afootnote{A H; \skt{विरु\appbrk{}द्धा\appbrk{}स\appbrk{}न\appbrk{}जा\appbrk{}न्रोगां प्रतिहन्ति}\ N}}\skt{ विरेचनं~। }\pend
\pstart\versepar \skt{वमनं }\edtext{\edtext{\skt{शमनम्वापि}}{\lemma{\skt{शमनम्}}\Afootnote{A H; \skt{शमनश्}\ N}}}{\lemma{\skt{वापि}}\Afootnote{A H; \skt{चापि}\ N}}\skt{ पूर्व्वम्वा हितसेवनं~॥}\pend

\pstart\versepar \ledsidenote{\passref{1.20.22}}\edtext{\edtext{\skt{सात्म्यतोऽल्पतया}}{\lemma{\skt{सात्म्यतो}}\Afootnote{A H; \skt{सात्य्म्यतो}\ N}}}{\lemma{\skt{ऽल्पतया}}\Afootnote{A H; \skt{ल्पतया}\ N}}\skt{ वापि }\edtext{\edtext{\skt{दीप्ता\-ग्ने\-स्त\-रु\-णस्य}}{\lemma{\skt{दीप्ताग्नेस्}}\Afootnote{H; \skt{दीप्ताग्ने}\ A N}}\skt{ च}}{\lemma{\skt{तरुणस्य च}}\Afootnote{N H; \skt{स्तरुणस्य च}\ A}}\skt{~। }\pend
\pstart\versepar \skt{}\edtext{\edtext{\skt{स्निग्धव्यायामिव}}{\lemma{\skt{स्निग्धव्यायाम्}}\Afootnote{H; \skt{स्निग्धव्याया}\ A; \skt{स्निग्धाव्याया}\ N}}}{\lemma{\skt{इव}}\Afootnote{H; \skt{मब}\ A; \skt{मव}\ N}}\skt{ }\edtext{\skt{लिना\-म्वि\-रु\-द्ध\-म्वि\-त\-थ\-म्भ\-वे\-दिति}}{\lemma{\skt{भवेदिति}}\Afootnote{N H; \skt{भवेत्}\ A}}\skt{~॥}\pend

\pstart\versepar \ledsidenote{\passref{1.20.22.add}}\edtext{\edtext{\edtext{\skt{क्षार\-म\-ग्नि\-ञ्ज\-ला\-युश्च}}{\lemma{\skt{क्षारम्}}\Afootnote{N\textsuperscript{pc} H\textsuperscript{pc}; \emph{om.}\ N\textsuperscript{ac} H\textsuperscript{ac}}}}{\lemma{\skt{अग्निञ्}}\Afootnote{H\textsuperscript{pc}; \emph{om.}\ N\textsuperscript{ac} H\textsuperscript{ac}; \skt{अग्नि}\ N\textsuperscript{pc}}}}{\lemma{\skt{जलायुश्च}}\Afootnote{N\textsuperscript{pc} H\textsuperscript{pc}; \emph{om.}\ N\textsuperscript{ac} H\textsuperscript{ac}}}\skt{ }\edtext{\skt{तथा}}{\lemma{\skt{तथा}}\Afootnote{N\textsuperscript{pc} H\textsuperscript{pc}; \emph{om.}\ N\textsuperscript{ac} H\textsuperscript{ac}}}\skt{ }\edtext{\skt{शोणितवर्ण्णनं}}{\lemma{\skt{शोणितवर्ण्णनं}}\Afootnote{H\textsuperscript{pc}; \emph{om.}\ N\textsuperscript{ac} H\textsuperscript{ac}; \skt{सोतवर्ण्णनं}\ N\textsuperscript{pc}}}\skt{~। }\pend
\pstart\versepar \skt{}\edtext{\edtext{\skt{दोष\-धा\-तु\-म\-ल\-ञ्चैव}}{\lemma{\skt{दोष\appbrk{}धा\appbrk{}तु\appbrk{}म\appbrk{}लञ्}}\Afootnote{H\textsuperscript{pc}; \emph{om.}\ N\textsuperscript{ac} H\textsuperscript{ac}; \skt{दोष\appbrk{}धा\appbrk{}तु\appbrk{}म\appbrk{}लश्}\ N\textsuperscript{pc}}}}{\lemma{\skt{चैव}}\Afootnote{N\textsuperscript{pc} H\textsuperscript{pc}; \emph{om.}\ N\textsuperscript{ac} H\textsuperscript{ac}}}\skt{ }\edtext{\edtext{\skt{कर्ण्ण\-ता\-ड\-न\-मेव}}{\lemma{\skt{कर्ण्णताडनम्}}\Afootnote{N\textsuperscript{pc} H\textsuperscript{pc}; \emph{om.}\ N\textsuperscript{ac} H\textsuperscript{ac}}}\skt{ च}}{\lemma{\skt{एव च}}\Afootnote{N\textsuperscript{pc} H\textsuperscript{pc}; \emph{om.}\ N\textsuperscript{ac} H\textsuperscript{ac}}}\skt{~॥ }\pend
\pstart\versepar \skt{}\edtext{\skt{आम\-प\-क्वे\-ष\-णी\-यञ्च}}{\lemma{\skt{आम\appbrk{}प\appbrk{}क्वे\appbrk{}ष\appbrk{}णी\appbrk{}यञ्च}}\Afootnote{H\textsuperscript{pc}; \emph{om.}\ N\textsuperscript{ac} H\textsuperscript{ac}; \skt{आस\appbrk{}प\appbrk{}न्के\appbrk{}ष\appbrk{}णी\appbrk{}यञ्च}\ N\textsuperscript{pc}}}\skt{ }\edtext{\edtext{\skt{आले\-प\-न\-वि\-धि\-न्तथा}}{\lemma{\skt{आले\appbrk{}प\appbrk{}न\appbrk{}वि\appbrk{}धिन्}}\Afootnote{N\textsuperscript{pc} H\textsuperscript{pc}; \emph{om.}\ N\textsuperscript{ac} H\textsuperscript{ac}}}}{\lemma{\skt{तथा}}\Afootnote{N\textsuperscript{pc} H\textsuperscript{pc}; \emph{om.}\ N\textsuperscript{ac} H\textsuperscript{ac}}}\skt{~। }\pend
\pstart\versepar \skt{}\edtext{\skt{व्रणि\-तो\-पा\-स\-नी\-यञ्च}}{\lemma{\skt{व्रणि\appbrk{}तो\appbrk{}पा\appbrk{}स\appbrk{}नी\appbrk{}यञ्च}}\Afootnote{H\textsuperscript{pc}; \emph{om.}\ N\textsuperscript{ac} H\textsuperscript{ac}; \skt{णितो\appbrk{}पा\appbrk{}स\appbrk{}नी\appbrk{}यञ्च}\ N\textsuperscript{pc}}}\skt{ }\edtext{\skt{विरुद्धान्नेन}}{\lemma{\skt{विरुद्धान्नेन}}\Afootnote{H\textsuperscript{pc}; \emph{om.}\ N\textsuperscript{ac} H\textsuperscript{ac}; \skt{विरुद्धान्नेव}\ N\textsuperscript{pc}}}\skt{ }\edtext{\skt{विंशतिः}}{\lemma{\skt{विंशतिः}}\Afootnote{N\textsuperscript{pc} H\textsuperscript{pc}; \emph{om.}\ N\textsuperscript{ac} H\textsuperscript{ac}}}\skt{~॥}\pend

\pstart \ledsidenote{\passref{1.20.trailer}}\edtext{\skt{॥ २०~॥}}{\lemma{\skt{॥ २०~॥}}\Afootnote{\skt{इति सुश्रु\appbrk{}त\appbrk{}सं\appbrk{}हि\appbrk{}तायां सूत्रस्थाने हिताहितीयो नाम विंशोऽध्यायः}\ A; \emph{om.}\ N\textsuperscript{ac}; \skt{॥ द्वितीयो दशः~॥}\ N\textsuperscript{pc}; \skt{द्वितीयो दश~॥ थं}\ H}}\pend

\endnumbering
\beginnumbering
\eledsection*{Adhyāya 21}
\pstart \ledsidenote{\passref{1.21.1}}\skt{॥ अथातो व्रणप्रश्नं }\edtext{\skt{व्याख्यास्यामः}}{\lemma{\skt{व्याख्यास्यामः}}\Afootnote{H; \skt{अध्यायं व्याख्यास्यामः}\ A N}}\skt{~॥}\pend

\pstart \ledsidenote{\passref{1.21.3}}\skt{वात\-पि\-त्त\-श्ले\-ष्माण एव }\edtext{\skt{देहे}}{\lemma{\skt{देहे}}\Afootnote{N H; \skt{देह}\ A}}\skt{ सम्भवहेतवो }\edtext{\skt{भवन्ति~॥ }\edtext{\edtext{\skt{तैरे\-वा\-स्या\-व्या\-प\-नै\-र\-धो\-म\-ध्यो\-र्द्ध\-स\-न्नि\-विष्टैः}}{\lemma{\skt{एवा\appbrk{}स्या\appbrk{}व्या\appbrk{}प\appbrk{}नैर्}}\Afootnote{H; \skt{एवाव्यापन्नैर्}\ A; \skt{एव व्यापन्नैर्}\ N}}}{\lemma{\skt{अधो\appbrk{}म\appbrk{}ध्यो\appbrk{}र्द्ध\appbrk{}स\appbrk{}न्नि\appbrk{}विष्टैः}}\Afootnote{H; \skt{अधो\appbrk{}म\appbrk{}ध्यो\appbrk{}र्ध्व\appbrk{}स\appbrk{}न्नि\appbrk{}विष्टैः}\ A; \skt{अधो\appbrk{}म\appbrk{}ध्यो\appbrk{}र्ध्व\appbrk{}स\appbrk{}न्नि\appbrk{}विष्टैः}\ N}}}{\lemma{\skt{भवन्ति~॥ तैर्}}\Afootnote{N H; \skt{तैर्}\ A}}\skt{ }\edtext{\edtext{\skt{शरी\-र\-मि\-द\-न्धा\-र्यते।}}{\lemma{\skt{शरीरम्}}\Afootnote{N H; \skt{शरीर}\ A}}}{\lemma{\skt{इदन्}}\Afootnote{H; \skt{मिदं}\ A; \skt{इद}\ N}}\skt{ अगारमिव }\edtext{\skt{स्थूणाभिरतश्च}}{\lemma{\skt{स्थूणाभिर्}}\Afootnote{N H; \skt{स्थूणाभिस्तिसृभिः}\ A}}\skt{ }\edtext{\edtext{\edtext{\edtext{\edtext{\skt{त्रिस्थू\-ण\-मि\-त्या\-हु\-रि\-त्येके।}}{\lemma{\skt{त्रिस्थूणम्}}\Afootnote{N H; \skt{त्रिस्थूण}\ A}}}{\lemma{\skt{इत्य्}}\Afootnote{N H; \emph{om.}\ A}}}{\lemma{\skt{आहुर्}}\Afootnote{N H; \emph{om.}\ A}}}{\lemma{\skt{इत्य्}}\Afootnote{N H; \skt{माहु}\ A}}\skt{ त}}{\lemma{\skt{एके। त}}\Afootnote{N H; \skt{रेके~। त}\ A}}\skt{ }\edtext{\skt{एव च}}{\lemma{\skt{एव च}}\Afootnote{A N H\textsuperscript{pc}; \skt{एव}\ H\textsuperscript{ac}}}\skt{ व्यापन्नाः प्रलयहेतवो }\edtext{\skt{भवन्ति}}{\lemma{\skt{भवन्ति}}\Afootnote{N H; \emph{om.}\ A}}\skt{ }\edtext{\skt{तदेभिः}}{\lemma{\skt{तदेभिः}}\Afootnote{N H; \skt{तदेभिर्}\ A}}\skt{ }\edtext{\skt{शोणितचतुर्त्थैः}}{\lemma{\skt{शोणितचतुर्त्थैः}}\Afootnote{N H; \skt{एव शोणितचतुर्थैः}\ A}}\skt{ सम्भ\-व\-स्थि\-ति\-प्र\-ल\-ये\-ष्व\-प्य\-वि\-र\-हितं शरीरम्भवति~॥}\pend

\pstart\versepar \ledsidenote{\passref{1.21.4}}\edtext{\skt{भवति}}{\lemma{\skt{भवति}}\Afootnote{A; \emph{om.}\ N H}}\skt{ }\edtext{\skt{चात्र}}{\lemma{\skt{चात्र}}\Afootnote{A; \emph{om.}\ N H}}\skt{~॥ }\pend
\pstart\versepar \skt{नर्त्ते देहः }\edtext{\skt{कफादस्ति न च}}{\lemma{\skt{कफादस्ति न च}}\Afootnote{N H; \skt{कफादस्ति न}\ A}}\skt{ }\edtext{\skt{पित्तान्न}}{\lemma{\skt{पित्तान्न}}\Afootnote{N H; \skt{पित्तान्न च}\ A}}\skt{ मारुतात्~। }\pend
\pstart\versepar \skt{शोणितादपि वा }\edtext{\skt{नित्य\-न्दे\-ह\-मे\-तैस्तु}}{\lemma{\skt{देहम्}}\Afootnote{N H; \skt{देह}\ A}}\skt{ धार्यते~॥}\pend

\pstart \ledsidenote{\passref{1.21.5}}\skt{तत्र वा }\edtext{\skt{गति\-ग\-न्धो\-पा\-दानो}}{\lemma{\skt{गति\appbrk{}ग\appbrk{}न्धो\appbrk{}पा\appbrk{}दानो}}\Afootnote{H; \skt{गतिगन्धनयोः, इति}\ A; \skt{गति\appbrk{}ग\appbrk{}न्धो\appbrk{}पा\appbrk{}दाने}\ N}}\skt{ धातुः~। तप सन्तापे~। श्लिष आलिङ्गने~। }\edtext{\skt{एषां}}{\lemma{\skt{एषां}}\Afootnote{N H; \skt{एतेषां}\ A}}\skt{ }\edtext{\skt{कृद्विहितैः}}{\lemma{\skt{कृद्विहितैः}}\Afootnote{A H; \skt{कृत्रविहितैः}\ N}}\skt{ }\edtext{\skt{प्रत्ययैर्व्वातं}}{\lemma{\skt{व्वातं}}\Afootnote{H; \skt{वातः}\ A; \skt{वात}\ N}}\skt{ }\edtext{\skt{पित्तं}}{\lemma{\skt{पित्तं}}\Afootnote{A H; \skt{पित्त}\ N}}\skt{ }\edtext{\skt{श्लेष्मा}}{\lemma{\skt{श्लेष्मा}}\Afootnote{H; \skt{श्लेष्मेति}\ A; \skt{श्लेष्मान}\ N}}\skt{ }\edtext{\skt{इति}}{\lemma{\skt{इति}}\Afootnote{N H; \skt{च}\ A}}\skt{ रूपाणि }\edtext{\skt{भवन्ति}}{\lemma{\skt{भवन्ति}}\Afootnote{A H; \skt{भवन्तिन्ति}\ N}}\skt{~॥}\pend

\pstart \ledsidenote{\passref{1.21.6}}\edtext{\skt{तेषां}}{\lemma{\skt{तेषां}}\Afootnote{N H; \skt{दोष}\ A}}\skt{ }\edtext{\edtext{\skt{स्थानान्यत}}{\lemma{\skt{स्थानान्य्}}\Afootnote{N\textsuperscript{pc} H; \skt{स्थानान्य}\ A; \skt{स्थानान्यि}\ N\textsuperscript{ac}}}}{\lemma{\skt{अत}}\Afootnote{N H; \skt{त}\ A}}\skt{ ऊर्द्ध्वम्वक्ष्यामः~॥ तत्र समासेन }\edtext{\skt{वायुः}}{\lemma{\skt{वायुः}}\Afootnote{N\textsuperscript{pc} H; \skt{वातः}\ A; \skt{वायु}\ N\textsuperscript{ac}}}\skt{ }\edtext{\skt{श्रोणी\-गु\-द\-सं\-श्रयः}}{\lemma{\skt{श्रोणी\appbrk{}गु\appbrk{}द\appbrk{}सं\appbrk{}श्रयः}}\Afootnote{N H; \skt{श्रोणि\appbrk{}गु\appbrk{}द\appbrk{}सं\appbrk{}श्रयः; तदुपर्यधो नाभेः पक्वाशयः}\ A}}\skt{~। }\edtext{\skt{पक्वा\-मा\-स\-य\-म\-ध्य\-स्थ\-म्पि\-त्तस्य}}{\lemma{\skt{पक्वा\appbrk{}मा\appbrk{}स\appbrk{}य\appbrk{}म\appbrk{}ध्य\appbrk{}स्थम्}}\Afootnote{N H; \skt{पक्वा\appbrk{}मा\appbrk{}श\appbrk{}य\appbrk{}मध्यं}\ A}}\skt{~। }\edtext{\skt{आमासयः}}{\lemma{\skt{आमासयः}}\Afootnote{N H; \skt{आमाशयः}\ A}}\skt{ श्लेष्मणः~॥}\pend

\pstart \ledsidenote{\passref{1.21.7}}\skt{अतः परम्पञ्चधा विभज्यन्ते~। तत्र वातस्य }\edtext{\skt{वातव्याधिके}}{\lemma{\skt{वातव्याधिके}}\Afootnote{N H; \skt{वातव्याधौ}\ A}}\skt{ वक्ष्यामः~। पित्तस्य यकृत्प्लीहानौ }\edtext{\edtext{\skt{हृद\-य\-न्दृ\-ष्टि\-स्त्व\-गिति}}{\lemma{\skt{दृष्टिस्}}\Afootnote{N H; \skt{दृष्टि}\ A}}}{\lemma{\skt{त्वगिति}}\Afootnote{N H; \skt{स्त्व}\ A}}\skt{~। }\edtext{\skt{पूर्व्वोक्तन्तु}}{\lemma{\skt{पूर्व्वोक्तन्तु}}\Afootnote{H; \skt{क्पूर्वोक्तं च}\ A; \skt{पूर्वोक्तञ्च}\ N}}\skt{ }\edtext{\edtext{\skt{श्लेष्मणस्तूरसः}}{\lemma{\skt{श्लेष्मणस्}}\Afootnote{N H; \skt{श्लेष्मण}\ A}}}{\lemma{\skt{तूरसः}}\Afootnote{N H; \skt{उरः शिरः}\ A}}\skt{ कण्ठः }\edtext{\skt{शिरः}}{\lemma{\skt{शिरः}}\Afootnote{N H; \skt{जिह्वामूलं}\ A}}\skt{ सन्धय इति~॥ }\edtext{\skt{एतानि}}{\lemma{\skt{एतानि}}\Afootnote{N H; \skt{पूर्वोक्तं च एतानि}\ A}}\skt{ खलु }\edtext{\edtext{\edtext{\skt{दोष\-स्था\-ना\-न्य\-व्या\-प\-न्ना\-ना\-म्भ\-वन्ति}}{\lemma{\skt{दोषस्थानान्य्}}\Afootnote{N H; \skt{दोषाणां स्थाना}\ A}}}{\lemma{\skt{अव्यापन्नानाम्}}\Afootnote{N H; \skt{न्यव्यापन्नानाम्}\ A}}}{\lemma{\skt{भवन्ति}}\Afootnote{N H; \emph{om.}\ A}}\skt{~। }\edtext{\edtext{\skt{पूर्व्वोक्तञ्च}}{\lemma{\skt{पूर्व्वोक्तञ्}}\Afootnote{N H; \emph{om.}\ A}}}{\lemma{\skt{च}}\Afootnote{N H; \emph{om.}\ A}}\skt{~॥}\pend

\pstart\versepar \ledsidenote{\passref{1.21.8}}\edtext{\skt{विस\-र्गा\-दा\-न\-वि\-क्षेपैः}}{\lemma{\skt{विस\appbrk{}र्गा\appbrk{}दा\appbrk{}न\appbrk{}वि\appbrk{}क्षेपैः}}\Afootnote{N H; \skt{भवति चात्र विस\appbrk{}र्गो\appbrk{}दा\appbrk{}न\appbrk{}वि\appbrk{}क्षेपैः}\ A}}\skt{ }\edtext{\skt{सोमसूर्यानिला}}{\lemma{\skt{सोमसूर्यानिला}}\Afootnote{A N\textsuperscript{pc} H; \skt{सोम\appbrk{}सू\appbrk{}र्या\appbrk{}द\appbrk{}निला}\ N\textsuperscript{ac}}}\skt{ यथा~। }\pend
\pstart\versepar \skt{}\edtext{\skt{धारयन्ति}}{\lemma{\skt{धारयन्ति}}\Afootnote{A N H\textsuperscript{pc}; \skt{धार न्ति}\ H\textsuperscript{ac}}}\skt{ जग\-द्दे\-ह\-ङ्क\-फ\-पि\-त्ता\-नि\-ला\-स्तथा~॥}\pend

\pstart \ledsidenote{\passref{1.21.9}}\skt{तत्र }\edtext{\skt{जिज्ञास्यते~॥ }\edtext{\edtext{\skt{किम्पि\-त्त\-व्य\-ति\-रे\-के\-णा\-न्योऽ\-ग्नि\-रु\-ताहो}}{\lemma{\skt{पित्त\appbrk{}व्य\appbrk{}ति\appbrk{}रे\appbrk{}के\appbrk{}णान्यो}}\Afootnote{N H; \skt{पित्त\appbrk{}व्य\appbrk{}ति\appbrk{}रे\appbrk{}का\appbrk{}दन्यो}\ A}}}{\lemma{\skt{ऽग्निरुताहो}}\Afootnote{H; \skt{ऽग्निः? आहो}\ A; \skt{ग्निरुताहो}\ N}}}{\lemma{\skt{जिज्ञास्यते~॥ किम्}}\Afootnote{N H; \skt{जिज्ञास्यं किं}\ A}}\skt{ }\edtext{\edtext{\edtext{\skt{पित्त\-मे\-वा\-ग्नि\-रिति}}{\lemma{\skt{पित्तम्}}\Afootnote{N H; \skt{स्वित्पित्त}\ A}}}{\lemma{\skt{एवाग्निर्}}\Afootnote{N H; \skt{मेवाग्नि}\ A}}}{\lemma{\skt{इति}}\Afootnote{N H; \skt{रिति}\ A}}\skt{~। अत्रोच्यते~। न खलु }\edtext{\edtext{\edtext{\skt{पित्त\-व्य\-ति\-रे\-के\-णा\-न्योऽ\-ग्नि\-रु\-प\-ल\-भ्यते}}{\lemma{\skt{पित्त\appbrk{}व्य\appbrk{}ति\appbrk{}रे\appbrk{}के\appbrk{}णान्यो}}\Afootnote{N H; \skt{पित्त\appbrk{}व्य\appbrk{}ति\appbrk{}रे\appbrk{}का\appbrk{}दन्यो}\ A}}}{\lemma{\skt{ऽग्निर्}}\Afootnote{N H; \skt{ऽग्नि}\ A}}}{\lemma{\skt{उपलभ्यते}}\Afootnote{N H; \skt{रुपलभ्यते}\ A}}\skt{~। आग्नेयत्वात्पित्ते }\edtext{\skt{दह\-न\-प\-च\-ना\-दिषु}}{\lemma{\skt{दह\appbrk{}न\appbrk{}प\appbrk{}च\appbrk{}ना\appbrk{}दिषु}}\Afootnote{N H; \skt{दह\appbrk{}न\appbrk{}प\appbrk{}च\appbrk{}ना\appbrk{}दिष्व}\ A}}\skt{ }\edtext{\edtext{\skt{प्रव\-र्त्त\-मा\-नोऽ\-ग्नि\-व\-दु\-प\-चारः}}{\lemma{\skt{प्रवर्त्तमानो}}\Afootnote{H; \skt{भिप्रवर्तमाने}\ A; \skt{प्रवर्त्तमाने}\ N}}}{\lemma{\skt{ऽग्नि\appbrk{}व\appbrk{}दु\appbrk{}प\appbrk{}चारः}}\Afootnote{H; \skt{ऽग्नि\appbrk{}व\appbrk{}दु\appbrk{}प\appbrk{}चारः}\ A; \skt{ग्निवदुपचारः}\ N}}\skt{ }\edtext{\skt{क्रियन्ते }\edtext{\skt{न्तरग्निरिति}}{\lemma{\skt{इति}}\Afootnote{N H; \skt{रिति}\ A}}}{\lemma{\skt{क्रियन्ते न्तरग्निर्}}\Afootnote{H; \skt{क्रियतेऽन्तरग्नि}\ A; \skt{क्रियन्ते रग्निर्}\ N}}\skt{~। }\edtext{\skt{क्षिणे}}{\lemma{\skt{क्षिणे}}\Afootnote{H; \skt{क्षीणे}\ A N}}\skt{ }\edtext{\skt{वाग्निगुणे}}{\lemma{\skt{वाग्निगुणे}}\Afootnote{N H; \skt{ह्यग्निगुणे}\ A}}\skt{ }\edtext{\skt{तत्स\-मा\-न\-द्र\-व्यो\-प\-यो\-गा\-द\-ति\-वृद्धे}}{\lemma{\skt{तत्स\appbrk{}मा\appbrk{}न\appbrk{}द्र\appbrk{}व्यो\appbrk{}प\appbrk{}यो\appbrk{}गा\appbrk{}द\appbrk{}ति\appbrk{}वृद्धे}}\Afootnote{A H; \skt{तत्स\appbrk{}मा\appbrk{}न\appbrk{}द्र\appbrk{}व्यो\appbrk{}प\appbrk{}यो\appbrk{}गा\appbrk{}द\appbrk{}ति\appbrk{}प्र\appbrk{}वृद्धे}\ N}}\skt{ शीत\-क्रि\-यो\-प\-यो\-गा\-दा\-ग\-माच्च पश्यामो न खलु }\edtext{\edtext{\edtext{\skt{पित्त\-व्य\-ति\-रे\-के\-णा\-न्योऽ\-ग्नि\-रिति}}{\lemma{\skt{पित्त\appbrk{}व्य\appbrk{}ति\appbrk{}रे\appbrk{}के\appbrk{}णान्यो}}\Afootnote{N H; \skt{पित्त\appbrk{}व्य\appbrk{}ति\appbrk{}रे\appbrk{}का\appbrk{}दन्यो}\ A}}}{\lemma{\skt{ऽग्निर्}}\Afootnote{H; \skt{ऽग्नि}\ A; \skt{ग्निर्}\ N}}}{\lemma{\skt{इति}}\Afootnote{N H; \skt{रिति}\ A}}\skt{~।}\pend

\pstart \ledsidenote{\passref{1.21.10}}\edtext{\skt{तत्तु}}{\lemma{\skt{तत्तु}}\Afootnote{H; \skt{तच्चा}\ A; \skt{तत्र}\ N}}\skt{ }\edtext{\skt{दृष्टहेतुकेन}}{\lemma{\skt{दृष्टहेतुकेन}}\Afootnote{A H; \skt{दृष्टहेतुके}\ N}}\skt{ विशेषेण }\edtext{\skt{पक्वा\-मा\-स\-य\-म\-ध्यस्थं}}{\lemma{\skt{पक्वा\appbrk{}मा\appbrk{}स\appbrk{}य\appbrk{}म\appbrk{}ध्यस्थं}}\Afootnote{N H; \skt{पक्वा\appbrk{}मा\appbrk{}श\appbrk{}य\appbrk{}म\appbrk{}ध्यस्थं}\ A}}\skt{ }\edtext{\edtext{\edtext{\skt{चतु\-र्व्वि\-ध\-म\-प्य\-न्न\-म्प\-चति}}{\lemma{\skt{चतुर्व्विधम्}}\Afootnote{N H; \skt{पित्तं चतुर्विध}\ A}}}{\lemma{\skt{अप्य्}}\Afootnote{N H; \skt{मन्न}\ A}}}{\lemma{\skt{अन्नम्}}\Afootnote{N H; \skt{पानं}\ A}}\skt{ विवेचयति च दोष\-र\-स\-मू\-त्र\-पु\-री\-षाणि~। तत्रस्थमेव }\edtext{\skt{चात्मशक्त्या}}{\lemma{\skt{चात्मशक्त्या}}\Afootnote{A N\textsuperscript{pc} H; \skt{चात्मक्त्या}\ N\textsuperscript{ac}}}\skt{ }\edtext{\skt{शेषा\-णा\-म्पि\-त्त\-स्था\-नानां}}{\lemma{\skt{पित्तस्थानानां}}\Afootnote{A H; \skt{पित्तस्थानाना}\ N}}\skt{ शरीरस्य चाग्नि\-क\-र्म्म\-णा\-नु\-ग्र\-ह\-ङ्क\-रोति~। तस्मिन्पित्ते }\edtext{\edtext{\skt{पाचकोऽग्निरिति}}{\lemma{\skt{ऽग्निर्}}\Afootnote{H; \skt{ऽग्नि}\ A; \skt{ग्निर्}\ N}}}{\lemma{\skt{इति}}\Afootnote{N H; \skt{रिति}\ A}}\skt{ संज्ञा~॥ }\edtext{\skt{यकृत्प्लीह्नोः}}{\lemma{\skt{यकृत्प्लीह्नोः}}\Afootnote{N H; \skt{यत्तु यकृत्प्लीह्नोः}\ A}}\skt{ पित्त\-न्त\-स्मि\-न्र\-ञ्जको }\edtext{\skt{ग्निरिति}}{\lemma{\skt{ग्निरिति}}\Afootnote{H; \skt{ऽग्निरिति}\ A; \skt{ग्नि}\ N}}\skt{ संज्ञा स }\edtext{\skt{रसस्य}}{\lemma{\skt{रसस्य}}\Afootnote{A H; \skt{रस्यस्य}\ N}}\skt{ }\edtext{\skt{रागकृदुक्तः}}{\lemma{\skt{रागकृदुक्तः}}\Afootnote{N H; \skt{रागकृदुक्तः}\ A}}\skt{~। }\edtext{\skt{यत्तु}}{\lemma{\skt{यत्तु}}\Afootnote{N H; \skt{यत्पित्तं}\ A}}\skt{ }\edtext{\edtext{\edtext{\edtext{\edtext{\skt{हृदि\-स्थ\-म्पि\-त्त\-न्त\-स्मि\-न्सा\-ध\-कोऽ\-ग्नि\-रिति}}{\lemma{\skt{हृदिस्थम्}}\Afootnote{N H; \skt{हृद}\ A}}}{\lemma{\skt{पित्तन्}}\Afootnote{N H; \skt{य}\ A}}}{\lemma{\skt{तस्मिन्साधको}}\Afootnote{H; \skt{स्थं तस्मिन्साधको}\ A; \skt{शक्तिसाधको}\ N}}}{\lemma{\skt{ऽग्निर्}}\Afootnote{H; \skt{ऽग्नि}\ A; \skt{ग्निर्}\ N}}}{\lemma{\skt{इति}}\Afootnote{N H; \skt{रिति}\ A}}\skt{ संज्ञा }\edtext{\skt{शोभि\-त\-प्रा\-र्थि\-त\-म\-नो\-र\-थ\-सा\-ध\-न\-कृ\-दुक्तः}}{\lemma{\skt{शोभि\appbrk{}त\appbrk{}प्रा\appbrk{}र्थि\appbrk{}त\appbrk{}म\appbrk{}नो\appbrk{}र\appbrk{}थ\appbrk{}सा\appbrk{}ध\appbrk{}न\appbrk{}कृ\appbrk{}दुक्तः}}\Afootnote{H; \skt{सोऽभि\appbrk{}प्रा\appbrk{}र्थि\appbrk{}त\appbrk{}म\appbrk{}नो\appbrk{}र\appbrk{}थ\appbrk{}सा\appbrk{}ध\appbrk{}न\appbrk{}कृ\appbrk{}दुक्तः}\ A; \skt{शोभि\appbrk{}प्रा\appbrk{}र्थि\appbrk{}त\appbrk{}म\appbrk{}नो\appbrk{}र\appbrk{}थ\appbrk{}सा\appbrk{}ध्य\appbrk{}कृ\appbrk{}दुक्तः}\ N}}\skt{~॥ }\edtext{\skt{यत्तु}}{\lemma{\skt{यत्तु}}\Afootnote{N H; \skt{य}\ A}}\skt{ }\edtext{\skt{दृष्ट्याम्पित्तं}}{\lemma{\skt{दृष्ट्याम्}}\Afootnote{N H; \skt{द्दृष्ट्यां}\ A}}\skt{ }\edtext{\edtext{\edtext{\skt{तस्मि\-न्ना\-लो\-च\-कोऽ\-ग्नि\-रिति}}{\lemma{\skt{तस्मिन्नालोचको}}\Afootnote{N H; \skt{तस्मिन्नालोचको}\ A}}}{\lemma{\skt{ऽग्निर्}}\Afootnote{H; \skt{ऽग्नि}\ A; \skt{ग्निर्}\ N}}}{\lemma{\skt{इति}}\Afootnote{N H; \skt{रिति}\ A}}\skt{ }\edtext{\skt{संज्ञा स}}{\lemma{\skt{संज्ञा स}}\Afootnote{A H; \skt{स}\ N}}\skt{ }\edtext{\skt{रूपग्रहणे}}{\lemma{\skt{रूपग्रहणे}}\Afootnote{N H; \skt{रूपग्रहणाऽ}\ A}}\skt{ धिकृतः~॥ यत्तु त्वचि }\edtext{\skt{पित्तन्तस्य}}{\lemma{\skt{तस्य}}\Afootnote{N H; \skt{तस्मि}\ A}}\skt{ }\edtext{\edtext{\edtext{\skt{भ्राजकोऽग्निरिति}}{\lemma{\skt{भ्राजको}}\Afootnote{N H; \skt{न्भ्राजको}\ A}}}{\lemma{\skt{ऽग्निर्}}\Afootnote{H; \skt{ऽग्नि}\ A; \skt{ग्निर्}\ N}}}{\lemma{\skt{इति}}\Afootnote{N H; \skt{रिति}\ A}}\skt{ }\edtext{\skt{संज्ञा }\edtext{\skt{सोऽभ्य\-ङ्ग\-प\-रि\-से\-का\-व\-गा\-हा\-ल\-य\-ना\-दीनां}}{\lemma{\skt{ऽभ्य\appbrk{}ङ्ग\appbrk{}प\appbrk{}रि\appbrk{}से\appbrk{}का\appbrk{}व\appbrk{}गा\appbrk{}हा\appbrk{}ल\appbrk{}य\appbrk{}ना\appbrk{}दीनां}}\Afootnote{\skt{ऽभ्य\appbrk{}ङ्ग\appbrk{}प\appbrk{}रि\appbrk{}षे\appbrk{}का\appbrk{}व\appbrk{}गा\appbrk{}हा\appbrk{}ले\appbrk{}प\appbrk{}ना\appbrk{}दीनां}\ A; \skt{ऽभ्य\appbrk{}ङ्ग\appbrk{}प\appbrk{}रि\appbrk{}षे\appbrk{}का\appbrk{}व\appbrk{}गा\appbrk{}हा\appbrk{}ले\appbrk{}प\appbrk{}ना\appbrk{}दीना}\ N; \skt{भ्यङ्ग\appbrk{}प\appbrk{}रि\appbrk{}से\appbrk{}का\appbrk{}व\appbrk{}गा\appbrk{}हा\appbrk{}ल\appbrk{}य\appbrk{}ना\appbrk{}दीनां}\ H}}}{\lemma{\skt{संज्ञा सो}}\Afootnote{A N; \skt{संज्ञा सोऽ}\ H}}\skt{ }\edtext{\skt{क्रियाद्रव्यादीनां पाचयिता}}{\lemma{\skt{क्रियाद्रव्यादीनां पाचयिता}}\Afootnote{H; \skt{क्रियाद्रव्याणां पक्ता छायानां च प्रकाशकः}\ A; \skt{क्रियाद्रव्याणां पाचयिता}\ N}}\skt{~॥॥}\pend

\pstart\versepar \ledsidenote{\passref{1.21.11}}\edtext{\skt{भवति}}{\lemma{\skt{भवति}}\Afootnote{A N; \skt{भ}\ H}}\skt{ }\edtext{\skt{चात्र}}{\lemma{\skt{चात्र}}\Afootnote{A N; \skt{भवति चात्र~॥}\ H}}\skt{~॥ }\pend
\pstart\versepar \skt{पित्तं }\edtext{\skt{तीक्ष्ण\-न्द्र\-व\-म्पू\-ति\-नी\-ल\-म्पी\-त\-न्त\-थैव}}{\lemma{\skt{पीतन्}}\Afootnote{A H; \skt{पीत्तं}\ N}}\skt{ च~। }\pend
\pstart\versepar \skt{उष्ण\-ङ्क\-टु\-र\-स\-ञ्चैव }\edtext{\edtext{\edtext{\skt{विद\-ग्ध\-न्त्व\-म्ल\-मि\-ष्यते}}{\lemma{\skt{विदग्धन्}}\Afootnote{A H; \skt{विदग्ध}\ N}}}{\lemma{\skt{त्वम्लम्}}\Afootnote{H; \skt{चाम्लम्}\ A; \skt{त्वम्ब्लम्}\ N}}}{\lemma{\skt{इष्यते}}\Afootnote{N H; \skt{एव च}\ A}}\skt{~॥}\pend

\pstart \ledsidenote{\passref{1.21.12}}\skt{अत ऊर्ध्वं }\edtext{\edtext{\skt{श्लेष्म\-स्था\-ना\-न्य\-नु\-व्या\-ख्या\-स्यामः}}{\lemma{\skt{श्लेष्मस्थानान्य्}}\Afootnote{N H; \skt{श्लेष्मस्थाना}\ A}}}{\lemma{\skt{अनुव्याख्यास्यामः}}\Afootnote{N H; \skt{न्यनुव्याख्यास्यामः}\ A}}\skt{~॥ }\edtext{\skt{तत्रामासयः}}{\lemma{\skt{तत्रामासयः}}\Afootnote{N H; \skt{तत्र, आमाशयः}\ A}}\skt{ }\edtext{\skt{पित्ता\-स\-य\-स्यो\-परि}}{\lemma{\skt{पित्ता\appbrk{}स\appbrk{}य\appbrk{}स्यो\appbrk{}परि}}\Afootnote{N H; \skt{पित्ता\appbrk{}श\appbrk{}य\appbrk{}स्यो\appbrk{}परि}\ A}}\skt{ }\edtext{\edtext{\skt{तत्प्र\-त्य\-नी\-क\-त्वा\-दू\-र्ध्व\-ग\-ति\-त्वा\-त्ते\-ज\-स\-श्चन्द्र}}{\lemma{\skt{तत्प्र\appbrk{}त्य\appbrk{}नी\appbrk{}क\appbrk{}त्वा\appbrk{}दू\appbrk{}र्ध्व\appbrk{}ग\appbrk{}ति\appbrk{}त्वा\appbrk{}त्ते\appbrk{}ज\appbrk{}सश्}}\Afootnote{N H; \skt{ष्टात्त\appbrk{}त्प्र\appbrk{}त्य\appbrk{}नी\appbrk{}क\appbrk{}त्वा\appbrk{}दू\appbrk{}र्ध्व\appbrk{}ग\appbrk{}ति\appbrk{}त्वा\appbrk{}त्ते\appbrk{}जसः}\ A}}}{\lemma{\skt{चन्द्र}}\Afootnote{A H; \skt{चन्द्रा}\ N}}\skt{ }\edtext{\skt{इवादित्यस्य स}}{\lemma{\skt{इवादित्यस्य स}}\Afootnote{H; \skt{इव आदित्यस्य}\ A; \skt{दित्यस्य स}\ N}}\skt{ चतु\-र्व्वि\-ध\-स्या\-हा\-र\-स्या\-धारः स च }\edtext{\edtext{\skt{तत्रौदकैर्गुणैः}}{\lemma{\skt{तत्रौदकैर्}}\Afootnote{N H; \skt{तत्रौदकै}\ A}}}{\lemma{\skt{गुणैः}}\Afootnote{N H; \skt{र्गुणै}\ A}}\skt{ }\edtext{\skt{प्रक्लिन्नो}}{\lemma{\skt{प्रक्लिन्नो}}\Afootnote{N H; \skt{राहारः प्रक्लिन्नो}\ A}}\skt{ भिन्न\-सं\-घा\-त\-स्सु\-ख\-जरो भवति~॥}\pend

\pstart\versepar \ledsidenote{\passref{1.21.13}}\skt{भवति चात्र~॥ }\pend
\pstart\versepar \skt{माधु\-र्या\-त्पि\-च्छि\-ल\-त्वाच्च }\edtext{\skt{प्रक्ले\-दि\-त्वा\-त्त\-थैव च}}{\lemma{\skt{प्रक्ले\appbrk{}दि\appbrk{}त्वा\appbrk{}त्त\appbrk{}थैव च}}\Afootnote{A H; \skt{प्रक्लि\appbrk{}न्न\appbrk{}त्वा\appbrk{}त्त\appbrk{}थैव च}\ N}}\skt{~। }\pend
\pstart\versepar \skt{}\edtext{\skt{आमासये}}{\lemma{\skt{आमासये}}\Afootnote{N H; \skt{आमाशये}\ A}}\skt{ सम्भवति श्लेष्मा मधुरशीतलः~॥}\edtext{}{\lemma{\skt{भवति चात्र~॥ माधु\appbrk{}र्या\appbrk{}त्पि\appbrk{}च्छि\appbrk{}ल\appbrk{}त्वाच्च}}\Afootnote{\skt{माधु\appbrk{}र्या\appbrk{}त्पि\appbrk{}च्छि\appbrk{}ल\appbrk{}त्वाच्च}\ A H N}}\pend

\pstart \ledsidenote{\passref{1.21.14}}\skt{स तत्र }\edtext{\skt{स्थितः}}{\lemma{\skt{स्थितः}}\Afootnote{N H; \skt{स्थ एव}\ A}}\skt{ स्वशक्त्या }\edtext{\skt{शेषाणां}}{\lemma{\skt{शेषाणां}}\Afootnote{A H; \skt{शेषाणा}\ N}}\skt{ }\edtext{\skt{श्लेष्मस्थानानां}}{\lemma{\skt{श्लेष्मस्थानानां}}\Afootnote{A H; \skt{श्लेष्मस्थानाना}\ N}}\skt{ }\edtext{\skt{शरी\-र\-स्यो\-द\-क\-क\-र्म्म\-णा\-नु\-ग्र\-ह\-ङ्क\-रोति}}{\lemma{\skt{शरी\appbrk{}र\appbrk{}स्यो\appbrk{}द\appbrk{}क\appbrk{}क\appbrk{}र्म्म\appbrk{}णा\appbrk{}नु\appbrk{}ग्र\appbrk{}हङ्}}\Afootnote{H; \skt{शरीरस्य चोद\appbrk{}क\appbrk{}क\appbrk{}र्म\appbrk{}णाऽ\appbrk{}नु\appbrk{}ग्रहं}\ A; \skt{शरीरस्य उद\appbrk{}क\appbrk{}क\appbrk{}र्मो\appbrk{}नु\appbrk{}ग्र\appbrk{}हङ्}\ N}}\skt{~॥ }\edtext{\skt{उरःस्थस्तु}}{\lemma{\skt{उरःस्थस्तु}}\Afootnote{H; \skt{उरःस्थ}\ A; \skt{उरस्थस्तु}\ N}}\skt{ }\edtext{\edtext{\skt{त्रिक\-स\-न्धा\-र\-ण\-मा\-त्म\-वी\-र्य्ये\-णा\-न्न\-र\-स\-स\-हि\-तेन}}{\lemma{\skt{त्रिक\appbrk{}स\appbrk{}न्धा\appbrk{}र\appbrk{}णम्}}\Afootnote{N H; \skt{स्त्रिकसन्धारण}\ A}}}{\lemma{\skt{आत्म\appbrk{}वी\appbrk{}र्य्ये\appbrk{}णा\appbrk{}न्न\appbrk{}र\appbrk{}स\appbrk{}स\appbrk{}हि\appbrk{}तेन}}\Afootnote{N H; \skt{मात्म\appbrk{}वी\appbrk{}र्ये\appbrk{}णा\appbrk{}न्न\appbrk{}र\appbrk{}स\appbrk{}स\appbrk{}हि\appbrk{}तेन}\ A}}\skt{ }\edtext{\skt{हृद\-या\-व\-ल\-म्ब\-न\-म्व\-र\-णञ्च}}{\lemma{\skt{वरणञ्च}}\Afootnote{N H; \emph{om.}\ A}}\skt{ करोति~। }\edtext{\skt{कण्ठस्थस्तु}}{\lemma{\skt{कण्ठस्थस्तु}}\Afootnote{N H\textsuperscript{pc}; \skt{जिह्वा\appbrk{}मू\appbrk{}ल\appbrk{}क\appbrk{}ण्ठस्थो}\ A; \skt{कण्ठस्तु}\ H\textsuperscript{ac}}}\skt{ जिह्वेन्द्रियस्य }\edtext{\skt{सौम्य\-त्वा\-त्स\-म्य\-ग्र\-स\-ज्ञाने}}{\lemma{\skt{सौम्य\appbrk{}त्वा\appbrk{}त्स\appbrk{}म्य\appbrk{}ग्र\appbrk{}स\appbrk{}ज्ञाने}}\Afootnote{A N; \skt{सौम्यत्वात् 𑑛 सम्यग्रसज्ञाने}\ H}}\skt{ }\edtext{\skt{विवर्त्तते}}{\lemma{\skt{विवर्त्तते}}\Afootnote{H; \skt{वर्तते}\ A; \skt{निवर्त्तते}\ N}}\skt{~॥ }\edtext{\skt{शिरःस्थस्तु}}{\lemma{\skt{शिरःस्थस्तु}}\Afootnote{H; \skt{शिरःस्थः}\ A; \skt{शिरस्थस्तु}\ N}}\skt{ }\edtext{\edtext{\skt{स्नेह\-स\-न्त\-र्प्प\-णा\-धि\-कृ\-त\-त्वा\-दि\-न्द्रि\-या\-णा\-मा\-त्म\-वी\-र्य्ये\-णा\-नु\-ग्रहं}}{\lemma{\skt{स्नेह\appbrk{}स\appbrk{}न्त\appbrk{}र्प्प\appbrk{}णा\appbrk{}धि\appbrk{}कृ\appbrk{}त\appbrk{}त्वा\appbrk{}दि\appbrk{}न्द्रि\appbrk{}या\appbrk{}णाम्}}\Afootnote{N H; \skt{स्नेह\appbrk{}स\appbrk{}न्त\appbrk{}र्प\appbrk{}णा\appbrk{}धि\appbrk{}कृ\appbrk{}त\appbrk{}त्वा\appbrk{}दि\appbrk{}न्द्रि\appbrk{}याणा}\ A}}}{\lemma{\skt{आत्म\appbrk{}वी\appbrk{}र्य्ये\appbrk{}णा\appbrk{}नु\appbrk{}ग्रहं}}\Afootnote{N H; \skt{मात्म\appbrk{}वी\appbrk{}र्ये\appbrk{}णा\appbrk{}नु\appbrk{}ग्रहं}\ A}}\skt{ }\edtext{\skt{करोति}}{\lemma{\skt{करोति}}\Afootnote{N H; \skt{करो}\ A}}\skt{~॥ }\edtext{\skt{सन्धिस्थस्तु}}{\lemma{\skt{सन्धिस्थस्तु}}\Afootnote{N H; \emph{om.}\ A}}\skt{ }\edtext{\skt{सन्धि\-सं\-श्ले\-षा\-त्सर्व्वं}}{\lemma{\skt{सन्धि\appbrk{}सं\appbrk{}श्ले\appbrk{}षा\appbrk{}त्सर्व्वं}}\Afootnote{H; \skt{त्; सन्धिस्थः श्लेष्मा सर्व}\ A; \emph{om.}\ N}}\skt{ }\edtext{\skt{सन्धि\-सं\-श्ले\-षा\-त्स\-र्व्व\-स\-न्ध्य\-ग्र\-नु\-ह\-ङ्क\-रोति}}{\lemma{\skt{सन्धि\appbrk{}सं\appbrk{}श्ले\appbrk{}षा\appbrk{}त्स\appbrk{}र्व्व\appbrk{}स\appbrk{}न्ध्य\appbrk{}ग्र\appbrk{}नु\appbrk{}हङ्}}\Afootnote{H; \skt{सन्धि\appbrk{}सं\appbrk{}श्ले\appbrk{}षा\appbrk{}त्स\appbrk{}र्व\appbrk{}स\appbrk{}न्ध्य\appbrk{}नु\appbrk{}ग्रहं}\ A; \skt{सन्धि\appbrk{}श्ले\appbrk{}षा\appbrk{}त्स\appbrk{}र्व्व\appbrk{}स\appbrk{}न्ध्य\appbrk{}नु\appbrk{}ग्र\appbrk{}हङ्}\ N}}\skt{~॥}\pend

\pstart\versepar \ledsidenote{\passref{1.21.15}}\edtext{\skt{भवति}}{\lemma{\skt{भवति}}\Afootnote{A; \emph{om.}\ N H}}\skt{ }\edtext{\skt{चात्र}}{\lemma{\skt{चात्र}}\Afootnote{A; \emph{om.}\ N H}}\skt{~॥ }\pend
\pstart\versepar \skt{श्लेष्मा श्वेतो }\edtext{\skt{गुरु}}{\lemma{\skt{गुरु}}\Afootnote{N H; \skt{गुरुः}\ A}}\skt{ स्निग्धः }\edtext{\skt{पिच्छिलः}}{\lemma{\skt{पिच्छिलः}}\Afootnote{A H; \skt{पिच्छिल}\ N}}\skt{ शीत एव च~। }\pend
\pstart\versepar \skt{}\edtext{\edtext{\skt{मधु\-र\-श्चा\-वि\-दग्धः}}{\lemma{\skt{मधुरश्}}\Afootnote{N H; \skt{मधुर}\ A}}}{\lemma{\skt{चाविदग्धः}}\Afootnote{N H; \skt{स्त्वविदग्धः}\ A}}\skt{ }\edtext{\skt{स्याद्विदग्धो}}{\lemma{\skt{स्याद्विदग्धो}}\Afootnote{N H; \skt{स्याद्विदग्धो}\ A}}\skt{ }\edtext{\skt{लवणः}}{\lemma{\skt{लवणः}}\Afootnote{H; \skt{लवण}\ A N}}\skt{ स्मृतः~॥}\pend

\pstart \ledsidenote{\passref{1.21.16}}\edtext{\skt{शोणि\-त\-स्था\-न\-प्रा\-ग\-भि\-हितं}}{\lemma{\skt{शोणि\appbrk{}त\appbrk{}स्था\appbrk{}न\appbrk{}प्रा\appbrk{}ग\appbrk{}भि\appbrk{}हितं}}\Afootnote{H; \skt{शोणितस्य स्थानं यकृत्प्लीहानौ, तच्च प्रागभिहितं; तत्रस्थमेव शेषाणां शोणि\appbrk{}त\appbrk{}स्था\appbrk{}ना\appbrk{}ना\appbrk{}म\appbrk{}नु\appbrk{}ग्रहं करोति}\ A; \skt{शोणितं स्थानं प्रागभिहितं}\ N}}\skt{~।}\pend

\pstart\versepar \ledsidenote{\passref{1.21.17}}\edtext{\skt{अनु\-ष्ण\-शी\-त\-म्म\-धुरं}}{\lemma{\skt{अनुष्णशीतम्}}\Afootnote{N H; \skt{भवति चात्र अनुष्णशीतं}\ A}}\skt{ स्निग्धं रक्तञ्च वर्णतः~। }\pend
\pstart\versepar \skt{शोणितङ्गुरु विस्रं }\edtext{\skt{स्याद्विदाहश्चास्य}}{\lemma{\skt{स्याद्विदाहश्}}\Afootnote{N H; \skt{स्याद्विदाहश्}\ A}}\skt{ पित्तवत्~॥}\pend

\pstart \ledsidenote{\passref{1.21.18}}\skt{एतानि खलु }\edtext{\skt{दोषस्थानान्यत्र}}{\lemma{\skt{अत्र}}\Afootnote{N H; \skt{एषु}\ A}}\skt{ सञ्चीयन्ते दोषाः~। प्राक्स\-ञ्च\-य\-हे\-तु\-रुक्तः~॥ }\edtext{\skt{सञ्चितानां}}{\lemma{\skt{सञ्चितानां}}\Afootnote{N H; \skt{तत्र सञ्चितानां}\ A}}\skt{ खलु दोषाणां }\edtext{\skt{स्तब्ध\-पू\-र्ण्ण\-को\-ष्ठता}}{\lemma{\skt{स्तब्ध\appbrk{}पू\appbrk{}र्ण्ण\appbrk{}को\appbrk{}ष्ठता}}\Afootnote{A H; \skt{तच्च\appbrk{}अ\appbrk{}पू\appbrk{}र्ण्ण\appbrk{}को\appbrk{}ष्ठता}\ N}}\skt{ पीतावभासता मन्दोष्मता चाङ्गानां }\edtext{\edtext{\edtext{\skt{गौर\-व\-मा\-ल\-स्य\-ञ्चेति}}{\lemma{\skt{गौरवम्}}\Afootnote{N H; \skt{गौरव}\ A}}}{\lemma{\skt{आलस्यञ्}}\Afootnote{N H; \skt{मालस्यं}\ A}}}{\lemma{\skt{चेति}}\Afootnote{N H; \skt{चय\appbrk{}का\appbrk{}र\appbrk{}ण\appbrk{}वि\appbrk{}द्वे\appbrk{}ष\appbrk{}श्चेति}\ A}}\skt{ लिङ्गानि भवन्ति~। तत्र }\edtext{\skt{प्रथमः}}{\lemma{\skt{प्रथमः}}\Afootnote{A H; \skt{प्रथम}\ N}}\skt{ क्रियाकालः~॥}\pend

\pstart \ledsidenote{\passref{1.21.19}}\skt{अत }\edtext{\skt{ऊर्ध्व\-म्प्र\-को\-प\-का\-र\-णानि}}{\lemma{\skt{प्रको\appbrk{}प\appbrk{}का\appbrk{}र\appbrk{}णानि}}\Afootnote{N H; \skt{प्रकोपणानि}\ A}}\skt{ वक्ष्यामः~। तत्र }\edtext{\skt{बल\-व\-द्वि\-ग्र\-हा\-ति\-व्य\-वा\-य\-व्या\-या\-म\-प्र\-त\-प\-न\-प्र\-धा\-व\-ना\-भि\-घा\-त\-ल\-ङ्घ\-न\-प्ल\-व\-न\-प्र\-त\-र\-ण\-जा\-ग\-र\-ण\-ह\-य\-ग\-ज\-र\-थ\-प्र\-या\-णा\-ति\-क\-टु\-क\-क\-षा\-य\-ति\-क्त\-ल\-घु\-रू\-क्ष\-शी\-त\-वी\-र्य्य\-शु\-ष्क\-शा\-क\-व\-ल्लू\-र\-को\-र\-दू\-ष\-श्या\-मा\-क\-नी\-वा\-र\-मु\-द्ग\-म\-सू\-रा\-ढ\-की\-ह\-रे\-णु\-नि\-ष्पा\-वा\-न\-श\-न\-वा\-त\-मू\-त्र\-पु\-री\-ष\-शु\-क्र\-व\-थू\-द्गा\-र\-च्छ\-र्द्दि\-बा\-ष्प\-वे\-ग\-वि\-घा\-ता\-दि\-भि\-र्व्वि\-शे\-षै\-र्व्वायुः}}{\lemma{\skt{बल\appbrk{}व\appbrk{}द्वि\appbrk{}ग्र\appbrk{}हा\appbrk{}ति\appbrk{}व्य\appbrk{}वा\appbrk{}य\appbrk{}व्या\appbrk{}या\appbrk{}म\appbrk{}प्र\appbrk{}त\appbrk{}प\appbrk{}न\appbrk{}प्र\appbrk{}धा\appbrk{}व\appbrk{}ना\appbrk{}भि\appbrk{}घा\appbrk{}त\appbrk{}ल\appbrk{}ङ्घ\appbrk{}न\appbrk{}प्ल\appbrk{}व\appbrk{}न\appbrk{}प्र\appbrk{}त\appbrk{}र\appbrk{}ण\appbrk{}जा\appbrk{}ग\appbrk{}र\appbrk{}ण\appbrk{}ह\appbrk{}य\appbrk{}ग\appbrk{}ज\appbrk{}र\appbrk{}थ\appbrk{}प्र\appbrk{}या\appbrk{}णा\appbrk{}ति\appbrk{}क\appbrk{}टु\appbrk{}क\appbrk{}क\appbrk{}षा\appbrk{}य\appbrk{}ति\appbrk{}क्त\appbrk{}ल\appbrk{}घु\appbrk{}रू\appbrk{}क्ष\appbrk{}शी\appbrk{}त\appbrk{}वी\appbrk{}र्य्य\appbrk{}शु\appbrk{}ष्क\appbrk{}शा\appbrk{}क\appbrk{}व\appbrk{}ल्लू\appbrk{}र\appbrk{}को\appbrk{}र\appbrk{}दू\appbrk{}ष\appbrk{}श्या\appbrk{}मा\appbrk{}क\appbrk{}नी\appbrk{}वा\appbrk{}र\appbrk{}मु\appbrk{}द्ग\appbrk{}म\appbrk{}सू\appbrk{}रा\appbrk{}ढ\appbrk{}की\appbrk{}ह\appbrk{}रे\appbrk{}णु\appbrk{}नि\appbrk{}ष्पा\appbrk{}वा\appbrk{}न\appbrk{}श\appbrk{}न\appbrk{}वा\appbrk{}त\appbrk{}मू\appbrk{}त्र\appbrk{}पु\appbrk{}री\appbrk{}ष\appbrk{}शु\appbrk{}क्र\appbrk{}व\appbrk{}थू\appbrk{}द्गा\appbrk{}र\appbrk{}च्छ\appbrk{}र्द्दि\appbrk{}बा\appbrk{}ष्प\appbrk{}वे\appbrk{}ग\appbrk{}वि\appbrk{}घा\appbrk{}ता\appbrk{}दि\appbrk{}भिर्}}\Afootnote{H; \skt{बल\appbrk{}व\appbrk{}द्वि\appbrk{}ग्र\appbrk{}हा\appbrk{}ति\appbrk{}व्या\appbrk{}या\appbrk{}म\appbrk{}व्य\appbrk{}वा\appbrk{}या\appbrk{}ध्य\appbrk{}य\appbrk{}न\appbrk{}प्र\appbrk{}प\appbrk{}त\appbrk{}न\appbrk{}प्र\appbrk{}धा\appbrk{}व\appbrk{}न\appbrk{}प्र\appbrk{}पी\appbrk{}ड\appbrk{}ना\appbrk{}भि\appbrk{}घा\appbrk{}त\appbrk{}ल\appbrk{}ङ्घ\appbrk{}न\appbrk{}प्ल\appbrk{}व\appbrk{}न\appbrk{}त\appbrk{}र\appbrk{}ण\appbrk{}रा\appbrk{}त्रि\appbrk{}जा\appbrk{}ग\appbrk{}र\appbrk{}ण\appbrk{}भा\appbrk{}र\appbrk{}ह\appbrk{}रण- गज\appbrk{}तु\appbrk{}र\appbrk{}ग\appbrk{}र\appbrk{}थ\appbrk{}प\appbrk{}दा\appbrk{}ति\appbrk{}च\appbrk{}र्या\appbrk{}क\appbrk{}टु\appbrk{}क\appbrk{}षा\appbrk{}य\appbrk{}ति\appbrk{}क्त\appbrk{}रू\appbrk{}क्ष\appbrk{}ल\appbrk{}घु\appbrk{}शी\appbrk{}त\appbrk{}वी\appbrk{}र्य\appbrk{}शु\appbrk{}ष्क\appbrk{}शा\appbrk{}क\appbrk{}व\appbrk{}ल्लू\appbrk{}र\appbrk{}व\appbrk{}र\appbrk{}को\appbrk{}द्दा\appbrk{}ल\appbrk{}क\appbrk{}को\appbrk{}र\appbrk{}दू\appbrk{}ष\appbrk{}श्या\appbrk{}मा\appbrk{}क\appbrk{}नी\appbrk{}वार- मुद्ग\appbrk{}म\appbrk{}सू\appbrk{}रा\appbrk{}ढ\appbrk{}की\appbrk{}ह\appbrk{}रे\appbrk{}णु\appbrk{}क\appbrk{}ला\appbrk{}य\appbrk{}नि\appbrk{}ष्पा\appbrk{}वा\appbrk{}न\appbrk{}श\appbrk{}न\appbrk{}वि\appbrk{}ष\appbrk{}मा\appbrk{}श\appbrk{}ना\appbrk{}ध्य\appbrk{}श\appbrk{}न\appbrk{}वा\appbrk{}त\appbrk{}मू\appbrk{}त्र\appbrk{}पु\appbrk{}री\appbrk{}ष\appbrk{}शु\appbrk{}क्र\appbrk{}च्छ\appbrk{}र्दि\appbrk{}क्ष\appbrk{}व\appbrk{}थू\appbrk{}द्गा\appbrk{}र\appbrk{}बा\appbrk{}ष्प\appbrk{}वे\appbrk{}ग\appbrk{}वि\appbrk{}घा\appbrk{}ता\appbrk{}दि\appbrk{}भिर्}\ A; \skt{बल\appbrk{}द्वि\appbrk{}ग्र\appbrk{}हा\appbrk{}ति\appbrk{}व्य\appbrk{}वा\appbrk{}य\appbrk{}व्या\appbrk{}या\appbrk{}म\appbrk{}प्र\appbrk{}त\appbrk{}प\appbrk{}न\appbrk{}प्र\appbrk{}धा\appbrk{}व\appbrk{}ना\appbrk{}भि\appbrk{}घा\appbrk{}त\appbrk{}ल\appbrk{}ङ्घ\appbrk{}न\appbrk{}प्ल\appbrk{}व\appbrk{}न\appbrk{}प्र\appbrk{}त\appbrk{}र्प\appbrk{}न\appbrk{}जा\appbrk{}ग\appbrk{}र\appbrk{}ण\appbrk{}ह\appbrk{}य\appbrk{}र\appbrk{}थ\appbrk{}ग\appbrk{}ज\appbrk{}प्र\appbrk{}या\appbrk{}णा\appbrk{}ति\appbrk{}क\appbrk{}टु\appbrk{}क\appbrk{}क\appbrk{}षा\appbrk{}य\appbrk{}ति\appbrk{}क्त\appbrk{}ल\appbrk{}घु\appbrk{}रू\appbrk{}क्ष\appbrk{}शी\appbrk{}त\appbrk{}वी\appbrk{}र्य\appbrk{}शु\appbrk{}ष्क\appbrk{}सा\appbrk{}क\appbrk{}व\appbrk{}ल्लू\appbrk{}र\appbrk{}को\appbrk{}र\appbrk{}दू\appbrk{}ष\appbrk{}स्या\appbrk{}मा\appbrk{}क\appbrk{}नी\appbrk{}वा\appbrk{}र\appbrk{}मु\appbrk{}द्ग\appbrk{}म\appbrk{}सू\appbrk{}रा\appbrk{}ढ\appbrk{}की\appbrk{}ह\appbrk{}रे\appbrk{}णु\appbrk{}नि\appbrk{}ष्पा\appbrk{}वा\appbrk{}न\appbrk{}स\appbrk{}ना\appbrk{}ध्या\appbrk{}स\appbrk{}न\appbrk{}वा\appbrk{}त\appbrk{}मू\appbrk{}त्र\appbrk{}पु\appbrk{}री\appbrk{}ष\appbrk{}शु\appbrk{}क्र\appbrk{}क्ष\appbrk{}व\appbrk{}थू\appbrk{}द्गा\appbrk{}र\appbrk{}च्छ\appbrk{}र्दि\appbrk{}बा\appbrk{}ष्प\appbrk{}वे\appbrk{}ग\appbrk{}वि\appbrk{}घा\appbrk{}ता\appbrk{}दिभिः}\ N}}\skt{ }\edtext{\edtext{\skt{प्रको\-प\-मा\-प\-द्यते}}{\lemma{\skt{प्रकोपम्}}\Afootnote{N H; \skt{प्रकोप}\ A}}}{\lemma{\skt{आपद्यते}}\Afootnote{N H; \skt{मापद्यते}\ A}}\skt{~॥}\pend

\pstart\versepar \ledsidenote{\passref{1.21.20}}\edtext{\skt{स शीता\-भ्र\-प्र\-या\-तेषु}}{\lemma{\skt{स शीता\appbrk{}भ्र\appbrk{}प्र\appbrk{}या\appbrk{}तेषु}}\Afootnote{N H; \skt{स शीता\appbrk{}भ्र\appbrk{}प्र\appbrk{}वा\appbrk{}तेषु}\ A}}\skt{ }\edtext{\skt{घर्म्मार्त्ते च}}{\lemma{\skt{घर्म्मार्त्ते च}}\Afootnote{N H; \skt{घर्मान्ते च}\ A}}\skt{ विशेषतः~। }\pend
\pstart\versepar \skt{}\edtext{\edtext{\skt{प्रत्यू\-ष\-स्य\-प\-राह्णे}}{\lemma{\skt{प्रत्यूषस्य्}}\Afootnote{N H; \skt{प्रत्यूष}\ A}}\skt{ च}}{\lemma{\skt{अपराह्णे च}}\Afootnote{H; \skt{स्यपराह्णे तु}\ A; \skt{आपराह्णे च}\ N}}\skt{ }\edtext{\skt{जीर्ण्णान्ते च}}{\lemma{\skt{जीर्ण्णान्ते च}}\Afootnote{N H; \skt{जीर्णेऽन्ने च}\ A}}\skt{ }\edtext{\skt{प्रकुप्यति}}{\lemma{\skt{प्रकुप्यति}}\Afootnote{A N H\textsuperscript{pc}; \skt{कुप्यति}\ H\textsuperscript{ac}}}\skt{~॥}\pend

\pstart \ledsidenote{\passref{1.21.21}}\edtext{\edtext{\skt{क्रोध\-शो\-क\-भ\-या\-या\-सो\-प\-वा\-स\-वि\-द\-ग्ध\-मै\-थु\-नो\-प\-ग\-म\-न\-क\-ट्व\-म्ल\-ल\-व\-णो\-ष्ण\-ती\-क्ष्ण\-रू\-क्ष\-ल\-घु\-वि\-दा\-हि\-ति\-ल\-तै\-ल\-पि\-ण्या\-क\-कु\-ल\-त्थ\-स\-र्ष\-पा\-त\-सी\-ह\-रि\-त\-शा\-क\-व\-रा\-ह\-म\-त्स्या\-जा\-वि\-क\-मां\-स\-द\-धि\-म\-स्तु\-त\-क्र\-सौ\-वी\-र\-क\-सु\-रा\-म्ल\-फ\-ल\-क\-ट्व\-र\-प्र\-भृ\-ति\-भि\-र्व्वि\-शेषैः}}{\lemma{\skt{क्रोध\appbrk{}शो\appbrk{}क\appbrk{}भ\appbrk{}या\appbrk{}या\appbrk{}सो\appbrk{}प\appbrk{}वा\appbrk{}स\appbrk{}वि\appbrk{}द\appbrk{}ग्ध\appbrk{}मै\appbrk{}थु\appbrk{}नो\appbrk{}प\appbrk{}ग\appbrk{}म\appbrk{}न\appbrk{}क\appbrk{}ट्व\appbrk{}म्ल\appbrk{}ल\appbrk{}व\appbrk{}णो\appbrk{}ष्ण\appbrk{}ती\appbrk{}क्ष्ण\appbrk{}रू\appbrk{}क्ष\appbrk{}ल\appbrk{}घु\appbrk{}वि\appbrk{}दा\appbrk{}हि\appbrk{}ति\appbrk{}ल\appbrk{}तै\appbrk{}ल\appbrk{}पि\appbrk{}ण्या\appbrk{}क\appbrk{}कु\appbrk{}ल\appbrk{}त्थ\appbrk{}स\appbrk{}र्ष\appbrk{}पा\appbrk{}त\appbrk{}सी\appbrk{}ह\appbrk{}रि\appbrk{}त\appbrk{}शा\appbrk{}क\appbrk{}व\appbrk{}रा\appbrk{}ह\appbrk{}म\appbrk{}त्स्या\appbrk{}जा\appbrk{}वि\appbrk{}क\appbrk{}मां\appbrk{}स\appbrk{}द\appbrk{}धि\appbrk{}म\appbrk{}स्तु\appbrk{}त\appbrk{}क्र\appbrk{}सौ\appbrk{}वी\appbrk{}र\appbrk{}क\appbrk{}सु\appbrk{}रा\appbrk{}म्ल\appbrk{}फ\appbrk{}ल\appbrk{}क\appbrk{}ट्व\appbrk{}र\appbrk{}प्र\appbrk{}भृ\appbrk{}ति\appbrk{}भिर्}}\Afootnote{H\textsuperscript{pc}; \skt{क्रोध\appbrk{}शो\appbrk{}क\appbrk{}भ\appbrk{}या\appbrk{}या\appbrk{}सो\appbrk{}प\appbrk{}वा\appbrk{}स\appbrk{}वि\appbrk{}द\appbrk{}ग्ध\appbrk{}मै\appbrk{}थु\appbrk{}नो\appbrk{}प\appbrk{}ग\appbrk{}म\appbrk{}न\appbrk{}क\appbrk{}ट्व\appbrk{}म्ल\appbrk{}ल\appbrk{}वण- तीक्ष्णो\appbrk{}ष्ण\appbrk{}ल\appbrk{}घु\appbrk{}वि\appbrk{}दा\appbrk{}हि\appbrk{}ति\appbrk{}ल\appbrk{}तै\appbrk{}ल\appbrk{}पि\appbrk{}ण्या\appbrk{}क\appbrk{}कु\appbrk{}ल\appbrk{}त्थ\appbrk{}स\appbrk{}र्ष\appbrk{}पा\appbrk{}त\appbrk{}सी\appbrk{}ह\appbrk{}रि\appbrk{}त\appbrk{}क\appbrk{}शाक- गोधा\appbrk{}म\appbrk{}त्स्या\appbrk{}जा\appbrk{}वि\appbrk{}क\appbrk{}मां\appbrk{}स\appbrk{}द\appbrk{}धि\appbrk{}त\appbrk{}क्र\appbrk{}कू\appbrk{}र्चि\appbrk{}का\appbrk{}म\appbrk{}स्तु\appbrk{}सौ\appbrk{}वी\appbrk{}र\appbrk{}क\appbrk{}सु\appbrk{}रा\appbrk{}वि\appbrk{}का\appbrk{}रा\appbrk{}म्ल\appbrk{}फ\appbrk{}ल\appbrk{}क\appbrk{}ट्व\appbrk{}र\appbrk{}प्र\appbrk{}भृ\appbrk{}तिभिः}\ A; \skt{क्रोध\appbrk{}शो\appbrk{}क\appbrk{}भ\appbrk{}या\appbrk{}या\appbrk{}सो\appbrk{}प\appbrk{}वा\appbrk{}स\appbrk{}वि\appbrk{}द\appbrk{}ग्ध\appbrk{}मै\appbrk{}थु\appbrk{}नो\appbrk{}प\appbrk{}ग\appbrk{}म\appbrk{}न\appbrk{}क\appbrk{}ट्व\appbrk{}म्ब्ल\appbrk{}ल\appbrk{}व\appbrk{}णो\appbrk{}ष्ण\appbrk{}ती\appbrk{}क्ष्ण\appbrk{}रू\appbrk{}ष्ण\appbrk{}ल\appbrk{}घु\appbrk{}वि\appbrk{}दा\appbrk{}हि\appbrk{}ति\appbrk{}ल\appbrk{}तै\appbrk{}ल\appbrk{}पि\appbrk{}ण्या\appbrk{}क\appbrk{}कु\appbrk{}ल\appbrk{}त्थ\appbrk{}स\appbrk{}र्ष\appbrk{}पा\appbrk{}त\appbrk{}सी\appbrk{}ह\appbrk{}रि\appbrk{}त\appbrk{}शा\appbrk{}क\appbrk{}व\appbrk{}रा\appbrk{}ह\appbrk{}म\appbrk{}त्स्या\appbrk{}जा\appbrk{}वि\appbrk{}क\appbrk{}मा\appbrk{}म्स\appbrk{}द\appbrk{}धि\appbrk{}म\appbrk{}स्तु\appbrk{}त\appbrk{}क्र\appbrk{}सौ\appbrk{}वी\appbrk{}र\appbrk{}क\appbrk{}सु\appbrk{}रा\appbrk{}म्व्ल\appbrk{}फ\appbrk{}ल\appbrk{}क\appbrk{}ट्व\appbrk{}र\appbrk{}प्र\appbrk{}भृ\appbrk{}ति\appbrk{}भिर्}\ N; \skt{क्रोध\appbrk{}शो\appbrk{}क\appbrk{}भ\appbrk{}या\appbrk{}या\appbrk{}सो\appbrk{}प\appbrk{}वा\appbrk{}स\appbrk{}वि\appbrk{}द\appbrk{}ग्ध\appbrk{}मै\appbrk{}थु\appbrk{}नो\appbrk{}प\appbrk{}ग\appbrk{}म\appbrk{}न\appbrk{}क\appbrk{}ट्व\appbrk{}म्ल\appbrk{}ल\appbrk{}व\appbrk{}णो\appbrk{}ष्ण\appbrk{}ती\appbrk{}क्ष्ण\appbrk{}रू\appbrk{}क्ष\appbrk{}ल\appbrk{}घु\appbrk{}वि\appbrk{}दा\appbrk{}हि\appbrk{}ति\appbrk{}ल\appbrk{}तै\appbrk{}ल\appbrk{}पि\appbrk{}ण्या\appbrk{}क\appbrk{}कु\appbrk{}ल\appbrk{}त्थ\appbrk{}स\appbrk{}र्ष\appbrk{}पा\appbrk{}त\appbrk{}सी\appbrk{}ह\appbrk{}रि\appbrk{}त\appbrk{}शा\appbrk{}क\appbrk{}व\appbrk{}रा\appbrk{}ह\appbrk{}म\appbrk{}त्स्या\appbrk{}जा\appbrk{}वि\appbrk{}क\appbrk{}मां\appbrk{}स\appbrk{}द\appbrk{}धि\appbrk{}म\appbrk{}स्तु\appbrk{}त\appbrk{}क्र\appbrk{}सौ\appbrk{}वी\appbrk{}र\appbrk{}क\appbrk{}सु\appbrk{}रा\appbrk{}म्ल\appbrk{}फ\appbrk{}ल\appbrk{}क\appbrk{}ट्व\appbrk{}र\appbrk{}प्र\appbrk{}भृ\appbrk{}तिभि ः र्}\ H\textsuperscript{ac}}}}{\lemma{\skt{व्विशेषैः}}\Afootnote{N H; \emph{om.}\ A}}\skt{ }\edtext{\edtext{\skt{पित्त\-म्प्र\-को\-प\-मा\-प\-द्यते}}{\lemma{\skt{प्रकोपम्}}\Afootnote{N H; \skt{प्रकोप}\ A}}}{\lemma{\skt{आपद्यते}}\Afootnote{N H; \skt{मापद्यते}\ A}}\skt{~॥}\pend

\pstart\versepar \ledsidenote{\passref{1.21.22}}\edtext{\skt{तदुष्णे}}{\lemma{\skt{तदुष्णे}}\Afootnote{N H; \skt{तदुष्णैर्}\ A}}\skt{ }\edtext{\skt{वोष्णकाले च}}{\lemma{\skt{वोष्णकाले च}}\Afootnote{N H; \skt{उष्णकाले च}\ A}}\skt{ घनान्ते च }\edtext{\skt{विषेशतः}}{\lemma{\skt{विषेशतः}}\Afootnote{N H; \skt{विशेषतः}\ A}}\skt{~। }\pend
\pstart\versepar \skt{मध्याह्ने चार्द्धरात्रे च }\edtext{\edtext{\skt{जीर्य्यत्यन्ने}}{\lemma{\skt{जीर्य्यत्य्}}\Afootnote{N H; \skt{जीर्य}\ A}}\skt{ च}}{\lemma{\skt{अन्ने च}}\Afootnote{N H; \skt{त्यन्ने च}\ A}}\skt{ कुप्यति~॥}\pend

\pstart \ledsidenote{\passref{1.21.23}}\edtext{\edtext{\skt{दिवा\-स्व\-प्न\-व्या\-या\-मा\-ल\-स्य\-म\-धु\-रा\-म्ल\-ल\-व\-ण\-स्नि\-ग्ध\-शी\-त\-गु\-रु\-पि\-च्छि\-ला\-भि\-ष्य\-न्दि\-हा\-य\-न\-क\-य\-व\-नै\-ष\-धो\-त्क\-ट\-मा\-ष\-गो\-धू\-म\-ति\-ल\-पि\-ष्ट\-वि\-कृ\-त\-दु\-ग्ध\-द\-धि\-कृ\-श\-र\-पा\-य\-से\-क्षु\-वि\-का\-रा\-नू\-पो\-द\-क\-मां\-स\-व\-सा\-बि\-स\-मृ\-णा\-ल\-क\-से\-रु\-क\-शृ\-ङ्गा\-ट\-क\-म\-धु\-र\-व\-ल्ली\-फ\-ल\-स\-म\-स\-ना\-ध्य\-श\-न\-प्र\-भृ\-ति\-भि\-र्व्वि\-शेषैः}}{\lemma{\skt{दिवा\appbrk{}स्व\appbrk{}प्न\appbrk{}व्या\appbrk{}या\appbrk{}मा\appbrk{}ल\appbrk{}स्य\appbrk{}म\appbrk{}धु\appbrk{}रा\appbrk{}म्ल\appbrk{}ल\appbrk{}व\appbrk{}ण\appbrk{}स्नि\appbrk{}ग्ध\appbrk{}शी\appbrk{}त\appbrk{}गु\appbrk{}रु\appbrk{}पि\appbrk{}च्छि\appbrk{}ला\appbrk{}भि\appbrk{}ष्य\appbrk{}न्दि\appbrk{}हा\appbrk{}य\appbrk{}न\appbrk{}क\appbrk{}य\appbrk{}व\appbrk{}नै\appbrk{}ष\appbrk{}धो\appbrk{}त्क\appbrk{}ट\appbrk{}मा\appbrk{}ष\appbrk{}गो\appbrk{}धू\appbrk{}म\appbrk{}ति\appbrk{}ल\appbrk{}पि\appbrk{}ष्ट\appbrk{}वि\appbrk{}कृ\appbrk{}त\appbrk{}दु\appbrk{}ग्ध\appbrk{}द\appbrk{}धि\appbrk{}कृ\appbrk{}श\appbrk{}र\appbrk{}पा\appbrk{}य\appbrk{}से\appbrk{}क्षु\appbrk{}वि\appbrk{}का\appbrk{}रा\appbrk{}नू\appbrk{}पो\appbrk{}द\appbrk{}क\appbrk{}मां\appbrk{}स\appbrk{}व\appbrk{}सा\appbrk{}बि\appbrk{}स\appbrk{}मृ\appbrk{}णा\appbrk{}ल\appbrk{}क\appbrk{}से\appbrk{}रु\appbrk{}क\appbrk{}शृ\appbrk{}ङ्गा\appbrk{}ट\appbrk{}क\appbrk{}म\appbrk{}धु\appbrk{}र\appbrk{}व\appbrk{}ल्ली\appbrk{}फ\appbrk{}ल\appbrk{}स\appbrk{}म\appbrk{}स\appbrk{}ना\appbrk{}ध्य\appbrk{}श\appbrk{}न\appbrk{}प्र\appbrk{}भृ\appbrk{}ति\appbrk{}भिर्}}\Afootnote{\skt{दिवा\appbrk{}स्व\appbrk{}प्ना\appbrk{}व्या\appbrk{}या\appbrk{}मा\appbrk{}ल\appbrk{}स्य\appbrk{}म\appbrk{}धु\appbrk{}रा\appbrk{}म्ल\appbrk{}ल\appbrk{}व\appbrk{}ण\appbrk{}शी\appbrk{}त\appbrk{}स्नि\appbrk{}ग्ध\appbrk{}गु\appbrk{}रु\appbrk{}पि\appbrk{}च्छि\appbrk{}ला\appbrk{}भि\appbrk{}ष्य\appbrk{}न्दि\appbrk{}हा\appbrk{}य\appbrk{}न\appbrk{}क\appbrk{}य\appbrk{}व\appbrk{}क\appbrk{}नै\appbrk{}ष\appbrk{}धे\appbrk{}त्कट- माष\appbrk{}म\appbrk{}हा\appbrk{}मा\appbrk{}ष\appbrk{}गो\appbrk{}धू\appbrk{}म\appbrk{}ति\appbrk{}ल\appbrk{}पि\appbrk{}ष्ट\appbrk{}वि\appbrk{}कृ\appbrk{}ति\appbrk{}द\appbrk{}धि\appbrk{}दु\appbrk{}ग्ध\appbrk{}कृ\appbrk{}श\appbrk{}रा\appbrk{}पा\appbrk{}य\appbrk{}से\appbrk{}क्षु\appbrk{}वि\appbrk{}का\appbrk{}रा\appbrk{}नू\appbrk{}पौ\appbrk{}द\appbrk{}क\appbrk{}मां\appbrk{}स\appbrk{}व\appbrk{}सा\appbrk{}बि\appbrk{}स\appbrk{}मृ\appbrk{}णाल- कसे\appbrk{}रु\appbrk{}क\appbrk{}शृ\appbrk{}ङ्गा\appbrk{}ट\appbrk{}क\appbrk{}म\appbrk{}धु\appbrk{}र\appbrk{}व\appbrk{}ल्ली\appbrk{}फ\appbrk{}ल\appbrk{}स\appbrk{}म\appbrk{}श\appbrk{}ना\appbrk{}ध्य\appbrk{}श\appbrk{}न\appbrk{}प्र\appbrk{}भृ\appbrk{}तिभिः}\ A; \skt{दिवा\appbrk{}स्व\appbrk{}प्न\appbrk{}व्या\appbrk{}या\appbrk{}मा\appbrk{}ल\appbrk{}स्य\appbrk{}म\appbrk{}धु\appbrk{}रा\appbrk{}म्व्ल\appbrk{}ल\appbrk{}व\appbrk{}ण\appbrk{}स्नि\appbrk{}ग्ध\appbrk{}शी\appbrk{}त्गु\appbrk{}रु\appbrk{}पि\appbrk{}च्छि\appbrk{}ला\appbrk{}भि\appbrk{}ष्य\appbrk{}न्द\appbrk{}हा\appbrk{}य\appbrk{}न\appbrk{}क\appbrk{}य\appbrk{}व\appbrk{}नै\appbrk{}ष\appbrk{}धो\appbrk{}त्क\appbrk{}ट\appbrk{}मा\appbrk{}ष\appbrk{}गो\appbrk{}धू\appbrk{}म\appbrk{}ति\appbrk{}ल\appbrk{}पि\appbrk{}ष्ट\appbrk{}वि\appbrk{}कृ\appbrk{}त\appbrk{}दु\appbrk{}ग्ध\appbrk{}द\appbrk{}धि\appbrk{}कृ\appbrk{}स\appbrk{}र\appbrk{}पा\appbrk{}य\appbrk{}से\appbrk{}क्षु\appbrk{}वि\appbrk{}का\appbrk{}रा\appbrk{}नू\appbrk{}पो\appbrk{}द\appbrk{}क\appbrk{}मा\appbrk{}म्स\appbrk{}व\appbrk{}सा\appbrk{}बि\appbrk{}स\appbrk{}मृ\appbrk{}णा\appbrk{}ल\appbrk{}क\appbrk{}से\appbrk{}रु\appbrk{}क\appbrk{}शृ\appbrk{}ङ्ग\appbrk{}ट\appbrk{}क\appbrk{}म\appbrk{}धु\appbrk{}र\appbrk{}व\appbrk{}ल्ली\appbrk{}फ\appbrk{}ल\appbrk{}स\appbrk{}म\appbrk{}श\appbrk{}ना\appbrk{}ध्य\appbrk{}स\appbrk{}न\appbrk{}प्र\appbrk{}भृ\appbrk{}ति\appbrk{}भिर्}\ N; \skt{दिवास्वप्नो व्याया\appbrk{}मा\appbrk{}ल\appbrk{}स्य\appbrk{}म\appbrk{}धु\appbrk{}रा\appbrk{}म्ल\appbrk{}ल\appbrk{}व\appbrk{}ण\appbrk{}स्नि\appbrk{}ग्ध\appbrk{}शी\appbrk{}त\appbrk{}पि\appbrk{}च्छि\appbrk{}ला\appbrk{}भि\appbrk{}ष्य\appbrk{}न्दि\appbrk{}हा\appbrk{}य\appbrk{}न\appbrk{}क\appbrk{}य\appbrk{}व\appbrk{}नै\appbrk{}ष\appbrk{}धो\appbrk{}त्क\appbrk{}ट\appbrk{}मा\appbrk{}ष\appbrk{}गो\appbrk{}धू\appbrk{}म\appbrk{}ति\appbrk{}ल\appbrk{}पि\appbrk{}ष्ट\appbrk{}वि\appbrk{}कृ\appbrk{}त\appbrk{}दु\appbrk{}ग्ध\appbrk{}द\appbrk{}धि\appbrk{}कृ\appbrk{}श\appbrk{}र\appbrk{}पा\appbrk{}य\appbrk{}से\appbrk{}क्षु\appbrk{}वि\appbrk{}का\appbrk{}रा\appbrk{}नू\appbrk{}पो\appbrk{}द\appbrk{}क\appbrk{}मां\appbrk{}स\appbrk{}व\appbrk{}सा\appbrk{}बि\appbrk{}स\appbrk{}मृ\appbrk{}णा\appbrk{}ल\appbrk{}क\appbrk{}से\appbrk{}रु\appbrk{}क\appbrk{}शृ\appbrk{}ङ्गा\appbrk{}ट\appbrk{}क\appbrk{}म\appbrk{}धु\appbrk{}र\appbrk{}व\appbrk{}ल्ली\appbrk{}फ\appbrk{}ल\appbrk{}स\appbrk{}म\appbrk{}स\appbrk{}ना\appbrk{}ध्य\appbrk{}श\appbrk{}न\appbrk{}प्र\appbrk{}भृ\appbrk{}ति\appbrk{}भिर्}\ H\textsuperscript{ac}; \skt{दिवास्वप्नो व्याया\appbrk{}मा\appbrk{}ल\appbrk{}स्य\appbrk{}म\appbrk{}धु\appbrk{}रा\appbrk{}म्ल\appbrk{}ल\appbrk{}व\appbrk{}ण\appbrk{}स्नि\appbrk{}ग्ध\appbrk{}शी\appbrk{}त\appbrk{}गु\appbrk{}रु\appbrk{}पि\appbrk{}च्छि\appbrk{}ला\appbrk{}भि\appbrk{}ष्य\appbrk{}न्दि\appbrk{}हा\appbrk{}य\appbrk{}न\appbrk{}क\appbrk{}य\appbrk{}व\appbrk{}नै\appbrk{}ष\appbrk{}धो\appbrk{}त्क\appbrk{}ट\appbrk{}मा\appbrk{}ष\appbrk{}गो\appbrk{}धू\appbrk{}म\appbrk{}ति\appbrk{}ल\appbrk{}पि\appbrk{}ष्ट\appbrk{}वि\appbrk{}कृ\appbrk{}त\appbrk{}दु\appbrk{}ग्ध\appbrk{}द\appbrk{}धि\appbrk{}कृ\appbrk{}श\appbrk{}र\appbrk{}पा\appbrk{}य\appbrk{}से\appbrk{}क्षु\appbrk{}वि\appbrk{}का\appbrk{}रा\appbrk{}नू\appbrk{}पो\appbrk{}द\appbrk{}क\appbrk{}मां\appbrk{}स\appbrk{}व\appbrk{}सा\appbrk{}बि\appbrk{}स\appbrk{}मृ\appbrk{}णा\appbrk{}ल\appbrk{}क\appbrk{}से\appbrk{}रु\appbrk{}क\appbrk{}शृ\appbrk{}ङ्गा\appbrk{}ट\appbrk{}क\appbrk{}म\appbrk{}धु\appbrk{}र\appbrk{}व\appbrk{}ल्ली\appbrk{}फ\appbrk{}ल\appbrk{}स\appbrk{}म\appbrk{}स\appbrk{}ना\appbrk{}ध्य\appbrk{}श\appbrk{}न\appbrk{}प्र\appbrk{}भृ\appbrk{}ति\appbrk{}भिर्}\ H\textsuperscript{pc}}}}{\lemma{\skt{व्विशेषैः}}\Afootnote{N H; \emph{om.}\ A}}\skt{ श्लेष्मा }\edtext{\edtext{\skt{प्रको\-प\-मा\-प\-द्यते}}{\lemma{\skt{प्रकोपम्}}\Afootnote{N H; \skt{प्रकोप}\ A}}}{\lemma{\skt{आपद्यते}}\Afootnote{N H; \skt{मापद्यते}\ A}}\skt{~॥}\pend

\pstart\versepar \ledsidenote{\passref{1.21.24}}\edtext{\skt{स शीते}}{\lemma{\skt{स शीते}}\Afootnote{N H; \skt{स शीतैः}\ A}}\skt{ शीतकाले च वसन्ते च विशेषतः~। }\pend
\pstart\versepar \skt{पूर्व्वाह्णे च प्रदोषे च भुक्तमात्रे प्रकुप्यति~॥}\pend

\pstart \ledsidenote{\passref{1.21.25}}\skt{पित्त\-प्र\-को\-प\-णै\-रेव चाभीक्ष्णं }\edtext{\edtext{\edtext{\edtext{\edtext{\skt{द्रव\-स्नि\-ग्ध\-गु\-रु\-भि\-श्चा\-हा\-रै\-र्द्दि\-वा\-स्व\-प्न\-क्रो\-धा\-न\-ला\-त\-प\-श्र\-मा\-भि\-घा\-ता\-जी\-र्ण्ण\-वि\-रु\-द्धा\-ध्य\-श\-न\-प्र\-भृ\-ति\-भि\-र्व्वि\-शे\-षै\-र्र\-क्त\-म्प्र\-को\-प\-मा\-प\-द्यते}}{\lemma{\skt{चाहारैर्}}\Afootnote{N H; \skt{आहारैर्}\ A}}}{\lemma{\skt{द्दिवा\appbrk{}स्व\appbrk{}प्न\appbrk{}क्रो\appbrk{}धा\appbrk{}न\appbrk{}ला\appbrk{}त\appbrk{}प\appbrk{}श्र\appbrk{}मा\appbrk{}भि\appbrk{}घा\appbrk{}ता\appbrk{}जी\appbrk{}र्ण्ण\appbrk{}वि\appbrk{}रु\appbrk{}द्धा\appbrk{}ध्य\appbrk{}श\appbrk{}न\appbrk{}प्र\appbrk{}भृ\appbrk{}ति\appbrk{}भिर्}}\Afootnote{N H; \skt{दिवा\appbrk{}स्व\appbrk{}प्न\appbrk{}क्रो\appbrk{}धा\appbrk{}न\appbrk{}ला\appbrk{}त\appbrk{}प\appbrk{}श्र\appbrk{}मा\appbrk{}भि\appbrk{}घा\appbrk{}ता\appbrk{}जी\appbrk{}र्ण\appbrk{}वि\appbrk{}रु\appbrk{}द्धा\appbrk{}ध्य\appbrk{}श\appbrk{}ना\appbrk{}दि\appbrk{}भिर्}\ A}}}{\lemma{\skt{रक्तम्}}\Afootnote{N H; \skt{असृ}\ A}}}{\lemma{\skt{प्रकोपम्}}\Afootnote{N H; \skt{क्प्रकोप}\ A}}}{\lemma{\skt{आपद्यते}}\Afootnote{N H; \skt{मापद्यते}\ A}}\skt{~॥}\pend

\pstart\versepar \ledsidenote{\passref{1.21.26}}\edtext{\skt{यस्माद्रक्तम्विना}}{\lemma{\skt{यस्माद्रक्तम्}}\Afootnote{H; \skt{यस्माद्रक्तं}\ A; \skt{यस्माद्रक्तो}\ N}}\skt{ }\edtext{\skt{दोषैर्न्न}}{\lemma{\skt{दोषैर्न्न}}\Afootnote{N H; \skt{दोषैर्न}\ A}}\skt{ कदा\-चि\-त्प्र\-कु\-प्यति~। }\pend
\pstart\versepar \skt{तस्मात्तस्य }\edtext{\skt{यथा\-दो\-ष\-ङ्का\-ल\-म्वि\-द्या\-त्प्र\-को\-पनैः}}{\lemma{\skt{विद्यात्प्रकोपनैः}}\Afootnote{H; \skt{विद्यत्प्रकोपणे}\ A; \skt{विद्यीत्प्रकोपनैः}\ N}}\skt{~॥}\pend

\pstart \ledsidenote{\passref{1.21.27}}\edtext{\skt{तेषान्तु}}{\lemma{\skt{तेषान्तु}}\Afootnote{N H; \skt{तेषां}\ A}}\skt{ }\edtext{\skt{प्रको\-पा\-त्को\-ष्ठ\-तो\-द\-स\-ञ्च\-र\-णा\-म्ली\-का\-दा\-ह\-पि\-पा\-सा\-न्न\-द्वे\-ष\-हृ\-द\-यो\-त्क्लेदा}}{\lemma{\skt{प्रको\appbrk{}पा\appbrk{}त्को\appbrk{}ष्ठ\appbrk{}तो\appbrk{}द\appbrk{}स\appbrk{}ञ्च\appbrk{}र\appbrk{}णा\appbrk{}म्ली\appbrk{}का\appbrk{}दा\appbrk{}ह\appbrk{}पि\appbrk{}पा\appbrk{}सा\appbrk{}न्न\appbrk{}द्वे\appbrk{}ष\appbrk{}हृ\appbrk{}द\appbrk{}यो\appbrk{}त्क्लेदा}}\Afootnote{H; \skt{प्रको\appbrk{}पा\appbrk{}त्को\appbrk{}ष्ठ\appbrk{}तो\appbrk{}द\appbrk{}स\appbrk{}ञ्च\appbrk{}र\appbrk{}णा\appbrk{}म्ली\appbrk{}का\appbrk{}पि\appbrk{}पा\appbrk{}सा\appbrk{}प\appbrk{}रि\appbrk{}दा\appbrk{}हा\appbrk{}न्न\appbrk{}द्वे\appbrk{}ष\appbrk{}हृ\appbrk{}द\appbrk{}यो\appbrk{}त्क्ले\appbrk{}दाश्}\ A; \skt{प्रको\appbrk{}पा\appbrk{}त्को\appbrk{}ष्ठ\appbrk{}तो\appbrk{}द\appbrk{}स\appbrk{}ञ्च\appbrk{}रा\appbrk{}णा\appbrk{}म्व्ली\appbrk{}का\appbrk{}दा\appbrk{}ह\appbrk{}पि\appbrk{}पा\appbrk{}सा\appbrk{}न्न\appbrk{}द्वे\appbrk{}ष\appbrk{}हृ\appbrk{}द\appbrk{}यो\appbrk{}त्क्लेदा}\ N}}\skt{ }\edtext{\skt{भवन्ति}}{\lemma{\skt{भवन्ति}}\Afootnote{N H; \skt{च जायन्ते}\ A}}\skt{~। तत्र द्वितीयः क्रियाकालः~॥}\pend

\pstart \ledsidenote{\passref{1.21.28}}\skt{अत ऊर्ध्वम्प्रसरं वक्ष्यामः~। }\edtext{\skt{तेषा\-मे\-भि\-रा\-त\-ङ्क\-वि\-शेषैः}}{\lemma{\skt{आतङ्कविशेषैः}}\Afootnote{A H; \skt{अतङ्कविशेषैः}\ N}}\skt{ }\edtext{\skt{प्रकु\-पि\-ता\-ना\-म्पि\-ष्ट\-कि\-ण्वो\-द\-क\-स\-म\-वाय}}{\lemma{\skt{पिष्ट\appbrk{}कि\appbrk{}ण्वो\appbrk{}द\appbrk{}क\appbrk{}स\appbrk{}म\appbrk{}वाय}}\Afootnote{N H; \skt{किण्वो\appbrk{}द\appbrk{}क\appbrk{}पि\appbrk{}ष्ट\appbrk{}स\appbrk{}म\appbrk{}वाय}\ A}}\skt{ }\edtext{\skt{इवा\-भ्यु\-द्ग\-ता\-ना\-म्प्र\-सरो}}{\lemma{\skt{इवा\appbrk{}भ्यु\appbrk{}द्ग\appbrk{}ता\appbrk{}नाम्}}\Afootnote{H; \skt{इवोद्रिक्तानां}\ A; \skt{इवाभ्युद्गताना}\ N}}\skt{ भवति~। }\edtext{\edtext{\edtext{\skt{तेषा\-म्वा\-यु\-र्ग्ग\-ति\-म\-त्त्वा\-त्प्र\-स\-र\-ण\-हेतुः}}{\lemma{\skt{तेषाम्}}\Afootnote{A H; \skt{तेषा}\ N}}}{\lemma{\skt{वायुर्}}\Afootnote{N H; \skt{वायु}\ A}}}{\lemma{\skt{ग्गति\appbrk{}म\appbrk{}त्त्वा\appbrk{}त्प्र\appbrk{}स\appbrk{}र\appbrk{}ण\appbrk{}हेतुः}}\Afootnote{H; \skt{र्गति\appbrk{}म\appbrk{}त्त्वा\appbrk{}त्प्र\appbrk{}स\appbrk{}र\appbrk{}ण\appbrk{}हेतुः}\ A; \skt{गति\appbrk{}म\appbrk{}त्वा\appbrk{}त्प्र\appbrk{}स\appbrk{}र\appbrk{}णा\appbrk{}हेतुः}\ N}}\skt{~। }\edtext{\skt{सत्यप्याचैतन्ये स हि}}{\lemma{\skt{आचैतन्ये स हि}}\Afootnote{N H; \skt{अचैतन्ये~। स हि}\ A}}\skt{ रजोभूयिष्ठः~। रजश्च }\edtext{\skt{प्रव\-र्त्त\-क\-म्भा\-वा\-ना\-म्यथा}}{\lemma{\skt{भावानाम्}}\Afootnote{N H; \skt{सर्वभावानाम्}\ A}}\skt{ }\edtext{\edtext{\skt{महा\-नु\-द\-क\-स\-ञ्च\-योऽ\-ति\-प्र\-वृ\-द्ध\-त्वात्}}{\lemma{\skt{महा\appbrk{}नु\appbrk{}द\appbrk{}क\appbrk{}स\appbrk{}ञ्चयो}}\Afootnote{A H; \skt{महा\appbrk{}नु\appbrk{}द\appbrk{}क\appbrk{}स\appbrk{}ञ्चया}\ N}}\skt{~। }\edtext{\edtext{\edtext{\skt{सेतु\-म\-व\-दा\-र्या\-प\-रै\-रु\-द\-कै\-र्व्व्या\-मिश्रः}}{\lemma{\skt{अव\appbrk{}दा\appbrk{}र्या\appbrk{}प\appbrk{}रैर्}}\Afootnote{N H; \skt{मवदार्यापरे}\ A}}}{\lemma{\skt{उदकैर्}}\Afootnote{N H; \skt{णोदकेन}\ A}}}{\lemma{\skt{व्व्यामिश्रः}}\Afootnote{A H; \skt{व्व्यामिश्र}\ N}}}{\lemma{\skt{ऽति\appbrk{}प्र\appbrk{}वृ\appbrk{}द्ध\appbrk{}त्वात्~। सेतुम्}}\Afootnote{H; \skt{ऽतिवृद्धः सेतु}\ A; \skt{दति\appbrk{}प्र\appbrk{}वृ\appbrk{}द्ध\appbrk{}त्वात्~। सेतुम्}\ N}}\skt{ सर्व्वतः }\edtext{\skt{प्रधा\-व\-त्ये\-व\-न्दोषाः}}{\lemma{\skt{दोषाः}}\Afootnote{N H; \skt{दोषाः}\ A}}\skt{~। }\edtext{\skt{एकैकशो}}{\lemma{\skt{एकैकशो}}\Afootnote{N H; \skt{कदाचिदेकशो}\ A}}\skt{ }\edtext{\skt{द्वन्द्वशः}}{\lemma{\skt{द्वन्द्वशः}}\Afootnote{N H; \skt{द्विशः}\ A}}\skt{ समस्ताः शोणितसहिता वा }\edtext{\skt{तत्र}}{\lemma{\skt{तत्र}}\Afootnote{H; \skt{नेकधा प्रसरन्ति~। तद्यथा}\ A; \skt{तत्र वात}\ N}}\skt{ }\edtext{\skt{वातः}}{\lemma{\skt{वातः}}\Afootnote{A H; \emph{om.}\ N}}\skt{ पित्तं श्लेष्मा }\edtext{\skt{शोणितं}}{\lemma{\skt{शोणितं}}\Afootnote{A N; \skt{शोणितं~। 𑑛}\ H}}\skt{~। वातपित्ते~। }\edtext{\skt{वातश्लेष्माणौ}}{\lemma{\skt{वातश्लेष्माणौ}}\Afootnote{A H; \skt{वातश्लेष्मणौ}\ N}}\skt{~। पित्तश्लेष्माणौ~। वातशोणिते~। }\edtext{\skt{पित्तशोणिते}}{\lemma{\skt{पित्तशोणिते}}\Afootnote{A N H\textsuperscript{pc}; \skt{पित्तशोणि पित्तशोणिते}\ H\textsuperscript{ac}}}\skt{~। श्लेष्मशोणिते~। वात\-पि\-त्त\-शो\-णि\-तानि~। वात\-श्ले\-ष्म\-शो\-णि\-तानि~। पित्त\-श्ले\-ष्म\-शो\-णि\-तानि~। वातपित्तकफाः~। वात\-पि\-त्त\-क\-फ\-शो\-णि\-ता\-नीति~। एव\-म्प\-ञ्च\-द\-शधा }\edtext{\skt{प्रसरन्ति}}{\lemma{\skt{प्रसरन्ति}}\Afootnote{A H; \skt{प्रसरं भवति}\ N}}\skt{~॥}\pend

\pstart\versepar \ledsidenote{\passref{1.21.29}}\edtext{\edtext{\skt{कृत्स्नेऽर्धेऽवयवे}}{\lemma{\skt{ऽर्धे}}\Afootnote{A H; \skt{र्धे}\ N}}}{\lemma{\skt{ऽवयवे}}\Afootnote{A H; \skt{वयवे}\ N}}\skt{ }\edtext{\skt{चापि}}{\lemma{\skt{चापि}}\Afootnote{N H; \skt{वाऽपि}\ A}}\skt{ }\edtext{\skt{देहस्य}}{\lemma{\skt{देहस्य}}\Afootnote{N H; \skt{यत्राङ्गे}\ A}}\skt{ कुपितो भृशम्~। }\pend
\pstart\versepar \skt{दोषो }\edtext{\edtext{\skt{विकारम्भजते}}{\lemma{\skt{विकारम्}}\Afootnote{A H; \skt{विकाराम्}\ N}}}{\lemma{\skt{भजते}}\Afootnote{N H; \skt{नभसि}\ A}}\skt{ }\edtext{\skt{मेघो}}{\lemma{\skt{मेघो}}\Afootnote{N H; \skt{मेघ}\ A}}\skt{ }\edtext{\edtext{\skt{वृष्टिमिवाम्बरे}}{\lemma{\skt{वृष्टिम्}}\Afootnote{H; \emph{om.}\ A; \skt{वृष्टिर्}\ N}}}{\lemma{\skt{इवाम्बरे}}\Afootnote{N H; \skt{वत्तत्र वर्षति}\ A}}\skt{~॥}\pend

\pstart\versepar \ledsidenote{\passref{1.21.30}}\skt{नात्य\-र्थ\-ङ्कु\-पि\-त\-श्चापि लीनो }\edtext{\skt{मार्ग्गेषु}}{\lemma{\skt{मार्ग्गेषु}}\Afootnote{A H; \skt{मार्ग्गे}\ N}}\skt{ }\edtext{\skt{तिष्ठति}}{\lemma{\skt{तिष्ठति}}\Afootnote{A H; \skt{ष्ब्तिष्ठति}\ N}}\skt{~। }\pend
\pstart\versepar \skt{}\edtext{\skt{निःप्रत्यनीकः}}{\lemma{\skt{निःप्रत्यनीकः}}\Afootnote{H; \skt{निष्प्रत्यनीकः}\ A; \skt{निःप्रत्यनीकाः}\ N}}\skt{ कालेन }\edtext{\edtext{\skt{हेतुमासाद्य}}{\lemma{\skt{हेतुम्}}\Afootnote{N H; \skt{हेतु}\ A}}}{\lemma{\skt{आसाद्य}}\Afootnote{N H; \skt{मासाद्य}\ A}}\skt{ कुप्यति~॥}\pend

\pstart \ledsidenote{\passref{1.21.31}}\edtext{\skt{नगतस्य}}{\lemma{\skt{नगतस्य}}\Afootnote{K; \skt{तत्र वायोः पित्त\appbrk{}स्था\appbrk{}न\appbrk{}ग\appbrk{}तस्य पित्त\appbrk{}व\appbrk{}त्प्र\appbrk{}ती\appbrk{}कारः, पित्तस्य च कफ\appbrk{}स्था\appbrk{}न\appbrk{}ग\appbrk{}तस्य कफवत्, कफस्य च वात\appbrk{}स्था\appbrk{}न\appbrk{}ग\appbrk{}तस्य}\ A; \skt{तत्र वायो पित्त\appbrk{}स्था\appbrk{}न\appbrk{}ग\appbrk{}तस्य पित्तवत्क्रिया~। पित्तस्य च कफ\appbrk{}स्था\appbrk{}न\appbrk{}ग\appbrk{}तस्य कफवत्~। कफस्य च वात\appbrk{}स्था\appbrk{}न\appbrk{}ग\appbrk{}तस्य}\ N\textsuperscript{ac}; \skt{तत्र वायोः पित्त\appbrk{}स्था\appbrk{}न\appbrk{}ग\appbrk{}तस्य पित्तवत्क्रिया पित्तस्य च कफ\appbrk{}स्था\appbrk{}न\appbrk{}ग\appbrk{}तस्य कफवत्कफस्य च वात\appbrk{}स्था\appbrk{}न\appbrk{}ग\appbrk{}तस्य}\ N\textsuperscript{pc} H}}\skt{ }\edtext{\skt{वात\-व\-त्क्रि\-या\-वि\-भागः}}{\lemma{\skt{वात\appbrk{}व\appbrk{}त्क्रि\appbrk{}या\appbrk{}वि\appbrk{}भागः}}\Afootnote{K N H; \skt{वातवत् ; एष क्रियाविभागः}\ A}}\skt{~॥}\pend

\pstart \ledsidenote{\passref{1.21.32}}\skt{एवं }\edtext{\skt{प्रक्षुभितानां}}{\lemma{\skt{प्रक्षुभितानां}}\Afootnote{K N H; \skt{प्रकुपितानां}\ A}}\skt{ प्रसरतां }\edtext{\edtext{\skt{विमा\-र्ग\-ग\-म\-न\-मा\-टोपो}}{\lemma{\skt{विमा\appbrk{}र्ग\appbrk{}ग\appbrk{}म\appbrk{}नम्}}\Afootnote{K N H; \skt{वायो\appbrk{}र्वि\appbrk{}मा\appbrk{}र्ग\appbrk{}ग\appbrk{}मना}\ A}}}{\lemma{\skt{आटोपो}}\Afootnote{K N H; \skt{टोपौ, ओष\appbrk{}चो\appbrk{}ष\appbrk{}प\appbrk{}रि\appbrk{}दाह}\ A}}\skt{ }\edtext{\edtext{\edtext{\skt{धूमा\-य\-न\-म\-रो\-च\-क\-श्छ\-र्द्दि\-रिति}}{\lemma{\skt{धूमायनम्}}\Afootnote{K N H; \skt{धूमायनानि पित्तस्य}\ A}}}{\lemma{\skt{अरोचकश्}}\Afootnote{K N H; \skt{अरो\appbrk{}च\appbrk{}का\appbrk{}वि\appbrk{}पा\appbrk{}का\appbrk{}ङ्ग\appbrk{}सा\appbrk{}दाश्}\ A}}}{\lemma{\skt{इति}}\Afootnote{K N H; \skt{चेति}\ A}}\skt{ }\edtext{\skt{लिङ्गानि}}{\lemma{\skt{लिङ्गानि}}\Afootnote{K N H; \skt{श्लेष्मणो लिङ्गानि}\ A}}\skt{ भवन्ति~। तत्र तृतीयः क्रियाकालः~॥}\pend

\pstart \ledsidenote{\passref{1.21.33}}\skt{अत }\edtext{\skt{ऊर्द्ध्वं}}{\lemma{\skt{ऊर्द्ध्वं}}\Afootnote{A K; \skt{ऊर्द्ध्व}\ N H}}\skt{ स्थान\-सं\-श्र\-य\-म्व\-क्ष्यामः~। एवं }\edtext{\skt{खलु}}{\lemma{\skt{खलु}}\Afootnote{K N H; \emph{om.}\ A}}\skt{ }\edtext{\edtext{\edtext{\skt{प्रसृ\-ता\-स्तां\-स्ता\-ञ्छ\-री\-र\-प्र\-दे\-शा\-ना\-गम्य}}{\lemma{\skt{प्रसृतास्}}\Afootnote{K N; \skt{प्रकुपि}\ A; \skt{प्रसृ}\ H}}}{\lemma{\skt{तांस्ताञ्}}\Afootnote{K; \skt{ता तांस्ता}\ A; \skt{तास्ताञ्}\ N; \skt{तस्तास्ताञ्}\ H}}\skt{ }\edtext{\skt{तांस्ता\-न्व्या\-धी\-ञ्ज\-न\-यन्ति}}{\lemma{\skt{तान्व्याधीञ्}}\Afootnote{K H; \skt{स्तान्व्याधी}\ A; \skt{तां व्याधीं}\ N}}}{\lemma{\skt{छरी\appbrk{}र\appbrk{}प्र\appbrk{}दे\appbrk{}शा\appbrk{}ना\appbrk{}गम्य तांस्}}\Afootnote{K N H; \skt{ञ्शरी\appbrk{}र\appbrk{}प्र\appbrk{}दे\appbrk{}शा\appbrk{}ना\appbrk{}गम्य तां}\ A}}\skt{~। तं }\edtext{\skt{प्रति}}{\lemma{\skt{प्रति}}\Afootnote{K N H; \emph{om.}\ A}}\skt{ }\edtext{\skt{प्रतिरोगं}}{\lemma{\skt{प्रतिरोगं}}\Afootnote{K N H; \skt{ति}\ A}}\skt{ }\edtext{\skt{वक्ष्यामः~। ते}}{\lemma{\skt{वक्ष्यामः~। ते}}\Afootnote{K\textsuperscript{pc}; \skt{ते}\ A; \skt{वक्ष्यामः}\ K\textsuperscript{ac} N; \skt{वक्ष्यामः~। ये}\ H}}\skt{ }\edtext{\edtext{\skt{यदोदर्धयः}}{\lemma{\skt{यदोदर्}}\Afootnote{K\textsuperscript{pc}; \skt{यदो}\ A; \skt{यदोद}\ K\textsuperscript{ac}; \skt{यदा ए}\ N; \skt{दोषास्ते}\ H}}}{\lemma{\skt{धयः}}\Afootnote{K; \skt{दर}\ A; \skt{व}\ N; \skt{ऽधः}\ H}}\skt{ }\edtext{\skt{सन्नि\-वे\-श\-ङ्कु\-र्वन्ति ते}}{\lemma{\skt{कुर्वन्ति ते}}\Afootnote{K N H; \skt{कुर्वन्ति तदा}\ A}}\skt{ }\edtext{\edtext{\skt{गुल्म\-वृ\-ध्यु\-द\-रा\-ग्नि\-स\-ङ्गा\-ना\-ह\-वि\-सू\-चि\-का\-ती\-सा\-र\-प्र\-भृ\-ती\-ञ्ज\-न\-यन्ति}}{\lemma{\skt{गुल्म\appbrk{}वृ\appbrk{}ध्यु\appbrk{}द\appbrk{}रा\appbrk{}ग्नि\appbrk{}स\appbrk{}ङ्गा\appbrk{}ना\appbrk{}ह\appbrk{}वि\appbrk{}सू\appbrk{}चि\appbrk{}का\appbrk{}ती\appbrk{}सा\appbrk{}र\appbrk{}प्र\appbrk{}भृ\appbrk{}तीञ्}}\Afootnote{K N; \skt{गुल्म विद्र\appbrk{}ध्यु\appbrk{}द\appbrk{}रा\appbrk{}ग्नि\appbrk{}स\appbrk{}ङ्गा\appbrk{}ना\appbrk{}ह\appbrk{}वि\appbrk{}सू\appbrk{}चि\appbrk{}का\appbrk{}ति\appbrk{}सा\appbrk{}र\appbrk{}प्र\appbrk{}भृ\appbrk{}तीञ्}\ A; \skt{गुल्म\appbrk{}वृ\appbrk{}द्ध्यु\appbrk{}द\appbrk{}रा\appbrk{}ग्नि\appbrk{}स\appbrk{}ङ्गा\appbrk{}ना\appbrk{}ह\appbrk{}वि\appbrk{}सू\appbrk{}चि\appbrk{}का\appbrk{}ती\appbrk{}सा\appbrk{}र\appbrk{}प्र\appbrk{}भृती}\ H}}}{\lemma{\skt{जनयन्ति}}\Afootnote{A K N; \skt{न्जनयन्ति}\ H}}\skt{~॥ बस्तिगताः }\edtext{\skt{प्रमे\-हा\-श्म\-री\-मू\-त्रा\-घा\-त\-मू\-त्र\-दो\-ष\-प्र\-भृ\-तीन्~॥ गुदगतास्तु}}{\lemma{\skt{प्रमे\appbrk{}हा\appbrk{}श्म\appbrk{}री\appbrk{}मू\appbrk{}त्रा\appbrk{}घा\appbrk{}त\appbrk{}मू\appbrk{}त्र\appbrk{}दो\appbrk{}ष\appbrk{}प्र\appbrk{}भृ\appbrk{}तीन्~॥ गुदगतास्तु}}\Afootnote{K H; \skt{प्रमे\appbrk{}हा\appbrk{}श्म\appbrk{}री\appbrk{}मू\appbrk{}त्रा\appbrk{}घा\appbrk{}त\appbrk{}मू\appbrk{}त्र\appbrk{}दोष}\ A; \skt{प्रमेहाष्मरी मूत्रा\appbrk{}घा\appbrk{}त\appbrk{}मू\appbrk{}त्र\appbrk{}दो\appbrk{}ष\appbrk{}प्र\appbrk{}भृ\appbrk{}तीन्~॥ गुदगतास्तु}\ N\textsuperscript{ac}; \skt{प्रमे\appbrk{}हा\appbrk{}ष्म\appbrk{}रा\appbrk{}मू\appbrk{}त्रा\appbrk{}घा\appbrk{}त\appbrk{}मू\appbrk{}त्र\appbrk{}दो\appbrk{}ष\appbrk{}प्र\appbrk{}भृ\appbrk{}तीन्~॥ गुदगतास्तु}\ N\textsuperscript{pc}}}\skt{ }\edtext{\skt{भग\-न्द\-रा\-र्स\-प्र\-भृ\-तीन्~॥ मेढ्रगतास्तु}}{\lemma{\skt{भग\appbrk{}न्द\appbrk{}रा\appbrk{}र्स\appbrk{}प्र\appbrk{}भृ\appbrk{}तीन्~॥ मेढ्रगतास्तु}}\Afootnote{K; \skt{प्रभृतीन्; मेढ्रगता नि}\ A; \skt{भग\appbrk{}न्द\appbrk{}रा\appbrk{}र्शः\appbrk{}प्र\appbrk{}भृ\appbrk{}ती\appbrk{}न्मे\appbrk{}ढ्र\appbrk{}ग\appbrk{}तास्तु}\ N H}}\skt{ }\edtext{\skt{परि\-व\-र्त्ति\-का\-प\-दं\-श\-शू\-क\-दो\-ष\-प्र\-भृ\-तीन्~॥ वृषणगता}}{\lemma{\skt{परि\appbrk{}व\appbrk{}र्त्ति\appbrk{}का\appbrk{}प\appbrk{}दं\appbrk{}श\appbrk{}शू\appbrk{}क\appbrk{}दो\appbrk{}ष\appbrk{}प्र\appbrk{}भृ\appbrk{}तीन्~॥ वृषणगता}}\Afootnote{K H; \skt{रुद्ध\appbrk{}प्र\appbrk{}क\appbrk{}शो\appbrk{}प\appbrk{}दं\appbrk{}श\appbrk{}शू\appbrk{}क\appbrk{}दो\appbrk{}ष\appbrk{}प्र\appbrk{}भृ\appbrk{}तीन्; गुदगता भग\appbrk{}न्द\appbrk{}रा\appbrk{}र्शः\appbrk{}प्र\appbrk{}भृ\appbrk{}तीन्; वृषणगता}\ A; \skt{परि\appbrk{}व\appbrk{}र्त्ति\appbrk{}को\appbrk{}प\appbrk{}दं\appbrk{}श\appbrk{}शू\appbrk{}क\appbrk{}दो\appbrk{}ष\appbrk{}प्र\appbrk{}भृ\appbrk{}तीन्~॥ वृषणगतास्तु}\ N}}\skt{ }\edtext{\skt{वृषणवृद्धीन्~॥ ऊर्द्ध्वजत्रुगता}}{\lemma{\skt{वृषणवृद्धीन्~॥ ऊर्द्ध्वजत्रुगता}}\Afootnote{K N H; \skt{वृद्धीः; ऊर्ध्वजत्रुगता}\ A}}\skt{ }\edtext{\skt{ऊर्द्ध्वगा}}{\lemma{\skt{ऊर्द्ध्वगा}}\Afootnote{K N H; \skt{स्तूर्ध्व}\ A}}\skt{ }\edtext{\skt{गल\-ग\-ण्डा\-प\-ची\-प्र\-भृ\-तीन्}}{\lemma{\skt{गल\appbrk{}ग\appbrk{}ण्डा\appbrk{}प\appbrk{}ची\appbrk{}प्र\appbrk{}भृ\appbrk{}तीन्}}\Afootnote{K N H; \skt{जान्}\ A}}\skt{~॥ }\edtext{\skt{त्वङ्मा\-न्स\-सो\-णि\-त\-गताः}}{\lemma{\skt{त्वङ्मा\appbrk{}न्स\appbrk{}सो\appbrk{}णि\appbrk{}त\appbrk{}गताः}}\Afootnote{K; \skt{त्वङ्मं\appbrk{}स\appbrk{}शो\appbrk{}णि\appbrk{}तस्थाः}\ A; \skt{त्वक्मा\appbrk{}न्स\appbrk{}शो\appbrk{}णि\appbrk{}त\appbrk{}गताः}\ N; \skt{त्वङ्मा\appbrk{}न्स\appbrk{}शो\appbrk{}णि\appbrk{}त\appbrk{}गताः}\ H}}\skt{ }\edtext{\skt{क्षुद्र\-रो\-गा\-न्कु\-ष्ठादीं}}{\lemma{\skt{क्षुद्र\appbrk{}रो\appbrk{}गा\appbrk{}न्कु\appbrk{}ष्ठादीं}}\Afootnote{K; \skt{क्षुद्र\appbrk{}रो\appbrk{}गा\appbrk{}न्कु\appbrk{}ष्ठानि}\ A; \skt{क्षुद्ररोगां कुष्ठादी}\ N; \skt{क्षुद्र\appbrk{}रो\appbrk{}गा\appbrk{}न्कु\appbrk{}ष्ठादी}\ H}}\skt{ }\edtext{\skt{विसर्पांश्च}}{\lemma{\skt{विसर्पांश्च}}\Afootnote{K; \skt{विसर्पांश्च}\ A; \skt{न्विसर्पाश्च}\ N; \skt{न्विसर्प्पाश्च}\ H}}\skt{~॥ }\edtext{\skt{मांसगता}}{\lemma{\skt{मांसगता}}\Afootnote{K N H; \skt{मेदोगता}\ A}}\skt{ }\edtext{\skt{ग्रन्थ्य\-प\-च्य\-र्बु\-द\-ग\-ल\-ग\-ण्डा\-ल\-जी\-प्र\-भृ\-तीन्~॥ अस्थिगता}}{\lemma{\skt{ग्रन्थ्य\appbrk{}प\appbrk{}च्य\appbrk{}र्बु\appbrk{}द\appbrk{}ग\appbrk{}ल\appbrk{}ग\appbrk{}ण्डा\appbrk{}ल\appbrk{}जी\appbrk{}प्र\appbrk{}भृ\appbrk{}तीन्~॥ अस्थिगता}}\Afootnote{A K H; \skt{ग्रन्थि\appbrk{}अ\appbrk{}प\appbrk{}च्य\appbrk{}र्बु\appbrk{}द\appbrk{}ग\appbrk{}ल\appbrk{}ग\appbrk{}ण्डा\appbrk{}ल\appbrk{}जी\appbrk{}प्र\appbrk{}भृ\appbrk{}तीन्~॥ अस्थिगता}\ N}}\skt{ }\edtext{\skt{विद्र\-ध्य\-नु\-श\-यी\-प्र\-भृ\-तीन्~॥ पादगताः}}{\lemma{\skt{विद्र\appbrk{}ध्य\appbrk{}नु\appbrk{}श\appbrk{}यी\appbrk{}प्र\appbrk{}भृ\appbrk{}तीन्~॥ पादगताः}}\Afootnote{A K H; \skt{विद्र\appbrk{}ध्य\appbrk{}नु\appbrk{}श\appbrk{}यी\appbrk{}प्र\appbrk{}भृ\appbrk{}त।न्। पादगताः}\ N}}\skt{ }\edtext{\skt{श्लीप\-द\-वा\-त\-सो\-णि\-त\-प्र\-भृ\-तीन्~॥ सर्वगताः}}{\lemma{\skt{श्लीप\appbrk{}द\appbrk{}वा\appbrk{}त\appbrk{}सो\appbrk{}णि\appbrk{}त\appbrk{}प्र\appbrk{}भृ\appbrk{}तीन्~॥ सर्वगताः}}\Afootnote{K N; \skt{श्लीप\appbrk{}द\appbrk{}वा\appbrk{}त\appbrk{}शो\appbrk{}णि\appbrk{}त\appbrk{}वा\appbrk{}त\appbrk{}क\appbrk{}ण्ट\appbrk{}क\appbrk{}प्र\appbrk{}भृ\appbrk{}तीन्; सर्वाङ्गगता}\ A; \skt{श्लीप\appbrk{}द\appbrk{}वा\appbrk{}त\appbrk{}शो\appbrk{}णि\appbrk{}त\appbrk{}प्र\appbrk{}भृ\appbrk{}तीन्~॥ सर्व्वगताः}\ H}}\skt{ }\edtext{\skt{ज्वर\-शु\-क्र\-दो\-ष\-स\-र्वा\-ङ्ग\-रो\-ग\-प्र\-भृ\-तीन्~॥ }\edtext{\edtext{\skt{एवमन्येष्वपि}}{\lemma{\skt{अन्येष्व्}}\Afootnote{K N H; \emph{om.}\ A}}}{\lemma{\skt{अपि}}\Afootnote{K N H; \emph{om.}\ A}}}{\lemma{\skt{ज्वर\appbrk{}शु\appbrk{}क्र\appbrk{}दो\appbrk{}ष\appbrk{}स\appbrk{}र्वा\appbrk{}ङ्ग\appbrk{}रो\appbrk{}ग\appbrk{}प्र\appbrk{}भृ\appbrk{}तीन्~॥ एवम्}}\Afootnote{K N H; \skt{ज्वर\appbrk{}स\appbrk{}र्वा\appbrk{}ङ्ग\appbrk{}रो\appbrk{}ग\appbrk{}प्रभृ}\ A}}\skt{ }\edtext{\skt{स्थानेषु}}{\lemma{\skt{स्थानेषु}}\Afootnote{K N H; \emph{om.}\ A}}\skt{ }\edtext{\skt{रोगाणां}}{\lemma{\skt{रोगाणां}}\Afootnote{K; \emph{om.}\ A; \skt{रोगाणा}\ N; \skt{रोगानां}\ H}}\skt{ }\edtext{\skt{दोषसन्निवेशं}}{\lemma{\skt{दोषसन्निवेशं}}\Afootnote{K H; \skt{तीन्; ते}\ A; \skt{दोषआसन्निवे}\ N\textsuperscript{ac}; \skt{दोषसन्निवे}\ N\textsuperscript{pc}}}\skt{ }\edtext{\skt{जानीयात्~॥ }\edtext{\edtext{\skt{तेषा\-मे\-व\-म\-भि\-स\-न्नि\-वि\-ष्टा\-ना\-म्पू\-र्व\-रू\-प\-प्रा\-दु\-र्भावो}}{\lemma{\skt{एवम्}}\Afootnote{K N H; \skt{एव}\ A}}}{\lemma{\skt{अभि\appbrk{}स\appbrk{}न्नि\appbrk{}वि\appbrk{}ष्टा\appbrk{}नाम्}}\Afootnote{K N H; \skt{मभि\appbrk{}स\appbrk{}न्नि\appbrk{}वि\appbrk{}ष्टानां}\ A}}}{\lemma{\skt{जानीयात्~॥ तेषाम्}}\Afootnote{K N H; \skt{षां}\ A}}\skt{ }\edtext{\skt{भवति}}{\lemma{\skt{भवति}}\Afootnote{K N H; \skt{तं प्रति}\ A}}\skt{~। }\edtext{\skt{तत्र चतुर्थः}}{\lemma{\skt{तत्र चतुर्थः}}\Afootnote{K N H; \skt{रोगं वक्ष्यामः~। तत्रSS.१.२०. पूर्व\appbrk{}रू\appbrk{}प\appbrk{}ग\appbrk{}तेषु चतुर्थः}\ A}}\skt{ क्रियाकालः~॥}\edtext{}{\lemma{\skt{जनयन्ति~। तं}}\Afootnote{K N H; \skt{न्जनयन्}\ A}}\pend

\pstart \ledsidenote{\passref{1.21.34}}\skt{अत ऊर्द्ध्वं }\edtext{\skt{व्याधि\-द\-र्स\-न\-म्व\-क्ष्यामः}}{\lemma{\skt{व्याधिदर्सनम्}}\Afootnote{K N; \skt{व्याधेर्दर्शनं}\ A; \skt{व्याधिदर्शनम्}\ H}}\skt{~॥ }\edtext{\skt{सोफा\-र्बु\-द\-ग्र\-न्थि\-वि\-द्र\-धि\-वि\-स\-र्पा\-दीनां}}{\lemma{\skt{सोफा\appbrk{}र्बु\appbrk{}द\appbrk{}ग्र\appbrk{}न्थि\appbrk{}वि\appbrk{}द्र\appbrk{}धि\appbrk{}वि\appbrk{}स\appbrk{}र्पा\appbrk{}दीनां}}\Afootnote{K; \skt{शोफा\appbrk{}र्बु\appbrk{}द\appbrk{}ग्र\appbrk{}न्थि\appbrk{}वि\appbrk{}द्र\appbrk{}धि\appbrk{}वि\appbrk{}स\appbrk{}र्प\appbrk{}प्र\appbrk{}भृ\appbrk{}तीनां}\ A; \skt{शोफा\appbrk{}बु\appbrk{}द\appbrk{}ग्र\appbrk{}न्थि\appbrk{}वि\appbrk{}द्र\appbrk{}धि\appbrk{}वि\appbrk{}स\appbrk{}र्प्पा\appbrk{}दीनां}\ N; \skt{शोफा\appbrk{}र्ब्बु\appbrk{}द\appbrk{}ग्र\appbrk{}न्थि\appbrk{}वि\appbrk{}द्र\appbrk{}धि\appbrk{}वि\appbrk{}स\appbrk{}र्प्पा\appbrk{}दीनां}\ H}}\skt{ प्रव्य\-क्त\-ल\-क्ष\-णता ज्वरा\-ती\-सा\-र\-प्र\-भृ\-ती\-नाञ्च~। तत्र }\edtext{\skt{पञ्च\-म\-क्रि\-या\-कालः}}{\lemma{\skt{पञ्च\appbrk{}म\appbrk{}क्रि\appbrk{}या\appbrk{}कालः}}\Afootnote{K; \skt{पञ्चमः क्रियाकालः}\ A H N}}\skt{~॥}\pend

\pstart \ledsidenote{\passref{1.21.35}}\skt{अत }\edtext{\edtext{\skt{ऊर्द्ध्व\-म\-स्या\-व\-दी\-र्ण्णस्य}}{\lemma{\skt{ऊर्द्ध्वम्}}\Afootnote{K N H; \skt{ऊर्ध्व}\ A}}}{\lemma{\skt{अस्यावदीर्ण्णस्य}}\Afootnote{K N H; \skt{मेते\appbrk{}षा\appbrk{}म\appbrk{}व\appbrk{}दी\appbrk{}र्णानां}\ A}}\skt{ }\edtext{\edtext{\skt{व्रण\-भा\-व\-मा\-प\-न्नस्य}}{\lemma{\skt{व्रणभावम्}}\Afootnote{K N H; \skt{व्रणभाव}\ A}}}{\lemma{\skt{आपन्नस्य}}\Afootnote{K N H; \skt{मापन्नानां}\ A}}\skt{ षष्ठः क्रियाकालः~॥ }\edtext{\skt{ज्वरा\-ती\-सा\-र\-प्र\-भृ\-ती\-नाञ्च}}{\lemma{\skt{ज्वरा\appbrk{}ती\appbrk{}सा\appbrk{}र\appbrk{}प्र\appbrk{}भृ\appbrk{}ती\appbrk{}नाञ्च}}\Afootnote{K N H; \skt{ज्वरा\appbrk{}ति\appbrk{}सा\appbrk{}र\appbrk{}प्र\appbrk{}भृ\appbrk{}तीनां च}\ A}}\skt{ }\edtext{\skt{दीर्घ\-का\-ला\-नु\-बन्धः}}{\lemma{\skt{दीर्घ\appbrk{}का\appbrk{}ला\appbrk{}नु\appbrk{}बन्धः}}\Afootnote{A K; \skt{दीर्घ\appbrk{}का\appbrk{}ला\appbrk{}नु\appbrk{}बद्धः}\ N H}}\skt{~। }\edtext{\skt{तत्रा\-प्र\-ति\-क्रि\-य\-माणो}}{\lemma{\skt{तत्रा\appbrk{}प्र\appbrk{}ति\appbrk{}क्रि\appbrk{}य\appbrk{}माणो}}\Afootnote{K N H; \skt{तत्रा\appbrk{}प्र\appbrk{}ति\appbrk{}क्रि\appbrk{}य\appbrk{}माणे}\ A}}\skt{ }\edtext{\skt{साध्यतामुपैति}}{\lemma{\skt{उपैति}}\Afootnote{K N H; \skt{उपयान्ति}\ A}}\skt{~॥}\pend

\pstart\versepar \ledsidenote{\passref{1.21.36}}\edtext{\skt{भवति}}{\lemma{\skt{भवति}}\Afootnote{H; \skt{भवन्ति}\ A; \emph{om.}\ K N}}\skt{ चात्र~॥ }\pend
\pstart\versepar \skt{सञ्चयञ्च प्रकोपञ्च प्रसरं स्थानसंश्रयं~। }\pend
\pstart\versepar \skt{व्यक्तिम्भेदञ्च यो वेत्ति दोषाणां स }\edtext{\skt{भवेद्भिषक्}}{\lemma{\skt{भवेद्भिषक्}}\Afootnote{K H; \skt{भवेद्भिषक्}\ A; \skt{भूवे भिषक्}\ N\textsuperscript{ac}; \skt{भवे भिषक्}\ N\textsuperscript{pc}}}\skt{~॥}\edtext{}{\lemma{\skt{चात्र~॥ सञ्चयञ्च}}\Afootnote{A; \skt{सञ्चयञ्च}\ K N H}}\pend

\pstart\versepar \ledsidenote{\passref{1.21.37}}\skt{सञ्चये }\edtext{\skt{पहृता}}{\lemma{\skt{पहृता}}\Afootnote{K; \skt{ऽपहृता}\ A H; \skt{प्रहृता}\ N}}\skt{ दोषा लभन्ते नोत्तरा गतीः~। }\pend
\pstart\versepar \skt{ते }\edtext{\skt{तूत्तरासु}}{\lemma{\skt{तूत्तरासु}}\Afootnote{A K H; \skt{ष्वुत्तरासु}\ N}}\skt{ गतिषु भवन्ति }\edtext{\skt{बलवत्तराः}}{\lemma{\skt{बलवत्तराः}}\Afootnote{A K N\textsuperscript{pc} H; \skt{बलवत्तरा}\ N\textsuperscript{ac}}}\skt{~॥}\edtext{}{\lemma{\skt{गतीः~। ते}}\Afootnote{A K N\textsuperscript{pc} H; \skt{गती~। ते}\ N\textsuperscript{ac}}}\pend

\pstart\versepar \ledsidenote{\passref{1.21.38}}\edtext{\skt{सर्वै\-र्भा\-वै\-स्तृ\-भि\-र्वापि}}{\lemma{\skt{तृभिर्}}\Afootnote{K; \skt{त्रिभिर्}\ A N H}}\skt{ द्वाभ्यामेकेन वा पुनः~। }\pend
\pstart\versepar \skt{संसर्गे }\edtext{\skt{कुपितः}}{\lemma{\skt{कुपितः}}\Afootnote{A K\textsuperscript{pc} H; \skt{कुपित ं}\ K\textsuperscript{ac}; \skt{कुपित}\ N}}\skt{ }\edtext{\skt{क्रुद्धं}}{\lemma{\skt{क्रुद्धं}}\Afootnote{A K N H\textsuperscript{pc}; \skt{क्रुद्धःं}\ H\textsuperscript{ac}}}\skt{ }\edtext{\skt{दोषं}}{\lemma{\skt{दोषं}}\Afootnote{A K H; \skt{दोषां}\ N}}\skt{ दोषो }\edtext{\skt{नुधावति}}{\lemma{\skt{नुधावति}}\Afootnote{K N; \skt{ऽनुधावति}\ A H}}\skt{~॥}\pend

\pstart\versepar \ledsidenote{\passref{1.21.39}}\skt{संसर्गे यो }\edtext{\skt{गरीयां}}{\lemma{\skt{गरीयां}}\Afootnote{K N; \skt{गरीया}\ A H}}\skt{ }\edtext{\skt{स्यादुपक्रम्यः स वै}}{\lemma{\skt{स्यादुपक्रम्यः स वै}}\Afootnote{K\textsuperscript{pc}; \skt{न्स्यादुपक्रम्यः स वै}\ A; \skt{स्यादुपक्रम्य स वै}\ K\textsuperscript{ac} N; \skt{न्स्यादुपक्रम्य स वै}\ H}}\skt{ भवेत्~। }\pend
\pstart\versepar \skt{शेष\-दो\-षा\-वि\-रो\-धेन सन्निपाते तथैव }\edtext{\skt{च}}{\lemma{\skt{च}}\Afootnote{A K\textsuperscript{pc} N H; \skt{चा}\ K\textsuperscript{ac}}}\skt{~॥}\edtext{}{\lemma{\skt{भवेत्~। शेष\appbrk{}दो\appbrk{}षा\appbrk{}वि\appbrk{}रो\appbrk{}धेन}}\Afootnote{A K H; \skt{भवेत्~। तोड\appbrk{}दो\appbrk{}षा\appbrk{}वि\appbrk{}रो\appbrk{}धेन}\ N}}\pend

\pstart\versepar \ledsidenote{\passref{1.21.40}}\edtext{\skt{व्रणे तु}}{\lemma{\skt{व्रणे तु}}\Afootnote{K N H; \skt{वृणोति}\ A}}\skt{ }\edtext{\skt{यस्माद्रूढे}}{\lemma{\skt{यस्माद्रूढे}}\Afootnote{A K H; \skt{यस्माद्रूढो}\ N}}\skt{ }\edtext{\skt{पि}}{\lemma{\skt{पि}}\Afootnote{K; \skt{ऽपि}\ A H N}}\skt{ }\edtext{\skt{व्रणवस्तु न}}{\lemma{\skt{व्रणवस्तु न}}\Afootnote{K N H; \skt{व्रणवस्तु न}\ A}}\skt{ नश्यति~। }\pend
\pstart\versepar \skt{}\edtext{\edtext{\edtext{\edtext{\skt{आदे\-ह\-धा\-र\-णा\-ज्ज\-न्तो\-र्व्र\-ण\-स्त\-स्मा\-न्नि\-रु\-च्यत}}{\lemma{\skt{आदे\appbrk{}ह\appbrk{}धा\appbrk{}र\appbrk{}णा\appbrk{}ज्ज\appbrk{}न्तोर्}}\Afootnote{K N H; \skt{आदे\appbrk{}ह\appbrk{}धा\appbrk{}र\appbrk{}णा\appbrk{}त्तस्मा}\ A}}}{\lemma{\skt{व्रणस्}}\Afootnote{K H; \skt{द्व्रण}\ A; \skt{व्रण}\ N}}}{\lemma{\skt{तस्मान्}}\Afootnote{K N H; \emph{om.}\ A}}}{\lemma{\skt{निरुच्यत}}\Afootnote{K N H; \emph{om.}\ A}}\skt{ }\edtext{\skt{इति}}{\lemma{\skt{इति}}\Afootnote{K N H; \skt{इत्युच्यते बुधैः}\ A}}\pend

\pstart \ledsidenote{\passref{1.21.trailer}}\edtext{\skt{॥ २१~॥~॥}}{\lemma{\skt{॥ २१~॥~॥}}\Afootnote{\skt{इति सुश्रु\appbrk{}त\appbrk{}सं\appbrk{}हि\appbrk{}तायां सूत्रस्थाने व्रणप्रश्नाध्यायो नामै\appbrk{}क\appbrk{}विं\appbrk{}शो\appbrk{}ध्यायः}\ A; \skt{२१~॥ {\symfont ❈}}\ K; \skt{२१~॥}\ N H}}\pend

\endnumbering
\beginnumbering
\eledsection*{Adhyāya 22}
\pstart \ledsidenote{\passref{1.22.1}}\skt{अथातो }\edtext{\edtext{\skt{व्रणा\-स्रा\-व\-वि\-ज्ञा\-नी\-य\-म्व्या\-ख्या\-स्यामः}}{\lemma{\skt{व्रणा\appbrk{}स्रा\appbrk{}व\appbrk{}वि\appbrk{}ज्ञा\appbrk{}नी\appbrk{}यम्}}\Afootnote{N; \skt{व्रणा\appbrk{}स्रा\appbrk{}व\appbrk{}वि\appbrk{}ज्ञा\appbrk{}नी\appbrk{}य\appbrk{}म\appbrk{}ध्यायं}\ A; \skt{व्रणा\appbrk{}स्रा\appbrk{}व\appbrk{}वि\appbrk{}ज्ञा\appbrk{}नी\appbrk{}यम्व्या…}\ K; \skt{व्रणा\appbrk{}श्रा\appbrk{}व\appbrk{}वि\appbrk{}ज्ञा\appbrk{}नी\appbrk{}यम्}\ H}}}{\lemma{\skt{व्याख्यास्यामः}}\Afootnote{H; \skt{व्याख्यास्यमः}\ A N; \emph{om.}\ K}}\skt{~॥}\pend

\pstart \ledsidenote{\passref{1.22.3}}\edtext{\edtext{\skt{त्वङ्मा\-न्स\-सि\-रा\-स्ना\-य्व\-स्थि\-स\-न्धि\-को\-ष्ठ\-म\-र्मा\-ण्यष्टौ}}{\lemma{\skt{त्वङ्मा\appbrk{}न्स\appbrk{}सि\appbrk{}रा\appbrk{}स्ना\appbrk{}य्व\appbrk{}स्थि\appbrk{}स\appbrk{}न्धि\appbrk{}को\appbrk{}ष्ठ\appbrk{}म\appbrk{}र्माण्य्}}\Afootnote{K H; \skt{त्वङ्मां\appbrk{}स\appbrk{}सि\appbrk{}रा\appbrk{}स्ना\appbrk{}य्व\appbrk{}स्थि\appbrk{}स\appbrk{}न्धि\appbrk{}को\appbrk{}ष्ठ\appbrk{}म\appbrk{}र्माणी}\ A; \skt{त्वङ्मा\appbrk{}न्स\appbrk{}शि\appbrk{}रा\appbrk{}स्ना\appbrk{}य्व\appbrk{}स्थि\appbrk{}स\appbrk{}न्धि\appbrk{}कौ\appbrk{}ष्ठ\appbrk{}म\appbrk{}र्माण्य्}\ N}}}{\lemma{\skt{अष्टौ}}\Afootnote{K N H; \skt{त्यष्टौ}\ A}}\skt{ व्रणवस्तूनि }\edtext{\skt{भवन्ति}}{\lemma{\skt{भवन्ति}}\Afootnote{K N H; \emph{om.}\ A}}\skt{~। अत्र सर्व\-व्र\-ण\-स\-न्नि\-वेशः~॥}\pend

\pstart \ledsidenote{\passref{1.22.4}}\edtext{\skt{तत्रा\-द्यै\-क\-व\-स्तु\-स\-न्नि\-वेशी}}{\lemma{\skt{तत्रा\appbrk{}द्यै\appbrk{}क\appbrk{}व\appbrk{}स्तु\appbrk{}स\appbrk{}न्नि\appbrk{}वेशी}}\Afootnote{K N H; \skt{तत्र, आद्यै\appbrk{}क\appbrk{}व\appbrk{}स्तु\appbrk{}स\appbrk{}न्नि\appbrk{}वेशी}\ A}}\skt{ त्वग्भेदी व्रणः }\edtext{\skt{सूपचरो}}{\lemma{\skt{सूपचरो}}\Afootnote{A N H; \skt{सूपचरः}\ K\textsuperscript{ac} K\textsuperscript{pc}}}\skt{ }\edtext{\skt{भवति}}{\lemma{\skt{भवति}}\Afootnote{K\textsuperscript{pc} H; \emph{om.}\ A K\textsuperscript{ac} N}}\skt{~। }\edtext{\edtext{\skt{शेषास्त्वास्रावा}}{\lemma{\skt{शेषास्}}\Afootnote{A K H; \skt{सेषा}\ N}}}{\lemma{\skt{त्वास्रावा}}\Afootnote{K H; \skt{स्व}\ A; \skt{सास्रार्व}\ N}}\skt{ }\edtext{\skt{विज्ञेयाः~। य}}{\lemma{\skt{विज्ञेयाः~। य}}\Afootnote{K N H; \emph{om.}\ A}}\skt{ }\edtext{\edtext{\edtext{\skt{एव\-मु\-त्ति\-ष्ठ\-न्त्य\-व\-दी\-व्यन्ते}}{\lemma{\skt{एवम्}}\Afootnote{K N H; \emph{om.}\ A}}}{\lemma{\skt{उत्तिष्ठन्त्य्}}\Afootnote{K H; \skt{य}\ A; \skt{उत्तिष्ठ}\ N}}\skt{ च}}{\lemma{\skt{अवदीव्यन्ते च}}\Afootnote{K H; \skt{मवदीर्यमाणा दुरुपचाराः}\ A; \skt{वदीर्यन्ते}\ N}}\skt{~।}\pend

\pstart \ledsidenote{\passref{1.22.5}}\edtext{\skt{आय\-त\-श्च\-तु\-रस्रो}}{\lemma{\skt{आयतश्}}\Afootnote{K H; \skt{तत्रायतश्}\ A; \skt{च आयतश्}\ N}}\skt{ }\edtext{\edtext{\skt{वृत्तस्त्रिपुटक}}{\lemma{\skt{वृत्तस्}}\Afootnote{K H; \skt{वृत्त}\ A N}}}{\lemma{\skt{त्रिपुटक}}\Afootnote{K N H; \skt{स्त्रिपुटक}\ A}}\skt{ इति व्रणा\-कृ\-ति\-स\-मासः~। }\edtext{\skt{विशेषतस्तु}}{\lemma{\skt{विशेषतस्तु}}\Afootnote{K H; \skt{शेषास्तु}\ A; \skt{विशेषतस्तु}\ N}}\skt{ }\edtext{\skt{विकृताकृ𑑛तयो दुरुपक्रमा}}{\lemma{\skt{विकृताकृ𑑛तयो}}\Afootnote{K; \skt{विकृताकृतयो}\ A N H}}\skt{ भवन्ति~॥}\pend

\pstart \ledsidenote{\passref{1.22.6}}\skt{सर्व एव व्रणाः क्षिप्रं }\edtext{\edtext{\skt{संरो\-ह\-न्त्या\-त्म\-वतां}}{\lemma{\skt{संरोहन्त्य्}}\Afootnote{K\textsuperscript{pc} H; \skt{संरोहन्}\ A; \skt{संरोहयन्त्य्}\ K\textsuperscript{ac}; \skt{संरोह}\ N}}}{\lemma{\skt{आत्मवतां}}\Afootnote{K H; \skt{त्यात्मवतां}\ A; \skt{त्यात्मवता}\ N}}\skt{ }\edtext{\skt{सुभि\-ष\-ग्भि\-श्चो\-प\-क्रान्ताः}}{\lemma{\skt{सुभिषग्भिश्}}\Afootnote{A K H; \skt{सुभिषस्भिश्}\ N}}\skt{~। }\edtext{\skt{अनात्मवतां}}{\lemma{\skt{अनात्मवतां}}\Afootnote{K H; \skt{अनात्मवता}\ A N}}\skt{ }\edtext{\skt{मज्ञैश्चोपक्रान्ताः}}{\lemma{\skt{मज्ञैश्}}\Afootnote{A K N; \skt{अज्ञैश्}\ H}}\skt{ प्रदुष्यन्ति }\edtext{\skt{प्रवृ\-द्व\-त्वा\-द्दो\-षाणां}}{\lemma{\skt{प्रवृ\appbrk{}द्व\appbrk{}त्वा\appbrk{}द्दो\appbrk{}षाणां}}\Afootnote{K; \skt{प्रवृ\appbrk{}द्ध\appbrk{}त्वा\appbrk{}द्दो\appbrk{}षा\appbrk{}णाम्}\ A N H}}\skt{~।}\pend

\pstart \ledsidenote{\passref{1.22.7}}\edtext{\skt{तत्रातिसम्वृतो}}{\lemma{\skt{तत्रातिसम्वृतो}}\Afootnote{A K H; \skt{तत्रानिसम्}\ N}}\skt{ }\edtext{\skt{विवृतः}}{\lemma{\skt{विवृतः}}\Afootnote{K H; \skt{ऽतिविवृतो}\ A; \skt{वृतो}\ N}}\skt{ }\edtext{\edtext{\edtext{\skt{कठि\-नो\-ति\-मा\-त्र\-म\-ति\-मृ\-दु\-रु\-त्स\-न्नो\-व\-सन्नः}}{\lemma{\skt{कठि\appbrk{}नो\appbrk{}ति\appbrk{}मा\appbrk{}त्रम्}}\Afootnote{K; \skt{ऽतिकठिनो}\ A; \emph{om.}\ N; \skt{कठि\appbrk{}नोऽ\appbrk{}ति\appbrk{}मा\appbrk{}त्रम्}\ H}}}{\lemma{\skt{अतिमृदुर्}}\Afootnote{K H; \skt{ऽतिमृदु}\ A; \emph{om.}\ N}}}{\lemma{\skt{उत्सन्नोवसन्नः}}\Afootnote{H; \skt{रुत्सन्नोऽवसन्नो}\ A; \skt{उत्सन्}\ K; \emph{om.}\ N}}\skt{ }\edtext{\skt{शीतोत्युष्णः}}{\lemma{\skt{शीतोत्युष्णः}}\Afootnote{\skt{ऽतिशीतोऽत्युष्णः}\ A; \emph{om.}\ K N; \skt{शीतोऽत्युष्णः}\ H}}\skt{ }\edtext{\edtext{\edtext{\skt{कृष्ण\-र\-क्त\-पी\-त\-शु\-क्ला\-दि\-व\-र्ण्णै\-ष्व\-प्य\-र्थ\-वर्ण्णः}}{\lemma{\skt{कृष्ण\appbrk{}र\appbrk{}क्त\appbrk{}पी\appbrk{}त\appbrk{}शु\appbrk{}क्ला\appbrk{}दि\appbrk{}व\appbrk{}र्ण्णैष्व्}}\Afootnote{\skt{नग\appbrk{}न्धा\appbrk{}त्य\appbrk{}र्थ\appbrk{}दा\appbrk{}ह\appbrk{}पा\appbrk{}क\appbrk{}रा\appbrk{}ग\appbrk{}वे\appbrk{}द\appbrk{}ना\appbrk{}वा\appbrk{}निति पित्तेन, शेषाः कफेन; उन्मार्गी मुखा\appbrk{}त्मु\appbrk{}खा\appbrk{}न्त\appbrk{}र\appbrk{}वान्, उत्सङ्गः कोटरः’ इति चक्रः; कृष्ण\appbrk{}र\appbrk{}क्त\appbrk{}पी\appbrk{}त\appbrk{}शु\appbrk{}क्ला\appbrk{}दीनां वर्णानाम्}\ A; \skt{नोवसन्नः शीतोत्युष्णः कृष्ण\appbrk{}र\appbrk{}क्त\appbrk{}पी\appbrk{}त\appbrk{}शु\appbrk{}क्ला\appbrk{}दि\appbrk{}व\appbrk{}र्ण्णैष्व्}\ K; \skt{विवृतः कठि\appbrk{}नो\appbrk{}ति\appbrk{}मा\appbrk{}त्र\appbrk{}म\appbrk{}ति\appbrk{}मृ\appbrk{}दु\appbrk{}रु\appbrk{}त्स\appbrk{}न्नो\appbrk{}व\appbrk{}सन्नः शीतोत्युष्णः कृष्ण\appbrk{}र\appbrk{}क्त\appbrk{}पी\appbrk{}त\appbrk{}शु\appbrk{}क्ला\appbrk{}दि\appbrk{}व\appbrk{}र्ण्णेषु}\ N; \skt{कृष्ण\appbrk{}र\appbrk{}क्त\appbrk{}पी\appbrk{}त\appbrk{}शु\appbrk{}क्ला\appbrk{}दि\appbrk{}व\appbrk{}र्ण्णेष्व्}\ H}}}{\lemma{\skt{अप्य्}}\Afootnote{K H; \skt{अन्य}\ A; \emph{om.}\ N}}}{\lemma{\skt{अर्थवर्ण्णः}}\Afootnote{K H; \skt{तमवर्णो}\ A; \emph{om.}\ N}}\skt{ }\edtext{\skt{पूति\-मां\-स\-सि\-रा\-स्ना\-यु\-प्र\-ति\-पूर्ण्णः}}{\lemma{\skt{पूति\appbrk{}मां\appbrk{}स\appbrk{}सि\appbrk{}रा\appbrk{}स्ना\appbrk{}यु\appbrk{}प्र\appbrk{}ति\appbrk{}पूर्ण्णः}}\Afootnote{H; \skt{भैरवः पूति\appbrk{}पू\appbrk{}य\appbrk{}मां\appbrk{}स\appbrk{}सि\appbrk{}रा\appbrk{}स्ना\appbrk{}यु\appbrk{}प्र\appbrk{}भृ\appbrk{}तिभिः पूर्णः}\ A; \skt{पूतिमांस𑑛 सिरा\appbrk{}स्ना\appbrk{}यु\appbrk{}प्र\appbrk{}ति\appbrk{}पूर्ण्णः}\ K; \skt{त्यर्थवर्ण्णः पूति\appbrk{}मा\appbrk{}न्स\appbrk{}शि\appbrk{}रा\appbrk{}श्ना\appbrk{}यु\appbrk{}प्र\appbrk{}ति\appbrk{}पूर्ण्णः}\ N}}\skt{ }\edtext{\skt{पूतिपूयाश्रावी}}{\lemma{\skt{पूतिपूयाश्रावी}}\Afootnote{K H; \skt{पूतिपूयास्रा}\ A; \skt{प्रतिपूयाश्रावी}\ N}}\skt{ }\edtext{\edtext{\edtext{\edtext{\skt{ऊन्मा\-र्ग्यु\-त्स\-ङ्ग्य\-म\-नो\-ज्ञ\-द\-र्श\-न\-ग\-न्धो\-त्य\-र्थ\-वे\-द\-ना\-वा\-न्दा\-ह\-पा\-क\-रा\-ग\-क\-ण्डू\-शो\-फ\-पि\-ट\-को\-प\-द्रु\-तो\-त्य\-र्थ\-दु\-ष्ट\-शो\-णि\-त\-स्रावी}}{\lemma{\skt{ऊन्मार्ग्य्}}\Afootnote{K H; \skt{व्युन्मार्ग्यु}\ A; \skt{उन्मा}\ N}}}{\lemma{\skt{उत्सङ्ग्य्}}\Afootnote{K H; \skt{त्सङ्}\ A; \skt{ग्युत्स}\ N}}}{\lemma{\skt{अम\appbrk{}नो\appbrk{}ज्ञ\appbrk{}द\appbrk{}र्श\appbrk{}न\appbrk{}ग\appbrk{}न्धो\appbrk{}त्य\appbrk{}र्थ\appbrk{}वे\appbrk{}द\appbrk{}ना\appbrk{}वान्}}\Afootnote{K H; \skt{ग्यम\appbrk{}नो\appbrk{}ज्ञ\appbrk{}द\appbrk{}र्श\appbrk{}न\appbrk{}ग\appbrk{}न्धोऽ\appbrk{}त्यर्थं वेदनावान्}\ A; \skt{ग्यम\appbrk{}नो\appbrk{}ज्ञ\appbrk{}द\appbrk{}र्श\appbrk{}न\appbrk{}गन्धो त्यर्थ\appbrk{}वे\appbrk{}द\appbrk{}ना\appbrk{}वान्}\ N}}}{\lemma{\skt{दाह\appbrk{}पा\appbrk{}क\appbrk{}रा\appbrk{}ग\appbrk{}क\appbrk{}ण्डू\appbrk{}शो\appbrk{}फ\appbrk{}पि\appbrk{}ट\appbrk{}को\appbrk{}प\appbrk{}द्रु\appbrk{}तो\appbrk{}त्य\appbrk{}र्थ\appbrk{}दु\appbrk{}ष्ट\appbrk{}शो\appbrk{}णि\appbrk{}त\appbrk{}स्रावी}}\Afootnote{K; \skt{दाह\appbrk{}पा\appbrk{}क\appbrk{}रा\appbrk{}ग\appbrk{}क\appbrk{}ण्डू\appbrk{}शो\appbrk{}फ\appbrk{}पि\appbrk{}ड\appbrk{}को\appbrk{}प\appbrk{}द्रु\appbrk{}तोऽ\appbrk{}त्यर्थं दुष्ट\appbrk{}शो\appbrk{}णि\appbrk{}ता\appbrk{}स्रावी}\ A; \skt{दाह\appbrk{}पा\appbrk{}क\appbrk{}श\appbrk{}ग\appbrk{}क\appbrk{}ण्डू\appbrk{}शो\appbrk{}फ\appbrk{}पि\appbrk{}ट\appbrk{}को\appbrk{}प\appbrk{}द्रु\appbrk{}त\appbrk{}त्य\appbrk{}र्थ\appbrk{}म्ष्ण\appbrk{}शो\appbrk{}णि\appbrk{}ता\appbrk{}श्रावी}\ N; \skt{दाह\appbrk{}पा\appbrk{}क\appbrk{}रा\appbrk{}ग\appbrk{}क\appbrk{}ण्डू\appbrk{}शो\appbrk{}फ\appbrk{}पि\appbrk{}ट\appbrk{}को\appbrk{}प\appbrk{}द्रु\appbrk{}तोऽ\appbrk{}त्य\appbrk{}र्थ\appbrk{}दु\appbrk{}ष्ट\appbrk{}शो\appbrk{}णि\appbrk{}त\appbrk{}स्रावी}\ H}}\skt{ }\edtext{\skt{दीर्घ\-का\-ला\-नु\-बन्धी}}{\lemma{\skt{दीर्घ\appbrk{}का\appbrk{}ला\appbrk{}नु\appbrk{}बन्धी}}\Afootnote{A K H; \skt{दीर्घ\appbrk{}का\appbrk{}ला\appbrk{}नु\appbrk{}वीन्धी}\ N}}\skt{ चेति }\edtext{\skt{दुष्ट\-व्र\-ण\-लि\-ङ्गानि}}{\lemma{\skt{दुष्ट\appbrk{}व्र\appbrk{}ण\appbrk{}लि\appbrk{}ङ्गानि}}\Afootnote{A K\textsuperscript{pc} N H; \skt{व्रणलिङ्गानि~॥}\ K\textsuperscript{ac}}}\skt{~॥ }\edtext{\skt{तत्र}}{\lemma{\skt{तत्र}}\Afootnote{K H; \skt{तस्य}\ A; \emph{om.}\ N}}\skt{ }\edtext{\edtext{\skt{दोषो\-च्छ्रा\-य\-म\-वेक्ष्य}}{\lemma{\skt{दोषोच्छ्रायम्}}\Afootnote{K N; \skt{दोषोच्छ्रायेण ष}\ A; \skt{दोषोच्छ्रयम्}\ H}}}{\lemma{\skt{अवेक्ष्य}}\Afootnote{K N H; \skt{ट्त्वं विभज्य}\ A}}\skt{ यथास्वं }\edtext{\skt{प्रकुर्वीत}}{\lemma{\skt{प्रकुर्वीत}}\Afootnote{K N H; \skt{प्रतीकारे प्रयतेत}\ A}}\skt{~॥}\pend

\pstart \ledsidenote{\passref{1.22.8a}}\edtext{\skt{अथ}}{\lemma{\skt{अथ}}\Afootnote{K N H; \skt{अत ऊर्ध्वं}\ A}}\skt{ }\edtext{\edtext{\skt{सर्वा\-स्रा\-वा\-म्व\-क्ष्यामः}}{\lemma{\skt{सर्वास्रावाम्}}\Afootnote{K H; \skt{सर्वस्रावा}\ A; \skt{सर्वाश्रावम्}\ N}}}{\lemma{\skt{वक्ष्यामः}}\Afootnote{K N H; \skt{न्वक्ष्यामः}\ A}}\skt{~। }\edtext{\skt{तत्र}}{\lemma{\skt{तत्र}}\Afootnote{A K H; \skt{तत्रा}\ N}}\skt{ घृष्टासु }\edtext{\skt{छिन्नासु वा}}{\lemma{\skt{छिन्नासु वा}}\Afootnote{A K; \skt{च्छिन्नासु वा}\ N; \skt{झिन्नासु वा}\ H}}\skt{ }\edtext{\skt{त्वक्षु}}{\lemma{\skt{त्वक्षु}}\Afootnote{A K H; \skt{त्वक्सु}\ N}}\skt{ }\edtext{\skt{त्वक्स्फुटिते}}{\lemma{\skt{त्वक्स्फुटिते}}\Afootnote{K\textsuperscript{pc} N H; \skt{स्फोटे}\ A; \skt{त्वस्फुटिते}\ K\textsuperscript{ac}}}\skt{ }\edtext{\skt{भिन्नेवदारिते वा}}{\lemma{\skt{भिन्नेवदारिते वा}}\Afootnote{K\textsuperscript{pc}; \skt{भिन्ने विदारिते वा}\ A; \skt{भिन्नेवदारि वा}\ K\textsuperscript{ac}; \skt{भिन्नेव दारितेषु वा}\ N; \skt{भिन्नेऽवदारिते वा}\ H}}\skt{ सलिलप्रकाशो }\edtext{\edtext{\skt{भवत्यास्रावः}}{\lemma{\skt{भवत्य्}}\Afootnote{K N H; \skt{भव}\ A}}}{\lemma{\skt{आस्रावः}}\Afootnote{K N H; \skt{त्यास्रावः}\ A}}\skt{ }\edtext{\skt{किञ्चिद्विस्रः}}{\lemma{\skt{किञ्चिद्विस्रः}}\Afootnote{N H; \skt{किञ्चिद्विस्रः}\ A; \skt{किञ्चिद्विस्रा ः}\ K\textsuperscript{ac}; \skt{किञ्चिद्विस्र ः}\ K\textsuperscript{pc}}}\skt{ }\edtext{\skt{पीतावभासश्च}}{\lemma{\skt{पीतावभासश्च}}\Afootnote{K N H; \skt{पीतावभासश्च}\ A}}\skt{~। }\edtext{\skt{मांसगते तु}}{\lemma{\skt{मांसगते तु}}\Afootnote{K H; \skt{मांसगतः}\ A; \skt{मान्सगते}\ N}}\skt{ सर्पिः }\edtext{\skt{प्रकाशः}}{\lemma{\skt{प्रकाशः}}\Afootnote{A K H; \skt{प्रकाश}\ N}}\skt{ सान्द्रः }\edtext{\skt{स्वेतः}}{\lemma{\skt{स्वेतः}}\Afootnote{K H; \skt{श्वेतः}\ A; \skt{श्वेत}\ N}}\skt{ पिच्छिलश्च~। }\edtext{\skt{सिरागतस्तु}}{\lemma{\skt{सिरागतस्तु}}\Afootnote{K N H; \skt{सिरागतः}\ A}}\skt{ }\edtext{\edtext{\skt{सद्यश्छिन्नासु}}{\lemma{\skt{सद्यश्}}\Afootnote{K H; \skt{सद्य}\ A N}}}{\lemma{\skt{छिन्नासु}}\Afootnote{K\textsuperscript{pc} H; \skt{श्छिन्नासु}\ A; \skt{छिन्ना}\ K\textsuperscript{ac}; \skt{स्थिन्नासु}\ N}}\skt{ सिरासु रक्तातिप्रवृत्तिः पक्वासु च }\edtext{\skt{तोय\-ना\-डी\-भि\-रि\-वा\-ग\-मनं}}{\lemma{\skt{इवागमनं}}\Afootnote{K\textsuperscript{pc} H; \skt{इव तोयागमनं}\ A; \skt{इवागम ं}\ K\textsuperscript{ac}; \skt{इवागमन्न}\ N}}\skt{ पूयस्य }\edtext{\skt{आश्रावश्चात्र}}{\lemma{\skt{आश्रावश्}}\Afootnote{K N H; \skt{आस्रावश्}\ A}}\skt{ }\edtext{\edtext{\skt{तनुर्विच्छिन्नः}}{\lemma{\skt{तनुर्}}\Afootnote{K H; \skt{तनु}\ A; \skt{तत्त्र}\ N}}}{\lemma{\skt{विच्छिन्नः}}\Afootnote{K N H; \skt{र्विच्छिन्नः}\ A}}\skt{ }\edtext{\skt{सफेनो}}{\lemma{\skt{सफेनो}}\Afootnote{K H; \skt{पिच्छिलोऽवलम्बी श्यावो}\ A; \skt{सफे}\ N}}\skt{ }\edtext{\skt{वसाप्रतिमः}}{\lemma{\skt{वसाप्रतिमः}}\Afootnote{K H; \skt{ऽव\appbrk{}श्या\appbrk{}य\appbrk{}प्र\appbrk{}ति\appbrk{}मश्}\ A; \skt{नोवृसाप्रतिमः}\ N}}\skt{ }\edtext{\skt{सरक्तश्च}}{\lemma{\skt{सरक्तश्च}}\Afootnote{K N H; \skt{च}\ A}}\skt{~। }\edtext{\skt{स्नायुगतस्तु}}{\lemma{\skt{स्नायुगतस्तु}}\Afootnote{K H; \skt{स्नायुगतः}\ A; \skt{स्नायुवातसू}\ N}}\skt{ स्निग्धो घनः }\edtext{\skt{सिंघ\-ण\-क\-प्र\-तिमः}}{\lemma{\skt{सिंघ\appbrk{}ण\appbrk{}क\appbrk{}प्र\appbrk{}तिमः}}\Afootnote{K; \skt{सिङ्घा\appbrk{}ण\appbrk{}क\appbrk{}प्र\appbrk{}तिमः}\ A H; \skt{सिङ्घा\appbrk{}न\appbrk{}क\appbrk{}प्र\appbrk{}तिमः}\ N}}\skt{ सरक्तश्च~। }\edtext{\skt{अस्थिगतस्तु}}{\lemma{\skt{अस्थिगतस्तु}}\Afootnote{K\textsuperscript{pc} N H; \skt{अस्थिगतो}\ A; \skt{अस्थितस्तु}\ K\textsuperscript{ac}}}\skt{ }\edtext{\edtext{\skt{अस्थिन्यभिहते}}{\lemma{\skt{अस्थिन्य्}}\Afootnote{K N H; \skt{ऽस्थ}\ A}}}{\lemma{\skt{अभिहते}}\Afootnote{K N H; \skt{न्यभिहते}\ A}}\skt{ }\edtext{\skt{स्फुटिते}}{\lemma{\skt{स्फुटिते}}\Afootnote{A K H; \skt{स्फुढिते}\ N}}\skt{ }\edtext{\skt{भिन्नेऽवदारिते वा}}{\lemma{\skt{ऽवदारिते वा}}\Afootnote{\skt{दोषावदारिते वा}\ A; \skt{वदारिते वा}\ K N H}}\skt{ }\edtext{\edtext{\skt{दोष\-भ\-क्षि\-त\-त्वा\-च्छु\-क्ति\-धौ\-त\-मि\-वा\-भाति}}{\lemma{\skt{दोष\appbrk{}भ\appbrk{}क्षि\appbrk{}त\appbrk{}त्वा\appbrk{}च्छु\appbrk{}क्ति\appbrk{}धौ\appbrk{}तम्}}\Afootnote{K N H; \skt{दोष\appbrk{}भ\appbrk{}क्षि\appbrk{}त\appbrk{}त्वा\appbrk{}दस्थि निःसारं शुक्तिधौत}\ A}}}{\lemma{\skt{इवाभाति}}\Afootnote{K N H; \skt{मिवाभाति}\ A}}\skt{ }\edtext{\skt{अस्थिनिःसारश्च}}{\lemma{\skt{अस्थिनिःसारश्च}}\Afootnote{K N H; \emph{om.}\ A}}\skt{ }\edtext{\skt{भवति}}{\lemma{\skt{भवति}}\Afootnote{K N; \emph{om.}\ A; \skt{भ}\ H}}\skt{~। }\edtext{\edtext{\skt{आस्रावश्चात्र}}{\lemma{\skt{आस्रावश्}}\Afootnote{A K N; \emph{om.}\ H}}}{\lemma{\skt{चात्र}}\Afootnote{A K N; \emph{om.}\ H}}\skt{ }\edtext{\edtext{\skt{तनुर्विच्छिन्नो}}{\lemma{\skt{तनुर्}}\Afootnote{K; \emph{om.}\ A H; \skt{तनि}\ N}}}{\lemma{\skt{विच्छिन्नो}}\Afootnote{K; \emph{om.}\ A H; \skt{विच्छिन्ना}\ N}}\skt{ }\edtext{\skt{मज्जामिश्रः}}{\lemma{\skt{मज्जामिश्रः}}\Afootnote{K; \skt{मज्जमिश्रः}\ A N; \emph{om.}\ H}}\skt{ }\edtext{\skt{सरु\-धि\-र\-स्नि\-ग्धश्च}}{\lemma{\skt{सरु\appbrk{}धि\appbrk{}र\appbrk{}स्नि\appbrk{}ग्धश्च}}\Afootnote{K N; \skt{सरुधिरः स्निग्धश्च}\ A; \emph{om.}\ H}}\skt{ }\edtext{\skt{सन्धिगतस्तु}}{\lemma{\skt{सन्धिगतस्तु}}\Afootnote{K\textsuperscript{pc} N; \skt{सन्धिगतः}\ A; \skt{सन्धितस्तु}\ K\textsuperscript{ac}; \emph{om.}\ H}}\skt{ }\edtext{\skt{पीड्यमानो न}}{\lemma{\skt{पीड्यमानो न}}\Afootnote{A K N; \emph{om.}\ H}}\skt{ }\edtext{\skt{प्रवर्त्तते}}{\lemma{\skt{प्रवर्त्तते}}\Afootnote{A K N; \emph{om.}\ H}}\skt{~। }\edtext{\skt{तथा\-कु\-ञ्च\-न\-प्र\-सा\-र\-णो\-न्ना\-म\-न\-वि\-ना\-म\-नो\-स्का\-स\-न\-प्र\-धा\-वनैः}}{\lemma{\skt{तथा\appbrk{}कु\appbrk{}ञ्च\appbrk{}न\appbrk{}प्र\appbrk{}सा\appbrk{}र\appbrk{}णो\appbrk{}न्ना\appbrk{}म\appbrk{}न\appbrk{}वि\appbrk{}ना\appbrk{}म\appbrk{}नो\appbrk{}स्का\appbrk{}स\appbrk{}न\appbrk{}प्र\appbrk{}धा\appbrk{}वनैः}}\Afootnote{K; \skt{आकु\appbrk{}ञ्च\appbrk{}न\appbrk{}प्र\appbrk{}सा\appbrk{}र\appbrk{}णो\appbrk{}न्न\appbrk{}म\appbrk{}न\appbrk{}वि\appbrk{}न\appbrk{}म\appbrk{}न\appbrk{}प्र\appbrk{}धा\appbrk{}व\appbrk{}नो\appbrk{}त्का\appbrk{}स\appbrk{}न\appbrk{}प्र\appbrk{}वा\appbrk{}ह\appbrk{}णैश्च}\ A; \skt{तथा कुञ्च\appbrk{}न\appbrk{}प्र\appbrk{}सा\appbrk{}रणो नाम न विन\appbrk{}म\appbrk{}नो\appbrk{}च्छा\appbrk{}स\appbrk{}न\appbrk{}प्र\appbrk{}धा\appbrk{}वनैः च}\ N; \emph{om.}\ H}}\skt{ }\edtext{\skt{स्रवति}}{\lemma{\skt{स्रवति}}\Afootnote{A K N; \skt{वति}\ H}}\skt{~। }\edtext{\skt{आस्रावश्चात्र}}{\lemma{\skt{आस्रावश्}}\Afootnote{A K N; \skt{आश्रावश्}\ H}}\skt{ }\edtext{\edtext{\skt{तनुर्विच्छिन्नः}}{\lemma{\skt{तनुर्}}\Afootnote{K H; \emph{om.}\ A; \skt{तनु}\ N}}}{\lemma{\skt{विच्छिन्नः}}\Afootnote{K N H; \emph{om.}\ A}}\skt{ }\edtext{\skt{पिच्छिलोवलम्बी}}{\lemma{\skt{पिच्छिलोवलम्बी}}\Afootnote{K; \skt{पिच्छिलोऽवलम्बी}\ A H; \skt{पिच्छिलैचलस्वी}\ N}}\skt{ }\edtext{\skt{सरु\-धि\-रो\-न्म\-थि\-तश्च}}{\lemma{\skt{सरु\appbrk{}धि\appbrk{}रो\appbrk{}न्म\appbrk{}थि\appbrk{}तश्च}}\Afootnote{A K N; \skt{सरु\appbrk{}धि\appbrk{}रो\appbrk{}त्म\appbrk{}थि\appbrk{}तश्च}\ H}}\skt{ }\edtext{\skt{भवति}}{\lemma{\skt{भवति}}\Afootnote{K N H; \emph{om.}\ A}}\skt{~। }\edtext{\skt{कोष्ठगतस्तु}}{\lemma{\skt{तु}}\Afootnote{K N H; \skt{ऽसृङ्}\ A}}\skt{ }\edtext{\skt{मूत्र\-पु\-री\-ष\-पू\-य\-रु\-धि\-रो\-द\-कानि}}{\lemma{\skt{मूत्र\appbrk{}पु\appbrk{}री\appbrk{}ष\appbrk{}पू\appbrk{}य\appbrk{}रु\appbrk{}धि\appbrk{}रो\appbrk{}द\appbrk{}कानि}}\Afootnote{K H; \skt{मूत्र\appbrk{}पु\appbrk{}री\appbrk{}ष\appbrk{}पू\appbrk{}यो\appbrk{}द\appbrk{}कानि}\ A; \skt{सूत्र\appbrk{}पु\appbrk{}री\appbrk{}ष\appbrk{}पू\appbrk{}य\appbrk{}रु\appbrk{}धि\appbrk{}रो\appbrk{}द\appbrk{}कानि}\ N}}\skt{ स्रवति~। }\edtext{\skt{मर्मगतस्तु}}{\lemma{\skt{मर्मगतस्तु}}\Afootnote{K H; \skt{मर्मगतस्}\ A; \skt{मम्मगतस्तु}\ N}}\skt{ }\edtext{\skt{नोच्यते}}{\lemma{\skt{नोच्यते}}\Afootnote{K N H; \emph{om.}\ A}}\skt{ }\edtext{\skt{त्वगा\-दि\-ष्वे\-वा\-व\-रु\-द्ध\-त्वात्}}{\lemma{\skt{एवा\appbrk{}व\appbrk{}रु\appbrk{}द्ध\appbrk{}त्वात्}}\Afootnote{K N H; \skt{अव\appbrk{}रु\appbrk{}द्ध\appbrk{}त्वा\appbrk{}न्नो\appbrk{}च्यते}\ A}}\skt{~॥}\pend

\pstart \ledsidenote{\passref{1.22.11}}\edtext{\skt{अथ}}{\lemma{\skt{अथ}}\Afootnote{K N H; \skt{अत ऊर्ध्वं}\ A}}\skt{ }\edtext{\skt{सर्व\-व्र\-ण\-वे\-द\-ना\-म्व\-क्ष्यामः}}{\lemma{\skt{सर्व\appbrk{}व्र\appbrk{}ण\appbrk{}वे\appbrk{}द\appbrk{}नाम्}}\Afootnote{K N H; \skt{सर्व\appbrk{}व्र\appbrk{}ण\appbrk{}वे\appbrk{}दना}\ A}}\skt{~॥ }\edtext{\skt{तोद\-न\-भे\-द\-न\-च्छे\-द\-न\-ता\-ड\-ना\-व\-म\-न्थ\-ना\-या\-म\-न\-वि\-क्षे\-प\-ण\-चु\-मु\-चु\-मा\-य\-न\-नि\-र्द्द\-र्श\-ना\-व\-भ\-ञ्ज\-न\-स्फो\-ट\-न\-वि\-दा\-र\-णो\-त्पा\-ट\-न\-क\-म्प\-न\-वि\-वि\-ध\-शू\-ल\-वि\-श्ले\-ष\-ण\-वि\-कि\-र\-ण\-पू\-र\-ण\-स्त\-म्भ\-न\-श्व\-प\-ना\-कु\-ञ्च\-ना\-ङ्कु\-शिकाः}}{\lemma{\skt{तोद\appbrk{}न\appbrk{}भे\appbrk{}द\appbrk{}न\appbrk{}च्छे\appbrk{}द\appbrk{}न\appbrk{}ता\appbrk{}ड\appbrk{}ना\appbrk{}व\appbrk{}म\appbrk{}न्थ\appbrk{}ना\appbrk{}या\appbrk{}म\appbrk{}न\appbrk{}वि\appbrk{}क्षे\appbrk{}प\appbrk{}ण\appbrk{}चु\appbrk{}मु\appbrk{}चु\appbrk{}मा\appbrk{}य\appbrk{}न\appbrk{}नि\appbrk{}र्द्द\appbrk{}र्श\appbrk{}ना\appbrk{}व\appbrk{}भ\appbrk{}ञ्ज\appbrk{}न\appbrk{}स्फो\appbrk{}ट\appbrk{}न\appbrk{}वि\appbrk{}दा\appbrk{}र\appbrk{}णो\appbrk{}त्पा\appbrk{}ट\appbrk{}न\appbrk{}क\appbrk{}म्प\appbrk{}न\appbrk{}वि\appbrk{}वि\appbrk{}ध\appbrk{}शू\appbrk{}ल\appbrk{}वि\appbrk{}श्ले\appbrk{}ष\appbrk{}ण\appbrk{}वि\appbrk{}कि\appbrk{}र\appbrk{}ण\appbrk{}पू\appbrk{}र\appbrk{}ण\appbrk{}स्त\appbrk{}म्भ\appbrk{}न\appbrk{}श्व\appbrk{}प\appbrk{}ना\appbrk{}कु\appbrk{}ञ्च\appbrk{}ना\appbrk{}ङ्कु\appbrk{}शिकाः}}\Afootnote{K; \skt{तोद\appbrk{}न\appbrk{}भे\appbrk{}द\appbrk{}न\appbrk{}ता\appbrk{}ड\appbrk{}न\appbrk{}च्छे\appbrk{}द\appbrk{}ना\appbrk{}या\appbrk{}म\appbrk{}न\appbrk{}म\appbrk{}न्थ\appbrk{}न\appbrk{}वि\appbrk{}क्षे\appbrk{}प\appbrk{}ण\appbrk{}चु\appbrk{}मु\appbrk{}चु\appbrk{}मा\appbrk{}य\appbrk{}न\appbrk{}नि\appbrk{}र्द\appbrk{}ह\appbrk{}ना\appbrk{}व\appbrk{}भ\appbrk{}ञ्ज\appbrk{}न\appbrk{}स्फो\appbrk{}टन विदा\appbrk{}र\appbrk{}णो\appbrk{}त्पा\appbrk{}ट\appbrk{}न\appbrk{}क\appbrk{}म्प\appbrk{}न\appbrk{}वि\appbrk{}वि\appbrk{}ध\appbrk{}शू\appbrk{}ल\appbrk{}वि\appbrk{}श्ले\appbrk{}ष\appbrk{}ण\appbrk{}वि\appbrk{}कि\appbrk{}र\appbrk{}ण\appbrk{}स्त\appbrk{}म्भ\appbrk{}न\appbrk{}पू\appbrk{}र\appbrk{}ण\appbrk{}स्व\appbrk{}प्ना\appbrk{}कु\appbrk{}ञ्ज\appbrk{}नाङ्कु शिकाः}\ A; \skt{तोद\appbrk{}न\appbrk{}भे\appbrk{}द\appbrk{}न\appbrk{}च्छे\appbrk{}द\appbrk{}न्ता\appbrk{}ड\appbrk{}ना\appbrk{}व\appbrk{}म\appbrk{}च्छ\appbrk{}ना\appbrk{}या\appbrk{}म\appbrk{}न\appbrk{}वि\appbrk{}क्षे\appbrk{}प\appbrk{}ण\appbrk{}चु\appbrk{}म\appbrk{}चु\appbrk{}मा\appbrk{}य\appbrk{}न\appbrk{}नि\appbrk{}र्द\appbrk{}र्श\appbrk{}ना\appbrk{}व\appbrk{}भ\appbrk{}ञ्ज\appbrk{}स्फो\appbrk{}ट\appbrk{}न\appbrk{}वि\appbrk{}दा\appbrk{}र\appbrk{}णो\appbrk{}त्पा\appbrk{}ट\appbrk{}न\appbrk{}क\appbrk{}न्प\appbrk{}न\appbrk{}वि\appbrk{}वि\appbrk{}ध\appbrk{}शू\appbrk{}ल\appbrk{}वि\appbrk{}श्ले\appbrk{}ष\appbrk{}ण\appbrk{}वि\appbrk{}कि\appbrk{}र\appbrk{}ण\appbrk{}पू\appbrk{}र\appbrk{}ण\appbrk{}स्त\appbrk{}म्भ\appbrk{}न\appbrk{}स्व\appbrk{}प\appbrk{}ना\appbrk{}कु\appbrk{}ञ्च\appbrk{}ना\appbrk{}त्कु\appbrk{}सि\appbrk{}कास्}\ N; \skt{तोद\appbrk{}न\appbrk{}भे\appbrk{}द\appbrk{}न\appbrk{}च्छे\appbrk{}द\appbrk{}न\appbrk{}ता\appbrk{}ड\appbrk{}ना\appbrk{}व\appbrk{}म\appbrk{}न्थ\appbrk{}ना\appbrk{}या\appbrk{}म\appbrk{}न\appbrk{}वि\appbrk{}क्षे\appbrk{}प\appbrk{}ण\appbrk{}चु\appbrk{}मु\appbrk{}चु\appbrk{}मा\appbrk{}य\appbrk{}न\appbrk{}नि\appbrk{}र्द्द\appbrk{}र्श\appbrk{}ना\appbrk{}व\appbrk{}भ\appbrk{}ञ्ज\appbrk{}न\appbrk{}स्फो\appbrk{}ट\appbrk{}न\appbrk{}वि\appbrk{}दा\appbrk{}र\appbrk{}णो\appbrk{}त्पा\appbrk{}ट\appbrk{}न\appbrk{}क\appbrk{}म्प\appbrk{}न\appbrk{}वि\appbrk{}वि\appbrk{}ध\appbrk{}शू\appbrk{}ल\appbrk{}वि\appbrk{}श्ले\appbrk{}ष\appbrk{}ण\appbrk{}वि\appbrk{}कि\appbrk{}र\appbrk{}ण\appbrk{}पू\appbrk{}र\appbrk{}ण\appbrk{}स्त\appbrk{}म्भ\appbrk{}न\appbrk{}स्व\appbrk{}प\appbrk{}ना\appbrk{}कु\appbrk{}ञ्च\appbrk{}ना\appbrk{}ङ्कु\appbrk{}सिकाः}\ H}}\skt{ सम्भवन्ति~। }\edtext{\skt{विविधा वा}}{\lemma{\skt{विविधा वा}}\Afootnote{K N H; \skt{अनि\appbrk{}मि\appbrk{}त्त\appbrk{}वि\appbrk{}वि\appbrk{}धवे}\ A}}\skt{ }\edtext{\skt{यत्र}}{\lemma{\skt{यत्र}}\Afootnote{K N H; \skt{दनाप्रादुर्भावो वा}\ A}}\skt{ }\edtext{\edtext{\skt{मुहु\-र्मु\-हु\-र्वे\-दना}}{\lemma{\skt{मुहुर्मुहुर्}}\Afootnote{K H; \skt{मुहुर्मुहु}\ A; \skt{मुहुमुहुर्}\ N}}}{\lemma{\skt{वेदना}}\Afootnote{K H; \skt{र्यत्रागच्छन्ति वेदना}\ A; \skt{यत्र वेदना}\ N}}\skt{ }\edtext{\skt{आगच्छन्ति}}{\lemma{\skt{आगच्छन्ति}}\Afootnote{K N H; \skt{वि}\ A}}\skt{ }\edtext{\edtext{\skt{तम्वातिकमिति}}{\lemma{\skt{तम्}}\Afootnote{K H; \skt{शेषास्तं}\ A; \skt{तस्वा}\ N}}}{\lemma{\skt{वातिकम्}}\Afootnote{A K H; \skt{तिकम्}\ N}}\skt{ }\edtext{\skt{विद्यात्~। ऊषा\-चो\-ष\-प\-रि\-दा\-ह\-धू\-मा\-य\-नानि}}{\lemma{\skt{विद्यात्~। ऊषा\appbrk{}चो\appbrk{}ष\appbrk{}प\appbrk{}रि\appbrk{}दा\appbrk{}ह\appbrk{}धू\appbrk{}मा\appbrk{}य\appbrk{}नानि}}\Afootnote{K; \skt{विद्यात्; ओष\appbrk{}चो\appbrk{}ष\appbrk{}प\appbrk{}रि\appbrk{}दा\appbrk{}ह\appbrk{}धू\appbrk{}मा\appbrk{}य\appbrk{}नानि}\ A; \skt{किन्द्यात्~॥ ऊषा\appbrk{}चो\appbrk{}ष\appbrk{}प\appbrk{}रि\appbrk{}दा\appbrk{}ह\appbrk{}धू\appbrk{}मा\appbrk{}य\appbrk{}नानि}\ N; \skt{विन्द्यात्~। ऊषा\appbrk{}चो\appbrk{}ष\appbrk{}प\appbrk{}रि\appbrk{}दा\appbrk{}ह\appbrk{}धू\appbrk{}मा\appbrk{}य\appbrk{}नानि}\ H}}\skt{ }\edtext{\skt{यत्रा\-ङ्गा\-रा\-व\-की\-र्ण्ण\-मिव}}{\lemma{\skt{यत्रा\appbrk{}ङ्गा\appbrk{}रा\appbrk{}व\appbrk{}की\appbrk{}र्ण्णम्}}\Afootnote{K H; \skt{यत्र गात्र\appbrk{}म\appbrk{}ङ्गा\appbrk{}रा\appbrk{}व\appbrk{}की\appbrk{}र्णम्}\ A; \skt{यन्त्रा\appbrk{}गा\appbrk{}रा\appbrk{}व\appbrk{}की\appbrk{}र्ण्णम्}\ N}}\skt{ }\edtext{\skt{वेदना}}{\lemma{\skt{वेदना}}\Afootnote{K N H; \emph{om.}\ A}}\skt{ }\edtext{\skt{सरुजं}}{\lemma{\skt{सरुजं}}\Afootnote{K N H; \emph{om.}\ A}}\skt{ }\edtext{\skt{तीक्ष्ण\-स\-म्पा\-ति\-प\-च्यते}}{\lemma{\skt{तीक्ष्ण\appbrk{}स\appbrk{}म्पा\appbrk{}ति\appbrk{}प\appbrk{}च्यते}}\Afootnote{K H; \skt{पच्यते}\ A; \skt{तीक्ष्ण\appbrk{}स\appbrk{}म्या\appbrk{}ति\appbrk{}प\appbrk{}च्युते}\ N}}\skt{ यत्र }\edtext{\edtext{\skt{चोष्मा\-भि\-र्वृ\-द्धि\-र्भ\-वति}}{\lemma{\skt{चोष्माभिर्}}\Afootnote{K H; \skt{चोष्माभि}\ A N}}}{\lemma{\skt{भवति}}\Afootnote{K N H; \emph{om.}\ A}}\skt{ क्षते क्षारा\-व\-सि\-क्त\-वच्च }\edtext{\skt{यत्र}}{\lemma{\skt{यत्र}}\Afootnote{K N H; \emph{om.}\ A}}\skt{ वेदनाविशेषाः }\edtext{\skt{घोतात }\edtext{\skt{म्पैत्तिकमिति}}{\lemma{\skt{पैत्तिकम्}}\Afootnote{A K\textsuperscript{pc} H; \skt{पैत्तिकी म्}\ K\textsuperscript{ac}; \skt{पैत्तिकिम्}\ N}}}{\lemma{\skt{घोतात म्}}\Afootnote{K\textsuperscript{pc} H; \skt{स्तं}\ A; \skt{घोताता म्}\ K\textsuperscript{ac}; \skt{घातात}\ N}}\skt{ }\edtext{\skt{विंद्यात्~। पित्त\-व\-द्र\-क्त\-स\-मु\-त्थ\-ञ्जा\-नी\-यात्}}{\lemma{\skt{विंद्यात्~। पित्त\appbrk{}व\appbrk{}द्र\appbrk{}क्त\appbrk{}स\appbrk{}मु\appbrk{}त्थञ्}}\Afootnote{K H; \skt{विद्यात्; पित्त\appbrk{}व\appbrk{}द्र\appbrk{}क्त\appbrk{}स\appbrk{}मुत्थं}\ A; \skt{विन्द्यात्~॥ पित्त\appbrk{}व\appbrk{}द्र\appbrk{}क्त\appbrk{}स\appbrk{}मुच्छं}\ N}}\skt{~। विशेषो }\edtext{\skt{रक्तदो}}{\lemma{\skt{रक्तदो}}\Afootnote{K\textsuperscript{pc} H; \emph{om.}\ A K\textsuperscript{ac}; \skt{स्कुतोदो}\ N}}\skt{ }\edtext{\skt{रक्तस्रावश्च}}{\lemma{\skt{रक्तस्रावश्च}}\Afootnote{K\textsuperscript{pc} N H; \emph{om.}\ A K\textsuperscript{ac}}}\skt{ }\edtext{\skt{भवति}}{\lemma{\skt{भवति}}\Afootnote{K\textsuperscript{pc} N H; \emph{om.}\ A K\textsuperscript{ac}}}\skt{~। }\edtext{\skt{कण्डूर्गुरुत्वं}}{\lemma{\skt{कण्डूर्गुरुत्वं}}\Afootnote{K\textsuperscript{pc} H; \skt{त्कण्डूर्गुरुत्वं}\ A K\textsuperscript{ac}; \skt{कण्डूगुरुत्वं}\ N}}\skt{ }\edtext{\skt{सुप्तता}}{\lemma{\skt{सुप्तता}}\Afootnote{K N H; \skt{सुप्तत्वम्}\ A}}\skt{ }\edtext{\skt{स्वेदो\-ल्प\-वे\-द\-नत्वं}}{\lemma{\skt{स्वेदो\appbrk{}ल्प\appbrk{}वे\appbrk{}द\appbrk{}नत्वं}}\Afootnote{K; \skt{उप\appbrk{}दे\appbrk{}होऽ\appbrk{}ल्प\appbrk{}वे\appbrk{}द\appbrk{}नत्वं}\ A; \skt{स्वेदोल्पवेदनं}\ N; \skt{स्वेदोऽ\appbrk{}ल्प\appbrk{}वे\appbrk{}द\appbrk{}नत्वं}\ H}}\skt{ स्तम्भः शैत्यञ्च यत्र तं }\edtext{\edtext{\skt{श्लैष्मिकमिति}}{\lemma{\skt{श्लैष्मिकम्}}\Afootnote{A K H; \skt{श्लेष्मिक}\ N}}}{\lemma{\skt{इति}}\Afootnote{A K H; \skt{मिति}\ N}}\skt{ }\edtext{\skt{विद्यात्~। यत्र}}{\lemma{\skt{विद्यात्~। यत्र}}\Afootnote{A K N; \skt{विन्द्यात्~। यत्र}\ H}}\skt{ }\edtext{\skt{सर्व\-वे\-द\-ना\-स\-मु\-त्प\-त्तिस्तं}}{\lemma{\skt{सर्व\appbrk{}वे\appbrk{}द\appbrk{}ना\appbrk{}स\appbrk{}मु\appbrk{}त्प\appbrk{}त्तिस्तं}}\Afootnote{K N H; \skt{सर्वासां वेद\appbrk{}ना\appbrk{}ना\appbrk{}मु\appbrk{}त्प\appbrk{}त्तिस्तं}\ A}}\skt{ सान्नि\-पा\-ति\-क\-मिति }\edtext{\skt{विद्यात्}}{\lemma{\skt{विद्यात्}}\Afootnote{A K N; \skt{विंद्यात्}\ H}}\skt{~॥}\edtext{}{\lemma{\skt{जानीयात्~। विशेषो}}\Afootnote{K\textsuperscript{pc} H; \skt{जानीया}\ A K\textsuperscript{ac}; \skt{जानीयात्~॥ विसेषो}\ N}}\pend

\pstart \ledsidenote{\passref{1.22.12}}\skt{अत }\edtext{\skt{ऊर्द्ध्वं}}{\lemma{\skt{ऊर्द्ध्वं}}\Afootnote{A K H; \skt{ऊर्द्ध्व}\ N}}\skt{ }\edtext{\edtext{\skt{सर्व\-व\-र्ण्णा\-म्व\-क्ष्यामः}}{\lemma{\skt{सर्ववर्ण्णाम्}}\Afootnote{K\textsuperscript{pc}; \skt{व्रणवर्णा}\ A; \skt{सवर्ण्णाम्}\ K\textsuperscript{ac}; \skt{सर्ववर्ण्णा}\ N; \skt{सर्व्ववर्ण्णम्}\ H}}}{\lemma{\skt{वक्ष्यामः}}\Afootnote{K H; \skt{न्वक्ष्यामः}\ A N}}\skt{~। }\edtext{\skt{भस्म\-क\-पो\-स्थि\-वर्ण्णः}}{\lemma{\skt{भस्म\appbrk{}क\appbrk{}पो\appbrk{}स्थि\appbrk{}वर्ण्णः}}\Afootnote{K; \skt{भस्म\appbrk{}क\appbrk{}पो\appbrk{}ता\appbrk{}स्थि\appbrk{}वर्णः}\ A N H}}\skt{ }\edtext{\skt{परु\-षो\-रु\-ण\-वर्ण्णः}}{\lemma{\skt{परु\appbrk{}षो\appbrk{}रु\appbrk{}ण\appbrk{}वर्ण्णः}}\Afootnote{K N; \skt{परुषोऽरुणः}\ A; \skt{परु\appbrk{}षोऽ\appbrk{}रु\appbrk{}ण\appbrk{}वर्ण्णः}\ H}}\skt{ कृष्ण इति }\edtext{\skt{मारुतजस्य}}{\lemma{\skt{मारुतजस्य}}\Afootnote{K\textsuperscript{pc} N H; \skt{मारुतजस्यः}\ A; \skt{मारुतस्य}\ K\textsuperscript{ac}}}\skt{~। }\edtext{\skt{नीलः}}{\lemma{\skt{नीलः}}\Afootnote{A K H; \skt{नील}\ N}}\skt{ }\edtext{\skt{श्यावो}}{\lemma{\skt{श्यावो}}\Afootnote{K H; \skt{पीतो}\ A; \skt{स्यावो}\ N}}\skt{ हरितः }\edtext{\skt{पीतः}}{\lemma{\skt{पीतः}}\Afootnote{K N H; \skt{श्यावः}\ A}}\skt{ कृष्णो }\edtext{\skt{रक्तः}}{\lemma{\skt{रक्तः}}\Afootnote{A K H; \skt{रक्त}\ N}}\skt{ }\edtext{\skt{पिङ्गल}}{\lemma{\skt{पिङ्गल}}\Afootnote{K N H; \skt{पिङ्गलः कपिल}\ A}}\skt{ इति }\edtext{\skt{पित्त\-र\-क्त\-स\-मु\-त्थयोः}}{\lemma{\skt{पित्त\appbrk{}र\appbrk{}क्त\appbrk{}स\appbrk{}मु\appbrk{}त्थयोः}}\Afootnote{K N H; \skt{रक्त\appbrk{}पि\appbrk{}त्त\appbrk{}स\appbrk{}मु\appbrk{}त्थयोः}\ A}}\skt{~। }\edtext{\skt{श्वेतस्निग्धः}}{\lemma{\skt{श्वेतस्निग्धः}}\Afootnote{K N; \skt{श्वेतः स्निग्धः}\ A H}}\skt{ पाण्डुरिति श्लेष्मजस्य~। सर्ववर्ण्णोपेतः }\edtext{\skt{सान्निपातिक}}{\lemma{\skt{सान्निपातिक}}\Afootnote{K\textsuperscript{pc}; \skt{सान्निपातिक इति}\ A; \skt{सान्निपातिकस्य}\ K\textsuperscript{ac}; \skt{सान्निपातिकः}\ N H}}\skt{~॥}\pend

\pstart\versepar \ledsidenote{\passref{1.22.13}}\skt{भवति चात्र श्लोकः~॥ }\pend
\pstart\versepar \skt{न केवलं व्रणेषूक्तो वेद\-ना\-व\-र्ण्ण\-स\-ङ्ग्रहः }\pend
\pstart\versepar \skt{सर्व\-शो\-फ\-वि\-का\-रेषु }\edtext{\skt{व्रण\-व\-ल्ल\-क्ष\-ये\-द्भि\-ष\-गिति}}{\lemma{\skt{व्रण\appbrk{}व\appbrk{}ल्ल\appbrk{}क्ष\appbrk{}ये\appbrk{}द्भि\appbrk{}ष\appbrk{}गिति}}\Afootnote{H; \skt{व्रण\appbrk{}व\appbrk{}ल्ल\appbrk{}क्ष\appbrk{}ये\appbrk{}द्भि\appbrk{}षक्}\ A; \skt{व्रण\appbrk{}व\appbrk{}ल्ल\appbrk{}क्ष\appbrk{}ये\appbrk{}द्भि\appbrk{}ष\appbrk{}गिति~॥ २२~॥ {\symfont ❈}}\ K; \skt{व्रणवल्लक्षये भिषगिति~॥}\ N}}\edtext{}{\lemma{\skt{श्लोकः~॥ न}}\Afootnote{K H; \skt{न}\ A N}}\pend

\pstart \ledsidenote{\passref{1.22.trailer}}\edtext{\skt{॥ २२~॥ {\symfont ❈}}}{\lemma{\skt{॥ २२~॥ {\symfont ❈}}}\Afootnote{\skt{इति सुश्रु\appbrk{}त\appbrk{}सं\appbrk{}हि\appbrk{}तायां सूत्रस्थाने व्रणा\appbrk{}स्रा\appbrk{}व\appbrk{}वि\appbrk{}ज्ञा\appbrk{}नीयो नाम द्वाविं\appbrk{}श\appbrk{}ति\appbrk{}त\appbrk{}मोऽ\appbrk{}ध्यायः}\ A}}\skt{~॥}\pend

\endnumbering
\beginnumbering
\eledsection*{Adhyāya 23}
\pstart \ledsidenote{\passref{1.23.1}}\skt{अथातः कृत्याकृत्यविधिं }\edtext{\skt{व्याख्यास्यामः}}{\lemma{\skt{व्याख्यास्यामः}}\Afootnote{H; \skt{अध्यायं व्याख्यास्यामः}\ A N; \emph{om.}\ K}}\skt{~॥}\pend

\pstart \ledsidenote{\passref{1.23.3}}\skt{तत्र }\edtext{\skt{वयस्थानां}}{\lemma{\skt{वयस्थानां}}\Afootnote{K N; \skt{वयःस्थानां}\ A H}}\skt{ दृढानां प्राणवतां सत्व\-व\-ता\-मा\-त्म\-व\-ताञ्च }\edtext{\skt{सुचिकित्स्या}}{\lemma{\skt{सुचिकित्स्या}}\Afootnote{A K H; \skt{सु त्स्या}\ N}}\skt{ व्रणा }\edtext{\skt{भवन्ति}}{\lemma{\skt{भवन्ति}}\Afootnote{K N H; \emph{om.}\ A}}\skt{~॥ }\edtext{\skt{एकैकस्मिन्वा}}{\lemma{\skt{एकैकस्मिन्वा}}\Afootnote{K N H; \skt{एकस्मिन्वा}\ A}}\skt{ पुरुषे }\edtext{\skt{यत्रै\-त\-द्गु\-ण\-प\-ञ्चकं}}{\lemma{\skt{यत्रै\appbrk{}त\appbrk{}द्गु\appbrk{}ण\appbrk{}प\appbrk{}ञ्चकं}}\Afootnote{K N H; \skt{यत्रैतद्गुण}\ A}}\skt{ }\edtext{\skt{भवति}}{\lemma{\skt{भवति}}\Afootnote{K N H; \skt{चतुष्टयं}\ A}}\skt{ तस्य }\edtext{\skt{खलु}}{\lemma{\skt{खलु}}\Afootnote{K N H; \skt{सुख}\ A}}\skt{ साधनीयतमाः तत्र }\edtext{\skt{वयस्थानां}}{\lemma{\skt{वयस्थानां}}\Afootnote{K N; \skt{वयःस्थानां}\ A H}}\skt{ प्रत्य\-ग्र\-धा\-तु\-त्वा\-दाशु }\edtext{\skt{व्रणसंरोहो}}{\lemma{\skt{व्रणसंरोहो}}\Afootnote{K H; \skt{व्रणा रोहन्}\ A; \skt{संरोहो}\ N}}\skt{ }\edtext{\skt{भवति}}{\lemma{\skt{भवति}}\Afootnote{K N H; \skt{ति}\ A}}\skt{~। दृढानां }\edtext{\skt{स्थिर\-ब\-हु\-मां\-स\-त्वा\-च्छ\-स्त्र\-म\-व\-चा\-र्य\-माणं}}{\lemma{\skt{अवचार्यमाणं}}\Afootnote{A N H; \skt{अ…}\ K}}\skt{ }\edtext{\skt{सिरा\-स्ना\-य्वा\-दी\-न्वि\-शे\-षान्न}}{\lemma{\skt{सिरा\appbrk{}स्ना\appbrk{}य्वा\appbrk{}दी\appbrk{}न्वि\appbrk{}शे\appbrk{}षान्न}}\Afootnote{N H; \skt{सिरा\appbrk{}स्ना\appbrk{}य्वा\appbrk{}दि\appbrk{}वि\appbrk{}शे\appbrk{}षान्न}\ A; \emph{om.}\ K}}\skt{ }\edtext{\skt{प्राप्नोति}}{\lemma{\skt{प्राप्नोति}}\Afootnote{A N H; \emph{om.}\ K}}\skt{ }\edtext{\edtext{\skt{प्राण\-व\-ता\-म्वे\-द\-ना\-भि\-घा\-ता\-हा\-र\-य\-न्त्र\-णा\-दि\-भिर्न्न}}{\lemma{\skt{प्राणवताम्}}\Afootnote{A N H; \emph{om.}\ K}}}{\lemma{\skt{वेद\appbrk{}ना\appbrk{}भि\appbrk{}घा\appbrk{}ता\appbrk{}हा\appbrk{}र\appbrk{}य\appbrk{}न्त्र\appbrk{}णा\appbrk{}दि\appbrk{}भिर्न्न}}\Afootnote{A H; \emph{om.}\ K; \skt{वेद\appbrk{}ना\appbrk{}वि\appbrk{}घा\appbrk{}ता\appbrk{}हा\appbrk{}र\appbrk{}य\appbrk{}न्त्र\appbrk{}णा\appbrk{}भिर्न}\ N}}\skt{ }\edtext{\edtext{\skt{ग्लानिर्ब्भवति}}{\lemma{\skt{ग्लानिर्}}\Afootnote{A N H; \emph{om.}\ K}}}{\lemma{\skt{ब्भवति}}\Afootnote{N H; \skt{उत्पद्यते}\ A; \emph{om.}\ K}}\skt{~। }\edtext{\edtext{\edtext{\skt{सत्व\-व\-ता\-न्दा\-रु\-णै\-रपि}}{\lemma{\skt{सत्ववतान्}}\Afootnote{A N H; \emph{om.}\ K}}}{\lemma{\skt{दारुणैर्}}\Afootnote{A N H; \emph{om.}\ K}}}{\lemma{\skt{अपि}}\Afootnote{A N H; \emph{om.}\ K}}\skt{ }\edtext{\skt{क्रियाविशेषैर्न्न}}{\lemma{\skt{क्रियाविशेषैर्न्न}}\Afootnote{A N H; \emph{om.}\ K}}\skt{ }\edtext{\skt{व्यथा}}{\lemma{\skt{व्यथा}}\Afootnote{A N H; \emph{om.}\ K}}\skt{ }\edtext{\skt{भवति}}{\lemma{\skt{भवति}}\Afootnote{N H; \skt{भ}\ A; \emph{om.}\ K}}\skt{~। }\edtext{\edtext{\edtext{\edtext{\edtext{\skt{आत्म\-व\-ता\-म्पु\-न\-र्म्मि\-ता\-हा\-र\-वि\-हा\-रा\-दि\-भि\-रु\-प\-दे\-शै\-र्न्ना\-न्यथा}}{\lemma{\skt{आत्मवताम्}}\Afootnote{N H; \skt{व}\ A; \emph{om.}\ K}}}{\lemma{\skt{पुनर्}}\Afootnote{N H; \emph{om.}\ A K}}}{\lemma{\skt{म्मिता\appbrk{}हा\appbrk{}र\appbrk{}वि\appbrk{}हा\appbrk{}रा\appbrk{}दि\appbrk{}भिर्}}\Afootnote{N H; \emph{om.}\ A K}}}{\lemma{\skt{उपदेशैर्}}\Afootnote{N H; \emph{om.}\ A K}}\skt{ मतिः}}{\lemma{\skt{न्नान्यथा मतिः}}\Afootnote{N H; \emph{om.}\ A K}}\skt{ }\edtext{\skt{प्रवर्त्तते}}{\lemma{\skt{प्रवर्त्तते}}\Afootnote{H; \skt{ति}\ A; \emph{om.}\ K; \skt{प्रतिप्रवर्त्तते}\ N}}\skt{~। }\edtext{\edtext{\skt{सर्व्वञ्चैषां}}{\lemma{\skt{सर्व्वञ्}}\Afootnote{N H; \emph{om.}\ A K}}}{\lemma{\skt{चैषां}}\Afootnote{N H; \emph{om.}\ A K}}\skt{ }\edtext{\edtext{\skt{साधुर्ब्भवति}}{\lemma{\skt{साधुर्}}\Afootnote{N H; \emph{om.}\ A K}}}{\lemma{\skt{ब्भवति}}\Afootnote{N H; \skt{तस्मादे}\ A; \emph{om.}\ K}}\skt{~। }\edtext{\skt{तस्मादेषां}}{\lemma{\skt{तस्मादेषां}}\Afootnote{H; \skt{तेषां}\ A; \emph{om.}\ K; \skt{तस्मा}\ N}}\skt{ }\edtext{\skt{सुख\-सा\-ध\-नी\-य\-तमाः}}{\lemma{\skt{सुख\appbrk{}सा\appbrk{}ध\appbrk{}नी\appbrk{}य\appbrk{}तमाः}}\Afootnote{A H; \emph{om.}\ K; \skt{त्सुख\appbrk{}सा\appbrk{}ध\appbrk{}नी\appbrk{}य\appbrk{}तमाः}\ N}}\skt{~॥}\pend

\pstart \ledsidenote{\passref{1.23.4}}\edtext{\skt{त एव}}{\lemma{\skt{त एव}}\Afootnote{A N H; \skt{{\symfont ‡}}\ K}}\skt{ }\edtext{\skt{विपरीतगुणाः}}{\lemma{\skt{विपरीतगुणाः}}\Afootnote{A N H; \emph{om.}\ K}}\skt{~। }\edtext{\edtext{\skt{वृद्ध\-कृ\-षा\-ल्प\-प्रा\-ण\-भी\-रु\-ष्व\-ना\-त्म\-वत्सु}}{\lemma{\skt{वृद्ध\appbrk{}कृ\appbrk{}षा\appbrk{}ल्प\appbrk{}प्रा\appbrk{}ण\appbrk{}भी\appbrk{}रुष्व्}}\Afootnote{H; \skt{वृद्ध\appbrk{}कृ\appbrk{}शा\appbrk{}ल्प\appbrk{}प्रा\appbrk{}ण\appbrk{}भी\appbrk{}रुषु}\ A; \emph{om.}\ K; \skt{वृद्ध\appbrk{}कृ\appbrk{}शा\appbrk{}ल्प\appbrk{}प्रा\appbrk{}ण\appbrk{}भीरु}\ N}}\skt{ च}}{\lemma{\skt{अनात्मवत्सु च}}\Afootnote{H; \emph{om.}\ A K; \skt{ष्वनात्मवत्सु च}\ N}}\skt{ }\edtext{\skt{द्रष्टव्याः}}{\lemma{\skt{द्रष्टव्याः}}\Afootnote{A N H; \emph{om.}\ K}}\skt{~॥}\pend

\pstart \ledsidenote{\passref{1.23.5}}\edtext{\skt{स्फिक्पा\-यु\-प्र\-ज\-न\-न\-ल\-ला\-ट\-ग\-ण्डौ\-ष्ठ\-फ\-ल\-को\-षो\-द\-र\-ज\-त्रु\-मु\-खा\-भ्य\-न्त\-र\-संस्थाः}}{\lemma{\skt{स्फिक्पा\appbrk{}यु\appbrk{}प्र\appbrk{}ज\appbrk{}न\appbrk{}न\appbrk{}ल\appbrk{}ला\appbrk{}ट\appbrk{}ग\appbrk{}ण्डौ\appbrk{}ष्ठ\appbrk{}फ\appbrk{}ल\appbrk{}को\appbrk{}षो\appbrk{}द\appbrk{}र\appbrk{}ज\appbrk{}त्रु\appbrk{}मु\appbrk{}खा\appbrk{}भ्य\appbrk{}न्त\appbrk{}र\appbrk{}संस्थाः}}\Afootnote{H; \skt{स्फिक्पा\appbrk{}यु\appbrk{}प्र\appbrk{}ज\appbrk{}न\appbrk{}न\appbrk{}ल\appbrk{}ला\appbrk{}ट\appbrk{}ग\appbrk{}ण्डौ\appbrk{}ष्ठ\appbrk{}पृ\appbrk{}ष्ठ\appbrk{}क\appbrk{}र्ण\appbrk{}फ\appbrk{}ल\appbrk{}को\appbrk{}षो\appbrk{}द\appbrk{}र\appbrk{}ज\appbrk{}त्रु\appbrk{}मु\appbrk{}खा\appbrk{}भ्य\appbrk{}न्त\appbrk{}र\appbrk{}संस्थाः}\ A; \skt{{\symfont ‡}}\ K; \skt{स्फिक्पा\appbrk{}यु\appbrk{}प्र\appbrk{}ज\appbrk{}न\appbrk{}न\appbrk{}ल\appbrk{}ला\appbrk{}ट\appbrk{}ग\appbrk{}ण्डो\appbrk{}ष्ठ\appbrk{}फ\appbrk{}ल\appbrk{}को\appbrk{}शो\appbrk{}द\appbrk{}र\appbrk{}ज\appbrk{}त्रु\appbrk{}मु\appbrk{}खा\appbrk{}भ्य\appbrk{}न्त\appbrk{}र\appbrk{}संस्थाः}\ N}}\skt{ }\edtext{\skt{सुखरोपणीया}}{\lemma{\skt{सुखरोपणीया}}\Afootnote{A N H; \emph{om.}\ K}}\skt{ }\edtext{\skt{व्रणाः}}{\lemma{\skt{व्रणाः}}\Afootnote{A N H; \emph{om.}\ K}}\skt{~॥}\pend

\pstart \ledsidenote{\passref{1.23.6}}\edtext{\skt{अक्षि\-द\-न्त\-ना\-सा\-अ\-पा\-ङ्ग\-श्रो\-त्र\-ना\-भि\-ज\-ठ\-र\-से\-व\-नी\-नि\-त\-म्ब\-पा\-र्श्व\-कु\-क्षि\-स्त\-न\-क\-क्ष\-स\-न्धि\-भा\-ग\-गताः}}{\lemma{\skt{अक्षि\appbrk{}द\appbrk{}न्त\appbrk{}ना\appbrk{}सा\appbrk{}अ\appbrk{}पा\appbrk{}ङ्ग\appbrk{}श्रो\appbrk{}त्र\appbrk{}ना\appbrk{}भि\appbrk{}ज\appbrk{}ठ\appbrk{}र\appbrk{}से\appbrk{}व\appbrk{}नी\appbrk{}नि\appbrk{}त\appbrk{}म्ब\appbrk{}पा\appbrk{}र्श्व\appbrk{}कु\appbrk{}क्षि\appbrk{}स्त\appbrk{}न\appbrk{}क\appbrk{}क्ष\appbrk{}स\appbrk{}न्धि\appbrk{}भा\appbrk{}ग\appbrk{}गताः}}\Afootnote{H; \skt{अक्षि\appbrk{}द\appbrk{}न्त\appbrk{}ना\appbrk{}सा\appbrk{}पा\appbrk{}ङ्ग\appbrk{}श्रो\appbrk{}त्र\appbrk{}ना\appbrk{}भि\appbrk{}ज\appbrk{}ठ\appbrk{}र\appbrk{}से\appbrk{}व\appbrk{}नी\appbrk{}नि\appbrk{}त\appbrk{}म्ब\appbrk{}पा\appbrk{}र्श्व\appbrk{}कु\appbrk{}क्षि\appbrk{}व\appbrk{}क्षः\appbrk{}क\appbrk{}क्षा\appbrk{}स्त\appbrk{}न\appbrk{}स\appbrk{}न्धि\appbrk{}भा\appbrk{}ग\appbrk{}गताः}\ A; \emph{om.}\ K; \skt{अक्षि\appbrk{}द\appbrk{}न्त\appbrk{}ना\appbrk{}सा\appbrk{}व्या\appbrk{}ङ्ग\appbrk{}स्रो\appbrk{}त्र\appbrk{}ना\appbrk{}भि\appbrk{}ज\appbrk{}ठ\appbrk{}र\appbrk{}से\appbrk{}व\appbrk{}नी\appbrk{}नि\appbrk{}त\appbrk{}म्ब\appbrk{}पा\appbrk{}स्व\appbrk{}कु\appbrk{}क्षि\appbrk{}स्त\appbrk{}न\appbrk{}क\appbrk{}क्ष\appbrk{}स\appbrk{}न्धि\appbrk{}भा\appbrk{}ग\appbrk{}गताः}\ N}}\skt{ }\edtext{\skt{सफे\-न\-पू\-या\-नि\-ल\-र\-क्त\-वा\-हि\-नो\-न्तः\-श\-ल्याश्च}}{\lemma{\skt{सफे\appbrk{}न\appbrk{}पू\appbrk{}या\appbrk{}नि\appbrk{}ल\appbrk{}र\appbrk{}क्त\appbrk{}वा\appbrk{}हि\appbrk{}नो\appbrk{}न्तः\appbrk{}श\appbrk{}ल्याश्च}}\Afootnote{H; \skt{सफे\appbrk{}न\appbrk{}पू\appbrk{}य\appbrk{}र\appbrk{}क्ता\appbrk{}नि\appbrk{}ल\appbrk{}वा\appbrk{}हि\appbrk{}नोऽ\appbrk{}न्तः\appbrk{}श\appbrk{}ल्याश्च}\ A; \emph{om.}\ K; \skt{सफे\appbrk{}न\appbrk{}पू\appbrk{}या\appbrk{}नि\appbrk{}ल\appbrk{}र\appbrk{}क्त\appbrk{}वा\appbrk{}हि\appbrk{}नोन्तः}\ N}}\skt{ }\edtext{\skt{दुश्चिकित्स्याः}}{\lemma{\skt{दुश्चिकित्स्याः}}\Afootnote{A H; \emph{om.}\ K N}}\skt{~॥ }\edtext{\skt{रोमा\-न्तो\-प\-न\-ख\-म\-न्म\-ज\-ङ्घा\-स्थि\-सं\-श्रि\-ताश्च}}{\lemma{\skt{रोमा\appbrk{}न्तो\appbrk{}प\appbrk{}न\appbrk{}ख\appbrk{}म\appbrk{}न्म\appbrk{}ज\appbrk{}ङ्घा\appbrk{}स्थि\appbrk{}सं\appbrk{}श्रि\appbrk{}ताश्च}}\Afootnote{H; \skt{अधो\appbrk{}भा\appbrk{}गा\appbrk{}श्चो\appbrk{}र्ध्व\appbrk{}भा\appbrk{}ग\appbrk{}नि\appbrk{}र्वा\appbrk{}हिणो , रोमा\appbrk{}न्तो\appbrk{}प\appbrk{}न\appbrk{}ख\appbrk{}म\appbrk{}र्म\appbrk{}ज\appbrk{}ङ्घा\appbrk{}स्थि\appbrk{}सं\appbrk{}श्रि\appbrk{}ताश्च}\ A; \skt{{\symfont ‡}}\ K; \skt{सल्यादुःचिकित्स्याः~। रोमा\appbrk{}न्तो\appbrk{}प\appbrk{}न\appbrk{}ख\appbrk{}म\appbrk{}र्म\appbrk{}जं\appbrk{}घा\appbrk{}स्थि\appbrk{}सं\appbrk{}स्रि\appbrk{}तास्च}\ N}}\skt{ }\edtext{\edtext{\skt{भगन्दरमपि}}{\lemma{\skt{भगन्दरम्}}\Afootnote{H; \skt{भगन्दर}\ A; \emph{om.}\ K; \skt{भगन्तरम्}\ N}}}{\lemma{\skt{अपि}}\Afootnote{N H; \skt{मपि}\ A; \emph{om.}\ K}}\skt{ }\edtext{\skt{चान्तर्म्मुखं}}{\lemma{\skt{चान्तर्म्मुखं}}\Afootnote{A N H; \emph{om.}\ K}}\skt{ }\edtext{\edtext{\skt{सेव\-नी\-क\-क्चु\-प\-स्थ\-सं\-श्रि\-त\-मिति}}{\lemma{\skt{सेव\appbrk{}नी\appbrk{}क\appbrk{}क्चु\appbrk{}प\appbrk{}स्थ\appbrk{}सं\appbrk{}श्रि\appbrk{}तम्}}\Afootnote{H; \skt{सेव\appbrk{}नी\appbrk{}कु\appbrk{}ट\appbrk{}का\appbrk{}स्थि\appbrk{}सं\appbrk{}श्रि\appbrk{}तम्}\ A; \emph{om.}\ K; \skt{सेव\appbrk{}नी\appbrk{}क\appbrk{}द्यु\appbrk{}पस्थं संस्रितम्}\ N}}}{\lemma{\skt{इति}}\Afootnote{H; \emph{om.}\ A K; \skt{इति~॥ भवति चात्र}\ N}}\skt{~॥}\pend

\pstart\versepar \ledsidenote{\passref{1.23.7}}\edtext{\edtext{\skt{कुष्ठि\-ना\-म्वि\-ष\-जु\-ष्टानां}}{\lemma{\skt{कुष्ठिनाम्}}\Afootnote{A H; \emph{om.}\ K; \skt{कुष्टिनां}\ N}}}{\lemma{\skt{विषजुष्टानां}}\Afootnote{A H; \emph{om.}\ K; \skt{विपदुष्टानां}\ N}}\skt{ }\edtext{\edtext{\skt{शोषि\-णा\-म्म\-धु\-मे\-हि\-नाम्}}{\lemma{\skt{शोषिणाम्}}\Afootnote{A N H; \emph{om.}\ K}}}{\lemma{\skt{मधुमेहिनाम्}}\Afootnote{A N H; \emph{om.}\ K}}\skt{~। }\pend
\pstart\versepar \skt{}\edtext{\skt{व्रणाः}}{\lemma{\skt{व्रणाः}}\Afootnote{A N; \emph{om.}\ K; \skt{व्रणाः 𑑛}\ H}}\skt{ }\edtext{\skt{कृच्छ्रेण}}{\lemma{\skt{कृच्छ्रेण}}\Afootnote{A N H; \emph{om.}\ K}}\skt{ }\edtext{\skt{सिध्यन्ति}}{\lemma{\skt{सिध्यन्ति}}\Afootnote{A N H; \emph{om.}\ K}}\skt{ }\edtext{\edtext{\skt{येषाञ्चापि}}{\lemma{\skt{येषाञ्}}\Afootnote{A N H; \emph{om.}\ K}}}{\lemma{\skt{चापि}}\Afootnote{A N H; \emph{om.}\ K}}\skt{ }\edtext{\skt{व्रणे}}{\lemma{\skt{व्रणे}}\Afootnote{A N H; \emph{om.}\ K}}\skt{ }\edtext{\skt{व्रणाः}}{\lemma{\skt{व्रणाः}}\Afootnote{A N H; \skt{{\symfont ‡}}\ K}}\skt{~॥}\pend

\pstart \ledsidenote{\passref{1.23.8}}\edtext{\skt{अव\-पा\-टि\-का\-नि\-रु\-द्ध\-प्र\-का\-स\-स\-न्नि\-रु\-द्ध\-गु\-द\-ज\-ठ\-र\-ग्र\-न्थि\-क्ष\-त\-क्रि\-मयः}}{\lemma{\skt{अव\appbrk{}पा\appbrk{}टि\appbrk{}का\appbrk{}नि\appbrk{}रु\appbrk{}द्ध\appbrk{}प्र\appbrk{}का\appbrk{}स\appbrk{}स\appbrk{}न्नि\appbrk{}रु\appbrk{}द्ध\appbrk{}गु\appbrk{}द\appbrk{}ज\appbrk{}ठ\appbrk{}र\appbrk{}ग्र\appbrk{}न्थि\appbrk{}क्ष\appbrk{}त\appbrk{}क्रि\appbrk{}मयः}}\Afootnote{H; \skt{अव\appbrk{}पा\appbrk{}टि\appbrk{}का\appbrk{}नि\appbrk{}रु\appbrk{}द्ध\appbrk{}प्र\appbrk{}क\appbrk{}श\appbrk{}स\appbrk{}न्नि\appbrk{}रु\appbrk{}द्ध\appbrk{}गु\appbrk{}द\appbrk{}ज\appbrk{}ठ\appbrk{}राणि, ग्रन्थि\appbrk{}क्ष\appbrk{}त\appbrk{}क्रि\appbrk{}मयः}\ A; \emph{om.}\ K; \skt{अव\appbrk{}पा\appbrk{}टि\appbrk{}का\appbrk{}नि\appbrk{}रु\appbrk{}द्ध\appbrk{}प्र\appbrk{}क\appbrk{}स\appbrk{}सं\appbrk{}नि\appbrk{}रु\appbrk{}द्ध\appbrk{}गु\appbrk{}द\appbrk{}ज\appbrk{}ठ\appbrk{}र\appbrk{}ग्र\appbrk{}न्थि\appbrk{}त्क\appbrk{}त\appbrk{}क्ष\appbrk{}मयः}\ N}}\skt{ }\edtext{\skt{प्रतिश्यायजाः}}{\lemma{\skt{प्रतिश्यायजाः}}\Afootnote{A N H; \emph{om.}\ K}}\skt{ }\edtext{\skt{कोष्ठजाश्च}}{\lemma{\skt{कोष्ठजाश्च}}\Afootnote{A N H; \emph{om.}\ K}}\skt{ }\edtext{\edtext{\skt{त्वग्दो\-षि\-णा\-म्प्र\-मे\-हि\-णाम्वा}}{\lemma{\skt{त्वग्दोषिणाम्}}\Afootnote{A N H; \emph{om.}\ K}}\skt{ ये}}{\lemma{\skt{प्रमेहिणाम्वा ये}}\Afootnote{A N H; \emph{om.}\ K}}\skt{ }\edtext{\skt{प्रतिक्षतेषु}}{\lemma{\skt{प्रतिक्षतेषु}}\Afootnote{N H; \skt{परिक्षतेषु}\ A; \emph{om.}\ K}}\skt{ }\edtext{\skt{दृश्यन्ते}}{\lemma{\skt{दृश्यन्ते}}\Afootnote{A N H; \emph{om.}\ K}}\skt{ }\edtext{\skt{शर्क्करा}}{\lemma{\skt{शर्क्करा}}\Afootnote{A N H; \emph{om.}\ K}}\skt{ }\edtext{\skt{सिकतामेही}}{\lemma{\skt{सिकतामेही}}\Afootnote{N H; \skt{सिकतामेहो}\ A; \emph{om.}\ K}}\skt{ }\edtext{\skt{वात\-कु\-ण्डा\-लि\-का\-ष्ठीला}}{\lemma{\skt{वात\appbrk{}कु\appbrk{}ण्डा\appbrk{}लि\appbrk{}का\appbrk{}ष्ठीला}}\Afootnote{H; \skt{वात\appbrk{}कु\appbrk{}ण्ड\appbrk{}लि\appbrk{}काऽ\appbrk{}ष्ठीला}\ A N; \emph{om.}\ K}}\skt{ }\edtext{\skt{दन्त\-श\-र्क्क\-रो\-प\-कु\-श\-क\-ण्ठ\-का\-शा\-लू\-क\-नि\-ष्को\-ष\-ण\-दू\-षिता}}{\lemma{\skt{दन्त\appbrk{}श\appbrk{}र्क्क\appbrk{}रो\appbrk{}प\appbrk{}कु\appbrk{}श\appbrk{}क\appbrk{}ण्ठ\appbrk{}का\appbrk{}शा\appbrk{}लू\appbrk{}क\appbrk{}नि\appbrk{}ष्को\appbrk{}ष\appbrk{}ण\appbrk{}दू\appbrk{}षिता}}\Afootnote{H; \skt{दन्त\appbrk{}श\appbrk{}र्क\appbrk{}रो\appbrk{}प\appbrk{}कुशः कण्ठशालूकं निष्को\appbrk{}ष\appbrk{}ण\appbrk{}दू\appbrk{}षि\appbrk{}ताश्च}\ A; \emph{om.}\ K; \skt{दन्त\appbrk{}श\appbrk{}र्क\appbrk{}रो\appbrk{}प\appbrk{}कु\appbrk{}श\appbrk{}क\appbrk{}ण्ठ\appbrk{}शा\appbrk{}लू\appbrk{}क\appbrk{}नि\appbrk{}ष्क\appbrk{}ष\appbrk{}ण\appbrk{}दृ\appbrk{}षिता}\ N}}\skt{ }\edtext{\skt{दन्तवेष्टा}}{\lemma{\skt{दन्तवेष्टा}}\Afootnote{A N H; \emph{om.}\ K}}\skt{ }\edtext{\skt{विस\-र्प्पा\-स्थि\-क्ष\-तो\-रः\-क्ष\-त\-यश्च}}{\lemma{\skt{विस\appbrk{}र्प्पा\appbrk{}स्थि\appbrk{}क्ष\appbrk{}तो\appbrk{}रः\appbrk{}क्ष\appbrk{}त\appbrk{}यश्च}}\Afootnote{H; \skt{विस\appbrk{}र्पा\appbrk{}स्थि\appbrk{}क्ष\appbrk{}तो\appbrk{}रः\appbrk{}क्ष\appbrk{}त\appbrk{}व्र\appbrk{}ण\appbrk{}ग्र\appbrk{}न्थि\appbrk{}प्र\appbrk{}भृ\appbrk{}त\appbrk{}यश्च}\ A; \skt{{\symfont ‡}}\ K; \skt{विष\appbrk{}स\appbrk{}र्पा\appbrk{}स्थि\appbrk{}क्ष\appbrk{}तो\appbrk{}रः\appbrk{}क्ष\appbrk{}त\appbrk{}यस्च}\ N}}\skt{ }\edtext{\skt{याप्याः}}{\lemma{\skt{याप्याः}}\Afootnote{A N H; \emph{om.}\ K}}\skt{~॥}\pend

\pstart\versepar \ledsidenote{\passref{1.23.9}}\skt{॥ भवन्ति चात्र~॥ }\pend
\pstart\versepar \skt{साध्या }\edtext{\edtext{\skt{याप्यत्वमायान्ति}}{\lemma{\skt{याप्यत्वम्}}\Afootnote{A N H; \emph{om.}\ K}}}{\lemma{\skt{आयान्ति}}\Afootnote{A N H; \emph{om.}\ K}}\skt{ }\edtext{\edtext{\edtext{\skt{याप्या\-श्चा\-सा\-ध्य\-ता\-मियुः}}{\lemma{\skt{याप्याश्}}\Afootnote{A N H; \emph{om.}\ K}}}{\lemma{\skt{चासाध्यताम्}}\Afootnote{A N H; \emph{om.}\ K}}}{\lemma{\skt{इयुः}}\Afootnote{N H; \skt{तथा}\ A; \emph{om.}\ K}}\skt{~। }\pend
\pstart\versepar \skt{}\edtext{\skt{असाध्याश्चाददुः }\edtext{\edtext{\skt{प्राणा\-न्न\-रा\-णा\-म\-क्रि\-या\-वतां}}{\lemma{\skt{नराणाम्}}\Afootnote{A N H; \emph{om.}\ K}}}{\lemma{\skt{अक्रियावतां}}\Afootnote{A N H; \skt{{\symfont ‡}}\ K}}}{\lemma{\skt{असाध्याश्चाददुः प्राणान्}}\Afootnote{H; \skt{घ्नन्ति प्राणानसाध्यास्तु}\ A; \emph{om.}\ K; \skt{असा\appbrk{}ध्या\appbrk{}स्चा\appbrk{}द\appbrk{}दु\appbrk{}प्रा\appbrk{}णान्}\ N}}\skt{~॥}\edtext{}{\lemma{\skt{॥ भवन्ति चात्र~॥ साध्या}}\Afootnote{N; \skt{साध्या}\ A H; \emph{om.}\ K}}\pend

\pstart\versepar \ledsidenote{\passref{1.23.10}}\edtext{\skt{एवं}}{\lemma{\skt{एवं}}\Afootnote{N H; \skt{याप}\ A; \emph{om.}\ K}}\skt{ }\edtext{\edtext{\skt{याप्य\-म्वि\-जा\-नी\-या\-त्क्रिया}}{\lemma{\skt{याप्यम्}}\Afootnote{N H; \skt{नीयं}\ A; \emph{om.}\ K}}}{\lemma{\skt{विजानीयात्क्रिया}}\Afootnote{A N H; \emph{om.}\ K}}\skt{ }\edtext{\skt{धारयते}}{\lemma{\skt{धारयते}}\Afootnote{A N H; \emph{om.}\ K}}\skt{ }\edtext{\skt{हितं}}{\lemma{\skt{हितं}}\Afootnote{N H; \skt{तु यम्}\ A; \emph{om.}\ K}}\skt{~। }\pend
\pstart\versepar \skt{}\edtext{\skt{क्रियायान्तु}}{\lemma{\skt{क्रियायान्तु}}\Afootnote{A N H; \emph{om.}\ K}}\skt{ }\edtext{\skt{निवृत्तायां}}{\lemma{\skt{निवृत्तायां}}\Afootnote{A N H; \emph{om.}\ K}}\skt{ }\edtext{\skt{सद्य}}{\lemma{\skt{सद्य}}\Afootnote{A N H; \emph{om.}\ K}}\skt{ }\edtext{\skt{एव}}{\lemma{\skt{एव}}\Afootnote{A N H; \skt{{\symfont ‡}}\ K}}\skt{ }\edtext{\skt{विनश्यति}}{\lemma{\skt{विनश्यति}}\Afootnote{A N H; \emph{om.}\ K}}\skt{~॥}\pend

\pstart\versepar \ledsidenote{\passref{1.23.11}}\edtext{\skt{प्राप्ता}}{\lemma{\skt{प्राप्ता}}\Afootnote{A N H; \emph{om.}\ K}}\skt{ }\edtext{\skt{क्रिया}}{\lemma{\skt{क्रिया}}\Afootnote{N H; \skt{क्रियां}\ A; \emph{om.}\ K}}\skt{ }\edtext{\skt{धारयते}}{\lemma{\skt{धारयते}}\Afootnote{N H; \skt{धारयति}\ A; \emph{om.}\ K}}\skt{ }\edtext{\edtext{\skt{याप्य\-व्या\-धि\-त\-मा\-तुरं}}{\lemma{\skt{याप्यव्याधितम्}}\Afootnote{A H; \emph{om.}\ K; \skt{याप्यव्याधिकम्}\ N}}}{\lemma{\skt{आतुरं}}\Afootnote{A N H; \skt{{\symfont ‡}}\ K}}\skt{~। }\pend
\pstart\versepar \skt{}\edtext{\edtext{\skt{प्रप\-त\-न्त\-मि\-वा\-गा\-र\-म्वि\-ष्कम्भः}}{\lemma{\skt{प्रप\appbrk{}त\appbrk{}न्त\appbrk{}मि\appbrk{}वा\appbrk{}गा\appbrk{}रम्}}\Afootnote{H; \skt{प्रप\appbrk{}ति\appbrk{}ष्य\appbrk{}दि\appbrk{}वा\appbrk{}गारं}\ A; \emph{om.}\ K; \skt{प्रप\appbrk{}न्ति\appbrk{}मि\appbrk{}वा\appbrk{}गार}\ N}}}{\lemma{\skt{विष्कम्भः}}\Afootnote{A N H; \emph{om.}\ K}}\skt{ }\edtext{\skt{साधुयोजितः}}{\lemma{\skt{साधुयोजितः}}\Afootnote{A N H; \emph{om.}\ K}}\skt{~॥}\pend

\pstart\versepar \ledsidenote{\passref{1.23.11.1}}\edtext{\skt{क्रियायान्तु}}{\lemma{\skt{क्रियायान्तु}}\Afootnote{N H; \emph{om.}\ K}}\skt{ }\edtext{\skt{निवृत्तायां}}{\lemma{\skt{निवृत्तायां}}\Afootnote{N H; \emph{om.}\ K}}\skt{ }\edtext{\skt{सद्य}}{\lemma{\skt{सद्य}}\Afootnote{N H; \emph{om.}\ K}}\skt{ }\edtext{\skt{एव}}{\lemma{\skt{एव}}\Afootnote{N H; \skt{{\symfont ‡}}\ K}}\skt{ }\edtext{\skt{विनश्यति}}{\lemma{\skt{विनश्यति}}\Afootnote{N H; \emph{om.}\ K}}\skt{~।}\footnoteA{This repeats 1.23.10cd.}\skt{ }\pend
\pstart\versepar \skt{}\edtext{\skt{विष्क\-म्भ\-प\-त\-ना\-द्य\-द्व\-द्गृ\-हस्य}}{\lemma{\skt{विष्क\appbrk{}म्भ\appbrk{}प\appbrk{}त\appbrk{}ना\appbrk{}द्य\appbrk{}द्व\appbrk{}द्गृ\appbrk{}हस्य}}\Afootnote{N H; \emph{om.}\ K}}\skt{ }\edtext{\skt{पतनं}}{\lemma{\skt{पतनं}}\Afootnote{N H; \emph{om.}\ K}}\skt{ }\edtext{\skt{ध्रुवं}}{\lemma{\skt{ध्रुवं}}\Afootnote{N H; \emph{om.}\ K}}\skt{~॥}\pend

\pstart \ledsidenote{\passref{1.23.12}}\edtext{\skt{अत}}{\lemma{\skt{अत}}\Afootnote{A N H; \emph{om.}\ K}}\skt{ }\edtext{\edtext{\skt{ऊर्द्ध्व\-म\-सा\-ध्या\-न्प\-दे\-क्ष्यामः}}{\lemma{\skt{ऊर्द्ध्वम्}}\Afootnote{A H; \emph{om.}\ K; \skt{ऊर्धम्}\ N}}}{\lemma{\skt{असा\appbrk{}ध्या\appbrk{}न्प\appbrk{}दे\appbrk{}क्ष्यामः}}\Afootnote{H; \skt{असाध्यान्वक्ष्यामः}\ A; \emph{om.}\ K; \skt{असा\appbrk{}ध्या\appbrk{}मु\appbrk{}पं\appbrk{}दे\appbrk{}क्ष्यामः}\ N}}\skt{~॥ }\edtext{\skt{मांस\-पि\-ण्ड\-व\-दु\-द्गताः}}{\lemma{\skt{मांस\appbrk{}पि\appbrk{}ण्ड\appbrk{}व\appbrk{}दु\appbrk{}द्गताः}}\Afootnote{N H\textsuperscript{pc}; \skt{मांस\appbrk{}पि\appbrk{}ण्ड\appbrk{}व\appbrk{}दु\appbrk{}द्गताः}\ A; \emph{om.}\ K; \skt{मांसपिण्डो वदुद्गताः}\ H\textsuperscript{ac}}}\skt{ }\edtext{\skt{प्रसे\-कि\-नो\-न्तः\-पूया}}{\lemma{\skt{प्रसे\appbrk{}कि\appbrk{}नो\appbrk{}न्तः\appbrk{}पूया}}\Afootnote{H; \skt{प्रसे\appbrk{}कि\appbrk{}नोऽ\appbrk{}न्तः\appbrk{}पूय}\ A; \emph{om.}\ K; \skt{प्रसे\appbrk{}कि\appbrk{}नो\appbrk{}न्तः\appbrk{}पूय}\ N}}\skt{ }\edtext{\skt{वेदनावन्तः}}{\lemma{\skt{वेदनावन्तः}}\Afootnote{A N H; \emph{om.}\ K}}\skt{~। }\edtext{\skt{अश्वा\-पा\-न\-व\-दु\-द्वृ\-त्तौष्ठाः}}{\lemma{\skt{अश्वा\appbrk{}पा\appbrk{}न\appbrk{}व\appbrk{}दु\appbrk{}द्वृ\appbrk{}त्तौष्ठाः}}\Afootnote{H; \skt{ऽश्वा\appbrk{}पा\appbrk{}न\appbrk{}व\appbrk{}दु\appbrk{}द्वृ\appbrk{}त्तौष्ठाः}\ A; \emph{om.}\ K; \skt{अश्वा\appbrk{}पा\appbrk{}न\appbrk{}व\appbrk{}दु\appbrk{}द्वृ\appbrk{}तोष्ठाः}\ N}}\skt{ }\edtext{\skt{केचित्कठिना}}{\lemma{\skt{केचित्कठिना}}\Afootnote{A N H; \emph{om.}\ K}}\skt{ }\edtext{\skt{गोशृ\-ङ्ग\-व\-दु\-द्गताः}}{\lemma{\skt{गोशृ\appbrk{}ङ्ग\appbrk{}व\appbrk{}दु\appbrk{}द्गताः}}\Afootnote{H; \skt{गोशृ\appbrk{}ङ्ग\appbrk{}व\appbrk{}दु\appbrk{}द्गत}\ A; \emph{om.}\ K; \skt{गोशृ\appbrk{}ङ्ग\appbrk{}व\appbrk{}दु\appbrk{}गताः}\ N}}\skt{~। }\edtext{\skt{मृदु\-मां\-स\-प्र\-रोहाः}}{\lemma{\skt{मृदु\appbrk{}मां\appbrk{}स\appbrk{}प्र\appbrk{}रोहाः}}\Afootnote{A N H; \emph{om.}\ K}}\skt{~। }\edtext{\skt{अपरे}}{\lemma{\skt{अपरे}}\Afootnote{A N H; \emph{om.}\ K}}\skt{ }\edtext{\skt{दुष्ट\-रु\-धि\-रा\-श्रा\-विणः}}{\lemma{\skt{दुष्ट\appbrk{}रु\appbrk{}धि\appbrk{}रा\appbrk{}श्रा\appbrk{}विणः}}\Afootnote{H; \skt{दुष्ट\appbrk{}रु\appbrk{}धि\appbrk{}रा\appbrk{}स्रा\appbrk{}विण}\ A; \emph{om.}\ K; \skt{दुष्ट\appbrk{}रु\appbrk{}धि\appbrk{}र\appbrk{}स्रा\appbrk{}विणः}\ N}}\skt{~। }\edtext{\skt{तन्तु\-पि\-च्छि\-ला\-स्रा\-विणो वा}}{\lemma{\skt{तन्तु\appbrk{}पि\appbrk{}च्छि\appbrk{}ला\appbrk{}स्रा\appbrk{}विणो वा}}\Afootnote{H; \skt{स्तनु\appbrk{}शी\appbrk{}त\appbrk{}पि\appbrk{}च्छि\appbrk{}ला\appbrk{}स्रा\appbrk{}विणो वा}\ A; \skt{{\symfont ‡}}\ K; \skt{तनु\appbrk{}पि\appbrk{}च्छि\appbrk{}ल\appbrk{}स्रा\appbrk{}विणो वा}\ N}}\skt{ }\edtext{\skt{मध्योन्नताः}}{\lemma{\skt{मध्योन्नताः}}\Afootnote{A H; \emph{om.}\ K; \skt{मध्तोन्नताः}\ N}}\skt{~। }\edtext{\skt{केचि\-द\-व\-स\-न्न\-शु\-षि\-र\-प\-र्यन्ताः}}{\lemma{\skt{केचि\appbrk{}द\appbrk{}व\appbrk{}स\appbrk{}न्न\appbrk{}शु\appbrk{}षि\appbrk{}र\appbrk{}प\appbrk{}र्यन्ताः}}\Afootnote{N H; \skt{केचि\appbrk{}द\appbrk{}व\appbrk{}स\appbrk{}न्न\appbrk{}शु\appbrk{}षि\appbrk{}र\appbrk{}प\appbrk{}र्यन्ताः}\ A; \emph{om.}\ K}}\skt{~। }\edtext{\skt{शण\-तू\-ल\-व\-त्स्ना\-यु\-जा\-ल\-वन्तो}}{\lemma{\skt{शण\appbrk{}तू\appbrk{}ल\appbrk{}व\appbrk{}त्स्ना\appbrk{}यु\appbrk{}जा\appbrk{}ल\appbrk{}वन्तो}}\Afootnote{A N H; \emph{om.}\ K}}\skt{ }\edtext{\skt{दुर्द्दर्शनाः}}{\lemma{\skt{दुर्द्दर्शनाः}}\Afootnote{A N H; \emph{om.}\ K}}\skt{~। }\edtext{\skt{वसा\-मे\-दो\-म\-ज्जा\-म\-स्तु\-लु\-ङ्गा\-श्रा\-वि\-णश्च}}{\lemma{\skt{वसा\appbrk{}मे\appbrk{}दो\appbrk{}म\appbrk{}ज्जा\appbrk{}म\appbrk{}स्तु\appbrk{}लु\appbrk{}ङ्गा\appbrk{}श्रा\appbrk{}वि\appbrk{}णश्च}}\Afootnote{H; \skt{वसा\appbrk{}मे\appbrk{}दो\appbrk{}म\appbrk{}ज्ज\appbrk{}म\appbrk{}स्तु\appbrk{}लु\appbrk{}ङ्ग\appbrk{}स्रा\appbrk{}वि\appbrk{}णश्च}\ A; \emph{om.}\ K; \skt{वसा\appbrk{}मे\appbrk{}दो\appbrk{}म\appbrk{}ज्जा\appbrk{}म\appbrk{}स्तु\appbrk{}लु\appbrk{}ङ्ग\appbrk{}स्रा\appbrk{}वि\appbrk{}णस्च}\ N}}\skt{~। }\edtext{\skt{दोषसमुत्था}}{\lemma{\skt{दोषसमुत्था}}\Afootnote{H; \skt{दोषसमुत्थाः}\ A N; \emph{om.}\ K}}\skt{ }\edtext{\skt{पीता\-सि\-त\-मू\-त्र\-पु\-री\-ष\-वा\-त\-वा\-हि\-नश्च}}{\lemma{\skt{पीता\appbrk{}सि\appbrk{}त\appbrk{}मू\appbrk{}त्र\appbrk{}पु\appbrk{}री\appbrk{}ष\appbrk{}वा\appbrk{}त\appbrk{}वा\appbrk{}हि\appbrk{}नश्च}}\Afootnote{A N H; \emph{om.}\ K}}\skt{ }\edtext{\skt{कोष्ठप्राप्ताः~। त}}{\lemma{\skt{कोष्ठप्राप्ताः~। त}}\Afootnote{N H; \skt{कोष्ठस्थाः, त}\ A; \emph{om.}\ K}}\skt{ }\edtext{\skt{एवोभयतो}}{\lemma{\skt{एवोभयतो}}\Afootnote{H; \skt{एवोभयतो}\ A N; \emph{om.}\ K}}\skt{ }\edtext{\skt{भागं}}{\lemma{\skt{भागं}}\Afootnote{N H; \skt{भाग}\ A; \emph{om.}\ K}}\skt{ }\edtext{\skt{व्रणमुखेषु}}{\lemma{\skt{व्रणमुखेषु}}\Afootnote{A N H; \emph{om.}\ K}}\skt{ }\edtext{\skt{पूय\-र\-क्त\-नि\-र्व्वा\-हिणः}}{\lemma{\skt{पूय\appbrk{}र\appbrk{}क्त\appbrk{}नि\appbrk{}र्व्वा\appbrk{}हिणः}}\Afootnote{A H; \emph{om.}\ K; \skt{पूय\appbrk{}र\appbrk{}क्त\appbrk{}नि\appbrk{}र्वा\appbrk{}र्हिणः}\ N}}\skt{ }\edtext{\skt{क्षीणमांसानाञ्च}}{\lemma{\skt{क्षीणमांसानाञ्च}}\Afootnote{A N H; \emph{om.}\ K}}\skt{ }\edtext{\skt{सर्व्वतो}}{\lemma{\skt{सर्व्वतो}}\Afootnote{N H; \skt{सर्वतो}\ A; \emph{om.}\ K}}\skt{ }\edtext{\edtext{\skt{गत\-य\-स्त्व\-णु\-मु\-खा\-न्मां\-स\-बु\-द्बु\-द\-वन्तः}}{\lemma{\skt{गतयस्}}\Afootnote{A N H; \emph{om.}\ K}}}{\lemma{\skt{त्वणु\appbrk{}मु\appbrk{}खा\appbrk{}न्मां\appbrk{}स\appbrk{}बु\appbrk{}द्बु\appbrk{}द\appbrk{}वन्तः}}\Afootnote{H; \skt{चाणुमुखा मांस\appbrk{}बु\appbrk{}द्बु\appbrk{}द\appbrk{}वन्तः}\ A; \emph{om.}\ K; \skt{त्वणुमुखा मांस\appbrk{}बु\appbrk{}द्बु\appbrk{}द\appbrk{}वन्तः}\ N}}\skt{ }\edtext{\skt{सश\-ब्द\-वा\-त\-वा\-हि\-नश्च}}{\lemma{\skt{सश\appbrk{}ब्द\appbrk{}वा\appbrk{}त\appbrk{}वा\appbrk{}हि\appbrk{}नश्च}}\Afootnote{A N H; \emph{om.}\ K}}\skt{ }\edtext{\skt{शिरःकण्ठस्थाः}}{\lemma{\skt{शिरःकण्ठस्थाः}}\Afootnote{A N H; \emph{om.}\ K}}\skt{~। }\edtext{\skt{क्षीणमांसानाञ्च}}{\lemma{\skt{क्षीणमांसानाञ्च}}\Afootnote{A N H; \emph{om.}\ K}}\skt{ }\edtext{\skt{पूय\-रु\-धि\-र\-नि\-र्व्वा\-पिनः}}{\lemma{\skt{पूय\appbrk{}रु\appbrk{}धि\appbrk{}र\appbrk{}नि\appbrk{}र्व्वा\appbrk{}पिनः}}\Afootnote{N H; \skt{पूय\appbrk{}र\appbrk{}क्त\appbrk{}नि\appbrk{}र्वा\appbrk{}हिणो}\ A; \emph{om.}\ K}}\skt{~। }\edtext{\skt{अरो\-च\-का\-वि\-पा\-क\-का\-स\-श्वा\-सो\-प\-द्र\-व\-युक्ताः}}{\lemma{\skt{अरो\appbrk{}च\appbrk{}का\appbrk{}वि\appbrk{}पा\appbrk{}क\appbrk{}का\appbrk{}स\appbrk{}श्वा\appbrk{}सो\appbrk{}प\appbrk{}द्र\appbrk{}व\appbrk{}युक्ताः}}\Afootnote{A N H; \emph{om.}\ K}}\skt{~। }\edtext{\skt{भिन्ने च}}{\lemma{\skt{भिन्ने च}}\Afootnote{N H; \skt{भिन्ने वा}\ A; \emph{om.}\ K}}\skt{ }\edtext{\skt{शिरःकपाल}}{\lemma{\skt{शिरःकपाल}}\Afootnote{H; \skt{शिरःकपाले}\ A N; \emph{om.}\ K}}\skt{ }\edtext{\skt{यत्र}}{\lemma{\skt{यत्र}}\Afootnote{A N H; \emph{om.}\ K}}\skt{ }\edtext{\edtext{\skt{मस्तु\-लु\-ङ्ग\-द\-र्श\-न\-न्त्रि\-दो\-ष\-लि\-ङ्ग\-प्रा\-दु\-र्भावः}}{\lemma{\skt{मस्तु\appbrk{}लु\appbrk{}ङ्ग\appbrk{}द\appbrk{}र्श\appbrk{}नन्}}\Afootnote{A N H; \emph{om.}\ K}}}{\lemma{\skt{त्रिदो\appbrk{}ष\appbrk{}लि\appbrk{}ङ्ग\appbrk{}प्रा\appbrk{}दु\appbrk{}र्भावः}}\Afootnote{A N H; \emph{om.}\ K}}\skt{~। }\edtext{\skt{कासश्वासौ वा}}{\lemma{\skt{कासश्वासौ वा}}\Afootnote{A N H; \emph{om.}\ K}}\skt{ }\edtext{\skt{यस्येति}}{\lemma{\skt{यस्येति}}\Afootnote{A N H; \emph{om.}\ K}}\skt{~॥}\pend

\pstart\versepar \ledsidenote{\passref{1.23.13}}\edtext{\skt{भवन्ति}}{\lemma{\skt{भवन्ति}}\Afootnote{H; \skt{भवति}\ A N; \emph{om.}\ K}}\skt{ }\edtext{\skt{चात्र}}{\lemma{\skt{चात्र}}\Afootnote{A N H; \emph{om.}\ K}}\skt{ }\edtext{\skt{श्लोकाः}}{\lemma{\skt{श्लोकाः}}\Afootnote{H; \emph{om.}\ A K; \skt{श्लोकः}\ N}}\skt{~॥ }\pend
\pstart\versepar \skt{}\edtext{\edtext{\skt{वसामेदोऽथ}}{\lemma{\skt{वसामेदो}}\Afootnote{H; \skt{वसां मेदो}\ A; \emph{om.}\ K; \skt{वसामेदो}\ N}}}{\lemma{\skt{ऽथ}}\Afootnote{A H; \skt{{\symfont ‡}}\ K; \skt{थ}\ N}}\skt{ }\edtext{\edtext{\skt{मज्जा\-न\-म्म\-स्तु\-लु\-ङ्गञ्च}}{\lemma{\skt{मज्जानम्}}\Afootnote{A N H; \emph{om.}\ K}}\skt{ यः}}{\lemma{\skt{मस्तुलुङ्गञ्च यः}}\Afootnote{A N H; \emph{om.}\ K}}\skt{ स्रवेत्~। }\pend
\pstart\versepar \skt{आगन्तुजो }\edtext{\skt{व्रणः}}{\lemma{\skt{व्रणः}}\Afootnote{A N H; \emph{om.}\ K}}\skt{ }\edtext{\skt{सिध्येन्न}}{\lemma{\skt{सिध्येन्न}}\Afootnote{A N H; \emph{om.}\ K}}\skt{ }\edtext{\skt{सिध्ये\-द्दो\-ष\-स\-म्भवः}}{\lemma{\skt{सिध्ये\appbrk{}द्दो\appbrk{}ष\appbrk{}स\appbrk{}म्भवः}}\Afootnote{A N; \emph{om.}\ K; \skt{सिध्ये\appbrk{}द्दो\appbrk{}ष\appbrk{}स\appbrk{}म्भवः~॥ 𑑛}\ H}}\skt{~॥}\edtext{}{\lemma{\skt{स्रवेत्~। आगन्तुजो}}\Afootnote{N H; \skt{स्रवेत्~। आगन्तुस्तु}\ A; \emph{om.}\ K}}\pend

\pstart\versepar \ledsidenote{\passref{1.23.14}}\edtext{\skt{अमर्म्मोपहिते}}{\lemma{\skt{अमर्म्मोपहिते}}\Afootnote{A H; \skt{{\symfont ‡}}\ K; \skt{अमन्मोपहिते}\ N}}\skt{ }\edtext{\skt{देशे}}{\lemma{\skt{देशे}}\Afootnote{A N H; \emph{om.}\ K}}\skt{ }\edtext{\skt{सिरा\-सं\-ध्य\-स्थि\-व\-र्ज्जिते}}{\lemma{\skt{सिरा\appbrk{}सं\appbrk{}ध्य\appbrk{}स्थि\appbrk{}व\appbrk{}र्ज्जिते}}\Afootnote{A H; \emph{om.}\ K; \skt{शिरा\appbrk{}सं\appbrk{}न्ध्य\appbrk{}स्थि\appbrk{}व\appbrk{}र्जिते}\ N}}\skt{~। }\pend
\pstart\versepar \skt{}\edtext{\skt{विकारो }\edtext{\skt{योऽनुपर्येति}}{\lemma{\skt{ऽनुपर्येति}}\Afootnote{A H; \emph{om.}\ K; \skt{नुपर्येति}\ N}}}{\lemma{\skt{विकारो यो}}\Afootnote{N H; \skt{विकारोयो}\ A; \emph{om.}\ K}}\skt{ }\edtext{\skt{तदसाध्यस्य}}{\lemma{\skt{तदसाध्यस्य}}\Afootnote{A N H; \emph{om.}\ K}}\skt{ }\edtext{\skt{लक्षणं}}{\lemma{\skt{लक्षणं}}\Afootnote{A N H; \emph{om.}\ K}}\skt{~॥}\pend

\pstart\versepar \ledsidenote{\passref{1.23.15}}\edtext{\skt{क्रमेणोपचयं}}{\lemma{\skt{क्रमेणोपचयं}}\Afootnote{A H; \emph{om.}\ K; \skt{क्रमेणोपयं}\ N}}\skt{ }\edtext{\skt{प्राप्तो}}{\lemma{\skt{प्राप्तो}}\Afootnote{\skt{प्राप्य}\ A; \emph{om.}\ K; \skt{प्राप्तो}\ N; \skt{प्राप्तो 𑑛}\ H}}\skt{ }\edtext{\skt{धातूननुगतः}}{\lemma{\skt{धातूननुगतः}}\Afootnote{A H; \emph{om.}\ K; \skt{धतूननुगतः}\ N}}\skt{ शनैः~। }\pend
\pstart\versepar \skt{न }\edtext{\edtext{\skt{शक्य\-मु\-न्मू\-ल\-यितुं}}{\lemma{\skt{शक्यम्}}\Afootnote{H; \skt{शक्य}\ A N; \emph{om.}\ K}}}{\lemma{\skt{उन्मूलयितुं}}\Afootnote{A H; \emph{om.}\ K; \skt{मुन्मूलयितुं}\ N}}\skt{ }\edtext{\skt{वृद्धो}}{\lemma{\skt{वृद्धो}}\Afootnote{A N H; \emph{om.}\ K}}\skt{ }\edtext{\skt{वृक्ष}}{\lemma{\skt{वृक्ष}}\Afootnote{A N H; \emph{om.}\ K}}\skt{ }\edtext{\skt{इवामयः}}{\lemma{\skt{इवामयः}}\Afootnote{A N H; \skt{{\symfont ‡}}\ K}}\skt{~॥}\edtext{}{\lemma{\skt{शनैः~। न}}\Afootnote{A N H; \emph{om.}\ K}}\pend

\pstart\versepar \ledsidenote{\passref{1.23.16}}\edtext{\skt{स स्थिर\-त्वा\-न्म\-ह\-त्त्वाच्च}}{\lemma{\skt{स स्थिर\appbrk{}त्वा\appbrk{}न्म\appbrk{}ह\appbrk{}त्त्वाच्च}}\Afootnote{A N H; \emph{om.}\ K}}\skt{ }\edtext{\skt{धात्व\-नु\-क्र\-म\-णेन च}}{\lemma{\skt{धात्व\appbrk{}नु\appbrk{}क्र\appbrk{}म\appbrk{}णेन च}}\Afootnote{A H; \skt{{\symfont ‡}}\ K; \skt{धात्व\appbrk{}र\appbrk{}नु\appbrk{}क्र\appbrk{}मे\appbrk{}णेन च}\ N}}\skt{~। }\pend
\pstart\versepar \skt{}\edtext{\edtext{\skt{निह\-न्त्यौ\-ष\-ध\-वी\-र्याणि}}{\lemma{\skt{निहन्त्य्}}\Afootnote{A N H; \emph{om.}\ K}}}{\lemma{\skt{औषधवीर्याणि}}\Afootnote{A N H; \emph{om.}\ K}}\skt{ }\edtext{\skt{बल\-व\-त्त्वा\-त्त\-थैव}}{\lemma{\skt{बल\appbrk{}व\appbrk{}त्त्वा\appbrk{}त्त\appbrk{}थैव}}\Afootnote{N H; \skt{मन्त्रान्दुष्टग्रहो य}\ A; \emph{om.}\ K}}\skt{ }\edtext{\skt{च}}{\lemma{\skt{च}}\Afootnote{N H; \skt{था}\ A; \emph{om.}\ K}}\skt{~॥}\pend

\pstart\versepar \ledsidenote{\passref{1.23.17}}\edtext{\skt{अतो यो}}{\lemma{\skt{अतो यो}}\Afootnote{A N H; \emph{om.}\ K}}\skt{ }\edtext{\skt{विपरीतः}}{\lemma{\skt{विपरीतः}}\Afootnote{A H; \emph{om.}\ K; \skt{विपऋतः}\ N}}\skt{ }\edtext{\skt{स्यात्सुखसाध्यः स}}{\lemma{\skt{स्यात्सुखसाध्यः स}}\Afootnote{A N H; \emph{om.}\ K}}\skt{ }\edtext{\skt{उच्यते}}{\lemma{\skt{उच्यते}}\Afootnote{A N H; \emph{om.}\ K}}\skt{~। }\pend
\pstart\versepar \skt{}\edtext{\skt{अबद्धमूलः}}{\lemma{\skt{अबद्धमूलः}}\Afootnote{A H; \skt{{\symfont ‡}}\ K; \skt{अलन्दमूल}\ N}}\skt{ }\edtext{\skt{क्षुपको}}{\lemma{\skt{क्षुपको}}\Afootnote{A N H; \emph{om.}\ K}}\skt{ }\edtext{\skt{यद्वदुत्पाटने}}{\lemma{\skt{यद्वदुत्पाटने}}\Afootnote{A N H; \emph{om.}\ K}}\skt{ }\edtext{\skt{सुखं}}{\lemma{\skt{सुखं}}\Afootnote{N H; \skt{सुखः}\ A; \emph{om.}\ K}}\skt{~॥}\pend

\pstart\versepar \ledsidenote{\passref{1.23.18}}\edtext{\edtext{\edtext{\skt{त्रिभि\-र्द्दो\-षै\-र\-ना\-क्रान्तः}}{\lemma{\skt{त्रिभिर्}}\Afootnote{A N H; \emph{om.}\ K}}}{\lemma{\skt{द्दोषैर्}}\Afootnote{A N H; \emph{om.}\ K}}}{\lemma{\skt{अनाक्रान्तः}}\Afootnote{A N H; \emph{om.}\ K}}\skt{ }\edtext{\skt{श्यावौष्ठः}}{\lemma{\skt{श्यावौष्ठः}}\Afootnote{A H; \emph{om.}\ K; \skt{स्यावोष्ठ}\ N}}\skt{ }\edtext{\skt{पिडकी}}{\lemma{\skt{पिडकी}}\Afootnote{A N H; \emph{om.}\ K}}\skt{ }\edtext{\skt{समः}}{\lemma{\skt{समः}}\Afootnote{A N H; \emph{om.}\ K}}\skt{~। }\pend
\pstart\versepar \skt{}\edtext{\skt{अवेदनो}}{\lemma{\skt{अवेदनो}}\Afootnote{A N H; \emph{om.}\ K}}\skt{ }\edtext{\skt{निरास्रावो}}{\lemma{\skt{निरास्रावो}}\Afootnote{A N H; \emph{om.}\ K}}\skt{ }\edtext{\skt{व्रणः}}{\lemma{\skt{व्रणः}}\Afootnote{A N H; \emph{om.}\ K}}\skt{ }\edtext{\skt{शुद्ध}}{\lemma{\skt{शुद्ध}}\Afootnote{A N H; \emph{om.}\ K}}\skt{ }\edtext{\skt{इति}}{\lemma{\skt{इति}}\Afootnote{N H; \skt{इहो}\ A; \skt{{\symfont ‡}}\ K}}\skt{ }\edtext{\skt{स्मृतः}}{\lemma{\skt{स्मृतः}}\Afootnote{H; \skt{च्यते}\ A; \emph{om.}\ K; \skt{स्मृत}\ N}}\skt{~॥}\pend

\pstart\versepar \ledsidenote{\passref{1.23.19}}\edtext{\skt{कपो\-त\-व\-र्ण्ण\-प्र\-तिमो}}{\lemma{\skt{कपो\appbrk{}त\appbrk{}व\appbrk{}र्ण्ण\appbrk{}प्र\appbrk{}तिमो}}\Afootnote{H; \skt{कपो\appbrk{}त\appbrk{}व\appbrk{}र्ण\appbrk{}प्र\appbrk{}तिमा}\ A N; \emph{om.}\ K}}\skt{ }\edtext{\skt{यस्यान्ते}}{\lemma{\skt{यस्यान्ते}}\Afootnote{H; \skt{यस्यान्ताः}\ A; \emph{om.}\ K; \skt{यस्यान्ता}\ N}}\skt{ }\edtext{\skt{क्लेदवर्ज्जिताः}}{\lemma{\skt{क्लेदवर्ज्जिताः}}\Afootnote{A H; \emph{om.}\ K; \skt{क्लदवर्जिताः}\ N}}\skt{~॥ }\pend
\pstart\versepar \skt{}\edtext{\skt{स्थिराश्च}}{\lemma{\skt{स्थिराश्च}}\Afootnote{N H; \skt{स्थिराश्चि}\ A; \emph{om.}\ K}}\skt{ }\edtext{\skt{पिटकावन्तो}}{\lemma{\skt{पिटकावन्तो}}\Afootnote{H; \skt{पिटिकावन्तो}\ A; \emph{om.}\ K; \skt{पिटकोवन्तो}\ N}}\skt{ }\edtext{\skt{रोहतीति }\edtext{\skt{तमादिशेत्}}{\lemma{\skt{आदिशेत्}}\Afootnote{A N H; \skt{{\symfont ‡}}\ K}}}{\lemma{\skt{रोहतीति तम्}}\Afootnote{A N H; \emph{om.}\ K}}\skt{~॥}\pend

\pstart\versepar \ledsidenote{\passref{1.23.20}}\edtext{\edtext{\edtext{\edtext{\edtext{\skt{रूढ\-व\-र्म्मा\-ण\-म\-ग्र\-न्थि\-म\-शू\-न\-म\-रु\-ज\-म्व्रणं}}{\lemma{\skt{रूढवर्म्माणम्}}\Afootnote{H; \skt{रूढवर्त्मानम्}\ A; \emph{om.}\ K; \skt{गूढवर्मान}\ N}}}{\lemma{\skt{अग्रन्थिम्}}\Afootnote{A H; \emph{om.}\ K; \skt{मग्रन्थि}\ N}}}{\lemma{\skt{अशूनम्}}\Afootnote{A H; \emph{om.}\ K; \skt{मसूनं}\ N}}}{\lemma{\skt{अरुजम्}}\Afootnote{A H; \emph{om.}\ K; \skt{मरुजं}\ N}}}{\lemma{\skt{व्रणं}}\Afootnote{A N H; \emph{om.}\ K}}\skt{~। }\pend
\pstart\versepar \skt{}\edtext{\skt{त्वक्सवर्ण्णं}}{\lemma{\skt{त्वक्सवर्ण्णं}}\Afootnote{A H; \emph{om.}\ K; \skt{त्वक्सवर्ण्ण}\ N}}\skt{ }\edtext{\skt{समतलं}}{\lemma{\skt{समतलं}}\Afootnote{A N H; \emph{om.}\ K}}\skt{ }\edtext{\edtext{\skt{सम्य\-ग्रू\-ढ\-म्वि\-नि\-र्द्दि\-शेत्}}{\lemma{\skt{सम्यग्रूढम्}}\Afootnote{A N H; \emph{om.}\ K}}}{\lemma{\skt{विनिर्द्दिशेत्}}\Afootnote{A N H; \skt{{\symfont ‡}}\ K}}\skt{~॥}\pend

\pstart\versepar \ledsidenote{\passref{1.23.21}}\edtext{\skt{दोष\-प्र\-को\-पा\-द्व्या\-या\-मा\-द\-भि\-घा\-ता\-द\-जी\-र्ण्णतः}}{\lemma{\skt{दोष\appbrk{}प्र\appbrk{}को\appbrk{}पा\appbrk{}द्व्या\appbrk{}या\appbrk{}मा\appbrk{}द\appbrk{}भि\appbrk{}घा\appbrk{}ता\appbrk{}द\appbrk{}जी\appbrk{}र्ण्णतः}}\Afootnote{N H; \skt{दोष\appbrk{}प्र\appbrk{}को\appbrk{}पा\appbrk{}द्व्या\appbrk{}या\appbrk{}मा\appbrk{}द\appbrk{}भि\appbrk{}घा\appbrk{}ता\appbrk{}द\appbrk{}जी\appbrk{}र्णतः}\ A; \emph{om.}\ K}}\skt{~। }\pend
\pstart\versepar \skt{}\edtext{\skt{हर्षा\-त्क्रो\-धा\-द्भ\-या\-द्वापि}}{\lemma{\skt{हर्षा\appbrk{}त्क्रो\appbrk{}धा\appbrk{}द्भ\appbrk{}या\appbrk{}द्वापि}}\Afootnote{A N H; \emph{om.}\ K}}\skt{ }\edtext{\skt{व्रणो}}{\lemma{\skt{व्रणो}}\Afootnote{A N H; \emph{om.}\ K}}\skt{ }\edtext{\edtext{\skt{रूढोऽपि}}{\lemma{\skt{रूढो}}\Afootnote{A N H; \emph{om.}\ K}}}{\lemma{\skt{ऽपि}}\Afootnote{A H; \skt{{\symfont ‡}}\ K; \skt{वि}\ N}}\skt{ }\edtext{\skt{दीर्यते}}{\lemma{\skt{दीर्यते}}\Afootnote{A N H; \emph{om.}\ K}}\skt{~॥ २३~॥}\pend

\endnumbering
\beginnumbering
\eledsection*{Adhyāya 24}
\pstart \ledsidenote{\passref{1.24.1}}\skt{अथातो व्याधि\-स\-मु\-द्दे\-शीयं }\edtext{\skt{व्याख्यास्यामः}}{\lemma{\skt{व्याख्यास्यामः}}\Afootnote{H; \skt{अध्यायं व्याख्यास्यामः}\ A; \skt{अध्याय व्याख्यामः}\ N}}\skt{~॥}\pend

\pstart \ledsidenote{\passref{1.24.3}}\edtext{\skt{द्विविधा}}{\lemma{\skt{द्विविधा}}\Afootnote{N H; \skt{द्विविधास्तु}\ A}}\skt{ व्याधयः शस्त्रसाध्याः स्नेहा\-दि\-क्रि\-या\-सा\-ध्याश्च~। तत्र शस्त्रसाध्येषु स्नेहादिक्रिया न प्रतिषिध्यते }\edtext{\skt{स्नेहा\-दि\-क्रि\-या\-सा\-ध्येषु}}{\lemma{\skt{स्नेहा\appbrk{}दि\appbrk{}क्रि\appbrk{}या\appbrk{}सा\appbrk{}ध्येषु}}\Afootnote{A N; \skt{स्नेहादिसाध्येषु}\ H}}\skt{ शस्त्रकर्म न क्रियते~॥}\pend

\pstart \ledsidenote{\passref{1.24.4}}\edtext{\skt{अस्मिंस्तु}}{\lemma{\skt{अस्मिंस्तु}}\Afootnote{N H; \skt{अस्मिन्पुनः}\ A}}\skt{ शास्त्रे }\edtext{\skt{सर्वत्र}}{\lemma{\skt{सर्वत्र}}\Afootnote{N H; \skt{सर्वतन्त्र}\ A}}\skt{ सामान्यात्सर्वेषां व्याधीनां }\edtext{\skt{यथा\-स्थौ\-ल्ये\-ना\-व\-रोधः}}{\lemma{\skt{यथा\appbrk{}स्थौ\appbrk{}ल्ये\appbrk{}ना\appbrk{}व\appbrk{}रोधः}}\Afootnote{N H; \skt{यथा\appbrk{}स्थू\appbrk{}ल\appbrk{}म\appbrk{}व\appbrk{}रोधः}\ A}}\skt{ क्रियते~। प्रागभिहितं }\edtext{\skt{तद्दुः\-ख\-सं\-यो\-गा\-द्व्या\-धि\-रिति}}{\lemma{\skt{तद्दुः\appbrk{}ख\appbrk{}सं\appbrk{}यो\appbrk{}गा\appbrk{}द्व्या\appbrk{}धिर्}}\Afootnote{N H; \skt{‘त\appbrk{}द्दुः\appbrk{}ख\appbrk{}सं\appbrk{}योगा व्याधय’ (सू. अ. १)}\ A}}\skt{~॥ }\edtext{\skt{तच्च}}{\lemma{\skt{तच्च}}\Afootnote{A H; \skt{तच्च}\ N}}\skt{ दुःखं त्रिविधं~। आध्या\-त्मि\-क\-मा\-धि\-भौ\-ति\-क\-मा\-धि\-दै\-वि\-क\-मिति~। }\edtext{\skt{तच्च}}{\lemma{\skt{तच्च}}\Afootnote{N H; \skt{तत्तु}\ A}}\skt{ }\edtext{\skt{दुःखं}}{\lemma{\skt{दुःखं}}\Afootnote{N H; \emph{om.}\ A}}\skt{ सप्तविधे }\edtext{\skt{व्याधा\-वु\-प\-नि\-प\-तति}}{\lemma{\skt{उपनिपतति}}\Afootnote{N H; \skt{उपनिपतति~। ते}\ A}}\skt{~। }\edtext{\skt{सप्तविधास्तु}}{\lemma{\skt{सप्तविधास्तु}}\Afootnote{N H; \skt{पुनः सप्तविधा}\ A}}\skt{ व्याधयः~। तद्यथा आदि\-ब\-ल\-प्र\-वृत्ताः~। जन्म\-ब\-ल\-प्र\-वृत्ताः~। }\edtext{\skt{दोष\-ब\-ल\-प्र\-वृत्ताः}}{\lemma{\skt{दोष\appbrk{}ब\appbrk{}ल\appbrk{}प्र\appbrk{}वृत्ताः}}\Afootnote{\skt{दोष\appbrk{}ब\appbrk{}ल\appbrk{}प्र\appbrk{}वृत्ताः, सङ्घा\appbrk{}त\appbrk{}ब\appbrk{}ल\appbrk{}प्र\appbrk{}वृत्ताः}\ A; \skt{दोष\appbrk{}ब\appbrk{}ल\appbrk{}प्र\appbrk{}वृत्ताः~। संघात}\ N; \skt{दोष\appbrk{}ब\appbrk{}ल\appbrk{}प्र\appbrk{}वृत्ताः~। काल\appbrk{}ब\appbrk{}ल\appbrk{}प्र\appbrk{}वृत्ताः}\ H}}\skt{~। }\edtext{\skt{संघा\-त\-ब\-ल\-प्र\-वृत्ताः}}{\lemma{\skt{संघा\appbrk{}त\appbrk{}ब\appbrk{}ल\appbrk{}प्र\appbrk{}वृत्ताः}}\Afootnote{H; \skt{काल\appbrk{}ब\appbrk{}ल\appbrk{}प्र\appbrk{}वृत्ताः}\ A; \skt{बलप्रवृत्ताः}\ N}}\skt{~। दैव\-ब\-ल\-प्र\-वृत्ताः~। }\edtext{\skt{स्वभा\-व\-ब\-ल\-प्र\-वृत्ताः}}{\lemma{\skt{स्वभा\appbrk{}व\appbrk{}ब\appbrk{}ल\appbrk{}प्र\appbrk{}वृत्ताः}}\Afootnote{A; \skt{स्वभा\appbrk{}व\appbrk{}ब\appbrk{}ल\appbrk{}प्र\appbrk{}वृत्ताः}\ N H}}\skt{~। }\edtext{\skt{इति}}{\lemma{\skt{इति}}\Afootnote{A; \skt{ति}\ N; \skt{चेति}\ H}}\skt{~॥}\pend

\pstart \ledsidenote{\passref{1.24.5}}\edtext{\skt{तत्रा\-दि\-ब\-ल\-प्र\-वृत्ता}}{\lemma{\skt{तत्रा\appbrk{}दि\appbrk{}ब\appbrk{}ल\appbrk{}प्र\appbrk{}वृत्ता}}\Afootnote{N H; \skt{तत्र, आदि\appbrk{}ब\appbrk{}ल\appbrk{}प्र\appbrk{}वृत्ता}\ A}}\skt{ }\edtext{\skt{नाम}}{\lemma{\skt{नाम}}\Afootnote{N; \skt{ये}\ A; \emph{om.}\ H}}\skt{ शुक्र\-शो\-णि\-त\-दो\-षा\-न्वयाः~। }\edtext{\skt{कुष्ठा\-र्शः\-प्र\-भृ\-तयः~। }\edtext{\skt{तेऽपि}}{\lemma{\skt{ऽपि}}\Afootnote{A; \skt{पि}\ N H}}}{\lemma{\skt{कुष्ठा\appbrk{}र्शः\appbrk{}प्र\appbrk{}भृ\appbrk{}तयः~। ते}}\Afootnote{A H; \skt{कुष्ठा\appbrk{}र्श\appbrk{}प्र\appbrk{}भृ\appbrk{}तयः~॥ ते}\ N}}\skt{ }\edtext{\skt{द्विविविधा}}{\lemma{\skt{द्विविविधा}}\Afootnote{N; \skt{द्विविधाः}\ A H}}\skt{ }\edtext{\skt{मातृ\-जा\-पि\-तृ\-जाश्च}}{\lemma{\skt{मातृ\appbrk{}जा\appbrk{}पि\appbrk{}तृ\appbrk{}जाश्च}}\Afootnote{N H; \skt{मातृजाः, पितृजाश्च}\ A}}\skt{~॥ जन्म\-ब\-ल\-प्र\-वृत्ता }\edtext{\skt{नाम ये}}{\lemma{\skt{नाम ये}}\Afootnote{N H; \skt{ये}\ A}}\skt{ }\edtext{\edtext{\skt{मातु\-र\-प\-चा\-रा\-त्प\-ङ्गु\-ज\-ड\-जा\-त्य\-न्ध\-मू\-क\-ब\-धि\-र\-मि\-न्मि\-न\-वा\-म\-न\-प्र\-भृ\-तयो}}{\lemma{\skt{मातुर्}}\Afootnote{N H; \skt{मातु}\ A}}}{\lemma{\skt{अप\appbrk{}चा\appbrk{}रा\appbrk{}त्प\appbrk{}ङ्गु\appbrk{}ज\appbrk{}ड\appbrk{}जा\appbrk{}त्य\appbrk{}न्ध\appbrk{}मू\appbrk{}क\appbrk{}ब\appbrk{}धि\appbrk{}र\appbrk{}मि\appbrk{}न्मि\appbrk{}न\appbrk{}वा\appbrk{}म\appbrk{}न\appbrk{}प्र\appbrk{}भृ\appbrk{}तयो}}\Afootnote{\skt{रप\appbrk{}चा\appbrk{}रा\appbrk{}त्प\appbrk{}ङ्गु\appbrk{}जा\appbrk{}त्य\appbrk{}न्ध\appbrk{}ब\appbrk{}धि\appbrk{}र\appbrk{}मू\appbrk{}क\appbrk{}मि\appbrk{}न्मि\appbrk{}न\appbrk{}वा\appbrk{}म\appbrk{}न\appbrk{}प्र\appbrk{}भृ\appbrk{}तयो}\ A; \skt{आपचारात्~। पङ्गु\appbrk{}ज\appbrk{}ड\appbrk{}जा\appbrk{}त्य\appbrk{}न्ध\appbrk{}मू\appbrk{}क\appbrk{}ध\appbrk{}ब\appbrk{}धि\appbrk{}र\appbrk{}मि\appbrk{}न्मि\appbrk{}न\appbrk{}वा\appbrk{}म\appbrk{}न\appbrk{}प्र\appbrk{}भृ\appbrk{}तयो}\ N; \skt{अप\appbrk{}चा\appbrk{}रा\appbrk{}त्प\appbrk{}ङ्गु\appbrk{}ज\appbrk{}ड\appbrk{}जा\appbrk{}त्य\appbrk{}न्ध\appbrk{}मू\appbrk{}क\appbrk{}व\appbrk{}धि\appbrk{}र\appbrk{}मि\appbrk{}र्म्मि\appbrk{}ण\appbrk{}वा\appbrk{}म\appbrk{}न\appbrk{}प्र\appbrk{}भृ\appbrk{}तयो}\ H}}\skt{ जायन्ते ते }\edtext{\skt{द्विविधा}}{\lemma{\skt{द्विविधा}}\Afootnote{N H; \skt{ऽपि द्विविधाः}\ A}}\skt{ रसकृता दौहृ\-दा\-प\-चा\-र\-कृ\-ताश्च~॥ दोष\-ब\-ल\-प्र\-वृत्ता }\edtext{\skt{नाम य}}{\lemma{\skt{नाम य}}\Afootnote{H; \skt{ये}\ A; \skt{नाम~। ये}\ N}}\skt{ }\edtext{\edtext{\skt{आत\-ङ्का\-प\-चा\-र\-कृ\-तास्ते}}{\lemma{\skt{आत\appbrk{}ङ्का\appbrk{}प\appbrk{}चा\appbrk{}र\appbrk{}कृ\appbrk{}तास्}}\Afootnote{N H; \skt{आत\appbrk{}ङ्क\appbrk{}स\appbrk{}मु\appbrk{}त्पन्ना मिथ्या\appbrk{}हा\appbrk{}रा\appbrk{}चा\appbrk{}र\appbrk{}कृ\appbrk{}ताश्च; तेऽपि द्विविधाः आमा\appbrk{}श\appbrk{}य\appbrk{}स\appbrk{}मुत्थाः}\ A}}}{\lemma{\skt{ते}}\Afootnote{N; \skt{पक्वा\appbrk{}श\appbrk{}य\appbrk{}स\appbrk{}मु\appbrk{}त्थाश्च; पुनश्च}\ A; \skt{त}\ H}}\skt{ }\edtext{\skt{द्विविधाः}}{\lemma{\skt{द्विविधाः}}\Afootnote{A H; \skt{द्विविधा}\ N}}\skt{~। शारीरा }\edtext{\skt{मानसाश्च}}{\lemma{\skt{मानसाश्च}}\Afootnote{N H; \skt{मानसाश्च~। त एते आध्यात्मिकाः}\ A}}\skt{~॥}\pend

\pstart \ledsidenote{\passref{1.24.7a}}\skt{काल\-ब\-ल\-प्र\-वृत्ता }\edtext{\skt{नाम ये}}{\lemma{\skt{नाम ये}}\Afootnote{N H; \skt{ये}\ A}}\skt{ }\edtext{\skt{शीतो\-ष्ण\-वा\-त\-व\-र्षा\-प्र\-भृ\-तिभिः}}{\lemma{\skt{शीतो\appbrk{}ष्ण\appbrk{}वा\appbrk{}त\appbrk{}व\appbrk{}र्षा\appbrk{}प्र\appbrk{}भृ\appbrk{}तिभिः}}\Afootnote{\skt{शीतो\appbrk{}ष्ण\appbrk{}वा\appbrk{}त\appbrk{}व\appbrk{}र्षा\appbrk{}त\appbrk{}प\appbrk{}प्र\appbrk{}भृति}\ A; \skt{शीतो\appbrk{}ष्ण\appbrk{}व\appbrk{}वा\appbrk{}त\appbrk{}व\appbrk{}र्षा\appbrk{}प्र\appbrk{}भृ\appbrk{}त्तिभिः}\ N; \skt{शीतो\appbrk{}ष्ण\appbrk{}वा\appbrk{}त\appbrk{}व\appbrk{}र्ष\appbrk{}प्र\appbrk{}भृ\appbrk{}तिभिः}\ H}}\skt{ }\edtext{\edtext{\skt{समुत्पन्नास्तेऽपि}}{\lemma{\skt{समुत्पन्नास्ते}}\Afootnote{H; \skt{निमित्ताः; ते}\ A; \skt{समुतपन्नास्ते}\ N}}}{\lemma{\skt{ऽपि}}\Afootnote{A; \emph{om.}\ N; \skt{पि}\ H}}\skt{ }\edtext{\skt{द्विविधा}}{\lemma{\skt{द्विविधा}}\Afootnote{A H; \skt{पि}\ N}}\skt{ }\edtext{\skt{व्यापन्ना}}{\lemma{\skt{व्यापन्ना}}\Afootnote{H; \skt{व्यापन्}\ A; \skt{व्यापन्नाश्}\ N}}\skt{ }\edtext{\edtext{\skt{अव्यापन्नाश्च}}{\lemma{\skt{अव्यापन्नाश्}}\Afootnote{\skt{नर्तुकृताः, अव्या\appbrk{}प\appbrk{}न्न\appbrk{}र्तु\appbrk{}कृ\appbrk{}ताश्}\ A; \skt{च~। रिक्व\appbrk{}कृ\appbrk{}र्त\appbrk{}का\appbrk{}ल\appbrk{}कृतः~। अत्र सर्व\appbrk{}व्या\appbrk{}ध्यु\appbrk{}प\appbrk{}रोधः}\ N; \skt{अव्यापन्नाश्च}\ H}}}{\lemma{\skt{च}}\Afootnote{A; \emph{om.}\ N H}}\skt{~॥}\footnoteA{In the 1931 vulgate edition, this passage comes after 1.24.6, before \emph{daivabalapravṛttā …}.}\pend

\pstart \ledsidenote{\passref{1.24.6}}\skt{संघा\-त\-ब\-ल\-प्र\-वृत्ता }\edtext{\skt{नाम य}}{\lemma{\skt{नाम य}}\Afootnote{N H; \skt{य}\ A}}\skt{ }\edtext{\skt{आगन्तव}}{\lemma{\skt{आगन्तव}}\Afootnote{H; \skt{आगन्तवो}\ A; \skt{आगन्तवे}\ N}}\skt{ }\edtext{\skt{आधिभौतिकास्ते}}{\lemma{\skt{आधिभौतिकास्ते}}\Afootnote{N; \emph{om.}\ A; \skt{आधिभौतिकास्त}\ H}}\skt{ }\edtext{\skt{द्विविधाः}}{\lemma{\skt{द्विविधाः}}\Afootnote{N H; \emph{om.}\ A}}\skt{~। दुर्बलस्य }\edtext{\skt{बल\-व\-द्वि\-ग्र\-हा\-च्छ\-स्त्रा\-दि\-कृ\-ताश्च}}{\lemma{\skt{बल\appbrk{}व\appbrk{}द्वि\appbrk{}ग्र\appbrk{}हा\appbrk{}च्छ\appbrk{}स्त्रा\appbrk{}दि\appbrk{}कृ\appbrk{}ताश्च}}\Afootnote{H; \skt{बल\appbrk{}व\appbrk{}द्वि\appbrk{}ग्र\appbrk{}हात्; तेऽपि द्विविधाः- शस्त्रकृता, व्यालकृताश्च~। एते आधिभौतिकाः}\ A; \skt{बल\appbrk{}व\appbrk{}द्वि\appbrk{}ग्र\appbrk{}हा\appbrk{}च्छ\appbrk{}त्रा\appbrk{}दि\appbrk{}कृ\appbrk{}ताश्च}\ N}}\skt{~।}\pend

\pstart \ledsidenote{\passref{1.24.7}}\edtext{\skt{दैव\-ब\-ल\-प्र\-वृत्ता}}{\lemma{\skt{दैव\appbrk{}ब\appbrk{}ल\appbrk{}प्र\appbrk{}वृत्ता}}\Afootnote{N H; \skt{दैव\appbrk{}ब\appbrk{}ल\appbrk{}प्र\appbrk{}वृता ये}\ A; \emph{om.}\ K}}\skt{ }\edtext{\skt{नाम य}}{\lemma{\skt{नाम य}}\Afootnote{N H; \skt{देव\appbrk{}द्रो\appbrk{}हा\appbrk{}द\appbrk{}भि\appbrk{}श\appbrk{}प्तका अथर्वण}\ A; \emph{om.}\ K}}\skt{ }\edtext{\skt{औपसर्गिका}}{\lemma{\skt{औपसर्गिका}}\Afootnote{N; \skt{कृता उपसर्गजाश्च; ते}\ A; \emph{om.}\ K; \skt{औपसर्ग्गा}\ H}}\skt{ }\edtext{\skt{द्विविधा}}{\lemma{\skt{द्विविधा}}\Afootnote{H; \skt{ऽपि द्विविधाः}\ A; \emph{om.}\ K; \skt{द्विवि}\ N}}\skt{ }\edtext{\skt{अभि\-चा\-रा\-भि\-शा\-पा\-भि\-ष\-ङ्गजाः}}{\lemma{\skt{अभि\appbrk{}चा\appbrk{}रा\appbrk{}भि\appbrk{}शा\appbrk{}पा\appbrk{}भि\appbrk{}ष\appbrk{}ङ्गजाः}}\Afootnote{H; \skt{विद्यु\appbrk{}द\appbrk{}श\appbrk{}नि\appbrk{}कृताः, पिशा\appbrk{}चा\appbrk{}दि\appbrk{}कृ\appbrk{}ताश्च; पुनश्च द्विविधाः- संसर्गजा , आकस्मिकाश्च}\ A; \emph{om.}\ K; \skt{धा~। आभि\appbrk{}चा\appbrk{}रा\appbrk{}भि\appbrk{}षा\appbrk{}पा\appbrk{}भि\appbrk{}ष\appbrk{}ङ्गजाः}\ N}}\skt{~॥ }\edtext{\skt{स्वभा\-व\-ब\-ल\-प्र\-वृत्ता}}{\lemma{\skt{स्वभा\appbrk{}व\appbrk{}ब\appbrk{}ल\appbrk{}प्र\appbrk{}वृत्ता}}\Afootnote{A N H; \emph{om.}\ K}}\skt{ }\edtext{\skt{नाम}}{\lemma{\skt{नाम}}\Afootnote{N H; \skt{ये}\ A; \emph{om.}\ K}}\skt{ }\edtext{\skt{क्षुत्पि\-पा\-सा\-ज\-रा\-मृ\-त्यु\-नि\-द्रा\-प्र\-भृ\-तय}}{\lemma{\skt{क्षुत्पि\appbrk{}पा\appbrk{}सा\appbrk{}ज\appbrk{}रा\appbrk{}मृ\appbrk{}त्यु\appbrk{}नि\appbrk{}द्रा\appbrk{}प्र\appbrk{}भृ\appbrk{}तय}}\Afootnote{H; \skt{क्षुत्पि\appbrk{}पा\appbrk{}सा\appbrk{}ज\appbrk{}रा\appbrk{}मृ\appbrk{}त्यु\appbrk{}नि\appbrk{}द्रा\appbrk{}प्र\appbrk{}भृ\appbrk{}तयः}\ A; \emph{om.}\ K; \skt{क्षुत्पि\appbrk{}पा\appbrk{}सा\appbrk{}ज\appbrk{}रा\appbrk{}मृ\appbrk{}त्यु\appbrk{}निद्राः प्रभृतय}\ N}}\skt{ }\edtext{\skt{इति~। }\edtext{\skt{तेऽपि}}{\lemma{\skt{ऽपि}}\Afootnote{A; \emph{om.}\ K; \skt{पि}\ N H}}}{\lemma{\skt{इति~। ते}}\Afootnote{N H; \skt{ते}\ A; \emph{om.}\ K}}\skt{ }\edtext{\skt{द्विविधा}}{\lemma{\skt{द्विविधा}}\Afootnote{N H; \skt{द्विविधाः- का}\ A; \emph{om.}\ K}}\skt{ }\edtext{\skt{रक्षकृता}}{\lemma{\skt{रक्षकृता}}\Afootnote{N; \skt{लजा, अकालजा}\ A; \emph{om.}\ K; \skt{रक्तकृता}\ H}}\skt{ }\edtext{\skt{अरक्षकृताः}}{\lemma{\skt{अरक्षकृताः}}\Afootnote{N; \skt{श्च; तत्र परि}\ A; \emph{om.}\ K; \skt{अरक्तकृताश्}\ H}}\skt{~। }\edtext{\skt{रक्षकृतः}}{\lemma{\skt{रक्षकृतः}}\Afootnote{\skt{रक्षणकृताः}\ A; \emph{om.}\ K; \skt{रक्षकृत}\ N; \skt{च~। रक्तकृतः}\ H}}\skt{ }\edtext{\skt{कालकृतः}}{\lemma{\skt{कालकृतः}}\Afootnote{H; \skt{कालजाः, अ}\ A; \emph{om.}\ K; \skt{कालकृत}\ N}}\skt{~। }\edtext{\edtext{\skt{अर\-क्ष\-कृ\-तोऽ\-का\-ल\-कृतः}}{\lemma{\skt{अरक्षकृतो}}\Afootnote{\skt{परि\appbrk{}र\appbrk{}क्ष\appbrk{}ण\appbrk{}कृता}\ A; \emph{om.}\ K; \skt{अरक्षकृत}\ N; \skt{अरक्तकृतः}\ H}}\skt{~॥ अत्र}}{\lemma{\skt{ऽकालकृतः~॥ अत्र}}\Afootnote{K H; \skt{अकालजाः~। एते आधिदैविकाः~। अत्र}\ A; \skt{अकालकृत~। अत्र}\ N}}\skt{ }\edtext{\skt{सर्वत्र व्याध्युपरोधः}}{\lemma{\skt{सर्वत्र व्याध्युपरोधः}}\Afootnote{N; \skt{सर्व\appbrk{}व्या\appbrk{}ध्य\appbrk{}व\appbrk{}रोधः}\ A; \skt{सर्व\appbrk{}व्या\appbrk{}ध्यु\appbrk{}प\appbrk{}रोधः}\ K H}}\skt{~।}\pend

\pstart \ledsidenote{\passref{1.24.8}}\skt{सर्वेषां च व्याधी\-ना\-म्वा\-त\-पि\-त्त\-श्ले\-ष्माण एव मूलं~। }\edtext{\skt{तल्लि\-ङ्ग\-त्वा\-द्दृ\-ष्ट\-फ\-ल\-त्वा\-दा\-ग\-माच्च}}{\lemma{\skt{तल्लि\appbrk{}ङ्ग\appbrk{}त्वा\appbrk{}द्दृ\appbrk{}ष्ट\appbrk{}फ\appbrk{}ल\appbrk{}त्वा\appbrk{}दा\appbrk{}ग\appbrk{}माच्च}}\Afootnote{A K H; \skt{तल्लि\appbrk{}ङ्ग\appbrk{}त्वा\appbrk{}द्दृ\appbrk{}ष्ट\appbrk{}फ\appbrk{}ल\appbrk{}त्वा\appbrk{}दा\appbrk{}ग\appbrk{}मच्च}\ N}}\skt{ }\edtext{\skt{पश्यामः}}{\lemma{\skt{पश्यामः}}\Afootnote{K N H; \emph{om.}\ A}}\skt{~। यथा हि }\edtext{\skt{कृत्स्नं}}{\lemma{\skt{कृत्स्नं}}\Afootnote{A K H; \skt{कृत्स्न}\ N}}\skt{ }\edtext{\skt{विकारजातं}}{\lemma{\skt{विकारजातं}}\Afootnote{A K H; \skt{विका\appbrk{}र\appbrk{}जा\appbrk{}त\appbrk{}जा\appbrk{}तम्}\ N}}\skt{ }\edtext{\skt{वैश्वरूप्येण}}{\lemma{\skt{वैश्वरूप्येण}}\Afootnote{K; \skt{विश्वरूपेणा}\ A; \skt{वैस्वरूपेण}\ N; \skt{वैस्वरूप्येण}\ H}}\skt{ }\edtext{\skt{व्यवस्थितं}}{\lemma{\skt{व्यवस्थितं}}\Afootnote{K N H; \skt{वस्थितं}\ A}}\skt{~। }\edtext{\edtext{\skt{सत्व\-र\-ज\-स्त\-मां\-स्य\-व्य\-ति\-रिच्य}}{\lemma{\skt{सत्व\appbrk{}र\appbrk{}ज\appbrk{}स्त\appbrk{}मांस्य्}}\Afootnote{A K H; \skt{सत्व\appbrk{}र\appbrk{}ज\appbrk{}त\appbrk{}मांसि}\ N}}}{\lemma{\skt{अव्यतिरिच्य}}\Afootnote{K H; \skt{न व्यतिरिच्यन्}\ A; \skt{अव्यरिच्य}\ N}}\skt{ }\edtext{\skt{वर्त्तन्ते}}{\lemma{\skt{वर्त्तन्ते}}\Afootnote{K\textsuperscript{pc} N H; \skt{ते}\ A; \skt{वर्त्तते}\ K\textsuperscript{ac}}}\skt{~। एवमेव कृत्स्नं विकारजातं }\edtext{\skt{वैश्व\-रू\-प्ये\-णा\-व\-स्थितं}}{\lemma{\skt{वैश्व\appbrk{}रू\appbrk{}प्ये\appbrk{}णा\appbrk{}व\appbrk{}स्थितं}}\Afootnote{K H; \skt{विश्व\appbrk{}रू\appbrk{}पे\appbrk{}णा\appbrk{}व\appbrk{}स्थित}\ A; \skt{वैस्व\appbrk{}रू\appbrk{}पे\appbrk{}णा\appbrk{}व\appbrk{}स्थितं}\ N}}\skt{~। }\edtext{\skt{अव्यतिरिच्य}}{\lemma{\skt{अव्यतिरिच्य}}\Afootnote{K H; \skt{मव्यतिरिच्य}\ A N}}\skt{ वात\-पि\-त्त\-श्ले\-ष्माणो वर्त्तन्ते~। }\edtext{\edtext{\edtext{\skt{दोष\-धा\-तु\-म\-ल\-सं\-स\-र्गा\-दा\-य\-त\-न\-वि\-शे\-षा\-न्नि\-मि\-त्त\-त\-श्चैषां}}{\lemma{\skt{दोष\appbrk{}धा\appbrk{}तु\appbrk{}म\appbrk{}ल\appbrk{}सं\appbrk{}स\appbrk{}र्गा\appbrk{}दा\appbrk{}य\appbrk{}त\appbrk{}न\appbrk{}वि\appbrk{}शे\appbrk{}षान्}}\Afootnote{K\textsuperscript{pc} N H\textsuperscript{pc}; \skt{दोष\appbrk{}धा\appbrk{}तु\appbrk{}म\appbrk{}ल\appbrk{}सं\appbrk{}स\appbrk{}र्गा\appbrk{}दा\appbrk{}य\appbrk{}त\appbrk{}न\appbrk{}वि\appbrk{}शे\appbrk{}षान्}\ A; \skt{दोष\appbrk{}धा\appbrk{}तु\appbrk{}म\appbrk{}ल\appbrk{}सं\appbrk{}स\appbrk{}र्गा\appbrk{}य\appbrk{}दा\appbrk{}य\appbrk{}त\appbrk{}न\appbrk{}वि\appbrk{}शे\appbrk{}षाःन्}\ K\textsuperscript{ac}; \skt{दोष\appbrk{}धा\appbrk{}तु\appbrk{}म\appbrk{}ल\appbrk{}सं\appbrk{}स\appbrk{}र्ग्गा\appbrk{}दा\appbrk{}य\appbrk{}न\appbrk{}वि\appbrk{}शे\appbrk{}षान्}\ H\textsuperscript{ac}}}}{\lemma{\skt{निमित्ततश्}}\Afootnote{N H; \skt{निमित्तत}\ A; \skt{निमित्तत…श्}\ K}}}{\lemma{\skt{चैषां}}\Afootnote{K N H; \skt{श्चैषां}\ A}}\skt{ विकल्पो }\edtext{\skt{भवति}}{\lemma{\skt{भवति}}\Afootnote{K N H; \emph{om.}\ A}}\skt{~। }\edtext{\skt{दोष\-दू\-षि\-ते\-ष्व\-त्य\-र्थञ्च}}{\lemma{\skt{अत्यर्थञ्च}}\Afootnote{K N H; \skt{अत्यर्थं}\ A}}\skt{ धातुषु }\edtext{\skt{संज्ञा}}{\lemma{\skt{संज्ञा}}\Afootnote{A K\textsuperscript{pc} N; \skt{जि ज्ञा}\ K\textsuperscript{ac} H}}\skt{ }\edtext{\skt{भवति}}{\lemma{\skt{भवति}}\Afootnote{K\textsuperscript{pc}; \emph{om.}\ A; \skt{सा भवति}\ K\textsuperscript{ac} H; \skt{भवति~॥ …}\ N}}\skt{~॥ }\edtext{\edtext{\skt{रसजोऽयं}}{\lemma{\skt{रसजो}}\Afootnote{A K\textsuperscript{pc} N H; \emph{om.}\ K\textsuperscript{ac}}}\skt{ }\edtext{\skt{रक्तजोऽयं}}{\lemma{\skt{ऽयं}}\Afootnote{A; \skt{यं}\ K N H}}}{\lemma{\skt{ऽयं रक्तजो}}\Afootnote{\skt{ऽयं, शोणितजो}\ A; \skt{सो जो यं रक्तजो}\ K\textsuperscript{ac}; \skt{यं रक्तजो}\ K\textsuperscript{pc} N H}}\skt{ }\edtext{\edtext{\skt{मांसजोऽयं}}{\lemma{\skt{मांसजो}}\Afootnote{A K H; \skt{माशजो}\ N}}}{\lemma{\skt{ऽयं}}\Afootnote{A; \skt{यं}\ K N H}}\skt{ }\edtext{\edtext{\skt{मेदोजोऽयं}}{\lemma{\skt{मेदोजो}}\Afootnote{A K H; \skt{मेदजो}\ N}}}{\lemma{\skt{ऽयं}}\Afootnote{A; \skt{यं}\ K N H}}\skt{ }\edtext{\skt{अस्थिजोऽयं}}{\lemma{\skt{ऽयं}}\Afootnote{A; \skt{यं}\ K N H}}\skt{ }\edtext{\edtext{\skt{मज्जजोऽयं}}{\lemma{\skt{मज्जजो}}\Afootnote{A K N; \skt{मज्जाजो}\ H}}}{\lemma{\skt{ऽयं}}\Afootnote{A; \skt{यं}\ K N H}}\skt{ }\edtext{\skt{शुक्रजोऽयं}}{\lemma{\skt{ऽयं}}\Afootnote{A; \skt{यं}\ K N H}}\skt{ व्याधिरिति~॥}\pend

\pstart \ledsidenote{\passref{1.24.9}}\edtext{\skt{अश्र\-द्धा\-रो\-च\-क\-पि\-पा\-सा\-ङ्ग\-म\-र्द\-ज्व\-र\-हृ\-ल्ला\-सा\-तृ\-प्ति\-गौ\-र\-व\-पा\-ण्डु\-रो\-ग\-स्रो\-तो\-रो\-ध\-का\-र्श्य\-वै\-र\-स्या\-ङ्ग\-सादाः}}{\lemma{\skt{अश्र\appbrk{}द्धा\appbrk{}रो\appbrk{}च\appbrk{}क\appbrk{}पि\appbrk{}पा\appbrk{}सा\appbrk{}ङ्ग\appbrk{}म\appbrk{}र्द\appbrk{}ज्व\appbrk{}र\appbrk{}हृ\appbrk{}ल्ला\appbrk{}सा\appbrk{}तृ\appbrk{}प्ति\appbrk{}गौ\appbrk{}र\appbrk{}व\appbrk{}पा\appbrk{}ण्डु\appbrk{}रो\appbrk{}ग\appbrk{}स्रो\appbrk{}तो\appbrk{}रो\appbrk{}ध\appbrk{}का\appbrk{}र्श्य\appbrk{}वै\appbrk{}र\appbrk{}स्या\appbrk{}ङ्ग\appbrk{}सादाः}}\Afootnote{H; \skt{तत्र, अन्ना\appbrk{}श्र\appbrk{}द्धा\appbrk{}रो\appbrk{}च\appbrk{}का\appbrk{}वि\appbrk{}पा\appbrk{}का\appbrk{}ङ्ग\appbrk{}म\appbrk{}र्द\appbrk{}ज्व\appbrk{}र\appbrk{}हृ\appbrk{}ल्ला\appbrk{}स\appbrk{}तृ\appbrk{}प्ति\appbrk{}गौ\appbrk{}र\appbrk{}व\appbrk{}हृ\appbrk{}त्पा\appbrk{}ण्डु\appbrk{}रो\appbrk{}ग\appbrk{}मा\appbrk{}र्गो\appbrk{}प\appbrk{}रो\appbrk{}ध\appbrk{}का\appbrk{}र्श्य\appbrk{}वै\appbrk{}र\appbrk{}स्या\appbrk{}ङ्ग\appbrk{}सादा}\ A; \skt{अश्र\appbrk{}द्धा\appbrk{}रो\appbrk{}च\appbrk{}क\appbrk{}पि\appbrk{}पा\appbrk{}सा\appbrk{}ङ्ग\appbrk{}म\appbrk{}र्द\appbrk{}ज्व\appbrk{}र\appbrk{}हृ\appbrk{}ल्लास तृप्ति\appbrk{}गौ\appbrk{}र\appbrk{}व\appbrk{}पा\appbrk{}ण्डु\appbrk{}रो\appbrk{}ग\appbrk{}स्रो\appbrk{}तो\appbrk{}रो\appbrk{}ध\appbrk{}का\appbrk{}र्श्य\appbrk{}वै\appbrk{}र\appbrk{}स्या\appbrk{}ङ्ग\appbrk{}सादाः}\ K\textsuperscript{ac}; \skt{अश्र\appbrk{}द्धा\appbrk{}रो\appbrk{}च\appbrk{}क\appbrk{}पि\appbrk{}पा\appbrk{}सा\appbrk{}ङ्ग\appbrk{}म\appbrk{}र्द\appbrk{}ज्व\appbrk{}र\appbrk{}हृ\appbrk{}ल्लासा तृप्ति\appbrk{}गौ\appbrk{}र\appbrk{}व\appbrk{}पा\appbrk{}ण्डु\appbrk{}रो\appbrk{}ग\appbrk{}स्रो\appbrk{}तो\appbrk{}रो\appbrk{}ध\appbrk{}का\appbrk{}र्श्य\appbrk{}वै\appbrk{}र\appbrk{}स्या\appbrk{}ङ्ग\appbrk{}सादाः}\ K\textsuperscript{pc}; \skt{अश्र\appbrk{}द्वा\appbrk{}रो\appbrk{}च\appbrk{}का\appbrk{}पि\appbrk{}पा\appbrk{}सा\appbrk{}ङ्ग\appbrk{}म\appbrk{}र्द्द\appbrk{}ज्व\appbrk{}र\appbrk{}हृ\appbrk{}ल्ला\appbrk{}स\appbrk{}तृ\appbrk{}प्ति\appbrk{}गौ\appbrk{}र\appbrk{}व\appbrk{}पा\appbrk{}ण्डु\appbrk{}रो\appbrk{}ग\appbrk{}श्रा\appbrk{}तो\appbrk{}रो\appbrk{}धो\appbrk{}का\appbrk{}र्स्य\appbrk{}वै\appbrk{}र\appbrk{}स्या\appbrk{}ङ्ग\appbrk{}स्पदः}\ N}}\skt{~। }\edtext{\skt{अका\-ल\-प\-लि\-त\-ति\-मि\-र\-द\-र्श\-न\-र\-स\-दो\-षजा}}{\lemma{\skt{अका\appbrk{}ल\appbrk{}प\appbrk{}लि\appbrk{}त\appbrk{}ति\appbrk{}मि\appbrk{}र\appbrk{}द\appbrk{}र्श\appbrk{}न\appbrk{}र\appbrk{}स\appbrk{}दो\appbrk{}षजा}}\Afootnote{K N H; \skt{काल\appbrk{}ज\appbrk{}व\appbrk{}ली\appbrk{}प\appbrk{}लि\appbrk{}त\appbrk{}द\appbrk{}र्श\appbrk{}न\appbrk{}प्र\appbrk{}भृ\appbrk{}तयो रसदोषजा}\ A}}\skt{ }\edtext{\skt{विकाराः}}{\lemma{\skt{विकाराः}}\Afootnote{K N H; \skt{विकाराः;SS}\ A}}\skt{~॥ }\edtext{\edtext{\skt{कुष्ठ\-वि\-स\-र्प\-पि\-ट\-का\-स्ति\-ल\-का\-ल\-क\-न\-च्छ\-व्य\-ङ्ग\-म\-स\-क\-नी\-लि\-का\-को\-ठ\-प्ली\-ह\-गु\-ल्म\-वि\-द्र\-ध्य\-र्शोऽ\-र्बु\-दा\-सृ\-ग्द\-र\-र\-क्त\-पि\-त्त\-प्र\-भृ\-तयो}}{\lemma{\skt{कुष्ठ\appbrk{}वि\appbrk{}स\appbrk{}र्प\appbrk{}पि\appbrk{}ट\appbrk{}कास्}}\Afootnote{K H; \skt{कुष्ठ\appbrk{}वि\appbrk{}स\appbrk{}र्प\appbrk{}पि\appbrk{}ड\appbrk{}का\appbrk{}म\appbrk{}श\appbrk{}क\appbrk{}नी\appbrk{}लिका}\ A; \skt{कुष्ठ\appbrk{}वि\appbrk{}स\appbrk{}र्प्प\appbrk{}वि\appbrk{}ढ\appbrk{}कास्}\ N}}}{\lemma{\skt{तिल\appbrk{}का\appbrk{}ल\appbrk{}क\appbrk{}न\appbrk{}च्छ\appbrk{}व्य\appbrk{}ङ्ग\appbrk{}म\appbrk{}स\appbrk{}क\appbrk{}नी\appbrk{}लि\appbrk{}का\appbrk{}को\appbrk{}ठ\appbrk{}प्ली\appbrk{}ह\appbrk{}गु\appbrk{}ल्म\appbrk{}वि\appbrk{}द्र\appbrk{}ध्य\appbrk{}र्शोऽ\appbrk{}र्बु\appbrk{}दा\appbrk{}सृ\appbrk{}ग्द\appbrk{}र\appbrk{}र\appbrk{}क्त\appbrk{}पि\appbrk{}त्त\appbrk{}प्र\appbrk{}भृ\appbrk{}तयो}}\Afootnote{\skt{तिल\appbrk{}का\appbrk{}ल\appbrk{}क\appbrk{}न्य\appbrk{}च्छ\appbrk{}व्य\appbrk{}ङ्गे\appbrk{}न्द्र\appbrk{}लु\appbrk{}प्त\appbrk{}प्ली\appbrk{}ह\appbrk{}वि\appbrk{}द्र\appbrk{}धि\appbrk{}गु\appbrk{}ल्म\appbrk{}वा\appbrk{}त\appbrk{}शो\appbrk{}णि\appbrk{}ता\appbrk{}र्शोऽ\appbrk{}र्बु\appbrk{}दा\appbrk{}ङ्ग\appbrk{}म\appbrk{}र्दा\appbrk{}सृ\appbrk{}ग्द\appbrk{}र\appbrk{}र\appbrk{}क्त\appbrk{}पि\appbrk{}त्त\appbrk{}प्र\appbrk{}भृ\appbrk{}तयो}\ A; \skt{तिल\appbrk{}का\appbrk{}ल\appbrk{}क\appbrk{}न\appbrk{}च्छ\appbrk{}व्य\appbrk{}ङ्ग\appbrk{}म\appbrk{}स\appbrk{}क\appbrk{}नी\appbrk{}लि\appbrk{}का\appbrk{}को\appbrk{}ठ\appbrk{}प्ली\appbrk{}ह\appbrk{}गु\appbrk{}ल्म\appbrk{}वि\appbrk{}द्र\appbrk{}ध्य\appbrk{}र्सो\appbrk{}र्बु\appbrk{}दा\appbrk{}सृ\appbrk{}ग्द\appbrk{}र\appbrk{}र\appbrk{}क्त\appbrk{}पि\appbrk{}त्त\appbrk{}प्र\appbrk{}भृ\appbrk{}तयो}\ K; \skt{तिल\appbrk{}का\appbrk{}ल\appbrk{}क\appbrk{}न\appbrk{}च्छ\appbrk{}व्य\appbrk{}ङ्श\appbrk{}म\appbrk{}स\appbrk{}क\appbrk{}नी\appbrk{}लि\appbrk{}का\appbrk{}को\appbrk{}ठ\appbrk{}प्ली\appbrk{}ह\appbrk{}गु\appbrk{}ल्म\appbrk{}वि\appbrk{}द्र\appbrk{}ध्य\appbrk{}र्श्शा\appbrk{}र्बु\appbrk{}दा\appbrk{}सू\appbrk{}गू\appbrk{}र\appbrk{}र\appbrk{}क्त\appbrk{}पि\appbrk{}त्त\appbrk{}प्र\appbrk{}भृ\appbrk{}तयो}\ N; \skt{तिल\appbrk{}का\appbrk{}ल\appbrk{}क\appbrk{}न\appbrk{}च्छ\appbrk{}व्य\appbrk{}ङ्ग\appbrk{}म\appbrk{}स\appbrk{}क\appbrk{}नी\appbrk{}लि\appbrk{}का\appbrk{}को\appbrk{}ठ\appbrk{}प्ली\appbrk{}ह\appbrk{}गु\appbrk{}ल्म\appbrk{}वि\appbrk{}द्र\appbrk{}ध्य\appbrk{}र्शोऽ\appbrk{}र्ब्बु\appbrk{}दा\appbrk{}सृ\appbrk{}ग्द\appbrk{}र\appbrk{}र\appbrk{}क्त\appbrk{}पित्ति प्रभृतयो}\ H\textsuperscript{ac}; \skt{तिल\appbrk{}का\appbrk{}ल\appbrk{}क\appbrk{}न\appbrk{}च्छ\appbrk{}व्य\appbrk{}ङ्ग\appbrk{}म\appbrk{}स\appbrk{}क\appbrk{}नी\appbrk{}लि\appbrk{}का\appbrk{}को\appbrk{}ठ\appbrk{}प्ली\appbrk{}ह\appbrk{}गु\appbrk{}ल्म\appbrk{}वि\appbrk{}द्र\appbrk{}ध्य\appbrk{}र्शोऽ\appbrk{}र्ब्बु\appbrk{}दा\appbrk{}सृ\appbrk{}ग्द\appbrk{}र\appbrk{}र\appbrk{}क्त\appbrk{}पित्त प्रभृतयो}\ H\textsuperscript{pc}}}\skt{ }\edtext{\skt{रक्तदोषात्~॥ }\edtext{\skt{गुद\-मु\-ख\-मे\-ढ्र\-पा\-का\-श्चा\-धि\-मां\-शा\-र्ब्बु\-दा\-र्सो\-प\-जि\-ह्वो\-प\-कु\-श\-ग\-ल\-शु\-ण्डि\-का\-मां\-श\-सं\-घा\-तो\-ष्ठ\-प्र\-को\-प\-ग\-ल\-ग\-ण्ड\-ग\-ण्ड\-मा\-ला\-प्र\-भृ\-तयो}}{\lemma{\skt{चाधि\appbrk{}मां\appbrk{}शा\appbrk{}र्ब्बु\appbrk{}दा\appbrk{}र्सो\appbrk{}प\appbrk{}जि\appbrk{}ह्वो\appbrk{}प\appbrk{}कु\appbrk{}श\appbrk{}ग\appbrk{}ल\appbrk{}शु\appbrk{}ण्डि\appbrk{}का\appbrk{}मां\appbrk{}श\appbrk{}सं\appbrk{}घा\appbrk{}तो\appbrk{}ष्ठ\appbrk{}प्र\appbrk{}को\appbrk{}प\appbrk{}ग\appbrk{}ल\appbrk{}ग\appbrk{}ण्ड\appbrk{}ग\appbrk{}ण्ड\appbrk{}मा\appbrk{}ला\appbrk{}प्र\appbrk{}भृ\appbrk{}तयो}}\Afootnote{K N; \skt{श्च; अधि\appbrk{}मां\appbrk{}सा\appbrk{}र्बु\appbrk{}दा\appbrk{}र्शोऽ\appbrk{}धि\appbrk{}जि\appbrk{}ह्वो\appbrk{}प\appbrk{}जि\appbrk{}ह्वो\appbrk{}प\appbrk{}कु\appbrk{}श\appbrk{}ग\appbrk{}ल\appbrk{}शु\appbrk{}ण्डि\appbrk{}का\appbrk{}ल\appbrk{}जी\appbrk{}मां\appbrk{}स\appbrk{}स\appbrk{}ङ्घा\appbrk{}तौ\appbrk{}ष्ठ\appbrk{}प्र\appbrk{}को\appbrk{}प\appbrk{}ग\appbrk{}ल\appbrk{}ग\appbrk{}ण्ड\appbrk{}ग\appbrk{}ण्ड\appbrk{}मा\appbrk{}ला\appbrk{}प्र\appbrk{}भृ\appbrk{}तयो}\ A; \skt{चाधि\appbrk{}मां\appbrk{}सा\appbrk{}र्ब्बु\appbrk{}दा\appbrk{}र्शोऽ\appbrk{}प\appbrk{}जि\appbrk{}ह्वा\appbrk{}प\appbrk{}कु\appbrk{}श\appbrk{}ग\appbrk{}ल\appbrk{}शु\appbrk{}ण्डि\appbrk{}का\appbrk{}मां\appbrk{}स\appbrk{}सं\appbrk{}घा\appbrk{}तो\appbrk{}ष्ठ\appbrk{}प्र\appbrk{}को\appbrk{}प\appbrk{}ग\appbrk{}ल\appbrk{}ग\appbrk{}ण्ड\appbrk{}ग\appbrk{}ण्ड\appbrk{}मा\appbrk{}ला\appbrk{}प्र\appbrk{}भृ\appbrk{}तयो}\ H}}}{\lemma{\skt{रक्तदोषात्~॥ गुद\appbrk{}मु\appbrk{}ख\appbrk{}मे\appbrk{}ढ्र\appbrk{}पा\appbrk{}काश्}}\Afootnote{K H; \skt{रक्तदोषजाः , गुद\appbrk{}मु\appbrk{}ख\appbrk{}मे\appbrk{}ढ्र\appbrk{}पाका}\ A; \skt{रक्तदोषा~॥ गुद\appbrk{}मु\appbrk{}ख\appbrk{}मे\appbrk{}ढ्र\appbrk{}पा\appbrk{}काश्}\ N}}\skt{ }\edtext{\skt{मांसदोषात्~॥ ग्रन्थि\-वृ\-द्धि\-ग\-ल\-ग\-ण्डा\-र्ब्बु\-दो\-ष्ठ\-प्र\-को\-प\-म\-धु\-मे\-हा\-ति\-स्थौ\-ल्य\-प्र\-भृ\-तयो}}{\lemma{\skt{मांसदोषात्~॥ ग्रन्थि\appbrk{}वृ\appbrk{}द्धि\appbrk{}ग\appbrk{}ल\appbrk{}ग\appbrk{}ण्डा\appbrk{}र्ब्बु\appbrk{}दो\appbrk{}ष्ठ\appbrk{}प्र\appbrk{}को\appbrk{}प\appbrk{}म\appbrk{}धु\appbrk{}मे\appbrk{}हा\appbrk{}ति\appbrk{}स्थौ\appbrk{}ल्य\appbrk{}प्र\appbrk{}भृ\appbrk{}तयो}}\Afootnote{K N H; \skt{मांसदोषजाः; ग्रन्थि\appbrk{}वृ\appbrk{}द्धि\appbrk{}ग\appbrk{}ल\appbrk{}ग\appbrk{}ण्डा\appbrk{}र्बु\appbrk{}द\appbrk{}मे\appbrk{}दो\appbrk{}जौ\appbrk{}ष्ठ\appbrk{}प्र\appbrk{}को\appbrk{}प\appbrk{}म\appbrk{}धु\appbrk{}मे\appbrk{}हा\appbrk{}ति\appbrk{}स्थौ\appbrk{}ल्या\appbrk{}ति\appbrk{}स्वे\appbrk{}द\appbrk{}प्र\appbrk{}भृ\appbrk{}तयो}\ A}}\skt{ }\edtext{\skt{मेदोदोषात्~॥ अध्य\-स्थि\-द\-न्ता\-स्थि\-तो\-द\-शू\-ला\-द\-योऽ\-स्थि\-दो\-षात्}}{\lemma{\skt{मेदोदोषात्~॥ अध्य\appbrk{}स्थि\appbrk{}द\appbrk{}न्ता\appbrk{}स्थि\appbrk{}तो\appbrk{}द\appbrk{}शू\appbrk{}ला\appbrk{}दयो}}\Afootnote{H; \skt{मेदोदोषजाः; अध्य\appbrk{}स्थ्य\appbrk{}धि\appbrk{}द\appbrk{}न्ता\appbrk{}स्थि\appbrk{}तो\appbrk{}द\appbrk{}शू\appbrk{}ल\appbrk{}कु\appbrk{}न\appbrk{}ख\appbrk{}प्र\appbrk{}भृ\appbrk{}तयो}\ A; \skt{मेदोदोषात्~॥ अध्य\appbrk{}स्थि\appbrk{}द\appbrk{}न्ता\appbrk{}स्थि\appbrk{}तो\appbrk{}द\appbrk{}शू\appbrk{}ला\appbrk{}दयो}\ K; \skt{मेददोषात्~॥ अध्य\appbrk{}स्थि\appbrk{}द\appbrk{}न्ता\appbrk{}स्थि\appbrk{}तो\appbrk{}द\appbrk{}शू\appbrk{}ला\appbrk{}दयो}\ N}}\skt{~॥ तमो\-द\-र्श\-मू\-र्च्छा\-भ्र\-म\-पा\-र्श्व\-गौ\-र\-व\-हृ\-च्छू\-ल\-स्थू\-ल\-मू\-लो\-रु\-ज\-म्भेति }\edtext{\skt{मज्जदोषात्~॥ }\edtext{\edtext{\skt{क्लैब्य\-म\-प्र\-ह\-र्षश्च}}{\lemma{\skt{अप्रहर्षश्}}\Afootnote{K N H; \skt{ब्याप्रहर्ष}\ A}}}{\lemma{\skt{च}}\Afootnote{K H; \skt{शुक्रा\appbrk{}श्म\appbrk{}री\appbrk{}शु\appbrk{}क्र\appbrk{}मे\appbrk{}ह\appbrk{}शु\appbrk{}क्र\appbrk{}दो\appbrk{}षा\appbrk{}दय}\ A; \skt{च शुक्रदोषा}\ N}}}{\lemma{\skt{मज्जदोषात्~॥ क्लैब्यम्}}\Afootnote{K N H; \skt{मज्जदोषजाः; क्लै}\ A}}\skt{ }\edtext{\skt{शुक्रदोषात्~॥ त्वग्दोषः}}{\lemma{\skt{शुक्रदोषात्~॥ त्वग्दोषः}}\Afootnote{K H; \skt{श्च तद्दोषाः; त्वग्दोषाः}\ A; \skt{त्~॥ त्वग्दोषस्}\ N}}\skt{ }\edtext{\edtext{\skt{सङ्गोऽ\-ति\-प्र\-वृ\-र्त्तिर्वा}}{\lemma{\skt{सङ्गो}}\Afootnote{A K H; \skt{सङ्गा}\ N}}}{\lemma{\skt{ऽतिप्रवृर्त्तिर्वा}}\Afootnote{\skt{ऽतिप्रवृत्तिर}\ A; \skt{तिप्रवृत्तिर्वा}\ K H; \skt{तिप्रवृर्त्तिर्वा}\ N}}\skt{ }\edtext{\skt{मलानां}}{\lemma{\skt{मलानां}}\Afootnote{K N H; \skt{यथाप्रवृत्तिर्वा}\ A}}\skt{ }\edtext{\skt{मला\-य\-त\-न\-दो\-षात्~॥ इन्द्रियाणां}}{\lemma{\skt{मला\appbrk{}य\appbrk{}त\appbrk{}न\appbrk{}दो\appbrk{}षात्~॥ इन्द्रियाणां}}\Afootnote{K N H\textsuperscript{pc}; \skt{मला\appbrk{}य\appbrk{}त\appbrk{}न\appbrk{}दोषाः; इन्द्रियाणा}\ A; \skt{मला\appbrk{}य\appbrk{}त\appbrk{}दो\appbrk{}षात्~॥ इन्द्रियाणाम्}\ H\textsuperscript{ac}}}\skt{ }\edtext{\skt{अयथा}}{\lemma{\skt{अयथा}}\Afootnote{K H; \skt{म}\ A; \skt{मयः}\ N}}\skt{ }\edtext{\edtext{\skt{प्रवृ\-त्ति\-र\-प्र\-वृ\-त्तिर्वा}}{\lemma{\skt{प्रवृत्तिर्}}\Afootnote{K N H; \skt{प्रवृत्ति}\ A}}}{\lemma{\skt{अप्रवृत्तिर्वा}}\Afootnote{K N H; \skt{रयथाप्रवृत्ति}\ A}}\skt{ }\edtext{\skt{इन्द्रि\-या\-य\-त\-न\-दो\-षा\-दि\-त्येष}}{\lemma{\skt{इन्द्रि\appbrk{}या\appbrk{}य\appbrk{}त\appbrk{}न\appbrk{}दो\appbrk{}षा\appbrk{}दित्य्}}\Afootnote{K N H; \skt{र्वेन्द्रि\appbrk{}या\appbrk{}य\appbrk{}त\appbrk{}न\appbrk{}दोषाः; इत्य्}\ A}}\skt{ }\edtext{\skt{समासो}}{\lemma{\skt{समासो}}\Afootnote{K N H; \skt{समास}\ A}}\skt{ }\edtext{\skt{विस्तरं}}{\lemma{\skt{विस्तरं}}\Afootnote{K; \skt{उक्तः; विस्तरं}\ A; \skt{विस्तरे}\ N; \skt{विस्तर}\ H}}\skt{ निमित्तानि }\edtext{\edtext{\edtext{\skt{चैषा\-म्प्र\-ति\-रो\-ग\-म्व\-क्ष्यामः}}{\lemma{\skt{चैषाम्}}\Afootnote{A K H; \skt{चैषा}\ N}}}{\lemma{\skt{प्रतिरोगम्}}\Afootnote{A K H; \skt{प्रतिर्योगं}\ N}}}{\lemma{\skt{वक्ष्यामः}}\Afootnote{A K H; \skt{न्वक्ष्यामः}\ N}}\skt{~॥}\edtext{}{\lemma{\skt{ऽस्थिदोषात्~॥ तमो\appbrk{}द\appbrk{}र्श\appbrk{}मू\appbrk{}र्च्छा\appbrk{}भ्र\appbrk{}म\appbrk{}पा\appbrk{}र्श्व\appbrk{}गौ\appbrk{}र\appbrk{}व\appbrk{}हृ\appbrk{}च्छू\appbrk{}ल\appbrk{}स्थू\appbrk{}ल\appbrk{}मू\appbrk{}लो\appbrk{}रु\appbrk{}ज\appbrk{}म्भेति}}\Afootnote{H; \skt{ऽस्थिदोषजाः; तमो\appbrk{}द\appbrk{}र्श\appbrk{}न\appbrk{}मू\appbrk{}र्च्छा\appbrk{}भ्र\appbrk{}म\appbrk{}प\appbrk{}र्व\appbrk{}स्थू\appbrk{}ल\appbrk{}मू\appbrk{}ला\appbrk{}रु\appbrk{}र्ज\appbrk{}न्म\appbrk{}ने\appbrk{}त्रा\appbrk{}भि\appbrk{}ष्य\appbrk{}न्द\appbrk{}प्र\appbrk{}भृ\appbrk{}तयो}\ A; \skt{स्थिदोषात्~॥ तमो\appbrk{}द\appbrk{}र्श\appbrk{}मू\appbrk{}र्च्छा\appbrk{}भ्र\appbrk{}म\appbrk{}पा\appbrk{}र्श्व\appbrk{}गौ\appbrk{}र\appbrk{}व\appbrk{}हृ\appbrk{}च्छू\appbrk{}ल\appbrk{}स्थू\appbrk{}ल\appbrk{}मू\appbrk{}लो\appbrk{}रुं\appbrk{}ज\appbrk{}म्भेति}\ K; \skt{स्थिदोषजात्~॥ तमो\appbrk{}द\appbrk{}र्श\appbrk{}न\appbrk{}मू\appbrk{}र्च्छा\appbrk{}भ्र\appbrk{}मो\appbrk{}पा\appbrk{}र्श्व\appbrk{}गौ\appbrk{}र\appbrk{}व\appbrk{}हृ\appbrk{}च्छू\appbrk{}ल\appbrk{}स्थू\appbrk{}ल\appbrk{}सू\appbrk{}ल\appbrk{}रु\appbrk{}जाश्ति}\ N}}\pend

\pstart\versepar \ledsidenote{\passref{1.24.10}}\skt{भवति चात्र~॥ }\pend
\pstart\versepar \skt{कुपितानां हि }\edtext{\skt{दोषाणां}}{\lemma{\skt{दोषाणां}}\Afootnote{A K H; \skt{दोषाणा}\ N}}\skt{ }\edtext{\skt{शरीरे}}{\lemma{\skt{शरीरे}}\Afootnote{A K\textsuperscript{pc} H; \skt{रीरे}\ K\textsuperscript{ac}; \skt{सरीरो}\ N}}\skt{ }\edtext{\skt{परिधावतां}}{\lemma{\skt{परिधावतां}}\Afootnote{A K N H\textsuperscript{pc}; \skt{परिधावत्म्}\ H\textsuperscript{ac}}}\skt{~। }\pend
\pstart\versepar \skt{यत्र }\edtext{\skt{सङ्गः}}{\lemma{\skt{सङ्गः}}\Afootnote{A K N; \skt{सङ्गाः}\ H}}\skt{ }\edtext{\skt{सवै\-गु\-ण्या\-द्व्या\-धि\-स्त\-त्रो\-प\-जा\-यते}}{\lemma{\skt{सवै\appbrk{}गु\appbrk{}ण्या\appbrk{}द्व्या\appbrk{}धिस्}}\Afootnote{K N H; \skt{खवै\appbrk{}गु\appbrk{}ण्या\appbrk{}द्व्या\appbrk{}धिस्}\ A}}\footnoteA{Ḍalhaṇa pointed out that others read \emph{svavaiguṇyād} (instead of \emph{khavaiguṇyād}).}\skt{~॥}\edtext{}{\lemma{\skt{भवति चात्र~॥ कुपितानां हि}}\Afootnote{A; \skt{भवति चात्र~॥ कुपितानां हि}\ K H; \skt{भवति चात्र~॥ कुपितानां हि हि}\ N}}\pend

\pstart \ledsidenote{\passref{1.24.11}}\edtext{\skt{भूयोऽत्र}}{\lemma{\skt{ऽत्र}}\Afootnote{A K N H\textsuperscript{pc}; \skt{ऽ}\ H\textsuperscript{ac}}}\skt{ }\edtext{\skt{जिज्ञासते~। }\edtext{\skt{किम्वा\-त\-पि\-त्त\-श्ले\-ष्माणां}}{\lemma{\skt{वात\appbrk{}पि\appbrk{}त्त\appbrk{}श्ले\appbrk{}ष्माणां}}\Afootnote{N H; \skt{वातादीनां}\ A; \skt{वात\appbrk{}पि\appbrk{}त्त\appbrk{}श्ले\appbrk{}ष्मणां}\ K}}}{\lemma{\skt{जिज्ञासते~। किम्}}\Afootnote{N; \skt{जिज्ञास्यं, किं}\ A; \skt{जिज्ञास्यते~। किम्}\ K H}}\skt{ }\edtext{\skt{ज्वरा\-ति\-सा\-रा\-दीनां च}}{\lemma{\skt{ज्वरा\appbrk{}ति\appbrk{}सा\appbrk{}रा\appbrk{}दीनां च}}\Afootnote{N; \skt{ज्वरादीनां च}\ A; \skt{ज्वरा\appbrk{}ती\appbrk{}सा\appbrk{}रा\appbrk{}दी\appbrk{}नाञ्च}\ K H}}\skt{ }\edtext{\skt{नित्यं}}{\lemma{\skt{नित्यं}}\Afootnote{K N H; \skt{नित्यः}\ A}}\skt{ }\edtext{\skt{संस्लेषः}}{\lemma{\skt{संस्लेषः}}\Afootnote{K N; \skt{संश्लेषः}\ A H}}\skt{ परिच्छेदो }\edtext{\skt{वेति}}{\lemma{\skt{वेति}}\Afootnote{K H; \skt{वा? इति}\ A; \skt{वेति~। य}\ N}}\skt{~। }\edtext{\skt{यदि}}{\lemma{\skt{यदि}}\Afootnote{A K H; \skt{दि हि}\ N}}\skt{ }\edtext{\skt{नित्यं}}{\lemma{\skt{नित्यं}}\Afootnote{K N H; \skt{नित्यः}\ A}}\skt{ }\edtext{\skt{संश्लेषः }\edtext{\skt{स्यान्नित्यातुरा}}{\lemma{\skt{नित्यातुरा}}\Afootnote{K H; \skt{त्तर्हि नित्यातुराः}\ A; \skt{स्यानित्यातुरा}\ N}}}{\lemma{\skt{संश्लेषः स्यान्}}\Afootnote{K H; \skt{संश्लेषः स्या}\ A; \skt{संस्लेषः}\ N}}\skt{ }\edtext{\skt{एव}}{\lemma{\skt{एव}}\Afootnote{K N H; \emph{om.}\ A}}\skt{ }\edtext{\skt{सर्वप्राणिनः स्युः}}{\lemma{\skt{सर्वप्राणिनः स्युः}}\Afootnote{K H; \skt{सर्व एव प्राणिनः स्युः}\ A; \skt{सर्वप्रानिनः स्यु}\ N}}\skt{~। अथापि }\edtext{\skt{अन्यथाभावो}}{\lemma{\skt{अन्यथाभावो}}\Afootnote{K N H; \skt{अन्यथा}\ A}}\skt{ }\edtext{\skt{वातादीनां}}{\lemma{\skt{वातादीनां}}\Afootnote{A K H; \skt{वातादीना}\ N}}\skt{ ज्वरादीनां }\edtext{\skt{चान्यत्र}}{\lemma{\skt{चान्यत्र}}\Afootnote{K N H; \skt{च, ‘अन्यत्र}\ A}}\skt{ }\edtext{\skt{वर्त्त\-मा\-न\-स्या\-न्यस्य}}{\lemma{\skt{वर्त्त\appbrk{}मा\appbrk{}न\appbrk{}स्या\appbrk{}न्यस्य}}\Afootnote{K N H; \skt{वर्त\appbrk{}मा\appbrk{}ना\appbrk{}ना\appbrk{}म\appbrk{}न्यत्र}\ A}}\skt{ लिङ्गं न }\edtext{\skt{भवतीति~। कृत्वा वातादयो}}{\lemma{\skt{भवतीति~। कृत्वा वातादयो}}\Afootnote{K H; \skt{भवति’ इति कृत्वा यदुच्यते वातादयो}\ A; \skt{भवतीति~। कृत्वा वातादयौ}\ N}}\skt{ }\edtext{\skt{ज्वरादीणां}}{\lemma{\skt{ज्वरादीणां}}\Afootnote{\skt{ज्वरादीनां}\ A K H; \skt{ज्वरादीयां}\ N}}\skt{ }\edtext{\skt{मूलानीति तन्न}}{\lemma{\skt{मूलानीति तन्न}}\Afootnote{A K\textsuperscript{pc} N H; \skt{मूलानीति तन्नः}\ K\textsuperscript{ac}}}\skt{~॥ अत्रोच्यते~॥ }\edtext{\skt{दोषान्प्रत्याख्याय}}{\lemma{\skt{दोषान्प्रत्याख्याय}}\Afootnote{A K H; \skt{दोषाप्रत्याख्याय}\ N}}\skt{ ज्वरादयो न भवन्ति~। }\edtext{\skt{अथ च}}{\lemma{\skt{अथ च}}\Afootnote{K H; \skt{अथ च न}\ A N}}\skt{ }\edtext{\skt{नित्यं सम्बन्धः}}{\lemma{\skt{नित्यं सम्बन्धः}}\Afootnote{K H; \skt{(नित्यः) सम्बन्धः}\ A; \skt{नित्यसम्बद्धः}\ N}}\skt{~। यथा हि }\edtext{\skt{विद्यु\-द्वा\-ता\-स\-नि\-व\-र्ष\-ण्या\-काशं}}{\lemma{\skt{विद्यु\appbrk{}द्वा\appbrk{}ता\appbrk{}स\appbrk{}नि\appbrk{}व\appbrk{}र्ष\appbrk{}ण्या\appbrk{}काशं}}\Afootnote{K; \skt{विद्यु\appbrk{}द्वा\appbrk{}ता\appbrk{}श\appbrk{}नि\appbrk{}व\appbrk{}र्षा\appbrk{}ण्या\appbrk{}काशं}\ A; \skt{विद्यु\appbrk{}द्वा\appbrk{}ता\appbrk{}स\appbrk{}नि\appbrk{}व\appbrk{}र्षा\appbrk{}ण्या\appbrk{}कास}\ N; \skt{विद्यु\appbrk{}द्वा\appbrk{}ता\appbrk{}श\appbrk{}नि\appbrk{}व\appbrk{}र्ष\appbrk{}ण्या\appbrk{}का\appbrk{}सम्}\ H}}\skt{ }\edtext{\skt{प्रत्याख्याय न}}{\lemma{\skt{प्रत्याख्याय न}}\Afootnote{A K\textsuperscript{pc} N H; \skt{प्रत्याख्याय}\ K\textsuperscript{ac}}}\skt{ भवन्ति~। }\edtext{\skt{सत्यप्याकासे}}{\lemma{\skt{सत्यप्याकासे}}\Afootnote{N; \skt{सत्यप्याकाशे}\ A K H}}\skt{ }\edtext{\skt{कदाचिच्च न}}{\lemma{\skt{कदाचिच्च न}}\Afootnote{K N H; \skt{कदाचिन्न}\ A}}\skt{ भवन्ति~। अथ च निमित्ततो }\edtext{\skt{भवन्ति}}{\lemma{\skt{भवन्ति}}\Afootnote{K N H; \skt{तत एवोत्पत्तिरिति}\ A}}\skt{~। }\edtext{\skt{तर\-ङ्ग\-बु\-द्बु\-दा\-द\-य\-श्चो\-द\-क\-वि\-शेषाः}}{\lemma{\skt{चोदकविशेषाः}}\Afootnote{A K H; \skt{चोदके विसेषाः}\ N}}\skt{~। }\edtext{\edtext{\skt{एवम्वातादीनां}}{\lemma{\skt{एवम्}}\Afootnote{K N H; \skt{एव}\ A}}}{\lemma{\skt{वातादीनां}}\Afootnote{A K H; \skt{वातादीना}\ N}}\skt{ }\edtext{\skt{ज्वरादीनाञ्च न}}{\lemma{\skt{ज्वरादीनाञ्च न}}\Afootnote{K N H; \skt{ज्वरादीनां च ना}\ A}}\skt{ }\edtext{\skt{नित्यसंश्लेषो न}}{\lemma{\skt{नित्यसंश्लेषो न}}\Afootnote{K N H; \skt{प्येवं संश्लेषो न}\ A}}\skt{ }\edtext{\skt{विच्छेदः}}{\lemma{\skt{विच्छेदः}}\Afootnote{K N H; \skt{परिच्छेदः}\ A}}\skt{ शाश्वतिकः~। अथ च }\edtext{\skt{निमित्ततः}}{\lemma{\skt{निमित्ततः}}\Afootnote{K N H; \skt{निमित्त}\ A}}\skt{~। }\edtext{\skt{तेभ्य}}{\lemma{\skt{तेभ्य}}\Afootnote{K N H; \emph{om.}\ A}}\skt{ }\edtext{\skt{एवोत्पत्तिरिति}}{\lemma{\skt{एवोत्पत्तिर्}}\Afootnote{N H; \skt{त एवोत्पत्तिर्}\ A; \skt{एवो प पत्तिर्}\ K\textsuperscript{ac}; \skt{एवो त्पत्तिर्}\ K\textsuperscript{pc}}}\skt{~॥}\pend

\pstart\versepar \ledsidenote{\passref{1.24.12}}\edtext{\skt{विका\-र\-म्प\-रि\-मा\-णञ्च}}{\lemma{\skt{विकारम्}}\Afootnote{K H; \skt{भवति चात्र विकार}\ A; \skt{विकार}\ N}}\skt{ संख्या चैषां पृथक्पृथक्~। }\pend
\pstart\versepar \skt{विस्तरेणोत्तरे तन्त्रे }\edtext{\skt{सर्वा}}{\lemma{\skt{सर्वा}}\Afootnote{A K N; \skt{सर्व्व}\ H}}\skt{ }\edtext{\skt{बाधां}}{\lemma{\skt{बाधां}}\Afootnote{K N H; \skt{बाधाश्}\ A}}\skt{ }\edtext{\skt{प्रचक्ष्मह}}{\lemma{\skt{प्रचक्ष्मह}}\Afootnote{K H; \skt{च वक्ष्यते}\ A; \skt{प्रवक्ष्युह}\ N}}\skt{ }\edtext{\skt{इति}}{\lemma{\skt{इति}}\Afootnote{K H; \emph{om.}\ A; \skt{रेति}\ N}}\skt{~॥~॥}\pend

\pstart \ledsidenote{\passref{1.24.trailer}}\edtext{\skt{॥ २४~॥}}{\lemma{\skt{॥ २४~॥}}\Afootnote{\skt{इति सुश्रु\appbrk{}त\appbrk{}सं\appbrk{}हि\appbrk{}तायां सूत्रस्थाने व्याधि\appbrk{}स\appbrk{}मु\appbrk{}द्दे\appbrk{}शीयो नाम चतुर्विंशोऽध्यायः}\ A; \skt{२४~॥ {\symfont ❈}}\ K; \skt{२४~॥}\ H}}\pend

\endnumbering
\beginnumbering
\eledsection*{Adhyāya 25}
\pstart \ledsidenote{\passref{1.25.1}}\edtext{\edtext{\skt{अथा\-तोऽ\-ष्ट\-वि\-ध\-श\-स्त्र\-क\-र्मी\-य\-म\-ध्यायं}}{\lemma{\skt{अथातो}}\Afootnote{A K H; \skt{अथातो}\ N}}}{\lemma{\skt{ऽष्ट\appbrk{}वि\appbrk{}ध\appbrk{}श\appbrk{}स्त्र\appbrk{}क\appbrk{}र्मी\appbrk{}य\appbrk{}म\appbrk{}ध्यायं}}\Afootnote{A; \skt{ष्टवि\appbrk{}ध\appbrk{}श\appbrk{}स्त्र\appbrk{}क\appbrk{}र्म\appbrk{}वि\appbrk{}धि\appbrk{}नि\appbrk{}श्च\appbrk{}यम्}\ K; \skt{ष्टवि\appbrk{}ध\appbrk{}स\appbrk{}स्त्र\appbrk{}क\appbrk{}र्म\appbrk{}वि\appbrk{}धि\appbrk{}नि\appbrk{}श्चयं}\ N; \skt{ऽष्ट\appbrk{}वि\appbrk{}धि\appbrk{}श\appbrk{}स्त्र\appbrk{}क\appbrk{}र्म्म\appbrk{}वि\appbrk{}धि\appbrk{}नि\appbrk{}श्च\appbrk{}यम्}\ H}}\skt{ व्याख्यास्यामः~॥}\pend

\pstart \ledsidenote{\passref{1.25.3-4}}\footnoteA{Prose from here to 12ab. These passages are in verse in the vulgate.}\edtext{\skt{छेद्यास्तु}}{\lemma{\skt{छेद्यास्तु}}\Afootnote{K N\textsuperscript{pc} H; \emph{om.}\ N\textsuperscript{ac}}}\skt{ }\edtext{\edtext{\edtext{\skt{भग\-न्द\-रा\-र्शो\-र्बु\-द\-दु\-ष्ट\-व्र\-ण\-ना\-डी\-च\-र्म\-की\-ल\-ति\-ल\-का\-ल\-क\-मे\-दो\-ग्र\-न्थि\-श्ले\-ष्म\-नि\-मि\-त्ता\-श्चा\-म\-स्तथा}}{\lemma{\skt{भग\appbrk{}न्द\appbrk{}रा\appbrk{}र्शो\appbrk{}र्बु\appbrk{}द\appbrk{}दु\appbrk{}ष्ट\appbrk{}व्र\appbrk{}ण\appbrk{}ना\appbrk{}डी\appbrk{}च\appbrk{}र्म\appbrk{}की\appbrk{}ल\appbrk{}ति\appbrk{}ल\appbrk{}का\appbrk{}ल\appbrk{}क\appbrk{}मे\appbrk{}दो\appbrk{}ग्र\appbrk{}न्थि\appbrk{}श्ले\appbrk{}ष्म\appbrk{}नि\appbrk{}मि\appbrk{}त्ताश्}}\Afootnote{K H; \emph{om.}\ N\textsuperscript{ac}; \skt{भग\appbrk{}न्द\appbrk{}रा\appbrk{}सौ\appbrk{}र्बु\appbrk{}द\appbrk{}दु\appbrk{}ष्ट\appbrk{}व्र\appbrk{}ण\appbrk{}ना\appbrk{}डी\appbrk{}च\appbrk{}र्म\appbrk{}की\appbrk{}ल\appbrk{}ति\appbrk{}ल\appbrk{}का\appbrk{}ल\appbrk{}क\appbrk{}मे\appbrk{}दो\appbrk{}ग्र\appbrk{}न्थि\appbrk{}स्ले\appbrk{}ष्म\appbrk{}नि\appbrk{}मि\appbrk{}त्तास्}\ N\textsuperscript{pc}}}}{\lemma{\skt{चामस्}}\Afootnote{K N\textsuperscript{pc} H; \emph{om.}\ N\textsuperscript{ac}}}}{\lemma{\skt{तथा}}\Afootnote{K H; \emph{om.}\ N\textsuperscript{ac}; \skt{लथा}\ N\textsuperscript{pc}}}\skt{ }\edtext{\edtext{\skt{मांस\-सि\-रा\-स्ना\-यु\-को\-थ\-मां\-स\-क\-न्दी\-मां\-सा\-स्थि\-ग\-त\-मपि}}{\lemma{\skt{मांस\appbrk{}सि\appbrk{}रा\appbrk{}स्ना\appbrk{}यु\appbrk{}को\appbrk{}थ\appbrk{}मां\appbrk{}स\appbrk{}क\appbrk{}न्दी\appbrk{}मां\appbrk{}सा\appbrk{}स्थि\appbrk{}ग\appbrk{}तम्}}\Afootnote{K H; \emph{om.}\ N\textsuperscript{ac}; \skt{मान्स\appbrk{}सि\appbrk{}रा\appbrk{}स्ना\appbrk{}यु\appbrk{}को\appbrk{}थ\appbrk{}मा\appbrk{}न्स\appbrk{}क\appbrk{}र्दी\appbrk{}मा\appbrk{}न्सा\appbrk{}स्थि\appbrk{}ग\appbrk{}तम्}\ N\textsuperscript{pc}}}\skt{ च}}{\lemma{\skt{अपि च}}\Afootnote{K N\textsuperscript{pc} H; \emph{om.}\ N\textsuperscript{ac}}}\skt{ }\edtext{\edtext{\skt{शल्यमुत्तवः}}{\lemma{\skt{शल्यम्}}\Afootnote{K N\textsuperscript{pc} H; \emph{om.}\ N\textsuperscript{ac}}}}{\lemma{\skt{उत्तवः}}\Afootnote{K N\textsuperscript{pc}; \emph{om.}\ N\textsuperscript{ac}; \skt{उत्सङ्गः}\ H}}\skt{ }\edtext{\skt{शतपोनको}}{\lemma{\skt{शतपोनको}}\Afootnote{K N\textsuperscript{pc} H; \emph{om.}\ N\textsuperscript{ac}}}\skt{ }\edtext{\edtext{\skt{जतु\-म\-णि\-र्व\-ल्मी\-क\-म\-धू\-षो\-धि\-मांसः}}{\lemma{\skt{जतुमणिर्}}\Afootnote{K N\textsuperscript{pc} H; \emph{om.}\ N\textsuperscript{ac}}}}{\lemma{\skt{वल्मी\appbrk{}क\appbrk{}म\appbrk{}धू\appbrk{}षो\appbrk{}धि\appbrk{}मांसः}}\Afootnote{K; \emph{om.}\ N\textsuperscript{ac}; \skt{वल्मी\appbrk{}क\appbrk{}म\appbrk{}ध्रु\appbrk{}षो\appbrk{}धि\appbrk{}मांसः}\ N\textsuperscript{pc}; \skt{व्वल्मी\appbrk{}क\appbrk{}म\appbrk{}धु\appbrk{}षो\appbrk{}धि\appbrk{}मान्स}\ H\textsuperscript{ac}; \skt{व्वल्मी\appbrk{}क\appbrk{}म\appbrk{}धु\appbrk{}षो\appbrk{}धि\appbrk{}मान्सः}\ H\textsuperscript{pc}}}\skt{ }\edtext{\skt{मांस\-सं\-घा\-त\-ग\-ल\-शु\-ण्डिका}}{\lemma{\skt{मांस\appbrk{}सं\appbrk{}घा\appbrk{}त\appbrk{}ग\appbrk{}ल\appbrk{}शु\appbrk{}ण्डिका}}\Afootnote{K H; \emph{om.}\ N\textsuperscript{ac}; \skt{मांसं\appbrk{}घा\appbrk{}त\appbrk{}ग\appbrk{}ल\appbrk{}शु\appbrk{}ण्डिका}\ N\textsuperscript{pc}}}\skt{ }\edtext{\edtext{\edtext{\skt{इत्येवमादयो}}{\lemma{\skt{इत्य्}}\Afootnote{K N\textsuperscript{pc} H; \emph{om.}\ N\textsuperscript{ac}}}}{\lemma{\skt{एवम्}}\Afootnote{K N\textsuperscript{pc} H; \emph{om.}\ N\textsuperscript{ac}}}}{\lemma{\skt{आदयो}}\Afootnote{K N\textsuperscript{pc} H; \emph{om.}\ N\textsuperscript{ac}}}\skt{ }\edtext{\skt{विकाराः}}{\lemma{\skt{विकाराः}}\Afootnote{K N\textsuperscript{pc} H; \emph{om.}\ N\textsuperscript{ac}}}\skt{~॥}\pend

\pstart \ledsidenote{\passref{1.25.5-8}}\edtext{\skt{भेद्यास्तु}}{\lemma{\skt{भेद्यास्तु}}\Afootnote{K N\textsuperscript{pc} H; \emph{om.}\ N\textsuperscript{ac}}}\skt{ }\edtext{\skt{सर्वजमृते}}{\lemma{\skt{सर्वजमृते}}\Afootnote{\skt{सर्वजमृते}\ K N\textsuperscript{pc}; \emph{om.}\ N\textsuperscript{ac}; \skt{सर्व्वजामृते}\ H}}\skt{ }\edtext{\skt{विद्रध्यनुशयी}}{\lemma{\skt{विद्रध्यनुशयी}}\Afootnote{K H; \emph{om.}\ N\textsuperscript{ac}; \skt{विद्रप्ये तुशयी}\ N\textsuperscript{pc}}}\skt{ }\edtext{\skt{प्रमेहपिटका}}{\lemma{\skt{प्रमेहपिटका}}\Afootnote{K; \emph{om.}\ N\textsuperscript{ac}; \skt{प्रमेहपिढका}\ N\textsuperscript{pc}; \skt{प्रमेहपिटकां}\ H}}\skt{ }\edtext{\skt{ग्रन्थयश्च}}{\lemma{\skt{ग्रन्थयश्च}}\Afootnote{K N\textsuperscript{pc} H; \emph{om.}\ N\textsuperscript{ac}}}\skt{ }\edtext{\edtext{\edtext{\edtext{\skt{विस\-र्पा\-श्चा\-दि\-त\-स्त्र\-य\-स्त्रयो}}{\lemma{\skt{विसर्पाश्}}\Afootnote{K N\textsuperscript{pc} H; \emph{om.}\ N\textsuperscript{ac}}}}{\lemma{\skt{चादितस्}}\Afootnote{K N\textsuperscript{pc} H; \emph{om.}\ N\textsuperscript{ac}}}}{\lemma{\skt{त्रयस्}}\Afootnote{K N\textsuperscript{pc} H; \emph{om.}\ N\textsuperscript{ac}}}}{\lemma{\skt{त्रयो}}\Afootnote{K N\textsuperscript{pc} H; \emph{om.}\ N\textsuperscript{ac}}}\skt{ }\edtext{\skt{वृद्धयश्च}}{\lemma{\skt{वृद्धयश्च}}\Afootnote{K N\textsuperscript{pc} H; \emph{om.}\ N\textsuperscript{ac}}}\skt{ }\edtext{\skt{विदा\-रि\-का\-व\-मन्थौ}}{\lemma{\skt{विदा\appbrk{}रि\appbrk{}का\appbrk{}व\appbrk{}मन्थौ}}\Afootnote{K N\textsuperscript{pc} H; \emph{om.}\ N\textsuperscript{ac}}}\skt{ }\edtext{\skt{पुष्क\-रि\-का\-ना\-ड्य\-स्त\-न\-रो\-गाश्च}}{\lemma{\skt{पुष्क\appbrk{}रि\appbrk{}का\appbrk{}ना\appbrk{}ड्य\appbrk{}स्त\appbrk{}न\appbrk{}रो\appbrk{}गाश्च}}\Afootnote{K H; \emph{om.}\ N\textsuperscript{ac}; \skt{पुष्क\appbrk{}रि\appbrk{}का\appbrk{}ना\appbrk{}न्य\appbrk{}स्तन रोगाश्च}\ N\textsuperscript{pc}}}\skt{ }\edtext{\skt{सहोपदंशैः}}{\lemma{\skt{सहोपदंशैः}}\Afootnote{K H; \emph{om.}\ N\textsuperscript{ac}; \skt{सहोपदसैः}\ N\textsuperscript{pc}}}\skt{ }\edtext{\edtext{\skt{शोफा\-श्चा\-स\-र्व\-सराः}}{\lemma{\skt{शोफाश्}}\Afootnote{K H; \emph{om.}\ N\textsuperscript{ac}; \skt{सोहा}\ N\textsuperscript{pc}}}}{\lemma{\skt{चासर्वसराः}}\Afootnote{K H; \emph{om.}\ N\textsuperscript{ac}; \skt{र्स्वसर्वसराः}\ N\textsuperscript{pc}}}\skt{ }\edtext{\skt{प्रायशः}}{\lemma{\skt{प्रायशः}}\Afootnote{K N\textsuperscript{pc} H; \emph{om.}\ N\textsuperscript{ac}}}\skt{ }\edtext{\edtext{\skt{क्षुद्ररोगास्तथा}}{\lemma{\skt{क्षुद्ररोगास्}}\Afootnote{K N\textsuperscript{pc} H; \emph{om.}\ N\textsuperscript{ac}}}}{\lemma{\skt{तथा}}\Afootnote{K N\textsuperscript{pc}; \emph{om.}\ N\textsuperscript{ac} H}}\skt{ }\edtext{\skt{तालु\-पु\-प्पु\-ट\-तु\-ण्डि\-के\-री\-गि\-ला\-यु\-प्र\-भृ\-तयो}}{\lemma{\skt{तालु\appbrk{}पु\appbrk{}प्पु\appbrk{}ट\appbrk{}तु\appbrk{}ण्डि\appbrk{}के\appbrk{}री\appbrk{}गि\appbrk{}ला\appbrk{}यु\appbrk{}प्र\appbrk{}भृ\appbrk{}तयो}}\Afootnote{K H; \emph{om.}\ N\textsuperscript{ac}; \skt{ताल\appbrk{}पु\appbrk{}प्पु\appbrk{}ढ\appbrk{}तु\appbrk{}ण्डि\appbrk{}के\appbrk{}री\appbrk{}गि\appbrk{}ला\appbrk{}यु\appbrk{}प्र\appbrk{}भृ\appbrk{}तयो}\ N\textsuperscript{pc}}}\skt{ }\edtext{\skt{न्तः}}{\lemma{\skt{न्तः}}\Afootnote{K N\textsuperscript{pc} H; \emph{om.}\ N\textsuperscript{ac}}}\skt{ }\edtext{\skt{पूयशल्याः}}{\lemma{\skt{पूयशल्याः}}\Afootnote{K N\textsuperscript{pc} H; \emph{om.}\ N\textsuperscript{ac}}}\skt{ }\edtext{\skt{पञ्च}}{\lemma{\skt{पञ्च}}\Afootnote{K N\textsuperscript{pc} H; \emph{om.}\ N\textsuperscript{ac}}}\skt{ }\edtext{\skt{मेदः समुत्था}}{\lemma{\skt{मेदः समुत्था}}\Afootnote{K\textsuperscript{pc} N\textsuperscript{pc} H; \skt{मेदसमुत्था}\ K\textsuperscript{ac}; \emph{om.}\ N\textsuperscript{ac}}}\skt{ }\edtext{\edtext{\skt{अश्म\-री\-हे\-तो\-र्ब\-स्तिश्च}}{\lemma{\skt{अश्मरीहेतोर्}}\Afootnote{K H; \emph{om.}\ N\textsuperscript{ac}; \skt{अस्मरीहेतो}\ N\textsuperscript{pc}}}}{\lemma{\skt{बस्तिश्च}}\Afootnote{K H; \emph{om.}\ N\textsuperscript{ac}; \skt{वस्तिकश्च}\ N\textsuperscript{pc}}}\skt{~॥}\pend

\pstart \ledsidenote{\passref{1.25.9}}\edtext{\edtext{\skt{लेख्याश्चतस्रो}}{\lemma{\skt{लेख्याश्}}\Afootnote{A K H; \emph{om.}\ N\textsuperscript{ac}; \skt{लेखाश्}\ N\textsuperscript{pc}}}}{\lemma{\skt{चतस्रो}}\Afootnote{A K N\textsuperscript{pc} H; \emph{om.}\ N\textsuperscript{ac}}}\skt{ }\edtext{\skt{रोहिण्यः}}{\lemma{\skt{रोहिण्यः}}\Afootnote{A K; \emph{om.}\ N\textsuperscript{ac}; \skt{रोहिण्या}\ N\textsuperscript{pc}; \skt{रोहिण्य}\ H}}\skt{ }\edtext{\skt{व्रण\-ने\-त्र\-व\-र्मा\-द्य\-धि\-जि\-ह्वो\-प\-जिह्वे}}{\lemma{\skt{व्रण\appbrk{}ने\appbrk{}त्र\appbrk{}व\appbrk{}र्मा\appbrk{}द्य\appbrk{}धि\appbrk{}जि\appbrk{}ह्वो\appbrk{}प\appbrk{}जिह्वे}}\Afootnote{K; \skt{किला\appbrk{}स\appbrk{}मु\appbrk{}प\appbrk{}जिह्वि}\ A; \emph{om.}\ N\textsuperscript{ac}; \skt{व्रण\appbrk{}मे\appbrk{}त्र\appbrk{}व\appbrk{}र्मा\appbrk{}ण्य\appbrk{}धि\appbrk{}जि\appbrk{}ह्वो\appbrk{}प\appbrk{}जिह्वे}\ N\textsuperscript{pc}; \skt{व्रण\appbrk{}ने\appbrk{}त्र\appbrk{}व\appbrk{}र्म्मा\appbrk{}ण्य\appbrk{}धि\appbrk{}जि\appbrk{}ह्वो\appbrk{}प\appbrk{}जिह्वे}\ H}}\skt{ }\edtext{\skt{मांसोच्छ्रयः}}{\lemma{\skt{मांसोच्छ्रयः}}\Afootnote{K N\textsuperscript{pc} H; \emph{om.}\ A N\textsuperscript{ac}}}\skt{ }\edtext{\skt{किलासो}}{\lemma{\skt{किलासो}}\Afootnote{K N\textsuperscript{pc} H; \skt{का~। मेदोजो}\ A; \emph{om.}\ N\textsuperscript{ac}}}\skt{ }\edtext{\skt{दन्तवैदर्भः}}{\lemma{\skt{दन्तवैदर्भः}}\Afootnote{A K N\textsuperscript{pc} H; \emph{om.}\ N\textsuperscript{ac}}}\skt{ }\edtext{\skt{पञ्च}}{\lemma{\skt{पञ्च}}\Afootnote{K N\textsuperscript{pc} H; \skt{ग्रन्}\ A; \emph{om.}\ N\textsuperscript{ac}}}\skt{ }\edtext{\skt{मेदोजाश्च}}{\lemma{\skt{मेदोजाश्च}}\Afootnote{K N\textsuperscript{pc} H; \skt{थिर्व\appbrk{}र्त्मा\appbrk{}धि\appbrk{}जि\appbrk{}ह्विका}\ A; \emph{om.}\ N\textsuperscript{ac}}}\skt{~॥}\pend

\pstart \ledsidenote{\passref{1.25.10cd}}\edtext{\edtext{\skt{वेध्यास्तु}}{\lemma{\skt{वेध्यास्}}\Afootnote{A; \emph{om.}\ K N\textsuperscript{ac} H; \skt{व्येध्यास्}\ N\textsuperscript{pc}}}}{\lemma{\skt{तु}}\Afootnote{N\textsuperscript{pc}; \skt{सिरा बहु}\ A; \skt{वेध्यास्तु}\ K H; \emph{om.}\ N\textsuperscript{ac}}}\skt{ }\edtext{\skt{सिरो\-द\-र\-वृ\-द्धि\-प्र\-भृ\-तयः}}{\lemma{\skt{सिरो\appbrk{}द\appbrk{}र\appbrk{}वृ\appbrk{}द्धि\appbrk{}प्र\appbrk{}भृ\appbrk{}तयः}}\Afootnote{K N\textsuperscript{pc} H; \skt{विधा मूत्र\appbrk{}वृ\appbrk{}द्धि\appbrk{}र्द\appbrk{}को\appbrk{}द\appbrk{}रम्}\ A; \emph{om.}\ N\textsuperscript{ac}}}\skt{~॥}\pend

\pstart \ledsidenote{\passref{1.25.11-12ab}}\edtext{\skt{एष्यास्तु}}{\lemma{\skt{एष्यास्तु}}\Afootnote{K H; \skt{एष्या}\ A; \emph{om.}\ N\textsuperscript{ac}; \skt{एषास्तु}\ N\textsuperscript{pc}}}\skt{ }\edtext{\skt{गतिमन्तो}}{\lemma{\skt{गतिमन्तो}}\Afootnote{K N\textsuperscript{pc} H; \emph{om.}\ A N\textsuperscript{ac}}}\skt{ }\edtext{\skt{व्रणा}}{\lemma{\skt{व्रणा}}\Afootnote{K N\textsuperscript{pc} H; \emph{om.}\ A N\textsuperscript{ac}}}\skt{ }\edtext{\skt{नाड्यः}}{\lemma{\skt{नाड्यः}}\Afootnote{A K H; \emph{om.}\ N\textsuperscript{ac}; \skt{नान्यः}\ N\textsuperscript{pc}}}\skt{ }\edtext{\skt{शल्यानि च}}{\lemma{\skt{शल्यानि च}}\Afootnote{K N\textsuperscript{pc} H; \skt{सशल्याश्च व्रणा}\ A; \emph{om.}\ N\textsuperscript{ac}}}\skt{~। }\edtext{\skt{आहर्त्तव्यास्तु}}{\lemma{\skt{आहर्त्तव्यास्तु}}\Afootnote{K N\textsuperscript{pc}; \skt{उन्मार्गिणश्च ये}\ A; \emph{om.}\ N\textsuperscript{ac} H}}\skt{ }\edtext{\skt{दन्तान्तः}}{\lemma{\skt{दन्तान्तः}}\Afootnote{K N\textsuperscript{pc}; \skt{आहार्याः}\ A; \emph{om.}\ N\textsuperscript{ac} H}}\skt{ }\edtext{\skt{शर्क\-रा\-श्म\-री\-मू\-ढ\-ग\-र्भ\-क\-र्ण्ण\-म\-ल\-श\-ल्यानि}}{\lemma{\skt{शर्क\appbrk{}रा\appbrk{}श्म\appbrk{}री\appbrk{}मू\appbrk{}ढ\appbrk{}ग\appbrk{}र्भ\appbrk{}क\appbrk{}र्ण्ण\appbrk{}म\appbrk{}ल\appbrk{}श\appbrk{}ल्यानि}}\Afootnote{K N\textsuperscript{pc}; \skt{शर्करास्तिस्रो दन्त\appbrk{}क\appbrk{}र्ण\appbrk{}म\appbrk{}लोऽ\appbrk{}श्मरी~॥११॥ शल्यानि}\ A; \emph{om.}\ N\textsuperscript{ac} H}}\skt{ }\edtext{\skt{पादशर्करा च}}{\lemma{\skt{पादशर्करा च}}\Afootnote{K N\textsuperscript{pc}; \skt{मूढगर्भाश्च वर्चश्च निचितं गुदे}\ A; \emph{om.}\ N\textsuperscript{ac} H}}\skt{~॥}\pend

\pstart \ledsidenote{\passref{1.25.12cd}}\edtext{\skt{विस्राव्यो}}{\lemma{\skt{विस्राव्यो}}\Afootnote{K N\textsuperscript{pc} H; \skt{स्राव्या}\ A; \emph{om.}\ N\textsuperscript{ac}}}\skt{ }\edtext{\skt{विद्रधिः}}{\lemma{\skt{विद्रधिः}}\Afootnote{K N\textsuperscript{pc} H; \skt{विद्र}\ A; \emph{om.}\ N\textsuperscript{ac}}}\skt{ }\edtext{\skt{सर्वो}}{\lemma{\skt{सर्वो}}\Afootnote{K N\textsuperscript{pc} H; \skt{धयः}\ A; \emph{om.}\ N\textsuperscript{ac}}}\skt{ }\edtext{\skt{भवेदन्यत्र}}{\lemma{\skt{भवेदन्यत्र}}\Afootnote{K H; \skt{पञ्च भवेयुः}\ A; \emph{om.}\ N\textsuperscript{ac}; \skt{भवोदन्यत्र}\ N\textsuperscript{pc}}}\skt{ }\edtext{\skt{सर्वजात्}}{\lemma{\skt{सर्वजात्}}\Afootnote{K H; \skt{सर्वजादृते}\ A; \emph{om.}\ N\textsuperscript{ac}; \skt{सर्वजात}\ N\textsuperscript{pc}}}\skt{~॥}\pend

\pstart\versepar \ledsidenote{\passref{1.25.13}}\edtext{\skt{कुष्ठानि}}{\lemma{\skt{कुष्ठानि}}\Afootnote{A K H; \skt{कुष्टानि}\ N}}\skt{ वायुः सरुजः शोफो }\edtext{\skt{यश्चैकदेशजः}}{\lemma{\skt{यश्चैकदेशजः}}\Afootnote{A; \skt{चाप्येकदेशजः}\ K N H}}\skt{~। }\pend
\pstart\versepar \skt{}\edtext{\skt{पाल्यामयाः}}{\lemma{\skt{पाल्यामयाः}}\Afootnote{A; \skt{पाल्यामया}\ K N H}}\skt{ श्लीपदानि }\edtext{\skt{विषजुष्टं च}}{\lemma{\skt{विषजुष्टं च}}\Afootnote{A; \skt{विषदुष्टश्च}\ K; \skt{विषदुष्टञ्च}\ N; \skt{विषजुष्टश्च}\ H}}\skt{ शोणितम्~॥}\pend

\pstart\versepar \ledsidenote{\passref{1.25.14}}\skt{अर्बुदानि विसर्पाश्च }\edtext{\edtext{\skt{ग्रन्थ\-य\-श्चा\-दि\-तस्तु}}{\lemma{\skt{ग्रन्थयश्}}\Afootnote{A K H; \skt{ग्रन्थयाश्}\ N}}\skt{ ये}}{\lemma{\skt{चादितस्तु ये}}\Afootnote{\skt{चादितश्च ते}\ A; \skt{चादितश्च ये}\ K N H}}\skt{~। }\pend
\pstart\versepar \skt{}\edtext{\skt{त्रय\-स्त्र\-य\-श्चो\-प\-दंशाः}}{\lemma{\skt{चोपदंशाः}}\Afootnote{A; \skt{ते विश्राव्या}\ K H; \skt{ते विश्राष्चा}\ N}}\skt{ }\edtext{\skt{स्तनरोगा}}{\lemma{\skt{स्तनरोगा}}\Afootnote{A; \skt{स्तनरोगास्}\ K N H}}\skt{ }\edtext{\skt{विदारिका}}{\lemma{\skt{विदारिका}}\Afootnote{A; \skt{तथैव च}\ K H; \skt{तञ्चैव च}\ N}}\skt{~॥}\pend

\pstart\versepar \ledsidenote{\passref{1.25.15}}\edtext{\skt{सुषिरो}}{\lemma{\skt{सुषिरो}}\Afootnote{\skt{सु(शु)षिरो}\ A; \skt{सौषिरः}\ K N; \skt{सौशिरः}\ H}}\skt{ }\edtext{\skt{गलशालूकं}}{\lemma{\skt{गलशालूकं}}\Afootnote{A; \skt{कण्ठशालूकाः}\ K H; \skt{कण्ठशालूकः}\ N}}\skt{ }\edtext{\skt{कण्टकाः}}{\lemma{\skt{कण्टकाः}}\Afootnote{A K H; \skt{कणूका}\ N}}\skt{ कृमिदन्तकः~। }\pend
\pstart\versepar \skt{दन्तवेष्टः सोपकुशः शीतादो }\edtext{\skt{दन्तपुप्पुटः}}{\lemma{\skt{दन्तपुप्पुटः}}\Afootnote{A K H; \skt{दन्तपुप्पटः}\ N}}\skt{~॥}\pend

\pstart\versepar \ledsidenote{\passref{1.25.16}}\edtext{\skt{पित्ता\-सृ\-क्क\-फ\-जा\-श्चौष्ठ्याः}}{\lemma{\skt{चौष्ठ्याः}}\Afootnote{A N; \skt{च्यौष्ठ्याः}\ K H}}\skt{ क्षुद्ररोगाश्च भूयशः~। }\pend
\pstart\versepar \skt{सीव्या }\edtext{\skt{मेदःसमुत्थाश्च}}{\lemma{\skt{मेदःसमुत्थाश्च}}\Afootnote{A; \skt{मेदःसमुत्थास्तु}\ K H; \skt{मेदसमुत्थास्तु}\ N}}\skt{ }\edtext{\skt{भिन्नाः}}{\lemma{\skt{भिन्नाः}}\Afootnote{A; \skt{भित्वा}\ K N H}}\skt{ }\edtext{\skt{सुलिखिता}}{\lemma{\skt{सुलिखिता}}\Afootnote{A; \skt{विलिखितं}\ K N H}}\skt{ }\edtext{\skt{गदाः}}{\lemma{\skt{गदाः}}\Afootnote{A; \skt{तथा}\ K N H}}\skt{~॥}\pend

\pstart\versepar \ledsidenote{\passref{1.25.17-28}}\edtext{\skt{कफ\-ग्र\-न्थि\-र\-ल्प\-पालिः}}{\lemma{\skt{कफग्रन्थिर्}}\Afootnote{K H; \skt{कप्णग्रन्थिर्}\ N}}\skt{ }\edtext{\skt{कर्ण्ण\-स\-द्यो\-व्र\-णश्च यः}}{\lemma{\skt{कर्ण्ण\appbrk{}स\appbrk{}द्यो\appbrk{}व्र\appbrk{}णश्च यः}}\Afootnote{K H; \skt{कर्ण्णसद्यो व्रणश्च यः}\ N}}\skt{~॥ }\pend
\pstart\versepar \skt{}\edtext{\skt{शिरो\-ल\-ला\-टा\-क्षि\-कृ\-ट\-क\-र्ण्णे\-ष्व\-ना\-सा\-ग\-ण्ड\-कृ\-का\-टिका}}{\lemma{\skt{शिरो\appbrk{}ल\appbrk{}ला\appbrk{}टा\appbrk{}क्षि\appbrk{}कृ\appbrk{}ट\appbrk{}क\appbrk{}र्ण्णेष्व्}}\Afootnote{K N; \skt{शिरो\appbrk{}ल\appbrk{}ला\appbrk{}टा\appbrk{}क्षि\appbrk{}कू\appbrk{}ट\appbrk{}क\appbrk{}र्ण्णेष्व्}\ H}}\skt{~। }\pend
\pstart\versepar \skt{}\edtext{\skt{बाहू\-द\-र\-स्फि\-क्षा\-यु\-प्र\-ज\-न\-न\-मु\-ष्का\-दिषु}}{\lemma{\skt{बाहू\appbrk{}द\appbrk{}र\appbrk{}स्फि\appbrk{}क्षा\appbrk{}यु\appbrk{}प्र\appbrk{}ज\appbrk{}न\appbrk{}न\appbrk{}मु\appbrk{}ष्का\appbrk{}दिषु}}\Afootnote{K; \skt{बाहू\appbrk{}द\appbrk{}र\appbrk{}स्फिं\appbrk{}क्पा\appbrk{}यु\appbrk{}प्र\appbrk{}ज\appbrk{}न\appbrk{}न\appbrk{}मु\appbrk{}ष्का\appbrk{}दिषु}\ N; \skt{बाहू\appbrk{}द\appbrk{}र\appbrk{}स्फि\appbrk{}ङ्पा\appbrk{}यु\appbrk{}प्र\appbrk{}ज\appbrk{}न\appbrk{}न\appbrk{}मु\appbrk{}ष्का\appbrk{}दिषु}\ H}}\skt{ प्रदे\-से\-ष्व\-च\-लेषु मांसवत्सु च }\edtext{\skt{शीव्येत्~। जनु\-कु\-र्प\-र\-जं\-घा\-दिषु}}{\lemma{\skt{शीव्येत्~। जनु\appbrk{}कु\appbrk{}र्प\appbrk{}र\appbrk{}जं\appbrk{}घा\appbrk{}दिषु}}\Afootnote{K; \skt{शीव्ये\appbrk{}त्जा\appbrk{}नु\appbrk{}कु\appbrk{}र्प\appbrk{}र\appbrk{}जं\appbrk{}घा\appbrk{}दिषु}\ N H}}\skt{ प्रच\-ले\-ष्व\-ल्प\-मां\-सेषु न }\edtext{\skt{सीव्येत्~। वायुनिर्वाहिनो}}{\lemma{\skt{सीव्येत्~। वायुनिर्वाहिनो}}\Afootnote{K H; \skt{सीव्येत्~। वायु\appbrk{}नि\appbrk{}र्वा\appbrk{}हि\appbrk{}नोन्}\ N}}\skt{ }\edtext{\skt{न्तलोहितशल्याः}}{\lemma{\skt{न्तलोहितशल्याः}}\Afootnote{K; \skt{तुलोहितशल्याः}\ N; \skt{ऽन्त\appbrk{}र्ल्लो\appbrk{}हि\appbrk{}त\appbrk{}शल्याः}\ H}}\skt{ सविषांश्च तत्र }\edtext{\skt{वासीव्यं}}{\lemma{\skt{वासीव्यं}}\Afootnote{\skt{सीव्यं}\ K\textsuperscript{ac}; \skt{वासीव्यं}\ K\textsuperscript{pc} N H}}\skt{ }\edtext{\skt{व्रण\-म\-भि\-स\-मीक्ष्य}}{\lemma{\skt{अभिसमीक्ष्य}}\Afootnote{K H; \skt{अभिसमीक्ष्य}\ N}}\skt{ चेला\-स्थि\-पां\-सु\-तृ\-ण\-रो\-म\-शु\-ष्क\-र\-क्ता\-दी\-न्य\-पो\-हो\-त्कृ\-त्या\-कृष्य यथास्थानं स्थाययित्वा स्नायु\-सू\-त्र\-वा\-लानां }\edtext{\skt{मन्यतमेन}}{\lemma{\skt{मन्यतमेन}}\Afootnote{K N; \skt{अन्यतमेन}\ H}}\skt{ }\edtext{\skt{सीव्येत्~। सणा\-श्म\-न्त\-क\-मू\-र्वा\-त\-शीनां वा}}{\lemma{\skt{सीव्येत्~। सणा\appbrk{}श्म\appbrk{}न्त\appbrk{}क\appbrk{}मू\appbrk{}र्वा\appbrk{}त\appbrk{}शीनां वा}}\Afootnote{K\textsuperscript{pc}; \skt{सीव्येत्~। सणा\appbrk{}श्म\appbrk{}न्त\appbrk{}क\appbrk{}मू\appbrk{}र्वा\appbrk{}त\appbrk{}शीनां व}\ K\textsuperscript{ac}; \skt{सीव्येत्~। सणा\appbrk{}ष्म\appbrk{}न्त\appbrk{}क\appbrk{}मू\appbrk{}र्वा\appbrk{}त\appbrk{}शीनां वा}\ N; \skt{सीव्येत्~। तृणा\appbrk{}ष्म\appbrk{}न्त\appbrk{}क\appbrk{}मू\appbrk{}र्वा\appbrk{}त\appbrk{}सीनां वा}\ H}}\skt{ वल्कलैः~। सूच्यस्तु }\edtext{\skt{तिस्र}}{\lemma{\skt{तिस्र}}\Afootnote{K H; \skt{तिसु}\ N}}\skt{ उपदिश्यन्ते~। द्व्यङ्गुला त्र्यङ्गुला }\edtext{\skt{धनुवक्रेति}}{\lemma{\skt{धनुवक्रेति}}\Afootnote{K; \skt{धनुवक्रेभि}\ N H}}\skt{~। तत्र }\edtext{\skt{मांस\-ले\-ष्व\-व\-का\-शेषु}}{\lemma{\skt{अवकाशेषु}}\Afootnote{K N; \skt{अवकासेषु}\ H}}\skt{ त्र्यस्राः }\edtext{\edtext{\skt{सत्वा\-स्थि\-स्व\-ल्प\-मां\-सेषु}}{\lemma{\skt{सत्वास्थिस्व्}}\Afootnote{K; \skt{सीन्व्यस्थि}\ N H}}\skt{ च}}{\lemma{\skt{अल्पमांसेषु च}}\Afootnote{K; \skt{स्चल्पमांसेषु च}\ N H}}\skt{ द्व्यङ्गुला वृत्ता }\edtext{\skt{यत्का\-मा\-स\-म\-यो\-र्म\-र्मसु}}{\lemma{\skt{यत्का\appbrk{}मा\appbrk{}स\appbrk{}म\appbrk{}योर्}}\Afootnote{K; \skt{पत्का\appbrk{}मा\appbrk{}स\appbrk{}म\appbrk{}योर्}\ N H}}\skt{ च }\edtext{\skt{धनुवत्चक्रार्धं}}{\lemma{\skt{धनुवत्चक्रार्धं}}\Afootnote{K; \skt{धनुवच्चक्रार्ध}\ N; \skt{धनुर्व्वक्रार्ध}\ H}}\skt{ तृतीयांगुलेति~॥ }\edtext{\skt{फुणि}}{\lemma{\skt{फुणि}}\Afootnote{K; \skt{रेधर्स्वाण}\ N; \skt{गोष्ठशि}\ H}}\skt{ }\edtext{\skt{कातू\-न\-से\-व\-न्या\-वें\-ल्लि\-तकं}}{\lemma{\skt{कातू\appbrk{}न\appbrk{}से\appbrk{}व\appbrk{}न्या\appbrk{}वें\appbrk{}ल्लि\appbrk{}तकं}}\Afootnote{K N; \skt{कातू\appbrk{}न\appbrk{}से\appbrk{}व\appbrk{}न्या\appbrk{}वे\appbrk{}ल्लि\appbrk{}तकं}\ H}}\skt{ रजग्रन्थिबन्धं चेति~। समासेन }\edtext{\skt{सेवनविकल्पाः}}{\lemma{\skt{सेवनविकल्पाः}}\Afootnote{K; \skt{सेचनविकल्पाः}\ N H}}\skt{~। तेषा\-न्ना\-म\-भि\-रे\-वा\-कृ\-तयः प्रायेण व्याख्याताः तेषां प्रहा\-र\-मा\-सा\-द्यो\-प\-योगां }\edtext{\skt{बुध्यावेक्ष्य न}}{\lemma{\skt{बुध्यावेक्ष्य न}}\Afootnote{K H; \skt{वुद्व्यावेक्ष्य न}\ N}}\skt{ चातिसन्निकृष्टां }\edtext{\edtext{\skt{विप्र\-कृ\-ष्टा\-म\-ल्प\-ग्रा\-हि\-णी\-म\-ति\-बहु}}{\lemma{\skt{अल्पग्राहिणीम्}}\Afootnote{K H; \skt{अल्प}\ N}}}{\lemma{\skt{अतिबहु}}\Afootnote{K H; \emph{om.}\ N}}\skt{ ग्राहिणीम्वा }\edtext{\skt{सूचीम्पातयेत्}}{\lemma{\skt{सूचीम्}}\Afootnote{K H; \skt{सूची}\ N}}\skt{~॥ }\pend
\pstart\versepar \skt{भवति चात्र~॥ }\pend
\pstart\versepar \skt{दूरे निपतिताबाधं }\edtext{\skt{सन्निकृष्टे}}{\lemma{\skt{सन्निकृष्टे}}\Afootnote{K; \skt{सन्निकृष्ट}\ N H}}\skt{ विलुञ्चनं~। }\pend
\pstart\versepar \skt{}\edtext{\skt{अल्प\-ग्र\-हा\-त्त\-द्व\-देव}}{\lemma{\skt{अल्प\appbrk{}ग्र\appbrk{}हा\appbrk{}त्त\appbrk{}द्व\appbrk{}देव}}\Afootnote{K; \skt{अल्प\appbrk{}ग्र\appbrk{}हा\appbrk{}न्त\appbrk{}द्व\appbrk{}देव}\ N H}}\skt{ व्रणे कुर्या\-द\-त\-स्त्य\-जेत्~॥ }\pend
\pstart\versepar \skt{सम्य\-क्सी\-वि\-त\-म\-वेक्ष्य }\edtext{\skt{मधु\-घृ\-त\-यु\-तै\-र\-ञ्ज\-न\-म\-धु\-क\-लो\-ध्र\-प्रि\-य\-ङ्गु\-श\-ल्ल\-की\-फ\-ल\-र\-सा\-ञ्ज\-न\-क्षौ\-म\-म\-सी\-चूर्ण्णैः}}{\lemma{\skt{अञ्ज\appbrk{}न\appbrk{}म\appbrk{}धु\appbrk{}क\appbrk{}लो\appbrk{}ध्र\appbrk{}प्रि\appbrk{}य\appbrk{}ङ्गु\appbrk{}श\appbrk{}ल्ल\appbrk{}की\appbrk{}फ\appbrk{}ल\appbrk{}र\appbrk{}सा\appbrk{}ञ्ज\appbrk{}न\appbrk{}क्षौ\appbrk{}म\appbrk{}म\appbrk{}सी\appbrk{}चूर्ण्णैः}}\Afootnote{K; \skt{अञ्ज\appbrk{}र\appbrk{}ने\appbrk{}म\appbrk{}धु\appbrk{}क\appbrk{}लो\appbrk{}ध्र\appbrk{}प्रि\appbrk{}य\appbrk{}ङ्गु\appbrk{}स\appbrk{}ल्ल\appbrk{}की\appbrk{}फ\appbrk{}ल\appbrk{}र\appbrk{}सा\appbrk{}ञ्ज\appbrk{}न\appbrk{}क्षौ\appbrk{}म\appbrk{}म\appbrk{}सी\appbrk{}चूर्ण्णैः}\ N H}}\skt{ प्रतिसार्य }\edtext{\skt{बन्धे\-नो\-प\-च\-रेत्}}{\lemma{\skt{बन्धे\appbrk{}नो\appbrk{}प\appbrk{}च\appbrk{}रेत्}}\Afootnote{K H; \skt{बन्धे\appbrk{}नो\appbrk{}प\appbrk{}च\appbrk{}रेत}\ N}}\skt{~॥}\footnoteA{This passage replaces vulgate 17-28..}\edtext{}{\lemma{\skt{पातयेत्~॥ भवति चात्र~॥ दूरे}}\Afootnote{\skt{पातयेत्दूरे}\ K H; \skt{प्रोतयेत्~। एवं वेद\appbrk{}ना\appbrk{}व\appbrk{}लु\appbrk{}ञ्च\appbrk{}न\appbrk{}म्भ\appbrk{}वति~॥ भवति चात्र श्लोकः~। दूरे}\ N}}\pend

\pstart\versepar \ledsidenote{\passref{1.25.29}}\edtext{\skt{भवन्ति}}{\lemma{\skt{भवन्ति}}\Afootnote{N H; \emph{om.}\ A K}}\skt{ }\edtext{\skt{चात्र}}{\lemma{\skt{चात्र}}\Afootnote{N H; \emph{om.}\ A K}}\skt{~॥ }\pend
\pstart\versepar \skt{}\edtext{\skt{एतदष्टविधं}}{\lemma{\skt{एतदष्टविधं}}\Afootnote{A K N; \skt{श्लोकाः~॥ एत\appbrk{}द\appbrk{}ष्ट\appbrk{}वि\appbrk{}धङ्}\ H}}\skt{ कर्म समासेन प्रकीर्तितम्~। }\pend
\pstart\versepar \skt{}\edtext{\skt{चिकित्सितेषु}}{\lemma{\skt{चिकित्सितेषु}}\Afootnote{A K\textsuperscript{pc} N H; \skt{चिकित्सिते}\ K\textsuperscript{ac}}}\skt{ कार्त्स्न्येन }\edtext{\skt{विस्तरस्तस्य}}{\lemma{\skt{विस्तरस्}}\Afootnote{A K H; \skt{विस्तरन्}\ N}}\skt{ वक्ष्यते~॥}\pend

\pstart\versepar \ledsidenote{\passref{1.25.30}}\edtext{\skt{हीनातिरिक्तं}}{\lemma{\skt{हीनातिरिक्तं}}\Afootnote{A K H; \skt{हीनातिरिक्त}\ N}}\skt{ तिर्यक्च }\edtext{\skt{गात्र\-च्छे\-द\-न\-मा\-त्मनः}}{\lemma{\skt{गात्र\appbrk{}च्छे\appbrk{}द\appbrk{}न\appbrk{}मा\appbrk{}त्मनः}}\Afootnote{\skt{गात्र\appbrk{}च्छे\appbrk{}द\appbrk{}न\appbrk{}मा\appbrk{}त्मनः}\ A H K N}}\skt{~। }\pend
\pstart\versepar \skt{}\edtext{\edtext{\edtext{\skt{एता\-श्च\-त\-स्रोऽ\-ष्ट\-विधे}}{\lemma{\skt{एताश्}}\Afootnote{K N H; \skt{एता}\ A}}}{\lemma{\skt{चतस्रो}}\Afootnote{K H; \skt{श्चतस्रो}\ A; \skt{चतस्रो}\ N}}}{\lemma{\skt{ऽष्टविधे}}\Afootnote{A K H; \skt{ष्टविधे}\ N}}\skt{ कर्मणि }\edtext{\skt{व्यापदः}}{\lemma{\skt{व्यापदः}}\Afootnote{A K H; \skt{व्यापद}\ N}}\skt{ स्मृताः~।}\pend

\pstart\versepar \ledsidenote{\passref{1.25.31}}\edtext{\skt{अज्ञा\-न\-लो\-भा\-हि\-त\-वा\-क्य\-योग}}{\lemma{\skt{अज्ञा\appbrk{}न\appbrk{}लो\appbrk{}भा\appbrk{}हि\appbrk{}त\appbrk{}वा\appbrk{}क्य\appbrk{}योग}}\Afootnote{A; \skt{अज्ञा\appbrk{}न\appbrk{}लो\appbrk{}भा\appbrk{}हि\appbrk{}त\appbrk{}वा\appbrk{}क्य\appbrk{}यो\appbrk{}गैर्}\ K N H}}\skt{ }\pend
\pstart\versepar \skt{भय\-प्र\-मो\-है\-र\-प\-रैश्च भावैः~। }\pend
\pstart\versepar \skt{यदा प्रयुञ्जीत भिषक्कुशस्त्रं }\pend
\pstart\versepar \skt{तदा स }\edtext{\skt{शेषान्कुरुते विकारान्}}{\lemma{\skt{शेषान्कुरुते विकारान्}}\Afootnote{A; \skt{भावैः~। वैद्यो यदा शस्त्र\appbrk{}म\appbrk{}भि\appbrk{}प्र\appbrk{}यु\appbrk{}ञ्जन्}\ K\textsuperscript{ac}; \emph{om.}\ K\textsuperscript{pc} N H}}\skt{~॥}\edtext{}{\lemma{\skt{यदा प्रयुञ्जीत भिषक्कुशस्त्रं तदा स}}\Afootnote{A; \skt{वैद्यो यदा शस्त्र\appbrk{}म\appbrk{}भि\appbrk{}प्र\appbrk{}यु\appbrk{}ञ्ज\appbrk{}न्स\appbrk{}शे\appbrk{}ष\appbrk{}दोषां कुरुते विकारान्~॥ एताश्चतस्रो ष्टविधे कर्म्मणि व्यापदः स्मृताः~॥ अज्ञा\appbrk{}न\appbrk{}लो\appbrk{}भा\appbrk{}हि\appbrk{}त\appbrk{}वा\appbrk{}क्य\appbrk{}यो\appbrk{}गै\appbrk{}र्भ\appbrk{}य\appbrk{}प्र\appbrk{}मो\appbrk{}है\appbrk{}र\appbrk{}प\appbrk{}रैश्च}\ K\textsuperscript{ac}; \skt{वैद्यो यदा शस्त्र\appbrk{}म\appbrk{}भि\appbrk{}प्र\appbrk{}यु\appbrk{}ञ्ज\appbrk{}न्स\appbrk{}शे\appbrk{}ष\appbrk{}दोषां कुरुते विकारान्}\ K\textsuperscript{pc}; \skt{वैद्यो यदा शस्त्र\appbrk{}म\appbrk{}भि\appbrk{}प्र\appbrk{}यु\appbrk{}ञ्जन्स शेषदोषां कुरुरुते विकारान्}\ N; \skt{वैद्यो यदा शस्त्र\appbrk{}म\appbrk{}भि\appbrk{}प्र\appbrk{}यु\appbrk{}ञ्ज\appbrk{}न्स\appbrk{}शे\appbrk{}ष\appbrk{}दो\appbrk{}षा\appbrk{}न्कु\appbrk{}रुते विकारान्}\ H}}\pend

\pstart\versepar \ledsidenote{\passref{1.25.32}}\skt{तं }\edtext{\skt{क्षार\-श\-स्त्रा\-ग्नि\-भि\-रौ\-ष\-धैश्च}}{\lemma{\skt{औषधैश्च}}\Afootnote{A K H; \skt{ओषधैश्च}\ N}}\skt{ }\pend
\pstart\versepar \skt{}\edtext{\skt{भूयोऽ\-भि\-यु\-ञ्जा\-न\-म\-यु\-क्ति\-यु\-क्तम्}}{\lemma{\skt{ऽभियुञ्जानम्}}\Afootnote{A; \skt{भियुञ्जन्तम्}\ K N; \skt{ऽभियुञ्जन्तम्}\ H}}\skt{~। }\pend
\pstart\versepar \skt{जिजी\-वि\-षु\-र्दू\-रत एव वैद्यं }\pend
\pstart\versepar \skt{विव\-र्ज\-ये\-दु\-ग्र\-वि\-षा\-हि\-तु\-ल्यम्~॥}\pend

\pstart\versepar \ledsidenote{\passref{1.25.33}}\skt{तदेव युक्तं त्वति\-म\-र्म\-स\-न्धीन् }\pend
\pstart\versepar \skt{हिंस्यात्सिराः }\edtext{\edtext{\skt{स्नायुमथास्थि}}{\lemma{\skt{स्नायुम्}}\Afootnote{K N H; \skt{स्नायु}\ A}}}{\lemma{\skt{अथास्थि}}\Afootnote{K N H; \skt{मथास्थि}\ A}}\skt{ चैव~। }\pend
\pstart\versepar \skt{मूर्खप्रयुक्तं }\edtext{\skt{पुरुषं}}{\lemma{\skt{पुरुषं}}\Afootnote{A K H; \skt{पुरुष}\ N}}\skt{ क्षणेन }\pend
\pstart\versepar \skt{}\edtext{\skt{प्राणै\-र्वि\-यु\-ञ्ज्या\-द\-थवा}}{\lemma{\skt{वियुञ्ज्यादथवा}}\Afootnote{A; \skt{वियुञ्ज्यादथ वा}\ K H; \skt{वियुंज्यदथवा}\ N}}\skt{ }\edtext{\skt{कथंचित्}}{\lemma{\skt{कथंचित्}}\Afootnote{\skt{कदाचित्}\ A N H K}}\skt{~॥}\edtext{}{\lemma{\skt{त्वति\appbrk{}म\appbrk{}र्म\appbrk{}स\appbrk{}न्धी\appbrk{}न्हिं\appbrk{}स्या\appbrk{}त्सिराः}}\Afootnote{A; \skt{त्वतिमर्मसन्धीं हिंस्यात्सिरा}\ K N; \skt{त्वचि मर्म्मसन्धीन्हिंसा सिरा}\ H}}\pend

\pstart\versepar \ledsidenote{\passref{1.25.34}}\skt{भ्रमः प्रलापः पतनं प्रमोहो }\pend
\pstart\versepar \skt{}\edtext{\skt{विचेष्टनं}}{\lemma{\skt{विचेष्टनं}}\Afootnote{A K H; \skt{विवेष्टन}\ N}}\skt{ }\edtext{\edtext{\skt{लीयनमुष्णता}}{\lemma{\skt{लीयनम्}}\Afootnote{K N H; \skt{संलय}\ A}}\skt{ च}}{\lemma{\skt{उष्णता च}}\Afootnote{K N H; \skt{नोष्णते च}\ A}}\skt{~। }\pend
\pstart\versepar \skt{}\edtext{\skt{सुप्ताङ्गता}}{\lemma{\skt{सुप्ताङ्गता}}\Afootnote{K N H; \skt{स्रस्ताङ्गता}\ A}}\skt{ मूर्छ\-न\-मू\-र्ध्व\-वातः~। }\pend
\pstart\versepar \skt{तीव्रा रुजो }\edtext{\skt{वातकृताश्च }\edtext{\skt{तास्ताः}}{\lemma{\skt{ताः}}\Afootnote{K N H; \skt{स्ताः}\ A}}}{\lemma{\skt{वातकृताश्च तास्}}\Afootnote{K N H; \skt{वातकृताश्च ता}\ A}}\skt{~॥}\pend

\pstart \ledsidenote{\passref{1.25.35}}\skt{मांसोदकाभं }\edtext{\skt{रुधिरं च}}{\lemma{\skt{रुधिरं च}}\Afootnote{A N; \skt{रुधितञ्च}\ K H}}\skt{ गच्छे\-त्स\-र्वे\-न्द्रि\-या\-र्थो\-प\-र\-म\-स्त\-थैव~। }\edtext{\skt{दशा\-र्ध\-सं\-ख्ये\-ष्वथ}}{\lemma{\skt{अथ}}\Afootnote{K N H; \skt{अपि}\ A}}\skt{ विक्षतेषु सामान्यतो मर्मसु }\edtext{\edtext{\skt{लिङ्गमुक्तम्}}{\lemma{\skt{लिङ्गम्}}\Afootnote{K N H; \skt{लिङ्ग}\ A}}}{\lemma{\skt{उक्तम्}}\Afootnote{K N H; \skt{मुक्तम्}\ A}}\skt{~॥}\pend

\pstart\versepar \ledsidenote{\passref{1.25.36}}\skt{सुरे\-न्द्र\-गो\-प\-प्र\-तिमं प्रभूतं }\pend
\pstart\versepar \skt{रक्तं }\edtext{\skt{स्रवेद्वै क्षततश्च}}{\lemma{\skt{स्रवेद्वै क्षततश्च}}\Afootnote{A; \skt{स्रवे\appbrk{}त्त\appbrk{}त्क्ष\appbrk{}त\appbrk{}जश्च}\ K H; \skt{श्रवे\appbrk{}त्त\appbrk{}त्क्ष\appbrk{}ण\appbrk{}जञ्च}\ N}}\skt{ वायुः~। }\pend
\pstart\versepar \skt{करोति }\edtext{\skt{रोगा\-न्वि\-वि\-धा\-न्य\-थो\-क्तांश्}}{\lemma{\skt{रोगा\appbrk{}न्वि\appbrk{}वि\appbrk{}धा\appbrk{}न्य\appbrk{}थो\appbrk{}क्तांश्}}\Afootnote{A; \skt{रोगां विविधां यथोक्तां}\ K N; \skt{रोगा\appbrk{}न्वि\appbrk{}वि\appbrk{}धा\appbrk{}न्य\appbrk{}थोक्ता}\ H}}\skt{ }\pend
\pstart\versepar \skt{}\edtext{\skt{छिन्नासु}}{\lemma{\skt{छिन्नासु}}\Afootnote{A; \skt{सिरासु}\ K; \skt{शिरासु}\ N; \skt{न्सिरासु}\ H}}\skt{ }\edtext{\edtext{\skt{भिन्नास्वथवा}}{\lemma{\skt{भिन्नास्व्}}\Afootnote{A; \skt{छिन्नास्व्}\ K H; \skt{च्छिन्ना}\ N}}\skt{ सिरासु}}{\lemma{\skt{अथवा सिरासु}}\Afootnote{A; \skt{अथ विक्षतासु}\ K H; \skt{श्वथ विक्षतासु}\ N}}\skt{~॥}\pend

\pstart\versepar \ledsidenote{\passref{1.25.37}}\edtext{\skt{कौब्ज्यं}}{\lemma{\skt{कौब्ज्यं}}\Afootnote{A; \skt{कौब्जं}\ K H; \skt{कौब्ज}\ N}}\skt{ शरी\-रा\-व\-य\-वा\-व\-सादः }\pend
\pstart\versepar \skt{}\edtext{\edtext{\edtext{\skt{क्रिया\-स्व\-श\-क्ति\-स्तु\-मुला}}{\lemma{\skt{क्रियास्व्}}\Afootnote{K H; \skt{क्रिया}\ A N}}}{\lemma{\skt{अशक्तिस्}}\Afootnote{K H; \skt{स्वशक्ति}\ A; \skt{स्वशक्तिस्}\ N}}}{\lemma{\skt{तुमुला}}\Afootnote{K N H; \skt{स्तुमुला}\ A}}\skt{ रुजश्च~। }\pend
\pstart\versepar \skt{}\edtext{\skt{चिराद्व्रणो}}{\lemma{\skt{चिराद्व्रणो}}\Afootnote{K H; \skt{चिराद्व्र्रणो}\ A; \skt{चिरात्व्रणा}\ N}}\skt{ रोहति यस्य चापि }\pend
\pstart\versepar \skt{तं स्नायुविद्धं मनुजं व्यवस्येत्~॥}\pend

\pstart\versepar \ledsidenote{\passref{1.25.38}}\edtext{\edtext{\skt{शोफा\-ति\-वृ\-द्धि\-स्तु\-मुला}}{\lemma{\skt{शोफातिवृद्धिस्}}\Afootnote{N; \skt{शोफातिवृद्धि}\ A; \skt{शोफाभिवृद्धिस्}\ K H}}}{\lemma{\skt{तुमुला}}\Afootnote{K N H; \skt{स्तुमुला}\ A}}\skt{ रुजश्च }\pend
\pstart\versepar \skt{बलक्षयः }\edtext{\skt{पर्वसु}}{\lemma{\skt{पर्वसु}}\Afootnote{A; \skt{सर्वत}\ K N H\textsuperscript{pc}; \skt{सर्व्व}\ H\textsuperscript{ac}}}\skt{ }\edtext{\skt{भेदशोफौ}}{\lemma{\skt{भेदशोफौ}}\Afootnote{A; \skt{एव शोफः}\ K N H}}\skt{~। }\pend
\pstart\versepar \skt{क्षतेषु }\edtext{\edtext{\skt{सन्धि\-ष्व\-च\-ला\-च\-लेषु}}{\lemma{\skt{सन्धिष्व्}}\Afootnote{K N H; \skt{सन्धि}\ A}}}{\lemma{\skt{अचलाचलेषु}}\Afootnote{K N H; \skt{ष्वचलाचलेषु}\ A}}\skt{ स्यात् }\pend
\pstart\versepar \skt{सन्धि\-क\-र्मो\-प\-र\-तिश्च लिङ्गम्~॥}\edtext{}{\lemma{\skt{स्यात्स\appbrk{}न्धि\appbrk{}क\appbrk{}र्मो\appbrk{}प\appbrk{}र\appbrk{}तिश्च}}\Afootnote{A; \skt{स्यात्स\appbrk{}र्व\appbrk{}क\appbrk{}र्मो\appbrk{}प\appbrk{}र\appbrk{}मश्च}\ K N H}}\pend

\pstart\versepar \ledsidenote{\passref{1.25.39}}\skt{घोरा रुजो यस्य निशादिनेषु }\pend
\pstart\versepar \skt{}\edtext{\edtext{\skt{सर्वास्ववस्थासु}}{\lemma{\skt{सर्वास्व्}}\Afootnote{K N H; \skt{सर्वा}\ A}}\skt{ न}}{\lemma{\skt{अवस्थासु न}}\Afootnote{\skt{स्ववस्थासु न}\ A; \skt{अवस्थासु च}\ K N H}}\skt{ }\edtext{\edtext{\skt{शान्तिरस्ति}}{\lemma{\skt{शान्तिर्}}\Afootnote{\skt{शान्ति}\ A; \skt{नैति}\ K N H}}}{\lemma{\skt{अस्ति}}\Afootnote{\skt{रस्ति}\ A; \emph{om.}\ K H; \skt{शान्तिम्}\ N}}\skt{~। }\pend
\pstart\versepar \skt{तृष्णाऽङ्गसादौ श्वयथुश्च रुक्षः }\pend
\pstart\versepar \skt{तमस्थिविद्धं मनुजं व्यवस्येत्~॥}\edtext{}{\lemma{\skt{तृष्णाऽङ्गसादौ श्वयथुश्च रुक्षः तम्}}\Afootnote{\skt{तृष्णाङ्गसादौ श्वयथुश्च रुक्च तम्}\ A; \skt{शान्ति\appbrk{}म्भि\appbrk{}ष\appbrk{}ग्वि\appbrk{}प\appbrk{}श्चि\appbrk{}द्वि\appbrk{}दि\appbrk{}ता\appbrk{}र्थ\appbrk{}सू\appbrk{}त्र\appbrk{}स्तम्}\ K H; \skt{भिष\appbrk{}ग्वि\appbrk{}प\appbrk{}श्चि\appbrk{}द्भि\appbrk{}वि\appbrk{}दि\appbrk{}ता\appbrk{}र्थ\appbrk{}सू\appbrk{}त्र\appbrk{}न्त\appbrk{}म्तम्}\ N}}\pend

\pstart\versepar \ledsidenote{\passref{1.25.40}}\edtext{\edtext{\skt{यथास्वमेतानि}}{\lemma{\skt{यथास्वम्}}\Afootnote{K N H; \skt{यथास्व}\ A}}}{\lemma{\skt{एतानि}}\Afootnote{K N H; \skt{मेतानि}\ A}}\skt{ }\edtext{\skt{विभावयेयुर्}}{\lemma{\skt{विभावयेयुर्}}\Afootnote{\skt{विभावयेच्च}\ A K N H}}\skt{ }\pend
\pstart\versepar \skt{लिङ्गानि }\edtext{\edtext{\skt{मर्म\-स्व\-भि\-ता\-डि\-तेषु}}{\lemma{\skt{मर्मस्व्}}\Afootnote{K H; \skt{मर्म}\ A N}}}{\lemma{\skt{अभिताडितेषु}}\Afootnote{K H; \skt{स्वभिताडितेषु}\ A N}}\skt{~। }\pend
\pstart\versepar \skt{}\edtext{\skt{स्पर्शं न}}{\lemma{\skt{स्पर्शं न}}\Afootnote{A; \skt{पाण्डुर्विवर्ण्णः स्पर्शन्न}\ K H; \skt{पाण्त्ताविवर्ण्णः स्पर्शन्न}\ N}}\skt{ जानाति विपाण्डुवर्णो }\pend
\pstart\versepar \skt{यो }\edtext{\skt{मांस\-म\-र्म\-ण्य\-भि\-ता\-डितः}}{\lemma{\skt{मांसमर्मण्य्}}\Afootnote{A K\textsuperscript{pc} N H; \skt{मांसमर्माण्य्}\ K\textsuperscript{ac}}}\skt{ }\edtext{\skt{स्यात्}}{\lemma{\skt{स्यात्}}\Afootnote{A; \skt{सः}\ K N H}}\skt{~॥}\edtext{}{\lemma{\skt{जानाति विपाण्डुवर्णो यो}}\Afootnote{A; \skt{वेति यो}\ K N H}}\pend

\pstart\versepar \ledsidenote{\passref{1.25.41}}\edtext{\skt{आत्मानमेवाथ जघन्यकारी}}{\lemma{\skt{आत्मानमेवाथ जघन्यकारी}}\Afootnote{A; \skt{आत्मा\appbrk{}न\appbrk{}मे\appbrk{}वा\appbrk{}ति\appbrk{}ज\appbrk{}घ\appbrk{}न्य\appbrk{}कारी}\ K H; \skt{आत्मा\appbrk{}न\appbrk{}मे\appbrk{}वा\appbrk{}ति\appbrk{}ज\appbrk{}घ\appbrk{}न्य\appbrk{}कारी}\ N}}\skt{ शस्त्रेण यो }\edtext{\skt{हन्ति हि}}{\lemma{\skt{हन्ति हि}}\Afootnote{A; \skt{हिंसति}\ K N H}}\skt{ कर्म }\edtext{\skt{कुर्वन्}}{\lemma{\skt{कुर्वन्}}\Afootnote{A K H; \skt{कुर्व}\ N}}\skt{~। }\pend
\pstart\versepar \skt{}\edtext{\skt{तमा\-त्मा\-वा\-ना\-त्म\-हनं}}{\lemma{\skt{तमा\appbrk{}त्मा\appbrk{}वा\appbrk{}ना\appbrk{}त्म\appbrk{}हनं}}\Afootnote{\skt{तमा\appbrk{}त्म\appbrk{}वा\appbrk{}ना\appbrk{}त्म\appbrk{}हनं}\ A; \skt{तमा\appbrk{}त्म\appbrk{}वा\appbrk{}ना\appbrk{}त्म\appbrk{}नु\appbrk{}दङ्}\ K H; \skt{त्~। तर्का\appbrk{}न्म\appbrk{}वा\appbrk{}ना\appbrk{}न्म\appbrk{}न्त्रदं}\ N}}\skt{ कुवैद्यं }\edtext{\skt{विव\-र्ज\-ये\-दा\-यु\-र\-भी\-प्स\-मानः}}{\lemma{\skt{विव\appbrk{}र्ज\appbrk{}ये\appbrk{}दा\appbrk{}यु\appbrk{}र\appbrk{}भी\appbrk{}प्स\appbrk{}मानः}}\Afootnote{A N; \skt{विव\appbrk{}र्ज\appbrk{}ये\appbrk{}दा\appbrk{}यु\appbrk{}र\appbrk{}भी\appbrk{}प्स\appbrk{}मानः}\ K H}}\skt{~॥}\pend

\pstart\versepar \ledsidenote{\passref{1.25.42}}\skt{तिर्यक्प्रणिहिते शस्त्रे दोषाः }\edtext{\skt{पुर्वमुदाहृताः}}{\lemma{\skt{पुर्वमुदाहृताः}}\Afootnote{\skt{पूर्वमुदाहृताः}\ A; \skt{पुर्वमुदाहृताः}\ K N H}}\skt{~। }\pend
\pstart\versepar \skt{}\edtext{\skt{तस्मा\-त्प\-रि\-ह\-र\-न्दो\-षा\-न्कु\-रु\-या\-च्छ\-स्त्र\-नि\-पा\-त\-नम्}}{\lemma{\skt{दोषा\appbrk{}न्कु\appbrk{}रु\appbrk{}या\appbrk{}च्छ\appbrk{}स्त्र\appbrk{}नि\appbrk{}पा\appbrk{}त\appbrk{}नम्}}\Afootnote{\skt{दोषा\appbrk{}न्कु\appbrk{}र्या\appbrk{}च्छ\appbrk{}स्त्र\appbrk{}नि\appbrk{}पा\appbrk{}त\appbrk{}नम्}\ A H; \skt{हं दोषां कुर्या\appbrk{}च्छ\appbrk{}स्त्र\appbrk{}नि\appbrk{}पा\appbrk{}तनं}\ K\textsuperscript{ac}; \skt{दोषां कुर्या\appbrk{}च्छ\appbrk{}स्त्र\appbrk{}नि\appbrk{}पा\appbrk{}तनं}\ K\textsuperscript{pc}; \skt{दोषां कुरु\appbrk{}या\appbrk{}च्छ\appbrk{}स्त्र\appbrk{}नि\appbrk{}पा\appbrk{}त\appbrk{}नम्}\ N}}\skt{~॥}\pend

\pstart\versepar \ledsidenote{\passref{1.25.43}}\skt{मातरं पितरं }\edtext{\skt{पुत्रा\-न्बा\-न्ध\-वा\-नपि}}{\lemma{\skt{पुत्रा\appbrk{}न्बा\appbrk{}न्ध\appbrk{}वा\appbrk{}नपि}}\Afootnote{A H; \skt{पुत्रा\appbrk{}म्बा\appbrk{}न्ध\appbrk{}वा\appbrk{}नपि}\ K; \skt{पुत्रं बान्धवानपि}\ N}}\skt{ }\edtext{\skt{चातुरः}}{\lemma{\skt{चातुरः}}\Afootnote{A K H; \skt{चान्तरः}\ N}}\skt{~। }\pend
\pstart\versepar \skt{}\edtext{\skt{अप्ये\-ता\-न\-भि\-श\-ङ्केत}}{\lemma{\skt{अप्ये\appbrk{}ता\appbrk{}न\appbrk{}भि\appbrk{}श\appbrk{}ङ्केत}}\Afootnote{A; \skt{अप्येताञ्छङ्कते नित्यं}\ K N H}}\skt{ वैद्ये }\edtext{\skt{विश्वासमेति}}{\lemma{\skt{विश्वासमेति}}\Afootnote{A; \skt{विश्वासमेति}\ K N H}}\skt{ च~॥}\pend

\pstart\versepar \ledsidenote{\passref{1.25.44}}\edtext{\skt{विसृ\-ज\-त्या\-त्म\-ना+\-आ\-त्मानं न}}{\lemma{\skt{विसृ\appbrk{}ज\appbrk{}त्या\appbrk{}त्म\appbrk{}ना+\appbrk{}आ\appbrk{}त्मानं न}}\Afootnote{\skt{विसृ\appbrk{}ज\appbrk{}त्या\appbrk{}त्म\appbrk{}नाऽऽ\appbrk{}त्मानं न}\ A; \skt{विसृ\appbrk{}ज\appbrk{}त्या\appbrk{}त्म\appbrk{}ना\appbrk{}त्मानं न}\ K; \skt{विसृ\appbrk{}जे\appbrk{}त्या\appbrk{}त्म\appbrk{}ना\appbrk{}त्मा\appbrk{}नन्न}\ N; \skt{विसृ\appbrk{}ज\appbrk{}त्या\appbrk{}त्म\appbrk{}ना\appbrk{}त्मान न}\ H}}\skt{ चैनं परिशङ्कते~। }\pend
\pstart\versepar \skt{}\edtext{\skt{तस्मा\-त्पु\-त्र\-व\-दे\-वैनं}}{\lemma{\skt{तस्मा\appbrk{}त्पु\appbrk{}त्र\appbrk{}व\appbrk{}दे\appbrk{}वैनं}}\Afootnote{A; \skt{तस्मा\appbrk{}त्पु\appbrk{}त्र\appbrk{}व\appbrk{}दे\appbrk{}वेह}\ K N H}}\skt{ पालयेदातुरं भिषक्~॥}\pend

\pstart\versepar \ledsidenote{\passref{1.25.46}}\skt{कर्मणा }\edtext{\skt{कश्चिदेकेन}}{\lemma{\skt{कश्चिदेकेन}}\Afootnote{\skt{कश्चिदेकेन}\ A K N H}}\skt{ द्वाभ्यां }\edtext{\skt{कश्चित्त्रिभिस्तथा}}{\lemma{\skt{कश्चित्त्रिभिस्}}\Afootnote{\skt{कश्चित्त्रिभिस्}\ A K H; \skt{कश्चित्त्रिभिस्}\ N}}\skt{~। }\pend
\pstart\versepar \skt{विकारः साध्यते कश्चि\-च्च\-तु\-र्भि\-रपि }\edtext{\skt{कर्मभिः}}{\lemma{\skt{कर्मभिः}}\Afootnote{A K N; \skt{कर्म्मभिः~॥ २५~॥ {\symfont ❈}}\ H}}\skt{~॥}\pend

\pstart \ledsidenote{\passref{1.25.trailer}}\edtext{\skt{इति}}{\lemma{\skt{इति}}\Afootnote{A; \emph{om.}\ N}}\skt{ }\edtext{\skt{सुश्रु\-त\-सं\-हि\-तायां}}{\lemma{\skt{सुश्रु\appbrk{}त\appbrk{}सं\appbrk{}हि\appbrk{}तायां}}\Afootnote{A; \emph{om.}\ N}}\skt{ }\edtext{\edtext{\skt{सूत्र\-स्था\-नेऽ\-ष्ट\-वि\-ध\-श\-स्त्र\-क\-र्मण्यो}}{\lemma{\skt{सूत्रस्थाने}}\Afootnote{A; \emph{om.}\ N}}}{\lemma{\skt{ऽष्ट\appbrk{}वि\appbrk{}ध\appbrk{}श\appbrk{}स्त्र\appbrk{}क\appbrk{}र्मण्यो}}\Afootnote{\skt{ऽष्ट\appbrk{}वि\appbrk{}ध\appbrk{}श\appbrk{}स्त्र\appbrk{}क\appbrk{}र्मीयो}\ A; \emph{om.}\ N}}\skt{ }\edtext{\skt{नाम}}{\lemma{\skt{नाम}}\Afootnote{A; \emph{om.}\ N}}\skt{ }\edtext{\edtext{\skt{पञ्चविंशोऽध्यायः}}{\lemma{\skt{पञ्चविंशो}}\Afootnote{A; \emph{om.}\ N}}}{\lemma{\skt{ऽध्यायः}}\Afootnote{A; \emph{om.}\ N}}\skt{~॥}\pend

\endnumbering
\beginnumbering
\eledsection*{Adhyāya 26}
\pstart \ledsidenote{\passref{1.26.1}}\skt{अथातः }\edtext{\edtext{\skt{प्रन\-ष्ट\-श\-ल्य\-वि\-ज्ञा\-नी\-य\-म\-ध्यायं}}{\lemma{\skt{प्रन\appbrk{}ष्ट\appbrk{}श\appbrk{}ल्य\appbrk{}वि\appbrk{}ज्ञा\appbrk{}नी\appbrk{}यम्}}\Afootnote{A K N; \skt{प्रण\appbrk{}ष्ट\appbrk{}श\appbrk{}ल्य\appbrk{}वि\appbrk{}ज्ञा\appbrk{}नीयं}\ H}}}{\lemma{\skt{अध्यायं}}\Afootnote{A; \emph{om.}\ K; \skt{व्याख्यास्यामः}\ N\textsuperscript{ac} H; \skt{व्याख्यास्यामः~॥ अथातः प्रन\appbrk{}ष्ट\appbrk{}श\appbrk{}ल्य\appbrk{}वि\appbrk{}ज्ञा\appbrk{}नीयं}\ N\textsuperscript{pc}}}\skt{ }\edtext{\skt{व्याख्यास्यामः}}{\lemma{\skt{व्याख्यास्यामः}}\Afootnote{A N\textsuperscript{pc}; \emph{om.}\ K N\textsuperscript{ac} H}}\skt{~॥}\pend

\pstart \ledsidenote{\passref{1.26.3}}\edtext{\skt{अथ}}{\lemma{\skt{अथ}}\Afootnote{K N H; \skt{‘शल’}\ A}}\skt{ }\edtext{\edtext{\edtext{\skt{शल्य\-हिं\-सा\-या\-न्धा\-तु\-स्तस्य}}{\lemma{\skt{शल्यहिंसायान्}}\Afootnote{H; \skt{‘श्वल’ आशुग}\ A; \skt{शल हिंसायां}\ K N}}}{\lemma{\skt{धातुस्}}\Afootnote{K H; \skt{मने धातूः}\ A; \skt{धतुः}\ N}}}{\lemma{\skt{तस्य}}\Afootnote{K N H; \skt{त}\ A}}\skt{ }\edtext{\skt{यत्प्रत्ययस्य}}{\lemma{\skt{यत्प्रत्ययस्य}}\Afootnote{K N H; \skt{योराद्यस्य}\ A}}\skt{ शल्यमिति }\edtext{\skt{रूपम्भवति}}{\lemma{\skt{भवति}}\Afootnote{K N H; \emph{om.}\ A}}\skt{~।}\pend

\pstart \ledsidenote{\passref{1.26.4}}\edtext{\skt{तद्द्विविधं}}{\lemma{\skt{तद्द्विविधं}}\Afootnote{A K H; \skt{तत्द्विविधिं}\ N}}\skt{ }\edtext{\skt{शारी\-र\-मा\-ग\-न्तुकं}}{\lemma{\skt{आगन्तुकं}}\Afootnote{A; \skt{आगन्तुकश्}\ K N H}}\skt{ च~॥}\pend

\pstart \ledsidenote{\passref{1.26.5}}\edtext{\skt{सर्व\-श\-री\-रा\-बा\-ध\-करं शल्यं तदि\-हो\-प\-दि\-श्यत}}{\lemma{\skt{सर्व\appbrk{}श\appbrk{}री\appbrk{}रा\appbrk{}बा\appbrk{}ध\appbrk{}करं शल्यं तदि\appbrk{}हो\appbrk{}प\appbrk{}दि\appbrk{}श्यत}}\Afootnote{A; \skt{शारी\appbrk{}रा\appbrk{}ग\appbrk{}न्त\appbrk{}वश्च सर्व\appbrk{}श\appbrk{}री\appbrk{}रा\appbrk{}बाधा इहोपदृश्यन्त}\ K; \skt{शरी\appbrk{}रा\appbrk{}ग\appbrk{}ब\appbrk{}न्धश्च सर्व\appbrk{}श\appbrk{}री\appbrk{}रा\appbrk{}बाधा इहोपदृश्यन्त}\ N; \skt{सर्व्व\appbrk{}श\appbrk{}री\appbrk{}रा\appbrk{}बाधा इहोपदृश्यंत}\ H}}\skt{ इत्यतः शल्यशास्त्रम्~॥}\pend

\pstart \ledsidenote{\passref{1.26.6}}\skt{तत्र }\edtext{\skt{शारीराणि}}{\lemma{\skt{शारीराणि}}\Afootnote{K N H; \skt{शारीरं}\ A}}\skt{ }\edtext{\skt{दन्त\-न\-ख\-रो\-मा\-दीनि}}{\lemma{\skt{दन्त\appbrk{}न\appbrk{}ख\appbrk{}रो\appbrk{}मा\appbrk{}दीनि}}\Afootnote{K N H; \skt{दन्त\appbrk{}रो\appbrk{}म\appbrk{}न\appbrk{}खादि}\ A}}\skt{~। धातवो }\edtext{\skt{न}}{\lemma{\skt{न}}\Afootnote{\skt{ऽन्न}\ A N H; \skt{न्नं}\ K}}\skt{ मला दोषाश्च दुष्टाः }\edtext{\skt{आगन्तूनि}}{\lemma{\skt{आगन्तूनि}}\Afootnote{K H; \skt{आगन्त्वपि}\ A; \skt{आगंत्रूनि}\ N}}\skt{ }\edtext{\skt{शरी\-र\-श\-ल्य\-व्य\-ति\-रे\-केण}}{\lemma{\skt{शरी\appbrk{}र\appbrk{}श\appbrk{}ल्य\appbrk{}व्य\appbrk{}ति\appbrk{}रे\appbrk{}केण}}\Afootnote{K H; \skt{शारी\appbrk{}र\appbrk{}श\appbrk{}ल्य\appbrk{}व्य\appbrk{}ति\appbrk{}रे\appbrk{}केण}\ A N}}\skt{ यावन्तो भावाः }\edtext{\skt{दुःख\-मु\-त्पा\-द\-यन्ति}}{\lemma{\skt{दुःखम्}}\Afootnote{A K\textsuperscript{pc} N H; \skt{दुःखं दुःखम्}\ K\textsuperscript{ac}}}\skt{~॥}\pend

\pstart\versepar \ledsidenote{\passref{1.26.6.add}}\edtext{\skt{भवति}}{\lemma{\skt{भवति}}\Afootnote{N; \emph{om.}\ K H}}\skt{ }\edtext{\skt{चात्र}}{\lemma{\skt{चात्र}}\Afootnote{N; \emph{om.}\ K H}}\skt{~॥ }\pend
\pstart\versepar \skt{शरीरे सर्वशल्यानां गतयः पञ्चधा }\edtext{\skt{स्मृताः}}{\lemma{\skt{स्मृताः}}\Afootnote{K H; \skt{स्मृताः}\ N}}\skt{~। }\pend
\pstart\versepar \skt{}\edtext{\edtext{\skt{ऋज्वा\-ग\-त\-म\-वा\-ञ्चीनं}}{\lemma{\skt{ऋज्वागतम्}}\Afootnote{K H; \skt{दज्धारात}\ N}}}{\lemma{\skt{अवाञ्चीनं}}\Afootnote{K H; \skt{मवाचीनं}\ N}}\skt{ }\edtext{\edtext{\skt{तिर्य\-गू\-र्ध्व\-म\-धो\-गतं}}{\lemma{\skt{तिर्यगूर्ध्वम्}}\Afootnote{K H; \skt{तिर्यगूर्द्व}\ N}}}{\lemma{\skt{अधोगतं}}\Afootnote{K H; \skt{मधोगतम्}\ N}}\skt{~॥}\pend

\pstart \ledsidenote{\passref{1.26.9}}\edtext{\skt{तानि तु}}{\lemma{\skt{तानि तु}}\Afootnote{K N H; \skt{तानि}\ A}}\skt{ }\edtext{\skt{यदा}}{\lemma{\skt{यदा}}\Afootnote{K H; \emph{om.}\ A; \skt{यादा}\ N}}\skt{ }\edtext{\skt{वेग\-क्ष\-या\-त्प्र\-ति\-घा\-ताद्वा}}{\lemma{\skt{वेग\appbrk{}क्ष\appbrk{}या\appbrk{}त्प्र\appbrk{}ति\appbrk{}घा\appbrk{}ताद्वा}}\Afootnote{A; \skt{वेग\appbrk{}क्ष\appbrk{}या\appbrk{}त्प्र\appbrk{}ती\appbrk{}घा\appbrk{}ताद्वा}\ K N H}}\skt{ त्वगादिषु }\edtext{\edtext{\skt{व्रण\-व\-स्तु\-ष्व\-व\-ति\-ष्ठन्ते}}{\lemma{\skt{व्रणवस्तुष्व्}}\Afootnote{K N H; \skt{व्रणवस्तु}\ A}}}{\lemma{\skt{अवतिष्ठन्ते}}\Afootnote{H; \skt{ष्वतिष्ठन्ते}\ A; \skt{अतिष्ठन्ते}\ K; \skt{तिष्ठन्ते}\ N}}\skt{ }\edtext{\skt{धम\-नी\-स्रो\-तोऽ\-स्थि\-वि\-व\-र\-पे\-शी\-प्र\-भृ\-तिषु वा}}{\lemma{\skt{धम\appbrk{}नी\appbrk{}स्रो\appbrk{}तोऽ\appbrk{}स्थि\appbrk{}वि\appbrk{}व\appbrk{}र\appbrk{}पे\appbrk{}शी\appbrk{}प्र\appbrk{}भृ\appbrk{}तिषु वा}}\Afootnote{A; \skt{धम\appbrk{}नी\appbrk{}स्रो\appbrk{}तो\appbrk{}स्थि\appbrk{}पे\appbrk{}सी\appbrk{}वि\appbrk{}व\appbrk{}र\appbrk{}प्र\appbrk{}भृ\appbrk{}तिषु वा}\ K H; \skt{धम\appbrk{}नी\appbrk{}स्रो\appbrk{}तो\appbrk{}स्थि\appbrk{}पे\appbrk{}शी\appbrk{}वि\appbrk{}व\appbrk{}र\appbrk{}प्र\appbrk{}भृ\appbrk{}तिषु वा}\ N}}\skt{ शरी\-र\-प्र\-दे\-शेषु~॥}\pend

\pstart \ledsidenote{\passref{1.26.10}}\edtext{\skt{तत्र शल्य\-ल\-क्ष\-ण\-मु\-च्य\-मा\-न\-मु\-प\-धा\-रय}}{\lemma{\skt{तत्र शल्यलक्षणम्}}\Afootnote{A; \skt{तत्र लक्षणम्}\ K H; \skt{तत्तु लक्षणम्}\ N}}\skt{~। तत्तु द्विविधं सामान्यं वैशेषिकं च~। श्यावं }\edtext{\skt{पिडकाचितं}}{\lemma{\skt{पिडकाचितं}}\Afootnote{A; \skt{पिटकाचितं}\ K H; \skt{पिढकाचित}\ N}}\skt{ शोफ\-वे\-द\-ना\-वन्तं }\edtext{\skt{मुहुर्मुहुः}}{\lemma{\skt{मुहुर्मुहुः}}\Afootnote{K N H; \skt{मुहुर्मुहुः}\ A}}\skt{ शोणितास्राविणं }\edtext{\skt{बुद्बु\-द\-व\-दु\-न्नतं}}{\lemma{\skt{बुद्बु\appbrk{}द\appbrk{}व\appbrk{}दु\appbrk{}न्नतं}}\Afootnote{A; \skt{बुद्बु\appbrk{}द\appbrk{}व\appbrk{}दु\appbrk{}द्ग\appbrk{}तम्}\ K H N}}\skt{ }\edtext{\skt{मृदुमांसं च}}{\lemma{\skt{मृदुमांसं च}}\Afootnote{A N H; \skt{मृदुमांसंञ्च}\ K}}\skt{ }\edtext{\skt{व्रणं}}{\lemma{\skt{व्रणं}}\Afootnote{A K H; \skt{व्रण}\ N}}\skt{ }\edtext{\edtext{\skt{जानी\-या\-त्स\-श\-ल्योऽ\-य\-मिति}}{\lemma{\skt{जानीयात्सशल्यो}}\Afootnote{A K H; \skt{जानी\appbrk{}या\appbrk{}त्स\appbrk{}श\appbrk{}ल्यम्}\ N}}}{\lemma{\skt{ऽयम्}}\Afootnote{A; \skt{यम्}\ K H; \emph{om.}\ N}}\skt{~। }\edtext{\skt{सामा\-न्य\-मे\-त\-ल्ल\-क्ष\-ण\-मु\-क्तम्}}{\lemma{\skt{सामा\appbrk{}न्य\appbrk{}मे\appbrk{}त\appbrk{}ल्ल\appbrk{}क्ष\appbrk{}णम्}}\Afootnote{A; \skt{सामान्यतो लक्षणम्}\ K N H}}\skt{~। }\edtext{\skt{वैशेषिकं तु}}{\lemma{\skt{वैशेषिकं तु}}\Afootnote{A; \skt{विशेषस्}\ K N; \skt{विशेषतस्}\ H}}\skt{ त्वग्गते }\edtext{\skt{विवर्णः}}{\lemma{\skt{विवर्णः}}\Afootnote{A H; \skt{विवर्ण्णाः}\ K N}}\skt{ शोफो }\edtext{\skt{भवत्यायतः}}{\lemma{\skt{भवत्य्}}\Afootnote{A N; \skt{भवन्त्य्}\ K H}}\skt{ कठिनश्च मांसगते }\edtext{\skt{शोफाभिवृद्धिः}}{\lemma{\skt{शोफाभिवृद्धिः}}\Afootnote{H; \skt{शोफाभि(ति)वृद्धिः}\ A; \skt{शोफाभिवृ}\ K N}}\skt{ }\edtext{\skt{शल्य\-मा\-र्गा\-नु\-प\-सं\-रोहः}}{\lemma{\skt{शल्य\appbrk{}मा\appbrk{}र्गा\appbrk{}नु\appbrk{}प\appbrk{}सं\appbrk{}रोहः}}\Afootnote{A; \skt{द्धिरसंरोहः}\ K N; \skt{असंरोहः}\ H}}\skt{ }\edtext{\skt{पीड\-ना\-स\-हि\-ष्णुता}}{\lemma{\skt{पीड\appbrk{}ना\appbrk{}स\appbrk{}हि\appbrk{}ष्णुता}}\Afootnote{A K\textsuperscript{pc} H; \skt{पीडनासहिता}\ K\textsuperscript{ac}; \skt{पीड्या\appbrk{}ना\appbrk{}श\appbrk{}ष्ण\appbrk{}हिता}\ N}}\skt{ }\edtext{\skt{चोषपाकौ च}}{\lemma{\skt{चोषपाकौ च}}\Afootnote{A; \skt{चोषप्रपाकौ च}\ K N H}}\skt{ }\edtext{\edtext{\edtext{\skt{पेश्य\-न्त\-र\-स्थेऽ\-प्ये\-त\-देव}}{\lemma{\skt{पेश्यन्तरस्थे}}\Afootnote{A; \skt{मांसा\appbrk{}ना\appbrk{}म्पे\appbrk{}श्य\appbrk{}न्त\appbrk{}र\appbrk{}गते}\ K H; \skt{मान्सानाम्पश्यन्त ग}\ N}}}{\lemma{\skt{ऽप्य्}}\Afootnote{A H; \skt{प्य्}\ K; \emph{om.}\ N}}}{\lemma{\skt{एतदेव}}\Afootnote{A; \skt{एवं}\ K; \skt{तेप्येवं}\ N; \skt{एव}\ H}}\skt{ }\edtext{\skt{चोषशोफवर्जं}}{\lemma{\skt{चोषशोफवर्जं}}\Afootnote{A; \skt{चोषशोषवर्ज्जं}\ K H; \skt{चोषसोषवर्जं}\ N}}\skt{ सिरागते }\edtext{\edtext{\skt{सिरा\-ध्मा\-न\-म्सि\-रा\-शूलं}}{\lemma{\skt{सिराध्मानम्}}\Afootnote{A; \emph{om.}\ K N; \skt{चाध्मानं}\ H}}}{\lemma{\skt{सिराशूलं}}\Afootnote{A; \skt{चाध्मानं}\ K N; \emph{om.}\ H}}\skt{ }\edtext{\skt{सिराशोफश्च}}{\lemma{\skt{सिराशोफश्च}}\Afootnote{A; \skt{शूलञ्च}\ K H; \skt{चलञ्च}\ N}}\skt{ स्नायुगते }\edtext{\skt{स्नायु\-जा\-लो\-त्क्षे\-पणं}}{\lemma{\skt{स्नायु\appbrk{}जा\appbrk{}लो\appbrk{}त्क्षे\appbrk{}पणं}}\Afootnote{A; \skt{स्नायु\appbrk{}जा\appbrk{}ला\appbrk{}व\appbrk{}क्षे\appbrk{}पणं}\ K; \skt{स्नायु\appbrk{}जा\appbrk{}ल\appbrk{}व\appbrk{}क्षे\appbrk{}प\appbrk{}णम्}\ N; \skt{स्नायु\appbrk{}जा\appbrk{}लाऽ\appbrk{}सु\appbrk{}क्षे\appbrk{}पणं}\ H}}\skt{ }\edtext{\edtext{\skt{संरम्भश्चोग्रा}}{\lemma{\skt{संरम्भश्}}\Afootnote{A; \emph{om.}\ K H; \skt{स्रोत्ये}\ N}}\skt{ रुक्च स्रोतोगते}}{\lemma{\skt{चोग्रा रुक्च स्रोतोगते}}\Afootnote{A; \skt{स्रोतोगते}\ K H; \skt{गतो}\ N}}\skt{ }\edtext{\skt{स्रोतसां}}{\lemma{\skt{स्रोतसां}}\Afootnote{A N H; \skt{स्रोतमीसांश्}\ K\textsuperscript{ac}; \skt{स्रोतसांश्}\ K\textsuperscript{pc}}}\skt{ }\edtext{\skt{स्वक\-र्म\-गु\-ण\-हानिः;}}{\lemma{\skt{स्वक\appbrk{}र्म\appbrk{}गु\appbrk{}ण\appbrk{}हानिः}}\Afootnote{A N; \skt{च कर्मगुणहानिः}\ K H}}\skt{ }\edtext{\skt{धमनीस्थे}}{\lemma{\skt{धमनीस्थे}}\Afootnote{A; \skt{धमनीगते}\ K N H}}\skt{ सफेनं }\edtext{\edtext{\skt{रक्त\-मी\-र\-य\-न्न\-निलः}}{\lemma{\skt{रक्तम्}}\Afootnote{A K\textsuperscript{pc} N H; \skt{सरक्तम्}\ K\textsuperscript{ac}}}}{\lemma{\skt{ईरयन्ननिलः}}\Afootnote{A; \skt{ईर\appbrk{}य\appbrk{}न्स\appbrk{}मी\appbrk{}रणः}\ K H N}}\skt{ सशब्दो }\edtext{\skt{निर्ग\-च्छ\-त्य\-ङ्ग\-मर्दः पिपासा}}{\lemma{\skt{अङ्गमर्दः पिपासा}}\Afootnote{A; \skt{अङ्ग\appbrk{}म\appbrk{}र्द\appbrk{}पि\appbrk{}पासा}\ K N H}}\skt{ हृल्लासश्च अस्थिगते विवि\-ध\-वे\-द\-ना\-प्रा\-दु\-र्भावः }\edtext{\skt{शोफश्च}}{\lemma{\skt{शोफश्च}}\Afootnote{A; \skt{सशोफश्च}\ K H; \skt{सशोफाश्च}\ N}}\skt{ }\edtext{\skt{अस्थि\-वि\-व\-र\-ग\-तेऽ\-स्थि\-पू\-र्ण\-ताऽ\-स्थि\-तोदः संहर्षो}}{\lemma{\skt{ऽस्थि\appbrk{}पू\appbrk{}र्ण\appbrk{}ताऽ\appbrk{}स्थि\appbrk{}तोदः संहर्षो}}\Afootnote{A; \skt{त्वस्थि\appbrk{}चू\appbrk{}र्ण्ण\appbrk{}पू\appbrk{}र्ण्णता संहर्षो}\ K; \skt{त्वस्थि\appbrk{}चू\appbrk{}र्ण्ण\appbrk{}पू\appbrk{}र्ण्णता सहर्षो}\ N; \skt{त्वस्थि\appbrk{}चू\appbrk{}र्ण्ण\appbrk{}पू\appbrk{}र्ण्णतां संहर्षो}\ H}}\skt{ }\edtext{\skt{बलवांश्च}}{\lemma{\skt{बलवांश्च}}\Afootnote{A; \skt{वेगवांश्च}\ K H; \skt{वेगावाश्च}\ N}}\skt{ }\edtext{\edtext{\skt{सन्धि\-ग\-तेऽ\-स्थि\-व\-च्चे\-ष्टो\-प\-र\-मश्च}}{\lemma{\skt{सन्धिगते}}\Afootnote{A K H; \skt{सन्धिग}\ N}}}{\lemma{\skt{ऽस्थि\appbrk{}व\appbrk{}च्चे\appbrk{}ष्टो\appbrk{}प\appbrk{}र\appbrk{}मश्च}}\Afootnote{A; \skt{प्येवं चेष्टोपरमश्च}\ K; \skt{तांस्ये\appbrk{}व\appbrk{}चे\appbrk{}ष्टो\appbrk{}प\appbrk{}र\appbrk{}मश्च}\ N; \skt{ऽप्येवं चेष्टोपरमश्च}\ H}}\skt{ }\edtext{\skt{कोष्ठगत}}{\lemma{\skt{कोष्ठगत}}\Afootnote{A; \skt{कोष्ठगते}\ K H N}}\skt{ }\edtext{\skt{आटोपानाहौ}}{\lemma{\skt{आटोपानाहौ}}\Afootnote{A; \skt{त्वाटोपाना}\ K; \skt{आटोपानाहो}\ N; \skt{त्वाटोपानाहौ}\ H}}\skt{ }\edtext{\skt{मूत्र\-पु\-री\-षा\-हा\-र\-द\-र्शनं च}}{\lemma{\skt{मूत्र\appbrk{}पु\appbrk{}री\appbrk{}षा\appbrk{}हा\appbrk{}र\appbrk{}द\appbrk{}र्शनं च}}\Afootnote{A; \skt{हो पुरी\appbrk{}शा\appbrk{}हा\appbrk{}र\appbrk{}द\appbrk{}र्श\appbrk{}नानि च}\ K; \skt{पुरी\appbrk{}षा\appbrk{}हा\appbrk{}र\appbrk{}द\appbrk{}र्श\appbrk{}नानि च}\ N H}}\skt{ }\edtext{\skt{व्रण\-मु\-खा\-त्म\-र्म\-गते}}{\lemma{\skt{व्रण\appbrk{}मु\appbrk{}खा\appbrk{}त्म\appbrk{}र्म\appbrk{}गते}}\Afootnote{A; \skt{व्रण\appbrk{}मु\appbrk{}खा\appbrk{}द्भ\appbrk{}वन्ति मर्मगते तु}\ K N H}}\skt{ }\edtext{\skt{मर्म\-वि\-द्धा\-व\-च्चे\-ष्टते}}{\lemma{\skt{मर्म\appbrk{}वि\appbrk{}द्धा\appbrk{}व\appbrk{}च्चे\appbrk{}ष्टते}}\Afootnote{A; \skt{मर्म\appbrk{}वि\appbrk{}द्ध\appbrk{}व\appbrk{}दा\appbrk{}च\appbrk{}क्षते}\ K H; \skt{मर्म\appbrk{}वि\appbrk{}द्ध\appbrk{}व\appbrk{}द\appbrk{}च\appbrk{}क्षते}\ N}}\skt{~। सूक्ष्मगतिषु }\edtext{\edtext{\skt{शल्येष्वेतान्येव}}{\lemma{\skt{शल्येष्व्}}\Afootnote{A K H; \skt{शल्ये}\ N}}}{\lemma{\skt{एतान्य्}}\Afootnote{A K H; \skt{ष्वेतान्य्}\ N}}\skt{ }\edtext{\skt{लक्ष\-णा\-न्य\-स्प\-ष्टानि}}{\lemma{\skt{अस्पष्टानि}}\Afootnote{A; \skt{विस्पष्टानि}\ K H; \skt{विस्पष्टितानि}\ N}}\skt{ भवन्ति~॥}\edtext{}{\lemma{\skt{उपधारय~। तत्तु द्विविधं सामान्यं वैशेषिकं च}}\Afootnote{\skt{उपधारय~। तद्द्विविधं -- सामान्यं, वैशेषिकं च}\ A; \skt{उपधारयस्व}\ K H; \skt{उपधारयस्र}\ N}}\pend

\pstart \ledsidenote{\passref{1.26.11}}\skt{महान्त्यल्पानि वा }\edtext{\edtext{\skt{शुद्ध\-दे\-हा\-ना\-म\-नु\-लो\-म\-स\-न्नि\-वि\-ष्टानि}}{\lemma{\skt{शुद्धदेहानाम्}}\Afootnote{A N H; \skt{शुद्धदे}\ K}}}{\lemma{\skt{अनु\appbrk{}लो\appbrk{}म\appbrk{}स\appbrk{}न्नि\appbrk{}वि\appbrk{}ष्टानि}}\Afootnote{A N H; \emph{om.}\ K}}\skt{ }\edtext{\skt{रोहन्ति}}{\lemma{\skt{रोहन्ति}}\Afootnote{A; \emph{om.}\ K; \skt{उपरोहन्ति}\ N H}}\skt{ }\edtext{\skt{विशेषतः}}{\lemma{\skt{विशेषतः}}\Afootnote{A N H; \emph{om.}\ K}}\skt{ }\edtext{\skt{कण्ठ\-स्रो\-तः\-सि\-रा\-त्व\-क्पे\-श्य\-स्थि\-वि\-व\-रेषु}}{\lemma{\skt{कण्ठ\appbrk{}स्रो\appbrk{}तः\appbrk{}सि\appbrk{}रा\appbrk{}त्व\appbrk{}क्पे\appbrk{}श्य\appbrk{}स्थि\appbrk{}वि\appbrk{}व\appbrk{}रेषु}}\Afootnote{\skt{कण्ठ\appbrk{}स्रो\appbrk{}तः\appbrk{}सि\appbrk{}रा\appbrk{}त्व\appbrk{}क्पे\appbrk{}श्य\appbrk{}स्थि\appbrk{}वि\appbrk{}व\appbrk{}रेषः}\ A; \emph{om.}\ K; \skt{कण्ठ\appbrk{}ग\appbrk{}त\appbrk{}स्रो\appbrk{}तः\appbrk{}सि\appbrk{}रा\appbrk{}त्व\appbrk{}क्पे\appbrk{}श्य\appbrk{}स्थि\appbrk{}वि\appbrk{}व\appbrk{}रेषु}\ N H}}\skt{ }\edtext{\skt{दोष\-प्र\-को\-प\-व्या\-या\-मा\-भि\-घा\-ता\-जी\-र्णेभ्यः प्रचलितानि }\edtext{\skt{पुनर्बाधन्ते}}{\lemma{\skt{बाधन्ते}}\Afootnote{A; \emph{om.}\ K; \skt{राबाधकं भवेत्}\ N; \skt{आबा\appbrk{}ध\appbrk{}क\appbrk{}म्भ\appbrk{}वेत्}\ H}}}{\lemma{\skt{दोष\appbrk{}प्र\appbrk{}को\appbrk{}प\appbrk{}व्या\appbrk{}या\appbrk{}मा\appbrk{}भि\appbrk{}घा\appbrk{}ता\appbrk{}जी\appbrk{}र्णेभ्यः प्रचलितानि पुनर्}}\Afootnote{A; \emph{om.}\ K; \skt{प्रको\appbrk{}प\appbrk{}व्या\appbrk{}या\appbrk{}मा\appbrk{}द\appbrk{}भि\appbrk{}घा\appbrk{}तात्~। क्षतादपि स्थानां तु चलितं शल्यं पुन}\ N; \skt{दोष\appbrk{}प्र\appbrk{}को\appbrk{}प\appbrk{}व्या\appbrk{}या\appbrk{}मा\appbrk{}द\appbrk{}भि\appbrk{}घा\appbrk{}तात्~। क्षतादपि स्थानात्तु चलितं~। शल्यम्पुनर्}\ H}}\skt{~॥}\pend

\pstart \ledsidenote{\passref{1.26.12}}\skt{तत्र }\edtext{\skt{त्वक्प्रनष्टे}}{\lemma{\skt{त्वक्प्रनष्टे}}\Afootnote{A N; \skt{त्वक्प्रणष्टे}\ H}}\skt{ }\edtext{\skt{स्निग्धस्विन्नायां}}{\lemma{\skt{स्निग्धस्विन्नायां}}\Afootnote{A; \emph{om.}\ N H}}\skt{ }\edtext{\skt{मृन्मा\-ष\-य\-व\-गो\-धू\-म\-गो\-म\-य\-मृ\-दि\-तायां}}{\lemma{\skt{मृन्मा\appbrk{}ष\appbrk{}य\appbrk{}व\appbrk{}गो\appbrk{}धू\appbrk{}म\appbrk{}गो\appbrk{}म\appbrk{}य\appbrk{}मृ\appbrk{}दि\appbrk{}तायां}}\Afootnote{A; \skt{मृद्गो\appbrk{}म\appbrk{}य\appbrk{}चू\appbrk{}र्ण्ण\appbrk{}प्र\appbrk{}मृ\appbrk{}दि\appbrk{}तायां}\ N H}}\skt{ त्वचि यत्र संरम्भो }\edtext{\skt{वेदना वा}}{\lemma{\skt{वेदना वा}}\Afootnote{A H; \skt{वेदनावान्}\ N}}\skt{ भवति }\edtext{\skt{तत्र}}{\lemma{\skt{तत्र}}\Afootnote{A H; \emph{om.}\ N}}\skt{ शल्यं }\edtext{\skt{विजा\-नी\-या\-त्स्त्या\-न\-घृ\-त\-मृ\-च्च\-न्द\-न\-क\-ल्कैर्वा}}{\lemma{\skt{विजा\appbrk{}नी\appbrk{}या\appbrk{}त्स्त्या\appbrk{}न\appbrk{}घृ\appbrk{}त\appbrk{}मृ\appbrk{}च्च\appbrk{}न्द\appbrk{}न\appbrk{}क\appbrk{}ल्कैर्वा}}\Afootnote{A; \skt{इति जानीयात्~। तत्र स्त्यान\appbrk{}घृ\appbrk{}त\appbrk{}मृ\appbrk{}त्च\appbrk{}न्द\appbrk{}न\appbrk{}क\appbrk{}ल्कैर्वा}\ N; \skt{इति जानीयात्~। तत्र स्त्यान\appbrk{}घृ\appbrk{}त\appbrk{}मृ\appbrk{}च्च\appbrk{}न्द\appbrk{}न\appbrk{}क\appbrk{}ल्कैर्व्वा}\ H}}\skt{ }\edtext{\skt{प्रदिग्धायां}}{\lemma{\skt{प्रदिग्धायां}}\Afootnote{A; \skt{प्रदिग्धे}\ N H}}\skt{ }\edtext{\skt{शल्योष्मणा+आशु}}{\lemma{\skt{शल्योष्मणा+आशु}}\Afootnote{\skt{शल्योष्मणाऽऽशु}\ A H; \skt{सल्योष्मणाशु}\ N}}\skt{ }\edtext{\skt{विसरति}}{\lemma{\skt{विसरति}}\Afootnote{A H; \skt{र्चिसरति}\ N}}\skt{ घृत\-मु\-प\-शु\-ष्यति }\edtext{\skt{वा लेपो}}{\lemma{\skt{वा लेपो}}\Afootnote{\skt{चालेपो}\ A H N}}\skt{ यत्र तत्र }\edtext{\skt{शल्यं}}{\lemma{\skt{शल्यं}}\Afootnote{A H; \skt{शल्य}\ N}}\skt{ }\edtext{\skt{विजा\-नी\-या\-त्मां\-स\-प्र\-नष्टे}}{\lemma{\skt{विजा\appbrk{}नी\appbrk{}या\appbrk{}त्मां\appbrk{}स\appbrk{}प्र\appbrk{}नष्टे}}\Afootnote{A; \skt{जानी\appbrk{}या\appbrk{}त्मा\appbrk{}न्स\appbrk{}प्र\appbrk{}णष्टे}\ N H}}\skt{ }\edtext{\skt{स्नेहस्वेदादिभिः}}{\lemma{\skt{स्नेहस्वेदादिभिः}}\Afootnote{A; \skt{स्नेहादिभिर्}\ N H}}\skt{ }\edtext{\edtext{\edtext{\skt{क्रिया\-वि\-शे\-षै\-र\-वि\-रु\-द्दै\-रा\-तु\-र\-मु\-प\-पा\-द\-ये\-त्क\-र्शि\-तस्य}}{\lemma{\skt{क्रियाविशेषैर्}}\Afootnote{A; \emph{om.}\ N H}}}{\lemma{\skt{अविरुद्दैर्}}\Afootnote{\skt{अविरुद्धैर्}\ A; \emph{om.}\ N H}}\skt{ तु}}{\lemma{\skt{उप\appbrk{}पा\appbrk{}द\appbrk{}ये\appbrk{}त्क\appbrk{}र्शि\appbrk{}तस्य तु}}\Afootnote{A; \skt{उपपादयेत्~। कर्षिततु}\ N; \skt{उपपादयेत्~। कर्षितस्य तु}\ H}}\skt{ }\edtext{\skt{शिथि\-ली\-भू\-त\-म\-न\-व\-बद्धं}}{\lemma{\skt{शिथिलीभूतम्}}\Afootnote{A; \skt{शिथिलीकृतम्}\ N H}}\skt{ }\edtext{\skt{क्षुभ्यमाणं}}{\lemma{\skt{क्षुभ्यमाणं}}\Afootnote{A N; \skt{क्षुभ्यमानं}\ H}}\skt{ यत्र }\edtext{\skt{संरम्भं}}{\lemma{\skt{संरम्भं}}\Afootnote{A; \skt{सरम्भी}\ N; \skt{संरम्भी}\ H}}\skt{ }\edtext{\skt{वेदनां वा}}{\lemma{\skt{वेदनां वा}}\Afootnote{A H; \skt{वेदनां}\ N}}\skt{ जनयति तत्र शल्यं }\edtext{\edtext{\edtext{\skt{विजा\-नी\-या\-त्को\-ष्ठा\-स्थि\-स\-न्धि\-पे\-शी\-वि\-व\-रे\-ष्व\-व\-स्थि\-त\-मे\-व\-मेव}}{\lemma{\skt{विजा\appbrk{}नी\appbrk{}या\appbrk{}त्को\appbrk{}ष्ठा\appbrk{}स्थि\appbrk{}स\appbrk{}न्धि\appbrk{}पे\appbrk{}शी\appbrk{}वि\appbrk{}व\appbrk{}रेष्व्}}\Afootnote{A; \skt{जानी\appbrk{}या\appbrk{}त्का\appbrk{}ष्ठा\appbrk{}स्थि\appbrk{}स\appbrk{}न्धि\appbrk{}पे\appbrk{}सी\appbrk{}वि\appbrk{}व\appbrk{}रेष्व्}\ N; \skt{जानी\appbrk{}या\appbrk{}त्को\appbrk{}ष्ठा\appbrk{}स्थि\appbrk{}स\appbrk{}न्धि\appbrk{}पे\appbrk{}शी\appbrk{}वि\appbrk{}व\appbrk{}रेष्व्}\ H}}}{\lemma{\skt{एवम्}}\Afootnote{A H; \emph{om.}\ N}}}{\lemma{\skt{एव}}\Afootnote{A N; \emph{om.}\ H}}\skt{ }\edtext{\skt{परीक्षेत}}{\lemma{\skt{परीक्षेत}}\Afootnote{A N; \skt{परीक्षे}\ H}}\skt{ }\edtext{\skt{सिरा\-ध\-म\-नी\-स्रो\-तः\-स्ना\-यु\-प्र\-नष्टे}}{\lemma{\skt{सिरा\appbrk{}ध\appbrk{}म\appbrk{}नी\appbrk{}स्रो\appbrk{}तः\appbrk{}स्ना\appbrk{}यु\appbrk{}प्र\appbrk{}नष्टे}}\Afootnote{A; \skt{शिरा\appbrk{}ध\appbrk{}म\appbrk{}नी\appbrk{}स्रो\appbrk{}तस्सु च~। स्नायुप्रणष्टे}\ N; \skt{त्~। सिरा\appbrk{}ध\appbrk{}म\appbrk{}नी\appbrk{}स्रो\appbrk{}तःसु च~॥ स्नायुप्रणष्टे}\ H}}\skt{ खण्ड\-च\-क्र\-सं\-युक्ते }\edtext{\skt{याने}}{\lemma{\skt{याने}}\Afootnote{A H; \skt{या}\ N}}\skt{ }\edtext{\edtext{\skt{व्याधि\-त\-मा\-रो\-प्याशु}}{\lemma{\skt{व्याधितम्}}\Afootnote{\skt{व्याधित}\ A; \skt{त्वातुररम्}\ N; \skt{त्वातुरम्}\ H}}}{\lemma{\skt{आरोप्याशु}}\Afootnote{N H; \skt{मारोप्याशु}\ A}}\skt{ }\edtext{\edtext{\skt{विषमेऽध्वनि}}{\lemma{\skt{विषमे}}\Afootnote{A H; \skt{र्विषमे}\ N}}}{\lemma{\skt{ऽध्वनि}}\Afootnote{A; \skt{ध्वनि}\ N H}}\skt{ यायाद्यत्र संरम्भो वेदना वा }\edtext{\skt{भवति}}{\lemma{\skt{भवति}}\Afootnote{A H; \skt{न्भवति}\ N}}\skt{ तत्र शल्यं }\edtext{\skt{जानी\-या\-त\-स्थि\-प्र\-नष्टे}}{\lemma{\skt{जानी\appbrk{}या\appbrk{}त\appbrk{}स्थि\appbrk{}प्र\appbrk{}नष्टे}}\Afootnote{N; \skt{विजानीयात् ; अस्थिप्रनष्टे}\ A; \skt{जानीयात्~॥ अस्थिप्रणष्टे}\ H}}\skt{ }\edtext{\skt{स्नेह\-स्वे\-दो\-प\-प\-न्ना\-न्य\-स्थीनि}}{\lemma{\skt{स्नेह\appbrk{}स्वे\appbrk{}दो\appbrk{}प\appbrk{}प\appbrk{}न्ना\appbrk{}न्य\appbrk{}स्थीनि}}\Afootnote{A N; \skt{स्नेह\appbrk{}स्वे\appbrk{}दो\appbrk{}प\appbrk{}प\appbrk{}न्ना\appbrk{}न्य\appbrk{}स्थीनि}\ H}}\skt{ बन्ध\-न\-पी\-ड\-नाभ्यां }\edtext{\skt{भृश\-मु\-पा\-च\-रे\-द्यत्र}}{\lemma{\skt{भृशम्}}\Afootnote{N H; \skt{भृश\appbrk{}मु\appbrk{}पा\appbrk{}च\appbrk{}रेत्, यत्र संरम्भो वेदना भवति तत्र शल्यं विजानीयात् ; सन्धिप्रनष्टे स्नेह\appbrk{}स्वे\appbrk{}दो\appbrk{}प\appbrk{}प\appbrk{}न्ना\appbrk{}न्स\appbrk{}न्धी\appbrk{}न्प्र\appbrk{}स\appbrk{}र\appbrk{}णा\appbrk{}कु\appbrk{}ञ्च\appbrk{}न\appbrk{}ब\appbrk{}न्ध\appbrk{}न\appbrk{}पी\appbrk{}डनैः भृशम्}\ A}}\skt{ संरम्भो वेदना वा }\edtext{\skt{भवति}}{\lemma{\skt{भवति}}\Afootnote{A H; \skt{न्भवति}\ N}}\skt{ तत्र शल्यं }\edtext{\skt{जानी\-या\-त्स\-न्धि\-प्र\-नष्टे}}{\lemma{\skt{जानी\appbrk{}या\appbrk{}त्स\appbrk{}न्धि\appbrk{}प्र\appbrk{}नष्टे}}\Afootnote{N; \skt{विजानीयात् , सन्धिप्रनष्टे}\ A; \skt{जानीयात्~॥ सन्धिप्रणष्टे}\ H}}\skt{ }\edtext{\edtext{\skt{स्नेह\-स्वे\-दो\-प\-प\-न्ना\-न्स\-न्धी\-न्प्र\-स\-र\-णा\-कु\-ञ्च\-न\-ब\-न्ध\-न\-पी\-ड\-नै\-र्भृ\-श\-मु\-पा\-च\-रे\-त्यत्र}}{\lemma{\skt{स्नेह\appbrk{}स्वे\appbrk{}दो\appbrk{}प\appbrk{}प\appbrk{}न्ना\appbrk{}न्स\appbrk{}न्धी\appbrk{}न्प्र\appbrk{}स\appbrk{}र\appbrk{}णा\appbrk{}कु\appbrk{}ञ्च\appbrk{}न\appbrk{}ब\appbrk{}न्ध\appbrk{}न\appbrk{}पी\appbrk{}ड\appbrk{}नैर्}}\Afootnote{\skt{स्नेह\appbrk{}स्वे\appbrk{}दो\appbrk{}प\appbrk{}प\appbrk{}न्ना\appbrk{}न्स\appbrk{}न्धी\appbrk{}न्प्र\appbrk{}स\appbrk{}र\appbrk{}णा\appbrk{}कु\appbrk{}ञ्च\appbrk{}न\appbrk{}ब\appbrk{}न्ध\appbrk{}न\appbrk{}पी\appbrk{}डनै}\ A; \skt{स्नेह\appbrk{}स्वे\appbrk{}दो\appbrk{}प\appbrk{}पन्नां सन्धी\appbrk{}प्र\appbrk{}सा\appbrk{}र\appbrk{}णा\appbrk{}कु\appbrk{}ञ्च\appbrk{}न\appbrk{}ब\appbrk{}न्ध\appbrk{}न\appbrk{}पी\appbrk{}डनैः}\ N; \skt{स्नेह\appbrk{}स्वे\appbrk{}दो\appbrk{}प\appbrk{}प\appbrk{}न्ना\appbrk{}न्स\appbrk{}न्धी\appbrk{}न्प्र\appbrk{}सा\appbrk{}र\appbrk{}णा\appbrk{}कु\appbrk{}ञ्च\appbrk{}न\appbrk{}ब\appbrk{}न्ध\appbrk{}पी\appbrk{}ड\appbrk{}नैर्}\ H}}}{\lemma{\skt{भृशम्}}\Afootnote{N H; \skt{र्भृशम्}\ A}}\skt{ संरम्भो वेदना वा }\edtext{\skt{भवति}}{\lemma{\skt{भवति}}\Afootnote{A H; \skt{न्भवति}\ N}}\skt{ तत्र शल्यं }\edtext{\skt{विजा\-नी\-या\-त्म\-र्म\-प्र\-नष्टे}}{\lemma{\skt{विजा\appbrk{}नी\appbrk{}या\appbrk{}त्म\appbrk{}र्म\appbrk{}प्र\appbrk{}नष्टे}}\Afootnote{A N; \skt{जानीयात्~॥ मर्म्मप्रणष्टे}\ H}}\skt{ }\edtext{\edtext{\skt{त्वन\-न्य\-भा\-वा\-न्म\-र्म\-णा\-मुक्तं}}{\lemma{\skt{त्वन\appbrk{}न्य\appbrk{}भा\appbrk{}वा\appbrk{}न्म\appbrk{}र्म\appbrk{}णाम्}}\Afootnote{A H; \skt{त्वन\appbrk{}न्य\appbrk{}भा\appbrk{}वा\appbrk{}त्म\appbrk{}र्मणा}\ N}}}{\lemma{\skt{उक्तं}}\Afootnote{A H; \skt{मुक्तं}\ N}}\skt{ परीक्षणं भवति~॥}\pend

\pstart \ledsidenote{\passref{1.26.13}}\edtext{\edtext{\skt{सामा\-न्य\-ल\-क्ष\-ण\-मपि}}{\lemma{\skt{सामा\appbrk{}न्य\appbrk{}ल\appbrk{}क्ष\appbrk{}णम्}}\Afootnote{N H; \skt{सामान्यलक्षम्}\ A}}\skt{ च}}{\lemma{\skt{अपि च}}\Afootnote{A; \skt{अपि}\ N H}}\skt{ }\edtext{\skt{हस्ति\-स्क\-न्धा\-श्व\-पृ\-ष्ठ\-प\-र्व\-त\-द्रु\-मा\-रो\-ह\-ण\-ध\-नु\-र्व्या\-या\-म\-द्रु\-त\-या\-न\-नि\-यु\-द्धा\-ध्व\-ग\-म\-न\-ल\-ङ्घ\-न\-प्र\-त\-र\-ण\-प्ल\-व\-न\-व्या\-या\-मै\-र्जृ\-म्भो\-द्गा\-र\-का\-स\-क्ष\-व\-थु\-ष्ठी\-व\-न\-ह\-स\-न\-प्रा\-णा\-या\-मै\-र्वा\-त\-मू\-त्र\-पु\-री\-ष\-शु\-क्रो\-त्स\-र्गैर्वा}}{\lemma{\skt{हस्ति\appbrk{}स्क\appbrk{}न्धा\appbrk{}श्व\appbrk{}पृ\appbrk{}ष्ठ\appbrk{}प\appbrk{}र्व\appbrk{}त\appbrk{}द्रु\appbrk{}मा\appbrk{}रो\appbrk{}ह\appbrk{}ण\appbrk{}ध\appbrk{}नु\appbrk{}र्व्या\appbrk{}या\appbrk{}म\appbrk{}द्रु\appbrk{}त\appbrk{}या\appbrk{}न\appbrk{}नि\appbrk{}यु\appbrk{}द्धा\appbrk{}ध्व\appbrk{}ग\appbrk{}म\appbrk{}न\appbrk{}ल\appbrk{}ङ्घ\appbrk{}न\appbrk{}प्र\appbrk{}त\appbrk{}र\appbrk{}ण\appbrk{}प्ल\appbrk{}व\appbrk{}न\appbrk{}व्या\appbrk{}या\appbrk{}मै\appbrk{}र्जृ\appbrk{}म्भो\appbrk{}द्गा\appbrk{}र\appbrk{}का\appbrk{}स\appbrk{}क्ष\appbrk{}व\appbrk{}थु\appbrk{}ष्ठी\appbrk{}व\appbrk{}न\appbrk{}ह\appbrk{}स\appbrk{}न\appbrk{}प्रा\appbrk{}णा\appbrk{}या\appbrk{}मैर्}}\Afootnote{\skt{हस्ति\appbrk{}स्क\appbrk{}न्धा\appbrk{}श्व\appbrk{}पृ\appbrk{}ष्ठ\appbrk{}प\appbrk{}र्व\appbrk{}त\appbrk{}द्रु\appbrk{}मा\appbrk{}रो\appbrk{}ह\appbrk{}ण\appbrk{}ध\appbrk{}नु\appbrk{}र्व्या\appbrk{}या\appbrk{}म\appbrk{}द्रु\appbrk{}त\appbrk{}या\appbrk{}न\appbrk{}नि\appbrk{}यु\appbrk{}द्धा\appbrk{}ध्व\appbrk{}ग\appbrk{}म\appbrk{}न\appbrk{}ल\appbrk{}ङ्घ\appbrk{}न\appbrk{}प्ल\appbrk{}व\appbrk{}न\appbrk{}प्र\appbrk{}त\appbrk{}रण- व्याया\appbrk{}मै\appbrk{}र्जृ\appbrk{}म्भो\appbrk{}द्गा\appbrk{}र\appbrk{}का\appbrk{}स\appbrk{}क्ष\appbrk{}व\appbrk{}थु\appbrk{}ष्ठी\appbrk{}व\appbrk{}न\appbrk{}ह\appbrk{}स\appbrk{}न\appbrk{}प्रा\appbrk{}णा\appbrk{}या\appbrk{}मैर्}\ A; \skt{हस्ति\appbrk{}स्क\appbrk{}न्धा\appbrk{}श्व\appbrk{}पृ\appbrk{}ष्ठ\appbrk{}प\appbrk{}र्व\appbrk{}त\appbrk{}द्रु\appbrk{}मा\appbrk{}रो\appbrk{}ह\appbrk{}ण\appbrk{}द्रु\appbrk{}त\appbrk{}या\appbrk{}न\appbrk{}ल\appbrk{}ङ्घ\appbrk{}न\appbrk{}प्ल\appbrk{}व\appbrk{}न\appbrk{}व्या\appbrk{}या\appbrk{}म\appbrk{}जृ\appbrk{}म्भो\appbrk{}द्ग\appbrk{}र\appbrk{}क्षा\appbrk{}व\appbrk{}षु\appbrk{}ष्ठी\appbrk{}व\appbrk{}न\appbrk{}ह\appbrk{}स\appbrk{}न\appbrk{}प्रा\appbrk{}णा\appbrk{}या\appbrk{}मैर्}\ N; \skt{हस्ति\appbrk{}स्क\appbrk{}न्धा\appbrk{}श्व\appbrk{}पृ\appbrk{}ष्ठ\appbrk{}प\appbrk{}र्व्व\appbrk{}त\appbrk{}द्रु\appbrk{}मा\appbrk{}रो\appbrk{}ह\appbrk{}ण\appbrk{}द्रु\appbrk{}त\appbrk{}या\appbrk{}न\appbrk{}लं\appbrk{}घ\appbrk{}न\appbrk{}प्ल\appbrk{}व\appbrk{}न\appbrk{}व्या\appbrk{}या\appbrk{}म\appbrk{}जृ\appbrk{}म्भो\appbrk{}ङ्गा\appbrk{}र\appbrk{}क्ष\appbrk{}व\appbrk{}थु\appbrk{}ष्ठी\appbrk{}व\appbrk{}न\appbrk{}ह\appbrk{}स\appbrk{}न\appbrk{}प्रा\appbrk{}णा\appbrk{}या\appbrk{}मैर्}\ H}}\skt{ यत्र }\edtext{\skt{संरम्भो}}{\lemma{\skt{संरम्भो}}\Afootnote{A; \emph{om.}\ N H}}\skt{ }\edtext{\skt{वेदना वा}}{\lemma{\skt{वेदना वा}}\Afootnote{A N; \skt{वेदना}\ H}}\skt{ भवति तत्र शल्यं }\edtext{\skt{जानीयात्}}{\lemma{\skt{जानीयात्}}\Afootnote{N H; \skt{विजानीयात्}\ A}}\skt{~॥}\pend

\pstart\versepar \ledsidenote{\passref{1.26.15}}\edtext{\skt{भवन्ति}}{\lemma{\skt{भवन्ति}}\Afootnote{N H; \emph{om.}\ A}}\skt{ }\edtext{\skt{चात्र}}{\lemma{\skt{चात्र}}\Afootnote{N H; \emph{om.}\ A}}\skt{~। }\pend
\pstart\versepar \skt{}\edtext{\skt{आतुरश्चापि}}{\lemma{\skt{आतुरश्}}\Afootnote{A N; \skt{श्लोकाः~॥ आतुरश्}\ H}}\skt{ यं देशमभीक्ष्णं परिरक्षति~। }\pend
\pstart\versepar \skt{संवाह्यमानो बहुशस्तत्र शल्यं विनिर्दिशेत्~॥}\pend

\pstart\versepar \ledsidenote{\passref{1.26.16}}\edtext{\edtext{\skt{अल्पा\-बा\-ध\-म\-शूनं}}{\lemma{\skt{अल्पाबाधम्}}\Afootnote{A; \skt{अल्पताम}\ N H}}\skt{ च}}{\lemma{\skt{अशूनं च}}\Afootnote{A; \skt{पिसूनञ्च}\ N; \skt{पशूनञ्च}\ H}}\skt{ }\edtext{\skt{नीरुजं}}{\lemma{\skt{नीरुजं}}\Afootnote{A; \skt{निरुजन्}\ N H}}\skt{ }\edtext{\skt{निरुपद्रवम्}}{\lemma{\skt{निरुपद्रवम्}}\Afootnote{A H; \skt{निरुपद्रवां}\ N}}\skt{~। }\pend
\pstart\versepar \skt{}\edtext{\skt{प्रसन्नं}}{\lemma{\skt{प्रसन्नं}}\Afootnote{A; \skt{प्रसन्न}\ N H}}\skt{ मृदुपर्यन्तं }\edtext{\skt{निरा\-घ\-ट्ट\-म\-नु\-न्न\-तम्}}{\lemma{\skt{निरा\appbrk{}घ\appbrk{}ट्ट\appbrk{}म\appbrk{}नु\appbrk{}न्न\appbrk{}तम्}}\Afootnote{A; \skt{निरा\appbrk{}घ\appbrk{}ट्ट\appbrk{}म\appbrk{}वे\appbrk{}द\appbrk{}नम्}\ N; \skt{निरा\appbrk{}घ\appbrk{}टू\appbrk{}म\appbrk{}वे\appbrk{}दनं}\ H}}\skt{~॥}\pend

\pstart\versepar \ledsidenote{\passref{1.26.17}}\skt{एषण्या सर्वतो दृष्ट्वा }\edtext{\skt{यथामार्गं}}{\lemma{\skt{यथामार्गं}}\Afootnote{A; \skt{यथालाभं}\ N H}}\skt{ चिकित्सकः~। }\pend
\pstart\versepar \skt{}\edtext{\edtext{\skt{प्रसा\-रा\-कु\-ञ्च\-ना\-न्नूनं}}{\lemma{\skt{प्रसा\appbrk{}रा\appbrk{}कु\appbrk{}ञ्च\appbrk{}नान्}}\Afootnote{A; \skt{प्रसराकुञ्चने}\ N; \skt{प्रसाराकुञ्चने}\ H}}}{\lemma{\skt{नूनं}}\Afootnote{A N; \skt{न्यूनन्}\ H}}\skt{ }\edtext{\skt{निःशल्यमिति}}{\lemma{\skt{निःशल्यम्}}\Afootnote{A H; \skt{निशल्यम्}\ N}}\skt{ }\edtext{\skt{निर्दिशेत्}}{\lemma{\skt{निर्दिशेत्}}\Afootnote{A H; \skt{निच्चिशेत्}\ N}}\skt{~॥}\pend

\pstart\versepar \ledsidenote{\passref{1.26.18}}\skt{अस्थ्यात्मकं }\edtext{\skt{भज्यते तु}}{\lemma{\skt{भज्यते तु}}\Afootnote{A; \skt{भिद्यते तु}\ N H}}\skt{ शल्यमन्तश्च शीर्यते~। }\pend
\pstart\versepar \skt{प्रायो निर्भुज्यते }\edtext{\skt{शार्ङ्गमायसं}}{\lemma{\skt{शार्ङ्गम्}}\Afootnote{A H; \skt{शाङ्गम्}\ N}}\skt{ चेति निश्चयः~॥}\pend

\pstart\versepar \ledsidenote{\passref{1.26.19}}\edtext{\skt{वार्क्ष\-वै\-ण\-व\-ता\-र्णानि}}{\lemma{\skt{वार्क्ष\appbrk{}वै\appbrk{}ण\appbrk{}व\appbrk{}ता\appbrk{}र्णानि}}\Afootnote{A; \skt{वार्क्षञ्च वैणवं चैव ताल चाप्य्}\ N; \skt{वार्क्षञ्च वैणवञ्चैव तालजञ्चाप्य्}\ H}}\skt{ }\edtext{\skt{निर्ह्रयन्ते तु नो}}{\lemma{\skt{निर्ह्रयन्ते तु नो}}\Afootnote{\skt{निर्ह्रियन्ते तु नो}\ A; \skt{अनिर्हृतम्}\ N H}}\skt{ }\edtext{\skt{यदि}}{\lemma{\skt{यदि}}\Afootnote{A; \emph{om.}\ N H}}\skt{~। }\pend
\pstart\versepar \skt{पचन्ति }\edtext{\skt{रक्तं}}{\lemma{\skt{रक्तं}}\Afootnote{A H; \skt{रक्त}\ N}}\skt{ मांसं च क्षिप्रमेतानि देहिनाम्~॥}\pend

\pstart\versepar \ledsidenote{\passref{1.26.20}}\edtext{\skt{कानकं}}{\lemma{\skt{कानकं}}\Afootnote{A; \skt{कनकं}\ N H}}\skt{ }\edtext{\skt{राजतं}}{\lemma{\skt{राजतं}}\Afootnote{A; \skt{रजतं}\ N H}}\skt{ ताम्रं }\edtext{\skt{रैतिकं}}{\lemma{\skt{रैतिकं}}\Afootnote{A; \skt{कृष्णाय}\ N H}}\skt{ }\edtext{\skt{त्रपुसीसकम्}}{\lemma{\skt{त्रपुसीसकम्}}\Afootnote{A N; \skt{स्त्रपुसीसकं}\ H}}\skt{~। }\pend
\pstart\versepar \skt{}\edtext{\skt{चिर\-स्था\-ना\-द्वि\-ली\-यन्ते}}{\lemma{\skt{चिर\appbrk{}स्था\appbrk{}ना\appbrk{}द्वि\appbrk{}ली\appbrk{}यन्ते}}\Afootnote{N H; \skt{चिर\appbrk{}स्था\appbrk{}ना\appbrk{}द्वि\appbrk{}ली\appbrk{}यन्ते}\ A}}\skt{ पित्त\-ते\-जः\-प्र\-ता\-प\-नात्~॥}\pend

\pstart\versepar \ledsidenote{\passref{1.26.21}}\skt{स्वभावशीता मृदवो ये }\edtext{\skt{चान्येऽपीदृशा}}{\lemma{\skt{ऽपीदृशा}}\Afootnote{A; \skt{प्येवमा}\ N; \skt{ऽप्येवमा}\ H}}\skt{ }\edtext{\skt{मताः}}{\lemma{\skt{मताः}}\Afootnote{A; \skt{दयः}\ N H}}\skt{~। }\pend
\pstart\versepar \skt{}\edtext{\skt{द्रवीभूताः}}{\lemma{\skt{द्रवीभूताः}}\Afootnote{A N; \skt{द्रवीभूता}\ H}}\skt{ }\edtext{\skt{शरी\-रेऽ\-स्मि\-न्ने\-कत्वं}}{\lemma{\skt{ऽस्मिन्नेकत्वं}}\Afootnote{A H; \skt{ऽस्मिन्नेकत्व}\ N}}\skt{ यान्ति }\edtext{\skt{धातुभिः}}{\lemma{\skt{धातुभिः}}\Afootnote{A H; \skt{देहिभिः}\ N}}\skt{~॥}\pend

\pstart\versepar \ledsidenote{\passref{1.26.22}}\edtext{\skt{विषा\-ण\-द\-न्त\-के\-शा\-स्थि\-वे\-णु\-दा\-रू\-प\-लानि तु}}{\lemma{\skt{विषा\appbrk{}ण\appbrk{}द\appbrk{}न्त\appbrk{}के\appbrk{}शा\appbrk{}स्थि\appbrk{}वे\appbrk{}णु\appbrk{}दा\appbrk{}रू\appbrk{}प\appbrk{}लानि तु}}\Afootnote{H; \skt{विषा\appbrk{}ण\appbrk{}द\appbrk{}न्त\appbrk{}के\appbrk{}शा\appbrk{}स्थि\appbrk{}वे\appbrk{}णु\appbrk{}दा\appbrk{}रु\appbrk{}प\appbrk{}लानि तु}\ A; \skt{विषा\appbrk{}ण\appbrk{}द\appbrk{}न्त\appbrk{}के\appbrk{}सा\appbrk{}स्थि\appbrk{}वे\appbrk{}णु\appbrk{}दा\appbrk{}रो\appbrk{}प\appbrk{}लानि तु}\ N}}\skt{~। }\pend
\pstart\versepar \skt{शल्यानि न विशीर्यन्ते शरीरे }\edtext{\skt{मृन्मयानि च}}{\lemma{\skt{मृन्मयानि च}}\Afootnote{A; \skt{मृन्मयानि चेति}\ N H}}\skt{~॥}\pend

\pstart \ledsidenote{\passref{1.26.trailer}}\edtext{\skt{॥ २६~॥}}{\lemma{\skt{॥ २६~॥}}\Afootnote{\skt{इति सुश्रु\appbrk{}त\appbrk{}सं\appbrk{}हि\appbrk{}तायां सूत्रस्थाने प्रन\appbrk{}ष्ट\appbrk{}श\appbrk{}ल्य\appbrk{}वि\appbrk{}ज्ञा\appbrk{}नीयो नाम षड्विं\appbrk{}श\appbrk{}ति\appbrk{}त\appbrk{}मोऽ\appbrk{}ध्यायः}\ A; \skt{२६~॥}\ N H}}\pend

\endnumbering
\beginnumbering
\eledsection*{Adhyāya 27}
\pstart \ledsidenote{\passref{1.27.1}}\skt{अथातः }\edtext{\edtext{\skt{शल्या\-प\-न\-य\-नी\-य\-म\-ध्यायं}}{\lemma{\skt{शल्या\appbrk{}प\appbrk{}न\appbrk{}य\appbrk{}नी\appbrk{}यम्}}\Afootnote{A; \skt{शल्यापनयनी}\ N H}}}{\lemma{\skt{अध्यायं}}\Afootnote{A; \skt{यं}\ N H}}\skt{ व्याख्यास्यामः~॥}\pend

\pstart \ledsidenote{\passref{1.27.4}}\skt{तत्र }\edtext{\skt{समा\-से\-ना\-न\-व\-ब\-द्ध\-श\-ल्यो\-द्ध\-र\-णार्थं}}{\lemma{\skt{समा\appbrk{}से\appbrk{}ना\appbrk{}न\appbrk{}व\appbrk{}ब\appbrk{}द्ध\appbrk{}श\appbrk{}ल्यो\appbrk{}द्ध\appbrk{}र\appbrk{}णार्थं}}\Afootnote{A; \skt{समासेन त्वव\appbrk{}ब\appbrk{}द्ध\appbrk{}श\appbrk{}ल्यो\appbrk{}द्ध\appbrk{}र\appbrk{}णा\appbrk{}र्थान्}\ N; \skt{समासेन त्वव\appbrk{}ब\appbrk{}द्ध\appbrk{}श\appbrk{}ल्यो\appbrk{}द्ध\appbrk{}र\appbrk{}णा\appbrk{}र्थन्}\ H}}\skt{ }\edtext{\skt{पञ्चदश}}{\lemma{\skt{पञ्चदश}}\Afootnote{A; \skt{नव}\ N H}}\skt{ हेतून्वक्ष्यामः~। }\edtext{\skt{तद्यथा}}{\lemma{\skt{तद्यथा}}\Afootnote{A N; \skt{तद्यथा}\ H}}\skt{ स्वभावः }\edtext{\skt{पाचनं भेदनं दारणं पीडनं}}{\lemma{\skt{पाचनं भेदनं दारणं पीडनं}}\Afootnote{A; \skt{पाचनम्}\ N H}}\skt{ }\edtext{\skt{प्रमार्जनं}}{\lemma{\skt{प्रमार्जनं}}\Afootnote{A H; \skt{प्रमार्जन}\ N}}\skt{ }\edtext{\skt{निर्ध्मापनं वमनं विरेचनं}}{\lemma{\skt{निर्ध्मापनं वमनं विरेचनं}}\Afootnote{A; \skt{निर्धमन}\ N; \skt{निर्द्ध}\ H\textsuperscript{ac}; \skt{निर्द्धमनम्}\ H\textsuperscript{pc}}}\skt{ }\edtext{\skt{प्रक्षालनं}}{\lemma{\skt{प्रक्षालनं}}\Afootnote{A; \emph{om.}\ N H\textsuperscript{pc}; \skt{मनम्}\ H\textsuperscript{ac}}}\skt{ }\edtext{\skt{प्रतिमर्शः}}{\lemma{\skt{प्रतिमर्शः}}\Afootnote{A; \skt{प्रतिमर्षः}\ N H}}\skt{ प्रवाहणं आचूषणं }\edtext{\skt{अयस्कान्तो}}{\lemma{\skt{अयस्कान्तो}}\Afootnote{A; \skt{अय\appbrk{}स्का\appbrk{}न्ता\appbrk{}क\appbrk{}र्षणं}\ N H\textsuperscript{ac}; \skt{अयस्कान्ता हर्षणम्}\ H\textsuperscript{pc}}}\skt{ }\edtext{\skt{हर्षश्चेति}}{\lemma{\skt{हर्षश्चेति}}\Afootnote{A; \skt{हर्षणमिति}\ N H\textsuperscript{ac}; \skt{इति}\ H\textsuperscript{pc}}}\skt{~॥}\pend

\pstart \ledsidenote{\passref{1.27.5}}\edtext{\skt{तत्रा\-श्रु\-क्ष\-व\-थू\-द्गा\-र\-का\-स\-मू\-त्र\-पु\-री\-षा\-निलैः}}{\lemma{\skt{तत्रा\appbrk{}श्रु\appbrk{}क्ष\appbrk{}व\appbrk{}थू\appbrk{}द्गा\appbrk{}र\appbrk{}का\appbrk{}स\appbrk{}मू\appbrk{}त्र\appbrk{}पु\appbrk{}री\appbrk{}षा\appbrk{}निलैः}}\Afootnote{A; \skt{तत्र श्रुक्ष\appbrk{}व\appbrk{}थु\appbrk{}का\appbrk{}सो\appbrk{}द्गा\appbrk{}र\appbrk{}मू\appbrk{}त्र\appbrk{}पु\appbrk{}री\appbrk{}षा\appbrk{}निलैः}\ N; \skt{तत्रा\appbrk{}श्रु\appbrk{}क्ष\appbrk{}व\appbrk{}थु\appbrk{}का\appbrk{}सो\appbrk{}द्गा\appbrk{}र\appbrk{}मू\appbrk{}त्र\appbrk{}पु\appbrk{}री\appbrk{}षा\appbrk{}निलैः}\ H}}\skt{ }\edtext{\skt{स्वभा\-व\-ब\-ल\-प्र\-वृ\-त्तै\-र्न\-य\-ना\-दिभ्यः}}{\lemma{\skt{स्वभा\appbrk{}व\appbrk{}ब\appbrk{}ल\appbrk{}प्र\appbrk{}वृ\appbrk{}त्तैर्}}\Afootnote{A; \skt{स्वभावप्रवृत्तै}\ N; \skt{स्वभा\appbrk{}व\appbrk{}प्र\appbrk{}वृ\appbrk{}त्तैर्}\ H}}\skt{ }\edtext{\skt{पतति}}{\lemma{\skt{पतति}}\Afootnote{A; \skt{प्रपतन्ति}\ N H}}\skt{ }\edtext{\skt{मांसावगाढं}}{\lemma{\skt{मांसावगाढं}}\Afootnote{A; \skt{एष स्वभावो नाम मान्सावगाढम्}\ N H}}\skt{ }\edtext{\skt{शल्य\-म\-वि\-द\-ह्य\-मानं}}{\lemma{\skt{शल्य\appbrk{}म\appbrk{}वि\appbrk{}द\appbrk{}ह्य\appbrk{}मानं}}\Afootnote{A; \skt{अम\appbrk{}वि\appbrk{}द\appbrk{}ह्य\appbrk{}मानं}\ N; \skt{अवि\appbrk{}द\appbrk{}ह्य\appbrk{}मा\appbrk{}नम्}\ H}}\skt{ पाचयित्वा }\edtext{\skt{प्रको\-था(\-आ.\-पा)\-त्तस्य पूय\-शो\-णि\-त\-वे\-गा\-द्गौ\-र\-वाद्वा}}{\lemma{\skt{प्रको\appbrk{}था(\appbrk{}आ.\appbrk{}पा)\appbrk{}त्तस्य पूय\appbrk{}शो\appbrk{}णि\appbrk{}त\appbrk{}वे\appbrk{}गा\appbrk{}द्गौ\appbrk{}र\appbrk{}वाद्वा}}\Afootnote{\skt{प्रकोथात्तस्य पूय\appbrk{}शो\appbrk{}णि\appbrk{}त\appbrk{}वे\appbrk{}गा\appbrk{}द्गौ\appbrk{}र\appbrk{}वाद्वा}\ A; \skt{प्रकोथयेत्~। पूयसोणि}\ N; \skt{प्रकोथयेत्~। पूय\appbrk{}शो\appbrk{}णि\appbrk{}त\appbrk{}वे\appbrk{}गा\appbrk{}ङ्गौ\appbrk{}र\appbrk{}वाद्वा}\ H}}\skt{ }\edtext{\skt{पतति}}{\lemma{\skt{पतति}}\Afootnote{A H; \skt{त}\ N}}\skt{~। }\edtext{\skt{पक्व\-म\-भि\-द्य\-मानं भेद\-ये\-द्दा\-र\-येद्वा~। भिन्न\-म\-नि\-र\-स्य\-मानं}}{\lemma{\skt{पक्व\appbrk{}म\appbrk{}भि\appbrk{}द्य\appbrk{}मानं भेद\appbrk{}ये\appbrk{}द्दा\appbrk{}र\appbrk{}येद्वा~। भिन्न\appbrk{}म\appbrk{}नि\appbrk{}र\appbrk{}स्य\appbrk{}मानं}}\Afootnote{A; \skt{वैस्वा\appbrk{}ङ्गै\appbrk{}र\appbrk{}चा\appbrk{}धा\appbrk{}यत ति ए}\ N; \skt{ए}\ H\textsuperscript{ac}; \skt{एतत्पाचनन्}\ H\textsuperscript{pc}}}\skt{ }\edtext{\skt{पीडनीयैः पीड\-ये\-त्पा\-णि\-भिर्वा~। }\edtext{\skt{अणू\-न्य\-क्ष\-श\-ल्यानि}}{\lemma{\skt{अक्षशल्यानि}}\Afootnote{\skt{न्यक्षशल्यानि}\ A; \skt{अक्षिशल्यानि}\ N H}}}{\lemma{\skt{पीडनीयैः पीड\appbrk{}ये\appbrk{}त्पा\appbrk{}णि\appbrk{}भिर्वा~। अणून्य्}}\Afootnote{\skt{पीडनीयैः पीड\appbrk{}ये\appbrk{}त्पा\appbrk{}णि\appbrk{}भिर्वा~। अणू}\ A; \skt{तत्पाचनन्नाम अणून्य्}\ N H\textsuperscript{ac}; \skt{नाम~॥ अणून्य्}\ H\textsuperscript{pc}}}\skt{ }\edtext{\skt{परि\-षे\-च\-ना\-ध्मा\-प\-नै\-र्बा\-ल\-व\-स्त्र\-पा\-णिभिः}}{\lemma{\skt{परि\appbrk{}षे\appbrk{}च\appbrk{}ना\appbrk{}ध्मा\appbrk{}प\appbrk{}नै\appbrk{}र्बा\appbrk{}ल\appbrk{}व\appbrk{}स्त्र\appbrk{}पा\appbrk{}णिभिः}}\Afootnote{A; \skt{परि\appbrk{}षे\appbrk{}च\appbrk{}नो\appbrk{}प\appbrk{}ध\appbrk{}म\appbrk{}न\appbrk{}बा\appbrk{}ल\appbrk{}व\appbrk{}स्त्र\appbrk{}पा\appbrk{}णिभिः}\ N; \skt{परि\appbrk{}षे\appbrk{}च\appbrk{}णो\appbrk{}प\appbrk{}ध\appbrk{}म\appbrk{}न\appbrk{}बा\appbrk{}ल\appbrk{}व\appbrk{}स्त्र\appbrk{}पा\appbrk{}णिभिः}\ H}}\skt{ }\edtext{\skt{प्रमार्जयेत्~। आहा\-र\-शे\-ष\-श्ले\-ष्म\-ही\-ना\-णु\-श\-ल्यानि श्वस\-नो\-त्का\-स\-न\-प्र\-ध\-म\-नै\-र्नि\-र्ध\-मेत्}}{\lemma{\skt{प्रमार्जयेत्~। आहा\appbrk{}र\appbrk{}शे\appbrk{}ष\appbrk{}श्ले\appbrk{}ष्म\appbrk{}ही\appbrk{}ना\appbrk{}णु\appbrk{}श\appbrk{}ल्यानि श्वस\appbrk{}नो\appbrk{}त्का\appbrk{}स\appbrk{}न\appbrk{}प्र\appbrk{}ध\appbrk{}म\appbrk{}नैर्}}\Afootnote{A; \skt{प्रमा\appbrk{}र्ज\appbrk{}ये\appbrk{}दे\appbrk{}त\appbrk{}त्प्र\appbrk{}मा\appbrk{}र्जनं नाम~॥ आहारशेषं कण्ठगतं सेष्मगतं सिंघाणकं निश्वा\appbrk{}सो\appbrk{}त्का\appbrk{}ष\appbrk{}नाभ्यां निर्धमेदेतं}\ N; \skt{प्रमा\appbrk{}र्ज्ज\appbrk{}ये\appbrk{}दे\appbrk{}त\appbrk{}त्प्र\appbrk{}मा\appbrk{}र्ज्ज\appbrk{}न\appbrk{}न्नाम~॥ आहारशेषं कण्ठगतं श्लेष्मगतं सिंघा\appbrk{}ण\appbrk{}क\appbrk{}न्नि\appbrk{}श्वा\appbrk{}सो\appbrk{}त्का\appbrk{}स\appbrk{}ना\appbrk{}भ्या\appbrk{}न्नि\appbrk{}र्द्ध\appbrk{}मे\appbrk{}दे\appbrk{}तन्}\ H}}\skt{~। अन्नशल्यानि }\edtext{\skt{वम\-ना\-ङ्गु\-लि\-प्र\-ति\-म\-र्श\-प्र\-भृ\-तिभिः}}{\lemma{\skt{वम\appbrk{}ना\appbrk{}ङ्गु\appbrk{}लि\appbrk{}प्र\appbrk{}ति\appbrk{}म\appbrk{}र्श\appbrk{}प्र\appbrk{}भृ\appbrk{}तिभिः}}\Afootnote{A; \skt{वम\appbrk{}ना\appbrk{}ङ्गु\appbrk{}लि\appbrk{}स्प\appbrk{}र्श\appbrk{}न\appbrk{}प्र\appbrk{}भृ\appbrk{}तिभिः}\ N H}}\skt{~। }\edtext{\skt{विरेचनैः पक्वा\-श\-य\-ग\-तानि~। व्रण\-दो\-षा\-श\-य\-ग\-तानि}}{\lemma{\skt{विरेचनैः पक्वा\appbrk{}श\appbrk{}य\appbrk{}ग\appbrk{}तानि~। व्रण\appbrk{}दो\appbrk{}षा\appbrk{}श\appbrk{}य\appbrk{}ग\appbrk{}तानि}}\Afootnote{A; \skt{प्रमृ\appbrk{}स्यो\appbrk{}द्ध\appbrk{}रे\appbrk{}देश}\ N; \skt{प्रमृ\appbrk{}स्यो\appbrk{}द्ध\appbrk{}रे\appbrk{}देष}\ H}}\skt{ }\edtext{\skt{प्रक्षालनैः~। वात\-मू\-त्र\-पु\-री\-ष\-ग\-र्भ\-स\-ङ्गेषु}}{\lemma{\skt{प्रक्षालनैः~। वात\appbrk{}मू\appbrk{}त्र\appbrk{}पु\appbrk{}री\appbrk{}ष\appbrk{}ग\appbrk{}र्भ\appbrk{}स\appbrk{}ङ्गेषु}}\Afootnote{A; \skt{प्रतिमर्शो नाम~॥ वात\appbrk{}मू\appbrk{}त्र\appbrk{}पु\appbrk{}री\appbrk{}ष\appbrk{}ग\appbrk{}र्भ\appbrk{}संगे}\ N; \skt{प्रतिमर्शो नाम~॥ वात\appbrk{}मू\appbrk{}त्र\appbrk{}पु\appbrk{}री\appbrk{}ष\appbrk{}ग\appbrk{}र्ब्भ\appbrk{}स\appbrk{}ङ्गेषु}\ H}}\skt{ }\edtext{\skt{प्रवाहणमुक्तं}}{\lemma{\skt{प्रवाहणमुक्तं}}\Afootnote{\skt{प्रवा\appbrk{}ह\appbrk{}ण\appbrk{}मु\appbrk{}क्तम्}\ A; \skt{विरेचनैर्वा प्रवा\appbrk{}ह\appbrk{}णे\appbrk{}नो\appbrk{}द्ध\appbrk{}रे\appbrk{}दे\appbrk{}त\appbrk{}त्प्र\appbrk{}वा\appbrk{}हनं नाम}\ N; \skt{विरेचनैर्व्वा प्रवा\appbrk{}ह\appbrk{}णे\appbrk{}नो\appbrk{}द्ध\appbrk{}रे\appbrk{}दे\appbrk{}त\appbrk{}त्प्र\appbrk{}वा\appbrk{}ह\appbrk{}ण\appbrk{}न्नाम}\ H}}\skt{~। }\edtext{\edtext{\skt{मारु\-तो\-द\-क\-स\-वि\-ष\-रु\-धि\-र\-दु\-ष्ट\-स्त\-न्ये\-ष्वा\-चू\-ष\-ण\-मा\-स्येन}}{\lemma{\skt{मारु\appbrk{}तो\appbrk{}द\appbrk{}क\appbrk{}स\appbrk{}वि\appbrk{}ष\appbrk{}रु\appbrk{}धि\appbrk{}र\appbrk{}दु\appbrk{}ष्ट\appbrk{}स्त\appbrk{}न्येष्व्}}\Afootnote{A; \skt{मारु\appbrk{}तो\appbrk{}द\appbrk{}क\appbrk{}स्त\appbrk{}न्य\appbrk{}रु\appbrk{}धि\appbrk{}र\appbrk{}वि\appbrk{}ष\appbrk{}श\appbrk{}ल्यान्य्}\ N H}}}{\lemma{\skt{आचू\appbrk{}ष\appbrk{}ण\appbrk{}मा\appbrk{}स्येन}}\Afootnote{A; \skt{आचू\appbrk{}ष\appbrk{}णे\appbrk{}नो\appbrk{}द्ध\appbrk{}रे\appbrk{}दा\appbrk{}स्येन}\ N H}}\skt{ }\edtext{\edtext{\skt{विषाणैर्वा}}{\lemma{\skt{विषाणैर्}}\Afootnote{A; \skt{विषाणेन वा एत\appbrk{}दा\appbrk{}चू\appbrk{}ष\appbrk{}णन्}\ N H}}}{\lemma{\skt{वा}}\Afootnote{A; \skt{नाम}\ N H}}\skt{~। }\edtext{\edtext{\skt{अनु\-लो\-म\-म\-न\-व\-ब\-द्ध\-म\-क\-र्ण\-म\-न\-ल्प\-व्र\-ण\-मु\-ख\-म\-य\-स्का\-न्तेन}}{\lemma{\skt{अकर्णम्}}\Afootnote{A H; \skt{आकर्ण्णम्}\ N}}\skt{~। }\edtext{\skt{हृद्य\-व\-स्थि\-त\-म\-ने\-क\-का\-र\-णो\-त्पन्नं}}{\lemma{\skt{अने\appbrk{}क\appbrk{}का\appbrk{}र\appbrk{}णो\appbrk{}त्पन्नं}}\Afootnote{A H; \skt{अने\appbrk{}क\appbrk{}का\appbrk{}र\appbrk{}णे\appbrk{}नो\appbrk{}त्पन्नं}\ N}}}{\lemma{\skt{अयस्कान्तेन~। हृद्य्}}\Afootnote{A; \skt{अय\appbrk{}स्का\appbrk{}न्ते\appbrk{}नो\appbrk{}द्ध\appbrk{}रे\appbrk{}थृद्य्}\ N H}}\skt{ शोकशल्यं }\edtext{\skt{हर्षेणेति}}{\lemma{\skt{हर्षेणेति}}\Afootnote{A; \skt{हर्षणेनेति}\ N H}}\skt{~॥}\edtext{}{\lemma{\skt{निर्धमेत्~। अन्नशल्यानि}}\Afootnote{A; \skt{निर्धमनं नाम अन्यशल्यानि तु}\ N H\textsuperscript{pc}; \skt{निर्द्धमनन्नाम~॥ अ न्न शल्यानि तु}\ H\textsuperscript{ac}}}\pend

\pstart \ledsidenote{\passref{1.27.6}}\skt{सर्वशल्यानां }\edtext{\skt{तु}}{\lemma{\skt{तु}}\Afootnote{A; \skt{शरीरगताना}\ N; \skt{शरी\appbrk{}र\appbrk{}ग\appbrk{}ता\appbrk{}नाम्}\ H}}\skt{ }\edtext{\skt{महतामणूनां वा द्वावे\-वा\-ह\-र\-ण\-हेतू}}{\lemma{\skt{अणूनां वा द्वाव्}}\Afootnote{A; \skt{अणूनाञ्च द्वाव्}\ N H}}\skt{ भवतः }\edtext{\edtext{\skt{प्रति\-लो\-मोऽ\-नु\-लो\-मश्च}}{\lemma{\skt{प्रतिलोमो}}\Afootnote{A N; \skt{प्रतिलोमा}\ H}}}{\lemma{\skt{ऽनुलोमश्}}\Afootnote{A; \skt{नुलोमश्}\ N H}}\skt{~॥}\pend

\pstart \ledsidenote{\passref{1.27.7}}\skt{तत्र }\edtext{\skt{प्रति\-लो\-म\-म\-र्वा\-ची\-न\-मा\-न\-ये\-त\-नु\-लोमं}}{\lemma{\skt{प्रति\appbrk{}लो\appbrk{}म\appbrk{}म\appbrk{}र्वा\appbrk{}ची\appbrk{}न\appbrk{}मा\appbrk{}न\appbrk{}ये\appbrk{}त\appbrk{}नु\appbrk{}लोमं}}\Afootnote{A; \skt{प्रति\appbrk{}लो\appbrk{}म\appbrk{}म\appbrk{}वां\appbrk{}ची\appbrk{}न\appbrk{}म\appbrk{}प\appbrk{}न\appbrk{}ये\appbrk{}त\appbrk{}नु\appbrk{}लोमं}\ N H}}\skt{ पराचीनम्~॥}\pend

\pstart \ledsidenote{\passref{1.27.8}}\edtext{\skt{उत्तुण्डितं}}{\lemma{\skt{उत्तुण्डितं}}\Afootnote{H; \skt{उत्तुण्दितं}\ A; \skt{उतुण्डितं}\ N}}\skt{ छित्वा }\edtext{\skt{निर्घा\-त\-ये\-च्छे\-द\-नी\-य\-मु\-खम्}}{\lemma{\skt{निर्घा\appbrk{}त\appbrk{}ये\appbrk{}च्छे\appbrk{}द\appbrk{}नी\appbrk{}य\appbrk{}मु\appbrk{}खम्}}\Afootnote{A; \skt{निर्घातयेत्}\ N H}}\skt{~॥}\pend

\pstart \ledsidenote{\passref{1.27.9}}\edtext{\skt{छेद\-नी\-य\-मु\-खा\-न्यपि}}{\lemma{\skt{छेद\appbrk{}नी\appbrk{}य\appbrk{}मु\appbrk{}खा\appbrk{}न्यपि}}\Afootnote{\skt{छेद\appbrk{}नी\appbrk{}य\appbrk{}मु\appbrk{}खा\appbrk{}न्यपि}\ A; \skt{च्छेद\appbrk{}नी\appbrk{}य\appbrk{}मु\appbrk{}खा\appbrk{}न्य\appbrk{}च्छे\appbrk{}द\appbrk{}नी\appbrk{}य\appbrk{}मु\appbrk{}खानि}\ N H}}\skt{ }\edtext{\skt{कुक्षि\-व\-क्षः\-क\-क्षा\-वं\-क्ष\-ण\-प\-र्शु\-का\-न्त\-र\-प\-ति\-तानि च}}{\lemma{\skt{कुक्षि\appbrk{}व\appbrk{}क्षः\appbrk{}क\appbrk{}क्षा\appbrk{}वं\appbrk{}क्ष\appbrk{}ण\appbrk{}प\appbrk{}र्शु\appbrk{}का\appbrk{}न्त\appbrk{}र\appbrk{}प\appbrk{}ति\appbrk{}तानि च}}\Afootnote{A; \skt{कक्ष\appbrk{}व\appbrk{}क्षः।\appbrk{}प\appbrk{}र\appbrk{}शु\appbrk{}का\appbrk{}न्त\appbrk{}र\appbrk{}म्प\appbrk{}ति\appbrk{}तानि च}\ N; \skt{कक्ष\appbrk{}व\appbrk{}क्षः\appbrk{}प\appbrk{}र\appbrk{}शु\appbrk{}का\appbrk{}न्त\appbrk{}र\appbrk{}प\appbrk{}ति\appbrk{}तानि च}\ H}}\skt{ }\edtext{\skt{हस्तशक्यं}}{\lemma{\skt{हस्तशक्यं}}\Afootnote{A; \emph{om.}\ N H}}\skt{ }\edtext{\skt{यथामार्गेण}}{\lemma{\skt{यथामार्गेण}}\Afootnote{A; \skt{यथामार्गे}\ N H}}\skt{ }\edtext{\skt{हस्ते\-नै\-वा\-प\-हर्तुं}}{\lemma{\skt{हस्ते\appbrk{}नै\appbrk{}वा\appbrk{}प\appbrk{}हर्तुं}}\Afootnote{A; \skt{णाहर्तुं}\ N H}}\skt{ प्रयतेत~॥}\pend

\pstart \ledsidenote{\passref{1.27.10}}\edtext{\skt{हस्ते\-ना\-प\-ह\-र्तु\-म\-शक्यं}}{\lemma{\skt{हस्ते\appbrk{}ना\appbrk{}प\appbrk{}ह\appbrk{}र्तुम्}}\Afootnote{\skt{हस्ते\appbrk{}नै\appbrk{}वा\appbrk{}प\appbrk{}ह\appbrk{}र्तुम्}\ A; \skt{हस्ते\appbrk{}ना\appbrk{}वा\appbrk{}ह\appbrk{}र्तुम्}\ N; \skt{हस्तेन वा हर्त्तुम्}\ H}}\skt{ }\edtext{\skt{विशस्य}}{\lemma{\skt{विशस्य}}\Afootnote{A N; \skt{विसस्य}\ H}}\skt{ }\edtext{\skt{शस्त्रेण}}{\lemma{\skt{शस्त्रेण}}\Afootnote{A; \skt{सस्त्रे}\ N; \skt{यन्त्रेण}\ H}}\skt{ }\edtext{\skt{यन्त्रे\-णा\-प\-ह\-रेत्}}{\lemma{\skt{यन्त्रे\appbrk{}णा\appbrk{}प\appbrk{}ह\appbrk{}रेत्}}\Afootnote{A; \skt{णाहरेत्}\ N; \skt{शस्त्रेणाहरेत्}\ H}}\skt{~॥}\pend

\pstart\versepar \ledsidenote{\passref{1.27.11}}\edtext{\skt{भवति}}{\lemma{\skt{भवति}}\Afootnote{A; \emph{om.}\ N H}}\skt{ }\edtext{\skt{चात्र}}{\lemma{\skt{चात्र}}\Afootnote{A; \skt{भ भवति चात्र}\ N H}}\skt{~। }\pend
\pstart\versepar \skt{शीतलेन }\edtext{\skt{जलेनैनं}}{\lemma{\skt{जलेनैनं}}\Afootnote{A H; \skt{जलेनेनम्}\ N}}\skt{ }\edtext{\skt{मूर्च्छ\-न्त\-म\-व\-से\-च\-येत्}}{\lemma{\skt{मूर्च्छन्तम्}}\Afootnote{N H; \skt{मूर्च्छन्त}\ A}}\skt{~। }\pend
\pstart\versepar \skt{संरक्षेदस्य मर्माणि मुहु\-रा\-श्वा\-स\-येच्च तम्~॥}\edtext{}{\lemma{\skt{अवसेचयेत्~। संरक्षेदस्य}}\Afootnote{\skt{मवसेचयेत्~। संरक्षेदस्य}\ A; \skt{अवसेचयेत्~। संरक्षे चास्य}\ N; \skt{अवसेचयेत्~। संरक्षेच्चास्य}\ H}}\pend

\pstart \ledsidenote{\passref{1.27.12}}\skt{ततः शल्यमुद्धृत्य }\edtext{\skt{निर्लोहितं}}{\lemma{\skt{निर्लोहितं}}\Afootnote{A H; \skt{निर्लोहित}\ N}}\skt{ व्रणं कृत्वा }\edtext{\skt{स्वेदा\-र्ह\-म\-ग्नि\-घृ\-त\-प्र\-भृ\-तिभिः}}{\lemma{\skt{अग्नि\appbrk{}घृ\appbrk{}त\appbrk{}प्र\appbrk{}भृ\appbrk{}तिभिः}}\Afootnote{A H; \skt{अग्निप्रभृतिभिः}\ N}}\skt{ }\edtext{\skt{संस्वेद्य विदह्य}}{\lemma{\skt{संस्वेद्य विदह्य}}\Afootnote{\skt{संस्वेद्यावदह्य}\ A; \skt{स्वेदयि}\ N H}}\skt{ }\edtext{\skt{प्रदिह्य}}{\lemma{\skt{प्रदिह्य}}\Afootnote{A; \skt{त्वा मधु}\ N H}}\skt{ }\edtext{\skt{सर्पिर्मधुभ्यां}}{\lemma{\skt{सर्पिर्मधुभ्यां}}\Afootnote{A; \skt{सर्पिर्भ्यां}\ N H}}\skt{ }\edtext{\edtext{\skt{बद्ध्वा+\-आ\-चा\-रि\-क\-मु\-प\-दि\-शेत्}}{\lemma{\skt{बद्ध्वा+\appbrk{}आ\appbrk{}चा\appbrk{}रि\appbrk{}कम्}}\Afootnote{\skt{बद्ध्वाऽऽचारिकम्}\ A N H}}\skt{~। (? सिरा\-स्ना\-यु\-वि\-लग्नं}}{\lemma{\skt{उपदिशेत्~। (? सिरा\appbrk{}स्ना\appbrk{}यु\appbrk{}वि\appbrk{}लग्नं}}\Afootnote{A; \skt{उपदिशेत्}\ N H}}\skt{ }\edtext{\edtext{\skt{शला\-का\-दि\-भि\-र्वि\-मो\-च्या\-प\-न\-ये\-त्श्व\-य\-थु\-ग्र\-स्त\-वा\-रङ्गं}}{\lemma{\skt{शलाकादिभिर्}}\Afootnote{A; \emph{om.}\ N H}}}{\lemma{\skt{विमो\appbrk{}च्या\appbrk{}प\appbrk{}न\appbrk{}ये\appbrk{}त्श्व\appbrk{}य\appbrk{}थु\appbrk{}ग्र\appbrk{}स्त\appbrk{}वा\appbrk{}रङ्गं}}\Afootnote{A; \emph{om.}\ N H}}\skt{ }\edtext{\skt{समवपीड्य}}{\lemma{\skt{समवपीड्य}}\Afootnote{A; \emph{om.}\ N H}}\skt{ }\edtext{\skt{श्वयथुं}}{\lemma{\skt{श्वयथुं}}\Afootnote{A; \emph{om.}\ N H}}\skt{ }\edtext{\skt{दुर्बलवारङ्गं}}{\lemma{\skt{दुर्बलवारङ्गं}}\Afootnote{A; \emph{om.}\ N H}}\skt{ }\edtext{\edtext{\skt{कुशादिभिर्बद्ध्वा}}{\lemma{\skt{कुशादिभिर्}}\Afootnote{A; \emph{om.}\ N H}}\skt{~।)}}{\lemma{\skt{बद्ध्वा~।)}}\Afootnote{A; \emph{om.}\ N H}}\pend

\pstart \ledsidenote{\passref{1.27.13}}\edtext{\skt{हृदयमभितो}}{\lemma{\skt{हृदयमभितो}}\Afootnote{A; \skt{हृदयगनि}\ N; \skt{हृदयमभि}\ H}}\skt{ }\edtext{\skt{वर्तमानं}}{\lemma{\skt{वर्तमानं}}\Afootnote{A H; \skt{वर्त्तमान}\ N}}\skt{ शल्यं }\edtext{\skt{शीत\-ज\-ला\-दि\-भि\-रु\-द्वे\-जि\-त\-स्या\-प\-ह\-रे\-द्य\-था\-मार्गं}}{\lemma{\skt{उद्वे\appbrk{}जि\appbrk{}त\appbrk{}स्या\appbrk{}प\appbrk{}ह\appbrk{}रे\appbrk{}द्य\appbrk{}था\appbrk{}मार्गं}}\Afootnote{A; \skt{उद्वे\appbrk{}जि\appbrk{}त\appbrk{}स्यो\appbrk{}द्वि\appbrk{}ग्न\appbrk{}हृ\appbrk{}द\appbrk{}य\appbrk{}स्या\appbrk{}प\appbrk{}ह\appbrk{}रे\appbrk{}द्य\appbrk{}था\appbrk{}मार्ग}\ N H}}\skt{ }\edtext{\edtext{\skt{दुरु\-प\-ह\-र\-म\-न्य\-तोऽ\-प\-बा\-ध्य\-मानं}}{\lemma{\skt{दुरु\appbrk{}प\appbrk{}ह\appbrk{}र\appbrk{}म\appbrk{}न्यतो}}\Afootnote{\skt{दुरु\appbrk{}प\appbrk{}ह\appbrk{}र\appbrk{}म\appbrk{}न्यतो}\ A; \skt{दुरा\appbrk{}ह\appbrk{}र\appbrk{}ण\appbrk{}मन्ये}\ N; \skt{दुरा\appbrk{}ह\appbrk{}र\appbrk{}ण\appbrk{}म\appbrk{}न्यतो}\ H}}}{\lemma{\skt{ऽपबाध्यमानं}}\Afootnote{A; \skt{न पराध्यमानं}\ N H}}\skt{ }\edtext{\skt{पाट\-यि\-त्वो\-द्ध\-रेत्}}{\lemma{\skt{पाट\appbrk{}यि\appbrk{}त्वो\appbrk{}द्ध\appbrk{}रेत्}}\Afootnote{A; \skt{पाट\appbrk{}यि\appbrk{}त्वा\appbrk{}प\appbrk{}ह\appbrk{}रेत्}\ N H}}\skt{~॥}\pend

\pstart \ledsidenote{\passref{1.27.14-16}}\edtext{\skt{अस्थि\-वि\-व\-र\-प्र\-वि\-ष्ट\-म\-स्थि\-वि\-दष्टं}}{\lemma{\skt{अस्थिविदष्टं}}\Afootnote{A H; \skt{अस्थिविदञ्}\ N}}\skt{ }\edtext{\skt{चावगृह्य}}{\lemma{\skt{चावगृह्य}}\Afootnote{N H; \skt{वाऽवगृह्य}\ A}}\skt{ }\edtext{\skt{पद्भ्यां}}{\lemma{\skt{पद्भ्यां}}\Afootnote{N H; \skt{पादाभ्यां}\ A}}\skt{ }\edtext{\skt{यन्त्रे\-णा\-प\-ह\-रेत्~। एवमशक्यं गुणेन परिवेष्ट्य}}{\lemma{\skt{यन्त्रे\appbrk{}णा\appbrk{}प\appbrk{}ह\appbrk{}रेत्~। एवमशक्यं गुणेन परिवेष्ट्य}}\Afootnote{H; \skt{यन्त्रे\appbrk{}णा\appbrk{}प\appbrk{}ह\appbrk{}रेत्, अशक्यमेवं वा बलवद्भिः सुप\appbrk{}रि\appbrk{}गृ\appbrk{}ही\appbrk{}तस्य यन्त्रेण ग्राहयित्वा शल्यवारङ्गं प्रविभुज्य धनु\appbrk{}र्गु\appbrk{}णै\appbrk{}र्ब\appbrk{}द्ध्वै\appbrk{}क\appbrk{}त\appbrk{}श्चास्य पञ्चा\appbrk{}ङ्ग्या\appbrk{}मु\appbrk{}प\appbrk{}सं\appbrk{}य\appbrk{}त\appbrk{}स्या\appbrk{}श्वस्य वक्त्रकविके}\ A; \skt{यंत्रे\appbrk{}णा\appbrk{}प\appbrk{}ह\appbrk{}रेत्~। एवमशक्यं गुणैर्न}\ N}}\skt{ }\edtext{\skt{बध्वा\-श्व\-व\-क्त्र\-क\-टके}}{\lemma{\skt{बध्वा\appbrk{}श्व\appbrk{}व\appbrk{}क्त्र\appbrk{}क\appbrk{}टके}}\Afootnote{H\textsuperscript{pc}; \skt{बध्नीयात्, अथैनं कशया ताड\appbrk{}ये\appbrk{}द्य\appbrk{}थो\appbrk{}न्ना\appbrk{}म\appbrk{}य\appbrk{}ञ्शिरो वेगेन शल्यमुद्धरति}\ A; \skt{परिवेष्ट्य बध्वा\appbrk{}श्व\appbrk{}व\appbrk{}क्त्र\appbrk{}कट}\ N; \skt{बध्वा\appbrk{}श्व\appbrk{}व\appbrk{}क्त्र\appbrk{}क\appbrk{}टके~। पञ्चाङ्ग्यां}\ H\textsuperscript{ac}}}\skt{~। }\edtext{\skt{पञ्चाङ्ग्यां}}{\lemma{\skt{पञ्चाङ्ग्यां}}\Afootnote{H\textsuperscript{pc}; \skt{दृढां वा}\ A; \emph{om.}\ N H\textsuperscript{ac}}}\skt{ }\edtext{\skt{वृक्षशाखायाम्वा}}{\lemma{\skt{वृक्षशाखायाम्वा}}\Afootnote{H; \skt{वृक्ष\appbrk{}शा\appbrk{}खा\appbrk{}म\appbrk{}व\appbrk{}नम्य तस्यां पूर्ववद्बद्ध्वो}\ A; \skt{के~। पञ्चा\appbrk{}वृ\appbrk{}क्ष\appbrk{}सा\appbrk{}खायां वा}\ N}}\skt{ }\edtext{\edtext{\edtext{\edtext{\skt{अच्छे\-द्य\-दे\-शो\-त्तु\-ण्डि\-त\-म\-ष्ठी\-ला\-स्म\-मु\-द्ग\-रा\-णा\-म\-न्य\-त\-मस्य}}{\lemma{\skt{अच्छे\appbrk{}द्य\appbrk{}दे\appbrk{}शो\appbrk{}त्तु\appbrk{}ण्डि\appbrk{}तम्}}\Afootnote{H; \skt{द्धरेत्~॥१४॥ अदेशोत्तुण्डित}\ A; \skt{अच्छे\appbrk{}द्य\appbrk{}दे\appbrk{}शो\appbrk{}त्तु\appbrk{}ण्डित}\ N}}}{\lemma{\skt{अष्ठीलास्मम्}}\Afootnote{H; \skt{मष्ठीलाश्म}\ A; \skt{मष्ठीलास्म}\ N}}}{\lemma{\skt{उद्गराणाम्}}\Afootnote{H; \skt{मुद्गराणा}\ A; \skt{मुद्गराणाम्}\ N}}}{\lemma{\skt{अन्यतमस्य}}\Afootnote{N H\textsuperscript{pc}; \skt{मन्यतमस्य}\ A; \skt{अन्यतस्य}\ H\textsuperscript{ac}}}\skt{ प्रहारेण विचाल्य }\edtext{\skt{यथामार्ग्गत}}{\lemma{\skt{यथामार्ग्गत}}\Afootnote{N H; \skt{यथामार्गम्}\ A}}\skt{ एव यन्त्रेण }\edtext{\edtext{\edtext{\skt{मृदि\-त\-क\-र्ण्ण\-म\-क\-र्ण्ण\-ञ्चेति}}{\lemma{\skt{मृदितकर्ण्णम्}}\Afootnote{N H; \skt{विमृदितकर्णा}\ A}}}{\lemma{\skt{अकर्ण्णञ्}}\Afootnote{N H; \skt{नि कर्ण}\ A}}}{\lemma{\skt{चेति}}\Afootnote{N H; \skt{वन्}\ A}}\skt{~। }\edtext{\skt{अना\-बा\-ध\-क\-र\-दे\-शो\-त्तु\-ण्डि\-तानि}}{\lemma{\skt{अना\appbrk{}बा\appbrk{}ध\appbrk{}क\appbrk{}र\appbrk{}दे\appbrk{}शो\appbrk{}त्तु\appbrk{}ण्डि\appbrk{}तानि}}\Afootnote{H; \skt{त्यना\appbrk{}बा\appbrk{}ध\appbrk{}क\appbrk{}र\appbrk{}दे\appbrk{}शो\appbrk{}त्तु\appbrk{}ण्डि\appbrk{}तानि}\ A; \skt{अना\appbrk{}बा\appbrk{}ध\appbrk{}क\appbrk{}र\appbrk{}सो\appbrk{}तु\appbrk{}ण्डि\appbrk{}तानि}\ N}}\skt{ पुरस्तादेव~।}\pend

\pstart \ledsidenote{\passref{1.27.17}}\edtext{\skt{जातुषे}}{\lemma{\skt{जातुषे}}\Afootnote{A; \emph{om.}\ K; \skt{जातुषे तु}\ N H}}\skt{ }\edtext{\skt{कण्ठसक्ते}}{\lemma{\skt{कण्ठसक्ते}}\Afootnote{N; \skt{कण्ठासक्ते}\ A; \emph{om.}\ K; \skt{कण्ठशक्ते}\ H}}\skt{ }\edtext{\skt{कण्ठनाडीं}}{\lemma{\skt{कण्ठनाडीं}}\Afootnote{N; \skt{कण्ठे नाडीं}\ A; \emph{om.}\ K; \skt{कण्ठनाडी}\ H}}\skt{ }\edtext{\skt{प्रवेश्याग्नितप्तां च}}{\lemma{\skt{प्रवेश्याग्नितप्तां च}}\Afootnote{A; \emph{om.}\ K; \skt{प्रवे\appbrk{}श्या\appbrk{}ग्नि\appbrk{}त\appbrk{}प्ताश्च}\ N; \skt{प्रवे\appbrk{}श्या\appbrk{}ग्नि\appbrk{}त\appbrk{}प्तया}\ H}}\skt{ }\edtext{\skt{शलाकां}}{\lemma{\skt{शलाकां}}\Afootnote{A H; \emph{om.}\ K; \skt{शलाका}\ N}}\skt{ }\edtext{\skt{तयावगृह्य}}{\lemma{\skt{तयावगृह्य}}\Afootnote{A; \emph{om.}\ K; \skt{तु या प्राप्ता}\ N; \skt{तया विगृह्य}\ H\textsuperscript{ac}; \skt{तया विप्राप्त}\ H\textsuperscript{pc}}}\skt{ }\edtext{\edtext{\skt{शीताभिरद्भिः}}{\lemma{\skt{शीताभिर्}}\Afootnote{A N H; \emph{om.}\ K}}}{\lemma{\skt{अद्भिः}}\Afootnote{A N H; \emph{om.}\ K}}\skt{ }\edtext{\skt{परिषिच्य}}{\lemma{\skt{परिषिच्य}}\Afootnote{A N; \skt{षिच्य}\ K; \skt{परिषिञ्च्य}\ H}}\skt{ }\edtext{\skt{स्थिरी\-भू\-ता\-म्श\-ल्य\-मु\-द्ध\-रेत्}}{\lemma{\skt{स्थिरीभूताम्}}\Afootnote{\skt{स्थिरीभूतं}\ A N H K}}\skt{~॥}\pend

\pstart \ledsidenote{\passref{1.27.18-19}}\edtext{\skt{अजातुषं}}{\lemma{\skt{अजातुषं}}\Afootnote{A K; \emph{om.}\ N H}}\skt{ }\edtext{\skt{जतुना लिप्तया}}{\lemma{\skt{जतुना लिप्तया}}\Afootnote{K; \skt{तु जतु\appbrk{}म\appbrk{}धू\appbrk{}च्छि\appbrk{}ष्ट\appbrk{}प्र\appbrk{}लि\appbrk{}प्तया शलाकया}\ A; \skt{अजातुषं जतुना लिप्तया}\ N; \skt{अजानुषं जतुना लिप्तया}\ H}}\skt{ }\edtext{\skt{पूर्व\-क\-ल्पे\-ना\-स्थि\-श\-ल्य\-म\-न्यद्वा}}{\lemma{\skt{पूर्व\appbrk{}क\appbrk{}ल्पे\appbrk{}ना\appbrk{}स्थि\appbrk{}श\appbrk{}ल्यम्}}\Afootnote{K N H; \skt{पूर्व\appbrk{}क\appbrk{}ल्पे\appbrk{}ने\appbrk{}त्येके~॥१८॥ अस्थिशल्यम्}\ A}}\skt{ }\edtext{\edtext{\skt{तिर्य\-क्क\-ण्ठा\-स\-क्त\-म\-वेक्ष्य}}{\lemma{\skt{तिर्य\appbrk{}क्क\appbrk{}ण्ठा\appbrk{}स\appbrk{}क्तम्}}\Afootnote{A; \skt{तिर्य सक्तम्}\ K; \skt{तिर्य शक्ताम्}\ N; \skt{तिर्य\appbrk{}क्क\appbrk{}ण्ठ\appbrk{}श\appbrk{}क्तम्}\ H}}}{\lemma{\skt{अवेक्ष्य}}\Afootnote{A; \skt{अवेक्ष्यञ्च}\ K N H}}\skt{ }\edtext{\skt{केशोण्डुकं}}{\lemma{\skt{केशोण्डुकं}}\Afootnote{A; \skt{केशोण्डुवं}\ K N; \skt{केशोण्डुक}\ H}}\skt{ }\edtext{\skt{दृढै\-क\-सू\-त्र\-बद्धं द्रव\-भ\-क्तो\-प\-हितं}}{\lemma{\skt{दृढै\appbrk{}क\appbrk{}सू\appbrk{}त्र\appbrk{}बद्धं द्रव\appbrk{}भ\appbrk{}क्तो\appbrk{}प\appbrk{}हितं}}\Afootnote{A; \skt{दृढै\appbrk{}क\appbrk{}दी\appbrk{}र्घ\appbrk{}सू\appbrk{}त्र\appbrk{}ब\appbrk{}न्धन कृत्वा द्रव\appbrk{}भ\appbrk{}क्तो\appbrk{}प\appbrk{}हितं}\ K; \skt{दृढै\appbrk{}क\appbrk{}दी\appbrk{}र्घा\appbrk{}सू\appbrk{}त्र\appbrk{}ब\appbrk{}न्धनं कृत्वा द्रव\appbrk{}भ\appbrk{}क्तो\appbrk{}प\appbrk{}हितं}\ N; \skt{दृढै\appbrk{}क\appbrk{}दी\appbrk{}र्घ\appbrk{}सू\appbrk{}त्र\appbrk{}ब\appbrk{}न्ध\appbrk{}न\appbrk{}ङ्कृत्वा द्रव\appbrk{}भ\appbrk{}क्तो\appbrk{}प\appbrk{}हितं}\ H}}\skt{ }\edtext{\skt{पाय\-ये\-दा\-क\-ण्ठा\-त्पू\-र्ण\-कोष्ठं च}}{\lemma{\skt{पाय\appbrk{}ये\appbrk{}दा\appbrk{}क\appbrk{}ण्ठा\appbrk{}त्पू\appbrk{}र्ण\appbrk{}कोष्ठं च}}\Afootnote{A; \skt{पाय}\ K\textsuperscript{ac}; \skt{पाय\appbrk{}ये\appbrk{}दा\appbrk{}क\appbrk{}ण्ठाच्च पूर्ण्णकोष्ठम्}\ K\textsuperscript{pc} H; \skt{पाय\appbrk{}य\appbrk{}दा\appbrk{}क\appbrk{}ण्ठाच्च पूर्ण्णकोष्ठं}\ N}}\skt{ }\edtext{\skt{वाम\-ये\-त्व\-म\-तश्च}}{\lemma{\skt{वाम\appbrk{}ये\appbrk{}त्व\appbrk{}म\appbrk{}तश्च}}\Afootnote{A K\textsuperscript{pc} N H; \skt{येद्वमतश्च}\ K\textsuperscript{ac}}}\skt{ }\edtext{\skt{शल्यै\-क\-दे\-श\-सक्तं}}{\lemma{\skt{शल्यै\appbrk{}क\appbrk{}दे\appbrk{}श\appbrk{}सक्तं}}\Afootnote{A K; \skt{शल्यैकदेश}\ N; \skt{शल्यै\appbrk{}क\appbrk{}दे\appbrk{}श\appbrk{}शक्तं}\ H}}\skt{ }\edtext{\skt{ज्ञात्वा}}{\lemma{\skt{ज्ञात्वा}}\Afootnote{A; \emph{om.}\ K H; \skt{क्तं}\ N}}\skt{ सूत्रं }\edtext{\skt{सहसा}}{\lemma{\skt{सहसा}}\Afootnote{A H; \skt{अन्यद्वा सहसा}\ K N}}\skt{ }\edtext{\skt{त्वाक्षिपेत्मृदुना वा}}{\lemma{\skt{त्वाक्षिपेत्मृदुना वा}}\Afootnote{A; \skt{त्वाक्षिपेत्मृदुना}\ K N H}}\skt{ }\edtext{\skt{दन्त\-धा\-व\-न\-कू\-र्च\-के\-ना\-प\-ह\-रे\-त्प्र\-णु\-दे\-द्वान्तः~। क्षतकण्ठाय च}}{\lemma{\skt{दन्त\appbrk{}धा\appbrk{}व\appbrk{}न\appbrk{}कू\appbrk{}र्च\appbrk{}के\appbrk{}ना\appbrk{}प\appbrk{}ह\appbrk{}रे\appbrk{}त्प्र\appbrk{}णु\appbrk{}दे\appbrk{}द्वान्तः~। क्षतकण्ठाय च}}\Afootnote{\skt{दन्त\appbrk{}धा\appbrk{}व\appbrk{}न\appbrk{}कू\appbrk{}र्च\appbrk{}के\appbrk{}ना\appbrk{}प\appbrk{}ह\appbrk{}रे\appbrk{}त्प्र\appbrk{}णु\appbrk{}दे\appbrk{}द्वाऽ\appbrk{}अन्तः~। क्षतकण्ठाय च}\ A; \skt{दन्त\appbrk{}धा\appbrk{}व\appbrk{}न\appbrk{}कू\appbrk{}र्च\appbrk{}के\appbrk{}ना\appbrk{}ह\appbrk{}रेत्~। प्रणु\appbrk{}दे\appbrk{}द्वा\appbrk{}न्त\appbrk{}क्ष\appbrk{}त\appbrk{}क\appbrk{}ण्ठाय च}\ K\textsuperscript{ac}; \skt{दन्त\appbrk{}धा\appbrk{}व\appbrk{}न\appbrk{}कू\appbrk{}र्च\appbrk{}के\appbrk{}ना\appbrk{}ह\appbrk{}रे\appbrk{}त्प्र\appbrk{}णु\appbrk{}दे\appbrk{}द्वा\appbrk{}न्तः\appbrk{}क्ष\appbrk{}त\appbrk{}क\appbrk{}ण्ठाय च}\ K\textsuperscript{pc} H; \skt{दन्त\appbrk{}धा\appbrk{}व\appbrk{}न\appbrk{}कू\appbrk{}र्च्च\appbrk{}के\appbrk{}ना\appbrk{}ह\appbrk{}रेत्~। प्रणु\appbrk{}दे\appbrk{}द्वा\appbrk{}न्तः\appbrk{}क्ष\appbrk{}त\appbrk{}क\appbrk{}ण्ठय च}\ N}}\skt{ }\edtext{\skt{मधुसर्पिषी लेढुं}}{\lemma{\skt{मधुसर्पिषी लेढुं}}\Afootnote{A; \skt{मधुसर्पिर्लेहं}\ K H; \skt{मधुसर्पिलोह}\ N}}\skt{ प्रय\-च्छे\-त्त्रि\-फ\-ला\-चूर्णं वा }\edtext{\skt{मधु\-श\-र्क\-रा\-वि\-मि\-श्रम्}}{\lemma{\skt{मधु\appbrk{}श\appbrk{}र्क\appbrk{}रा\appbrk{}वि\appbrk{}मि\appbrk{}श्रम्}}\Afootnote{A; \skt{मधु\appbrk{}श\appbrk{}र्क्क\appbrk{}रा\appbrk{}मिश्रं}\ K N H}}\skt{~॥}\pend

\pstart \ledsidenote{\passref{1.27.20}}\edtext{\edtext{\edtext{\skt{उद\-क\-पू\-र्णो\-द\-र\-म\-वा\-क्शि\-र\-स\-म\-व\-पी\-ड\-ये\-द्धु\-नी\-या\-द्वा\-म\-येद्वा}}{\lemma{\skt{उद\appbrk{}क\appbrk{}पू\appbrk{}र्णो\appbrk{}द\appbrk{}रम्}}\Afootnote{A; \skt{उदकपूर्ण्णम्}\ K H N}}}{\lemma{\skt{अवाक्शिरसम्}}\Afootnote{A; \skt{अवाञ्च्छिरसम्}\ K N H}}}{\lemma{\skt{अव\appbrk{}पी\appbrk{}ड\appbrk{}ये\appbrk{}द्धु\appbrk{}नी\appbrk{}या\appbrk{}द्वा\appbrk{}म\appbrk{}येद्वा}}\Afootnote{A; \skt{अव\appbrk{}पी\appbrk{}ड\appbrk{}ये\appbrk{}त्धु\appbrk{}नु\appbrk{}या\appbrk{}द्वा\appbrk{}म\appbrk{}येच्च}\ K N H}}\skt{ }\edtext{\skt{भस्मराशौ वा निखनेदा मुखात्}}{\lemma{\skt{भस्मराशौ वा निखनेदा मुखात्}}\Afootnote{A; \emph{om.}\ K N H}}\skt{~॥}\pend

\pstart \ledsidenote{\passref{1.27.21}}\skt{ग्रासशल्ये तु }\edtext{\skt{कण्ठासक्ते}}{\lemma{\skt{कण्ठासक्ते}}\Afootnote{A; \skt{कण्ठसक्ते}\ K; \skt{कण्ठशक्ते}\ N H}}\skt{ }\edtext{\edtext{\skt{निःश\-ङ्क\-म\-न\-व\-बुद्धं}}{\lemma{\skt{निःशङ्कम्}}\Afootnote{A K N; \skt{निःशङ्काम्}\ H}}}{\lemma{\skt{अनवबुद्धं}}\Afootnote{A; \skt{अनवबद्धं}\ K N H}}\skt{ स्कन्धे }\edtext{\skt{मुष्टि\-ना\-भि\-ह\-न्या\-त्स्नेहं मद्यं पानीयं वा}}{\lemma{\skt{मुष्टि\appbrk{}ना\appbrk{}भि\appbrk{}ह\appbrk{}न्या\appbrk{}त्स्नेहं मद्यं पानीयं वा}}\Afootnote{\skt{मुष्टि\appbrk{}नाऽ\appbrk{}भि\appbrk{}ह\appbrk{}न्यात्, स्नेहं मद्यं पानीयं वा}\ A; \skt{मुष्टिनो निह\appbrk{}न्या\appbrk{}त्स्ने\appbrk{}ह\appbrk{}सु\appbrk{}खो\appbrk{}दकं वा}\ K; \skt{मुष्ठिनो हरेत्~। स्नेह\appbrk{}सु\appbrk{}खो\appbrk{}द\appbrk{}कम्वा}\ N; \skt{मुष्ठिना निहन्यात्~। स्नेह\appbrk{}सु\appbrk{}खो\appbrk{}द\appbrk{}कम्वा}\ H}}\skt{ पाययेत्~॥}\pend

\pstart\versepar \ledsidenote{\passref{1.27.23}}\edtext{\skt{भवन्ति}}{\lemma{\skt{भवन्ति}}\Afootnote{A N H; \emph{om.}\ K}}\skt{ चात्र~। }\pend
\pstart\versepar \skt{शल्या\-कृ\-ति\-वि\-शे\-षांश्च }\edtext{\edtext{\skt{स्थानान्यावेक्ष्य}}{\lemma{\skt{स्थानान्य्}}\Afootnote{A; \skt{स्थानञ्}\ K N H}}}{\lemma{\skt{आवेक्ष्य}}\Afootnote{A; \skt{चावेक्ष्य}\ K N H}}\skt{ }\edtext{\skt{बुद्धिमान्}}{\lemma{\skt{बुद्धिमान्}}\Afootnote{A; \skt{नैकधा}\ K N H}}\skt{~। }\pend
\pstart\versepar \skt{तथा }\edtext{\skt{यन्त्रपृथक्त्वं च}}{\lemma{\skt{यन्त्रपृथक्त्वं च}}\Afootnote{A K H; \skt{यत्रपृथक्त्वञ्च}\ N}}\skt{ सम्य\-क्श\-ल्य\-म\-था\-ह\-रेत्~॥}\edtext{}{\lemma{\skt{चात्र~। शल्या\appbrk{}कृ\appbrk{}ति\appbrk{}वि\appbrk{}शे\appbrk{}षांश्च}}\Afootnote{A; \skt{शल्या\appbrk{}कृ\appbrk{}ति\appbrk{}वि\appbrk{}शे\appbrk{}षांश्च}\ K H; \skt{शल्या\appbrk{}कृ\appbrk{}ति\appbrk{}वि\appbrk{}शे\appbrk{}षाञ्च}\ N}}\pend

\pstart\versepar \ledsidenote{\passref{1.27.24}}\edtext{\skt{कर्णवन्ति तु}}{\lemma{\skt{कर्णवन्ति तु}}\Afootnote{A; \skt{कर्ण्णवन्ति च}\ K N H}}\skt{ शल्यानि }\edtext{\skt{दुःखाहार्याणि यानि च}}{\lemma{\skt{दुःखाहार्याणि यानि च}}\Afootnote{A; \skt{दुःखान्याहरणे विदुः}\ K N H}}\skt{~। }\pend
\pstart\versepar \skt{आददीत भिषक्तस्मात्तानि युक्त्या समाहितः~॥}\pend

\pstart\versepar \ledsidenote{\passref{1.27.25}}\skt{एतैरुपायैः शल्यं तु नैव }\edtext{\skt{निर्यात्यते}}{\lemma{\skt{निर्यात्यते}}\Afootnote{A; \skt{नि ते}\ K\textsuperscript{ac}; \skt{नि हार्य ते}\ K\textsuperscript{pc}; \skt{निर्हार्यते}\ N H}}\skt{ यदि~। }\pend
\pstart\versepar \skt{मत्या }\edtext{\skt{निपुणया}}{\lemma{\skt{निपुणया}}\Afootnote{A K N; \skt{निपुनया}\ H}}\skt{ वैद्यो यन्त्रयोगैश्च निर्हरेत्~॥}\pend

\pstart\versepar \ledsidenote{\passref{1.27.26}}\edtext{\skt{शोथपाकौ}}{\lemma{\skt{शोथपाकौ}}\Afootnote{A; \skt{शोफपाकौ}\ K N H}}\skt{ }\edtext{\skt{रुजश्चोग्राः}}{\lemma{\skt{चोग्राः}}\Afootnote{A K H; \skt{चोग्रा}\ N}}\skt{ }\edtext{\edtext{\skt{कुर्या\-च्छ\-ल्य\-म\-नि\-र्हृ\-तम्}}{\lemma{\skt{कुर्याच्छल्यम्}}\Afootnote{K N H; \skt{कुर्याच्छल्य}\ A}}}{\lemma{\skt{अनिर्हृतम्}}\Afootnote{K\textsuperscript{pc} H; \skt{मनाहृतम्}\ A; \skt{अनिहृतम्}\ K\textsuperscript{ac}; \skt{इनिर्हृतम्}\ N}}\skt{~। }\pend
\pstart\versepar \skt{वैकल्यं मरणं चाऽपि }\edtext{\skt{तस्मा\-द्य\-त्ना\-द्वि\-नि\-र्ह\-रेत्}}{\lemma{\skt{तस्मा\appbrk{}द्य\appbrk{}त्ना\appbrk{}द्वि\appbrk{}नि\appbrk{}र्ह\appbrk{}रेत्}}\Afootnote{\skt{तस्मा\appbrk{}द्य\appbrk{}त्ना\appbrk{}द्वि\appbrk{}नि\appbrk{}र्ह\appbrk{}रेत्}\ A; \skt{तस्मा\appbrk{}द्य\appbrk{}त्ना\appbrk{}द्वि\appbrk{}नि\appbrk{}र्ह\appbrk{}रे\appbrk{}दिति~॥ २७~॥ {\symfont ❈}}\ K; \skt{तस्मा\appbrk{}द्य\appbrk{}त्ना\appbrk{}द्वि\appbrk{}नि\appbrk{}र्ह\appbrk{}रे\appbrk{}दिति}\ N H}}\skt{~॥}\pend

\pstart \ledsidenote{\passref{1.27.trailer}}\skt{इति सुश्रु\-त\-सं\-हि\-तायां सूत्रस्थाने शल्या\-प\-न\-य\-नीयो नाम सप्त\-विं\-श\-ति\-त\-मोऽ\-ध्यायः~॥}\pend

\endnumbering
\beginnumbering
\eledsection*{Adhyāya 28}
\pstart \ledsidenote{\passref{1.28.1}}\edtext{\skt{अथातो}}{\lemma{\skt{अथातो}}\Afootnote{A K; \emph{om.}\ N H}}\skt{ }\edtext{\skt{विप\-री\-त\-व्र\-ण\-वि\-ज्ञा\-नीयं}}{\lemma{\skt{विप\appbrk{}री\appbrk{}त\appbrk{}व्र\appbrk{}ण\appbrk{}वि\appbrk{}ज्ञा\appbrk{}नीयं}}\Afootnote{K; \skt{विप\appbrk{}री\appbrk{}ता\appbrk{}वि\appbrk{}प\appbrk{}री\appbrk{}त\appbrk{}व्र\appbrk{}ण\appbrk{}वि\appbrk{}ज्ञा\appbrk{}नी\appbrk{}यम्}\ A; \skt{अथातो विप\appbrk{}री\appbrk{}त\appbrk{}व्र\appbrk{}ण\appbrk{}वि\appbrk{}ज्ञा\appbrk{}नीयं}\ N H}}\skt{ }\edtext{\skt{व्याख्यास्यामः}}{\lemma{\skt{व्याख्यास्यामः}}\Afootnote{N H; \skt{अध्यायं व्याख्यास्यामः}\ A; \emph{om.}\ K}}\skt{~॥}\pend

\pstart\versepar \ledsidenote{\passref{1.28.3}}\skt{फला\-ग्नि\-ज\-ल\-वृ\-ष्टीनां पुष्पधूमाम्बुदा यथा~। }\pend
\pstart\versepar \skt{ख्यापयन्ति भविष्यत्वं }\edtext{\skt{रिष्टानि}}{\lemma{\skt{रिष्टानि}}\Afootnote{K\textsuperscript{pc} N H; \skt{तथा रिष्टानि}\ A; \emph{om.}\ K\textsuperscript{ac}}}\skt{ }\edtext{\edtext{\skt{मरणन्तथा}}{\lemma{\skt{मरणन्}}\Afootnote{K N H; \skt{पञ्चताम्}\ A}}}{\lemma{\skt{तथा}}\Afootnote{K N H; \emph{om.}\ A}}\skt{~॥}\pend

\pstart\versepar \ledsidenote{\passref{1.28.4}}\skt{तानि }\edtext{\skt{सौक्ष्म्यात्प्रमादाद्वा}}{\lemma{\skt{सौक्ष्म्यात्प्रमादाद्वा}}\Afootnote{A K H; \skt{सौक्ष्म्यात्प्रमदाद्वा}\ N}}\skt{ तथैवाशु व्यतिक्रमात्~। }\pend
\pstart\versepar \skt{गृह्यन्ते नोद्ग\-ता\-न्य\-ज्ञै\-र्मु\-मू\-र्षोर्न त्वसंभवात्~॥}\pend

\pstart\versepar \ledsidenote{\passref{1.28.5}}\edtext{\skt{ध्रुवं}}{\lemma{\skt{ध्रुवं}}\Afootnote{A; \skt{ध्रुवं हि}\ K N H}}\skt{ }\edtext{\skt{तु मरणं}}{\lemma{\skt{तु मरणं}}\Afootnote{A; \skt{रिष्टे मरणं}\ K H N}}\skt{ }\edtext{\skt{रिष्टे ब्राह्म\-णै\-स्त\-त्कि\-ला\-मलैः}}{\lemma{\skt{रिष्टे ब्राह्मणैस्}}\Afootnote{A; \skt{ब्राह्मणैस्}\ K; \skt{ब्राह्मण्यैस्}\ N H}}\skt{~। }\pend
\pstart\versepar \skt{}\edtext{\skt{रसा\-य\-न\-त\-पो\-ज\-प्य\-त\-त्प\-रैर्वा}}{\lemma{\skt{रसा\appbrk{}य\appbrk{}न\appbrk{}त\appbrk{}पो\appbrk{}ज\appbrk{}प्य\appbrk{}त\appbrk{}त्प\appbrk{}रैर्वा}}\Afootnote{A K H; \skt{रसा\appbrk{}य\appbrk{}न\appbrk{}त\appbrk{}पो\appbrk{}जा\appbrk{}प्य\appbrk{}त\appbrk{}त्प\appbrk{}रैर्वा}\ N}}\skt{ निवार्यते~॥}\pend

\pstart\versepar \ledsidenote{\passref{1.28.6}}\skt{नक्षत्रपीडा बहुधा यथा }\edtext{\skt{कालाद्विपच्यते}}{\lemma{\skt{कालाद्विपच्यते}}\Afootnote{K\textsuperscript{pc} H; \skt{कालं विपच्यते}\ A; \skt{कालापच्यते}\ K\textsuperscript{ac}; \skt{कालाद्विपर्च्यते}\ N}}\skt{~। }\pend
\pstart\versepar \skt{}\edtext{\skt{तथै\-वा\-रि\-ष्ट\-पाकं}}{\lemma{\skt{तथै\appbrk{}वा\appbrk{}रि\appbrk{}ष्ट\appbrk{}पाकं}}\Afootnote{A; \skt{रिष्टपाकन्}\ K H N}}\skt{ }\edtext{\skt{च ब्रुवते}}{\lemma{\skt{च ब्रुवते}}\Afootnote{A; \skt{तथा ब्रुवते}\ K\textsuperscript{ac}; \skt{तथा केचित्ब्रुवते}\ K\textsuperscript{pc} N H}}\skt{ बहवो जनाः~॥}\pend

\pstart\versepar \ledsidenote{\passref{1.28.7}}\edtext{\skt{असि\-द्धि\-मा\-प्नु\-या\-ल्लोके}}{\lemma{\skt{आप्नुयाल्लोके}}\Afootnote{A; \skt{प्राप्नुयाल्लोके}\ K N H}}\skt{ }\edtext{\skt{प्रति\-कु\-र्व\-न्ग\-ता\-युषः}}{\lemma{\skt{प्रति\appbrk{}कु\appbrk{}र्व\appbrk{}न्ग\appbrk{}ता\appbrk{}युषः}}\Afootnote{A K; \skt{प्रति\appbrk{}कु\appbrk{}र्व\appbrk{}ङ्ग\appbrk{}ता\appbrk{}युषः}\ N H\textsuperscript{pc}; \skt{प्रति\appbrk{}कु\appbrk{}र्व्व\appbrk{}ङ्ग\appbrk{}ग\appbrk{}ता\appbrk{}युषः}\ H\textsuperscript{ac}}}\skt{~। }\pend
\pstart\versepar \skt{}\edtext{\skt{अतोऽरिष्टनि}}{\lemma{\skt{अतोऽरिष्टनि}}\Afootnote{\skt{अतोऽरिष्टानि}\ A; \skt{तस्माद्यत्नेन रिष्टानि}\ K N H}}\skt{ }\edtext{\skt{यत्नेन लक्षयेत्कुशलो}}{\lemma{\skt{यत्नेन लक्षयेत्कुशलो}}\Afootnote{A; \skt{लक्षयेत्कुशलो}\ K H N}}\skt{ भिषक्~॥}\pend

\pstart\versepar \ledsidenote{\passref{1.28.8}}\skt{गन्ध\-व\-र्ण\-र\-सा\-दीनां विशेषाणां स्वभावतः~। }\pend
\pstart\versepar \skt{}\edtext{\skt{वैकृतं}}{\lemma{\skt{वैकृतं}}\Afootnote{A K H; \skt{वैकृत्य}\ N}}\skt{ यत्तदाचष्टे }\edtext{\skt{व्रणिनः}}{\lemma{\skt{व्रणिनः}}\Afootnote{A K H; \skt{व्रणिन}\ N}}\skt{ }\edtext{\skt{पक्वलक्षणम्}}{\lemma{\skt{पक्वलक्षणम्}}\Afootnote{A K H; \skt{पक्व\appbrk{}प\appbrk{}ल\appbrk{}क्ष\appbrk{}णम्}\ N}}\skt{~॥}\pend

\pstart\versepar \ledsidenote{\passref{1.28.9}}\edtext{\skt{कटुस्तीक्ष्णश्च}}{\lemma{\skt{तीक्ष्णश्च}}\Afootnote{A; \skt{तीव्रश्च}\ K N H}}\skt{ विस्रश्च गन्धस्तु पवनादिभिः~। }\pend
\pstart\versepar \skt{लोहगन्धिस्तु रक्तेन }\edtext{\skt{व्यामिश्रः}}{\lemma{\skt{व्यामिश्रः}}\Afootnote{A K H; \skt{व्यामिश्र}\ N}}\skt{ }\edtext{\skt{सान्निपातिकः}}{\lemma{\skt{सान्निपातिकः}}\Afootnote{A K\textsuperscript{pc} N H; \skt{सान्निपा कः}\ K\textsuperscript{ac}}}\skt{~॥}\pend

\pstart\versepar \ledsidenote{\passref{1.28.10}}\skt{लाजा\-त\-सी\-तै\-ल\-समाः }\edtext{\skt{किंचिद्विस्राश्च}}{\lemma{\skt{किंचिद्विस्राश्च}}\Afootnote{A; \skt{किञ्चिद्विस्राश्च}\ K H; \skt{किञ्चिद्विश्राश्च}\ N}}\skt{ गन्धतः~। }\pend
\pstart\versepar \skt{}\edtext{\skt{ज्ञेयाः}}{\lemma{\skt{ज्ञेयाः}}\Afootnote{A K H; \skt{ज्ञेया}\ N}}\skt{ }\edtext{\skt{प्रकृतिगन्धाः }\edtext{\edtext{\skt{स्युर\-तोऽ\-न्य\-द्ग\-न्ध\-वै\-कृ\-तम्}}{\lemma{\skt{अतो}}\Afootnote{A K\textsuperscript{pc} N; \emph{om.}\ K\textsuperscript{ac}; \skt{ततो}\ H}}}{\lemma{\skt{ऽन्य\appbrk{}द्ग\appbrk{}न्ध\appbrk{}वै\appbrk{}कृ\appbrk{}तम्}}\Afootnote{H; \skt{ऽन्य\appbrk{}द्ग\appbrk{}न्ध\appbrk{}वै\appbrk{}कृ\appbrk{}तम्}\ A; \skt{तो न्यद्ग\appbrk{}न्ध\appbrk{}वै\appbrk{}कृ\appbrk{}तम्}\ K\textsuperscript{ac}; \skt{न्यद्ग\appbrk{}न्ध\appbrk{}वै\appbrk{}कृ\appbrk{}तम्}\ K\textsuperscript{pc} N}}}{\lemma{\skt{प्रकृतिगन्धाः स्युर्}}\Afootnote{A K\textsuperscript{pc}; \skt{प्रकृतिगन्धा}\ K\textsuperscript{ac}; \skt{प्रकृतिगन्धा स्युर्}\ N; \skt{प्रकृतिगन्धास्तु}\ H}}\skt{~॥}\pend

\pstart\versepar \ledsidenote{\passref{1.28.11}}\edtext{\skt{मध्या\-गु\-र्वा\-ज्य\-सु\-म\-नः\-प\-द्म\-च\-न्द\-न\-च\-म्पकैः}}{\lemma{\skt{मध्या\appbrk{}गु\appbrk{}र्वा\appbrk{}ज्य\appbrk{}सु\appbrk{}म\appbrk{}नः\appbrk{}प\appbrk{}द्म\appbrk{}च\appbrk{}न्द\appbrk{}न\appbrk{}च\appbrk{}म्पकैः}}\Afootnote{\skt{मद्या\appbrk{}गु\appbrk{}र्वा\appbrk{}ज्य\appbrk{}सु\appbrk{}म\appbrk{}ना\appbrk{}प\appbrk{}द्म\appbrk{}च\appbrk{}न्द\appbrk{}न\appbrk{}च\appbrk{}म्पकैः}\ A; \skt{मद्याज्ययोः सुमनसां पद्म\appbrk{}च\appbrk{}न्द\appbrk{}न\appbrk{}यो\appbrk{}रपि}\ K H N}}\skt{~। }\pend
\pstart\versepar \skt{}\edtext{\skt{सगन्धा}}{\lemma{\skt{सगन्धा}}\Afootnote{A; \skt{सुगन्धा}\ K N H}}\skt{ दिव्यगन्धाश्च मुमूर्षूणां व्रणाः स्मृताः~॥}\pend

\pstart\versepar \ledsidenote{\passref{1.28.12}}\edtext{\skt{श्ववा\-जि\-मू\-षि\-क\-ध्वा\-ङ्क्ष\-पू\-ति\-व\-ल्लू\-र\-म\-त्कुणैः}}{\lemma{\skt{श्ववा\appbrk{}जि\appbrk{}मू\appbrk{}षि\appbrk{}क\appbrk{}ध्वा\appbrk{}ङ्क्ष\appbrk{}पू\appbrk{}ति\appbrk{}व\appbrk{}ल्लू\appbrk{}र\appbrk{}म\appbrk{}त्कुणैः}}\Afootnote{A; \skt{श्ववा\appbrk{}जि\appbrk{}मू\appbrk{}षि\appbrk{}का\appbrk{}ध्वां\appbrk{}क्ष\appbrk{}पू\appbrk{}ति\appbrk{}व\appbrk{}ल्लू\appbrk{}र\appbrk{}म\appbrk{}त्कुणैः}\ K H; \skt{श्ववा\appbrk{}जि\appbrk{}मू\appbrk{}षि\appbrk{}का\appbrk{}ध्वां\appbrk{}क्ष\appbrk{}पू\appbrk{}ति\appbrk{}व\appbrk{}ल्लू\appbrk{}र\appbrk{}मं\appbrk{}कुणैः}\ N}}\skt{~। }\pend
\pstart\versepar \skt{सगन्धाः पङ्कगन्धाश्च भूमिगन्धाश्च गर्हिताः~॥}\pend

\pstart\versepar \ledsidenote{\passref{1.28.13}}\edtext{\skt{कुङ्कु\-म\-ध्या\-म\-कं\-गुष्ठः सवर्ण्णाः}}{\lemma{\skt{कुङ्कु\appbrk{}म\appbrk{}ध्या\appbrk{}म\appbrk{}कं\appbrk{}गुष्ठः सवर्ण्णाः}}\Afootnote{H; \skt{कुङ्कु\appbrk{}म\appbrk{}ध्या\appbrk{}म\appbrk{}क\appbrk{}ङ्कु\appbrk{}ष्ठ\appbrk{}स\appbrk{}वर्णाः}\ A; \skt{कुंकु\appbrk{}म\appbrk{}ध्या\appbrk{}म\appbrk{}कां\appbrk{}गु\appbrk{}ष्ठ\appbrk{}स\appbrk{}वर्ण्णाः}\ K\textsuperscript{ac}; \skt{कुंकु\appbrk{}म\appbrk{}ध्या\appbrk{}म\appbrk{}कं\appbrk{}गु\appbrk{}ष्ठ\appbrk{}स\appbrk{}वर्ण्णाः}\ K\textsuperscript{pc} N}}\skt{ पित्तजाश्च ये~। }\pend
\pstart\versepar \skt{न }\edtext{\skt{दह्यन्ते}}{\lemma{\skt{दह्यन्ते}}\Afootnote{\skt{दह्यन्ते न}\ A K N H}}\skt{ }\edtext{\skt{शुष्यन्ति}}{\lemma{\skt{शुष्यन्ति}}\Afootnote{K N H; \skt{चूष्यन्ते भिषक्तान्परि}\ A}}\skt{ }\edtext{\skt{वर्ज्ज\-ये\-त्ता\-न्वि\-च\-क्षणः}}{\lemma{\skt{वर्ज्ज\appbrk{}ये\appbrk{}त्ता\appbrk{}न्वि\appbrk{}च\appbrk{}क्षणः}}\Afootnote{K H; \skt{वर्जयेत्}\ A; \skt{वर्जये ताम्विचक्षणः}\ N}}\skt{~॥}\edtext{}{\lemma{\skt{पित्तजाश्च ये~। न}}\Afootnote{K N; \skt{पित्तकोपतः~। न}\ A; \skt{पित्तजाश्च ये}\ H}}\pend

\pstart\versepar \ledsidenote{\passref{1.28.14}}\edtext{\skt{कण्डूमन्तः}}{\lemma{\skt{कण्डूमन्तः}}\Afootnote{K N H; \skt{अण्डूमन्तः}\ A}}\skt{ स्थिराः }\edtext{\skt{श्वेताः}}{\lemma{\skt{श्वेताः}}\Afootnote{A; \skt{स्निग्धाः}\ K N H}}\skt{ }\edtext{\skt{स्निग्धाः}}{\lemma{\skt{स्निग्धाः}}\Afootnote{A; \skt{श्वेताः}\ K; \skt{स्वेताः}\ N H}}\skt{ }\edtext{\skt{कफनिमित्ततः}}{\lemma{\skt{कफनिमित्ततः}}\Afootnote{A; \skt{कफनिमित्तजाः}\ K N H}}\skt{~। }\pend
\pstart\versepar \skt{}\edtext{\skt{दूयन्ते}}{\lemma{\skt{दूयन्ते}}\Afootnote{A; \skt{दह्यन्ते}\ K N H}}\skt{ }\edtext{\skt{वाऽपि दह्यन्ते}}{\lemma{\skt{वाऽपि दह्यन्ते}}\Afootnote{A; \skt{वेदनावन्तो}\ K N H}}\skt{ }\edtext{\skt{भिष\-क्ता\-न्प\-रि\-व\-र्ज\-येत्}}{\lemma{\skt{भिष\appbrk{}क्ता\appbrk{}न्प\appbrk{}रि\appbrk{}व\appbrk{}र्ज\appbrk{}येत्}}\Afootnote{A; \skt{भिषक्तानपि वर्जयेत्}\ K N H}}\skt{~॥}\pend

\pstart\versepar \ledsidenote{\passref{1.28.15}}\edtext{\skt{कृष्णास्तु ये}}{\lemma{\skt{कृष्णास्तु ये}}\Afootnote{A; \skt{कृष्णा ये}\ K N H}}\skt{ }\edtext{\skt{तनुस्रावा}}{\lemma{\skt{तनुस्रावा}}\Afootnote{A; \skt{त्वरुणास्रावा}\ K H; \skt{त्वरुणाश्रा}\ N}}\skt{ }\edtext{\skt{वातजा}}{\lemma{\skt{वातजा}}\Afootnote{A K; \skt{वाताजा}\ N H}}\skt{ मर्मतापिनः~। }\pend
\pstart\versepar \skt{स्वल्पामपि न कुर्वन्ति रुजं }\edtext{\skt{तान्प\-रि\-व\-र्ज\-येत्}}{\lemma{\skt{तान्प\appbrk{}रि\appbrk{}व\appbrk{}र्ज\appbrk{}येत्}}\Afootnote{A; \skt{तानपि वर्जयेत्}\ K N H}}\skt{~॥}\pend

\pstart\versepar \ledsidenote{\passref{1.28.16}}\skt{क्ष्वेडन्ति }\edtext{\skt{घुर्घुरायन्ते}}{\lemma{\skt{घुर्घुरायन्ते}}\Afootnote{A; \skt{घुर्घुरायन्ति}\ K N H}}\skt{ ज्वलन्तीव च ये व्रणाः~। }\pend
\pstart\versepar \skt{}\edtext{\skt{त्वङ्मांसस्थाश्च}}{\lemma{\skt{त्वङ्मांसस्थाश्च}}\Afootnote{A K H; \skt{त्वङ्मान्सस्था}\ N}}\skt{ पवनं सशब्दं विसृजन्ति ये~॥}\pend

\pstart\versepar \ledsidenote{\passref{1.28.17}}\edtext{\skt{ये च}}{\lemma{\skt{ये च}}\Afootnote{A H; \skt{ये त्व}\ K; \skt{ये}\ N}}\skt{ }\edtext{\skt{मर्मस्वसंभूता}}{\lemma{\skt{मर्मस्वसंभूता}}\Afootnote{A K H; \skt{त्व्मर्मस्वसंभूता}\ N}}\skt{ भव\-न्त्य\-त्य\-र्थ\-वे\-दनाः~। }\pend
\pstart\versepar \skt{दह्यन्ते }\edtext{\skt{चान्तरत्यर्थं}}{\lemma{\skt{चान्तरत्यर्थं}}\Afootnote{A K N; \skt{चान्तरात्यर्थम्}\ H}}\skt{ }\edtext{\skt{बहिः}}{\lemma{\skt{बहिः}}\Afootnote{A K H; \skt{ब्रहि}\ N}}\skt{ शीताश्च ये व्रणाः~॥}\pend

\pstart\versepar \ledsidenote{\passref{1.28.18}}\skt{दह्यन्ते }\edtext{\skt{बहिरत्यर्थं}}{\lemma{\skt{अत्यर्थं}}\Afootnote{A K H; \skt{अत्यर्थ}\ N}}\skt{ }\edtext{\edtext{\skt{भवन्त्यन्तश्च}}{\lemma{\skt{भवन्त्य्}}\Afootnote{K N H; \skt{भवन्}\ A}}}{\lemma{\skt{अन्तश्च}}\Afootnote{K N H; \skt{त्यन्तश्च}\ A}}\skt{ शीतलाः~। }\pend
\pstart\versepar \skt{}\edtext{\skt{शक्तिध्वजरथा}}{\lemma{\skt{शक्तिध्वजरथा}}\Afootnote{\skt{शक्तिध्वजरथाः}\ A; \skt{शक्ति\appbrk{}कु\appbrk{}न्त\appbrk{}ध्व\appbrk{}ज\appbrk{}रथा}\ K H; \skt{शक्ति\appbrk{}कु\appbrk{}न्त\appbrk{}ध्व\appbrk{}जा\appbrk{}रथा}\ N}}\skt{ }\edtext{\skt{कुन्त\-वा\-जि\-वा\-र\-ण\-गो\-वृषाः}}{\lemma{\skt{कुन्त\appbrk{}वा\appbrk{}जि\appbrk{}वा\appbrk{}र\appbrk{}ण\appbrk{}गो\appbrk{}वृषाः}}\Afootnote{A; \skt{वाजि\appbrk{}वा\appbrk{}र\appbrk{}ण\appbrk{}प\appbrk{}क्षिणः}\ K N\textsuperscript{pc} H; \skt{वाजिवारणा पक्षिणः}\ N\textsuperscript{ac}}}\skt{~॥}\pend

\pstart\versepar \ledsidenote{\passref{1.28.19}}\skt{येषु }\edtext{\skt{चाप्य\-व\-भा\-से\-र\-न्प्रा\-सा\-दा\-कृ\-त\-य\-स्तथा}}{\lemma{\skt{अव\appbrk{}भा\appbrk{}से\appbrk{}र\appbrk{}न्प्रा\appbrk{}सा\appbrk{}दा\appbrk{}कृ\appbrk{}त\appbrk{}यस्}}\Afootnote{A; \skt{अवभासेयुः प्रासा\appbrk{}दा\appbrk{}कृ\appbrk{}त\appbrk{}यस्}\ K; \skt{अवभासेयुः प्रासा\appbrk{}द\appbrk{}कृ\appbrk{}त\appbrk{}यस्}\ N; \skt{एव भाषेयुः प्रासा\appbrk{}दा\appbrk{}कृ\appbrk{}त\appbrk{}यस्}\ H}}\skt{~। }\pend
\pstart\versepar \skt{चूर्णावकीर्णा इव ये भान्ति }\edtext{\skt{वा न च चूर्णिताः}}{\lemma{\skt{वा न च चूर्णिताः}}\Afootnote{\skt{चानवचूर्णिताः}\ A K N H}}\skt{~॥}\pend

\pstart\versepar \ledsidenote{\passref{1.28.20}}\skt{प्राण\-मां\-स\-क्ष\-य\-श्वा\-स\-का\-सा\-रो\-च\-क\-पी\-डिताः~। }\pend
\pstart\versepar \skt{प्रवृ\-द्ध\-पू\-य\-रु\-धिरा व्रणा येषां च मर्मसु~॥}\pend

\pstart\versepar \ledsidenote{\passref{1.28.21}}\edtext{\skt{क्रियाभिः}}{\lemma{\skt{क्रियाभिः}}\Afootnote{A K H; \skt{क्रियाभि}\ N}}\skt{ सम्यगारब्धा न सिध्यन्ति च ये व्रणाः~। }\pend
\pstart\versepar \skt{}\edtext{\skt{वर्ज\-ये\-त्ता\-न्भि\-ष\-क्प्राज्ञः}}{\lemma{\skt{वर्ज\appbrk{}ये\appbrk{}त्ता\appbrk{}न्भि\appbrk{}ष\appbrk{}क्प्राज्ञः}}\Afootnote{\skt{वर्जयेत्तानपि प्राज्ञः}\ A; \skt{वर्जयेदपि तां प्राज्ञः}\ K N H}}\skt{ }\edtext{\skt{संरक्षन्नात्मनो}}{\lemma{\skt{संरक्षन्नात्मनो}}\Afootnote{A; \skt{संरक्षन्नात्मनो}\ K H; \skt{सरक्षनात्मनो}\ N}}\skt{ यशः~॥}\pend

\pstart \ledsidenote{\passref{1.28.trailer}}\edtext{\skt{इति सुश्रु\-त\-सं\-हि\-तायां सूत्रस्थाने विप\-री\-ता\-वि\-प\-री\-त\-व्र\-ण\-वि\-ज्ञा\-नीयो नामा\-ष्ट\-विं\-श\-ति\-त\-मोऽ\-ध्यायः}}{\lemma{\skt{इति सुश्रु\appbrk{}त\appbrk{}सं\appbrk{}हि\appbrk{}तायां सूत्रस्थाने विप\appbrk{}री\appbrk{}ता\appbrk{}वि\appbrk{}प\appbrk{}री\appbrk{}त\appbrk{}व्र\appbrk{}ण\appbrk{}वि\appbrk{}ज्ञा\appbrk{}नीयो नामा\appbrk{}ष्ट\appbrk{}विं\appbrk{}श\appbrk{}ति\appbrk{}त\appbrk{}मोऽ\appbrk{}ध्यायः}}\Afootnote{\skt{इति सुश्रु\appbrk{}त\appbrk{}सं\appbrk{}हि\appbrk{}तायां सूत्रस्थाने विप\appbrk{}री\appbrk{}ता\appbrk{}वि\appbrk{}प\appbrk{}री\appbrk{}त\appbrk{}व्र\appbrk{}ण\appbrk{}वि\appbrk{}ज्ञा\appbrk{}नीयो नामा\appbrk{}ष्टा\appbrk{}विं\appbrk{}श\appbrk{}ति\appbrk{}त\appbrk{}मोऽ\appbrk{}ध्यायः}\ A; \skt{२८~॥ {\symfont ❈}}\ K; \skt{॥ २८~॥}\ N; \skt{२८~॥}\ H}}\skt{~॥}\pend

\endnumbering
\beginnumbering
\eledsection*{Adhyāya 29}
\pstart \ledsidenote{\passref{1.29.1}}\edtext{\skt{अथातो}}{\lemma{\skt{अथातो}}\Afootnote{A K H; \skt{अथात}\ N}}\skt{ }\edtext{\edtext{\skt{विप\-री\-ता\-वि\-प\-री\-त\-स्व\-प्न\-नि\-द\-र्श\-नी\-य\-म\-ध्यायं}}{\lemma{\skt{विप\appbrk{}री\appbrk{}ता\appbrk{}वि\appbrk{}प\appbrk{}री\appbrk{}त\appbrk{}स्व\appbrk{}प्न\appbrk{}नि\appbrk{}द\appbrk{}र्श\appbrk{}नी\appbrk{}यम्}}\Afootnote{A; \skt{विप\appbrk{}री\appbrk{}त\appbrk{}दू\appbrk{}त\appbrk{}स्व\appbrk{}प्न\appbrk{}द\appbrk{}र्श\appbrk{}नी\appbrk{}यम्}\ K N; \skt{विप\appbrk{}री\appbrk{}त\appbrk{}दू\appbrk{}त\appbrk{}स्व\appbrk{}प्न\appbrk{}द\appbrk{}र्शनी}\ H}}}{\lemma{\skt{अध्यायं}}\Afootnote{A; \emph{om.}\ K; \skt{व्याख्यामः}\ N; \skt{यम्}\ H}}\skt{ }\edtext{\skt{व्याख्यास्यामः}}{\lemma{\skt{व्याख्यास्यामः}}\Afootnote{A H; \emph{om.}\ K N}}\skt{~॥}\pend

\pstart\versepar \ledsidenote{\passref{1.29.3}}\skt{दूत\-द\-र्श\-न\-सं\-भाषा }\edtext{\skt{वेषा\-श्चे\-ष्टि\-त\-मेव}}{\lemma{\skt{वेषाश्चेष्टितम्}}\Afootnote{A; \skt{वेशचेष्टितम्}\ K H; \skt{वेषचेष्टितम्}\ N}}\skt{ च~। }\pend
\pstart\versepar \skt{}\edtext{\skt{ऋक्षं वेला}}{\lemma{\skt{ऋक्षं वेला}}\Afootnote{A; \skt{ऋक्षवेला}\ K H N}}\skt{ तिथिश्चैव }\edtext{\skt{निमित्तं}}{\lemma{\skt{निमित्तं}}\Afootnote{A H; \skt{निमित्त}\ K N}}\skt{ }\edtext{\skt{शकुनोऽनिलः}}{\lemma{\skt{ऽनिलः}}\Afootnote{A N H; \skt{निलः}\ K}}\skt{~॥}\pend

\pstart\versepar \ledsidenote{\passref{1.29.4}}\edtext{\skt{देशो}}{\lemma{\skt{देशो}}\Afootnote{A; \skt{गेहे}\ K N H}}\skt{ वैद्यस्य वाग्देहमनसां च विचेष्टितम्~। }\pend
\pstart\versepar \skt{कथ\-य\-न्त्या\-तु\-र\-गतं शुभं वा यदि वाऽशुभम्~॥}\pend

\pstart\versepar \ledsidenote{\passref{1.29.5}}\edtext{\skt{पाख\-ण्डा\-श्र\-म\-व\-र्णानां}}{\lemma{\skt{पाख\appbrk{}ण्डा\appbrk{}श्र\appbrk{}म\appbrk{}व\appbrk{}र्णानां}}\Afootnote{\skt{पाख(\appbrk{}ष)\appbrk{}ण्डा\appbrk{}श्र\appbrk{}म\appbrk{}व\appbrk{}र्णानां}\ A; \skt{पाष\appbrk{}ण्डा\appbrk{}श्र\appbrk{}म\appbrk{}व\appbrk{}र्ण्णानां}\ K N H}}\skt{ }\edtext{\skt{सपक्षाः}}{\lemma{\skt{सपक्षाः}}\Afootnote{A; \skt{सवर्ण्णाः}\ K H; \skt{सवर्ण्ण}\ N}}\skt{ कर्मसिद्धये~। }\pend
\pstart\versepar \skt{त एव }\edtext{\skt{विपरीताः }\edtext{\skt{स्युर्दूताः}}{\lemma{\skt{दूताः}}\Afootnote{A K N\textsuperscript{pc} H; \skt{दूता}\ N\textsuperscript{ac}}}}{\lemma{\skt{विपरीताः स्युर्}}\Afootnote{A H; \skt{विपरीताः स्यु}\ K; \skt{विपरीता स्युर्}\ N}}\skt{ कर्मविपत्तये~॥}\pend

\pstart\versepar \ledsidenote{\passref{1.29.6}}\edtext{\skt{नपुंसकं स्त्री}}{\lemma{\skt{नपुंसकं स्त्री}}\Afootnote{A; \skt{नपुंसकः स्त्री}\ K N; \skt{नपुंसक स्त्री}\ H}}\skt{ बहवो }\edtext{\skt{नैककार्या असूयकाः}}{\lemma{\skt{नैककार्या असूयकाः}}\Afootnote{A; \skt{न्यथा कार्यानुसूयकाः}\ K N H}}\skt{~। }\pend
\pstart\versepar \skt{}\edtext{\skt{गर्द\-भो\-ष्ट्र\-र\-थ\-प्राप्ताः प्राप्ता वा स्युः}}{\lemma{\skt{गर्द\appbrk{}भो\appbrk{}ष्ट्र\appbrk{}र\appbrk{}थ\appbrk{}प्राप्ताः प्राप्ता वा स्युः}}\Afootnote{A; \skt{गर्द\appbrk{}भो\appbrk{}ष्ट्र\appbrk{}र\appbrk{}था\appbrk{}रूढा रुदन्त्यः सन्ध्यअयोस्तथा}\ K; \skt{गर्द\appbrk{}भो\appbrk{}ष्ट्र\appbrk{}र\appbrk{}था\appbrk{}रूढा रुदन्त्यः सन्ध्ययोस्तथा}\ N H}}\skt{ }\edtext{\skt{परम्पराः}}{\lemma{\skt{परम्पराः}}\Afootnote{A; \emph{om.}\ K N H}}\skt{~॥}\pend

\pstart\versepar \ledsidenote{\passref{1.29.7}}\skt{वैद्यं य उपसर्पन्ति }\edtext{\skt{दूतास्ते}}{\lemma{\skt{दूतास्ते}}\Afootnote{A; \skt{दूतास्ते}\ K N H}}\skt{ चापि गर्हिताः~। }\pend
\pstart\versepar \skt{}\edtext{\skt{पाश\-द\-ण्डा\-यु\-ध\-धराः}}{\lemma{\skt{पाश\appbrk{}द\appbrk{}ण्डा\appbrk{}यु\appbrk{}ध\appbrk{}धराः}}\Afootnote{A H; \skt{पाष\appbrk{}द\appbrk{}ण्डा\appbrk{}यु\appbrk{}ध\appbrk{}धरा}\ K; \skt{पाष\appbrk{}द\appbrk{}ण्डा\appbrk{}यु\appbrk{}ध\appbrk{}धराः}\ N}}\skt{ }\edtext{\skt{पाण्डु\-रे\-त\-र\-वा\-ससः}}{\lemma{\skt{पाण्डु\appbrk{}रे\appbrk{}त\appbrk{}र\appbrk{}वा\appbrk{}ससः}}\Afootnote{A; \skt{शुक्लै\appbrk{}त\appbrk{}र\appbrk{}नि\appbrk{}वा\appbrk{}सनः}\ K N H}}\skt{~॥}\pend

\pstart\versepar \ledsidenote{\passref{1.29.8}}\edtext{\skt{आर्द्र\-जी\-र्णा\-प\-स\-व्यै\-क\-म\-लि\-नो\-द्ध्व\-स्त\-वा\-ससः}}{\lemma{\skt{आर्द्र\appbrk{}जी\appbrk{}र्णा\appbrk{}प\appbrk{}स\appbrk{}व्यै\appbrk{}क\appbrk{}म\appbrk{}लि\appbrk{}नो\appbrk{}द्ध्व\appbrk{}स्त\appbrk{}वा\appbrk{}ससः}}\Afootnote{A; \skt{आर्द्र\appbrk{}जी\appbrk{}र्ण्णा\appbrk{}व\appbrk{}स\appbrk{}व्ये\appbrk{}क\appbrk{}म\appbrk{}लि\appbrk{}नोध्व वाससः}\ K\textsuperscript{ac}; \skt{आर्द्र\appbrk{}जी\appbrk{}र्ण्णा\appbrk{}व\appbrk{}स\appbrk{}व्ये\appbrk{}क\appbrk{}म\appbrk{}लि\appbrk{}नोध्व त वाससः}\ K\textsuperscript{pc} N; \skt{आर्द्द्र\appbrk{}जी\appbrk{}र्ण्णा\appbrk{}व\appbrk{}स\appbrk{}व्ये\appbrk{}क\appbrk{}म\appbrk{}लि\appbrk{}नो\appbrk{}ध्व\appbrk{}स्त\appbrk{}वा\appbrk{}ससः}\ H}}\skt{~। }\pend
\pstart\versepar \skt{}\edtext{\skt{न्यूनाधिकाङ्गा उद्विग्ना}}{\lemma{\skt{न्यूनाधिकाङ्गा उद्विग्ना}}\Afootnote{A; \skt{न्यूना\appbrk{}धि\appbrk{}का\appbrk{}ङ्गा\appbrk{}ति\appbrk{}कृष्णा}\ K H N}}\skt{ विकृता रौद्ररूपिणः~॥}\pend

\pstart\versepar \ledsidenote{\passref{1.29.9}}\edtext{\edtext{\skt{रूक्ष\-नि\-ष्ठु\-र\-वा\-दा\-श्चा\-प्य\-मा\-ङ्ग\-ल्या\-भि\-धा\-यिनः}}{\lemma{\skt{रूक्ष\appbrk{}नि\appbrk{}ष्ठु\appbrk{}र\appbrk{}वा\appbrk{}दा\appbrk{}श्चाप्य्}}\Afootnote{\skt{रूक्ष\appbrk{}नि\appbrk{}ष्ठु\appbrk{}र\appbrk{}व\appbrk{}क्तार}\ A; \skt{रूक्ष\appbrk{}नि\appbrk{}ष्ठु\appbrk{}र\appbrk{}व\appbrk{}क्ता\appbrk{}रस्}\ K N H}}}{\lemma{\skt{अमा\appbrk{}ङ्ग\appbrk{}ल्या\appbrk{}भि\appbrk{}धा\appbrk{}यिनः}}\Afootnote{\skt{स्त्वम\appbrk{}ङ्ग\appbrk{}ल्या\appbrk{}भि\appbrk{}धा\appbrk{}यिनः}\ A; \skt{त्वम\appbrk{}ङ्ग\appbrk{}ल्या\appbrk{}भि\appbrk{}धा\appbrk{}यिनः}\ K N H}}\skt{~। }\pend
\pstart\versepar \skt{}\edtext{\edtext{\skt{छिन्द\-न्त\-स्तृ\-ण\-का\-ष्ठानि}}{\lemma{\skt{छिन्दन्तस्}}\Afootnote{K H; \skt{छिन्दन्त}\ A; \skt{च्छिदन्त}\ N}}}{\lemma{\skt{तृणकाष्ठानि}}\Afootnote{K N H; \skt{स्तृणकाष्ठानि}\ A}}\skt{ स्पृशन्तो }\edtext{\skt{नासिकां स्तनम्}}{\lemma{\skt{नासिकां स्तनम्}}\Afootnote{A; \skt{नासिकास्तनम्}\ K H; \skt{नासिकास्तनः}\ N}}\skt{~॥}\pend

\pstart\versepar \ledsidenote{\passref{1.29.10}}\edtext{\skt{वस्त्रा\-न्ता\-ना\-मि\-का\-के\-श\-न\-ख\-रो\-म\-द\-शा\-स्पृशः}}{\lemma{\skt{वस्त्रा\appbrk{}न्ता\appbrk{}ना\appbrk{}मि\appbrk{}का\appbrk{}के\appbrk{}श\appbrk{}न\appbrk{}ख\appbrk{}रो\appbrk{}म\appbrk{}द\appbrk{}शा\appbrk{}स्पृशः}}\Afootnote{A; \skt{वस्त्रा\appbrk{}न्ता\appbrk{}ना\appbrk{}मि\appbrk{}का\appbrk{}के\appbrk{}श\appbrk{}न\appbrk{}ख\appbrk{}रोम ं स्पृशन्~। लिखन्तो वा महीं किंचिन्मुञ्चन्तो लोष्टभेदिनः}\ K\textsuperscript{ac}; \skt{वस्त्रा\appbrk{}न्ता\appbrk{}ना\appbrk{}मि\appbrk{}का\appbrk{}के\appbrk{}श\appbrk{}न\appbrk{}ख\appbrk{}रोम दशां स्पृशन्लिखन्तो वा महीं किंचिन्मुञ्चन्तो लोष्टभेदिनः}\ K\textsuperscript{pc} H; \skt{वस्त्रा\appbrk{}न्ता\appbrk{}ना\appbrk{}मि\appbrk{}का\appbrk{}के\appbrk{}श\appbrk{}न\appbrk{}ख\appbrk{}रो\appbrk{}म\appbrk{}द\appbrk{}शा\appbrk{}न्स्पृ\appbrk{}शन्~। लिखन्तो वा मही\appbrk{}ङ्कि\appbrk{}ञ्चि\appbrk{}त्मु\appbrk{}ञ्चन्तो लोष्टभेदिनः}\ N}}\skt{~। }\pend
\pstart\versepar \skt{}\edtext{\skt{स्रोतो\-व\-रो\-ध\-हृ\-द्ग\-ण्ड\-मू\-र्धो\-रः\-कु\-क्षि\-पा\-णयः}}{\lemma{\skt{स्रोतो\appbrk{}व\appbrk{}रो\appbrk{}ध\appbrk{}हृ\appbrk{}द्ग\appbrk{}ण्ड\appbrk{}मू\appbrk{}र्धो\appbrk{}रः\appbrk{}कु\appbrk{}क्षि\appbrk{}पा\appbrk{}णयः}}\Afootnote{A; \emph{om.}\ K N H}}\skt{~॥}\pend

\pstart\versepar \ledsidenote{\passref{1.29.11}}\skt{कपा\-लो\-प\-ल\-भ\-स्मा\-स्थि\-तु\-षा\-ङ्गा\-र\-क\-राश्च ये~। }\pend
\pstart\versepar \skt{विलिखन्तो महीं किंचिन्मुञ्चन्तो लोष्टभेदिनः~॥}\pend

\pstart\versepar \ledsidenote{\passref{1.29.12}}\edtext{\skt{तैल\-क\-र्द\-म\-दि\-ग्धाङ्गा}}{\lemma{\skt{तैल\appbrk{}क\appbrk{}र्द\appbrk{}म\appbrk{}दि\appbrk{}ग्धाङ्गा}}\Afootnote{A K H; \skt{तैल\appbrk{}क\appbrk{}र्द\appbrk{}म\appbrk{}लि\appbrk{}प्ताङ्गा}\ N}}\skt{ रक्त\-स्र\-ग\-नु\-ले\-पनाः~। }\pend
\pstart\versepar \skt{फलं पक्वमसारं वा }\edtext{\skt{गृहीत्वाऽन्यच्च}}{\lemma{\skt{गृहीत्वाऽन्यच्च}}\Afootnote{A K\textsuperscript{pc} H; \skt{गृहीत्वा}\ K\textsuperscript{ac}; \skt{गृहीत्यदप}\ N}}\skt{ तद्विधम्~॥}\pend

\pstart\versepar \ledsidenote{\passref{1.29.13}}\edtext{\skt{नखैर्नखान्तरं}}{\lemma{\skt{नखान्तरं}}\Afootnote{A; \skt{नखाङ्करञ्}\ K N H}}\skt{ }\edtext{\skt{वाऽपि}}{\lemma{\skt{वाऽपि}}\Afootnote{A; \skt{चापि}\ K N H}}\skt{ करेण चरणं तथा~। }\pend
\pstart\versepar \skt{}\edtext{\skt{उपा\-न\-च्च\-र्म\-हस्ता वा}}{\lemma{\skt{उपा\appbrk{}न\appbrk{}च्च\appbrk{}र्म\appbrk{}हस्ता वा}}\Afootnote{A K H; \skt{उपा\appbrk{}न\appbrk{}च\appbrk{}र्म\appbrk{}हस्ता वा}\ N}}\skt{ }\edtext{\skt{विकृ\-त\-व्या\-धि\-पी\-डिताः}}{\lemma{\skt{विकृ\appbrk{}त\appbrk{}व्या\appbrk{}धि\appbrk{}पी\appbrk{}डिताः}}\Afootnote{A; \skt{विकृताः व्याधिपीडिताः}\ K N H}}\skt{~॥}\pend

\pstart\versepar \ledsidenote{\passref{1.29.14}}\skt{वामाचारा }\edtext{\skt{रुदन्तश्च}}{\lemma{\skt{रुदन्तश्च}}\Afootnote{A; \skt{रुदन्तो}\ K\textsuperscript{ac}; \skt{रुदन्तो वा}\ K\textsuperscript{pc} N H}}\skt{ श्वासिनो }\edtext{\skt{विकृतेक्षणाः}}{\lemma{\skt{विकृतेक्षणाः}}\Afootnote{A K; \skt{विवृत्तेक्षणाः}\ N H}}\skt{~। }\pend
\pstart\versepar \skt{याम्यां }\edtext{\skt{दिशं}}{\lemma{\skt{दिशं}}\Afootnote{K N H; \skt{दिशि}\ A}}\skt{ प्राञ्जलयो }\edtext{\skt{विषमैकपदे}}{\lemma{\skt{विषमैकपदे}}\Afootnote{A; \skt{विषमैकपदैः}\ K H; \skt{विषमेकपदे}\ N}}\skt{ }\edtext{\skt{स्थिताः}}{\lemma{\skt{स्थिताः}}\Afootnote{A N H; \skt{स्थिता}\ K}}\skt{~॥}\pend

\pstart\versepar \ledsidenote{\passref{1.29.15abcd}}\skt{वैद्यं य उपसर्पन्ति दूतास्ते चापि गर्हिताः~। }\pend
\pstart\versepar \skt{दक्षिणाभिमुखं देशे }\edtext{\skt{मलिने}}{\lemma{\skt{मलिने}}\Afootnote{K N H; \emph{om.}\ A}}\skt{ }\edtext{\skt{क्रूरकर्मिणं~। \textasciicircum{}}}{\lemma{\skt{क्रूरकर्मिणं~। \textasciicircum{}}}\Afootnote{\skt{त्वशुचौ वा हुताशनम्}\ A; \skt{क्रूरकर्म्मिणं}\ K H N}}\pend

\pstart\versepar \ledsidenote{\passref{1.29.16}}\edtext{\skt{नग्नं भूमौ}}{\lemma{\skt{नग्नं भूमौ}}\Afootnote{A; \skt{भूमौ}\ K N H}}\skt{ शयानं }\edtext{\skt{वा}}{\lemma{\skt{वा}}\Afootnote{A; \skt{नग्नं वा}\ K N H}}\skt{ वेगोत्सर्गेषु वाऽशुचिम्~। }\pend
\pstart\versepar \skt{}\edtext{\skt{प्रकी\-र्णा\-के\-श\-म\-भ्यक्तं}}{\lemma{\skt{प्रकीर्णाकेशम्}}\Afootnote{\skt{प्रकीर्णकेशम्}\ A K N H}}\skt{ स्विन्नं विक्लवमेव वा~॥}\pend

\pstart\versepar \ledsidenote{\passref{1.29.17}}\skt{वैद्यं य उपसर्पन्ति दूतास्ते चापि }\edtext{\skt{गर्हिताः}}{\lemma{\skt{गर्हिताः}}\Afootnote{A; \skt{गर्हिताः}\ K N H}}\skt{~। }\pend
\pstart\versepar \skt{}\edtext{\skt{वैद्यस्य पैत्र्ये}}{\lemma{\skt{वैद्यस्य पैत्र्ये}}\Afootnote{A; \skt{पित्र्ये}\ K N H}}\skt{ }\edtext{\skt{दैवे वा कार्ये}}{\lemma{\skt{दैवे वा कार्ये}}\Afootnote{A; \skt{वा बहुकार्ये वा}\ K N H}}\skt{ }\edtext{\skt{चोत्पातदर्शने}}{\lemma{\skt{चोत्पातदर्शने}}\Afootnote{A; \skt{तथा चोत्पातदर्शने}\ K N H}}\skt{~॥}\pend

\pstart\versepar \ledsidenote{\passref{1.29.18}}\edtext{\skt{मध्याह्ने}}{\lemma{\skt{मध्याह्ने}}\Afootnote{A N H; \skt{मध्या…}\ K}}\skt{ }\edtext{\skt{चार्धरात्रे वा}}{\lemma{\skt{चार्धरात्रे वा}}\Afootnote{A N H; \skt{…र्धरात्रे वा}\ K}}\skt{ सन्ध्ययोः }\edtext{\skt{कृत्तिकासु वा}}{\lemma{\skt{कृत्तिकासु वा}}\Afootnote{K N H; \skt{कृत्तिकासु च}\ A}}\skt{~। }\pend
\pstart\versepar \skt{}\edtext{\skt{आर्द्रा\-श्ले\-षा\-म\-घा\-मू\-ला\-पू\-र्वासु}}{\lemma{\skt{आर्द्रा\appbrk{}श्ले\appbrk{}षा\appbrk{}म\appbrk{}घा\appbrk{}मू\appbrk{}ला\appbrk{}पू\appbrk{}र्वासु}}\Afootnote{K N\textsuperscript{pc} H; \skt{आर्द्रा\appbrk{}श्ले\appbrk{}षा\appbrk{}म\appbrk{}घा\appbrk{}मू\appbrk{}ल\appbrk{}पू\appbrk{}र्वासु}\ A; \emph{om.}\ N\textsuperscript{ac}}}\skt{ }\edtext{\skt{भरणीषु}}{\lemma{\skt{भरणीषु}}\Afootnote{A K N\textsuperscript{pc} H; \emph{om.}\ N\textsuperscript{ac}}}\skt{ }\edtext{\skt{वा}}{\lemma{\skt{वा}}\Afootnote{K N\textsuperscript{pc} H; \skt{च}\ A; \emph{om.}\ N\textsuperscript{ac}}}\skt{~॥}\footnoteA{This passage and 1.29.19ab were known to the Chinese translators An Shigao (ca. 148-180 CE) and Dharmarakṣa (fl. 284 CE) (Lu 2025: 7). The expression \emph{caturthyāṃ} is present in the vulgate only. The Chinese sources all have "fourth."}\pend

\pstart\versepar \ledsidenote{\passref{1.29.19}}\edtext{\edtext{\skt{नवम्यामथ}}{\lemma{\skt{नवम्याम्}}\Afootnote{K N H; \skt{चतुर्थ्यां वा नवम्यां}\ A}}}{\lemma{\skt{अथ}}\Afootnote{K N H; \emph{om.}\ A}}\skt{ }\edtext{\skt{षष्ठ्याम्वा}}{\lemma{\skt{षष्ठ्याम्वा}}\Afootnote{K N H; \skt{वा}\ A}}\skt{ }\edtext{\skt{पक्ष\-स\-न्धि\-दि\-नेषु वा}}{\lemma{\skt{पक्ष\appbrk{}स\appbrk{}न्धि\appbrk{}दि\appbrk{}नेषु वा}}\Afootnote{K N H; \skt{षष्ठ्यां सन्धिदिनेषु च}\ A}}\skt{~। }\pend
\pstart\versepar \skt{वैद्यं य उपसर्पन्ति दूतास्ते चापि गर्हिताः~।}\pend

\pstart\versepar \ledsidenote{\passref{1.29.20}}\skt{स्विन्नाभितप्ता मध्याह्ने }\edtext{\skt{ज्वलनस्य}}{\lemma{\skt{ज्वलनस्य}}\Afootnote{A K H; \skt{ज्वलन}\ N}}\skt{ समीपतः~। }\pend
\pstart\versepar \skt{गर्हिताः पित्तरोगेषु दूता }\edtext{\skt{वैद्यम्समागताः}}{\lemma{\skt{समागताः}}\Afootnote{K N H; \skt{उपागताः}\ A}}\skt{~॥}\pend

\pstart\versepar \ledsidenote{\passref{1.29.21}}\skt{त एव कफरोगेषु }\edtext{\skt{कर्मसिद्धिकराः}}{\lemma{\skt{कर्मसिद्धिकराः}}\Afootnote{A K H; \skt{कर्मसिद्धिकरा}\ N}}\skt{ स्मृताः~। }\pend
\pstart\versepar \skt{एतेन }\edtext{\skt{शेषं}}{\lemma{\skt{शेषं}}\Afootnote{A N H; \skt{शेष}\ K}}\skt{ व्याख्यातं }\edtext{\skt{स्वबुध्या}}{\lemma{\skt{स्वबुध्या}}\Afootnote{K N H; \skt{बुद्ध्वा}\ A}}\skt{ }\edtext{\skt{विभजेत्भिषक्}}{\lemma{\skt{विभजेत्भिषक्}}\Afootnote{K N H; \skt{संविभजेत्तु तत्}\ A}}\skt{~॥}\pend

\pstart\versepar \ledsidenote{\passref{1.29.22}}\skt{रक्त\-पि\-त्ता\-ति\-सा\-रेषु प्रमेहेषु तथैव च~। }\pend
\pstart\versepar \skt{}\edtext{\skt{प्रशस्तो}}{\lemma{\skt{प्रशस्तो}}\Afootnote{K N\textsuperscript{pc} H; \skt{प्रशस्तो जलरोधेषु दूत}\ A; \skt{आर्द्द्रा\appbrk{}श्ले\appbrk{}षा\appbrk{}म\appbrk{}घा\appbrk{}मू\appbrk{}ला\appbrk{}पू\appbrk{}र्व्वासु भरणीषु वा~।प्रशस्तो}\ N\textsuperscript{ac}}}\skt{ }\edtext{\skt{जलसंरोधो}}{\lemma{\skt{जलसंरोधो}}\Afootnote{K N H; \emph{om.}\ A}}\skt{ }\edtext{\skt{वैद्या\-तु\-र\-स\-मा\-गमः}}{\lemma{\skt{वैद्या\appbrk{}तु\appbrk{}र\appbrk{}स\appbrk{}मा\appbrk{}गमः}}\Afootnote{K N H; \skt{वैद्यसमागमः}\ A}}\skt{~॥}\pend

\pstart\versepar \ledsidenote{\passref{1.29.23}}\edtext{\skt{विज्ञायैवं}}{\lemma{\skt{विज्ञायैवं}}\Afootnote{A; \skt{एवम्}\ K N H}}\skt{ विभागं }\edtext{\skt{तु}}{\lemma{\skt{तु}}\Afootnote{A; \skt{विज्ञाय}\ K N H}}\skt{ शेषं बुध्येत पण्डितः~। }\pend
\pstart\versepar \skt{}\edtext{\skt{शुक्लवासाः}}{\lemma{\skt{शुक्लवासाः}}\Afootnote{A K\textsuperscript{pc} N H; \skt{शुक्लवासा}\ K\textsuperscript{ac}}}\skt{ }\edtext{\skt{शुचिर्गौरः}}{\lemma{\skt{शुचिर्गौरः}}\Afootnote{A; \skt{शुचिगौरः}\ K N H}}\skt{ }\edtext{\skt{श्यामो वा}}{\lemma{\skt{श्यामो वा}}\Afootnote{A K N H\textsuperscript{pc}; \skt{स्याश्यामो वा}\ H\textsuperscript{ac}}}\skt{ प्रियदर्शनः~॥}\pend

\pstart\versepar \ledsidenote{\passref{1.29.24}}\skt{स्वस्याम्जातौ स्वगोत्रो वा }\edtext{\skt{दूतः}}{\lemma{\skt{दूतः}}\Afootnote{A K N\textsuperscript{pc} H; \skt{दूत}\ N\textsuperscript{ac}}}\skt{ }\edtext{\skt{कार्यकरः}}{\lemma{\skt{कार्यकरः}}\Afootnote{A K; \skt{कार्येः कर}\ N\textsuperscript{ac}; \skt{कार्येः करः}\ N\textsuperscript{pc}; \skt{कार्जकरः}\ H}}\skt{ स्मृतः~। }\pend
\pstart\versepar \skt{}\edtext{\skt{गोया\-ने\-ना\-ग\-त\-स्तुष्टः}}{\lemma{\skt{गोया\appbrk{}ने\appbrk{}ना\appbrk{}ग\appbrk{}तस्}}\Afootnote{A K N H\textsuperscript{pc}; \skt{गोया\appbrk{}ने\appbrk{}ग\appbrk{}ना\appbrk{}ग\appbrk{}तस्}\ H\textsuperscript{ac}}}\skt{ }\edtext{\skt{पादाभ्यां}}{\lemma{\skt{पादाभ्यां}}\Afootnote{A; \skt{पद्भ्याम्वा}\ K N H}}\skt{ }\edtext{\skt{शुभचेष्टितः}}{\lemma{\skt{शुभचेष्टितः}}\Afootnote{A; \skt{क्लिष्टचेष्टितः}\ K N H}}\skt{~॥}\pend

\pstart\versepar \ledsidenote{\passref{1.29.25}}\edtext{\skt{स्मृति\-मा\-न्वि\-धि\-का\-लज्ञः}}{\lemma{\skt{स्मृति\appbrk{}मा\appbrk{}न्वि\appbrk{}धि\appbrk{}का\appbrk{}लज्ञः}}\Afootnote{A; \skt{धृतिमांश्च विधिज्ञश्च कालज्ञः}\ K N H}}\skt{ स्वतन्त्रः प्रतिपत्तिमान्~। }\pend
\pstart\versepar \skt{अलङ्कृतो }\edtext{\skt{मङ्गलवान्दूतः}}{\lemma{\skt{मङ्गलवान्दूतः}}\Afootnote{A; \skt{मङ्गलवाग्दूतः}\ K N H}}\skt{ कार्यकरः स्मृतः~॥}\edtext{}{\lemma{\skt{स्वतन्त्रः प्रतिपत्तिमान्~। अलङ्कृतो}}\Afootnote{A; \skt{प्रतिपत्तिमां अलङ्कृतो}\ K N; \skt{प्रतिपत्तिमान्~॥ अलङ्कृतो}\ H}}\pend

\pstart\versepar \ledsidenote{\passref{1.29.26}}\skt{स्वस्थं }\edtext{\edtext{\skt{प्राङ्मुखमासीनं}}{\lemma{\skt{प्राङ्मुखम्}}\Afootnote{A N H; \skt{प्राङ्मुख}\ K}}}{\lemma{\skt{आसीनं}}\Afootnote{A N H; \skt{मासीनं}\ K}}\skt{ समे देशे शुचौ शुचिम्~। }\pend
\pstart\versepar \skt{}\edtext{\skt{उपसर्पति यो}}{\lemma{\skt{उपसर्पति यो}}\Afootnote{A; \skt{उपगच्छन्ति ये}\ K N H}}\skt{ }\edtext{\skt{वैद्यं स च}}{\lemma{\skt{वैद्यं स च}}\Afootnote{A; \skt{वैद्यं ते च}\ K N H}}\skt{ }\edtext{\skt{कार्यकरः}}{\lemma{\skt{कार्यकरः}}\Afootnote{A; \skt{कार्यकराः}\ K N H}}\skt{ }\edtext{\skt{स्मृतः}}{\lemma{\skt{स्मृतः}}\Afootnote{A; \skt{स्मृताः}\ K N H}}\skt{~॥}\pend

\pstart\versepar \ledsidenote{\passref{1.29.27}}\edtext{\skt{मांसो\-द\-कु\-म्भा\-त\-प\-त्र\-वि\-प्र\-वा\-र\-ण\-गो\-वृषाः}}{\lemma{\skt{मांसो\appbrk{}द\appbrk{}कु\appbrk{}म्भा\appbrk{}त\appbrk{}प\appbrk{}त्र\appbrk{}वि\appbrk{}प्र\appbrk{}वा\appbrk{}र\appbrk{}ण\appbrk{}गो\appbrk{}वृषाः}}\Afootnote{A; \skt{मांस\appbrk{}कु\appbrk{}म्भो\appbrk{}द\appbrk{}क\appbrk{}च्छ\appbrk{}त्त्र\appbrk{}वि\appbrk{}प्र\appbrk{}वा\appbrk{}र\appbrk{}ण\appbrk{}गो\appbrk{}वृषाः}\ K N H\textsuperscript{pc}; \skt{मांस\appbrk{}कु\appbrk{}म्भो\appbrk{}द\appbrk{}क\appbrk{}च्छ\appbrk{}त्त्र\appbrk{}वि\appbrk{}वा\appbrk{}र\appbrk{}ण\appbrk{}गो\appbrk{}वृषाः}\ H\textsuperscript{ac}}}\skt{~। }\pend
\pstart\versepar \skt{शुक्लवर्णाश्च पूज्यन्ते प्रस्थाने }\edtext{\skt{दर्शनं गताः}}{\lemma{\skt{दर्शनं गताः}}\Afootnote{A; \skt{दर्शनागताः}\ K N H}}\skt{~॥}\pend

\pstart\versepar \ledsidenote{\passref{1.29.28}}\edtext{\skt{स्त्री पुत्रिणी}}{\lemma{\skt{स्त्री पुत्रिणी}}\Afootnote{A N H; \skt{स्त्री पुत्रिणीं}\ K}}\skt{ }\edtext{\skt{सवत्सा }\edtext{\skt{गौर्व\-र्ध\-मा\-न\-म\-ल\-ङ्कृता}}{\lemma{\skt{अलङ्कृता}}\Afootnote{A H; \skt{अलङ्कृताः}\ K; \skt{अलकृताः}\ N}}}{\lemma{\skt{सवत्सा गौर्}}\Afootnote{A K\textsuperscript{pc} N H; \skt{सत्सा गौर्}\ K\textsuperscript{ac}}}\skt{~। }\pend
\pstart\versepar \skt{कन्या मत्स्याः }\edtext{\skt{फलं}}{\lemma{\skt{फलं}}\Afootnote{A K\textsuperscript{pc} N H; \skt{फल}\ K\textsuperscript{ac}}}\skt{ }\edtext{\skt{चामं}}{\lemma{\skt{चामं}}\Afootnote{A K\textsuperscript{pc} N H; \skt{ञ्मं}\ K\textsuperscript{ac}}}\skt{ }\edtext{\skt{स्वस्तिकं}}{\lemma{\skt{स्वस्तिकं}}\Afootnote{A; \skt{स्वस्तीका}\ K N H}}\skt{ मोदका दधि~॥}\pend

\pstart\versepar \ledsidenote{\passref{1.29.29}}\skt{हिर\-ण्या\-क्ष\-त\-पात्रं वा रत्नानि सुमनो }\edtext{\skt{नृपः}}{\lemma{\skt{नृपः}}\Afootnote{A N H; \skt{नृप}\ K}}\skt{~। }\pend
\pstart\versepar \skt{}\edtext{\skt{अप्रशान्तोऽनलो}}{\lemma{\skt{अप्रशान्तोऽनलो}}\Afootnote{A; \skt{अप्रशान्तबलो}\ K N H}}\skt{ वाजी }\edtext{\skt{हंसश्चाषः}}{\lemma{\skt{चाषः}}\Afootnote{\skt{चापः}\ A; \skt{चाशः}\ K; \skt{चासः}\ N H}}\skt{ शिखी तथा~॥}\pend

\pstart\versepar \ledsidenote{\passref{1.29.30}}\edtext{\skt{ब्रह्म\-दु\-न्दु\-भि\-पु\-ण्याहः}}{\lemma{\skt{ब्रह्म\appbrk{}दु\appbrk{}न्दु\appbrk{}भि\appbrk{}पु\appbrk{}ण्याहः}}\Afootnote{K N H; \skt{ब्रह्म\appbrk{}दु\appbrk{}न्दु\appbrk{}भि\appbrk{}जी\appbrk{}मूत}\ A}}\skt{ शङ्ख\-वे\-णु\-र\-थ\-स्वनाः~। }\pend
\pstart\versepar \skt{}\edtext{\skt{सिंह\-गो\-वृ\-ष\-ना\-दाश्च}}{\lemma{\skt{सिंह\appbrk{}गो\appbrk{}वृ\appbrk{}ष\appbrk{}ना\appbrk{}दाश्च}}\Afootnote{\skt{सिंह\appbrk{}गो\appbrk{}वृ\appbrk{}ष\appbrk{}ना\appbrk{}दाश्च ह्रे}\ A; \skt{सिंह\appbrk{}मे\appbrk{}घ\appbrk{}नि\appbrk{}ना\appbrk{}दैश्च}\ K\textsuperscript{ac}; \skt{सिंह\appbrk{}मे\appbrk{}घ\appbrk{}नि\appbrk{}ना\appbrk{}दै\appbrk{}आश्च}\ K\textsuperscript{pc}; \skt{सिंह\appbrk{}मे\appbrk{}घ\appbrk{}नि\appbrk{}दा\appbrk{}नाश्च}\ N; \skt{सिंह\appbrk{}मे\appbrk{}घ\appbrk{}नि\appbrk{}ना\appbrk{}दाश्च}\ H}}\skt{ }\edtext{\skt{ह्रेषितं}}{\lemma{\skt{ह्रेषितं}}\Afootnote{\skt{हेषितं}\ A K H N}}\skt{ }\edtext{\skt{गजबृंहितम्}}{\lemma{\skt{गजबृंहितम्}}\Afootnote{A K H; \skt{गजबृहितं}\ N}}\skt{~॥}\pend

\pstart\versepar \ledsidenote{\passref{1.29.31}}\edtext{\skt{शस्तं}}{\lemma{\skt{शस्तं}}\Afootnote{A K\textsuperscript{pc} N H; \skt{प्रशस्तं}\ K\textsuperscript{ac}}}\skt{ }\edtext{\skt{हंसरुतं}}{\lemma{\skt{हंसरुतं}}\Afootnote{A K\textsuperscript{pc} N H; \skt{हंस}\ K\textsuperscript{ac}}}\skt{ }\edtext{\skt{नॄणां}}{\lemma{\skt{नॄणां}}\Afootnote{K\textsuperscript{pc} N H; \skt{नॄणां}\ A; \emph{om.}\ K\textsuperscript{ac}}}\skt{ कौशिकं चैव वामतः~। }\pend
\pstart\versepar \skt{प्रस्थाने यायिनः श्रेष्ठा वाचश्च }\edtext{\skt{हृदयङ्गमाः}}{\lemma{\skt{हृदयङ्गमाः}}\Afootnote{A; \skt{हृदयं प्रियाः~। आस्थिता वा न वेश्म\appbrk{}ध्व\appbrk{}ज\appbrk{}तो\appbrk{}र\appbrk{}ण\appbrk{}वे\appbrk{}दिकाः}\ K\textsuperscript{ac}; \skt{हृदयं प्रियाः आस्थिता वा न भो वेश्म\appbrk{}ध्व\appbrk{}ज\appbrk{}तो\appbrk{}र\appbrk{}ण\appbrk{}वे\appbrk{}दिकाः}\ K\textsuperscript{pc} H\textsuperscript{pc}; \skt{हृदयप्रियाः~। आस्थिता वा न वेश्म\appbrk{}ध्व\appbrk{}ज\appbrk{}तो\appbrk{}र\appbrk{}ण\appbrk{}वे\appbrk{}दिकाः}\ N; \skt{हृदयं प्रियाः~॥ आस्थिता वा न वो वेश्म\appbrk{}ध्व\appbrk{}ज\appbrk{}तो\appbrk{}र\appbrk{}ण\appbrk{}वे\appbrk{}दिकाः}\ H\textsuperscript{ac}}}\skt{~॥}\edtext{}{\lemma{\skt{कौशिकं चैव वामतः~। प्रस्थाने यायिनः श्रेष्ठा वाचश्च}}\Afootnote{A; \skt{तं नॄणाम्वाचश्च}\ K\textsuperscript{ac}; \skt{वाचश्च}\ K\textsuperscript{pc} N H}}\pend

\pstart\versepar \ledsidenote{\passref{1.29.33}}\skt{दिक्षु }\edtext{\skt{शान्तासु}}{\lemma{\skt{शान्तासु}}\Afootnote{A; \skt{सु}\ K\textsuperscript{ac}; \skt{शस्ता सु}\ K\textsuperscript{pc} N H}}\skt{ वक्तारो मधुरं }\edtext{\skt{पृष्ठतोऽनुगाः}}{\lemma{\skt{ऽनुगाः}}\Afootnote{A; \skt{नुगाः}\ K N H}}\skt{~। }\pend
\pstart\versepar \skt{वामा वा दक्षिणा वाऽपि }\edtext{\skt{शकुनाः}}{\lemma{\skt{शकुनाः}}\Afootnote{A K H; \skt{शकुना}\ N}}\skt{ कर्मसिद्धये~॥}\pend

\pstart\versepar \ledsidenote{\passref{1.29.34}}\edtext{\edtext{\skt{शुष्केऽ\-श\-नि\-ह\-तेऽ\-पत्रे}}{\lemma{\skt{ऽशनिहते}}\Afootnote{A; \skt{सनिहते}\ K N; \skt{ऽसनिहते}\ H}}}{\lemma{\skt{ऽपत्रे}}\Afootnote{A H; \skt{पत्रे}\ K N}}\skt{ वल्लीनद्धे सकण्टके~। }\pend
\pstart\versepar \skt{}\edtext{\edtext{\skt{वृक्षेऽ\-थ\-वाऽ\-श्म\-भ\-स्मा\-स्थि\-वि\-ट्तु\-षा\-ङ्गा\-र\-पां\-शुषु}}{\lemma{\skt{वृक्षे}}\Afootnote{A K H; \skt{वृक्षो}\ N}}}{\lemma{\skt{ऽथ\appbrk{}वाऽ\appbrk{}श्म\appbrk{}भ\appbrk{}स्मा\appbrk{}स्थि\appbrk{}वि\appbrk{}ट्तु\appbrk{}षा\appbrk{}ङ्गा\appbrk{}र\appbrk{}पां\appbrk{}शुषु}}\Afootnote{A; \skt{थ वाश्म\appbrk{}भ\appbrk{}स्मा\appbrk{}स्थि\appbrk{}भू\appbrk{}तु\appbrk{}षां\appbrk{}गा\appbrk{}र\appbrk{}पां\appbrk{}शुषु}\ K H; \skt{थ वाश्म\appbrk{}भ\appbrk{}स्मा\appbrk{}स्थि\appbrk{}भू\appbrk{}तु\appbrk{}षा\appbrk{}ङ्गा\appbrk{}र\appbrk{}पा\appbrk{}शुषु}\ N}}\skt{~॥}\pend

\pstart\versepar \ledsidenote{\passref{1.29.35}}\edtext{\skt{चैत्य\-व\-ल्मी\-क\-वि\-ष\-म\-स्थिता}}{\lemma{\skt{चैत्य\appbrk{}व\appbrk{}ल्मी\appbrk{}क\appbrk{}वि\appbrk{}ष\appbrk{}म\appbrk{}स्थिता}}\Afootnote{K N H; \skt{चैत्य\appbrk{}व\appbrk{}ल्मी\appbrk{}क\appbrk{}वि\appbrk{}ष\appbrk{}म\appbrk{}स्थिता}\ A}}\skt{ }\edtext{\skt{दीप्तखरस्वराः}}{\lemma{\skt{दीप्तखरस्वराः}}\Afootnote{A; \skt{दीप्तखरस्वनाः}\ K H; \skt{दीप्तखरस्वणाः}\ N}}\skt{~। }\pend
\pstart\versepar \skt{पुरतो दिक्षु दीप्तासु वक्तारो नार्थसाधकाः~॥}\pend

\pstart\versepar \ledsidenote{\passref{1.29.36}}\edtext{\skt{पुन्नामानः}}{\lemma{\skt{पुन्नामानः}}\Afootnote{A; \skt{पुंसंज्ञाः}\ K N H}}\skt{ }\edtext{\skt{खगा}}{\lemma{\skt{खगा}}\Afootnote{A; \skt{पक्षिणो}\ K N H}}\skt{ वामाः स्त्रीसंज्ञा }\edtext{\skt{दक्षिणाः}}{\lemma{\skt{दक्षिणाः}}\Afootnote{A K\textsuperscript{pc} N H; \skt{दक्षिणा}\ K\textsuperscript{ac}}}\skt{ शुभाः~। }\pend
\pstart\versepar \skt{दक्षि\-णा\-द्वा\-म\-ग\-मनं }\edtext{\skt{प्रशस्तं}}{\lemma{\skt{प्रशस्तं}}\Afootnote{A K\textsuperscript{pc} H; \skt{प्रशस्तंः}\ K\textsuperscript{ac}; \skt{प्रशस्तञ्च}\ N}}\skt{ श्वशृगालयोः~। }\pend
\pstart\versepar \skt{}\edtext{\skt{वामं नकुलचाषाणां}}{\lemma{\skt{वामं नकुलचाषाणां}}\Afootnote{A; \skt{चाश\appbrk{}कौ\appbrk{}शि\appbrk{}क\appbrk{}यो\appbrk{}रे\appbrk{}वन्}\ K; \skt{चास\appbrk{}कौ\appbrk{}शि\appbrk{}क\appbrk{}यो\appbrk{}रे\appbrk{}वन्}\ N H}}\skt{ नोभयं शशसर्पयोः~॥}\pend

\pstart \ledsidenote{\passref{1.29.37cd}}\skt{दर्शनं वा रुतं }\edtext{\skt{चापि न}}{\lemma{\skt{चापि न}}\Afootnote{\skt{वाऽपि न}\ A K N H}}\skt{ गोधा\-कृ\-क\-ला\-सयोः~॥}\pend

\pstart\versepar \ledsidenote{\passref{1.29.38}}\edtext{\skt{दूतै\-र\-नि\-ष्टै\-स्तु\-ल्या\-ना\-म\-शस्तं}}{\lemma{\skt{अनिष्टैस्}}\Afootnote{A K; \skt{अणिष्टै}\ N; \skt{अरिष्टैस्}\ H}}\skt{ दर्शनं नृणाम्~। }\pend
\pstart\versepar \skt{}\edtext{\skt{कुल\-त्थ\-ति\-ल\-का\-र्पा\-स\-तु\-ष\-पा\-षा\-ण\-भ\-स्म\-नाम्}}{\lemma{\skt{कुल\appbrk{}त्थ\appbrk{}ति\appbrk{}ल\appbrk{}का\appbrk{}र्पा\appbrk{}स\appbrk{}तु\appbrk{}ष\appbrk{}पा\appbrk{}षा\appbrk{}ण\appbrk{}भ\appbrk{}स्म\appbrk{}नाम्}}\Afootnote{A; \skt{कुल\appbrk{}त्थ\appbrk{}ति\appbrk{}ल\appbrk{}क\appbrk{}र्पा\appbrk{}स\appbrk{}तु\appbrk{}ष\appbrk{}पा\appbrk{}षा\appbrk{}ण\appbrk{}भ\appbrk{}स्म\appbrk{}नान्}\ K; \skt{कुल\appbrk{}त्थ\appbrk{}ति\appbrk{}ल\appbrk{}क\appbrk{}र्पा\appbrk{}स\appbrk{}तु\appbrk{}ष\appbrk{}पा\appbrk{}षा\appbrk{}ण\appbrk{}भ\appbrk{}स्म\appbrk{}नाम्}\ N H}}\skt{~॥}\pend

\pstart\versepar \ledsidenote{\passref{1.29.39}}\skt{पात्रं नेष्टं तथा\-ङ्गा\-र\-तै\-ल\-क\-र्द\-म\-पू\-रि\-तम्~। }\pend
\pstart\versepar \skt{}\edtext{\skt{प्रस\-न्ने\-त\-र\-म\-द्यानां}}{\lemma{\skt{प्रस\appbrk{}न्ने\appbrk{}त\appbrk{}र\appbrk{}म\appbrk{}द्यानां}}\Afootnote{A; \skt{विना सु रया मद्यानां}\ K\textsuperscript{ac} H; \skt{विना सु राम्वा मद्यानां}\ K\textsuperscript{pc} N}}\skt{ पूर्णं वा }\edtext{\skt{रक्तसर्षपैः}}{\lemma{\skt{रक्तसर्षपैः}}\Afootnote{A K H; \skt{रक्तसर्सपैः}\ N}}\skt{~॥}\pend

\pstart\versepar \ledsidenote{\passref{1.29.40}}\edtext{\skt{शव\-का\-ष्ठ\-प\-ला\-शानां}}{\lemma{\skt{शव\appbrk{}का\appbrk{}ष्ठ\appbrk{}प\appbrk{}ला\appbrk{}शानां}}\Afootnote{A N H; \skt{शव\appbrk{}का\appbrk{}ष्ठ\appbrk{}प\appbrk{}ला\appbrk{}षानां}\ K}}\skt{ }\edtext{\skt{शुष्काणां}}{\lemma{\skt{शुष्काणां}}\Afootnote{A; \skt{रूढानाम्}\ K N H}}\skt{ पथि सङ्गमाः~। }\pend
\pstart\versepar \skt{नेष्यन्ते पति\-ता\-न्त\-स्थ\-दी\-ना\-न्ध\-रि\-प\-व\-स्तथा~॥}\pend

\pstart\versepar \ledsidenote{\passref{1.29.41}}\edtext{\skt{मृदुः}}{\lemma{\skt{मृदुः}}\Afootnote{A K\textsuperscript{pc} H; \skt{मृदु}\ K\textsuperscript{ac} N}}\skt{ }\edtext{\skt{शीतोऽनुकूलश्च}}{\lemma{\skt{ऽनुकूलश्च}}\Afootnote{A H; \skt{नुकूलश्च}\ K N}}\skt{ सुगन्धिश्चानिलः शुभः~। }\pend
\pstart\versepar \skt{}\edtext{\skt{खरो\-ष्णोऽ\-नि\-ष्ट\-ग\-न्धश्च}}{\lemma{\skt{ऽनिष्टगन्धश्च}}\Afootnote{A; \skt{निष्ठगन्धश्च}\ K H\textsuperscript{ac}; \skt{निष्टगन्धश्च}\ N; \skt{नि म्ब गन्धश्च}\ H\textsuperscript{pc}}}\skt{ }\edtext{\skt{प्रतिलोमश्च}}{\lemma{\skt{प्रतिलोमश्च}}\Afootnote{A; \skt{प्रतिलोमो न}\ K N H}}\skt{ }\edtext{\skt{गर्हितः}}{\lemma{\skt{गर्हितः}}\Afootnote{A; \skt{शस्यते}\ K H; \skt{सस्यते}\ N}}\skt{~॥}\pend

\pstart\versepar \ledsidenote{\passref{1.29.42}}\edtext{\skt{ग्रन्थ्यर्बुदादिषु}}{\lemma{\skt{ग्रन्थ्यर्बुदादिषु}}\Afootnote{A K\textsuperscript{pc} N H; \emph{om.}\ K\textsuperscript{ac}}}\skt{ }\edtext{\skt{सदा}}{\lemma{\skt{सदा}}\Afootnote{A K\textsuperscript{pc} N H; \emph{om.}\ K\textsuperscript{ac}}}\skt{ }\edtext{\skt{छेदशब्दस्तु}}{\lemma{\skt{छेदशब्दस्तु}}\Afootnote{A K\textsuperscript{pc} N H; \emph{om.}\ K\textsuperscript{ac}}}\skt{ }\edtext{\skt{पूजितः}}{\lemma{\skt{पूजितः}}\Afootnote{A K\textsuperscript{pc} N H; \emph{om.}\ K\textsuperscript{ac}}}\skt{~। }\pend
\pstart\versepar \skt{विद्र\-ध्यु\-द\-र\-गु\-ल्मेषु }\edtext{\skt{भेद\-श\-ब्द\-स्त\-थैव}}{\lemma{\skt{भेदशब्दस्}}\Afootnote{A K H; \skt{भेदशब्द}\ N}}\skt{ च~॥}\pend

\pstart\versepar \ledsidenote{\passref{1.29.43}}\skt{रक्त\-पि\-त्ता\-ति\-सा\-रेषु }\edtext{\skt{रुद्धशब्दः}}{\lemma{\skt{रुद्धशब्दः}}\Afootnote{A K H; \skt{रुद्धशब्द}\ N}}\skt{ प्रशस्यते~। }\pend
\pstart\versepar \skt{}\edtext{\skt{एवं व्याधिविशेषेण निमि\-त्त\-मु\-प\-धा\-र\-येत्}}{\lemma{\skt{एवं व्याधिविशेषेण निमि\appbrk{}त्त\appbrk{}मु\appbrk{}प\appbrk{}धा\appbrk{}र\appbrk{}येत्}}\Afootnote{A; \skt{दौर्मनस्यञ्च वैद्यस्य यात्रायां नैव पूजितं}\ K H; \skt{दौर्मन्यस्यञ्च वैद्यस्य यात्रायां नैव पूजितं}\ N}}\skt{~॥}\pend

\pstart\versepar \ledsidenote{\passref{1.29.45}}\skt{प्रतिषिद्धं तथा भग्नं क्षुतं स्खलि\-त\-मा\-ह\-तम्~। }\pend
\pstart\versepar \skt{दौर्मनस्यं च वैद्यस्य यात्रायां न प्रशस्यते~॥}\pend

\pstart\versepar \ledsidenote{\passref{1.29.46}}\edtext{\edtext{\edtext{\skt{प्रवे\-शेऽ\-प्ये\-त\-दु\-द्दे\-शा\-द\-वेक्ष्यं}}{\lemma{\skt{प्रवेशे}}\Afootnote{A K\textsuperscript{pc} N H; \skt{प्रवेश्य्\_\_ए}\ K\textsuperscript{ac}}}}{\lemma{\skt{ऽप्य्}}\Afootnote{A H; \skt{प्य्}\ K N}}\skt{ च}}{\lemma{\skt{एत\appbrk{}दु\appbrk{}द्दे\appbrk{}शा\appbrk{}द\appbrk{}वेक्ष्यं च}}\Afootnote{A; \skt{एव\appbrk{}मे\appbrk{}त\appbrk{}त्स्या\appbrk{}द\appbrk{}वेक्ष्य च}\ K N H}}\skt{ }\edtext{\skt{तथा+आतुरे}}{\lemma{\skt{तथा+आतुरे}}\Afootnote{\skt{तथाऽऽतुरे}\ A; \skt{यथातुरं}\ K N H}}\skt{~। }\pend
\pstart\versepar \skt{}\edtext{\skt{प्रतिद्वारं}}{\lemma{\skt{प्रतिद्वारं}}\Afootnote{A; \skt{प्रतिद्वारे}\ K N H}}\skt{ गृहे }\edtext{\skt{वाऽस्य}}{\lemma{\skt{वाऽस्य}}\Afootnote{\skt{चास्य}\ A K N H}}\skt{ }\edtext{\skt{पुनरेतन्न}}{\lemma{\skt{पुनरेतन्न}}\Afootnote{A; \skt{इदम्भूयो न}\ K H N}}\skt{ }\edtext{\skt{गण्यते}}{\lemma{\skt{गण्यते}}\Afootnote{A; \skt{पूज्यते}\ K N H}}\skt{~॥}\pend

\pstart\versepar \ledsidenote{\passref{1.29.49}}\skt{भाण्डानां संकरस्थानां }\edtext{\skt{स्थानात्संचरणं}}{\lemma{\skt{स्थानात्संचरणं}}\Afootnote{\skt{स्थानात्सञ्चारणं}\ A; \skt{स्थान\appbrk{}स\appbrk{}ञ्च\appbrk{}र\appbrk{}णन्}\ K N H}}\skt{ तथा~। }\pend
\pstart\versepar \skt{निखातोत्पाटनं भङ्गः }\edtext{\skt{पतनं}}{\lemma{\skt{पतनं}}\Afootnote{A; \skt{शस्त्राणान्}\ K H; \skt{सस्त्राणा}\ N}}\skt{ निर्गमस्तथा~॥}\pend

\pstart\versepar \ledsidenote{\passref{1.29.50}}\skt{वैद्या\-स\-ना\-व\-सादो वा }\edtext{\skt{रोगी}}{\lemma{\skt{रोगी}}\Afootnote{A; \skt{आतुरो}\ K N H}}\skt{ }\edtext{\skt{वा स्यादधोमुखः}}{\lemma{\skt{वा स्यादधोमुखः}}\Afootnote{A; \skt{वाप्यधोमुखः}\ K N H}}\skt{~। }\pend
\pstart\versepar \skt{वैद्यं }\edtext{\skt{संभाषमाणोऽङ्गं}}{\lemma{\skt{संभाषमाणोऽङ्गं}}\Afootnote{A; \skt{संभाषते ङ्गानि}\ K N; \skt{सम्भाषतेऽङ्गानि}\ H}}\skt{ }\edtext{\skt{कुड्यमा स्तरणानि}}{\lemma{\skt{कुड्यमा स्तरणानि}}\Afootnote{A; \skt{कुड्य\appbrk{}मा\appbrk{}स्त\appbrk{}र\appbrk{}णानि}\ K H; \skt{कूड्य\appbrk{}मा\appbrk{}स्त\appbrk{}र\appbrk{}णानि}\ N}}\skt{ वा~॥}\pend

\pstart\versepar \ledsidenote{\passref{1.29.51}}\edtext{\skt{अनुसृज्य}}{\lemma{\skt{अनुसृज्य}}\Afootnote{K N H; \skt{प्रमृज्याद्वा}\ A}}\skt{ }\edtext{\skt{धुनुयाद्वा}}{\lemma{\skt{धुनुयाद्वा}}\Afootnote{K N H; \skt{धुनीयाद्वा}\ A}}\skt{ }\edtext{\skt{करौ पृष्ठं शिरस्तथा}}{\lemma{\skt{करौ पृष्ठं शिरस्}}\Afootnote{A; \skt{प्रगृ\appbrk{}ही\appbrk{}त\appbrk{}शि\appbrk{}रस्}\ K N H}}\skt{~। }\pend
\pstart\versepar \skt{हस्तं }\edtext{\skt{चाकृष्य}}{\lemma{\skt{चाकृष्य}}\Afootnote{A; \skt{वाकृष्य}\ K N H}}\skt{ वैद्यस्य न्यसेच्छिरसि चोरसि~॥}\pend

\pstart \ledsidenote{\passref{1.29.52cd}}\skt{न स सिध्यति }\edtext{\skt{वैद्यो वा}}{\lemma{\skt{वैद्यो वा}}\Afootnote{A; \skt{वैद्यो पि}\ K N H}}\skt{ गृहे यस्य न पूज्यते~॥}\pend

\pstart \ledsidenote{\passref{1.29.52ab}}\edtext{\skt{यो }\edtext{\skt{वैद्यमुन्मुखः}}{\lemma{\skt{उन्मुखः}}\Afootnote{A; \skt{मुहुर्मुहुः}\ K; \skt{मुहुर्मुहुः}\ N H}}}{\lemma{\skt{यो वैद्यम्}}\Afootnote{A; \skt{वैद्यम्}\ K N H}}\skt{ }\edtext{\skt{पृच्छेदुन्मार्ष्टि}}{\lemma{\skt{पृच्छेदुन्मार्ष्टि}}\Afootnote{\skt{पश्यन्नुन्मार्ष्टि}\ A; \skt{पृच्छन्मार्ष्टि वा}\ K H; \skt{पृच्छन्माष्टिर्वा}\ N}}\skt{ }\edtext{\skt{स्वाङ्गमातुरः}}{\lemma{\skt{स्वाङ्गम्}}\Afootnote{A K H; \skt{चाङ्गम्}\ N}}\skt{~।}\pend

\pstart\versepar \ledsidenote{\passref{1.29.53}}\edtext{\skt{भवने पूज्यते}}{\lemma{\skt{भवने पूज्यते}}\Afootnote{A; \skt{भूयः सम्पूज्यते}\ K H; \skt{भूय संपूज्यते}\ N}}\skt{ }\edtext{\skt{वाऽपि}}{\lemma{\skt{वाऽपि}}\Afootnote{A; \skt{यस्य}\ K H; \emph{om.}\ N}}\skt{ }\edtext{\skt{यस्य}}{\lemma{\skt{यस्य}}\Afootnote{A; \skt{गृहे}\ K N H}}\skt{ वैद्यः स सिध्यति~। }\pend
\pstart\versepar \skt{शुभं शुभेषु }\edtext{\skt{दूतादिष्वशुभं}}{\lemma{\skt{दूतादिष्व्}}\Afootnote{A K H; \skt{दूतानिष्व्}\ N}}\skt{ ह्यशुभेषु च~॥}\pend

\pstart\versepar \ledsidenote{\passref{1.29.54}}\skt{आतुरस्य ध्रुवं तस्मा\-द्दू\-ता\-दी\-न्ल\-क्ष\-ये\-द्भि\-षक्~। }\pend
\pstart\versepar \skt{स्वप्नानतः प्रवक्ष्यामि मरणाय शुभाय च~॥}\edtext{}{\lemma{\skt{तस्मा\appbrk{}द्दू\appbrk{}ता\appbrk{}दी\appbrk{}न्ल\appbrk{}क्ष\appbrk{}ये\appbrk{}द्भि\appbrk{}षक्~। स्वप्नानतः}}\Afootnote{A; \skt{तस्माद्दूतादीं सम्परीक्षयेत्~॥ स्वप्नानतः}\ K; \skt{तस्मा\appbrk{}द्दू\appbrk{}ता\appbrk{}दी\appbrk{}न्स\appbrk{}म्प\appbrk{}री\appbrk{}क्ष\appbrk{}येत्~॥ स्वप्नानतः}\ N H}}\pend

\pstart\versepar \ledsidenote{\passref{1.29.55}}\edtext{\skt{सुहृदो यां श्च}}{\lemma{\skt{सुहृदो यां श्च}}\Afootnote{A; \skt{पश्यन्ति सुहृदो यांस्}\ K H; \skt{पश्यन्तु सुहृदो यांस्}\ N}}\skt{ }\edtext{\skt{पश्यन्ति}}{\lemma{\skt{पश्यन्ति}}\Afootnote{A; \skt{तु}\ K N H}}\skt{ }\edtext{\skt{व्याधितो वा}}{\lemma{\skt{व्याधितो वा}}\Afootnote{A; \skt{स्वप्नां}\ K; \skt{स्वप्ना}\ N H}}\skt{ }\edtext{\skt{स्वयं}}{\lemma{\skt{स्वयं}}\Afootnote{A K; \skt{न्स्वयम्}\ N H}}\skt{ तथा~। }\pend
\pstart\versepar \skt{स्नेहा\-भ्य\-क्त\-श\-री\-रस्तु }\edtext{\skt{कर\-भ\-व्या\-ल\-ग\-र्दभैः}}{\lemma{\skt{कर\appbrk{}भ\appbrk{}व्या\appbrk{}ल\appbrk{}ग\appbrk{}र्दभैः}}\Afootnote{A; \skt{कर\appbrk{}भ\appbrk{}व्या\appbrk{}ड\appbrk{}ग\appbrk{}र्दभैः}\ K N H}}\skt{~॥}\edtext{}{\lemma{\skt{तथा~। स्नेहा\appbrk{}भ्य\appbrk{}क्त\appbrk{}श\appbrk{}री\appbrk{}रस्तु}}\Afootnote{A; \skt{अथापि वा स्नेहा\appbrk{}भ्य\appbrk{}क्त\appbrk{}श\appbrk{}री\appbrk{}रस्य}\ K N H}}\pend

\pstart\versepar \ledsidenote{\passref{1.29.56}}\edtext{\skt{वरा\-है\-र्म\-हि\-षै\-र्वाऽपि यो}}{\lemma{\skt{वाऽपि यो}}\Afootnote{A K H; \skt{वापि}\ N}}\skt{ }\edtext{\skt{याया\-द्द\-क्षि\-णा\-मुखः}}{\lemma{\skt{याया\appbrk{}द्द\appbrk{}क्षि\appbrk{}णा\appbrk{}मुखः}}\Afootnote{\skt{याया\appbrk{}द्द\appbrk{}क्षि\appbrk{}णा\appbrk{}मुखः}\ A; \skt{याया\appbrk{}द्द\appbrk{}क्षि\appbrk{}णा\appbrk{}मुखं}\ K N H}}\skt{~। }\pend
\pstart\versepar \skt{}\edtext{\skt{रक्ताम्बरधरा}}{\lemma{\skt{रक्ताम्बरधरा}}\Afootnote{A; \skt{कृष्णा रक्ताम्बरधरा}\ K N H}}\skt{ }\edtext{\skt{कृष्णा हसन्ती}}{\lemma{\skt{कृष्णा हसन्ती}}\Afootnote{A; \skt{हसन्ती}\ K H N}}\skt{ }\edtext{\skt{मुक्तमूर्धजा}}{\lemma{\skt{मुक्तमूर्धजा}}\Afootnote{A K\textsuperscript{pc}; \skt{मुक्तमूर्धजाः}\ K\textsuperscript{ac} N H}}\skt{~॥}\pend

\pstart\versepar \ledsidenote{\passref{1.29.57}}\edtext{\skt{यं वा}}{\lemma{\skt{यं वा}}\Afootnote{K N H; \skt{यं चा}\ A}}\skt{ कर्षति बद्ध्वा स्त्री }\edtext{\skt{नृत्यन्ती}}{\lemma{\skt{नृत्यन्ती}}\Afootnote{A; \skt{हसन्तन्}\ K N H}}\skt{ दक्षिणामुखम्~। }\pend
\pstart\versepar \skt{}\edtext{\skt{अन्ता\-व\-सा\-यि\-भिर्यो}}{\lemma{\skt{अन्ता\appbrk{}व\appbrk{}सा\appbrk{}यि\appbrk{}भिर्यो}}\Afootnote{A; \skt{अन्ते निवासिभिर्यो}\ K N H}}\skt{ }\edtext{\skt{वा+आकृष्यते दक्षिणामुखः}}{\lemma{\skt{वा+आकृष्यते दक्षिणामुखः}}\Afootnote{\skt{वाऽकृष्यते दक्षिणामुखः}\ A; \skt{वा प्याकृ\appbrk{}ष्ये\appbrk{}द्द\appbrk{}क्षि\appbrk{}णा\appbrk{}मुखं}\ K H N}}\skt{~॥}\pend

\pstart\versepar \ledsidenote{\passref{1.29.58}}\edtext{\skt{परिष्वजेरन्यं}}{\lemma{\skt{परिष्वजेरन्यं}}\Afootnote{A; \skt{परिष्वजेयुर्}\ K\textsuperscript{ac}; \skt{परिष्वजेयुर्यं}\ K\textsuperscript{pc} H; \skt{परिष्वज्येये युञ्}\ N}}\skt{ }\edtext{\skt{वाऽपि}}{\lemma{\skt{वाऽपि}}\Afootnote{A; \skt{पि}\ K\textsuperscript{ac}; \skt{चा पि}\ K\textsuperscript{pc} N H}}\skt{ }\edtext{\skt{प्रेताः}}{\lemma{\skt{प्रेताः}}\Afootnote{A K\textsuperscript{pc} N H; \skt{प्रेता}\ K\textsuperscript{ac}}}\skt{ }\edtext{\edtext{\skt{प्रव्रजितास्तथा}}{\lemma{\skt{प्रव्रजितास्}}\Afootnote{K\textsuperscript{pc} N H; \skt{प्रव्रजिता(न)स्}\ A; \skt{प्रव्रजिता न्}\ K\textsuperscript{ac}}}}{\lemma{\skt{तथा}}\Afootnote{A K\textsuperscript{pc} N H; \skt{त}\ K\textsuperscript{ac}}}\skt{~। }\pend
\pstart\versepar \skt{}\edtext{\edtext{\skt{मुहुराघ्रायते}}{\lemma{\skt{मुहुर्}}\Afootnote{A; \skt{था}\ K\textsuperscript{ac}; \emph{om.}\ K\textsuperscript{pc} N H}}\skt{ }\edtext{\skt{यस्तु}}{\lemma{\skt{तु}}\Afootnote{A; \skt{मुहुः}\ K N H}}}{\lemma{\skt{आघ्रायते यस्}}\Afootnote{A; \skt{आघ्राय्यते यश्च}\ K H; \skt{आघ्रायते यश्च}\ N}}\skt{ श्वाप\-दै\-र्वि\-कृ\-ता\-ननैः~॥}\pend

\pstart\versepar \ledsidenote{\passref{1.29.59}}\skt{पिबेन्मधु च तैलं च यो वा }\edtext{\skt{पङ्केऽवसीदति}}{\lemma{\skt{ऽवसीदति}}\Afootnote{A; \skt{वसीदति}\ K N H}}\skt{~। }\pend
\pstart\versepar \skt{पङ्क\-प्र\-दि\-ग्ध\-गात्रो वा }\edtext{\skt{प्रनृ\-त्ये\-त्प्र\-ह\-से\-त्तथा}}{\lemma{\skt{प्रनृ\appbrk{}त्ये\appbrk{}त्प्र\appbrk{}ह\appbrk{}से\appbrk{}त्तथा}}\Afootnote{A; \skt{नृत्येद्वाथ हसेत वा}\ K H; \skt{नृत्येद्वाथ हसेन वा}\ N}}\skt{~॥}\pend

\pstart\versepar \ledsidenote{\passref{1.29.60}}\skt{निरम्बरश्च यो }\edtext{\skt{रक्तां}}{\lemma{\skt{रक्तां}}\Afootnote{A K N; \skt{रक्तं}\ H}}\skt{ }\edtext{\skt{धारयेच्छिरसि}}{\lemma{\skt{धारयेच्छिरसि}}\Afootnote{\skt{धारयेच्छिरसा}\ A K N H}}\skt{ स्रजम्~। }\pend
\pstart\versepar \skt{यस्य वंशो नलो वाऽपि }\edtext{\skt{तालो}}{\lemma{\skt{तालो}}\Afootnote{A; \skt{वृक्षो}\ K N H}}\skt{ वोरसि जायते~॥}\pend

\pstart \ledsidenote{\passref{1.29.60.1}}\skt{मस्तकाद्यश्च तालो वा उच्छ्रिता }\edtext{\skt{वीणुवीरुधः}}{\lemma{\skt{वीणुवीरुधः}}\Afootnote{N; \skt{वेणुवीरुधः}\ K H}}\skt{~॥}\pend

\pstart\versepar \ledsidenote{\passref{1.29.61}}\skt{यं वा मत्स्यो ग्रसेद्यो वा }\edtext{\skt{ज्वलनं}}{\lemma{\skt{ज्वलनं}}\Afootnote{K H; \skt{जननीं}\ A; \skt{ज्वल}\ N}}\skt{ }\edtext{\skt{प्रविशेन्नरः}}{\lemma{\skt{प्रविशेन्}}\Afootnote{A H; \skt{प्रविसे}\ K N}}\skt{~। }\pend
\pstart\versepar \skt{}\edtext{\skt{पर्व\-ता\-ग्रा\-त्प\-तेद्यो वा}}{\lemma{\skt{पर्व\appbrk{}ता\appbrk{}ग्रा\appbrk{}त्प\appbrk{}तेद्यो वा}}\Afootnote{A H; \skt{पर्व\appbrk{}ता\appbrk{}ण्ग्रा\appbrk{}त्प\appbrk{}तेद्यो वा}\ K N}}\skt{ श्वभ्रे वा तमसावृते~॥}\pend

\pstart\versepar \ledsidenote{\passref{1.29.62}}\edtext{\skt{ह्रियते}}{\lemma{\skt{ह्रियते}}\Afootnote{K N H; \skt{ह्रियेत}\ A}}\skt{ }\edtext{\skt{स्रोतसा यो वा यो वा}}{\lemma{\skt{स्रोतसा यो वा यो वा}}\Afootnote{A K\textsuperscript{pc} N H; \skt{स्रोतसा यो वा}\ K\textsuperscript{ac}}}\skt{ }\edtext{\skt{मौण्ड्य\-म\-वा\-प्नु\-यात्}}{\lemma{\skt{मौण्ड्यम्}}\Afootnote{K H; \skt{मौण्ड्य}\ A; \skt{मोण्ड्यम्}\ N}}\skt{~। }\pend
\pstart\versepar \skt{पराजीयेत }\edtext{\skt{बध्येत}}{\lemma{\skt{बध्येत}}\Afootnote{A; \skt{युद्धैर्वा}\ K N H}}\skt{ }\edtext{\skt{काका\-द्यै\-र्वाऽ\-भि\-भू\-यते}}{\lemma{\skt{काकाद्यैर्}}\Afootnote{A N H; \skt{…काद्यैर्}\ K}}\skt{~॥}\edtext{}{\lemma{\skt{अवाप्नुयात्~। पराजीयेत}}\Afootnote{K N; \skt{मवाप्नुयात्~। पराजीयेत}\ A; \skt{अवाप्नुयात्~। पराजयेत}\ H}}\pend

\pstart\versepar \ledsidenote{\passref{1.29.63}}\skt{पतनं तारकादीनां }\edtext{\skt{प्रणाशं}}{\lemma{\skt{प्रणाशं}}\Afootnote{A; \skt{प्रणाशो}\ K H; \skt{प्रणासो}\ N}}\skt{ दीपचक्षुषोः~। }\pend
\pstart\versepar \skt{यः }\edtext{\skt{पश्येद्देवतानां च}}{\lemma{\skt{पश्येद्देवतानां च}}\Afootnote{\skt{पश्येद्देवतानां}\ A K N H}}\skt{ }\edtext{\skt{(आ.वा)}}{\lemma{\skt{(आ.वा)}}\Afootnote{\skt{वा}\ A K N H}}\skt{ }\edtext{\skt{प्रक\-म्प\-म\-व\-ने\-स्तथा}}{\lemma{\skt{अवनेस्}}\Afootnote{A; \skt{पतनन्}\ K N H}}\skt{~॥}\pend

\pstart\versepar \ledsidenote{\passref{1.29.64}}\skt{यस्य छर्दिर्विरेको वा }\edtext{\skt{दशनाः}}{\lemma{\skt{दशनाः}}\Afootnote{A; \skt{दशना}\ K; \skt{दसना}\ N; \skt{दसनाः}\ H}}\skt{ }\edtext{\skt{प्रपतन्ति वा}}{\lemma{\skt{प्रपतन्ति वा}}\Afootnote{A K H; \skt{त्प्रपतन्ति वा}\ N}}\skt{~। }\pend
\pstart\versepar \skt{}\edtext{\skt{शाल्मलीं}}{\lemma{\skt{शाल्मलीं}}\Afootnote{A; \skt{शाल्मलिं}\ K N H}}\skt{ किंशुकं }\edtext{\skt{यूपं}}{\lemma{\skt{यूपं}}\Afootnote{A K H; \skt{पूयम्}\ N}}\skt{ वल्मीकं पारिभद्रकम्~॥}\pend

\pstart\versepar \ledsidenote{\passref{1.29.65}}\edtext{\skt{पुष्पाढ्यं}}{\lemma{\skt{पुष्पाढ्यं}}\Afootnote{A K H; \skt{पुष्पाद्यं}\ N}}\skt{ कोविदारं वा }\edtext{\skt{चितां वा }\edtext{\skt{योऽधिरोहति}}{\lemma{\skt{ऽधिरोहति}}\Afootnote{A H; \skt{धिरोहति}\ K N}}}{\lemma{\skt{चितां वा यो}}\Afootnote{A K N; \skt{चितम्वा यो}\ H}}\skt{~। }\pend
\pstart\versepar \skt{}\edtext{\skt{कार्पा\-स\-तै\-ल\-पि\-ण्या\-क\-लो\-हानि}}{\lemma{\skt{कार्पा\appbrk{}स\appbrk{}तै\appbrk{}ल\appbrk{}पि\appbrk{}ण्या\appbrk{}क\appbrk{}लो\appbrk{}हानि}}\Afootnote{A; \skt{कर्पा\appbrk{}स\appbrk{}तै\appbrk{}ल\appbrk{}पि\appbrk{}ण्या\appbrk{}क\appbrk{}लो\appbrk{}हानि}\ K N H}}\skt{ लवणं तिलान्~॥}\pend

\pstart \ledsidenote{\passref{1.29.66cd}}\skt{स्वस्थः स लभते व्याधिं व्याधितो }\edtext{\edtext{\skt{मृत्युमृच्छति}}{\lemma{\skt{मृत्युम्}}\Afootnote{K H; \skt{मृत्यु}\ A N}}}{\lemma{\skt{ऋच्छति}}\Afootnote{\skt{मृच्छति}\ A; \skt{अर्च्छति}\ K H; \skt{गच्छति}\ N}}\skt{~॥}\pend

\pstart\versepar \ledsidenote{\passref{1.29.67}}\skt{यथास्वं प्रकृतिस्वप्नो विस्मृतो विहतस्तथा~। }\pend
\pstart\versepar \skt{चिन्ताकृतो दिवा }\edtext{\skt{दृष्तो}}{\lemma{\skt{दृष्तो}}\Afootnote{\skt{दृष्टो}\ A; \skt{यश्च}\ K N H}}\skt{ भव\-न्त्य\-फ\-ल\-दास्तु ते~॥}\pend

\pstart\versepar \ledsidenote{\passref{1.29.68}}\skt{ज्वरितानां शुना सख्यं कपिसख्यं तु शोषिणाम्~। }\pend
\pstart\versepar \skt{उन्मादे राक्षसैः प्रेतैरपस्मारे }\edtext{\skt{प्रवर्तनम्}}{\lemma{\skt{प्रवर्तनम्}}\Afootnote{A; \skt{तु नर्तनं}\ K N H}}\skt{~॥}\pend

\pstart\versepar \ledsidenote{\passref{1.29.69}}\skt{मेहातिसारिणां }\edtext{\skt{तोयपानं स्नेहस्य}}{\lemma{\skt{तोयपानं स्नेहस्य}}\Afootnote{A; \skt{तोय\appbrk{}स्ने\appbrk{}ह\appbrk{}पा\appbrk{}नन्तु}\ K H N}}\skt{ }\edtext{\skt{कुष्ठिनाम्}}{\lemma{\skt{कुष्ठिनाम्}}\Afootnote{A; \skt{कुष्ठिनः}\ K N H}}\skt{~। }\pend
\pstart\versepar \skt{}\edtext{\skt{गुल्मेषु}}{\lemma{\skt{गुल्मेषु}}\Afootnote{A; \skt{गुल्मे तु}\ K N H}}\skt{ स्थावरोत्पत्तिः कोष्ठे मूर्ध्नि शिरोरुजि~॥}\pend

\pstart\versepar \ledsidenote{\passref{1.29.70}}\edtext{\edtext{\skt{पिपा\-सा\-श्वा\-स\-यो\-रध्वा}}{\lemma{\skt{पिपा\appbrk{}सा\appbrk{}श्वा\appbrk{}स\appbrk{}योर्}}\Afootnote{K N H; \emph{om.}\ A}}}{\lemma{\skt{अध्वा}}\Afootnote{K N H; \emph{om.}\ A}}\skt{ }\edtext{\skt{च्छर्द्यां}}{\lemma{\skt{च्छर्द्यां}}\Afootnote{K N H; \emph{om.}\ A}}\skt{ }\edtext{\skt{शष्कुलिभक्षणं}}{\lemma{\skt{शष्कुलिभक्षणं}}\Afootnote{K N H; \skt{शष्कुलीभक्षणं}\ A}}\skt{~। }\pend
\pstart\versepar \skt{}\edtext{\skt{हारिद्रं}}{\lemma{\skt{हारिद्रं}}\Afootnote{N H; \skt{छर्द्यामध्वा श्वासपिपासयोः~। हारिद्रं}\ A; \skt{हरिद्रं}\ K}}\skt{ भोजनं वापि }\edtext{\skt{यद्भ\-वे\-त्पा\-ण्डु\-रो\-गिणः}}{\lemma{\skt{यद्भ\appbrk{}वे\appbrk{}त्पा\appbrk{}ण्डु\appbrk{}रो\appbrk{}गिणः}}\Afootnote{K N H; \skt{यस्य स्यात्पाण्डुरोगिणः}\ A}}\skt{~॥}\pend

\pstart \ledsidenote{\passref{1.29.71ab}}\skt{रक्तपित्ती }\edtext{\skt{पिबेद्यश्च}}{\lemma{\skt{पिबेद्यश्च}}\Afootnote{K H; \skt{पिबेद्यस्तु}\ A; \skt{भवेद्यश्च}\ N}}\skt{ शोणितं स }\edtext{\skt{विनश्यति}}{\lemma{\skt{विनश्यति}}\Afootnote{A K N\textsuperscript{ac} H; \skt{विनति}\ N\textsuperscript{pc}}}\skt{~।}\footnoteA{There follow ten verses not present in the vulgate.}\pend

\pstart\versepar \ledsidenote{\passref{1.29.71.add1}}\skt{यो नरो बहुभिर्मुष्णैः }\edtext{\skt{क्रुद्धैश्च}}{\lemma{\skt{क्रुद्धैश्च}}\Afootnote{N H; \skt{क्रुद्धेश्च}\ K}}\skt{ प्रतिबोध्यते~। }\pend
\pstart\versepar \skt{हन्यते बहुबिः शस्त्रैर्ब्राह्मणं यो न पश्यति~॥}\pend

\pstart\versepar \ledsidenote{\passref{1.29.71.add2}}\skt{दन्त\-प्र\-क्षा\-ल\-न\-स्नानं स्वप्ने }\edtext{\skt{केशविलेखनं}}{\lemma{\skt{केशविलेखनं}}\Afootnote{K H; \skt{क्लेशविलेखनं}\ N}}\skt{~। }\pend
\pstart\versepar \skt{}\edtext{\edtext{\skt{रोगा\-ग\-म\-नि\-मि\-त्त\-न्त\-द्ध\-न्व\-न्त\-रि\-वचो}}{\lemma{\skt{रोगा\appbrk{}ग\appbrk{}म\appbrk{}नि\appbrk{}मि\appbrk{}त्तन्}}\Afootnote{H\textsuperscript{pc}; \skt{रोगा\appbrk{}ग\appbrk{}म\appbrk{}नि\appbrk{}मि\appbrk{}त्तंन्}\ K; \skt{रागा\appbrk{}ग\appbrk{}म\appbrk{}नि\appbrk{}मि\appbrk{}त्तन्}\ N; \skt{रोगागमनिमि}\ H\textsuperscript{ac}}}}{\lemma{\skt{तद्ध\appbrk{}न्व\appbrk{}न्त\appbrk{}रि\appbrk{}वचो}}\Afootnote{K H; \skt{तद्ध\appbrk{}न्व\appbrk{}न्त\appbrk{}रि\appbrk{}र्वचो}\ N}}\skt{ यथा~॥}\pend

\pstart\versepar \ledsidenote{\passref{1.29.71.add3}}\edtext{\edtext{\skt{अप\-र्वा\-णी\-क्षु\-र्ग\-ण्डानि}}{\lemma{\skt{अपर्वाणीक्षुर्}}\Afootnote{K; \skt{अपर्वाशीक्षु}\ N H}}}{\lemma{\skt{गण्डानि}}\Afootnote{K; \skt{खण्टानि}\ N; \skt{गश्टानि}\ H\textsuperscript{ac}; \skt{गदश्टानि}\ H\textsuperscript{pc}}}\skt{ गुडम्वा योऽतिभक्षयेत्~। }\pend
\pstart\versepar \skt{प्रतिबुद्धो }\edtext{\edtext{\skt{विजा\-नी\-या\-द\-ङ्ग\-व्या\-धि\-मु\-प\-स्थितं}}{\lemma{\skt{विजा\appbrk{}नी\appbrk{}या\appbrk{}द\appbrk{}ङ्ग\appbrk{}व्या\appbrk{}धिम्}}\Afootnote{H; \skt{विजानीयादङ्ग}\ K\textsuperscript{ac}; \skt{विजा\appbrk{}नी\appbrk{}या\appbrk{}द\appbrk{}ङ्ग\appbrk{}व्याधि}\ K\textsuperscript{pc}; \skt{विजा\appbrk{}नी\appbrk{}या\appbrk{}द\appbrk{}र्ङ्ग\appbrk{}व्या\appbrk{}धिम्}\ N}}}{\lemma{\skt{उपस्थितं}}\Afootnote{N H; \skt{अपस्थितस्थितं}\ K\textsuperscript{ac}; \skt{अपस्थितं}\ K\textsuperscript{pc}}}\skt{~॥}\edtext{}{\lemma{\skt{ऽतिभक्षयेत्~। प्रतिबुद्धो}}\Afootnote{K N; \skt{विभक्षयेत्~। प्रतिबुद्धो}\ H}}\pend

\pstart\versepar \ledsidenote{\passref{1.29.71.add4}}\skt{अयुक्तं }\edtext{\edtext{\skt{यानमारुह्य}}{\lemma{\skt{यानम्}}\Afootnote{K N; \skt{यः न}\ H}}}{\lemma{\skt{आरुह्य}}\Afootnote{K N; \skt{सारुब्धा}\ H}}\skt{ शुष्कवृक्षञ्च सर्वतः~। }\pend
\pstart\versepar \skt{प्लवते }\edtext{\skt{तोयमध्ये वा}}{\lemma{\skt{तोयमध्ये वा}}\Afootnote{K H; \skt{तोयमध्ये}\ N}}\skt{ }\edtext{\skt{नासरोगं}}{\lemma{\skt{नासरोगं}}\Afootnote{N; \skt{स रोगं}\ K H}}\skt{ प्राप्नुयान्महत्॥}\pend

\pstart\versepar \ledsidenote{\passref{1.29.71.add5}}\skt{तिलकल्केन मांसेन वसया कर्दमेन वा। }\pend
\pstart\versepar \skt{गात्रं यस्योपलिप्येत स }\edtext{\edtext{\skt{रोग\-म्प्रा\-प्नु\-या\-न्म\-हत्}}{\lemma{\skt{रोगम्}}\Afootnote{K H; \skt{रोग}\ N}}}{\lemma{\skt{प्राप्नुयान्महत्}}\Afootnote{K N; \skt{प्राप्नुयात्महत्}\ H}}\skt{~॥}\pend

\pstart\versepar \ledsidenote{\passref{1.29.71.add6}}\skt{चाण्डालैर्वा श्वपाकैर्वा }\edtext{\skt{चोरैः}}{\lemma{\skt{चोरैः}}\Afootnote{K N; \skt{चौरैः}\ H}}\skt{ }\edtext{\skt{प्रव्रजितैस्तथा}}{\lemma{\skt{प्रव्रजितैस्}}\Afootnote{K N; \skt{प्रव्रजितेस्}\ H}}\skt{~। }\pend
\pstart\versepar \skt{}\edtext{\skt{अन्यैश्च}}{\lemma{\skt{अन्यैश्च}}\Afootnote{K H; \skt{अन्येश्च}\ N}}\skt{ }\edtext{\skt{नीचबीभत्सैः}}{\lemma{\skt{नीचबीभत्सैः}}\Afootnote{N H; \skt{नीचभीभत्सैः}\ K}}\skt{ प्रहारो न प्रशस्यते।}\pend

\pstart\versepar \ledsidenote{\passref{1.29.71.add7}}\skt{मद्यस्यान्नस्य मांसस्य }\edtext{\skt{तिलतैलस्य}}{\lemma{\skt{तिलतैलस्य}}\Afootnote{K N; \skt{तिलतेलस्य}\ H}}\skt{ सर्प्पिषः~। }\pend
\pstart\versepar \skt{स्वप्ने किञ्चिद्यदि प्राप्य पश्ये\-द्व्या\-धि\-मु\-प\-स्थितं~॥}\pend

\pstart\versepar \ledsidenote{\passref{1.29.71.add8}}\skt{रक्तवर्ण्णन्तु यत्कि\-ञ्चि\-त्स\-र्व्व\-न्तन्न प्रशस्यते~। }\pend
\pstart\versepar \skt{}\edtext{\skt{रक्टं}}{\lemma{\skt{रक्टं}}\Afootnote{K H; \skt{रक्ट}\ N}}\skt{ }\edtext{\skt{यदामकं}}{\lemma{\skt{यदामकं}}\Afootnote{K H; \skt{यादामकं}\ N}}\skt{ मांसञ्चन्दनञ्च प्रशस्यते~॥}\pend

\pstart\versepar \ledsidenote{\passref{1.29.71.add9}}\skt{यद्वा शुक्लेन वर्णेन सर्व\-मे\-त\-त्प्र\-श\-स्यते~। }\pend
\pstart\versepar \skt{}\edtext{\edtext{\skt{कर्पासमस्थ्यथो}}{\lemma{\skt{अस्थ्य्}}\Afootnote{K; \skt{अप्य्}\ N H}}}{\lemma{\skt{अथो}}\Afootnote{K H; \skt{अधो}\ N}}\skt{ भस्म दधि चैवात्र वर्जितं~।}\pend

\pstart\versepar \ledsidenote{\passref{1.29.71.add10}}\skt{यत्तु कृष्णेन वर्णेन सर्व्वन्तन्न प्रशस्यते~। }\pend
\pstart\versepar \skt{भूमिलाभो हिरण्यञ्च हस्ती चैव प्रशस्यते~॥}\pend

\pstart\versepar \ledsidenote{\passref{1.29.75}}\skt{अत }\edtext{\skt{ऊर्ध्वम्प्रवक्ष्यामि}}{\lemma{\skt{ऊर्ध्वम्}}\Afootnote{A K N; \skt{ऊर्द्वम्}\ H}}\skt{ }\edtext{\skt{प्रशस्तं}}{\lemma{\skt{प्रशस्तं}}\Afootnote{A K H; \skt{प्रसस्त}\ N}}\skt{ स्वप्नदर्शनं~। }\pend
\pstart\versepar \skt{}\edtext{\edtext{\skt{देवा\-न्द्वि\-जा\-न्गो\-वृ\-ष\-भा\-न्जी\-वतः}}{\lemma{\skt{देवान्}}\Afootnote{A K H; \skt{देवा}\ N}}}{\lemma{\skt{द्विजा\appbrk{}न्गो\appbrk{}वृ\appbrk{}ष\appbrk{}भा\appbrk{}न्जी\appbrk{}वतः}}\Afootnote{A H; \skt{द्विजं गोवृषभां जीवतः}\ K; \skt{द्विजं गोवृषभां जीवितः}\ N}}\skt{ सुहृदो }\edtext{\skt{नृपान्}}{\lemma{\skt{नृपान्}}\Afootnote{A H; \skt{नृपां}\ K N}}\skt{~॥}\pend

\pstart\versepar \ledsidenote{\passref{1.29.76}}\edtext{\edtext{\skt{समृद्धमग्निं}}{\lemma{\skt{समृद्धम्}}\Afootnote{K N H; \skt{समिद्धम्}\ A}}}{\lemma{\skt{अग्निं}}\Afootnote{A K H; \skt{अग्निम~।}\ N}}\skt{ साधूंश्च निर्मलानि जलानि च~। }\pend
\pstart\versepar \skt{}\edtext{\skt{पश्ये\-त्क\-ल्या\-ण\-भा\-वाय}}{\lemma{\skt{पश्ये\appbrk{}त्क\appbrk{}ल्या\appbrk{}ण\appbrk{}भा\appbrk{}वाय}}\Afootnote{K N H; \skt{पश्ये\appbrk{}त्क\appbrk{}ल्या\appbrk{}ण\appbrk{}ला\appbrk{}भाय}\ A}}\skt{ व्याधे\-र\-प\-ग\-माय च~॥}\pend

\pstart\versepar \ledsidenote{\passref{1.29.77}}\skt{मांसं }\edtext{\skt{मत्स्यं}}{\lemma{\skt{मत्स्यं}}\Afootnote{K N H; \skt{मत्स्या}\ A}}\skt{ }\edtext{\skt{स्रजः}}{\lemma{\skt{स्रजः}}\Afootnote{K N H; \skt{न्स्रजः}\ A}}\skt{ }\edtext{\skt{श्वेता}}{\lemma{\skt{श्वेता}}\Afootnote{A K N; \skt{स्वेता}\ H}}\skt{ वासांसि च फलानि च~। }\pend
\pstart\versepar \skt{}\edtext{\skt{लभते}}{\lemma{\skt{लभते}}\Afootnote{K N H; \skt{लभेत}\ A}}\skt{ धनलाभाय व्याधे\-र\-प\-ग\-माय च~॥}\pend

\pstart\versepar \ledsidenote{\passref{1.29.79}}\skt{नदी\-न\-द\-स\-मु\-द्रांश्च }\edtext{\skt{क्षुभि\-ता\-न्क\-लु\-षो\-द\-कान्}}{\lemma{\skt{क्षुभि\appbrk{}ता\appbrk{}न्क\appbrk{}लु\appbrk{}षो\appbrk{}द\appbrk{}कान्}}\Afootnote{A; \skt{क्षुभितां निर्मलोदकां}\ K H; \skt{क्षुभिता निर्मलोदकां}\ N}}\skt{~। }\pend
\pstart\versepar \skt{तरे\-त्क\-ल्या\-ण\-ला\-भाय व्याधे\-र\-प\-ग\-माय च~॥}\pend

\pstart\versepar \ledsidenote{\passref{1.29.78}}\edtext{\edtext{\skt{प्रासा\-दा\-न्स\-फ\-ला\-न्वृ\-क्षा\-न्वा\-र\-णा\-न्प\-र्व\-ता\-न्तथा}}{\lemma{\skt{प्रासा\appbrk{}दा\appbrk{}न्स\appbrk{}फ\appbrk{}ला\appbrk{}न्वृ\appbrk{}क्षा\appbrk{}न्वा\appbrk{}र\appbrk{}णा\appbrk{}न्प\appbrk{}र्व\appbrk{}तान्}}\Afootnote{H; \skt{महा\appbrk{}प्रा\appbrk{}सा\appbrk{}द\appbrk{}स\appbrk{}फ\appbrk{}ल\appbrk{}वृ\appbrk{}क्ष\appbrk{}वा\appbrk{}र\appbrk{}ण\appbrk{}पर्व}\ A; \skt{प्रासादां सफलां वृक्षां वारणां पर्वतांस्}\ K N}}}{\lemma{\skt{तथा}}\Afootnote{K N H; \skt{ता}\ A}}\skt{~। }\pend
\pstart\versepar \skt{}\edtext{\skt{आरु\-हे\-द्द्र\-व्य\-ला\-भाय}}{\lemma{\skt{आरु\appbrk{}हे\appbrk{}द्द्र\appbrk{}व्य\appbrk{}ला\appbrk{}भाय}}\Afootnote{K N H; \skt{न्~। आरो\appbrk{}हे\appbrk{}द्द्र\appbrk{}व्य\appbrk{}ला\appbrk{}भाय}\ A}}\skt{ व्याधे\-र\-प\-ग\-माय च~॥}\pend

\pstart\versepar \ledsidenote{\passref{1.29.81}}\edtext{\skt{ईदृग्विधां}}{\lemma{\skt{ईदृग्विधां}}\Afootnote{K N H; \skt{एवं}\ A}}\skt{ }\edtext{\skt{शुभान्स्वप्नान्यो वा}}{\lemma{\skt{शुभान्स्वप्नान्यो वा}}\Afootnote{\skt{रूपा\appbrk{}ञ्शु\appbrk{}भा\appbrk{}न्स्व\appbrk{}प्नान्यः}\ A; \skt{छुभान्यो न्वै}\ K; \skt{च्छुभांन्यान्वै}\ N; \skt{छुभान्यो वा}\ H}}\skt{ स्वप्ना\-न्प\-श्ये\-त्स\-दा\-तुरः~। }\pend
\pstart\versepar \skt{स दीर्घायुरिति }\edtext{\skt{ज्ञेयस्तस्मै}}{\lemma{\skt{ज्ञेयस्}}\Afootnote{A; \skt{ख्यातस्}\ K N H}}\skt{ कर्म }\edtext{\skt{समाचरेत्}}{\lemma{\skt{समाचरेत्}}\Afootnote{A H; \skt{समाचरेदिति~॥ २९~॥ {\symfont ❈}}\ K; \skt{समाचरेदिति~॥ २९}\ N}}\skt{~॥}\edtext{}{\lemma{\skt{स्वप्ना\appbrk{}न्प\appbrk{}श्ये\appbrk{}त्स\appbrk{}दा\appbrk{}तुरः~। स}}\Afootnote{K H; \skt{पश्येद्व्याधितो नरः~। स}\ A; \skt{स्वप्नां पस्येत्सदातुरः~। स}\ N}}\pend

\pstart \ledsidenote{\passref{1.29.trailer}}\skt{इति सुश्रु\-त\-सं\-हि\-तायां सूत्रस्थाने विप\-री\-ता\-वि\-प\-री\-त\-स्व\-प्न\-नि\-द\-र्श\-नीयो नामै\-को\-न\-त्रिं\-श\-त्त\-मोऽ\-ध्यायः~॥}\pend

\endnumbering
\beginnumbering
\eledsection*{Adhyāya 30}
\pstart \ledsidenote{\passref{1.30.1}}\skt{अथातः }\edtext{\skt{पञ्चे\-न्द्रि\-या\-र्थ\-वि\-प्र\-ति\-प\-त्ति\-म\-ध्यायं}}{\lemma{\skt{अध्यायं}}\Afootnote{A N; \emph{om.}\ K H}}\skt{ }\edtext{\skt{व्याख्यास्यामः}}{\lemma{\skt{व्याख्यास्यामः}}\Afootnote{A N H; \emph{om.}\ K}}\skt{~॥}\pend

\pstart\versepar \ledsidenote{\passref{1.30.3}}\edtext{\skt{शरी\-र\-शी\-ल\-यो\-र्यस्य}}{\lemma{\skt{शरी\appbrk{}र\appbrk{}शी\appbrk{}ल\appbrk{}योर्}}\Afootnote{A K H; \skt{शरीरशीलयो}\ N}}\skt{ }\edtext{\skt{प्रकृ\-ते\-र्वि\-कृ\-ति\-र्भ\-वेत्}}{\lemma{\skt{प्रकृतेर्}}\Afootnote{A; \skt{प्रकृतिर्}\ K N H}}\skt{~। }\pend
\pstart\versepar \skt{तत्त्व रिष्टं समासेन व्यासतस्तु निबोध मे~॥}\edtext{}{\lemma{\skt{भवेत्~। तत्त्व}}\Afootnote{A; \skt{भवेत्~। तद}\ K\textsuperscript{ac}; \skt{भवेत्~। तदा}\ K\textsuperscript{pc}; \skt{भवेत्तदा}\ N H}}\pend

\pstart\versepar \ledsidenote{\passref{1.30.4}}\skt{शृणोति }\edtext{\skt{विविधान्शब्दान्यो}}{\lemma{\skt{विविधान्शब्दान्यो}}\Afootnote{K H; \skt{विविधाञ्शब्दान्यो}\ A; \skt{विविधां शब्दां यो}\ N}}\skt{ दिव्या\-ना\-म\-भा\-वतः~। }\pend
\pstart\versepar \skt{समु\-द्र\-पु\-र\-मे\-घा\-ना\-म\-सं\-पत्तौ च }\edtext{\skt{निःस्वनान्}}{\lemma{\skt{निःस्वनान्}}\Afootnote{A; \skt{तां स्वनां}\ K N; \skt{तान्स्वनान्}\ H}}\skt{~॥}\edtext{}{\lemma{\skt{दिव्या\appbrk{}ना\appbrk{}म\appbrk{}भा\appbrk{}वतः~। समु\appbrk{}द्र\appbrk{}पु\appbrk{}र\appbrk{}मे\appbrk{}घा\appbrk{}नाम्}}\Afootnote{A; \skt{विद्यानसतो बहू\appbrk{}न्स\appbrk{}मु\appbrk{}द्र\appbrk{}पु\appbrk{}र\appbrk{}मे\appbrk{}घा\appbrk{}नाम्}\ K N; \skt{दिव्यानसतो बहून्~। समु\appbrk{}द्र\appbrk{}पु\appbrk{}र\appbrk{}मे\appbrk{}घा\appbrk{}नाम्}\ H}}\pend

\pstart\versepar \ledsidenote{\passref{1.30.5}}\edtext{\skt{तान्स्व\-ना\-न्ना\-व\-गृ\-ह्णाति}}{\lemma{\skt{तान्स्व\appbrk{}ना\appbrk{}न्ना\appbrk{}व\appbrk{}गृ\appbrk{}ह्णाति}}\Afootnote{A; \skt{तां स्वनाम्वा न गृह्णीते}\ K N; \skt{तान्स्वनाम्वा न गृह्णीते}\ H}}\skt{ }\edtext{\skt{मन्यते}}{\lemma{\skt{मन्यते}}\Afootnote{A; \skt{गृह्णीते}\ K N H}}\skt{ चान्यशब्दवत्~। }\pend
\pstart\versepar \skt{ग्राम्या\-र\-ण्य\-स्व\-नां\-श्चापि }\edtext{\skt{विप\-री\-ता\-न्शृ\-णोति च}}{\lemma{\skt{विप\appbrk{}री\appbrk{}ता\appbrk{}न्शृ\appbrk{}णोति च}}\Afootnote{\skt{विप\appbrk{}री\appbrk{}ता\appbrk{}ञ्शृ\appbrk{}णोति च}\ A; \skt{गृह्णाति विपरीतवत्}\ K N H}}\skt{~॥}\edtext{}{\lemma{\skt{चान्यशब्दवत्~। ग्राम्या\appbrk{}र\appbrk{}ण्य\appbrk{}स्व\appbrk{}नांश्}}\Afootnote{A; \skt{चान्य\appbrk{}श\appbrk{}ब्द\appbrk{}व\appbrk{}त्ग्रा\appbrk{}मा\appbrk{}र\appbrk{}ण्य\appbrk{}स्व\appbrk{}नांश्}\ K N H}}\pend

\pstart\versepar \ledsidenote{\passref{1.30.6}}\edtext{\skt{द्विषच्छब्देषु}}{\lemma{\skt{द्विषच्छब्देषु}}\Afootnote{A; \skt{द्विषच्छब्देन}\ K N H}}\skt{ रमते }\edtext{\skt{सुहृच्छब्देषु}}{\lemma{\skt{सुहृच्छब्देषु}}\Afootnote{A; \skt{सुहृच्छब्देन}\ K N H}}\skt{ कुप्यति~। }\pend
\pstart\versepar \skt{न शृणोति }\edtext{\skt{च योऽकस्मात्तं}}{\lemma{\skt{च योऽकस्मात्तं}}\Afootnote{A; \skt{यच्चाकस्मान्न}\ K N H}}\skt{ }\edtext{\skt{ब्रुवन्ति}}{\lemma{\skt{ब्रुवन्ति}}\Afootnote{A; \skt{गृह्णीते}\ K N H}}\skt{ }\edtext{\skt{गतायुषम्}}{\lemma{\skt{गतायुषम्}}\Afootnote{A; \skt{गतासुन्तं प्रचक्षते}\ K N H}}\skt{~॥}\edtext{}{\lemma{\skt{कुप्यति~। न शृणोति}}\Afootnote{A; \skt{कुप्यति}\ K N H}}\pend

\pstart\versepar \ledsidenote{\passref{1.30.7}}\edtext{\edtext{\skt{यत्तूष्णमिति}}{\lemma{\skt{यत्तूष्णम्}}\Afootnote{K N H; \skt{यस्तूष्णम्}\ A}}}{\lemma{\skt{इति}}\Afootnote{K N H; \skt{इव}\ A}}\skt{ गृह्णाति }\edtext{\edtext{\skt{शीतमुष्णं}}{\lemma{\skt{शीतम्}}\Afootnote{K N H; \skt{शीत}\ A}}\skt{ च}}{\lemma{\skt{उष्णं च}}\Afootnote{K N H; \skt{मुष्णं च}\ A}}\skt{ शीतवत्~। }\pend
\pstart\versepar \skt{संजा\-त\-शी\-त\-पि\-डको यश्च }\edtext{\skt{दाहेन}}{\lemma{\skt{दाहेन}}\Afootnote{A K H; \skt{दाह्येन}\ N}}\skt{ पीड्यते~॥}\edtext{}{\lemma{\skt{शीतवत्~। संजा\appbrk{}त\appbrk{}शी\appbrk{}त\appbrk{}पि\appbrk{}डको यश्च}}\Afootnote{A; \skt{शीतवत्~। सञ्जा\appbrk{}त\appbrk{}शी\appbrk{}त\appbrk{}पि\appbrk{}टको यश्च}\ K; \skt{शीतवत्~। सञ्जा\appbrk{}त\appbrk{}शी\appbrk{}त\appbrk{}पि\appbrk{}ढकौ यश्च}\ N; \skt{शीतवत्~। सञ्जा\appbrk{}त\appbrk{}शी\appbrk{}त\appbrk{}पि\appbrk{}टकौ यश्च}\ H}}\pend

\pstart\versepar \ledsidenote{\passref{1.30.8}}\edtext{\skt{उष्ण\-गा\-त्रोऽ\-ति\-गा\-त्रञ्च यः}}{\lemma{\skt{ऽतिगात्रञ्च यः}}\Afootnote{K N H; \skt{ऽतिमात्रं च यः}\ A}}\skt{ शीतेन प्रवेपते~। }\pend
\pstart\versepar \skt{}\edtext{\edtext{\skt{प्रहा\-रा\-न्ना\-भि\-जा\-नाति}}{\lemma{\skt{प्रहारान्}}\Afootnote{A K; \skt{प्रहारन्}\ N; \skt{प्रहारा}\ H}}\skt{ }\edtext{\skt{योऽङ्ग\-च्छे\-द\-म\-थापि}}{\lemma{\skt{ऽङ्गच्छेदम्}}\Afootnote{\skt{ऽङ्गच्छेद}\ A; \skt{गच्छे त}\ K\textsuperscript{ac} N H; \skt{गच्छे त्तु}\ K\textsuperscript{pc}}}}{\lemma{\skt{नाभिजानाति यो}}\Afootnote{A; \skt{नाभिजानाति स}\ K N; \skt{त्नाभिजानाति स}\ H}}\skt{ वा~॥}\edtext{}{\lemma{\skt{अथापि वा}}\Afootnote{\skt{मथापि वा}\ A; \skt{यमालयम्}\ K N H}}\pend

\pstart\versepar \ledsidenote{\passref{1.30.9}}\edtext{\skt{पांशु\-नै\-वा\-व\-की\-र्णानि यश्च}}{\lemma{\skt{पांशु\appbrk{}नै\appbrk{}वा\appbrk{}व\appbrk{}की\appbrk{}र्णानि यश्च}}\Afootnote{K N H; \skt{पांशु\appbrk{}ने\appbrk{}वा\appbrk{}व\appbrk{}की\appbrk{}र्णानि यश्च}\ A}}\skt{ गात्राणि मन्यते~। }\pend
\pstart\versepar \skt{वर्णान्यता वा }\edtext{\skt{राज्यो वा}}{\lemma{\skt{राज्यो वा}}\Afootnote{A K H; \skt{राज्ये वा}\ N}}\skt{ यस्य गात्रे भवन्ति हि~॥}\pend

\pstart\versepar \ledsidenote{\passref{1.30.10}}\edtext{\skt{स्नातानुलिप्तं यं}}{\lemma{\skt{स्नातानुलिप्तं यं}}\Afootnote{A K H; \skt{स्नातानुलिप्त यञ्}\ N}}\skt{ चापि भजन्ते नीलमक्षिकाः~। }\pend
\pstart\versepar \skt{}\edtext{\skt{सुगन्धिर्वाऽति}}{\lemma{\skt{सुगन्धिर्वाऽति}}\Afootnote{A; \skt{सुगन्धो वाति}\ K N H}}\skt{ }\edtext{\skt{योऽकस्मात्तं}}{\lemma{\skt{योऽकस्मात्तं}}\Afootnote{A; \skt{चाकस्मात्तं}\ K H; \skt{चाकस्मात्त}\ N}}\skt{ ब्रुवन्ति गतायुषम्~॥}\pend

\pstart\versepar \ledsidenote{\passref{1.30.11}}\skt{विपरीतेन गृह्णाति }\edtext{\skt{रसा\-न्य\-श्चो\-प\-यो\-जि\-तान्}}{\lemma{\skt{रसान्यश्}}\Afootnote{A H; \skt{रसां यश्}\ K N}}\skt{~। }\pend
\pstart\versepar \skt{उपयुक्ताः }\edtext{\skt{क्रमाद्यस्य रसा दोषाभिवृद्धये}}{\lemma{\skt{क्रमाद्यस्य रसा दोषाभिवृद्धये}}\Afootnote{A; \skt{च रसान्यस्य दोषाभिवृद्धये}\ K H; \skt{च रसां यस्य दोषाभिवृद्धये}\ N}}\skt{~॥}\edtext{}{\lemma{\skt{चोपयोजितान्~। उपयुक्ताः}}\Afootnote{A; \skt{चोपयोजितां क्रमोपयुक्तांश्}\ K N; \skt{चोपयोजितान्~। क्रमोपयुक्ताश्}\ H}}\pend

\pstart\versepar \ledsidenote{\passref{1.30.12}}\skt{यस्य दोषाग्निसाम्यं च }\edtext{\skt{कुर्यु\-र्मि\-थ्यो\-प\-यो\-जिताः}}{\lemma{\skt{कुर्युर्}}\Afootnote{A K H; \skt{कुर्यु}\ N}}\skt{~। }\pend
\pstart\versepar \skt{यो वा }\edtext{\skt{रसान्न}}{\lemma{\skt{रसान्न}}\Afootnote{A; \skt{रसं न}\ K N H}}\skt{ }\edtext{\skt{संवेत्ति}}{\lemma{\skt{संवेत्ति}}\Afootnote{A; \skt{सम्वेत्ति तं}\ K H; \skt{सम्वेति तम्}\ N}}\skt{ }\edtext{\skt{गतासुं तं}}{\lemma{\skt{गतासुं तं}}\Afootnote{A; \skt{ब्रुव गतायुषं}\ K\textsuperscript{ac}; \skt{ब्रुवन्ति गतायुषं}\ K\textsuperscript{pc} H N}}\skt{ }\edtext{\skt{प्रचक्षते}}{\lemma{\skt{प्रचक्षते}}\Afootnote{A; \emph{om.}\ K N H}}\skt{~॥}\pend

\pstart\versepar \ledsidenote{\passref{1.30.13}}\edtext{\skt{सुगन्धं}}{\lemma{\skt{सुगन्धं}}\Afootnote{A; \skt{सुरभिं दुरभिं}\ K H; \skt{सुर\appbrk{}भी\appbrk{}न्दु\appbrk{}र\appbrk{}भिम्}\ N}}\skt{ }\edtext{\skt{वेत्ति}}{\lemma{\skt{वेत्ति}}\Afootnote{A K H; \skt{वेति}\ N}}\skt{ }\edtext{\skt{दुर्गन्धं}}{\lemma{\skt{दुर्गन्धं}}\Afootnote{A; \skt{दुरभिं}\ K H; \skt{दुरभि}\ N}}\skt{ }\edtext{\skt{दुर्गन्धस्य}}{\lemma{\skt{दुर्गन्धस्य}}\Afootnote{A; \skt{सुरभीति च~। यो}\ K H; \skt{न्सुरभीति च~। यो}\ N}}\skt{ }\edtext{\skt{सुगन्धिताम्}}{\lemma{\skt{सुगन्धिताम्}}\Afootnote{A; \skt{वा गन्धन्न}\ K N H}}\skt{~। }\pend
\pstart\versepar \skt{}\edtext{\skt{गृह्णीते}}{\lemma{\skt{गृह्णीते}}\Afootnote{A; \skt{गृह्णाति}\ K N H}}\skt{ }\edtext{\skt{वाऽन्यथा गन्धं}}{\lemma{\skt{वाऽन्यथा गन्धं}}\Afootnote{\skt{वाऽन्यथा गन्धं}\ A; \emph{om.}\ K H; \skt{शान्ते च}\ N}}\skt{ }\edtext{\skt{शान्ते}}{\lemma{\skt{शान्ते}}\Afootnote{A K H; \emph{om.}\ N}}\skt{ दीपे च }\edtext{\skt{नीरुजः}}{\lemma{\skt{नीरुजः}}\Afootnote{A; \skt{मानवः}\ K N H}}\skt{~॥}\pend

\pstart \ledsidenote{\passref{1.30.15cd}}\skt{दिवा ज्योतींषि }\edtext{\skt{यश्चा ज्वलितानीव}}{\lemma{\skt{चा ज्वलितानीव}}\Afootnote{N; \skt{चापि ज्वलितानीव}\ A K H}}\skt{ पश्यति~।}\pend

\pstart\versepar \ledsidenote{\passref{1.30.16}}\skt{रात्रौ }\edtext{\skt{सूर्यं}}{\lemma{\skt{सूर्यं}}\Afootnote{A; \skt{चन्द्रं}\ K N H}}\skt{ ज्वलन्तं वा }\edtext{\skt{दिवा वा}}{\lemma{\skt{दिवा वा}}\Afootnote{A; \skt{सूर्यम्वा}\ K N H}}\skt{ चन्द्रवर्चसम्~। }\pend
\pstart\versepar \skt{अमेघोपप्लवे यश्च }\edtext{\skt{शक्र\-चा\-प\-त\-डि\-द्गु\-णान्}}{\lemma{\skt{शक्र\appbrk{}चा\appbrk{}प\appbrk{}त\appbrk{}डि\appbrk{}द्गु\appbrk{}णान्}}\Afootnote{A H; \skt{शक्र\appbrk{}चा\appbrk{}प\appbrk{}त\appbrk{}डि\appbrk{}द्गुणां}\ K; \skt{शक्र\appbrk{}चा\appbrk{}प\appbrk{}त\appbrk{}डि\appbrk{}गुणां}\ N}}\skt{~॥}\pend

\pstart\versepar \ledsidenote{\passref{1.30.17}}\edtext{\edtext{\skt{तडि\-त्त्व\-तोऽ\-सि\-तान्यो}}{\lemma{\skt{तडित्त्वतो}}\Afootnote{A; \skt{तडिद्वतो}\ K N H}}\skt{ वा}}{\lemma{\skt{ऽसितान्यो वा}}\Afootnote{A; \skt{गुणान्यच्च}\ K N H}}\skt{ निर्मले गगने घनान्~। }\pend
\pstart\versepar \skt{विमा\-न\-या\-न\-प्रा\-सा\-दै\-र्यश्च }\edtext{\skt{संकु\-ल\-म\-म्ब\-रम्}}{\lemma{\skt{संकु\appbrk{}ल\appbrk{}म\appbrk{}म्ब\appbrk{}रम्}}\Afootnote{A; \skt{संकुलमम्बरं}\ K N H}}\skt{~॥}\edtext{}{\lemma{\skt{घनान्~। विमा\appbrk{}न\appbrk{}या\appbrk{}न\appbrk{}प्रा\appbrk{}सा\appbrk{}दै\appbrk{}र्यश्च}}\Afootnote{A H; \skt{घनां विमा\appbrk{}न\appbrk{}या\appbrk{}न\appbrk{}प्रा\appbrk{}सा\appbrk{}दै\appbrk{}र्यश्च}\ K N}}\pend

\pstart \ledsidenote{\passref{1.30.18ab}}\edtext{\edtext{\skt{यश्चाप्यनिर्म्मलं}}{\lemma{\skt{यश्चाप्य्}}\Afootnote{K N H; \skt{यश्चा}\ A}}}{\lemma{\skt{अनिर्म्मलं}}\Afootnote{K N H; \skt{निलं}\ A}}\skt{ }\edtext{\edtext{\skt{मूर्त्तिमन्तरिक्षे}}{\lemma{\skt{मूर्त्तिम्}}\Afootnote{K N H; \skt{मूर्ति}\ A}}}{\lemma{\skt{अन्तरिक्षे}}\Afootnote{K N H; \skt{मन्तमन्तरिक्षं}\ A}}\skt{ }\edtext{\skt{प्रपश्यति}}{\lemma{\skt{प्रपश्यति}}\Afootnote{K N H; \skt{च पश्यति}\ A}}\skt{~॥}\pend

\pstart\versepar \ledsidenote{\passref{1.30.18.add}}\skt{द्वन्द्वानि }\edtext{\skt{शीतोष्णञ्च}}{\lemma{\skt{शीतोष्णञ्च}}\Afootnote{\skt{शीतमुष्कञ्च}\ H}}\skt{ कालावस्था दिशस्तथा~। }\pend
\pstart\versepar \skt{}\edtext{\skt{विपरीतेन}}{\lemma{\skt{विपरीतेन}}\Afootnote{H\textsuperscript{pc}; \skt{विपरीते}\ H\textsuperscript{ac}}}\skt{ गृह्णाति }\edtext{\skt{भावानान्यश्च यो}}{\lemma{\skt{भावानान्यश्च यो}}\Afootnote{\skt{भावानान्योश्च यो}\ H}}\skt{ नरः॥}\pend

\pstart \ledsidenote{\passref{1.30.18cd}}\skt{धूम\-नी\-हा\-र\-वा\-सो\-भि\-रा\-वृ\-ता\-मिव मेदिनीम्~॥}\pend

\pstart\versepar \ledsidenote{\passref{1.30.19}}\skt{प्रदीप्तमिव लोकं च यो वा प्लुतमिवाम्भसा~। }\pend
\pstart\versepar \skt{भूमि\-म\-ष्टा\-प\-दा\-कारां लेखाभिर्यश्च पश्यति~॥}\pend

\pstart\versepar \ledsidenote{\passref{1.30.23}}\skt{यो वा मयूरकण्ठाभं विधू\-म\-म्व\-ह्नि\-मी\-क्षते~। }\pend
\pstart\versepar \skt{}\edtext{\edtext{\skt{ध्रुव\-मा\-का\-श\-गङ्गां}}{\lemma{\skt{ध्रुवम्}}\Afootnote{K N H; \emph{om.}\ A}}\skt{ वा तं}}{\lemma{\skt{आकाशगङ्गां वा तं}}\Afootnote{K N H; \emph{om.}\ A}}\skt{ }\edtext{\skt{ब्रुवन्ति}}{\lemma{\skt{ब्रुवन्ति}}\Afootnote{K N H; \emph{om.}\ A}}\skt{ }\edtext{\skt{गतायुषं}}{\lemma{\skt{गतायुषं}}\Afootnote{K N H; \skt{आतुरस्य भवेन्मृत्युः स्वस्थो व्याधि\appbrk{}म\appbrk{}वा\appbrk{}प्नु\appbrk{}यात्}\ A}}\pend

\pstart\versepar \ledsidenote{\passref{1.30.21}}\edtext{\skt{ज्योत्स्ना\-द\-र्शो\-ष्ण\-तो\-येषु}}{\lemma{\skt{ज्योत्स्ना\appbrk{}द\appbrk{}र्शो\appbrk{}ष्ण\appbrk{}तो\appbrk{}येषु}}\Afootnote{A; \skt{। ज्योत्स्ना\appbrk{}द\appbrk{}र्श\appbrk{}तो\appbrk{}येषु}\ K\textsuperscript{ac}; \skt{यो ज्योत्स्ना\appbrk{}द\appbrk{}र्श\appbrk{}तो\appbrk{}येषु श्}\ K\textsuperscript{pc}; \skt{यो ज्योत्स्ना\appbrk{}द\appbrk{}र्श\appbrk{}तो\appbrk{}येषु}\ N H}}\skt{ }\edtext{\skt{छायां यश्च न}}{\lemma{\skt{छायां यश्च न}}\Afootnote{A; \skt{यश्च च्छायाञ्च}\ K\textsuperscript{ac}; \skt{च च्छायाञ्चन}\ K\textsuperscript{pc}; \skt{स्वच्छायां च न}\ N H}}\skt{ पश्यति~। }\pend
\pstart\versepar \skt{}\edtext{\skt{पश्य\-त्ये\-का\-ङ्ग\-हीनां वा}}{\lemma{\skt{एकाङ्गहीनां वा}}\Afootnote{A K H; \skt{एकाङ्गहीनम्वा}\ N}}\skt{ विकृतां }\edtext{\skt{वाऽन्यसत्त्वजाम्}}{\lemma{\skt{वाऽन्यसत्त्वजाम्}}\Afootnote{A K N; \skt{वाऽन्यसत्वजां}\ H}}\skt{~॥}\pend

\pstart\versepar \ledsidenote{\passref{1.30.22}}\skt{श्वका\-क\-क\-ङ्क\-गृ\-ध्राणां प्रेतानां }\edtext{\skt{यक्षरक्षसाम्}}{\lemma{\skt{यक्षरक्षसाम्}}\Afootnote{A H; \skt{यक्षराक्षसां}\ K N}}\skt{~। }\pend
\pstart\versepar \skt{पिशा\-चो\-र\-ग\-ना\-गानां भूतानां }\edtext{\skt{विकृतामपि}}{\lemma{\skt{विकृतामपि}}\Afootnote{A; \skt{विकृतानपि}\ K N H}}\skt{~॥}\pend

\pstart\versepar \ledsidenote{\passref{1.30.22.add1}}\skt{स्वस्थः स लभते व्याधिं व्याधितो }\edtext{\edtext{\skt{मृत्युमर्च्छति}}{\lemma{\skt{मृत्युम्}}\Afootnote{K H; \skt{मृत्यु}\ N}}}{\lemma{\skt{अर्च्छति}}\Afootnote{\skt{अर्च्छति~॥ {\symfont ❈}}\ K; \skt{गच्छति}\ N; \skt{अर्च्छति~॥ …}\ H}}\skt{~॥~॥ }\pend
\pstart\versepar \skt{व्रण\-प्र\-श्न\-म्व्र\-णा\-स्रावं कृत्या\-कृ\-त्य\-वि\-धि\-न्तथा~।}\pend

\pstart\versepar \ledsidenote{\passref{1.30.22.add2}}\skt{व्याध्यु\-द्दे\-शी\-य\-म\-ध्यायं शस्त्र\-क\-र्म्मा\-ष्ट\-क\-न्तथा~॥ }\pend
\pstart\versepar \skt{प्रण\-ष्ट\-श\-ल्य\-वि\-ज्ञानं }\edtext{\edtext{\skt{शल्या\-प\-न\-य\-न\-मेव}}{\lemma{\skt{शल्या\appbrk{}प\appbrk{}न\appbrk{}य\appbrk{}नम्}}\Afootnote{K\textsuperscript{pc} N H; \skt{शल्यापनयन}\ K\textsuperscript{ac}}}}{\lemma{\skt{एव}}\Afootnote{K\textsuperscript{pc} N H; \skt{व}\ K\textsuperscript{ac}}}\skt{ च~।}\pend

\pstart\versepar \ledsidenote{\passref{1.30.22.add3}}\skt{विप\-री\-त\-व्र\-ण\-ज्ञानं दूत\-स्व\-प्न\-वि\-प\-र्ययं~॥ }\pend
\pstart\versepar \skt{पञ्चे\-न्द्रि\-या\-र्थ\-वि\-भ्रा\-न्ति\-म्प्रोक्तं वै तृतीयो }\edtext{\skt{दश}}{\lemma{\skt{दश}}\Afootnote{K H; \skt{दशः}\ N}}\skt{~॥}\pend

\pstart \ledsidenote{\passref{1.30.trailer}}\edtext{\skt{॥ ३०~॥}}{\lemma{\skt{॥ ३०~॥}}\Afootnote{N H; \skt{इति सुश्रु\appbrk{}त\appbrk{}सं\appbrk{}हि\appbrk{}तायां सूत्रस्थाने पञ्चे\appbrk{}न्द्रि\appbrk{}या\appbrk{}र्थ\appbrk{}वि\appbrk{}प्र\appbrk{}ति\appbrk{}प\appbrk{}त्ति\appbrk{}र्नाम त्रिंशोऽध्यायः}\ A; \skt{{\symfont ❈}}\ K}}\pend

\endnumbering
\beginnumbering
\eledsection*{Adhyāya 31}
\pstart \ledsidenote{\passref{1.31.1}}\edtext{\skt{अथा\-त\-श्च्छा\-या\-वि\-प्र\-ति\-प\-त्ति\-म्व्या\-ख्या\-स्यामः}}{\lemma{\skt{व्याख्यास्यामः}}\Afootnote{H; \skt{अध्यायं व्याख्यास्यामः}\ A N; \emph{om.}\ K}}\skt{~॥}\pend

\pstart\versepar \ledsidenote{\passref{1.31.3}}\skt{श्यावा लोहितिका }\edtext{\skt{नीलाः}}{\lemma{\skt{नीलाः}}\Afootnote{K\textsuperscript{pc} H; \skt{नीला}\ A K\textsuperscript{ac}; \skt{नीः}\ N}}\skt{ पीतिका }\edtext{\skt{वापि}}{\lemma{\skt{वापि}}\Afootnote{A K N; \skt{चापि}\ H}}\skt{ }\edtext{\skt{देहिनां}}{\lemma{\skt{देहिनां}}\Afootnote{K N H; \skt{मानवम्}\ A}}\skt{~। }\pend
\pstart\versepar \skt{अभिद्रवन्ति }\edtext{\skt{यञ्छाया स}}{\lemma{\skt{छाया स}}\Afootnote{K N H; \skt{छायाः स}\ A}}\skt{ परा\-सु\-र\-सं\-शयं~।}\pend

\pstart\versepar \ledsidenote{\passref{1.31.4}}\edtext{\skt{ह्रीरपक्रामति}}{\lemma{\skt{ह्रीरपक्रामति}}\Afootnote{K N H; \skt{ह्रीरपक्रमते}\ A}}\skt{ }\edtext{\skt{यतः}}{\lemma{\skt{यतः}}\Afootnote{K N H; \skt{यस्य}\ A}}\skt{ }\edtext{\skt{कान्ति\-स्मृ\-ति\-धृ\-ति\-श्रियः}}{\lemma{\skt{कान्ति\appbrk{}स्मृ\appbrk{}ति\appbrk{}धृ\appbrk{}ति\appbrk{}श्रियः}}\Afootnote{K N H; \skt{प्रभा\appbrk{}स्मृ\appbrk{}ति\appbrk{}धृ\appbrk{}ति\appbrk{}श्रियः}\ A}}\skt{~। }\pend
\pstart\versepar \skt{अकस्माद्यं }\edtext{\skt{भजन्ते च स}}{\lemma{\skt{भजन्ते च स}}\Afootnote{K N H; \skt{भजन्ते वा स}\ A}}\skt{ }\edtext{\skt{परा\-सु\-र\-सं\-श\-यम्}}{\lemma{\skt{परासुर्}}\Afootnote{K N H; \skt{गतासुर्}\ A}}\skt{~।}\pend

\pstart\versepar \ledsidenote{\passref{1.31.5}}\edtext{\skt{यस्याधरौष्ठः}}{\lemma{\skt{यस्याधरौष्ठः}}\Afootnote{A K H; \skt{यस्याधरौष्ठं}\ N}}\skt{ पतितः }\edtext{\edtext{\skt{क्षिप्त\-श्चो\-र्ध\-न्त\-थो\-त्तरः}}{\lemma{\skt{क्षिप्तश्}}\Afootnote{A K N; \skt{क्षिप्तञ्}\ H}}}{\lemma{\skt{चोर्धन्}}\Afootnote{K N H; \skt{चोर्ध्वं}\ A}}\skt{~। }\pend
\pstart\versepar \skt{उभौ वा }\edtext{\skt{जाम्बवाभासौ स}}{\lemma{\skt{जाम्बवाभासौ स}}\Afootnote{K N H; \skt{जाम्बवाभासौ}\ A}}\skt{ }\edtext{\edtext{\skt{परा\-सु\-र\-सं\-श\-यम्}}{\lemma{\skt{परासुर्}}\Afootnote{K N H; \skt{दु}\ A}}}{\lemma{\skt{असंशयम्}}\Afootnote{K N H; \skt{र्लभं तस्य जीवितम्}\ A}}\skt{~।}\pend

\pstart\versepar \ledsidenote{\passref{1.31.6}}\skt{आरक्ता दशना यस्य }\edtext{\skt{श्यावा वा स्युः}}{\lemma{\skt{श्यावा वा स्युः}}\Afootnote{A K H; \skt{स्यावा व स्युः}\ N}}\skt{ पतन्ति वा~। }\pend
\pstart\versepar \skt{}\edtext{\skt{खाञ्जनप्रतिभो}}{\lemma{\skt{खाञ्जनप्रतिभो}}\Afootnote{K\textsuperscript{pc}; \skt{खञ्जनप्रतिमा}\ A; \skt{ख्ञ्जनप्रतिभो}\ K\textsuperscript{ac}; \skt{खाञ्जनप्रतिभा}\ N; \skt{खञ्जनप्रतिभा}\ H}}\skt{ वापि तं गता\-यु\-ष\-मा\-दि\-शेत्~।}\pend

\pstart\versepar \ledsidenote{\passref{1.31.7}}\skt{कृष्णा }\edtext{\skt{स्निग्धावलिप्ता वा}}{\lemma{\skt{स्निग्धावलिप्ता वा}}\Afootnote{K H; \skt{स्तब्धाऽवलिप्ता वा}\ A; \skt{स्निग्धावलिप्ता}\ N}}\skt{ जिह्वा }\edtext{\skt{शूना च}}{\lemma{\skt{शूना च}}\Afootnote{A K H; \skt{गुना च}\ N}}\skt{ यस्य वै~। }\pend
\pstart\versepar \skt{}\edtext{\skt{कर्कशा च}}{\lemma{\skt{कर्कशा च}}\Afootnote{K N H; \skt{कर्कशा वा}\ A}}\skt{ भवेद्यस्य सो चिरा\-द्वि\-ज\-हा\-त्य\-सून्~।}\pend

\pstart\versepar \ledsidenote{\passref{1.31.8}}\skt{कुटिला स्फुटिता वापि }\edtext{\skt{शूना वा}}{\lemma{\skt{शूना वा}}\Afootnote{K N H; \skt{शुष्का वा}\ A}}\skt{ यस्य नासिका~। }\pend
\pstart\versepar \skt{}\edtext{\skt{भग्ना वा}}{\lemma{\skt{भग्ना वा}}\Afootnote{K N H; \skt{अवस्फूर्जति मग्ना वा}\ A}}\skt{ }\edtext{\skt{स्फुरते}}{\lemma{\skt{स्फुरते}}\Afootnote{K N H; \emph{om.}\ A}}\skt{ }\edtext{\skt{वापि स}}{\lemma{\skt{वापि स}}\Afootnote{K N; \emph{om.}\ A; \skt{चापि स}\ H}}\skt{ }\edtext{\edtext{\skt{परा\-सु\-र\-सं\-शयं}}{\lemma{\skt{परासुर्}}\Afootnote{K N H; \emph{om.}\ A}}}{\lemma{\skt{असंशयं}}\Afootnote{K H; \skt{न स जीवति मानवः}\ A; \skt{असंसयः}\ N}}\skt{~।}\pend

\pstart\versepar \ledsidenote{\passref{1.31.9}}\skt{संक्षिप्ते विषमे स्तब्धे रक्ते }\edtext{\skt{सुप्ते च}}{\lemma{\skt{सुप्ते च}}\Afootnote{K N H; \skt{स्रस्ते च}\ A}}\skt{ लोचने~। }\pend
\pstart\versepar \skt{}\edtext{\skt{यस्य}}{\lemma{\skt{यस्य}}\Afootnote{K N H; \emph{om.}\ A}}\skt{ }\edtext{\skt{स्यातां}}{\lemma{\skt{स्यातां}}\Afootnote{K N H; \skt{स्यातां वा}\ A}}\skt{ }\edtext{\skt{श्रुते}}{\lemma{\skt{श्रुते}}\Afootnote{K N H; \skt{प्रस्रुते}\ A}}\skt{ }\edtext{\skt{वापि स}}{\lemma{\skt{वापि स}}\Afootnote{K N H; \skt{यस्य स}\ A}}\skt{ }\edtext{\edtext{\skt{परा\-सु\-र\-सं\-शयं}}{\lemma{\skt{परासुर्}}\Afootnote{K H; \skt{गतायु}\ A; \skt{परासु}\ N}}}{\lemma{\skt{असंशयं}}\Afootnote{K N H; \skt{र्नरो ध्रुवम्}\ A}}\skt{~।}\pend

\pstart\versepar \ledsidenote{\passref{1.31.11}}\edtext{\skt{न धारयति यः}}{\lemma{\skt{न धारयति यः}}\Afootnote{K\textsuperscript{pc} N H\textsuperscript{pc}; \emph{om.}\ A; \skt{न धारयन्ति यः}\ K\textsuperscript{ac} H\textsuperscript{ac}}}\skt{ }\edtext{\edtext{\edtext{\skt{शीर्ष\-न्ना\-ह\-र\-न्त्य\-न्न\-मा\-स्यगं}}{\lemma{\skt{शीर्षन्}}\Afootnote{K N H; \emph{om.}\ A}}}{\lemma{\skt{नाहरन्त्य्}}\Afootnote{K N H\textsuperscript{ac}; \skt{नाहरत्य्}\ A H\textsuperscript{pc}}}}{\lemma{\skt{आस्यगं}}\Afootnote{K N H; \skt{आस्यस्थं न}\ A}}\skt{~। }\pend
\pstart\versepar \skt{}\edtext{\skt{एका\-ग्र\-दृ\-ष्टि\-र्मू\-ढात्मा}}{\lemma{\skt{एकाग्रदृष्टिर्}}\Afootnote{K H; \skt{धारयति यः शिरः~। एकाग्रदृष्टिर्}\ A; \skt{एकाग्रदृष्टि}\ N}}\skt{ स }\edtext{\edtext{\skt{परा\-सु\-र\-सं\-शयं}}{\lemma{\skt{परासुर्}}\Afootnote{K H; \skt{द्यः प्राणान्जहाति सः}\ A; \skt{परासु य}\ N}}}{\lemma{\skt{असंशयं}}\Afootnote{K H; \emph{om.}\ A; \skt{संशयं}\ N}}\skt{~।}\pend

\pstart\versepar \ledsidenote{\passref{1.31.12}}\skt{बलवान्दुर्बलो वापि सम्मोहं यो }\edtext{\skt{धिगच्छति}}{\lemma{\skt{धिगच्छति}}\Afootnote{K N H; \skt{ऽधिगच्छति}\ A}}\skt{~। }\pend
\pstart\versepar \skt{उत्थाप्यमानो }\edtext{\skt{बहुशः स}}{\lemma{\skt{बहुशः स}}\Afootnote{K N H; \skt{बहुशस्तं}\ A}}\skt{ }\edtext{\edtext{\skt{परा\-सु\-र\-सं\-शयं}}{\lemma{\skt{परासुर्}}\Afootnote{K N H; \skt{प}\ A}}}{\lemma{\skt{असंशयं}}\Afootnote{K H; \skt{क्वं भिषगादिशेत्}\ A; \skt{असं शयः}\ N}}\skt{~।}\pend

\pstart\versepar \ledsidenote{\passref{1.31.13}}\skt{उत्तानः सर्वदा शेते~। पादौ }\edtext{\skt{विकुरुते च यः}}{\lemma{\skt{विकुरुते च यः}}\Afootnote{A K\textsuperscript{pc} N H; \skt{वि ः}\ K\textsuperscript{ac}}}\skt{~। }\pend
\pstart\versepar \skt{}\edtext{\skt{विप्र\-सा\-र\-ण\-शीलो वा स}}{\lemma{\skt{विप्र\appbrk{}सा\appbrk{}र\appbrk{}ण\appbrk{}शीलो वा स}}\Afootnote{K N H; \skt{विप्र\appbrk{}सा\appbrk{}र\appbrk{}ण\appbrk{}शीलो वा}\ A}}\skt{ }\edtext{\edtext{\skt{परा\-सु\-र\-सं\-शयं}}{\lemma{\skt{परासुर्}}\Afootnote{K N H; \skt{न}\ A}}}{\lemma{\skt{असंशयं}}\Afootnote{K H; \skt{स जीवति मानवः}\ A; \skt{असंशयः}\ N}}\skt{~।}\pend

\pstart\versepar \ledsidenote{\passref{1.31.14}}\skt{शीत\-पा\-द\-क\-रो\-च्छ्वासः }\edtext{\skt{छिन्नश्वासश्च यो}}{\lemma{\skt{छिन्नश्वासश्च यो}}\Afootnote{K H; \skt{छिन्नोच्छ्वासश्च यो}\ A; \skt{च्छिन्नास्वासश्च यो}\ N}}\skt{ भवेत्~। }\pend
\pstart\versepar \skt{काकोच्छ्वासश्च यो }\edtext{\skt{मर्त्यः स}}{\lemma{\skt{मर्त्यः स}}\Afootnote{K N H; \skt{मर्त्यस्}\ A}}\skt{ }\edtext{\edtext{\skt{परा\-सु\-र\-सं\-शयं}}{\lemma{\skt{परासुर्}}\Afootnote{K N H; \emph{om.}\ A}}}{\lemma{\skt{असंशयं}}\Afootnote{K N H; \skt{तं धीरः परिवर्जयेत्}\ A}}\skt{~।}\pend

\pstart\versepar \ledsidenote{\passref{1.31.15}}\skt{निद्रा न च्छिद्यते यस्य यो वा जागर्त्ति सर्वदा~। }\pend
\pstart\versepar \skt{मुह्येद्वा वक्तुकामश्च प्रत्याख्येयः स जानता~।}\pend

\pstart\versepar \ledsidenote{\passref{1.31.16}}\edtext{\edtext{\skt{परि\-लि\-हे\-दु\-त्त\-रो\-ष्ठ\-मु\-ङ्ग\-रांश्च}}{\lemma{\skt{परि\appbrk{}लि\appbrk{}हे\appbrk{}दु\appbrk{}त्त\appbrk{}रो\appbrk{}ष्ठम्}}\Afootnote{K N H; \skt{उत्तरौष्ठं}\ A}}}{\lemma{\skt{उङ्गरांश्च}}\Afootnote{K N H; \skt{च यो लिह्यादुत्कारांश्च}\ A}}\skt{ }\edtext{\skt{करोति वा}}{\lemma{\skt{करोति वा}}\Afootnote{K N H; \skt{करोति यः}\ A}}\skt{~। }\pend
\pstart\versepar \skt{प्रेतैर्वा भाषते }\edtext{\skt{सार्धं स}}{\lemma{\skt{सार्धं स}}\Afootnote{K N H; \skt{सार्धं प्रे}\ A}}\skt{ }\edtext{\edtext{\skt{परा\-सु\-र\-सं\-शयं}}{\lemma{\skt{परासुर्}}\Afootnote{K N H; \skt{तरू}\ A}}}{\lemma{\skt{असंशयं}}\Afootnote{K H; \skt{पं तमादिशेत्}\ A; \skt{असंशयः}\ N}}\skt{~।}\pend

\pstart\versepar \ledsidenote{\passref{1.31.17}}\skt{खेभ्यः सरोमकूपेभ्यो यस्य रक्तं प्रवर्त्तते~। }\pend
\pstart\versepar \skt{पुरु\-ष\-स्या\-वि\-षा\-र्त्तस्य स }\edtext{\edtext{\skt{परा\-सु\-र\-सं\-शयं}}{\lemma{\skt{परासुर्}}\Afootnote{K N H; \emph{om.}\ A}}}{\lemma{\skt{असंशयं}}\Afootnote{K N H; \skt{द्यो जह्यात्स जीवितम्}\ A}}\skt{~।}\pend

\pstart\versepar \ledsidenote{\passref{1.31.18}}\skt{वाताष्ठीला तु हृदये }\edtext{\skt{यस्यो\-र्ध्व\-म\-नु\-या\-यिनी}}{\lemma{\skt{यस्योर्ध्वम्}}\Afootnote{K; \skt{यस्योर्ध्वम्}\ A; \skt{यस्योर्द्धम्}\ N H}}\skt{~। }\pend
\pstart\versepar \skt{रुजा\-न्न\-वि\-द्वे\-ष\-करी स }\edtext{\skt{परा\-सु\-र\-सं\-शयं}}{\lemma{\skt{असंशयं}}\Afootnote{A K N; \skt{असंशयः}\ H}}\skt{~।}\pend

\pstart\versepar \ledsidenote{\passref{1.31.19}}\skt{अन\-न्यो\-प\-द्र\-व\-कृतः }\edtext{\skt{पादः}}{\lemma{\skt{पादः}}\Afootnote{K N H; \skt{शोफः}\ A}}\skt{ }\edtext{\skt{शोफः}}{\lemma{\skt{शोफः}}\Afootnote{K\textsuperscript{pc}; \skt{पाद}\ A; \skt{शोफ}\ K\textsuperscript{ac} N H}}\skt{ समुत्थितः }\pend
\pstart\versepar \skt{पुरुषं हन्ति }\edtext{\skt{नारीञ्च}}{\lemma{\skt{नारीञ्च}}\Afootnote{K N H; \skt{नारीं तु}\ A}}\skt{ मुखजो गुह्यजो द्वयं~।}\pend

\pstart\versepar \ledsidenote{\passref{1.31.20}}\edtext{\skt{अतीसारो}}{\lemma{\skt{अतीसारो}}\Afootnote{K N H; \skt{अतिसारो}\ A}}\skt{ ज्वरो }\edtext{\skt{ध्मानं}}{\lemma{\skt{ध्मानं}}\Afootnote{K N H; \skt{हिक्का}\ A}}\skt{ छर्द्दिः }\edtext{\skt{शूनाङ्गमेढ्रता}}{\lemma{\skt{शूनाङ्गमेढ्रता}}\Afootnote{K N; \skt{शूनाण्डमेढ्रता}\ A; \skt{सूनाङ्गमेढ्रता}\ H}}\skt{~। }\pend
\pstart\versepar \skt{}\edtext{\skt{कासिन}}{\lemma{\skt{कासिन}}\Afootnote{K N; \skt{श्वासिनः}\ A; \skt{कासिनः}\ H}}\skt{ }\edtext{\skt{स्वासिनो}}{\lemma{\skt{स्वासिनो}}\Afootnote{K N; \skt{कासिनो}\ A; \skt{श्वासिनो}\ H}}\skt{ वापि यस्य तं क्षीणमादिशेत्~।}\pend

\pstart\versepar \ledsidenote{\passref{1.31.21}}\skt{स्वेदो दाहश्च }\edtext{\skt{बलवां}}{\lemma{\skt{बलवां}}\Afootnote{K N; \skt{बलवा}\ A H}}\skt{ }\edtext{\skt{हिक्का}}{\lemma{\skt{हिक्का}}\Afootnote{K N; \skt{न्हिक्का}\ A H}}\skt{ श्वासश्च मानवं~। }\pend
\pstart\versepar \skt{बलवन्तमपि }\edtext{\skt{प्राणैर्वियुञ्जन्ति न}}{\lemma{\skt{वियुञ्जन्ति न}}\Afootnote{A K H; \skt{वियुजन्ति न}\ N}}\skt{ }\edtext{\skt{संशयं}}{\lemma{\skt{संशयं}}\Afootnote{K N H; \skt{संशयः}\ A}}\skt{~।}\pend

\pstart\versepar \ledsidenote{\passref{1.31.22}}\skt{श्यावा जिह्वा भवेद्यस्य सव्यं चाक्षि निमज्जति }\pend
\pstart\versepar \skt{मुखं च जायते }\edtext{\skt{पूतिर्यस्य}}{\lemma{\skt{पूतिर्}}\Afootnote{K N H; \skt{पूति}\ A}}\skt{ तं परिवर्जयेत्~।}\pend

\pstart\versepar \ledsidenote{\passref{1.31.23}}\edtext{\skt{नेत्रे}}{\lemma{\skt{नेत्रे}}\Afootnote{K N H; \emph{om.}\ A}}\skt{ }\edtext{\skt{चाम्रेण}}{\lemma{\skt{चाम्रेण}}\Afootnote{K H; \skt{वक्त्रमा}\ A; \skt{चास्येन}\ N}}\skt{ }\edtext{\skt{पूर्येते}}{\lemma{\skt{पूर्येते}}\Afootnote{K N H; \skt{पूर्यते}\ A}}\skt{ }\edtext{\skt{स्विद्येते}}{\lemma{\skt{स्विद्येते}}\Afootnote{K N H; \skt{ऽश्रूणां स्विद्यतश्}\ A}}\skt{ चरणौ }\edtext{\skt{तथा}}{\lemma{\skt{तथा}}\Afootnote{K N H; \skt{उभौ}\ A}}\skt{~। }\pend
\pstart\versepar \skt{चक्षुश्चाकुलतां याति यमराष्ट्रं गमिष्यतः~।}\pend

\pstart\versepar \ledsidenote{\passref{1.31.24}}\edtext{\skt{अतिमात्रं}}{\lemma{\skt{अतिमात्रं}}\Afootnote{A K H; \skt{अतिमात्र}\ N}}\skt{ लघूनि }\edtext{\skt{स्युर्गा\-त्रा\-ण्य\-ति\-गु\-रूणि वा}}{\lemma{\skt{अतिगुरूणि वा}}\Afootnote{K N H; \skt{गुरुकाणि वा}\ A}}\skt{~। }\pend
\pstart\versepar \skt{यस्याकस्मात्स विज्ञेयो गन्ता }\edtext{\skt{चैव}}{\lemma{\skt{चैव}}\Afootnote{K N H; \skt{वैव}\ A}}\skt{ }\edtext{\skt{यमालयं}}{\lemma{\skt{यमालयं}}\Afootnote{K N H; \skt{स्वतालयम्}\ A}}\skt{~।}\pend

\pstart\versepar \ledsidenote{\passref{1.31.25}}\edtext{\skt{पङ्क\-म\-त्स्य\-व\-सा\-तै\-ल\-घृ\-त\-ग\-न्धाश्च ये}}{\lemma{\skt{पङ्क\appbrk{}म\appbrk{}त्स्य\appbrk{}व\appbrk{}सा\appbrk{}तै\appbrk{}ल\appbrk{}घृ\appbrk{}त\appbrk{}ग\appbrk{}न्धाश्च ये}}\Afootnote{K; \skt{पङ्क\appbrk{}म\appbrk{}त्स्य\appbrk{}व\appbrk{}सा\appbrk{}तै\appbrk{}ल\appbrk{}घृ\appbrk{}त\appbrk{}ग\appbrk{}न्धांश्च ये}\ A N H\textsuperscript{pc}; \skt{पङ्क\appbrk{}म\appbrk{}त्स्य\appbrk{}व\appbrk{}सा\appbrk{}तै\appbrk{}लं\appbrk{}घृ\appbrk{}त\appbrk{}ग\appbrk{}न्धाश्च ये}\ H\textsuperscript{ac}}}\skt{ }\edtext{\skt{नराः}}{\lemma{\skt{नराः}}\Afootnote{K N; \skt{नराः}\ A H}}\skt{~। }\pend
\pstart\versepar \skt{}\edtext{\skt{पिष्टगन्धांश्च ये}}{\lemma{\skt{पिष्टगन्धांश्च ये}}\Afootnote{K; \skt{मृष्टगन्धांश्च ये}\ A; \skt{पिष्टगन्धांश्च यो}\ N\textsuperscript{ac}; \skt{पिष्टगन्धांश्च ये}\ N\textsuperscript{pc}; \skt{विष्टगन्धाश्च ये}\ H}}\skt{ वान्ति }\edtext{\skt{गतास्ते}}{\lemma{\skt{गतास्ते}}\Afootnote{K N H; \skt{गन्तार}\ A}}\skt{ }\edtext{\skt{यमसादनं}}{\lemma{\skt{यमसादनं}}\Afootnote{K N H; \skt{स्ते यमालयम्}\ A}}\skt{~।}\pend

\pstart\versepar \ledsidenote{\passref{1.31.26}}\skt{यूका ललाटमायान्ति }\edtext{\skt{बलिन्नाश्नन्ति}}{\lemma{\skt{नाश्नन्ति}}\Afootnote{A K H; \skt{नाश्नान्ति}\ N}}\skt{ वायसाः~। }\pend
\pstart\versepar \skt{येषाञ्चापि }\edtext{\skt{रतिन्नास्ति}}{\lemma{\skt{रतिन्}}\Afootnote{K N; \skt{रतिर्}\ A H}}\skt{ }\edtext{\skt{गतास्ते}}{\lemma{\skt{गतास्ते}}\Afootnote{K N H; \skt{यातार}\ A}}\skt{ }\edtext{\skt{यमसादनं}}{\lemma{\skt{यमसादनं}}\Afootnote{K N H; \skt{स्ते यमालयम्}\ A}}\skt{~।}\pend

\pstart\versepar \ledsidenote{\passref{1.31.27}}\edtext{\skt{ज्वरा\-ती\-सा\-र\-शोफा }\edtext{\skt{स्युर्य\-स्या\-न्यो\-न्या\-व\-सा\-दिनः}}{\lemma{\skt{यस्या\appbrk{}न्यो\appbrk{}न्या\appbrk{}व\appbrk{}सा\appbrk{}दिनः}}\Afootnote{A K N; \skt{यस्या\appbrk{}न्यो\appbrk{}न्या\appbrk{}व\appbrk{}सा\appbrk{}दिनं}\ H}}}{\lemma{\skt{ज्वरा\appbrk{}ती\appbrk{}सा\appbrk{}र\appbrk{}शोफा स्युर्}}\Afootnote{K N; \skt{ज्वरा\appbrk{}ति\appbrk{}सा\appbrk{}र\appbrk{}शोफाः स्युर्}\ A; \skt{ज्वरा\appbrk{}ती\appbrk{}सा\appbrk{}र\appbrk{}शोफाः स्युर्}\ H}}\skt{~। }\pend
\pstart\versepar \skt{}\edtext{\skt{प्रक्षीणञ्च}}{\lemma{\skt{प्रक्षीणञ्च}}\Afootnote{K N H; \skt{प्रक्षीणब}\ A}}\skt{ }\edtext{\skt{बलं}}{\lemma{\skt{बलं}}\Afootnote{K N H; \skt{लमां}\ A}}\skt{ }\edtext{\skt{यस्य न स}}{\lemma{\skt{यस्य न स}}\Afootnote{K N H; \skt{सस्य नासौ}\ A}}\skt{ }\edtext{\edtext{\skt{शक्यश्चिकित्सितुं}}{\lemma{\skt{शक्यश्}}\Afootnote{K H; \skt{शक्य}\ A N}}}{\lemma{\skt{चिकित्सितुं}}\Afootnote{K N H; \skt{श्चिकित्सितम्}\ A}}\skt{~।}\pend

\pstart\versepar \ledsidenote{\passref{1.31.28}}\skt{क्षीणस्य यस्य क्षुत्तृष्णे }\edtext{\edtext{\skt{हृद्यै\-र्मृ\-ष्टै\-र्हि\-तै\-स्तथा}}{\lemma{\skt{मृष्टैर्}}\Afootnote{K N H; \skt{मिष्टैर्}\ A}}}{\lemma{\skt{तथा}}\Afootnote{K N H; \skt{तथा~। न}\ A}}\skt{~। }\pend
\pstart\versepar \skt{}\edtext{\skt{अन्नपानैर्न्न}}{\lemma{\skt{अन्नपानैर्न्न}}\Afootnote{K N H; \skt{शाम्य\appbrk{}तोऽ\appbrk{}न्न\appbrk{}पा\appbrk{}नैश्च}\ A}}\skt{ }\edtext{\skt{सास्येते}}{\lemma{\skt{सास्येते}}\Afootnote{K; \emph{om.}\ A; \skt{शास्येत}\ N H}}\skt{ तस्य मृत्युरुपस्थितः~।}\pend

\pstart\versepar \ledsidenote{\passref{1.31.29}}\skt{प्रवाहिका शिरःशूलं }\edtext{\skt{कोष्ठशूलश्च}}{\lemma{\skt{कोष्ठशूलश्च}}\Afootnote{K N; \skt{कोष्ठशूलं च}\ A H}}\skt{ }\edtext{\skt{दारुण}}{\lemma{\skt{दारुण}}\Afootnote{K\textsuperscript{pc}; \skt{दारुणम्}\ A N H; \skt{दारुणः}\ K\textsuperscript{ac}}}\skt{~। }\pend
\pstart\versepar \skt{पिपासा बलहानिश्च तस्य मृत्युरुपस्थितः~।}\pend

\pstart\versepar \ledsidenote{\passref{1.31.30}}\skt{विष\-मे\-णो\-प\-चा\-रेण कर्मभिश्च पुराकृतैः~। }\pend
\pstart\versepar \skt{}\edtext{\skt{अनित्यत्वाच्च}}{\lemma{\skt{अनित्यत्वाच्च}}\Afootnote{K H; \skt{अनित्यत्वाच्च}\ A; \skt{अनित्याच्च}\ N}}\skt{ जन्तूनां जीवितं निधनम्व्रजेत्~।}\pend

\pstart\versepar \ledsidenote{\passref{1.31.31}}\skt{प्रेता }\edtext{\skt{दूताः}}{\lemma{\skt{दूताः}}\Afootnote{K H; \skt{भूताः}\ A; \skt{दूता}\ N}}\skt{ पिशाचाश्च रक्षांसि विविधानि च~। }\pend
\pstart\versepar \skt{}\edtext{\edtext{\skt{मर\-णी\-य\-न्न\-र\-न्नि\-त्य\-मु\-प\-स\-र्पन्ति}}{\lemma{\skt{मरणीयन्}}\Afootnote{K N H; \skt{मरणाभि}\ A}}}{\lemma{\skt{नरन्}}\Afootnote{K N H; \skt{मुखं}\ A}}\skt{ }\edtext{\skt{सर्वदा}}{\lemma{\skt{सर्वदा}}\Afootnote{K N H; \skt{मानवम्}\ A}}\pend

\pstart\versepar \ledsidenote{\passref{1.31.32}}\skt{तानि }\edtext{\skt{भेषजवीर्याणि}}{\lemma{\skt{भेषजवीर्याणि}}\Afootnote{A K N; \skt{भैषजवीर्याणि}\ H}}\skt{ }\edtext{\skt{प्रतिनिघ्नन्ति}}{\lemma{\skt{प्रतिनिघ्नन्ति}}\Afootnote{K N H; \skt{प्रतिघ्नन्ति जि}\ A}}\skt{ }\edtext{\skt{सर्वदा}}{\lemma{\skt{सर्वदा}}\Afootnote{K N H; \skt{घांसया}\ A}}\skt{~। }\pend
\pstart\versepar \skt{}\edtext{\skt{तस्मान्मोहाः}}{\lemma{\skt{तस्मान्मोहाः}}\Afootnote{K; \skt{तस्मान्मोघाः}\ A; \skt{तस्मात्मोहा}\ N; \skt{तस्मात्मोहाः}\ H}}\skt{ }\edtext{\skt{क्रियास्सर्वा}}{\lemma{\skt{क्रियास्}}\Afootnote{A K H; \skt{क्रिया}\ N}}\skt{ }\edtext{\skt{भवन्तीह}}{\lemma{\skt{भवन्तीह}}\Afootnote{K N H; \skt{भवन्त्येव}\ A}}\skt{ }\edtext{\skt{गतायुषमिति}}{\lemma{\skt{गतायुषमिति}}\Afootnote{K; \skt{गतायुषाम्}\ A; \skt{गतायुषम्}\ N H}}\skt{~॥}\pend

\pstart \ledsidenote{\passref{1.31.trailer}}\edtext{\skt{॥ ३१~॥}}{\lemma{\skt{॥ ३१~॥}}\Afootnote{\skt{इति सुश्रु\appbrk{}त\appbrk{}सं\appbrk{}हि\appbrk{}तायां सूत्रस्थाने छाया\appbrk{}वि\appbrk{}प्र\appbrk{}ति\appbrk{}प\appbrk{}त्ति\appbrk{}र्ना\appbrk{}मै\appbrk{}क\appbrk{}त्रिं\appbrk{}श\appbrk{}त्त\appbrk{}मोऽ\appbrk{}ध्यायः}\ A; \skt{३१~॥ {\symfont ❈}}\ K; \skt{इति}\ N H}}\pend

\endnumbering
